% Auto-generated from data/segments.json + layout.json (+ transforms) — do not edit by hand.
% volume: 04Vin04
% mode: reading
% segments: 2759
% transform_rules: 0
% toc_marks: 442

% Volume layout defaults (layout.json ``layout``).
\csromanlayoutapply{1}{1.5}{21.6pt}{5pt}{6.3pt}{65pt}{2.5em}%

% Reading physical-page layout (layout.json ``page_layout_reading_mode``).
\csromanreadingpagelayoutvolume{1}{1.5}{21.6pt}{5pt}{6.3pt}{65pt}{2.5em}%
\csromanreadingpagelayoutenable

\pitaka{\textbf{วินยปิฏก}}
\csromanmatikaanchor{mtk.0}
\csromantocmarknopage{chapter}{จูฬวคฺคปาฬิ}{mtk.0}
\bookmark[dest={mtk.0},level=0]{จูฬวคฺคปาฬิ}
\gambhira{\textbf{จูฬวคฺคปาฬิ}}
\namakkaram{นโม ตสฺส ภควโต อรหโต สมฺมาสมฺพุทฺธสฺส.}
\csromanmatikaanchor{mtk.1}
\csromantocmarknopage{chapter}{1. กมฺมกฺขนฺธก}{mtk.1}
\bookmark[dest={mtk.1},level=0]{1. กมฺมกฺขนฺธก}
\chapterhead{1. กมฺมกฺขนฺธก}
\csromanmatikaanchor{mtk.2}
\csromantocmarkref{section}{1. ตชฺชนียกมฺม}{mtk.2}
\bookmark[dest={mtk.2},level=1]{1. ตชฺชนียกมฺม}
\csromanheader{1}{1. ตชฺชนียกมฺม}
\proseitem{1}{\csromanfolio{1}เตน สมเยน พุทฺโธ ภควา สาวตฺถิยํ วิหรติ เชตวเน อนาถปิณฺฑิกสฺส อาราเม.\csromanspacer{} เตน โข ปน สมเยน ปณฺฑุกโลหิตกา ภิกฺขู อตฺตนา ภณฺฑนการกา กลหการกา วิวาทการกา ภสฺสการกา สํเฆ อธิกรณการกา, เยปิ จญฺเญ ภิกฺขู ภณฺฑนการกา กลหการกา วิวาทการกา ภสฺสการกา สํเฆ อธิกรณการกา, เต อุปสงฺกมิตฺวา เอวํ วทนฺติ “มา โข ตุเมฺห อายสฺมนฺโต เอโส อเชสิ, พลวาพลวํ ปฏิมนฺเตถ, ตุเมฺห เตน ปณฺฑิตตรา จ พฺยตฺตตรา จ พหุสฺสุตตรา จ อลมตฺตตรา จ\footnote{อลมตฺถตรา จ (สฺยา, ก)}, มา จสฺส ภายิตฺถ, มยมฺปิ ตุมฺหากํ ปกฺขา ภวิสฺสามา”ติ.\csromanspacer{} เตน อนุปฺปนฺนานิ เจว ภณฺฑนานิ อุปฺปชฺชนฺติ, อุปฺปนฺนานิ จ ภณฺฑนานิ ภิยฺโยภาวาย เวปุลฺลาย สํวตฺตนฺติ.\csromanspacer{} เย เต ภิกฺขู อปฺปิจฺฉา ฯเปฯ เต อุชฺฌายนฺติ ขิยฺยนฺติ วิปาเจนฺติ “กถํ หิ นาม ปณฺฑุกโลหิตกา ภิกฺขู อตฺตนา ภณฺฑนการกา กลหการกา วิวาทการกา ภสฺสการกา สํเฆ อธิกรณการกา, เยปิ จญฺเญ ภิกฺขู ภณฺฑนการกา กลหการกา วิวาทการกา ภสฺสการกา สํเฆ อธิกรณการกา, เต อุปสงฺกมิตฺวา เอวํ วกฺขนฺติ ‘มา โข ตุเมฺห อายสฺมนฺโต เอโส อเชสิ, \csromanfolio{2}พลวาพลวํ ปฏิมนฺเตถ, ตุเมฺห เตน ปณฺฑิตตรา จ พฺยตฺตตรา จ พหุสฺสุตตรา จ อลมตฺตตรา จ, มา จสฺส ภายิตฺถ, มยมฺปิ ตุมฺหากํ ปกฺขา ภวิสฺสามา’ติ, เตน อนุปฺปนฺนานิ เจว ภณฺฑนานิ อุปฺปชฺชนฺติ, อุปฺปนฺนานิ จ ภณฺฑนานิ ภิยฺโยภาวาย เวปุลฺลาย สํวตฺตนฺตี”ติ.}

\proseitem{2}{อถ โข เต ภิกฺขู ภควโต เอตมตฺถํ อาโรเจสุํ.\csromanspacer{} อถ โข ภควา เอตสฺมิํ นิทาเน เอตสฺมิํ ปกรเณ ภิกฺขุสํฆํ สนฺนิปาตาเปตฺวา ภิกฺขู ปฏิปุจฺฉิ “สจฺจํ กิร ภิกฺขเว ปณฺฑุกโลหิตกา ภิกฺขู อตฺตนา ภณฺฑนการกา กลหการกา วิวาทการกา ภสฺสการกา สํเฆ อธิกรณการกา, เยปิ จญฺเญ ภิกฺขู ภณฺฑนการกา กลหการกา วิวาทการกา ภสฺสการกา สํเฆ อธิกรณการกา, เต อุปสงฺกมิตฺวา เอวํ วทนฺติ ‘มา โข ตุเมฺห อายสฺมนฺโต เอโส อเชสิ, พลวาพลวํ ปฏิมนฺเตถ, ตุเมฺห เตน ปณฺฑิตตรา จ พฺยตฺตตรา จ พหุสฺสุตตรา จ อลมตฺตตรา จ, มา จสฺส ภายิตฺถ, มยมฺปิ ตุมฺหากํ ปกฺขา ภวิสฺสามา’ติ, เตน อนุปฺปนฺนานิ เจว ภณฺฑนานิ อุปฺปชฺชนฺติ, อุปฺปนฺนานิ จ ภณฺฑนานิ ภิยฺโยภาวาย เวปุลฺลาย สํวตฺตนฺตี”ติ.\csromanspacer{} สจฺจํ ภควาติ.\csromanspacer{} วิครหิ พุทฺโธ ภควา “อนนุจฺฉวิกํ ภิกฺขเว เตสํ โมฆปุริสานํ อนนุโลมิกํ อปฺปติรูปํ อสฺสามณกํ อกปฺปิยํ อกรณียํ, กถํ หิ นาม เต ภิกฺขเว โมฆปุริสา อตฺตนา ภณฺฑนการกา ฯเปฯ สํเฆ อธิกรณการกา, เยปิ จญฺเญ ภิกฺขู ภณฺฑนการกา ฯเปฯ สํเฆ อธิกรณการกา, เต อุปสงฺกมิตฺวา เอวํ วกฺขนฺติ ‘มา โข ตุเมฺห อายสฺมนฺโต เอโส อเชสิ, พลวาพลวํ ปฏิมนฺเตถ, ตุเมฺห เตน ปณฺฑิตตรา จ พฺยตฺตตรา จ พหุสฺสุตตรา จ อลมตฺตตรา จ, มา จสฺส ภายิตฺถ, มยมฺปิ ตุมฺหากํ ปกฺขา ภวิสฺสามา’ติ, เตน อนุปฺปนฺนานิ เจว ภณฺฑนานิ อุปฺปชฺชนฺติ, อุปฺปนฺนานิ จ ภณฺฑนานิ ภิยฺโยภาวาย เวปุลฺลาย สํวตฺตนฺติ.\csromanspacer{} เนตํ ภิกฺขเว อปฺปสนฺนานํ วา ปสาทาย ปสนฺนานํ วา ภิยฺโยภาวาย, อถ เขฺวตํ ภิกฺขเว อปฺปสนฺนานญฺเจว อปฺปสาทาย ปสนฺนานญฺจ เอกจฺจานํ อญฺญถตฺตายา”ติ.}

\prose{อถ โข ภควา เต\footnote{ปณฺฑุกโลหิตเก (สฺยา)} ภิกฺขู อเนกปริยาเยน วิครหิตฺวา ทุพฺภรตาย ทุปฺโปสตาย มหิจฺฉตาย อสนฺตุฏฺฐิตาย\footnote{อสนฺตุฏฺฐตาย (สฺยา), อสนฺตุฏฺฐิยา (สี)} สงฺคณิกาย \csromanfolio{3}โกสชฺชสฺส อวณฺณํ ภาสิตฺวา อเนกปริยาเยน สุภรตาย สุโปสตาย อปฺปิจฺฉสฺส สนฺตุฏฺฐสฺส สลฺเลขสฺส ธุตสฺส ปาสาทิกสฺส อปจยสฺส วีริยารมฺภสฺส\footnote{วิริยารมฺภสฺส (สี), วีริยารพฺภสฺส (ก)} วณฺณํ ภาสิตฺวา ภิกฺขูนํ ตทนุจฺฉวิกํ ตทนุโลมิกํ ธมฺมิํ กถํ กตฺวา ภิกฺขู อามนฺเตสิ\csromandash{}เตน หิ ภิกฺขเว สํโฆ ปณฺฑุกโลหิตกานํ ภิกฺขูนํ ตชฺชนียกมฺมํ กโรตุ.\csromanspacer{} เอวญฺจ ปน ภิกฺขเว กาตพฺพํ, ปฐมํ ปณฺฑุกโลหิตกา ภิกฺขู โจเทตพฺพา, โจเทตฺวา สาเรตพฺพา, สาเรตฺวา อาปตฺติํ\footnote{อาปตฺติ (สี, สฺยา)} อาโรเปตพฺพา, อาปตฺติํ อาโรเปตฺวา พฺยตฺเตน ภิกฺขุนา ปฏิพเลน สํโฆ ญาเปตพฺโพ\csromandash{}}

\proseitem{3}{“สุณาตุ เม ภนฺเต สํโฆ, อิเม ปณฺฑุกโลหิตกา ภิกฺขู อตฺตนา ภณฺฑนการกา กลหการกา วิวาทการกา ภสฺสการกา สํเฆ อธิกรณการกา, เยปิ จญฺเญ ภิกฺขู ภณฺฑนการกา กลหการกา วิวาทการกา ภสฺสการกา สํเฆ อธิกรณการกา, เต อุปสงฺกมิตฺวา เอวํ วทนฺติ ‘มา โข ตุเมฺห อายสฺมนฺโต เอโส อเชสิ, พลวาพลวํ ปฏิมนฺเตถ, ตุเมฺห เตน ปณฺฑิตตรา จ พฺยตฺตตฺตตรา จ พหุสฺสุตตรา จ อลมตฺตตรา จ, มา จสฺส ภายิตฺถ, มยมฺปิ ตุมฺหากํ ปกฺขา ภวิสฺสามา’ติ, เตน อนุปฺปนฺนานิ เจว ภณฺฑนานิ อุปฺปชฺชนฺติ, อุปฺปนฺนานิ จ ภณฺฑนานิ ภิยฺโยภาวาย เวปุลฺลาย สํวตฺตนฺติ, ยทิ สํฆสฺส ปตฺตกลฺลํ, สํโฆ ปณฺฑุกโลหิตกานํ ภิกฺขูนํ ตชฺชนียกมฺมํ กเรยฺย, เอสา ญตฺติ.}

\prose{สุณาตุ เม ภนฺเต สํโฆ, อิเม ปณฺฑุกโลหิตกา ภิกฺขู อตฺตนา ภณฺฑนการกา ฯเปฯ สํเฆ อธิกรณการกา, เยปิ จญฺเญ ภิกฺขู ภณฺฑนการกา ฯเปฯ สํเฆ อธิกรณการกา, เต อุปสงฺกมิตฺวา เอวํ วทนฺติ ‘มา โข ตุเมฺห อายสฺมนฺโต เอโส อเชสิ, พลวาพลวํ ปฏิมนฺเตถ, ตุเมฺห เตน ปณฺฑิตตรา จ พฺยตฺตตรา จ พหุสฺสุตตรา จ อลมตฺตตรา จ, มา จสฺส ภายิตฺถ, มยมฺปิ ตุมฺหากํ ปกฺขา ภวิสฺสามา’ติ, เตน อนุปฺปนฺนานิ เจว ภณฺฑนานิ อุปฺปชฺชนฺติ, อุปฺปนฺนานิ จ ภณฺฑนานิ ภิยฺโยภาวาย เวปุลฺลาย สํวตฺตนฺติ, สํโฆ ปณฺฑุกโลหิตกานํ. \csromanfolio{4}ภิกฺขูนํ ตชฺชนียกมฺมํ กโรติ, ยสฺสายสฺมโต ขมติ ปณฺฑุกโลหิตกานํ ภิกฺขูนํ ตชฺชนียสฺส กมฺมสฺส กรณํ, โส ตุณฺหสฺส, ยสฺส นกฺขมติ, โส ภาเสยฺย.}

\prose{ทุติยมฺปิ เอตมตฺถํ วทามิ.\csromanspacer{} สุณาตุ เม ภนฺเต สํโฆ, อิเม ปณฺฑุกโลหิตกา ภิกฺขู อตฺตนา ภณฺฑนการกา ฯเปฯ สํเฆ อธิกรณการกา, เยปิ จญฺเญ ภิกฺขู ภณฺฑนการกา ฯเปฯ สํเฆ อธิกรณการกา, เต อุปสงฺกมิตฺวา เอวํ วทนฺติ ‘มา โข ตุเมฺห อายสฺมนฺโต เอโส อเชสิ, พลวาพลวํ ปฏิมนฺเตถ, ตุเมฺห เตน ปณฺฑิตตรา จ พฺยตฺตตรา จ พหุสฺสุตตรา จ อลมตฺตตรา จ, มา จสฺส ภายิตฺถ, มยมฺปิ ตุมฺหากํ ปกฺขา ภวิสฺสามา’ติ, เตน อนุปฺปนฺนานิ เจว ภณฺฑนานิ อุปฺปชฺชนฺติ, อุปฺปนฺนานิ จ ภณฺฑนานิ ภิยฺโยภาวาย เวปุลฺลาย สํวตฺตนฺติ, สํโฆ ปณฺฑุกโลหิตกานํ ภิกฺขูนํ ตชฺชนียกมฺมํ กโรติ, ยสฺสายสฺมโต ขมติ ปณฺฑุกโลหิตกานํ ภิกฺขูนํ ตชฺชนียสฺส กมฺมสฺส กรณํ, โส ตุณฺหสฺส, ยสฺส นกฺขมติ, โส ภาเสยฺย.}

\prose{ตติยมฺปิ เอตมตฺถํ วทามิ.\csromanspacer{} สุณาตุ เม ภนฺเต สํโฆ, อิเม ปณฺฑุกโลหิตกา ภิกฺขู อตฺตนา ภณฺฑนการกา ฯเปฯ สํเฆ อธิกรณการกา, เยปิ จญฺเญ ภิกฺขู ภณฺฑนการกา ฯเปฯ สํเฆ อธิกรณการกา, เต อุปสงฺกมิตฺวา เอวํ วทนฺติ ‘มา โข ตุเมฺห อายสฺมนฺโต เอโส อเชสิ, พลวาพลวํ ปฏิมนฺเตถ, ตุเมฺห เตน ปณฺฑิตตรา จ พฺยตฺตตรา จ พหุสฺสุตตรา จ อลมตฺตตรา จ, มา จสฺส ภายิตฺถ, มยมฺปิ ตุมฺหากํ ปกฺขา ภวิสฺสามา’ติ, เตน อนุปฺปนฺนานิ เจว ภณฺฑนานิ อุปฺปชฺชนฺติ, อุปฺปนฺนานิ จ ภณฺฑนานิ ภิยฺโยภาวาย เวปุลฺลาย สํวตฺตนฺติ, สํโฆ ปณฺฑุกโลหิตกานํ ภิกฺขูนํ ตชฺชนียกมฺมํ กโรติ, ยสฺสายสฺมโต ขมติ ปณฺฑุกโลหิตกานํ ภิกฺขูนํ ตชฺชนียสฺส กมฺมสฺส กรณํ, โส ตุณฺหสฺส, ยสฺส นกฺขมติ, โส ภาเสยฺย.}

\prose{กตํ สํเฆน ปณฺฑุกโลหิตกานํ ภิกฺขูนํ ตชฺชนียกมฺมํ, ขมติ สํฆสฺส, ตสฺมา ตุณฺหี, เอวเมตํ ธารยามี”ติ.}

\csromanmatikaanchor{mtk.3}
\csromantocmarkref{subsection}{อธมฺมกมฺมทฺวาทสก}{mtk.3}
\bookmark[dest={mtk.3},level=2]{อธมฺมกมฺมทฺวาทสก}
\csromanheader{2}{\textbf{อธมฺมกมฺมทฺวาทสก}}
\proseitem{4}{\csromanfolio{5}ตีหิ ภิกฺขเว องฺเคหิ สมนฺนาคตํ ตชฺชนียกมฺมํ อธมฺมกมฺมญฺจ\footnote{อธมฺมกมฺมญฺเจว (สฺยา)} โหติ อวินยกมฺมญฺจ ทุวูปสนฺตญฺจ.\csromanspacer{} อสมฺมุขา กตํ โหติ, อปฺปฏิปุจฺฉา กตํ โหติ, อปฺปฏิญฺญาย กตํ โหติ.\csromanspacer{} อิเมหิ โข ภิกฺขเว ตีหงฺเคหิ สมนฺนาคตํ ตชฺชนียกมฺมํ อธมฺมกมฺมญฺจ โหติ อวินยกมฺมญฺจ ทุวูปสนฺตญฺจ. (1)}

\prose{อปเรหิปิ ภิกฺขเว ตีหงฺเคหิ สมนฺนาคตํ ตชฺชนียกมฺมํ อธมฺมกมฺมญฺจ โหติ อวินยกมฺมญฺจ ทุวูปสนฺตญฺจ.\csromanspacer{} อนาปตฺติยา กตํ โหติ, อเทสนาคามินิยา อาปตฺติยา กตํ โหติ, เทสิตาย อาปตฺติยา กตํ โหติ.\csromanspacer{} อิเมหิ โข ภิกฺขเว ตีหงฺเคหิ สมนฺนาคตํ ตชฺชนียกมฺมํ อธมฺมกมฺมญฺจ โหติ อวินยกมฺมญฺจ ทุวูปสนฺตญฺจ. (2)}

\prose{อปเรหิปิ ภิกฺขเว ตีหงฺเคหิ สมนฺนาคตํ ตชฺชนียกมฺมํ อธมฺมกมฺมญฺจ โหติ อวินยกมฺมญฺจ ทุวูปสนฺตญฺจ.\csromanspacer{} อโจเทตฺวา กตํ โหติ, อสาเรตฺวา กตํ โหติ, อาปตฺติํ อนาโรเปตฺวา กตํ โหติ.\csromanspacer{} อิเมหิ โข ภิกฺขเว ตีหงฺเคหิ สมนฺนาคตํ ตชฺชนียกมฺมํ อธมฺมกมฺมญฺจ โหติ อวินยกมฺมญฺจ ทุวูปสนฺตญฺจ. (3)}

\prose{อปเรหิปิ ภิกฺขเว ตีหงฺเคหิ สมนฺนาคตํ ตชฺชนียกมฺมํ อธมฺมกมฺมญฺจ โหติ อวินยกมฺมญฺจ ทุวูปสนฺตญฺจ.\csromanspacer{} อสมฺมุขากตํ โหติ, อธมฺเมน กตํ โหติ, วคฺเคน กตํ โหติ.\csromanspacer{} อิเมหิ โข ภิกฺขเว ตีหงฺเคหิ สมนฺนาคตํ ตชฺชนียกมฺมํ อธมฺมกมฺมญฺจ โหติ อวินยกมฺมญฺจ ทุวูปสนฺตญฺจ. (4)}

\prose{อปเรหิปิ ภิกฺขเว ตีหงฺเคหิ สมนฺนาคตํ ตชฺชนียกมฺมํ อธมฺมกมฺมญฺจ โหติ อวินยกมฺมญฺจ ทุวูปสนฺตญฺจ.\csromanspacer{} อปฺปฏิปุจฺฉา กตํ โหติ, อธมฺเมน กตํ โหติ, วคฺเคน กตํ โหติ.\csromanspacer{} อิเมหิ โข ภิกฺขเว ตีหงฺเคหิ สมนฺนาคตํ ตชฺชนียกมฺมํ อธมฺมกมฺมญฺจ โหติ อวินยกมฺมญฺจ ทุวูปสนฺตญฺจ. (5)}

\prose{อปเรหิปิ ภิกฺขเว ตีหงฺเคหิ สมนฺนาคตํ ตชฺชนียกมฺมํ อธมฺมกมฺมญฺจ โหติ อวินยกมฺมญฺจ ทุวูปสนฺตญฺจ.\csromanspacer{} อปฺปฏิญฺญาย กตํ โหติ, อธมฺเมน กตํ โหติ, วคฺเคน กตํ โหติ.\csromanspacer{} อิเมหิ โข ภิกฺขเว ตีหงฺเคหิ สมนฺนาคตํ ตชฺชนียกมฺมํ อธมฺมกมฺมญฺจ โหติ อวินยกมฺมญฺจ ทุวูปสนฺตญฺจ. (6)}

\prose{\csromanfolio{6}อปเรหิปิ ภิกฺขเว ตีหงฺเคหิ สมนฺนาคตํ ตชฺชนียกมฺมํ อธมฺมกมฺมญฺจ โหติ อวินยกมฺมญฺจ ทุวูปสนฺตญฺจ.\csromanspacer{} อนาปตฺติยา กตํ โหติ, อธมฺเมน กตํ โหติ, วคฺเคน กตํ โหติ.\csromanspacer{} อิเมหิ โข ภิกฺขเว ตีหงฺเคหิ สมนฺนาคตํ ตชฺชนียกมฺมํ อธมฺมกมฺมญฺจ โหติ อวินยกมฺมญฺจ ทุวูปสนฺตญฺจ. (7)}

\prose{อปเรหิปิ ภิกฺขเว ตีหงฺเคหิ สมนฺนาคตํ ตชฺชนียกมฺมํ อธมฺมกมฺมญฺจ โหติ อวินยกมฺมญฺจ ทุวูปสนฺตญฺจ.\csromanspacer{} อเทสนาคามินิยา อาปตฺติยา กตํ โหติ, อธมฺเมน กตํ โหติ, วคฺเคน กตํ โหติ.\csromanspacer{} อิเมหิ โข ภิกฺขเว ตีหงฺเคหิ สมนฺนาคตํ ตชฺชนียกมฺมํ อธมฺมกมฺมญฺจ โหติ อวินยกมฺมญฺจ ทุวูปสนฺตญฺจ. (8)}

\prose{อปเรหิปิ ภิกฺขเว ตีหงฺเคหิ สมนฺนาคตํ ตชฺชนียกมฺมํ อธมฺมกมฺมญฺจ โหติ อวินยกมฺมญฺจ ทุวูปสนฺตญฺจ.\csromanspacer{} เทสิตาย อาปตฺติยา กตํ โหติ, อธมฺเมน กตํ โหติ, วคฺเคน กตํ โหติ.\csromanspacer{} อิเมหิ โข ภิกฺขเว ตีหงฺเคหิ สมนฺนาคตํ ตชฺชนียกมฺมํ อธมฺมกมฺมญฺจ โหติ อวินยกมฺมญฺจ ทุวูปสนฺตญฺจ. (9)}

\prose{อปเรหิปิ ภิกฺขเว ตีหงฺเคหิ สมนฺนาคตํ ตชฺชนียกมฺมํ อธมฺมกมฺมญฺจ โหติ อวินยกมฺมญฺจ ทุวูปสนฺตญฺจ.\csromanspacer{} อโจเทตฺวา กตํ โหติ, อธมฺเมน กตํ โหติ, วคฺเคน กตํ โหติ.\csromanspacer{} อิเมหิ โข ภิกฺขเว ตีหงฺเคหิ สมนฺนาคตํ ตชฺชนียกมฺมํ อธมฺมกมฺมญฺจ โหติ อวินยกมฺมญฺจ ทุวูปสนฺตญฺจ. (10)}

\prose{อปเรหิปิ ภิกฺขเว ตีหงฺเคหิ สมนฺนาคตํ ตชฺชนียกมฺมํ อธมฺมกมฺมญฺจ โหติ อวินยกมฺมญฺจ ทุวูปสนฺตญฺจ.\csromanspacer{} อสาเรตฺวา กตํ โหติ, อธมฺเมน กตํ โหติ, วคฺเคน กตํ โหติ.\csromanspacer{} อิเมหิ โข ภิกฺขเว ตีหงฺเคหิ สมนฺนาคตํ ตชฺชนียกมฺมํ อธมฺมกมฺมญฺจ โหติ อวินยกมฺมญฺจ ทุวูปสนฺตญฺจ. (11)}

\prose{อปเรหิปิ ภิกฺขเว ตีหงฺเคหิ สมนฺนาคตํ ตชฺชนียกมฺมํ อธมฺมกมฺมญฺจ โหติ อวินยกมฺมญฺจ ทุวูปสนฺตญฺจ.\csromanspacer{} อาปตฺติํ อนาโรเปตฺวา กตํ โหติ, อธมฺเมน กตํ โหติ, วคฺเคน กตํ โหติ.\csromanspacer{} อิเมหิ โข \csromanfolio{7}ภิกฺขเว ตีหงฺเคหิ สมนฺนาคตํ ตชฺชนียกมฺมํ อธมฺมกมฺมญฺจ โหติ อวินยกมฺมญฺจ ทุวูปสนฺตญฺจ. (12)}

\nitthitam{อธมฺมกมฺมทฺวาทสกํ นิฏฺฐิตํ.}
\csromanmatikaanchor{mtk.4}
\csromantocmarkref{subsection}{ธมฺมกมฺมทฺวาทสก}{mtk.4}
\bookmark[dest={mtk.4},level=2]{ธมฺมกมฺมทฺวาทสก}
\csromanheader{2}{\textbf{ธมฺมกมฺมทฺวาทสก}}
\proseitem{5}{ตีหิ ภิกฺขเว องฺเคหิ สมนฺนาคตํ ตชฺชนียกมฺมํ ธมฺมกมฺมญฺจ โหติ วินยกมฺมญฺจ สุวูปสนฺตญฺจ.\csromanspacer{} สมฺมุขา กตํ โหติ, ปฏิปุจฺฉา กตํ โหติ, ปฏิญฺญาย กตํ โหติ.\csromanspacer{} อิเมหิ โข ภิกฺขเว ตีหงฺเคหิ สมนฺนาคตํ ตชฺชนียกมฺมํ ธมฺมกมฺมญฺจ โหติ วินยกมฺมญฺจ สุวูปสนฺตญฺจ. (1)}

\prose{อปเรหิปิ ภิกฺขเว ตีหงฺเคหิ สมนฺนาคตํ ตชฺชนียกมฺมํ ธมฺมกมฺมญฺจ โหติ วินยกมฺมญฺจ สุวูปสนฺตญฺจ.\csromanspacer{} อาปตฺติยา กตํ โหติ, เทสนาคามินิยา อาปตฺติยา กตํ โหติ, อเทสิตาย อาปตฺติยา กตํ โหติ.\csromanspacer{} อิเมหิ โข ภิกฺขเว ตีหงฺเคหิ สมนฺนาคตํ ตชฺชนียกมฺมํ ธมฺมกมฺมญฺจ โหติ วินยกมฺมญฺจ สุวูปสนฺตญฺจ. (2)}

\prose{อปเรหิปิ ภิกฺขเว ตีหงฺเคหิ สมนฺนาคตํ ตชฺชนียกมฺมํ ธมฺมกมฺมญฺจ โหติ วินยกมฺมญฺจ สุวูปสนฺตญฺจ.\csromanspacer{} โจเทตฺวา กตํ โหติ, สาเรตฺวา กตํ โหติ, อาปตฺติํ อาโรเปตฺวา กตํ โหติ.\csromanspacer{} อิเมหิ โข ภิกฺขเว ตีหงฺเคหิ สมนฺนาคตํ ตชฺชนียกมฺมํ ธมฺมกมฺมญฺจ โหติ วินยกมฺมญฺจ สุวูปสนฺตญฺจ. (3)}

\prose{อปเรหิปิ ภิกฺขเว ตีหงฺเคหิ สมนฺนาคตํ ตชฺชนียกมฺมํ ธมฺมกมฺมญฺจ โหติ วินยกมฺมญฺจ สุวูปสนฺตญฺจ.\csromanspacer{} สมฺมุขา กตํ โหติ, ธมฺเมน กตํ โหติ, สมคฺเคน กตํ โหติ.\csromanspacer{} อิเมหิ โข ภิกฺขเว ตีหงฺเคหิ สมนฺนาคตํ ตชฺชนียกมฺมํ ธมฺมกมฺมญฺจ โหติ วินยกมฺมญฺจ สุวูปสนฺตญฺจ. (4)}

\prose{อปเรหิปิ ภิกฺขเว ตีหงฺเคหิ สมนฺนาคตํ ตชฺชนียกมฺมํ ธมฺมกมฺมญฺจ โหติ วินยกมฺมญฺจ สุวูปสนฺตญฺจ.\csromanspacer{} ปฏิปุจฺฉา กตํ โหติ, ธมฺเมน กตํ โหติ, สมคฺเคน กตํ โหติ.\csromanspacer{} อิเมหิ โข ภิกฺขเว ตีหงฺเคหิ สมนฺนาคตํ ตชฺชนียกมฺมํ ธมฺมกมฺมญฺจ โหติ วินยกมฺมญฺจ สุวูปสนฺตญฺจ. (5)}

\prose{\csromanfolio{8}อปเรหิปิ ภิกฺขเว ตีหงฺเคหิ สมนฺนาคตํ ตชฺชนียกมฺมํ ธมฺมกมฺมญฺจ โหติ วินยกมฺมญฺจ สุวูปสนฺตญฺจ.\csromanspacer{} ปฏิญฺญาย กตํ โหติ, ธมฺเมน กตํ โหติ, สมคฺเคน กตํ โหติ.\csromanspacer{} อิเมหิ โข ภิกฺขเว ตีหงฺเคหิ สมนฺนาคตํ ตชฺชนียกมฺมํ ธมฺมกมฺมญฺจ โหติ วินยกมฺมญฺจ สุวูปสนฺตญฺจ. (6)}

\prose{อปเรหิปิ ภิกฺขเว ตีหงฺเคหิ สมนฺนาคตํ ตชฺชนียกมฺมํ ธมฺมกมฺมญฺจ โหติ วินยกมฺมญฺจ สุวูปสนฺตญฺจ.\csromanspacer{} อาปตฺติยา กตํ โหติ, ธมฺเมน กตํ โหติ, สมคฺเคน กตํ โหติ.\csromanspacer{} อิเมหิ โข ภิกฺขเว ตีหงฺเคหิ สมนฺนาคตํ ตชฺชนียกมฺมํ ธมฺมกมฺมญฺจ โหติ วินยกมฺมญฺจ สุวูปสนฺตญฺจ. (7)}

\prose{อปเรหิปิ ภิกฺขเว ตีหงฺเคหิ สมนฺนาคตํ ตชฺชนียกมฺมํ ธมฺมกมฺมญฺจ โหติ วินยกมฺมญฺจ สุวูปสนฺตญฺจ.\csromanspacer{} เทสนาคามินิยา อาปตฺติยา กตํ โหติ, ธมฺเมน กตํ โหติ, สมคฺเคน กตํ โหติ.\csromanspacer{} อิเมหิ โข ภิกฺขเว ตีหงฺเคหิ สมนฺนาคตํ ตชฺชนียกมฺมํ ธมฺมกมฺมญฺจ โหติ วินยกมฺมญฺจ สุวูปสนฺตญฺจ. (8)}

\prose{อปเรหิปิ ภิกฺขเว ตีหงฺเคหิ สมนฺนาคตํ ตชฺชนียกมฺมํ ธมฺมกมฺมญฺจ โหติ วินยกมฺมญฺจ สุวูปสนฺตญฺจ.\csromanspacer{} อเทสิตาย อาปตฺติยา กตํ โหติ, ธมฺเมน กตํ โหติ, สมคฺเคน กตํ โหติ.\csromanspacer{} อิเมหิ โข ภิกฺขเว ตีหงฺเคหิ สมนฺนาคตํ ตชฺชนียกมฺมํ ธมฺมกมฺมญฺจ โหติ วินยกมฺมญฺจ สุวูปสนฺตญฺจ. (9)}

\prose{อปเรหิปิ ภิกฺขเว ตีหงฺเคหิ สมนฺนาคตํ ตชฺชนียกมฺมํ ธมฺมกมฺมญฺจ โหติ วินยกมฺมญฺจ สุวูปสนฺตญฺจ.\csromanspacer{} โจเทตฺวา กตํ โหติ, ธมฺเมน กตํ โหติ, สมคฺเคน กตํ โหติ.\csromanspacer{} อิเมหิ โข ภิกฺขเว ตีหงฺเคหิ สมนฺนาคตํ ตชฺชนียกมฺมํ ธมฺมกมฺมญฺจ โหติ วินยกมฺมญฺจ สุวูปสนฺตญฺจ. (10)}

\prose{อปเรหิปิ ภิกฺขเว ตีหงฺเคหิ สมนฺนาคตํ ตชฺชนียกมฺมํ ธมฺมกมฺมญฺจ โหติ วินยกมฺมญฺจ สุวูปสนฺตญฺจ.\csromanspacer{} สาเรตฺวา กตํ โหติ, ธมฺเมน กตํ โหติ, สมคฺเคน กตํ โหติ.\csromanspacer{} อิเมหิ โข ภิกฺขเว ตีหงฺเคหิ สมนฺนาคตํ ตชฺชนียกมฺมํ ธมฺมกมฺมญฺจ โหติ วินยกมฺมญฺจ สุวูปสนฺตญฺจ. (11)}

\prose{อปเรหิปิ ภิกฺขเว ตีหงฺเคหิ สมนฺนาคตํ ตชฺชนียกมฺมํ ธมฺมกมฺมญฺจ โหติ วินยกมฺมญฺจ สุวูปสนฺตญฺจ.\csromanspacer{} อาปตฺติํ อาโรเปตฺวา กตํ โหติ, ธมฺเมน กตํ โหติ, สมคฺเคน กตํ โหติ.\csromanspacer{} อิเมหิ โข ภิกฺขเว \csromanfolio{9}ตีหงฺเคหิ สมนฺนาคตํ ตชฺชนียกมฺมํ ธมฺมกมฺมญฺจ โหติ วินยกมฺมญฺจ สุวูปสนฺตญฺจ. (12)}

\nitthitam{ธมฺมกมฺมทฺวาทสกํ นิฏฺฐิตํ.}
\csromanmatikaanchor{mtk.5}
\csromantocmarkref{subsection}{อากงฺขมานฉกฺก}{mtk.5}
\bookmark[dest={mtk.5},level=2]{อากงฺขมานฉกฺก}
\csromanheader{2}{\textbf{อากงฺขมานฉกฺก}}
\proseitem{6}{\csromansymbolfootnote{*}{วิ 5. 219 ปิฏฺเฐปิ.}ตีหิ ภิกฺขเว องฺเคหิ สมนฺนาคตสฺส ภิกฺขุโน อากงฺขมาโน สํโฆ ตชฺชนียกมฺมํ กเรยฺย, ภณฺฑนการโก โหติ กลหการโก วิวาทการโก ภสฺสการโก สํเฆ อธิกรณการโก, พาโล โหติ อพฺยตฺโต อาปตฺติพหุโล อนปทาโน, คิหิสํสฏฺโฐ วิหรติ อนนุโลมิเกหิ คิหิสํสคฺเคหิ.\csromanspacer{} อิเมหิ โข ภิกฺขเว ตีหงฺเคหิ สมนฺนาคตสฺส ภิกฺขุโน อากงฺขมาโน สํโฆ ตชฺชนียกมฺมํ กเรยฺย. (1)}

\prose{อปเรหิปิ ภิกฺขเว ตีหงฺเคหิ สมนฺนาคตสฺส ภิกฺขุโน อากงฺขมาโน สํโฆ ตชฺชนียกมฺมํ กเรยฺย.\csromanspacer{} อธิสีเล สีลวิปนฺโน โหติ, อชฺฌาจาเร อาจารวิปนฺโน โหติ, อติทิฏฺฐิยา ทิฏฺฐิวิปนฺโน โหติ.\csromanspacer{} อิเมหิ โข ภิกฺขเว ตีหงฺเคหิ สมนฺนาคตสฺส ภิกฺขุโน อากงฺขมาโน สํโฆ ตชฺชนียกมฺมํ กเรยฺย. (2)}

\prose{อปเรหิปิ ภิกฺขเว ตีหงฺเคหิ สมนฺนาคตสฺส ภิกฺขุโน อากงฺขมาโน สํโฆ ตชฺชนียกมฺมํ กเรยฺย.\csromanspacer{} พุทฺธสฺส อวณฺณํ ภาสติ, ธมฺมสฺส อวณฺณํ ภาสติ, สํฆสฺส อวณฺณํ ภาสติ.\csromanspacer{} อิเมหิ โข ภิกฺขเว ตีหงฺเคหิ สมนฺนาคตสฺส ภิกฺขุโน อากงฺขมาโน สํโฆ ตชฺชนียกมฺมํ กเรยฺย. (3)}

\prose{ติณฺณํ ภิกฺขเว ภิกฺขูนํ อากงฺขมาโน สํโฆ ตชฺชนียกมฺมํ กเรยฺย.\csromanspacer{} เอโก ภณฺฑนการโก โหติ กลหการโก วิวาทการโก ภสฺสการโก สํเฆ อธิกรณการโก, เอโก พาโล โหติ อพฺยตฺโต อาปตฺติพหุโล อนปทาโน, เอโก คิหิสํสฏฺโฐ วิหรติ อนนุโลมิเกหิ คิหิสํสคฺเคหิ.\csromanspacer{} อิเมสํ โข ภิกฺขเว ติณฺณํ ภิกฺขูนํ อากงฺขมาโน สํโฆ ตชฺชนียกมฺมํ กเรยฺย. (4)}

\prose{\csromanfolio{10}อปเรสมฺปิ ภิกฺขเว ติณฺณํ ภิกฺขูนํ อากงฺขมาโน สํโฆ ตชฺชนียกมฺมํ กเรยฺย.\csromanspacer{} เอโก อธิสีเล สีลวิปนฺโน โหติ, เอโก อชฺฌาจาเร อาจารวิปนฺโน โหติ, เอโก อติทิฏฺฐิยา ทิฏฺฐิวิปนฺโน โหติ.\csromanspacer{} อิเมสํ โข ภิกฺขเว ติณฺณํ ภิกฺขูนํ อากงฺขมาโน สํโฆ ตชฺชนียกมฺมํ กเรยฺย. (5)}

\prose{อปเรสมฺปิ ภิกฺขเว ติณฺณํ ภิกฺขูนํ อากงฺขมาโน สํโฆ ตชฺชนียกมฺมํ กเรยฺย.\csromanspacer{} เอโก พุทฺธสฺส อวณฺณํ ภาสติ, เอโก ธมฺมสฺส อวณฺณํ ภาสติ, เอโก สํฆสฺส อวณฺณํ ภาสติ.\csromanspacer{} อิเมสํ โข ภิกฺขเว ติณฺณํ ภิกฺขูนํ อากงฺขมาโน สํโฆ ตชฺชนียกมฺมํ กเรยฺย. (6)}

\nitthitam{อากงฺขมานฉกฺกํ นิฏฺฐิตํ.}
\csromanmatikaanchor{mtk.6}
\csromantocmarkref{subsection}{อฏฺฐารสวตฺต}{mtk.6}
\bookmark[dest={mtk.6},level=2]{อฏฺฐารสวตฺต}
\csromanheader{2}{\textbf{อฏฺฐารสวตฺต}}
\proseitem{7}{ตชฺชนียกมฺมกเตน ภิกฺขเว ภิกฺขุนา สมฺมา วตฺติตพฺพํ.\csromanspacer{} ตตฺรายํ สมฺมาวตฺตนา\csromandash{}น อุปสมฺปาเทตพฺพํ, น นิสฺสโย ทาตพฺโพ, น สามเณโร อุปฏฺฐาเปตพฺโพ, น ภิกฺขุโนวาทกสมฺมุติ\footnote{สมฺมติ (สฺยา)} สาทิตพฺพา, สมฺมเตนปิ ภิกฺขุนิโย น โอวทิตพฺพา, ยาย อาปตฺติยา สํเฆน ตชฺชนียกมฺมํ กตํ โหติ, สา อาปตฺติ น อาปชฺชิตพฺพา, อญฺญา วา ตาทิสิกา, ตโต วา ปาปิฏฺฐตรา, กมฺมํ น ครหิตพฺพํ, กมฺมิกา น ครหิตพฺพา, น ปกตตฺตสฺส ภิกฺขุโน อุโปสโถ ฐเปตพฺโพ, น ปวารณา ฐเปตพฺพา, น สวจนียํ กาตพฺพํ, น อนุวาโท ปฏฺฐเปตพฺโพ, น โอกาโส กาเรตพฺโพ, น โจเทตพฺโพ, น สาเรตพฺโพ, น ภิกฺขูหิ\footnote{น ภิกฺขู ภิกฺขูหิ (สฺยา)} สมฺปโยเชตพฺพนฺติ.}

\nitthitam{ตชฺชนียกมฺเม อฏฺฐารสวตฺตํ นิฏฺฐิตํ.}
\csromanmatikaanchor{mtk.7}
\csromantocmarkref{subsection}{นปฺปฏิปฺปสฺสมฺเภตพฺพ-อฏฺฐารสก}{mtk.7}
\bookmark[dest={mtk.7},level=2]{นปฺปฏิปฺปสฺสมฺเภตพฺพ-อฏฺฐารสก}
\csromanheader{2}{\textbf{นปฺปฏิปฺปสฺสมฺเภตพฺพ-อฏฺฐารสก}}
\proseitem{8}{อถ โข สํโฆ ปณฺฑุกโลหิตกานํ ภิกฺขูนํ ตชฺชนียกมฺมํ อกาสิ.\csromanspacer{} เต สํเฆน ตชฺชนียกมฺมกตา สมฺมา วตฺตนฺติ, โลมํ \csromanfolio{11}ปาเตนฺติ, เนตฺถารํ วตฺตนฺติ, ภิกฺขู อุปสงฺกมิตฺวา เอวํ วทนฺติ “มยํ อาวุโส สํเฆน ตชฺชนียกมฺมกตา สมฺมา วตฺตาม, โลมํ ปาเตม, เนตฺถารํ วตฺตาม, กถํ นุ โข อเมฺหหิ ปฏิปชฺชิตพฺพนฺ”ติ.\csromanspacer{} ภิกฺขู ภควโต เอตมตฺถํ อาโรเจสุํ ฯเปฯ.\csromanspacer{} เตน หิ ภิกฺขเว สํโฆ ปณฺฑุกโลหิตกานํ ภิกฺขูนํ ตชฺชนียกมฺมํ ปฏิปฺปสฺสมฺเภตุ.}

\prose{ปญฺจหิ ภิกฺขเว องฺเคหิ สมนฺนาคตสฺส ภิกฺขุโน ตชฺชนียกมฺมํ นปฺปฏิปฺปสฺสมฺเภตพฺพํ.\csromanspacer{} อุปสมฺปาเทติ, นิสฺสยํ เทติ, สามเณรํ อุปฏฺฐาเปติ, ภิกฺขูโนวาทกสมฺมุติํ สาทิยติ, สมฺมโตปิ ภิกฺขุนิโย โอวทติ.\csromanspacer{} อิเมหิ โข ภิกฺขเว ปญฺจหงฺเคหิ สมนฺนาคตสฺส ภิกฺขุโน ตชฺชนียกมฺมํ นปฺปฏิปฺปสฺสมฺเภตพฺพํ.}

\prose{\csromansymbolfootnote{*}{วิ 5. 317 ปิฏฺเฐปิ.}อปเรหิปิ ภิกฺขเว ปญฺจหงฺเคหิ สมนฺนาคตสฺส ภิกฺขุโน ตชฺชนียกมฺมํ นปฺปฏิปฺปสฺสมฺเภตพฺพํ.\csromanspacer{} ยาย อาปตฺติยา สํเฆน ตชฺชนียกมฺมํ กตํ โหติ, ตํ อาปตฺติํ อาปชฺชติ, อญฺญํ วา ตาทิสิกํ, ตโต วา ปาปิฏฺฐตรํ, กมฺมํ ครหติ, กมฺมิเก ครหติ.\csromanspacer{} อิเมหิ โข ภิกฺขเว ปญฺจหงฺเคหิ สมนฺนาคตสฺส ภิกฺขุโน ตชฺชนียกมฺมํ นปฺปฏิปฺปสฺสมฺเภตพฺพํ.}

\prose{อฏฺฐหิ ภิกฺขเว องฺเคหิ สมนฺนาคตสฺส ภิกฺขุโน ตชฺชนียกมฺมํ นปฺปฏิปฺปสฺสมฺเภตพฺพํ.\csromanspacer{} ปกตตฺตสฺส ภิกฺขุโน อุโปสถํ ฐเปติ, ปวารณํ ฐเปติ, สวจนียํ กโรติ, อนุวาทํ ปฏฺฐเปติ, โอกาสํ กาเรติ, โจเทติ, สาเรติ, ภิกฺขูหิ สมฺปโยเชติ\footnote{ภิกฺขู ภิกฺขูหิ สมฺปโยเชติ (สฺยา)}.\csromanspacer{} อิเมหิ โข ภิกฺขเว อฏฺฐหงฺเคหิ สมนฺนาคตสฺส ภิกฺขุโน ตชฺชนียกมฺมํ นปฺปฏิปฺปสฺสมฺเภตพฺพํ.}

\nitthitam{นปฺปฏิปฺปสฺสมฺเภตพฺพ-อฏฺฐารสกํ นิฏฺฐิตํ.}
\csromanmatikaanchor{mtk.8}
\csromantocmarkref{subsection}{ปฏิปฺปสฺสมฺเภตพฺพ-อฏฺฐารสก}{mtk.8}
\bookmark[dest={mtk.8},level=2]{ปฏิปฺปสฺสมฺเภตพฺพ-อฏฺฐารสก}
\csromanheader{2}{\textbf{ปฏิปฺปสฺสมฺเภตพฺพ-อฏฺฐารสก}}
\proseitem{9}{ปญฺจหิ ภิกฺขเว องฺเคหิ สมนฺนาคตสฺส ภิกฺขุโน ตชฺชนียกมฺมํ ปฏิปฺปสฺสมฺเภตพฺพํ.\csromanspacer{} น อุปสมฺปาเทติ, น นิสฺสยํ เทติ, น สามเณรํ อุปฏฺฐาเปติ, น ภิกฺขุโนวาทกสมฺมุติํ สาทิยติ, สมฺมโตปิ ภิกฺขุนิโย น โอวทติ.\csromanspacer{} อิเมหิ โข ภิกฺขเว ปญฺจหงฺเคหิ สมนฺนาคตสฺส ภิกฺขุโน ตชฺชนียกมฺมํ ปฏิปฺปสฺสมฺเภตพฺพํ.}

\prose{\csromanfolio{12}อปเรหิปิ ภิกฺขเว ปญฺจหงฺเคหิ สมนฺนาคตสฺส ภิกฺขุโน ตชฺชนียกมฺมํ ปฏิปฺปสฺสมฺเภตพฺพํ.\csromanspacer{} ยาย อาปตฺติยา สํเฆน ตชฺชนียกมฺมํ กตํ โหติ, ตํ อาปตฺติํ น อาปชฺชติ, อญฺญํ วา ตาทิสิกํ, ตโต วา ปาปิฏฺฐตรํ, กมฺมํ น ครหติ, กมฺมิเก น ครหติ.\csromanspacer{} อิเมหิ โข ภิกฺขเว ปญฺจหงฺเคหิ สมนฺนาคตสฺส ภิกฺขุโน ตชฺชนียกมฺมํ ปฏิปฺปสฺสมฺเภตพฺพํ.}

\prose{อฏฺฐหิ ภิกฺขเว องฺเคหิ สมนฺนาคตสฺส ภิกฺขุโน ตชฺชนียกมฺมํ ปฏิปฺปสฺสมฺเภตพฺพํ.\csromanspacer{} น ปกตตฺตสฺส ภิกฺขุโน อุโปสถํ ฐเปติ, น ปวารณํ ฐเปติ, น สวจนียํ กโรติ, น อนุวาทํ ปฏฺฐเปติ, น โอกาสํ กาเรติ, น โจเทติ, น สาเรติ, น ภิกฺขูหิ สมฺปโยเชติ.\csromanspacer{} อิเมหิ โข ภิกฺขเว อฏฺฐงฺเคหิ สมนฺนาคตสฺส ภิกฺขุโน ตชฺชนียกมฺมํ ปฏิปฺปสฺสมฺเภตพฺพํ.}

\nitthitam{ปฏิปฺปสฺสมฺเภตพฺพ-อฏฺฐารสกํ นิฏฺฐิตํ.}
\proseitem{10}{เอวญฺจ ปน ภิกฺขเว ปฏิปฺปสฺสมฺเภตพฺพํ, เตหิ ภิกฺขเว ปณฺฑุกโลหิตเกหิ ภิกฺขูหิ สํฆํ อุปสงฺกมิตฺวา เอกํสํ อุตฺตราสงฺคํ กริตฺวา วุฑฺฒานํ ภิกฺขูนํ ปาเท วนฺทิตฺวา อุกฺกุฏิกํ นิสีทิตฺวา อญฺชลิํ ปคฺคเหตฺวา เอวมสฺส วจนีโย “มยํ ภนฺเต สํเฆน ตชฺชนียกมฺมกตา สมฺมา วตฺตาม, โลมํ ปาเตม, เนตฺถารํ วตฺตาม, ตชฺชนียสฺส กมฺมสฺส ปฏิปฺปสฺสทฺธิํ ยาจามา”ติ.\csromanspacer{} ทุติยมฺปิ ยาจิตพฺพา.\csromanspacer{} ตติยมฺปิ ยาจิตพฺพา.\csromanspacer{} พฺยตฺเตน ภิกฺขุนา ปฏิพเลน สํโฆ ญาเปตพฺโพ\csromandash{}}

\prose{“สุณาตุ เม ภนฺเต สํโฆ, อิเม ปณฺฑุกโลหิตกา ภิกฺขู สํเฆน ตชฺชนียกมฺมกตา สมฺมา วตฺตนฺติ, โลมํ ปาเตนฺติ, เนตฺถารํ วตฺตนฺติ, ตชฺชนียสฺส กมฺมสฺส ปฏิปฺปสฺสทฺธิํ ยาจนฺติ, ยทิ สํฆสฺส ปตฺตกลฺลํ, สํโฆ ปณฺฑุกโลหิตกานํ ภิกฺขูนํ ตชฺชนียกมฺมํ ปฏิปฺปสฺสมฺเภยฺย, เอสา ญตฺติ.}

\prose{สุณาตุ เม ภนฺเต สํโฆ, อิเม ปณฺฑุกโลหิตกา ภิกฺขู สํเฆน ตชฺชนียกมฺมกตา สมฺมา วตฺตนฺติ, โลมํ ปาเตนฺติ, เนตฺถารํ วตฺตนฺติ, ตชฺชนียสฺส กมฺมสฺส ปฏิปฺปสฺสทฺธิํ ยาจนฺติ, สํโฆ ปณฺฑุกโลหิตกานํ ภิกฺขูนํ ตชฺชนียกมฺมํ ปฏิปฺปสฺสมฺเภติ, ยสฺสายสฺมโต ขมติ ปณฺฑุกโลหิตกานํ ภิกฺขูนํ ตชฺชนียสฺส กมฺมสฺส ปฏิปฺปสฺสทฺธิ, โส ตุณฺหสฺส, ยสฺส นกฺขมติ, โส ภาเสยฺย.}

\prose{\csromanfolio{13}ทุติยมฺปิ เอตมตฺถํ วทามิ.\csromanspacer{} สุณาตุ เม ภนฺเต สํโฆ, อิเม ปณฺฑุกโลหิตกา ภิกฺขู สํเฆน ตชฺชนียกมฺมกตา สมฺมา วตฺตนฺติ, โลมํ ปาเตนฺติ, เนตฺถารํ วตฺตนฺติ, ตชฺชนียสฺส กมฺมสฺส ปฏิปฺปสฺสทฺธิํ ยาจนฺติ, สํโฆ ปณฺฑุกโลหิตกานํ ภิกฺขูนํ ตชฺชนียกมฺมํ ปฏิปฺปสฺสมฺเภติ, ยสฺสายสฺมโต ขมติ ปณฺฑุกโลหิตกานํ ภิกฺขูนํ ตชฺชนียสฺส กมฺมสฺส ปฏิปฺปสฺสทฺธิ, โส ตุณฺหสฺส, ยสฺส นกฺขมติ, โส ภาเสยฺย.}

\prose{ตติยมฺปิ เอตมตฺถํ วทามิ.\csromanspacer{} สุณาตุ เม ภนฺเต สํโฆ, อิเม ปณฺฑุกโลหิตกา ภิกฺขู สํเฆน ตชฺชนียกมฺมกตา สมฺมา วตฺตนฺติ, โลมํ ปาเตนฺติ, เนตฺถารํ วตฺตนฺติ, ตชฺชนียสฺส กมฺมสฺส ปฏิปฺปสฺสทฺธิํ ยาจนฺติ, สํโฆ ปณฺฑุกโลหิตกานํ ภิกฺขูนํ ตชฺชนียกมฺมํ ปฏิปฺปสฺสมฺเภติ, ยสฺสายสฺมโต ขมติ ปณฺฑุกโลหิตกานํ ภิกฺขูนํ ตชฺชนียสฺส กมฺมสฺส ปฏิปฺปสฺสทฺธิ, โส ตุณฺหสฺส, ยสฺส นกฺขมติ, โส ภาเสยฺย.}

\prose{ปฏิปฺปสฺสทฺธํ สํเฆน ปณฺฑุกโลหิตกานํ ภิกฺขูนํ ตชฺชนียกมฺมํ, ขมติ สํฆสฺส, ตสฺมา ตุณฺหี, เอวเมตํ ธารยามี”ติ.}

\nitthitam{ตชฺชนียกมฺมํ นิฏฺฐิตํ ปฐมํ.}
\csromanmatikaanchor{mtk.9}
\csromantocmarkref{section}{2. นิยสฺสกมฺม}{mtk.9}
\bookmark[dest={mtk.9},level=1]{2. นิยสฺสกมฺม}
\csromanheader{1}{2. นิยสฺสกมฺม}
\proseitem{11}{เตน โข ปน สมเยน อายสฺมา เสยฺยสโก พาโล โหติ อพฺยตฺโต อาปตฺติพหุโล อนปทาโน, คิหิสํสฏฺโฐ วิหรติ อนนุโลมิเกหิ คิหิสํสคฺเคหิ, อปิสฺสุ ภิกฺขู ปกตา\footnote{ปกตตฺตา (สี, สฺยา)} ปริวาสํ เทนฺตา มูลาย ปฏิกสฺสนฺตา มานตฺตํ เทนฺตา อพฺเภนฺตา.\csromanspacer{} เย เต ภิกฺขู อปฺปิจฺฉา ฯเปฯ เต อุชฺฌายนฺติ ขิยฺยนฺติ วิปาเจนฺติ “กถํ หิ นาม อายสฺมา เสยฺยสโก พาโล ภวิสฺสติ อพฺยตฺโต อาปตฺติพหุโล อนปทาโน, คิหิสํสฏฺโฐ วิหริสฺสติ อนนุโลมิเกหิ คิหิสํสคฺเคหิ, อปิสฺสุ ภิกฺขู ปกตา ปริวาสํ เทนฺตา มูลาย ปฏิกสฺสนฺตา มานตฺตํ เทนฺตา อพฺเภนฺตา”ติ.\csromanspacer{} อถ โข เต ภิกฺขู ภควโต เอตมตฺถํ อาโรเจสุํ.}

\prose{\csromanfolio{14}อถ โข ภควา เอตสฺมิํ นิทาเน เอตสฺมิํ ปกรเณ ภิกฺขุสํฆํ สนฺนิปาตาเปตฺวา ภิกฺขู ปฏิปุจฺฉิ “สจฺจํ กิร ภิกฺขเว เสยฺยสโก ภิกฺขุ พาโล อพฺยตฺโต อาปตฺติพหุโล อนปทาโน, คิหิสํสฏฺโฐ วิหรติ อนนุโลมิเกหิ คิหิสํสคฺเคหิ, อปิสฺสุ ภิกฺขู ปกตา ปริวาสํ เทนฺตา มูลาย ปฏิกสฺสนฺตา มานตฺตํ เทนฺตา อพฺเภนฺตา”ติ.\csromanspacer{} สจฺจํ ภควาติ.\csromanspacer{} วิครหิ พุทฺโธ ภควา อนนุจฺฉวิกํ ภิกฺขเว ตสฺส โมฆปุริสสฺส อนนุโลมิกํ อปฺปติรูปํ อสฺสามณกํ อกปฺปิยํ อกรณียํ.\csromanspacer{} กถํ หิ นาม โส ภิกฺขเว โมฆปุริโส พาโล ภวิสฺสติ อพฺยตฺโต อาปตฺติพหุโล อนปทาโน, คิหิสํสฏฺโฐ วิหริสฺสติ อนนุโลมิเกหิ คิหิสํสคฺเคหิ, อปิสฺสุ ภิกฺขู ปกตา ปริวาสํ เทนฺตา มูลาย ปฏิกสฺสนฺตา มานตฺตํ เทนฺตา อพฺเภนฺตา.\csromanspacer{} เนตํ ภิกฺขเว อปฺปสนฺนานํ วา ปสาทาย ฯเปฯ วิครหิตฺวา ฯเปฯ ธมฺมิํ กถํ กตฺวา ภิกฺขู อามนฺเตสิ\csromandash{}เตน หิ ภิกฺขเว สํโฆ เสยฺยสกสฺส ภิกฺขุโน นิยสฺสกมฺมํ\footnote{นิยสกมฺมํ (ก)} กโรตุ “นิสฺสาย เต วตฺถพฺพนฺ”ติ.\csromanspacer{} เอวญฺจ ปน ภิกฺขเว กาตพฺพํ, ปฐมํ เสยฺยสโก ภิกฺขุ โจเทตพฺโพ, โจเทตฺวา สาเรตพฺโพ, สาเรตฺวา อาปตฺติํ อาโรเปตพฺโพ\footnote{อาปตฺติ อาโรเปตพฺพา (สี, สฺยา)}, อาปตฺติํ อาโรเปตฺวา พฺยตฺเตน ภิกฺขุนา ปฏิพเลน สํโฆ ญาเปตพฺโพ\csromandash{}}

\proseitem{12}{“สุณาตุ เม ภนฺเต สํโฆ, อยํ เสยฺยสโก ภิกฺขุ, พาโล อพฺยตฺโต อาปตฺติพหุโล อนปทาโน, คิหิสํสฏฺโฐ วิหรติ อนนุโลมิเกหิ คิหิสํสคฺเคหิ, อปิสฺสุ ภิกฺขู ปกตา ปริวาสํ เทนฺตา มูลาย ปฏิกสฺสนฺตา มานตฺตํ เทนฺตา อพฺเภนฺตา, ยทิ สํฆสฺส ปตฺตกลฺลํ, สํโฆ เสยฺยสกสฺส ภิกฺขุโน นิยสฺสกมฺมํ กเรยฺย ‘นิสฺสาย เต วตฺถพฺพนฺ’ติ, เอสา ญตฺติ.}

\prose{สุณาตุ เม ภนฺเต สํโฆ, อยํ เสยฺยสโก ภิกฺขุ พาโล อพฺยตฺโต อาปตฺติพหุโล อนปทาโน, คิหิสํสฏฺโฐ วิหรติ อนนุโลมิเกหิ คิหิสํสคฺเคหิ, อปิสฺสุ ภิกฺขู ปกตา ปริวาสํ เทนฺตา มูลาย ปฏิกสฺสนฺตา มานตฺตํ เทนฺตา อพฺเภนฺตา, สํโฆ เสยฺยสกสฺส ภิกฺขุโน นิยสฺสกมฺมํ กโรติ ‘นิสฺสาย เต วตฺถพฺพนฺ’ติ, ยสฺสายสฺมโต ขมติ เสยฺยสกสฺส ภิกฺขุโน นิยสฺสสฺส กมฺมสฺส กรณํ ‘นิสฺสาย เต วตฺถพฺพนฺ’ติ, โส ตุณฺหสฺส, ยสฺส นกฺขมติ, โส ภาเสยฺย.}

\prose{\csromanfolio{15}ทุติยมฺปิ เอตมตฺถํ วทามิ ฯเปฯ.\csromanspacer{} ตติยมฺปิ เอตมตฺถํ วทามิ.\csromanspacer{} สุณาตุ เม ภนฺเต สํโฆ, อยํ เสยฺยสโก ภิกฺขุ พาโล อพฺยตฺโต อาปตฺติพหุโล อนปทาโน, คิหิสํสฏฺโฐ วิหรติ อนนุโลมิเกหิ คิหิสํสคฺเคหิ, อปิสฺสุ ภิกฺขู ปกตา ปริวาสํ เทนฺตา มูลาย ปฏิกสฺสนฺตา มานตฺตํ เทนฺตา อพฺเภนฺตา, สํโฆ เสยฺยสกสฺส ภิกฺขุโน นิยสฺสกมฺมํ กโรติ ‘นิสฺสาย เต วตฺถพฺพนฺ’ติ, ยสฺสายสฺมโต ขมติ เสยฺยสกสฺส ภิกฺขุโน นิยสฺสสฺส กมฺมสฺส กรณํ ‘นิสฺสาย เต วตฺถพฺพนฺ’ติ, โส ตุณฺหสฺส, ยสฺส นกฺขมติ, โส ภาเสยฺย.}

\prose{กตํ สํเฆน เสยฺยสกสฺส ภิกฺขุโน นิยสฺสกมฺมํ ‘นิสฺสาย เต วตฺถพฺพนฺ’ติ, ขมติ สํฆสฺส, ตสฺมา ตุณฺหี, เอวเมตํ ธารยามี”ติ.}

\csromanmatikaanchor{mtk.10}
\csromantocmarkref{subsection}{อธมฺมกมฺมทฺวาทสก}{mtk.10}
\bookmark[dest={mtk.10},level=2]{อธมฺมกมฺมทฺวาทสก}
\csromanheader{2}{\textbf{อธมฺมกมฺมทฺวาทสก}}
\proseitem{13}{ตีหิ ภิกฺขเว องฺเคหิ สมนฺนาคตํ นิยสฺสกมฺมํ อธมฺมกมฺมญฺจ โหติ อวินยกมฺมญฺจ ทุวูปสนฺตญฺจ.\csromanspacer{} อสมฺมุขา กตํ โหติ, อปฺปฏิปุจฺฉา กตํ โหติ, อปฺปฏิญฺญาย กตํ โหติ.\csromanspacer{} อิเมหิ โข ภิกฺขเว ตีหงฺเคหิ สมนฺนาคตํ นิยสฺสกมฺมํ อธมฺมกมฺมญฺจ โหติ อวินยกมฺมญฺจ ทุวูปสนฺตญฺจ. (1)}

\prose{อปเรหิปิ ภิกฺขเว ตีหงฺเคหิ สมนฺนาคตํ นิยสฺสกมฺมํ อธมฺมกมฺมญฺจ โหติ อวินยกมฺมญฺจ ทุวูปสนฺตญฺจ.\csromanspacer{} อนาปตฺติยา กตํ โหติ, อเทสนาคามินิยา อาปตฺติยา กตํ โหติ, เทสิตาย อาปตฺติยา กตํ โหติ.\csromanspacer{} อิเมหิ โข ภิกฺขเว ตีหงฺเคหิ สมนฺนาคตํ นิยสฺสกมฺมํ อธมฺมกมฺมญฺจ โหติ อวินยกมฺมญฺจ ทุวูปสนฺตญฺจ. (2)}

\prose{อปเรหิปิ ภิกฺขเว ตีหงฺเคหิ สมนฺนาคตํ นิยสฺสกมฺมํ อธมฺมกมฺมญฺจ โหติ อวินยกมฺมญฺจ ทุวูปสนฺตญฺจ.\csromanspacer{} อโจเทตฺวา กตํ โหติ, อสาเรตฺวา กตํ โหติ, อาปตฺติํ อนาโรเปตฺวา กตํ โหติ.\csromanspacer{} อิเมหิ โข ภิกฺขเว ตีหงฺเคหิ สมนฺนาคตํ นิยสฺสกมฺมํ อธมฺมกมฺมญฺจ โหติ อวินยกมฺมญฺจ ทุวูปสนฺตญฺจ. (3)}

\prose{อปเรหิปิ ภิกฺขเว ฯเปฯ.\csromanspacer{} อสมฺมุขา กตํ โหติ, อธมฺเมน กตํ โหติ, วคฺเคน กตํ โหติ.\csromanspacer{} อิเมหิ โข ภิกฺขเว ฯเปฯ. (4)}

\prose{\csromanfolio{16}อปเรหิปิ ภิกฺขเว ฯเปฯ.\csromanspacer{} อปฺปฏิปุจฺฉา กตํ โหติ, อธมฺเมน กตํ โหติ, วคฺเคน กตํ โหติ.\csromanspacer{} อิเมหิ โข ภิกฺขเว ฯเปฯ. (5)}

\prose{อปเรหิปิ ภิกฺขเว ฯเปฯ.\csromanspacer{} อปฺปฏิญฺญาย กตํ โหติ, อธมฺเมน กตํ โหติ, วคฺเคน กตํ โหติ.\csromanspacer{} อิเมหิ โข ภิกฺขเว ฯเปฯ. (6)}

\prose{อปเรหิปิ ภิกฺขเว ฯเปฯ.\csromanspacer{} อนาปตฺติยา กตํ โหติ, อธมฺเมน กตํ โหติ, วคฺเคน กตํ โหติ.\csromanspacer{} อิเมหิ โข ภิกฺขเว ฯเปฯ. (7)}

\prose{อปเรหิปิ ภิกฺขเว ฯเปฯ.\csromanspacer{} อเทสนาคามินิยา อาปตฺติยา กตํ โหติ, อธมฺเมน กตํ โหติ, วคฺเคน กตํ โหติ.\csromanspacer{} อิเมหิ โข ภิกฺขเว ฯเปฯ. (8)}

\prose{อปเรหิปิ ภิกฺขเว ฯเปฯ.\csromanspacer{} เทสิตาย อาปตฺติยา กตํ โหติ, อธมฺเมน กตํ โหติ, วคฺเคน กตํ โหติ.\csromanspacer{} อิเมหิ โข ภิกฺขเว ฯเปฯ. (9)}

\prose{อปเรหิปิ ภิกฺขเว ฯเปฯ.\csromanspacer{} อโจเทตฺวา กตํ โหติ, อธมฺเมน กตํ โหติ, วคฺเคน กตํ โหติ.\csromanspacer{} อิเมหิ โข ภิกฺขเว ฯเปฯ. (10)}

\prose{อปเรหิปิ ภิกฺขเว ฯเปฯ.\csromanspacer{} อสาเรตฺวา กตํ โหติ, อธมฺเมน กตํ โหติ, วคฺเคน กตํ โหติ.\csromanspacer{} อิเมหิ โข ภิกฺขเว ฯเปฯ. (11)}

\prose{อปเรหิปิ ภิกฺขเว ตีหงฺเคหิ สมนฺนาคตํ นิยสฺสกมฺมํ อธมฺมกมฺมญฺจ โหติ อวินยกมฺมญฺจ ทุวูปสนฺตญฺจ.\csromanspacer{} อาปตฺติํ อนาโรเปตฺวา กตํ โหติ, อธมฺเมน กตํ โหติ, วคฺเคน กตํ โหติ.\csromanspacer{} อิเมหิ โข ภิกฺขเว ตีหงฺเคหิ สมนฺนาคตํ นิยสฺสกมฺมํ อธมฺมกมฺมญฺจ โหติ อวินยกมฺมญฺจ ทุวูปสนฺตญฺจ. (12)}

\vspace{\tipitakaparskip}
\csromancenter{อธมฺมกมฺมทฺวาทสกํ ทิฏฺฐิตํ.}
\csromanmatikaanchor{mtk.11}
\csromantocmarkref{subsection}{ธมฺมกมฺมทฺวาทสก}{mtk.11}
\bookmark[dest={mtk.11},level=2]{ธมฺมกมฺมทฺวาทสก}
\csromanheader{2}{\textbf{ธมฺมกมฺมทฺวาทสก}}
\proseitem{14}{ตีหิ ภิกฺขเว องฺเคหิ สมนฺนาคตํ นิยสฺสกมฺมํ ธมฺมกมฺมญฺจ โหติ วินยกมฺมญฺจ สุวูปสนฺตญฺจ.\csromanspacer{} สมฺมุขา กตํ โหติ, ปฏิปุจฺฉา กตํ โหติ, ปฏิญฺญาย กตํ โหติ.\csromanspacer{} อิเมหิ โข ภิกฺขเว ตีหงฺเคหิ สมนฺนาคตํ นิยสฺสกมฺมํ ธมฺมกมฺมญฺจ โหติ วินยกมฺมญฺจ สุวูปสนฺตญฺจ. (1)}

\prose{\csromanfolio{17}อปเรหิปิ ภิกฺขเว ตีหงฺเคหิ สมนฺนาคตํ นิยสฺสกมฺมํ ธมฺมกมฺมญฺจ โหติ วินยกมฺมญฺจ สุวูปสนฺตญฺจ.\csromanspacer{} อาปตฺติยา กตํ โหติ, เทสนาคามินิยา อาปตฺติยา กตํ โหติ, อเทสิตาย อาปตฺติยา กตํ โหติ.\csromanspacer{} อิเมหิ โข ภิกฺขเว ตีหงฺเคหิ สมนฺนาคตํ นิยสฺสกมฺมํ ธมฺมกมฺมญฺจ โหติ วินยกมฺมญฺจ สุวูปสนฺตญฺจ. (2)}

\prose{อปเรหิปิ ภิกฺขเว ตีหงฺเคหิ สมนฺนาคตํ นิยสฺสกมฺมํ ธมฺมกมฺมญฺจ โหติ วินยกมฺมญฺจ สุวูปสนฺตญฺจ.\csromanspacer{} โจเทตฺวา กตํ โหติ, สาเรตฺวา กตํ โหติ, อาปตฺติํ อาโรเปตฺวา กตํ โหติ.\csromanspacer{} อิเมหิ โข ภิกฺขเว ตีหงฺเคหิ สมนฺนาคตํ นิยสฺสกมฺมํ ธมฺมกมฺมญฺจ โหติ วินยกมฺมญฺจ สุวูปสนฺตญฺจ. (3)}

\prose{อปเรหิปิ ภิกฺขเว ฯเปฯ.\csromanspacer{} สมฺมุขา กตํ โหติ, ธมฺเมน กตํ โหติ, สมคฺเคน กตํ โหติ.\csromanspacer{} อิเมหิ โข ภิกฺขเว ฯเปฯ. (4)}

\prose{อปเรหิปิ ภิกฺขเว ฯเปฯ.\csromanspacer{} ปฏิปุจฺฉา กตํ โหติ, ธมฺเมน กตํ โหติ, สมคฺเคน กตํ โหติ.\csromanspacer{} อิเมหิ โข ภิกฺขเว ฯเปฯ. (5)}

\prose{อปเรหิปิ ภิกฺขเว ฯเปฯ.\csromanspacer{} ปฏิญฺญายกตํ โหติ, ธมฺเมน กตํ โหติ, สมคฺเคน กตํ โหติ.\csromanspacer{} อิเมหิ โข ภิกฺขเว ฯเปฯ. (6)}

\prose{อปเรหิปิ ภิกฺขเว ฯเปฯ.\csromanspacer{} อาปตฺติยา กตํ โหติ, ธมฺเมน กตํ โหติ, สมคฺเคน กตํ โหติ.\csromanspacer{} อิเมหิ โข ภิกฺขเว ฯเปฯ. (7)}

\prose{อปเรหิปิ ภิกฺขเว ฯเปฯ.\csromanspacer{} เทสนาคามินิยา อาปตฺติยา กตํ โหติ, ธมฺเมน กตํ โหติ, สมคฺเคน กตํ โหติ.\csromanspacer{} อิเมหิ โข ภิกฺขเว ฯเปฯ. (8)}

\prose{อปเรหิปิ ภิกฺขเว ฯเปฯ.\csromanspacer{} อเทสิตาย อาปตฺติยา กตํ โหติ, ธมฺเมน กตํ โหติ, สมคฺเคน กตํ โหติ.\csromanspacer{} อิเมหิ โข ภิกฺขเว ฯเปฯ. (9)}

\prose{อปเรหิปิ ภิกฺขเว ฯเปฯ.\csromanspacer{} โจเทตฺวา กตํ โหติ, ธมฺเมน กตํ โหติ, สมคฺเคน กตํ โหติ.\csromanspacer{} อิเมหิ โข ภิกฺขเว ฯเปฯ. (10)}

\prose{อปเรหิปิ ภิกฺขเว ฯเปฯ.\csromanspacer{} สาเรตฺวา กตํ โหติ, ธมฺเมน กตํ โหติ, สมคฺเคน กตํ โหติ.\csromanspacer{} อิเมหิ โข ภิกฺขเว ฯเปฯ. (11)}

\prose{\csromanfolio{18}อปเรหิปิ ภิกฺขเว ตีหงฺเคหิ สมนฺนาคตํ นิยสฺสกมฺมํ ธมฺมกมฺมญฺจ โหติ วินยกมฺมญฺจ สุวูปสนฺตญฺจ.\csromanspacer{} อาปตฺติํ อาโรเปตฺวา กตํ โหติ, ธมฺเมน กตํ โหติ, สมคฺเคน กตํ โหติ.\csromanspacer{} อิเมหิ โข ภิกฺขเว ตีหงฺเคหิ สมนฺนาคตํ นิยสฺสกมฺมํ ธมฺมกมฺมญฺจ โหติ วินยกมฺมญฺจ สุวูปสนฺตญฺจ. (12)}

\nitthitam{ธมฺมกมฺมทฺวาทสกํ นิฏฺฐิตํ.}
\csromanmatikaanchor{mtk.12}
\csromantocmarkref{subsection}{อากงฺขมานฉกฺก}{mtk.12}
\bookmark[dest={mtk.12},level=2]{อากงฺขมานฉกฺก}
\csromanheader{2}{\textbf{อากงฺขมานฉกฺก}}
\proseitem{15}{\csromansymbolfootnote{*}{วิ 5. 219 ปิฏฺเฐปิ.}ตีหิ ภิกฺขเว องฺเคหิ สมนฺนาคตสฺส ภิกฺขุโน อากงฺขมาโน สํโฆ นิยสฺสกมฺมํ กเรยฺย.\csromanspacer{} ภณฺฑนการโก โหติ กลหการโก วิวาทการโก ภสฺสการโก สํเฆ อธิกรณการโก, พาโล โหติ อพฺยตฺโต อาปตฺติพหุโล อนปทาโน, คิหิสํสฏฺโฐ วิหรติ อนนุโลมิเกหิ คิหิสํสคฺเคหิ.\csromanspacer{} อิเมหิ โข ภิกฺขเว ตีหงฺเคหิ สมนฺนาคตสฺส ภิกฺขุโน อากงฺขมาโน สํโฆ นิยสฺสกมฺมํ กเรยฺย. (1)}

\prose{อปเรหิปิ ภิกฺขเว ตีหงฺเคหิ สมนฺนาคตสฺส ภิกฺขุโน อากงฺขมาโน สํโฆ นิยสฺสกมฺมํ กเรยฺย.\csromanspacer{} อธิสีเล สีลวิปนฺโน โหติ, อชฺฌาจาเร อาจารวิปนฺโน โหติ, อติทิฏฺฐิยา ทิฏฺฐิวิปนฺโน โหติ.\csromanspacer{} อิเมหิ โข ภิกฺขเว ตีหงฺเคหิ สมนฺนาคตสฺส ภิกฺขุโน อากงฺขมาโน สํโฆ นิยสฺสกมฺมํ กเรยฺย. (2)}

\prose{อปเรหิปิ ภิกฺขเว ตีหงฺเคหิ สมนฺนาคตสฺส ภิกฺขุโน อากงฺขมาโน สํโฆ นิยสฺสกมฺมํ กเรยฺย.\csromanspacer{} พุทฺธสฺส อวณฺณํ ภาสติ, ธมฺมสฺส อวณฺณํ ภาสติ, สํฆสฺส อวณฺณํ ภาสติ.\csromanspacer{} อิเมหิ โข ภิกฺขเว ตีหงฺเคหิ สมนฺนาคตสฺส ภิกฺขุโน อากงฺขมาโน สํโฆ นิยสฺสกมฺมํ กเรยฺย. (3)}

\prose{ติณฺณํ ภิกฺขเว ภิกฺขูนํ อากงฺขมาโน สํโฆ นิยสฺสกมฺมํ กเรยฺย.\csromanspacer{} เอโก ภณฺฑนการโก โหติ กลหการโก วิวาทการโก ภสฺสการโก สํเฆ อธิกรณการโก, เอโก พาโล โหติ อพฺยตฺโต อาปตฺติพหุโล อนปทาโน, เอโก คิหิสํสฏฺโฐ วิหรติ \csromanfolio{19}อนนุโลมิเกหิ คิหิสํสคฺเคหิ.\csromanspacer{} อิเมสํ โข ภิกฺขเว ติณฺณํ ภิกฺขูนํ อากงฺขมาโน สํโฆ นิยสฺสกมฺมํ กเรยฺย. (4)}

\prose{อปเรสมฺปิ ภิกฺขเว ติณฺณํ ภิกฺขูนํ อากงฺขมาโน สํโฆ นิยสฺสกมฺมํ กเรยฺย.\csromanspacer{} เอโก อธิสีเล สีลวิปนฺโน โหติ, เอโก อชฺฌาจาเร อาจารวิปนฺโน โหติ, เอโก อติทิฏฺฐิยา ทิฏฺฐิวิปนฺโน โหติ.\csromanspacer{} อิเมสํ โข ภิกฺขเว ติณฺณํ ภิกฺขูนํ อากงฺขมาโน สํโฆ นิยสฺสกมฺมํ กเรยฺย. (5)}

\prose{อปเรสมฺปิ ภิกฺขเว ติณฺณํ ภิกฺขูนํ อากงฺขมาโน สํโฆ นิยสฺสกมฺมํ กเรยฺย.\csromanspacer{} เอโก พุทฺธสฺส อวณฺณํ ภาสติ, เอโก ธมฺมสฺส อวณฺณํ ภาสติ, เอโก สํฆสฺส อวณฺณํ ภาสติ.\csromanspacer{} อิเมสํ โข ภิกฺขเว ติณฺณํ ภิกฺขูนํ อากงฺขมาโน สํโฆ นิยสฺสกมฺมํ กเรยฺย. (6)}

\nitthitam{อากงฺขมานฉกฺกํ นิฏฺฐิตํ.}
\csromanmatikaanchor{mtk.13}
\csromantocmarkref{subsection}{อฏฺฐารสวตฺต}{mtk.13}
\bookmark[dest={mtk.13},level=2]{อฏฺฐารสวตฺต}
\csromanheader{2}{\textbf{อฏฺฐารสวตฺต}}
\proseitem{16}{นิยสฺสกมฺมกเตน ภิกฺขเว ภิกฺขุนา สมฺมา วตฺติตพฺพํ.\csromanspacer{} ตตฺรายํ สมฺมาวตฺตนา\csromandash{}น อุปสมฺปาเทตพฺพํ, น นิสฺสโย ทาตพฺโพ, น สามเณโร อุปฏฺฐาเปตพฺโพ, น ภิกฺขุโนวาทกสมฺมุติ สาทิตพฺพา, สมฺมเตนปิ ภิกฺขุนิโย น โอวทิตพฺพา, ยาย อาปตฺติยา สํเฆน นิยสฺสกมฺมํ กตํ โหติ, สา อาปตฺติ น อาปชฺชิตพฺพา, อญฺญา วา ตาทิสิกา, ตโต วา ปาปิฏฺฐตรา, กมฺมํ น ครหิตพฺพํ, กมฺมิกา น ครหิตพฺพา, น ปกตตฺตสฺส ภิกฺขุโน อุโปสโถ ฐเปตพฺโพ, น ปวารณา ฐเปตพฺพา, น สวจนียํ กาตพฺพํ, น อนุวาโท ปฏฺฐเปตพฺโพ, น โอกาโส การาเปตพฺโพ, น โจเทตพฺโพ, น สาเรตพฺโพ, น ภิกฺขูหิ สมฺปโยเชตพฺพนฺติ.}

\nitthitam{นิยสฺสกมฺเม อฏฺฐารสวตฺตํ นิฏฺฐิตํ.}
\proseitem{17}{อถ โข สํโฆ เสยฺยสกสฺส ภิกฺขุโน นิยสฺสกมฺมํ อกาสิ “นิสฺสาย เต วตฺถพฺพนฺ”ติ.\csromanspacer{} โส สํเฆน นิยสฺสกมฺมกโต \csromanfolio{20}กลฺยาณมิตฺเต เสวมาโน ภชมาโน ปยิรุปาสมาโน อุทฺทิสาเปนฺโต ปริปุจฺฉนฺโต พหุสฺสุโต โหติ อาคตาคโม ธมฺมธโร วินยธโร มาติกาธโร ปณฺฑิโต วิยตฺโต เมธาวี ลชฺชี กุกฺกุจฺจโก สิกฺขากาโม, สมฺมา วตฺตติ, โลมํ ปาเตติ, เนตฺถารํ วตฺตติ, ภิกฺขู อุปสงฺกมิตฺวา เอวํ วเทติ “อหํ อาวุโส สํเฆน นิยสฺสกมฺมกโต สมฺมา วตฺตามิ, โลมํ ปาเตมิ, เนตฺถารํ วตฺตามิ, กถํ นุ โข มยา ปฏิปชฺชิตพฺพนฺ”ติ.\csromanspacer{} ภิกฺขู ภควโต เอตมตฺถํ อาโรเจสุํ ฯเปฯ.\csromanspacer{} เตน หิ ภิกฺขเว สํโฆ เสยฺยสกสฺส ภิกฺขุโน นิยสฺสกมฺมํ ปฏิปฺปสฺสมฺเภตุ.}

\csromanmatikaanchor{mtk.14}
\csromantocmarkref{subsection}{นปฺปฏิปฺปสฺสมฺเภตพฺพ-อฏฺฐารสก}{mtk.14}
\bookmark[dest={mtk.14},level=2]{นปฺปฏิปฺปสฺสมฺเภตพฺพ-อฏฺฐารสก}
\csromanheader{2}{\textbf{นปฺปฏิปฺปสฺสมฺเภตพฺพ-อฏฺฐารสก}}
\proseitem{18}{ปญฺจหิ ภิกฺขเว องฺเคหิ สมนฺนาคตสฺส ภิกฺขุโน นิยสฺสกมฺมํ นปฺปฏิปฺปสฺสมฺเภตพฺพํ.\csromanspacer{} อุปสมฺปาเทติ, นิสฺสยํ เทติ, สามเณรํ อุปฏฺฐาเปติ, ภิกฺขุโนวาทกสมฺมุติํ สาทิยติ, สมฺมโตปิ ภิกฺขุนิโย โอวทติ.\csromanspacer{} อิเมหิ โข ภิกฺขเว ปญฺจหงฺเคหิ สมนฺนาคตสฺส ภิกฺขุโน นิยสฺสกมฺมํ นปฺปฏิปฺปสฺสมฺเภตพฺพํ.}

\prose{\csromansymbolfootnote{*}{วิ 5. 317 ปิฏฺเฐปิ.}อปเรหิปิ ภิกฺขเว ปญฺจหงฺเคหิ สมนฺนาคตสฺส ภิกฺขุโน นิยสฺสกมฺมํ นปฺปฏิปฺปสฺสมฺเภตพฺพํ.\csromanspacer{} ยาย อาปตฺติยา สํเฆน นิยสฺสกมฺมํ กตํ โหติ, ตํ อาปตฺติํ อาปชฺชติ, อญฺญํ วา ตาทิสิกํ, ตโต วา ปาปิฏฺฐตรํ, กมฺมํ ครหติ, กมฺมิเก ครหติ.\csromanspacer{} อิเมหิ โข ภิกฺขเว ปญฺจหงฺเคหิ สมนฺนาคตสฺส ภิกฺขุโน นิยสฺสกมฺมํ นปฺปฏิปฺปสฺสมฺเภตพฺพํ.}

\prose{อฏฺฐหิ ภิกฺขเว องฺเคหิ สมนฺนาคตสฺส ภิกฺขุโน นิยสฺสกมฺมํ นปฺปฏิปฺปสฺสมฺเภตพฺพํ.\csromanspacer{} ปกตตฺตสฺส ภิกฺขุโน อุโปสถํ ฐเปติ, ปวารณํ ฐเปติ, สวจนียํ กโรติ, อนุวาทํ ปฏฺฐเปติ, โอกาสํ กาเรติ, โจเทติ, สาเรติ, ภิกฺขูหิ สมฺปโยเชติ.\csromanspacer{} อิเมหิ โข ภิกฺขเว อฏฺฐหงฺเคหิ สมนฺนาคตสฺส ภิกฺขุโน นิยสฺสกมฺมํ นปฺปฏิปฺปสฺสมฺเภตพฺพํ.}

\nitthitam{\textbf{นปฺปฏิปฺปสฺสมฺเภตพฺพ-อฏฺฐารสก}ํ นิฏฺฐิตํ.}
\csromanmatikaanchor{mtk.15}
\csromantocmarkref{subsection}{ปฏิปฺปสฺสมฺเภตพฺพ-อฏฺฐารสก}{mtk.15}
\bookmark[dest={mtk.15},level=2]{ปฏิปฺปสฺสมฺเภตพฺพ-อฏฺฐารสก}
\csromanheader{2}{\textbf{ปฏิปฺปสฺสมฺเภตพฺพ-อฏฺฐารสก}}
\proseitem{19}{\csromanfolio{21}ปญฺจหิ ภิกฺขเว องฺเคหิ สมนฺนาคตสฺส ภิกฺขุโน นิยสฺสกมฺมํ ปฏิปฺปสฺสมฺเภตพฺพํ.\csromanspacer{} น อุปสมฺปาเทติ, น นิสฺสยํ เทติ, น สามเณรํ อุปฏฺฐาเปติ, น ภิกฺขุโนวาทกสมฺมุติํ สาทิยติ, สมฺมโตปิ ภิกฺขุนิโย น โอวทติ.\csromanspacer{} อิเมหิ โข ภิกฺขเว ปญฺจหงฺเคหิ สมนฺนาคตสฺส ภิกฺขุโน นิยสฺสกมฺมํ ปฏิปฺปสฺสมฺเภตพฺพํ.}

\prose{อปเรหิปิ ภิกฺขเว ปญฺจหงฺเคหิ สมนฺนาคตสฺส ภิกฺขุโน นิยสฺสกมฺมํ ปฏิปฺปสฺสมฺเภตพฺพํ.\csromanspacer{} ยาย อาปตฺติยา สํเฆน นิยสฺสกมฺมํ กตํ โหติ, ตํ อาปตฺติํ น อาปชฺชติ, อญฺญํ วา ตาทิสิกํ, ตโต วา ปาปิฏฺฐตรํ, กมฺมํ น ครหติ, กมฺมิเก น ครหติ.\csromanspacer{} อิเมหิ โข ภิกฺขเว ปญฺจหงฺเคหิ สมนฺนาคตสฺส ภิกฺขุโน นิยสฺสกมฺมํ ปฏิปฺปสฺสมฺเภตพฺพํ.}

\prose{อฏฺฐหิ ภิกฺขเว องฺเคหิ สมนฺนาคตสฺส ภิกฺขุโน นิยสฺสกมฺมํ ปฏิปฺปสฺสมฺเภตพฺพํ.\csromanspacer{} น ปกตตฺตสฺส ภิกฺขุโน อุโปสถํ ฐเปติ, น ปวารณํ ฐเปติ, น สวจนียํ กโรติ, น อนุวาทํ ปฏฺฐเปติ, น โอกาสํ กาเรติ, น โจเทติ, น สาเรติ, น ภิกฺขูหิ สมฺปโยเชติ.\csromanspacer{} อิเมหิ โข ภิกฺขเว อฏฺฐหงฺเคหิ สมนฺนาคตสฺส ภิกฺขุโน นิยสฺสกมฺมํ ปฏิปฺปสฺสมฺเภตพฺพํ.}

\nitthitam{\textbf{ปฏิปฺปสฺสมฺเภตพฺพ-อฏฺฐารสก}ํ นิฏฺฐิตํ.}
\proseitem{20}{เอวญฺจ ปน ภิกฺขเว ปฏิปฺปสฺสมฺเภตพฺพํ, เตน ภิกฺขเว เสยฺยสเกน ภิกฺขุนา สํฆํ อุปสงฺกมิตฺวา เอกํสํ อุตฺตราสงฺคํ กริตฺวา วุฑฺฒานํ ภิกฺขูนํ ปาเท วนฺทิตฺวา อุกฺกุฏิกํ นิสีทิตฺวา อญฺชลิํ ปคฺคเหตฺวา เอวมสฺส วจนีโย “อหํ ภนฺเต สํเฆน นิยสฺสกมฺมกโต สมฺมา วตฺตามิ, โลมํ ปาเตมิ, เนตฺถารํ วตฺตามิ, นิยสฺสสฺส กมฺมสฺส ปฏิปฺปสฺสทฺธิํ ยาจามี”ติ.\csromanspacer{} ทุติยมฺปิ ยาจิตพฺพา.\csromanspacer{} ตติยมฺปิ ยาจิตพฺพา.\csromanspacer{} พฺยตฺเตน ภิกฺขุนา ปฏิพเลน สํโฆ ญาเปตพฺโพ\csromandash{}}

\prose{“สุณาตุ เม ภนฺเต สํโฆ, อยํ เสยฺยสโก ภิกฺขุ สํเฆน นิยสฺสกมฺมกโต สมฺมา วตฺตติ, โลมํ ปาเตติ, เนตฺถารํ วตฺตติ, นิยสฺสสฺส กมฺมสฺส ปฏิปฺปสฺสทฺธิํ ยาจติ, ยทิ สํฆสฺส ปตฺตกลฺลํ, สํโฆ เสยฺยสกสฺส ภิกฺขุโน นิยสฺสกมฺมํ ปฏิปฺปสฺสมฺเภยฺย, เอสา ญตฺติ.}

\prose{\csromanfolio{22}สุณาตุ เม ภนฺเต สํโฆ, อยํ เสยฺยสโก ภิกฺขุ สํเฆน นิยสฺสกมฺมกโต สมฺมา วตฺตติ, โลมํ ปาเตติ, เนตฺถารํ วตฺตติ, นิยสฺสสฺส กมฺมสฺส ปฏิปฺปสฺสทฺธิํ ยาจติ, สํโฆ เสยฺยสกสฺส ภิกฺขุโน นิยสฺสกมฺมํ ปฏิปฺปสฺสมฺเภติ, ปสฺสายสฺมโต ขมติ เสยฺยสกสฺส ภิกฺขุโน นิยสฺสสฺส กมฺมสฺส ปฏิปฺปสฺสทฺธิ, โส ตุณฺหสฺส, ยสฺส นกฺขมติ, โส ภาเสยฺย.}

\prose{ทุติยมฺปิ เอตมตฺถํ วทามิ.\csromanspacer{} สุณาตุ เม ภนฺเต สํโฆ, อยํ เสยฺยสโก ภิกฺขุ สํเฆน นิยสฺสกมฺมกโต สมฺมา วตฺตติ, โลมํ ปาเตติ, เนตฺถารํ วตฺตติ, นิยสฺสสฺส กมฺมสฺส ปฏิปฺปสฺสทฺธิํ ยาจติ, สํโฆ เสยฺยสกสฺส ภิกฺขุโน นิยสฺสกมฺมํ ปฏิปฺปสฺสมฺเภติ, ยสฺสายสฺมโต ขมติ เสยฺยสกสฺส ภิกฺขุโน นิยสฺสสฺส กมฺมสฺส ปฏิปฺปสฺสทฺธิ, โส ตุณฺหสฺส, ยสฺส นกฺขมติ, โส ภาเสยฺย.}

\prose{ตติยมฺปิ เอตมตฺถํ วทามิ.\csromanspacer{} สุณาตุ เม ภนฺเต สํโฆ, อยํ เสยฺยสโก ภิกฺขุ สํเฆน นิยสฺสกมฺมกโต สมฺมา วตฺตติ, โลมํ ปาเตติ, เนตฺถารํ วตฺตติ, นิยสฺสสฺส กมฺมสฺส ปฏิปฺปสฺสทฺธิํ ยาจติ, สํโฆ เสยฺยสกสฺส ภิกฺขุโน นิยสฺสกมฺมํ ปฏิปฺปสฺสมฺเภติ, ยสฺสายสฺมโต ขมติ เสยฺยสกสฺส ภิกฺขุโน นิยสฺสสฺส กมฺมสฺส ปฏิปฺปสฺสทฺธิ, โส ตุณฺหสฺส, ยสฺส นกฺขมติ, โส ภาเสยฺย.}

\prose{ปฏิปฺปสฺสทฺธํ สํเฆน เสยฺยสกสฺส ภิกฺขุโน นิยสฺสกมฺมํ, ขมติ สํฆสฺส, ตสฺมา ตุณฺหี, เอวเมตํ ธารยามี”ติ.}

\nitthitam{นิยสฺสกมฺมํ นิฏฺฐิตํ ทุติยํ.}
\csromanmatikaanchor{mtk.16}
\csromantocmarkref{section}{3. ปพฺพาชนียกมฺม}{mtk.16}
\bookmark[dest={mtk.16},level=1]{3. ปพฺพาชนียกมฺม}
\csromanheader{1}{3. ปพฺพาชนียกมฺม}
\proseitem{21}{\csromansymbolfootnote{*}{อิทํ วตฺถุ วิ 1. 274 ปิฏฺฐาทีสุปิ อาคตํ.}เตน โข ปน สมเยน อสฺสชิปุนพฺพสุกา นาม\footnote{นาม ภิกฺขู (ก)} กีฏาคิริสฺมิํ อาวาสิกา โหนฺติ อลชฺชิโน ปาปภิกฺขู, เต + เอวรูปํ อนาจารํ อาจรนฺติ, มาลาวจฺฉํ โรเปนฺติปิ โรปาเปนฺติปิ, สิญฺจนฺติปิ สิญฺจาเปนฺติปิ, โอจินนฺติปิ โอจินาเปนฺติปิ, คนฺเถนฺติปิ คนฺถาเปนฺติปิ, เอกโตวณฺฏิกมาลํ \csromanfolio{23}กโรนฺติปิ การาเปนฺติปิ, อุภโตวณฺฏิกมาลํ กโรนฺติปิ การาเปนฺติปิ, มญฺชริกํ กโรนฺติปิ การาเปนฺติปิ, วิธูติกํ\footnote{วิธุติกํ (สฺยา)} กโรนฺติปิ การาเปนฺติปิ, วฏํสกํ กโรนฺติปิ การาเปนฺติปิ, อาเวฬํ กโรนฺติปิ การาเปนฺติปิ, อุรจฺฉทํ กโรนฺติปิ การาเปนฺติปิ.\csromanspacer{} เต กุลิตฺถีนํ กุลธีตานํ กุลกุมารีนํ กุลสุณฺหานํ กุลทาสีนํ เอกโตวณฺฏิกมาลํ หรนฺติปิ หราเปนฺติปิ, อุภโตวณฺฏิกมาลํ หรนฺติปิ หราเปนฺติปิ, มญฺชริกํ หรนฺติปิ หราเปนฺติปิ, วิธูติกํ หรนฺติปิ หราเปนฺติปิ, วฏํสกํ หรนฺติปิ หราเปนฺติปิ, อาเวฬํ หรนฺติปิ หราเปนฺติปิ, อุรจฺฉทํ หรนฺติปิ หราเปนฺติปิ.\csromanspacer{} เต กุลิตฺถีหิ กุลธีตาหิ กุลกุมารีหิ กุลสุณฺหาหิ กุลทาสีหิ สทฺธิํ เอกภาชเนปิ ภุญฺชนฺติ, เอกถาลเกปิ ปิวนฺติ, เอกาสเนปิ นิสีทนฺติ, เอกมญฺเจปิ ตุวฏฺเฏนฺติ, เอกตฺถรณาปิ ตุวฏฺเฏนฺติ, เอกปาวุรณาปิ ตุวฏฺเฏนฺติ, เอกตฺถรณปาวุรณาปิ ตุวฏฺเฏนฺติ, วิกาเลปิ ภุญฺชนฺติ, มชฺชมฺปิ ปิวนฺติ, มาลาคนฺธวิเลปนมฺปิ ธาเรนฺติ, นจฺจนฺติปิ คายนฺติปิ วาเทนฺติปิ ลาเสนฺติปิ, นจฺจนฺติยาปิ นจฺจนฺติ นจฺจนฺติยาปิ คายนฺติ นจฺจนฺติยาปิ วาเทนฺติ นจฺจนฺติยาปิ ลาเสนฺติ, คายนฺติยาปิ นจฺจนฺติ คายนฺติยาปิ คายนฺติ คายนฺติยาปิ วาเทนฺติ คายนฺติยาปิ ลาเสนฺติ, วาเทนฺติยาปิ นจฺจนฺติ วาเทนฺติยาปิ คายนฺติ วาเทนฺติยาปิ วาเทนฺติ วาเทนฺติยาปิ ลาเสนฺติ, ลาเสนฺติยาปิ นจฺจนฺติ ลาเสนฺติยาปิ คายนฺติ ลาเสนฺติยาปิ วาเทนฺติ ลาเสนฺติยาปิ ลาเสนฺติ.\csromanspacer{} อฏฺฐปเทปิ กีฬนฺติ, ทสปเทปิ กีฬนฺติ, อากาเสปิ กีฬนฺติ, ปริหารปเถปิ กีฬนฺติ, สนฺติกายปิ กีฬนฺติ, ขลิกายปิ กีฬนฺติ, ฆฏิกายปิ กีฬนฺติ, สลากหตฺเถนปิ กีฬนฺติ, อกฺเขนปิ กีฬนฺติ, ปงฺคจีเรนปิ กีฬนฺติ, วงฺกเกนปิ กีฬนฺติ, โมกฺขจิกายปิ กีฬนฺติ, จิงฺคุลเกนปิ กีฬนฺติ, ปตฺตาฬฺหเกนปิ กีฬนฺติ, รถเกนปิ กีฬนฺติ, ธนุเกนปิ กีฬนฺติ, อกฺขริกายปิ กีฬนฺติ, มเนสิกายปิ กีฬนฺติ, ยถาวชฺเชนปิ กีฬนฺติ, หตฺถิสฺมิมฺปิ สิกฺขนฺติ, อสฺสสฺมิมฺปิ สิกฺขนฺติ, รถสฺมิมฺปิ สิกฺขนฺติ, ธนุสฺมิมฺปิ สิกฺขนฺติ, ถรุสฺมิมฺปิ สิกฺขนฺติ, หตฺถิสฺสปิ ปุรโต ธาวนฺติ, อสฺสสฺสปิ ปุรโต ธาวนฺติ, รถสฺสปิ ปุรโต\footnote{ปุรโต ธาวนฺติ (สฺยา)} ธาวนฺติปิ อาธาวนฺติปิ, อุสฺเสเฬนฺติปิ, อปฺโผเฏนฺติปิ, นิพฺพุชฺฌนฺติปิ, มุฏฺฐีหิปิ ยุชฺฌนฺติ, รงฺคมชฺเฌปิ สงฺฆาฏิํ ปตฺถริตฺวา นจฺจกิํ\footnote{นจฺจนฺติํ (สี, สฺยา)} เอวํ วทนฺติ “อิธ ภคินิ นจฺจสฺสู”ติ, นลาฏิกมฺปิ เทนฺติ, วิวิธมฺปิ อนาจารํ อาจรนฺติ.}

\proseitem{22}{\csromanfolio{24}เตน โข ปน สมเยน อญฺญตโร ภิกฺขุ กาสีสุ วสฺสํวุฏฺโฐ\footnote{วสฺสํวุตฺโถ (สี, สฺยา)} สาวตฺถิํ คจฺฉนฺโต ภควนฺตํ ทสฺสนาย เยน กีฏาคิริ ตทวสริ.\csromanspacer{} อถ โข โส ภิกฺขุ ปุพฺพณฺหสมยํ นิวาเสตฺวา ปตฺตจีวรมาทาย กีฏาคิริํ ปิณฺฑาย ปาวิสิ ปาสาทิเกน อภิกฺกนฺเตน ปฏิกฺกนฺเตน อาโลกิเตน วิโลกิเตน สมิญฺชิเตน\footnote{สมฺมิญฺชิเตน (สี, สฺยา, กํ)} ปสาริเตน โอกฺขิตฺตจกฺขุ อิริยาปถสมฺปนฺโน.\csromanspacer{} มนุสฺสา ตํ ภิกฺขุํ ปสฺสิตฺวา เอวมาหํสุ “กฺวายํ อพลพโล วิย มนฺทมนฺโท วิย ภากุฏิกภากุฏิโก วิย, โก อิมสฺส อุปคตสฺส ปิณฺฑกมฺปิ ทสฺสติ, อมฺหากํ ปน อยฺยา อสฺสชิปุนพฺพสุกา สณฺหา สขิลา สุขสมฺภาสา มิหิตปุพฺพงฺคมา เอหิสฺวาคตวาทิโน อพฺภากุฏิกา อุตฺตานมุขา ปุพฺพภาสิโน, เตสํ โข นาม ปิณฺโฑ ทาตพฺโพ”ติ.}

\prose{อทฺทสา โข อญฺญตโร อุปาสโก ตํ ภิกฺขุํ กีฏาคิริสฺมิํ ปิณฺฑาย จรนฺตํ, ทิสฺวาน เยน โส ภิกฺขุ เตนุปสงฺกมิ, อุปสงฺกมิตฺวา ตํ ภิกฺขุํ อภิวาเทตฺวา เอตทโวจ “อปิ ภนฺเต ปิณฺโฑ ลพฺภตี”ติ.\csromanspacer{} น โข อาวุโส ปิณฺโฑ ลพฺภตีติ.\csromanspacer{} เอหิ ภนฺเต, ฆรํ คมิสฺสามาติ.\csromanspacer{} อถ โข โส อุปาสโก ตํ ภิกฺขุํ ฆรํ เนตฺวา โภเชตฺวา เอตทโวจ “กหํ ภนฺเต อยฺโย คมิสฺสตี”ติ.\csromanspacer{} สาวตฺถิํ โข อหํ อาวุโส คมิสฺสามิ ภควนฺตํ ทสฺสนายาติ.\csromanspacer{} เตน หิ ภนฺเต มม วจเนน ภควโต ปาเท สิรสา วนฺท, เอวญฺจ วเทหิ “ทุฏฺโฐ ภนฺเต กีฏาคิริสฺมิํ อาวาโส, อสฺสชิปุนพฺพสุกา นาม กีฏาคิริสฺมิํ อาวาสิกา อลชฺชิโน ปาปภิกฺขู, เต เอวรูปํ อนาจารํ อาจรนฺติ, มาลาวจฺฉํ โรเปนฺติปิ โรปาเปนฺติปิ, สิญฺจนฺติปิ สิญฺจาเปนฺติปิ, โอจินนฺติปิ โอจินาเปนฺติปิ, คนฺเถนฺติปิ คนฺถาเปนฺติปิ, เอกโตวณฺฏิกมาลํ กโรนฺติปิ การาเปนฺติปิ, อุภโตวณฺฏิกมาลํ กโรนฺติปิ การาเปนฺติปิ, มญฺชริกํ กโรนฺติปิ การาเปนฺติปิ, วิธูติกํ กโรนฺติปิ การาเปนฺติปิ, วฏํสกํ กโรนฺติปิ การาเปนฺติปิ, อาเวฬํ กโรนฺติปิ การาเปนฺติปิ, อุรจฺฉทํ กโรนฺติปิ การาเปนฺติปิ.\csromanspacer{} เต กุลิตฺถีนํ กุลธีตานํ กุลกุมารีนํ กุลสุณฺหานํ กุลทาสีนํ เอกโตวณฺฏิกมาลํ หรนฺติปิ หราเปนฺติปิ, อุภโตวณฺฏิกมาลํ หรนฺติปิ หราเปนฺติปิ, มญฺชริกํ หรนฺติปิ หราเปนฺติปิ, \csromanfolio{25}วิธูติกํ หรนฺติปิ หราเปนฺติปิ, วฏํสกํ หรนฺติปิ หราเปนฺติปิ, อาเวฬํ หรนฺติปิ หราเปนฺติปิ, อุรจฺฉทํ หรนฺติปิ หราเปนฺติปิ.\csromanspacer{} เต กุลิตฺถีหิ กุลธีตาหิ กุลกุมารีหิ กุลสุณฺหาหิ กุลทาสีหิ สทฺธิํ เอกภาชเนปิ ภุญฺชนฺติ, เอกถาลเกปิ ปิวนฺติ, เอกาสเนปิ นิสีทนฺติ, เอกมญฺเจปิ ตุวฏฺเฏนฺติ, เอกตฺถรณาปิ ตุวฏฺเฏนฺติ, เอกปาวุรณาปิ ตุวฏฺเฏนฺติ, เอกตฺถรณปาวุรณาปิ ตุวฏฺเฏนฺติ, วิกาเลปิ ภุญฺชนฺติ, มชฺชมฺปิ ปิวนฺติ, มาลาคนฺธวิเลปนมฺปิ ธาเรนฺติ, นจฺจนฺติปิ คายนฺติปิ วาเทนฺติปิ ลาเสนฺติปิ, นจฺจนฺติยาปิ นจฺจนฺติ นจฺจนฺติยาปิ คายนฺติ นจฺจนฺติยาปิ วาเทนฺติ นจฺจนฺติยาปิ ลาเสนฺติ ฯเปฯ (จกฺกํ กาตพฺพํ).\csromanspacer{} ลาเสนฺติยาปิ นจฺจนฺติ ลาเสนฺติยาปิ คายนฺติ ลาเสนฺติยาปิ วาเทนฺติ ลาเสนฺติยาปิ ลาเสนฺติ, อฏฺฐปเทปิ กีฬนฺติ, ทสปเทปิ กีฬนฺติ, อากาเสปิ กีฬนฺติ, ปริหารปเถปิ กีฬนฺติ, สนฺติกายปิ กีฬนฺติ, ขลิกายปิ กีฬนฺติ, ฆฏิกายปิ กีฬนฺติ, สลากหตฺเถนปิ กีฬนฺติ, อกฺเขนปิ กีฬนฺติ, ปงฺคจีเรนปิ กีฬนฺติ, วงฺกเกนปิ กีฬนฺติ, โมกฺขจิกายปิ กีฬนฺติ, จิงฺคุลเกนปิ กีฬนฺติ, ปตฺตาฬฺหเกนปิ กีฬนฺติ, รถเกนปิ กีฬนฺติ, ธนุเกนปิ กีฬนฺติ, อกฺขริกายปิ กีฬนฺติ, มเนสิกายปิ กีฬนฺติ, ยถาวชฺเชนปิ กีฬนฺติ, หตฺถิสฺมิมฺปิ สิกฺขนฺติ, อสฺสสฺมิมฺปิ สิกฺขนฺติ, รถสฺมิมฺปิ สิกฺขนฺติ, ธนุสฺมิมฺปิ สิกฺขนฺติ, ถรุสฺมิมฺปิ สิกฺขนฺติ, หตฺถิสฺสปิ ปุรโต ธาวนฺติ, อสฺสสฺสปิ ปุรโต ธาวนฺติ, รถสฺสปิ ปุรโต ธาวนฺติปิ อาธาวนฺติปิ, อุสฺเสเฬนฺติปิ, อปฺโผเฏนฺติปิ, นิพฺพุชฺฌนฺติปิ, มุฏฺฐีหิปิ ยุชฺฌนฺติ, รงฺคมชฺเฌปิ สงฺฆาฏิํ ปตฺถริตฺวา นจฺจกิํ เอวํ วทนฺติ “อิธ ภคินิ นจฺจสฺสู”ติ, นลาฏิกมฺปิ เทนฺติ, วิวิธมฺปิ อนาจารํ อาจรนฺติ.\csromanspacer{} เยปิ เต ภนฺเต มนุสฺสา ปุพฺเพ สทฺธา อเหสุํ ปสนฺนา, เตปิ เอตรหิ อสฺสทฺธา อปฺปสนฺนา, ยานิปิ ตานิ สํฆสฺส ปุพฺเพ ทานปถานิ, ตานิปิ เอตรหิ อุปจฺฉินฺนานิ, ริญฺจนฺติ เปสลา ภิกฺขู, นิวสนฺติ ปาปภิกฺขู, สาธุ ภนฺเต ภควา กีฏาคิริํ ภิกฺขู ปหิเณยฺย, ยถายํ กีฏาคิริสฺมิํ อาวาโส สณฺฐเหยฺยา”ติ.}

\prose{“เอวมาวุโส”ติ โข โส ภิกฺขุ ตสฺส อุปาสกสฺส ปฏิสฺสุณิตฺวา อุฏฺฐายาสนา เยน สาวตฺถิ เตน ปกฺกามิ, อนุปุพฺเพน เยน สาวตฺถิ เชตวนํ อนาถปิณฺฑิกสฺส อาราโม, เยน ภควา เตนุปสงฺกมิ, อุปสงฺกมิตฺวา ภควนฺตํ อภิวาเทตฺวา เอกมนฺตํ นิสีทิ, อาจิณฺณํ โข ปเนตํ พุทฺธานํ ภควนฺตานํ อาคนฺตุเกหิ ภิกฺขูหิ สทฺธิํ ปฏิสมฺโมทิตุํ. \csromanfolio{26}อถ โข ภควา ตํ ภิกฺขุํ เอตทโวจ “กจฺจิ ภิกฺขุ ขมนียํ, กจฺจิ ยาปนียํ, กจฺจิสิ อปฺปกิลมเถน อทฺธานํ อาคโต, กุโต จ ตฺวํ ภิกฺขุ อาคจฺฉสี”ติ.\csromanspacer{} ขมนียํ ภควา, ยาปนียํ ภควา, อปฺปกิลมเถน จ อหํ ภนฺเต อทฺธานํ อาคโต, อิธาหํ ภนฺเต กาสีสุ วสฺสํวุฏฺโฐ สาวตฺถิํ อาคจฺฉนฺโต ภควนฺตํ ทสฺสนาย เยน กีฏาคิริ ตทวสริํ.\csromanspacer{} อถ ขฺวาหํ ภนฺเต ปุพฺพณฺหสมยํ นิวาเสตฺวา ปตฺตจีวรมาทาย กีฏาคิริํ ปิณฺฑาย ปาวิสิํ, อทฺทสา โข มํ ภนฺเต อญฺญตโร อุปาสโก กีฏาคิริสฺมิํ ปิณฺฑาย จรนฺตํ, ทิสฺวาน เยนาหํ เตนุปสงฺกมิ, อุปสงฺกมิตฺวา มํ อภิวาเทตฺวา เอตทโวจ “อปิ ภนฺเต ปิณฺโฑ ลพฺภตี”ติ, น โข อาวุโส ปิณฺโฑ ลพฺภตีติ, เอหิ ภนฺเต, ฆรํ คมิสฺสามาติ, อถ โข ภนฺเต โส อุปาสโก มํ ฆรํ เนตฺวา โภเชตฺวา เอตทโวจ “กหํ ภนฺเต อยฺโย คมิสฺสาตี”ติ, สาวตฺถิํ โข อหํ อาวุโส คมิสฺสามิ ภควนฺตํ ทสฺสนายาติ, เตน หิ ภนฺเต มม วจเนน ภควโต ปาเท สิรสา วนฺท, เอวญฺจ วเทหิ “ทุฏฺโฐ ภนฺเต กีฏาคิริสฺมิํ อาวาโส, อสฺสชิปุนพฺพสุกา นาม กีฏาคิริสฺมิํ อาวาสิกา อลชฺชิโน ปาปภิกฺขู, เต เอวรูปํ อนาจารํ อาจรนฺติ, มาลาวจฺฉํ โรเปนฺติปิ โรปาเปนฺติปิ ฯเปฯ วิวิธมฺปิ อนาจารํ อาจรนฺติ.\csromanspacer{} เยปิ เต ภนฺเต มนุสฺสา ปุพฺเพ สทฺธา อเหสุํ ปสนฺนา, เตปิ เอตรหิ อสฺสทฺธา อปฺปสนฺนา, ยานิปิ ตานิ สํฆสฺส ปุพฺเพ ทานปถานิ, ตานิปิ เอตรหิ อุปจฺฉินฺนานิ, ริญฺจนฺติ เปสลา ภิกฺขู, นิวสนฺติ ปาปภิกฺขู, สาธุ ภนฺเต ภควา กีฏาคิริํ ภิกฺขู ปหิเณยฺย, ยถายํ กีฏาคิริสฺมิํ อาวาโส สณฺฐเหยฺยา”ติ, ตโต อหํ ภควา อาคจฺฉามีติ.}

\proseitem{23}{อถ โข ภควา เอตสฺมิํ นิทาเน เอตสฺมิํ ปกรเณ ภิกฺขุสํฆํ สนฺนิปาตาเปตฺวา ภิกฺขู ปฏิปุจฺฉิ “สจฺจํ กิร ภิกฺขเว อสฺสชิปุนพฺพสุกา นาม กีฏาคิริสฺมิํ อาวาสิกา อลชฺชิโน ปาปภิกฺขู, เต เอวรูปํ อนาจารํ อาจรนฺติ, มาลาวจฺฉํ โรเปนฺติปิ ฯเปฯ วิวิธมฺปิ อนาจารํ อาจรนฺติ.\csromanspacer{} เยปิ เต มนุสฺสา ปุพฺเพ สทฺธา อเหสุํ ปสนฺนา, เตปิ เอตรหิ อสฺสทฺธา อปฺปสนฺนา, ยานิปิ ตานิ สํฆสฺส ปุพฺเพ ทานปถานิ, ตานิปิ เอตรหิ อุปจฺฉินฺนานิ, ริญฺจนฺติ เปสลา ภิกฺขู, นิวสนฺติ ปาปภิกฺขู”ติ.\csromanspacer{} สจฺจํ ภควาติ.\csromanspacer{} วิครหิ พุทฺโธ ภควา อนนุจฺฉวิกํ ฯเปฯ.\csromanspacer{} กถํ หิ นาม เต \csromanfolio{27}ภิกฺขเว โมฆปุริสา เอวรุปํ อนาจารํ อาจริสฺสนฺติ, มาลาวจฺฉํ โรเปสฺสนฺติปิ โรปาเปสฺสนฺติปิ, สิญฺจิสฺสนฺติปิ สิญฺจาเปสฺสนฺติปิ, โอจินิสฺสนฺติปิ โอจินาเปสฺสนฺติปิ, คนฺเถสฺสนฺติปิ คนฺถาเปสฺสนฺติปิ, เอกโตวณฺฏิกมาลํ กริสฺสนฺติปิ การาเปสฺสนฺติปิ, อุภโตวณฺฏิกมาลํ กริสฺสนฺติปิ การาเปสฺสนฺติปิ, มญฺชริกํ กริสฺสนฺติปิ การาเปสฺสนฺติปิ, วิธูติกํ กริสฺสนฺติปิ การาเปสฺสนฺติปิ, วฏํสกํ กริสฺสนฺติปิ การาเปสฺสนฺติปิ, อาเวฬํ กริสฺสนฺติปิ การาเปสฺสนฺติปิ, อุรจฺฉทํ กริสฺสนฺติปิ การาเปสฺสนฺติปิ.\csromanspacer{} เต กุลิตฺถีนํ กุลธีตานํ กุลกุมารีนํ กุลสุณฺหานํ กุลทาสีนํ เอกโตวณฺฏิกมาลํ หริสฺสนฺติปิ หราเปสฺสนฺติปิ, อุภโตวณฺฏิกมาลํ หริสฺสนฺติปิ หราเปสฺสนฺติปิ, มญฺชริกํ หริสฺสนฺติปิ หราเปสฺสนฺติปิ, วิธูติกํ หริสฺสนฺติปิ หราเปสฺสนฺติปิ, วฏํสกํ หริสฺสนฺติปิ หราเปสฺสนฺติปิ, อาเวฬํ หริสฺสนฺติปิ หราเปสฺสนฺติปิ, อุรจฺฉทํ หริสฺสนฺติปิ หราเปสฺสนฺติปิ.\csromanspacer{} เต กุลิตฺถีหิ กุลธีตาหิ กุลกุมารีหิ กุลสุณฺหาหิ กุลทาสีหิ สทฺธิํ เอกภาชเนปิ ภุญฺชิสฺสนฺติ, เอกถาลเกปิ ปิวิสฺสนฺติ, เอกาสเนปิ นิสีทิสฺสนฺติ, เอกมญฺเจปิ ตุวฏฺฏิสฺสนฺติ, เอกตฺถรณาปิ ตุวฏฺฏิสฺสนฺติ, เอกปาวุรณาปิ ตุวฏฺฏิสฺสนฺติ, เอกตฺถรณปาวุรณาปิ ตุวฏฺฏิสฺสนฺติ, วิกาเลปิ ภุญฺชิสฺสนฺติ, มชฺชมฺปิ ปิวิสฺสนฺติ, มาลาคนฺธวิเลปนมฺปิ ธาเรสฺสนฺติ, นจฺจิสฺสนฺติปิ คายิสฺสนฺติปิ วาเทสฺสนฺติปิ ลาเสสฺสนฺติปิ, นจฺจนฺติยาปิ นจฺจิสฺสนฺติ นจฺจนฺติยาปิ คายิสฺสนฺติ นจฺจนฺติยาปิ วาเทสฺสนฺติ นจฺจนฺติยาปิ ลาเสสฺสนฺติ, คายนฺติยาปิ นจฺจิสฺสนฺติ คายนฺติยาปิ คายิสฺสนฺติ คายนฺติยาปิ วาเทสฺสนฺติ คายนฺติยาปิ ลาเสสฺสนฺติ, วาเทนฺติยาปิ นจฺจิสฺสนฺติ วาเทนฺติยาปิ คายิสฺสนฺติ วาเทนฺติยาปิ วาเทสฺสนฺติ วาเทนฺติยาปิ ลาเสสฺสนฺติ, ลาเสนฺติยาปิ นจฺจิสฺสนฺติ ลาเสนฺติยาปิ คายิสฺสนฺติ, ลาเสนฺติยาปิ วาเทสฺสนฺติ, ลาเสนฺติยาปิ ลาเสสฺสนฺติ, อฏฺฐปเทปิ กีฬิสฺสนฺติ, ทสปเทปิ กีฬิสฺสนฺติ, อากาเสปิ กีฬิสฺสนฺติ, ปริหารปเถปิ กีฬิสฺสนฺติ, สนฺติกายปิ กีฬิสฺสนฺติ, ขลิกายปิ กีฬิสฺสนฺติ, ฆฏิกายปิ กีฬิสฺสนฺติ, สลากหตฺเถนปิ กีฬิสฺสนฺติ, อกฺเขนปิ กีฬิสฺสนฺติ, ปงฺคจีเรนปิ กีฬิสฺสนฺติ, วงฺกเกนปิ กีฬิสฺสนฺติ, โมกฺขจิกายปิ กีฬิสฺสนฺติ, จิงฺคุลเกนปิ กีฬิสฺสนฺติ, ปตฺตาฬฺหเกนปิ กีฬิสฺสนฺติ, รถเกนปิ กีฬิสฺสนฺติ, ธนุเกนปิ กีฬิสฺสนฺติ, อกฺขริกายปิ กีฬิสฺสนฺติ, มเนสิกายปิ กีฬิสฺสนฺติ, ยถาวชฺเชนปิ กีฬิสฺสนฺติ, หตฺถิสฺมิมฺปิ สิกฺขิสฺสนฺติ, อสฺสสฺมิมฺปิ สิกฺขิสฺสนฺติ, รถสฺมิมฺปิ สิกฺขิสฺสนฺติ, ธนุสฺมิมฺปิ สิกฺขิสฺสนฺติ, ถรุสฺมิมฺปิ สิกฺขิสฺสนฺติ, หตฺถิสฺสปิ ปุรโต \csromanfolio{28}ธาวิสฺสนฺติ, อสฺสสฺสปิ ปุรโต มาวิสฺสนฺติ, รถสฺสปิ ปุรโต\footnote{ปุรโต ธาวิสฺสนฺติ (สฺยา)} ธาวิสฺสนฺติปิ อาธาวิสฺสนฺติปิ, อุสฺเสเฬสฺสนฺติปิ, อปฺโผเฏสฺสนฺติปิ, นิพฺพุชฺฌิสฺสนฺติปิ, มุฏฺฐีหิปิ ยุชฺฌิสฺสนฺติ, รงฺคมชฺเฌปิ สงฺฆาฏิํ ปตฺถริตฺวา นจฺจกิํ เอวํ วกฺขนฺติ\footnote{วทิสฺสนฺติ (ก)} “อิธ ภคินิ นจฺจสฺสู”ติ, นลาฏิกมฺปิ ทสฺสนฺติ, วิวิธมฺปิ อนาจารํ อาจาริสฺสนฺติ, เนตํ ภิกฺขเว อปฺปสนฺนานํ วา ปสาทาย ฯเปฯ วิครหิตฺวา ฯเปฯ ธมฺมิํ กถํ กตฺวา สาริปุตฺตโมคฺคลฺลาเน อามนฺเตสิ “คจฺฉถ ตุเมฺห สาริปุตฺตา, กีฏาคิริํ คนฺตฺวา อสฺสชิปุนพฺพสุกานํ ภิกฺขูนํ กีฏาคิริสฺมา ปพฺพาชนียกมฺมํ กโรถ, ตุมฺหากํ เอเต สทฺธิวิหาริโน”ติ.}

\prose{กถํ มยํ ภนฺเต อสฺสชิปุนพฺพสุกานํ ภิกฺขูนํ กีฏาคิริสฺมา ปพฺพาชนียกมฺมํ กโรม, จณฺฑา เต ภิกฺขู ผรุสาติ.\csromanspacer{} เตน หิ ตุเมฺห สาริปุตฺตา พหุเกหิ ภิกฺขูหิ สทฺธิํ คจฺฉถาติ. “เอวํ ภนฺเต”ติ โข สาริปุตฺตโมคฺคลฺลานา ภควโต ปจฺจสฺโสสุํ.\csromanspacer{} เอวญฺจ ปน ภิกฺขเว กาตพฺพํ, ปฐมํ อสฺสชิปุนพฺพสุกา ภิกฺขู โจเทตพฺพา, โจเทตฺวา สาเรตพฺพา, สาเรตฺวา อาปตฺติํ อาโรเปตพฺพา, อาปตฺติํ อาโรเปตฺวา พฺยตฺเตน ภิกฺขุนา ปฏิพเลน สํโฆ ญาเปตพฺโพ\csromandash{}}

\proseitem{24}{“สุณาตุ เม ภนฺเต สํโฆ, อิเม อสฺสชิปุนพฺพสุกา ภิกฺขู กุลทูสกา ปาปสมาจารา, อิเมสํ ปาปกา สมาจารา ทิสฺสนฺติ เจว สุยฺยนฺติ จ, กุลานิ จ อิเมหิ ทุฏฺฐานิ ทิสฺสนฺติ เจว สุยฺยนฺติ จ, ยทิ สํฆสฺส ปตฺตกลฺลํ, สํโฆ อสฺสชิปุนพฺพสุกานํ ภิกฺขูนํ กีฏาคิริสฺมา ปพฺพาชนียกมฺมํ กเรยฺย ‘น อสฺสชิปุนพฺพสุเกหิ ภิกฺขูหิ กีฏาคิริสฺมิํ วตฺถพฺพนฺ’ติ, เอสา ญตฺติ.}

\prose{สุณาตุ เม ภนฺเต สํโฆ, อิเม อสฺสชิปุนพฺพสุกา ภิกฺขู กุลทูสกา ปาปสมาจารา, อิเมสํ ปาปกา สมาจารา ทิสฺสนฺติ เจว สุยฺยนฺติ จ, กุลานิ จ อิเมหิ ทุฏฺฐานิ ทิสฺสนฺติ เจว สุยฺยนฺติ จ, สํโฆ อสฺสชิปุนพฺพสุกานํ ภิกฺขูนํ กีฏาคิริสฺมา ปพฺพาชนียกมฺมํ กโรติ ‘น อสฺสชิปุนพฺพสุเกหิ ภิกฺขูหิ กีฏาคิริสฺมิํ วตฺถพฺพนฺ’ติ, ยสฺสายสฺมโต ขมติ อสฺสชิปุนพฺพสุกานํ ภิกฺขูนํ กีฏาคิริสฺมา ปพฺพาชนียสฺส กมฺมสฺส กรณํ ‘น อสฺสชิปุนพฺพสุเกหิ ภิกฺขูหิ กีฏาคิริสฺมิํ วตฺถพฺพนฺ’ติ, โส ตุณฺหสฺส, ยสฺส นกฺขมติ, โส ภาเสยฺย.}

\prose{\csromanfolio{29}ทุติยมฺปิ เอตมตฺถํ วทามิ ฯเปฯ.\csromanspacer{} ตติยมฺปิ เอตมตฺถํ วทามิ.\csromanspacer{} สุณาตุ เม ภนฺเต สํโฆ, อิเม อสฺสชิปุนพฺพสุกา ภิกฺขู กุลทูสกา ปาปสมาจารา, อิเมสํ ปาปกา สมาจารา ทิสฺสนฺติ เจว สุยฺยนฺติ จ, กุลานิ จ อิเมหิ ทุฏฺฐานิ ทิสฺสนฺติ เจว สุยฺยนฺติ จ, สํโฆ อสฺสชิปุนพฺพสุกานํ ภิกฺขูนํ กีฏาคิริสฺมา ปพฺพาชนียกมฺมํ กโรติ ‘น อสฺสชิปุนพฺพสุเกหิ ภิกฺขูหิ กีฏาคิริสฺมิํ วตฺตพฺพนฺ’ติ, ยสฺสายสฺมโต ขมติ อสฺสชิปุนพฺพสุกานํ ภิกฺขูนํ กีฏาคิริสฺมา ปพฺพาชนียสฺส กมฺมสฺส กรณํ ‘น อสฺสชิปุนพฺพสุเกหิ ภิกฺขูหิ กีฏาคิริสฺมิํ วตฺตพฺพนฺ’ติ, โส ตุณฺหสฺส, ยสฺส นกฺขมติ, โส ภาเสยฺย.}

\prose{กตํ สํเฆน อสฺสชิปุนพฺพสุกานํ ภิกฺขูนํ กีฏาคิริสฺมา ปพฺพาชนียกมฺมํ ‘น อสฺสชิปุนพฺพสุเกหิ ภิกฺขูหิ กีฏาคิริสฺมิํ วตฺตพฺพนฺ’ติ, ขมติ สํฆสฺส, ตสฺมา ตุณฺหี, เอวเมตํ ธารยามี”ติ.}

\csromanmatikaanchor{mtk.17}
\csromantocmarkref{subsection}{อธมฺมกมฺมทฺวาทสก}{mtk.17}
\bookmark[dest={mtk.17},level=2]{อธมฺมกมฺมทฺวาทสก}
\csromanheader{2}{\textbf{อธมฺมกมฺมทฺวาทสก}}
\proseitem{25}{ตีหิ ภิกฺขเว องฺเคหิ สมนฺนาคตํ ปพฺพาชนียกมฺมํ อธมฺมกมฺมญฺจ โหติ อวินยกมฺมญฺจ ทุวูปสนฺตญฺจ.\csromanspacer{} อสมฺมุขา กตํ โหติ, อปฺปฏิปุจฺฉา กตํ โหติ, อปฺปฏิญฺญาย กตํ โหติ.\csromanspacer{} อิเมหิ โข ภิกฺขเว ตีหงฺเคหิ สมนฺนาคตํ ปพฺพาชนียกมฺมํ อธมฺมกมฺมญฺจ โหติ อวินยกมฺมญฺจ ทุวูปสนฺตญฺจ.}

\prose{อปเรหิปิ ภิกฺขเว ตีหงฺเคหิ สมนฺนาคตํ ปพฺพาชนียกมฺมํ อธมฺมกมฺมญฺจ โหติ อวินยกมฺมญฺจ ทุวูปสนฺตญฺจ.\csromanspacer{} อนาปตฺติยา กตํ โหติ, อเทสนาคามินิยา อาปตฺติยา กตํ โหติ, เทสิตาย อาปตฺติยา กตํ โหติ ฯเปฯ.\csromanspacer{} อโจเทตฺวา กตํ โหติ, อสาเรตฺวา กตํ โหติ, อาปตฺติํ อนาโรเปตฺวา กตํ โหติ ฯเปฯ.\csromanspacer{} อสมฺมุขา กตํ โหติ, อธมฺเมน กตํ โหติ, วคฺเคน กตํ โหติ ฯเปฯ.\csromanspacer{} อปฺปฏิปุจฺฉา กตํ โหติ, อธมฺเมน กตํ โหติ, วคฺเคน กตํ โหติ ฯเปฯ.\csromanspacer{} อปฺปฏิญฺญาย กตํ โหติ, อธมฺเมน กตํ โหติ, วคฺเคน กตํ โหติ ฯเปฯ.\csromanspacer{} อนาปตฺติยา กตํ โหติ, อธมฺเมน กตํ โหติ, วคฺเคน กตํ โหติ ฯเปฯ.\csromanspacer{} อเทสนาคามินิยา อาปตฺติยา กตํ โหติ, อธมฺเมน กตํ โหติ, วคฺเคน กตํ โหติ ฯเปฯ.\csromanspacer{} เทสิตาย อาปตฺติยา กตํ โหติ, อธมฺเมน กตํ โหติ, วคฺเคน กตํ โหติ ฯเปฯ.\csromanspacer{} อโจเทตฺวา กตํ โหติ, อธมฺเมน กตํ โหติ, วคฺเคน กตํ \csromanfolio{30}โหติ ฯเปฯ.\csromanspacer{} อสาเรตฺวา กตํ โหติ, อธมฺเมน กตํ โหติ, วคฺเคน กตํ โหติ ฯเปฯ.\csromanspacer{} อาปตฺติํ อนาโรเปตฺวา กตํ โหติ, อธมฺเมน กตํ โหติ, วคฺเคน กตํ โหติ.\csromanspacer{} อิเมหิ โข ภิกฺขเว ตีหงฺเคหิ สมนฺนาคตํ ปพฺพาชนียกมฺมํ อธมฺมกมฺมญฺจ โหติ อวินยกมฺมญฺจ ทุวูปสนฺตญฺจ.}

\nitthitam{อธมฺมกมฺมทฺวาทสกํ นิฏฺฐิตํ.}
\csromanmatikaanchor{mtk.18}
\csromantocmarkref{subsection}{ธมฺมกมฺมทฺวาทสก}{mtk.18}
\bookmark[dest={mtk.18},level=2]{ธมฺมกมฺมทฺวาทสก}
\csromanheader{2}{\textbf{ธมฺมกมฺมทฺวาทสก}}
\proseitem{26}{ตีหิ ภิกฺขเว องฺเคหิ สมนฺนาคตํ ปพฺพาชนียกมฺมํ ธมฺมกมฺมญฺจ โหติ วินยกมฺมญฺจ สุวูปสนฺตญฺจ.\csromanspacer{} สมฺมุขา กตํ โหติ, ปฏิปุจฺฉา กตํ โหติ, ปฏิญฺญาย กตํ โหติ.\csromanspacer{} อิเมหิ โข ภิกฺขเว ตีหงฺเคหิ สมนฺนาคตํ ปพฺพาชนียกมฺมํ ธมฺมกมฺมญฺจ โหติ วินยกมฺมญฺจ สุวูปสนฺตญฺจ.}

\prose{อปเรหิปิ ภิกฺขเว ตีหงฺเคหิ สมนฺนาคตํ ปพฺพาชนียกมฺมํ ธมฺมกมฺมญฺจ โหติ วินยกมฺมญฺจ สุวูปสนฺตญฺจ.\csromanspacer{} อาปตฺติยา กตํ โหติ, เทสนาคามินิยา อาปตฺติยา กตํ โหติ, อเทสิตาย อาปตฺติยา กตํ โหติ ฯเปฯ.\csromanspacer{} โจเทตฺวา กตํ โหติ, สาเรตฺวา กตํ โหติ, อาปตฺติํ อาโรเปตฺวา กตํ โหติ ฯเปฯ.\csromanspacer{} สมฺมุขา กตํ โหติ, ธมฺเมน กตํ โหติ, สมคฺเคน กตํ โหติ ฯเปฯ.\csromanspacer{} ปฏิปุจฺฉา กตํ โหติ, ธมฺเมน กตํ โหติ, สมคฺเคน กตํ โหติ ฯเปฯ.\csromanspacer{} ปฏิญฺญาย กตํ โหติ, ธมฺเมน กตํ โหติ, สมคฺเคน กตํ โหติ ฯเปฯ.\csromanspacer{} อาปตฺติยา กตํ โหติ, ธมฺเมน กตํ โหติ, สมคฺเคน กตํ โหติ ฯเปฯ.\csromanspacer{} เทสนาคามินิยา อาปตฺติยา กตํ โหติ, ธมฺเมน กตํ โหติ, สมคฺเคน กตํ โหติ ฯเปฯ.\csromanspacer{} อเทสิตาย อาปตฺติยา กตํ โหติ, ธมฺเมน กตํ โหติ, สมคฺเคน กตํ โหติ ฯเปฯ.\csromanspacer{} โจเทตฺวา กตํ โหติ, ธมฺเมน กตํ โหติ, สมคฺเคน กตํ โหติ ฯเปฯ.\csromanspacer{} สาเรตฺวา กตํ โหติ, ธมฺเมน กตํ โหติ, สมคฺเคน กตํ โหติ ฯเปฯ.\csromanspacer{} อาปตฺติํ อาโรเปตฺวา กตํ โหติ, ธมฺเมน \csromanfolio{31}กตํ โหติ, สมคฺเคน กตํ โหติ.\csromanspacer{} อิเมหิ โข ภิกฺขเว ตีหงฺเคหิ สมนฺนาคตํ ปพฺพาชนียกมฺมํ ธมฺมกมฺมญฺจ โหติ วินยกมฺมญฺจ สุวูปสนฺตญฺจ.}

\nitthitam{ธมฺมกมฺมทฺวาทสกํ นิฏฺฐิตํ.}
\csromanmatikaanchor{mtk.19}
\csromantocmarkref{subsection}{อากงฺขมานจุทฺทสก}{mtk.19}
\bookmark[dest={mtk.19},level=2]{อากงฺขมานจุทฺทสก}
\csromanheader{2}{\textbf{อากงฺขมานจุทฺทสก}}
\proseitem{27}{\csromansymbolfootnote{*}{วิ 5. 219 ปิฏฺเฐปิ.}ตีหิ ภิกฺขเว องฺเคหิ สมนฺนาคตสฺส ภิกฺขุโน อากงฺขมาโน สํโฆ ปพฺพาชนียกมฺมํ กเรยฺย.\csromanspacer{} ภณฺฑนการโก โหติ กลหการโก วิวาทการโก ภสฺสการโก สํเฆ อธิกรณการโก, พาโล โหติ อพฺยตฺโต อาปตฺติพหุโล อนปทาโน, คิหิสํสฏฺโฐ วิหรติ อนนุโลมิเกหิ คิหิสํสคฺเคหิ.\csromanspacer{} อิเมหิ โข ภิกฺขเว ตีหงฺเคหิ สมนฺนาคตสฺส ภิกฺขุโน อากงฺขมาโน สํโฆ ปพฺพาชนียกมฺมํ กเรยฺย. (1)}

\prose{อปเรหิปิ ภิกฺขเว ตีหงฺเคหิ สมนฺนาคตสฺส ภิกฺขุโน อากงฺขมาโน สํโฆ ปพฺพาชนียกมฺมํ กเรยฺย.\csromanspacer{} อธิสีเล สีลวิปนฺโน โหติ, อชฺฌาจาเร อาจารวิปนฺโน โหติ, อติทิฏฺฐิยา ทิฏฺฐิวิปนฺโน โหติ.\csromanspacer{} อิเมหิ โข ภิกฺขเว ตีหงฺเคหิ สมนฺนาคตสฺส ภิกฺขุโน อากงฺขมาโน สํโฆ ปพฺพาชนียกมฺมํ กเรยฺย. (2)}

\prose{อปเรหิปิ ภิกฺขเว ตีหงฺเคหิ สมนฺนาคตสฺส ภิกฺขุโน อากงฺขมาโน สํโฆ ปพฺพาชนียกมฺมํ กเรยฺย.\csromanspacer{} พุทฺธสฺส อวณฺณํ ภาสติ, ธมฺมสฺส อวณฺณํ ภาสติ, สํฆสฺส อวณฺณํ ภาสติ.\csromanspacer{} อิเมหิ โข ภิกฺขเว ตีหงฺเคหิ สมนฺนาคตสฺส ภิกฺขุโน อากงฺขมาโน สํโฆ ปพฺพาชนียกมฺมํ กเรยฺย. (3)}

\prose{อปเรหิปิ ภิกฺขเว ตีหงฺเคหิ สมนฺนาคตสฺส ภิกฺขุโน อากงฺขมาโน สํโฆ ปพฺพาชนียกมฺมํ กเรยฺย.\csromanspacer{} กายิเกน ทเวน สมนฺนาคโต โหติ, วาจสิเกน ทเวน สมนฺนาคโต โหติ, กายิกวาจสิเกน ทเวน สมนฺนาคโต โหติ.\csromanspacer{} อิเมหิ โข ภิกฺขเว ตีหงฺเคหิ สมนฺนาคตสฺส ภิกฺขุโน อากงฺขมาโน สํโฆ ปพฺพาชนียกมฺมํ กเรยฺย. (4)}

\prose{\csromanfolio{32}อปเรหิปิ ภิกฺขเว ตีหงฺเคหิ สมนฺนาคตสฺส ภิกฺขุโน อากงฺขมาโน สํโฆ ปพฺพาชนียกมฺมํ กเรยฺย.\csromanspacer{} กายิเกน อนาจาเรน สมนฺนาคโต โหติ, วาจสิเกน อนาจาเรน สมนฺนาคโต โหติ, กายิกวาจสิเกน อนาจาเรน สมนฺนาคโต โหติ.\csromanspacer{} อิเมหิ โข ภิกฺขเว ตีหงฺเคหิ สมนฺนาคตสฺส ภิกฺขุโน อากงฺขมาโน สํโฆ ปพฺพาชนียกมฺมํ กเรยฺย. (5)}

\prose{อปเรหิปิ ภิกฺขเว ตีหงฺเคหิ สมนฺนาคตสฺส ภิกฺขุโน อากงฺขมาโน สํโฆ ปพฺพาชนียกมฺมํ กเรยฺย.\csromanspacer{} กายิเกน อุปฆาติเกน สมนฺนาคโต โหติ, วาจสิเกน อุปฆาติเกน สมนฺนาคโต โหติ, กายิกวาจสิเกน อุปฆาติเกน สมนฺนาคโต โหติ.\csromanspacer{} อิเมหิ โข ภิกฺขเว ตีหงฺเคหิ สมนฺนาคตสฺส ภิกฺขุโน อากงฺขมาโน สํโฆ ปพฺพาชนียกมฺมํ กเรยฺย. (6)}

\prose{อปเรหิปิ ภิกฺขเว ตีหงฺเคหิ สมนฺนาคตสฺส ภิกฺขุโน อากงฺขมาโน สํโฆ ปพฺพาชนียกมฺมํ กเรยฺย.\csromanspacer{} กายิเกน มิจฺฉาชีเวน สมนฺนาคโต โหติ, วาจสิเกน มิจฺฉาชีเวน สมนฺนาคโต โหติ, กายิกวาจสิเกน มิจฺฉาชีเวน สมนฺนาคโต โหติ.\csromanspacer{} อิเมหิ โข ภิกฺขเว ตีหงฺเคหิ สมนฺนาคตสฺส ภิกฺขุโน อากงฺขมาโน สํโฆ ปพฺพาชนียกมฺมํ กเรยฺย. (7)}

\prose{ติณฺณํ ภิกฺขเว ภิกฺขูนํ อากงฺขมาโน สํโฆ ปพฺพาชนียกมฺมํ กเรยฺย.\csromanspacer{} เอโก ภณฺฑนการโก โหติ กลหการโก วิวาทการโก ภสฺสการโก สํเฆ อธิกรณการโก, เอโก พาโล โหติ อพฺยตฺโต อาปตฺติพหุโล อนปทาโน, เอโก คิหิสํสฏฺโฐ วิหรติ อนนุโลมิเกหิ คิหิสํสคฺเคหิ.\csromanspacer{} อิเมสํ โข ภิกฺขเว ติณฺณํ ภิกฺขูนํ อากงฺขมาโน สํโฆ ปพฺพาชนียกมฺมํ กเรยฺย. (8)}

\prose{อปเรสมฺปิ ภิกฺขเว ติณฺณํ ภิกฺขูนํ อากงฺขมาโน สํโฆ ปพฺพาชนียกมฺมํ กเรยฺย.\csromanspacer{} เอโก อธิสีเล สีลวิปนฺโน โหติ, เอโก อชฺฌาจาเร อาจารวิปนฺโน โหติ, เอโก อติทิฏฺฐิยา ทิฏฺฐิวิปนฺโน โหติ.\csromanspacer{} อิเมสํ โข ภิกฺขเว ติณฺณํ ภิกฺขูนํ อากงฺขมาโน สํโฆ ปพฺพาชนียกมฺมํ กเรยฺย. (9)}

\prose{\csromanfolio{33}อปเรสมฺปิ ภิกฺขเว ติณฺณํ ภิกฺขูนํ อากงฺขมาโน สํโฆ ปพฺพาชนียกมฺมํ กเรยฺย.\csromanspacer{} เอโก พุทฺธสฺส อวณฺณํ ภาสติ, เอโก ธมฺมสฺส อวณฺณํ ภาสติ, เอโก สํฆสฺส อวณฺณํ ภาสติ.\csromanspacer{} อิเมสํ โข ภิกฺขเว ติณฺณํ ภิกฺขูนํ อากงฺขมาโน สํโฆ ปพฺพาชนียกมฺมํ กเรยฺย. (10)}

\prose{อปเรสมฺปิ ภิกฺขเว ติณฺณํ ภิกฺขูนํ อากงฺขมาโน สํโฆ ปพฺพาชนียกมฺมํ กเรยฺย.\csromanspacer{} เอโก กายิเกน ทเวน สมนฺนาคโต โหติ, เอโก วาจสิเกน ทเวน สมนฺนาคโต โหติ, เอโก กายิกวาจสิเกน ทเวน สมนฺนาคโต โหติ.\csromanspacer{} อิเมสํ โข ภิกฺขเว ติณฺณํ ภิกฺขูนํ อากงฺขมาโน สํโฆ ปพฺพาชนียกมฺมํ กเรยฺย. (11)}

\prose{อปเรสมฺปิ ภิกฺขเว ติณฺณํ ภิกฺขูนํ อากงฺขมาโน สํโฆ ปพฺพาชนียกมฺมํ กเรยฺย.\csromanspacer{} เอโก กายิเกน อนาจาเรน สมนฺนาคโต โหติ, เอโก วาจสิเกน อนาจาเรน สมนฺนาคโต โหติ, เอโก กายิกวาจสิเกน อนาจาเรน สมนฺนาคโต โหติ.\csromanspacer{} อิเมสํ โข ภิกฺขเว ติณฺณํ ภิกฺขูนํ อากงฺขมาโน สํโฆ ปพฺพาชนียกมฺมํ กเรยฺย. (12)}

\prose{อปเรสมฺปิ ภิกฺขเว ติณฺณํ ภิกฺขูนํ อากงฺขมาโน สํโฆ ปพฺพาชนียกมฺมํ กเรยฺย.\csromanspacer{} เอโก กายิเกน อุปฆาติเกน สมนฺนาคโต โหติ, เอโก วาจสิเกน อุปฆาติเกน สมนฺนาคโต โหติ.\csromanspacer{} อิเมสํ โข ภิกฺขเว ติณฺณํ ภิกฺขูนํ อากงฺขมาโน สํโฆ ปพฺพาชนียกมฺมํ กเรยฺย. (13)}

\prose{อปเรสมฺปิ ภิกฺขเว ติณฺณํ ภิกฺขูนํ อากงฺขมาโน สํโฆ ปพฺพาชนิยกมฺมํ กเรยฺย.\csromanspacer{} เอโก กายิเกน มิจฺฉาชีเวน สมนฺนาคโต โหติ, เอโก วาจสิเกน มิจฺฉาชีเวน สมนฺนาคโต โหติ, เอโก กายิกวาจสิเกน มิจฺฉาชีเวน สมนฺนาคโต โหติ.\csromanspacer{} อิเมสํ โข ภิกฺขเว ติณฺณํ ภิกฺขูนํ อากงฺขมาโน สํโฆ ปพฺพาชนียกมฺมํ กเรยฺย. (14)}

\nitthitam{อากงฺขมานจุทฺทสกํ นิฏฺฐิตํ.}
\csromanmatikaanchor{mtk.20}
\csromantocmarkref{subsection}{อฏฺฐารสวตฺต}{mtk.20}
\bookmark[dest={mtk.20},level=2]{อฏฺฐารสวตฺต}
\csromanheader{2}{\textbf{อฏฺฐารสวตฺต}}
\proseitem{28}{\csromanfolio{34}ปพฺพาชนียกมฺมกเตน ภิกฺขเว ภิกฺขุนา สมฺมา วตฺติตพฺพํ.\csromanspacer{} ตตฺรายํ สมฺมาวตฺตนา\csromandash{}น อุปสมฺปาเทตพฺพํ, น นิสฺสโย ทาตพฺโพ, น สามเณโร อุปฏฺฐาเปตพฺโพ, น ภิกฺขุโนวาทกสมฺมุติ สาทิตพฺพา, สมฺมเตนปิ ภิกฺขุนิโย น โอวทิตพฺพา, ยาย อาปตฺติยา สํเฆน ปพฺพาชนียกมฺมํ กตํ โหติ, สา อาปตฺติ น อาปชฺชิตพฺพา, อญฺญา วา ตาทิสิกา, ตโต วา ปาปิฏฺฐตรา, กมฺมํ น ครหิตพฺพํ, กมฺมิกา น ครหิตพฺพา, น ปกตตฺตสฺส ภิกฺขุโน อุโปสโถ ฐเปตพฺโพ, น ปวารณา ฐเปตพฺพา, น สวจนียํ กาตพฺพํ, น อนุวาโท อฏฺฐเปตพฺโพ, น โอกาโส กาเรตพฺโพ, น โจเทตพฺโพ, น สาเรตพฺโพ, น ภิกฺขูหิ สมฺปโยเชตพฺพนฺติ.}

\nitthitam{ปพฺพาชนียกมฺเม อฏฺฐารสวตฺตํ นิฏฺฐิตํ.}
\proseitem{29}{อถ โข สาริปุตฺตโมคฺคลฺลานปฺปมุโข ภิกฺขุสํโฆ กีฏาคิริํ คนฺตฺวา อสฺสชิปุนพฺพสุกานํ ภิกฺขูนํ กีฏาคิริสฺมา ปพฺพาชนียกมฺมํ อกาสิ “น อสฺสชิปุนพฺพสุเกหิ ภิกฺขูหิ กีฏาคิริสฺมิํ วตฺถพฺพนฺ”ติ.\csromanspacer{} เต สํเฆน ปพฺพาชนียกมฺมกตา น สมฺมา วตฺตนฺติ, น โลมํ ปาเตนฺติ, น เนตฺถารํ วตฺตนฺติ, น ภิกฺขู ขมาเปนฺติ, อกฺโกสนฺติ, ปริภาสนฺติ, ฉนฺทคามิตา โทสคามิตา โมหคามิตา ภยคามิตา ปาเปนฺติ, ปกฺกมนฺติปิ วิพฺภมนฺติปิ.\csromanspacer{} เย เต ภิกฺขู อปฺปิจฺฉา ฯเปฯ เต อุชฺฌายนฺติ ขิยฺยนฺติ วิปาเจนฺติ “กถํ หิ นาม อสฺสชิปุนพฺพสุกา ภิกฺขู สํเฆน ปพฺพาชนียกมฺมกตา น สมฺมา วตฺติสฺสนฺติ, น โลมํ ปาเตสฺสนฺติ, น เนตฺถารํ วตฺติสฺสนฺติ, น ภิกฺขู ขมาเปสฺสนฺติ, อกฺโกสิสฺสนฺติ, ปริภาสิสฺสนฺติ, ฉนฺทคามิตา โทสคามิตา โมหคามิตา ภยคามิตา ปาเปสฺสนฺติ, ปกฺกมิสฺสนฺติปิ วิพฺภมิสฺสนฺติปี”ติ.\csromanspacer{} อถ โข เต ภิกฺขู ภควโต เอตมตฺถํ อาโรเจสุํ.}

\prose{อถ โข ภควา เอตสฺมิํ นิทาเน เอตสฺมิํ ปกรเณ ภิกฺขุสํฆํ สนฺนิปาตาเปตฺวา ภิกฺขู ปฏิปุจฺฉิ “สจฺจํ กิร ภิกฺขเว อสฺสชิปุนพฺพสุกา ภิกฺขู สํเฆน ปพฺพาชนียกมฺมกตา น สมฺมา วตฺตนฺติ, น โลมํ ปาเตนฺติ, น เนตฺถารํ วตฺตนฺติ, น ภิกฺขู ขมาเปนฺติ, อกฺโกสนฺติ, ปริภาสนฺติ, ฉนฺทคามิตา โทสคามิตา โมหคามิตา ภยคามิตา ปาเปนฺติ, ปกฺกมนฺติปิ วิพฺภมนฺติปี”ติ.\csromanspacer{} สจฺจํ ภควาติ.\csromanspacer{} วิครหิ พุทฺโธ ภควา \csromanfolio{35}อนนุจฺฉวิกํ ฯเปฯ.\csromanspacer{} กถํ หิ นาม เต ภิกฺขเว โมฆปุริสา สํเฆน ปพฺพาชนียกมฺมกตา น สมฺมา วตฺติสฺสนฺติ, น โลมํ ปาเตสฺสนฺติ, น เนตฺถารํ วตฺติสฺสนฺติ, น ภิกฺขู ขมาเปสฺสนฺติ, อกฺโกสิสฺสนฺติ, ปริภาสิสฺสนฺติ, ฉนฺทคามิตา โทสคามิตา โมหคามิตา ภยคามิตา ปาเปสฺสนฺติ, ปกฺกมิสฺสนฺติปิ วิพฺภมิสฺสนฺติปิ.\csromanspacer{} เนตํ ภิกฺขเว อปฺปสนฺนานํ วา ปสาทาย ฯเปฯ วิครหิตฺวา ฯเปฯ ธมฺมิํ กถํ กตฺวา ภิกฺขู อามนฺเตสิ\csromandash{}เตน หิ ภิกฺขเว สํโฆ ปพฺพาชนียกมฺมํ ปฏิปฺปสฺสมฺเภตุ.}

\csromanmatikaanchor{mtk.21}
\csromantocmarkref{subsection}{นปฺปฏิปฺปสฺสมฺเภตพฺพ-อฏฺฐารสก}{mtk.21}
\bookmark[dest={mtk.21},level=2]{นปฺปฏิปฺปสฺสมฺเภตพฺพ-อฏฺฐารสก}
\csromanheader{2}{\textbf{นปฺปฏิปฺปสฺสมฺเภตพฺพ-อฏฺฐารสก}}
\proseitem{30}{ปญฺจหิ ภิกฺขเว องฺเคหิ สมนฺนาคตสฺส ภิกฺขุโน ปพฺพาชนียกมฺมํ นปฺปฏิปฺปสฺสมฺเภตพฺพํ.\csromanspacer{} อุปสมฺปาเทติ, นิสฺสยํ เทติ, สามเณรํ อุปฏฺฐาเปติ, ภิกฺขุโนวาทกสมฺมุติํ สาทิยติ, สมฺมโตปิ ภิกฺขุนิโย โอวทติ.\csromanspacer{} อิเมหิ โข ภิกฺขเว ปญฺจหงฺเคหิ สมนฺนาคตสฺส ภิกฺขุโน ปพฺพาชนียกมฺมํ นปฺปฏิปฺปสฺสมฺเภตพฺพํ.}

\prose{\csromansymbolfootnote{*}{วิ 5. 317 ปิฏฺเฐปิ.}อปเรหิปิ ภิกฺขเว ปญฺจหงฺเคหิ สมนฺนาคตสฺส ภิกฺขุโน ปพฺพาชนีย กมฺมํ นปฺปฏิปฺปสฺสมฺเภตพฺพํ.\csromanspacer{} ยาย อาปตฺติยา สํเฆน ปพฺพาชนียกมฺมํ กตํ โหติ, ตํ อาปตฺติํ อาปชฺชติ, อญฺญํ วา ตาทิสิกํ, ตโต วา ปาปิฏฺฐตรํ, กมฺมํ ครหติ, กมฺมิเก ครหติ.\csromanspacer{} อิเมหิ โข ภิกฺขเว ปญฺจหงฺเคหิ สมนฺนาคตสฺส ภิกฺขุโน ปพฺพาชนียกมฺมํ นปฺปฏิปฺปสฺสมฺเภตพฺพํ.}

\prose{อฏฺฐหิ ภิกฺขเว องฺเคหิ สมนฺนาคตสฺส ภิกฺขุโน ปพฺพาชนียกมฺมํ นปฺปฏิปฺปสฺสมฺเภตพฺพํ.\csromanspacer{} ปกตตฺตสฺส ภิกฺขุโน อุโปสถํ ฐเปติ, ปวารณํ ฐเปติ, สวจนียํ กโรติ, อนุวาทํ ปฏฺฐเปติ, โอกาสํ กาเรติ, โจเทติ, สาเรติ, ภิกฺขูหิ สมฺปโยเชติ.\csromanspacer{} อิเมหิ โข ภิกฺขเว อฏฺฐหงฺเคหิ สมนฺนาคตสฺส ภิกฺขุโน ปพฺพาชนียกมฺมํ นปฺปฏิปฺปสฺสมฺเภตพฺพํ.}

\nitthitam{ปพฺพาชนียกมฺเม นปฺปฏิปฺปสฺสมฺเภตพฺพ-อฏฺฐารสกํ นิฏฺฐิตํ.}
\csromanmatikaanchor{mtk.22}
\csromantocmarkref{subsection}{ปฏิปฺปสฺสมฺเภตพฺพ-อฏฺฐารสก}{mtk.22}
\bookmark[dest={mtk.22},level=2]{ปฏิปฺปสฺสมฺเภตพฺพ-อฏฺฐารสก}
\csromanheader{2}{\textbf{ปฏิปฺปสฺสมฺเภตพฺพ-อฏฺฐารสก}}
\proseitem{31}{ปญฺจหิ ภิกฺขเว องฺเคหิ สมนฺนาคตสฺส ภิกฺขุโน ปพฺพาชนียกมฺมํ ปฏิปฺปสฺสมฺเภตพฺพํ.\csromanspacer{} น อุปสมฺปาเทติ, น นิสฺสยํ เทติ, น สามเณรํ \csromanfolio{36}อุปฏฺฐาเปติ, น ภิกฺขุโนวาทกสมฺมุติํ สาทิยติ, สมฺมโตปิ ภิกฺขุนิโย น โอวทติ.\csromanspacer{} อิเมหิ โข ภิกฺขเว ปญฺจหงฺเคหิ สมนฺนาคตสฺส ภิกฺขุโน ปพฺพาชนียกมฺมํ ปฏิปฺปสฺสมฺเภตพฺพํ.}

\prose{อปเรหิปิ ภิกฺขเว ปญฺจหงฺเคหิ สมนฺนาคตสฺส ภิกฺขุโน ปพฺพาชนียกมฺมํ ปฏิปฺปสฺสมฺเภตพฺพํ.\csromanspacer{} ยาย อาปตฺติยา สํเฆน ปพฺพาชนียกมฺมํ กตํ โหติ, ตํ อาปตฺติํ น อาปชฺชติ, อญฺญํ วา ตาทิสิกํ, ตโต วา ปาปิฏฺฐตรํ, กมฺมํ น ครหติ, กมฺมิเก น ครหติ.\csromanspacer{} อิเมหิ โข ภิกฺขเว ปญฺจหงฺเคหิ สมนฺนาคตสฺส ภิกฺขุโน ปพฺพาชนียกมฺมํ ปฏิปฺปสฺสมฺเภตพฺพํ.}

\prose{อฏฺฐหิ ภิกฺขเว องฺเคหิ สมนฺนาคตสฺส ภิกฺขุโน ปพฺพาชนียกมฺมํ ปฏิปฺปสฺสมฺเภตพฺพํ.\csromanspacer{} น ปกตตฺตสฺส ภิกฺขุโน อุโปสถํ ฐเปติ, น ปวารณํ ฐเปติ, น สวจนียํ กโรติ, น อนุวาทํ ปฏฺฐเปติ, น โอกาสํ กาเรติ, น โจเทติ, น สาเรติ, น ภิกฺขูหิ สมฺปโยเชติ.\csromanspacer{} อิเมหิ โข ภิกฺขเว อฏฺฐหงฺเคหิ สมนฺนาคตสฺส ภิกฺขุโน ปพฺพาชนียกมฺมํ ปฏิปฺปสฺสมฺเภตพฺพํ.}

\nitthitam{ปพฺพาชนียกมฺเม ปฏิปฺปสฺสมฺเภตพฺพ-อฏฺฐารสกํ นิฏฺฐิตํ.}
\proseitem{32}{เอวญฺจ ปน ภิกฺขเว ปฏิปฺปสฺสมฺเภตพฺพํ, เตน ภิกฺขเว ปพฺพาชนียกมฺมกเตน ภิกฺขุนา สํฆํ อุปสงฺกมิตฺวา เอกํสํ อุตฺตราสงฺคํ กริตฺวา วุฑฺฒานํ ภิกฺขูนํ ปาเท วนฺทิตฺวา อุกฺกุฏิกํ นิสีทิตฺวา อญฺชลิํ ปคฺคเหตฺวา เอวมสฺส วจนีโย “อหํ ภนฺเต สํเฆน ปพฺพาชนียกมฺมกโต สมฺมา วตฺตามิ, โลมํ ปาเตมิ, เนตฺถารํ วตฺตามิ, ปพฺพาชนียสฺส กมฺมสฺส ปฏิปฺปสฺสทฺธิํ ยาจามี”ติ.\csromanspacer{} ทุติยมฺปิ ยาจิตพฺพา.\csromanspacer{} ตติยมฺปิ ยาจิตพฺพา.\csromanspacer{} พฺยตฺเตน ภิกฺขุนา ปฏิพเลน สํโฆ ญาเปตพฺโพ\csromandash{}}

\prose{“สุณาตุ เม ภนฺเต สํโฆ, อยํ อิตฺตนฺนาโม ภิกฺขุ สํเฆน ปพฺพาชนียกมฺมกโต สมฺมา วตฺตติ, โลมํ ปาเตติ, เนตฺถารํ วตฺตติ, ปพฺพาชนียสฺส กมฺมสฺส ปฏิปฺปสฺสทฺธิํ ยาจติ, ยทิ สํฆสฺส ปตฺตกลฺลํ, สํโฆ อิตฺถนฺนามสฺส ภิกฺขุโน ปพฺพาชนียกมฺมํ ปฏิปฺปสฺสมฺเภยฺย, เอสา ญตฺติ.}

\prose{สุณาตุ เม ภนฺเต สํโฆ, อยํ อิตฺถนฺนาโม ภิกฺขุ สํเฆน ปพฺพาชนียกมฺมกโต สมฺมา วตฺตติ, โลมํ ปาเตติ, เนตฺถารํ วตฺตติ, \csromanfolio{37}ปพฺพาชนียสฺส กมฺมสฺส ปฏิปฺปสฺสทฺธิํ ยาจติ, สํโฆ อิตฺถนฺนามสฺส ภิกฺขุโน ปพฺพาชนียกมฺมํ ปฏิปฺปสฺสมฺเภติ, ยสฺสายสฺมโต ขมติ อิตฺถนฺนามสฺส ภิกฺขุโน ปพฺพาชนียสฺส กมฺมสฺส ปฏิปฺปสฺสทฺธิ, โส ตุณฺหสฺส, ยสฺส นกฺขมติ, โส ภาเสยฺย.}

\prose{ทุติยมฺปิ เอตมตฺถํ วทามิ ฯเปฯ.\csromanspacer{} ตติยมฺปิ เอตมตฺถํ วทามิ.\csromanspacer{} สุณาตุ เม ภนฺเต สํโฆ, อยํ อิตฺถนฺนาโม ภิกฺขุ สํเฆน ปพฺพาชนียกมฺมกโต สมฺมา วตฺตติ, โลมํ ปาเตติ, เนตฺถารํ วตฺตติ, ปพฺพาชนียสฺส กมฺมสฺส ปฏิปฺปสฺสทฺธิํ ยาจติ, สํโฆ อิตฺถนฺนามสฺส ภิกฺขุโน ปพฺพาชนียกมฺมํ ปฏิปฺปสฺสมฺเภติ, ยสฺสายสฺมโต ขมติ อิตฺถนฺนามสฺส ภิกฺขุโน ปพฺพาชนียสฺส กมฺมสฺส ปฏิปฺปสฺสทฺธิ, โส ตุณฺหสฺส, ยสฺส นกฺขมติ, โส ภาเสยฺย.}

\prose{ปฏิปฺปสฺสทฺธํ สํเฆน อิตฺถนฺนามสฺส ภิกฺขุโน ปพฺพาชนียกมฺมํ, ขมติ สํฆสฺส, ตสฺมา ตุณฺหี, เอวเมตํ ธารยามี”ติ.}

\nitthitam{ปพฺพาชนียกมฺมํ นิฏฺฐิตํ ตติยํ.}
\csromanmatikaanchor{mtk.23}
\csromantocmarkref{section}{4. ปฏิสารณียกมฺม}{mtk.23}
\bookmark[dest={mtk.23},level=1]{4. ปฏิสารณียกมฺม}
\csromanheader{1}{4. ปฏิสารณียกมฺม}
\proseitem{33}{เตน โข ปน สมเยน อายสฺมา สุธมฺโม มจฺฉิกาสณฺเฑ จิตฺตสฺส คหปติโน อาวาสิโก โหติ นวกมฺมิโก ธุวภตฺติโก.\csromanspacer{} ยทา จิตฺโต คหปติ สํฆํ วา คณํ วา ปุคฺคลํ วา นิมนฺเตตุกาโม โหติ, ตทา น อายสฺมนฺตํ สุธมฺมํ อนปโลเกตฺวา สํฆํ วา คณํ วา ปุคฺคลํ วา นิมนฺเตติ.}

\prose{เตน โข ปน สมเยน สมฺพหุลา เถรา ภิกฺขู อายสฺมา จ สาริปุตฺโต อายสฺมา จ มหาโมคฺคลฺลาโน อายสฺมา จ มหากจฺจาโน อายสฺมา จ มหาโกฏฺฐิโก อายสฺมา จ มหากปฺปิโน อายสฺมา จ มหาจุนฺโท อายสฺมา จ อนุรุทฺโธ อายสฺมา จ เรวโต อายสฺมา จ อุปาลิ อายสฺมา จ อานนฺโท อายสฺมา จ ราหุโล กาสีสุ จาริกํ จรมานา เยน มจฺฉิกาสณฺโฑ ตทวสรุํ.}

\prose{อสฺโสสิ โข จิตฺโต คหปติ “เถรา กิร ภิกฺขู มจฺฉิกาสณฺฑํ อนุปฺปตฺตา”ติ.\csromanspacer{} อถ โข จิตฺโต คหปติ เยน เถรา ภิกฺขู \csromanfolio{38}เตนุปสงฺกมิ, อุปสงฺกมิตฺวา เถเร ภิกฺขู อภิวาเทตฺวา เอกมนฺตํ นิสีทิ, เอกมนฺตํ นิสินฺนํ โข จิตฺตํ คหปติํ อายสฺมา สาริปุตฺโต ธมฺมิยา กถาย สนฺทสฺเสสิ สมาทเปสิ สมุตฺเตเชสิ สมฺปหํเสสิ.}

\prose{อถ โข จิตฺโต คหปติ อายสฺมตา สาริปุตฺเตน ธมฺมิยา กถาย สนฺทสฺสิโต สมาทปิโต สมุตฺเตชิโต สมฺปหํสิโต เถเร ภิกฺขู เอตทโวจ “อธิวาเสนฺตุ เม ภนฺเต เถรา สฺวาตนาย อาคนฺตุกภตฺตนฺ”ติ.\csromanspacer{} อธิวาเสสุํ โข เถรา ภิกฺขู\footnote{อธิวาเสสุํ โข เต เถรา ภิกฺขู (สฺยา)} ตุณฺหีภาเวน.}

\prose{อถ โข จิตฺโต คหปติ เถรานํ ภิกฺขูนํ อธิวาสนํ วิทิตฺวา อุฏฺฐายาสนา เถเร ภิกฺขู อภิวาเทตฺวา ปทกฺขิณํ กตฺวา เยนายสฺมา สุธมฺโม เตนุปสงฺกมิ, อุปสงฺกมิตฺวา อายสฺมนฺตํ สุธมฺมํ อภิวาเทตฺวา เอกมนฺตํ อฏฺฐาสิ, เอกมนฺตํ ฐิโต โข จิตฺโต คหปติ อายสฺมนฺตํ สุธมฺมํ เอตทโวจ “อธิวาเสตุ เม ภนฺเต อยฺโย สุธมฺโม สฺวาตนาย ภตฺตํ สทฺธิํ เถเรหี”ติ.\csromanspacer{} อถ โข อายสฺมา สุธมฺโม “ปุพฺเพ ขฺวายํ จิตฺโต คหปติ ยทา สํฆํ วา คณํ วา ปุคฺคลํ วา นิมนฺเตตุกาโม, น มํ อนปโลเกตฺวา สํฆํ วา คณํ วา ปุคฺคลํ วา นิมนฺเตติ, โสทานิ มํ อนปโลเกตฺวา เถเร ภิกฺขู นิมนฺเตสิ, ทุฏฺโฐทานายํ จิตฺโต คหปติ อนเปกฺโข วิรตฺตรูโป มยี”ติ จิตฺตํ คหปติํ เอตทโวจ “อลํ คหปติ, นาธิวาเสมี”ติ.\csromanspacer{} ทุติยมฺปิ โข ฯเปฯ.\csromanspacer{} ตติยมฺปิ โข จิตฺโต คหปติ อายสฺมนฺตํ สุธมฺมํ เอตทโวจ “อธิวาเสตุ เม ภนฺเต อยฺโย สุธมฺโม สฺวาตนาย ภตฺตํ สทฺธิํ เถเรหี”ติ.\csromanspacer{} อลํ คหปติ, นาธิวาเสมีติ.\csromanspacer{} อถ โข จิตฺโต คหปติ “กิํ เม กริสฺสติ อยฺโย สุธมฺโม อธิวาเสนฺโต วา อนธิวาเสนฺโต วา”ติ อายสฺมนฺตํ สุธมฺมํ อภิวาเทตฺวา ปทกฺขิณํ กตฺวา ปกฺกามิ.}

\proseitem{34}{อถ โข จิตฺโต คหปติ ตสฺสา รตฺติยา อจฺจเยน เถรานํ ภิกฺขูนํ ปณีตํ ขาทนียํ โภชนียํ ปฏิยาทาเปสิ.\csromanspacer{} อถ โข อายสฺมา สุธมฺโม “ยํนูนาหํ จิตฺตสฺส คหปติโน เถรานํ ภิกฺขูนํ ปฏิยตฺตํ ปสฺเสยฺยนฺ”ติ ปุพฺพณฺหสมยํ นิวาเสตฺวา ปตฺตจีวรมาทาย เยน จิตฺตสฺส คหปติโน นิเวสนํ เตนุปสงฺกมิ, อุปสงฺกมิตฺวา ปญฺญตฺเต อาสเน นิสีทิ.\csromanspacer{} อถ โข จิตฺโต คหปติ เยนายสฺมา สุธมฺโม เตนุปสงฺกมิ, \csromanfolio{39}อุปสงฺกมิตฺวา อายสฺมนฺตํ สุธมฺมํ อภิวาเทตฺวา เอกมนฺตํ นิสีทิ, เอกมนฺตํ นิสินฺนํ โข จิตฺตํ คหปติํ อายสฺมา สุธมฺโม เอตทโวจ “ปหูตํ โข เต อิทํ คหปติ ขาทนียํ โภชนียํ ปฏิยตฺตํ, เอกา จ โข อิธ นตฺถิ ยทิทํ ติลสํคุฬิกา”ติ.\csromanspacer{} พหุมฺหิ วต ภนฺเต รตเน พุทฺธวจเน\footnote{ภนฺเต พุทฺธวจเน (สฺยา)} วิชฺชมาเน อยฺเยน สุธมฺเมน ยเทว กิญฺจิ ภาสิตํ ยทิทํ ติลสํคุฬิกาติ.\csromanspacer{} ภูตปุพฺพํ ภนฺเต ทกฺขิณาปถกา วาณิชา ปุรตฺถิมํ ชนปทํ อคมํสุ วาณิชฺชาย, เต ตโต กุกฺกุฏิํ อาเนสุํ.\csromanspacer{} อถ โข สา ภนฺเต กุกฺกุฏี กาเกน สทฺธิํ สํวาสํ กปฺเปสิ, สา โปตกํ ชเนสิ, ยทา โข โส ภนฺเต กุกฺกุฏโปตโก กากวสฺสํ วสฺสิตุกาโม โหติ, “กากกุกฺกุฏี”ติ วสฺสติ.\csromanspacer{} ยทา กุกฺกุฏิวสฺสํ วสฺสิตุกาโม โหติ, “กุกฺกุฏิกากา”ติ วสฺสติ.\csromanspacer{} เอวเมว โข ภนฺเต พหุมฺหิ รตเน พุทฺธวจเน วิชฺชมาเน อยฺเยน สุธมฺเมน ยเทว กิญฺจิ ภาสิตํ ยทิทํ ติลสํคุฬิกาติ.\csromanspacer{} อกฺโกสสิ มํ ตฺวํ คหปติ, ปริภาสสิ มํ ตฺวํ คหปติ, เอโส เต คหปติ อาวาโส, ปกฺกมิสฺสามีติ.\csromanspacer{} นาหํ ภนฺเต อยฺยํ สุธมฺมํ อกฺโกสามิ ปริภาสามิ\footnote{น ปริภาสามิ (สี, สฺยา)}, วสตุ ภนฺเต อยฺโย สุธมฺโม มจฺฉิกาสณฺเฑ รมณียํ อมฺพาฏกวนํ, อหํ อยฺยสฺส สุธมฺมสฺส อุสฺสุกฺกํ กริสฺสามิ จีวรปิณฺฑปาตเสนาสนคิลานปฺปจฺจยเภสชฺชปริกฺขารานนฺติ.\csromanspacer{} ทุติยมฺปิ โข ฯเปฯ.\csromanspacer{} ตติยมฺปิ โข อายสฺมา สุธมฺโม จิตฺตํ คหปติํ เอตทโวจ “อกฺโกสสิ มํ ตฺวํ คหปติ, ปริภาสสิ มํ ตฺวํ คหปติ, เอโส เต คหปติ อาวาโส, ปกฺกมิสฺสามี”ติ.\csromanspacer{} กหํ ภนฺเต อยฺโย สุธมฺโม คมิสฺสตีติ.\csromanspacer{} สาวตฺถิํ โข อหํ คหปติ คมิสฺสามิ ภควนฺตํ ทสฺสนายาติ.\csromanspacer{} เตน หิ ภนฺเต ยญฺจ อตฺตนา ภณิตํ, ยญฺจ มยา ภณิตํ, ตํ สพฺพํ ภควโต อาโรเจหิ, อนจฺฉริยํ โข ปเนตํ ภนฺเต ยํ อยฺโย สุธมฺโม ปุนเทว มจฺฉิกาสณฺฑํ ปจฺจาคจฺเฉยฺยาติ.}

\proseitem{35}{อถ โข อายสฺมา สุธมฺโม เสนาสนํ สํสาเมตฺวา ปตฺตจีวรมาทาย เยน สาวตฺถิ เตน ปกฺกามิ, อนุปุพฺเพน เยน สาวตฺถิ เชตวนํ อนาถปิณฺฑิกสฺส อาราโม, เยน ภควา เตนุปสงฺกมิ, อุปสงฺกมิตฺวา ภควนฺตํ อภิวาเทตฺวา เอกมนฺตํ นิสีทิ, เอกมนฺตํ นิสินฺโน โข อายสฺมา สุธมฺโม ยญฺจ อตฺตนา ภณิตํ, ยญฺจ จิตฺเตน คหปตินา ภณิตํ, ตํ สพฺพํ ภควโต อาโรเจสิ.}

\prose{\csromanfolio{40}วิครหิ พุทฺโธ ภควา อนนุจฺฉวิกํ โมฆปุริส อนนุโลมิกํ อปฺปติรูปํ อสฺสามณกํ อกปฺปิยํ อกรณียํ.\csromanspacer{} กถํ หิ นาม ตฺวํ โมฆปุริส จิตฺตํ คหปติํ สทฺธํ ปสนฺนํ ทายกํ การกํ สํฆุปฏฺฐากํ หีเนน ขุํเสสาสิ, หีเนน วมฺเภสฺสสิ.\csromanspacer{} เนตํ โมฆปุริส อปฺปสนฺนานํ วา ปสาทาย ฯเปฯ วิครหิตฺวา ฯเปฯ ธมฺมิํ กถํ กตฺวา ภิกฺขู อามนฺเตสิ\csromandash{}เตน หิ ภิกฺขเว สํโฆ สุธมฺมสฺส ภิกฺขุโน ปฏิสารณียกมฺมํ กโรตุ “จิตฺโต เต คหปติ ขมาเปตพฺโพ”ติ.\csromanspacer{} เอวญฺจ ปน ภิกฺขเว กาตพฺพํ, ปฐมํ สุธมฺโม ภิกฺขุ โจเทตพฺโพ, โจเทตฺวา สาเรตพฺโพ, สาเรตฺวา อาปตฺติํ อาโรเปตพฺโพ, อาปตฺติํ อาโรเปตฺวา พฺยตฺเตน ภิกฺขุนา ปฏิพเลน สํโฆ ญาเปตพฺโพ\csromandash{}}

\proseitem{36}{“สุณาตุ เม ภนฺเต สํโฆ, อยํ สุธมฺโม ภิกฺขุ จิตฺตํ คหปติํ สทฺธํ ปสนฺนํ ทายกํ การกํ สํฆุปฏฺฐากํ หีเนน ขุํเสติ, หีเนน วมฺเภติ, ยทิ สํฆสฺส ปตฺตกลฺลํ, สํโฆ สุธมฺมสฺส ภิกฺขุโน ปฏิสารณียกมฺมํ กเรยฺย ‘จิตฺโต เต คตปติ ขมาเปตพฺโพ’ติ, เอสา ญตฺติ.}

\prose{สุณาตุ เม ภนฺเต สํโฆ, อยํ สุธมฺโม ภิกฺขุ จิตฺตํ คหปติํ สทฺธํ ปสนฺนํ ทายกํ การกํ สํฆุปฏฺฐากํ หีเนน ขุํเสติ, หีเนน วมฺเภติ, สํโฆ สุธมฺมสฺส ภิกฺขุโน ปฏิสารณียกมฺมํ กโรติ ‘จิตฺโต เต คหปติ ขมาเปตพฺโพ’ติ, ยสฺสายสฺมโต ขมติ สุธมฺมสฺส ภิกฺขุโน ปฏิสารณียสฺส กมฺมสฺส กรณํ ‘จิตฺโต เต คหปติ ขมาเปตพฺโพ’ติ, โส ตุณฺหสฺส, ยสฺส นกฺขมติ, โส ภาเสยฺย.}

\prose{ทุติยมฺปิ เอตมตฺถํ วทามิ ฯเปฯ.\csromanspacer{} ตติยมฺปิ เอตมตฺถํ วทามิ.\csromanspacer{} สุณาตุ เม ภนฺเต สํโฆ, อยํ สุธมฺโม ภิกฺขุ จิตฺตํ คหปติํ สทฺธํ ปสนฺนํ ทายกํ การกํ สํฆุปฏฺฐากํ หีเนน ขุํเสติ, หีเนน วมฺเภติ, สํโฆ สุธมฺมสฺส ภิกฺขุโน ปฏิสารณียกมฺมํ กโรติ ‘จิตฺโต เต คหปติ ขมาเปตพฺโพ’ติ, ยสฺสายสฺมโต ขมติ สุธมฺมสฺส ภิกฺขุโน ปฏิสารณียสฺส กมฺมสฺส กรณํ ‘จิตฺโต เต คหปติ ขมาเปตพฺโพ’ติ, โส ตุณฺหสฺส, ยสฺส นกฺขมติ, โส ภาเสยฺย.}

\prose{กตํ สํเฆน สุธมฺมสฺส ภิกฺขุโน ปฏิสารณียกมฺมํ ‘จิตฺโต เต คหปติ ขมาเปตพฺโพ’ติ, ขมติ สํฆสฺส, ตสฺมา ตุณฺหี, เอวเมตํ ธารยามี”ติ.}

\csromanmatikaanchor{mtk.24}
\csromantocmarkref{subsection}{อธมฺมกมฺมทฺวาทสก}{mtk.24}
\bookmark[dest={mtk.24},level=2]{อธมฺมกมฺมทฺวาทสก}
\csromanheader{2}{\textbf{อธมฺมกมฺมทฺวาทสก}}
\proseitem{37}{\csromanfolio{41}ตีหิ ภิกฺขเว องฺเคหิ สมนฺนาคตํ ปฏิสารณียกมฺมํ อธมฺมกมฺมญฺจ โหติ อวินยกมฺมญฺจ ทุวูปสนฺตญฺจ.\csromanspacer{} อสมฺมุขา กตํ โหติ, อปฺปฏิปุจฺฉา กตํ โหติ, อปฺปฏิญฺญาย กตํ โหติ.\csromanspacer{} อิเมหิ โข ภิกฺขเว ตีหงฺเคหิ สมนฺนาคตํ ปฏิสารณียกมฺมํ อธมฺมกมฺมญฺจ โหติ อวินยกมฺมญฺจ ทุวูปสนฺตญฺจ.}

\prose{อปเรหิปิ ภิกฺขเว ตีหงฺเคหิ สมนฺนาคตํ ปฏิสารณียกมฺมํ อธมฺมกมฺมญฺจ โหติ อวินยกมฺมญฺจ ทุวูปสนฺตญฺจ.\csromanspacer{} อนาปตฺติยา กตํ โหติ, อเทสนาคามินิยา อาปตฺติยา กตํ โหติ, เทสิตาย อาปตฺติยา กตํ โหติ ฯเปฯ.\csromanspacer{} อโจเทตฺวา กตํ โหติ, อสาเรตฺวา กตํ โหติ, อาปตฺติํ อนาโรเปตฺวา กตํ โหติ ฯเปฯ.\csromanspacer{} อสมฺมุขา กตํ โหติ, อธมฺเมน กตํ โหติ, วคฺเคน กตํ โหติ ฯเปฯ.\csromanspacer{} อปฺปฏิปุจฺฉา กตํ โหติ, อธมฺเมน กตํ โหติ, วคฺเคน กตํ โหติ ฯเปฯ.\csromanspacer{} อปฺปฏิญฺญาย กตํ โหติ, อธมฺเมน กตํ โหติ, วคฺเคน กตํ โหติ ฯเปฯ.\csromanspacer{} อนาปตฺติยา กตํ โหติ, อธมฺเมน กตํ โหติ, วคฺเคน กตํ โหติ ฯเปฯ.\csromanspacer{} อเทสนาคามินิยา อาปตฺติยา กตํ โหติ, อธมฺเมน กตํ โหติ, วคฺเคน กตํ โหติ ฯเปฯ.\csromanspacer{} เทสิตาย อาปตฺติยา กตํ โหติ, อธมฺเมน กตํ โหติ, วคฺเคน กตํ โหติ ฯเปฯ.\csromanspacer{} อโจเทตฺวา กตํ โหติ, อธมฺเมน กตํ โหติ, วคฺเคน กตํ โหติ ฯเปฯ.\csromanspacer{} อสาเรตฺวา กตํ โหติ, อธมฺเมน กตํ โหติ, วคฺเคน กตํ โหติ ฯเปฯ.\csromanspacer{} อาปตฺติํ อนาโรเปตฺวา กตํ โหติ, อธมฺเมน กตํ โหติ, วคฺเคน กตํ โหติ.\csromanspacer{} อิเมหิ โข ภิกฺขเว ตีหงฺเคหิ สมนฺนาคตํ ปฏิสารณียกมฺมํ อธมฺมกมฺมญฺจ โหติ อวินยกมฺมญฺจ ทุวูปสนฺตญฺจ.}

\nitthitam{ปฏิสารณียกมฺเม อธมฺมกมฺมทฺวาทสกํ นิฏฺฐิตํ.}
\csromanmatikaanchor{mtk.25}
\csromantocmarkref{subsection}{ธมฺมกมฺมทฺวาทสก}{mtk.25}
\bookmark[dest={mtk.25},level=2]{ธมฺมกมฺมทฺวาทสก}
\csromanheader{2}{\textbf{ธมฺมกมฺมทฺวาทสก}}
\proseitem{38}{ตีหิ ภิกฺขเว องฺเคหิ สมนฺนาคตํ ปฏิสารณียกมฺมํ ธมฺมกมฺมญฺจ โหติ วินยกมฺมญฺจ สุวูปสนฺตญฺจ.\csromanspacer{} สมฺมุขา กตํ โหติ, ปฏิปุจฺฉา กตํ โหติ, ปฏิญฺญาย กตํ โหติ.\csromanspacer{} อิเมหิ โข ภิกฺขเว ตีหงฺเคหิ สมนฺนาคตํ ปฏิสารณียกมฺมํ ธมฺมกมฺมญฺจ โหติ วินยกมฺมญฺจ สุวูปสนฺตญฺจ.}

\prose{\csromanfolio{42}อปเรหิปิ ภิกฺขเว ตีหงฺเคหิ สมนฺนาคตํ ปฏิสารณียกมฺมํ ธมฺมกมฺมญฺจ โหติ วินยกมฺมญฺจ สุวูปสนฺตญฺจ.\csromanspacer{} อาปตฺติยา กตํ โหติ, เทสนาคามินิยา อาปตฺติยา กตํ โหติ, อเทสิตาย อาปตฺติยา กตํ โหติ ฯเปฯ.\csromanspacer{} โจเทตฺวา กตํ โหติ, สาเรตฺวา กตํ โหติ, อาปตฺติํ อาโรเปตฺวา กตํ โหติ ฯเปฯ.\csromanspacer{} สมฺมุขา กตํ โหติ, ธมฺเมน กตํ โหติ, สมคฺเคน กตํ โหติ ฯเปฯ.\csromanspacer{} ปฏิปุจฺฉา กตํ โหติ, ธมฺเมน กตํ โหติ, สมคฺเคน กตํ โหติ ฯเปฯ.\csromanspacer{} ปฏิญฺญาย กตํ โหติ, ธมฺเมน กตํ โหติ, สมคฺเคน กตํ โหติ ฯเปฯ.\csromanspacer{} อาปตฺติยา กตํ โหติ, ธมฺเมน กตํ โหติ, สมคฺเคน กตํ โหติ ฯเปฯ.\csromanspacer{} เทสนาคามินิยา อาปตฺติยา กตํ โหติ, ธมฺเมน กตํ โหติ, สมคฺเคน กตํ โหติ ฯเปฯ.\csromanspacer{} อเทสิตาย อาปตฺติยา กตํ โหติ, ธมฺเมน กตํ โหติ, สมคฺเคน กตํ โหติ ฯเปฯ.\csromanspacer{} โจเทตฺวา กตํ โหติ, ธมฺเมน กตํ โหติ, สมคฺเคน กตํ โหติ ฯเปฯ.\csromanspacer{} สาเรตฺวา กตํ โหติ, ธมฺเมน กตํ โหติ, สมคฺเคน กตํ โหติ ฯเปฯ.\csromanspacer{} อาปตฺติํ อาโรเปตฺวา กตํ โหติ, ธมฺเมน กตํ โหติ, สมคฺเคน กตํ โหติ.\csromanspacer{} อิเมหิ โข ภิกฺขเว ตีหงฺเคหิ สมนฺนาคตํ ปฏิสารณียกมฺมํ ธมฺมกมฺมญฺจ โหติ วินยกมฺมญฺจ สุวูปสนฺตญฺจ.}

\nitthitam{ปฏิสารณียกมฺเม ธมฺมกมฺมทฺวาทสกํ นิฏฺฐิตํ.}
\csromanmatikaanchor{mtk.26}
\csromantocmarkref{subsection}{อากงฺขมานจตุกฺก}{mtk.26}
\bookmark[dest={mtk.26},level=2]{อากงฺขมานจตุกฺก}
\csromanheader{2}{\textbf{อากงฺขมานจตุกฺก}}
\proseitem{39}{ปญฺจหิ ภิกฺขเว องฺเคหิ สมนฺนาคตสฺส ภิกฺขุโน อากงฺขมาโน สํโฆ ปฏิสารณียกมฺมํ กเรยฺย.\csromanspacer{} คิหีนํ อลาภาย ปริสกฺกติ คิหีนํ อนตฺถาย ปริสกฺกติ, คิหีนํ อนาวาสาย\footnote{อวาสาย (สี)} ปริสกฺกติ, คิหี อกฺโกสติ ปริภาสติ, คิหี คิหีหิ เภเทติ.\csromanspacer{} อิเมหิ โข ภิกฺขเว ปญฺจหงฺเคหิ สมนฺนาคตสฺส ภิกฺขุโน อากงฺขมาโน สํโฆ ปฏิสารณียกมฺมํ กเรยฺย. (1)}

\prose{อปเรหิปิ ภิกฺขเว ปญฺจหงฺเคหิ สมนฺนาคตสฺส ภิกฺขุโน อากงฺขมาโน สํโฆ ปฏิสารณียกมฺมํ กเรยฺย.\csromanspacer{} คิหีนํ พุทฺธสฺส อวณฺณํ \csromanfolio{43}ภาสติ, คิหีนํ ธมฺมสฺส อวณฺณํ ภาสติ, คิหีนํ สํฆสฺส อวณฺณํ ภาสติ, คิหี หีเนน ขุํเสติ หีเนน วมฺเภติ, คิหีนํ ธมฺมิกํ ปฏิสฺสวํ น สจฺจาเปติ.\csromanspacer{} อิเมหิ โข ภิกฺขเว ปญฺจหงฺเคหิ สมนฺนาคตสฺส ภิกฺขุโน อากงฺขมาโน สํโฆ ปฏิสารณียกมฺมํ กเรยฺย. (2)}

\prose{ปญฺจนฺนํ ภิกฺขเว ภิกฺขูนํ อากงฺขมาโน สํโฆ ปฏิสารณียกมฺมํ กเรยฺย.\csromanspacer{} เอโก คิหีนํ อลาภาย ปริสกฺกติ, เอโก คิหีนํ อนตฺถาย ปริสกฺกติ, เอโก คิหีนํ อนาวาสาย ปริสกฺกติ, เอโก คิหี อกฺโกสติ ปริภาสติ, เอโก คิหี คิหีหิ เภเทติ.\csromanspacer{} อิเมสํ โข ภิกฺขเว ปญฺจนฺนํ ภิกฺขูนํ อากงฺขมาโน สํโฆ ปฏิสารณียกมฺมํ กเรยฺย. (3)}

\prose{อปเรสมฺปิ ภิกฺขเว ปญฺจนฺนํ ภิกฺขูนํ อากงฺขมาโน สํโฆ ปฏิสารณียกมฺมํ กเรยฺย.\csromanspacer{} เอโก คิหีนํ พุทฺธสฺส อวณฺณํ ภาสติ, เอโก คิหีนํ ธมฺมสฺส อวณฺณํ ภาสติ, เอโก คิหีนํ สํฆสฺส อวณฺณํ ภาสติ, เอโก คิหี หีเนน ขุํเสติ หีเนน วมฺเภติ, เอโก คิหีนํ ธมฺมิกํ ปฏิสฺสวํ น สจฺจาเปติ.\csromanspacer{} อิเมสํ โข ภิกฺขเว ปญฺจนฺนํ ภิกฺขูนํ อากงฺขมาโน สํโฆ ปฏิสารณียกมฺมํ กเรยฺย. (4)}

\nitthitam{อากงฺขมานจตุกฺกํ นิฏฺฐิตํ.}
\csromanmatikaanchor{mtk.27}
\csromantocmarkref{subsection}{อฏฺฐารสวตฺต}{mtk.27}
\bookmark[dest={mtk.27},level=2]{อฏฺฐารสวตฺต}
\csromanheader{2}{\textbf{อฏฺฐารสวตฺต}}
\proseitem{40}{ปฏิสารณียกมฺมกเตน ภิกฺขเว ภิกฺขุนา สมฺมา วตฺติตพฺพํ.\csromanspacer{} ตตฺรายํ สมฺมาวตฺตนา\csromandash{}น อุปสมฺปาเทตพฺพํ, น นิสฺสโย ทาตพฺโพ, น สามเณโร อุปฏฺฐาเปตพฺโพ, น ภิกฺขุโนวาทกสมฺมุติ สาทิตพฺพา, สมฺมเตนปิ ภิกฺขุนิโย น โอวทิตพฺพา, ยาย อาปตฺติยา สํเฆน ปฏิสารณียกมฺมํ กตํ โหติ, สา อาปตฺติ น อาปชฺชิตพฺพา, อญฺญา วา ตาทิสิกา, ตโต วา ปาปิฏฺฐตรา, กมฺมํ น ครหิตพฺพํ, กมฺมิกา น ครหิตพฺพา, น ปกตตฺตสฺส ภิกฺขุโน อุโปสโถ ฐเปตพฺโพ, น ปวารณา ฐเปตพฺพา, น สวจนียํ กาตพฺพํ, น อนุวาโท ปฏฺฐเปตพฺโพ, น โอกาโส กาเรตพฺโพ, น โจเทตพฺโพ, น สาเรตพฺโพ, น ภิกฺขูหิ สมฺปโยเชตพฺพนฺติ.}

\nitthitam{ปฏิสารณียกมฺเม อฏฺฐารสวตฺตํ นิฏฺฐิตํ.}
\proseitem{41}{\csromanfolio{44}อถ โข สํโฆ สุธมฺมสฺส ภิกฺขุโน ปฏิสารณียกมฺมํ อกาสิ ‘จิตฺโต เต คหปติ ขมาเปตพฺโพ’ติ.\csromanspacer{} โส สํเฆน ปฏิสารณียกมฺมกโต มจฺฉิกาสณฺฑํ คนฺตฺวา มงฺกุภูโต นาสกฺขิ จิตฺตํ คหปติํ ขมาเปตุํ, ปุนเทว สาวตฺถิํ ปจฺจาคญฺฉิ.\csromanspacer{} ภิกฺขู เอวมาหํสุ “ขมาปิโต ตยา อาวุโส สุธมฺม จิตฺโต คหปตี”ติ.\csromanspacer{} อิธาหํ อาวุโส มจฺฉิกาสณฺฑํ คนฺตฺวา มงฺกุภูโต นาสกฺขิํ จิตฺตํ คหปติํ ขมาเปตุนฺติ.\csromanspacer{} ภิกฺขู ภควโต เอตมตฺถํ อาโรเจสุํ ฯเปฯ.\csromanspacer{} เตน หิ ภิกฺขเว สํโฆ สุธมฺมสฺส ภิกฺขุโน อนุทูตํ เทตุ จิตฺตํ คหปติํ ขมาเปตุํ.\csromanspacer{} เอวญฺจ ปน ภิกฺขเว ทาตพฺโพ, ปฐมํ ภิกฺขุ ยาจิตพฺโพ, ยาจิตฺวา พฺยตฺเตน ภิกฺขุนา ปฏิพเลน สํโฆ ญาเปตพฺโพ\csromandash{}}

\prose{“สุณาตุ เม ภนฺเต สํโฆ, ยทิ สํฆสฺส ปตฺตกลฺลํ, สํโฆ อิตฺถนฺนามํ ภิกฺขุํ สุธมฺมสฺส ภิกฺขุโน อนุทูตํ ทเทยฺย จิตฺตํ คหปติํ ขมาเปตุํ, เอสา ญตฺติ.}

\prose{สุณาตุ เม ภนฺเต สํโฆ, สํโฆ อิตฺถนฺนามํ ภิกฺขุํ สุธมฺมสฺส ภิกฺขุโน อนุทูตํ เทติ จิตฺตํ คหปติํ ขมาเปตุํ, ยสฺสายสฺมโต ขมติ อิตฺถนฺนามสฺส ภิกฺขุโน สุธมฺมสฺส ภิกฺขุโน อนุทูตสฺส ทานํ จิตฺตํ คหปติํ ขมาเปตุํ, โส ตุณฺหสฺส, ยสฺส นกฺขมติ, โส ภาเสยฺย.}

\prose{ทินฺโน สํเฆน อิตฺถนฺนาโม ภิกฺขุ สุธมฺมสฺส ภิกฺขุโน อนุทูโต จิตฺตํ คหปติํ ขมาเปตุํ, ขมติ สํฆสฺส, ตสฺมา ตุณฺหี, เอวเมตํ ธารยามี”ติ.}

\proseitem{42}{เตน ภิกฺขเว สุธมฺเมน ภิกฺขุนา อนุทูเตน ภิกฺขุนา สทฺธิํ มจฺฉิกาสณฺฑํ คนฺตฺวา จิตฺโต คหปติ ขมาเปตพฺโพ “ขม คหปติ, ปสาเทมิ ตนฺ”ติ.\csromanspacer{} เอวญฺเจ วุจฺจมาโน ขมติ, อิจฺเจตํ กุสลํ, โน เจ ขมติ, อนุทูเตน ภิกฺขุนา วตฺตพฺโพ “ขม คหปติ อิมสฺส ภิกฺขุโน, ปสาเทติ ตนฺ”ติ.\csromanspacer{} เอวญฺเจ วุจฺจมาโน ขมติ, อิจฺเจตํ กุสลํ, โน เจ ขมติ, อนุทูเตน ภิกฺขุนา วตฺตพฺโพ “ขม คหปติ อิมสฺส ภิกฺขุโน, อหํ ตํ ปสาเทมี”ติ.\csromanspacer{} เอวญฺเจ วุจฺจมาโน ขมติ, อิจฺเจตํ กุสลํ, โน เจ ขมติ, อนุทูเตน ภิกฺขุนา วตฺตพฺโพ “ขม คหปติ อิมสฺส ภิกฺขุโน สํฆสฺส วจเนนา”ติ.\csromanspacer{} เอวญฺเจ วุจฺจมาโน ขมติ, อิจฺเจตํ กุสลํ, โน \csromanfolio{45}เจ ขมติ, อนุทูเตน ภิกฺขุนา สุธมฺโม ภิกฺขุ\footnote{สุธมฺมํ ภิกฺขุ ... สา อาปตฺติ เทสาเปตพฺพาติ (สี, สฺยา)} จิตฺตสฺส คหปติโน ทสฺสนูปเจรํ อวิชหาเปตฺวา สวนูปจารํ อวิชหาเปตฺวา เอกํสํ อุตฺตราสงฺคํ การาเปตฺวา อุกฺกุฏิกํ นิสีทาเปตฺวา อญฺชลิํ ปคฺคณฺหาเปตฺวา ตํ อาปตฺติํ เทสาเปตพฺโพติ1.}

\prose{อถ โข อายสฺมา สุธมฺโม อนุทูเตน ภิกฺขุนา สทฺธิํ มจฺฉิกาสณฺฑํ คนฺตฺวา จิตฺตํ คหปติํ ขมาเปสิ, โส สมฺมา วตฺตติ, โลมํ ปาเตติ, เนตฺถารํ วตฺตติ, ภิกฺขู อุปสงฺกํตฺวา เอวํ วเทติ “อหํ อาวุโส สํเฆน ปฏิสารณียกมฺมกโต สมฺมา วตฺตามิ, โลมํ ปาเตมิ, เนตฺถารํ วตฺตามิ, กถํ นุ โข มยา ปฏิปชฺชิตพฺพนฺ”ติ.\csromanspacer{} ภิกฺขู ภควโต เอตมตฺถํ อาโรเจสุํ ฯเปฯ.\csromanspacer{} เตน หิ ภิกฺขเว สํโฆ สุธมฺมสฺส ภิกฺขุโน ปฏิสารณียกมฺมํ ปฏิปฺปสฺสมฺเภตุ.}

\csromanmatikaanchor{mtk.28}
\csromantocmarkref{subsection}{นปฺปฏิปฺปสฺสมฺเภตพฺพ-อฏฺฐารสก}{mtk.28}
\bookmark[dest={mtk.28},level=2]{นปฺปฏิปฺปสฺสมฺเภตพฺพ-อฏฺฐารสก}
\csromanheader{2}{\textbf{นปฺปฏิปฺปสฺสมฺเภตพฺพ-อฏฺฐารสก}}
\proseitem{43}{ปญฺจหิ ภิกฺขเว องฺเคหิ สมนฺนาคตสฺส ภิกฺขุโน ปฏิสารณียกมฺมํ นปฺปฏิปฺปสฺสมฺเภตพฺพํ.\csromanspacer{} อุปสมฺปาเทติ, นิสฺสยํ เทติ, สามเณรํ อุปฏฺฐาเปติ, ภิกฺขุโนวาทกสมฺมุติํ สาทิยติ, สมฺมโตปิ ภิกฺขุนิโย โอวทติ.\csromanspacer{} อิเมหิ โข ภิกฺขเว ปญฺจหงฺเคหิ สมนฺนาคตสฺส ภิกฺขุโน ปฏิสารณียกมฺมํ นปฺปฏิปฺปสฺสมฺเภตพฺพํ.}

\prose{\csromansymbolfootnote{*}{วิ 5. 317 ปิฏฺเฐปิ.}อปเรหิปิ ภิกฺขเว ปญฺจหงฺเคหิ สมนฺนาคตสฺส ภิกฺขุโน ปฏิสารณียกมฺมํ นปฺปฏิปฺปสฺสมฺเภตพฺพํ.\csromanspacer{} ยาย อาปตฺติยา สํเฆน ปฏิสารณียกมฺมํ กตํ โหติ, ตํ อาปตฺติํ อาปชฺชติ, อญฺญํ วา ตาทิสิกํ, ตโต วา ปาปิฏฺฐตรํ, กมฺมํ ครหติ, กมฺมิเก ครหติ.\csromanspacer{} อิเมหิ โข ภิกฺขเว ปญฺจหงฺเคหิ สมนฺนาคตสฺส ภิกฺขุโน ปฏิสารณียกมฺมํ นปฺปฏิปฺปสฺสมฺเภตพฺพํ.}

\prose{อฏฺฐหิ ภิกฺขเว องฺเคหิ สมนฺนาคตสฺส ภิกฺขุโน ปฏิสารณียกมฺมํ นปฺปฏิปฺปสฺสมฺเภตพฺพํ.\csromanspacer{} ปกตตฺตสฺส ภิกฺขุโน อุโปสถํ ฐเปติ, ปวารณํ ฐเปติ, สวจนียํ กโรติ, อนุวาทํ ปฏฺฐเปติ, โอกาสํ กาเรติ, \csromanfolio{46}โจเทติ, สาเรติ, ภิกฺขูหิ สมฺปโยเชติ.\csromanspacer{} อิเมหิ โข ภิกฺขเว อฏฺฐหงฺเคหิ สมนฺนาคตสฺส ภิกฺขุโน ปฏิสารณียกมฺมํ นปฺปฏิปฺปสฺสมฺเภตพฺพํ.\csromanspacer{} ปฏิสารณียกมฺเม นปฺปฏิปฺปสฺสมฺเภตพฺพ-อฏฺฐารสกํ นิฏฺฐิตํ.}

\csromanmatikaanchor{mtk.29}
\csromantocmarkref{subsection}{ปฏิปฺปสฺสมฺเภตพฺพ-อฏฺฐารสก}{mtk.29}
\bookmark[dest={mtk.29},level=2]{ปฏิปฺปสฺสมฺเภตพฺพ-อฏฺฐารสก}
\csromanheader{2}{\textbf{ปฏิปฺปสฺสมฺเภตพฺพ-อฏฺฐารสก}}
\proseitem{44}{ปญฺจหิ ภิกฺขเว องฺเคหิ สมนฺนาคตสฺส ภิกฺขุโน ปฏิสารณียกมฺมํ ปฏิปฺปสฺสมฺเภตพฺพํ.\csromanspacer{} น อุปสมฺปาเทติ, น นิสฺสยํ เทติ, น สามเณรํ อุปฏฺฐาเปติ, น ภิกฺขุโนวาทกสมฺมุติํ สาทิยติ, สมฺมโตปิ ภิกฺขุนิโย น โอวทติ.\csromanspacer{} อิเมหิ โข ภิกฺขเว ปญฺจหงฺเคหิ สมนฺนาคตสฺส ภิกฺขุโน ปฏิสารณียกมฺมํ ปฏิปฺปสฺสมฺเภตพฺพํ.}

\prose{อปเรหิปิ ภิกฺขเว ปญฺจหงฺเคหิ สมนฺนาคตสฺส ภิกฺขุโน ปฏิสารณียกมฺมํ ปฏิปฺปสฺสมฺเภตพฺพํ.\csromanspacer{} ยาย อาปตฺติยา สํเฆน ปฏิสารณียกมฺมํ กตํ โหติ, ตํ อาปตฺติํ น อาปชฺชติ, อญฺญํ วา ตาทิสิกํ, ตโต วา ปาปิฏฺฐตรํ, กมฺมํ น ครหติ, กมฺมิเก น ครหติ.\csromanspacer{} อิเมหิ โข ภิกฺขเว ปญฺจหงฺเคหิ สมนฺนาคตสฺส ภิกฺขุโน ปฏิสารณียกมฺมํ ปฏิปฺปสฺสมฺเภตพฺพํ.}

\prose{อฏฺฐหิ ภิกฺขเว องฺเคหิ สมนฺนาคตสฺส ภิกฺขุโน ปฏิสารณียกมฺมํ ปฏิปฺปสฺสมฺเภตพฺพํ.\csromanspacer{} น ปกตตฺตสฺส ภิกฺขุโน อุโปสถํ ฐเปติ, น ปวารณํ ฐเปติ, น สวจนียํ กโรติ, น อนุวาทํ ปฏฺฐเปติ, น โอกาสํ กาเรติ, น โจเทติ, น สาเรติ, น ภิกฺขูหิ สมฺปโยเชติ.\csromanspacer{} อิเมหิ โข ภิกฺขเว อฏฺฐหงฺเคหิ สมนฺนาคตสฺส ภิกฺขุโน ปฏิสารณียกมฺมํ ปฏิปฺปสฺสมฺเภตพฺพํ.\csromanspacer{} ปฏิสารณียกมฺเม ปฏิปฺปสฺสมฺเภตพฺพ-อฏฺฐารสกํ นิฏฺฐิตํ.}

\proseitem{45}{เอวญฺจ ปน ภิกฺขเว ปฏิปฺปสฺสมฺเภตพฺพํ, เตน\footnote{เตน หิ (สฺยา, ก)} ภิกฺขเว สุธมฺเมน ภิกฺขุนา สํฆํ อุปสงฺกมิตฺวา เอกํสํ อุตฺตราสงฺคํ กริตฺวา วุฑฺฒานํ ภิกฺขูนํ ปาเท วนฺทิตฺวา อุกฺกุฏิกํ นิสีทิตฺวา อญฺชลิํ ปคฺคเหตฺวา เอวมสฺส วจนีโย “อหํ ภนฺเต สํเฆน ปฏิสารณียกมฺมกโต สมฺมา วตฺตามิ, โลมํ ปาเตมิ, เนตฺถารํ วตฺตามิ, ปฏิสารณียสฺส กมฺมสฺส ปฏิปฺปสฺสทฺธิํ \csromanfolio{47}ยาจามี”ติ.\csromanspacer{} ทุติยมฺปิ ยาจิตพฺพา.\csromanspacer{} ตติยมฺปิ ยาจิตพฺพา.\csromanspacer{} พฺยตฺเตน ภิกฺขุนา ปฏิพเลน สํโฆ ญาเปตพฺโพ\csromandash{}}

\prose{“สุณาตุ เม ภนฺเต สํโฆ, อยํ สุธมฺโม ภิกฺขุ สํเฆน ปฏิสารณียกมฺมกโต สมฺมา วตฺตติ, โลมํ ปาเตติ, เนตฺถารํ วตฺตติ, ปฏิสารณียสฺส กมฺมสฺส ปฏิปฺปสฺสทฺธิํ ยาจติ, ยทิ สํฆสฺส ปตฺตกลฺลํ, สํโฆ สุธมฺมสฺส ภิกฺขุโน ปฏิสารณียกมฺมํ ปฏิปฺปสฺสมฺเภยฺย, เอสา ญตฺติ.}

\prose{สุณาตุ เม ภนฺเต สํโฆ, อยํ สุธมฺโม ภิกฺขุ สํเฆน ปฏิสารณียกมฺมกโต สมฺมา วตฺตติ, โลมํ ปาเตติ, เนตฺถารํ วตฺตติ, ปฏิสารณียสฺส กมฺมสฺส ปฏิปฺปสฺสทฺธิํ ยาจติ, สํโฆ สุธมฺมสฺส ภิกฺขุโน ปฏิสารณียกมฺมํ ปฏิปฺปสฺสมฺเภติ, ยสฺสายสฺมโต ขมติ สุธมฺมสฺส ภิกฺขุโน ปฏิสารณียสฺส กมฺมสฺส ปฏิปฺปสฺสทฺธิ, โส ตุณฺหสฺส, ยสฺส นกฺขมติ, โส ภาเสยฺย.}

\prose{ทุติยมฺปิ เอตมตฺถํ วทามิ ฯเปฯ.\csromanspacer{} ตติยมฺปิ เอตมตฺถํ วทามิ.\csromanspacer{} สุณาตุ เม ภนฺเต สํโฆ, อยํ สุธมฺโม ภิกฺขุ สํเฆน ปฏิสารณียกมฺมกโต สมฺมา วตฺตติ, โลมํ ปาเตติ, เนตฺถารํ วตฺตติ, ปฏิสารณียสฺส กมฺมสฺส ปฏิปฺปสฺสทฺธิํ ยาจติ, สํโฆ สุธมฺมสฺส ภิกฺขุโน ปฏิสารณียกมฺมํ ปฏิปฺปสฺสมฺเภติ, ยสฺสายสฺมโต ขมติ สุธมฺมสฺส ภิกฺขุโน ปฏิสารณียสฺส กมฺมสฺส ปฏิปฺปสฺสทฺธิ, โส ตุณฺหสฺส, ยสฺส นกฺขมติ, โส ภาเสยฺย.}

\prose{ปฏิปฺปสฺสทฺธํ สํเฆน สุธมฺมสฺส ภิกฺขุโน ปฏิสารณียกมฺมํ, ขมติ สํฆสฺส, ตสฺมา ตุณฺหี, เอวเมตํ ธารยามี”ติ.}

\nitthitam{ปฏิสารณียกมฺมํ นิฏฺฐิตํ จตุตฺถํ.}
\csromanmatikaanchor{mtk.30}
\csromantocmarkref{section}{5. อาปตฺติยา อทสฺสเน อุกฺเขปนียกมฺม}{mtk.30}
\bookmark[dest={mtk.30},level=1]{5. อาปตฺติยา อทสฺสเน อุกฺเขปนียกมฺม}
\csromanheader{1}{5. อาปตฺติยา อทสฺสเน อุกฺเขปนียกมฺม}
\proseitem{46}{เตน สมเยน พุทฺโธ ภควา โกสมฺพิยํ วิหรติ โฆสิตาราเม.\csromanspacer{} เตน โข ปน สมเยน อายสฺมา ฉนฺโน อาปตฺติํ อาปชฺชิตฺวา น อิจฺฉติ อาปตฺติํ ปสฺสิตุํ.\csromanspacer{} เย เต ภิกฺขู อปฺปิจฺฉา ฯเปฯ เต \csromanfolio{48}อุชฺฌายนฺติ ขิยฺยนฺติ วิปาเจนฺติ “กถํ หิ นาม อายสฺมา ฉนฺโน อาปตฺติํ อาปชฺชิตฺวา น อิจฺฉิสฺสติ อาปตฺติํ ปสฺสิตุนฺ”ติ.\csromanspacer{} อถ โข เต ภิกฺขู ภควโต เอตมตฺถํ อาโรเจสุํ.}

\prose{อถ โข ภควา เอตสฺมิํ นิทาเน เอตสฺมิํ ปกรเณ ภิกฺขุสํฆํ สนฺนิปาตาเปตฺวา ภิกฺขู ปฏิปุจฺฉิ “สจฺจํ กิร ภิกฺขเว ฉนฺโน ภิกฺขุ อาปตฺติํ อาปชฺชิตฺวา น อิจฺฉติ อาปตฺติํ ปสฺสิตุนฺ”ติ.\csromanspacer{} สจฺจํ ภควาติ.\csromanspacer{} วิครหิ พุทฺโธ ภควา อนนุจฺฉวิกํ ฯเปฯ.\csromanspacer{} กถํ หิ นาม โส ภิกฺขเว โมฆปุริโส อาปตฺติํ อาปชฺชิตฺวา น อิจฺฉิสฺสติ อาปตฺติํ ปสฺสิตุํ, เนตํ ภิกฺขเว อปฺปสนฺนานํ วา ปสฺสทาย ฯเปฯ วิครหิตฺวา ฯเปฯ ธมฺมิํ กถํ กตฺวา ภิกฺขู อามนฺเตสิ\csromandash{}เตน หิ ภิกฺขเว สํโฆ ฉนฺนสฺส ภิกฺขุโน อาปตฺติยา อทสฺสเน อุกฺเขปนียกมฺมํ กโรตุ อสมฺโภคํ สํเฆน.\csromanspacer{} เอวญฺจ ปน ภิกฺขเว กาตพฺพํ, ปฐมํ ฉนฺโน ภิกฺขุ โจเทตพฺโพ, โจเทตฺวา สาเรตพฺโพ, สาเรตฺวา อาปตฺติํ อาโรเปตพฺโพ, อาปตฺติํ อาโรเปตฺวา พฺยตฺเตน ภิกฺขุนา ปฏิพเลน สํโฆ ญาเปตพฺโพ\csromandash{}}

\proseitem{47}{“สุณาตุ เม ภนฺเต สํโฆ, อยํ ฉนฺโน ภิกฺขุ อาปตฺติํ อาปชฺชิตฺวา น อิจฺฉติ อาปตฺติํ ปสฺสิตุํ, ยทิ สํฆสฺส ปตฺตกลฺลํ, สํโฆ ฉนฺนสฺส ภิกฺขุโน อาปตฺติยา อทสฺสเน อุกฺเขปนียกมฺมํ กเรยฺย อสมฺโภคํ สํเฆน, เอสา ญตฺติ.}

\prose{สุณาตุ เม ภนฺเต สํโฆ, อยํ ฉนฺโน ภิกฺขุ อาปตฺติํ อาปชฺชิตฺวา น อิจฺฉติ อาปตฺติํ ปสฺสิตุํ, สํโฆ ฉนฺนสฺส ภิกฺขุโน อาปตฺติยา อทสฺสเน อุกฺเขปนียกมฺมํ กโรติ อสมฺโภคํ สํเฆน, ยสฺสายสฺมโต ขมติ ฉนฺนสฺส ภิกฺขุโน อาปตฺติยา อทสฺสเน อุกฺเขปนียสฺส กมฺมสฺส กรณํ อสมฺโภคํ สํเฆน, โส ตุณฺหสฺส, ยสฺส นกฺขมติ, โส ภาเสยฺย.}

\prose{ทุติยมฺปิ เอตมตฺถํ วทามิ ฯเปฯ.\csromanspacer{} ตติยมฺปิ เอตมตฺถํ วทามิ.\csromanspacer{} สุณาตุ เม ภนฺเต สํโฆ, อยํ ฉนฺโน ภิกฺขุ อาปตฺติํ อาปชฺชิตฺวา น อิจฺฉติ อาปตฺติํ ปสฺสิตุํ.\csromanspacer{} สํโฆ ฉนฺนสฺส ภิกฺขุโน อาปตฺติยา อทสฺสเน อุกฺเขปนียกมฺมํ กโรติ อสมฺโภคํ สํเฆน, ยสฺสายสฺมโต ขมติ ฉนฺนสฺส ภิกฺขุโน อาปตฺติยา อทสฺสเน อุกฺเขปนียสฺส กมฺมสฺส \csromanfolio{49}กรณํ อสมฺโภคํ สํเฆน, โส ตุณฺหสฺส, ยสฺส นกฺขมติ, โส ภาเสยฺย.}

\prose{กตํ สํเฆน ฉนฺนสฺส ภิกฺขุโน อาปตฺติยา อทสฺสเน อุกฺเขปนียกมฺมํ อสมฺโภคํ สํเฆน, ขมติ สํฆสฺส, ตสฺมา ตุณฺหี, เอวเมตํ ธารยามี”ติ.}

\prose{อาวาสปรมฺปรญฺจ ภิกฺขเว สํสถ “ฉนฺโน ภิกฺขุ สํเฆน อาปตฺติยา\footnote{ฉนฺโน ภิกฺขุ อาปตฺติยา (สี, ก)} อทสฺสเน อุกฺเขปนียกมฺมกโต อสมฺโภคํ สํเฆนา”ติ.}

\csromanmatikaanchor{mtk.31}
\csromantocmarkref{subsection}{อธมฺมกมฺมทฺวาทสก}{mtk.31}
\bookmark[dest={mtk.31},level=2]{อธมฺมกมฺมทฺวาทสก}
\csromanheader{2}{\textbf{อธมฺมกมฺมทฺวาทสก}}
\proseitem{48}{ตีหิ ภิกฺขเว องฺเคหิ สมนฺนาคตํ อาปตฺติยา อทสฺสเน อุกฺเขปนียกมฺมํ อธมฺมกมฺมญฺจ โหติ อวินยกมฺมญฺจ ทุวูปสนฺตญฺจ.\csromanspacer{} อสมฺมุขา กตํ โหติ, อปฺปฏิปุจฺฉา กตํ โหติ, อปฺปฏิญฺญาย กตํ โหติ.\csromanspacer{} อิเมหิ โข ภิกฺขเว ตีหงฺเคหิ สมนฺนาคตํ อาปตฺติยา อทสาเน อุกฺเขปนียกมฺมํ อธมฺมกมฺมญฺจ โหติ อวินยกมฺมญฺจ ทุวูปสนฺตญฺจ.}

\prose{อปเรหิปิ ภิกฺขเว ตีหงฺเคหิ สมนฺนาคตํ อาปตฺติยา อทสฺสเน อุกฺเขปนียกมฺมํ อธมฺมกมฺมญฺจ โหติ อวินยกมฺมญฺจ ทุวูปสนฺตญฺจ.\csromanspacer{} อนาปตฺติยา กตํ โหติ, อเทสนาคามินิยา อาปตฺติยา กตํ โหติ, เทสิตาย อาปตฺติยา กตํ โหติ ฯเปฯ.\csromanspacer{} อโจเทตฺวา กตํ โหติ, อสาเรตฺวา กตํ โหติ, อาปตฺติํ อนาโรเปตฺวา กตํ โหติ ฯเปฯ.\csromanspacer{} อสมฺมุขา กตํ โหติ, อธมฺเมน กตํ โหติ, วคฺเคน กตํ โหติ ฯเปฯ.\csromanspacer{} อปฺปฏิปุจฺฉา กตํ โหติ, อธมฺเมน กตํ โหติ, วคฺเคน กตํ โหติ ฯเปฯ.\csromanspacer{} อปฺปฏิญฺญาย กตํ โหติ, อธมฺเมน กตํ โหติ, วคฺเคน กตํ โหติ ฯเปฯ.\csromanspacer{} อนาปตฺติยา กตํ โหติ, อธมฺเมน กตํ โหติ, วคฺเคน กตํ โหติ ฯเปฯ.\csromanspacer{} อเทสนาคามินิยา อาปตฺติยา กตํ โหติ, อธมฺเมน กตํ โหติ, วคฺเคน กตํ โหติ ฯเปฯ.\csromanspacer{} เทสิตาย อาปตฺติยา กตํ โหติ, อธมฺเมน กตํ โหติ, วคฺเคน กตํ โหติ ฯเปฯ.\csromanspacer{} อโจเทตฺวา กตํ โหติ, อธมฺเมน กตํ โหติ, วคฺเคน กตํ โหติ ฯเปฯ.\csromanspacer{} อสาเรตฺวา กตํ โหติ, อธมฺเมน กตํ โหติ, วคฺเคน กตํ โหติ ฯเปฯ.\csromanspacer{} อาปตฺติํ \csromanfolio{50}อนาโรเปตฺวา กตํ โหติ, อธมฺเมน กตํ โหติ, วคฺเคน กตํ โหติ.\csromanspacer{} อิเมหิ โข ภิกฺขเว ตีหงฺเคหิ สมนฺนาคตํ อาปตฺติยา อทสฺสเน อุกฺเขปนียกมฺมํ อธมฺมกมฺมญฺจ โหติ อวินยกมฺมญฺจ ทุวูปสนฺตญฺจ.}

\nitthitam{อาปตฺติยา อทสฺสเน อุกฺเขปนียกมฺเม อธมฺมกมฺมทฺวาทสกํ นิฏฺฐิตํ.}
\csromanmatikaanchor{mtk.32}
\csromantocmarkref{subsection}{ธมฺมกมฺมทฺวาทสก}{mtk.32}
\bookmark[dest={mtk.32},level=2]{ธมฺมกมฺมทฺวาทสก}
\csromanheader{2}{\textbf{ธมฺมกมฺมทฺวาทสก}}
\proseitem{49}{ตีหิ ภิกฺขเว องฺเคหิ สมนฺนาคตํ อาปตฺติยา อทสฺสเน อุกฺเขปนียกมฺมํ ธมฺมกมฺมญฺจ โหติ วินยกมฺมญฺจ สุวูปสนฺตญฺจ.\csromanspacer{} สมฺมุขา กตํ โหติ, ปฏิปุจฺฉา กตํ โหติ, ปฏิญฺญาย กตํ โหติ.\csromanspacer{} อิเมหิ โข ภิกฺขเว ตีหงฺเคหิ สมนฺนาคตํ อาปตฺติยา อทสฺสเน อุกฺเขปนียกมฺมํ ธมฺมกมฺมญฺจ โหติ วินยกมฺมญฺจ สุวูปสนฺตญฺจ.}

\prose{อปเรหิปิ ภิกฺขเว ตีหงฺเคหิ สมนฺนาคตํ ฯเปฯ.\csromanspacer{} อาปตฺติยา กตํ โหติ, เทสนาคามินิยา อาปตฺติยา กตํ โหติ, อเทสิตาย อาปตฺติยา กตํ โหติ ฯเปฯ.\csromanspacer{} โจเทตฺวา กตํ โหติ, สาเรตฺวา กตํ โหติ, อาปตฺติํ อาโรเปตฺวา กตํ โหติ ฯเปฯ.\csromanspacer{} สมฺมุขา กตํ โหติ, ธมฺเมน กตํ โหติ, สมคฺเคน กตํ โหติ ฯเปฯ.\csromanspacer{} ปฏิปุจฺฉา กตํ โหติ, ธมฺเมน กตํ โหติ, สมคฺเคน กตํ โหติ ฯเปฯ.\csromanspacer{} ปฏิญฺญาย กตํ โหติ, ธมฺเมน กตํ โหติ, สมคฺเคน กตํ โหติ ฯเปฯ.\csromanspacer{} อาปตฺติยา กตํ โหติ, ธมฺเมน กตํ โหติ, สมคฺเคน กตํ โหติ ฯเปฯ.\csromanspacer{} เทสนาคามินิยา อาปตฺติยา กตํ โหติ, ธมฺเมน กตํ โหติ, สมคฺเคน กตํ โหติ ฯเปฯ.\csromanspacer{} อเทสิตาย อาปตฺติยา กตํ โหติ, ธมฺเมน กตํ โหติ, สมคฺเคน กตํ โหติ ฯเปฯ.\csromanspacer{} โจเทตฺวา กตํ โหติ, ธมฺเมน กตํ โหติ, สมคฺเคน กตํ โหติ ฯเปฯ.\csromanspacer{} สาเรตฺวา กตํ โหติ, ธมฺเมน กตํ โหติ, สมคฺเคน กตํ โหติ ฯเปฯ.\csromanspacer{} อาปตฺติํ อาโรเปตฺวา กตํ โหติ, ธมฺเมน กตํ โหติ, สมคฺเคน กตํ โหติ.\csromanspacer{} อิเมหิ โข ภิกฺขเว ตีหงฺเคหิ สมนฺนาคตํ อาปตฺติยา อทสฺสเน อุกฺเขปนียกมฺมํ ธมฺมกมฺมญฺจ โหติ วินยกมฺมญฺจ สุวูปสนฺตญฺจ.}

\nitthitam{อาปตฺติยา อทสฺสเน อุกฺเขปนียกมฺเม ธมฺมกมฺมทฺวาทสกํ นิฏฺฐิตํ.}
\csromanmatikaanchor{mtk.33}
\csromantocmarkref{subsection}{อากงฺขมานฉกฺก}{mtk.33}
\bookmark[dest={mtk.33},level=2]{อากงฺขมานฉกฺก}
\csromanheader{2}{\textbf{อากงฺขมานฉกฺก}}
\proseitem{50}{\csromanfolio{51}\csromansymbolfootnote{*}{วิ 5. 220 ปิฏฺเฐปิ.}ตีหิ ภิกฺขเว องฺเคหิ สมนฺนาคตสฺส ภิกฺขุโน อากงฺขมาโน สํโฆ อาปตฺติยา อทสฺสเน อุกฺเขปนียกมฺมํ กเรยฺย.\csromanspacer{} ภณฺฑนการโก โหติ กลหการโก วิวาทการโก ภสฺสการโก สํเฆ อธิกรณการโก, พาโล โหติ อพฺยตฺโต อาปตฺติพหุโล อนปทาโน, คิหิสํสฏฺโฐ วิหรติ อนนุโลมิเกหิ คิหิสํสคฺเคหิ.\csromanspacer{} อิเมหิ โข ภิกฺขเว ตีหงฺเคหิ สมนฺนาคตสฺส ภิกฺขุโน อากงฺขมาโน สํโฆ อาปตฺติยา อทสฺสเน อุกฺเขปนียกมฺมํ กเรยฺย.}

\prose{อปเรหิปิ ภิกฺขเว ตีหงฺเคหิ สมนฺนาคตสฺส ภิกฺขุโน อากงฺขมาโน สํโฆ อาปตฺติยา อทสฺสเน อุกฺเขปนียกมฺมํ กเรยฺย.\csromanspacer{} อธิสีเล สีลวิปนฺโน โหติ, อชฺฌาจาเร อาจารวิปนฺโน โหติ, อติทิฏฺฐิยา ทิฏฺฐิวิปนฺโน โหติ.\csromanspacer{} อิเมหิ โข ภิกฺขเว ตีหงฺเคหิ สมนฺนาคตสฺส ภิกฺขุโน อากงฺขมาโน สํโฆ อาปตฺติยา อทสฺสเน อุกฺเขปนียกมฺมํ กเรยฺย.}

\prose{อปเรหิปิ ภิกฺขเว ตีหงฺเคหิ สมนฺนาคตสฺส ภิกฺขุโน อากงฺขมาโน สํโฆ อาปตฺติยา อทสฺสเน อุกฺเขปนียกมฺมํ กเรยฺย.\csromanspacer{} พุทฺธสฺส อวณฺณํ ภาสติ, ธมฺมสฺส อวณฺณํ ภาสติ, สํฆสฺส อวณฺณํ ภาสติ.\csromanspacer{} อิเมหิ โข ภิกฺขเว ตีหงฺเคหิ สมนฺนาคตสฺส ภิกฺขุโน อากงฺขมาโน สํโฆ อาปตฺติยา อทสฺสเน อุกฺเขปนียกมฺมํ กเรยฺย.}

\prose{ติณฺณํ ภิกฺขเว ภิกฺขูนํ อากงฺขมาโน สํโฆ อาปตฺติยา อทสฺสเน อุกฺเขปนียกมฺมํ กเรยฺย.\csromanspacer{} เอโก ภณฺฑนการโก โหติ กลหการโก วิวาทการโก ภสฺสการโก สํเฆ อธิกรณการโก, เอโก พาโล โหติ อพฺยตฺโต อาปตฺติพหุโล อนปทาโน, เอโก คิหิสํสฏฺโฐ วิหรติ อนนุโลมิเกหิ คิหิสํสคฺเคหิ.\csromanspacer{} อิเมสํ โข ภิกฺขเว ติณฺณํ ภิกฺขูนํ อากงฺขมาโน สํโฆ อาปตฺติยา อทสฺสเน อุกฺเขปนียกมฺมํ กเรยฺย.}

\prose{อปเรสมฺปิ ภิกฺขเว ติณฺณํ ภิกฺขูนํ อากงฺขมาโน สํโฆ อาปตฺติยา อทสฺสเน อุกฺเขปนียกมฺมํ กเรยฺย.\csromanspacer{} เอโก อธิสีเล สีลวิปนฺโน \csromanfolio{52}โหติ, เอโก อชฺฌาจาเร อาจารวิปนฺโน โหติ, เอโก อติทิฏฺฐิยา ทิฏฺฐิวิปนฺโน โหติ.\csromanspacer{} อิเมสํ โข ภิกฺขเว ติณฺณํ ภิกฺขูนํ อากงฺขมาโน สํโฆ อาปตฺติยา อทสฺสเน อุกฺเขปนียกมฺมํ กเรยฺย.}

\prose{อปเรสมฺปิ ภิกฺขเว ติณฺณํ ภิกฺขูนํ อากงฺขมาโน สํโฆ อาปตฺติยา อทสฺสเน อุกฺเขปนียกมฺมํ กเรยฺย.\csromanspacer{} เอโก พุทฺธสฺส อวณฺณํ ภาสติ, เอโก ธมฺมสฺส อวณฺณํ ภาสติ, เอโก สํฆสฺส อวณฺณํ ภาสติ.\csromanspacer{} อิเมสํ โข ภิกฺขเว ติณฺณํ ภิกฺขูนํ อากงฺขมาโน สํโฆ อาปตฺติยา อทสฺสเน อุกฺเขปนียกมฺมํ กเรยฺย.}

\csromanheader{2}{อาปตฺติยา อทสฺสเน อุกฺเขปนียกมฺเม}
\nitthitam{อากงฺขมานฉกฺกํ นิฏฺฐิตํ.}
\csromanmatikaanchor{mtk.34}
\csromantocmarkref{subsection}{เตจตฺตาลีสวตฺต}{mtk.34}
\bookmark[dest={mtk.34},level=2]{เตจตฺตาลีสวตฺต}
\csromanheader{2}{\textbf{เตจตฺตาลีสวตฺต}}
\vspace{\tipitakaparskip}
\csromancenter{อาปตฺติยา อทสฺสเน อุกฺเขปนียกมฺมกเตน ภิกฺขเว ภิกฺขุนา สมฺมา วตฺติตพฺพํ.\csromanspacer{} ตตฺรายํ สมฺมาวตฺตนา\csromandash{}น อุปสมฺปาเทตพฺพํ, น นิสฺสโย ทาตพฺโพ, น สามเณโร อุปฏฺฐาเปตพฺโพ, น ภิกฺขุโนวาทกสมฺมุติ สาทิตพฺพา, สมฺมเตนปิ ภิกฺขุนิโย น โอวทิตพฺพา, ยาย อาปตฺติยา สํเฆน อาปตฺติยา อทสฺสเน อุกฺเขปนียกมฺมํ กตํ โหติ, สา อาปตฺติ น อาปชฺชิตพฺพา, อญฺญา วา ตาทิสิกา, ตโต วา ปาปิฏฺฐตรา, กมฺมํ น ครหิตพฺพํ, กมฺมิกา น ครหิตพฺพา, น ปกตตฺตสฺส ภิกฺขุโน อภิวาทนํ ปจฺจุฏฺฐานํ อญฺชลิกมฺมํ สามีจิกมฺมํ อาสนาภิหาโร เสยฺยาภิหาโร ปาโททกํ ปาทปีฐํ ปาทกถลิกํ ปตฺตจีวรปฺปฏิคฺคหณํ นหาเน ปิฏฺฐิปริกมฺมํ สาทิตพฺพํ, น ปกตตฺโต ภิกฺขุ สีลวิปตฺติยา อนุทฺธํเสตพฺโพ, น ปตฺติยา อนุทฺธํเสตพฺโพ, น ทิฏฺฐิวิปตฺติยา อนุทฺธํเสตพฺโพ, น อาชีววิปตฺติยา ทฺธํเสตพฺโพ, น ภิกฺขุ ภิกฺขูหิ เภเทตพฺโพ\footnote{น ภิกฺขู ภิกฺขูหิ เภเทตพฺพา (สฺยา)}, น คิหิทฺธโช ธาเรตพฺโพ, น ติตฺถิยทฺธโช ธาเรตพฺโพ, น ติตฺถิยา เสวิตพฺพา, ภิกฺขู เสวิตพฺพา,ภิกฺขุสิกฺขาย สิกฺขิตพฺพํ\footnote{ภิกฺขุสิกฺขา สิกฺขิตพฺพา (สฺยา)}, น ปกตตฺเตน ภิกฺขุนา สทฺธิํ เอกจฺฉนฺเน อาวาเส วตฺถพฺพํ, น เอกจฺฉนฺเน}
\prosecont{\csromanfolio{53}อนาวาเส วตฺถพฺพํ, น เอกจฺฉนฺเน อาวาเส วา อนาวาเส วา วตฺถพฺพํ, ปกตตฺตํ ภิกฺขุํ ทิสฺวา อาสนา วุฏฺฐาตพฺพํ, น ปกตตฺโต ภิกฺขุ อาสาเทตพฺโพ อนฺโต วา พหิ วา, น ปกตตฺตสฺส ภิกฺขุโน อุโปสโถ ฐเปตพฺโพ, น ปวารณา ฐเปตพฺพา, น สวจนียํ กาตพฺพํ, น อนุวาโท ปฏฺฐเปตพฺโพ, น โอกาโส กาเรตพฺโพ, น โจเทตพฺโพ, น สาเรตพฺโพ, น ภิกฺขูหิ สมฺปโยเชตพฺพนฺติ.}

\csromanheader{2}{อาปตฺติยา อทสฺสเน อุกฺเขปนียกมฺเม}
\nitthitam{เตจตฺตาลีสวตฺตํ นิฏฺฐิตํ.}
\proseitem{52}{อถ โข สํโฆ ฉนฺนสฺส ภิกฺขุโน อาปตฺติยา อทสฺสเน อุกฺเขปนียกมฺมํ อกาสิ อสมฺโภคํ สํเฆน.\csromanspacer{} โส สํเฆน อาปตฺติยา อทสฺสเน อุกฺเขปนียกมฺมกโต ตมฺหา อาวาสา อญฺญํ อาวาสํ อคมาสิ.\csromanspacer{} ตตฺถ ภิกฺขู เนว อภิวาเทสุํ, น ปจฺจุฏฺเฐสุํ, น อญฺชลิกมฺมํ น สามีจิกมฺมํ อกํสุ, น สกฺกริํสุ, น ครุํ กริํสุ\footnote{น ครุกริํสุ (ก)}, น มาเนสุํ, น ปูเชสุํ.\csromanspacer{} โส ภิกฺขูหิ อสกฺกริยมาโน อครุกริยมาโน อมานิยมาโน อปูชิยมาโน อสกฺการปกโต ตมฺหาปิ อาวาสา อญฺญํ อาวาสํ อคมาสิ.\csromanspacer{} ตตฺถปิ ภิกฺขู เนว อภิวาเทสุํ, น ปจฺจุฏฺเฐสุํ, น อญฺชลิกมฺมํ น สามีจิกมฺมํ อกํสุ, น สกฺกริํสุ, น ครุํ กริํสุ, น มาเนสุํ, น ปูเชสุํ.\csromanspacer{} โส ภิกฺขูหิ อสกฺกริยมาโน อครุกริยมาโน อมานิยมาโน อปูชิยมาโน อสกฺการปกโต ตมฺหาปิ อาวาสา อญฺญํ อาวาสํ อคมาสิ.\csromanspacer{} ตตฺถปิ ภิกฺขู เนว อภิวาเทสุํ, น ปจฺจุฏฺเฐสุํ, น อญฺชลิกมฺมํ น สามีจิกมฺมํ อกํสุ, น สกฺกริํสุ, น ครุํ กริํสุ, น มาเนสุํ, น ปูเชสุํ.\csromanspacer{} โส ภิกฺขูหิ อสกฺกริยมาโน อครุกริยมาโน อมานิยมาโน อปูชิยมาโน อสกฺการปกโต ปุนเทว โกสมฺพิํ ปจฺจาคญฺฉิ.\csromanspacer{} โส สมฺมา วตฺตติ, โลมํ ปาเตติ, เนตฺถารํ วตฺตติ, ภิกฺขู อุปสงฺกมิตฺวา เอวํ วเทติ “อหํ อาวุโส สํเฆน อาปตฺติยา อทสฺสเน อุกฺเขปนียกมฺมกโต สมฺมา วตฺตามิ, โลมํ ปาเตมิ, เนตฺถารํ วตฺตามิ, กถํ นุ โข มยา ปฏิปชฺชิตพฺพนฺ”ติ.\csromanspacer{} ภิกฺขู ภควโต เอตมตฺถํ อาโรเจสุํ ฯเปฯ.\csromanspacer{} เตน หิ ภิกฺขเว สํโฆ ฉนฺนสฺส ภิกฺขุโน อาปตฺติยา อทสฺสเน อุกฺเขปนียกมฺมํ ปฏิปฺปสฺสมฺเภตุ.}

\csromanmatikaanchor{mtk.35}
\csromantocmarkref{subsection}{นปฺปฏิปฺปสฺสมฺเภตพฺพเตจตฺตาลีสก}{mtk.35}
\bookmark[dest={mtk.35},level=2]{นปฺปฏิปฺปสฺสมฺเภตพฺพเตจตฺตาลีสก}
\csromanheader{2}{1. \textbf{นปฺปฏิปฺปสฺสมฺเภตพฺพเตจตฺตาลีสก}}
\proseitem{53}{\csromanfolio{54}ปญฺจหิ ภิกฺขเว องฺเคหิ สมนฺนาคตสฺส ภิกฺขุโน อาปตฺติยา อทสฺสเน อุกฺเขปนียกมฺมํ นปฺปฏิปฺปสฺสมฺเภตพฺพํ.\csromanspacer{} อุปสมฺปาเทติ, นิสฺสยํ เทติ, สามเณรํ อุปฏฺฐาเปติ, ภิกฺขุโนวาทกสมฺมุติํ สาทิยติ, สมฺมโตปิ ภิกฺขุนิโย โอวทติ.\csromanspacer{} อิเมหิ โข ภิกฺขเว ปญฺจหงฺเคหิ สมนฺนาคตสฺส ภิกฺขุโน อาปตฺติยา อทสฺสเน อุกฺเขปนียกมฺมํ นปฺปฏิปฺปสฺสมฺเภตพฺพํ. (5)}

\prose{\csromansymbolfootnote{*}{วิ 5. 317 ปิฏฺเฐปิ.}อปเรหิปิ ภิกฺขเว ปญฺจหงฺเคหิ สมนฺนาคตสฺส ภิกฺขุโน อาปตฺติยา อทสฺสเน อุกฺเขปนียกมฺมํ นปฺปฏิปฺปสฺสมฺเภตพฺพํ.\csromanspacer{} ยาย อาปตฺติยา สํเฆน อาปตฺติยา อทสฺสเน อุกฺเขปนียกมฺมํ กตํ โหติ, ตํ อาปตฺติํ อาปชฺชติ, อญฺญํ วา ตาทิสิกํ, ตโต วา ปาปิฏฺฐตรํ, กมฺมํ ครหติ, กมฺมิเก ครหติ.\csromanspacer{} อิเมหิ โข ภิกฺขเว ปญฺจหงฺเคหิ สมนฺนาคตสฺส ภิกฺขุโน อาปตฺติยา อทสฺสเน อุกฺเขปนียกมฺมํ นปฺปฏิปฺปสฺสมฺเภตพฺพํ. (\textbf{1}0)}

\prose{อปเรหิปิ ภิกฺขเว ปญฺจหงฺเคหิ สมนฺนาคตสฺส ภิกฺขุโน อาปตฺติยา อทสฺสเน อุกฺเขปนียกมฺมํ นปฺปฏิปฺปสฺสมฺเภตพฺพํ.\csromanspacer{} ปกตตฺตสฺส ภิกฺขุโน อภิวาทนํ ปจฺจุฏฺฐานํ อญฺชลิกมฺมํ สามีจิกมฺมํ อาสนาภิหารํ สาทิยติ.\csromanspacer{} อิเมหิ โข ภิกฺขเว ปญฺจหงฺเคหิ สมนฺนาคตสฺส ภิกฺขุโน อาปตฺติยา อทสฺสเน อุกฺเขปนียกมฺมํ นปฺปฏิปฺปสฺสมฺเภตพฺพํ. (\textbf{1}5)}

\prose{อปเรหิปิ ภิกฺขเว ปญฺจหงฺเคหิ สมนฺนาคตสฺส ภิกฺขุโน อาปตฺติยา อทสฺสเน อุกฺเขปนียกมฺมํ นปฺปฏิปฺปสฺสมฺเภตพฺพํ.\csromanspacer{} ปกตตฺตสฺส ภิกฺขุโน เสยฺยาภิหารํ ปาโททกํ ปาทปิฐํ ปาทกถลิกํ ปตฺตจีวรปฺปฏิคฺคหณํ นหาเน ปิฏฺฐิปริกมฺมํ สาทิยติ.\csromanspacer{} อิเมหิ โข ภิกฺขเว ปญฺจหงฺเคหิ สมนฺนาคตสฺส ภิกฺขุโน อาปตฺติยา อทสฺสเน อุกฺเขปนียกมฺมํ นปฺปฏิปฺปสฺสมฺเภตพฺพํ. (20)}

\prose{อปเรหิปิ ภิกฺขเว ปญฺจหงฺเคหิ สมนฺนาคตสฺส ภิกฺขุโน อาปตฺติยา อทสฺสเน อุกฺเขปนียกมฺมํ นปฺปฏิปฺปสฺสมฺเภตพฺพํ.\csromanspacer{} ปกตตฺตํ ภิกฺขุํ สีลวิปตฺติยา อนุทฺธํเสติ, อาจารวิปตฺติยา อนุทฺธํเสติ, ทิฏฺฐิวิปตฺติยา อนุทฺธํเสติ, อาชีววิปตฺติยา อนุทฺธํเสติ, ภิกฺขุํ\footnote{ภิกฺขู (สี, สฺยา)} ภิกฺขูหิ เภเทติ.}

\prose{\csromanfolio{55}อิเมหิ โข ภิกฺขเว ปญฺจหงฺเคหิ สมนฺนาคตสฺส ภิกฺขุโน อาปตฺติยา อทสฺสเน อุกฺเขปนียกมฺมํ นปฺปฏิปฺปสฺสมฺเภตพฺพํ. (25)}

\prose{อปเรหิปิ ภิกฺขเว ปญฺจหงฺเคหิ สมนฺนาคตสฺส ภิกฺขุโน อาปตฺติยา อทสฺสเน อุกฺเขปนียกมฺมํ นปฺปฏิปฺปสฺสมฺเภตพฺพํ.\csromanspacer{} คิหิทฺธชํ ธาเรติ, ติตฺถิยทฺธชํ ธาเรติ, ติตฺถิเย เสวติ, ภิกฺขู น เสวติ, ภิกฺขุสิกฺขาย น สิกฺขติ.\csromanspacer{} อิเมหิ โข ภิกฺขเว ปญฺจหงฺเคหิ สมนฺนาคตสฺส ภิกฺขุโน อาปตฺติยา อทสฺสเน อุกฺเขปนียกมฺมํ นปฺปฏิปฺปสฺสมฺเภตพฺพํ. (30)}

\prose{อปเรหิปิ ภิกฺขเว ปญฺจหงฺเคหิ สมนฺนาคตสฺส ภิกฺขุโน อาปตฺติยา อทสฺสเน อุกฺเขปนียกมฺมํ นปฺปฏิปฺปสฺสมฺเภตพฺพํ.\csromanspacer{} ปกตตฺเตน ภิกฺขุนา สทฺธิํ เอกจฺฉนฺเน อาวาเส วสติ, เอกจฺฉนฺเน อนาวาเส วสติ, เอกจฺฉนฺเน อาวาเส วา อนาวาเส วา วสติ, ปกตตฺตํ ภิกฺขุํ ทิสฺวา อาสนา น วุฏฺฐาติ, ปกตตฺตํ ภิกฺขุํ อาสาเทติ อนฺโต วา พหิ วา.\csromanspacer{} อิเมหิ โข ภิกฺขเว ปญฺจหงฺเคหิ สมนฺนาคตสฺส ภิกฺขุโน อาปตฺติยา อทสฺสเน อุกฺเขปนียกมฺมํ นปฺปฏิปฺปสฺสมฺเภตพฺพํ. (35)}

\prose{อฏฺฐหิ ภิกฺขเว องฺเคหิ สมนฺนาคตสฺส ภิกฺขุโน อาปตฺติยา อทสฺสเน อุกฺเขปนียกมฺมํ นปฺปฏิปฺปสฺสมฺเภตพฺพํ.\csromanspacer{} ปกตตฺตสฺส ภิกฺขุโน อุโปสถํ ฐเปติ, ปวารณํ ฐเปติ, สวจนียํ กโรติ, อนุวาทํ ปฏฺฐเปติ, โอกาสํ กาเรติ, โจเทติ, สาเรติ, ภิกฺขูหิ สมฺปโยเชติ.\csromanspacer{} อิเมหิ โข ภิกฺขเว อฏฺฐหงฺเคหิ สมนฺนาคตสฺส ภิกฺขุโน อาปตฺติยา อทสฺสเน อุกฺเขปนียกมฺมํ นปฺปฏิปฺปสฺสมฺเภตพฺพํ. (43)}

\csromanheader{2}{อาปตฺติยา อทสฺสเน อุกฺเขปนียกมฺเม}
\nitthitam{นปฺปฏิปฺปสฺสมฺเภตพฺพเตจตฺตาลีสกํ นิฏฺฐิตํ.}
\csromanmatikaanchor{mtk.36}
\csromantocmarkref{subsection}{ปฏิปฺปสฺสมฺเภตพฺพเตจตฺตาลีสก}{mtk.36}
\bookmark[dest={mtk.36},level=2]{ปฏิปฺปสฺสมฺเภตพฺพเตจตฺตาลีสก}
\csromanheader{2}{\textbf{ปฏิปฺปสฺสมฺเภตพฺพเตจตฺตาลีสก}}
\proseitem{54}{ปญฺจหิ ภิกฺขเว องฺเคหิ สมนฺนาคตสฺส ภิกฺขุโน อาปตฺติยา อทสฺสเน อุกฺเขปนียกมฺมํ ปฏิปฺปสฺสมฺเภตพฺพํ.\csromanspacer{} น อุปสมฺปาเทติ, น นิสฺสยํ เทติ, น สามเณรํ อุปฏฺฐาเปติ, น ภิกฺขุโนวาทกสมฺมุติํ สาทิยติ, สมฺมโตปิ ภิกฺขุนิโย น โอวทติ.\csromanspacer{} อิเมหิ โข ภิกฺขเว ปญฺจหงฺเคหิ \csromanfolio{56}สมนฺนาคตสฺส ภิกฺขุโน อาปตฺติยา อทสฺสเน อุกฺเขปนียกมฺมํ ปฏิปฺปสฺสมฺเภตพฺพํ. (5)}

\prose{อปเรหิปิ ภิกฺขเว ปญฺจหงฺเคหิ สมนฺนาคตสฺส ภิกฺขุโน อาปตฺติยา อทสฺสเน อุกฺเขปนียกมฺมํ ปฏิปฺปสฺสมฺเภตพฺพํ.\csromanspacer{} ยาย อาปตฺติยา สํเฆน อาปตฺติยา อทสฺสเน อุกฺเขปนียกมฺมํ กตํ โหติ, ตํ อาปตฺติํ น อาปชฺชติ, อญฺญํ วา ตาทิสิกํ, ตโต วา ปาปิฏฺฐตรํ, กมฺมํ น ครหติ, กมฺมิเก น ครหติ.\csromanspacer{} อิเมหิ โข ภิกฺขเว ปญฺจหงฺเคหิ สมนฺนาคตสฺส ภิกฺขุโน อาปตฺติยา อทสฺสเน อุกฺเขปนียกมฺมํ ปฏิปฺปสฺสมฺเภตพฺพํ. (10)}

\prose{อปเรหิปิ ภิกฺขเว ปญฺจหงฺเคหิ สมนฺนาคตสฺส ภิกฺขุโน อาปตฺติยา อทสฺสเน อุกฺเขปนียกมฺมํ ปฏิปฺปสฺสมฺเภตพฺพํ.\csromanspacer{} น ปกตตฺตสฺส ภิกฺขุโน อภิวาทนํ ปจฺจุฏฺฐานํ อญฺชลิกมฺมํ สามีจิกมฺมํ อาสนาภิหารํ สาทิยติ.\csromanspacer{} อิเมหิ โข ภิ-กฺขเว ปญฺจหงฺเคหิ สมนฺนาคตสฺส ภิกฺขุโน อาปตฺติยา อทสฺสเน อุกฺเขปนียกมฺมํ ปฏิปฺปสฺสมฺเภตพฺพํ. (15)}

\prose{อปเรหิปิ ภิกฺขเว ปญฺจหงฺเคหิ สมนฺนาคตสฺส ภิกฺขุโน อาปตฺติยา อทสฺสเน อุกฺเขปนียกมฺมํ ปฏิปฺปสฺสมฺเภตพฺพํ.\csromanspacer{} น ปกตตฺตสฺส ภิกฺขุโน เสยฺยาภิหารํ ปาโททกํ ปาทปีฐํ ปาทกถลิกํ ปตฺตจีวรปฺปฏิคฺคหณํ นหาเน ปิฏฺฐิปริกมฺมํ สาทิยติ.\csromanspacer{} อิเมหิ โข ภิกฺขเว ปญฺจหงฺเคหิ สมนฺนาคตสฺส ภิกฺขุโน อาปตฺติยา อทสฺสเน อุกฺเขปนียกมฺมํ ปฏิปฺปสฺสมฺเภตพฺพํ. (20)}

\prose{อปเรหิปิ ภิกฺขเว ปญฺจหงฺเคหิ สมนฺนาคตสฺส ภิกฺขุโน อาปตฺติยา อทสฺสเน อุกฺเขปนียกมฺมํ ปฏิปฺปสฺสมฺเภตพฺพํ.\csromanspacer{} น ปกตตฺตํ ภิกฺขุํ สีลวิปตฺติยา อนุทฺธํเสติ, น อาจารวิปตฺติยา อนุทฺธํเสติ, น ทิฏฺฐิวิปตฺติยา อนุทฺธํเสติ, น อาชีววิปตฺติยา อนุทฺธํเสติ, น ภิกฺขุํ ภิกฺขูหิ เภเทติ.\csromanspacer{} อิเมหิ โข ภิกฺขเว ปญฺจหงฺเคหิ สมนฺนาคตสฺส ภิกฺขุโน อาปตฺติยา อทสฺสเน อุกฺเขปนียกมฺมํ ปฏิปฺปสฺสมฺเภตพฺพํ. (25)}

\prose{อปเรหิปิ ภิกฺขเว ปญฺจหงฺเคหิ สมนฺนาคตสฺส ภิกฺขุโน อาปตฺติยา อทสฺสเน อุกฺเขปนียกมฺมํ ปฏิปฺปสฺสมฺเภตพฺพํ.\csromanspacer{} น คิหิทฺธชํ ธาเรติ, น ติตฺถิยทฺธชํ ธาเรติ, น ติตฺถิเย เสวติ, ภิกฺขู เสวติ, \csromanfolio{57}ภิกฺขุสิกฺขาย สิกฺขติ.\csromanspacer{} อิเมหิ โข ภิกฺขเว ปญฺจหงฺเคหิ สมนฺนาคตสฺส ภิกฺขุโน อาปตฺติยา อทสฺสเน อุกฺเขปนียกมฺมํ ปฏิปฺปสฺสมฺเภตพฺพํ. (30)}

\prose{อปเรหิปิ ภิกฺขเว ปญฺจหงฺเคหิ สมนฺนาคตสฺสภิกฺขุโน อาปตฺติยา อทสฺสเน อุกฺเขปนียกมฺมํ ปฏิปฺปสฺสมฺเภตพฺพํ.\csromanspacer{} น ปกตตฺเตน ภิกฺขุนา สทฺธิํ เอกจฺฉนฺเน อาวาเส วสติ, น เอกจฺฉนฺเน อนาวาเส วสติ, น เอกจฺฉนฺเน อาวาเส วา อนาวาเส วา วสติ, ปกตตฺตํ ภิกฺขุํ ทิสฺวา อาสนา วุฏฺฐาติ, น ปกตตฺตํ ภิกฺขุํ อาสาเทติ อนฺโต วา พหิ วา.\csromanspacer{} อิเมหิ โข ภิกฺขเว ปญฺจหงฺเคหิ สมนฺนาคตสฺส ภิกฺขุโน อาปตฺติยา อทสฺสเน อุกฺเขปนียกมฺมํ ปฏิปฺปสฺสมฺเภตพฺพํ. (35)}

\prose{อฏฺฐหิ ภิกฺขเว องฺเคหิ สมนฺนาคตสฺส ภิกฺขุโน อาปตฺติยา อทสฺสเน อุกฺเขปนียกมฺมํ ปฏิปฺปสฺสมฺเภตพฺพํ.\csromanspacer{} น ปกตตฺตสฺส ภิกฺขุโน อุโปสถํ ฐเปติ, น ปวารณํ ฐเปติ, น สวจนียํ กโรติ, น อนุวาทํ ปฏฺฐเปติ, น โอกาสํ กาเรติ, น โจเทติ, น สาเรติ, น ภิกฺขูหิ สมฺปโยเชติ.\csromanspacer{} อิเมหิ โข ภิกฺขเว อฏฺฐหงฺเคหิ สมนฺนาคตสฺส ภิกฺขุโน อาปตฺติยา อทสฺสเน อุกฺเขปนียกมฺมํ ปฏิปฺปสฺสมฺเภตพฺพํ. (43)}

\csromanheader{2}{อาปตฺติยา อทสฺสเน อุกฺเขปนียกมฺเม}
\nitthitam{ปฏิปฺปสฺสมฺเภตพฺพเตจตฺตาลีสกํ นิฏฺฐิตํ.}
\proseitem{55}{เอวญฺจ ปน ภิกฺขเว ปฏิปฺปสฺสมฺเภตพฺพํ, เตน ภิกฺขเว ฉนฺเนน ภิกฺขุนา สํฆํ อุปสงฺกมิตฺวา เอกํสํ อุตฺตราสงฺคํ กริตฺวา วุฑฺฒานํ ภิกฺขูนํ ปาเท วนฺทิตฺวา อุกฺกุฏิกํ นิสีทิตฺวา อญฺชลิํ ปคฺคเหตฺวา เอวมสฺส วจนีโย “อหํ ภนฺเต สํเฆน อาปตฺติยา อทสฺสเน อุกฺเขปนียกมฺมกโต สมฺมา วตฺตามิ, โลมํ ปาเตมิ, เนตฺถารํ วตฺตามิ, อาปตฺติยา อทสฺสเน อุกฺเขปนียสฺส กมฺมสฺส ปฏิปฺปสฺสทฺธิํ ยาจามี”ติ.\csromanspacer{} ทุติยมฺปิ ยาจิตพฺพา.\csromanspacer{} ตติยมฺปิ ยาจิตพฺพา.\csromanspacer{} พฺยตฺเตน ภิกฺขุนา ปฏิพเลน สํโฆ ญาเปตพฺโพ\csromandash{}}

\prose{“สุณาตุ เม ภนฺเต สํโฆ, อยํ ฉนฺโน ภิกฺขุ สํเฆน อาปตฺติยา อทสฺสเน อุกฺเขปนียกมฺมกโต สมฺมา วตฺตติ, โลมํ ปาเตติ, เนตฺถารํ วตฺตติ, อาปตฺติยา อทสฺสเน อุกฺเขปนียสฺส กมฺมสฺส ปฏิปฺปสฺสทฺธิํ \csromanfolio{58}ยาจติ, ยทิ สํฆสฺส ปตฺตกลฺลํ, สํโฆ ฉนฺนสฺส ภิกฺขุโน อาปตฺติยา อทสฺสเน อุกฺเขปนียกมฺมํ ปฏิปฺปสฺสมฺเภยฺย, เอสา ญตฺติ.}

\prose{สุณาตุ เม ภนฺเต สํโฆ, อยํ ฉนฺโน ภิกฺขุ สํเฆน อาปตฺติยา อทสฺสเน อุกฺเขปนียกมฺมกโต สมฺมา วตฺตติ, โลมํ ปาเตติ, เนตฺถารํ, วตฺตติ, อาปตฺติยา อทสฺสเน อุกฺเขปนียสฺส กมฺมสฺส ปฏิปฺปสฺสทฺธิํ ยาจติ, สํโฆ ฉนฺนสฺส ภิกฺขุโน อาปตฺติยา อทสฺสเน อุกฺเขปนียกมฺมํ ปฏิปฺปสฺสมฺเภติ, ยสฺสายสฺมโต ขมติ ฉนฺนสฺส ภิกฺขุโน อาปตฺติยา อทสฺสเน อุกฺเขปนียสฺส กมฺมสฺส ปฏิปฺปสฺสทฺธิ, โส ตุณฺหสฺส, ยสฺส นกฺขมติ, โส ภาเสยฺย.}

\prose{ทุติยมฺปิ เอตมตฺถํ วทามิ ฯเปฯ.\csromanspacer{} ตติยมฺปิ เอตมตฺถํ วทามิ.\csromanspacer{} สุณาตุ เม ภนฺเต สํโฆ, อยํ ฉนฺโน ภิกฺขุ สํเฆน อาปตฺติยา อทสฺสเน อุกฺเขปนียกมฺมกโต สมฺมา วตฺตติ, โลมํ ปาเตติ, เนตฺถารํ วตฺตติ, อาปตฺติยา อทสฺสเน อุกฺเขปนียสฺส กมฺมสฺส ปฏิปฺปสฺสทฺธิํ ยาจติ, สํโฆ ฉนฺนสฺส ภิกฺขุโน อาปตฺติยา อทสฺสเน อุกฺเขปนียกมฺมํ ปฏิปฺปสฺสมฺเภติ, ยสฺสายสฺมโต ขมติ ฉนฺนสฺส ภิกฺขุโน อาปตฺติยา อทสฺสเน อุกฺเขปนียสฺส กมฺมสฺส ปฏิปฺปสฺสทฺธิ, โส ตุณฺหสฺส, ยสฺส นกฺขมติ, โส ภาเสยฺย.}

\prose{ปฏิปฺปสฺสทฺธํ สํเฆน ฉนฺนสฺส ภิกฺขุโน อาปตฺติยา อทสฺสเน อุกฺเขปนียกมฺมํ, ขมติ สํฆสฺส, ตสฺมา ตุณฺหี, เอวเมตํ ธารยามี”ติ.}

\nitthitam{อาปตฺติยา อทสฺสเน อุกฺเขปนียกมฺมํ นิฏฺฐิตํ ปญฺจมํ.}
\csromanmatikaanchor{mtk.37}
\csromantocmarkref{section}{6. อาปตฺติยา อปฺปฏิกมฺเม อุกฺเขปนียกมฺม ...}{mtk.37}
\bookmark[dest={mtk.37},level=1]{6. อาปตฺติยา อปฺปฏิกมฺเม อุกฺเขปนียกมฺม ...}
\csromanheader{1}{6. อาปตฺติยา อปฺปฏิกมฺเม อุกฺเขปนียกมฺม}
\proseitem{56}{เตน สมเยน พุทฺโธ ภควา โกสมฺพิยํ วิหรติ โฆสิตาราเม.\csromanspacer{} เตน โข ปน สมเยน อายสฺมา ฉนฺโน อาปตฺติํ อาปชฺชิตฺวา น อิจฺฉติ อาปตฺติํ ปฏิกาตุํ.\csromanspacer{} เย เต ภิกฺขู อปฺปิจฺฉา ฯเปฯ เต อุชฺฌายนฺติ ขิยฺยนฺติ วิปาเจนฺติ “กถํ หิ นาม อายสฺมา ฉนฺโน อาปตฺติํ อาปชฺชิตฺวา น อิจฺฉิสฺสติ อาปตฺติํ ปฏิกาตุนฺ”ติ.\csromanspacer{} อถ โข เต ภิกฺขู ภควโต เอตมตฺถํ อาโรเจสุํ.}

\prose{\csromanfolio{59}อถ โข ภควา เอตสฺมิํ นิทาเน เอตสฺมิํ ปกรเณ ภิกฺขุสํฆํ สนฺนิปาตาเปตฺวา ภิกฺขู ปฏิปุจฺฉิ “สจฺจํ กิร ภิกฺขเว ฉนฺโน ภิกฺขุ อาปตฺติํ อาปชฺชิตฺวา น อิจฺฉติ อาปตฺติํ ปฏิกาตุนฺ”ติ.\csromanspacer{} สจฺจํ ภควาติ.\csromanspacer{} วิครหิ พุทฺโธ ภควา อนนุจฺฉวิกํ ฯเปฯ.\csromanspacer{} กถํ หิ นาม โส ภิกฺขเว โมฆปุริโส อาปตฺติํ อาปชฺชิตฺวา น อิจฺฉิสฺสติ อาปตฺติํ ปฏิกาตุํ, เนตํ ภิกฺขเว อปฺปสนฺนานํ วา ปสาทาย ฯเปฯ วิครหิตฺวา ฯเปฯ ธมฺมิํ กถํ กตฺวา ภิกฺขู อามนฺเตสิ\csromandash{}เตน หิ ภิกฺขเว สํโฆ ฉนฺนสฺส ภิกฺขุโน อาปตฺติยา อปฺปฏิกมฺเม อุกฺเขปนียกมฺมํ กโรตุ อสมฺโภคํ สํเฆน, เอวญฺจ ปน ภิกฺขเว กาตพฺพํ, ปฐมํ ฉนฺโน ภิกฺขุ โจเทตพฺโพ, โจเทตฺวา สาเรตพฺโพ, สาเรตฺวา อาปตฺติํ อาโรเปตพฺโพ, อาปตฺติํ อาโรเปตฺวา พฺยตฺเตน ภิกฺขุนา ปฏิพเลน สํโฆ ญาเปตพฺโพ\csromandash{}}

\prose{“สุณาตุ เม ภนฺเต สํโฆ, อยํ ฉนฺโน ภิกฺขุ อาปตฺติํ อาปชฺชิตฺวา น อิจฺฉติ อาปตฺติํ ปฏิกาตุํ, ยทิ สํฆสฺส ปตฺตกลฺลํ, สํโฆ ฉนฺนสฺส ภิกฺขุโน อาปตฺติยา อปฺปฏิกมฺเม อุกฺเขปนียกมฺมํ กเรยฺย อสมฺโภคํ สํเฆน, เอสา ญตฺติ.}

\prose{สุณาตุ เม ภนฺเต สํโฆ, อยํ ฉนฺโน ภิกฺขุ อาปตฺติํ อาปชฺชิตฺวา น อิจฺฉติ อาปตฺติํ ปฏิกาตุํ, สํโฆ ฉนฺนสฺส ภิกฺขุโน อาปตฺติยา อปฺปฏิกมฺเม อุกฺเขปนียกมฺมํ กโรติ อสมฺโภคํ สํเฆน, ยสฺสายสฺมโต ขมติ ฉนฺนสฺส ภิกฺขุโน อาปตฺติยา อปฺปฏิกมฺเม อุกฺเขปนียสฺส กมฺมสฺส กรณํ อสมฺโภคํ สํเฆน, โส ตุณฺหสฺส, ยสฺส นกฺขมติ, โส ภาเสยฺย.}

\prose{ทุติยมฺปิ เอตมตฺถํ วทามิ ฯเปฯ.\csromanspacer{} ตติยมฺปิ เอตมตฺถํ วทามิ.\csromanspacer{} สุณาตุ เม ภนฺเต สํโฆ, อยํ ฉนฺโน ภิกฺขุ อาปตฺติํ อาปชฺชิตฺวา น อิจฺฉติ อาปตฺติํ ปฏิกาตุํ, สํโฆ ฉนฺนสฺส ภิกฺขุโน อาปตฺติยา อปฺปฏิกมฺเม อุกฺเขปนียกมฺมํ กโรติ อสมฺโภคํ สํเฆน, ยสฺสายสฺมโต ขมติ ฉนฺนสฺส ภิกฺขุโน อาปตฺติยา อปฺปฏิกมฺเม อุกฺเขปนียสฺส กมฺมสฺส กรณํ อสมฺโภคํ สํเฆน, โส ตุณฺหสฺส, ยสฺส นกฺขมติ, โส ภาเสยฺย.}

\prose{กตํ สํเฆน ฉนฺนสฺส ภิกฺขุโน อาปตฺติยา อปฺปฏิกมฺเม อุกฺเขปนียกมฺมํ อสมฺโภคํ สํเฆน, ขมติ สํฆสฺส, ตสฺมา ตุณฺหี, เอวเมตํ ธารยามี”ติ.}

\prose{\csromanfolio{60}อาวาสปรมฺปรญฺจ ภิกฺขเว สํสถ “ฉนฺโน ภิกฺขุ สํเฆน อาปตฺติยา\footnote{ฉนฺโน ภิกฺขุ อาปตฺติยา (สี, ก)} อปฺปฏิกมฺเม อุกฺเขปนียกมฺมกโต อสมฺโภคํ สํเฆนา”ติ.}

\csromanmatikaanchor{mtk.38}
\csromantocmarkref{subsection}{อธมฺมกมฺมทฺวาทสก}{mtk.38}
\bookmark[dest={mtk.38},level=2]{อธมฺมกมฺมทฺวาทสก}
\csromanheader{2}{\textbf{อธมฺมกมฺมทฺวาทสก}}
\proseitem{57}{ตีหิ ภิกฺขเว องฺเคหิ สมนฺนาคตํ อาปตฺติยา อปฺปฏิกมฺเม อุกฺเขปนียกมฺมํ อธมฺมกมฺมญฺจ โหติ อวินยกมฺมญฺจ ทุวูปสนฺตญฺจ.\csromanspacer{} อสมฺมุขา กตํ โหติ, อปฺปฏิปุจฺฉา กตํ โหติ, อปฺปฏิญฺญาย กตํ โหติ.\csromanspacer{} อิเมหิ โข ภิกฺขเว ตีหงฺเคหิ สมนฺนาคตํ อาปตฺติยา อปฺปฏิกมฺเม อุกฺเขปนียกมฺมํ อธมฺมกมฺมญฺจ โหติ อวินยกมฺมญฺจ ทุวูปสนฺตญฺจ.}

\prose{อปเรหิปิ ภิกฺขเว ตีหงฺเคหิ สมนฺนาคตํ อาปตฺติยา อปฺปฏิกมฺเม อุกฺเขปนียกมฺมํ อธมฺมกมฺมญฺจ โหติ อวินยกมฺมญฺจ ทุวูปสนฺตญฺจ.\csromanspacer{} อนาปตฺติยา กตํ โหติ, อเทสนาคามินิยา อาปตฺติยา กตํ โหติ, เทสิตาย อาปตฺติยา กตํ โหติ ฯเปฯ.\csromanspacer{} อโจเทตฺวา กตํ โหติ, อสาเรตฺวา กตํ โหติ, อาปตฺติํ อนาโรเปตฺวา กตํ โหติ ฯเปฯ.\csromanspacer{} อสมฺมุขา กตํ โหติ, อธมฺเมน กตํ โหติ, วคฺเคน กตํ โหติ ฯเปฯ.\csromanspacer{} อปฺปฏิปุจฺฉา กตํ โหติ, อธมฺเมน กตํ โหติ, วคฺเคน กตํ โหติ ฯเปฯ.\csromanspacer{} อปฺปฏิญฺญาย กตํ โหติ, อธมฺเมน กตํ โหติ, วคฺเคน กตํ โหติ ฯเปฯ.\csromanspacer{} อนาปตฺติยา กตํ โหติ, อธมฺเมน กตํ โหติ, วคฺเคน กตํ โหติ ฯเปฯ.\csromanspacer{} อเทสนาคามินิยา อาปตฺติยา กตํ โหติ, อธมฺเมน กตํ โหติ, วคฺเคน กตํ โหติ ฯเปฯ.\csromanspacer{} เทสิตาย อาปตฺติยา กตํ โหติ, อธมฺเมน กตํ โหติ, วคฺเคน กตํ โหติ ฯเปฯ.\csromanspacer{} อโจเทตฺวา กตํ โหติ, อธมฺเมน กตํ โหติ, วคฺเคน กตํ โหติ ฯเปฯ.\csromanspacer{} อสาเรตฺวา กตํ โหติ, อธมฺเมน กตํ โหติ, วคฺเคน กตํ โหติ ฯเปฯ.\csromanspacer{} อาปตฺติํ อนาโรเปตฺวา กตํ โหติ, อธมฺเมน กตํ โหติ, วคฺเคน กตํ โหติ.\csromanspacer{} อิเมหิ โข ภิกฺขเว หีหงฺเคหิ สมนฺนาคตํ อาปตฺติยา อปฺปฏิกมฺเม อุกฺเขปนียกมฺมํ อธมฺมกมฺมญฺจ โหติ อวินยกมฺมญฺจ ทุวูปสนฺตญฺจ.}

\csromanheader{2}{อาปตฺติยา อปฺปฏิกมฺเม อุกฺเขปนียกมฺเม}
\nitthitam{\textbf{อธมฺมกมฺมทฺวาทสก}ํ นิฏฺฐิตํ.}
\csromanmatikaanchor{mtk.39}
\csromantocmarkref{subsection}{ธมฺมกมฺมทฺวาทสก}{mtk.39}
\bookmark[dest={mtk.39},level=2]{ธมฺมกมฺมทฺวาทสก}
\csromanheader{2}{\textbf{ธมฺมกมฺมทฺวาทสก}}
\proseitem{58}{\csromanfolio{61}ตีหิ ภิกฺขเว องฺเคหิ สมนฺนาคตํ อาปตฺติยา อปฺปฏิกมฺเม อุกฺเขปนียกมฺมํ ธมฺมกมฺมญฺจ โหติ วินยกมฺมญฺจ สุวูปสนฺตญฺจ.\csromanspacer{} สมฺมุขา กตํ โหติ, ปฏิปุจฺฉา กตํ โหติ, ปฏิญฺญาย กตํ โหติ.\csromanspacer{} อิเมหิ โข ภิกฺขเว ตีหงฺเคหิ สมนฺนาคตํ อาปตฺติยา อปฺปฏิกมฺเม อุกฺเขปนียกมฺมํ ธมฺมกมฺมญฺจ โหติ วินยกมฺมญฺจ สุวูปสนฺตญฺจ.}

\prose{อปเรหิปิ ภิกฺขเว ตีหงฺเคหิ สมนฺนาคตํ อาปตฺติยา อปฺปฏิกมฺเม อุกฺเขปนียกมฺมํ ธมฺมกมฺมญฺจ โหติ วินยกมฺมญฺจ สุวูปสนฺตญฺจ.\csromanspacer{} อาปตฺติยา กตํ โหติ, เทสนาคามินิยา อาปตฺติยา กตํ โหติ, อเทสิตาย อาปตฺติยา กตํ โหติ ฯเปฯ.\csromanspacer{} โจเทตฺวา กตํ โหติ, สาเรตฺวา กตํ โหติ, อาปตฺติํ อาโรเปตฺวา กตํ โหติ ฯเปฯ.\csromanspacer{} สมฺมุขา กตํ โหติ, ธมฺเมน กตํ โหติ.\csromanspacer{} สมคฺเคน กตํ โหติ ฯเปฯ.\csromanspacer{} ปฏิปุจฺฉา กตํ โหติ, ธมฺเมน กตํ โหติ, สมคฺเคน กตํ โหติ ฯเปฯ.\csromanspacer{} ปฏิญฺญาย กตํ โหติ, ธมฺเมน กตํ โหติ, สมคฺเคน กตํ โหติ ฯเปฯ.\csromanspacer{} อาปตฺติยา กตํ โหติ, ธมฺเมน กตํ โหติ, สมคฺเคน กตํ โหติ ฯเปฯ.\csromanspacer{} เทสนาคามินิยา อาปตฺติยา กตํ โหติ, ธมฺเมน กตํ โหติ, สมคฺเคน กตํ โหติ ฯเปฯ.\csromanspacer{} อเทสิตาย อาปตฺติยา กตํ โหติ, ธมฺเมน กตํ โหติ, สมคฺเคน กตํ โหติ ฯเปฯ.\csromanspacer{} โจเทตฺวา กตํ โหติ, ธมฺเมน กตํ โหติ, สมคฺเคน กตํ โหติ ฯเปฯ.\csromanspacer{} สาเรตฺวา กตํ โหติ, ธมฺเมน กตํ โหติ, สมคฺเคน กตํ โหติ ฯเปฯ.\csromanspacer{} อาปตฺติํ อาโรเปตฺวา กตํ โหติ, ธมฺเมน กตํ โหติ, สมคฺเคน กตํ โหติ.\csromanspacer{} อิเมหิ โข ภิกฺขเว ตีหงฺเคหิ สมนฺนาคตํ อาปตฺติยา อปฺปฏิกมฺเม อุกฺเขปนียกมฺมํ ธมฺมกมฺมญฺจ โหติ วินยกมฺมญฺจ สุวูปสนฺตญฺจ.}

\csromanheader{2}{อาปตฺติยา อปฺปฏิกมฺเม อุกฺเขปนียกมฺเม}
\nitthitam{\textbf{ธมฺมกมฺมทฺวาทสก}ํ นิฏฺฐิตํ.}
\csromanmatikaanchor{mtk.40}
\csromantocmarkref{subsection}{อากงฺขมานฉกฺก}{mtk.40}
\bookmark[dest={mtk.40},level=2]{อากงฺขมานฉกฺก}
\csromanheader{2}{\textbf{อากงฺขมานฉกฺก}}
\proseitem{59}{\csromanfolio{62}\csromansymbolfootnote{*}{วิ 5. 220 ปิฏฺเฐปิ.}ตีหิ ภิกฺขเว องฺเคหิ สมนฺนาคตสฺส ภิกฺขุโน อากงฺขมาโน สํโฆ อาปตฺติยา อปฺปฏิกมฺเม อุกฺเขปนียกมฺมํ กเรยฺย.\csromanspacer{} ภณฺฑนการโก โหติ กลหการโก วิวาทการโก ภสฺสการโก สํเฆ อธิกรณการโก, พาโล โหติ อพฺยตฺโต อาปตฺติพหุโล อนปทาโน, คิหิสํสฏฺโฐ วิหรติ อนนุโลมิเกหิ คิหิสํสคฺเคหิ.\csromanspacer{} อิเมหิ โข ภิกฺขเว ตีหงฺเคหิ สมนฺนาคตสฺส ภิกฺขุโน อากงฺขมาโน สํโฆ อาปตฺติยา อปฺปฏิกมฺเม อุกฺเขปนียกมฺมํ กเรยฺย.}

\prose{อปเรหิปิ ภิกฺขเว ตีหงฺเคหิ สมนฺนาคตสฺส ภิกฺขุโน อากงฺขมาโน สํโฆ อาปตฺติยา อปฺปฏิกมฺเม อุกฺเขปนียกมฺมํ กเรยฺย.\csromanspacer{} อธิสีเล สีลวิปนฺโน โหติ, อชฺฌาจาเร อาจารวิปนฺโน โหติ, อติทิฏฺฐิยา ทิฏฺฐิวิปนฺโน โหติ.\csromanspacer{} อิเมหิ โข ภิกฺขเว ตีหงฺเคหิ สมนฺนาคตสฺส ภิกฺขุโน อากงฺขมาโน สํโฆ อาปตฺติยา อปฺปฏิกมฺเม อุกฺเขปนียกมฺมํ กเรยฺย.}

\prose{อปเรหิปิ ภิกฺขเว ตีหงฺเคหิ สมนฺนาคตสฺส ภิกฺขุโน อากงฺขมาโน สํโฆ อาปตฺติยา อปฺปฏิกมฺเม อุกฺเขปนียกมฺมํ กเรยฺย.\csromanspacer{} พุทฺธสฺส อวณฺณํ ภาสติ, ธมฺมสฺส อวณฺณํ ภาสติ, สํฆสฺส อวณฺณํ ภาสติ.\csromanspacer{} อิเมหิ โข ภิกฺขเว ตีหงฺเคหิ สมนฺนาคตสฺส ภิกฺขุโน อากงฺขมาโน สํโฆ อาปตฺติยา อปฺปฏิกมฺเม อุกฺเขปนียกมฺมํ กเรยฺย.}

\prose{ติณฺณํ ภิกฺขเว ภิกฺขูนํ อากงฺขมาโน สํโฆ อาปตฺติยา อปฺปฏิกมฺเม อุกฺเขปนียกมฺมํ กเรยฺย.\csromanspacer{} เอโก ภณฺฑนการโก โหติ กลหการโก วิวาทการโก ภสฺสการโก สํเฆ อธิกรณการโก, เอโก พาโล โหติ อพฺยตฺโต อาปตฺติพหุโล อนปทาโน, เอโก คิหิสํสฏฺโฐ วิหรติ อนนุโลมิเกหิ คิหิ สํสคฺเคหิ.\csromanspacer{} อิเมสํ โข ภิกฺขเว ติณฺณํ ภิกฺขูนํ อากงฺขมาโน สํโฆ อาปตฺติยา อปฺปฏิกมฺเม อุกฺเขปนียกมฺมํ กเรยฺย.}

\prose{อปเรสมฺปิ ภิกฺขเว ติณฺณํ ภิกฺขูนํ อากงฺขมาโน สํโฆ อาปตฺติยา อปฺปฏิกมฺเม อุกฺเขปนียกมฺมํ กเรยฺย.\csromanspacer{} เอโก อธิสีเล \csromanfolio{63}สีลวิปนฺโน โหติ.\csromanspacer{} เอโก อชฺฌาจาเร อาจารวิปนฺโน โหติ, เอโก อติทิฏฺฐิยา ทิฏฺฐิวิปนฺโน โหติ.\csromanspacer{} อิเมสํ โข ภิกฺขเว ติณฺณํ ภิกฺขูนํ อากงฺขมาโน สํโฆ อาปตฺติยา อปฺปฏิกมฺเม อุกฺเขปนียกมฺมํ กเรยฺย.}

\prose{อปเรสมฺปิ ภิกฺขเว ติณฺณํ ภิกฺขูนํ อากงฺขมาโน สํโฆ อาปตฺติยา อปฺปฏิกมฺเม อุกฺเขปนียกมฺมํ กเรยฺย.\csromanspacer{} เอโก พุทฺธสฺส อวณฺณํ ภาสติ, เอโก ธมฺมสฺส อวณฺณํ ภาสติ, เอโก สํฆสฺส อวณฺณํ ภาสติ.\csromanspacer{} อิเมสํ โข ภิกฺขเว ติณฺณํ ภิกฺขูนํ อากงฺขมาโน สํโฆ อาปตฺติยา อปฺปฏิกมฺเม อุกฺเขปนียกมฺมํ กเรยฺย.}

\csromanheader{2}{อาปตฺติยา อปฺปฏิกมฺเม อุกฺเขปนียกมฺเม}
\nitthitam{อากงฺขมานฉกฺกํ นิฏฺฐิตํ.}
\csromanmatikaanchor{mtk.41}
\csromantocmarkref{subsection}{เตจตฺตาลีสวตฺต}{mtk.41}
\bookmark[dest={mtk.41},level=2]{เตจตฺตาลีสวตฺต}
\csromanheader{2}{\textbf{เตจตฺตาลีสวตฺต}}
\vspace{\tipitakaparskip}
\csromancenter{อาปตฺติยา อปฺปฏิกมฺเม อุกฺเขปนียกมฺมกเตน ภิกฺขเว ภิกฺขุนา สมฺมา วตฺติตพฺพํ.\csromanspacer{} ตตฺรยํ สมฺมา วตฺตนา\csromandash{}น อุปสมฺปาเทตพฺพํ, น นิสฺสโย ทาตพฺโพ, น สามเณโร อุปฏฺฐาเปตพฺโพ, น ภิกฺขุโนวาทกสมฺมุติ สาทิตพฺพา, สมฺมเตนปิ ภิกฺขุนิโย น โอวทิตพฺพา, ยาย อาปตฺติยา สํเฆน อาปตฺติยา อปฺปฏิกมฺเม อุกฺเขปนียกมฺมํ กตํ โหติ, สา อาปตฺติ น อาปชฺชิตพฺพา, อญฺญา วา ตาทิสิกา, ตโต วา ปาปิฏฺฐตรา, กมฺมํ น ครหิตพฺพํ, กมฺมิกา น ครหิตพฺพา, น ปกตตฺตสฺส ภิกฺขุโน อภิวาทนํ ปจฺจุฏฺฐานํ อญฺชลิกมฺมํ สามิจิกมฺมํ อาสนาภิหาโร เสยฺยาภิหาโร ปาโททกํ ปาทปีฐํ ปาทกถลิกํ ปตฺตจีวรปฺปฏิคฺคหณํ นหาเน ปิฏฺฐิปริกมฺมํ สาทิตพฺพํ, น ปกตตฺโต ภิกฺขุ สีลวิปตฺติยา อนุทฺธํเสตพฺโพ, น อาจารวิปตฺติยา อนุทฺธํเสตพฺโพ, น ทิฏฺฐิวิปตฺติยา อนุทฺธํเสตพฺโพ, น อาชีววิปตฺติยา อนุทฺธํเสตพฺโพ, น ภิกฺขุ ภิกฺขูหิ เภเทตพฺโพ, น คิหิทฺธโช ธาเรตพฺโพ, น ติตฺถิยทฺธโช ธาเรตพฺโพ, น ติตฺถิยา เสวิตพฺพา, ภิกฺขู เสวิตพฺพา, ภิกฺขุสิกฺขาย สิกฺขิตพฺพํ, น ปกตตฺเตน ภิกฺขุนา สทฺธิํ เอกจฺฉนฺเน อาวาเส วตฺถพฺพํ, น เอกจฺฉนฺเน อนาวาเส วตฺถพฺพํ, น เอกจฺฉนฺเน อาวาเส}
\prosecont{\csromanfolio{64}วา อนาวาเส วา วตฺถพฺพํ, ปกตตฺตํ ภิกฺขุํ ทิสฺวา อาสนา วุฏฺฐาตพฺพํ, น ปกตตฺโต ภิกฺขุ อาสาเทตพฺโพ อนฺโต วา พหิ วา, น ปกตตฺตสฺส ภิกฺขุโน อุโปสโถ ฐเปตพฺโพ, น ปวารณา ฐเปตพฺพา, น สวจนียํ กาตพฺพํ, น อนุวาโท ปฏฺฐเปตพฺโพ, น โอกาโส กาเรตพฺโพ, น โจเทตพฺโพ, น สาเรตพฺโพ, น ภิกฺขูหิ สมฺปโยเชตพฺพนฺติ.}

\csromanheader{2}{อาปตฺติยา อปฺปฏิกมฺเม อุกฺเขปนียกมฺเม}
\nitthitam{เตจตฺตาลีสวตฺตํ นิฏฺฐิตํ.}
\proseitem{61}{อถ โข สํโฆ ฉนฺนสฺส ภิกฺขุโน อาปตฺติยา อปฺปฏิกมฺเม อุกฺเขปนียกมฺมํ อกาสิ อสมฺโภคํ สํเฆน.\csromanspacer{} โส สํเฆน อาปตฺติยา อปฺปฏิกมฺเม อุกฺเขปนียกมฺมกโต ตมฺหา อาวาสา อญฺญํ อาวาสํ อคมาสิ.\csromanspacer{} ตตฺถ ภิกฺขู เนว อภิวาเทสุํ, น ปจฺจุฏฺเฐสุํ, น อญฺชลิกมฺมํ, น สามีจิกมฺมํ อกํสุ, น สกฺกริสุ, น ครุํ กริํสุ, น มาเนสุํ, น ปูเชสุํ.\csromanspacer{} โส ภิกฺขูหิ อสกฺกริยมาโน อครุกริยมาโน อมานิยมาโน อปูชิยมาโน อสกฺการปกโต ตมฺหาปิ อาวาสา อญฺญํ อาวาสํ อคมาสิ.\csromanspacer{} ตตฺถปิ ภิกฺขู เนว อภิวาเทสุํ, น ปจฺจุฏฺเฐสุํ, น อญฺชลิกมฺมํ, น สามีจิกมฺมํ อกํสุ, น สกฺกริํสุ, น ครุํ กริํสุ, น มาเนสุํ, น ปูเชสุํ.\csromanspacer{} โส ภิกฺขูหิ อสกฺกริยมาโน อครุกริยมาโน อมานิยมาโน อปูชิยมาโน อสกฺการปกโต ตมฺหาปิ อาวาสา อญฺญํ อาวาสํ อคมาสิ.\csromanspacer{} ตตฺถปิ ภิกฺขู เนว อภิวาเทสุํ, น ปจฺจุฏฺเฐสุํ, น อญฺชลิกมฺมํ, น สามีจิกมฺมํ อกํสุ, น สกฺกริํสุ, น ครุํ กริํสุ, น มาเนสุํ, น ปูเชสุํ.\csromanspacer{} โส ภิกฺขูหิ อสกฺกริยมาโน อครุกริยมาโน อมานิยมาโน อปูชิยมาโน อสกฺการปกโต ปุนเทว โกสมฺพิํ ปจฺจาคญฺฉิ.\csromanspacer{} โส สมฺมา วตฺตติ, โลมํ ปาเตติ, เนตฺถารํ วตฺตติ, ภิกฺขู อุปสงฺกมิตฺวา เอวํ วเทติ “อหํ อาวุโส สํเฆน อาปตฺติยา อปฺปฏิกมฺเม อุกฺเขปนียกมฺมกโต สมฺมา วตฺตามิ, โลมํ ปาเตมิ, เนตฺถารํ วตฺตามิ, กถํ นุ ขา มยา ปฏิปชฺชิตพฺพนฺ”ติ.\csromanspacer{} ภิกฺขู ภควโต เอตมตฺถํ อาโรเจสุํ ฯเปฯ.\csromanspacer{} เตน หิ ภิกฺขเว สํโฆ ฉนฺนสฺส ภิกฺขุโน อาปตฺติยา อปฺปฏิกมฺเม อุกฺเขปนียกมฺมํ ปฏิปฺปสฺสมฺเภตุ.}

\csromanmatikaanchor{mtk.42}
\csromantocmarkref{subsection}{นปฺปฏิปฺปสฺสมฺเภตพฺพเตจตฺตาลีสก}{mtk.42}
\bookmark[dest={mtk.42},level=2]{นปฺปฏิปฺปสฺสมฺเภตพฺพเตจตฺตาลีสก}
\csromanheader{2}{\textbf{นปฺปฏิปฺปสฺสมฺเภตพฺพเตจตฺตาลีสก}}
\proseitem{62}{\csromanfolio{65}ปญฺจหิ ภิกฺขเว องฺเคหิ สมนฺนาคตสฺส ภิกฺขุโน อาปตฺติยา อปฺปฏิกมฺเม อุกฺเขปนียกมฺมํ นปฺปฏิปฺปสฺสมฺเภตพฺพํ.\csromanspacer{} อุปสมฺปาเทติ, นิสฺสยํ เทติ, สามเณรํ อุปฏฺฐาเปติ, ภิกฺขุโนวาทกสมฺมุติํ สาทิยติ, สมฺมโตปิ ภิกฺขุนิโย โอวทติ.\csromanspacer{} อิเมหิ โข ภิกฺขเว ปญฺจหงฺเคหิ สมนฺนาคตสฺส ภิกฺขุโน อาปตฺติยา อปฺปฏิกมฺเม อุกฺเขปนียกมฺมํ นปฺปฏิปฺปสฺสมฺเภตพฺพํ.}

\prose{\csromansymbolfootnote{*}{วิ 5. 317 ปิฏฺเฐปิ.}อปเรหิปิ ภิกฺขเว ปญฺจหงฺเคหิ สมนฺนาคตสฺส ภิกฺขุโน อาปตฺติยา อปฺปฏิกมฺเม อุกฺเขปนียกมฺมํ นปฺปฏิปฺปสฺสมฺเภตพฺพํ.\csromanspacer{} ยาย อาปตฺติยา สํเฆน อาปตฺติยา อปฺปฏิกมฺเม อุกฺเขปนียกมฺมํ กตํ โหติ, ตํ อาปตฺติํ อาปชฺชติ, อญฺญํ วา ตาทิสิกํ, ตโต วา ปาปิฏฺฐตรํ, กมฺมํ ครหติ, กมฺมิเก ครหติ ฯเปฯ.\csromanspacer{} ปกตตฺตสฺส ภิกฺขุโน อภิวาทนํ ปจฺจุฏฺฐานํ อญฺชลิกมฺมํ สามีจิกมฺมํ อาสนาภิหารํ สาทิยติ ฯเปฯ.\csromanspacer{} ปกตตฺตสฺส ภิกฺขุโน เสยฺยาภิหารํ ปาโททกํ ปาทปีฐํ ปาทกถลิกํ ปตฺตจีวรปฺปฏิคฺคหณํ นหาเน ปิฏฺฐิปริกมฺมํ สาทิยติ ฯเปฯ.\csromanspacer{} ปกตตฺตํ ภิกฺขุํ สีลวิปตฺติยา อนุทฺธํเสติ, อาจารวิปตฺติยา อนุทฺธํเสติ, ทิฏฺฐิวิปตฺติยา อนุทฺธํเสติ, อาชีววิปตฺติยา อนุทฺธํเสติ, ภิกฺขุํ ภิกฺขูหิ เภเทติ ฯเปฯ.\csromanspacer{} คิหิทฺธชํ ธาเรติ, ติตฺถิยทฺธชํ ธาเรติ, ติตฺถิเย เสวติ, ภิกฺขู น เสวติ, ภิกฺขุสิกฺขาย น สิกฺขติ ฯเปฯ.\csromanspacer{} ปกตตฺเตน ภิกฺขุนา สทฺธิํ เอกจฺฉนฺเน อาวาเส วสติ, เอกจฺฉนฺเน อนาวาเส วสติ, เอกจฺฉนฺเน อาวาเส วา อนาวาเส วา วสติ, ปกตตฺตํ ภิกฺขุํ ทิสฺวา อาสนา น วุฏฺฐาติ, ปกตตฺตํ ภิกฺขุํ อาสาเทติ อนฺโต วา พหิ วา.\csromanspacer{} อิเมหิ โข ภิกฺขเว ปญฺจหงฺเคหิ สมนฺนาคตสฺส ภิกฺขุโน อาปตฺติยา อปฺปฏิกมฺเม อุกฺเขปนียกมฺมํ นปฺปฏิปฺปสฺสมฺเภตพฺพํ.}

\prose{อฏฺฐหิ ภิกฺขเว องฺเคหิ สมนฺนาคตสฺส ภิกฺขุโน อาปตฺติยา อปฺปฏิกมฺเม อุกฺเขปนียกมฺมํ นปฺปฏิปฺปสฺสมฺเภตพฺพํ.\csromanspacer{} ปกตตฺตสฺส ภิกฺขุโน อุโปสถํ ฐเปติ, ปวารณํ ฐเปติ, สวจนียํ กโรติ, อนุวาทํ ปฏฺฐเปติ, โอกาสํ กาเรติ, โจเทติ, สาเรติ, ภิกฺขูหิ \csromanfolio{66}สมฺปโยเชติ.\csromanspacer{} อิเมหิ โข ภิกฺขเว อฏฺฐหงฺเคหิ สมนฺนาคตสฺส ภิกฺขุโน อาปตฺติยา อปฺปฏิกมฺเม อุกฺเขปนียกมฺมํ นปฺปฏิปฺปสฺสมฺเภตพฺพํ.}

\csromanheader{2}{อาปตฺติยา อปฺปฏิกมฺเม อุกฺเขปนียกมฺเม}
\nitthitam{นปฺปฏิปฺปสฺสมฺเภตพฺพเตจตฺตาลีสกํ นิฏฺฐิตํ.}
\csromanmatikaanchor{mtk.43}
\csromantocmarkref{subsection}{ปฏิปฺปสฺสมฺเภตพฺพเตจตฺตาลีสก}{mtk.43}
\bookmark[dest={mtk.43},level=2]{ปฏิปฺปสฺสมฺเภตพฺพเตจตฺตาลีสก}
\csromanheader{2}{\textbf{ปฏิปฺปสฺสมฺเภตพฺพเตจตฺตาลีสก}}
\proseitem{63}{ปญฺจหิ ภิกฺขเว องฺเคหิ สมนฺนาคตสฺส ภิกฺขุโน อาปตฺติยา อปฺปฏิกมฺเม อุกฺเขปนียกมฺมํ ปฏิปฺปสฺสมฺเภตพฺพํ.\csromanspacer{} น อุปสมฺปาเทติ, น นิสฺสยํ เทติ, น สามเณรํ อุปฏฺฐาเปติ, น ภิกฺขุโนวาทกสมฺมุติํ สาทิยติ, สมฺมโตปิ ภิกฺขุนิโย น โอวทติ.\csromanspacer{} อิเมหิ โข ภิกฺขเว ปญฺจหงฺเคหิ สมนฺนาคตสฺส ภิกฺขุโน อาปตฺติยา อปฺปฏิกมฺเม อุกฺเขปนียกมฺมํ ปฏิปฺปสฺสมฺเภตพฺพํ.}

\prose{อปเรหิปิ ภิกฺขเว ปญฺจหงฺเคหิ สมนฺนาคตสฺส ภิกฺขุโน อาปตฺติยา อปฺปฏิกมฺเม อุกฺเขปนียกมฺมํ ปฏิปฺปสฺสมฺเภตพฺพํ.\csromanspacer{} ยาย อาปตฺติยา สํเฆน อาปตฺติยา อปฺปฏิกมฺเม อุกฺเขปนียกมฺมํ กตํ โหติ, ตํ อาปตฺติํ น อาปชฺชติ, อญฺญํ วา ตาทิสิกํ, ตโต วา ปาปิฏฺฐตรํ, กมฺมํ น ครหติ, กมฺมิเก น ครหติ ฯเปฯ.\csromanspacer{} น ปกตตฺตสฺส ภิกฺขุโน อภิวาทนํ ปจฺจุฏฺฐานํ อญฺชลิกมฺมํ สามีจิกมฺมํ อาสนาภิหารํ สาทิยติ ฯเปฯ.\csromanspacer{} น ปกตตฺตสฺส ภิกฺขุโน เสยฺยาภิหารํ ปาโททกํ ปาทปีฐํ ปาทกถลิกํ ปตฺตจีวรปฺปฏิคฺคหณํ นหาเน ปิฏฺฐิปริกมฺมํ สาทิยติ ฯเปฯ.\csromanspacer{} น ปกตตฺตํ ภิกฺขุํ สีลวิปตฺติยา อนุทฺธํเสติ, น อาจารวิปตฺติยา อนุทฺธํเสติ, น ทิฏฺฐิวิปตฺติยา อนุทฺธํเสติ, น อาชีววิปตฺติยา อนุทฺธํเสติ, น ภิกฺขุํ ภิกฺขูหิ เภเทติ ฯเปฯ.\csromanspacer{} น คิหิทฺธชํ ธาเรติ, น ติตฺถิยทฺธชํ ธาเรติ, น ติตฺถิเย เสวติ, ภิกฺขู เสวติ, ภิกฺขุสิกฺขาย สิกฺขติ ฯเปฯ.\csromanspacer{} น ปกตตฺเตน ภิกฺขุนา สทฺธิํ เอกจฺฉนฺเน อาวาเส วสติ, น เอกจฺฉนฺเน อนาวาเส วสติ, น เอกจฺฉนฺเน อาวาเส วา อนาวาเส วา วสติ, ปกตตฺตํ ภิกฺขุํ ทิสฺวา อาสนา วุฏฺฐาติ, น ปกตตฺตํ ภิกฺขุํ อาสาเทติ อนฺโต วา พหิ วา.\csromanspacer{} อิเมหิ โข ภิกฺขเว ปญฺจหงฺเคหิ สมนฺนาคตสฺส ภิกฺขุโน อาปตฺติยา อปฺปฏิกมฺเม อุกฺเขปนียกมฺมํ ปฏิปฺปสฺสมฺเภตพฺพํ.}

\prose{\csromanfolio{67}อฏฺฐหิ ภิกฺขเว องฺเคหิ สมนฺนาคตสฺส ภิกฺขุโน อาปตฺติยา อปฺปฏิกมฺเม อุกฺเขปนียกมฺมํ ปฏิปฺปสฺสมฺเภตพฺพํ.\csromanspacer{} น ปกตตฺตสฺส ภิกฺขุโน อุโปสถํ ฐเปติ, น ปวารณํ ฐเปติ, น สวจนียํ กโรติ, น อนุวาทํ ปฏฺฐเปติ, น โอกาสํ กาเรติ, น โจเทติ, น สาเรติ, น ภิกฺขูหิ สมฺปโยเชติ.\csromanspacer{} อิเมหิ โข ภิกฺขเว อฏฺฐหงฺเคหิ สมนฺนาคตสฺส ภิกฺขุโน อาปตฺติยา อปฺปฏิกมฺเม อุกฺเขปนียกมฺมํ ปฏิปฺปสฺสมฺเภตพฺพํ.}

\csromanheader{2}{อาปตฺติยา อปฺปฏิกมฺเม อุกฺเขปนียกมฺเม}
\nitthitam{ปฏิปฺปสฺสมฺเภตพฺพเตจตฺตาลีสกํ นิฏฺฐิตํ.}
\proseitem{64}{เอวญฺจ ปน ภิกฺขเว ปฏิปฺปสฺสมฺเภตพฺพํ, เตน ภิกฺขเว ฉนฺเนน ภิกฺขุนา สํฆํ อุปสงฺกมิตฺวา เอกํสํ อุตฺตราสงฺคํ กริตฺวา วุฑฺฒานํ ภิกฺขูนํ ปาเท วนฺทิตฺวา อุกฺกุฏิกํ นิสีทิตฺวา อญฺชลิํ ปคฺคเหตฺวา เอวมสฺส วจนีโย “อหํ ภนฺเต สํเฆน อาปตฺติยา อปฺปฏิกมฺเม อุกฺเขปนียกมฺมกโต สมฺมา วตฺตามิ, โลมํ ปาเตมิ, เนตฺถารํ วตฺตามิ, อาปตฺติยา อปฺปฏิกมฺเม อุกฺเขปนียสฺส กมฺมสฺส ปฏิปฺปสฺสทฺธิํ ยาจามี”ติ.\csromanspacer{} ทุติยมฺปิ ยาจิตพฺพา.\csromanspacer{} ตติยมฺปิ ยาจิตพฺพา.\csromanspacer{} พฺยตฺเตน ภิกฺขุนา ปฏิพเลน สํโฆ ญาเปตพฺโพ\csromandash{}}

\prose{“สุณาตุ เม ภนฺเต สํโฆ, อยํ ฉนฺโน ภิกฺขุ สํเฆน อาปตฺติยา อปฺปฏิกมฺเม อุกฺเขปนียกมฺมกโต สมฺมา วตฺตติ, โลมํ ปาเตติ, เนตฺถารํ วตฺตติ, อาปตฺติยา อปฺปฏิกมฺเม อุกฺเขปนียสฺส กมฺมสฺส ปฏิปฺปสฺสทฺธิํ ยาจติ, ยทิ สํฆสฺส ปตฺตกลฺลํ, สํโฆ ฉนฺนสฺส ภิกฺขุโน อาปตฺติยา อปฺปฏิกมฺเม อุกฺเขปนียกมฺมํ ปฏิปฺปสฺสมฺเภยฺย, เอสา ญตฺติ.}

\prose{สุณาตุ เม ภนฺเต สํโฆ, อยํ ฉนฺโน ภิกฺขุ สํเฆน อาปตฺติยา อปฺปฏิกมฺเม อุกฺเขปนียกมฺมกโต สมฺมา วตฺตติ, โลมํ ปาเตติ, เนตฺถารํ วตฺตติ, อาปตฺติยา อปฺปฏิกมฺเม อุกฺเขปนียสฺส กมฺมสฺส ปฏิปฺปสฺสทฺธิํ ยาจติ, สํโฆ ฉนฺนสฺส ภิกฺขุโน อาปตฺติยา อปฺปฏิกมฺเม อุกฺเขปนียกมฺมํ ปฏิปฺปสฺสมฺเภติ, ยสฺสายสฺมโต ขมติ ฉนฺนสฺส ภิกฺขุโน อาปตฺติยา อปฺปฏิกมฺเม อุกฺเขปนียสฺส กมฺมสฺส ปฏิปฺปสฺสทฺธิ, โส ตุณฺหสฺส, ยสฺส นกฺขมติ, โส ภาเสยฺย.}

\prose{\csromanfolio{68}ทุติยมฺปิ เอตมตฺถํ วทามิ ฯเปฯ.\csromanspacer{} ตติยมฺปิ เอตมตฺถํ วทามิ.\csromanspacer{} สุณาตุ เม ภนฺเต สํโฆ, อยํ ฉนฺโน ภิกฺขุ สํเฆน อาปตฺติยา อปฺปฏิกมฺเม อุกฺเขปนียกมฺมกโต สมฺมา วตฺตติ, โลมํ ปาเตติ, เนตฺถารํ วตฺตติ, อาปตฺติยา อปฺปฏิกมฺเม อุกฺเขปนียสฺส กมฺมสฺส ปฏิปฺปสฺสทฺธิํ ยาจติ, สํโฆ ฉนฺนสฺส ภิกฺขุโน อาปตฺติยา อปฺปฏิกมฺเม อุกฺเขปนียกมฺมํ ปฏิปฺปสฺสมฺเภติ, ยสฺสายสฺมโต ขมติ ฉนฺนสฺส ภิกฺขุโน อาปตฺติยา อปฺปฏิกมฺเม อุกฺเขปนียสฺส กมฺมสฺส ปฏิปฺปสฺสทฺธิ, โส ตุณฺหสฺส, ยสฺส นกฺขมติ, โส ภาเสยฺย.}

\prose{ปฏิปฺปสฺสทฺธํ สํเฆน ฉนฺนสฺส ภิกฺขุโน อาปตฺติยา อปฺปฏิกมฺเม อุกฺเขปนียกมฺมํ, ขมติ สํฆสฺส, ตสฺมา ตุณฺหี, เอวเมตํ ธารยามี”ติ.}

\nitthitam{อาปตฺติยา อปฺปฏิกมฺเม อุกฺเขปนียกมฺมํ นิฏฺฐิตํ ฉฏฺฐํ.}
\csromanmatikaanchor{mtk.44}
\csromantocmarkref{section}{7. ปาปิกาย ทิฏฺฐิยา อปฺปฏินิสฺสคฺเค อุกฺเขปนียกมฺม}{mtk.44}
\bookmark[dest={mtk.44},level=1]{7. ปาปิกาย ทิฏฺฐิยา อปฺปฏินิสฺสคฺเค อุกฺเขปนียกมฺม}
\csromanheader{1}{7. ปาปิกาย ทิฏฺฐิยา อปฺปฏินิสฺสคฺเค อุกฺเขปนียกมฺม}
\proseitem{65}{\csromansymbolfootnote{*}{อิทํ วตฺถุ วิ 2. 175; ม 1. 182 ปิฏฺฐาทีสุปิ อาคตํ.}เตน สมเยน พุทฺโธ ภควา สาวตฺถิยํ วิหรติ เชตวเน อนาถปิณฺฑิกสฺส อาราเม.\csromanspacer{} เตน โข ปน สมเยน อริฏฺฐสฺส นาม ภิกฺขุโน คทฺธพาธิปุพฺพสฺส\footnote{คนฺธพาธิปุพฺพสฺส (ก)} เอวรูปํ ปาปกํ ทิฏฺฐิคตํ อุปฺปนฺนํ โหติ “ตถาหํ ภควตา ธมฺมํ เทสิตํ อาชานามิ, ยถา เยเม อนฺตรายิกา ธมฺมา วุตฺตา ภควตา, เต ปฏิเสวโต นาลํ อนฺตรายายา”ติ.\csromanspacer{} อสฺโสสุํ โข สมฺพหุลา ภิกฺขู “อริฏฺฐสฺส นาม กิร ภิกฺขุโน คทฺธพาธิปุพฺพสฺส เอวรูปํ ปาปกํ ทิฏฺฐิคตํ อุปฺปนฺนํ ‘ตถาหํ ภควตา ธมฺมํ เทสิตํ อาชานามิ, ยถา เยเม อนฺตรายิกา ธมฺมา วุตฺตา ภควตา, เต ปฏิเสวโต นาลํ อนฺตรายายา’ติ”.\csromanspacer{} อถ โข เต ภิกฺขู เยน อริฏฺโฐ ภิกฺขุ คทฺธพาธิปุพฺโพ เตนุปสงฺกมิํสุ, อุปสงฺกมิตฺวา อริฏฺฐํ ภิกฺขุํ คทฺธพาธิปุพฺพํ เอตทโวจุํ “สจฺจํ กิร เต อาวุโส อริฏฺฐ เอวรูปํ ปาปกํ ทิฏฺฐิคตํ อุปฺปนฺนํ ‘ตถาหํ ภควตา ธมฺมํ เทสิตํ อาชานามิ, ยถา เยเม อนฺตรายิกา ธมฺมา วุตฺตา ภควตา, เต ปฏิเสวโต นาลํ อนฺตรายายา’ติ”.\csromanspacer{} เอวํพฺยา โข อหํ อาวุโส ภควตา ธมฺมํ เทสิตํ อาชานามิ, ยถา \csromanfolio{69}เยเม อนฺตรายิกา ธมฺมา วุตฺตา ภควตา, เต ปฏิเสวโต นาลํ อนฺตรายายาติ.}

\prose{มาวุโส อริฏฺฐ เอวํ อวจ, มา ภควนฺตํ อพฺภาจิกฺขิ, น หิ สาธุ ภควโต อพฺภกฺขานํ\footnote{อพฺภาจิกฺขนํ (ก)}, น หิ ภควา เอวํ วเทยฺย, อเนกปริยาเยนาวุโส อริฏฺฐ อนฺตรายิกา ธมฺมา อนฺตรายิกา วุตฺตา ภควตา, อลญฺจ ปน เต ปฏิเสวโต อนฺตรายาย.\csromanspacer{} อปฺปสฺสาทา กามา วุตฺตา ภควตา พหุทุกฺขา พหุปายาสา\footnote{พหูปายายา (สี, สฺยา)}, อาทีนโว เอตฺถ ภิยฺโย.\csromanspacer{} อฏฺฐิกงฺกลูปมา กามา วุตฺตา ภควตา พหุทุกฺขา พหุปายาสา, อาทีนโว เอตฺถ ภิยฺโย.\csromanspacer{} มํสเปสูปมา กามา วุตฺตา ภควตา ฯเปฯ.\csromanspacer{} ติณุกฺกูปมา กามา วุตฺตา ภควตา.\csromanspacer{} องฺคารกาสูปมา กามา วุตฺตา ภควตา.\csromanspacer{} สุปินกูปมา กามา วุตฺตา ภควตา.\csromanspacer{} ยาจิตกูปมา กามา วุตฺตา ภควตา.\csromanspacer{} รุกฺขผลูปมา กามา วุตฺตา ภควตา.\csromanspacer{} อสิสูนูปมา กามา วุตฺตา ภควตา.\csromanspacer{} สตฺติสูลูปมา กามา วุตฺตา ภควตา.\csromanspacer{} สปฺปสิรูปมา กามา วุตฺตา ภควตา พหุทุกฺขา พหุปายาสา, อาทีนโว เอตฺถ ภิยฺโยติ.}

\prose{เอวมฺปิ โข อริฏฺโฐ ภิกฺขุ คทฺธพาธิปุพฺโพ เตหิ ภิกฺขูหิ วุจฺจมาโน ตเถว ตํ ปาปกํ ทิฏฺฐิคตํ ถามสา ปรามาสา อภินิวิสฺส โวหรติ “เอวํพฺยา โข อหํ อาวุโส ภควตา ธมฺมํ เทสิตํ อาชานามิ, ยถา เยเม อนฺตรายิกา ธมฺมา วุตฺตา ภควตา, เต ปฏิเสวโต นาลํ อนฺตรายายา”ติ.\csromanspacer{} ยโต จ โข เต ภิกฺขู นาสกฺขิํสุ อริฏฺฐํ ภิกฺขุํ คทฺธพาธิปุพฺพํ เอตสฺมา ปาปกา ทิฏฺฐิคตา วิเวเจตุํ, อถ โข เต ภิกฺขู เยน ภควา เตนุปสงฺกมิํสุ, อุปสงฺกมิตฺวา ภควโต เอตมตฺถํ อาโรเจสุํ.\csromanspacer{} อถ โข ภควา เอตสฺมิํ นิทาเน เอตสฺมิํ ปกรเณ ภิกฺขุสํฆํ สนฺนิปาตาเปตฺวา อริฏฺฐํ ภิกฺขุํ คทฺธพาธิปุพฺพํ ปฏิปุจฺฉิ “สจฺจํ กิร เต อริฏฺฐ เอวรูปํ ปาปกํ ทิฏฺฐิคตํ อุปฺปนฺนํ ‘ตถาหํ ภควตา ธมฺมํ เทสิตํ อาชานามิ, ยถา เยเม อนฺตรายิกา ธมฺมา วุตฺตา ภควตา, เต ปฏิเสวโต นาลํ อนฺตรายายา’ติ”.\csromanspacer{} เอวํพฺยา โข อหํ ภนฺเต ภควตา ธมฺมํ เทสิตํ อาชานามิ, ยถา เยเม อนฺตรายิกา ธมฺมา วุตฺตา ภควตา, เต ปฏิเสวโต นาลํ อนฺตรายายาติ.}

\prose{\csromanfolio{70}กสฺส นุ โข นาม ตฺวํ โมฆปุริส มยา เอวํ ธมฺมํ เทสิตํ อาชานาสิ.\csromanspacer{} นนุ มยา โมฆปุริส อเนกปริยาเยน อนฺตรายิกา ธมฺมา อนฺตรายิกา วุตฺตา, อลญฺจ ปน เต ปฏิเสวโต อนฺตรายาย.\csromanspacer{} อปฺปสฺสาทา กามา วุตฺตา มยา พหุทุกฺขา พหุปายาสา, อาทีนโว เอตฺถ ภิยฺโย.\csromanspacer{} อฏฺฐิกงฺกลูปมา กามา วุตฺตา มยา พหุทุกฺขา พหุปายาสา, อาทีนโว เอตฺถ ภิยฺโย.\csromanspacer{} มํสเปสูปมา กามา วุตฺตา มยา ฯเปฯ.\csromanspacer{} ติณุกฺกูปมา กามา วุตฺตา มยา.\csromanspacer{} องฺคารกาสูปมา กามา วุตฺตา มยา.\csromanspacer{} สุปินกูปมา กามา วุตฺตา มยา.\csromanspacer{} ยาจิตกูปมา กามา วุตฺตา มยา.\csromanspacer{} รูกฺขผลูปมา กามา วุตฺตา มยา.\csromanspacer{} อสิสูนูปมา กามา วุตฺตา มยา.\csromanspacer{} สตฺติสูลูปมา กามา วุตฺตา มยา.\csromanspacer{} สปฺปสิรูปมา กามา วุตฺถา มยา พหุทุกฺขา พหุปายาสา, อาทีนโว เอตฺถ ภิยฺโย.\csromanspacer{} อถ จ ปน ตฺวํ โมฆปุริส อตฺตนา ทุคฺคหิเตน\footnote{ทุคฺคหิเตน ทิฏฺฐิคเตน (สฺยา)} อเมฺห เจว อพฺภาจิกฺขสิ, อตฺตานญฺจ ขณสิ, พหุญฺจ อปุญฺญํ ปสวสิ, ตญฺหิ เต โมฆปุริส ภวิสฺสติ ทีฆรตฺตํ อหิตาย ทุกฺขาย.\csromanspacer{} เนตํ โมฆปุริส อปฺปสนฺนานํ วา ปสาทาย ฯเปฯ วิครหิตฺวา ฯเปฯ ธมฺมิํ กถํ กตฺวา ภิกฺขู อามนฺเตสิ\csromandash{}เตน หิ ภิกฺขเว สํโฆ อริฏฺฐสฺส ภิกฺขุโน คทฺธพาธิปุพฺพสฺส ปาปิกาย ทิฏฺฐิยา อปฺปฏินิสฺสคฺเค อุกฺเขปนียกมฺมํ กโรตุ อสมฺโภคํ สํเฆน.\csromanspacer{} เอวญฺจ ปน ภิกฺขเว กาตพฺพํ, ปฐมํ อริฏฺโฐ ภิกฺขุ โจเทตพฺโพ, โจเทตฺวา สาเรตพฺโพ, สาเรตฺวา อาปตฺติํ อาโรเปตพฺโพ, อาปตฺติํ อาโรเปตฺวา พฺยตฺเตน ภิกฺขุนา ปฏิพเลน สํโฆ ญาเปตพฺโพ\csromandash{}}

\proseitem{66}{“สุณาตุ เม ภนฺเต สํโฆ, อริฏฺฐสฺส ภิกฺขุโน คทฺธพาธิปุพฺพสฺส เอวรูปํ ปาปกํ ทิฏฺฐิคตํ อุปฺปนฺนํ ‘ตถาหํ ภควตา ธมฺมํ เทสิตํ อาชานามิ, ยถา เยเม อนฺตรายิกา ธมฺมา วุตฺตา ภควตา, เต ปฏิเสวโต นาลํ อนฺตรายายา’ติ.\csromanspacer{} โส ตํ ทิฏฺฐิํ น ปฏินิสฺสชฺชติ, ยทิ สํฆสฺส ปตฺตกลฺลํ, สํโฆ อริฏฺฐสฺส ภิกฺขุโน คทฺธพาธิปุพฺพสฺส ปาปิกาย ทิฏฺฐิยา อปฺปฏินิสฺสคฺเค อุกฺเขปนียกมฺมํ กเรยฺย อสมฺโภคํ สํเฆน, เอสา ญตฺติ.}

\prose{สุณาตุ เม ภนฺเต สํโฆ, อริฏฺฐสฺส ภิกฺขุโน คทฺธพาธิปุพฺพสฺส เอวรูปํ ปาปกํ ทิฏฺฐิคตํ อุปฺปนฺนํ ‘ตถาหํ ภควตา ธมฺมํ เทสิตํ อาชานามิ, \csromanfolio{71}ยถา เยเม อนฺตรายิกา ธมฺมา วุตฺตา ภควตา, เต ปฏิเสวโต นาลํ อนฺตรายายา’ติ.\csromanspacer{} โส ตํ ทิฏฺฐิํ น ปฏินิสฺสชฺชติ, สํโฆ อริฏฺฐสฺส ภิกฺขุโน คทฺธพาธิปุพฺพสฺส ปาปิกาย ทิฏฺฐิยา อปฺปฏินิสฺสคฺเค อุกฺเขปนียกมฺมํ กโรติ อสมฺโภคํ สํเฆน, ยสฺสายสฺมโต ขมติ อริฏฺฐสฺส ภิกฺขุโน คทฺธพาธิปุพฺพสฺส ปาปิกาย ทิฏฺฐิยา อปฺปฏินิสฺสคฺเค อุกฺเขปนียสฺส กมฺมสฺส กรณํ อสมฺโภคํ สํเฆน, โส ตุณฺหสฺส, ยสฺส นกฺขมติ, โส ภาเสยฺย.}

\prose{ทุติยมฺปิ เอตมตฺวํ วทามิ ฯเปฯ.\csromanspacer{} ตติยมฺปิ เอตมตฺถํ วทามิ.\csromanspacer{} สุณาตุ เม ภนฺเต สํโฆ, อริฏฺฐสฺส ภิกฺขุโน คทฺธพาธิปุพฺพสฺส เอวรูปํ ปาปกํ ทิฏฺฐิคตํ อุปฺปนฺนํ ‘ตถาหํ ภควตา ธมฺมํ เทสิตํ อาชานามิ, ยถา เยเม อนฺตรายิกา ธมฺมา วุตฺตา ภควตา, เต ปฏิเสวโต นาลํ อนฺตรายายา’ติ, โส ตํ ทิฏฺฐิํ น ปฏินิสฺสชฺชติ, สํโฆ อริฏฺฐสฺส ภิกฺขุโน คทฺธพาธิปุพฺพสฺส ปาปิกาย ทิฏฺฐิยา อปฺปฏินิสฺสคฺเค อุกฺเขปนียกมฺมํ กโรติ อสมฺโภคํ สํเฆน, ยสฺสายสฺมโต ขมติ อริฏฺฐสฺส ภิกฺขุโน คทฺธพาธิปุพฺพสฺส ปาปิกาย ทิฏฺฐิยา อปฺปฏินิสฺสคฺเค อุกฺเขปนียสฺส กมฺมสฺส กรณํ อสมฺโภคํ สํเฆน, โส ตุณฺหสฺส, ยสฺส นกฺขมติ, โส ภาเสยฺย.}

\prose{กตํ สํเฆน อริฏฺฐสฺส ภิกฺขุโน คทฺธพาธิปุพฺพสฺส ปาปิกาย ทิฏฺฐิยา อปฺปฏินิสฺสคฺเค อุกฺเขปนียกมฺมํ อสมฺโภคํ สํเฆน, ขมติ สํฆสฺส, ตสฺมา ตุณฺหี, เอวเมตํ ธารยามี”ติ.}

\prose{อาวาสปรมฺปรญฺจ ภิกฺขเว สํสถ “อริฏฺโฐ ภิกฺขุ คทฺธพาธิปุพฺโพ สํเฆน ปาปิกาย\footnote{คทฺธพาธิปุพฺโพ ปาปิกาย (สี, ก)} ทิฏฺฐิยา อปฺปฏินิสฺสคฺเค อุกฺเขปนียกมฺมกโต อสมฺโภคํ สํเฆนา”ติ.}

\csromanmatikaanchor{mtk.45}
\csromantocmarkref{subsection}{อธมฺมกมฺมทฺวาทสก}{mtk.45}
\bookmark[dest={mtk.45},level=2]{อธมฺมกมฺมทฺวาทสก}
\csromanheader{2}{\textbf{อธมฺมกมฺมทฺวาทสก}}
\proseitem{67}{ตีหิ ภิกฺขเว องฺเคหิ สมนฺนาคตํ ปาปิกาย ทิฏฺฐิยา อปฺปฏินิสฺสคฺเค อุกฺเขปนียกมฺมํ อธมฺมกมฺมญฺจ โหติ อวินยกมฺมญฺจ ทุวูปสนฺตญฺจ.\csromanspacer{} อสมฺมุขา กตํ โหติ, อปฺปฏิปุจฺฉา กตํ โหติ, อปฺปฏิญฺญาย กตํ โหติ.\csromanspacer{} อิเมหิ โข ภิกฺขเว ตีหงฺเคหิ สมนฺนาคตํ \csromanfolio{72}ปาปิกาย ทิฏฺฐิยา อปฺปฏินิสฺสคฺเค อุกฺเขปนียกมฺมํ อธมฺมกมฺมญฺจ โหติ อวินยกมฺมญฺจ ทุวูปสนฺตญฺจ.}

\prose{อปเรหิปิ ภิกฺขเว ตีหงฺเคหิ สมนฺนาคตํ ปาปิกาย ทิฏฺฐิยา อปฺปฏินิสฺสคฺเค อุกฺเขปนียกมฺมํ อธมฺมกมฺมญฺจ โหติ อวินยกมฺมญฺจ ทุวูปสนฺตญฺจ.\csromanspacer{} อนาปตฺติยา กตํ โหติ, อเทสนาคามินิยา อาปตฺติยา กตํ โหติ, เทสิตาย อาปตฺติยา กตํ โหติ ฯเปฯ.\csromanspacer{} อโจเทตฺวา กตํ โหติ, อสาเรตฺวา กตํ โหติ, อาปตฺติํ อนาโรเปตฺวา กตํ โหติ ฯเปฯ.\csromanspacer{} อสมฺมุขา กตํ โหติ, อธมฺเมน กตํ โหติ, วคฺเคน กตํ โหติ ฯเปฯ.\csromanspacer{} อปฺปฏิปุจฺฉา กตํ โหติ, อธมฺเมน กตํ โหติ, วคฺเคน กตํ โหติ ฯเปฯ.\csromanspacer{} อปฺปฏิญฺญาย กตํ โหติ, อธมฺเมน กตํ โหติ, วคฺเคน กตํ โหติ ฯเปฯ.\csromanspacer{} อนาปตฺติยา กตํ โหติ, อธมฺเมน กตํ โหติ, วคฺเคน กตํ โหติ ฯเปฯ.\csromanspacer{} อเทสนาคามินิยา อาปตฺติยา กตํ โหติ, อธมฺเมน กตํ โหติ, วคฺเคน กตํ โหติ ฯเปฯ.\csromanspacer{} เทสิตาย อาปตฺติยา กตํ โหติ, อธมฺเมน กตํ โหติ, วคฺเคน กตํ โหติ ฯเปฯ.\csromanspacer{} อโจเทตฺวา กตํ โหติ, อธมฺเมน กตํ โหติ, วคฺเคน กตํ โหติ ฯเปฯ.\csromanspacer{} อาสาเรตฺวา กตํ โหติ, อธมฺเมน กตํ โหติ, วคฺเคน กตํ โหติ ฯเปฯ.\csromanspacer{} อาปตฺติํ อนาโรเปตฺวา กตํ โหติ, อธมฺเมน กตํ โหติ, วคฺเคน กตํ โหติ.\csromanspacer{} อิเมหิ โข ภิกฺขเว ตีหงฺเคหิ สมนฺนาคตํ ปาปิกาย ทิฏฺฐิยา อปฺปฏินิสฺสคฺเค อุกฺเขปนียกมฺมํ อธมฺมกมฺมญฺจ โหติ อวินยกมฺมญฺจ ทุวูปสนฺตญฺจ.}

\csromanheader{2}{ปาปิกาย ทิฏฺฐิยา อปฺปฏินิสฺสคฺเค อุกฺเขปนียกมฺเม}
\nitthitam{อธมฺมกมฺมทฺวาทสกํ นิฏฺฐิตํ.}
\csromanmatikaanchor{mtk.46}
\csromantocmarkref{subsection}{ธมฺมกมฺมทฺวาทสก}{mtk.46}
\bookmark[dest={mtk.46},level=2]{ธมฺมกมฺมทฺวาทสก}
\csromanheader{2}{\textbf{ธมฺมกมฺมทฺวาทสก}}
\proseitem{68}{ตีหิ ภิกฺขเว องฺเคหิ สมนฺนาคตํ ปาปิกาย ทิฏฺฐิยา อปฺปฏินิสฺสคฺเค อุกฺเขปนียกมฺมํ ธมฺมกมฺมญฺจ โหติ วินยกมฺมญฺจ สุวูปสนฺตญฺจ.\csromanspacer{} สมฺมุขา กตํ โหติ, ปฏิปุจฺฉา กตํ โหติ, ปฏิญฺญาย กตํ โหติ.\csromanspacer{} อิเมหิ โข ภิกฺขเว ตีหงฺเคหิ สมนฺนาคตํ ปาปิกาย ทิฏฺฐิยา อปฺปฏินิสฺสคฺเค อุกฺเขปนียกมฺมํ ธมฺมกมฺมญฺจ โหติ วินยกมฺมญฺจ สุวูปสนฺตญฺจ.}

\prose{\csromanfolio{73}อปเรหิปิ ภิกฺขเว ตีหงฺเคหิ สมนฺนาคตํ ปาปิกาย ทิฏฺฐิยา อปฺปฏินิสฺสคฺเค อุกฺเขปนียกมฺมํ ธมฺมกมฺมญฺจ โหติ วินยกมฺมญฺจ สุวูปสนฺตญฺจ.\csromanspacer{} อาปตฺติยา กตํ โหติ, เทสนาคามินิยา อาปตฺติยา กตํ โหติ, อเทสิตาย อาปตฺติยา กตํ โหติ ฯเปฯ.\csromanspacer{} โจเทตฺวา กตํ โหติ, สาเรตฺวา กตํ โหติ, อาปตฺติํ อาโรเปตฺวา กตํ โหติ ฯเปฯ.\csromanspacer{} สมฺมุขา กตํ โหติ, ธมฺเมน กตํ โหติ, สมคฺเคน กตํ โหติ ฯเปฯ.\csromanspacer{} ปฏิปุจฺฉา กตํ โหติ, ธมฺเมน กตํ โหติ, สมคฺเคน กตํ โหติ ฯเปฯ.\csromanspacer{} ปฏิญฺญาย กตํ โหติ, ธมฺเมน กตํ โหติ, สมคฺเคน กตํ โหติ ฯเปฯ.\csromanspacer{} อาปตฺติยา กตํ โหติ, ธมฺเมน กตํ โหติ, สมคฺเคน กตํ โหติ ฯเปฯ.\csromanspacer{} เทสนาคามินิยา อาปตฺติยา กตํ โหติ, ธมฺเมน กตํ โหติ, สมคฺเคน กตํ โหติ ฯเปฯ.\csromanspacer{} อเทสิตาย อาปตฺติยา กตํ โหติ, ธมฺเมน กตํ โหติ, สมคฺเคน กตํ โหติ ฯเปฯ.\csromanspacer{} โจเทตฺวา กตํ โหติ, ธมฺเมน กตํ โหติ, สมคฺเคน กตํ โหติ ฯเปฯ.\csromanspacer{} สาเรตฺวา กตํ โหติ, ธมฺเมน กตํ โหติ, สมคฺเคน กตํ โหติ ฯเปฯ.\csromanspacer{} อาปตฺติํ อาโรเปตฺวา กตํ โหติ, ธมฺเมน กตํ โหติ, สมคฺเคน กตํ โหติ.\csromanspacer{} อิเมหิ โข ภิกฺขเว ตีหงฺเคหิ สมนฺนาคตํ ปาปิกาย ทิฏฺฐิยา อปฺปฏินิสฺสคฺเค อุกฺเขปนียกมฺมํ ธมฺมกมฺมญฺจ โหติ วินยกมฺมญฺจ สุวูปสนฺตญฺจ.\csromanspacer{} ปาปิกาย ทิฏฺฐายา อปฺปฏินิสฺสคฺเค อุกฺเขปนียกมฺเม ธมฺมกมฺมทฺวาทสกํ นิฏฺฐิตํ.}

\csromanmatikaanchor{mtk.47}
\csromantocmarkref{subsection}{อากงฺขมานฉกฺก}{mtk.47}
\bookmark[dest={mtk.47},level=2]{อากงฺขมานฉกฺก}
\csromanheader{2}{\textbf{อากงฺขมานฉกฺก}}
\proseitem{69}{\csromansymbolfootnote{*}{วิ 5. 220 ปิฏฺเฐปิ.}ตีหิ ภิกฺขเว องฺเคหิ สมนฺนาคตสฺส ภิกฺขุโน อากงฺขมาโน สํโฆ ปาปิกาย ทิฏฺฐิยา อปฺปฏินิสฺสคฺเค อุกฺเขปนียกมฺมํ กเรยฺย.\csromanspacer{} ภณฺฑนการโก โหติ กลหการโก วิวาทการโก ภสฺสการโก สํเฆ อธิกรณการโก, พาโล โหติ อพฺยตฺโต อาปตฺติพหุโล อนปทาโน, คิหิสํสฏฺโฐ วิหรติ อนนุโลมิเกหิ คิหิสํสคฺเคหิ.\csromanspacer{} อิเมหิ โข ภิกฺขเว ตีหงฺเคหิ สมนฺนาคตสฺส ภิกฺขุโน อากงฺขมาโน สํโฆ ปาปิกาย ทิฏฺฐิยา อปฺปฏินิสฺสคฺเค อุกฺเขปนียกมฺมํ กเรยฺย. (1)}

\prose{\csromanfolio{74}อปเรหิปิ ภิกฺขเว ตีหงฺเคหิ สมนฺนาคตสฺส ภิกฺขุโน อากงฺขมาโน สํโฆ ปาปิกาย ทิฏฺฐิยา อปฺปฏินิสฺสคฺเค อุกฺเขปนียกมฺมํ กเรยฺย.\csromanspacer{} อธิสีเล สีลวิปนฺโน โหติ, อชฺฌาจาเร อาจารวิปนฺโน โหติ, อติทิฏฺฐิยา ทิฏฺฐิวิปนฺโน โหติ.\csromanspacer{} อิเมหิ โข ภิกฺขเว ตีหงฺเคหิ สมนฺนาคตสฺส ภิกฺขุโน อากงฺขมาโน สํโฆ ปาปิกาย ทิฏฺฐิยา อปฺปฏินิสฺสคฺเค อุกฺเขปนียกมฺมํ กเรยฺย. (2)}

\prose{อปเรหิปิ ภิกฺขเว ตีหงฺเคหิ สมนฺนาคตสฺส ภิกฺขุโน อากงฺขมาโน สํโฆ ปาปิกาย ทิฏฺฐิยา อปฺปฏินิสฺสคฺเค อุกฺเขปนียกมฺมํ กเรยฺย.\csromanspacer{} พุทฺธสฺส อวณฺณํ ภาสติ, ธมฺมสฺส อวณฺณํ ภาสติ, สํฆสฺส อวณฺณํ ภาสติ.\csromanspacer{} อิเมหิ โข ภิกฺขเว ตีหงฺเคหิ สมนฺนาคตสฺส ภิกฺขุโน อากงฺขมาโน สํโฆ ปาปิกาย ทิฏฺฐิยา อปฺปฏินิสฺสคฺเค อุกฺเขปนียกมฺมํ กเรยฺย. (3)}

\prose{ติณฺณํ ภิกฺขเว ภิกฺขูนํ อากงฺขมาโน สํโฆ ปาปิกาย ทิฏฺฐิยา อปฺปฏินิสฺสคฺเค อุกฺเขปนียกมฺมํ กเรยฺย.\csromanspacer{} เอโก ภณฺฑนการโก โหติ กลหการโก วิวาทการโก ภสฺสการโก สํเฆ อธิกรณการโก, เอโก พาโล โหติ อพฺยตฺโต อาปตฺติพหุโล อนปทาโน, เอโก คิหิสํสฏฺโฐ วิหรติ อนนุโลมิเกหิ คิหิสํสคฺเคหิ.\csromanspacer{} อิเมสํ โข ภิกฺขเว ติณฺณํ ภิกฺขูนํ อากงฺขมาโน สํโฆ ปาปิกาย ทิฏฺฐิยา อปฺปฏินิสฺสคฺเค อุกฺเขปนียกมฺมํ กเรยฺย. (4)}

\prose{อปเรสมฺปิ ภิกฺขเว ติณฺณํ ภิกฺขูนํ อากงฺขมาโน สํโฆ ปาปิกาย ทิฏฺฐิยา อปฺปฏินิสฺสคฺเค อุกฺเขปนียกมฺมํ กเรยฺย.\csromanspacer{} เอโก อธิสีเล สีลวิปนฺโน โหติ, เอโก อชฺฌาจาเร อาจารวิปนฺโน โหติ, เอโก อติทิฏฺฐิยา ทิฏฺฐิวิปนฺโน โหติ.\csromanspacer{} อิเมสํ โข ภิกฺขเว ติณฺณํ ภิกฺขูนํ อากงฺขมาโน สํโฆ ปาปิกาย ทิฏฺฐิยา อปฺปฏินิสฺสคฺเค อุกฺเขปนียกมฺมํ กเรยฺย. (5)}

\prose{อปเรสมฺปิ ภิกฺขเว ติณฺณํ ภิกฺขูนํ อากงฺขมาโน สํโฆ ปาปิกาย ทิฏฺฐิยา อปฺปฏินิสฺสคฺเค อุกฺเขปนียกมฺมํ กเรยฺย.\csromanspacer{} เอโก พุทฺธสฺส อวณฺณํ ภสติ, เอโก ธมฺมสฺส อวณฺณํ ภาสติ, เอโก สํฆสฺส อวณฺณํ \csromanfolio{75}ภาสติ.\csromanspacer{} อิเมสํ โข ภิกฺขเว ติณฺณํ ภิกฺขูนํ อากงฺขมาโน สํโฆ ปาปิกาย ทิฏฺฐิยา อปฺปฏินิสฺสคฺเค อุกฺเขปนียกมฺมํ กเรยฺย. (6) ปาปิกาย ทิฏฺฐิยา อปฺปฏินิสฺสคฺเค อุกฺเขปนียกมฺเม อกงฺขมานฉกฺกํ นิฏฺฐิตํ.}

\csromanmatikaanchor{mtk.48}
\csromantocmarkref{subsection}{เตจตฺตาลีสวตฺต}{mtk.48}
\bookmark[dest={mtk.48},level=2]{เตจตฺตาลีสวตฺต}
\csromanheader{2}{\textbf{เตจตฺตาลีสวตฺต}}
\vspace{\tipitakaparskip}
\csromancenter{ปาปิกาย ทิฏฺฐิยา อปฺปฏินิสฺสคฺเค อุกฺเขปนียกมฺมกเตน ภิกฺขเว ภิกฺขุนา สมฺมา วตฺติตพฺพํ.\csromanspacer{} ตตฺรายํ สมฺมาวตฺตนา\csromandash{}น อุปสมฺปาเทตพฺพํ, น นิสฺสโย ทาตพฺโพ, น สามเณโร อุปฏฺฐาเปตพฺโพ, น ภิกฺขุโนวาทกสมฺมุติ สาทิตพฺพา, สมฺมเตนปิ ภิกฺขุนิโย น โอวทิตพฺพา, ยาย อาปตฺติยา สํเฆน ปาปิกาย ทิฏฺฐิยา อปฺปฏินิสฺสคฺเค อุกฺเขปนียกมฺมํ กตํ โหติ, สา อาปตฺติ น อาปชฺชิตพฺพา, อญฺญา วา ตาทิสิกา, ตโต วา ปาปิฏฺฐตรา, กมฺมํ น ครหิตพฺพํ, กมฺมิกา น ครหิตพฺพา - ป-.\csromanspacer{} น ปกตตฺตสฺส ภิกฺขุโน อุโปสโถ ฐเปตพฺโพ, น ปวารณา ฐเปตพฺพา, น สวจนียํ กาตพฺพํ, น อนุวาโท ปฏฺฐเปตพฺโพ, น โอกาโส กาเรตพฺโพ, น โจเทตพฺโพ, น สาเรตพฺโพ, น ภิกฺขูหิ สมฺปโยเชตพฺพนฺติ.\csromanspacer{} ปาปิกาย ทิฏฺฐิยา อปฺปฏินิสฺสคฺเค อุกฺเขปนียกมฺเม \textbf{เตจตฺตาลีสวตฺต}ํ นิฏฺฐิตํ.}
\proseitem{71}{อถ โข สํโฆ อริฏฺฐสฺส ภิกฺขุโน คทฺธพาธิปุพฺพสฺส ปาปิกาย ทิฏฺฐิยา อปฺปฏินิสฺสคฺเค อุกฺเขปนียกมฺมํ อกาสิ อสมฺโภคํ สํเฆน.\csromanspacer{} โส สํเฆน ปาปิกาย ทิฏฺฐิยา อปฺปฏินิสฺสคฺเค อุกฺเขปนียกมฺมกโต วิพฺภมิ.\csromanspacer{} เย เต ภิกฺขู อปฺปิจฺฉา ฯเปฯ เต อุชฺฌายนฺติ ขิยฺยนฺติ วิปาเจนฺติ “กถํ หิ นาม อริฏฺโฐ ภิกฺขุ คทฺธพาธิปุพฺโพ สํเฆน ปาปิกาย ทิฏฺฐิยา อปฺปฏินิสฺสคฺเค อุกฺเขปนียกมฺมกโต วิพฺภมิสฺสตี”ติ.\csromanspacer{} อถ โข เต ภิกฺขู ภควโต เอตมตฺถํ อาโรเจสุํ.}

\prose{อถ โข ภควา เอตสฺมิํ นิทาเน เอตสฺมิํ ปกรเณ ภิกฺขุสํฆํ สนฺนิปาตาเปตฺวา ภิกฺขู ปฏิปุจฺฉิ “สจฺจํ กิร ภิกฺขเว อริฏฺโฐ ภิกฺขุ \csromanfolio{76}คทฺธพาธิปุพฺโพ สํเฆน ปาปิกาย ทิฏฺฐิยา อปฺปฏินิสฺสคฺเค อุกฺเขปนียกมฺมกโต วิพฺภมตี”ติ\footnote{วิพฺภมีติ (สี, ก)}.\csromanspacer{} สจฺจํ ภควาติ.\csromanspacer{} วิครหิ พุทฺโธ ภควา อนนุจฺฉวิกํ ฯเปฯ.\csromanspacer{} กถํ หิ นาม โส ภิกฺขเว โมฆปุริโส สํเฆน ปาปิกาย ทิฏฺฐิยา อปฺปฏินิสฺสคฺเค อุกฺเขปนียกมฺมกโต วิพฺภมิสฺสติ.\csromanspacer{} เนตํ ภิกฺขเว อปฺปสนฺนานํ วา ปสาทาย ฯเปฯ วิครหิตฺวา ฯเปฯ ธมฺมิํ กถํ กตฺวา ภิกฺขู อามนฺเตสิ\csromandash{} เตน หิ ภิกฺขเว สํโฆ ปาปิกาย ทิฏฺฐิยา อปฺปฏินิสฺสคฺเค อุกฺเขปนียกมฺมํ ปฏิปฺปสฺสมฺเภตุ.\csromanspacer{} \textbf{นปฺปฏิปฺปสฺสมฺเภตพฺพเตจตฺตาลีสก}}

\proseitem{72}{ปญฺจหิ ภิกฺขเว องฺเคหิ สมนฺนาคตสฺส ภิกฺขุโน ปาปิกาย ทิฏฺฐิยา อปฺปฏินิสฺสคฺเค อุกฺเขปนียกมฺมํ นปฺปฏิปฺปสฺสมฺเภตพฺพํ.\csromanspacer{} อุปสมฺปาเทติ, นิสฺสยํ เทติ, สามเณรํ อุปฏฺฐาเปติ, ภิกฺขุโนวาทกสมฺมุติํ สาทิยติ, สมฺมโตปิ ภิกฺขุนิโย โอวทติ.\csromanspacer{} อิเมหิ โข ภิกฺขเว ปญฺจหงฺเคหิ สมนฺนาคตสฺส ภิกฺขุโน ปาปิกาย ทิฏฺฐิยา อปฺปฏินิสฺสคฺเค อุกฺเขปนียกมฺมํ นปฺปฏิปฺปสฺสมฺเภตพฺพํ.}

\prose{\csromansymbolfootnote{*}{วิ 5. 317 ปิฏฺเฐปิ.}อปเรหิปิ ภิกฺขเว ปญฺจหงฺเคหิ สมนฺนาคตสฺส ภิกฺขุโน ปาปิกาย ทิฏฺฐิยา อปฺปฏินิสฺสคฺเค อุกฺเขปนียกมฺมํ นปฺปฏิปฺปสฺสมฺเภ- ตพฺพํ.\csromanspacer{} ยาย อาปตฺติยา สํเฆน ปาปิกาย ทิฏฺฐิยา อปฺปฏินิสฺสคฺเค อุกฺเขปนียกมฺมํ กตํ โหติ, ตํ อาปตฺติํ อาปชฺชติ, อญฺญํ วา ตาทิสิกํ, ตโต วา ปาปิฏฺฐตรํ, กมฺมํ ครหติ, กมฺมิเก ครหติ.\csromanspacer{} อิเมหิ โข ภิกฺขเว ปญฺจหงฺเคหิ สมนฺนาคตสฺส ภิกฺขุโน ปาปิกาย ทิฏฺฐิยา อปฺปฏินิสฺสคฺเค อุกฺเขปนียกมฺมํ นปฺปฏิปฺปสฺสมฺเภตพฺพํ ฯเปฯ.}

\prose{อฏฺฐหิ ภิกฺขเว องฺเคหิ สมนฺนาคตสฺส ภิกฺขุโน ปาปิกาย ทิฏฺฐิยา อปฺปฏินิสฺสคฺเค อุกฺเขปนียกมฺมํ นปฺปฏิปฺปสฺสมฺเภตพฺพํ.\csromanspacer{} ปกตตฺตสฺส ภิกฺขุโน อุโปสถํ ฐเปติ, ปวารณํ ฐเปติ, สวจนียํ กโรติ, อนุวาทํ ปฏฺฐเปติ, โอกาสํ กาเรติ, โจเทติ, สาเรติ, ภิกฺขูหิ สมฺปโยเชติ.\csromanspacer{} อิเมหิ โข ภิกฺขเว อฏฺฐหงฺเคหิ สมนฺนาคตสฺส ภิกฺขุโน ปาปิกาย ทิฏฺฐิยา อปฺปฏินิสฺสคฺเค อุกฺเขปนียกมฺมํ นปฺปฏิปฺปสฺสมฺเภตพฺพํ.\csromanspacer{} ปาปิกาย ทิฏฺฐิยา อปฺปฏินิสฺสคฺเค อุกฺเขปนียกมฺเม \textbf{นปฺปฏิปฺปสฺสมฺเภตพฺพเตจตฺตาลีสก}ํ นิฏฺฐิตํ.}

\csromanmatikaanchor{mtk.49}
\csromanmatikaanchor{mtk.50}
\csromantocmarkref{subsection}{นปฺปฏิปฺปสฺสมฺเภตพฺพเตจตฺตาลีสก}{mtk.49}
\bookmark[dest={mtk.49},level=2]{นปฺปฏิปฺปสฺสมฺเภตพฺพเตจตฺตาลีสก}
\csromantocmarkref{subsection}{ปฏิปฺปสฺสมฺเภตพฺพเตจตฺตาลีสก}{mtk.50}
\bookmark[dest={mtk.50},level=2]{ปฏิปฺปสฺสมฺเภตพฺพเตจตฺตาลีสก}
\csromanheader{2}{\textbf{ปฏิปฺปสฺสมฺเภตพฺพเตจตฺตาลีสก}}
\proseitem{73}{\csromanfolio{77}ปญฺจหิ ภิกฺขเว องฺเคหิ สมนฺนาคตสฺส ภิกฺขุโน ปาปิกาย ทิฏฺฐิยา อปฺปฏินิสฺสคฺเค อุกฺเขปนียกมฺมํ ปฏิปฺปสฺสมฺเภตพฺพํ.\csromanspacer{} น อุปสมฺปาเทติ, น นิสฺสยํ เทติ, น สามเณรํ อุปฏฺฐาเปติ, น ภิกฺขุโนวาทกสมฺมุติํ สาทิยติ, สมฺมโตปิ ภิกฺขุนิโย น โอวทติ.\csromanspacer{} อิเมหิ โข ภิกฺขเว ปญฺจหงฺเคหิ สมนฺนาคตสฺส ภิกฺขุโน ปาปิกาย ทิฏฺฐิยา อปฺปฏินิสฺสคฺเค อุกฺเขปนียกมฺมํ ปฏิปฺปสฺสมฺเภตพฺพํ.}

\prose{อปเรหิปิ ภิกฺขเว ปญฺจหงฺเคหิ สมนฺนาคตสฺส ภิกฺขุโน ปาปิกาย ทิฏฺฐิยา อปฺปฏินิสฺสคฺเค อุกฺเขปนียกมฺมํ ปฏิปฺปสฺสมฺเภตพฺพํ.\csromanspacer{} ยาย อาปตฺติยา สํเฆน ปาปิกาย ทิฏฺฐิยา อปฺปฏินิสฺสคฺเค อุกฺเขปนียกมฺมํ กตํ โหติ, ตํ อาปตฺติํ น อาปชฺชติ, อญฺญํ วา ตาทิสิกํ, ตโต วา ปาปิฏฺฐตรํ, กมฺมํ น ครหติ, กมฺมิเก น ครหติ.\csromanspacer{} อิเมหิ โข ภิกฺขเว ปญฺจหงฺเคหิ สมนฺนาคตสฺส ภิกฺขุโน ปาปิกาย ทิฏฺฐิยา อปฺปฏินิสฺสคฺเค อุกฺเขปนียกมฺมํ ปฏิปฺปสฺสมฺเภตพฺพํ ฯเปฯ.}

\prose{อฏฺฐหิ ภิกฺขเว องฺเคหิ สมนฺนาคตสฺส ภิกฺขุโน ปาปิกาย ทิฏฺฐิยา อปฺปฏินิสฺสคฺเค อุกฺเขปนียกมฺมํ ปฏิปฺปสฺสมฺเภตพฺพํ.\csromanspacer{} น ปกตตฺตสฺส ภิกฺขุโน อุโปสถํ ฐเปติ, น ปวารณํ ฐเปติ, น สวจนียํ กโรติ, น อนุวาทํ ปฏฺฐเปติ, น โอกาสํ กาเรติ, น โจเทติ, น สาเรติ, น ภิกฺขูหิ สมฺปโยเชติ.\csromanspacer{} อิเมหิ โข ภิกฺขเว อฏฺฐหงฺเคหิ สมนฺนาคตสฺส ภิกฺขุโน ปาปิกาย ทิฏฺฐิยา อปฺปฏินิสฺสคฺเค อุกฺเขปนียกมฺมํ ปฏิปฺปสฺสมฺเภตพฺพํ.\csromanspacer{} ปาปิกาย ทิฏฺฐิยา อปฺปฏินิสฺสคฺเค อุกฺเขปนียกมฺเม \textbf{ปฏิปฺปสฺสมฺเภตพฺพเตจตฺตาลีสก}ํ นิฏฺฐิตํ.}

\proseitem{74}{เอวญฺจ ปน ภิกฺขเว ปฏิปฺปสฺสมฺเภตพฺพํ, เตน ภิกฺขเว ปาปิกาย ทิฏฺฐิยา อปฺปฏินิสฺสคฺเค อุกฺเขปนียกมฺมกเตน ภิกฺขุนา สํฆํ อุปสงฺกมิตฺวา เอกํสํ อุตฺตราสงฺคํ กริตฺวา วุฑฺฒานํ ภิกฺขูนํ ปาเท วนฺทิตฺวา อุกฺกุฏิกํ นิสีทิตฺวา อญฺชลิํ ปคฺคเหตฺวา เอวมสฺส วจนีโย “อหํ ภนฺเต สํเฆน ปาปิกาย ทิฏฺฐิยา อปฺปฏินิสฺสคฺเค อุกฺเขปนียกมฺมกโต สมฺมา วตฺตามิ, โลมํ ปาเตมิ, เนตฺถารํ วตฺตามิ, ปาปิกาย ทิฏฺฐิยา อปฺปฏินิสฺสคฺเค \csromanfolio{78}อุกฺเขปนียสฺส กมฺมสฺส ปฏิปฺปสฺสทฺธิํ ยาจามี”ติ.\csromanspacer{} ทุติยมฺปิ ยาจิตพฺพา.\csromanspacer{} ตติยมฺปิ ยาจิตพฺพา.\csromanspacer{} พฺยตฺเตน ภิกฺขุนา ปฏิพเลน สํโฆ ญาเปตพฺโพ\csromandash{}}

\prose{“สุณาตุ เม ภนฺเต สํโฆ, อยํ อิตฺถนฺนาโม ภิกฺขุ สํเฆน ปาปิกาย ทิฏฺฐิยา อปฺปฏินิสฺสคฺเค อุกฺเขปนียกมฺมกโต สมฺมา วตฺตติ, โลมํ ปาเตติ, เนตฺถารํ วตฺตติ, ปาปิกาย ทิฏฺฐิยา อปฺปฏินิสฺสคฺเค อุกฺเขปนียสฺส กมฺมสฺส ปฏิปฺปสฺสทฺธิํ ยาจติ, ยทิ สํฆสฺส ปตฺตกลฺลํ, สํโฆ อิตฺถนฺนามสฺส ภิกฺขุโน ปาปิกาย ทิฏฺฐิยา อปฺปฏินิสฺสคฺเค อุกฺเขปนียกมฺมํ ปฏิปฺปสฺสมฺเภยฺย, เอสา ญตฺติ.}

\prose{สุณาตุ เม ภนฺเต สํโฆ, อยํ อิตฺถนฺนาโม ภิกฺขุ สํเฆน ปาปิกาย ทิฏฺฐิยา อปฺปฏินิสฺสคฺเค อุกฺเขปนียกมฺมกโต สมฺมา วตฺตติ, โลมํ ปาเตติ, เนตฺถารํ วตฺตติ, ปาปิกาย ทิฏฺฐิยา อปฺปฏินิสฺสคฺเค อุกฺเขปนียสฺส กมฺมสฺส ปฏิปฺปสฺสทฺธิํ ยาจติ, สํโฆ อิตฺถนฺนามสฺส ภิกฺขุโน ปาปิกาย ทิฏฺฐิยา อปฺปฏินิสฺสคฺเค อุกฺเขปนียกมฺมํ ปฏิปฺปสฺสมฺเภติ, ยสฺสายสฺมโต ขมติ อิตฺถนฺนามสฺส ภิกฺขุโน ปาปิกาย ทิฏฺฐิยา อปฺปฏินิสฺสคฺเค อุกฺเขปนียสฺส กมฺมสฺส ปฏิปฺปสฺสทฺธิ, โส ตุณฺหสฺส, ยสฺส นกฺขมติ, โส ภาเสยฺย.}

\prose{ทุติยมฺปิ เอตมตฺถํ วทามิ ฯเปฯ.\csromanspacer{} ตติยมฺปิ เอตมตฺถํ วทามิ.\csromanspacer{} สุณาตุ เม ภนฺเต สํโฆ, อยํ อิตฺถนฺนาโม ภิกฺขุ สํเฆน ปาปิกาย ทิฏฺฐิยา อปฺปฏินิสฺสคฺเค อุกฺเขปนียกมฺมกโต สมฺมา วตฺตติ, โลมํ ปาเตติ, เนตฺถารํ วตฺตติ, ปาปิกาย ทิฏฺฐิยา อปฺปฏินิสฺสคฺเค อุกฺเขปนียสฺส กมฺมสฺส ปฏิปฺปสฺสทฺธิํ ยาจติ, สํโฆ อิตฺถนฺนามสฺส ภิกฺขุโน ปาปิกาย ทิฏฺฐิยา อปฺปฏินิสฺสคฺเค อุกฺเขปนียกมฺมํ ปฏิปฺปสฺสมฺเภติ, ยสฺสายสฺมโต ขมติ อิตฺถนฺนามสฺส ภิกฺขุโน ปาปิกาย ทิฏฺฐิยา อปฺปฏินิสฺสคฺเค อุกฺเขปนียสฺส กมฺมสฺส ปฏิปฺปสฺสทฺธิ, โส ตุณฺหสฺส, ยสฺส นกฺขมติ, โส ภาเสยฺย.}

\prose{ปฏิปฺปสฺสทฺธํ สํเฆน อิตฺถนฺนามสฺส ภิกฺขุโน ปาปิกาย ทิฏฺฐิยา อปฺปฏินิสฺสคฺเค อุกฺเขปนียกมฺมํ, ขมติ สํฆสฺส, ตสฺมา ตุณฺหี, เอวเมตํ ธารยามี”ติ.\csromanspacer{} ปาปิกาย ทิฏฺฐิยา อปฺปฏินิสฺสคฺเค อุกฺเขปนียกมฺมํ นิฏฺฐิตํ สตฺตมํ.\csromanspacer{} กมฺมกฺขนฺธโก ปฐโม.\csromanspacer{} อิมมฺหิ ขนฺธเก วตฺถู สตฺต.}

\titlehead{\textbf{ตสฺสุทฺทานํ}}
\csromanmatikaanchor{mtk.51}
\csromantocmarkref{section}{อุทฺทานคาถา}{mtk.51}
\bookmark[dest={mtk.51},level=1]{อุทฺทานคาถา}
\csromangathameasure{ปณฺฑุโลหิตกา ภิกฺขู, สยํ ภณฺฑนการกา. \\ ตาทิเส อุปสงฺกมฺม, อุสฺสหิํสุ จ ภณฺฑเน. \\ อนุปฺปนฺนาปิ ชายนฺติ, อุปฺปนฺนานิปิ วฑฺฒเร2. \\ อปฺปิจฺฉา เปสลา ภิกฺขู, อุชฺฌายนฺติ ปทสฺสโต3. \\ สทฺธมฺมฏฺฐิติโก พุทฺโธ, สยมฺภู อคฺคปุคฺคโล. \\ อาณาเปสิ ตชฺชนีย, กมฺมํ สาวตฺถิยํ ชิโน. \\ อสมฺมุขาปฺปฏิปุจฺฉา, ปฺปฏิญฺญาย กตญฺจ ยํ. \\ อนาปตฺติ อเทสเน, เทสิตาย กตญฺจ ยํ. \\ อโจเทตฺวา อสาเรตฺวา, อนาโรเปตฺวา จ ยํ กตํ. \\ อสมฺมุขา อธมฺเมน, วคฺเคน จาปิ ยํ กตํ. \\ อปฺปฏิปุจฺฉา อธมฺเมน, ปุน วคฺเคน4 ยํ กตํ. \\ อปฺปฏิญฺญาย อธมฺเมน, วคฺเคน จาปิ4 ยํ กตํ. \\ อนาปตฺติ อธมฺเมน, วคฺเคน จาปิ4 ยํ กตํ. \\ อเทสนาคามินิยา, อธมฺเมวคฺคเมว จ. \\ เทสิตาย อธมฺเมน, วคฺเคนาปิ ตเถว จ. \\ อโจเทตฺวา อธมฺเมน, วคฺเคนาปิ ตเถว จ. \\ อสาเรตฺวา อธมฺเมน, วคฺเคนาปิ ตเถว จ. \\ อนาโรเปตฺวา อธมฺเมน, วคฺเคนาปิ ตเถว จ. \\ กณฺหวารนเยเนว, สุกฺกวารํ วิชานิยา. \\ สํโฆ อากงฺขมาโน จ, ยสฺส ตชฺชนียํ กเร. \\ ภณฺฑนํ พาโล สํสฏฺโฐ, อธิสีเล อชฺฌาจาเร. \\ อติทิฏฺฐิวิปนฺนสฺส, สํโฆ ตชฺชนิยํ กเร. \\ พุทฺธธมฺมสฺส สํฆสฺส, อวณฺณํ โย จ ภาสติ. \\ ติณฺณนฺนมฺปิ จ ภิกฺขูนํ, สํโฆ ตชฺชนิยํ กเร. \\ ภณฺฑนํ การโก เอโก, พาโล สํสคฺคนิสฺสิโต. \\ อธิสีเล อชฺฌาจาเร, ตเถว อติทิฏฺฐิยา. \\ พุทฺธธมฺมสฺส สํฆสฺส, อวณฺณํ โย จ ภาสติ. \\ ตชฺชนียกมฺมกโต, เอวํ สมฺมานุวตฺตนา. \\ อุปสมฺปทนิสฺสยา, สามเณรํ อุปฏฺฐนา. \\ โอวาทสมฺมเตนาปิ, น กเร ตชฺชนีกโต. \\ นาปชฺเช ตญฺจ อาปตฺติํ, ตาทิสญฺจ ตโต ปรํ. \\ กมฺมญฺจ กมฺมิเก จาปิ, น ครเห ตถาวิโธ. \\ อุโปสถํ ปวารณํ, ปกตตฺตสฺส นฏฺฐเป. \\ สวจนิํ อนุวาโท, โอกาโส โจทเนน จ. \\ สารณํ สมฺปโยคญฺจ, น กเรยฺย ตถาวิโธ. \\ อุปสมฺปทนิสฺสยา, สามเณรํ อุปฏฺฐนา. \\ โอวาทสมฺมเตนาปิ, ปญฺจหงฺเคหิ น สมฺมติ. \\ ตญฺจาปชฺชติ อาปตฺติํ, ตาทิสญฺจ ตโต ปรํ. \\ กมฺมญฺจ กมฺมิเก จาปิ, ครหนฺโต น สมฺมติ. \\ อุโปสถํ ปวารณํ, สวจนียา จ โนวาโท. \\ โอกาโส โจทนญฺเจว, สารณา สมฺปโยชนา. \\ อิเมหฏฺฐงฺเคหิ โย ยุตฺโต, ตชฺชนานุปสมฺมติ. \\ กณฺหวารนเยเนว, สุกฺกวารํ วิชานิยา. \\ พาโล อาปตฺติพหุโล, สํสฏฺโฐปิ จ เสยฺยโส. \\ นิยสฺสกมฺมํ สมฺพุทฺโธ, อาณาเปสิ มหามุนิ. \\ กีฏาคิริสฺมิํ เทฺว ภิกฺขู, อสฺสชิปุนพฺพสุกา. \\ อนาจารญฺจ วิวิธํ, อาจริํสุ อสญฺญตา. \\ ปพฺพาชนียํ สมฺพุทฺโธ, กมฺมํ สาวตฺถิยํ ชิโน. \\ มจฺฉิกาสณฺเฑ สุธมฺโม, จิตฺตสฺสาวาสิโก อหุ. \\ ชาติวาเทน ขุํเสติ, สุธมฺโม จิตฺตุปาสกํ. \\ ปฏิสารณียกมฺมํ, อาณาเปสิ ตถาคโต. \\ โกสมฺพิยํ ฉนฺนํ ภิกฺขุํ, นิจฺฉนฺตาปตฺติํ ปสฺสิตุํ. \\ อทสฺสเน อุกฺขิปิตุํ, อาณาเปสิ ชินุตฺตโม. \\ ฉนฺโน ตํเยว อาปตฺติํ, ปฏิกาตุํ น อิจฺฉติ. \\ อุกฺเขปนาปฺปฏิกมฺเม, อาณาเปสิ วินายโก. \\ ปาปทิฏฺฐิ อริฏฺฐสฺส, อาสิ อญฺญาณนิสฺสิตา. \\ ทิฏฺฐิยาปฺปฏินิสฺสคฺเค, อุกฺเขปํ ชินภาสิตํ. \\ นิยสฺสกมฺมํ ปพฺพชฺชํ, ตเถว ปฏิสารณี. \\ อทสฺสนาปฺปฏิกมฺเม, อนิสฺสคฺเค จ ทิฏฺฐิยา. \\ ทวานาจารูปฆาติ, มิจฺฉา-อาชีวเมว จ. \\ ปพฺพาชนียกมฺมมฺหิ, อติเรกปทา อิเม. \\ อลาภาวณฺณา เทฺว ปญฺจ, เทฺว ปญฺจกาติ นามกา. \\ ปฏิสารณียกมฺมมฺหิ, อติเรกปทา อิเม. \\ ตชฺชนียํ นิยสฺสญฺจ, ทุเว กมฺมาปิ สาทิสา4. \\ ปพฺพชฺชา ปฏิสารี จ, อตฺถิ ปทาติริตฺตตา. \\ ตโย อุกฺเขปนา กมฺมา, สทิสา เต วิภตฺติโต. \\ ตชฺชนียนเยนาปิ, เสสกมฺมํ วิชานิยาติ.}
\csromangathagroup{\csromangathastanza{\csromanfolio{79}ปณฺฑุโลหิตกา ภิกฺขู, สยํ ภณฺฑนการกา. \\ ตาทิเส อุปสงฺกมฺม, อุสฺสหิํสุ จ ภณฺฑเน.}\csromangathastanza{อนุปฺปนฺนาปิ ชายนฺติ\footnote{อนุปฺปนฺนานิ ชายนฺติ (สี, สฺยา) 2. อุปฺปนฺนานิ ปวฑฺฒเร (สี), อุปฺปนฺนาปิ ปวฑฺฒนฺติ (ก) 3. ปรีสโต (สฺยา), ปรสฺสโต (สี)}, อุปฺปนฺนานิปิ วฑฺฒเร2. \\ อปฺปิจฺฉา เปสลา ภิกฺขู, อุชฺฌายนฺติ ปทสฺสโต3.}\csromangathastanza{สทฺธมฺมฏฺฐิติโก พุทฺโธ, สยมฺภู อคฺคปุคฺคโล. \\ อาณาเปสิ ตชฺชนีย, กมฺมํ สาวตฺถิยํ ชิโน.}\csromangathastanza{อสมฺมุขาปฺปฏิปุจฺฉา, ปฺปฏิญฺญาย กตญฺจ ยํ. \\ อนาปตฺติ อเทสเน, เทสิตาย กตญฺจ ยํ.}\csromangathastanza{อโจเทตฺวา อสาเรตฺวา, อนาโรเปตฺวา จ ยํ กตํ. \\ อสมฺมุขา อธมฺเมน, วคฺเคน จาปิ\footnote{วคฺเคนาปิ จ (สี, สฺยา)} ยํ กตํ.}\csromangathastanza{อปฺปฏิปุจฺฉา อธมฺเมน, ปุน วคฺเคน4 ยํ กตํ. \\ อปฺปฏิญฺญาย อธมฺเมน, วคฺเคน จาปิ4 ยํ กตํ.}\csromangathastanza{อนาปตฺติ\footnote{อนาปตฺติยา (สี, สฺยา)} อธมฺเมน, วคฺเคน จาปิ4 ยํ กตํ. \\ อเทสนาคามินิยา, อธมฺเมวคฺคเมว จ.}\csromangathastanza{เทสิตาย อธมฺเมน, วคฺเคนาปิ ตเถว จ. \\ อโจเทตฺวา อธมฺเมน, วคฺเคนาปิ ตเถว จ.}\csromangathastanza{อสาเรตฺวา อธมฺเมน, วคฺเคนาปิ ตเถว จ. \\ อนาโรเปตฺวา อธมฺเมน, วคฺเคนาปิ ตเถว จ.}\csromangathastanza{กณฺหวารนเยเนว, สุกฺกวารํ วิชานิยา. \\ สํโฆ อากงฺขมาโน จ, ยสฺส ตชฺชนียํ กเร.}\csromangathastanza{ภณฺฑนํ พาโล สํสฏฺโฐ, อธิสีเล อชฺฌาจาเร. \\ อติทิฏฺฐิวิปนฺนสฺส, สํโฆ ตชฺชนิยํ กเร.}\csromangathastanza{\csromanfolio{80}พุทฺธธมฺมสฺส สํฆสฺส, อวณฺณํ โย จ ภาสติ. \\ ติณฺณนฺนมฺปิ จ ภิกฺขูนํ, สํโฆ ตชฺชนิยํ กเร.}\csromangathastanza{ภณฺฑนํ การโก เอโก, พาโล สํสคฺคนิสฺสิโต. \\ อธิสีเล อชฺฌาจาเร, ตเถว อติทิฏฺฐิยา.}\csromangathastanza{พุทฺธธมฺมสฺส สํฆสฺส, อวณฺณํ โย จ ภาสติ. \\ ตชฺชนียกมฺมกโต, เอวํ สมฺมานุวตฺตนา.}\csromangathastanza{อุปสมฺปทนิสฺสยา, สามเณรํ อุปฏฺฐนา. \\ โอวาทสมฺมเตนาปิ, น กเร ตชฺชนีกโต.}\csromangathastanza{นาปชฺเช ตญฺจ อาปตฺติํ, ตาทิสญฺจ ตโต ปรํ. \\ กมฺมญฺจ กมฺมิเก จาปิ, น ครเห ตถาวิโธ.}\csromangathastanza{อุโปสถํ ปวารณํ, ปกตตฺตสฺส นฏฺฐเป. \\ สวจนิํ\footnote{น สวจนิยํ (สี, สฺยา)} อนุวาโท, โอกาโส โจทเนน จ.}\csromangathastanza{สารณํ สมฺปโยคญฺจ, น กเรยฺย ตถาวิโธ. \\ อุปสมฺปทนิสฺสยา, สามเณรํ อุปฏฺฐนา.}\csromangathastanza{โอวาทสมฺมเตนาปิ, ปญฺจหงฺเคหิ\footnote{ปญฺจ-องฺโค (ก)} น สมฺมติ. \\ ตญฺจาปชฺชติ อาปตฺติํ, ตาทิสญฺจ ตโต ปรํ.}\csromangathastanza{กมฺมญฺจ กมฺมิเก จาปิ, ครหนฺโต น สมฺมติ. \\ อุโปสถํ ปวารณํ, สวจนียา จ โนวาโท.}\csromangathastanza{โอกาโส โจทนญฺเจว, สารณา สมฺปโยชนา. \\ อิเมหฏฺฐงฺเคหิ โย ยุตฺโต, ตชฺชนานุปสมฺมติ.}\csromangathastanza{กณฺหวารนเยเนว, สุกฺกวารํ วิชานิยา. \\ พาโล อาปตฺติพหุโล, สํสฏฺโฐปิ จ เสยฺยโส.}\csromangathastanza{นิยสฺสกมฺมํ สมฺพุทฺโธ, อาณาเปสิ มหามุนิ. \\ กีฏาคิริสฺมิํ เทฺว ภิกฺขู, อสฺสชิปุนพฺพสุกา.}\csromangathastanza{\csromanfolio{81}อนาจารญฺจ วิวิธํ, อาจริํสุ อสญฺญตา. \\ ปพฺพาชนียํ สมฺพุทฺโธ, กมฺมํ สาวตฺถิยํ ชิโน.}\csromangathastanza{มจฺฉิกาสณฺเฑ สุธมฺโม, จิตฺตสฺสาวาสิโก อหุ. \\ ชาติวาเทน ขุํเสติ, สุธมฺโม จิตฺตุปาสกํ.}\csromangathastanza{ปฏิสารณียกมฺมํ, อาณาเปสิ ตถาคโต. \\ โกสมฺพิยํ ฉนฺนํ ภิกฺขุํ, นิจฺฉนฺตาปตฺติํ ปสฺสิตุํ.}\csromangathastanza{อทสฺสเน อุกฺขิปิตุํ, อาณาเปสิ ชินุตฺตโม. \\ ฉนฺโน ตํเยว อาปตฺติํ, ปฏิกาตุํ น อิจฺฉติ.}\csromangathastanza{อุกฺเขปนาปฺปฏิกมฺเม, อาณาเปสิ วินายโก. \\ ปาปทิฏฺฐิ อริฏฺฐสฺส, อาสิ อญฺญาณนิสฺสิตา.}\csromangathastanza{ทิฏฺฐิยาปฺปฏินิสฺสคฺเค\footnote{ทิฏฺฐิ-อปฺปฏินิสฺสคฺเค (ก)}, อุกฺเขปํ ชินภาสิตํ. \\ นิยสฺสกมฺมํ ปพฺพชฺชํ\footnote{ทิฏฺฐิ-อปฺปฏินิสฺสคฺเค (ก)}, ตเถว ปฏิสารณี.}\csromangathastanza{อทสฺสนาปฺปฏิกมฺเม, อนิสฺสคฺเค จ ทิฏฺฐิยา. \\ ทวานาจารูปฆาติ, มิจฺฉา-อาชีวเมว จ.}\csromangathastanza{ปพฺพาชนียกมฺมมฺหิ, อติเรกปทา อิเม. \\ อลาภาวณฺณา เทฺว ปญฺจ, เทฺว ปญฺจกาติ นามกา\footnote{เทฺว ปญฺจโกติ นามโก (ก) 4. กมฺเมสุ สทิสํ (ก)}.}\csromangathastanza{ปฏิสารณียกมฺมมฺหิ, อติเรกปทา อิเม. \\ ตชฺชนียํ นิยสฺสญฺจ, ทุเว กมฺมาปิ สาทิสา4.}\csromangathastanza{ปพฺพชฺชา\footnote{ปพฺพาชา (ก)} ปฏิสารี จ, อตฺถิ ปทาติริตฺตตา. \\ ตโย อุกฺเขปนา กมฺมา, สทิสา เต วิภตฺติโต. \\ ตชฺชนียนเยนาปิ, เสสกมฺมํ วิชานิยาติ.}}

\nitthitam{\textbf{กมฺมกฺขนฺธโก นิฏฺฐิโต.}}
\csromanmatikaanchor{mtk.52}
\csromantocmarknopage{chapter}{2. ปาริวาสิกกฺขนฺธก}{mtk.52}
\bookmark[dest={mtk.52},level=0]{2. ปาริวาสิกกฺขนฺธก}
\chapterheadpageread{2. ปาริวาสิกกฺขนฺธก}
\csromanmatikaanchor{mtk.53}
\csromantocmarkref{section}{1. ปาริวาสิกวตฺต}{mtk.53}
\bookmark[dest={mtk.53},level=1]{1. ปาริวาสิกวตฺต}
\csromanheader{1}{1. ปาริวาสิกวตฺต}
\proseitem{75}{\csromanfolio{82}เตน สมเยน พุทฺโธ ภควา สาวตฺถิยํ วิหรติ เชตวเน อนาถปิณฺฑิกสฺส อาราเม.\csromanspacer{} เตน โข ปน สมเยน ปาริวาสิกา ภิกฺขู สาทิยนฺติ ปกตตฺตานํ ภิกฺขูนํ ภิวาทนํ ปจฺจุฏฺฐานํ อญฺชลิกมฺมํ สามีจิกมฺมํ อาสนาภิหารํ เสยฺยาภิหารํ ปาโททกํ ปาทปีฐํ ปาทกถลิกํ ปตฺตจีวรปฏิคฺคหณํ นหาเน ปฏฺฐิปริกมฺมํ.\csromanspacer{} เย เต ภิกฺขู อปฺปิจฺฉา ฯเปฯ เต อุชฺฌายนฺติ ขิยฺยนฺติ วิปาเจนฺติ “กถํ หิ นาม ปาริวาสิกา ภิกฺขู สาทิยิสฺสนฺติ ปกตตฺตานํ ภิกฺขูนํ อภิวาทนํ ปจฺจุฏฺฐานํ อญฺชลิกมฺมํ สามีจิกมฺมํ อาสนาภิหารํ เสยฺยาภิหารํ ปาโททกํ ปาทปีฐํ ปาทกถลิกํ ปตฺตจีวรปฏิคฺคหณํ นหาเน ปิฏฺฐิปริกมฺมนฺ”ติ.\csromanspacer{} อถ โข เต ภิกฺขู ภควโต เอตมตฺถํ อาโรเจสุํ.}

\prose{อถ โข ภควา เอตสฺมิํ นิทาเน เอตสฺมิํ ปกรเณ ภิกฺขุสํฆํ สนฺนิปาตาเปตฺวา ภิกฺขู ปฏิปุจฺฉิ “สจฺจํ กิร ภิกฺขเว ปาริวาสิกา ภิกฺขู สาทิยนฺติ ปกตตฺตานํ ภิกฺขูนํ อภิวาทนํ ปจฺจุฏฺฐานํ อญฺชลิกมฺมํ สามีจิกมฺมํ อาสนาภิหารํ เสยฺยาภิหารํ ปาโททกํ ปาทปีฐํ ปาทกถลิกํ ปตฺตจีวรปฏิคฺคหณํ นหาเน ปิฏฺฐิปริกมฺมนฺ”ติ.\csromanspacer{} สจฺจํ ภควาติ.\csromanspacer{} วิครหิ พุทฺโธ ภควา อนนุจฺฉวิกํ ฯเปฯ.\csromanspacer{} กถํ หิ นาม ภิกฺขเว ปาริวาสิกา ภิกฺขู สาทิยิสฺสนฺติ ปกตตฺตานํ ภิกฺขูนํ อภิวาทนํ ปจฺจุฏฺฐานํ อญฺชลิกมฺมํ สามีจิกมฺมํ อาสนาภิหารํ เสยฺยาภิหารํ ปาโททกํ ปาทปีฐํ ปาทกถลิกํ ปตฺตจีวรปฏิคฺคหณํ นหาเน ปิฏฺฐิปริกมฺมํ.\csromanspacer{} เนตํ ภิกฺขเว อปฺปสนฺนานํ วา ปสาทาย ฯเปฯ วิครหิตฺวา ฯเปฯ ธมฺมิํ กถํ กตฺวา ภิกฺขู อามนฺเตสิ\csromandash{}น ภิกฺขเว ปาริวาสิเกน ภิกฺขุนา สาทิตพฺพํ ปกตตฺตานํ ภิกฺขูนํ อภิวาทนํ ปจฺจุฏฺฐานํ อญฺชลิกมฺมํ สามีจิกมฺมํ อาสนาภิหาโร เสยฺยาภิหาโร ปาโททกํ ปาทปีฐํ ปาทกถลิกํ ปตฺตจีวรปฏิคฺคหณํ นหาเน ปิฏฺฐิปริกมฺมํ, โย สาทิเยยฺย, อาปตฺติ ทุกฺกฏสฺส.}

\prose{อนชานามิ ภิกฺขเว ปาริวาสิกานํ ภิกฺขูนํ มิถุ ยถาวุฑฺฒํ อภิวาทนํ ปจฺจุฏฺฐานํ อญฺชลิกมฺมํ สามีจิกมฺมํ อาสนาภิหารํ เสยฺยาภิหารํ ปาโททกํ ปาทปีฐํ ปาทกถลิกํ ปตฺตจีวรปฏิคฺคหณํ นหาเน ปิฏฺฐิปริกมฺมํ.}

\prose{\csromanfolio{83}อนุชานามิ ภิกฺขเว ปาริวาสิกานํ ภิกฺขูนํ ปญฺจ ยถาวุฑฺฒํ อุโปสถํ ปวารณํ วสฺสิกสาฏิกํ โอโณชนํ ภตฺตํ.\csromanspacer{} เตน หิ ภิกฺขเว ปาริวาสิกานํ ภิกฺขูนํ วตฺตํ ปญฺญเปสฺสามิ, ยถา ปาริวาสิเกหิ ภิกฺขูหิ วตฺติตพฺพํ.}

\proseitem{76}{ปาริวาสิเกน ภิกฺขเว ภิกฺขุนา สมฺมา วตฺติตพฺพํ.\csromanspacer{} ตตฺรายํ สมฺมาวตฺตนา\csromandash{}}

\prose{น อุปสมฺปาเทตพฺพํ, น นิสฺสโย ทาตพฺโพ, น สามเณโร อุปฏฺฐาเปตพฺโพ, น ภิกฺขุโนวาทกสมฺมุติ สาทิตพฺพา, สมฺมเตนปิ ภิกฺขุนิโย น โอวทิตพฺพา, ยาย อาปตฺติยา สํเฆน ปริวาโส ทินฺโน โหติ, สา อาปตฺติ น อาปชฺชิตพฺพา, อญฺญา วา ตาทิสิกา, ตโต วา ปาปิฏฺฐตรา, กมฺมํ น ครหิตพฺพํ, กมฺมิกา น ครหิตพฺพา, น ปกตตฺตสฺส ภิกฺขุโน อุโปสโถ ฐเปตพฺโพ, น ปวารณา ฐเปตพฺพา, น สวจนียํ กาตพฺพํ, น อนุวาโท ปฏฺฐเปตพฺโพ, น โอกาโส กาเรตพฺโพ, น โจเทตพฺโพ, น สาเรตพฺโพ, น ภิกฺขูหิ สมฺปโยเชตพฺพํ.}

\prose{น ภิกฺขเว ปาริวาสิเกน ภิกฺขุนา ปกตตฺตสฺส ภิกฺขุโน ปุรโต คนฺตพฺพํ, น ปุรโต นิสีทิตพฺพํ, โย โหติ สํฆสฺส อาสนปริยนฺโต เสยฺยาปริยนฺโต วิหารปริยนฺโต, โส ตสฺส ปทาตพฺโพ, เตน จ โส สาทิตพฺโพ.}

\prose{น ภิกฺขเว ปาริวาสิเกน ภิกฺขุนา ปกตตฺเตน ภิกฺขุนา ปุเรสมเณน วา ปจฺฉาสมเณน วา กุลานิ อุปสงฺกมิตพฺพานิ, น อารญฺญิกงฺคํ สมาทาตพฺพํ\footnote{สมาทิตพฺพํ (ก)}, น ปิณฺฑปาติกงฺคํ สมาทาตพฺพํ, น จ ตปฺปจฺจยา ปิณฺฑปาโต นีหราเปตพฺโพ “มา มํ ชานิํสู”ติ.}

\prose{ปาริวาสิเกน ภิกฺขเว ภิกฺขุนา อาคนฺตุเกน อาโรเจตพฺพํ, อาคนฺตุกสฺส อาโรเจตพฺพํ, อุโปสเถ อาโรเจตพฺพํ, ปวารณาย อาโรเจตพฺพํ, สเจ คิลาโน โหติ, ทูเตนปิ อาโรเจตพฺพํ.}

\prose{น ภิกฺขเว ปาริวาสิเกน ภิกฺขุนา สภิกฺขุกา อาวาสา อภิกฺขุโก อาวาโส คนฺตพฺโพ อญฺญตฺร ปกตตฺเตน อญฺญตฺร อนฺตรายา.}

\prose{\csromanfolio{84}น ภิกฺขเว ปาริวาสิเกน ภิกฺขุนา สภิกฺขุกา อาวาสา อภิกฺขุโก อนาวาโส คนฺตพฺโพ อญฺญตฺร ปกตตฺเตน อญฺญตฺร อนฺตรายา.\csromanspacer{} น ภิกฺขเว ปาริวาสิเกน ภิกฺขุนา สภิกฺขุกา อาวาสา อภิกฺขุโก อาวาโส วา อนาวาโส วา คนฺตพฺโพ อญฺญตฺร ปกตตฺเตน อญฺญตฺร อนฺตรายา.}

\prose{น ภิกฺขเว ปาริวาสิเกน ภิกฺขุนา สภิกฺขุกา อนาวาสา อภิกฺขุโก อาวาโส คนฺตพฺโพ อญฺญตฺร ปกตตฺเตน อญฺญตฺร อนฺตรายา.\csromanspacer{} น ภิกฺขเว ปาริวาสิเกน ภิกฺขุนา สภิกฺขุกา อนาวาสา อภิกฺขุโก อนาวาโส คนฺตพฺโพ อญฺญตฺร ปกตตฺเตน อญฺญตฺร อนฺตรายา.\csromanspacer{} น ภิกฺขเว ปาริวาสิเกน ภิกฺขุนา สภิกฺขุกา อนาวาสา อภิกฺขุโก อาวาโส วา อนาวาโส วา คนฺตพฺโพ อญฺญตฺร ปกตตฺเตน อญฺญตฺร อนฺตรายา.}

\prose{น ภิกฺขเว ปาริวาสิเกน ภิกฺขุนา สภิกฺขุกา อาวาสา วา อนาวาสา วา อภิกฺขุโก อาวาโส คนฺตพฺโพ อญฺญตฺร ปกตตฺเตน อญฺญตฺร อนฺตรายา.\csromanspacer{} น ภิกฺขเว ปาริวาสิเกน ภิกฺขุนา สภิกฺขุกา อาวาสา วา อนาวาสา วา อภิกฺขุโก อนาวาโส คนฺตพฺโพ อญฺญตฺร ปกตตฺเตน อญฺญตฺร อนฺตรายา.\csromanspacer{} น ภิกฺขเว ปาริวาสิเกน ภิกฺขุนา สภิกฺขุกา อาวาสา วา อนาวาสา วา อภิกฺขุโก อาวาโส วา อนาวาโส วา คนฺตพฺโพ อญฺญตฺร ปกตตฺเตน อญฺญตฺร อนฺตรายา.}

\prose{น ภิกฺขเว ปาริวาสิเกน ภิกฺขุนา สภิกฺขุกา อาวาสา สภิกฺขุโก อาวาโส คนฺตพฺโพ ยตฺถสฺสุ ภิกฺขู นานาสํวาสกา อญฺญตฺร ปกตตฺเตน อญฺญตฺร อนฺตรายา.\csromanspacer{} น ภิกฺขเว ปาริวาสิเกน ภิกฺขุนา สภิกฺขุกา อาวาสา สภิกฺขุโก อนาวาโส คนฺตพฺโพ ยตฺถสฺสุ ภิกฺขู นานาสํวาสกา อญฺญตฺร ปกตตฺเตน อญฺญตฺร อนฺตรายา.\csromanspacer{} น ภิกฺขเว ปาริวาสิเกน ภิกฺขุนา สภิกฺขุกา อาวาสา สภิกฺขุโก อาวาโส วา อนาวาโส วา คนฺตพฺโพ ยตฺถสฺสุ ภิกฺขู นานาสํวาสกา อญฺญตฺร ปกตตฺเตน อญฺญตฺร อนฺตรายา.}

\prose{น ภิกฺขเว ปาริวาสิเกน ภิกฺขุนา สภิกฺขุกา อนาวาสา สภิกฺขุโก อาวาโส คนฺตพฺโพ ยตฺถสฺสุ ภิกฺขู นานาสํวาสกา \csromanfolio{85}อญฺญตฺร ปกตตฺเตน อญฺญตฺร อนฺตรายา.\csromanspacer{} น ภิกฺขเว ปาริวาสิเกน ภิกฺขุนา สภิกฺขุกา อนาวาสา สภิกฺขุโก อนาวาโส คนฺตพฺโพ ยตฺถสฺสุ ภิกฺขู นานาสํวาสกา อญฺญตฺร ปกตตฺเตน อญฺญตฺร อนฺตรายา.\csromanspacer{} น ภิกฺขเว ปาริวาสิเกน ภิกฺขุนา สภิกฺขุกา อนาวาสา สภิกฺขุโก อาวาโส วา อนาวาโส วา คนฺตพฺโพ ยตฺถสฺสุ ภิกฺขู นานาสํวาสกา อญฺญตฺร ปกตตฺเตน อญฺญตฺร อนฺตรายา.}

\prose{น ภิกฺขเว ปาริวาสิเกน ภิกฺขุนา สภิกฺขุกา อาวาสา วา อนาวาสา วา สภิกฺขุโก อาวาโส คนฺตพฺโพ ยตฺถสฺสุ ภิกฺขู นานาสํวาสกา อญฺญตฺร ปกตตฺเตน อญฺญตฺร อนฺตรายา.\csromanspacer{} น ภิกฺขเว ปาริวาสิเกน ภิกฺขุนา สภิกฺขุกา อาวาสา วา อนาวาสา วา สภิกฺขุโก อนาวาโส คนฺตพฺโพ ยตฺถสฺสุ ภิกฺขู นานาสํวาสกา อญฺญตฺร ปกตตฺเตน อญฺญตฺร อนฺตรายา.\csromanspacer{} น ภิกฺขเว ปาริวาสิเกน ภิกฺขุนา สภิกฺขุกา อาวาสา วา อนาวาสา วา สภิกฺขุโก อาวาโส วา อนาวาโส วา คนฺตพฺโพ ยตฺถสฺสุ ภิกฺขู นานาสํวาสกา อญฺญตฺร ปกตตฺเตน อญฺญตฺร อนฺตรายา.}

\proseitem{80}{คนฺตพฺโพ ภิกฺขเว ปาริวาสิเกน ภิกฺขุนา สภิกฺขุกา อาวาสา สภิกฺขุโก อาวาโส ยตฺถสฺสุ ภิกฺขู สมานสํวาสกา ยํ ชญฺญา “สกฺโกมิ อชฺเชว คนฺตุนฺ”ติ.}

\prose{คนฺตพฺโพ ภิกฺขเว ปาริวาสิเกน ภิกฺขุนา สภิกฺขุกา อาวาสา สภิกฺขุโก อนาวาโส ยตฺถสฺสุ ภิกฺขู สมานสํวาสกา ยํ ชญฺญา “สกฺโกมิ อชฺเชว คนฺตุนฺ”ติ.}

\prose{คนฺตพฺโพ ภิกฺขเว ปาริวาสิเกน ภิกฺขุนา สภิกฺขุกา อาวาสา สภิกฺขุโก อาวาโส วา อนาวาโส วา ยตฺถสฺสุ ภิกฺขู สมานสํวาสกา ยํ ชญฺญา “สกฺโกมิ อชฺเชว คนฺตุนฺ”ติ.}

\prose{คนฺตพฺโพ ภิกฺขเว ปาริวาสิเกน ภิกฺขุนา สภิกฺขุกา อนาวาสา สภิกฺขุโก อาวาโส ยตฺถสฺสุ ภิกฺขู สมานสํวาสกา ยํ ชญฺญา “สกฺโกมิ อชฺเชว คนฺตุนฺ”ติ.}

\prose{\csromanfolio{86}คนฺตพฺโพ ภิกฺขเว ปาริวาสิเกน ภิกฺขุนา สภิกฺขุกา อนาวาสา สภิกฺขุโก อนาวาโส ยตฺถสฺสุ ภิกฺขู สมานสํวาสกา ยํ ชญฺญา “สกฺโกมิ อชฺเชว คนฺตุนฺ”ติ.}

\prose{คนฺตพฺโพ ภิกฺขเว ปาริวาสิเกน ภิกฺขุนา สภิกฺขุกา อนาวาสา สภิกฺขุโก อาวาโส วา อนาวาโส วา ยตฺถสฺสุ ภิกฺขู สมานสํวาสกา ยํ ชญฺญา “สกฺโกมิ อชฺเชว คนฺตุนฺ”ติ.}

\prose{คนฺตพฺโพ ภิกฺขเว ปาริวาสิเกน ภิกฺขุนา สภิกฺขุกา อาวาสา วา อนาวาสา วา สภิกฺขุโก อาวาโส ยตฺถสฺสุ ภิกฺขู สมานสํวาสกา ยํ ชญฺญา “สกฺโกมิ อชฺเชว คนฺตุนฺ”ติ.}

\prose{คนฺตพฺโพ ภิกฺขเว ปาริวาสิเกน ภิกฺขุนา สภิกฺขุกา อาวาสา วา อนาวาสา วา สภิกฺขุโก อนาวาโส ยตฺถสฺสุ ภิกฺขู สมานสํวาสกา ยํ ชญฺญา “สกฺโกมิ อชฺเชว คนฺตุนฺ”ติ.}

\prose{คนฺตพฺโพ ภิกฺขเว ปาริวาสิเกน ภิกฺขุนา สภิกฺขุกา อาวาสา วา อนาวาสา วา สภิกฺขุโก อาวาโส วา อนาวาโส วา ยตฺถสฺสุ ภิกฺขู สมานสํวาสกา ยํ ชญฺญา “สกฺโกมิ อชฺเชว คนฺตุนฺ”ติ.}

\proseitem{81}{น ภิกฺขเว ปาริวาสิเกน ภิกฺขุนา ปกตตฺเตน ภิกฺขุนา สทฺธิํ เอกจฺฉนฺเน อาวาเส วตฺตพฺพํ, น เอกจฺฉนฺเน อนาวาเส วตฺถพฺพํ, น เอกจฺฉนฺเน อาวาเส วา อนาวาเส วา วตฺถพฺพํ, ปกตตฺตํ ภิกฺขุํ ทิสฺวา อาสนา วุฏฺฐาตพฺพํ, ปกตตฺโต ภิกฺขุ อาสเนน นิมนฺเต ตพฺโพ, น ปกตตฺเตน ภิกฺขุนา สทฺธิํ เอกาสเน นิสีทิตพฺพํ, น นีเจ อาสเน นิสินฺเน อุจฺเจ อาสเน นิสีทิตพฺพํ, น ฉมายํ นิสินฺเน อาสมน นิสีทิตพฺพํ, น เอกจงฺกเม จงฺกมิตพฺพํ, น นีเจ จงฺกเม จงฺกมนฺเต อุจฺเจ จงฺกเม จงฺกมิตพฺพํ, น ฉมายํ จงฺกมนฺเต\footnote{จงฺกมนฺตํ (อฏฺฐกถายํ สํวณฺเณตพฺพปาโฐ)} จงฺเกเม จงฺกมิตพฺพํ.}

\proseitem{82}{น ภิกฺขเว ปาริวาสิเกน ภิกฺขุนา ปาริวาสิเกน วุฑฺฒตเรน ภิกฺขุนา สทฺธิํ ฯเปฯ มูลายปฏิกสฺสนารเหน ภิกฺขุนา สทฺธิํ ฯเปฯ มานตฺตารเหน ภิกฺขุนา สทฺธิํ ฯเปฯ มานตฺตจาริเกน ภิกฺขุนา สทฺธิํ ฯเปฯ อพฺภานารเหน ภิกฺขุนา สทฺธิํ เอกจฺฉนฺเน อาวาเส วตฺถพฺพํ, น เอกจฺฉนฺเน อนาวาเส วตฺถพฺพํ, น เอกจฺฉนฺเน อาวาเส วา อนาวาเส วา วตฺถพฺพํ, น เอกาสเน \csromanfolio{87}นิสีทิตพฺพํ, น นีเจ อาสเน นิสินฺเน อุจฺเจ อาสเน นิสีทิตพฺพํ, น ฉมายํ นิสินฺเน อาสเน นิสีทิตพฺพํ, น เอกจงฺกเม จงฺกมิตพฺพํ, น นีเจ จงฺกเม จงฺกมนฺเต อุจฺเจ จงฺกเม จงฺกมิตพฺพํ, น ฉมายํ จงฺกมนฺเต จงฺกเม จงฺกมิตพฺพํ.}

\prose{\csromansymbolfootnote{*}{วิ 3. 442 ปิฏฺเฐ.}ปาริวาสิกจตุตฺโถ เจ ภิกฺขเว ปริวาสํ ทเทยฺย, มูลาย ปฏิกสฺเสยฺย, มานตฺตํ ทเทยฺย, ตํวีโส อพฺเภยฺย, อกมฺมํ\footnote{อกมฺมํ ตํ (สฺยา)} น จ กรณียนฺติ.}

\nitthitam{จตุนฺนวุติปาริวาสิกวตฺตํ นิฏฺฐิตํ.}
\proseitem{83}{อถ โข อายสฺมา อุปาลิ เยน ภควา เตนุปสงฺกมิ, อุปสงฺกมิตฺวา ภควนฺตํ อภิวาเทตฺวา เอกมนฺตํ นิสีทิ, เอกมนฺตํ นิสินฺโน โข อายสฺมา อุปาลิ ภควนฺตํ เอตทโวจ “กติ นุ โข ภนฺเต ปาริวาสิกสฺส ภิกฺขุโน รตฺติจฺเฉทา”ติ.\csromanspacer{} ตโย โข อุปาลิ ปาริวาสิกสฺส ภิกฺขุโน รตฺติจฺเฉทา, สหวาโส วิปฺปวาโส อนาโรจนา, อิเม โข อุปาลิ ตโย ปาริวาสิกสฺส ภิกฺขุโน รตฺติจฺเฉทาติ.}

\proseitem{84}{เตน โข ปน สมเยน สาวตฺถิยํ มหาภิกฺขุสํโฆ สนฺนิปติโต โหติ, น สกฺโกนฺติ ปาริวาสิกา ภิกฺขู ปริวาสํ โสเธตุํ, ภควโต\footnote{เต ภิกฺขู ภควโต (สฺยา, เอวมุปริปิ)} เอตมตฺถํ อาโรเจสุํ.\csromanspacer{} อนุชานามิ ภิกฺขเว ปริวาสํ นิกฺขิปิตุํ.\csromanspacer{} เอวญฺจ ปน ภิกฺขเว นิกฺขิปิตพฺโพ, เตน ปาริวาสิเกน ภิกฺขุนา เอกํ ภิกฺขุํ อุปสงฺกมิตฺวา เอกํสํ อุตฺตราสงฺคํ กริตฺวา อุกฺกุฏิกํ นิสีทิตฺวา อญฺชลิํ ปคฺคเหตฺวา เอวมสฺส วจนีโย “ปริวาสํ นิกฺขิปามี”ติ, นิกฺขิตฺโต โหติ ปริวาโส. “วตฺตํ นิกฺขิปามี”ติ, นิกฺขิตฺโต โหติ ปริวาโส.}

\proseitem{85}{เตน โข ปน สมเยน สาวตฺถิยา ภิกฺขู ตหํ ตหํ ปกฺกมิํสุ, สกฺโกนฺติ ปาริวาสิกา ภิกฺขู ปริวาสํ โสเธสุํ, ภควโต เอตมตฺถํ อาโรเจสุํ.\csromanspacer{} อนุชานามิ ภิกฺขเว ปริวาสํ สมาทิยิตุํ\footnote{สมาทาตุํ (สฺยา, กํ)}.\csromanspacer{} เอวญฺจ ปน ภิกฺขเว สมาทิยิตพฺโพ\footnote{สมาทิตพฺโพ (สี, สฺยา, กํ)}, เตน ปาริวาสิเกน ภิกฺขุนา เอกํ ภิกฺขุํ \csromanfolio{88}อุปสงฺกมิตฺวา เอกํสํ อุตฺตราสงฺคํ กริตฺวา อุกฺกุฏิกํ นิสีทิตฺวา อญฺชลิํ ปคฺคเหตฺวา เอวมสฺส วจนีโย “ปริวาสํ สมาทิยามี”ติ, สมาทินฺโน โหติ ปริวาโส. “วตฺตํ สมาทิยามี”ติ, สมาทินฺโน โหติ ปริวาโส.}

\nitthitam{ปาริวาสิกวตฺตํ นิฏฺฐิตํ.}
\csromanmatikaanchor{mtk.54}
\csromantocmarkref{section}{2. มูลายปฏิกสฺสนารหวตฺต}{mtk.54}
\bookmark[dest={mtk.54},level=1]{2. มูลายปฏิกสฺสนารหวตฺต}
\csromanheader{1}{2. มูลายปฏิกสฺสนารหวตฺต}
\proseitem{86}{เตน โข ปน สมเยน มูลายปฏิกสฺสนารหา ภิกฺขู สาทิยนฺติ ปกตตฺตานํ ภิกฺขูนํ อภิวาทนํ ปจฺจุฏฺฐานํ อญฺชลิกมฺมํ สามีจิกมฺมํ อาสนาภิหารํ เสยฺยาภิหารํ ปาโททกํ ปาทปีฐํ ปาทกถลิกํ ปตฺตจีวรปฏิคฺคหณํ นหาเน ปิฏฺฐิปริกมฺมํ.\csromanspacer{} เย เต ภิกฺขู อปฺปิจฺฉา ฯเปฯ เต อุชฺฌายนฺติ ขิยฺยนฺติ วิปาเจนฺติ “กถํ หิ นาม มูลายปฏิกสฺสนารหา ภิกฺขู สาทิยิสฺสนฺติ ปกตตฺตานํ ภิกฺขูนํ อภิวาทนํ ปจฺจุฏฺฐานํ ฯเปฯ นหาเน ปิฏฺฐิปริกมฺมนฺ”ติ.\csromanspacer{} อถ โข เต ภิกฺขู ภควโต เอตมตฺถํ อาโรเจสุํ.}

\prose{อถ โข ภควา เอตสฺมิํ นิทาเน เอตสฺมิํ ปกรเณ ภิกฺขุสํฆํ สนฺนิปาตาเปตฺวา ภิกฺขู ปฏิปุจฺฉิ “สจฺจํ กิร ภิกฺขเว มูลายปฏิกสฺสนารหา ภิกฺขู สาทิยนฺติ ปกตตฺตานํ ภิกฺขูนํ อภิวาทนํ ปจฺจุฏฺฐานํ ฯเปฯ นหาเน ปิฏฺฐิปริกมฺมนฺ”ติ.\csromanspacer{} สจฺจํ ภควาติ.\csromanspacer{} วิครหิ พุทฺโธ ภควา อนนุจฺฉวิกํ ฯเปฯ.\csromanspacer{} กถํ หิ นาม ภิกฺขเว มูลายปฏิกสฺสนารหา ภิกฺขู สาทิยิสฺสนฺติ ปกตตฺตานํ ภิกฺขูนํ อภิวาทนํ ปจฺจุฏฺฐานํ ฯเปฯ นหาเน ปิฏฺฐิปริกมฺมํ.\csromanspacer{} เนตํ ภิกฺขเว อปฺปสนฺนานํ วา ปสฺสาทาย ฯเปฯ วิครหิตฺวา ฯเปฯ ธมฺมิํ กถํ กตฺวา ภิกฺขู อามนฺเตสิ\csromandash{}น ภิกฺขเว มูลายปฏิกสฺสนารเหน ภิกฺขุนา สาทิตพฺพํ ปกตตฺตานํ ภิกฺขูนํ อภิวาทนํ ปจฺจุฏฺฐานํ อญฺชลิกมฺมํ สามีจิกมฺมํ อาสนาภิหาโร เสยฺยาภิหาโร ปาโททกํ ปาทปีฐํ ปาทกถลิกํ ปตฺตจีวรปฏิคฺคหณํ นหาเน ปิฏฺฐิปริกมฺมํ, โย สาทิเยยฺย, อาปตฺติ ทุกฺกฏสฺส.\csromanspacer{} อนุชานามิ ภิกฺขเว มูลายปฏิกสฺสนารหานํ ภิกฺขูนํ มิถุ ยถาวุฑฺฒํ อภิวาทนํ ปจฺจุฏฺฐานํ ฯเปฯ นหาเน ปิฏฺฐิปริกมฺมํ.\csromanspacer{} อนุชานามิ ภิกฺขเว มูลายปฏิกสฺสนารหานํ ภิกฺขูนํ ปญฺจ ยถาวุฑฺฒํ \csromanfolio{89}อุโปสถํ ปวารณํ วสฺสิกสาฏิกํ โอโณชนํ ภตฺตํ.\csromanspacer{} เตน หิ ภิกฺขเว มูลายปฏิกสฺสนารหานํ ภิกฺขูนํ วตฺตํ ปญฺญเปสฺสามิ, ยถา มูลายปฏิกสฺสนารเหหิ ภิกฺขูหิ วตฺติตพฺพํ.}

\proseitem{87}{มูลายปฏิกสฺสนารเหน ภิกฺขเว ภิกฺขุนา สมฺมา วตฺติตพฺพํ.\csromanspacer{} ตตฺรายํ สมฺมาวตฺตนา\csromandash{}}

\prose{น อุปสมฺปาเทตพฺพํ, น นิสฺสโย ทาตพฺโพ, น สามเณโร อุปฏฺฐาเปตพฺโพ, น ภิกฺขุโนวาทกสมฺมุติ สาทิตพฺพา, สมฺมเตนปิ ภิกฺขุนิโย น โอวทิตพฺพา, ยาย อาปตฺติยา สํเฆน มูลาย ปฏิกสฺสนารโห กโต โหติ, สา อาปตฺติ น อาปชฺชิตพฺพา, อญฺญา วา ตาทิสิกา, ตโต วา ปาปิฏฺฐตรา, กมฺมํ น ครหิตพฺพํ, กมฺมิกา น ครหิตพฺพา, น ปกตตฺตสฺส ภิกฺขุโน อุโปสโถ ฐเปตพฺโพ, น ปวารณา ฐเปตพฺพา, น สวจนียํ กาตพฺพํ, น อนุวาโท ปฏฺฐเปตพฺโพ, น โอกาโส กาเรตพฺโพ, น โจเทตพฺโพ, น สาเรตพฺโพ, น ภิกฺขูหิ สมฺปโยเชตพฺพํ.}

\prose{น ภิกฺขเว มูลายปฏิกสฺสนารเหน ภิกฺขุนา ปกตตฺตสฺส ภิกฺขุโน ปุรโต คนฺตพฺพํ, น ปุรโต นิสีทิตพฺพํ, โย โหติ สํฆสฺส อาสนปริยนฺโต เสยฺยาปริยนฺโต วิหารปริยนฺโต, โส ตสฺส ปทาตพฺโพ, เตน จ โส สาทิตพฺโพ.}

\prose{น ภิกฺขเว มูลายปฏิกสฺสนารเหน ภิกฺขุนา ปกตตฺเตน ภิกฺขุนา ปุเรสมเณน วา ปจฺฉาสมเณน วา กุลานิ อุปสงฺกมิตพฺพานิ, น อารญฺญิกงฺคํ สมาทาตพฺพํ, น ปิณฺฑปาติกงฺคํ สมาทาตพฺพํ, น จ ตปฺปจฺจยา ปิณฺฑปาโต นีหราเปตพฺโพ “มา มํ ชานิํสู”ติ.}

\prose{น ภิกฺขเว มูลายปฏิกสฺสนารเหน ภิกฺขุนา สภิกฺขุกา อาวาสา อภิกฺขุโก อาวาโส คนฺตพฺโพ อญฺญตฺร ปกตตฺเตน อญฺญตฺร อนฺตรายา.}

\prose{น ภิกฺขเว มูลายปฏิกสฺสนารเหน ภิกฺขุนา สภิกฺขุกา อาวาสา อภิกฺขุโก อนาวาโส คนฺตพฺโพ อญฺญตฺร ปกตตฺเตน อญฺญตฺร อนฺตรายา.}

\prose{\csromanfolio{90}น ภิกฺขเว มูลายปฏิกสฺสนารเหน ภิกฺขุนา สภิกฺขุกา อาวาสา อภิกฺขุโก อาวาโส วา อนาวาโส วา คนฺตพฺโพ อญฺญตฺร ปกตตฺเตน อญฺญตฺร อนฺตรายา.}

\prose{น ภิกฺขเว มูลายปฏิกสฺสนารเหน ภิกฺขุนา สภิกฺขุกา อนาวาสา อภิกฺขุโก อาวาโส คนฺตพฺโพ ฯเปฯ อภิกฺขุโก อนาวาโส คนฺตพฺโพ ฯเปฯ อภิกฺขุโก อาวาโส วา อนาวาโส วา คนฺตพฺโพ อญฺญตฺร ปกตฺเตน อญฺญตฺร อนฺตรายา.}

\prose{น ภิกฺขเว มูลายปฏิกสฺสนารเหน ภิกฺขุนา สภิกฺขุกา อาวาสา วา อนาวาสา วา อภิกฺขุโก อาวาโส คนฺตพฺโพ ฯเปฯ อภิกฺขุโก อนาวาโส คนฺตพฺโพ ฯเปฯ อภิกฺขุโก อาวาโส วา อนาวาโส วา คนฺตพฺโพ อญฺญตฺร ปกตตฺเตน อญฺญตฺร อนฺตรายา.}

\prose{น ภิกฺขเว มูลายปฏิกสฺสนารเหน ภิกฺขุนา สภิกฺขุกา อาวาสา สภิกฺขุโก อาวาโส คนฺตพฺโพ ฯเปฯ สภิกฺขุโก อนาวาโส คนฺตพฺโพ ฯเปฯ สภิกฺขุโก อาวาโส วา อนาวาโส วา คนฺตพฺโพ ยตฺถสฺสุ ภิกฺขู นานาสํวาสกา อญฺญตฺร ปกตตฺเตน อญฺญตฺร อนฺตรายา.}

\prose{น ภิกฺขเว มูลายปฏิกสฺสนารเหน ภิกฺขุนา สภิกฺขุกา อานาวาสา สภิกฺขุโก อาวาโส คนฺตพฺโพ ฯเปฯ สภิกฺขุโก อนาวาโส คนฺตพฺโพ ฯเปฯ สภิกฺขุโก อาวาโส วา อนาวาโส วา คนฺตพฺโพ ยตฺถสฺสุ ภิกฺขู นานาสํวาสกา อญฺญตฺร ปกตตฺเตน อญฺญตฺร อนฺตรายา.}

\prose{น ภิกฺขเว มูลายปฏิกสฺสนารเหน ภิกฺขุนา สภิกฺขุกา อาวาสา วา อนาวาสา วา สภิกฺขุโก อาวาโส คนฺตพฺโพ ฯเปฯ สภิกฺขุโก อนาวาโส คนฺตพฺโพ ฯเปฯ สภิกฺขุโก อาวาโส วา อนาวาโส วา คนฺตพฺโพ ยตฺถสฺสุ ภิกฺขู นานาสํวาสกา อญฺญตฺร ปกตตฺเตน อญฺญตฺร อนฺตรายา.}

\prose{คนฺตพฺโพ ภิกฺขเว มูลายปฏิกสฺสนารเหน ภิกฺขุนา สภิกฺขุกา อาวาสา สภิกฺขุโก อาวาโส ฯเปฯ สภิกฺขุโก อนาวาโส ฯเปฯ สภิกฺขุโก อาวาโส วา อนาวาโส วา ยตฺถสฺสุ ภิกฺขู สมานสํวาสกา ยํ ชญฺญา “สกฺโกมิ อชฺเชว คนฺตุนฺ”ติ.}

\prose{\csromanfolio{91}คนฺตพฺโพ ภิกฺขเว มูลายปฏิกสฺสนารเหน ภิกฺขุนา สภิกฺขุกา อนาวาสา สภิกฺขุโก อาวาโส ฯเปฯ สภิกฺขุโก อนาวาโส ฯเปฯ สภิกฺขุโก อาวาโส วา อนาวาโส วา ยตฺถสฺสุ ภิกฺขู สมานสํวาสกา ยํ ชญฺญา “สกฺโกมิ อชฺเชว คนฺตุนฺ”ติ.}

\prose{คนฺตพฺโพ ภิกฺขเว มูลายปฏิกสฺสนารเหน ภิกฺขุนา สภิกฺขุกา อาวาสา วา อนาวาสา วา สภิกฺขุโก อาวาโส ฯเปฯ สภิกฺขุโก อนาวาโส ฯเปฯ สภิกฺขุโก อาวาโส วา อนาวาโส วา ยตฺถสฺสุ ภิกฺขู สมานสํวาสกา ยํ ชญฺญา “สกฺโกมิ อชฺเชว คนฺตุนฺ”ติ.}

\prose{น ภิกฺขเว มูลายปฏิกสฺสนารเหน ภิกฺขุนา ปกตตฺเตน ภิกฺขุนา สทฺธิํ เอกจฺฉนฺเน อาวาเส วตฺถพฺพํ, น เอกจฺฉนฺเน อนาวาเส วตฺถพฺพํ, น เอกจฺฉนฺเน อาวาเส วา อนาวาเส วา วตฺถพฺพํ, ปกตตฺตํ ภิกฺขุํ ทิสฺวา อาสนา วุฏฺฐาตพฺพํ, ปกตตฺโต ภิกฺขุ อาสเนน นิมนฺเตตพฺโพ, น ปกตตฺเตน ภิกฺขุนา สทฺธิํ เอกาสเน นิสีทิตพฺพํ, น นีเจ อาสเน นิสินฺเน อุจฺเจ อาสเน นิสีทิตพฺพํ, น ฉมายํ นิสินฺเน อาสเน นิสีทิตพฺพํ, น เอกจงฺกเม จงฺกมิตพฺพํ, น นีเจ จงฺกเม จงฺกมนฺเต อุจฺเจ จงฺกเม จงฺกมิตพฺพํ, น ฉมายํ จงฺกมนฺเต จงฺกเม จงฺกมิตพฺพํ.}

\prose{น ภิกฺขเว มูลายปฏิกสฺสนารเหน ภิกฺขุนา ปาริวาสิเกน ภิกฺขุนา สทฺธิํ ฯเปฯ มูลายปฏิกสฺสนารเหน วุฑฺฒตเรน ภิกฺขุนา สทฺธิํ ฯเปฯ มานตฺตารเหน ภิกฺขุนา สทฺธิํ ฯเปฯ มานตฺตจาริเกน ภิกฺขุนา สทฺธิํ ฯเปฯ อพฺภานารเหน ภิกฺขุนา สทฺธิํ เอกจฺฉนฺเน อาวาเส วตฺถพฺพํ, น เอกจฺเฉนฺเน อนาวาเส วตฺตพฺพํ, น เอกจฺฉนฺเน อาวาเส วา อนาวาเส วา วตฺถพฺพํ, น เอกาสเน นิสีทิตพฺพํ, น นีเจ อาสเน นิสินฺเน อุจฺเจ อาสเน นิสีทิตพฺพํ, น ฉมายํ นิสินฺเน อาสเน นิสีทิตพฺพํ, น เอกจงฺกเม จงฺกมิตพฺพํ, น นีเจ จงฺกเม จงฺกมนฺเต อุจฺเจ จงฺกเม จงฺกมิตพฺพํ, น ฉมายํ จงฺกมนฺเต จงฺกเม จงฺกมิตพฺพํ.}

\prose{\csromansymbolfootnote{*}{วิ 3. 442 ปิฏฺเฐ.}มูลายปฏิกสฺสนารหจตุตฺโถ เจ ภิกฺขเว ปริวาสํ ทเทยฺย, มูลาย ปฏิกสฺเสยฺย, มานตฺตํ ทเทยฺย, ตํวีโส อพฺเภยฺย, อกมฺมํ น จ กรณียนฺติ.}

\nitthitam{มูลายปฏิกสฺสนารหวตฺตํ นิฏฺฐิตํ.}
\csromanmatikaanchor{mtk.55}
\csromantocmarkref{section}{3. มานตฺตารหวตฺต}{mtk.55}
\bookmark[dest={mtk.55},level=1]{3. มานตฺตารหวตฺต}
\csromanheader{1}{3. มานตฺตารหวตฺต}
\proseitem{88}{\csromanfolio{92}เตน โข ปน สมเยน มานตฺตารหา ภิกฺขู สาทิยนฺติ ปกตตฺตานํ ภิกฺขูนํ อภิวาทนํ ปจฺจุฏฺฐานํ อญฺชลิกมฺมํ สามีจิกมฺมํ อาสนาภิหารํ เสยฺยาภิหารํ ปาโททกํ ปาทปีฐํ ปาทกถลิกํ ปตฺตจีวรปฏิคฺคหณํ นหาเน ปิฏฺฐิปริกมฺมํ.\csromanspacer{} เย เต ภิกฺขู อปฺปิจฺฉา ฯเปฯ เต อุชฺฌายนฺติ ขิยฺยนฺติ วิปาเจนฺติ “กถํ หิ นาม มานตฺตารหา ภิกฺขู สาทิยิสฺสนฺติ ปกตตฺตานํ ภิกฺขูนํ อภิวาทนํ ปจฺจุฏฺฐานํ อญฺชลิกมฺมํ สามีจิกมฺมํ อาสนาภิหารํ เสยฺยาภิหารํ ปาโททกํ ปาทปีฐํ ปาทกถลิกํ ปตฺตจีวรปฏิคฺคหณํ นหาเน ปิฏฺฐิปริกมฺมนฺ”ติ.\csromanspacer{} อถ โข เต ภิกฺขู ภควโต เอตมตฺถํ อาโรเจสุํ ฯเปฯ.\csromanspacer{} สจฺจํ กิร ภิกฺขเว มานตฺตารหา ภิกฺขู สาทิยนฺติ ปกตตฺตานํ ภิกฺขูนํ อภิวาทนํ ปจฺจุฏฺฐานํ ฯเปฯ นหาเน ปิฏฺฐิปริกมฺมนฺติ.\csromanspacer{} สจฺจํ ภควาติ.\csromanspacer{} วิครหิ พุทฺโธ ภควา อนนุจฺฉวิกํ ฯเปฯ.\csromanspacer{} กถํ หิ นาม ภิกฺขเว มานตฺตารหา ภิกฺขู สาทิยิสฺสนฺติ ปกตตฺตานํ ภิกฺขูนํ อภิวาทนํ ฯเปฯ นหาเน ปิฏฺฐิปริกมฺมํ.\csromanspacer{} เนตํ ภิกฺขเว อปฺปสนฺนานํ วา ปสาทาย ฯเปฯ วิครหิตฺวา ฯเปฯ ธมฺมิํ กถํ กตฺวา ภิกฺขู อามนฺเตสิ\csromandash{}น ภิกฺขเว มานตฺตารเหน ภิกฺขุนา สาทิตพฺพํ ปกตตฺตานํ ภิกฺขูนํ อภิวาทนํ ปจฺจุฏฺฐานํ ฯเปฯ นหาเน ปิฏฺฐิปริกมฺมํ, โย สาทิเยยฺย, อาปตฺติ ทุกฺกฏสฺส.\csromanspacer{} อนุชานามิ ภิกฺขเว มานตฺตารหานํ ภิกฺขูนํ มิถุ ยถาวุฑฺฒํ อภิวาทนํ ปจฺจุฏฺฐานํ ฯเปฯ นหาเน ปิฏฺฐิปริกมฺมํ.\csromanspacer{} อนุชานามิ ภิกฺขเว มานตฺตารหานํ ภิกฺขูนํ ปญฺจ ยถาวุฑฺฒํ อุโปสถํ ปวารณํ วสฺสิกสาฏิกํ โอโณชนํ ภตฺตํ.\csromanspacer{} เตน หิ ภิกฺขเว มานตฺตารหานํ ภิกฺขูนํ วตฺตํ ปญฺญเปสฺสามิ, ยถา มานตฺตารเหหิ ภิกฺขูหิ วตฺติตพฺพํ.}

\proseitem{89}{มานตฺตารเหน ภิกฺขเว ภิกฺขุนา สมฺมา วตฺติตพฺพํ.\csromanspacer{} ตตฺรายํ สมฺมาวตฺตนา\csromandash{}}

\prose{น อุปสมฺปาเทตพฺพํ ฯเปฯ (ยถา มูลาย ปฏิกสฺสนา, ตถา วิตฺถาเรตพฺพํ.) น ภิกฺขูหิ สมฺปโยเชตพฺพํ.}

\prose{น ภิกฺขเว มานตฺตารเหน ภิกฺขุนา ปกตตฺตสฺส ภิกฺขุโน ปุรโต คนฺตพฺพํ, น ปุรโต นิสีทิตพฺพํ, โย โหติ สํฆสฺส อาสนปริยนฺโต เสยฺยาปริยนฺโต วิหารปริยนฺโต, โส ตสฺส ปทาตพฺโพ, เตน จ โส สาทิตพฺโพ.}

\prose{\csromanfolio{93}น ภิกฺขเว มานตฺตารเหน ภิกฺขุนา ปกตตฺเตน ภิกฺขุนา ปุเรสมเณน วา ปจฺฉาสมเณน วา กุลานิ อุปสงฺกมิตพฺพานิ, น อารญฺญิกงฺคํ สมาทาตพฺพํ, น ปิณฺฑปาติกงฺคํ สมาทาตพฺพํ, น จ ตปฺปจฺจยา ปิณฺฑปาโต นีหราเปตพฺโพ “มา มํ ชานิํสู”ติ.}

\prose{น ภิกฺขเว มานตฺตารเหน ภิกฺขุนา สภิกฺขุกา อาวาสา อภิกฺขุโก อาวาโส คนฺตพฺโพ อญฺญตฺร ปกตตฺเตน อญฺญตฺร อนฺตรายา.}

\prose{น ภิกฺขเว มานตฺตารเหน ภิกฺขุนา สภิกฺขุกา อาวาสา อภิกฺขุโก อนาวาโส คนฺตพฺโพ อญฺญตฺร ปกตตฺเตน อญฺญตฺร อนฺตรายา.}

\prose{น ภิกฺขเว มานตฺตารเหน ภิกฺขุนา สภิกฺขุกา อาวาสา อภิกฺขุโก อาวาโส วา อนาวาโส วา คนฺตพฺโพ อญฺญตฺร ปกตตฺเตน อญฺญตฺร อนฺตรายา.}

\prose{น ภิกฺขเว มานตฺตารเหน ภิกฺขุนา สภิกฺขุกา อนาวาสา อภิกฺขุโก อาวาโส คนฺตพฺโพ ฯเปฯ อภิกฺขุโก อนาวาโส คนฺตพฺโพ ฯเปฯ อภิกฺขุโก อาวาโส วา อนาวาโส วา คนฺตพฺโพ อญฺญตฺร ปกตตฺเตน อญฺญตฺร อนฺตรายา.}

\prose{น ภิกฺขเว มานตฺตารเหน ภิกฺขุนา สภิกฺขุกา อาวาสา วา อนาวาสา วา อภิกฺขุโก อาวาโส คนฺตพฺโพ ฯเปฯ อภิกฺขุโก อนาวาโส คนฺตพฺโพ ฯเปฯ อภิกฺขุโก อาวาโส วา อนาวาโส วา คนฺตพฺโพ อญฺญตฺร ปกตตฺเตน อญฺญตฺร อนฺตรายา.}

\prose{น ภิกฺขเว มานตฺตารเหน ภิกฺขุนา สภิกฺขุกา อาวาสา สภิกฺขุโก อาวาโส คนฺตพฺโพ ฯเปฯ สภิกฺขุโก อนาวาโส คนฺตพฺโพ ฯเปฯ สภิกฺขุโก อาวาโส วา อนาวาโส วา คนฺตพฺโพ ยตฺถสฺสุ ภิกฺขู นานาสํวาสกา อญฺญตฺร ปกตตฺเตน อญฺญตฺร อนฺตรายา.}

\prose{น ภิกฺขเว มานตฺตารเหน ภิกฺขุนา สภิกฺขุกา อนาวาสา สภิกฺขุโก อาวาโส คนฺตพฺโพ ฯเปฯ สภิกฺขุโก อนาวาโส \csromanfolio{94}คนฺตพฺโพ ฯเปฯ สภิกฺขุโก อาวาโส วา อนาวาโส วา คนฺตพฺโพ ยตฺถสฺสุ ภิกฺขู นานาสํวาสกา อญฺญตฺร ปกตตฺเตน อญฺญตฺร อนฺตรายา.}

\prose{น ภิกฺขเว มานตฺตารเหน ภิกฺขุนา สภิกฺขุกา อาวาสา วา อนาวาสา วา สภิกฺขุโก อาวาโส คนฺตพฺโพ ฯเปฯ สภิกฺขุโก อนาวาโส คนฺตพฺโพ ฯเปฯ สภิกฺขุโก อาวาโส วา อนาวาโส วา คนฺตพฺโพ ยตฺถสฺสุ ภิกฺขู นานาสํวาสกา อญฺญตฺร ปกตตฺเตน อญฺญตฺร อนฺตรายา.}

\prose{คนฺตพฺโพ ภิกฺขเว มานตฺตารเหน ภิกฺขุนา สภิกฺขุกา อาวาสา สภิกฺขุโก อาวาโส ฯเปฯ สภิกฺขุโก อนาวาโส ฯเปฯ สภิกฺขุโก อาวาโส วา อนาวาโส วา ยตฺถสฺสุ ภิกฺขู สมานสํวาสกา ยํ ชญฺญา “สกฺโกมิ อชฺเชว คนฺตุนฺ”ติ.}

\prose{คนฺตพฺโพ ภิกฺขเว มานตฺตารเหน ภิกฺขุนา สภิกฺขุกา อนาวาสา สภิกฺขุโก อาวาโส ฯเปฯ สภิกฺขุโก อนาวาโส ฯเปฯ สภิกฺขุโก อาวาโส วา อนาวาโส วา ยตฺถสฺสุ ภิกฺขู สมานสํวาสกา ยํ ชญฺญา “สกฺโกมิ อชฺเชว คนฺตุนฺ”ติ.}

\prose{คนฺตพฺโพ ภิกฺขเว มานตฺตารเหน ภิกฺขุนา สภิกฺขุกา อาวาสา วา อนาวาสา วา สภิกฺขุโก อาวาโส ฯเปฯ สภิกฺขุโก อนาวาโส ฯเปฯ สภิกฺขุโก อาวาโส วา อนาวาโส วา ยตฺถสฺสุ ภิกฺขู สมานสํวาสกา ยํ ชญฺญา “สกฺโกมิ อชฺเชว คนฺตุนฺ”ติ.}

\prose{น ภิกฺขเว มานตฺตารเหน ภิกฺขุนา ปกตตฺเตน ภิกฺขุนา สทฺธิํ เอกจฺฉนฺเน อาวาเส วตฺถพฺพํ, น เอกจฺฉนฺเน อนาวาเส วตฺถพฺพํ, น เอกจฺฉนฺเน อาวาเส วา อนาวาเส วา วตฺถพฺพํ, ปกตตฺตํ ภิกฺขุํ ทิสฺวา อาสนา วุฏฺฐาตพฺพํ, ปกตตฺโต ภิกฺขุ อาสเนน นิมนฺเตตพฺโพ, น ปกตตฺเตน ภิกฺขุนา สทฺธิํ เอกาสเน นิสีทิตพฺพํ, น นีเจ อาสเน นิสินฺเน อุจฺเจ อาสเน นิสีทิตพฺพํ, น ฉมายํ นิสินฺเน อาสเน นิสีทิตพฺพํ, น เอกจงฺกเม จงฺกมิตพฺพํ, น นีเจ จงฺกเม จงฺกมนฺเต อุจฺเจ จงฺกเม จงฺกมิตพฺพํ, น ฉมายํ จงฺกมนฺเต จงฺกเม จงฺกมิตพฺพํ.}

\prose{\csromanfolio{95}น ภิกฺขเว มานตฺตารเหน ภิกฺขุนา ปาริวาสิเกน ภิกฺขุนา สทฺธิํ ฯเปฯ มูลายปฏิกสฺสนารเหน ภิกฺขุนา สทฺธิํ ฯเปฯ มานตฺตารเหน วุฑฺฒตเรน ภิกฺขุนา สทฺธิํ ฯเปฯ มานตฺตจาริเกน ภิกฺขุนา สทฺธิํ ฯเปฯ อพฺภานาเรเหน ภิกฺขุนา สทฺธิํ เอกจฺฉนฺเน อาวาเส วตฺถพฺพํ, น เอกจฺฉนฺเน อนาวาเส วตฺถพฺพํ, น เอกจฺฉนฺเน อาวาเส วา อนาวาเส วา วตฺถพฺพํ, น เอกาสเน นิสีทิตพฺพํ, น นีเจ อาสเน นิสินฺเน อุจฺเจ อาสเน นิสีทิตพฺพํ, น ฉมายํ นิสินฺเน อาสเน นิสีทิตพฺพํ, น เอกจงฺกเม จงฺกมิตพฺพํ, น นีเจ จงฺกเม จงฺกมนฺเต อุจฺเจ จงฺกเม จงฺกมิตพฺพํ, น ฉมายํ จงฺกมนฺเต จงฺกเม จงฺกมิตพฺพํ.}

\prose{\csromansymbolfootnote{*}{วิ 3. 442 ปิฏฺเฐ.}มานตฺตารหจตุตฺโถ เจ ภิกฺขเว ปริวาสํ ทเทยฺย, มูลาย ปฏิกสฺเสยฺย, มานตฺตํ ทเทยฺย, ตํวีโส อพฺเภยฺย, อกมฺมํ น จ กรณียนฺติ.}

\nitthitam{มานตฺตารหวตฺตํ นิฏฺฐิตํ.}
\csromanmatikaanchor{mtk.56}
\csromantocmarkref{section}{4. มานตฺตจาริกวตฺต}{mtk.56}
\bookmark[dest={mtk.56},level=1]{4. มานตฺตจาริกวตฺต}
\csromanheader{1}{4. มานตฺตจาริกวตฺต}
\proseitem{90}{เตน โข ปน สมเยน มานตฺตจาริกา ภิกฺขู สาทิยนฺติ ปกตตฺตานํ ภิกฺขูนํ อภิวาทนํ ปจฺจุฏฺฐานํ อญฺชลิกมฺมํ สามีจิกมฺมํ อาสนาภิหารํ เสยฺยาภิหารํ ปาโททกํ ปาทปีฐํ ปาทกถลิกํ ปตฺตจีวรปฏิคฺคหณํ นหาเน ปิฏฺฐิปริกมฺมํ.\csromanspacer{} เย เต ภิกฺขู อปฺปิจฺฉา ฯเปฯ เต อุชฺฌายนฺติ ขิยฺยนฺติ วิปาเจนฺติ “กถํ หิ นาม มานตฺตจาริกา ภิกฺขู สาทิยิสฺสนฺติ ปกตตฺตานํ ภิกฺขูนํ อภิวาทนํ ปจฺจุฏฺฐานํ ฯเปฯ นหาเน ปิฏฺฐิปริกมฺมนฺ”ติ.\csromanspacer{} อถ โข เต ภิกฺขู ภควโต เอตมตฺถํ อาโรเจสุํ.}

\prose{อถ โข ภควา เอตสฺมิํ นิทาเน เอตสฺมิํ ปกรเณ ภิกฺขุสํฆํ สนฺนิปาตาเปตฺวา ภิกฺขู ปฏิปุจฺฉิ “สจฺจํ กิร ภิกฺขเว มานตฺตจาริกา ภิกฺขู สาทิยนฺติ ปกตตฺตานํ ภิกฺขูนํ อภิวาทนํ ปจฺจุฏฺฐานํ ฯเปฯ นหาเน ปิฏฺฐิปริกมฺมนฺ”ติ.\csromanspacer{} สจฺจํ ภควาติ.\csromanspacer{} วิครหิ พุทฺโธ ภควา อนนุจฺฉวิกํ ฯเปฯ.\csromanspacer{} กถํ หิ นาม ภิกฺขเว มานตฺตจาริกา ภิกฺขู สาทิยิสฺสนฺติ ปกตตฺตานํ ภิกฺขูนํ อภิวาทนํ ฯเปฯ นหาเน ปิฏฺฐิปริกมฺมํ.\csromanspacer{} เนตํ ภิกฺขเว อปฺปสนฺนานํ วา ปสาทาย ฯเปฯ วิครหิตฺวา ฯเปฯ ธมฺมิํ กถํ กตฺวา ภิกฺขู อามนฺเตสิ\csromandash{}}

\prose{\csromanfolio{96}น ภิกฺขเว มานตฺตจาริเกน ภิกฺขุนา สาทิตพฺพํ ปกตตฺตานํ ภิกฺขูนํ อภิวาทนํ ปจฺจุฏฺฐานํ อญฺชลิกมฺมํ สามีจิกมฺมํ อาสนาภิหาโร เสยฺยาภิหาโร ปาโททกํ ปาทปีฐํ ปาทกถลิกํ ปตฺตจีวรปฏิคฺคหณํ นหาเน ปิฏฺฐิปริกมฺมํ, โย สาทิเยยฺย, อาปตฺติ ทุกฺกฏสฺส.\csromanspacer{} อนุชานามิ ภิกฺขเว มานตฺตจาริกานํ ภิกฺขูนํ มิถุ ยถาวุฑฺฒํ อภิวาทนํ ปจฺจุฏฺฐานํ ฯเปฯ นหาเน ปิฏฺฐิปริกมฺมํ.\csromanspacer{} อนุชานามิ ภิกฺขเว มานตฺตจริกานํ ภิกฺขูนํ ปญฺจ ยถาวุฑฺฒํ อุโปสถํ ปวารณํ วสฺสิกสาฏิกํ โอโณชนํ ภตฺตํ.\csromanspacer{} เตน หิ ภิกฺขเว มานตฺตจาริกานํ ภิกฺขูนํ วตฺตํ ปญฺญเปสฺสามิ, ยถา มานตฺตจาริเกหิ ภิกฺขูหิ วตฺติตพฺพํ.}

\proseitem{91}{มานตฺตจาริเกน ภิกฺขเว ภิกฺขุนา สมฺมา วตฺติตพฺพํ.\csromanspacer{} ตตฺรายํ สมฺมาวตฺตนา\csromandash{}}

\prose{น อุปสมฺปาเทตพฺพํ, น นิสฺสโย ทาตพฺโพ, น สามเณโร อุปฏฺฐาเปตพฺโพ, น ภิกฺขุโนวาทกสมฺมุติ สาทิตพฺพา, สมฺมเตนปิ ภิกฺขุนิโย น โอวทิตพฺพา, ยาย อาปตฺติยา สํเฆน มานตฺตํ ทินฺนํ โหติ, สา อาปตฺติ น อาปชฺชิตพฺพา, อญฺญา วา ตาทิสิกา, ตโต วา ปาปิฏฺฐตรา, กมฺมํ น ครหิตพฺพํ, กมฺมิกา น ครหิตพฺพา, น ปกตตฺตสฺส ภิกฺขุโน อุโปสโถ ฐเปตพฺโพ, น ปวารณา ฐเปตพฺพา, น สวจนียํ กาตพฺพํ, น อนุวาโท ปฏฺฐเปตพฺโพ, น โอกาโส กาเรตพฺโพ น โจเทตพฺโพ, น สาเรตพฺโพ, น ภิกฺขูหิ สมฺปโยเชตพฺพํ.}

\prose{น ภิกฺขเว มานตฺตจาริเกน ภิกฺขุนา ปกตตฺตสฺส ภิกฺขุโน ปุรโต คนฺตพฺพํ, น ปุรโต นิสีทิตพฺพํ, โย โหติ สํฆสฺส อาสนปริยนฺโต เสยฺยาปริยนฺโต วิหารปริยนฺโต, โส ตสฺส ปทาตพฺโพ, เตน จ โส สาทิตพฺโพ.}

\prose{น ภิกฺขเว มานตฺตจาริเกน ภิกฺขุนา ปกตตฺเตน ภิกฺขุนา ปุเรสมเณน วา ปจฺฉาสมเณน วา กุลานิ อุปสงฺกมิตพฺพานิ, น อารญฺญิกงฺคํ สมาทาตพฺพํ, น ปิณฺฑปาติกงฺคํ สมาทาตพฺพํ, น จ ตปฺปจฺจยา ปิณฺฑปาโต นีหราเปตพฺโพ “มา มํ ชานิํสู”ติ.}

\prose{มานตฺตจาริเกน ภิกฺขเว ภิกฺขุนา อาคนฺตุเกน อาโรเจตพฺพํ, อาคนฺตุกสฺส อาโรเจตพฺพํ, อุโปสเถ อาโรเจตพฺพํ, ปวารณาย \csromanfolio{97}อาโรเจตพฺพํ, เทวสิกํ อาโรเจตพฺพํ, สเจ คิลาโน โหติ, ทูเตนปิ อาโรเจตพฺพํ.}

\prose{น ภิกฺขเว มานตฺตจาริเกน ภิกฺขุนา สภิกฺขุกา อาวาสา อภิกฺขุโก อาวาโส คนฺตพฺโพ อญฺญตฺร สํเฆน อญฺญตฺร อนฺตรายา.}

\prose{น ภิกฺขเว มานตฺตจาริเกน ภิกฺขุนา สภิกฺขุกา อาวาสา อภิกฺขุโก อนาวาโส คนฺตพฺโพ อญฺญตฺร สํเฆน อญฺญตฺร อนฺตรายา.}

\prose{น ภิกฺขเว มานตฺตจาริเกน ภิกฺขุนา สภิกฺขุกา อาวาสา อภิกฺขุโก อาวาโส วา อนาวาโส วา คนฺตพฺโพ อญฺญตฺร สํเฆน อญฺญตฺร อนฺตรายา.}

\prose{น ภิกฺขเว มานตฺตจาริเกน ภิกฺขุนา สภิกฺขุกา อนาวาสา อภิกฺขุโก อาวาโส คนฺตพฺโพ ฯเปฯ อภิกฺขุโก อนาวาโส คนฺตพฺโพ ฯเปฯ อภิกฺขุโก อาวาโส วา อนาวาโส วา คนฺตพฺโพ อญฺญตฺร สํเฆน อญฺญตฺร อนฺตรายา.}

\prose{น ภิกฺขเว มานตฺตจาริเกน ภิกฺขุนา สภิกฺขุกา อาวาสา วา อนาวาสา วา อภิกฺขุโก อาวาโส คนฺตพฺโพ ฯเปฯ อภิกฺขุโก อนาวาโส คนฺตพฺโพ ฯเปฯ อภิกฺขุโก อาวาโส วา อนาวาโส วา คนฺตพฺโพ อญฺญตฺร สํเฆน อญฺญตฺร อนฺตรายา.}

\prose{น ภิกฺขเว มานตฺตจาริเกน ภิกฺขุนา สภิกฺขุกา อาวาสา สภิกฺขุโก อาวาโส คนฺตพฺโพ ฯเปฯ สภิกฺขุโก อานาวาโส คนฺตพฺโพ ฯเปฯ สภิกฺขุโก อาวาโส วา อนาวาโส วา คนฺตพฺโพ ยตฺถสฺสุ ภิกฺขู นานาสํวาสกา อญฺญตฺร สํเฆน อญฺญตฺร อนฺตรายา.}

\prose{น ภิกฺขเว มานตฺตจาริเกน ภิกฺขุนา สภิกฺขุกา อนาวาสา สภิกฺขุโก อาปาโส คนฺตพฺโพ ฯเปฯ สภิกฺขุโก อนาวาโส คนฺตพฺโพ ฯเปฯ สภิกฺขุโก อาวาโส วา อนาวาโส วา คนฺตพฺโพ ยตฺถสฺสุ ภิกฺขู นานาสํวาสกา อญฺญตฺร สํเฆน อญฺญตฺร อนฺตรายา.}

\prose{น ภิกฺขเว มานตฺตจาริเกน ภิกฺขุนา สภิกฺขุกา อาวาสา วา อนาวาสา วา สภิกฺขุโก อาวาโส คนฺตพฺโพ ฯเปฯ สภิกฺขุโก อนาวาโส คนฺตพฺโพ ฯเปฯ สภิกฺขุโก อาวาโส วา อนาวาโส วา \csromanfolio{98}คนฺตพฺโพ ยตฺถสฺสุ ภิกฺขู นานาสํวาสกา อญฺญตฺร สํเฆน อญฺญตฺร อนฺตรายา.}

\prose{คนฺตพฺโพ ภิกฺขเว มานตฺตจาริเกน ภิกฺขุนา สภิกฺขุกา อาวาสา สภิกฺขุโก อาวาโส ฯเปฯ สภิกฺขุโก อนาวาโส ฯเปฯ สภิกฺขุโก อาวาโส วา อนาวาโส วา ยตฺถสฺสุ ภิกฺขู สมานสํวาสกา ยํ ชญฺญา “สกฺโกมิ อชฺเชว คนฺตุนฺ”ติ.}

\prose{คนฺตพฺโพ ภิกฺขเว มานตฺตจาริเกน ภิกฺขุนา สภิกฺขุกา อนาวาสา สภิกฺขุโก อาวาโส ฯเปฯ สภิกฺขุโก อนาวาโส ฯเปฯ สภิกฺขุโก อาวาโส วา อนาวาโส วา ยตฺถสฺสุ ภิกฺขู สมานสํวาสกา ยํ ชญฺญา “สกฺโกมิ อชฺเชว คนฺตุนฺ”ติ.}

\prose{คนฺตพฺโพ ภิกฺขเว มานตฺตจาริเกน ภิกฺขุนา สภิกฺขุกา อาวาสา วา อนาวาสา วา สภิกฺขุโก อาวาโส ฯเปฯ สภิกฺขุโก อนาวาโส ฯเปฯ สภิกฺขุโก อาวาโส วา อนาวาโส วา ยตฺถสฺสุ ภิกฺขู สมานสํวาสกา ยํ ชญฺญา “สกฺโกมิ อชฺเชว คนฺตุนฺ”ติ.}

\prose{น ภิกฺขเว มานตฺตจาริเกน ภิกฺขุนา ปกตตฺเตน ภิกฺขุนา สทฺธิํ เอกจฺฉนฺเน อาวาเส วตฺถพฺพํ, น เอกจฺฉนฺเน อนาวาเส วตฺถพฺพํ, น เอกจฺฉนฺเน อาวาเส วา อนาวาเส วา วตฺถพฺพํ, ปกตตฺตํ ภิกฺขุํ ทิสฺวา อาสนา วุฏฺฐาตพฺพํ, ปกตตฺโต ภิกฺขุ อาสเนน นิมนฺเตตพฺโพ, น ปกตตฺเตน ภิกฺขุนา สทฺธิํ เอกาสเน นิสีทิตพฺพํ, น นีเจ อาสเน นิสินฺเน อุจฺเจ อาสเน นิสีทิตพฺพํ, น ฉมายํ นิสินฺเน อาสเน นิสีทิตพฺพํ, น เอกจงฺกเม จงฺกมิตพฺพํ, น นีเจ จงฺกเม จงฺกมนฺเต อุจฺเจ จงฺกเม จงฺกมิตพฺพํ, น ฉมายํ จงฺกมนฺเต จงฺกมิตพฺพํ.}

\prose{น ภิกฺขเว มานตฺตจาริเกน ภิกฺขุนา ปาริวาสิเกน ภิกฺขุนา สทฺธิํ ฯเปฯ มูลายปฏิกสฺสนารเหน ภิกฺขุนา สทฺธิํ ฯเปฯ มานตฺตารเหน ภิกฺขุนา สทฺธิํ ฯเปฯ มานตฺตจาริเกน วุฑฺฒตเรน ภิกฺขุนา สทฺธิํ ฯเปฯ อพฺภานารเหน ภิกฺขุนา สทฺธิํ เอกจฺฉนฺเน อาวาเส วตฺถพฺพํ, น เอกจฺฉนฺเน อนาวาเส วตฺถพฺพํ, น เอกจฺฉนฺเน อาวาเส วา อนาวเส วา วตฺถพฺพํ, น เอกาสเน นิสีทิตพฺพํ, น นีเจ อาสเน นิสินฺเน อุจฺเจ อาสเน นิสีทิตพฺพํ, น ฉมายํ นิสินฺเน อาสเน นิสีทิตพฺพํ, น เอกจงฺกเม \csromanfolio{99}จงฺกมิตพฺพํ, น นีเจ จงฺกเม จงฺกมนฺเต อุจฺเจ จงฺกเม จงฺกมิตพฺพํ, น ฉมายํ จงฺกมนฺเต จงฺกเม จงฺกมิตพฺพํ.}

\prose{\csromansymbolfootnote{*}{วิ 3. 442 ปิฏฺเฐ.}มานตฺตจาริกจตุตฺโถ เจ ภิกฺขเว ปริวาสํ ทเทยฺย, มูลาย ปฏิกสฺเสยฺย, มานตฺตํ ทเทยฺย, ตํวีโส อพฺเภยฺย, อกมฺมํ น จ กรณียนฺติ.}

\proseitem{92}{อถ โข อายสฺมา อุปาลิ เยน ภควา เตนุปสงฺกมิ, อุปสงฺกมิตฺวา ภควนฺตํ อภิวาเทตฺวา เอกมนฺตํ นิสีทิ, เอกมนฺตํ นิสินฺโน โข อายสฺมา อุปาลิ ภควนฺตํ เอตทโวจ “กติ นุ โข ภนฺเต มานตฺตจาริกสฺส ภิกฺขุโน รตฺติจฺเฉทา”ติ.\csromanspacer{} จตฺตาโร โข อุปาลิ มานตฺตจาริกสฺส ภิกฺขุโน รตฺติจฺเฉทา, สหวาโส วิปฺปวาโส อนาโรจนา อูเน คเณ จรณํ\footnote{จรณนฺติ (ก)}, อิเม โข อุปาลิ จตฺตาโร มานตฺตจาริกสฺส ภิกฺขุโน รตฺติจฺเฉทาติ.}

\proseitem{93}{เตน โข ปน สมเยน สาวตฺถิยํ มหา ภิกฺขุสํโฆ สนฺนิปติโต โหติ, น สกฺโกนฺติ มานตฺตจาริกา ภิกฺขู มานตฺตํ โสเธตุํ, ภควโต เอตมตฺถํ อาโรเจสุํ.\csromanspacer{} อนุชานามิ ภิกฺขเว มานตฺตํ นิกฺขิปิตุํ.\csromanspacer{} เอวญฺจ ปน ภิกฺขเว นิกฺขิปิตพฺพํ, เตน มานตฺตจาริเกน ภิกฺขุนา เอกํ ภิกฺขุํ อุปสงฺกมิตฺวา เอกํสํ อุตฺตราสงฺคํ กริตฺวา อุกฺกุฏิกํ นิสีทิตฺวา อญฺชลิํ ปคฺคเหตฺวา เอวมสฺส วจนีโย “มานตฺตํ นิกฺขิปามี”ติ, นิกฺขิตฺตํ โหติ มานตฺตํ. “วตฺตํ นิกฺขิปามี”ติ, นิกฺขิตฺตํ โหติ มานตฺตนฺติ.}

\proseitem{94}{เตน โข ปน สมเยน สาวตฺถิยา ภิกฺขู ตหํ ตหํ ปกฺกมิํสุ, สกฺโกนฺติ มานตฺตจาริกา ภิกฺขู มานตฺตํ โสเธตุํ, ภควโต เอตมตฺถํ อาโรเจสุํ.\csromanspacer{} อนุชานามิ ภิกฺขเว มานตฺตํ สมาทิยิตุํ.\csromanspacer{} เอวญฺจ ปน ภิกฺขเว สมาทิยิตพฺพํ, เตน มานตฺตจาริเกน ภิกฺขุนา เอกํ ภิกฺขุํ อุปสงฺกมิตฺวา เอกํสํ อุตฺตราสงฺคํ กริตฺวา อุกฺกุฏิกํ นิสีทิตฺวา อญฺชลิํ ปคฺคเหตฺวา เอวมสฺส วจนีโย “มานตฺตํ สมาทิยามี”ติ, สมาทินฺนํ โหติ มานตฺตํ. “วตฺตํ สมาทิยามี”ติ, สมาทินฺนํ โหติ มานตฺตนฺติ.}

\nitthitam{มานตฺตจาริกวตฺตํ นิฏฺฐิตํ.}
\csromanheader{2}{5. อพฺภานารตวตฺต}
\csromanmatikaanchor{mtk.57}
\csromantocmarkref{section}{5. อพฺภานารหวตฺต}{mtk.57}
\bookmark[dest={mtk.57},level=1]{5. อพฺภานารหวตฺต}
\proseitem{95}{\csromanfolio{100}เตน โข ปน สมเยน อพฺภานารหาภิกฺขู สาทิยนฺติ ปกตตฺตานํ ภิกฺขูนํ อภิวาทนํ ปจฺจุฏฺฐานํ ฯเปฯ นหาเน ปิฏฺฐิปริกมฺมํ.\csromanspacer{} เย เต ภิกฺขู อปฺปิจฺฉา ฯเปฯ เต อุชฺฌายนฺติ ขิยฺยนฺติ วิปาเจนฺติ “กถํ หิ นาม อพฺภานารหา ภิกฺขู สาทิยิสฺสนฺติ ปกตตฺตานํ ภิกฺขูนํ อภิวาทนํ ปจฺจุฏฺฐานํ ฯเปฯ นหาเน ปิฏฺฐิปริกมฺมนฺ”ติ.\csromanspacer{} อถ โข เต ภิกฺขู ภควโต เอตมตฺถํ อาโรเจสุํ.}

\prose{อถ โข ภควา เอตสฺมิํ นิทาเน เอตสฺมิํ ปกรเณ ภิกฺขุสํฆํ สนฺนิปาตาเปตฺวา ภิกฺขู ปฏิปุจฺฉิ “สจฺจํ กิร ภิกฺขเว อพฺภานารหา ภิกฺขู สาทิยนฺติ ปกตตฺตานํ ภิกฺขูนํ อภิวาทนํ ปจฺจุฏฺฐานํ ฯเปฯ นหาเน ปิฏฺฐิปริกมฺมนฺ”ติ, สจฺจํ ภควาติ.\csromanspacer{} วิครหิ พุทฺโธ ภควา อนนุจฺฉวิกํ ฯเปฯ.\csromanspacer{} กถํ หิ นาม ภิกฺขเว อพฺภานารหา ภิกฺขู สาทิยิสฺสนฺติ ปกตตฺถานํ ภิกฺขูนํ อภิวาทนํ ปจฺจุฏฺฐานํ ฯเปฯ นหาเน ปิฏฺฐิปริกมฺมํ.\csromanspacer{} เนตํ ภิกฺขเว อปฺปสนฺนานํ วา ปสาทาย ฯเปฯ วิครหิตฺวา ฯเปฯ ธมฺมิํ กถํ กตฺวา ภิกฺขู อามนฺเตสิ\csromandash{} น ภิกฺขเว อพฺภานารเหน ภิกฺขุนา สาทิตพฺพํ ปกตตฺตานํ ภิกฺขูนํ อภิวาทนํ ปจฺจุฏฺฐานํ ฯเปฯ นหาเน ปิฏฺฐิปริกมฺมํ.\csromanspacer{} โย สาทิเยยฺย, อาปตฺติ ทุกฺกฏสฺส.\csromanspacer{} อนุชานามิ ภิกฺขเว อพฺภานารหานํ ภิกฺขูนํ มิถุ ยถาวุฑฺฒํ อภิวาทนํ ปจฺจุฏฺฐานํ ฯเปฯ นหาเน ปิฏฺฐิปริกมฺมํ, อนุชานามิ ภิกฺขเว อพฺภานารหานํ ภิกฺขูนํ ปญฺจ ยถาวุฑฺฒํ อุโปสถํ ปวารณํ วสฺสิกสาฏิกํ โอโณชนํ ภตฺตํ.\csromanspacer{} เตน หิ ภิกฺขเว อพฺภานารหานํ ภิกฺขูนํ ปตฺตํ ปญฺญ ปสฺสามิ, ยถา อพฺภานารเหหิ ภิกฺขูหิ วตฺติตพฺพํ.}

\proseitem{96}{อพฺภานารเหน ภิกฺขเว ภิกฺขุนา สมฺมา วตฺติตพฺพํ.\csromanspacer{} ตตฺรายํ สมฺมาวตฺตนา\csromandash{}}

\prose{น อุปสมฺปาเทตพฺพํ ฯเปฯ (ยถา เหฏฺฐา, ตถา วิตฺถาเรตพฺพํ,) น ภิกฺขูหิ สมฺปโยเชตพฺพํ.}

\prose{น ภิกฺขเว อพฺภานารเหน ภิกฺขุนา ปกตตฺตสฺส ภิกฺขุโน ปุรโต คนฺตพฺพํ, น ปุรโต นิสีทิตพฺพํ, โย โหติ สํฆสฺส อาสนปริยนฺโต เสยฺยาปริยนฺโต วิหารปริยนฺโต, โส ตสฺส ปทาตพฺโพ, เตน จ โส สาทิตพฺโพ.}

\prose{\csromanfolio{101}น ภิกฺขเว อพฺภานารเหน ภิกฺขุนา ปกตตฺเตน ภิกฺขุนา ปุเรสมเณน วา ปจฺฉาสมเณน วา กุลานิ อุปสงฺกมิตพฺพานิ, น อารญฺญิกงฺคํ สมาทาตพฺพํ, น ปิณฺฑปาติกงฺคํ สมาทาตพฺพํ, น จ ตปฺปจฺจยา ปิณฺฑปาโต นีหราเปตพฺโพ “มา มํ ชานิํสู”ติ.}

\prose{น ภิกฺขเว อพฺภานารเหน ภิกฺขุนา สภิกฺขุกา อาวาสา อภิกฺขุโก อาวาโส คนฺตพฺโพ อญฺญตฺร ปกตตฺเตน อญฺญตฺร อนฺตรายา.}

\prose{น ภิกฺขเว อพฺภานารเหน ภิกฺขุนา สภิกฺขุกา อาวาสา อภิกฺขุโก อนาวาโส คนฺตพฺโพ อญฺญตฺร ปกตตฺเตน อญฺญตฺร อนฺตรายา ฯเปฯ (ยถา เหฏฺฐา, ตถา วิตฺถาเรตพฺพา.)}

\prose{คนฺตพฺโพ ภิกฺขเว อพฺภานารเหน ภิกฺขุนา สภิกฺขุกา อาวาสา ฯเปฯ อนาวาสา ฯเปฯ อาวาสา วา อนาวาสา วา สภิกฺขุโก อาวาโส ฯเปฯ สภิกฺขุโก อนาวาโส ฯเปฯ สภิกฺขุโก อาวาโส วา อนาวาโส วา ยตฺถสฺสุ ภิกฺขู สมานสํวาสกา ยํ ชญฺญา “สกฺโกมิ อชฺเชว คนฺตุนฺ”ติ.}

\prose{น ภิกฺขเว อพฺภานารเหน ภิกฺขุนา ปกตตฺเตน ภิกฺขุนา สทฺธิํ เอกจฺฉนฺเน อาวาเส วตฺถพฺพํ, น เอกจฺฉนฺเน อนาวาเส วตฺถพฺพํ, น เอกจฺฉนฺเน อาวาเส วา อนาวาเส วา วตฺถพฺพํ, ปกตตฺตํ ภิกฺขุํ ทิสฺวา อาสนา วุฏฺฐาตพฺพํ, ปกตตฺโต ภิกฺขุ อาสเนน นิมนฺเตตพฺโพ, น ปกตตฺเตน ภิกฺขุนา สทฺธิํ เอกาสเน นิสีทิตพฺพํ, น นีเจ อาสเน นิสินฺเน อุจฺเจ อาสเน นิสีทิตพฺพํ, น ฉมายํ นิสินฺเน อาสเน นิสีทิตพฺพํ, น เอกจงฺกเม จงฺกมิตพฺพํ, น นีเจ จงฺกเม จงฺกมนฺเต อุจฺเจ จงฺกเม จงฺกมิตพฺพํ, น ฉมายํ จงฺกมนฺเต จงฺกเม จงฺกมิตพฺพํ.}

\csromanmatikaanchor{mtk.58}
\csromantocmarkref{section}{อุทฺทานคาถา}{mtk.58}
\bookmark[dest={mtk.58},level=1]{อุทฺทานคาถา}
\prose{น ภิกฺขเว อพฺภานารเหน ภิกฺขุนา ปาริวาสิเกน ภิกฺขุนา สทฺธิํ ฯเปฯ มูลายปฏิกสฺสนารเหน ภิกฺขุนา สทฺธิํ ฯเปฯ มานตฺตารเหน ภิกฺขุนา สทฺธิํ ฯเปฯ มานตฺตจาริเกน ภิกฺขุนา สทฺธิํ ฯเปฯ อพฺภานารเหน วุฑฺฒตเรน ภิกฺขุนา สทฺธิํ เอกจฺฉนฺเน อาวาเส วตฺถพฺพํ, น เอกจฺฉนฺเน อนาวาเส วตฺถพฺพํ, น เอกจฺฉนฺเน อาวาเส วา อนาวาเส วา วตฺถพฺพํ, น เอกาสเน นิสีทิตพฺพํ, น นีเจ อาสเน นิสินฺเน อุจฺเจ อาสเน นิสีทิตพฺพํ, น ฉมายํ นิสินฺเน อาสเน นิสีทิตพฺพํ, น เอกจงฺกเม \csromanfolio{102}จงฺกมิตพฺพํ, น นีเจ จงฺกเม จงฺกมนฺเต อุจฺเจ จงฺกเม จงฺกมิตพฺพํ, น ฉมายํ จงฺกมนฺเต จงฺกเม จงฺกมิตพฺพํ.}

\prose{\csromansymbolfootnote{*}{วิ 3. 442 ปิฏฺเฐ.}อพฺภานารหจตุตฺโถ เจ ภิกฺขเว ปริวาสํ ทเทยฺย, มูลาย ปฏิกสฺเสยฺย, มานตฺตํ ทเทยฺย, ตํวีโส อพฺเภยฺย, อกมฺมํ น จ กรณียนฺติ.}

\vspace{\tipitakaparskip}
\csromancenter{อพฺภานารหวตฺตํ นิฏฺฐิตํ.\csromanspacer{} ปาริวาสิกกฺขนฺธโก ทุติโย.\csromanspacer{} อิมมฺหิ ขนฺธเก วตฺถู ปญฺจ.}
\titlehead{\textbf{ตสฺสุทฺทานํ}}
\csromangathameasure{ปาริวาสิกา สาเทนฺติ, ปกตตฺตาน ภิกฺขุนํ. \\ อภิวาทนํ ปจฺจุฏฺฐานํ, อญฺชลิํ จ สามีจิยํ. \\ อาสนํ เสยฺยาภิหารํ, ปาโท ปีฐํ กถลิกํ. \\ ปตฺตํ นหาเน ปริกมฺมํ, อุชฺฌายนฺติ จ เปสลา. \\ ทุกฺกฏํ สาทิยนฺตสฺส, มิถุ ปญฺจ ยถาวุฑฺฒํ. \\ อุโปสถํ ปวารณํ, วสฺสิโกโณชโภชนํ. \\ สมฺมา จ วตฺตนา ตตฺถ, ปกตตฺตสฺส คจฺฉนฺตํ. \\ โย จ โหติ ปริยนฺโต, ปุเร ปจฺฉา ตเถว จ. \\ อารญฺญปิณฺฑนีหาโร, อาคนฺตุเก อุโปสเถ. \\ ปวารณาย ทูเตน, คนฺตพฺโพ จ สภิกฺขุโก. \\ เอกจฺฉนฺเน จ วุฏฺฐานํ, ตเถว จ นิมนฺตเย. \\ อาสเน นีเจ จงฺกเม, ฉมายํ จงฺกเมน จ. \\ วุฑฺฒตเรน อกมฺมํ, รตฺติจฺเฉทา จ โสธนา. \\ นิกฺขิปนํ สมาทานํ, วตฺตํว ปาริวาสิเก. \\ มูลาย มานตฺตารหา, ตถา มานตฺตจาริตา. \\ อพฺภานารเห นโย จาปิ, สมฺเภทํ นยโต ปุน. \\ ปาริวาสิเกสุ ตโย, จตุ มานตฺตจาริเก. \\ น สเมนฺติ รตฺติจฺเฉเทสุ มานตฺเตสุ จ เทวสิ. เทฺว กมฺมา สทิสา เสสา ตโย กมฺมา สมาสมาติ.}
\csromangathagroup{\csromangathastanza{ปาริวาสิกา สาเทนฺติ, ปกตตฺตาน ภิกฺขุนํ. \\ อภิวาทนํ ปจฺจุฏฺฐานํ, อญฺชลิํ จ สามีจิยํ.}\csromangathastanza{อาสนํ เสยฺยาภิหารํ, ปาโท ปีฐํ กถลิกํ. \\ ปตฺตํ นหาเน ปริกมฺมํ, อุชฺฌายนฺติ จ เปสลา.}\csromangathastanza{ทุกฺกฏํ สาทิยนฺตสฺส, มิถุ ปญฺจ ยถาวุฑฺฒํ\footnote{ปุนาปเร (ก)}. \\ อุโปสถํ ปวารณํ, วสฺสิโกโณชโภชนํ.}\csromangathastanza{สมฺมา จ วตฺตนา ตตฺถ, ปกตตฺตสฺส คจฺฉนฺตํ. \\ โย จ โหติ ปริยนฺโต, ปุเร ปจฺฉา ตเถว จ\footnote{น ปุเร ปจฺฉาสมเณน (สี, สฺยา)}.}\csromangathastanza{อารญฺญปิณฺฑนีหาโร, อาคนฺตุเก อุโปสเถ. \\ ปวารณาย ทูเตน, คนฺตพฺโพ จ สภิกฺขุโก.}\csromangathastanza{เอกจฺฉนฺเน จ วุฏฺฐานํ, ตเถว จ นิมนฺตเย. \\ อาสเน นีเจ จงฺกเม, ฉมายํ จงฺกเมน จ.}\csromangathastanza{วุฑฺฒตเรน อกมฺมํ, รตฺติจฺเฉทา จ โสธนา. \\ นิกฺขิปนํ สมาทานํ, วตฺตํว ปาริวาสิเก\footnote{รตฺติ วา ปาริวาสิเก (ก), ญาตพฺพํ ปาริวาสิกา (สี, สฺยา)}.}\csromangathastanza{\csromanfolio{103}มูลาย มานตฺตารหา, ตถา มานตฺตจาริตา. \\ อพฺภานารเห นโย จาปิ, สมฺเภทํ นยโต\footnote{สมฺเภทนยโต (สฺยา)} ปุน.}\csromangathastanza{ปาริวาสิเกสุ ตโย, จตุ มานตฺตจาริเก. \\ น สเมนฺติ รตฺติจฺเฉเทสุ\footnote{รตฺติจฺเฉเท (อิติปิ), รตฺติจฺเฉทา (สฺยา)} มานตฺเตสุ จ เทวสิ. เทฺว กมฺมา สทิสา เสสา ตโย กมฺมา สมาสมาติ\footnote{รตฺติจฺเฉเท (อิติปิ), รตฺติจฺเฉทา (สฺยา)}.}}

\nitthitam{\textbf{ปาริวาสิกกฺขนฺธโก นิฏฺฐิโต.}}
\csromanmatikaanchor{mtk.59}
\csromantocmarknopage{chapter}{3. สมุจฺจยกฺขนฺธก}{mtk.59}
\bookmark[dest={mtk.59},level=0]{3. สมุจฺจยกฺขนฺธก}
\chapterheadpageread{3. สมุจฺจยกฺขนฺธก}
\csromanmatikaanchor{mtk.60}
\csromantocmarkref{section}{1. สุกฺกวิสฺสฏฺฐิ}{mtk.60}
\bookmark[dest={mtk.60},level=1]{1. สุกฺกวิสฺสฏฺฐิ}
\csromanheader{1}{1. สุกฺกวิสฺสฏฺฐิ}
\proseitem{97}{\csromanfolio{104}เตน สมเยน พุทฺโธ ภควา สาวตฺถิยํ วิหรติ เชตวเน อนาถปิณฺฑิกสฺส อาราเม.\csromanspacer{} เตน โข ปน สมเยน อายสฺมา อุทายี เอกํ อาปตฺติํ อาปนฺโน โหติ สญฺเจตนิกํ สุกฺกวิสฺสฏฺฐิํ อปฺปฏิจฺฉนฺนํ, โส ภิกฺขูนํ อาโรเจสิ “อหํ อาวุโส เอกํ อาปตฺติํ อาปชฺชิํ สญฺเจตนิกํ สุกฺกวิสฺสฏฺฐิํ อปฺปฏิจฺฉนฺนํ, กถํ นุ โข มยา ปฏิปชฺชิตพฺพนฺ”ติ.\csromanspacer{} ภควโต\footnote{เต ภิกฺขู ภควโต (สฺยา)} เอตมตฺถํ อาโรเจสุํ.\csromanspacer{} เตน หิ ภิกฺขเว สํโฆ อุทายิสฺส ภิกฺขุโน เอกิสฺสา อาปตฺติยา สญฺเจตนิกาย สุกฺกวิสฺสฏฺฐิยา อปฺปฏิจฺฉนฺนาย ฉารตฺตํ มานตฺตํ เทตุ.\csromanspacer{} เอวญฺจ ปน ภิกฺขเว ทาตพฺพํ\csromandash{}}

\csromanmatikaanchor{mtk.61}
\csromantocmarkref{subsection}{อปฺปฏิจฺฉนฺนมานตฺต}{mtk.61}
\bookmark[dest={mtk.61},level=2]{อปฺปฏิจฺฉนฺนมานตฺต}
\csromanheader{2}{\textbf{อปฺปฏิจฺฉนฺนมานตฺต}}
\prose{98 เตน ภิกฺขเว อุทายินา ภิกฺขุนา สํฆํ อุปสงฺกมิตฺวา เอกํสํ อุตฺตราสงฺคํ กริตฺวา วุฑฺฒานํ ภิกฺขูนํ ปาเท วนฺทิตฺวา อุกฺกุฏิกํ นิสีทิตฺวา อญฺชลิํ ปคฺคเหตฺวา เอวมสฺส วจนีโย “อหํ ภนฺเต เอกํ อาปตฺติํ อาปชฺชิํ สญฺเจตนิกํ สุกฺกวิสฺสฏฺฐิํ อปฺปฏิจฺฉนฺนํ, โสหํ ภนฺเต สํฆํ เอกิสฺสา อาปตฺติยา สญฺเจตนิกาย สุกฺกวิสฺสฏฺฐิยา อปฺปฏิจฺฉนฺนาย ฉารตฺตํ มานตฺตํ ยาจามิ.\csromanspacer{} อหํ ภนฺเต เอกํ อาปตฺติํ อาปชฺชิํ สญฺเจตนิกํ สุกฺกวิสฺสฏฺฐิํ อปฺปฏิจฺฉนฺนํ, ทุติยมฺปิ โสหํ\footnote{ทุติยมฺปิ (สี, ก)} ภนฺเต สํฆํ เอกิสฺสา อาปตฺติยา สญฺเจตนิกาย สุกฺกวิสฺสฏฺฐิยา อปฺปฏิจฺฉนฺนาย ฉารตฺตํ มานตฺตํ ยาจามิ.\csromanspacer{} อหํ ภนฺเต เอกํ อาปตฺติํ อาปชฺชิํ สญฺเจตนิกํ สุกฺกวิสฺสฏฺฐิํ อปฺปฏิจฺฉนฺนํ, ตติยมฺปิ โสหํ\footnote{ตติยมฺปิ (สี, ก)} ภนฺเต สํฆํ เอกิสฺสา อาปตฺติยา สญฺเจตนิกาย สุกฺกวิสฺสฏฺฐิยา อปฺปฏิจฺฉนฺนาย ฉารตฺตํ มานตฺตํ ยาจามี”ติ.\csromanspacer{} พฺยตฺเตน ภิกฺขุนา ปฏิพเลน สํโฆ ญาเปตพฺโพ\csromandash{}}

\proseitem{99}{“สุณาตุ เม ภนฺเต สํโฆ, อยํ อุทายี ภิกฺขุ เอกํ อาปตฺติํ อาปชฺชิ สญฺเจตนิกํ สุกฺกวิสฺสฏฺฐิํ อปฺปฏิจฺฉนฺนํ, โส สํฆํ เอกิสฺสา อาปตฺติยา สญฺเจตนิกาย สุกฺกวิสฺสฏฺฐิยา อปฺปฏิจฺฉนฺนาย ฉารตฺตํ มานตฺตํ ยาจติ, ยทิ สํฆสฺส ปตฺตกลฺลํ, สํโฆ อุทายิสฺส ภิกฺขุโน \csromanfolio{105}เอกิสฺสา อาปตฺติยา สญฺเจตนิกาย สุกฺกวิสฺสฏฺฐิยา อปฺปฏิจฺฉนฺนาย ฉารตฺตํ มานตฺตํ ทเทยฺย, เอสา ญตฺติ.}

\prose{สุณาตุ เม ภนฺเต สํโฆ, อยํ อุทายี ภิกฺขุ เอกํ อาปตฺติํ อาปชฺชิ สญฺเจตนิกํ สุกฺกวิสฺสฏฺฐิํ อปฺปฏิจฺฉนฺนํ, โส สํฆํ เอกิสฺสา อาปตฺติยา สญฺเจตนิกาย สุกฺกวิสฺสฏฺฐิยา อปฺปฏิจฺฉนฺนาย ฉารตฺตํ มานตฺตํ ยาจติ, สํโฆ อุทายิสฺส ภิกฺขุโน เอกิสฺสา อาปตฺติยา สญฺเจตนิกาย สุกฺกวิสฺสฏฺฐิยา อปฺปฏิจฺฉนฺนาย ฉารตฺตํ มานตฺตํ เทติ, ยสฺสายสฺมโต ขมติ อุทายิสฺส ภิกฺขุโน เอกิสฺสา อาปตฺติยา สญฺเจตนิกาย สุกฺกวิสฺสฏฺฐิยา อปฺปฏิจฺฉนฺนาย ฉารตฺตํ มานตฺตสฺส ทานํ, โส ตุณฺหสฺส, ยสฺส นกฺขมติ, โส ภาเสยฺย.}

\prose{ทุติยมฺปิ เอตมตฺถํ วทามิ.\csromanspacer{} สุณาตุ เม ภนฺเต สํโฆ, อยํ อุทายี ภิกฺขุ เอกํ อาปตฺติํ อาปชฺชิ สญฺเจตนิกํ สุกฺกวิสฺสฏฺฐิํ อปฺปฏิจฺฉนฺนํ, โส สํฆํ เอกิสฺสา อาปตฺติยา สญฺเจตนิกาย สุกฺกวิสฺสฏฺฐิยา อปฺปฏิจฺฉนฺนาย ฉารตฺตํ มานตฺตํ ยาจติ, สํโฆ อุทายิสฺส ภิกฺขุโน เอกิสฺสา อาปตฺติยา สญฺเจตนิกาย สุกฺกวิสฺสฏฺฐิยา อปฺปฏิจฺฉนฺนาย ฉารตฺตํ มานตฺตํ เทติ, ยสฺสายสฺมโต ขมติ อุทายิสฺส ภิกฺขุโน เอกิสฺสา อาปตฺติยา สญฺเจตนิกาย สุกฺกวิสฺสฏฺฐิยา อปฺปฏิจฺฉนฺนาย ฉารตฺตํ มานตฺตสฺส ทานํ, โส ตุณฺหสฺส, ยสฺส นกฺขมติ, โส ภาเสยฺย.}

\prose{ตติยมฺปิ เอตมตฺถํ วทามิ.\csromanspacer{} สุณาตุ เม ภนฺเต สํโฆ, อยํ อุทายี ภิกฺขุ เอกํ อาปตฺติํ อาปชฺชิ สญฺเจตนิกํ สุกฺกวิสฺสฏฺฐิํ อปฺปฏิจฺฉนฺนํ, โส สํฆํ เอกิสฺสา อาปตฺติยา สญฺเจตนิกาย สุกฺกวิสฺสฏฺฐิยา อปฺปฏิจฺฉนฺนาย ฉารตฺตํ มานตฺตํ ยาจติ, สํโฆ อุทายิสฺส ภิกฺขุโน เอกิสฺสา อาปตฺติยา สญฺเจตนิกาย สุกฺกวิสฺสฏฺฐิยา อปฺปฏิจฺฉนฺนาย ฉารตฺตํ มานตฺตํ เทติ, ยสฺสายสฺมโต ขมติ อุทายิสฺส ภิกฺขุโน เอกิสฺสา อาปตฺติยา สญฺเจตนิกาย สุกฺกวิสฺสฏฺฐิยา อปฺปฏิจฺฉนฺนาย ฉารตฺตํ มานตฺตสฺส ทานํ, โส ตุณฺหสฺส, ยสฺส นกฺขมติ, โส ภาเสยฺย.}

\prose{ทินฺนํ สํเฆน อุทายิสฺส ภิกฺขุโน เอกิสฺสา อาปตฺติยา สญฺเจตนิกาย สุกฺกวิสฺสฏฺฐิยา อปฺปฏิจฺฉนฺนาย ฉารตฺตํ มานตฺตํ, ขมติ สํฆสฺส, ตสฺมา ตุณฺหี, เอวเมตํ ธารยามี”ติ.}

\csromanmatikaanchor{mtk.62}
\csromantocmarkref{subsection}{อปฺปฏิจฺฉนฺน-อพฺภาน}{mtk.62}
\bookmark[dest={mtk.62},level=2]{อปฺปฏิจฺฉนฺน-อพฺภาน}
\csromanheader{2}{\textbf{อปฺปฏิจฺฉนฺน-อพฺภาน}}
\proseitem{100}{\csromanfolio{106}โส จิณฺณมานตฺโต ภิกฺขูนํ อาโรเจสิ “อหํ อาวุโส เอกํ อาปตฺติํ อาปชฺชิ สญฺเจตนิกํ สุกฺกวิสฺสฏฺฐิํ อปฺปฏิจฺฉนฺนํ, โสหํ สํฆํ เอกิสฺสา อาปตฺติยา สญฺเจตนิกาย สุกฺกวิสฺสฏฺฐิยา อปฺปฏิจฺฉนฺนาย ฉารตฺตํ มานตฺตํ ยาจิํ, ตสฺส เม สํโฆ เอกิสฺสา อาปตฺติยา สญฺเจตนิกาย สุกฺกวิสฺสฏฺฐิยา อปฺปฏิจฺฉนฺนาย ฉารตฺตํ มานตฺตํ อทาสิ, โสหํ จิณฺณมานตฺโต, กถํ นุ โข มยา ปฏิปชฺชิตพฺพนฺ”ติ.\csromanspacer{} ภควโต เอตมตฺถํ อาโรเจสุํ.\csromanspacer{} เตน หิ ภิกฺขเว สํโฆ อุทายิํ ภิกฺขุํ อพฺเภตุ.\csromanspacer{} เอวญฺจ ปน ภิกฺขเว อพฺเภตพฺโพ\csromandash{}}

\prose{เตน ภิกฺขเว อุทายินา ภิกฺขุนา สํฆํ อุปสงฺกมิตฺวา เอกํสํ อุตฺตราสงฺคํ กริตฺวา วุฑฺฒานํ ภิกฺขูนํ ปาเท วนฺทิตฺวา อุกฺกุฏิกํ นิสีทิตฺวา อญฺชลิํ ปคฺคเหตฺวา เอวมสฺส วจนีโย “อหํ ภนฺเต เอกํ อาปตฺติํ อาปชฺชิ สญฺเจตนิกํ สุกฺกวิสฺสฏฺฐิํ อปฺปฏิจฺฉนฺนํ, โสหํ สํฆํ เอกิสฺสา อาปตฺติยา สญฺเจตนิกาย สุกฺกวิสฺสฏฺฐิยา อปฺปฏิจฺฉนฺนาย ฉารตฺตํ มานตฺตํ ยาจิํ, ตสฺส เม สํโฆ เอกิสฺสา อาปตฺติยา สญฺเจตนิกาย สุกฺกวิสฺสฏฺฐิยา อปฺปฏิจฺฉนฺนาย ฉารตฺตํ มานตฺตํ อทาสิ, โสหํ ภนฺเต จิณฺณมานตฺโต สํฆํ อพฺภานํ ยาจามิ.}

\prose{อหํ ภนฺเต เอกํ อาปตฺติํ อาปชฺชิํ สญฺเจตนิกํ สุกฺกวิสฺสฏฺฐิํ อปฺปฏิจฺฉนฺนํ โสหํ สํฆํ เอกิสฺสา อาปตฺติยา สญฺเจตนิกาย สุกฺกวิสฺสฏฺฐิยา อปฺปฏิจฺฉนฺนาย ฉารตฺตํ มานตฺตํ ยาจิํ, ตสฺส เม สํโฆ เอกิสฺสา อาปตฺติยา สญฺเจตนิกาย สุกฺกวิสฺสฏฺฐิยา อปฺปฏิจฺฉนฺนาย ฉารตฺตํ มานตฺตํ อทาสิ, โสหํ จิณฺณมานตฺโต ทุติยมฺปิ ภนฺเต สํฆํ อพฺภานํ ยาจามิ.}

\prose{อหํ ภนฺเต เอกํ อาปตฺติํ อาปชฺชิํ สญฺเจตนิกํ สุกฺกวิสฺสฏฺฐิํ อปฺปฏิจฺฉนฺนํ, โสหํ สํฆํ เอกิสฺสา อาปตฺติยา สญฺเจตนิกาย สุกฺกวิสฺสฏฺฐิยา อปฺปฏิจฺฉนฺนาย ฉารตฺตํ มานตฺตํ ยาจิํ, ตสฺส เม สํโฆ เอกิสฺสา อาปตฺติยา สญฺเจตนิกาย สุกฺกวิสฺสฏฺฐิยา อปฺปฏิจฺฉนฺนาย ฉารตฺตํ มานตฺตํ อทาสิ, โสหํ จิณฺณมานตฺโต ตติยมฺปิ ภนฺเต สํฆํ อพฺภานํ ยาจามี”ติ.\csromanspacer{} พฺยตฺเตน ภิกฺขุนา ปฏิพเลน สํโฆ ญาเปตพฺโพ\csromandash{}}

\proseitem{101}{“สุณาตุ เม ภนฺเต สํโฆ, อยํ อุทายี ภิกฺขุ เอกํ อาปตฺติํ อาปชฺชิ สญฺเจตนิกํ สุกฺกวิสฺสฏฺฐิํ อปฺปฏิจฺฉนฺนํ, โส สํฆํ เอกิสฺสา \csromanfolio{107}อาปตฺติยา สญฺเจตนิกาย สุกฺกวิสฺสฏฺฐิยา อปฺปฏิจฺฉนฺนาย ฉารตฺตํ มานตฺตํ ยาจิ, สํโฆ อุทายิสฺส ภิกฺขุโน เอกิสฺสา อาปตฺติยา สญฺเจตนิกายา สุกฺกวิสฺสฏฺฐิยา อปฺปฏิจฺฉนฺนาย ฉารตฺตํ มานตฺตํ อทาสิ, โส จิณฺณมานตฺโต สํฆํ อพฺภานํ ยาจติ, ยทิ สํฆสฺส ปตฺตกลฺลํ, สํโฆ อุทายิํ ภิกฺขุํ อพฺเภยฺย, เอสา ญตฺติ.}

\prose{สุณาตุ เม ภนฺเต สํโฆ, อยํ อุทายี ภิกฺขุ เอกํ อาปตฺติํ อาปชฺชิ สญฺเจตนิกํ สุกฺกวิสฺสฏฺฐิํ อปฺปฏิจฺฉนฺนํ, โส สํฆํ เอกิสฺสา อาปตฺติยา สญฺเจตนิกาย สุกฺกวิสฺสฏฺฐิยา อปฺปฏิจฺฉนฺนาย ฉารตฺตํ มานตฺตํ ยาจิ, สํโฆ อุทายิสฺส ภิกฺขุโน เอกิสฺสา อาปตฺติยา สญฺเจตนิกาย สุกฺกวิสฺสฏฺฐิยา อปฺปฏิจฺฉนฺนาย ฉารตฺตํ มานตฺตํ อทาสิ, โส จิณฺณมานตฺโต สํฆํ อพฺภานํ ยาจติ, สํโฆ อุทายิํ ภิกฺขุํ อพฺเภติ, ยสฺสายสฺมโต ขมติ อุทายิสฺส ภิกฺขุโน อพฺภานํ, โส ตุณฺหสฺส, ยสฺส นกฺขมติ, โส ภาเสยฺย.}

\prose{ทุติยมฺปิ เอตมตฺถํ วทามิ.\csromanspacer{} สุณาตุ เม ภนฺเต สํโฆ, อยํ อุทายี ภิกฺขุ เอกํ อาปตฺติํ อาปชฺชิ สญฺเจตนิกํ สุกฺกวิสฺสฏฺฐิํ อปฺปฏิจฺฉนฺนํ, โส สํฆํ เอกิสฺสา อาปตฺติยา สญฺเจตนิกาย สุกฺกวิสฺสฏฺฐิยา อปฺปฏิจฺฉนฺนาย ฉารตฺตํ มานตฺตํ ยาจิ, สํโฆ อุทายิสฺส ภิกฺขุโน เอกิสฺสา อาปตฺติยา สญฺเจตนิกาย สุกฺกวิสฺสฏฺฐิยา อปฺปฏิจฺฉนฺนาย ฉารตฺตํ มานตฺตํ อทาสิ, โส จิณฺณมานตฺโต สํฆํ อพฺภานํ ยาจติ, สํโฆ อุทายิํ ภิกฺขุํ อพฺเภติ, ยสฺสายสฺมโต ขมติ อุทายิสฺส ภิกฺขุโน อพฺภานํ, โส ตุณฺหสฺส, ยสฺส นกฺขมติ, โส ภาเสยฺย.}

\prose{ตติยมฺปิ เอตมตฺถํ วทามิ.\csromanspacer{} สุณาตุ เม ภนฺเต สํโฆ, อยํ อุทายี ภิกฺขุ เอกํ อาปตฺติํ อาปชฺชิ สญฺเจตนิกํ สุกฺกวิสฺสฏฺฐิํ อปฺปฏิจฺฉนฺนํ, โส สํฆํ เอกิสฺสา อาปตฺติยา สญฺเจตนิกาย สุกฺกวิสฺสฏฺฐิยา อปฺปฏิจฺฉนฺนาย ฉารตฺตํ มานตฺตํ ยาจิ, สํโฆ อุทายิสฺส ภิกฺขุโน เอกิสฺสา อาปตฺติยา สญฺเจตนิกาย สุกฺกวิสฺสฏฺฐิยา อปฺปฏิจฺฉนฺนาย ฉารตฺตํ มานตฺตํ อทาสิ, โส จิณฺณมานตฺโต สํฆํ อพฺภานํ ยาจติ, สํโฆ อุทายิํ ภิกฺขุํ อพฺเภติ, ยสฺสายสฺมโต ขมติ อุทายิสฺส ภิกฺขุโน อพฺภานํ, โส ตุณฺหสฺส, ยสฺส นกฺขมติ, โส ภาเสยฺย.}

\prose{\csromanfolio{108}อพฺภิโต สํเฆน อุทายี ภิกฺขุ, ขมติ สํฆสฺส, ตสฺมา ตุณฺหี เอวเมตํ ธารยามี”ติ.}

\csromanmatikaanchor{mtk.63}
\csromantocmarkref{subsection}{เอกาหปฺปฏิจฺฉนฺนปริวาส}{mtk.63}
\bookmark[dest={mtk.63},level=2]{เอกาหปฺปฏิจฺฉนฺนปริวาส}
\csromanheader{2}{\textbf{เอกาหปฺปฏิจฺฉนฺนปริวาส}}
\proseitem{102}{เตน โข ปน สมเยน อายสฺมา อุทายี เอกํ อาปตฺติํ อาปนฺโน โหติ สญฺเจตนิกํ สุกฺกวิสฺสฏฺฐิํ เอกาหปฺปฏิจฺฉนฺนํ, โส ภิกฺขูนํ อาโรเจสิ “อหํ อาวุโส เอกํ อาปตฺติํ อาปชฺชิํ สญฺเจตนิกํ สุกฺกวิสฺสฏฺฐิํ เอกาหปฺปฏิจฺฉนฺนํ, กถํ นุ โข มยา ปฏิปชฺชิตพฺพนฺ”ติ.\csromanspacer{} ภควโต เอตมตฺถํ อาโรเจสุํ.\csromanspacer{} เตน หิ ภิกฺขเว สํโฆ อุทายิสฺส ภิกฺขุโน เอกิสฺสา อาปตฺติยา สญฺเจตนิกาย สุกฺกวิสฺสฏฺฐิยา เอกาหปฺปฏิจฺฉนฺนาย เอกาหปริวาสํ เทตุ.\csromanspacer{} เอวญฺจ ปน ภิกฺขเว ทาตพฺโพ\csromandash{}}

\prose{เตน ภิกฺขเว อุทายินา ภิกฺขุนา สํฆํ อุปสงฺกมิตฺวา เอกํสํ อุตฺตราสงฺคํ กริตฺวา วุฑฺฒานํ ภิกฺขูนํ ปาเท วนฺทิตฺวา อุกฺกุฏิกํ นิสีทิตฺวา อญฺชลิํ ปคฺคเหตฺวา เอวมสฺส วจนีโย “อหํ ภนฺเต เอกํ อาปตฺติํ อาปชฺชิํ สญฺเจตนิกํ สุกฺกวิสฺสฏฺฐิํ เอกาหปฺปฏิจฺฉนฺนํ, โสหํ ภนฺเต สํฆํ เอกิสฺสา อาปตฺติยา สญฺเจตนิกาย สุกฺกวิสฺสฏฺฐิยา เอกาหปฺปฏิจฺฉนฺนาย เอกาหปริวาสํ ยาจามี”ติ.\csromanspacer{} ทุติยมฺปิ ยาจิตพฺโพ.\csromanspacer{} ตติยมฺปิ ยาจิตพฺโพ.\csromanspacer{} พฺยตฺเตน ภิกฺขุนา ปฏิพเลน สํโฆ ญาเปตพฺโพ\csromandash{}}

\proseitem{103}{“สุณาตุ เม ภนฺเต สํโฆ, อยํ อุทายี ภิกฺขุ เอกํ อาปตฺติํ อาปชฺชิ สญฺเจตนิกํ สุกฺกวิสฺสฏฺฐิํ เอกาหปฺปฏิจฺฉนฺนํ, โส สํฆํ เอกิสฺสา อาปตฺติยา สญฺเจตนิกาย สุกฺกวิสฺสฏฺฐิยา เอกาหปฺปฏิจฺฉนฺนาย เอกาหปริวาสํ ยาจติ, ยทิ สํฆสฺส ปตฺตกลฺลํ, สํโฆ อุทายิสฺส ภิกฺขุโน เอกิสฺสา อาปตฺติยา สญฺเจตนิกาย สุกฺกวิสฺสฏฺฐิยา เอกาหปฺปฏิจฺฉนฺนาย เอกาหปริวาสํ ทเทยฺย, เอสา ญตฺติ.}

\prose{สุณาตุ เม ภนฺเต สํโฆ, อยํ อุทายี ภิกฺขุ เอกํ อาปตฺติํ อาปชฺชิ สญฺเจตนิกํ สุกฺกวิสฺสฏฺฐิํ เอกาหปฺปฏิจฺฉนฺนํ, โส สํฆํ เอกิสฺสา อาปตฺติยา สญฺเจตนิกาย สุกฺกวิสฺสฏฺฐิยา เอกาหปฺปฏิจฺฉนฺนาย เอกาหปริวาสํ ยาจติ, สํโฆ อุทายิสฺส ภิกฺขุโน เอกิสฺสา อาปตฺติยา สญฺเจตนิกาย สุกฺกวิสฺสฏฺฐิยา เอกาหปฺปฏิจฺฉนฺนาย เอกาหปริวาสํ \csromanfolio{109}เทหิ, ยสฺสายสฺมโต ขมติ อุทายิสฺส ภิกฺขุโน เอกิสฺสา อาปตฺติยา สญฺเจตนิกาย สุกฺกวิสฺสฏฺฐิยา เอกาหปฺปฏิจฺฉนฺนาย เอกาหปริวาสสฺส ทานํ, โส ตุณฺหสฺส, ยสฺส นกฺขมติ, โส ภาเสยฺย.}

\prose{ทุติยมฺปิ เอตมตฺถํ วทามิ ฯเปฯ.\csromanspacer{} ตติยมฺปิ เอตมตฺถํ วทามิ ฯเปฯ.}

\prose{ทินฺโน สํเฆน อุทายิสฺส ภิกฺขุโน เอกิสฺสา อาปตฺติยา สญฺเจตนิกาย สุกฺกวิสฺสฏฺฐิยา เอกาหปฺปฏิจฺฉนฺนาย เอกาหปริวาโส, ขมติ สํฆสฺส, ตสฺมา ตุณฺหี, เอวเมตํ ธารยามี”ติ.}

\csromanmatikaanchor{mtk.64}
\csromantocmarkref{subsection}{เอกาหปฺปฏิจฺฉนฺนมานตฺต}{mtk.64}
\bookmark[dest={mtk.64},level=2]{เอกาหปฺปฏิจฺฉนฺนมานตฺต}
\csromanheader{2}{\textbf{เอกาหปฺปฏิจฺฉนฺนมานตฺต}}
\proseitem{104}{โส ปริวุตฺถปริวาโส ภิกฺขูนํ อาโรเจสิ “อหํ อาวุโส เอกํ อาปตฺติํ อาปชฺชิํ สญฺเจตนิกํ สุกฺกวิสฺสฏฺฐิํ เอกาหปฺปฏิจฺฉนฺนํ, โสหํ สํฆํ เอกิสฺสา อาปตฺติยา สญฺเจตนิกาย สุกฺกวิสฺสฏฺฐิยา เอกาหปฺปฏิจฺฉนฺนาย เอกาหปริวาสํ ยาจิํ, ตสฺส เม สํโฆ เอกิสฺสา อาปตฺติยา สญฺเจตนิกาย สุกฺกวิสฺสฏฺฐิยา เอกาหปฺปฏิจฺฉนฺนาย เอกาหปริวาสํ อทาสิ, โสหํ ปริวุตฺถปริวาโส, กถํ นุ โข มยา ปฏิปชฺชิตพฺพนฺ”ติ.\csromanspacer{} ภควโต เอตมตฺถํ อาโรเจสุํ.\csromanspacer{} เตน หิ ภิกฺขเว สํโฆ อุทายิสฺส ภิกฺขุโน เอกิสฺสา อาปตฺติยา สญฺเจตนิกาย สุกฺกวิสฺสฏฺฐิยา เอกาหปฺปฏิจฺฉนฺนาย ฉารตฺตํ มานตฺตํ เทตุ.\csromanspacer{} เอวญฺจ ปน ภิกฺขเว ทาตพฺพํ\csromandash{}}

\prose{เตน ภิกฺขเว อุทายินา ภิกฺขุนา สํฆํ อุปสงฺกมิตฺวา ฯเปฯ เอวมสฺส วจนีโย “อหํ ภนฺเต เอกํ อาปตฺติํ อาปชฺชิํ สญฺเจตนิกํ สุกฺกวิสฺสฏฺฐิํ เอกาหปฺปฏิจฺฉนฺนํ, โสหํ สํฆํ เอกิสฺสา อาปตฺติยา สญฺเจตนิกาย สุกฺกวิสฺสฏฺฐิยา เอกาหปฺปฏิจฺฉนฺนาย เอกาหปริวาสํ ยาจิํ, ตสฺส เม สํโฆ เอกิสฺสา อาปตฺติยา สญฺเจตนิกาย สุกฺกวิสฺสฏฺฐิยา เอกาหปฺปฏิจฺฉนฺนาย เอกาหปริวาสํ อทาสิ, โสหํ ภนฺเต ปริวุตฺถปริวาโส สํฆํ เอกิสฺสา อาปตฺติยา สญฺเจตนิกาย สุกฺกวิสฺสฏฺฐิยา เอกาหปฺปฏิจฺฉนฺนาย ฉารตฺตํ มานตฺตํ ยาจามี”ติ.\csromanspacer{} ทุติยมฺปิ ยาจิตพฺพํ\footnote{ยาจิตพฺโพ (สี, เอวมุปริปิ)}.\csromanspacer{} ตติยมฺปิ ยาจิตพฺพํ1.\csromanspacer{} พฺยตฺเตน ภิกฺขุนา ปฏิพเลน สํโฆ ญาเปตพฺโพ\csromandash{}}

\proseitem{105}{\csromanfolio{110}“สุณาตุ เม ภนฺเต สํโฆ, อยํ อุทายี ภิกฺขุ เอกํ อาปตฺติํ อาปชฺชิ สญฺเจตนิกํ สุกฺกวิสฺสฏฺฐิํ เอกาหปฺปฏิจฺฉนฺนํ, โส สํฆํ เอกิสฺสา อาปตฺติยา สญฺเจตนิกายา สุกฺกวิสฺสฏฺฐิยา เอกาหปฺปฏิจฺฉนฺนาย เอกาหปริวาสํ ยาจิ, สํโฆ อุทายิสฺส ภิกฺขุโน เอกิสฺสา อาปตฺติยา สญฺเจตนิกาย สุกฺกวิสฺสฏฺฐิยา เอกาหปฺปฏิจฺฉนฺนาย เอกาหปริวาสํ อทาสิ, โส ปริวุตฺถปริวาโส สํฆํ เอกิสฺสา อาปตฺติยา สญฺเจตนิกาย สุกฺกวิสฺสฏฺฐิยา เอกาหปฺปฏิจฺฉนฺนาย ฉารตฺตํ มานตฺตํ ยาจติ, ยทิ สํฆสฺส ปตฺตกลฺลํ, สํโฆ อุทายิสฺส ภิกฺขุโน เอกิสฺสา อาปตฺติยา สญฺเจตนิกาย สุกฺกวิสฺสฏฺฐิยา เอกาหปฺปฏิจฺฉนฺนาย ฉารตฺตํ มานตฺตํ ทเทยฺย, เอสา ญตฺติ.}

\prose{สุณาตุ เม ภนฺเต สํโฆ, อยํ อุทายี ภิกฺขุ เอกํ อาปตฺติํ อาปชฺชิ สญฺเจตนิกํ สุกฺกวิสฺสฏฺฐิํ เอกาหปฺปฏิจฺฉนฺนํ, โส สํฆํ เอกิสฺสา อาปตฺติยา สญฺเจตนิกาย สุกฺกวิสฺสฏฺฐิยา เอกาหปฺปฏิจฺฉนฺนาย เอกาหปริวาสํ ยาจิ, สํโฆ อุทายิสฺส ภิกฺขุโน เอกิสฺสา อาปตฺติยา สญฺเจตนิกาย สุกฺกวิสฺสฏฺฐิยา เอกาหปฺปฏิจฺฉนฺนาย เอกาหปริวาสํ อทาสิ, โส ปริวุตฺถปริวาโส สํฆํ เอกิสฺสา อาปตฺติยา สญฺเจตนิกาย สุกฺกวิสฺสฏฺฐิยา เอกาหปฺปฏิจฺฉนฺนาย ฉารตฺตํ มานตฺตํ ยาจติ, สํโฆ อุทายิสฺส ภิกฺขุโน เอกิสฺสา อาปตฺติยา สญฺเจตนิกาย สุกฺกวิสฺสฏฺฐิยา เอกาหปฺปฏิจฺฉนฺนาย ฉารตฺตํ มานตฺตํ เทติ, ยสฺสายสฺมโต ขมติ อุทายิสฺส ภิกฺขุโน เอกิสฺสา อาปตฺติยา สญฺเจตนิกาย สุกฺกวิสฺสฏฺฐิยา เอกาหปฺปฏิจฺฉนฺนาย ฉารตฺต มานตฺตสฺส ทานํ, โส ตุณฺหสฺส, ยสฺส นกฺขมติ, โส ภาเสยฺย.}

\prose{ทุติยมฺปิ เอตมตฺตํ วทามิ ฯเปฯ.\csromanspacer{} ตติยมฺปิ เอตมตฺตํ วทามิ ฯเปฯ.}

\prose{ทินฺนํ สํเฆน อุทายิสฺส ภิกฺขุโน เอกิสฺสา อาปตฺติยา สญฺเจตนิกาย สุกฺกวิสฺสฏฺฐิยา เอกาหปฺปฏิจฺฉนฺนาย ฉารตฺตํ มานตฺตํ, ขมติ สํฆสฺส, ตสฺมา ตุณฺหี, เอวเมตํ ธารยามี”ติ.}

\csromanmatikaanchor{mtk.65}
\csromantocmarkref{subsection}{เอกาหปฺปฏิจฺฉนฺน-อพฺภาน}{mtk.65}
\bookmark[dest={mtk.65},level=2]{เอกาหปฺปฏิจฺฉนฺน-อพฺภาน}
\csromanheader{2}{\textbf{เอกาหปฺปฏิจฺฉนฺน-อพฺภาน}}
\proseitem{106}{โส จิณฺณมานตฺโต ภิกฺขูนํ อาโรเจสิ “อหํ อาวุโส เอกํ อาปตฺติํ อาปชฺชิํ สญฺเจตนิกํ สุกฺกวิสฺสฏฺฐิํ เอกาหปฺปฏิจฺฉนํ, \csromanfolio{111}โสหํ สํฆํ เอกิสฺสา อาปตฺติยา สญฺเจตนิกาย สุกฺกวิสฺสฏฺฐิยา เอกาหปฺปฏิจฺฉนฺนาย เอกาหปริวาสํ ยาจิํ, ตสฺส เม สํโฆ เอกิสฺสา อาปตฺติยา สญฺเจตนิกาย สุกฺกวิสฺสฏฺฐิยา เอกาหปฺปฏิจฺฉนฺนาย เอกาหปริวาสํ อทาสิ, โสหํ ปริวุตฺถปริวาโส สํฆํ เอกิสฺสา อาปตฺติยา สญฺเจตนิกาย สุกฺกวิสฺสฏฺฐิยา เอกาหปฺปฏิจฺฉนฺนาย ฉารตฺตํ มานตฺตํ ยาจิํ, ตสฺส เม สํโฆ เอกิสฺสา อาปตฺติยา สญฺเจตนิกาย สุกฺกวิสฺสฏฺฐิยา เอกาหปฺปฏิจฺฉนฺนาย ฉารตฺตํ มานตฺตํ อทาสิ, โสหํ จิณฺณมานตฺโต, กถํ นุ โข มยา ปฏิปชฺชิตพฺพนฺ”ติ.\csromanspacer{} ภควโต เอตมตฺถํ อาโรเจสุํ.\csromanspacer{} เตน หิ ภิกฺขเว สํโฆ อุทายิํ ภิกฺขุํ อพฺเภตุ.\csromanspacer{} เอวญฺจ ปน ภิกฺขเว อพฺเภตพฺโพ\csromandash{}}

\prose{เตน ภิกฺขเว อุทายินา ภิกฺขุนา สํฆํ อุปสงฺกมิตฺวา ฯเปฯ เอวมสฺส วจนีโย “อหํ ภนฺเต เอกํ อาปตฺติํ อาปชฺชิํ สญฺเจตนิกํ สุกฺกวิสฺสฏฺฐิํ เอกาหปฺปฏิจฺฉนฺนํ, โสหํ สํฆํ เอกิสฺสา อาปตฺติยา สญฺเจตนิกาย สุกฺกวิสฺสฏฺฐิยา เอกาหปฺปฏิจฺฉนฺนาย เอกาหปริวาสํ ยาจิํ, ตสฺส เม สํโฆ เอกิสฺสา อาปตฺติยา สญฺเจตนิกาย สุกฺกวิสฺสฏฺฐิยา เอกาหปฺปฏิจฺฉนฺนาย เอกาหปริวาสํ อทาสิ, โสหํ ปริวุตฺถปริวาโส สํฆํ เอกิสฺสา อาปตฺติยา สญฺเจตนิกาย สุกฺกวิสฺสฏฺฐิยา เอกาหปฺปฏิจฺฉนฺนาย ฉารตฺตํ มานตฺตํ ยาจิํ, ตสฺส เม สํโฆ เอกิสฺสา อาปตฺติยา สญฺเจตนิกาย สุกฺกวิสฺสฏฺฐิยา เอกาหปฺปฏิจฺฉนฺนาย ฉารตฺตํ มานตฺตํ อทาสิ, โสหํ ภนฺเต จิณฺณมานตฺโต สํฆํ อพฺภานํ ยาจามี”ติ.\csromanspacer{} ทุติยมฺปิ ยาจิตพฺพํ.\csromanspacer{} ตติยมฺปิ ยาจิตพฺพํ.\csromanspacer{} พฺยตฺเตน ภิกฺขุนา ปฏิพเลน สํโฆ ญาเปตพฺโพ\csromandash{}}

\proseitem{107}{“สุณาตุ เม ภนฺเต สํโฆ, อยํ อุทายี ภิกฺขุ เอกํ อาปตฺติํ อาปชฺชิ สญฺเจตนิกํ สุกฺกวิสฺสฏฺฐิ เอกาหปฺปฏิจฺฉนฺนํ, โส สํฆํ เอกิสฺสา อาปตฺติยา สญฺเจตนิกาย สุกฺกวิสฺสฏฺฐิยา เอกาหปฺปฏิจฺฉนฺนาย เอกาหปริวาสํ ยาจิ, สํโฆ อุทายิสฺส ภิกฺขุโน เอกิสฺสา อาปตฺติยา สญฺเจตนิกาย สุกฺกวิสฺสฏฺฐิยา เอกาหปฺปฏิจฺฉนฺนาย เอกาหปริวาสํ อทาสิ, โส ปริวุตฺถปริวาโส สํฆํ เอกิสฺสา อาปตฺติยา สญฺเจตนิกาย สุกฺกวิสฺสฏฺฐิยา เอกาหปฺปฏิจฺฉนฺนาย ฉารตฺตํ มานตฺตํ ยาจิ, สํโฆ อุทายิสฺส ภิกฺขุโน เอกิสฺสา อาปตฺติยา สญฺเจตนิกาย สุกฺกวิสฺสฏฺฐิยา เอกาหปฺปฏิจฺฉนฺนาย ฉารตฺตํ มานตฺตํ อทาสิ, โส \csromanfolio{112}จิณฺณมานตฺโต สํฆํ อพฺภานํ ยาจติ, ยทิ สํฆสฺส ปตฺตกลฺลํ, สํโฆ อุทายิํ ภิกฺขุํ อพฺเภยฺย, เอสา ญตฺติ.}

\prose{สุณาตุ เม ภนฺเต สํโฆ, อยํ อุทายี ภิกฺขุ เอกํ อาปตฺติํ อาปชฺชิ สญฺเจตนิกํ สุกฺกวิสฺสฏฺฐิํ เอกาหปฺปฏิจฺฉนฺนํ, โส สํฆํ เอกิสฺสา อาปตฺติยา สญฺเจตนิกาย สุกฺกวิสฺสฏฺฐิยา เอกาหปฺปฏิจฺฉนฺนาย เอกาหปริวาสํ ยาจิ, สํโฆ อุทายิสฺส ภิกฺขุโน เอกิสฺสา อาปตฺติยา สญฺเจตนิกาย สุกฺกวิสฺสฏฺฐิยา เอกาหปฺปฏิจฺฉนฺนาย เอกาหปริวาสํ อทาสิ, โส ปริวุตฺถปริวาโส สํฆํ เอกิสฺสา อาปตฺติยา สญฺเจตนิกาย สุกฺกวิสฺสฏฺฐิยา เอกาหปฺปฏิจฺฉนฺนาย ฉารตฺตํ มานตฺตํ ยาจิ, สํโฆ อุทายิสฺส ภิกฺขุโน เอติสฺสา อาปตฺติยา สญฺเจตนิกาย สุกฺกวิสฺสฏฺฐิยา เอกาหปฺปฏิจฺฉนฺนาย ฉารตฺตํ มานตฺตํ อทาสิ, โส จิณฺณมานตฺโต สํฆํ อพฺภานํ ยาจติ สํโฆ อุทายิํ ภิกฺขุํ อพฺเภติ, ยสฺสายสฺมโต ขมติ อุทายิสฺส ภิกฺขุโน อพฺภานํ, โส ตุณฺหสฺส, ยสฺส นกฺขมติ, โส ภาเสยฺย.}

\prose{ทุติยมฺปิ เอตมตฺถํ วทามิ ฯเปฯ.\csromanspacer{} ตติยมฺปิ เอตมตฺถํ วทามิ ฯเปฯ.}

\prose{อพฺภิโต สํเฆน อุทายี ภิกฺขุ, ขมติ สํฆสฺส, ตสฺมา ตุณฺหี เอวเมตํ ธารยามี”ติ.}

\csromanmatikaanchor{mtk.66}
\csromantocmarkref{subsection}{ปญฺจาหปฺปฏิจฺฉนฺนปริวาส}{mtk.66}
\bookmark[dest={mtk.66},level=2]{ปญฺจาหปฺปฏิจฺฉนฺนปริวาส}
\csromanheader{2}{\textbf{ปญฺจาหปฺปฏิจฺฉนฺนปริวาส}}
\proseitem{108}{เตน โข ปน สมเยน อายสฺมา อุทายี เอกํ อาปตฺติํ อาปนฺโน โหติ สญฺเจตนิกํ สุกฺกวิสฺสฏฺฐิํ ทฺวีหปฺปฏิจฺฉนฺนํ ฯเปฯ ตีหปฺปฏิจฺฉนฺนํ ฯเปฯ จตูหปฺปฏิจฺฉนฺนํ ฯเปฯ ปญฺจาหปฺปฏิจฺฉนฺนํ.\csromanspacer{} โส ภิกฺขูนํ อาโรเจสิ “อหํ อาวุโส เอกํ อาปตฺติํ อาปชฺชิํ สญฺเจตนิกํ สุกฺกวิสฺสฏฺฐิํ ปญฺจาหปฺปฏิจฺฉนฺนํ, กถํ นุ โข มยา ปฏิปชฺชิตพฺพนฺ”ติ.\csromanspacer{} ภควโต เอตมตฺถํ อาโรเจสุํ.\csromanspacer{} เตน หิ ภิกฺขเว สํโฆ อุทายิสฺส ภิกฺขุโน เอกิสฺสา อาปตฺติยา สญฺเจตนิกาย สุกฺกวิสฺสฏฺฐิยา ปญฺจาหปฺปฏิจฺฉนฺนาย ปญฺจาหปริวาสํ เทตุ.\csromanspacer{} เอวญฺจ ปน ภิกฺขเว ทาตพฺโพ\csromandash{}}

\csromanmatikaanchor{mtk.67}
\csromantocmarkref{subsection}{ปาริวาสิกมูลายปฏิกสฺสนา}{mtk.67}
\bookmark[dest={mtk.67},level=2]{ปาริวาสิกมูลายปฏิกสฺสนา}
\prose{เตน ภิกฺขเว อุทายินา ภิกฺขุนา สํฆํ อุปสงฺกมิตฺวา เอกํสํ อุตฺตราสงฺคํ กริตฺวา วุฑฺฒานํ ภิกฺขูนํ ปาเท วนฺทิตฺวา อุกฺกุฏิกํ นิสีทิตฺวา อญฺชลิํ ปคฺคเหตฺวา เอวมสฺส วจนีโย “อหํ ภนฺเต เอกํ อาปตฺติํ \csromanfolio{113}อาปชฺชิํ สญฺเจตนิกํ สุกฺกวิสฺสฏฺฐิํ ปญฺจาหปฺปฏิจฺฉนฺนํ, โสหํ ภนฺเต สํฆํ เอกิสฺสา อาปตฺติยา สญฺเจตนิกาย สุกฺกวิสฺสฏฺฐิยา ปญฺจาหปฺปฏิจฺฉนฺนาย ปญฺจาหปริวาสํ ยาจามี”ติ.\csromanspacer{} ทุติยมฺปิ ยาจิตพฺโพ.\csromanspacer{} ตติยมฺปิ ยาจิตพฺโพ.\csromanspacer{} พฺยตฺเตน ภิกฺขุนา ปฏิพเลน สํโฆ ญาเปตพฺโพ\csromandash{}}

\proseitem{109}{“สุณาตุ เม ภนฺเต สํโฆ, อยํ อุทายี ภิกฺขุ เอกํ อาปตฺติํ อาปชฺชิ สญฺเจตนิกํ สุกฺกวิสฺสฏฺฐิํ ปญฺจาหปฺปฏิจฺฉนฺนํ, โส สํฆํ เอกิสฺสา อาปตฺติยา สญฺเจตนิกาย สุกฺกวิสฺสฏฺฐิยา ปญฺจาหปฺปฏิจฺฉนฺนาย ปญฺจาหปริวาสํ ยาจติ, ยทิสํฆสฺส ปตฺตกลฺลํ, สํโฆ อุทายิสฺส ภิกฺขุโน เอกิสฺสา อาปตฺติยา สญฺเจตนิกาย สุกฺกวิสฺสฏฺฐิยา ปญฺจาหปฺปฏิจฺฉนฺนาย ปญฺจาหปริวาสํ ทเทยฺย, เอสา ญตฺติ.}

\prose{สุณาตุ เม ภนฺเต สํโฆ, อยํ อุทายี ภิกฺขุ เอกํ อาปตฺติํ อาปชฺชิ สญฺเจตนิกํ สุกฺกวิสฺสฏฺฐิํ ปญฺจาหปฺปฏิจฺฉนฺนํ, โส สํฆํ เอกิสฺสา อาปตฺติยา สญฺเจตนิกาย สุกฺกวิสฺสฏฺฐิยา ปญฺจาหปฺปฏิจฺฉนฺนาย ปญฺจาหปริวาสํ ยาจติ, สํโฆ อุทายิสฺส ภิกฺขุโน เอกิสฺสา อาปตฺติยา สญฺเจตนิกาย สุกฺกวิสฺสฏฺฐิยา ปญฺจาหปฺปฏิจฺฉนฺนาย ปญฺจาหปริวาสํ เทติ, ยสฺสายสฺมโต ขมติ อุทายิสฺส ภิกฺขุโน เอกิสฺสา อาปตฺติยา สญฺเจตนิกาย สุกฺกวิสฺสฏฺฐิยา ปญฺจาหปฺปฏิจฺฉนฺนาย ปญฺจาหปริวาสสฺส ทานํ, โส ตุณฺหสฺส, ยสฺส นกฺขมติ, โส ภาเสยฺย.}

\prose{ทุติยมฺปิ เอตมตฺถํ วทามิ ฯเปฯ.\csromanspacer{} ตติยมฺปิ เอตมตฺถํ วทามิ ฯเปฯ.}

\prose{ทินฺโน สํเฆน อุทายิสฺส ภิกฺขุโน เอกิสฺสา อาปตฺติยา สญฺเจตนิกาย สุกฺกวิสฺสฏฺฐิยา ปญฺจาหปฺปฏิจฺฉนฺนาย ปญฺจาหปริวาโส, ขมติ สํฆสฺส, ตสฺมา ตุณฺหี, เอวเมตํ ธารยามี”ติ.}

\csromanheader{2}{\textbf{ปริวาสิกมูลายปฏิกสฺสนา}}
\proseitem{110}{โส ปริวสนฺโต อนฺตรา เอกํ อาปตฺติํ อาปชฺชิ สญฺเจตนิกํ สุกฺกวิสฺสฏฺฐิํ อปฺปฏิจฺฉนฺนํ, โส ภิกฺขูนํ อาโรเจสิ “อหํ อาวุโส เอกํ อาปตฺติํ อาปชฺชิํ สญฺเจตนิกํ สุกฺกวิสฺสฏฺฐิํ ปญฺจาหปฺปฏิจฺฉนฺนํ, โสหํ สํฆํ เอกิสฺสา อาปตฺติยา สญฺเจตนิกาย สุกฺกวิสฺสฏฺฐิยา ปญฺจาหปฺปฏิจฺฉนฺนาย ปญฺจาหปริวาสํ ยาจิํ, ตสฺส เม สํโฆ เอกิสฺสา \csromanfolio{114}อาปตฺติยา สญฺเจตนิกาย สุกฺกวิสฺสฏฺฐิยา ปญฺจาหปฺปฏิจฺฉนฺนาย ปญฺจาหปริวาสํ อทาสิ, โสหํ ปริวสนฺโต อนฺตรา เอกํ อาปตฺติํ อาปชฺชิํ สญฺเจตนิกํ สุกฺกวิสฺสฏฺฐิํ อปฺปฏิจฺฉนฺนํ, กถํ นุ โข มยา ปฏิปชฺชิตพฺพนฺ”ติ.\csromanspacer{} ภควโต เอตมตฺถํ อาโรเจสุํ.\csromanspacer{} เตน หิ ภิกฺขเว สํโฆ อุทายิํ ภิกฺขุํ อนฺตรา เอกิสฺสา อาปตฺติยา สญฺเจตนิกาย สุกฺกวิสฺสฏฺฐิยา อปฺปฏิจฺฉนฺนาย มูลาย ปฏิกสฺสตุ.\csromanspacer{} เอวญฺจ ปน ภิกฺขเว มูลาย ปฏิกสฺสิตพฺโพ\csromandash{}}

\prose{เตน ภิกฺขเว อุทายินา ภิกฺขุนา สํฆํ อุปสงฺกมิตฺวา ฯเปฯ เอวมสฺส วจนีโย “อหํ ภนฺเต เอกํ อาปตฺติํ อาปชฺชิํ สญฺเจตนิกํ สุกฺกวิสฺสฏฺฐิํ ปญฺจาหปฺปฏิจฺฉนฺนํ, โสหํ สํฆํ เอกิสฺสา อาปตฺติยา สญฺเจตนิกาย สุกฺกวิสฺสฏฺฐิยา ปญฺจาหปฺปฏิจฺฉนฺนาย ปญฺจาหปริวาสํ ยาจิํ, ตสฺส เม สํโฆ เอกิสฺสา อาปตฺติยา สญฺเจตนิกาย สุกฺกวิสฺสฏฺฐิยา ปญฺจาหปฺปฏิจฺฉนฺนาย ปญฺจาหปริวาสํ อทาสิ, โสหํ ปริวสนฺโต อนฺตรา เอกํ อาปตฺติํ อาปชฺชิํ สญฺเจตนิกํ สุกฺกวิสฺสฏฺฐิํ อปฺปฏิจฺฉนฺนํ, โสหํ ภนฺเต สํฆํ อนฺตรา เอกิสฺสา อาปตฺติยา สญฺเจตนิกาย สุกฺกวิสฺสฏฺฐิยา อปฺปฏิจฺฉนฺนาย มูลาย ปฏิกสฺสนํ ยาจามี”ติ.\csromanspacer{} ทุติยมฺปิ ยาจิตพฺพา\footnote{ยาจิตพฺโพ (สี, เอวมุปริวิ)}.\csromanspacer{} ตติยมฺปิ ยาจิตพฺพา2.\csromanspacer{} พฺยตฺเตน ภิกฺขุนา ปฏิพเลน สํโฆ ญาเปตพฺโพ\csromandash{}}

\proseitem{111}{“สุณาตุ เม ภนฺเต สํโฆ, อยํ อุทายี ภิกฺขุ เอกํ อาปตฺติํ อาปชฺชิ สญฺเจตนิกํ สุกฺกวิสฺสฏฺฐิํ ปญฺจาหปฺปฏิจฺฉนฺนํ, โส สํฆํ เอกิสฺสา อาปตฺติยา สญฺเจตนิกาย สุกฺกวิสฺสฏฺฐิยา ปญฺจาหปฺปฏิจฺฉนฺนาย ปญฺจาหปริวาสํ ยาจิ, สํโฆ อุทายิสฺส ภิกฺขุโน เอกิสฺสา อาปตฺติยา สญฺเจตนิกาย สุกฺกวิสฺสฏฺฐิยา ปญฺจาหปฺปฏิจฺฉนฺนาย ปญฺจาหปริวาสํ อทาสิ, โส ปริวสนฺโต อนฺตรา เอกํ อาปตฺติํ อาปชฺชิ สญฺเจตนิกํ สุกฺกวิสฺสฏฺฐิํ อปฺปฏิจฺฉนฺนํ, โส สํฆํ อนฺตรา เอกิสฺสา อาปตฺติยา สญฺเจตนิกาย สุกฺกวิสฺสฏฺฐิยา อปฺปฏิจฺฉนฺนาย มูลายปฏิกสฺสนํ ยาจติ, ยทิ สํฆสฺส ปตฺตกลฺลํ, สํโฆ อุทายิํ ภิกฺขุํ อนฺตรา เอกิสฺสา อาปตฺติยา สญฺเจตนิกาย สุกฺกวิสฺสฏฺฐิยา อปฺปฏิจฺฉนฺนาย มูลาย ปฏิกสฺเสยฺย, เอสา ญตฺติ.}

\prose{\csromanfolio{115}สุณาตุ เม ภนฺเต สํโฆ, อยํ อุทายี ภิกฺขุ เอกํ อาปตฺติํ อาปชฺชิ สญฺเจตนิกํ สุกฺกวิสฺสฏฺฐิํ ปญฺจาหปฺปฏิจฺฉนฺนํ, โส สํฆํ เอกิสฺสา อาปตฺติยา สญฺเจตนิกาย สุกฺกวิสฺสฏฺฐิยา ปญฺจาหปฺปฏิจฺฉนฺนาย ปญฺจาหปริวาสํ ยาจิ, สํโฆ อุทายิสฺส ภิกฺขุโน เอกิสฺสา อาปตฺติยา สญฺเจตนิกาย สุกฺกวิสฺสฏฺฐิยา ปญฺจาหปฺปฏิจฺฉนฺนาย ปญฺจาหปริวาสํ อทาสิ, โส ปริวสนฺโต อนฺตรา เอกํ อาปตฺติํ อาปชฺชิ สญฺเจตนิกํ สุกฺกวิสฺสฏฺฐิํ อปฺปฏิจฺฉนฺนํ, โส สํฆํ อนฺตรา เอกิสฺสา อาปตฺติยา สญฺเจตนิกาย สุกฺกวิสฺสฏฺฐิยา อปฺปฏิจฺฉนฺนาย มูลายปฏิกสฺสนํ ยาจติ, สํโฆ อุทายิํ ภิกฺขุํ อนฺตรา เอกิสฺสา อาปตฺติยา สญฺเจตนิกาย สุกฺกวิสฺสฏฺฐิยา อปฺปฏิจฺฉนฺนาย มูลาย ปฏิกสฺสติ, ยสฺสายสฺมโต ขมติ อุทายิสฺส ภิกฺขุโน อนฺตรา เอกิสฺสา อาปตฺติยา สญฺเจตนิกาย สุกฺกวิสฺสฏฺฐิยา อปฺปฏิจฺฉนฺนาย มูลายปฏิกสฺสนา, โส ตุณฺหสฺส, ยสฺส นกฺขมติ, โส ภาเสยฺย.}

\prose{ทุติยมฺปิ เอตมตฺถํ วทามิ ฯเปฯ.\csromanspacer{} ตติยมฺปิ เอตมตฺถํ วทามิ ฯเปฯ.}

\prose{ปฏิกสฺสิโต สํเฆน อุทายี ภิกฺขุ อนฺตรา เอกิสฺสา อาปตฺติยา สญฺเจตนิกาย สุกฺกวิสฺสฏฺฐิยา อปฺปฏิจฺฉนฺนาย มูลายปฏิกสฺสนา\footnote{มูลาย (สฺยา)}, ขมติ สํฆสฺส, ตสฺมา ตุณฺหี, เอวเมตํ ธารยามี”ติ.}

\csromanmatikaanchor{mtk.68}
\csromantocmarkref{subsection}{มานตฺตารหมูลายปฏิกสฺสนา}{mtk.68}
\bookmark[dest={mtk.68},level=2]{มานตฺตารหมูลายปฏิกสฺสนา}
\csromanheader{2}{\textbf{มานตฺตารหมูลายปฏิกสฺสนา}}
\proseitem{112}{โส ปริวุตฺถปริวาโส มานตฺตารโห อนฺตรา เอกํ อาปตฺติํ อาปชฺชิ สญฺเจตนิกํ สุกฺกวิสฺสฏฺฐิํ อปฺปฏิจฺฉนฺนํ, โส ภิกฺขูนํ อาโรเจสิ “อหํ อาวุโส เอกํ อาปตฺติํ อาปชฺชิํ สญฺเจตนิกํ สุกฺกวิสฺสฏฺฐิํ ปญฺจาหปฺปฏิจฺฉนฺนํ โสหํ สํฆํ เอกิสฺสา อาปตฺติยา สญฺเจตนิกาย สุกฺกวิสฺสฏฺฐิยา ปญฺจาหปฺปฏิจฺฉนฺนาย ปญฺจาหปริวาสํ ยาจิํ, ตสฺส เม สํโฆ เอกิสฺสา อาปตฺติยา สญฺเจตนิกาย สุกฺกวิสฺสฏฺฐิยา ปญฺจาหปฺปฏิจฺฉนฺนาย ปญฺจาหปริวาสํ อทาสิ, โสหํ ปริวสนฺโต อนฺตรา เอกํ อาปตฺติํ อาปชฺชิํ สญฺเจตนิกํ สุกฺกวิสฺสฏฺฐิํ อปฺปฏิจฺฉนฺนํ, โสหํ สํฆํ อนฺตรา เอกิสฺสา อาปตฺติยา สญฺเจตนิกาย สุกฺกวิสฺสฏฺฐิยา อปฺปฏิจฺฉนฺนาย มูลายปฏิกสฺสนํ ยาจิํ, ตํ มํ สํโฆ อนฺตรา เอกิสฺสา อาปตฺติยา \csromanfolio{116}สญฺเจตนิกาย สุกฺกวิสฺสฏฺฐิยา อปฺปฏิจฺฉนฺนาย มูลาย ปฏิกสฺสิ, โสหํ ปริวุตฺถปริวาโส มานตฺตารโห อนฺตรา เอกํ อาปตฺติํ อาปชฺชิํ สญฺเจตนิกํ สุกฺกวิสฺสฏฺฐิํ อปฺปฏิจฺฉนฺนํ, กถํ นุ โข มายา ปฏิปชฺชิตพฺพนฺ”ติ.\csromanspacer{} ภควโต เอตมตฺถํ อาโรเจสุํ.\csromanspacer{} เตน หิ ภิกฺขเว สํโฆ อุทายิํ ภิกฺขุํ อนฺตรา เอกิสฺสา อาปตฺติยา สญฺเจตนิกาย สุกฺกวิสฺสฏฺฐิยา อปฺปฏิจฺฉนฺนาย มูลาย ปฏิกสฺสตุ.\csromanspacer{} เอวญฺจ ปน ภิกฺขเว มูลาย ปฏิกสฺสิตพฺโพ\csromandash{}}

\prose{เตน ภิกฺขเว อุทายินา ภิกฺขุนา สํฆํ อุปสงฺกมิตฺวา ฯเปฯ เอวมสฺส วจนีโย “อหํ ภนฺเต เอกํ อาปตฺติํ อาปชฺชิํ สญฺเจตนิกํ สุกฺกวิสฺสฏฺฐิํ ปญฺจาหปฺปฏิจฺฉนฺนํ ฯเปฯ โสหํ ปริวุตฺถปริวาโส มานตฺตารโห อนฺตรา เอกํ อาปตฺติํ อาปชฺชิํ สญฺเจตนิกํ สุกฺกวิสฺสฏฺฐิํ อปฺปฏิจฺฉนฺนํ โสหํ สํฆํ อนฺตรา เอกิสฺสา อาปตฺติยา สญฺเจตนิกาย สุกฺกวิสฺสฏฺฐิยา อปฺปฏิจฺฉนฺนาย มูลายปฏิกสฺสนํ ยาจามี”ติ.\csromanspacer{} ทุติยมฺปิ ยาจิตพฺพา.\csromanspacer{} ตติยมฺปิ ยาจิตพฺพา.\csromanspacer{} พฺยตฺเตน ภิกฺขุนา ปฏิพเลน สํโฆ ญาเปตพฺโพ\csromandash{}}

\proseitem{113}{“สุณาตุ เม ภนฺเต สํโฆ, อยํ อุทายี ภิกฺขุ เอกํ อาปตฺติํ อาปชฺชิ สญฺเจตนิกํ สุกฺกวิสฺสฏฺฐิํ ปญฺจาหปฺปฏิจฺฉนฺนํ ฯเปฯ โส ปริวุตฺถปริวาโส มานตฺตารโห อนฺตรา เอกํ อาปตฺติํ อาปชฺชิ สญฺเจตนิกํ สุกฺกวิสฺสฏฺฐิํ อปฺปฏิจฺฉนฺนํ, โส สํฆํ อนฺตรา เอกิสฺสา อาปตฺติยา สญฺเจตนิกาย สุกฺกวิสฺสฏฺฐิยา อปฺปฏิจฺฉนฺนาย มูลายปฏิกสฺสนํ ยาจติ, ยทิ สํฆสฺส ปตฺตกลฺลํ, สํโฆ อุทายิํ ภิกฺขุํ อนฺตรา เอกิสฺสา อาปตฺติยา สญฺเจตนิกาย สุกฺกวิสฺสฏฺฐิยา อปฺปฏิจฺฉนฺนาย มูลาย ปฏิกสฺเสยฺย, เอสา ญตฺติ.}

\prose{สุณาตุ เม ภนฺเต สํโฆ, อยํ อุทายี ภิกฺขุ เอกํ อาปตฺติํ อาปชฺชิ สญฺเจตนิกํ สุกฺกวิสฺสฏฺฐิํ ปญฺจาหปฺปฏิจฺฉนฺนํ ฯเปฯ โส ปริวุตฺถปริวาโส มานตฺตารโห อนฺตรา เอกํ อาปตฺติํ อาปชฺชิ สญฺเจตนิกํ สุกฺกวิสฺสฏฺฐิํ อปฺปฏิจฺฉนฺนํ, โส สํฆํ อนฺตรา เอกิสฺสา อาปตฺติยา สญฺเจตนิกาย สุกฺกวิสฺสฏฺฐิยา อปฺปฏิจฺฉนฺนาย มูลายปฏิกสฺสนํ ยาจติ, สํโฆ อุทายิํ ภิกฺขุํ อนฺตรา เอกิสฺสา อาปตฺติยา สญฺเจตนิกาย สุกฺกวิสฺสฏฺฐิยา อปฺปฏิจฺฉนฺนาย มูลาย ปฏิกสฺสติ, ยสฺสายสฺมโต ขมติ อุทายิสฺส ภิกฺขุโน อนฺตรา เอกิสฺสา อาปตฺติยา สญฺเจตนิกาย \csromanfolio{117}สุกฺกวิสฺสฏฺฐิยา อปฺปฏิจฺฉนฺนาย มูลายปฏิกสฺสนา, โส ตุณฺหสฺส, ยสฺส นกฺขมติ, โส ภาเสยฺย.}

\prose{ทุติยมฺปิ เอตมตฺถํ วทามิ ฯเปฯ.\csromanspacer{} ตติยมฺปิ เอตมตฺถํ วทามิ ฯเปฯ.}

\prose{ปฏิกสฺสิโต สํเฆน อุทายี ภิกฺขุ อนฺตรา เอกิสฺสา อาปตฺติยา สญฺเจตนิกาย สุกฺกวิสฺสฏฺฐิยา อปฺปฏิจฺฉนฺนาย มูลายปฏิกสฺสนา, ขมติ สํฆสฺส, ตสฺมา ตุณฺหี, เอวเมตํ ธารยามี”ติ.}

\csromanmatikaanchor{mtk.69}
\csromantocmarkref{subsection}{ติกาปตฺติมานตฺต}{mtk.69}
\bookmark[dest={mtk.69},level=2]{ติกาปตฺติมานตฺต}
\csromanheader{2}{\textbf{ติกาปตฺติมานตฺต}}
\proseitem{114}{โส ปริวุตฺถปริวาโส ภิกฺขูนํ อาโรเจสิ “อหํ อาวุโส เอกํ อาปตฺติํ อาปชฺชิํ สญฺเจตนิกํ สุกฺกวิสฺสฏฺฐิํ ปญฺจาหปฺปฏิจฺฉนฺนํ ฯเปฯ โสหํ ปริวุตฺถปริวาโส, กถํ นุ โข มยา ปฏิปชฺชิตพฺพนฺ”ติ.\csromanspacer{} ภควโต เอตมตฺถํ อาโรเจสุํ.\csromanspacer{} เตน หิ ภิกฺขเว สํโฆ อุทายิสฺส ภิกฺขุโน ติสฺสนฺนํ อาปตฺตีนํ ฉารตฺตํ มานตฺตํ เทตุ.\csromanspacer{} เอวญฺจ ปน ภิกฺขเว ทาตพฺพํ\csromandash{}}

\prose{เตน ภิกฺขเว อุทายินา ภิกฺขุนา สํฆํ อุปสงฺกมิตฺวา ฯเปฯ เอวมสฺส วจนีโย “อหํ ภนฺเต เอกํ อาปตฺติํ อาปชฺชิํ สญฺเจตนิกํ สุกฺกวิสฺสฏฺฐิํ ปญฺจาหปฺปฏิจฺฉนฺนํ, โสหํ สํฆํ เอกิสฺสา อาปตฺติยา สญฺเจตนิกาย สุกฺกวิสฺสฏฺฐิยา ปญฺจาหปฺปฏิจฺฉนฺนาย ปญฺจาหปริวาสํ ยาจิํ, ตสฺส เม สํโฆ เอกิสฺสา อาปตฺติยา สญฺเจตนิกาย สุกฺกวิสฺสฏฺฐิยา ปญฺจาหปฺปฏิจฺฉนฺนาย ปญฺจาหปริวาสํ อทาสิ ฯเปฯ โสหํ ภนฺเต ปริวุตฺถปริวาโส สํฆํ ติสฺสนฺนํ อาปตฺตีนํ ฉารตฺถํ มานตฺตํ ยาจามี”ติ.\csromanspacer{} ทุติยมฺปิ ยาจิตพฺพํ.\csromanspacer{} ตติยมฺปิ ยาจิตพฺพํ.\csromanspacer{} พฺยตฺเตน ภิกฺขุนา ปฏิพเลน สํโฆ ญาเปตพฺโพ\csromandash{}}

\proseitem{115}{“สุณาตุ เม ภนฺเต สํโฆ, อยํ อุทายี ภิกฺขุ เอกํ อาปตฺติํ อาปชฺชิ สญฺเจตนิกํ สุกฺกวิสฺสฏฺฐิํ ปญฺจาหปฺปฏิจฺฉนฺนํ ฯเปฯ.\csromanspacer{} โส ปริวุตฺถปริวาโส สํฆํ ติสฺสนฺนํ อาปตฺตีนํ ฉารตฺตํ มานตฺตํ ยาจติ, ยทิ สํฆสฺส ปตฺตกลฺลํ, สํโฆ อุทายิสฺส ภิกฺขุโน ติสฺสนฺนํ อาปตฺตีนํ ฉารตฺตํ มานตฺตํ ทเทยฺย, เอสา ญตฺติ.}

\prose{สุณาตุ เม ภนฺเต สํโฆ, อยํ อุทายี ภิกฺขุ เอกํ อาปตฺติํ อาปชฺชิ สญฺเจตนิกํ สุกฺกวิสฺสฏฺฐิํ ปญฺจาหปฺปฏิจฺฉนฺนํ ฯเปฯ.\csromanspacer{} โส ปริวุตฺถปริวาโส \csromanfolio{118}สํฆํ ติสฺสนฺนํ อาปตฺตีนํ ฉารตฺตํ มานตฺตํ ยาจติ, สํโฆ อุทายิสฺส ภิกฺขุโน ติสฺสนฺนํ อาปตฺตีนํ ฉารตฺตํ มานตฺตํ เทติ, ยสฺสายสฺมโต ขมติ อุทายิสฺส ภิกฺขุโน ติสฺสนฺนํ อาปตฺตีนํ ฉารตฺตํ มานตฺตสฺส ทานํ, โส ตุณฺหสฺส, ยสฺส นกฺขมติ, โส ภาเสยฺย.}

\prose{ทุติยมฺปิ เอตมตฺถํ วทามิ ฯเปฯ.\csromanspacer{} ตติยมฺปิ เอตมตฺถํ วทามิ ฯเปฯ.}

\prose{ทินฺนํ สํเฆน อุทายิสฺส ภิกฺขุโน ติสฺสนฺนํ อาปตฺตีนํ ฉารตฺตํ มานตฺตํ, ขมติ สํฆสฺส, ตสฺมา ตุณฺหี, เอวเมตํ ธารยามี”ติ.}

\csromanmatikaanchor{mtk.70}
\csromantocmarkref{subsection}{มานตฺตจาริกมูลายปฏิกสฺสนา}{mtk.70}
\bookmark[dest={mtk.70},level=2]{มานตฺตจาริกมูลายปฏิกสฺสนา}
\csromanheader{2}{\textbf{มานตฺตจาริกมูลายปฏิกสฺสนา}}
\proseitem{116}{โส มานตฺตํ จรนฺโต อนฺตรา เอกํ อาปตฺติํ อาปชฺชิ สญฺเจตนิกํ สุกฺกวิสฺสฏฺฐิํ อปฺปฏิจฺฉนฺนํ, โส ภิกฺขูนํ อาโรเจสิ “อหํ อาวุโส เอกํ อาปตฺติํ อาปชฺชิํ สญฺเจตนิกํ สุกฺกวิสฺสฏฺฐิํ ปญฺจาหปฺปฏิจฺฉนฺนํ ฯเปฯ โสหํ มานตฺตํ จรนฺโต อนฺตรา เอกํ อาปตฺติํ อาปชฺชิํ สญฺเจตนิกํ สุกฺกวิสฺสฏฺฐิํ อปฺปฏิจฺฉนฺนํ, กถํ นุ โข มยา ปฏิปชฺชิตพฺพนฺ”ติ.\csromanspacer{} ภควโต เอตมตฺถํ อาโรเจสุํ.\csromanspacer{} เตน หิ ภิกฺขเว สํโฆ อุทายิํ ภิกฺขุํ อนฺตรา เอกิสฺสา อาปตฺติยา สญฺเจตนิกาย สุกฺกวิสฺสฏฺฐิยา อปฺปฏิจฺฉนฺนาย มูลาย ปฏิกสฺสิตฺวา ฉารตฺตํ มานตฺตํ เทตุ.\csromanspacer{} เอวญฺจ ปน ภิกฺขเว มูลาย ปฏิกสฺสิตพฺโพ, เตน ภิกฺขเว อุทายินา ภิกฺขุนา สํฆํ อุปสงฺกมิตฺวา ฯเปฯ เอวมสฺส วจนีโย “อหํ ภนฺเต เอกํ อาปตฺติํ อาปชฺชิํ สญฺเจตนิกํ สุกฺกวิสฺสฏฺฐิํ ปญฺจาหปฺปฏิจฺฉนฺนํ ฯเปฯ โสหํ มานตฺตํ จรนฺโต อนฺตรา เอกํ อาปตฺติํ อาปชฺชิํ สญฺเจตนิกํ สุกฺกวิสฺสฏฺฐิํ อปฺปฏิจฺฉนฺนํ โสหํ ภนฺเต สํฆํ อนฺตรา เอกิสฺสา อาปตฺติยา สญฺเจตนิกาย สุกฺกวิสฺสฏฺฐิยา อปฺปฏิจฺฉนฺนาย มูลายปฏิกสฺสนํ ยาจามี”ติ.\csromanspacer{} ทุติยมฺปิ ยาจิตพฺพา.\csromanspacer{} ตติยมฺปิ ยาจิตพฺพา.\csromanspacer{} พฺยตฺเตน ภิกฺขุนา ปฏิพเลน สํโฆ ญาเปตพฺโพ\csromandash{}}

\proseitem{117}{“สุณาตุ เม ภนฺเต สํโฆ, อยํ อุทายี ภิกฺขุ ฯเปฯ มูลายปฏิกสฺสนํ ยาจติ, ยทิ สํฆสฺส ปตฺตกลฺลํ, สํโฆ อุทายิํ ภิกฺขุํ อนฺตรา เอกิสฺสา อาปตฺติยา สญฺเจตนิกาย สุกฺกวิสฺสฏฺฐิยา อปฺปฏิจฺฉนฺนาย มูลาย ปฏิกสฺเสยฺย, เอสา ญตฺติ ฯเปฯ.\csromanspacer{} ปฏิกสฺสิโต สํเฆน อุทายี ภิกฺขุ อนฺตรา เอกิสฺสา อาปตฺติยา สญฺเจตนิกาย \csromanfolio{119}สุกฺกวิสฺสฏฺฐิยา อปฺปฏิจฺฉนฺนาย มูลายปฏิกสฺสนา, ขมติ สํฆสฺส, ตสฺมา ตุณฺหี, เอวเมตํ ธารยามี”ติ.}

\prose{เอวญฺจ ปน ภิกฺขเว ฉารตฺตํ มานตฺตํ ทาตพฺพํ, เตน ภิกฺขเว อุทายินา ภิกฺขุนา ฯเปฯ เอวมสฺส วจนีโย “อหํ ภนฺเต เอกํ อาปตฺติํ อาปชฺชิํ สญฺเจตนิกํ สุกฺกวิสฺสฏฺฐิํ ปญฺจาหปฺปฏิจฺฉนฺนํ ฯเปฯ โสหํ มานตฺตํ จรนฺโต อนฺตรา เอกํ อาปตฺติํ อาปชฺชิํ สญฺเจตนิกํ สุกฺกวิสฺสฏฺฐิํ อปฺปฏิจฺฉนฺนํ, โสหํ สํฆํ อนฺตรา เอกิสฺสา อาปตฺติยา สญฺเจตนิกาย สุกฺกวิสฺสฏฺฐิยา อปฺปฏิจฺฉนฺนาย มูลายปฏิกสฺสนํ ยาจิํ, ตํ มํ สํโฆ อนฺตรา เอกิสฺสา อาปตฺติยา สญฺเจตนิกายา สุกฺกวิสฺสฏฺฐิยา อปฺปฏิจฺฉนฺนาย มูลาย ปฏิกสฺสิ.\csromanspacer{} โสหํ ภนฺเต สํฆํ อนฺตรา เอกิสฺสา อาปตฺติยา สญฺเจตนิกาย สุกฺกวิสฺสฏฺฐิยา อปฺปฏิจฺฉนฺนาย ฉารตฺตํ มานตฺตํ ยาจามี”ติ.\csromanspacer{} ทุติยมฺปิ ยาจิตพฺพํ.\csromanspacer{} ตติยมฺปิ ยาจิตพฺพํ.\csromanspacer{} พฺยตฺเตน ภิกฺขุนา ปฏิพเลน สํโฆ ญาเปตพฺโพ\csromandash{}}

\prose{“สุณาตุ เม ภนฺเต สํโฆ อยํ อุทายี ภิกฺขุ ฯเปฯ ฉารตฺตํ มานตฺตํ ยาจติ, ยทิ สํฆสฺส ปตฺตกลฺลํ, สํโฆ อุทายิสฺส ภิกฺขุโน อนฺตรา เอกิสฺสา อาปตฺติยา สญฺเจตนิกาย สุกฺกวิสฺสฏฺฐิยา อปฺปฏิจฺฉนฺนาย ฉารตฺตํ มานตฺตํ ทเทยฺย, เอสา ญตฺติ ฯเปฯ.\csromanspacer{} ทินฺนํ สํเฆน อุทายิสฺส ภิกฺขุโน อนฺตรา เอกิสฺสา อาปตฺติยา สญฺเจตนิกาย สุกฺกวิสฺสฏฺฐิยา อปฺปฏิจฺฉนฺนาย ฉารตฺตํ มานตฺตํ, ขมติ สํฆสฺส, ตสฺมา ตุณฺหี, เอวเมตํ ธารยามี”ติ.}

\csromanmatikaanchor{mtk.71}
\csromantocmarkref{subsection}{อพฺภานารหมูลายปฏิกสฺสนา}{mtk.71}
\bookmark[dest={mtk.71},level=2]{อพฺภานารหมูลายปฏิกสฺสนา}
\csromanheader{2}{\textbf{อพฺภานารหมูลายปฏิกสฺสนา}}
\proseitem{118}{โส จิณฺณมานตฺโต อพฺภานารโห อนฺตรา เอกํ อาปตฺติํ อาปชฺชิ สญฺเจตนิกํ สุกฺกวิสฺสฏฺฐิํ อปฺปฏิจฺฉนฺนํ, โส ภิกฺขูนํ อาโรเจสิ “อหํ อาวุโส เอกํ อาปตฺติํ อาปชฺชิํ สญฺเจตนิกํ สุกฺกวิสฺสฏฺฐิํ ปญฺจาหปฺปฏิจฺฉนฺนํ ฯเปฯ โสหํ จิณฺณมานตฺโต อพฺภานารโห อนฺตรา เอกํ อาปตฺติํ อาปชฺชิํ สญฺเจตนิกํ สุกฺกวิสฺสฏฺฐิํ อปฺปฏิจฺฉนฺนํ, กถํ นุ โข มยา ปฏิปชฺชิตพฺพนฺ”ติ.\csromanspacer{} ภควโต เอตมตฺถํ อาโรเจสุํ.\csromanspacer{} เตน หิ ภิกฺขเว สํโฆ อุทายิํ ภิกฺขุํ อนฺตรา เอกิสฺสา อาปตฺติยา สญฺเจตนิกาย สุกฺกวิสฺสฏฺฐิยา อปฺปฏิจฺฉนฺนาย มูลาย ปฏิกสฺสิตฺวา ฉารตฺตํ มานตฺตํ เทตุ.\csromanspacer{} เอวญฺจ ปน ภิกฺขเว มูลาย ปฏิกสฺสิตพฺโพ ฯเปฯ.}

\prose{\csromanfolio{120}เอวญฺจ ปน ภิกฺขเว ฉารตฺตํ มานตฺตํ ทาตพฺพํ ฯเปฯ.}

\prose{ทินฺนํ สํเฆน อุทายิสฺส ภิกฺขุโน อนฺตรา เอกิสฺสา อาปตฺติยา สญฺเจตนิกาย สุกฺกวิสฺสฏฺฐิยา อปฺปฏิจฺฉนฺนาย ฉารตฺตํ มานตฺตํ, ขมติ สํฆสฺส, ตสฺมา ตุณฺหี, เอวเมตํ ธารยามี”ติ.}

\csromanmatikaanchor{mtk.72}
\csromantocmarkref{subsection}{มูลายปฏิกสฺสิต-อพฺภาน}{mtk.72}
\bookmark[dest={mtk.72},level=2]{มูลายปฏิกสฺสิต-อพฺภาน}
\csromanheader{2}{\textbf{มูลายปฏิกสฺสิต-อพฺภาน}}
\proseitem{119}{โส จิณฺณมานตฺโต ภิกฺขูนํ อาโรเจสิ “อหํ อาวุโส เอกํ อาปตฺติํ อาปชฺชิํ สญฺเจตนิกํ สุกฺกวิสฺสฏฺฐิํ ปญฺจาหปฺปฏิจฺฉนฺนํ ฯเปฯ โสหํ จิณฺณมานตฺโต, กถํ นุ โข มยา ปฏิปชฺชิตพฺพนฺ”ติ.\csromanspacer{} ภควโต เอตมตฺถํ อาโรเจสุํ.\csromanspacer{} เตน หิ ภิกฺขเว สํโฆ อุทายิํ ภิกฺขุํ อพฺเภตุ.\csromanspacer{} เอวญฺจ ปน ภิกฺขเว อพฺเภตพฺโพ\csromandash{}}

\prose{เตน ภิกฺขเว อุทายินา ภิกฺขุนา สํฆํ อุปสงฺกมิตฺวา ฯเปฯ เอวมสฺส วจนีโย “อหํ ภนฺเต เอกํ อาปตฺติํ อาปชฺชิํ สญฺเจตนิกํ สุกฺกวิสฺสฏฺฐิํ ปญฺจาหปฺปฏิจฺฉนฺนํ, โสหํ สํฆํ เอกิสฺสา อาปตฺติยา สญฺเจตนิกาย สุกฺกวิสฺสฏฺฐิยา ปญฺจาหปฺปฏิจฺฉนฺนาย ปญฺจาหปริวาสํ ยาจิํ, ตสฺส เม สํโฆ เอกิสฺสา อาปตฺติยา สญฺเจตนิกาย สุกฺกวิสฺสฏฺฐิยา ปญฺจาหปฺปฏิจฺฉนฺนาย ปญฺจาหปริวาสํ อทาสิ.\csromanspacer{} โสหํ ปริวสนฺโต อนฺตรา เอกํ อาปตฺติํ อาปชฺชิํ สญฺเจตนิกํ สุกฺกวิสฺสฏฺฐิํ อปฺปฏิจฺฉนฺนํ, โสหํ สํฆํ อนฺตรา เอกิสฺสา อาปตฺติยา สญฺเจตนิกาย สุกฺกวิสฺสฏฺฐิยา อปฺปฏิจฺฉนฺนาย มูลายปฏิกสฺสนํ ยาจิํ, ตํ มํ สํโฆ อนฺตรา เอกิสฺสา อาปตฺติยา สญฺเจตนิกาย สุกฺกวิสฺสฏฺฐิยา อปฺปฏิจฺฉนฺนาย มูลาย ปฏิกสฺสิ.\csromanspacer{} โสหํ ปริวุตฺถปริวาโส มานตฺตารโห อนฺตรา เอกํ อาปตฺติํ อาปชฺชิํ สญฺเจตนิกํ สุกฺกวิสฺสฏฺฐิํ อปฺปฏิจฺฉนฺนํ, โสหํ สํฆํ อนฺตรา เอกิสฺสา อาปตฺติยา สญฺเจตนิกาย สุกฺกวิสฺสฏฺฐิยา อปฺปฏิจฺฉนฺนาย มูลายปฏิกสฺสนํ ยาจิํ, ตํ มํ สํโฆ อนฺตรา เอกิสฺสา อาปตฺติยา สญฺเจตนิกาย สุกฺกวิสฺสฏฺฐิยา อปฺปฏิจฺฉนฺนาย มูลาย ปฏิกสฺสิ.\csromanspacer{} โสหํ ปริวุตฺถวาโส สํฆํ ติสฺสนฺนํ อาปตฺตีนํ ฉารตฺตํ มานตฺตํ ยาจิํ, ตสฺส เม สํโฆ ติสฺสนฺนํ อาปตฺตีนํ ฉารตฺตํ มานตฺตํ อทาสิ.\csromanspacer{} โสหํ มานตฺตํ จรนฺโต อนฺตรา เอกํ อาปตฺติํ อาปชฺชิํ สญฺเจตนิกํ สุกฺกวิสฺสฏฺฐิํ อปฺปฏิจฺฉนฺนํ, โสหํ สํฆํ อนฺตรา เอกิสฺสา อาปตฺติยา สญฺเจตนิกาย สุกฺกวิสฺสฏฺฐิยา อปฺปฏิจฺฉนฺนาย มูลายปฏิกสฺสนํ ยาจิํ, ตํ มํ \csromanfolio{121}สํโฆ อนฺตรา เอกิสฺสา อาปตฺติยา สญฺเจตนิกาย สุกฺกวิสฺสฏฺฐิยา อปฺปฏิจฺฉนฺนาย มูลาย ปฏิกสฺสิ, โสหํ สํฆํ อนฺตรา เอกิสฺสา อาปตฺติยา สญฺเจตนิกาย สุกฺกวิสฺสฏฺฐิยา อปฺปฏิจฺฉนฺนาย ฉารตฺตํ มานตฺตํ ยาจิํ, ตสฺส เม สํโฆ อนฺตรา เอกิสฺสา อาปตฺติยา สญฺเจตนิกาย สุกฺกวิสฺสฏฺฐิยา อปฺปฏิจฺฉนฺนาย ฉารตฺตํ มานตฺตํ อทาสิ.\csromanspacer{} โสหํ จิณฺณมานตฺโต อพฺภานารโห อนฺตรา เอกํ อาปตฺติํ อาปชฺชิํ สญฺเจตนิกํ สุกฺกวิสฺสฏฺฐิํ อปฺปฏิจฺฉนฺนํ, โสหํ สํฆํ อนฺตรา เอกิสฺสา อาปตฺติยา สญฺเจตนิกาย สุกฺกวิสฺสฏฺฐิยา อปฺปฏิจฺฉนฺนาย มูลายปฏิกสฺสนํ ยาจิํ, ตํ มํ สํโฆ อนฺตรา เอกิสฺสา อาปตฺติยา สญฺเจตนิกาย สุกฺกวิสฺสฏฺฐิยา อปฺปฏิจฺฉนฺนาย มูลาย ปฏิกสฺสิ, โสหํ สํฆํ อนฺตรา เอกิสฺสา อาปตฺติยา สญฺเจตนิกาย สุกฺกวิสฺสฏฺฐิยา อปฺปฏิจฺฉนฺนาย ฉารตฺตํ มานตฺตํ ยาจิํ, ตสฺส เม สํโฆ อนฺตรา เอกิสฺสา อาปตฺติยา สญฺเจตนิกาย สุกฺกวิสฺสฏฺฐิยา อปฺปฏิจฺฉนฺนาย ฉารตฺตํ มานตฺตํ อทาสิ, โสหํ ภนฺเต จิณฺณมานตฺโต สํฆํ อพฺภานํ ยาจามี”ติ.}

\prose{ทุติยมฺปิ ยาจิตพฺพํ.\csromanspacer{} ตติยมฺปิ ยาจิตพฺพํ.\csromanspacer{} พฺยตฺเตน ภิกฺขุนา ปฏิพเลน สํโฆ ญาเปตพฺโพ\csromandash{}}

\proseitem{120}{“สุณาตุ เม ภนฺเต สํโฆ, อยํ อุทายี ภิกฺขุ เอกํ อาปตฺติํ อาปชฺชิ สญฺเจตนิกํ สุกฺกวิสฺสฏฺฐิํ ปญฺจาหปฺปฏิจฺฉนฺนํ, โส สํฆํ เอกิสฺสา อาปตฺติยา สญฺเจตนิกาย สุกฺกวิสฺสฏฺฐิยา ปญฺจาหปฺปฏิจฺฉนฺนาย ปญฺจาหปริวาสํ ยาจิ, สํโฆ อุทายิสฺส ภิกฺขุโน เอกิสฺสา อาปตฺติยา สญฺเจตนิกาย สุกฺกวิสฺสฏฺฐิยา ปญฺจาหปฺปฏิจฺฉนฺนาย ปญฺจาหปริวาสํ อทาสิ.\csromanspacer{} โส ปริวสนฺโต อนฺตรา เอกํ อาปตฺติํ อาปชฺชิ สญฺเจตนิกํ สุกฺกวิสฺสฏฺฐิํ อปฺปฏิจฺฉนฺนํ, โส สํฆํ อนฺตรา เอกิสฺสา อาปตฺติยา สญฺเจตนิกาย สุกฺกวิสฺสฏฺฐิยา อปฺปฏิจฺฉนฺนาย มูลายปฏิกสฺสนํ ยาจิ, สํโฆ อุทายิํ ภิกฺขุํ อนฺตรา เอกิสฺสา อาปตฺติยา สญฺเจตนิกาย สุกฺกวิสฺสฏฺฐิยา อปฺปฏิจฺฉนฺนาย มูลาย ปฏิกสฺสิ.\csromanspacer{} โส ปริวุตฺถปริวาโส มานตฺตารโห อนฺตรา เอกํ อาปตฺติํ อาปชฺชิ สญฺเจตนิกํ สุกฺกวิสฺสฏฺฐิํ อปฺปฏิจฺฉนฺนํ, โส สํฆํ อนฺตรา เอกิสฺสา อาปตฺติยา สญฺเจตนิกาย สุกฺกวิสฺสฏฺฐิยา อปฺปฏิจฺฉนฺนาย มูลายปฏิกสฺสนํ ยาจิ, สํโฆ อุทายิํ ภิกฺขุํ อนฺตรา เอกิสฺสา อาปตฺติยา \csromanfolio{122}สญฺเจตนิกาย สุกฺกวิสฺสฏฺฐิยา อปฺปฏิจฺฉนฺนาย มูลาย ปฏิกสฺสิ.\csromanspacer{} โส ปริวุตฺถปริวาโส สํฆํ ติสฺสนฺนํ อาปตฺตีนํ ฉารตฺตํ มานตฺตํ ยาจิ, สํโฆ อุทายิสฺส ภิกฺขุโน ติสฺสนฺนํ อาปตฺตีนํ ฉารตฺตํ มานตฺตํ อทาสิ.\csromanspacer{} โส มานตฺตํ จรนฺโต อนฺตรา เอกํ อาปตฺติํ อาปชฺชิ สญฺเจตนิกํ สุกฺกวิสฺสฏฺฐิํ อปฺปฏิจฺฉนฺนํ, โส สํฆํ อนฺตรา เอกิสฺสา อาปตฺติยา สญฺเจตนิกาย สุกฺกวิสฺสฏฺฐิยา อปฺปฏิจฺฉนฺนาย มูลายปฏิกสฺสนํ ยาจิ, สํโฆ อุทายิํ ภิกฺขุํ อนฺตรา เอกิสฺสา อาปตฺติยา สญฺเจตนิกาย สุกฺกวิสฺสฏฺฐิยา อปฺปฏิจฺฉนฺนาย มูลาย ปฏิกสฺสิ.\csromanspacer{} โส สํฆํ อนฺตรา เอกิสฺสา อาปตฺติยา สญฺเจตนิกาย สุกฺกวิสฺสฏฺฐิยา อปฺปฏิจฺฉนฺนาย ฉารตฺตํ มานตฺตํ ยาจิ, สํโฆ อุทายิสฺส ภิกฺขุโน อนฺตรา เอกิสฺสา อาปตฺติยา สญฺเจตนิกาย สุกฺกวิสฺสฏฺฐิยา อปฺปฏิจฺฉนฺนาย ฉารตฺตํ มานตฺตํ อทาสิ.\csromanspacer{} โส จิณฺณมานตฺโต อพฺภานารโห อนฺตรา เอกํ อาปตฺติํ อาปชฺชิ สญฺเจตนิกํ สุกฺกวิสฺสฏฺฐิํ อปฺปฏิจฺฉนฺนํ, โส สํฆํ อนฺตรา เอกิสฺสา อาปตฺติยา สญฺเจตนิกาย สุกฺกวิสฺสฏฺฐิยา อปฺปฏิจฺฉนฺนาย มูลาย ปฏิกสฺสนํ ยาจิ, สํโฆ อุทายิํ ภิกฺขุํ อนฺตรา เอกิสฺสา อาปตฺติยา สญฺเจตนิกาย สุกฺกวิสฺสฏฺฐิยา อปฺปฏิจฺฉนฺนาย มูลาย ปฏิกสฺสิ.\csromanspacer{} โส สํฆํ อนฺตรา เอกิสฺสา อาปตฺติยา สญฺเจตนิกาย สุกฺกวิสฺสฏฺฐิยา อปฺปฏิจฺฉนฺนาย ฉารตฺตํ มานตฺตํ ยาจิ, สํโฆ อุทายิสฺส ภิกฺขุโน อนฺตรา เอกิสฺสา อาปตฺติยา สญฺเจตนิกาย สุกฺกวิสฺสฏฺฐิยา อปฺปฏิจฺฉนฺนาย ฉารตฺตํ มานตฺตํ อทาสิ, โส จิณฺณมานตฺโต สํฆํ อพฺภานํ ยาจติ.\csromanspacer{} ยทิ สํฆสฺส ปตฺตกลฺลํ, สํโฆ อุทายิํ ภิกฺขุํ อพฺเภยฺย, เอสา ญตฺติ.}

\prose{สุณาตุ เม ภนฺเต สํโฆ, อยํ อุทายี ภิกฺขุ เอกํ อาปตฺติํ อาปชฺชิ สญฺเจตนิกํ สุกฺกวิสฺสฏฺฐิํ ปญฺจาหปฺปฏิจฺฉนฺนํ ฯเปฯ โส จิณฺณมานตฺโต สํฆํ อพฺภานํ ยาจติ.\csromanspacer{} สํโฆ อุทายิํ ภิกฺขุํ อพฺเภติ, ยสฺสายสฺมโต ขมติ อุทายิสฺส ภิกฺขุโน อพฺภานํ, โส ตุณฺหสฺส, ยสฺส นกฺขมติ โส ภาเสยฺย.}

\prose{ทุติยมฺปิ เอตมตฺถํ วทามิ ฯเปฯ.\csromanspacer{} ตติยมฺปิ เอตมตฺถํ วทามิ ฯเปฯ.}

\prose{อพฺภิโต สํเฆน อุทายี ภิกฺขุ, ขมติ สํฆสฺส, ตสฺมา ตุณฺหี, เอวเมตํ ธารยามี”ติ.}

\csromanheader{2}{\textbf{ปกฺขฏิจฺฉนฺนปริวาส}}
\csromanmatikaanchor{mtk.73}
\csromantocmarkref{subsection}{ปกฺขปฺปฏิจฺฉนฺนปริวาส}{mtk.73}
\bookmark[dest={mtk.73},level=2]{ปกฺขปฺปฏิจฺฉนฺนปริวาส}
\proseitem{121}{\csromanfolio{123}เตน โข ปน สมเยน อายสฺมา อุทายี เอกํ อาปตฺติํ อาปนฺโน โหติ สญฺเจตนิกํ สุกฺกวิสฺสฏฺฐิํ ปกฺขปฺปฏิจฺฉนฺนํ, โส ภิกฺขูนํ อาโรเจสิ “อหํ อาวุโส เอกํ อาปตฺติํ อาปชฺชิํ สญฺเจตนิกํ สุกฺกวิสฺสฏฺฐิํ ปกฺขปฺปฏิจฺฉนฺนํ, กถํ นุ โข มยา ปฏิปชฺชิตพฺพนฺ”ติ.\csromanspacer{} ภควโต เอตมตฺถํ อาโรเจสุํ.\csromanspacer{} เตน หิ ภิกฺขเว สํโฆ อุทายิสฺส ภิกฺขุโน เอกิสฺสา อาปตฺติยา สญฺเจตนิกาย สุกฺกวิสฺสฏฺฐิยา ปกฺขปฺปฏิจฺฉนฺนาย ปกฺขปริวาสํ เทตุ, เอวญฺจ ปน ภิกฺขเว ทาตพฺโพ\csromandash{}}

\prose{เตน ภิกฺขเว อุทายินา ภิกฺขุนา สํฆํ อุปสงฺกมิตฺวา ฯเปฯ เอวมสฺส วจนีโย “อหํ ภนฺเต เอกํ อาปตฺติํ อาปชฺชิํ สญฺเจตนิกํ สุกฺกวิสฺสฏฺฐิํ ปกฺขปฺปฏิจฺฉนฺนํ, โสหํ ภนฺเต สํฆํ เอกิสฺสา อาปตฺติยา สญฺเจตนิกาย สุกฺกวิสฺสฏฺฐิยา ปกฺขปฺปฏิจฺฉนฺนาย ปกฺขปริวาสํ ยาจามี”ติ.\csromanspacer{} ทุติยมฺปิ ยาจิตพฺโพ.\csromanspacer{} ตติยมฺปิ ยาจิตพฺโพ.\csromanspacer{} พฺยตฺเตน ภิกฺขุนา ปฏิพเลน สํโฆ ญาเปตพฺโพ\csromandash{}}

\proseitem{122}{“สุณาตุ เม ภนฺเต สํโฆ, อยํ อุทายี ภิกฺขุ เอกํ อาปตฺติํ อาปชฺชิ สญฺเจตนิกํ สุกฺกวิสฺสฏฺฐิํ ปกฺขปฺปฏิจฺฉนฺนํ, โส สํฆํ เอกิสฺสา อาปตฺติยา สญฺเจตนิกาย สุกฺกวิสฺสฏฺฐิยา ปกฺขปฺปฏิจฺฉนฺนาย ปกฺขปริวาสํ ยาจติ, ยทิ สํฆสฺส ปตฺตกลฺลํ, สํโฆ อุทายิสฺส ภิกฺขุโน เอกิสฺสา อาปตฺติยา สญฺเจตนิกาย สุกฺกวิสฺสฏฺฐิยา ปกฺขปฺปฏิจฺฉนฺนาย ปกฺขปริวาสํ ทเทยฺย, เอสา ญตฺติ.}

\prose{สุณาตุ เม ภนฺเต สํโฆ, อยํ อุทายี ภิกฺขุ เอกํ อาปตฺติํ อาปชฺชิ สญฺเจตนิกํ สุกฺกวิสฺสฏฺฐิํ ปกฺขปฺปฏิจฺฉนฺนํ, โส สํฆํ เอกิสฺสา อาปตฺติยา สญฺเจตนิกาย สุกฺกวิสฺสฏฺฐิยา ปกฺขปฺปฏิจฺฉนฺนาย ปกฺขปริวาสํ ยาจติ, สํโฆ อุทายิสฺส ภิกฺขุโน เอกิสฺสา อาปตฺติยา สญฺเจตนิกาย สุกฺกวิสฺสฏฺฐิยา ปกฺขปฺปฏิจฺฉนฺนาย ปกฺขปริวาสํ เทติ, ยสฺสายสฺมโต ขมติ อุทายิสฺส ภิกฺขุโน เอกิสฺสา อาปตฺติยา สญฺเจตนิกาย สุกฺกวิสฺสฏฺฐิยา ปกฺขปฺปฏิจฺฉนฺนาย ปกฺขปริวาสสฺส ทานํ, โส ตุณฺหสฺส, ยสฺส นกฺขมติ, โส ภาเสยฺย.}

\prose{ทุติยมฺปิ เอตมตฺถํ วทามิ ฯเปฯ.\csromanspacer{} ตติยมฺปิ เอตมตฺถํ วทามิ ฯเปฯ.}

\prose{\csromanfolio{124}ทินฺโน สํเฆน อุทายิสฺส ภิกฺขุโน เอกิสฺสา อาปตฺติยา สญฺเจตนิกาย สุกฺกวิสฺสฏฺฐิยา ปกฺขปฺปฏิจฺฉนฺนาย ปกฺขปริวาโส, ขมติ สํฆสฺส, ตสฺมา ตุณฺหี, เอวเมตํ ธารยามี”ติ.}

\csromanmatikaanchor{mtk.74}
\csromantocmarkref{subsection}{ปกฺขปาริวาสิกมูลายปฏิกสฺสนา}{mtk.74}
\bookmark[dest={mtk.74},level=2]{ปกฺขปาริวาสิกมูลายปฏิกสฺสนา}
\csromanheader{2}{\textbf{ปกฺขปาริวาสิกมูลายปฏิกสฺสนา}}
\proseitem{123}{โส ปริวสนฺโต อนฺตรา เอกํ อาปตฺติํ อาปชฺชิ สญฺเจตนิกํ สุกฺกวิสฺสฏฺฐิํ ปญฺจาหปฺปฏิจฺฉนฺนํ, โส ภิกฺขูนํ อาโรเจสิ “อหํ อาวุโส เอกํ อาปตฺติํ อาปชฺชิํ สญฺเจตนิกํ สุกฺกวิสฺสฏฺฐิํ ปกฺขปฺปฏิจฺฉนฺนํ, โสหํ สํฆํ เอกิสฺสา อาปตฺติยา สญฺเจตนิกาย สุกฺกวิสฺสฏฺฐิยา ปกฺขปฺปฏิจฺฉนฺนาย ปกฺขปริวาสํ ยาจิํ, ตสฺส เม สํโฆ เอกิสฺสา อาปตฺติยา สญฺเจตนิกาย สุกฺกวิสฺสฏฺฐิยา ปกฺขปฺปฏิจฺฉนฺนาย ปกฺขปริวาสํ อทาสิ, โสหํ ปริวสนฺโต อนฺตรา เอกํ อาปตฺติํ อาปชฺชิํ สญฺเจตนิกํ สุกฺกวิสฺสฏฺฐิํ ปญฺจาหปฺปฏิจฺฉนฺนํ, กถํ นุ โข มยา ปฏิปชฺชิตพฺพนฺ”ติ.\csromanspacer{} ภควโต เอตมตฺถํ อาโรเจสุํ.\csromanspacer{} เตน หิ ภิกฺขเว สํโฆ อุทายิํ ภิกฺขุํ อนฺตรา เอกิสฺสา อาปตฺติยา สญฺเจตนิกาย สุกฺกวิสฺสฏฺฐิยา ปญฺจาหปฺปฏิจฺฉนฺนาย มูลาย ปฏิกสฺสิตฺวา ปุริมาย อาปตฺติยา สโมธานปริวาสํ เทตุ, เอวญฺจ ปน ภิกฺขเว มูลาย ปฏิกสฺสิตพฺโพ\csromandash{}}

\prose{เตน ภิกฺขเว อุทายินา ภิกฺขุนา สํฆํ อุปสงฺกมิตฺวา ฯเปฯ เอวมสฺส วจนีโย “อหํ ภนฺเต เอกํ อาปตฺติํ อาปชฺชิํ สญฺเจตนิกํ สุกฺกวิสฺสฏฺฐิํ ปกฺขปฺปฏิจฺฉนฺนํ, โสหํ สํฆํ เอกิสฺสา อาปตฺติยา สญฺเจตนิกาย สุกฺกวิสฺสฏฺฐิยา ปกฺขปฺปฏิจฺฉนฺนาย ปกฺขปริวาสํ ยาจิํ, ตสฺส เม สํโฆ เอกิสฺสา อาปตฺติยา สญฺเจตนิกาย สุกฺกวิสฺสฏฺฐิยา ปกฺขปฺปฏิจฺฉนฺนาย ปกฺขปริวาสํ อทาสิ.\csromanspacer{} โสหํ ปริวสนฺโต อนฺตรา เอกํ อาปตฺติํ อาปชฺชิํ สญฺเจตนิกํ สุกฺกวิสฺสฏฺฐิํ ปญฺจาหปฺปฏิจฺฉนฺนํ, โสหํ ภนฺเต สํฆํ อนฺตรา เอกิสฺสา อาปตฺติยา สญฺเจตนิกาย สุกฺกวิสฺสฏฺฐิยา ปญฺจาหปฺปฏิจฺฉนฺนาย มูลายปฏิกสฺสนํ ยาจามี”ติ.\csromanspacer{} ทุติยมฺปิ ยาจิตพฺพา.\csromanspacer{} ตติยมฺปิ ยาจิตพฺพา.\csromanspacer{} พฺยตฺเตน ภิกฺขุนา ปฏิพเลน สํโฆ ญาเปตพฺโพ\csromandash{}}

\proseitem{124}{“สุณาตุ เม ภนฺเต สํโฆ, อยํ อุทายี ภิกฺขุ เอกํ อาปตฺติํ อาปชฺชิ สญฺเจตนิกํ สุกฺกวิสฺสฏฺฐิํ ปกฺขปฺปฏิจฺฉนฺนํ, โส สํฆํ เอกิสฺสา อาปตฺติยา สญฺเจตนิกาย สุกฺกวิสฺสฏฺฐิยา ปกฺขปฺปฏิจฺฉนฺนาย ปกฺขปริวาสํ \csromanfolio{125}ยาจิ, สํโฆ อุทายิสฺส ภิกฺขุโน เอกิสฺสา อาปตฺติยา สญฺเจตนิกาย สุกฺกวิสฺสฏฺฐิยา ปกฺขปฺปฏิจฺฉนฺนาย ปกฺขปริวาสํ อทาสิ.\csromanspacer{} โส ปริวสนฺโต อนฺตรา เอกํ อาปตฺติํ อาปชฺชิ สญฺเจตนิกํ สุกฺกวิสฺสฏฺฐิํ ปญฺจาหปฺปฏิจฺฉนฺนํ, โส สํฆํ อนฺตรา เอกิสฺสา อาปตฺติยา สญฺเจตนิกาย สุกฺกวิสฺสฏฺฐิยา ปญฺจาหปฺปฏิจฺฉนฺนาย มูลายปฏิกสฺสนํ ยาจติ.\csromanspacer{} ยทิ สํฆสฺส ปตฺตกลฺลํ, สํโฆ อุทายิํ ภิกฺขุํ อนฺตรา เอกิสฺสา อาปตฺติยา สญฺเจตนิกาย สุกฺกวิสฺสฏฺฐิยา ปญฺจาหปฺปฏิจฺฉนฺนาย มูลาย ปฏิกสฺเสยฺย, เอสา ญตฺติ.}

\prose{สุณาตุ เม ภนฺเต สํโฆ, อยํ อุทายี ภิกฺขุ เอกํ อาปตฺติํ อาปชฺชิ สญฺเจตนิกํ สุกฺกวิสฺสฏฺฐิํ ปกฺขปฺปฏิจฺฉนฺนํ, โส สํฆํ เอกิสฺสา อาปตฺติยา สญฺเจตนิกาย สุกฺกวิสฺสฏฺฐิยา ปกฺขปฺปฏิจฺฉนฺนาย ปกฺขปริวาสํ ยาจิ, สํโฆ อุทายิสฺส ภิกฺขุโน เอกิสฺสา อาปตฺติยา สญฺเจตนิกาย สุกฺกวิสฺสฏฺฐิยา ปกฺขปฺปฏิจฺฉนฺนาย ปกฺขปริวาสํ อทาสิ.\csromanspacer{} โส ปริวสนฺโต อนฺตรา เอกํ อาปตฺติํ อาปชฺชิ สญฺเจตนิกํ สุกฺกวิสฺสฏฺฐิํ ปญฺจาหปฺปฏิจฺฉนฺนํ, โส สํฆํ อนฺตรา เอกิสฺสา อาปตฺติยา สญฺเจตนิกาย สุกฺกวิสฺสฏฺฐิยา ปญฺจาหปฺปฏิจฺฉนฺนาย มูลายปฏิกสฺสนํ ยาจติ.\csromanspacer{} สํโฆ อุทายิํ ภิกฺขุํ อนฺตรา เอกิสฺสา อาปตฺติยา สญฺเจตนิกาย สุกฺกวิสฺสฏฺฐิยา ปญฺจาหปฺปฏิจฺฉนฺนาย มูลาย ปฏิกสฺสติ, ยสฺสายสฺมโต ขมติ อุทายิสฺส ภิกฺขุโน อนฺตรา เอกิสฺสา อาปตฺติยา สญฺเจตนิกาย สุกฺกวิสฺสฏฺฐิยา ปญฺจาหปฺปฏิจฺฉนฺนาย มูลายปฏิกสฺสนา, โส ตุณฺหสฺส, ยสฺส นกฺขมติ, โส ภาเสยฺย.}

\prose{ทุติยมฺปิ เอตมตฺถํ วาทมิ ฯเปฯ.\csromanspacer{} ตติยมฺปิ เอตมตฺถํ วทามิ ฯเปฯ.}

\prose{ปฏิกสฺสิโต สํเฆน อุทายี ภิกฺขุ อนฺตรา เอกิสฺสา อาปตฺติยา สญฺเจตนิกาย สุกฺกวิสฺสฏฺฐิยา ปญฺจาหปฺปฏิจฺฉนฺนาย มูลายปฏิกสฺสนา, ขมติ สํฆสฺส, ตสฺมา ตุณฺหี, เอวเมตํ ธารยามี”ติ.}

\csromanmatikaanchor{mtk.75}
\csromantocmarkref{subsection}{สโมธานปริวาส}{mtk.75}
\bookmark[dest={mtk.75},level=2]{สโมธานปริวาส}
\csromanheader{2}{\textbf{สโมธานปริวาส}}
\proseitem{125}{เอวญฺจ ปน ภิกฺขเว ปุริมาย อาปตฺติยา สโมธานปริวาโส ทาตพฺโพ\csromandash{}เตน ภิกฺขเว อุทายินา ภิกฺขุนา สํฆํ อุปสงฺกมิตฺวา ฯเปฯ เอวมสฺส วจนีโย “อหํ ภนฺเต เอกํ อาปตฺติํ อาปชฺชิํ สญฺเจตนิกํ \csromanfolio{126}สุกฺกวิสฺสฏฺฐิํ ปกฺขปฺปฏิจฺฉนฺนํ, โสหํ สํฆํ เอกิสฺสา อาปตฺติยา สญฺเจตนิกาย สุกฺกวิสฺสฏฺฐิยา ปกฺขปฺปฏิจฺฉนฺนาย ปกฺขปริวาสํ ยาจิํ, ตสฺส เม สํโฆ เอกิสฺสา อาปตฺติยา สญฺเจตนิกาย สุกฺกวิสฺสฏฺฐิยา ปกฺขปฺปฏิจฺฉนฺนาย ปกฺขปริวาสํ อทาสิ.\csromanspacer{} โสหํ ปริวสนฺโต อนฺตรา เอกํ อาปตฺติํ อาปชฺชิํ สญฺเจตนิกํ สุกฺกวิสฺสฏฺฐิํ ปญฺจาหปฺปฏิจฺฉนฺนํ, โสหํ สํฆํ อนฺตรา เอกิสฺสา อาปตฺติยา สญฺเจตนิกาย สุกฺกวิสฺสฏฺฐิยา ปญฺจาหปฺปฏิจฺฉนฺนาย มูลายปฏิกสฺสนํ ยาจิํ, ตํ มํ สํโฆ อนฺตรา เอกิสฺสา อาปตฺติยา สญฺเจตนิกาย สุกฺกวิสฺสฏฺฐิยา ปญฺจาหปฺปฏิจฺฉนฺนาย มูลาย ปฏิกสฺสิ, โสหํ ภนฺเต สํฆํ อนฺตรา เอกิสฺสา อาปตฺติยา สญฺเจตนิกาย สุกฺกวิสฺสฏฺฐิยา ปญฺจาหปฺปฏิจฺฉนฺนาย ปุริมาย อาปตฺติยา สโมธานปริวาสํ ยาจามี”ติ.}

\prose{ทุติยมฺปิ ยาจิตพฺโพ.\csromanspacer{} ตติยมฺปิ ยาจิตพฺโพ.\csromanspacer{} พฺยตฺเตน ภิกฺขุนา ปฏิพเลน สํโฆ ญาเปตพฺโพ\csromandash{}}

\proseitem{126}{“สุณาตุ เม ภนฺเต สํโฆ, อยํ อุทายี ภิกฺขุ เอกํ อาปตฺติํ อาปชฺชิ สญฺเจตนิกํ สุกฺกวิสฺสฏฺฐิํ ปกฺขปฺปฏิจฺฉนฺนํ, โส สํฆํ เอกิสฺสา อาปตฺติยา สญฺเจตนิกาย สุกฺกวิสฺสฏฺฐิยา ปกฺขปฺปฏิจฺฉนฺนาย ปกฺขปริวาสํ ยาจิ, สํโฆ อุทายิสฺส ภิกฺขุโน เอกิสฺสา อาปตฺติยา สญฺเจตนิกาย สุกฺกวิสฺสฏฺฐิยา ปกฺขปฺปฏิจฺฉนฺนาย ปกฺขปริวาสํ อทาสิ.\csromanspacer{} โส ปริวสนฺโต อนฺตรา เอกํ อาปตฺติํ อาปชฺชิ สญฺเจตนิกํ สุกฺกวิสฺสฏฺฐิํ ปญฺจาหปฺปฏิจฺฉนฺนํ, โส สํฆํ อนฺตรา เอกิสฺสา อาปตฺติยา สญฺเจตนิกาย สุกฺกวิสฺสฏฺฐิยา ปญฺจาหปฺปฏิจฺฉนฺนาย มูลายปฏิกสฺสนํ ยาจิ, สํโฆ อุทายิํ ภิกฺขุํ อนฺตรา เอกิสฺสา อาปตฺติยา สญฺเจตนิกาย สุกฺกวิสฺสฏฺฐิยา ปญฺจาหปฺปฏิจฺฉนฺนาย มูลาย ปฏิกสฺสิ, โส สํฆํ อนฺตรา เอกิสฺสา อาปตฺติยา สญฺเจตนิกาย สุกฺกวิสฺสฏฺฐิยา ปญฺจาหปฺปฏิจฺฉนฺนาย ปุริมาย อาปตฺติยา สโมธานปริวาสํ ยาจติ.\csromanspacer{} ยทิ สํฆสฺส ปตฺตกลฺลํ, สํโฆ อุทายิสฺส ภิกฺขุโน อนฺตรา เอกิสฺสา อาปตฺติยา สญฺเจตนิกาย สุกฺกวิสฺสฏฺฐิยา ปญฺจาหปฺปฏิจฺฉนฺนาย ปุริมาย อาปตฺติํยา สโมธานปริวาสํ ทเทยฺย, เอสา ญตฺติ.}

\prose{สุณาตุ เม ภนฺเต สํโฆ, อยํ อุทายี ภิกฺขุ เอกํ อาปตฺติํ อาปชฺชิ สญฺเจตนิกํ สุกฺกวิสฺสฏฺฐิํ ปกฺขปฺปฏิจฺฉนฺนํ, โส สํฆํ เอกิสฺสา \csromanfolio{127}อาปตฺติยา สญฺเจตนิกาย สุกฺกวิสฺสฏฺฐิยา ปกฺกปฺปฏิจฺฉนฺนาย ปกฺขปริวาสํ ยาจิ, สํโฆ อุทายิสฺส ภิกฺขุโน เอกิสฺสา อาปตฺติยา สญฺเจตนิกาย สุกฺกวิสฺสฏฺฐิยา ปกฺขปฺปฏิจฺฉนฺนาย ปกฺขปริวาสํ อทาสิ.\csromanspacer{} โส ปริวสนฺโต อนฺตรา เอกํ อาปตฺติํ อาปชฺชิ สญฺเจตนิกํ สุกฺกวิสฺสฏฺฐิํ ปญฺจาหปฺปฏิจฺฉนฺนํ, โส สํฆํ อนฺตรา เอกิสฺสา อาปตฺติยา สญฺเจตนิกาย สุกฺกวิสฺสฏฺฐิยา ปญฺจาหปฺปฏิจฺฉนฺนาย มูลายปฏิกสฺสนํ ยาจิ, สํโฆ อุทายิํ ภิกฺขุํ อนฺตรา เอกิสฺสา อาปตฺติยา สญฺเจตนิกาย สุกฺกวิสฺสฏฺฐิยา ปญฺจาหปฺปฏิจฺฉนฺนาย มูลาย ปฏิกสฺสิ, โส สํฆํ อนฺตรา เอกิสฺสา อาปตฺติยา สญฺเจตนิกาย สุกฺกวิสฺสฏฺฐิยา ปญฺจาหปฺปฏิจฺฉนฺนาย ปุริมาย อาปตฺติยา สโมธานปริวาสํ ยาจติ.\csromanspacer{} สํโฆ อุทายิสฺส ภิกฺขุโน อนฺตรา เอกิสฺสา อาปตฺติยา สญฺเจตนิกาย สุกฺกวิสฺสฏฺฐิยา ปญฺจาหปฺปฏิจฺฉนฺนาย ปุริมาย อาปตฺติยา สโมธานปริวาสํ เทติ, ยสฺสายสฺมโต ขม-ติ อุทายิสฺส ภิกฺขุโน อนฺตรา เอกิสฺสา อาปตฺติยา สญฺเจตนิกาย สุกฺกวิสฺสฏฺฐิยา ปญฺจาหปฺปฏิจฺฉนฺนาย ปุริมาย อาปตฺติยา สโมธานปริวาสสฺส ทานํ, โส ตุณฺหสฺส, ยสฺส นกฺขมติ, โส ภ ภาเสยฺย.}

\prose{ทุติยมฺปิ เอตมตฺถํ วทามิ ฯเปฯ.\csromanspacer{} ตติยมฺปิ เอตมตฺถํ วทามิ ฯเปฯ.}

\prose{ทินฺโน สํเฆน อุทายิสฺส ภิกฺขุโน อนฺตรา เอกิสฺสา อาปตฺติยา สญฺเจตนิกาย สุกฺกวิสฺสฏฺฐิยา ปญฺจาหปฺปฏิจฺฉนฺนาย ปุริมาย อาปตฺติยา สโมธานปริวาโส, ขมติ สํฆสฺส, ตสฺมา ตุณฺหี, เอวเมตํ ธารยามี”ติ.}

\csromanmatikaanchor{mtk.76}
\csromantocmarkref{subsection}{มานตฺตารหมูลายปฏิกสฺสนาทิ}{mtk.76}
\bookmark[dest={mtk.76},level=2]{มานตฺตารหมูลายปฏิกสฺสนาทิ}
\csromanheader{2}{\textbf{มานตฺตารหมูลายปฏิกสฺสนาทิ}}
\proseitem{127}{โส ปริวุตฺถปริวาโส มานตฺตารโห อนฺตรา เอกํ อาปตฺติํ อาปชฺชิ สญฺเจตนิกํ สุกฺกวิสฺสฏฺฐิํ ปญฺจาหปฺปฏิจฺฉนฺนํ.\csromanspacer{} โส ภิกฺขูนํ อาโรเจสิ “อหํ อาวุโส เอกํ อาปตฺติํ อาปชฺชิํ สญฺเจตนิกํ สุกฺกวิสฺสฏฺฐิํ ปกฺขปฺปฏิจฺฉนฺนํ ฯเปฯ โสหํ ปริวุตฺถปริวาโส มานตฺตารโห อนฺตรา เอกํ อาปตฺติํ อาปชฺชิํ สญฺเจตนิกํ สุกฺกวิสฺสฏฺฐิํ ปญฺจาหปฺปฏิจฺฉนฺนํ, กถํ นุ โข มยา ปฏิปชฺชิตพฺพนฺ”ติ.\csromanspacer{} ภควโต เอตมตฺถํ อาโรเจสุํ.\csromanspacer{} เตน หิ ภิกฺขเว สํโฆ อุทายิํ ภิกฺขุํ อนฺตรา เอกิสฺสา อาปตฺติยา สญฺเจตนิกาย สุกฺกวิสฺสฏฺฐิยา ปญฺจาหปฺปฏิจฺฉนฺนาย มูลาย ปฏิกสฺสิตฺวา \csromanfolio{128}ปุริมาย อาปตฺติยา สโมธานปริวาสํ เทตุ.\csromanspacer{} เอวญฺจ ปน ภิกฺขเว มูลาย ปฏิกสฺสิตพฺโพ ฯเปฯ.}

\prose{เอวญฺจ ปน ภิกฺขเว ปุริมาย อาปตฺติยา สโมธานปริวาโส ทาตพฺโพ ฯเปฯ เทติ ฯเปฯ.}

\prose{ทินฺโน สํเฆน อุทายิสฺส ภิกฺขุโน อนฺตรา เอกิสฺสา อาปตฺติยา สญฺเจตนิกาย สุกฺกวิสฺสฏฺฐิยา ปญฺจาหปฺปฏิจฺฉนฺนาย ปุริมาย อาปตฺติยา สโมธานปริวาโส, ขมติ สํฆสฺส, ตสฺมา ตุณฺหี, เอวเมตํ ธารยามี”ติ.}

\csromanmatikaanchor{mtk.77}
\csromantocmarkref{subsection}{ติกาปตฺติมานตฺต}{mtk.77}
\bookmark[dest={mtk.77},level=2]{ติกาปตฺติมานตฺต}
\csromanheader{2}{\textbf{ติกาปตฺติมานตฺต}}
\proseitem{128}{โส ปริวุตฺถปริวาโส ภิกฺขูนํ อาโรโจสิ “อหํ อาวุโส เอกํ อาปตฺติํ อาปชฺชิํ สญฺเจตนิกํ สุกฺกวิสฺสฏฺฐิํ ปกฺขปฺปฏิจฺฉนฺนํ ฯเปฯ โสหํ ปริวุตฺถปริวาโส, กถํ นุ โข มยา ปฏิปชฺชิตพฺพนฺ”ติ.\csromanspacer{} ภควโต เอตมตฺถํ อาโรเจสุํ.\csromanspacer{} เตน หิ ภิกฺขเว สํโฆ อุทายิสฺส ภิกฺขุโน ติสฺสนฺนํ อาปตฺตีนํ ฉารตฺตํ มานตฺตํ เทตุ.\csromanspacer{} เอวญฺจ ปน ภิกฺขเว ทาตพฺพํ\csromandash{}}

\prose{เตน ภิกฺขเว อุทายิน ภิกฺขุนา สํฆํ อุปสงฺกมิตฺวา ฯเปฯ เอวมสฺส วจนีโย “อหํ ภนฺเต เอกํ อาปตฺติํ อาปชฺชิํ สญฺเจตนิกํ สุกฺกวิสฺสฏฺฐิํ ปกฺขปฺปฏิจฺฉนฺนํ ฯเปฯ โสหํ ภนฺเต ปริวุตฺถปริวาโส สํฆํ ติสฺสนฺนํ อาปตฺตีนํ ฉารตฺตํ มานตฺตํ ยาจามี”ติ.\csromanspacer{} ทุติยมฺปิ ยาจิตพฺพํ.\csromanspacer{} ตติยมฺปิ ยาจิตพฺพํ.\csromanspacer{} พฺยตฺเตน ภิกฺขุนา ปฏิพเลน สํโฆ ญาเปตพฺโพ\csromandash{}}

\proseitem{129}{“สุณาตุ เม ภนฺเต สํโฆ, อยํ อุทายี ภิกฺขุ เอกํ อาปตฺติํ อาปชฺชิ สญฺเจตนิกํ สุกฺกวิสฺสฏฺฐิํ ปกฺขปฺปฏิจฺฉนฺนํ ฯเปฯ โส ปริวุตฺถปริวาโส สํฆํ ติสฺสนฺนํ อาปตฺตีนํ ฉารตฺตํ มานตฺตํ ยาจติ.\csromanspacer{} ยทิ สํฆสฺส ปตฺตกลฺลํ, สํโฆ อุทายิสฺส ภิกฺขุโน ติสฺสนฺนํ อาปตฺตีนํ ฉารตฺตํ มานตฺตํ ทเทยฺย, เอสา ญตฺติ.}

\prose{สุณาตุ เม ภนฺเต สํโฆ, อยํ อุทายี ภิกฺขุ เอกํ อาปตฺติํ อาปชฺชิ สญฺเจตนิกํ สุกฺกวิสฺสฏฺฐิํ ปกฺขปฺปฏิจฺฉนฺนํ ฯเปฯ โส ปริวุตฺถปริวาโส สํฆํ ติสฺสนฺนํ อาปตฺตีนํ ฉารตฺตํ มานตฺตํ ยาจติ.\csromanspacer{} สํโฆ อุทายิสฺส ภิกฺขุโน ติสฺสนฺนํ อาปตฺตีนํ ฉารตฺตํ มานตฺตํ เทติ, ยสฺสายสฺมโต \csromanfolio{129}ขมติ อุทายิสฺส ภิกฺขุโน ติสฺสนฺนํ อาปตฺตีนํ ฉารตฺตํ มานตฺตํ ทานํ, โส ตุณฺหสฺส, ยสฺส นกฺขมติ, โส ภาเสยฺย.}

\prose{ทุติยมฺปิ เอตมตฺถํ วทามิ ฯเปฯ.\csromanspacer{} ตติยมฺปิ เอตมตฺถํ วทามิ ฯเปฯ.}

\prose{ทินฺนํ สํเฆน อุทายิสฺส ภิกฺขุโน ติสฺสนฺนํ อาปตฺตีนํ ฉารตฺตํ มานตฺตํ, ขมติ สํฆสฺส, ตสฺมา ตุณฺหี, เอวเมตํ ธารยามี”ติ.}

\csromanmatikaanchor{mtk.78}
\csromantocmarkref{subsection}{มานตฺตจาริกมูลายปฏิกสฺสนาทิ}{mtk.78}
\bookmark[dest={mtk.78},level=2]{มานตฺตจาริกมูลายปฏิกสฺสนาทิ}
\csromanheader{2}{\textbf{มานตฺตจาริกมูลายปฏิกสฺสนาทิ}}
\proseitem{130}{โส มานตฺตํ จรนฺโต อนฺตรา เอกํ อาปตฺติํ อาปชฺชิ สญฺเจตนิกํ สุกฺกวิสฺสฏฺฐิํ ปญฺจาหปฺปฏิจฺฉนฺนํ.\csromanspacer{} โส ภิกฺขูนํ อาโรเจสิ “อหํ อาวุโส เอกํ อาปตฺติํ อาปชฺชิํ สญฺเจตนิกํ สุกฺกวิสฺสฏฺฐิํ ปกฺขปฺปฏิจฺฉนฺนํ ฯเปฯ โสหํ มานตฺตํ จรนฺโต อนฺตรา เอกํ อาปตฺติํ อาปชฺชิํ สญฺเจตนิกํ สุกฺกวิสฺสฏฺฐิํ ปญฺจาหปฺปฏิจฺฉนฺนํ, กถํ นุ โข มยา ปฏิปชฺชิตพฺพนฺ”ติ.\csromanspacer{} ภควโต เอตมตฺถํ อาโรเจสุํ.\csromanspacer{} เตน หิ ภิกฺขเว สํโฆ อุทายิํ ภิกฺขุํ อนฺตรา เอกิสฺสา อาปตฺติยา สญฺเจตนิกาย สุกฺกวิสฺสฏฺฐิยา ปญฺจหปฺปฏิจฺฉนฺนาย มูลาย ปฏิกสฺสิตฺวา ปุริมาย อาปตฺติยา สโมธานปริวาสํ ทตฺวา ฉารตฺตํ มานตฺตํ เทตุ.\csromanspacer{} เอวญฺจ ปน ภิกฺขเว มูลาย ปฏิกสฺสิตพฺโพ ฯเปฯ.}

\prose{เอวญฺจ ปน ภิกฺขเว ปุริมาย อาปตฺติยา สโมธานปริวาโส ทาตพฺโพ ฯเปฯ.}

\prose{เอวญฺจ ปน ภิกฺขเว ฉารตฺตํ มานตฺตํ ทาตพฺพํ ฯเปฯ เทติ ฯเปฯ.}

\prose{ทินฺนํ สํเฆน อุทายิสฺส ภิกฺขุโน อนฺตรา เอกิสฺสา อาปตฺติยา สญฺเจตนิกาย สุกฺกวิสฺสฏฺฐิยา ปญฺจาหปฺปฏิจฺฉนฺนาย ฉารตฺตํ มานตฺตํ, ขมติ สํฆสฺส, ตสฺมา ตุณฺหี, เอวเมตํ ธารยามี”ติ.}

\csromanmatikaanchor{mtk.79}
\csromantocmarkref{subsection}{อพฺภานารหมูลายปฏิกสฺสนาทิ}{mtk.79}
\bookmark[dest={mtk.79},level=2]{อพฺภานารหมูลายปฏิกสฺสนาทิ}
\csromanheader{2}{\textbf{อพฺภานารหมูลายปฏิกสฺสนาทิ}}
\proseitem{131}{โส จิณฺณมานตฺโต อพฺภานารโห อนฺตรา เอกํ อาปตฺติํ อาปชฺชิ สญฺเจตนิกํ สุกฺกวิสฺสฏฺฐิํ ปญฺจาหปฺปฏิจฺฉนฺนํ.\csromanspacer{} โส ภิกฺขูนํ อาโรเจสิ “อหํ อาวุโส เอกํ อาปตฺติํ อาปชฺชิํ สญฺเจตนิกํ สุกฺกวิสฺสฏฺฐิํ ปกฺขปฺปฏิจฺฉนฺนํ ฯเปฯ โสหํ จิณฺณมานตฺโต อพฺภานารโห อนฺตรา เอกํ อาปตฺติํ อาปชฺชิํ สญฺเจตนิกํ สุกฺกวิสฺสฏฺฐิํ ปญฺจาหปฺปฏิจฺฉนฺนํ, กถํ นุ โข \csromanfolio{130}มยา ปฏิปชฺชิตพฺพนฺ”ติ.\csromanspacer{} ภควโต เอตมตฺถํ อาโรเจสุํ.\csromanspacer{} เตน หิ ภิกฺขเว สํโฆ อุทายิํ ภิกฺขุํ อนฺตรา เอกิสฺสา อาปตฺติยา สญฺเจตนิกาย สุกฺกวิสฺสฏฺฐิยา ปญฺจาหปฺปฏิจฺฉนฺนาย มูลาย ปฏิกสฺสิตฺวา ปุริมาย อาปตฺติยา สโมธานปริวาสํ ทตฺวา ฉารตฺตํ มานตฺตํ เทตุ.\csromanspacer{} เอวญฺจ ปน ภิกฺขเว มูลาย ปฏิกสฺสิตพฺโพ ฯเปฯ.}

\prose{เอวญฺจ ปน ภิกฺขเว ปุริมาย อาปตฺติยา สโมธานปริวาโส ทาตพฺโพ ฯเปฯ.}

\prose{เอวญฺจ ปน ภิกฺขเว ฉารตฺตํ มานตฺตํ ทาตพฺพํ ฯเปฯ เทติ ฯเปฯ.}

\prose{ทินฺนํ สํเฆน อุทายิสฺส ภิกฺขุโน อนฺตรา เอกิสฺสา อาปตฺติยา สญฺเจตนิกาย สุกฺกวิสฺสฏฺฐิยา ปญฺจาหปฺปฏิจฺฉนฺนาย ฉารตฺตํ มานตฺตํ, ขมติ สํฆสฺส, ตสฺมา ตุณฺหี, เอวเมตํ ธารยามี”ติ.}

\csromanmatikaanchor{mtk.80}
\csromantocmarkref{subsection}{ปกฺขปฺปฏิจฺฉนฺน-อพฺภาน}{mtk.80}
\bookmark[dest={mtk.80},level=2]{ปกฺขปฺปฏิจฺฉนฺน-อพฺภาน}
\csromanheader{2}{\textbf{ปกฺขปฺปฏิจฺฉนฺน-อพฺภาน}}
\proseitem{132}{โส จิณฺณมานตฺโต ภิกฺขูนํ อาโรเจสิ “อหํ อาวุโส เอกํ อาปตฺติํ อาปชฺชิํ สญฺเจตนิกํ สุกฺกวิสฺสฏฺฐิํ ปกฺขปฺปฏิจฺฉนฺนํ ฯเปฯ โสหํ จิณฺณมานตฺโต, กถํ นุ โข มยา ปฏิปชฺชิตพฺพนฺ”ติ.\csromanspacer{} ภควโต เอตมตฺถํ อาโรเจสุํ.\csromanspacer{} เตน หิ ภิกฺขเว สํโฆ อุทายิํ ภิกฺขุํ อพฺเภตุ.\csromanspacer{} เอวญฺจ ปน ภิกฺขเว อพฺเภตพฺโพ\csromandash{}}

\prose{เตน ภิกฺขเว อุทายินา ภิกฺขุนา สํฆํ อุปสงฺกมิตฺวา เอกํสํ อุตฺตราสงฺคํ กริตฺวา วุฑฺฒานํ ภิกฺขูนํ ปาเท วนฺทิตฺวา อุกฺกุฏิกํ นิสีทิตฺวา อญฺชลิํ ปคฺคเหตฺวา เอวมสฺส วจนีโย “อหํ ภนฺเต เอกํ อาปตฺติํ อาปชฺชิํ สญฺเจตนิกํ สุกฺกวิสฺสฏฺฐิํ ปกฺขปฺปฏิจฺฉนฺนํ, โสหํ สํฆํ เอกิสฺสา อาปตฺติยา สญฺเจตนิกาย สุกฺกวิสฺสฏฺฐิยา ปกฺขปฺปฏิจฺฉนฺนาย ปกฺขปริวาสํ ยาจิํ, ตสฺส เม สํโฆ เอกิสฺสา อาปตฺติยา สญฺเจตนิกาย สุกฺกวิสฺสฏฺฐิยา ปกฺขปฺปฏิจฺฉนฺนาย ปกฺขปริวาสํ อทาสิ.\csromanspacer{} โสหํ ปริวสนฺโต อนฺตรา เอกํ อาปตฺติํ อาปชฺชิํ สญฺเจตนิกํ สุกฺกวิสฺสฏฺฐิํ ปญฺจาหปฺปฏิจฺฉนฺนํ, โสหํ สํฆํ อนฺตรา เอกิสฺสา อาปตฺติยา สญฺเจตนิกาย สุกฺกวิสฺสฏฺฐิยา ปญฺจาหปฺปฏิจฺฉนฺนาย มูลายปฏิกสฺสนํ ยาจิํ, ตํ มํ สํโฆ อนฺตรา เอกิสฺสา อาปตฺติยา สญฺเจตนิกาย สุกฺกวิสฺสฏฺฐิยา ปญฺจาหปฺปติจฺฉนฺนาย มูลาย ปฏิกสฺสิ, โสหํ สํฆํ อนฺตรา เอกิสฺสา \csromanfolio{131}อาปตฺติยา สญฺเจตนิกาย สุกฺกวิสฺสฏฺฐิยา ปญฺจาหปฺปฏิจฺฉนฺนาย ปุริมาย อาปตฺติยา สโมธานปริวาสํ ยาจิํ, ตสฺส เม สํโฆ อนฺตรา เอกิสฺสา อาปตฺติยา สญฺเจตนิกาย สุกฺกวิสฺสฏฺฐิยา ปญฺจาหปฺปฏิจฺฉนฺนาย ปุริมาย อาปตฺติยา สโมธานปริวาสํ อทาสิ.\csromanspacer{} โสหํ ปริวุตฺถปริวาโส มานตฺตารโห อนฺตรา เอกํ อาปตฺติํ อาปชฺชิํ สญฺเจตนิกํ สุกฺกวิสฺสฏฺฐิํ ปญฺจาหปฺปฏิจฺฉนฺนํ, โสหํ สํฆํ อนฺตรา เอกิสฺสา อาปตฺติยา สญฺเจตนิกาย สุกฺกวิสฺสฏฺฐิยา ปญฺจาหปฺปฏิจฺฉนฺนาย มูลายปฏิกสฺสนํ ยาจิํ, ตํ มํ สํโฆ อนฺตรา เอกิสฺสา อาปตฺติยา สญฺเจตนิกาย สุกฺกวิสฺสฏฺฐิยา ปญฺจาหปฺปฏิจฺฉนฺนาย มูลาย ปฏิกสฺสิ, โสหํ สํฆํ อนฺตรา เอกิสฺสา อาปตฺติยา สญฺเจตนิกาย สุกฺกวิสฺสฏฺฐิยา ปญฺจาหปฺปฏิจฺฉนฺนาย ปุริมาย อาปตฺติยา สโมธานปริวาสํ ยาจิํ, ตสฺส เม สํโฆ อนฺตรา เอกิสฺสา อาปตฺติยา สญฺเจตนิกาย สุกฺกวิสฺสฏฺฐิยา ปญฺจาหปฺปฏิจฺฉนฺนาย ปุริมาย อาปตฺติยา สโมธานปริวาสํ อทาสิ.\csromanspacer{} โสหํ ปริวุตฺถปริวาโส สํฆํ ติสฺสนฺนํ อาปตฺตีนํ ฉารตฺตํ มานตฺตํ ยาจิํ, ตสฺส เม สํโฆ ติสฺสนฺนํ อาปตฺตีนํ ฉารตฺตํ มานตฺตํ อทาสิ.\csromanspacer{} โสหํ มานตฺตํ จรนฺโต อนฺตรา เอกํ อาปตฺติํ อาปชฺชิํ สญฺเจตนิกํ สุกฺกวิสฺสฏฺฐิํ ปญฺจาหปฺปฏิจฺฉนฺนํ, โสหํ สํฆํ อนฺตรา เอกิสฺสา อาปตฺติยา สญฺเจตนิกาย สุกฺกวิสฺสฏฺฐิยา ปญฺจาหปฺปฏิจฺฉนฺนาย มูลายปฏิกสฺสนํ ยาจิํ, ตํ มํ สํโฆ อนฺตรา เอกิสฺสา อาปตฺติยา สญฺเจตนิกาย สุกฺกวิสฺสฏฺฐิยา ปญฺจาหปฺปฏิจฺฉนฺนาย มูลาย ปฏิกสฺสิ, โสหํ สํฆํ อนฺตรา เอกิสฺสา อาปตฺติยา สญฺเจกนิกาย สุกฺกวิสฺสฏฺฐิยา ปญฺจาหปฺปฏิจฺฉนฺนาย ปุริมาย อาปตฺติยา สโมธานปริวาสํ ยาจิํ, ตสฺส เม สํโฆ อนฺตรา เอกิสฺสา อาปตฺติยา สญฺเจตนิกาย สุกฺกวิสฺสฏฺฐิยา ปญฺจาหปฺปฏิจฺฉนฺนาย ปุริมาย อาปตฺติยา สโมธานปริวาสํ อทาสิ.\csromanspacer{} โสหํ ปริวุตฺถปริวาโส สํฆํ อนฺตรา เอกิสฺสา อาปตฺติยา สญฺเจตนิกาย สุกฺกวิสฺสฏฺฐิยา ปญฺจาหปฺปฏิจฺฉนฺนาย ฉารตฺตํ มานตฺตํ ยาจิํ, ตสฺส เม สํโฆ อนฺตรา เอกิสฺสา อาปตฺติยา สญฺเจตนิกาย สุกฺกวิสฺสฏฺฐิยา ปญฺจาหปฺปฏิจฺฉนฺนาย ฉารตฺตํ มานตฺตํ อทาสิ.\csromanspacer{} โสหํ จิณฺณมานตฺโต อพฺภานารโห อนฺตรา เอกํ อาปตฺติํ อาปชฺชิํ สญฺเจตนิกํ สุกฺกวิสฺสฏฺฐิํ ปญฺจาหปฺปฏิจฺฉนฺนํ, โสหํ สํฆํ อนฺตรา เอกิสฺสา อาปตฺติยา สญฺเจตนิกาย สุกฺกวิสฺสฏฺฐิยา ปญฺจาหปฺปฏิจฺฉนฺนาย มูลายปฏิกสฺสนํ ยาจิํ, ตํ มํ สํโฆ อนฺตรา เอกิสฺสา อาปตฺติยา สญฺเจตนิกาย สุกฺกวิสฺสฏฺฐิยา \csromanfolio{132}ปญฺจาหปฺปฏิจฺฉนฺนาย มูลาย ปฏิกสฺสิ, โสหํ สํฆํ อนฺตรา เอกิสฺสา อาปตฺติยา สญฺเจตนิกาย สุกฺกวิสฺสฏฺฐิยา ปญฺจาหปฺปฏิจฺฉนฺนาย ปุริมาย อาปตฺติยา สโมธานปริวาสํ ยาจิํ, ตสฺส เม สํโฆ อนฺตรา เอกิสฺสา อาปตฺติยา สญฺเจตนิกาย สุกฺกวิสฺสฏฺฐิยา ปญฺจาหปฺปฏิจฺฉนฺนาย ปุริมาย อาปตฺติยา สโมธานปริวาสํ อทาสิ, โสหํ ปริวุตฺถปริวาโส สํฆํ อนฺตรา เอกิสฺสา อาปตฺติยา สญฺเจตนิกาย สุกฺกวิสฺสฏฺฐิยา ปญฺจาหปฺปฏิจฺฉนฺนาย ฉารตฺตํ มานตฺตํ ยาจิํ, ตสฺส เม สํโฆ อนฺตรา เอกิสฺสา อาปตฺติยา สญฺเจตนิกาย สุกฺกวิสฺสฏฺฐิยา ปญฺจาหปฺปฏิจฺฉนฺนาย ฉารตฺตํ มานตฺตํ อทาสิ, โสหํ ภนฺเต จิณฺณมานตฺโต สํฆํ อพฺภานํ ยาจามี”ติ.}

\prose{ทุติยมฺปิ ยาจิตพฺพํ.\csromanspacer{} ตติยมฺปิ ยาจิตพฺพํ.\csromanspacer{} พฺยตฺเตน ภิกฺขุนา ปฏิพเลน สํโฆ ญาเปตพฺโพ\csromandash{}}

\proseitem{133}{“สุณาตุ เม ภนฺเต สํโฆ, อยํ อุทายี ภิกฺขุ เอกํ อาปตฺติํ อาปชฺชิ สญฺเจตนิกํ สุกฺกวิสฺสฏฺฐิํ ปกฺขปฺปฏิจฺฉนฺนํ, โส สํฆํ เอกิสฺสา อาปตฺติยา สญฺเจตนิกาย สุกฺกวิสฺสฏฺฐิยา ปกฺขปฺปฏิจฺฉนฺนาย ปกฺขปริวาสํ ยาจิ, สํโฆ อุทายิสฺส ภิกฺขุโน เอกิสฺสา อาปตฺติยา สญฺเจตนิกาย สุกฺกวิสฺสฏฺฐิยา ปกฺขปฺปฏิจฺฉนฺนาย ปกฺขปริวาสํ อทาสิ.\csromanspacer{} โส ปริวสนฺโต อนฺตรา เอกํ อาปตฺติํ อาปชฺชิ สญฺเจตนิกํ สุกฺกวิสฺสฏฺฐิํ ปญฺจาหปฺปฏิจฺฉนฺนํ, โส สํฆํ อนฺตรา เอกิสฺสา อาปตฺติยา สญฺเจตนิกาย สุกฺกวิสฺสฏฺฐิยา ปญฺจาหปฺปฏิจฺฉนฺนาย มูลายปฏิกสฺสนํ ยาจิ, สํโฆ อุทายิํ ภิกฺขุํ อนฺตรา เอกิสฺสา อาปตฺติยา สญฺเจตนิกาย สุกฺกวิสฺสฏฺฐิยา ปญฺจาหปฺปฏิจฺฉนฺนาย มูลาย ปฏิกสฺสิ, โส สํฆํ อนฺตรา เอกิสฺสา อาปตฺติยา สญฺเจตนิกาย สุกฺกวิสฺสฏฺฐิยา ปญฺจาหปฺปฏิจฺฉนฺนาย ปุริมาย อาปตฺติยา สโมธานปริวาสํ ยาจิ, สํโฆ อุทายิสฺส ภิกฺขุโน อนฺตรา เอกิสฺสา อาปตฺติยา สญฺเจตนิกาย สุกฺกวิสฺสฏฺฐิยา ปญฺจาหปฺปฏิจฺฉนฺนาย ปุริมาย อาปตฺติยา สโมธานปริวาสํ อทาสิ.\csromanspacer{} โส ปริวุตฺถปริวาโส มานตฺตารโห อนฺตรา เอกํ อาปตฺติํ อาปชฺชิ สญฺเจตนิกํ สุกฺกวิสฺสฏฺฐิํ ปญฺจหปฺปฏิจฺฉนฺนํ, โส สํฆํ อนฺตรา เอกิสฺสา อาปตฺติยา สญฺเจตนิกาย สุกฺกวิสฺสฏฺฐิยา ปญฺจาหปฺปฏิจฺฉนฺนาย มูลายปฏิกสฺสนํ ยาจิ, สํโฆ อุทายิํ ภิกฺขุํ อนฺตรา เอกิสฺสา อาปตฺติยา สญฺเจตนิกาย สุกฺกวิสฺสฏฺฐิยา ปญฺจหปฺปฏิจฺฉนฺนาย มูลาย ปฏิกสฺสิ, โส สํฆํ อนฺตรา \csromanfolio{133}เอกิสฺสา อาปตฺติยา สญฺเจตนิกาย สุกฺกวิสฺสฏฺฐิยา ปญฺจาหปฺปฏิจฺฉนฺนาย ปุริมาย อาปตฺติยา สโมธานปริวาสํ ยาจิ, สํโฆ อุทายิสฺส ภิกฺขุโน อนฺตรา เอกิสฺสา อาปตฺติยา สญฺเจตนิกาย สุกฺกวิสฺสฏฺฐิยา ปญฺจาหปฺปฏิจฺฉนฺนาย ปุริมาย อาปตฺติยา สโมธานปริวาสํ อทาสิ.\csromanspacer{} โส ปริวุตฺถปริวาโส สํฆํ ติสฺสนฺนํ อาปตฺตีนํ ฉารตฺตํ มานตฺตํ ยจิ, สํโฆ อุทายิสฺส ภิกฺขุโน ติสฺสนฺนํ อาปตฺตีนํ ฉารตฺตํ มานตฺตํ อทาสิ.\csromanspacer{} โส มานตฺตํ จรนฺโต อนฺตรา เอกํ อาปตฺติํ อาปชฺชิ สญฺเจตนิกํ สุกฺกวิสฺสฏฺฐิํ ปญฺจาหปฺปฏิจฺฉนฺนํ, โส สํฆํ อนฺตรา เอกิสฺสา อาปตฺติยา สญฺเจตนิกาย สุกฺกวิสฺสฏฺฐิยา ปญฺจาหปฺปฏิจฺฉนฺนาย มูลายปฏิกสฺสนํ ยาจิ, สํโฆ อุทายิํ ภิกฺขุํ อนฺตรา เอกิสฺสา อาปตฺติยา สญฺเจตนิกาย สุกฺกวิสฺสฏฺฐิยา ปญฺจาหปฺปฏิจฺฉนฺนาย มูลาย ปฏิกสฺสิ, โส สํฆํ อนฺตรา เอกิสฺสา อาปตฺติยา สญฺเจตนิกาย สุกฺกวิสฺสฏฺฐิยา ปญฺจาหปฺปฏิจฺฉนฺนาย ปุริมาย อาปตฺติยา สโมธานปริวาสํ ยาจิ, สํโฆ อุทายิสฺส ภิกฺขุโน อนฺตรา เอกิสฺสา อาปตฺติยา สญฺเจตนิกาย สุกฺกวิสฺสฏฺฐิยา ปญฺจาหปฺปฏิจฺฉนฺนาย ปุริมาย อาปตฺติยา สโมธานปริวาสํ อทาสิ, โส ปริวุตฺถปริวาโส สํฆํ\footnote{โส สํฆํ (ก)} อนฺตรา เอกิสฺสา อาปตฺติยา สญฺเจตนิกาย สุกฺกวิสฺสฏฺฐิยา ปญฺจาหปฺปฏิจฺฉนฺนาย ฉารตฺตํ มานตฺตํ ยาจิ, สํโฆ อุทายิสฺส ภิกฺขุโน อนฺตรา เอกิสฺสา อาปตฺติยา สญฺเจตนิกาย สุกฺกวิสฺสฏฺฐิยา ปญฺจาหปฺปฏิจฺฉนฺนาย ฉารตฺตํ มานตฺตํ อทาสิ.\csromanspacer{} โส จิณฺณมานตฺโต อพฺภานารโห อนฺตรา เอกํ อาปตฺติํ อาปชฺชิ สญฺเจตนิกํ สุกฺกวิสฺสฏฺฐิํ ปญฺจาหปฺปฏิจฺฉนฺนํ, โส สํฆํ อนฺตรา เอกิสฺสา อาปตฺติยา สญฺเจตนิกาย สุกฺกวิสฺสฏฺฐิยา ปญฺจาหปฺปฏิจฺฉนฺนาย มูลายปฏิกสฺสนํ ยาจิ, สํโฆ อุทายิํ ภิกฺขุํ อนฺตรา เอกิสฺสา อาปตฺติยา สญฺเจตนิกาย สุกฺกวิสฺสฏฺฐิยา ปญฺจาหปฺปฏิจฺฉนฺนาย มูลาย ปฏิกสฺสิ, โส สํฆํ อนฺตรา เอกิสฺสา อาปตฺติยา สญฺเจตนิกาย สุกฺกวิสฺสฏฺฐิยา ปญฺจาหปฺปฏิจฺฉนฺนาย ปุริมาย อาปตฺติยา สโมธานปริวาสํ ยาจิ, สํโฆ อุทายิสฺส ภิกฺขุโน อนฺตรา เอกิสฺสา อาปตฺติยา สญฺเจตนิกาย สุกฺกวิสฺสฏฺฐิยา ปญฺจาหปฺปฏิจฺฉนฺนาย ปุริมาย อาปตฺติยา สโมธานปริวาสํ อทาสิ, โส ปริวุตฺถปริวาโส สํฆํ อนฺตรา สโมธานปริวาสํ อทาสิ, โส ปริวุตฺถปริวาโส สํฆํ อนฺตรา เอกิสฺสา อาปตฺติยา สญฺเจตนิกาย สุกฺกวิสฺสฏฺฐิยา ปญฺจาหปฺปฏิจฺฉนฺนาย \csromanfolio{134}ฉารตฺตํ มานตฺตํ ยาจิ, สํโฆ อุทายิสฺส ภิกฺขุโน อนฺตรา เอกิสฺสา อาปตฺติยา สญฺเจตนิกาย สุกฺกวิสฺสฏฺฐิยา ปญฺจาหปฺปฏิจฺฉนฺนาย ฉารตฺตํ มานตฺตํ อทาสิ, โส จิณฺณมานตฺโต สํฆํ อพฺภานํ ยาจติ.\csromanspacer{} ยทิ สํฆสฺส ปตฺตกลฺลํ, สํโฆ อุทายิํ อพฺเภยฺย, เอสา ญตฺติ.}

\prose{สุณาตุ เม ภนฺเต สํโฆ, อยํ อุทายี ภิกฺขุ เอกํ อาปตฺติํ อาปชฺชิ สญฺเจตนิกํ สุกฺกวิสฺสฏฺฐิํ ปกฺขปฺปฏิจฺฉนฺนํ ฯเปฯ โส จิณฺณมานตฺโต สํฆํ อพฺภานํ ยาจติ.\csromanspacer{} สํโฆ อุทายิํ ภิกฺขุํ อพฺเภติ, ยสฺสายสฺมโต ขมติ อุทายิสฺส ภิกฺขุโน อพฺภานํ, โส ตุณฺหสฺส, ยสฺส นกฺขมติ, โส ภาเสยฺย.}

\prose{ทุติยมฺปิ เอตมตฺตํ วทามิ ฯเปฯ.\csromanspacer{} ตติยมฺปิ เอตมตฺถํ วทามิ ฯเปฯ.}

\prose{อพฺภิโต สํเฆน อุทายี ภิกฺขุ, ขมติ สํฆสฺส, ตสฺมา ตุณฺหี, เอวเมตํ ธารยามี”ติ.}

\vspace{\tipitakaparskip}
\csromancenter{สุกฺกวิสฺสฏฺฐิ สมตฺตา.}
\csromanmatikaanchor{mtk.81}
\csromantocmarkref{section}{2. ปริวาส}{mtk.81}
\bookmark[dest={mtk.81},level=1]{2. ปริวาส}
\csromanheader{1}{2. ปริวาส}
\csromanmatikaanchor{mtk.82}
\csromantocmarkref{subsection}{อคฺฆสโมธานปริวาส}{mtk.82}
\bookmark[dest={mtk.82},level=2]{อคฺฆสโมธานปริวาส}
\csromanheader{2}{\textbf{อคฺฆสโมธานปริวาส}}
\proseitem{134}{เตน โข ปน สมเยน อญฺญตโร ภิกฺขุ สมฺพหุลา สํฆาทิเสสา อาปตฺติโย อาปนฺโน โหติ, เอกา อาปตฺติ เอกาหปฺปฏิจฺฉนฺนา, เอกา อาปตฺติ ทฺวีหปฺปฏิจฺฉนฺนา, เอกา อาปตฺติ ตีหปฺปฏิจฺฉนฺนา, เอกา อาปตฺติ จตูหปฺปฏิจฺฉนฺนา, เอกา อาปตฺติ ปญฺจาหปฺปฏิจฺฉนฺนา, เอกา อาปตฺติ ฉาหปฺปฏิจฺฉนฺนา, เอกา อาปตฺติ สตฺตาหปฺปฏิจฺฉนฺนา, เอกา อาปตฺติ อฏฺฐาหปฺปฏิจฺฉนฺนา, เอกา อาปตฺติ นวาหปฺปฏิจฺฉนฺนา, เอกา อาปตฺติ ทสาหปฺปฏิจฺฉนฺนา.\csromanspacer{} โส ภิกฺขูนํ อาโรเจสิ “อหํ อาวุโส สมฺพหุลา สํฆาทิเสสา อาปตฺติโย อาปชฺชิํ, เอกา อาปตฺติ เอกาหปฺปฏิจฺฉนฺนา ฯเปฯ เอกา อาปตฺติ ทสาหปฺปฏิจฺฉนฺนา, กถํ นุ โข มยา ปฏิปชฺชิตพฺพนฺ”ติ.\csromanspacer{} ภควโต เอตมตฺถํ อาโรเจสุํ.\csromanspacer{} เตน หิ ภิกฺขเว สํโฆ ตสฺส ภิกฺขุโน ตาสํ อาปตฺตีนํ ยา อาปตฺติ ทสาหปฺปฏิจฺฉนฺนา, ตสฺสา อคฺเฆน สโมธานปริวาสํ เทตุ.\csromanspacer{} เอวญฺจ ปน ภิกฺขเว ทาตพฺโพ\csromandash{}}

\prose{\csromanfolio{135}เตน ภิกฺขเว ภิกฺขุนา สํฆํ อุปสงฺกมิตฺวา ฯเปฯ เอวมสฺส วจนีโย “อหํ ภนฺเต สมฺพหุลา สํฆาทิเสสา อาปตฺติโย อาปชฺชิํ, เอกา อาปตฺติ เอกาหปฺปฏิจฺฉนฺนา ฯเปฯ เอกา อาปตฺติ ทสาหปฺปฏิจฺฉนฺนา, โสหํ ภนฺเต สํฆํ ตาสํ อาปตฺตีนํ ยา อาปตฺติ ทสาหปฺปฏิจฺฉนฺนา, ตสฺสา อคฺเฆน สโมธานปริวาสํ ยาจามี”ติ.\csromanspacer{} ทุติยมฺปิ ยาจิตพฺโพ.\csromanspacer{} ตติยมฺปิ ยาจิตพฺโพ.\csromanspacer{} พฺยตฺเตน ภิกฺขุนา ปฏิพเลน สํโฆ ญาเปตพฺโพ\csromandash{}}

\proseitem{135}{“สุณาตุ เม ภนฺเต สํโฆ, อยํ อิตฺถนฺนาโม ภิกฺขุ สมฺพหุลา สํฆาทิเสสา อาปตฺติโย อาปชฺชิ, เอกา อาปตฺติ เอกาหปฺปฏิจฺฉนฺนา ฯเปฯ เอกา อาปตฺติ ทสาหปฺปฏิจฺฉนฺนา, โส สํฆํ ตาสํ อาปตฺตีนํ ยา อาปตฺติ ทสาหปฺปฏิจฺฉนฺนา, ตสฺสา อคฺเฆน สโมธานปริวาสํ ยาจติ.\csromanspacer{} ยทิ สํฆสฺส ปตฺตกลฺลํ, สํโฆ อิตฺถนฺนามสฺส ภิกฺขุโน ตาสํ อาปตฺตีนํ ยา อาปตฺติ ทสาหปฺปฏิจฺฉนฺนา, ตสฺสา อคฺเฆน สโมธานปริวาสํ ทเทยฺย, เอสา ญตฺติ.}

\prose{สุณาตุ เม ภนฺเต สํโฆ, อยํ อิตฺถนฺนาโม ภิกฺขุ สมฺพหุลา สํฆาทิเสสา อาปตฺติโย อาปชฺชิ, เอกา อาปตฺติ เอกาหปฺปฏิจฺฉนฺนา ฯเปฯ เอกา อาปตฺติ ทสาหปฺปฏิจฺฉนฺนา, โส สํฆํ ตาสํ อาปตฺตีนํ ยา อาปตฺติ ทสาหปฺปฏิจฺฉนฺนา, ตสฺสา อคฺเฆน สโมธานปริวาสํ ยาจติ.\csromanspacer{} สํโฆ อิตฺถนฺนามสฺส ภิกฺขุโน ตาสํ อาปตฺตีนํ ยา อาปตฺติ ทสาหปฺปฏิจฺฉนฺนา, ตสฺสา อคฺเฆน สโมธานปริวาสํ เทติ, ยสฺสายสฺมโต ขมติ อิตฺถนฺนามสฺส ภิกฺขุโน ตาสํ อาปตฺตีนํ ยา อาปตฺติ ทสาหปฺปฏิจฺฉนฺนา, ตสฺสา อคฺเฆน สโมธานปริวาสสฺส ทานํ, โส ตุณฺหสฺส, ยสฺส นกฺขมติ, โส ภาเสยฺย.}

\prose{ทุติยมฺปิ เอตมตฺถํ วทามิ ฯเปฯ.\csromanspacer{} ตติยมฺปิ เอตมตฺถํ วทามิ ฯเปฯ.}

\prose{ทินฺโน สํเฆน อิตฺถนฺนามสฺส ภิกฺขุโน ตาสํ อาปตฺตีนํ ยา อาปตฺติ ทสาหปฺปฏิจฺฉนฺนา, ตสฺสา อคฺเฆน สโมธานปริวาโส, ขมติ สํฆสฺส, ตสฺมา ตุณฺหี, เอวเมตํ ธารยามี”ติ.}

\csromanmatikaanchor{mtk.83}
\csromantocmarkref{subsection}{สพฺพจิรปฺปฏิจฺฉนฺน-อคฺฆสโมธาน}{mtk.83}
\bookmark[dest={mtk.83},level=2]{สพฺพจิรปฺปฏิจฺฉนฺน-อคฺฆสโมธาน}
\csromanheader{2}{\textbf{สพฺพจิรปฺปฏิจฺฉนฺน-อคฺฆสโมธาน}}
\proseitem{136}{เตน โข ปน สมเยน อญฺญตโร ภิกฺขุ สมฺพหุลา สํฆาทิเสสา อาปตฺติโย อาปนฺโน โหติ, เอกา อาปตฺติ \csromanfolio{136}เอกาหปฺปฏิจฺฉนฺนา, เทฺว อาปตฺติโย ทฺวีหปฺปฏิจฺฉนฺนาโย\footnote{ทฺวีหปฺปฏิจฺฉนฺนา (ก, เอวํ ยาวทสาหปฺปฏิจฺฉนฺนา)}, ติสฺโส อาปตฺติโย ตีหปฺปฏิจฺฉนฺนาโย, จตสฺโส อาปตฺติโย จตูหปฺปฏิจฺฉนฺนาโย, ปญฺจ อาปตฺติโย ปญฺจาหปฺปฏิจฺฉนฺนาโย, ฉ อาปตฺติโย ฉาหปฺปฏิจฺฉนฺนาโย, สตฺต อาปตฺติโย สตฺตาหปฺปฏิจฺฉนฺนาโย, อฏฺฐ อาปตฺติโย อฏฺฐาหปฺปฏิจฺฉนฺนาโย, นว อาปตฺติโย นวาหปฺปฏิจฺฉนฺนาโย, ทส อาปตฺติโย ทสาหปฺปฏิจฺฉนฺนาโย, โส ภิกฺขูนํ อาโรเจสิ “อหํ อาวุโส สมฺพหุลา สํฆาทิเสสา อาปตฺติโย อาปชฺชิํ, เอกา อาปตฺติ เอกาหปฺปฏิจฺฉนฺนา ฯเปฯ ทส อาปตฺติโย ทสาหปฺปฏิจฺฉนฺนาโย, กถํ นุ โข มยา ปฏิปชฺชิตพฺพนฺ”ติ.\csromanspacer{} ภควโต เอตมตฺถํ อาโรเจสุํ.\csromanspacer{} เตน หิ ภิกฺขเว สํโฆ ตสฺส ภิกฺขุโน ตาสํ อาปตฺตีนํ ยา อาปตฺติโย สพฺพจิรปฺปฏิจฺฉนฺนาโย, ตาสํ อคฺเฆน สโมธานปริวาสํ เทตุ.\csromanspacer{} เอวญฺจ ปน ภิกฺขเว ทาตพฺโพ\csromandash{}}

\prose{เตน ภิกฺขเว ภิกฺขุนา สํฆํ อุปสงฺกมิตฺวา ฯเปฯ เอวมสฺส วจนีโย “อหํ ภนฺเต สมฺพหุลา สํฆาทิเสสา อาปตฺติโย อาปชฺชิํ, เอกา อาปตฺติ เอกาหปฺปฏิจฺฉนฺนา ฯเปฯ ทส อาปตฺติโย ทสาหปฺปฏิจฺฉนฺนาโย, โสหํ ภนฺเต สํฆํ ตาสํ อาปตฺตีนํ ยา อาปตฺติโย สพฺพจิรปฺปฏิจฺฉานฺนาโย, ตาสํ อคฺเฆน สโมธานปริวาสํ ยาจามี”ติ.\csromanspacer{} ทุติยมฺปิ ยาจิตพฺโพ.\csromanspacer{} ตติยมฺปิ ยาจิตพฺโพ.\csromanspacer{} พฺยตฺเตน ภิกฺขุนา ปฏิพเลน ปฏิพเลน สํโฆ ญาเปตพฺโพ\csromandash{}}

\proseitem{137}{“สุณาตุ เม ภนฺเต สํโฆ, อยํ อิตฺถนฺนาโม ภิกฺขุ สมฺพหุลา สํฆาทิเสสา อาปตฺติโย อาปชฺชิ, เอกา อาปตฺติ เอกาหปฺปฏิจฺฉนฺนา ฯเปฯ ทส อาปตฺติโย ทสาหปฺปฏิจฺฉนฺนาโย, โส สํฆํ ตาสํ อาปตฺตีนํ ยา อาปตฺติโย สพฺพจิรปฺปฏิจฺฉนฺนาโย, ตาสํ อคฺเฆน สโมธานปริวาสํ ยาจติ.\csromanspacer{} ยทิ สํฆสฺส ปตฺตกลฺลํ, สํโฆ อิตฺถนฺนามสฺส ภิกฺขุโน ตาสํ อาปตฺตีนํ ยา อาปตฺติโย สพฺพจิรปฺปฏิจฺฉนฺนาโย, ตาสํ อคฺเฆน สโมธานปริวาสํ ทเทยฺย, เอสา ญตฺติ.}

\prose{สุณาตุ เม ภนฺเต สํโฆ, อยํ อิตฺถนฺนาโม ภิกฺขุ สมฺพหุลา สํฆาทิเสสา อาปตฺติโย อาปชฺชิ, เอกา อาปตฺติ เอกาหปฺปฏิจฺฉนฺนา.\csromanspacer{} ทส อาปตฺติโย ทสาหปฺปฏิจฺฉนฺนาโย, โส สํฆํ ตาสํ อาปตฺตีนํ \csromanfolio{137}ยา อาปตฺติโย สพฺพจิรปฺปฏิจฺฉนฺนาโย, ตาสํ อคฺเฆน สโมธานปริวาสํ ยาจติ.\csromanspacer{} สํโฆ อิตฺถนฺนามสฺส ภิกฺขุโน ตาสํ อาปตฺตีนํ ยา อาปตฺติโย สพฺพจิรปฺปฏิจฺฉนฺนาโย, ตาสํ อคฺเฆน สโมธานปริวาสํ เทติ, ยสฺสายสฺมโต ขมติ อิตฺถนฺนามสฺส ภิกฺขุโน ตาสํ อาปตฺตีนํ ยา อาปตฺติโย สพฺพจิรปฺปฏิจฺฉนฺนาโย, ตาสํ อคฺเฆน สโมธานปริวาสสฺส ทานํ, โส ตุณฺหสฺส, ยสฺส นกฺขมติ, โส ภาเสยฺย.}

\prose{ทุติยมฺปิ เอตมตฺถํ วทามิ ฯเปฯ.\csromanspacer{} ตติยมฺปิ เอตมตฺถํ วทามิ ฯเปฯ.}

\prose{ทินฺโน สํเฆน อิตฺถนฺนามสฺส ภิกฺขุโน ตาสํ อาปตฺตีนํ ยา อาปตฺติโย สพฺพจิรปฺปฏิจฺฉนฺนาโย, ตาสํ อคฺเฆน สโมธานปริวาโส, ขมติ สํฆสฺส, ตสฺมา ตุณฺหี, เอวเมตํ ธารยามี”ติ.}

\csromanmatikaanchor{mtk.84}
\csromantocmarkref{subsection}{เทฺวมาสปริวาส}{mtk.84}
\bookmark[dest={mtk.84},level=2]{เทฺวมาสปริวาส}
\csromanheader{2}{\textbf{เทฺวมาสปริวาส}}
\proseitem{138}{เตน โข ปน สมเยน อญฺญตโร ภิกฺขุ เทฺว สํฆาทิเสสา อาปตฺติโย อาปนฺโน โหติ เทฺวมาสปฺปฏิจฺฉนฺนาโย.\csromanspacer{} ตสฺส เอตทโหสิ “อหํ โข เทฺว สํฆาทิเสสา อาปตฺติโย อาปชฺชิํ เทฺวมาสปฺปฏิจฺฉนฺนาโย, ยํนูนาหํ สํฆํ เอกิสฺสา อาปตฺติยา เทฺวมาสปฺปฏิจฺฉนฺนาย เทฺวมาสปริวาสํ ยาเจยฺยนฺ”ติ.\csromanspacer{} โส สํฆํ เอกิสฺสา อาปตฺติยา เทฺวมาสปฺปฏิจฺฉนฺนาย เทฺวมาสปริวาสํ ยาจิ, ตสฺส สํโฆ เอกิสฺสา อาปตฺติยา เทฺวมาสปฺปฏิจฺฉนฺนาย เทฺวมาสปริวาสํ อทาสิ.\csromanspacer{} ตสฺส ปริวสนฺตสฺส ลชฺชีธมฺโม โอกฺกมิ\csromandash{}อหํ โข เทฺว สํฆาทิเสสา อาปตฺติโย อาปชฺชิํ เทฺวมาสปฺปฏิจฺฉนฺนาโย, ตสฺส เม เอตทโหสิ “อหํ โข เทฺว สํฆาทิเสสา อาปตฺติโย อาปชฺชิํ เทฺวมาสปฺปฏิจฺฉนฺนาโย, ยํนูนาหํ สํฆํ เอกิสฺสา อาปตฺติยา เทฺวมาสปฺปฏิจฺฉนฺนาย เทฺวมาสปริวาสํ ยาเจยฺยนฺ”ติ, โสหํ สํฆํ เอกิสฺสา อาปตฺติยา เทฺวมาสปฺปฏิจฺฉนฺนาย เทฺวมาสปริวาสํ ยาจิํ, ตสฺส เม สํโฆ เอกิสฺสา อาปตฺติยา เทฺวมาสปฺปฏิจฺฉนฺนาย เทฺวมาสปริวาสํ อทาสิ, ตสฺส เม ปริวสนฺตสฺส ลชฺชีธมฺโม โอกฺกมิ “ยํนูนาหํ สํฆํ อิตริสฺสาปิ อาปตฺติยา เทฺวมาสปฺปฏิจฺฉนฺนาย เทฺวมาสปริวาสํ ยาเจยฺยนฺ”ติ.}

\prose{โส ภิกฺขูนํ อาโรเจสิ “อหํ อาวุโส เทฺว สํฆาทิเสสา อาปตฺติโย อาปชฺชิํ เทฺวมาสปฺปฏิจฺฉนฺนาโย, ตสฺส เม เอตทโหสิ \csromanfolio{138}“อหํ โข เทฺว สํฆาทิเสสา อาปตฺติโย อาปชฺชิํ เทฺวมาสปฺปฏิจฺฉนฺนาโย, ยํนูนาหํ สํฆํ เอกิสฺสา อาปตฺติยา เทฺวมาสปฺปฏิจฺฉนฺนาย เทฺวมาสปริวาสํ ยาเจยฺยนฺ”ติ.\csromanspacer{} โสหํ สํฆํ เอกิสฺสา อาปตฺติยา เทฺวมาสปฺปฏิจฺฉนฺนาย เทฺวมาสปริวาสํ ยาจิํ, ตสฺส เม สํโฆ เอกิสฺสา อาปตฺติยา เทฺวมาสปฺปฏิจฺฉนฺนาย เทฺวมาสปริวาสํ อทาสิ, ตสฺส เม ปริวสนฺตสฺส ลชฺชีธมฺโม โอกฺกมิ\csromandash{}อหํ โข เทฺว สํฆาทิเสสา อาปตฺติโย อาปชฺชิํ เทฺวมาสปฺปฏิจฺฉนฺนาโย, ตสฺส เม เอตทโหสิ “อหํ โข เทฺว สํฆาทิเสสา อาปตฺติโย อาปชฺชิํ เทฺวมาสปฺปฏิจฺฉนฺนาโย, ยํนูนาหํ สํฆํ เอกิสฺสา อาปตฺติยา เทฺวมาสปฺปฏิจฺฉนฺนาย เทฺวมาสปริวาสํ ยาเจยฺยนฺ”ติ, โสหํ สํฆํ เอกิสฺสา อาปตฺติยา เทฺวมาสปฺปฏิจฺฉนฺนาย เทฺวมาสปริวาสํ ยาจิํ, ตสฺส เม สํโฆ เอกิสฺสา อาปตฺติยา เทฺวมาสปฺปฏิจฺฉนฺนาย เทฺวมาสปริวาสํ อทาสิ, ตสฺส เม ปริวสนฺตสฺส ลชฺชีธมฺโม โอกฺกมิ “ยํนูนาหํ สํฆํ อิตริสฺสาปิ อาปตฺติยา เทฺวมาสปฺปฏิจฺฉนฺนาย เทฺวมาสปริวาสํ ยาเจยฺยนฺ”ติ.\csromanspacer{} กถํ นุ โข มยา ปฏิปชฺชิตพฺพนฺติ.\csromanspacer{} ภควโต เอตมตฺถํ อาโรเจสุํ.\csromanspacer{} เตน หิ ภิกฺขเว สํโฆ ตสฺส ภิกฺขุโน อิตริสฺสาปิ อาปตฺติยา เทฺวมาสปฺปฏิจฺฉนฺนาย เทฺวมาสปริวาสํ เทตุ.\csromanspacer{} เอวญฺจ ปน ภิกฺขเว ทาตพฺโพ\csromandash{}}

\prose{เตน ภิกฺขเว ภิกฺขุนา สํฆํ อุปสงฺกมิตฺวา เอกํสํ อุตฺตราสงฺคํ กริตฺวา วุฑฺฒานํ ภิกฺขูนํ ปาเท วนฺทิตฺวา อุกฺกุฏิกํ นิสีทิตฺวา อญฺชลิํ ปคฺคเหตฺวา เอวมสฺส วจนีโย “อหํ ภนฺเต เทฺว สํฆาทิเสสา อาปตฺติโย อาปชฺชิํ เทฺวมาสปฺปฏิจฺฉนฺนาโย, ตสฺส เม เอตทโหสิ ‘อหํ โข เทฺว สํฆาทิเสสา อาปตฺติโย อาปชฺชิํ เทฺวมาสปฺปฏิจฺฉนฺนาโย, ยํนูนาหํ สํฆํ เอกิสฺสา อาปตฺติยา เทฺวมาสปฺปฏิจฺฉนฺนาย เทฺวมาสปริวาสํ ยาเจยฺยนฺ’ติ, โสหํ สํฆํ เอกิสฺสา อาปตฺติยา เทฺวมาสปฺปฏิจฺฉนฺนาย เทฺวมาสปริวาสํ ยาจิํ, ตสฺส เม สํโฆ เอกิสฺสา อาปตฺติยา เทฺวมาสปฺปฏิจฺฉนฺนาย เทฺวมาสปริวาสํ อทาสิ, ตสฺส เม ปริวสนฺตสฺส ลชฺชีธมฺโม โอกฺกมิ\csromandash{}อหํ โข เทฺว สํฆาทิเสสา อาปตฺติโย อาปชฺชิํ เทฺวมาสปฺปฏิจฺฉนฺนาโย, ตสฺส เม เอตทโหสิ ‘อหํ โข เทฺว สํฆาทิเสสา อาปตฺติโย อาปชฺชิํ เทฺวมาสปฺปฏิจฺฉนฺนาโย, ยํนูนาหํ สํฆํ เอกิสฺสา อาปตฺติยา เทฺวมาสปฺปฏิจฺฉนฺนาย เทฺวมาสปริวาสํ ยาเจยฺยนฺ’ติ, โสหํ สํฆํ เอกิสฺสา อาปตฺติยา เทฺวมาสปฺปฏิจฺฉนฺนาย \csromanfolio{139}เทฺวมาสปริวาสํ ยาจิํ, ตสฺส เม สํโฆ เอกิสฺสา อาปตฺติยา เทฺวมาสปฺปฏิจฺฉนฺนาย เทฺวมาสปริวาสํ อทาสิ, ตสฺส เม ปริวสนฺตสฺส ลชฺชีธมฺโม โอกฺกมิ ‘ยํนูนาหํ สํฆํ อิตริสฺสาปิ อาปตฺติยา เทฺวมาสปฺปฏิจฺฉนฺนาย เทฺวมาสปริวาสํ ยาเจยฺยนฺ’ติ, โสหํ ภนฺเต สํฆํ อิตริสฺสาปิ อาปตฺติยา เทฺวมาสปฺปฏิจฺฉนฺนาย เทฺวมาสปริวาสํ ยาจามี”ติ.}

\prose{ทุติยมฺปิ ยาจิตพฺโพ.\csromanspacer{} ตติยมฺปิ ยาจิตพฺโพ.\csromanspacer{} พฺยตฺเตน ภิกฺขุนา ปฏิพเลน สํโฆ ญาเปตพฺโพ\csromandash{}}

\proseitem{139}{“สุณาตุ เม ภนฺเต สํโฆ, อยํ อิตฺถนฺนาโม ภิกฺขุ เทฺว สํฆาทิเสสา อาปตฺติโย อาปชฺชิ เทฺวมาสปฺปฏิจฺฉนฺนาโย, ตสฺส เอตทโหสิ ‘อหํ โข เทฺว สํฆาทิเสสา อาปตฺติโย อาปชฺชิํ เทฺวมาสปฺปฏิจฺฉนฺนาโย, ยํนูนาหํ สํฆํ เอกิสฺสา อาปตฺติยา เทฺวมาสปฺปฏิจฺฉนฺนาย เทฺวมาสปริวาสํ ยาเจยฺยนฺ’ติ, โส สํฆํ เอกิสฺสา อาปตฺติยา เทฺวมาสปฺปฏิจฺฉนฺนาย เทฺวมาสปริวาสํ ยาจิ, ตสฺส สํโฆ เอกิสฺสา อาปตฺติยา เทฺวมาสปฺปฏิจฺฉนฺนาย เทฺวมาสปริวาสํ อทาสิ, ตสฺส ปริวสนฺตสฺส ลชฺชีธมฺโม โอกฺกมิ\csromandash{}อหํ โข เทฺว สํฆาทิเสสา อาปตฺติโย อาปชฺชิํ เทฺวมาสปฺปฏิจฺฉนฺนาโย, ตสฺส เม เอตทโหสิ ‘อหํ โข เทฺว สํฆาทิเสสา อาปตฺติโย อาปชฺชิํ เทฺวมาสปฺปฏิจฺฉนฺนาโย, ยํนูนาหํ สํฆํ เอกิสฺสา อาปตฺติยา เทฺวมาสปฺปฏิจฺฉนฺนาย เทฺวมาสปริวาสํ ยาเจยฺยนฺ’ติ, โสหํ สํฆํ เอกิสฺสา อาปตฺติยา เทฺวมาสปฺปฏิจฺฉนฺนาย เทฺวมาสปริวาสํ ยาจิํ, ตสฺส เม สํโฆ เอกิสฺสา อาปตฺติยา เทฺวมาสปฺปฏิจฺฉนฺนาย เทฺวมาสปริวาสํ อทาสิ, ตสฺส เม ปริวสนฺตสฺส ลชฺชีธมฺโม โอกฺกมิ ‘ยํนูนาหํ สํฆํ อิตริสฺสาปิ อาปตฺติยา เทฺวมาสปฺปฏิจฺฉนฺนาย เทฺวมาสปริวาสํ ยาเจยฺยนฺ’ติ, โส สํฆํ อิตริสฺสาปิ อาปตฺติยา เทฺวมาสปฺปฏิจฺฉนฺนาย เทฺวมาสปริวาสํ ยาจติ, ยทิ สํฆสฺส ปตฺตกลฺลํ, สํโฆ อิตฺถนฺนามสฺส ภิกฺขุโน อิตริสฺสาปิ อาปตฺติยา เทฺวมาสปฺปฏิจฺฉนฺนาย เทฺวมาสปริวาสํ เทยฺย, เอสา ญตฺติ.}

\prose{สุณาตุ เม ภนฺเต สํโฆ, อยํ อิตฺถนฺนาโม ภิกฺขุ เทฺว สํฆาทิเสสา อาปตฺติโย อาปชฺชิ เทฺวมาสปฺปฏิจฺฉนฺนาโย, ตสฺส เอตทโหสิ ‘อหํ โข เทฺว สํฆาทิเสสา อาปตฺติโย อาปชฺชิํ เทฺวมาสปฺปฏิจฺฉนฺนาโย, ยํนูนาหํ สํฆํ เอกิสฺสา อาปตฺติยา \csromanfolio{140}เทฺวมาสปฺปฏิจฺฉนฺนาย เทฺวมาสปริวาสํ ยาเจยฺยนฺ’ติ, โส สํฆํ เอกิสฺสา อาปตฺติยา เทฺวมาสปฺปฏิจฺฉนฺนาย เทฺวมาสปริวาสํ ยาจิ, ตสฺส สํโฆ เอกิสฺสา อาปตฺติยา เทฺวมาสปฺปฏิจฺฉนฺนาย เทฺวมาสปริวาสํ อทาสิ, ตสฺส ปริวสนฺตสฺส ลชฺชีธมฺโม โอกฺกมิ\csromandash{}อหํ โข เทฺว สํฆาทิเสสา อาปตฺติโย อาปชฺชิํ เทฺวมาสปฺปฏิจฺฉนฺนาโย, ตสฺส เม เอตทโหสิ ‘อหํ โข เทฺว สํฆาทิเสสา อาปตฺติโย อาปชฺชิํ เทฺวมาสปฺปฏิจฺฉนฺนาโย, ยํนูนาหํ สํฆํ เอกิสฺสา อาปตฺติยา เทฺวมาสปฺปฏิจฺฉนฺนาย เทฺวสปริวาสํ ยาเจยฺยนฺ’ติ, โสหํ สํฆํ เอกิสฺสา อาปตฺติยา เทฺวมาสปฺปฏิจฺฉนฺนาย เทฺวมาสปริวาสํ ยาจิํ, ตสฺส เม สํโฆ เอกิสฺสา อาปตฺติยา เทฺวมาสปฺปฏิจฺฉนฺนาย เทฺวมาสปริวาสํ อทาสิ, ตสฺส เม ปริวสนฺตสฺส ลชฺชีธมฺโม โอกฺกมิ ‘ยํนูนาหํ สํฆํ อิตริสฺสาปิ อาปตฺติยา เทฺวมาสปฺปฏิจฺฉนฺนาย เทฺวมาสปริวาสํ ยาเจยฺยนฺ’ติ, โส สํฆํ อิตริสฺสาปิ อาปตฺติยา เทฺวมาสปฺปฏิจฺฉนฺนาย เทฺวมาสปริวาสํ ยาจติ, สํโฆ อิตฺถนฺนามสฺส ภิกฺขุโน อิตริสฺสาปิ อาปตฺติยา เทฺวมาสปฺปฏิจฺฉนฺนาย เทฺวมาสปริวาสํ เทติ, ยสฺสายสฺมโต ขมติ อิตฺถนฺนามสฺส ภิกฺขุโน อิตริสฺสาปิ อาปตฺติยา เทฺวมาสปฺปฏิจฺฉนฺนาย เทฺวมาสปริวาสสฺส ทานํ, โส ตุณฺหสฺส, ยสฺส นกฺขมติ, โส ภาเสยฺย.}

\prose{ทุติยมฺปิ เอตมตฺถํ วทามิ ฯเปฯ.\csromanspacer{} ตติยมฺปิ เอตมตฺถํ วทามิ ฯเปฯ.}

\prose{ทินฺโน สํเฆน อิตฺถนฺนามสฺส ภิกฺขุโน อิตริสฺสาปิ อาปตฺติยา เทฺวมาสปฺปฏิจฺฉนฺนาย เทฺวมสปริวาโส, ขมติ สํฆสฺส, ตสฺมา ตุณฺหี, เอวเมตํ ธารยามี”ติ.\csromanspacer{} เตน ภิกฺขเว ภิกฺขุนา ตทุปาทาย เทฺว มาสา ปริวสิตพฺพา.}

\csromanmatikaanchor{mtk.85}
\csromantocmarkref{subsection}{เทฺว มาสา ปริวสิตพฺพวิธิ}{mtk.85}
\bookmark[dest={mtk.85},level=2]{เทฺว มาสา ปริวสิตพฺพวิธิ}
\csromanheader{2}{\textbf{เทฺว มาสา ปริวสิตพฺพวิธิ}}
\proseitem{140}{อิธ ปน ภิกฺขเว ภิกฺขุ เทฺว สํฆาทิเสสา อาปตฺติโย อาปชฺชติ เทฺวมาสปฺปฏิจฺฉนฺนาโย, ตสฺส เอวํ โหติ “อหํ โข เทฺว สํฆาทิเสสา อาปตฺติโย อาปชฺชิํ เทฺวมาสปฺปฏิจฺฉนฺนาโย, ยํนุนาหํ สํฆํ เอกิสฺสา อาปตฺติยา เทฺวมาสปฺปฏิจฺฉนฺนาย เทฺวมาสปริวาสํ ยาเจยฺยนฺ”ติ, โส สํฆํ เอกิสฺสา อาปตฺติยา เทฺวมาสปฺปฏิจฺฉนฺนาย เทฺวมาสปริวาสํ ยาจติ, ตสฺส สํโฆ เอกิสฺสา อาปตฺติยา เทฺวมาสปฺปฏิจฺฉนฺนาย เทฺวมาสปริวาสํ เทติ, ตสฺส ปริวสนฺตสฺส ลชฺชีธมฺโม \csromanfolio{141}โอกฺกมิ\csromandash{}อหํ โข เทฺว สํฆาทิเสสา อาปตฺติโย อาปชฺชิํ เทฺวมาสปฺปฏิจฺฉนฺนาโย, ตสฺส เม เอตทโหสิ “อหํ โข เทฺว สํฆาทิเสสา อาปตฺติโย อาปชฺชิํ เทฺวมาสปฺปฏิจฺฉนฺนาโย, ยํนูนาหํ สํฆํ เอกิสฺสา อาปตฺติยา เทฺวมาสปฏิจฺฉนฺนาย เทฺวมาสปริวาสํ ยาเจยฺยนฺ”ติ, โสหํ สํฆํ เอกิสฺสา อาปตฺติยา เทฺวมาสปฺปฏิจฺฉนฺนาย เทฺวมาสปริวาสํ ยาจิํ, ตสฺส เม สํโฆ เอกิสฺสา อาปตฺติยา เทฺวมาสปฺปฏิจฺฉนฺนาย เทฺวมาสปริวาสํ อทาสิ, ตสฺส เม ปริวสนฺตสฺส ลชฺชีธมฺโม โอกฺกมิ ‘ยํนุนาหํ สํฆํ อิตริสฺสาปิ อาปตฺติยา เทฺวมาสปฺปฏิจฺฉนฺนาย เทฺวมาสปริวาสํ ยาเจยฺยนฺ”ติ, โส สํฆํ อิตริสฺสาปิ อาปตฺติยา เทฺวมาสปฺปฏิจฺฉนฺนาย เทฺวมาสปริวาสํ ยาจติ, ตสฺส สํโฆ อิตริสฺสาปิ อาปตฺติยา เทฺวมาสปฺปฏิจฺฉนฺนาย เทฺวมาสปริวาสํ เทติ.\csromanspacer{} เตน ภิกฺขเว ภิกฺขุนา ตทุปาทาย เทฺว มาสา ปริวสิตพฺพา.}

\proseitem{141}{อิธ ปน ภิกฺขเว ภิกฺขุ เทฺว สํฆาทิเสสา อาปตฺติโย อาปชฺชติ เทฺวมาสปฺปฏิจฺฉนฺนาโย, เอกํ อาปตฺติํ ชานาติ, เอกํ อาปตฺติํ น ชานาติ.\csromanspacer{} โส สํฆํ ยํ อาปตฺติํ ชานาติ, ตสฺสา อาปตฺติยา เทฺวมาสปฺปฏิจฺฉนฺนาย เทฺวมาสปริวาสํ ยาจติ, ตสฺส สํโฆ ตสฺสา อาปตฺติยา เทฺวมาสปฺปฏิจฺฉนฺนาย เทฺวมาสปริวาสํ เทติ, โส ปริวสนฺโต อิตรมฺปิ อาปตฺติํ ชานาติ, ตสฺส เอวํ โหติ “อหํ โข เทฺว สํฆาทิเสสา อาปตฺติโย อาปชฺชิํ เทฺวมาสปฺปฏิจฺฉนฺนาโย, เอกํ อาปตฺติ ชานิํ, เอกํ อาปตฺติํ น ชานิํ, โสหํ สํฆํ ยํ อาปตฺติํ ชานิํ, ตสฺสา อาปตฺติยา เทฺวมาสปฺปฏิจฺฉนฺนาย เทฺวมาสปริวาสํ ยาจิํ, ตสฺส เม สํโฆ ตสฺสา อาปตฺติยา เทฺวมาสปฺปฏิจฺฉนฺนาย เทฺวมาสปริวาสํ อทาสิ, โสหํ ปริวสนฺโต อิตรมฺปิ อาปตฺติํ ชานามิ, ยํนุนาหํ สํฆํ อิตริสฺสาปิ อาปตฺติยา เทฺวมาสปฺปฏิจฺฉนฺนาย เทฺวมาสปริวาสํ ยาเจยฺยนฺ”ติ, โส สํฆํ อิตริสฺสาปิ อาปตฺติยา เทฺวมาสปฺปฏิจฺฉนฺนาย เทฺวมาสปริวาสํ ยาจติ, ตสฺส สํโฆ อิตริสฺสาปิ อาปตฺติยา เทฺวมาสปฺปฏิจฺฉนฺนาย เทฺวมาสปริวาสํ เทติ.\csromanspacer{} เตน ภิกฺขเว ภิกฺขุนา ตทุปาทาย เทฺวมาสา ปริวสิตพฺพา.}

\proseitem{142}{อิธ ปน ภิกฺขเว ภิกฺขุ เทฺว สํฆาทิเสสา อาปตฺติโย อาปชฺชติ เทฺวมาสปฺปฏิจฺฉนฺนาโย, เอกํ อาปตฺติํ สรติ, เอกํ อาปตฺติํ นสฺสารติ.\csromanspacer{} โส สํฆํ ยํ อาปตฺติํ สรติ, ตสฺสา อาปตฺติยา \csromanfolio{142}เทฺวมาสปฺปฏิจฺฉนฺนาย เทฺวมาสปริวาสํ ยาจติ, ตสฺส สํโฆ ตสฺสา อาปตฺติยา เทฺวมาสปฺปฏิจฺฉนฺนาย เทฺวมาสปริวาสํ เทติ, โส ปริวสนฺโต อิตรมฺปิ อาปตฺติํ สรติ, ตสฺส เอวํ โหติ “อหํ โข เทฺว สํฆาทิเสสา อาปตฺติโย อาปชฺชิํ เทฺวมาสปฺปฏิจฺฉนฺนาโย, เอกํ อาปตฺติํ สริํ, เอกํ อาปตฺติํ นสฺสริํ, โสหํ สํฆํ ยํ อาปตฺติํ สริํ, ตสฺสา อาปตฺติยา เทฺวมาสปฺปฏิจฺฉนฺนาย เทฺวมาสปริวาสํ ยาจิํ, ตสฺส เม สํโฆ ตสฺสา อาปตฺติยา เทฺวมาสปฺปฏิจฺฉนฺนาย เทฺวมาสปริวาสํ อทาสิ, โสหํ ปริวสนฺโต อิตรมฺปิ อาปตฺติํ สรามิ, ยํนูนาหํ สํฆํ อิตริสฺสาปิ อาปตฺติยา เทฺวมาสปฺปฏิจฺฉนฺนาย เทฺวมาสปริวาสํ ยาเจยฺยนฺ”ติ, โส สํฆํ อิตริสฺสาปิ อาปตฺติยา เทฺวมาสปฺปฏิจฺฉนฺนาย เทฺวมาสปริวาสํ ยาจติ, ตสฺส สํโฆ อิตริสฺสาปิ อาปตฺติยา เทฺวมาสปฺปฏิจฺฉนฺนาย เทฺวมาสปริวาสํ เทติ.\csromanspacer{} เตน ภิกฺขเว ภิกฺขุนา ตทุปาทาย เทฺว มาสา ปริวสิตพฺพา.}

\proseitem{143}{อิธ ปน ภิกฺขเว ภิกฺขุ เทฺว สํฆาทิเสสา อาปตฺติโย อาปชฺชติ เทฺวมาสปฺปฏิจฺฉนฺนาโย, เอกาย อาปตฺติยา นิพฺเพมติโก, เอกาย อาปตฺติยา เวมติโก.\csromanspacer{} โส สํฆํ ยาย อาปตฺติยา นิพฺเพมติโก, ตสฺสา อาปตฺติยา เทฺวมาสปฺปฏิจฺฉนฺนาย เทฺวมาสปริวาสํ ยาจติ, ตสฺส สํโฆ ตสฺสา อาปตฺติยา เทฺวมาสปฺปฏิจฺฉนฺนาย เทฺวมาสปริวาสํ เทติ, โส ปริวสนฺโต อิตริสฺสาปิ อาปตฺติยา นิพฺเพมติโก โหติ, ตสฺส เอวํ โหติ “อหํ โข เทฺว สํฆาทิเสสา อาปตฺติโย อาปชฺชิํ เทฺวมาสปฺปฏิจฺฉนฺนาโย, เอกาย อาปตฺติยา นิพฺเพมติโก, เอกาย อาปตฺติยา เวมติโก, โสหํ สํฆํ ยาย อาปตฺติยา นิพฺเพมติโก, ตสฺสา อาปตฺติยา เทฺวมาสปฺปฏิจฺฉนฺนาย เทฺวมาสปริวาสํ ยาจิํ, ตสฺส เม สํโฆ ตสฺสา อาปตฺติยา เทฺวมาสปฺปฏิจฺฉนฺนาย เทฺวมาสปริวาสํ อทาสิ, โสหํ ปริวสนฺโต อิตริสฺสาปิ อาปตฺติยา นิพฺเพมติโก, ยํนูนาหํ สํฆํ อิตริสฺสาปิ อาปตฺติยา เทฺวมาสปฺปฏิจฺฉนฺนาย เทฺวมาสปริวาสํ ยาเจยฺยนฺ”ติ, โส สํฆํ อิตริสฺสาปิ อาปตฺติยา เทฺวมาสปฺปฏิจฺฉนฺนาย เทฺวมาสปริวาสํ ยาจติ, ตสฺส สํโฆ อิตริสฺสาปิ อาปตฺติยา เทฺวมาสปฺปฏิจฺฉนฺนาย เทฺวมาสปริวาสํ เทติ.\csromanspacer{} เตน ภิกฺขเว ภิกฺขุนา ตทุปาทาย เทฺว มาสา ปริวสิตพฺพา.}

\proseitem{144}{\csromanfolio{143}อิธ ปน ภิกฺขเว ภิกฺขุ เทฺว สํฆาทิเสสา อาปตฺติโย อาปชฺชติ เทฺวมาสปฺปฏิจฺฉนฺนาโย, เอกา อาปตฺติ ชานปฺปฏิจฺฉนฺนา, เอกา อาปตฺติ อชานปฺปฏิจฺฉนฺนา.\csromanspacer{} โส สํฆํ ตาสํ อาปตฺตีนํ เทฺวมาสปฺปฏิจฺฉนฺนานํ เทฺวมาสปริวาสํ ยาจติ, ตสฺส สํโฆ ตาสํ อาปตฺตีนํ เทฺวมาสปฺปฏิจฺฉนฺนานํ เทฺวมาสปริวาสํ เทติ, ตสฺส ปริวสนฺตสฺส อญฺโญ ภิกฺขุ อาคจฺฉติ พหุสฺสุโต อาคตาคโม ธมฺมธโร วินยธโร มาติกาธโร ปณฺฑิโต วิยตฺโต เมธาวี ลชฺชี กุกฺกุจฺจโก สิกฺขากาโม, โส เอวํ วเทติ “กิํ อยํ อาวุโส ภิกฺขุ อาปนฺโน กิสฺสายํ ภิกฺขุ ปริวสตี”ติ.\csromanspacer{} เต เอวํ วเทนฺติ “อยํ อาวุโส ภิกฺขุ เทฺว สํฆาทิเสสา อาปตฺติโย อาปชฺชิ เทฺวมาสปฺปฏิจฺฉนฺนาโย, เอกา อาปตฺติ ชานปฺปฏิจฺฉนฺนา, เอกา อาปตฺติ อชานปฺปฏิจฺฉนฺนา, โส สํฆํ ตาสํ อาปตฺตีนํ เทฺวมาสปฺปฏิจฺฉนฺนานํ เทฺวมาสปริวาสํ ยาจติ, ตสฺส สํโฆ ตาสํ อาปตฺตีนํ เทฺวมาสปฺปฏิจฺฉนฺนานํ เทฺวมาสปริวาสํ อทาสิ, ตาโย อยํ อาวุโส ภิกฺขุ อาปนฺโน, ตาสายํ ภิกฺขุ ปริวสตี”ติ.\csromanspacer{} โส เอวํ วเทติ “ยายํ อาวุโส อาปตฺติ ชานปฺปฏิจฺฉนฺนา, ธมฺมิกํ ตสฺสา อาปตฺติยา ปริวาสทานํ, ธมฺมตา รุหติ.\csromanspacer{} ยา จ ขฺวายํ อาวุโส อาปตฺติ อชานปฺปฏิจฺฉนฺนา, อธมฺมิกํ ตสฺสา อาปตฺติยา ปริวาสทานํ, อธมฺมตฺตา น รุหติ.\csromanspacer{} เอกิสฺสา อาวุโส อาปตฺติยา ภิกฺขุ มานตฺตารโห”ติ.}

\proseitem{145}{อิธ ปน ภิกฺขเว ภิกฺขุ เทฺว สํฆาทิเสสา อาปตฺติโย อาปชฺชติ เทฺวมาสปฺปฏิจฺฉนฺนาโย, เอกา อาปตฺติ สรมานปฺปฏิจฺฉนฺนา, เอกา อาปตฺติ อสฺสรมานปฺปฏิจฺฉนฺนา.\csromanspacer{} โส สํฆํ ตาสํ อาปตฺตีนํ เทฺวมาสปฺปฏิจฺฉนฺนานํ เทฺวมาสปริวาสํ ยาจติ, ตสฺส สํโฆ ตาสํ อาปตฺตีนํ เทฺวมาสปฺปฏิจฺฉนฺนานํ เทฺวมาสปริวาสํ เทติ, ตสฺส ปริวสนฺตสฺส อญฺโญ ภิกฺขุ อาคจฺฉติ พหุสฺสุโต อาคตาคโม ธมฺมธโร วินยธโร มาติกาธโร ปณฺฑิโต วิยตฺโต เมธาวี ลชฺชี กุกฺกุจฺจโก สิกฺขากาโม, โส เอวํ วเทติ “กิํ อยํ อาวุโส ภิกฺขุ อาปนฺโน, กิสฺสายํ ภิกฺขุ ปริวสตี”ติ.\csromanspacer{} เต เอวํ วเทนฺติ “อยํ อาวุโส ภิกฺขุ เทฺว สํฆาทิเสสา อาปตฺติโย อาปชฺชิ เทฺวมาสปฺปฏิจฺฉนฺนาโย, เอกา อาปตฺติ สรมานปฺปฏิจฺฉนฺนา, เอกา อาปตฺติ อสฺสรมานปฺปฏิจฺฉนฺนา, โส สํฆํ ตาสํ อาปตฺตีนํ เทฺวมาสปฺปฏิจฺฉนฺนานํ เทฺวมาสปริวาสํ ยาจิ, ตสฺส สํโฆ ตาสํ อาปตฺตีนํ เทฺวมาสปฺปฏิจฺฉนฺนานํ เทฺวมาสปริวาสํ อทาสิ, \csromanfolio{144}ตาโย อยํ อาวุโส ภิกฺขุ อาปนฺโน, ตาสายํ ภิกฺขุ ปริวสตี”ติ.\csromanspacer{} โส เอวํ วเทติ “ยายํ อาวุโส อาปตฺติ สรมานปฺปฏิจฺฉนฺนา, ธมฺมิกํ ตสฺสา อาปตฺติยา ปริวาสทานํ, ธมฺมตฺตา รุหติ.\csromanspacer{} ยา จ ขฺวายํ อาวุโส อาปตฺติ อสฺสรมานปฺปฏิจฺฉนฺนา, อธมฺมิกํ ตสฺส อาปตฺติยา ปริวาสทานํ, อธมฺมตฺตา น รุหติ.\csromanspacer{} เอกิสฺสา อาวุโส อาปตฺติยา ภิกฺขุ มานตฺตารโห”ติ.}

\proseitem{146}{อิธ ปน ภิกฺขเว ภิกฺขุ เทฺว สํฆาทิเสสา อาปตฺติโย อาปชฺชติ เทฺวมาสปฺปฏิจฺฉนฺนาโย, เอกา อาปตฺติ นิพฺเพมติกปฺปฏิจฺฉนฺนา, เอกา อาปตฺติ เวมติกปฺปฏิจฺฉนฺนา.\csromanspacer{} โส สํฆํ ตาสํ อาปตฺตีนํ เทฺวมาสปฺปฏิจฺฉนฺนานํ เทฺวมาสปริวาสํ ยาจติ, ตสฺส สํโฆ ตาสํ อาปตฺตีนํ เทฺวมาสปฺปฏิจฺฉนฺนานํ เทฺวมาสปริวาสํ เทติ, ตสฺส ปริวสนฺตสฺส อญฺโญ ภิกฺขุ อาคจฺฉติ พหุสฺสุโต อาคตาคโม ธมฺมธโร วินยธโร มาติกาธโร ปณฺฑิโต ปิยตฺโต เมธาวี ลชฺชี กุกฺกุจฺจโก สิกฺขากาโม, โส เอวํ วเทติ “กิํ อยํ อาวุโส ภิกฺขุ อาปนฺโน, กิสฺสายํ ติกฺขุ ปริวสตี”ติ.\csromanspacer{} เต เอวํ วเทนฺติ “อยํ อาวุโส ภิกฺขุ เทฺว สํฆาทิเสสา อาปตฺติโย อาปชฺชิ เทฺวมาสปฺปฏิจฺฉนฺนาโย, เอกา อาปตฺติ นิพฺเพมติกปฺปฏิจฺฉนฺนา, เอกา อาปตฺติ เวมติกปฺปฏิจฺฉนฺนา, โส สํฆํ ตาสํ อาปตฺตีนํ เทฺวมาสปฺปฏิจฺฉนฺนานํ เทฺวมาสปริวาสํ ยาจิ, ตสฺส สํโฆ ตาสํ อาปตฺตีนํ เทฺวมาสปฺปฏิจฺฉนฺนานํ เทฺวมาสปริวาสํ อทาสิ, ตาโย อยํ อาวุโส ภิกฺขุ อาปนฺโน, ตาสายํ ภิกฺขุ ปริวสตี”ติ.\csromanspacer{} โส เอวํ วเทติ “ยายํ อาวุโส อาปตฺติ นิพฺเพมติกปฺปฏิจฺฉนฺนา, ธมฺมิกํ ตสฺสา อาปตฺติยา ปริวาสทานํ, ธมฺมตฺตา รุหติ.\csromanspacer{} ยา จ ขฺวายํ อาวุโส อาปตฺติ เวมติกปฺปฏิจฺฉนฺนา, อธมฺมิกํ ตสฺสา อาปตฺติยา ปริวาสทานํ, อธมฺมตฺตา น รุหติ.\csromanspacer{} เอกิสฺสา อาวุโส อาปตฺติยา ภิกฺขุ มานตฺตารโห”ติ.}

\proseitem{147}{เตน โข ปน สมเยน อญฺญตโร ภิกฺขุ เทฺว สํฆาทิเสสา อาปตฺติโย อาปนฺโน โหติ เทฺวมาสปฺปฏิจฺฉนฺนาโย.\csromanspacer{} ตสฺส เอตทโหสิ “อหํ โข เทฺว สํฆาทิเสสา อาปตฺติโย อาปชฺชิํ เทฺวมาสปฺปฏิจฺฉนฺนาโย, ยํนูนาหํ สํฆํ ทฺวินฺนํ อาปตฺตีนํ เทฺวมาสปฺปฏิจฺฉนฺนานํ เอกมาสปริวาสํ ยาเจยฺยนฺ”ติ.\csromanspacer{} โส สํฆํ ทฺวินฺนํ อาปตฺตีนํ เทฺวมาสปฺปฏิจฺฉนฺนานํ \csromanfolio{145}เอกมาสปริวาสํ ยาจิ, ตสฺส สํโฆ ทฺวินฺนํ อาปตฺตีนํ เทฺวมาสปฺปฏิจฺฉนฺนานํ เอกมาสปริวาสํ อทาสิ, ตสฺส ปริวสนฺตสฺส ลชฺชีธมฺโม โอกฺกมิ\csromandash{}อหํ โข เทฺว สํฆาทิเสสา อาปตฺติโย อาปชฺชิํ เทฺวมาสปฺปฏิจฺฉนฺนาโย, ตสฺส เม เอตทโหสิ “อหํ โข เทฺว สํฆาทิเสสา อาปตฺติโย อาปชฺชิํ เทฺวมาสปฺปฏิจฺฉนฺนาโย, ยํนูนาหํ สํฆํ ทฺวินฺนํ อาปตฺตีนํ เทฺวมาสปฺปฏิจฺฉนฺนานํ เอกมาสปริวาสํ ยาเจยฺยนฺ”ติ, โสหํ สํฆํ ทฺวินฺนํ อาปตฺตีนํ เทฺวมาสปฺปฏิจฺฉนฺนานํ เอกมาสปริวาสํ ยาจิํ, ตสฺส เม สํโฆ ทฺวินฺนํ อาปตฺตีนํ เทฺวมาสปฺปฏิจฺฉนฺนานํ เอกมาสปริวาสํ อทาสิ, ตสฺส เม ปริวสนฺตสฺส ลชฺชีธมฺโม โอกฺกมิ “ยํนูนาหํ สํฆํ ทฺวินฺนํ อาปตฺตีนํ เทฺวมาสปฺปฏิจฺฉนฺนานํ อิตรมฺปิ มาสํ ปริวาสํ\footnote{อิตรมฺปิ มาสปริวาสํ (สฺยา, ก, เอวมุปริปิ)} ยาเจยฺยนฺ”ติ.}

\prose{โส ภิกฺขูนํ อาโรเจสิ “อหํ โข อาวุโส เทฺว สํฆาทิเสสา อาปตฺติโย อาปชฺชิํ เทฺวมาสปฺปฏิจฺฉนฺนาโย, ตสฺส เม เอตทโหสิ ‘อหํ โข เทฺว สํฆาทิเสสา อาปตฺติโย อาปชฺชิํ เทฺวมาสปฺปฏิจฺฉนฺนาโย, ยํนูนาหํ สํฆํ ทฺวินฺนํ อาปตฺตีนํ เทฺวมาสปฺปฏิจฺฉนฺนานํ เอกมาสปริวาสํ ยาเจยฺยนฺ’ติ, โสหํ สํฆํ ทฺวินฺนํ อาปตฺตีนํ เทฺวมาสปฺปฏิจฺฉนฺนานํ เอกมาสปริวาสํ ยาจิํ, ตสฺส เม สํโฆ ทฺวินฺนํ อาปตฺตีนํ เทฺวมาสปฺปฏิจฺฉนฺนานํ เอกมาสปริวาสํ อทาสิ, ตสฺส เม ปริวสนฺตสฺส ลชฺชีธมฺโม โอกฺกมิ\csromandash{}อหํ โข เทฺว สํฆาทิเสสา อาปตฺติโย อาปชฺชิํ เทฺวมาสปฺปฏิจฺฉนฺนาโย, ตสฺส เม เอตทโหสิ ‘อหํ โข เทฺว สํฆาทิเสสา อาปตฺติโย อาปชฺชิํ เทฺวมาสปฺปฏิจฺฉนฺนาโย, ยํนูนาหํ สํฆํ ทฺวินฺนํ อาปตฺตีนํ เทฺวมาสปฺปฏิจฺฉนฺนานํ เอกมาสปริวาสํ ยาเจยฺยนฺ’ติ, โสหํ สํฆํ ทฺวินฺนํ อาปตฺตีนํ เทฺวมาสปฺปฏิจฺฉนฺนานํ เอกมาสปริวาสํ ยาจิํ, ตสฺส เม สํโฆ ทฺวินฺนํ อาปตฺตีนํ เทฺวมาสปฺปฏิจฺฉนฺนานํ เอกมาสปริวาสํ อทาสิ, ตสฺส เม ปริวสนฺตสฺส ลชฺชีธมฺโม โอกฺกมิ ‘ยํนูนาหํ สํฆํ ทฺวินฺนํ อาปตฺตีนํ เทฺวมาสปฺปฏิจฺฉนฺนานํ อิตรมฺปิ มาสํ ปริวาสํ ยาเจยฺยนฺ’ติ.\csromanspacer{} กถํ นุ โข มยา ปฏิปชฺชิตพฺพนฺ”ติ.\csromanspacer{} ภควโต เอตมตฺถํ อาโรเจสุํ.\csromanspacer{} เตน หิ ภิกฺขเว สํโฆ ตสฺส ภิกฺขุโน ทฺวินฺนํ อาปตฺตีนํ เทฺวมาสปฺปฏิจฺฉนฺนานํ อิตรมฺปิ มาสํ ปริวาสํ เทตุ.\csromanspacer{} เอวญฺจ ปน ภิกฺขเว ทาตพฺโพ\csromandash{}}

\prose{เตน ภิกฺขเว ภิกฺขุนา สํฆํ อุปสงฺกมิตฺวา ฯเปฯ เอวมสฺส วจนีโย “อหํ ภนฺเต เทฺว สํฆาทิเสสา อาปตฺติโย อาปชฺชิํ \csromanfolio{146}เทฺวมาสปฺปฏิจฺฉนฺนาโย, ตสฺส เม เอตทโหสิ ‘อหํ โข เทฺว สํฆาทิเสสา อาปตฺติโย อาปชฺชิํ เทฺวมาสปฺปฏิจฺฉนฺนาโย, ยํนูนาหํ สํฆํ ทฺวินฺนํ อาปตฺตีนํ เทฺวมาสปฺปฏิจฺฉนฺนานํ เอกมาสปริวาสํ ยาเจยฺยนฺ’ติ, โสหํ สํฆํ ทฺวินฺนํ อาปตฺตีนํ เทฺวมาสปฺปฏิจฺฉนฺนานํ เอกมาสปริวาสํ ยาจิํ, ตสฺส เม สํโฆ ทฺวินฺนํ อาปตฺตีนํ เทฺวมาสปฺปฏิจฺฉนฺนานํ เอกมาสปริวาสํ อทาสิ, ตสฺส เม ปริวสนฺตสฺส ลชฺชีธมฺโม โอกฺกมิ\csromandash{} อหํ โข เทฺว สํฆาทิเสสา อาปตฺติโย อาปชฺชิํ เทฺวมาสปฺปฏิจฺฉนฺนาโย, ตสฺส เม เอตทโหสิ ‘อหํ โข เทฺว สํฆาทิเสสา อาปตฺติโย อาปชฺชิํ เทฺวมาสปฺปฏิจฺฉนฺนาโย ยํนูนาหํ สํฆํ ทฺวินฺนํ อาปตฺตีนํ เทฺวมาสปฺปฏิจฺฉนฺนานํ เอกมาสปริวาสํ ยาเจยฺยนฺ’ติ, โสหํ สํฆํ ทฺวินฺนํ อาปตฺตีนํ เทฺวมาสปฺปฏิจฺฉนฺนานํ เอกมาสปริวาสํ ยาจิํ, ตสฺส เม สํโฆ ทฺวินฺนํ อาปตฺตีนํ เทฺวมาสปฺปฏิจฺฉนฺนานํ เอกมาสปริวาสํ อทาสิ, ตสฺส เม ปริวสนฺตสฺส ลชฺชีธมฺโม โอกฺกมิ ‘ยํนูนาหํ สํฆํ ทฺวินฺนํ อาปตฺตีนํ เทฺวมาสปฺปฏิจฺฉนฺนานํ อิตรมฺปิ มาสํ ปริวาสํ ยาเจยฺยนฺ’ติ, โสหํ ภนฺเต สํฆํ ทฺวินฺนํ อาปตฺตีนํ เทฺวมาสปฺปฏิจฺฉนฺนานํ อิตรมฺปิ มาสํ ปริวาสํ ยาจามี”ติ.\csromanspacer{} ทุติยมฺปิ ยาจิตพฺโพ.\csromanspacer{} ตติยมฺปิ ยาจิตพฺโพ.\csromanspacer{} พฺยตฺเตน ภิกฺขุนา ปฏิพเลน สํโฆ ญาเปตพฺโพ\csromandash{}}

\proseitem{148}{“สุณาตุ เม ภนฺเต สํโฆ, อยํ อิตฺถนฺนาโม ภิกฺขุ เทฺว สํฆาทิเสสา อาปตฺติโย อาปชฺชิ เทฺวมาสปฺปฏิจฺฉนฺนาโย, ตสฺส เอตทโหสิ ‘อหํ โข เทฺว สํฆาทิเสสา อาปตฺติโย อาปชฺชิํ เทฺวมาสปฺปฏิจฺฉนฺนาโย, ยํนูนาหํ สํฆํ ทฺวินฺนํ อาปตฺตีนํ เทฺวมาสปฺปฏิจฺฉนฺนานํ เอกมาสปริวาสํ ยาเจยฺยนฺ’ติ, โส สํฆํ ทฺวินฺนํ อาปตฺตีนํ เทฺว มาสปฺปฏิจฺฉนฺนานํ เอกมาสปริวาสํ ยาจิ, ตสฺส สํโฆ ทฺวินฺนํ อาปตฺตีนํ เทฺวมาสปฺปฏิจฺฉนฺนานํ เอกมาสปริวาสํ อทาสิ, ตสฺส ปริวสนฺตสฺส ลชฺชีธมฺโม โอกฺกมิ\csromandash{}อหํ โข เทฺว สํฆาทิเสสา อาปตฺติโย อาปชฺชิํ เทฺวมาสปฺปฏิจฺฉนฺนาโย, ตสฺส เม เอตทโหสิ ‘อหํ โข เทฺว สํฆาทิเสสา อาปตฺติโย อาปชฺชิํ เทฺวมาสปฺปฏิจฺฉนฺนาโย, ยํนูนาหํ สํฆํ ทฺวินฺนํ อาปตฺตีนํ เทฺวมาสปฺปฏิจฺฉนฺนานํ เอกมาสปริวาสํ ยาเจยฺยนฺ’ติ, โสหํ สํฆํ ทฺวินฺนํ อาปตฺตีนํ เทฺวมาสปฺปฏิจฺฉนฺนานํ เอกมาสปริวาสํ ยาจิํ, ตสฺส เม สํโฆ ทฺวินฺนํ อาปตฺตีนํ เทฺวมาสปฺปฏิจฺฉนฺนานํ เอกมาสปริวาสํ อทาสิ, ตสฺส เม ปริวสนฺตสฺส ลชฺชีธมฺโม โอกฺกมิ ‘ยํนูนาหํ สํฆํ ทฺวินฺนํ อาปตฺตีนํ เทฺวมาสปฺปฏิจฺฉนฺนานํ อิตรมฺปิ มาสํ ปริวาสํ ยาเจยฺยนฺ’ติ, โส สํฆํ ทฺวินฺนํ \csromanfolio{147}อาปตฺตีนํ เทฺวมาสปฺปฏิจฺฉนฺนานํ อิตรมฺปิ มาสํ ปริวาสํ ยาจติ, ยทิ สํฆสฺส ปตฺตกลฺลํ, สํโฆ อิตฺถนฺนามสฺส ภิกฺขุโน ทฺวินฺนํ อาปตฺตีนํ เทฺวมาสปฺปฏิจฺฉนฺนานํ อิตรมฺปิ มาสํ ปริวาสํ ทเทยฺย, เอสา ญตฺติ.}

\prose{สุณาตุ เม ภนฺเต สํโฆ, อยํ อิตฺถนฺนาโม ภิกฺขุ เทฺว สํฆาทิเสสา อาปตฺติโย อาปชฺชิ เทฺวมาสปฺปฏิจฺฉนฺนาโย, ตสฺส เอตทโหสิ ‘อหํ โข เทฺว สํฆาทิเสสา อาปตฺติโย อาปชฺชิํ เทฺวมาสปฺปฏิจฺฉนฺนาโย, ยํนูนาหํ สํฆํ ทฺวินฺนํ อาปตฺตีนํ เทฺวมาสปฺปฏิจฺฉนฺนานํ เอกมาสปริวาสํ ยาเจยฺยนฺ’ติ, โส สํฆํ ทฺวินฺนํ อาปตฺตีนํ เทฺวมาสปฺปฏิจฺฉนฺนานํ เอกมาสปริวาสํ ยาจิ, ตสฺส สํโฆ ทฺวินฺนํ อาปตฺตีนํ เทฺวมาสปฺปฏิจฺฉนฺนานํ เอกมาสปริวาสํ อทาสิ, ตสฺส ปริวสนฺตสฺส ลชฺชีธมฺโม โอกฺกมิ\csromandash{}อหํ โข เทฺว สํฆาทิเสสา อาปตฺติโย อาปชฺชิํ เทฺวมาสปฺปฏิจฺฉนฺนาโย, กสฺส เม เอตทโหสิ ‘อหํ โข เทฺว สํฆาทิเสสา อาปตฺติโย อาปชฺชิํ เทฺวมาสปฺปฏิจฺฉนฺนาโย, ยํนูนาหํ สํฆํ ทฺวินฺนํ อาปตฺตีนํ เทฺวมาสปฺปฏิจฺฉนฺนานํ เอกมาสปริวาสํ ยาเจยฺยนฺ’ติ, โสหํ สํฆํ ทฺวินฺนํ อาปตฺตีนํ เทฺวมาสปฺปฏิจฺฉนฺนานํ เอกมาสปริวาสํ ยาจิํ, ตสฺส เม สํโฆ ทฺวินฺนํ อาปตฺตีนํ เทฺวมาสปฺปฏิจฺฉนฺนานํ เอกมาสปริวาสํ อทาสิ, ตสฺส เม ปริวสนฺตสฺส ลชฺชีธมฺโม โอกฺกมิ ‘ยํนูนาหํ สํฆํ ทฺวินฺนํ อาปตฺตีนํ เทฺวมาสปฺปฏิจฺฉนฺนานํ อิตรมฺปิ มาสํ ปริวาสํ ยาเจยฺยนฺ’ติ, โส สํฆํ ทฺวินฺนํ อาปตฺตีนํ เทฺวมาสปฺปฏิจฺฉนฺนานํ อิตรมฺปิ มาสํ ปริวาสํ ยาจติ, สํโฆ อิตฺถนฺนามสฺส ภิกฺขุโน ทฺวินฺนํ อาปตฺตีนํ เทฺวมาสปฺปฏิจฺฉนฺนานํ อิตรมฺปิ มาสํ ปริวาสํ เทติ, ยสฺสายสฺมโต ขมติ อิตฺถนฺนามสฺส ภิกฺขุโน ทฺวินฺนํ อาปตฺตีนํ เทฺวมาสปฺปฏิจฺฉนฺนานํ อิตรมฺปิ มาสํ ปริวาสสฺส\footnote{อิตรมฺปิ มาสปริวาสสฺส (ก), อิตรสฺสปิ มาสปริวาสสฺส (สฺยา)} ทานํ, โส ตุณฺหสฺส, ยสฺส นกฺขมติ, โส ภาเสยฺย.}

\prose{ทุติยมฺปิ เอตมตฺถํ วทามิ ฯเปฯ.\csromanspacer{} ตติยมฺปิ เอตมตฺถํ วทามิ ฯเปฯ.}

\prose{ทินฺโน สํเฆน อิตฺถนฺนามสฺส ภิกฺขุโน ทฺวินฺนํ อาปตฺตีนํ เทฺวมาสปฺปฏิจฺฉนฺนานํ อิตรมฺปิ มาสํ ปริวาโส\footnote{อิตรมฺปิ มาสปริวาโส (ก), อิตโรปิ มาสปริวาโส (สฺยา)}, ขมติ สํฆสฺส, ตสฺมา ตุณฺหี, เอวเมตํ ธารยามี”ติ.}

\prose{เตน ภิกฺขเว ภิกฺขุนา ปุริมํ อุปาทาย เทฺว มาสา ปริวสิตพฺพา.}

\proseitem{149}{\csromanfolio{148}อิธ ปน ภิกฺขเว ภิกฺขุ เทฺว สํฆาทิเสสา อาปตฺติโย อาปชฺชติ เทฺวมาสปฺปฏิจฺฉนฺนาโย.\csromanspacer{} ตสฺส เอวํ โหติ “อหํ โข เทฺว สํฆาทิเสสา อาปตฺติโย อาปชฺชิํ เทฺวมาสปฺปฏิจฺฉนฺนาโย, ยํนูนาหํ สํฆํ ทฺวินฺนํ อาปตฺตีนํ เทฺวมาสปฺปฏิจฺฉนฺนานํ เอกมาสปริวาสํ ยาเจยฺยนฺ”ติ, โส สํฆํ ทฺวินฺนํ อาปตฺตีนํ เทฺวมาสปฺปฏิจฺฉนฺนานํ เอกมาสปริวาสํ ยาจติ, ตสฺส สํโฆ ทฺวินฺนํ อาปตฺตีนํ เทฺวมาสปฺปฏิจฺฉนฺนานํ เอกมาสปริวาสํ เทติ, ตสฺส ปริวสนฺตสฺส ลชฺชีธมฺโม โอกฺกมิ\csromandash{}อหํ โข เทฺว สํฆาทิเสสา อาปตฺติโย อาปชฺชิํ เทฺวมาสปฺปฏิจฺฉนฺนาโย, ตสฺส เม เอตทโหสิ “อหํ โข เทฺว สํฆาทิเสสา อาปตฺติโย อาปชฺชิํ เทฺวมาสปฺปฏิจฺฉนฺนาโย, ยํนูนาหํ สํฆํ ทฺวินฺนํ อาปตฺตีนํ เทฺวมาสปฺปฏิจฺฉนฺนานํ เอกมาสปริวาสํ ยาเจยฺยนฺ”ติ, โสหํ สํฆํ ทฺวินฺนํ อาปตฺตีนํ เทฺวมาสปฺปฏิจฺฉนฺนานํ เอกมาสปริวาสํ ยาจิํ, ตสฺส เม สํโฆ ทฺวินฺนํ อาปตฺตีนํ เทฺวมาสปฺปฏิจฺฉนฺนานํ เอกมาสปริวาสํ อทาสิ, ตสฺส เม ปริวสนฺตสฺส ลชฺชีธมฺโม โอกฺกมิ “ยํนูนาหํ สํฆํ ทฺวินฺนํ อาปตฺตีนํ เทฺวมาสปฺปฏิจฺฉนฺนานํ อิตรมฺปิ มาสํ ปริวาสํ ยาเจยฺยนฺ”ติ, โส สํฆํ ทฺวินฺนํ อาปตฺตีนํ เทฺวมาสปฺปฏิจฺฉนฺนานํ อิตรมฺปิ มาสํ ปริวาสํ ยาจติ, ตสฺส สํโฆ ทฺวินฺนํ อาปตฺตีนํ เทฺวมาสปฺปฏิจฺฉนฺนานํ อิตรมฺปิ มาสํ ปริวาสํ เทติ.\csromanspacer{} เตน ภิกฺขเว ภิกฺขุนา ปุริมํ อุปาทาย เทฺว มาสา ปริวสิตพฺพา.}

\proseitem{150}{อิธ ปน ภิกฺขเว ภิกฺขุ เทฺว สํฆาทิเสสา อาปตฺติโย อาปชฺชติ เทฺวมาสปฺปฏิจฺฉนฺนาโย, เอกํ มาสํ ชานาติ, เอกํ มาสํ น ชานติ.\csromanspacer{} โส สํฆํ ทฺวินฺนํ อาปตฺตีนํ เทฺวมาสปฺปฏิจฺฉนฺนานํ ยํ มาสํ ชานาติ, ตํ มาสํ ปริวาสํ ยาจติ, ตสฺส สํโฆ ทฺวินฺนํ อาปตฺตีนํ เทฺวมาสปฺปฏิจฺฉนฺนานํ ยํ มาสํ ชานาติ, ตํ มาสํ ปริวาสํ เทติ, โส ปริวสนฺโต อิตรมฺปิ มาสํ ชานาติ, ตสฺส เอวํ โหติ ‘อหํ โข เทฺว สํฆาทิเสสา อาปตฺติโย อาปชฺชิํ เทฺวมาสปฺปฏิจฺฉนฺนาโย, เอกํ มาสํ ชานิํ, เอกํ มาสํ น ชานิํ, โสหํ สํฆํ ทฺวินฺนํ อาปตฺตีนํ เทฺวมาสปฺปฏิจฺฉนฺนานํ ยํ มาสํ ชานิํ, ตํ มาสํ ปริวาสํ ยาจิํ, ตสฺส เม สํโฆ ทฺวินฺนํ อาปตฺตีนํ เทฺวมาสปฺปฏิจฺฉนฺนานํ ยํ มาสํ ชานิํ, ตํ มาสํ ปริวาสํ อทาสิ, โสหํ ปริวสนฺโต อิตรมฺปิ มาสํ ชานามิ, ยํนูนาหํ สํฆํ ทฺวินฺนํ อาปตฺตีนํ เทฺวมาสปฺปฏิจฺฉนฺนานํ อิตรมฺปิ มาสํ ปริวาสํ ยาเจยฺยนฺ’ติ, โส สํฆํ ทฺวินฺนํ อาปตฺตีนํ เทฺวมาสปฺปฏิจฺฉนฺนานํ อิตรมฺปิ มาสํ ปริวาสํ \csromanfolio{149}ยาจติ, ตสฺส สํโฆ ทฺวินฺนํ อาปตฺตีนํ เทฺวมาสปฺปฏิจฺฉนฺนานํ อิตรมฺปิ มาสํ ปริวาสํ เทติ.\csromanspacer{} เตน ภิกฺขเว ภิกฺขุนา ปุริมํ อุปาทาย เทฺว มาสา ปริวสิตพฺพา.}

\proseitem{151}{อิธ ปน ภิกฺขเว ภิกฺขุ เทฺว สํฆาทิเสสา อาปตฺติโย อาปชฺชติ เทฺวมาสปฺปฏิจฺฉนฺนาโย, เอกํ มาสํ สรติ, เอกํ มาสํ นสฺสรติ.\csromanspacer{} โส สํฆํ ทฺวินฺนํ อาปตฺตีนํ เทฺวมาสปฺปฏิจฺฉนฺนานํ ยํ มาสํ สรติ, ตํ มาสํ ปริวาสํ ยาจติ, ตสฺส สํโฆ ทฺวินฺนํ อาปตฺตีนํ เทฺวมาสปฺปฏิจฺฉนฺนานํ ยํ มาสํ สรติ, ตํ มาสํ ปริวาสํ เทติ, โส ปริวสนฺโต อิตรมฺปิ มาสํ สรติ, ตสฺส เอวํ โหติ, “อหํ โข เทฺว สํฆาทิเสสา อาปตฺติโย อาปชฺชิํ เทฺวมาสปฺปฏิจฺฉนฺนาโย, เอกํ มาสํ สริํ, เอกํ มาสํ นสฺสริํ, โสหํ สํฆํ ทฺวินฺนํ อาปตฺตีนํ เทฺวมาสปฺปฏิจฺฉนฺนานํ ยํ มาสํ สริํ, ตํ มาสํ ปริวาสํ ยาจิํ, ตสฺส เม สํโฆ ทฺวินฺนํ อาปตฺตีนํ เทฺวมาสปฺปฏิจฺฉนฺนานํ ยํ มาสํ สริํ, ตํ มาสํ ปริวาสํ อทาสิ, โสหํ ปริวสนฺโต อิตรมฺปิ มาสํ สรามิ, ยํนูนาหํ สํฆํ ทฺวินฺนํ อาปตฺตีนํ เทฺวมาสปฺปฏิจฺฉนฺนานํ อิตรมฺปิ มาสํ ปริวาสํ ยาเจยฺยนฺ”ติ, โส สํฆํ ทฺวินฺนํ อาปตฺตีนํ เทฺวมาสปฺปฏิจฺฉนฺนานํ อิตรมฺปิ มาสํ ปริวาสํ ยาจติ, ตสฺส สํโฆ ทฺวินฺนํ อาปตฺตีนํ เทฺวมาสปฺปฏิจฺฉนฺนานํ อิตรมฺปิ มาสํ ปริวาสํ เทติ.\csromanspacer{} เตน ภิกฺขเว ภิกฺขุนา ปุริมํ อุปาทาย เทฺว มาสา ปริวสิตพฺพา.}

\proseitem{152}{อิธ ปน ภิกฺขเว ภิกฺขุ เทฺว สํฆาทิเสสา อาปตฺติโย อาปชฺชติ เทฺวมาสปฺปฏิจฺฉนฺนาโย, เอกํ มาสํ นิพฺเพมติโก, เอกํ มาสํ เวมติโก.\csromanspacer{} โส สํฆํ ทฺวินฺนํ อาปตฺตีนํ เทฺวมาสปฺปฏิจฺฉนฺนานํ ยํ มาสํ นิพฺเพมติโก, ตํ มาสํ ปริวาสํ ยาจติ, ตสฺส สํโฆ ทฺวินฺนํ อาปตฺตีนํ เทฺวมาสปฺปฏิจฺฉนฺนานํ ยํ มาสํ นิพฺเพมติโก, ตํ มาสํ ปริวาสํ เทติ, โส ปริวสนฺโต อิตรมฺปิ มาสํ นิพฺเพมติโก โหติ, ตสฺส เอวํ โหติ “อหํ โข เทฺว สํฆาทิเสสา อาปตฺติโย อาปชฺชิํ เทฺวมาสปฺปฏิจฺฉนฺนาโย, เอกํ มาสํ นิพฺเพมติโก, เอกํ มาสํ เวมติโก, โสหํ สํฆํ ทฺวินฺนํ อาปตฺตีนํ เทฺวมาสปฺปฏิจฺฉนฺนานํ ยํ มาสํ นิพฺเพมติโก, ตํ มาสํ ปริวาสํ ยาจิํ, ตสฺส เม สํโฆ ทฺวินฺนํ อาปตฺตีนํ เทฺวมาสปฺปฏิจฺฉนฺนานํ ยํ มาสํ นิพฺเพมติโก, ตํ มาสํ ปริวาสํ อทาสิ, โสหํ ปริวสนฺโต อิตรมฺปิ มาสํ นิพฺเพมติโก, ยํนูนาหํ สํฆํ ทฺวินฺนํ อาปตฺตีนํ \csromanfolio{150}เทฺวมาสปฺปฏิจฺฉนฺนานํ อิตรมฺปิ มาสํ ปริวาสํ ยาเจยฺยนฺ”ติ, โส สํฆํ ทฺวินฺนํ อาปตฺตีนํ เทฺวมาสปฺปฏิจฺฉนฺนานํ อิตรมฺปิ มาสํ ปริวาสํ ยาจติ, ตสฺส สํโฆ ทฺวินฺนํ อาปตฺตีนํ เทฺวมาสปฺปฏิจฺฉนฺนานํ อิตรมฺปิ มาสํ ปริวาสํ เทติ.\csromanspacer{} เตน ภิกฺขเว ภิกฺขุนา ปุริมํ อุปาทาย เทฺว มาสา ปริวสิตพฺพา.}

\proseitem{153}{อิธ ปน ภิกฺขเว ภิกฺขุ เทฺว สํฆาทิเสสา อาปตฺติโย อาปชฺชติ เทฺวมาสปฺปฏิจฺฉนฺนาโย, เอโก มาโส ชานปฺปฏิจฺฉนฺโน, เอโก มาโส อชานปฺปฏิจฺฉนฺโน.\csromanspacer{} โส สํฆํ ทฺวินฺนํ อาปตฺตีนํ เทฺวมาสปฺปฏิจฺฉนฺนานํ เทฺวมาสปริวาสํ ยาจติ, ตสฺส สํโฆ ทฺวินฺนํ อาปตฺตีนํ เทฺวมาสปฺปฏิจฺฉนฺนานํ เทฺวมาสปริวาสํ เทติ, ตสฺส ปริวสนฺตสฺส อญฺโญ ภิกฺขุ อาคจฺฉติ พหุสฺสุโต อาคตาคโม ธมฺมธโร วินยธโร มาติกาธโร ปณฺฑิโต วิยตฺโต เมธาวี ลชฺชี กุกฺกุจฺจโก สิกฺขากาโม, โส เอวํ วเทติ “กิํ อยํ อาวุโส ภิกฺขุ อาปนฺโน, กิสฺสายํ ภิกฺขุ ปริวสตี”ติ.\csromanspacer{} เต เอวํ วเทนฺติ “อยํ อาวุโส ภิกฺขุ เทฺว สํฆาทิเสสา อาปตฺติโย อาปชฺชิ เทฺวมาสปฺปฏิจฺฉนฺนาโย, เอโก มาโส ชานปฺปฏิจฺฉนฺโน, เอโก มาโส อชานปฺปฏิจฺฉนฺโน.\csromanspacer{} โส สํฆํ ทฺวินฺนํ อาปตฺตีนํ เทฺวมาสปฺปฏิจฺฉนฺนานํ เทฺวมาสปริวาสํ ยาจิ, ตสฺส สํโฆ ทฺวินฺนํ อาปตฺตีนํ เทฺวมาสปฺปฏิจฺฉนฺนานํ เทฺวมาสปริวาสํ อทาสิ, ตาโย อยํ อาวุโส ภิกฺขุ อาปนฺโน, ตาสายํ ภิกฺขุ ปริวสตี”ติ.\csromanspacer{} โส เอวํ วเทติ “ยฺวายํ อาวุโส มาโส ชานปฺปฏิจฺฉนฺโน, ธมฺมิกํ ตสฺส มาสสฺส ปริวาสทานํ, ธมฺมตฺตา รุหติ.\csromanspacer{} โย จ ขฺวายํ อาวุโส มาโส อชานปฺปฏิจฺฉนฺโน, อธมฺมิกํ ตสฺส มาสสฺส ปริวาสทานํ, อธมฺมตฺตา น รุหติ, เอกสฺส อาวุโส มาสสฺส ภิกฺขุ มานตฺตารโห”ติ.}

\proseitem{154}{อิธ ปน ภิกฺขเว ภิกฺขุ เทฺว สํฆาทิเสสา อาปตฺติโย อาปชฺชติ เทฺวมาสปฺปฏิจฺฉนฺนาโย, เอโก มาโส สรมานปฺปฏิจฺฉนฺโน, เอโก มาโส อสฺสรมานปฺปฏิจฺฉนฺโน.\csromanspacer{} โส สํฆํ ทฺวินฺนํ อาปตฺตีนํ เทฺวมาสปฺปฏิจฺฉนฺนานํ เทฺวมาสปริวาสํ ยาจติ, ตสฺส สํโฆ ทฺวินฺนํ อาปตฺตีนํ เทฺวมาสปฺปฏิจฺฉนฺนานํ เทฺวมาสปริวาสํ เทติ, ตสฺส ปริวสนฺตสฺส อญฺโญ ภิกฺขุ อาคจฺฉติ พหุสฺสุโต อาคตาคโม ธมฺมธโร วินยธโร มาติกาธโร ปณฺฑิโต วิยตฺโต เมธาวี ลชฺชี กุกฺกุจฺจโก \csromanfolio{151}สิกฺขากาโม, โส เอวํ วเทติ “กิํ อยํ อาวุโส ภิกฺขุ อาปนฺโน, กิสฺสายํ ภิกฺขุ ปริวสตี”ติ.\csromanspacer{} เต เอวํ วเทนฺติ “อยํ อาวุโส ภิกฺขุ เทฺว สํฆาทิเสสา อาปตฺติโย อาปชฺชิ เทฺวมาสปฺปฏิจฺฉนฺนาโย, เอโก มาโส สรมานปฺปฏิจฺฉนฺโน, เอโก มาโส อสฺสรมานปฺปฏิจฺฉนฺโน, โส สํฆํ ทฺวินฺนํ อาปตฺตีนํ เทฺวมาสปฺปฏิจฺฉนฺนานํ เทฺวมาสปริวาทํ ยาจิ, ตสฺส สํโฆ ทฺวินฺนํ อาปตฺตีนํ เทฺวมาสปฺปฏิจฺฉนฺนานํ เทฺวมาสปริวาสํ อทาสิ, ตาโย อยํ อาวุโส ภิกฺขุ อาปนฺโน, ตาสายํ ภิกฺขุ ปริวสตี”ติ.\csromanspacer{} โส เอวํ วเทติ “ยฺวายํ อาวุโส มาโส สรมานปฺปฏิจฺฉนฺโน, ธมฺมิกํ ตสฺส มาสสฺส ปริวาสทานํ, ธมฺมตฺตา รุหติ.\csromanspacer{} โย จ ขฺวายํ อาวุโส มาโส อสฺสรมานปฺปฏิจฺฉนฺโน, อธมฺมิกํ ตสฺส มาสสฺส ปริวาสทานํ, อธมฺมตฺตา น รุหติ, เอกสฺส อาวุโส มาสสฺส ภิกฺขุ มานตฺตารโห”ติ.}

\proseitem{155}{อิธ ปน ภิกฺขเว ภิกฺขุ เทฺว สํฆาทิเสสา อาปตฺติโย อาปชฺชติ เทฺวมาสปฺปฏิจฺฉนฺนาโย, เอโก มาโส นิพฺเพมติกปฺปฏิจฺฉนฺโน, เอโก มาโส เวมติกปฺปฏิจฺฉนฺโน.\csromanspacer{} โส สํฆํ ทฺวินฺนํ อาปตฺตีนํ เทฺวมาสปฺปฏิจฺฉนฺนานํ เทฺวมาสปริวาสํ ยาจติ, ตสฺส สํโฆ ทฺวินฺนํ อาปตฺตีนํ เทฺวมาสปฺปฏิจฺฉนฺนานํ เทฺวมาสปริวาสํ เทติ, ตสฺส ปริวสนฺตสฺส อญฺโญ ภิกฺขุ อาคจฺฉติ พหุสฺสุโต อาคตาคโม ธมฺมธโร วินยธโร มาติกาธโร ปณฺฑิโต วิยตฺโต เมธาวี ลชฺชี กุกฺกุจฺจโก สิกฺขากาโม, โส เอวํ วเทติ “กิํ อยํ อาวุโส ภิกฺขุ อาปนฺโน, กิสฺสายํ ภิกฺขุ ปริวสตี”ติ.\csromanspacer{} เต เอวํ วเทนฺติ “อยํ อาวุโส ภิกฺขุ เทฺว สํฆาทิเสสา อาปตฺติโย อาปชฺชิ เทฺวมาสปฺปฏิจฺฉนฺนาโย, เอโก มาโส นิพฺเพมติกปฺปฏิจฺฉนฺโน, เอโก มาโส เวมติกปฺปฏิจฺฉนฺโน, โส สํฆํ ทฺวินฺนํ อาปตฺตีนํ เทฺวมาสปฺปฏิจฺฉนฺนานํ เทฺวมาสปริวาสํ ยาจิ, ตสฺส สํโฆ ทฺวินฺนํ อาปตฺตีนํ เทฺวมาสปฺปฏิจฺฉนฺนานํ เทฺวมาสปริวาสํ อทาสิ, ตาโย อยํ อาวุโส ภิกฺขุ อาปนฺโน, ตาสายํ ภิกฺขุ ปริวสตี”ติ.\csromanspacer{} โส เอวํ วเทติ “ยฺวายํ อาวุโส มาโส นิพฺเพมติกปฺปฏิจฺฉนฺโน, ธมฺมิกํ ตสฺส มาสสฺส ปริวาสทานํ, ธมฺมตฺตา รุหติ.\csromanspacer{} โย จ ขฺวายํ อาวุโส มาโส เวมติกปฺปฏิจฺฉนฺโน, อธมฺมิกํ ตสฺส มาสสฺส ปริวาสทานํ, อธมฺมตฺตา น รุหติ, เอกสฺส อาวุโส มาสสฺส ภิกฺขุ มานตฺตารโห”ติ.}

\csromanmatikaanchor{mtk.86}
\csromantocmarkref{subsection}{สุทฺธนฺตปริวาส}{mtk.86}
\bookmark[dest={mtk.86},level=2]{สุทฺธนฺตปริวาส}
\csromanheader{2}{\textbf{สุทฺธนฺตปริวาส}}
\proseitem{156}{\csromanfolio{152}เตน โข ปน สมเยน อญฺญตโร ภิกฺขุ สมฺพหุลา สํฆาทิเสสา อาปตฺติโย อาปนฺโน โหติ, โส อาปตฺติปริยนฺตํ น ชานาติ, รตฺติปริยนฺตํ น ชานาติ, อาปตฺติปริยนฺตํ นสฺสรติ, รตฺติปริยนฺตํ นสฺสรติ, อาปตฺติปริยนฺเต เวมติโก, รตฺติปริยนฺเต เวมติโก, โส ภิกฺขูนํ อาโรเจสิ “อหํ อาวุโส สมฺพหุลา สํฆาทิเสสา อาปตฺติโย อาปชฺชิํ, อาปตฺติปริยนฺตํ น ชานามิ, รตฺติปริยนฺตํ น ชานามิ, อาปตฺติปริยนฺตํ นสฺสรามิ, รตฺติปริยนฺตํ นสฺสรามิ, อาปตฺติปริยนฺเต เวมติโก, รตฺติปริยนฺเต เวมติโก, กถํ นุ โข มยา ปฏิปชฺชิตพฺพนฺ”ติ.\csromanspacer{} ภควโต เอตมตฺถํ อาโรเจสุํ.\csromanspacer{} เตน หิ ภิกฺขเว สํโฆ ตสฺส ภิกฺขุโน ตาสํ อาปตฺตีนํ สุทฺธนฺตปริวาสํ เทตุ.\csromanspacer{} เอวญฺจ ปน ภิกฺขเว ทาตพฺโพ\csromandash{}}

\prose{เตน ภิกฺขเว ภิกฺขุนา สํฆํ อุปสงฺกมิตฺวา ฯเปฯ เอวมสฺส วจนีโย “อหํ ภนฺเต สมฺพหุลา สํฆาทิเสสา อาปตฺติโย อาปชฺชิํ, อาปตฺติปริยนฺตํ น ชานามิ, รตฺติปริยนฺตํ น ชานามิ, อาปตฺติปริยนฺตํ นสฺสรามิ, รตฺติปริยนฺตํ นสฺสรามิ, อาปตฺติปริยนฺเต เวมติโก, รตฺติปริยนฺเต เวมติโก, โสหํ ภนฺเต สํฆํ ตาสํ อาปตฺตีนํ สุทฺธนฺตปริวาสํ ยาจามี”ติ.\csromanspacer{} ทุติยมฺปิ ยาจิตพฺโพ.\csromanspacer{} ตติยมฺปิ ยาจิตพฺโพ.\csromanspacer{} พฺยตฺเตน ภิกฺขุนา ปฏิพเลน สํโฆ ญาเปตพฺโพ\csromandash{}}

\proseitem{157}{“สุณาตุ เม ภนฺเต สํโฆ, อยํ อิตฺถนฺนาโม ภิกฺขุ สมฺพหุลา สํฆาทิเสสา อาปตฺติโย อาปชฺชิ, อาปตฺติปริยนฺตํ น ชานาติ, รตฺติปริยนฺตํ น ชานาติ, อาปตฺติปริยนฺตํ นสฺสรติ, รตฺติปริยนฺตํ นสฺสรติ, อาปตฺติปริยนฺเต เวมติโก, รตฺติปริยนฺเต เวมติโก โส สํฆํ ตาสํ อาปตฺตีนํ สุทฺธนฺตปริวาสํ ยาจติ.\csromanspacer{} ยทิ สํฆสฺส ปตฺตกลฺลํ, สํโฆ อิตฺถนฺนามสฺส ภิกฺขุโน ตาสํ อาปตฺตีนํ สุทฺธนฺตปริวาสํ ทเทยฺย, เอสา ญตฺติ.}

\prose{สุณาตุ เม ภนฺเต สํโฆ, อยํ อิตฺถนฺนาโม ภิกฺขุ สมฺพหุลา สํฆาทิเสสา อาปตฺติโย อาปชฺชิ, อาปตฺติปริยนฺตํ น ชานาติ, รตฺติปริยนฺตํ น ชานาติ, อาปตฺติปริยนฺตํ นสฺสรติ, รตฺติปริยนฺตํ นสฺสรติ, อาปตฺติปริยนฺเต เวมติโก, รตฺติปริยนฺเต เวมติโก, โส สํฆํ \csromanfolio{153}ตาสํ อาปตฺตีนํ สุทฺธนฺตปริวาสํ ยาจติ.\csromanspacer{} สํโฆ อิตฺถนฺนามสฺส ภิกฺขุโน ตาสํ อาปตฺตีนํ สุทฺธนฺตปริวาสํ เทติ, ยสฺสยสฺมโต ขมติ อิตฺถนฺนามสฺส ภิกฺขุโน ตาสํ อาปตฺตีนํ สุทฺธนฺตปริวาสสฺส ทานํ, โส ตุณฺหสฺส, ยสฺส นกฺขมติ, โส ภาเสยฺย.}

\prose{ทุติยมฺปิ เอตมตฺถํ วทามิ ฯเปฯ.\csromanspacer{} ตติยมฺปิ เอตมตฺถํ วทามิ ฯเปฯ.}

\prose{ทินฺโน สํเฆน อิตฺถนฺนามสฺส ภิกฺขุโน ตาสํ อาปตฺตีนํ สุทฺธนฺตปริวาโส, ขมติ สํฆสฺส, ตสฺมา ตุณฺหี, เอวเมตํ ธารยามี”ติ.}

\proseitem{158}{เอวํ โข ภิกฺขเว สุทฺธนฺตปริวาโส ทาตพฺโพ, เอวํ ปริวาโส ทาตพฺโพ.\csromanspacer{} กถญฺจ ภิกฺขเว สุทฺธนฺตปริวาโส ทาตพฺโพ.\csromanspacer{} อาปตฺติปริยนฺตํ น ชานาติ, รตฺติปริยนฺตํ น ชานาติ, อาปตฺติปริยนฺตํ นสฺสรติ, รตฺติปริยนฺตํ นสฺสรติ, อาปตฺติปริยนฺเต เวมติโก, รตฺติปริยนฺเต เวมติโก.\csromanspacer{} สุทฺธนฺตปริวาโส ทาตพฺโพ.}

\prose{อาปตฺติปริยนฺตํ ชานาติ, รตฺติปริยนฺตํ น ชานาติ, อาปตฺติปริยนฺตํ สรติ, รตฺติปริยนฺตํ นสฺสรติ, อาปตฺติปริยนฺเต นิพฺเพมติโก, รตฺติปริยนฺเต เวมาติโก.\csromanspacer{} สุทฺธนฺตปริวาโส ทาตพฺโพ.}

\prose{อาปตฺติปริยนฺตํ เอกจฺจํ ชานาติ, เอกจฺจํ น ชานาติ, รตฺติปริยนฺตํ น ชานาติ, อาปตฺติปริยนฺตํ เอกจฺจํ สรติ, เอกจฺจํ นสฺสรติ, รตฺติปริยนฺตํ นสฺสรติ, อาปตฺติปริยนฺเต เอกจฺเจ เวมติโก, เอกจฺเจ นิพฺเพมติโก, รตฺติปริยนฺเต เวมติโก.\csromanspacer{} สุทฺธนฺตปริวาโส ทาตพฺโพ.}

\prose{อาปตฺติปริยนฺตํ น ชานาติ, รตฺติปริยนฺตํ เอกจฺจํ ชานาติ, เอกจฺจํ น ชานาติ, อาปตฺติปริยนฺตํ นสฺสรติ, รตฺติปริยนฺตํ เอกจฺจํ สรติ, เอกจฺจํ นสฺสรติ, อาปตฺติปริยนฺเต เวมติโก, รตฺติปริยนฺเต เอกจฺเจ เวมติโก, เอกจฺเจ นิพฺเพมติโก.\csromanspacer{} สุทฺธนฺตปริวาโส ทาตพฺโพ.}

\prose{อาปตฺติปริยนฺตํ ชานาติ, รตฺติปริยนฺตํ เอกจฺจํ ชานาติ, เอกจฺจํ น ชานาติ, อาปตฺติปริยนฺตํ สรติ, รตฺติปริยนฺตํ เอกจฺจํ สรติ, เอกจฺจํ นสฺสรติ, อาปตฺติปริยนฺเต นิพฺเพมติโก, รตฺติปริยนฺเต เอกจฺเจ เวมติโก, เอกจฺเจ นิพฺเพมติโก.\csromanspacer{} สุทฺธนฺตปริวาโส ทาตพฺโพ.}

\prose{อาปตฺติปริยนฺตํ เอกจฺจํ ชานาติ, เอกจฺจํ น ชานาติ, รตฺติปริยนฺตํ เอกจฺจํ ชานาติ, เอกจฺจํ น ชานาติ, อาปตฺติปริยนฺตํ เอกจฺจํ สรติ, เอกจฺจํ \csromanfolio{154}นสฺสรติ, รตฺติปริยนฺตํ เอกจฺจํ สรติ, เอกจฺจํ นสฺสรติ, อาปตฺติปริยนฺเต เอกจฺเจ เวมติโก, เอกจฺเจ นิพฺเพมติโก, รตฺติปริยนฺเต เอกจฺเจ เวมติโก, เอกจฺเจ นิพฺเพมติโก.\csromanspacer{} สุทฺธนฺตปริวาโส ทาตพฺโพ.\csromanspacer{} เอวํ โข ภิกฺขเว สุทฺธนฺตปริวาโส ทาตพฺโพ.}

\proseitem{159}{กถญฺจ ภิกฺขเว ปริวาโส ทาตพฺโพ, อาปตฺติปริยนฺตํ ชานาติ, รตฺติปริยนฺตํ ชานาติ, อาปตฺติปริยนฺตํ สรติ, รตฺติปริยนฺตํ สรติ, อาปตฺติปริยนฺเต นิพฺเพมติโก, รตฺติปริยนฺเต นิพฺเพมติโก.\csromanspacer{} ปริวาโส ทาตพฺโพ.}

\prose{อาปตฺติปริยนฺตํ น ชานาติ, รตฺติปริยนฺตํ ชานาติ, อาปตฺติปริยนฺตํ นสฺสรติ, รตฺติปริยนฺตํ สรติ, อาปตฺติปริยนฺเต เวมติโก, รตฺติปริยนฺเต นิพฺเพมติโก.\csromanspacer{} ปริวาโส ทาตพฺโพ.}

\prose{อาปตฺติปริยนฺตํ เอกจฺจํ ชานาติ, เอกจฺจํ น ชานาติ, รตฺติปริยนฺตํ ชานาติ, อาปตฺติปริยนฺตํ เอกจฺจํ สรติ, เอกจฺจํ นสฺสรติ, รตฺติปริยนฺตํ สรติ, อาปตฺติปริยนฺเต เอกจฺเจ เวมติโก, เอกจฺเจ นิพฺเพมติโก, รตฺติปริยนฺเต นิพฺเพมติโก.\csromanspacer{} ปริวาโส ทาตพฺโพ.\csromanspacer{} เอวํ โข ภิกฺขเว ปริวาโส ทาตพฺโพ.}

\nitthitam{ปริวาโส นิฏฺฐิโต.}
\csromanmatikaanchor{mtk.87}
\csromantocmarkref{section}{3. จตฺตาลีสก}{mtk.87}
\bookmark[dest={mtk.87},level=1]{3. จตฺตาลีสก}
\csromanheader{1}{3. จตฺตาลีสก}
\proseitem{160}{เตน โข ปน สมเยน อญฺญตโร ภิกฺขุ ปริวสนฺโต วิพฺภมิ.\csromanspacer{} โส ปุน ปจฺจาคนฺตฺวา ภิกฺขู อุปสมฺปทํ ยาจิ.\csromanspacer{} ภควโต เอตมตฺถํ อาโรเจสุํ.}

\prose{อิธ ปน ภิกฺขเว ภิกฺขุ ปริวสนฺโต วิพฺภมติ, วิพฺภนฺตกสฺส ภิกฺขเว ปริวาโส น รุหติ.\csromanspacer{} โส เจ ปุน อุปสมฺปชฺชติ, ตสฺส ตเทว ปุริมํ ปริวาสทานํ, โย ปริวาโส ทินฺโน สุทินฺโน, โย ปริวุตฺโถ สุปริวุตฺโถ, อวเสโส ปริวสิตพฺโพ. (1)}

\prose{อิธ ปน ภิกฺขเว ภิกฺขุ ปริวสนฺโต สามเณโร โหติ, สามเณรสฺส ภิกฺขเว ปริวาโส น รุหติ.\csromanspacer{} โส เจ ปุน อุปสมฺปชฺชติ, \csromanfolio{155}ตสฺส ตเทว ปุริมํ ปริวาสทานํ, โย ปริวาโส ทินฺโน สุทินฺโน, โย ปริวุตฺโถ สุปริวุตฺโถ, อวเสโส ปริวสิตพฺโพ. (2)}

\prose{อิธ ปน ภิกฺขเว ภิกฺขุ ปริวสนฺโต อุมฺมตฺตโก โหติ, อุมฺมตฺตกสฺส ภิกฺขเว ปริวาโส น รุหติ.\csromanspacer{} โส เจ ปุน อนุมฺมตฺตโก โหติ, ตสฺส ตเทว ปุริมํ ปริวาสทานํ, โย ปริวาโส ทินฺโน สุทินฺโน, โย ปริวุตฺโถ สุปริวุตฺโถ, อวเสโส ปริวสิตพฺโพ. (3)}

\prose{อิธ ปน ภิกฺขเว ภิกฺขุ ปริวสนฺโต ขิตฺตจิตฺโต โหติ, ขิตฺตจิตฺตสฺส ภิกฺขเว ปริวาโส น รุหติ.\csromanspacer{} โส เจ ปุน อขิตฺตจิตฺโต โหติ, ตสฺส ตเทว ปุริมํ ปริวาสทานํ, โย ปริวาโส ทินฺโน สุทินฺโน, โย ปริวุตฺโถ สุปริวุตฺโถ, อวเสโส ปริวสิตพฺโพ. (4)}

\prose{อิธ ปน ภิกฺขเว ภิกฺขุ ปริวสนฺโต เวทนาฏฺโฏ โหติ, เวทนาฏฺฏสฺส ภิกฺขเว ปริวาโส น รุหติ.\csromanspacer{} โส เจ ปุน อเวทนาฏฺโฏ โหติ, ตสฺส ตเทว ปุริมํ ปริวาสทานํ, โย ปริวาโส ทินฺโน สุทินฺโน, โย ปริวุตฺโถ สุปริวุตฺโถ, อวเสโส ปริวสิตพฺโพ. (5)}

\prose{อิธ ปน ภิกฺขเว ภิกฺขุ ปริวสนฺโต อาปตฺติยา อทสฺสเน อุกฺขิปิยฺยติ\footnote{อุกฺขิปิยติ (สฺยา), อุกฺขิปียติ (ก)}, อุกฺขิตฺตกสฺส ภิกฺขเว ปริวาโส น รุหติ.\csromanspacer{} โส เจ ปุน โอสาริยฺยติ\footnote{โอสาริยติ (สฺยา), โอสารียติ (ก)}, ตสฺส ตเทว ปุริมํ ปริวาสทานํ, โย ปริวาโส ทินฺโน สุทินฺโน, โย ปริวุตฺโถ สุปริวุตฺโถ, อวเสโส ปริวสิตพฺโพ. (6)}

\prose{อิธ ปน ภิกฺขเว ภิกฺขุ ปริวสนฺโต อาปตฺติยา อปฺปฏิกมฺเม อุกฺขิปิยฺยติ, อุกฺขิตฺตกสฺส ภิกฺขเว ปริวาโส น รุหติ.\csromanspacer{} โส เจ ปุน โอสาริยฺยติ, ตสฺส ตเทว ปุริมํ ปริวาสทานํ, โย ปริวาโส ทินฺโน สุทินฺโน, โย ปริวุตฺโถ สุปริวุตฺโถ, อวเสโส ปริวสิตพฺโพ. (7)}

\prose{อิธ ปน ภิกฺขเว ภิกฺขุ ปริวสนฺโต ปาปิกาย ทิฏฺฐิยา อปฺปฏินิสฺสคฺเค อุกฺขิปิยฺยติ, อุกฺขิตฺตกสฺส ภิกฺขเว ปริวาโส น รุหติ.\csromanspacer{} โส เจ ปุน โอสาริยฺยติ, ตสฺส ตเทว ปุริมํ ปริวาสทานํ, โย ปริวาโส ทินฺโน สุทินฺโน, โย ปริวุตฺโถ สุปริวุตฺโถ, อวเสโส ปริวสิตพฺโพ. (8)}

\proseitem{161}{\csromanfolio{156}อิธ ปน ภิกฺขเว ภิกฺขุ มูลายปฏิกสฺสนารโห วิพฺภมติ, วิพฺภนฺตกสฺส ภิกฺขเว มูลายปฏิกสฺสนา น รุหติ.\csromanspacer{} โส เจ ปุน อุปสมฺปชฺชติ, ตสฺส ตเทว ปุริมํ ปริวาสทานํ, โย ปริวาโส ทินฺโน สุทินฺโน, โย ปริวุตฺโถ สุปริวุตฺโถ, โส ภิกฺขุ มูลาย ปฏิกสฺสิตพฺโพ. (9)}

\prose{อิธ ปน ภิกฺขเว ภิกฺขุ มูลายปฏิกสฺสนารโห สามเณโร โหติ ฯเปฯ อุมฺมตฺตโก โหติ ฯเปฯ ขิตฺตจิตฺโต โหติ ฯเปฯ เวทนาฏฺโฏ โหติ ฯเปฯ อาปตฺติยา อทสฺสเน อุกฺขิปิยฺยติ ฯเปฯ อาปตฺติยา อปฺปฏิกมฺเม อุกฺขิปิยฺยติ ฯเปฯ ปาปิกาย ทิฏฺฐิยา อปฺปฏินิสฺสคฺเค อุกฺขิปิยฺยติ, อุกฺขิตฺตกสฺส ภิกฺขเว มูลาย ปฏิกสฺสนา น รุหติ.\csromanspacer{} โส เจ ปุน โอสาริยฺยติ, ตสฺส ตเทว ปุริมํ ปริวาสทานํ, โย ปริวาโส ทินฺโน สุทินฺโน, โย ปริวุตฺโถ สุปริวุตฺโถ, โส ภิกฺขุ มูลาย ปฏิกสฺสิตพฺโพ. (16)}

\proseitem{162}{อิธ ปน ภิกฺขเว ภิกฺขุ มานตฺตารโห วิพฺภมติ, วิพฺภนฺตกสฺส ภิกฺขเว มานตฺตทานํ น รุหติ.\csromanspacer{} โส เจ ปุน อุปสมฺปชฺชติ, ตสฺส ตเทว ปุริมํ ปริวาสทานํ, โย ปริวาโส ทินฺโน สุทินฺโน, โย ปริวุตฺโถ สุปริวุตฺโถ, ตสฺส ภิกฺขุโน มานตฺตํ ทาตพฺพํ. (17)}

\prose{อิธ ปน ภิกฺขเว ภิกฺขุ มานตฺตารโห สามเณโร โหติ ฯเปฯ อุมฺมตฺตโก โหติ ฯเปฯ ขิตฺตจิตฺโต โหติ ฯเปฯ เวทนาฏฺโฏ โหติ ฯเปฯ อาปตฺติยา อทสฺสเน อุกฺขิปิยฺยติ ฯเปฯ อาปตฺติยา อปฺปฏิกมฺเม อุกฺขิปิยฺยติ ฯเปฯ ปาปิกาย ทิฏฺฐิยา อปฺปฏินิสฺสคฺเค อุกฺขิปิยฺยติ, อุกฺขิตฺตกสฺส ภิกฺขเว มานตฺตทานํ น รุหติ.\csromanspacer{} โส เจ ปุน โอสาริยฺยติ, ตสฺส ตเทว ปุริมํ ปริวาสทานํ, โย ปริวาโส ทินฺโน สุทินฺโน, โย ปริวุตฺโถ สุปริวุตฺโถ, ตสฺส ภิกฺขุโน มานตฺตํ ทาตพฺพํ. (24)}

\proseitem{163}{อิธ ปน ภิกฺขเว ภิกฺขุ มานตฺตํ จรนฺโต วิพฺภมติ, วิพฺภนฺตกสฺส ภิกฺขเว มานตฺตจริยา น รุหติ.\csromanspacer{} โส เจ ปุน อุปสมฺปชฺชติ, ตสฺส ตเทว ปุริมํ ปริวาสทานํ, โย ปริวาโส ทินฺโน สุทินฺโน, โย ปริวุตฺโถ สุปริวุตฺโถ, ยํ มานตฺตํ ทินฺนํ สุทินฺนํ, ยํ มานตฺตํ จิณฺณํ สุจิณฺณํ, อวเสสํ จริตพฺพํ. (25)}

\csromanmatikaanchor{mtk.88}
\csromantocmarkref{section}{4. ฉติํสก}{mtk.88}
\bookmark[dest={mtk.88},level=1]{4. ฉติํสก}
\prose{\csromanfolio{157}อิธ ปน ภิกฺขเว ภิกฺขุ มานตฺตํ จรนฺโต สามเณโร โหติ ฯเปฯ อุมฺมตฺตโก โหติ ฯเปฯ ขิตฺตจิตฺโต โหติ ฯเปฯ เวทนาฏฺโฏ โหติ ฯเปฯ อาปตฺติยา อทสฺสเน อุกฺขิปิยฺยติ ฯเปฯ อาปตฺติยา อปฺปฏิกมฺเม อุกฺขิปิยฺยติ ฯเปฯ ปาปิกาย ทิฏฺฐิยา อปฺปฏินิสฺสคฺเค อุกฺขิปิยฺยติ, อุกฺขิตฺตกสฺส ภิกฺขเว มานตฺตจริยา น รุหติ.\csromanspacer{} โส เจ ปุน โอสาริยฺยติ, ตสฺส ตเทว ปุริมํ ปริวาสทานํ, โย ปริวาโส ทินฺโน สุทินฺโน, โย ปริวุตฺโถ สุปริวุตฺโถ, ยํ มานตฺตํ ทินฺนํ สุทินฺนํ, ยํ มานตฺตํ จิณฺณํ สุจิณฺณํ, อวเสสํ จริตพฺพํ. (32)}

\proseitem{164}{อิธ ปน ภิกฺขเว ภิกฺขุ อพฺภานารโห วิพฺภมติ, วิพฺภนฺตกสฺส ภิกฺขเว อพฺภานํ น รุหติ.\csromanspacer{} โส เจ ปุน อุปสมฺปชฺชติ, ตสฺส ตเทว ปุริมํ ปริวาสทานํ, โย ปริวาโส ทินฺโน สุทินฺโน, โย ปริวุตฺโถ สุปริวุตฺโถ, ยํ มานตฺตํ ทินฺนํ สุทินฺนํ, ยํ มานตฺตํ จิณฺณํ สุจิณฺณํ, โส ภิกฺขุ อพฺเภตพฺโพ. (33)}

\prose{อิธ ปน ภิกฺขเว ภิกฺขุ อพฺภานารโห สามเณโร โหติ ฯเปฯ อุมฺมตฺตโก โหติ ฯเปฯ ขิตฺตจิตฺโต โหติ ฯเปฯ เวทนาฏฺโฏ โหติ ฯเปฯ อาปตฺติยา อทสฺสเน อุกฺขิปิยฺยติ ฯเปฯ อาปตฺติยา อปฺปฏิกมฺเม อุกฺขิปิยฺยติ ฯเปฯ ปาปิกาย ทิฏฺฐิยา อปฺปฏินิสฺสคฺเค อุกฺขิปิยฺยติ, อุกฺขิตฺตกสฺส ภิกฺขเว อพฺภานํ น รุหติ.\csromanspacer{} โส เจ ปุน โอสาริยฺยติ, ตสฺส ตเทว ปุริมํ ปริวาสทานํ, โย ปริวาโส ทินฺโน สุทินฺโน, โย ปริวุตฺโถ สุปริวุตฺโถ, ยํ มานตฺตํ ทินฺนํ สุทินฺนํ, ยํ มานตฺตํ จิณฺณํ สุจิณฺณํ, โส ภิกฺขุ อพฺเภตพฺโพ. (40)}

\vspace{\tipitakaparskip}
\csromancenter{จตฺตาลีสกํ สมตฺตํ.}
\csromanheader{2}{4. ฉตฺติํสก}
\proseitem{165}{อิธ ปน ภิกฺขเว ภิกฺขุ ปริวสนฺโต อนฺตรา สมฺพหุลา สํฆาทิเสสา อาปตฺติโย อาปชฺชติ ปริมาณา\footnote{ปริมาณาโย (สี, สฺยา)} อปฺปฏิจฺฉนฺนาโย, โส ภิกฺขุ มูลาย ปฏิกสฺสิตพฺโพ. (1)}

\prose{\csromanfolio{158}อิธ ปน ภิกฺขเว ภิกฺขุ ปริวสนฺโต อนฺตรา สมฺพหุลา สํฆาทิเสสา อาปตฺติโย อาปชฺชติ ปริมาณา ปฏิจฺฉนฺนาโย, โส ภิกฺขุ มูลาย ปฏิกสฺสิตพฺโพ, ยถาปฏิจฺฉนฺนานญฺจสฺส อาปตฺตีนํ ปุริมาย อาปตฺติยา สโมธานปริวาโส ทาตพฺโพ. (2)}

\prose{อิธ ปน ภิกฺขเว ภิกฺขุ ปริวสนฺโต อนฺตรา สมฺพหุลา สํฆาทิเสสา อาปตฺติโย อาปชฺชติ ปริมาณา ปฏิจฺฉนฺนาโยปิ อปฺปฏิจฺฉนฺนาโยปิ, โส ภิกฺขุ มูลาย ปฏิกสฺสิตพฺโพ, ยถาปฏิจฺฉนฺนานญฺจสฺส อาปตฺตีนํ ปุริมาย อาปตฺติยา สโมธานปริวาโส ทาตพฺโพ. (3)}

\prose{อิธ ปน ภิกฺขเว ภิกฺขุ ปริวสนฺโต อนฺตรา สมฺพหุลา สํฆาทิเสสา อาปตฺติโย อาปชฺชติ อปริมาณา\footnote{อปริมาณาโย (สี, สฺยา)} อปฺปฏิจฺฉนฺนาโย ฯเปฯ อปริมาณา ปฏิจฺฉนฺนาโย ฯเปฯ อปริมาณา ปฏิจฺฉนฺนาโยปิ อปฺปฏิจฺฉนฺนาโยปิ ฯเปฯ ปริมาณาโยปิ อปริมาณาโยปิ อปฺปฏิจฺฉนฺนาโย ฯเปฯ ปริมาณาโยปิ อปริมาณาโยปิ ปฏิจฺฉนฺนาโย ฯเปฯ ปริมาณาโยปิ อปริมาณาโยปิ ปฏิจฺฉนฺนาโยปิ อปฺปฏิจฺฉนฺนาโยปิ, โส ภิกฺขุ มูลาย ปฏิกสฺสิตพฺโพ.\csromanspacer{} ยถาปฏิจฺฉนฺนานญฺจสฺส อาปตฺตีนํ ปุริมาย อาปตฺติยา สโมธานปริวาโส ทาตพฺโพ. (9)}

\prose{อิธ ปน ภิกฺขเว ภิกฺขุ มานตฺตารโห ฯเปฯ มานตฺตํ จรนฺโต ฯเปฯ (ยถา ปริวาสํ, ตถา วิตฺถาเรตพฺพํ.) อพฺภานารโห อนฺตรา สมฺพหุลา สํฆาทิเสสา อาปตฺติโย อาปชฺชติ ปริมาณา อปฺปฏิจฺฉนฺนาโย ฯเปฯ ปริมาณา ปฏิจฺฉนฺนาโย ฯเปฯ ปริมาณา ปฏิจฺฉนฺนาโยปิ อปฺปฏิจฺฉนฺนาโยปิ ฯเปฯ อปริมาณา อปฺปฏิจฺฉนฺนาโย ฯเปฯ อปริมาณา ปฏิจฺฉนฺนาโย ฯเปฯ อปริมาณา ปฏิจฺฉนฺนาโยปิ อปฺปฏิจฺฉนฺนาโยปิ ฯเปฯ ปริมาณาโยปิ อปริมาณาโยปิ อปฺปฏิจฺฉนฺนาโย ฯเปฯ ปริมาณาโยปิ อปริมาณาโยปิ ปฏิจฺฉนฺนาโย ฯเปฯ ปริมาณาโยปิ อปริมาณาโยปิ ปฏิจฺฉนฺนาโยปิ อปฺปฏิจฺฉนฺนาโยปิ, โส ภิกฺขุ มูลาย ปฏิกสฺสิตพฺโพ.\csromanspacer{} ยถาปฏิจฺฉนฺนานญฺจสฺส อาปตฺตีนํ ปุริมาย อาปตฺติยา สโมธานปริวาโส ทาตพฺโพ. (36)}

\vspace{\tipitakaparskip}
\csromancenter{ฉตฺติํสกํ สมตฺตํ.}
\csromanmatikaanchor{mtk.89}
\csromantocmarkref{section}{5. มานตฺตสต}{mtk.89}
\bookmark[dest={mtk.89},level=1]{5. มานตฺตสต}
\csromanheader{1}{5. มานตฺตสต}
\proseitem{166}{\csromanfolio{159}อิธ ปน ภิกฺขเว ภิกฺขุ สมฺพหุลา สํฆาทิเสสา อาปตฺติโย อาปชฺชิตฺวา อปฺปฏิจฺฉาเทตฺวา วิพฺภมติ.\csromanspacer{} โส ปุน\footnote{โส เจ ปุน (ก)} อุปสมฺปนฺโน ตา อาปตฺติโย นจฺฉาเทติ, ตสฺส ภิกฺขเว ภิกฺขุโน มานตฺตํ ทาตพฺพํ. (1)}

\prose{อิธ ปน ภิกฺขเว ภิกฺขุ สมฺพหุลา สํฆาทิเสสา อาปตฺติโย อาปชฺชิตฺวา อปฺปฏิจฺฉาเทตฺวา วิพฺภมติ.\csromanspacer{} โส ปุน อุปสมฺปนฺโน ตา อาปตฺติโย ฉาเทติ, ตสฺส ภิกฺขเว ภิกฺขุโน ปจฺฉิมสฺมิํ อาปตฺติกฺขนฺเธ ยถาปฏิจฺฉนฺเน ปริวาสํ ทตฺวา มานตฺตํ ทาตพฺพํ. (2)}

\prose{อิธ ปน ภิกฺขเว ภิกฺขุ สมฺพหุลา สํฆาทิเสสา อาปตฺติโย อาปชฺชิตฺวา ปฏิจฺฉาเทตฺวา วิพฺภมติ.\csromanspacer{} โส ปุน อุปสมฺปนฺโน ตา อาปตฺติโย นจฺฉาเทติ, ตสฺส ภิกฺขเว ภิกฺขุโน ปุริมสฺมิํ อาปตฺติกฺขนฺเธ ยถาปฏิจฺฉนฺเน ปริวาสํ ทตฺวา มานตฺตํ ทาตพฺพํ. (3)}

\prose{อิธ ปน ภิกฺขเว ภิกฺขุ สมฺพหุลา สํฆาทิเสสา อาปตฺติโย อาปชฺชิตฺวา ปฏิจฺฉาเทตฺวา วิพฺภมติ.\csromanspacer{} โส ปุน อุปสมฺปนฺโน ตา อาปตฺติโย ฉาเทติ, ตสฺส ภิกฺขเว ภิกฺขุโน ปุริมสฺมิํ จ ปจฺฉิมสฺมิํ จ อาปตฺติกฺขนฺเธ ยถาปฏิจฺฉนฺเน ปริวาสํ ทตฺวา มานตฺตํ ทาตพฺพํ. (4)}

\proseitem{167}{อิธ ปน ภิกฺขเว ภิกฺขุ สมฺพหุลา สํฆาทิเสสา อาปตฺติโย อาปชฺชติ, ตสฺส โหนฺติ อาปตฺติโย ปฏิจฺฉนฺนาโยปิ อปฺปฏิจฺฉนฺนาโยปิ, โส วิพฺภมิตฺวา ปุน อุปสมฺปนฺโน ยา อาปตฺติโย ปุพฺเพ ฉาเทสิ, ตา อาปตฺติโย ปจฺฉา นจฺฉาเทติ, ยา อาปตฺติโย ปุพฺเพ นจฺฉาเทสิ, ตา อาปตฺติโย ปจฺฉา นจฺฉาเทติ.\csromanspacer{} ตสฺส ภิกฺขเว ภิกฺขุโน ปุริมสฺมิํ อาปตฺติกฺขนฺเธ ยถาปฏิจฺฉนฺเน ปริวาสํ ทตฺวา มานตฺตํ ทาตพฺพํ. (5)}

\prose{อิธ ปน ภิกฺขเว ภิกฺขุ สมฺพหุลา สํฆาทิเสสา อาปตฺติโย อาปชฺชติ, ตสฺส โหนฺติ อาปตฺติโย ปฏิจฺฉนฺนาโยปิ อปฺปฏิจฺฉนฺนาโยปิ.\csromanspacer{} โส วิพฺภมิตฺวา ปุน อุปสมฺปนฺโน ยา อาปตฺติโย ปุพฺเพ ฉาเทสิ, ตา อาปตฺติโย ปจฺฉา นจฺฉาเทติ.\csromanspacer{} ยา อาปตฺติโย ปุพฺเพ นจฺฉาเทสิ, ตา อาปตฺติโย ปจฺฉา ฉาเทติ.\csromanspacer{} ตสฺส ภิกฺขเว ภิกฺขุโน ปุริมสฺมิํ จ ปจฺฉิมสฺมิํ จ อาปตฺติกฺขนฺเธ ยถาปฏิจฺฉนฺเน ปริวาสํ ทตฺวา มานตฺตํ ทาตพฺพํ. (6)}

\prose{\csromanfolio{160}อิธ ปน ภิกฺขเว ภิกฺขุ สมฺพหุลา สํฆาทิเสสา อาปตฺติโย อาปชฺชติ, ตสฺส โหนฺติ อาปตฺติโย ปฏิจฺฉนฺนาโยปิ อปฺปฏิจฺฉนฺนาโยปิ.\csromanspacer{} โส วิพฺภมิตฺวา ปุน อุปสมฺปนฺโน ยา อาปตฺติโย ปุพฺเพ ฉาเทสิ, ตา อาปตฺติโย ปจฺฉา ฉาเทติ.\csromanspacer{} ยา อาปตฺติโย ปุพฺเพ นจฺฉาเทสิ, ตา อาปตฺติโย ปจฺฉา นจฺฉาเทติ.\csromanspacer{} ตสฺส ภิกฺขเว ภิกฺขุโน ปุริมสฺมิํ จ ปจฺฉิมสฺมิํ จ อาปตฺติกฺขนฺเธ ยถาปฏิจฺฉนฺเน ปริวาสํ ทตฺวา มานตฺตํ ทาตพฺพํ. (7)}

\prose{อิธ ปน ภิกฺขเว ภิกฺขุ สมฺพหุลา สํฆาทิเสสา อาปตฺติโย อาปชฺชติ, ตสฺส โหนฺติ อาปตฺติโย ปฏิจฺฉนฺนาโยปิ อปฺปฏิจฺฉนฺนาโยปิ.\csromanspacer{} โส วิพฺภมิตฺวา ปุน อุปสมฺปนฺโน ยา อาปตฺติโย ปุพฺเพ ฉาเทสิ, ตา อาปตฺติโย ปจฺฉา ฉาเทติ.\csromanspacer{} ยา อาปตฺติโย ปุพฺเพ นจฺฉาเทสิ, ตา อาปตฺติโย ปจฺฉา ฉาเทติ.\csromanspacer{} ตสฺส ภิกฺขเว ภิกฺขุโน ปุริมสฺมิํ จ ปจฺฉิมสฺมิํ จ อาปตฺติกฺขนฺเธ ยถาปฏิจฺฉนฺเน ปริวาสํ ทตฺวา มานตฺตํ ทาตพฺพํ. (8)}

\proseitem{168}{อิธ ปน ภิกฺขเว ภิกฺขุ สมฺพหุลา สํฆาทิเสสา อาปตฺติโย อาปชฺชติ, เอกจฺจา อาปตฺติโย ชานาติ, เอกจฺจา อาปตฺติโย น ชานาติ.\csromanspacer{} ยา อาปตฺติโย ชานาติ, ตา อาปตฺติโย ฉาเทติ.\csromanspacer{} ยา อาปตฺติโย น ชานาติ, ตา อาปตฺติโย นจฺฉาเทติ.\csromanspacer{} โส วิพฺภมิตฺวา ปุน อุปสมฺปนฺโน ยา อาปตฺติโย ปุพฺเพ ชานิตฺวา ฉาเทสิ, ตา อาปตฺติโย ปจฺฉา ชานิตฺวา นจฺฉาเทติ.\csromanspacer{} ยา อาปตฺติโย ปุพฺเพ อชานิตฺวา นจฺฉาเทสิ, ตา อาปตฺติโย ปจฺฉา ชานิตฺวา นจฺฉาเทติ.\csromanspacer{} ตสฺส ภิกฺขเว ภิกฺขุโน ปุริมสฺมิํ อาปตฺติกฺขนฺเธ ยถาปฏิจฺฉนฺเน ปริวาสํ ทตฺวา มานตฺตํ ทาตพฺพํ. (9)}

\prose{อิธ ปน ภิกฺขเว ภิกฺขุ สมฺพหุลา สํฆาทิเสสา อาปตฺติโย อาปชฺชติ, เอกจฺจา อาปตฺติโย ชานาติ, เอกจฺจา อาปตฺติโย น ชานาติ.\csromanspacer{} ยา อาปตฺติโย ชานาติ, ตา อาปตฺติโย ฉาเทติ.\csromanspacer{} ยา อาปตฺติโย น ชานาติ, ตา อาปตฺติโย นจฺฉาเทติ.\csromanspacer{} โส วิพฺภมิตฺวา ปุน อุปสมฺปนฺโน ยา อาปตฺติโย ปุพฺเพ ชานิตฺวา ฉาเทสิ, ตา อาปตฺติโย ปจฺฉา ชานิตฺวา นจฺฉาเทติ.\csromanspacer{} ยา อาปตฺติโย ปุพฺเพ อชานิตฺวา นจฺฉาเทสิ, ตา อาปตฺติโย ปจฺฉา ชานิตฺวา ฉาเทติ. \csromanfolio{161}ตสฺส ภิกฺขเว ภิกฺขุโน ปุริมสฺมิํ จ ปจฺฉิมสฺมิํ จ อาปตฺติกฺขนฺเธ ยถาปฏิจฺฉนฺเน ปริวาสํ ทตฺวา มานตฺตํ ทาตพฺพํ. (10)}

\prose{อิธ ปน ภิกฺขเว ภิกฺขุ สมฺพหุลา สํฆาทิเสสา อาปตฺติโย อาปชฺชติ, เอกจฺจา อาปตฺติโย ชานาติ, เอกจฺจา อาปตฺติโย น ชานาติ.\csromanspacer{} ยา อาปตฺติโย ชานาติ, ตา อาปตฺติโย ฉาเทติ.\csromanspacer{} ยา อาปตฺติโย น ชานาติ, ตา อาปตฺติโย นจฺฉาเทติ.\csromanspacer{} โส วิพฺภมิตฺวา ปุน อุปสมฺปนฺโน ยา อาปตฺติโย ปุพฺเพ ชานิตฺวา ฉาเทสิ, ตา อาปตฺติโย ปจฺฉา ชานิตฺวา ฉาเทติ.\csromanspacer{} ยา อาปตฺติโย ปุพฺเพ อชานิตฺวา นจฺฉาเทสิ, ตา อาปตฺติโย ปจฺฉา ชานิตฺวา นจฺฉาเทติ.\csromanspacer{} ตสฺส ภิกฺขเว ภิกฺขุโน ปุริมสฺมิํ จ ปจฺฉิมสฺมิํ จ อาปตฺติกฺขนฺเธ ยถาปฏิจฺฉนฺเน ปริวาสํ ทตฺวา มานตฺตํ ทาตพฺพํ. (11)}

\prose{อิธ ปน ภิกฺขเว ภิกฺขุ สมฺพหุลา สํฆาทิเสสา อาปตฺติโย อาปชฺชติ, เอกจฺจา อาปตฺติโย ชานาติ, เอกจฺจา อาปตฺติโย น ชานาติ.\csromanspacer{} ยา อาปตฺติโย ชานาติ, ตา อาปตฺติโย ฉาเทติ.\csromanspacer{} ยา อาปตฺติโย น ชานาติ, ตา อาปตฺติโย นจฺฉาเทติ.\csromanspacer{} โส วิพฺภมิตฺวา ปุน อุปสมฺปนฺโน ยา อาปตฺติโย ปุพฺเพ ชานิตฺวา ฉาเทสิ, ตา อาปตฺติโย ปจฺฉา ชานิตฺวา ฉาเทติ.\csromanspacer{} ยา อาปตฺติโย ปุพฺเพ อชานิตฺวา นจฺฉาเทสิ, ตา อาปตฺติโย ปจฺฉา ชานิตฺวา ฉาเทติ.\csromanspacer{} ตสฺส ภิกฺขเว ภิกฺขุโน ปุริมสฺมิํ จ ปจฺฉิมสฺมิํ จ อาปตฺติกฺขนฺเธ ยถาปฏิจฺฉนฺเน ปริวาสํ ทตฺวา มานตฺตํ ทาตพฺพํ. (12)}

\proseitem{196}{อิธ ปน ภิกฺขเว ภิกฺขุ สมฺพหุลา สํฆาทิเสสา อาปตฺติโย อาปชฺชติ, เอกจฺจา อาปตฺติโย สรติ, เอกจฺจา อาปตฺติโย นสฺสรติ.\csromanspacer{} ยา อาปตฺติโย สรติ, ตา อาปตฺติโย ฉาเทติ.\csromanspacer{} ยา อาปตฺติโย นสฺสรติ, ตา อาปตฺติโย นจฺฉาเทติ.\csromanspacer{} โส วิพฺภมิตฺวา ปุน อุปสมฺปนฺโน ยา อาปตฺติโย ปุพฺเพ สริตฺวา ฉาเทสิ, ตา อาปตฺติโย ปจฺฉา สริตฺวา นจฺฉาเทติ.\csromanspacer{} ยา อาปตฺติโย ปุพฺเพ อสฺสริตฺวา นจฺฉาเทสิ, ตา อาปตฺติโย ปจฺฉา สริตฺวา นจฺฉาเทติ.\csromanspacer{} ตสฺส ภิกฺขเว ภิกฺขุโน ปุริมสฺมิํ อาปตฺติกฺขนฺเธ ยถาปฏิจฺฉนฺเน ปริวาสํ ทตฺวา มานตฺตํ ทาตพฺพํ. (13)}

\prose{อิธ ปน ภิกฺขเว ภิกฺขุ สมฺพหุลา สํฆาทิเสสา อาปตฺติโย อาปชฺชติ, เอกจฺจา อาปตฺติโย สรติ, เอกจฺจา อาปตฺติโย นสฺสรติ. \csromanfolio{162}ยา อาปตฺติโย สรติ, ตา อาปตฺติโย ฉาเทติ.\csromanspacer{} ยา อาปตฺติโย นสฺสรติ, ตา อาปตฺติโย นจฺฉาเทติ.\csromanspacer{} โส วิพฺภมิตฺวา ปุน อุปสมฺปนฺโน ยา อาปตฺติโย ปุพฺเพ สริตฺวา ฉาเทสิ, ตา อาปตฺติโย ปจฺฉา สริตฺวา นจฺฉาเทติ.\csromanspacer{} ยา อาปตฺติโย ปุพฺเพ อสฺสริตฺวา นจฺฉาเทสิ, ตา อาปตฺติโย ปจฺฉา สริตฺวา ฉาเทติ.\csromanspacer{} ตสฺส ภิกฺขเว ภิกฺขุโน ปุริมสฺมิํ จ ปจฺฉิมสฺมิํ จ อาปตฺติกฺขนฺเธ ยถาปฏิจฺเฉนฺเน ปริวาสํ ทตฺวา มานตฺตํ ทาตพฺพํ. (14)}

\prose{อิธ ปน ภิกฺขเว ภิกฺขุ สมฺพหุลา สํฆาทิเสสา อาปตฺติโย อาปชฺชติ, เอกจฺจา อาปตฺติโย สรติ, เอกจฺจา อาปตฺติโย นสฺสรติ.\csromanspacer{} ยา อาปตฺติโย สรติ, ตา อาปตฺติโย ฉาเทติ.\csromanspacer{} ยา อาปตฺติโย นสฺสรติ, ตา อาปตฺติโย นจฺฉาเทติ.\csromanspacer{} โส วิพฺภมิตฺวา ปุน อุปสมฺปนฺโน ยา อาปตฺติโย ปุพฺเพ สริตฺวา ฉาเทสิ, ตา อาปตฺติโย ปจฺฉา สริตฺวา ฉาเทติ.\csromanspacer{} ยา อาปตฺติโย ปุพฺเพ อสฺสริตฺวา นจฺฉาเทสิ, ตา อาปตฺติโย ปจฺฉา สริตฺวา นจฺฉาเทติ.\csromanspacer{} ตสฺส ภิกฺขเว ภิกฺขุโน ปุริมสฺมิํ จ ปจฺฉิมสฺมิํ จ อาปตฺติกฺขนฺเธ ยถาปฏิจฺฉนฺเน ปริวาสํ ทตฺวา มานตฺตํ ทาตพฺพํ. (15)}

\prose{อิธ ปน ภิกฺขเว ภิกฺขุ สมฺพหุลา สํฆาทิเสสา อาปตฺติโย อาปชฺชติ, เอกจฺจา อาปตฺติโย สรติ, เอกจฺจา อาปตฺติโย นสฺสรติ.\csromanspacer{} ยา อาปตฺติโย สรติ, ตา อาปตฺติโย ฉาเทติ.\csromanspacer{} ยา อาปตฺติโย นสฺสรติ, ตา อาปตฺติโย นจฺฉาเทติ.\csromanspacer{} โส วิพฺภมิตฺวา ปุน อุปสมฺปนฺโน ยา อาปตฺติโย ปุพฺเพ สริตฺวา ฉาเทสิ, ตา อาปตฺติโย ปจฺฉา สริตฺวา ฉาเทติ.\csromanspacer{} ยา อาปติโย ปุพฺเพ อสฺสริตฺวา นจฺฉาเทสิ, ตา อาปตฺติโย ปจฺฉา สริตฺวา ฉาเทติ.\csromanspacer{} ตสฺส ภิกฺขเว ภิกฺขุโน ปุริมสฺมิํ จ ปจฺฉิมสฺมิํ จ อาปตฺติกฺขนฺเธ ยถาปฏิจฺฉนฺเน ปริวาสํ ทตฺวา มานตฺตํ ทาตพฺพํ. (16)}

\proseitem{170}{อิธ ปน ภิกฺขเว ภิกฺขุ สมฺพหุลา สํฆาทิเสสา อาปตฺติโย อาปชฺชติ, เอกจฺจาสุ อาปตฺตีสุ นิพฺเพมติโก, เอกจฺจาสุ อาปตฺตีสุ เวมติโก.\csromanspacer{} ยาสุ อาปตฺตีสุ นิพฺเพมติโก, ตา อาปตฺติโย ฉาเทติ.\csromanspacer{} ยาสุ อาปตฺตีสุ เวมติโก, ตา อาปตฺติโย นจฺฉาเทติ.\csromanspacer{} โส วิพฺภมิตฺวา ปุน อุปสมฺปนฺโน ยา อาปตฺติโย ปุพฺเพ นิพฺเพมติโก ฉาเทสิ, ตา อาปตฺติโย ปจฺฉา นิพฺเพมติโก นจฺฉาเทติ.\csromanspacer{} ยา \csromanfolio{163}อาปตฺติโย ปุพฺเพ เวมติโก นจฺฉาเทสิ, ตา อาปตฺติโย ปจฺฉา นิพฺเพมติโก นจฺฉาเทติ.\csromanspacer{} ตสฺส ภิกฺขเว ภิกฺขุโน ปุริมสฺมิํ อาปตฺติกฺขนฺเธ ยถาปฏิจฺฉนฺเน ปริวาสํ ทตฺวา มานตฺตํ ทาตพฺพํ. (17)}

\prose{อิธ ปน ภิกฺขเว ภิกฺขุ สมฺพหุลา สํฆาทิเสสา อาปตฺติโย อาปชฺชติ, เอกจฺจาสุ อาปตฺตีสุ นิพฺเพมติโก, เอกจฺจาสุ อาปตฺตีสุ เวมติโก.\csromanspacer{} ยาสุ อาปตฺตีสุ นิพฺเพมติโก, ตา อาปตฺติโย ฉาเทติ.\csromanspacer{} ยาสุ อาปตฺตีสุ เวมติโก, ตา อาปตฺติโย นจฺฉาเทติ.\csromanspacer{} โส วิพฺภมิตฺวา ปุน อุปสมฺปนฺโน ยา อาปตฺติโย ปุพฺเพ นิพฺเพมติโก ฉาเทสิ, ตา อาปตฺติโย ปจฺฉา นิพฺเพมติโก นจฺฉาเทติ.\csromanspacer{} ยา อาปตฺติโย ปุพฺเพ เวมติโก นจฺฉาเทสิ, ตา อาปตฺติโย ปจฺฉา นิพฺเพมติโก ฉาเทติ.\csromanspacer{} ตสฺส ภิกฺขเว ภิกฺขุโน ปุริมสฺมิํ จ ปจฺฉิมสฺมิํ จ อาปตฺติกฺขนฺเธ ยถาปฏิจฺฉนฺเน ปริวาสํ ทตฺวา มานตฺตํ ทาตพฺพํ. (18)}

\prose{อิธ ปน ภิกฺขเว ภิกฺขุ สมฺพหุลา สํฆาทิเสสา อาปตฺติโย อาปชฺชติ, เอกจฺจาสุ อาปตฺตีสุ นิพฺเพมติโก, เอกจฺจาสุ อาปตฺตีสุ เวมติโก.\csromanspacer{} ยาสุ อาปตฺตีสุ นิพฺเพมติโก, ตา อาปตฺติโย ฉาเทติ.\csromanspacer{} ยาสุ อาปตฺตีสุ เวมติโก, ตา อาปตฺติโย นจฺฉาเทติ.\csromanspacer{} โส วิพฺภมิตฺวา ปุน อุปสมฺปนฺโน ยา อาปตฺติโย ปุพฺเพ นิพฺเพมติโก ฉาเทสิ, ตา อาปตฺติโย ปจฺฉา นิพฺเพมติโก ฉาเทติ.\csromanspacer{} ยา อาปตฺติโย ปุพฺเพ เวมติโก นจฺฉาเทสิ, ตา อาปตฺติโย ปจฺฉา นิพฺเพมติโก นจฺฉาเทติ.\csromanspacer{} ตสฺส ภิกฺขเว ภิกฺขุโน ปุริมสฺมิํ จ ปจฺฉิมสฺมิํ จ อาปตฺติกฺขนฺเธ ยถาปฏิจฺฉนฺเน ปริวาสํ ทตฺวา มานตฺตํ ทาตพฺพํ. (19)}

\prose{อิธ ปน ภิกฺขเว ภิกฺขุ สมฺพหุลา สํฆาทิเสสา อาปตฺติโย อาปชฺชติ, เอกจฺจาสุ อาปตฺตีสุ นิพฺเพมติโก, เอกจฺจาสุ อาปตฺตีสุ เวมติโก.\csromanspacer{} ยาสุ อาปตฺตีสุ นิพฺเพมติโก, ตา อาปตฺติโย ฉาเทติ.\csromanspacer{} ยาสุ อาปตฺตีสุ เวมติโก, ตา อาปตฺติโย นจฺฉาเทติ.\csromanspacer{} โส วิพฺภมิตฺวา ปุน อุปสมฺปนฺโน ยา อาปตฺติโย ปุพฺเพ นิพฺเพมติโก ฉาเทสิ, ตา อาปตฺติโย ปจฺฉา นิพฺเพมติโก ฉาเทติ.\csromanspacer{} ยา อาปตฺติโย ปุพฺเพ เวมติโก นจฺฉาเทสิ, ตา อาปตฺติโย ปจฺฉา นิพฺเพมติโก ฉาเทติ.\csromanspacer{} ตสฺส ภิกฺขเว ภิกฺขุโน ปุริมสฺมิํ จ ปจฺฉิมสฺมิํ จ อาปตฺติกฺขนฺเธ ยถาปฏิจฺฉนฺเน ปริวาสํ ทตฺวา มานตฺตํ ทาตพฺพํ. (20)}

\proseitem{171}{\csromanfolio{164}อิธ ปน ภิกฺขเว ภิกฺขุ สมฺพหุลา สํฆาทิเสสา อาปตฺติโย อาปชฺชิตฺวา อปฺปฏิจฺฉาเทตฺวา สามเณโร โหติ ฯเปฯ อุมฺมตฺตโก โหติ ฯเปฯ ขิตฺตจิตฺโต โหติ ฯเปฯ (ยถา เหฏฺฐา, ตถา วิตฺถาเรตพฺพํ.) เวทนาฏฺโฏ โหติ ฯเปฯ ตสฺส โหนฺติ อาปตฺติโย ปฏิจฺฉนฺนาโยปิ อปฺปฏิจฺฉนฺนาโยปิ ฯเปฯ เอกจฺจา อาปตฺติโย ชานาติ, เอกจฺจา อาปตฺติโย น ชานาติ ฯเปฯ เอกจฺจา อาปตฺติโย สรติ, เอกจฺจา อาปตฺติโย นสฺสรติ ฯเปฯ เอกจฺจาสุ อาปตฺตีสุ นิพฺเพมติโก, เอกจฺจาสุ อาปตฺตีสุ เวมติโก.\csromanspacer{} ยาสุ อาปตฺตีสุ นิพฺเพมติโก, ตา อาปตฺติโย ฉาเทติ.\csromanspacer{} ยาสุ อาปตฺตีสุ เวมติโก, ตา อาปตฺติโย นจฺฉาเทติ.\csromanspacer{} โส เวทนาฏฺโฏ โหติ.\csromanspacer{} โส ปุน อเวทนาฏฺโฏ หุตฺวา ยา อาปตฺติโย ปุพฺเพ นิพฺเพมติโก ฉาเทสิ, ตา อาปตฺติโย ปจฺฉา นิพฺเพมติโก นจฺฉาเทติ.\csromanspacer{} ยา อาปตฺติโย ปุพฺเพ เวมติโก นจฺฉาเทสิ, ตา อาปตฺติโย ปจฺฉา นิพฺเพมติโก นจฺฉาเทติ ฯเปฯ.\csromanspacer{} ยา อาปตฺติโย ปุพฺเพ นิพฺเพมติโก ฉาเทสิ, ตา อาปตฺติโย ปจฺฉา นิพฺเพมติโก นจฺฉาเทติ.\csromanspacer{} ยา อาปตฺติโย ปุพฺเพ เวมติโก นจฺฉาเทสิ, ตา อาปตฺติโย ปจฺฉา นิพฺเพมติโก ฉาเทติ ฯเปฯ.\csromanspacer{} ยา อาปตฺติโย ปุพฺเพ นิพฺเพมติโก ฉาเทสิ, ตา อาปตฺติโย ปจฺฉา นิพฺเพมติโก ฉาเทติ.\csromanspacer{} ยา อาปตฺติโย ปุพฺเพ เวมติโก นจฺฉาเทสิ, ตา อาปตฺติโย ปจฺฉา นิพฺเพมติโก นจฺฉาเทติ ฯเปฯ.\csromanspacer{} ยา อาปตฺติโย ปุพฺเพ นิพฺเพมติโก ฉาเทสิ, ตา อาปตฺติโย ปจฺฉา นิพฺเพมติโก ฉาเทติ.\csromanspacer{} ยา อาปตฺติโย ปุพฺเพ เวมติโก นจฺฉาเทสิ, ตา อาปตฺติโย ปจฺฉา นิพฺเพมติโก ฉาเทติ.\csromanspacer{} ตสฺส ภิกฺขเว ภิกฺขุโน ปุริมสฺมิํ จ ปจฺฉิมสฺมิํ จ อาปตฺติกฺขนฺเธ ยถาปฏิจฺฉนฺเน ปริวาสํ ทตฺวา มานตฺตํ ทาตพฺพํ. (100)}

\nitthitam{มานตฺตสตํ นิฏฺฐิตํ.}
\csromanmatikaanchor{mtk.90}
\csromantocmarkref{section}{6. สมูลายสโมธานปริวาสจตุสฺสต}{mtk.90}
\bookmark[dest={mtk.90},level=1]{6. สมูลายสโมธานปริวาสจตุสฺสต}
\csromanheader{1}{6. สมูลายสโมธานปริวาสจตุสฺสต}
\proseitem{172}{อิธ ปน ภิกฺขเว ภิกฺขุ ปริวสนฺโต อนฺตรา สมฺพหุลา สํฆาทิเสสา อาปตฺติโย อาปชฺชิตฺวา อปฺปฏิจฺฉาเทตฺวา วิพฺภมติ, \csromanfolio{165}โส ปุน อุปสมฺปนฺโน ตา อาปตฺติโย นจฺฉาเทติ, โส ภิกฺขุ มูลาย ปฏิกสฺสิตพฺโพ. (1)}

\prose{อิธ ปน ภิกฺขเว ภิกฺขุ ปริวสนฺโต อนฺตรา สมฺพหุลา สํฆาทิเสสา อาปตฺติโย อาปชฺชิตฺวา อปฺปฏิจฺฉาเทตฺวา วิพฺภมติ, โส ปุน อุปสมฺปนฺโน ตา อาปตฺติโย ฉาเทติ, โส ภิกฺขุ มูลาย ปฏิกสฺสิตพฺโพ, ยถาปฏิจฺฉนฺนานญฺจสฺส อาปตฺตีนํ ปุริมาย อาปตฺติยา สโมธานปริวาโส ทาตพฺโพ. (2)}

\prose{อิธ ปน ภิกฺขเว ภิกฺขุ ปริวสนฺโต อนฺตรา สมฺพหุลา สํฆาทิเสสา อาปตฺติโย อาปชฺชิตฺวา ปฏิจฺฉาเทตฺวา วิพฺภมติ, โส ปุน อุปสมฺปนฺโน ตา อาปตฺติโย นจฺฉาเทติ, โส ภิกฺขุ มูลาย ปฏิกสฺสิตพฺโพ, ยถาปฏิจฺฉนฺนานญฺจสฺส อาปตฺตีนํ ปุริมาย อาปตฺติยา สโมธานปริวาโส ทาตพฺโพ. (3)}

\prose{อิธ ปน ภิกฺขเว ภิกฺขุ ปริวสนฺโต อนฺตรา สมฺพหุลา สํฆาทิเสสา อาปตฺติโย อาปชฺชิตฺวา ปฏิจฺฉาเทตฺวา วิพฺภมติ, โส ปุน อุปสมฺปนฺโน ตา อาปตฺติโย ฉาเทติ, โส ภิกฺขุ มูลาย ปฏิกสฺสิตพฺโพ, ยถาปฏิจฺฉนฺนานญฺจสฺส อาปตฺตีนํ ปุริมาย อาปตฺติยา สโมธานปริวาโส ทาตพฺโพ. (4)}

\proseitem{173}{อิธ ปน ภิกฺขเว ภิกฺขุ ปริวสนฺโต อนฺตรา สมฺพหุลา สํฆาทิเสสา อาปตฺติโย อาปชฺชติ, ตสฺส โหนฺติ อาปตฺติโย ปฏิจฺฉนฺนาโยปิ อปฺปฏิจฺฉนฺนาโยปิ.\csromanspacer{} โส วิพฺภมิตฺวา ปุน อุปสมฺปนฺโน ยา อาปตฺติโย ปุพฺเพ ฉาเทสิ, ตา อาปตฺติโย ปจฺฉา นจฺฉาเทติ.\csromanspacer{} ยา อาปตฺติโย ปุพฺเพ นจฺฉาเทสิ, ตา อาปตฺติโย ปจฺฉา นจฺฉาเทติ.\csromanspacer{} โส ภิกฺขุ มูลาย ปฏิกสฺสิตพฺโพ, ยถาปฏิจฺฉนฺนานญฺจสฺส อาปตฺตีนํ ปุริมาย อาปตฺติยา สโมธานปริวาโส ทาตพฺโพ. (5)}

\prose{อิธ ปน ภิกฺขเว ภิกฺขุ ปริวสนฺโต อนฺตรา สมฺพหุลา สํฆาทิเสสา อาปตฺติโย อาปชฺชติ, ตสฺส โหนฺติ อาปตฺติโย ปฏิจฺฉนฺนาโยปิ อปฺปฏิจฺฉนฺนาโยปิ.\csromanspacer{} โส วิพฺภมิตฺวา ปุน อุปสมฺปนฺโน ยา อาปตฺติโย ปุพฺเพ ฉาเทสิ, ตา อาปตฺติโย ปจฺฉา นจฺฉาเทติ. \csromanfolio{166}ยา อาปตฺติโย ปุพฺเพ นจฺฉาเทสิ, ตา อาปตฺติโย ปจฺฉา ฉาเทติ.\csromanspacer{} โส ภิกฺขุ มูลาย ปฏิกสฺสิตพฺโพ, ยถาปฏิจฺฉนฺนานญฺจสฺส อาปตฺตีนํ ปุริมาย อาปตฺติยา สโมธานปริวาโส ทาตพฺโพ. (6)}

\prose{อิธ ปน ภิกฺขเว ภิกฺขุ ปริวสนฺโต อนฺตรา สมฺพหุลา สํฆาทิเสสา อาปตฺติโย อาปชฺชติ, ตสฺส โหนฺติ อาปตฺติโย ปฏิจฺฉนฺนาโยปิ อปฺปฏิจฺฉนฺนาโยปิ.\csromanspacer{} โส วิพฺภมิตฺวา ปุน อุปสมฺปนฺโน ยา อาปตฺติโย ปุพฺเพ ฉาเทสิ, ตา อาปตฺติโย ปจฺฉา ฉาเทติ.\csromanspacer{} ยา อาปตฺติโย ปุพฺเพ นจฺฉาเทสิ, ตา อาปตฺติโย ปจฺฉา นจฺฉาเทติ.\csromanspacer{} โส ภิกฺขุ มูลาย ปฏิกสฺสิตพฺโพ, ยถาปฏิจฺฉนฺนานญฺจสฺส อาปตฺตีนํ ปุริมาย อาปตฺติยา สโมธานปริวาโส ทาตพฺโพ. (7)}

\prose{อิธ ปน ภิกฺขเว ภิกฺขุ ปริวสนฺโต อนฺตรา สมฺพหุลา สํฆาทิเสสา อาปตฺติโย อาปชฺชติ, ตสฺส โหนฺติ อาปตฺติโย ปฏิจฺฉนฺนาโยปิ อปฺปฏิจฺฉนฺนาโยปิ.\csromanspacer{} โส วิพฺภมิตฺวา ปุน อุปสมฺปนฺโน ยา อาปตฺติโย ปุพฺเพ ฉาเทสิ, ตา อาปตฺติโย ปจฺฉา ฉาเทติ.\csromanspacer{} ยา อาปตฺติโย ปุพฺเพ นจฺฉาเทสิ, ตา อาปตฺติโย ปจฺฉา ฉาเทติ.\csromanspacer{} โส ภิกฺขุ มูลาย ปฏิกสฺสิตพฺโพ, ยถาปฏิจฺฉนฺนานญฺจสฺส อาปตฺตีนํ ปุริมาย อาปตฺติยา สโมธานปริวาโส ทาตพฺโพ. (8)}

\proseitem{174}{อิธ ปน ภิกฺขเว ภิกฺขุ ปริวสนฺโต อนฺตรา สมฺพหุลา สํฆาทิเสสา อาปตฺติโย อาปชฺชติ, เอกจฺจา อาปตฺติโย ชานาติ, เอกจฺจา อาปตฺติโย น ชานาติ.\csromanspacer{} ยา อาปตฺติโย ชานาติ, ตา อาปตฺติโย ฉาเทติ.\csromanspacer{} ยา อาปตฺติโย น ชานาติ, ตา อาปตฺติโย นจฺฉาเทติ.\csromanspacer{} โส วิพฺภมิตฺวา ปุน อุปสมฺปนฺโน ยา อาปตฺติโย ปุพฺเพ ชานิตฺวา ฉาเทสิ, ตา อาปตฺติโย ปจฺฉา ชานิตฺวา นจฺฉาเทติ.\csromanspacer{} ยา อาปตฺติโย ปุพฺเพ อชานิตฺวา นจฺฉาเทสิ, ตา อาปตฺติโย ปจฺฉา ชานิตฺวา นจฺฉาเทติ.\csromanspacer{} โส ภิกฺขุ มูลาย ปฏิกสฺสิตพฺโพ, ยถาปฏิจฺฉนฺนานญฺจสฺส อาปตฺตีนํ ปุริมาย อาปตฺติยา สโมธานปริวาโส ทาตพฺโพ. (9)}

\prose{อิธ ปน ภิกฺขเว ภิกฺขุ ปริวสนฺโต อนฺตรา สมฺพหุลา สํฆาทิเสสา อาปตฺติโย อาปชฺชติ, เอกจฺจา อาปตฺติโย ชานาติ, เอกจฺจา \csromanfolio{167}อาปตฺติโย น ชานาติ.\csromanspacer{} ยา อาปตฺติโย ชานาติ, ตา อาปตฺติโย ฉาเทติ.\csromanspacer{} ยา อาปตฺติโย น ชานาติ, ตา อาปตฺติโย นจฺฉาเทติ.\csromanspacer{} โส วิพฺภมิตฺวา ปุน อุปสมฺปนฺโน ยา อาปตฺติโย ปุพฺเพ ชานิตฺวา ฉาเทสิ, ตา อาปตฺติโย ปจฺฉา ชานิตฺวา นจฺฉาเทติ.\csromanspacer{} ยา อาปตฺติโย ปุพฺเพ อชานิตฺวา นจฺฉาเทสิ, ตา อาปตฺติโย ปจฺฉา ชานิตฺวา ฉาเทติ.\csromanspacer{} โส ภิกฺขุ มูลาย ปฏิกสฺสิตพฺโพ, ยถาปฏิจฺฉนฺนานญฺจสฺส อาปตฺตีนํ ปุริมาย อาปตฺติยา สโมธานปริวาโส ทาตพฺโพ. (10)}

\prose{อิธ ปน ภิกฺขุ ปริวสนฺโต อนฺตรา สมฺพหุลา สํฆาทิเสสา อาปตฺติโย อาปชฺชติ, เอกจฺจา อาปตฺติโย ชานาติ, เอกจฺจา อาปตฺติโย น ชานาติ.\csromanspacer{} ยา อาปตฺติโย ชานาติ, ตา อาปตฺติโย ฉาเทติ.\csromanspacer{} ยา อาปตฺติโย น ชานาติ, ตา อาปตฺติโย นจฺฉาเทติ.\csromanspacer{} โส วิพฺภมิตฺวา ปุน อุปสมฺปนฺโน ยา อาปตฺติโย ปุพฺเพ ชานิตฺวา ฉาเทสิ, ตา อาปตฺติโย ปจฺฉา ชานิตฺวา ฉาเทติ.\csromanspacer{} ยา อาปตฺติโย ปุพฺเพ อชานิตฺวา นจฺฉาเทสิ, ตา อาปตฺติโย ปจฺฉา ชานิตฺวา นจฺฉาเทติ.\csromanspacer{} โส ภิกฺขุ มูลาย ปฏิกสฺสิตพฺโพ, ยถาปฏิจฺฉนฺนานญฺจสฺส อาปตฺตีนํ ปุริมาย อาปตฺติยา สโมธานปริวาโส ทาตพฺโพ. (11)}

\prose{อิธ ปน ภิกฺขเว ภิกฺขุ ปริวสนฺโต อนฺตรา สมฺพหุลา สํฆาทิเสสา อาปตฺติโย อาปชฺชติ, เอกจฺจา อาปตฺติโย ชานาติ, เอกจฺจา อาปตฺติโย น ชานาติ.\csromanspacer{} ยา อาปตฺติโย ชานาติ, ตา อาปตฺติโย ฉาเทติ.\csromanspacer{} ยา อาปตฺติโย น ชานาติ, ตา อาปตฺติโย นจฺฉาเทติ.\csromanspacer{} โส วิพฺภมิตฺวา ปุน อุปสมฺปนฺโน ยา อาปตฺติโย ปุพฺเพ ชานิตฺวา ฉาเทสิ, ตา อาปตฺติโย ปจฺฉา ชานิตฺวา ฉาเทติ.\csromanspacer{} ยา อาปตฺติโย ปุพฺเพ อชานิตฺวา นจฺฉาเทสิ, ตา อาปตฺติโย ปจฺฉา ชานิตฺวา ฉาเทติ.\csromanspacer{} โส ภิกฺขุ มูลาย ปฏิกสฺสิตพฺโพ, ยถาปฏิจฺฉนฺนานญฺจสฺส อาปตฺตีนํ ปุริมาย อาปตฺติยา สโมธานปริวาโส ทาตพฺโพ. (12)}

\proseitem{175}{อิธ ปน ภิกฺขเว ภิกฺขุ ปริวสนฺโต อนฺตรา สมฺพหุลา สํฆาทิเสสา อาปตฺติโย อาปชฺชติ, เอกจฺจา อาปตฺติโย สรติ เอกจฺจา อาปตฺติโย นสฺสรติ.\csromanspacer{} ยา อาปตฺติโย สรติ, ตา \csromanfolio{168}อาปตฺติโย ฉาเทติ.\csromanspacer{} ยา อาปตฺติโย นสฺสรติ, ตา อาปตฺติโย นจฺฉาเทติ.\csromanspacer{} โส วิพฺภมิตฺวา ปุน อุปสมฺปนฺโน ยา อาปตฺติโย ปุพฺเพ สริตฺวา ฉาเทสิ, ตา อาปตฺติโย ปจฺฉา สริตฺวา นจฺฉาเทติ.\csromanspacer{} ยา อาปตฺติโย ปุพฺเพ อสฺสริตฺวา นจฺฉาเทสิ, ตา อาปตฺติโย ปจฺฉา สริตฺวา นจฺฉาเทติ.\csromanspacer{} โส ภิกฺขุ มูลาย ปฏิกสฺสิตพฺโพ, ยถา ปฏิจฺฉนฺนานญฺจสฺส อาปตฺตีนํ ปุริมาย อาปตฺติยา สโมธานปริวาโส ทาตพฺโพ. (13)}

\prose{อิธ ปน ภิกฺขเว ภิกฺขุ ปริวสนฺโต อนฺตรา สมฺพหุลา สํฆาทิเสสา อาปตฺติโย อาปชฺชติ, เอกจฺจา อาปตฺติโย สรติ, เอกจฺจา อาปตฺติโย นสฺสรติ.\csromanspacer{} ยา อาปตฺติโย สรติ, ตา อาปตฺติโย ฉาเทติ.\csromanspacer{} ยา อาปตฺติโย นสฺสรติ, ตา อาปตฺติโย นจฺฉาเทติ.\csromanspacer{} โส วิพฺภมิตฺวา ปุน อุปสมฺปนฺโน ยา อาปตฺติโย ปุพฺเพ สริตฺวา ฉาเทสิ, ตา อาปตฺติโย ปจฺฉา สริตฺวา นจฺฉาเทติ.\csromanspacer{} ยา อาปตฺติโย ปุพฺเพ อสฺสริตฺวา นจฺฉาเทสิ, ตา อาปตฺติโย ปจฺฉา สริตฺวา ฉาเทติ.\csromanspacer{} โส ภิกฺขุ มูลาย ปฏิกสฺสิตพฺโพ, ยถาปฏิจฺฉนฺนานญฺจสฺส อาปตฺตีนํ ปุริมาย อาปตฺติยา สโมธานปริวาโส ทาตพฺโพ. (14)}

\prose{อิธ ปน ภิกฺขเว ภิกฺขุ ปริวสนฺโต อนฺตรา สมฺพหุลา สํฆาทิเสสา อาปตฺติโย อาปชฺชติ, เอกจฺจา อาปตฺติโย สรติ, เอกจฺจา อาปตฺติโย นสฺสรติ.\csromanspacer{} ยา อาปตฺติโย สรติ, ตา อาปตฺติโย ฉาเทติ.\csromanspacer{} ยา อาปตฺติโย นสฺสรติ, ตา อาปตฺติโย นจฺฉาเทติ.\csromanspacer{} โส วิพฺภมิตฺวา ปุน อุปสมฺปนฺโน ยา อาปตฺติโย ปุพฺเพ สริตฺวา ฉาเทสิ, ตา อาปตฺติโย ปจฺฉา สริตฺวา ฉาเทติ.\csromanspacer{} ยา อาปตฺติโย ปุพฺเพ อสฺสริตฺวา นจฺฉาเทสิ, ตา อาปตฺติโย ปจฺฉา สริตฺวา นจฺฉาเทติ.\csromanspacer{} โส ภิกฺขุ มูลาย ปฏิกสฺสิตพฺโพ, ยถาปฏิจฺฉนฺนานญฺจสฺส อาปตฺตีนํ ปุริมาย อาปตฺติยา สโมธานปริวาโส ทาตพฺโพ. (15)}

\prose{อิธ ปน ภิกฺขเว ภิกฺขุ ปริวสนฺโต อนฺตรา สมฺพหุลา สํฆาทิเสสา อาปตฺติโย อาปชฺชติ, เอกจฺจา อาปตฺติโย สรติ, เอกจฺจา อาปตฺติโย นสฺสรติ.\csromanspacer{} ยา อาปตฺติโย สรติ, ตา อาปตฺติโย ฉาเทติ.\csromanspacer{} ยา อาปตฺติโย นสฺสรติ, ตา อาปตฺติโย นจฺฉาเทติ.\csromanspacer{} โส วิพฺภมิตฺวา ปุน อุปสมฺปนฺโน ยา อาปตฺติโย ปุพฺเพ สริตฺวา \csromanfolio{169}ฉาเทสิ, ตา อาปตฺติโย ปจฺฉา สริตฺวา ฉาเทติ.\csromanspacer{} ยา อาปตฺติโย ปุพฺเพ อสฺสริตฺวา นจฺฉาเทสิ, ตา อาปตฺติโย ปจฺฉา สริตฺวา ฉาเทติ.\csromanspacer{} โส ภิกฺขุ มูลาย ปฏิกสฺสิตพฺโพ, ยถาปฏิจฺฉนฺนานญฺจสฺส อาปตฺตีนํ ปุริมาย อาปตฺติยา สโมธานปริวาโส ทาตพฺโพ. (16)}

\proseitem{176}{อิธ ปน ภิกฺขเว ภิกฺขุ ปริวสนฺโต อนฺตรา สมฺพหุลา สํฆาทิเสสา อาปตฺติโย อาปชฺชติ, เอกจฺจาสุ อาปตฺตีสุ นิพฺเพมติโก, เอกจฺจาสุ อาปตฺตีสุ เวมติโก.\csromanspacer{} ยาสุ อาปตฺตีสุ นิพฺเพมติโก, ตา อาปตฺติโย ฉาเทติ.\csromanspacer{} ยาสุ อาปตฺตีสุ เวมติโก, ตา อาปตฺติโย นจฺฉาเทติ.\csromanspacer{} โส วิพฺภมิตฺวา ปุน อุปสมฺปนฺโน ยา อาปตฺติโย ปุพฺเพ นิพฺเพมติโก ฉาเทสิ, ตา อาปตฺติโย ปจฺฉา นิพฺเพมติโก นจฺฉาเทติ.\csromanspacer{} ยา อาปตฺติโย ปุพฺเพ เวมติโก นจฺฉาเทสิ, ตา อาปตฺติโย ปจฺฉา นิพฺเพมติโก นจฺฉาเทติ.\csromanspacer{} โส ภิกฺขุ มูลาย ปฏิกสฺสิตพฺโพ, ยถาปฏิจฺฉนฺนานญฺจสฺส อาปตฺตีนํ ปุริมาย อาปตฺติยา สโมธานปริวาโส ทาตพฺโพ. (17)}

\prose{อิธ ปน ภิกฺขเว ภิกฺขุ ปริวสนฺโต อนฺตรา สมฺพหุลา สํฆาทิเสสา อาปตฺติโย อาปชฺชติ, เอกจฺจาสุ อาปตฺตีสุ นิพฺเพมติโก, เอกจฺจาสุ อาปตฺตีสุ เวมติโก.\csromanspacer{} ยาสุ อาปตฺตีสุ นิพฺเพมติโก, ตา อาปตฺติโย ฉาเทติ.\csromanspacer{} ยาสุ อาปตฺตีสุ เวมติโก, ตา อาปตฺติโย นจฺฉาเทติ.\csromanspacer{} โส วิพฺภมิตฺวา ปุน อุปสมฺปนฺโน ยา อาปตฺติโย ปุพฺเพ นิพฺเพมติโก ฉาเทสิ, ตา อาปตฺติโย ปจฺฉา นิพฺเพมติโก นจฺฉาเทติ.\csromanspacer{} ยา อาปตฺติโย ปุพฺเพ เวมติโก นจฺฉาเทสิ, ตา อาปตฺติโย ปจฺฉา นิพฺเพมติโก ฉาเทติ.\csromanspacer{} โส ภิกฺขุ มูลาย ปฏิกสฺสิตพฺโพ, ยถาปฏิจฺฉนฺนานญฺจสฺส อาปตฺตีนํ ปุริมาย อาปตฺติยา สโมธานปริวาโส ทาตพฺโพ. (18)}

\prose{อิธ ปน ภิกฺขเว ภิกฺขุ ปริวสนฺโต อนฺตรา สมฺพหุลา สํฆาทิเสสา อาปตฺติโย อาปชฺชติ, เอกจฺจาสุ อาปตฺตีสุ นิพฺเพมติโก, เอกจฺจาสุ อาปตฺตีสุ เวมติโก.\csromanspacer{} ยาสุ อาปตฺตีสุ นิพฺเพมติโก, ตา อาปตฺติโย ฉาเทติ.\csromanspacer{} ยาสุ อาปตฺตีสุ เวมติโก, ตา อาปตฺติโย นจฺฉาเทติ.\csromanspacer{} โส วิพฺภมิตฺวา ปุน อุปสมฺปนฺโน ยา \csromanfolio{170}อาปตฺติโย ปุพฺเพ นิพฺเพมติโก ฉาเทสิ, ตา อาปตฺติโย ปจฺฉา นิพฺเพมติโก ฉาเทติ.\csromanspacer{} ยา อาปตฺติโย ปุพฺเพ เวมติโก นจฺฉาเทสิ, ตา อาปตฺติโย ปจฺฉา นิพฺเพมติโก นจฺฉาเทติ.\csromanspacer{} โส ภิกฺขุ มูลาย ปฏิกสฺสิตพฺโพ, ยถาปฏิจฺฉนฺนานญฺจสฺส อาปตฺตีนํ ปุริมาย อาปตฺติยา สโมธานปริวาโส ทาตพฺโพ. (19)}

\prose{อิธ ปน ภิกฺขเว ภิกฺขุ ปริวสนฺโต อนฺตรา สมฺพหุลา สํฆาทิเสสา อาปตฺติโย อาปชฺชติ, เอกจฺจาสุ อาปตฺตีสุ นิพฺเพมติโก, เอกจฺจาสุ อาปตฺตีสุ เวมติโก.\csromanspacer{} ยาสุ อาปตฺตีสุ นิพฺเพมติโก, ตา อาปตฺติโย ฉาเทติ.\csromanspacer{} ยาสุ อาปตฺตีสุ เวมติโก, ตา อาปตฺติโย นจฺฉาเทติ.\csromanspacer{} โส วิพฺภมิตฺวา\footnote{โส ภิกฺขุ วิพฺภมิตฺวา (ก)} ปุน อุปสมฺปนฺโน ยา อาปตฺติโย ปุพฺเพ นิพฺเพมติโก ฉาเทสิ, ตา อาปตฺติโย ปจฺฉา นิพฺเพมติโก ฉาเทติ.\csromanspacer{} ยา อาปตฺติโย ปุพฺเพ เวมติโก นจฺฉาเทสิ, ตา อาปตฺติโย ปจฺฉา นิพฺเพมติโก ฉาเทติ.\csromanspacer{} โส ภิกฺขุ มูลาย ปฏิกสฺสิตพฺโพ, ยถาปฏิจฺฉนฺนานญฺจสฺส อาปตฺตีนํ ปุริมาย อาปตฺติยา สโมธานปริวาโส ทาตพฺโพ. (20)}

\proseitem{177}{อิธ ปน ภิกฺขเว ภิกฺขุ ปริวสนฺโต อนฺตรา สมฺพหุลา สํฆาทิเสสา อาปตฺติโย อาปชฺชิตฺวา อปฺปฏิจฺฉาเทตฺวา สามเณโร โหติ ฯเปฯ อุมฺมตฺตโก โหติ ฯเปฯ ขิตฺตจิตฺโต โหติ ฯเปฯ เวทนาฏฺโฏ โหติ ฯเปฯ ตสฺส โหนฺติ อาปตฺติโย ปฏิจฺฉนฺนาโยปิ อปฺปฏิจฺฉนฺนาโยปิ. (ยถา เหฏฺฐา วิตฺถาริตํ, ตถา วิตฺถาเรตพฺพํ) ฯเปฯ เอกจฺจา อาปตฺติโย ชานาติ, เอกจฺจา อาปตฺติโย น ชานาติ ฯเปฯ เอกจฺจา อาปตฺติโย สรติ, เอกจฺจา อาปตฺติโย นสฺสรติ ฯเปฯ เอกจฺจาสุ อาปตฺตีสุ นิพฺเพมติโก, เอกจฺจาสุ อาปตฺตีสุ เวมติโก.\csromanspacer{} ยาสุ อาปตฺตีสุ นิพฺเพมติโก, ตา อาปตฺติโย ฉาเทติ.\csromanspacer{} ยาสุ อาปตฺตีสุ เวมติโก, ตา อาปตฺติโย นจฺฉาเทติ.\csromanspacer{} โส เวทนาฏฺโฏ โหติ, โส ปุน อเวทนาฏฺโฏ หุตฺวา ยา อาปตฺติโย ปุพฺเพ นิพฺเพมติโก ฉาเทสิ, ตา อาปตฺติโย ปจฺฉา นิพฺเพมติโก นจฺฉาเทติ.\csromanspacer{} ยา อาปตฺติโย ปุพฺเพ เวมติโก นจฺฉาเทสิ, ตา อาปตฺติโย ปจฺฉา นิพฺเพมติโก นจฺฉาเทติ ฯเปฯ.\csromanspacer{} ยา อาปตฺติโย ปุพฺเพ นิพฺเพมติโก ฉาเทสิ, ตา อาปตฺติโย ปจฺฉา นิพฺเพมติโก \csromanfolio{171}นจฺฉาเทติ.\csromanspacer{} ยา อาปตฺติโย ปุพฺเพ เวมติโก นจฺฉาเทสิ, ตา อาปตฺติโย ปจฺฉา นิพฺเพมติโก ฉาเทติ ฯเปฯ.\csromanspacer{} ยา อาปตฺติโย ปุพฺเพ นิพฺเพมติโก ฉาเทสิ, ตา อาปตฺติโย ปจฺฉา นิพฺเพมติโก ฉาเทติ.\csromanspacer{} ยา อาปตฺติโย ปุพฺเพ เวมติโก นจฺฉาเทสิ, ตา อาปตฺติโย ปจฺฉา นิพฺเพมติโก นจฺฉาเทติ ฯเปฯ.\csromanspacer{} ยา อาปตฺติโย ปุพฺเพ นิพฺเพมติโก ฉาเทสิ, ตา อาปตฺติโย ปจฺฉา นิพฺเพมติโก ฉาเทติ.\csromanspacer{} ยา อาปตฺติโย ปุพฺเพ เวมติโก นจฺฉาเทสิ, ตา อาปตฺติโย ปจฺฉา นิพฺเพมติโก ฉาเทติ.\csromanspacer{} โส ภิกฺขุ มูลาย ปฏิกสฺสิตพฺโพ, ยถาปฏิจฺฉนฺนานญฺจสฺส อาปตฺตีนํ ปุริมาย อาปตฺติยา สโมธานปริวาโส ทาตพฺโพ. (100)}

\proseitem{178}{อิธ ปน ภิกฺขเว ภิกฺขุ มานตฺตารโห ฯเปฯ มานตฺตํ จรนฺโต ฯเปฯ อพฺภานารโห อนฺตรา สมฺพหุลา สํฆาทิเสสา อาปตฺติโย อาปชฺชิตฺวา อปฺปฏิจฺฉาเทตฺวา วิพฺภมติ ฯเปฯ (มานตฺตารโห จ มานตฺตจารี จ อพฺภานารโห จ ยถา ปริวาโส วิตฺถาริโต, ตถา วิตฺถาเรตพฺโพ.) (320)}

\proseitem{179}{อิธ ปน ภิกฺขเว ภิกฺขุ อพฺภานารโห อนฺตรา สมฺพหุลา สํฆาทิเสสา อาปตฺติโย อาปชฺชิตฺวา อปฺปฏิจฺฉาเทตฺวา สามเณโร โหติ ฯเปฯ อุมฺมตฺตโก โหติ ฯเปฯ ขิตฺตจิตฺโต โหติ ฯเปฯ เวทนาฏฺโฏ โหติ ฯเปฯ ตสฺส โหนฺติ อาปตฺติโย ปฏิจฺฉนฺนาโยปิ อปฺปฏิจฺฉนฺนาโยปิ ฯเปฯ เอกจฺจา อาปตฺติโย ชานาติ, เอกจฺจา อาปตฺติโย น ชานาติ ฯเปฯ เอกจฺจา อาปตฺติโย สรติ, เอกจฺจา อาปตฺติโย นสฺสรติ ฯเปฯ เอกจฺจาสุ อาปตฺตีสุ นิพฺเพมติโก, เอกจฺจาสุ อาปตฺตีสุ เวมติโก.\csromanspacer{} ยาสุ อาปตฺตีสุ นิพฺเพมติโก, ตา อาปตฺติโย ฉาเทติ.\csromanspacer{} ยาสุ อาปตฺตีสุ เวมติโก, ตา อาปตฺติโย นจฺฉาเทติ.\csromanspacer{} โส เวทนาฏฺโฏ โหติ, โส ปุน อเวทนาฏฺโฏ หุตฺวา ยา อาปตฺติโย ปุพฺเพ นิพฺเพมติโก ฉาเทสิ, ตา อาปตฺติโย ปจฺฉา นิพฺเพมติโก นจฺฉาเทติ.\csromanspacer{} ยา อาปตฺติโย ปุพฺเพ เวมติโก นจฺฉาเทสิ, ตา อาปตฺติโย ปจฺฉา นิพฺเพมติโก นจฺฉาเทติ ฯเปฯ.\csromanspacer{} ยา อาปตฺติโย ปุพฺเพ นิพฺเพมติโก ฉาเทสิ, ตา อาปตฺติโย ปจฺฉา นิพฺเพมติโก นจฺฉาเทติ.\csromanspacer{} ยา อาปตฺติโย ปุพฺเพ เวมติโก \csromanfolio{172}นจฺฉาเทสิ, ตา อาปตฺติโย ปจฺฉา นิพฺเพมติโก ฉาเทติ ฯเปฯ.\csromanspacer{} ยา อาปตฺติโย ปุพฺเพ นิพฺเพมติโก ฉาเทสิ, ตา อาปตฺติโย ปจฺฉา นิพฺเพมติโก ฉาเทติ.\csromanspacer{} ยา อาปตฺติโย ปุพฺเพ เวมติโก นจฺฉาเทสิ, ตา อาปตฺติโย ปจฺฉา นิพฺเพมติโก นจฺฉาเทติ ฯเปฯ.\csromanspacer{} ยา อาปตฺติโย ปุพฺเพ นิพฺเพมติโก ฉาเทสิ, ตา อาปตฺติโย ปจฺฉา นิพฺเพมติโก ฉาเทติ.\csromanspacer{} ยา อาปตฺติโย ปุพฺเพ เวมติโก นจฺฉาเทสิ, ตา อาปตฺติโย ปจฺฉา นิพฺเพมติโก ฉาเทติ.\csromanspacer{} โส ภิกฺขุ มูลาย ปฏิกสฺสิตพฺโพ, ยถาปฏิจฺฉนฺนานญฺจสฺส อาปตฺตีนํ ปุริมาย อาปตฺติยา สโมธานปริวาโส ทาตพฺโพ. (400)}

\nitthitam{สมูลายสโมธานปริวาสจตุสฺสตํ นิฏฺฐิตํ.}
\csromanmatikaanchor{mtk.91}
\csromantocmarkref{section}{7. ปริมาณาทิวาร-อฏฺฐก}{mtk.91}
\bookmark[dest={mtk.91},level=1]{7. ปริมาณาทิวาร-อฏฺฐก}
\csromanheader{1}{7. ปริมาณาทิวาร-อฏฺฐก}
\proseitem{180}{อิธ ปน ภิกฺขเว ภิกฺขุ สมฺพหุลา สํฆาทิเสสา อาปตฺติโย อาปชฺชิตฺวา ปริมาณา อปฺปฏิจฺฉาเทตฺวา ฯเปฯ อปริมาณา อปฺปฏิจฺฉาเทตฺวา ฯเปฯ เอกนามา อปฺปฏิจฺฉาเทตฺวา ฯเปฯ นานานามา อปฺปฏิจฺฉาเทตฺวา ฯเปฯ สภาคา อปฺปฏิจฺฉาเทตฺวา ฯเปฯ วิสภาคา อปฺปฏิจฺฉาเทตฺวา ฯเปฯ ววตฺถิตา อปฺปฏิจฺฉาเทตฺวา ฯเปฯ สมฺภินฺนา อปฺปฏิจฺฉาเทตฺวา วิพฺภมติ ฯเปฯ. (ยถา เหฏฺฐา, ตถา วิตฺถาเรตพฺพํ.)}

\nitthitam{ปริมาณาทิวาร-อฏฺฐกํ นิฏฺฐิตํ.}
\csromanmatikaanchor{mtk.92}
\csromantocmarkref{section}{8. เทฺวภิกฺขุวาร-เอกาทสก}{mtk.92}
\bookmark[dest={mtk.92},level=1]{8. เทฺวภิกฺขุวาร-เอกาทสก}
\csromanheader{1}{8. เทฺวภิกฺขุวาร-เอกาทสก}
\proseitem{181}{เทฺว ภิกฺขู สํฆาทิเสสํ อาปนฺนา โหนฺติ, เต สํฆาทิเสเส สํฆาทิเสสทิฏฺฐิโน โหนฺติ, เอโก ฉาเทติ, เอโก นจฺฉาเทติ.\csromanspacer{} โย ฉาเทติ, โส ทุกฺกฏํ เทสาเปตพฺโพ, ยถาปฏิจฺฉนฺเน\footnote{ยถาปฏิจฺฉนฺนานํ (สี)} จสฺส ปริวาสํ ทตฺวา อุภินฺนมฺปิ มานตฺตํ ทาตพฺพํ. (1)}

\prose{\csromanfolio{173}เทฺว ภิกฺขู สํฆาทิเสสํ อาปนฺนา โหนฺติ, เต สํฆาทิเสเส เวมติกา โหนฺติ, เอโก ฉาเทติ, เอโก นจฺฉาเทติ.\csromanspacer{} โย ฉาเทติ, โส ทุกฺกฏํ เทสาเปตพฺโพ, ยถาปฏิจฺฉนฺเน จสฺส ปริวาสํ ทตฺวา อุภินฺนมฺปิ มานตฺตํ ทาตพฺพํ. (2)}

\prose{เทฺว ภิกฺขู สํฆาทิเสสํ อาปนฺนา โหนฺติ, เต สํฆาทิเสเส มิสฺสกทิฏฺฐิโน โหนฺติ, เอโก ฉาเทติ, เอโก นจฺฉาเทติ.\csromanspacer{} โย ฉาเทติ, โส ทุกฺกฏํ เทสาเปตพฺโพ, ยถาปฏิจฺฉนฺเน จสฺส ปริวาสํ ทตฺวา อุภินฺนมฺปิ มานตฺตํ ทาตพฺพํ. (3)}

\prose{เทฺว ภิกฺขู มิสฺสกํ อาปนฺนา โหนฺติ, เต มิสฺสเก สํฆาทิเสสทิฏฺฐิโน โหนฺติ, เอโก ฉาเทติ, เอโก นจฺฉาเทติ.\csromanspacer{} โย ฉาเทติ, โส ทุกฺกฏํ เทสาเปตพฺโพ, ยถาปฏิจฺฉนฺเน จสฺส ปริวาสํ ทตฺวา อุภินฺนมฺปิ มานตฺตํ ทาตพฺพํ. (4)}

\prose{เทฺว ภิกฺขู มิสฺสกํ อาปนฺนา โหนฺติ, เต มิสฺสเก มิสฺสกทิฏฺฐิโน โหนฺติ, เอโก ฉาเทติ, เอโก นจฺฉาเทติ.\csromanspacer{} โย ฉาเทติ, โส ทุกฺกฏํ เทสาเปตพฺโพ, ยถาปฏิจฺฉนฺเน จสฺส ปริวาสํ ทตฺวา อุภินฺนมฺปิ มานตฺตํ ทาตพฺพํ. (5)}

\prose{เทฺว ภิกฺขู สุทฺธกํ อาปนฺนา โหนฺติ, เต สุทฺธเก สํฆาทิเสสทิฏฺฐิโน โหนฺติ, เอโก ฉาเทติ, เอโก นจฺฉาเทติ.\csromanspacer{} โย ฉาเทติ, โส ทุกฺกฏํ เทสาเปตพฺโพ, อุโภปิ ยถาธมฺมํ การาเปตพฺพา. (6)}

\prose{เทฺว ภิกฺขู สุทฺธกํ อาปนฺนา โหนฺติ, เต สุทฺธเก สุทฺธกทิฏฺฐิโน โหนฺติ, เอโก ฉาเทติ, เอโก นจฺฉาเทติ.\csromanspacer{} โย ฉาเทติ, โส ทุกฺกฏํ เทสาเปตพฺโพ, อุโภปิ ยถาธมฺมํ การาเปตพฺพา. (7)}

\prose{เทฺว ภิกฺขู สํฆาทิเสสํ อาปนฺนา โหนฺติ, เต สํฆาทิเสเส สํฆาทิเสสทิฏฺฐิโน โหนฺติ, เอกสฺส โหติ “อาโรเจสฺสามี”ติ, เอกสฺส โหติ “น อาโรเจสฺสามี”ติ, โส ปฐมมฺปิ ยามํ ฉาเทติ, ทุติยมฺปิ ยามํ ฉาเทติ, ตติยมฺปิ ยามํ ฉาเทติ, อุฏฺฐิเต อรุเณ ฉนฺนา โหติ อาปตฺติ.\csromanspacer{} โย ฉาเทติ, โส ทุกฺกฏํ เทสาเปตพฺโพ, ยถาปฏิจฺฉนฺเน จสฺส ปริวาสํ ทตฺวา อุภินฺนมฺปิ มานตฺตํ ทาตพฺพํ. (8)}

\prose{\csromanfolio{174}เทฺว ภิกฺขู สํฆาทิเสสํ อาปนฺนา โหนฺติ, เต สํฆาทิเสเส สํฆาทิเสสทิฏฺฐิโน โหนฺติ, เต คจฺฉนฺติ “อาโรเจสฺสามา”ติ, เอกสฺส อนฺตรามคฺเค มกฺขธมฺโม อุปฺปชฺชติ “น อาโรเจสฺสามี”ติ, โส ปฐมมฺปิ ยามํ ฉาเทติ, ทุติยมฺปิ ยามํ ฉาเทติ, ตติยมฺปิ ยามํ ฉาเทติ, อุฏฺฐิเต อรุเณ ฉนฺนา โหติ อาปตฺติ.\csromanspacer{} โย ฉาเทติ, โส ทุกฺกฏํ เทสาเปตพฺโพ, ยถาปฏิจฺฉนฺเน จสฺส ปริวาสํ ทตฺวา อุภินฺนมฺปิ มานตฺตํ ทาตพฺพํ. (9)}

\prose{เทฺว ภิกฺขู สํฆาทิเสสํ อาปนฺนา โหนฺติ, เต สํฆาทิเสเส สํฆาทิเสสทิฏฺฐิโน โหนฺติ, เต อุมฺมตฺตกา โหนฺติ, เต ปจฺฉา อนุมฺมตฺตกา หุตฺวา เอโก ฉาเทติ, เอโก นจฺฉาเทติ.\csromanspacer{} โย ฉาเทติ, โส ทุกฺกฏํ เทสาเปตพฺโพ, ยถาปฏิจฺฉนฺเน จสฺส ปริวาสํ ทตฺวา อุภินฺนมฺปิ มานตฺตํ ทาตพฺพํ. (10)}

\prose{เทฺว ภิกฺขู สํฆาทิเสสํ อาปนฺนา โหนฺติ, เต ปาติโมกฺเข อุทฺทิสฺสมาเน เอวํ วทนฺติ “อิทาเนว โข มยํ ชานาม อยมฺปิ กิร ธมฺโม สุตฺตาคโต สุตฺตปริยาปนฺโน อนฺวทฺธมาสํ อุทฺเทสํ อาคจฺฉตี”ติ, เต สํฆาทิเสเส สํฆาทิเสสทิฏฺฐิโน โหนฺติ, เอโก ฉาเทติ, เอโก นจฺฉาเทติ.\csromanspacer{} โย ฉาเทติ, โส ทุกฺกฏํ เทสาเปตพฺโพ, ยถาปฏิจฺฉนฺเน จสฺส ปริวาสํ ทตฺวา อุภินฺนมฺปิ มานตฺตํ ทาตพฺพํ. (11)}

\nitthitam{เทฺวภิกฺขุวาร-เอกาทสกํ นิฏฺฐิตํ.}
\csromanmatikaanchor{mtk.93}
\csromantocmarkref{section}{9. มูลาย-อวิสุทฺธินวก}{mtk.93}
\bookmark[dest={mtk.93},level=1]{9. มูลาย-อวิสุทฺธินวก}
\csromanheader{1}{9. มูลาย-อวิสุทฺธินวก}
\proseitem{182}{อิธ ปน ภิกฺขเว ภิกฺขุ สมฺพหุลา สํฆาทิเสสา อาปตฺติโย อาปชฺชติ ปริมาณมฺปิ อปริมาณมฺปิ เอกนามมฺปิ นานานามปฺปิ สภาคมฺปิ วิสภาคมฺปิ ววตฺถิตมฺปิ สมฺภินฺนมฺปิ.\csromanspacer{} โส สํฆํ ตาสํ อาปตฺตีนํ สโมธานปริวาสํ ยาจติ, ตสฺส สํโฆ ตาสํ อาปตฺตีนํ สโมธานปริวาสํ เทติ, โส ปริวสนฺโต อนฺตรา สมฺพหุลา สํฆาทิเสสา อาปตฺติโย อาปชฺชติ ปริมาณาโย อปฺปฏิจฺฉนฺนาโย.\csromanspacer{} โส สํฆํ อนฺตรา-อาปตฺตีนํ มูลายปฏิกสฺสนํ ยาจติ, ตํ สํโฆ อนฺตรา-อาปตฺตีนํ มูลาย \csromanfolio{175}ปฏิกสฺสติ ธมฺมิเกน กมฺเมน อกุปฺเปน ฐานารเหน, ธมฺเมน สโมธานปริวาสํ เทติ, อธมฺเมน มานตฺตํ เทติ, อธมฺเมน อพฺเภติ.\csromanspacer{} โส ภิกฺขเว ภิกฺขุ อวิสุทฺโธ ตาหิ อาปตฺตีหิ. (1)}

\prose{อิธ ปน ภิกฺขเว ภิกฺขุ สมฺพหุลา สํฆาทิเสสา อาปตฺติโย อาปชฺชติ ปริมาณมฺปิ อปริมาณมฺปิ เอกนามมฺปิ นานานามมฺปิ สภาคมฺปิ วิสภาคมฺปิ ววตฺถิตมฺปิ สมฺภินฺนมฺปิ.\csromanspacer{} โส สํฆํ ตาสํ อาปตฺตีนํ สโมธานปริวาสํ ยาจติ, ตสฺส สํโฆ ตาสํ อาปตฺตีนํ สโมธานปริวาสํ เทติ, โส ปริวสนฺโต อนฺตรา สมฺพหุลา สํฆาทิเสสา อาปตฺติโย อาปชฺชติ ปริมาณาโย ปฏิจฺฉนฺนาโย.\csromanspacer{} โส สํฆํ อนฺตรา-อาปตฺตีนํ มูลายปฏิกสฺสนํ ยาจติ, ตํ สํโฆ อนฺตรา-อาปตฺตีนํ มูลายปฏิกสฺสติ ธมฺมิเกน กมฺเมน อกุปฺเปน ฐานารเหน, ธมฺเมน สโมธานปริวาสํ เทติ, อธมฺเมน มานตฺตํ เทติ, อธมฺเมน อพฺเภติ.\csromanspacer{} โส ภิกฺขเว ภิกฺขุ อวิสุทฺโธ ตาหิ อาปตฺตีหิ. (2)}

\prose{อิธ ปน ภิกฺขเว ภิกฺขุ สมฺพหุลา สํฆาทิเสสา อาปตฺติโย อาปชฺชติ ปริมาณมฺปิ อปริมาณมฺปิ เอกนามมฺปิ นานานามมฺปิ สภาคมฺปิ วิสภาคมฺปิ ววตฺถิตมฺปิ สมฺภินฺนมฺปิ.\csromanspacer{} โส สํฆํ ตาสํ อาปตฺตีนํ สโมธานปริวาสํ ยาจติ, ตสฺส สํโฆ ตาสํ อาปตฺตีนํ สโมธานปริวาสํ เทติ, โส ปริวสนฺโต อนฺตรา สมฺพหุลา สํฆาทิเสสา อาปตฺติโย อาปชฺชติ ปริมาณาโย ปฏิจฺฉนฺนาโยปิ อปฺปฏิจฺฉนฺนาโยปิ.\csromanspacer{} โส สํฆํ อนฺตรา-อาปตฺตีนํ มูลายปฏิกสฺสนํ ยาจติ, ตํ สํโฆ อนฺตรา-อาปตฺตีนํ มูลาย ปฏิกสฺสติ ธมฺมิเกน กมฺเมน อกุปฺเปน ฐานารเหน, ธมฺเมน สโมธานปริวาสํ เทติ, อธมฺเมน มานตฺตํ เทติ, อธมฺเมน อพฺเภติ.\csromanspacer{} โส ภิกฺขเว ภิกฺขุ อวิสุทฺโธ ตาหิ อาปตฺตีหิ. (3)}

\prose{อิธ ปน ภิกฺขเว ภิกฺขุ สมฺพหุลา สํฆาทิเสสา อาปตฺติโย อาปชฺชติ ปริมาณมฺปิ อปริมาณมฺปิ เอกนามมฺปิ นานานามมฺปิ สภาคมฺปิ วิสภาคมฺปิ ววตฺถิตมฺปิ สมฺภินฺนมฺปิ.\csromanspacer{} โส สํฆํ ตาสํ อาปตฺตีนํ สโมธานปริวาสํ ยาจติ, ตสฺส สํโฆ ตาสํ อาปตฺตีนํ สโมธานปริวาสํ เทติ, โส ปริวสนฺโต อนฺตรา สมฺพหุลา สํฆาทิเสสา อาปตฺติโย อาปชฺชติ อปริมาณาโย อปฺปฏิจฺฉนฺนาโย ฯเปฯ อปริมาณาโย ปฏิจฺฉนฺนาโย ฯเปฯ อปริมาณาโย ปฏิจฺฉนฺนาโยปิ \csromanfolio{176}อปฺปฏิจฺฉนฺนาโยปิ ฯเปฯ ปริมาณาโยปิ อปริมาณาโยปิ อปฺปฏิจฺฉนฺนาโย.\csromanspacer{} โส สํฆํ อนฺตรา อาปตฺตีนํ มูลายปฏิกสฺสนํ ยาจติ, ตํ สํโฆ อนฺตรา-อาปตฺตีนํ มูลาย ปฏิกสฺสติ ธมฺมิเกน กมฺเมน อกุปฺเปน ฐานารเหน, ธมฺเมน สโมธานปริวาสํ เทติ, อธมฺเมน มานตฺตํ เทติ, อธมฺเมน อพฺเภติ.\csromanspacer{} โส ภิกฺขเว ภิกฺขุ อวิสุทฺโธ ตาหิ อาปตฺตีหิ. (4 - 7)}

\prose{อิธ ปน ภิกฺขเว ภิกฺขุ สมฺพหุลา สํฆาทิเสสา อาปตฺติโย อาปชฺชติ ปริมาณมฺปิ อปริมาณมฺปิ เอกนามมฺปิ นานานามมฺปิ สภาคมฺปิ วิสภาคมฺปิ ววตฺถิตมฺปิ สมฺภินฺนมฺปิ.\csromanspacer{} โส สํฆํ ตาสํ อาปตฺตีนํ สโมธานปริวาสํ ยาจติ, ตสฺส สํโฆ ตาสํ อาปตฺตีนํ สโมธานปริวาสํ เทติ, โส ปริวสนฺโต อนฺตรา สมฺพหุลา สํฆาทิเสสา อาปตฺติโย อาปชฺชติ ปริมาณาโยปิ อปริมาณาโยปิ ปฏิจฺฉนฺนาโย.\csromanspacer{} โส สํฆํ อนฺตรา-อาปตฺตีนํ มูลายปฏิกสฺสนํ ยาจติ, ตํ สํโฆ อนฺตรา-อาปตฺตีนํ มูลาย ปฏิกสฺสติ ธมฺมิเกน กมฺเมน อกุปฺเปน ฐานารเหน, ธมฺเมน สโมธานปริวาสํ เทติ, อธมฺเมน มานตฺตํ เทติ, อธมฺเมน อพฺเภติ.\csromanspacer{} โส ภิกฺขเว ภิกฺขุ อวิสุทฺโธ ตาหิ อาปตฺตีหิ. (8)}

\prose{อิธ ปน ภิกฺขเว ภิกฺขุ สมฺพหุลา สํฆาทิเสสา อาปตฺติโย อาปชฺชติ ปริมาณมฺปิ อปริมาณมฺปิ เอกนามมฺปิ นานานามมฺปิ สภาคมฺปิ วิสภาคมฺปิ ววตฺถิตมฺปิ สมฺภินฺนมฺปิ.\csromanspacer{} โส สํฆํ ตาสํ อาปตฺตีนํ สโมธานปริวาสํ ยาจติ, ตสฺส สํโฆ ตาสํ อาปตฺตีนํ สโมธานปริวาสํ เทติ, โส ปริวสนฺโต อนฺตรา สมฺพหุลา สํฆาทิเสสา อาปตฺติโย อาปชฺชติ ปริมาณาโยปิ อปริมาณาโยปิ ปฏิจฺฉนฺนาโยปิ อปฺปฏิจฺฉนฺนาโยปิ.\csromanspacer{} โส สํฆํ อนฺตรา-อาปตฺตีนํ มูลายปฏิกสฺสนํ ยาจติ, ตํ สํโฆ อนฺตรา-อาปตฺตีนํ มูลาย ปฏิกสฺสติ ธมฺมิเกน กมฺเมน อกุปฺเปน ฐานารเหน, ธมฺเมน สโมธานปริวาสํ เทติ, อธมฺเมน มานตฺตํ เทติ, อธมฺเมน อพฺเภติ.\csromanspacer{} โส ภิกฺขเว ภิกฺขุ อวิสุทฺโธ ตาหิ อาปตฺตีหิ. (9)}

\nitthitam{มูลาย-อวิสุทฺธินวกํ\footnote{สมูลาวิสุทฺธินวกํ (ก)} นิฏฺฐิตํ.}
\csromanmatikaanchor{mtk.94}
\csromantocmarkref{section}{10. ทุติยนวก}{mtk.94}
\bookmark[dest={mtk.94},level=1]{10. ทุติยนวก}
\csromanheader{1}{10. ทุติยนวก\footnote{อิทํ นวกํ โปราณโปตฺถเกสุ อวิสุกฺขิวเสเนว อาคตํ, วุจฺจมานตติยนวเกน จ สํสฏฺฐํ. ตํ ปฏิวิโสธเกหิ อสํสฏฺฐํ กตฺวา วิสุํ ปติฏฺฐาปิตํ. สีหฬ-สฺยาม-โปตฺถเกสุ ปน ตํ วิสุทฺธิวเสเนว อาคตํ. ตํ ปเนวํ เวทิตพฺพํ\csromandash{}}}
\csromanmatikaanchor{mtk.95}
\csromantocmarkref{section}{10. มูลายวิสุทฺธินวก (สี, สฺยา)}{mtk.95}
\bookmark[dest={mtk.95},level=1]{10. มูลายวิสุทฺธินวก (สี, สฺยา)}
\proseitem{183}{\csromanfolio{177}อิธ ปน ภิกฺขเว ภิกฺขุ สมฺพหุลา สํฆาทิเสสา อาปตฺติโย อาปชฺชติ ปริมาณมฺปิ อปริมาณมฺปิ ฯเปฯ ววตฺถิตมฺปิ สมฺภินฺนมฺปิ.\csromanspacer{} โส สํฆํ ตาสํ อาปตฺตีนํ สโมธานปริวาสํ ยาจติ, ตสฺส สํโฆ ตาสํ อาปตฺตีนํ สโมธานปริวาสํ เทติ, โส ปริวสนฺโต อนฺตรา สมฺพหุลา สํฆาทิเสสา อาปตฺติโย อาปชฺชติ ปริมาณา อปฺปฏิจฺฉนฺนาโย.\csromanspacer{} โส สํฆํ อนฺตรา-อาปตฺตีนํ มูลายปฏิกสฺสนํ ยาจติ, ตํ สํโฆ อนฺตรา-อาปตฺตีนํ มูลาย ปฏิกสฺสติ อธมฺมิเกน กมฺเมน กุปฺเปน อฏฺฐานารเหน, อธมฺเมน สโมธานปริวาสํ เทติ, ธมฺเมน มานตฺตํ เทติ, ธมฺเมน อพฺเภติ.\csromanspacer{} โส ภิกฺขเว ภิกฺขุ อวิสุทฺโธ ตาหิ อาปตฺตีหิ. (\textbf{1})}

\csromanmatikaanchor{mtk.96}
\csromantocmarkref{section}{10. มูลายวิสุทฺธินวก (สี, สฺยา)}{mtk.96}
\bookmark[dest={mtk.96},level=1]{10. มูลายวิสุทฺธินวก (สี, สฺยา)}
\prose{อิธ ปน ภิกฺขเว ภิกฺขุ สมฺพหุลา สํฆาทิเสสา อาปตฺติโย อาปชฺชติ ปริมาณมฺปิ อปริมาณมฺปิ ฯเปฯ ววตฺถิตมฺปิ สมฺภินฺนมฺปิ.\csromanspacer{} โส สํฆํ ตาสํ อาปตฺตีนํ สโมธานปริวาสํ ยาจติ, ตสฺส สํโฆ ตาสํ อาปตฺตีนํ สโมธานปริวาสํ เทติ, โส ปริวสนฺโต อนฺตรา สมฺพหุลา สํฆาทิเสสา อาปตฺติโย อาปชฺชติ ปริมาณา ปฏิจฺฉนฺนาโย.\csromanspacer{} โส สํฆํ อนฺตรา-อาปตฺตีนํ มูลายปฏิกสฺสนํ ยาจติ, ตํ สํโฆ อนฺตรา-อาปตฺตีนํ มูลาย ปฏิกสฺสติ อธมฺมิเกน กมฺเมน กุปฺเปน \csromanfolio{178}อฏฺฐานารเหน, อธมฺเมน สโมธานปริวาสํ เทติ, ธมฺเมน มานตฺตํ เทติ, ธมฺเมน อพฺเภติ.\csromanspacer{} โส ภิกฺขเว ภิกฺขุ อวิสุทฺโธ ตาหิ อาปตฺตีหิ. (2)}

\prose{อิธ ปน ภิกฺขเว ภิกฺขุ สมฺพหุลา สํฆาทิเสสา อาปตฺติโย อาปชฺชติ ปริมาณมฺปิ อาปริมาณมฺปิ ฯเปฯ ววตฺถิตมฺปิ สมฺภินฺนมฺปิ.\csromanspacer{} โส สํฆํ ตาสํ อาปตฺตีนํ สโมธานปริวาสํ ยาจติ, ตสฺส สํโฆ ตาสํ อาปตฺตีนํ สโมธานปริวาสํ เทติ, โส ปริวสนฺโต อนฺตรา สมฺพหุลา สํฆาทิเสสา อาปตฺติโย อาปชฺชติ ปริมาณา ปฏิจฺฉนฺนาโยปิ อปฺปฏิจฺฉนฺนาโยปิ.\csromanspacer{} โส สํฆํ อนฺตรา-อาปตฺตีนํ มูลายปฏิกสฺสนํ ยาจติ, ตํ สํโฆ อนฺตรา-อาปตฺตีนํ มูลาย ปฏิกสฺสติ อธมฺมิเกน กมฺเมน กุปฺเปน อฏฺฐานารเหน, อธมฺเมน สโมธานปริวาสํ เทติ, ธมฺเมน มานตฺตํ เทติ, ธมฺเมน อพฺเภติ.\csromanspacer{} โส ภิกฺขเว ภิกฺขุ อวิสุทฺโธ ตาหิ อาปตฺตีหิ. (3)}

\csromanmatikaanchor{mtk.97}
\csromantocmarkref{section}{10. มูลายวิสุทฺธินวก (สี, สฺยา)}{mtk.97}
\bookmark[dest={mtk.97},level=1]{10. มูลายวิสุทฺธินวก (สี, สฺยา)}
\prose{อิธ ปน ภิกฺขเว ภิกฺขุ สมฺพหุลา สํฆาทิเสสา อาปตฺติโย อาปชฺชติ ปริมาณมฺปิ อปริมาณมฺปิ เอกนามมฺปิ นานานามมฺปิ สภาคมฺปิ วิสภาคมฺปิ ววตฺถิตมฺปิ สมฺภินฺนมฺปิ.\csromanspacer{} โส สํฆํ ตาสํ อาปตฺตีนํ สโมธานปริวาสํ ยาจติ, ตสฺส สํโฆ ตาสํ อาปตฺตีนํ สโมธานปริวาสํ \csromanfolio{179}เทติ, โส ปริวสนฺโต อนฺตรา สมฺพหุลา สํฆาทิเสสา อาปตฺติโย อาปชฺชติ อปริมาณา อปฺปฏิจฺฉนฺนาโย ฯเปฯ อปริมาณา ปฏิจฺฉนฺนาโย ฯเปฯ อปริมาณา ปฏิจฺฉนฺนาโยปิ อปฺปฏิจฺฉนฺนาโยปิ ฯเปฯ ปริมาณาโยปิ อปริมาณาโยปิ อปฺปฏิจฺฉนฺนาโย.\csromanspacer{} โส สํฆํ อนฺตรา อาปตฺตีนํ มูลายปฏิกสฺสนํ ยาจติ, ตํ สํโฆ อนฺตรา-อาปตฺตีนํ มูลาย ปฏิกสฺสติ อธมฺมิเกน กมฺเมน กุปฺเปน อฏฺฐานารเหน, อธมฺเมน สโมธานปริวาสํ เทติ, ธมฺเมน มานตฺตํ เทติ, ธมฺเมน อพฺเภติ.\csromanspacer{} โส ภิกฺขเว ภิกฺขุ อวิสุทฺโธ ตาหิ อาปตฺตีหิ. (4-7)}

\prose{อิธ ปน ภิกฺขเว ภิกฺขุ สมฺพหุลา สํฆาทิเสสา อาปตฺติโย อาปชฺชติ ปริมาณมฺปิ อาปริมาณมฺปิ ฯเปฯ ววตฺถิตมฺปิ สมฺภินฺนมฺปิ.\csromanspacer{} โส สํฆํ ตาสํ อาปตฺตีนํ สโมธานปริวาสํ ยาจติ, ตสฺส สํโฆ ตาสํ อาปตฺตีนํ สโมธานปริวาสํ เทติ, โส ปริวสนฺโต อนฺตรา สมฺพหุลา สํฆาทิเสสา อาปตฺติโย อาปชฺชติ ปริมาณาโยปิ อปริมาณาโยปิ ปฏิจฺฉนฺนาโย.\csromanspacer{} โส สํฆํ อนฺตรา-อาปตฺตีนํ มูลายปฏิกสฺสนํ ยาจติ, ตํ สํโฆ อนฺตรา-อาปตฺตีนํ มูลาย ปฏิกสฺสติ อธมฺมิเกน กมฺเมน กุปฺเปน อฏฺฐานารเหน, อธมฺเมน สโมธานปริวาสํ เทติ, ธมฺเมน มานตฺตํ เทติ, ธมฺเมน อพฺเภติ.\csromanspacer{} โส ภิกฺขเว ภิกฺขุ อวิสุทฺโธ ตาหิ อาปตฺตีหิ. (8)}

\csromanmatikaanchor{mtk.98}
\csromanmatikaanchor{mtk.99}
\csromantocmarkref{section}{11. ตติยนวก}{mtk.98}
\bookmark[dest={mtk.98},level=1]{11. ตติยนวก}
\csromantocmarkref{section}{10. มูลายวิสุทฺธินวก (สี, สฺยา)}{mtk.99}
\bookmark[dest={mtk.99},level=1]{10. มูลายวิสุทฺธินวก (สี, สฺยา)}
\prose{\csromanfolio{180}อิธ ปน ภิกฺขเว ภิกฺขุ สมฺพหุลา สํฆาทิเสสา อาปตฺติโย อาปชฺชติ ปริมาณมฺปิ อปริมาณมฺปิ ฯเปฯ ววตฺถิตมฺปิ สมฺภินฺนมฺปิ.\csromanspacer{} โส สํฆํ ตาสํ อาปตฺตีนํ สโมธานปริวาสํ ยาจติ, ตสฺส สํโฆ ตาสํ อาปตฺตีนํ สโมธานปริวาสํ เทติ, โส ปริวสนฺโต อนฺตรา สมฺพหุลา สํฆาทิเสสา อาปตฺติโย อาปชฺชติ ปริมาณาโยปิ อปริมาณาโยปิ ปฏิจฺฉนฺนาโยปิ อปฺปฏิจฺฉนฺนาโยปิ.\csromanspacer{} โส สํฆํ อนฺตรา-อาปตฺตีนํ มูฏิกสฺสนํ ยาจติ, ตํ สํโฆ อนฺตรา-อาปตฺตีนํ มูลาย ปฏิกสฺสติ อธมฺมิเกน กมฺเมน กุปฺเปน อฏฺฐานารเหน, อธมฺเมน สโมธานปริวาสํ เทติ, ธมฺเมน มานตฺตํ เทติ, ธมฺเมน อพฺเภติ.\csromanspacer{} โส ภิกฺขเว ภิกฺขุ อวิสุทฺโธ ตาหิ อาปตฺตีหิ. (9)}

\nitthitam{ทุติยนวกํ นิฏฺฐิตํ.}
\csromanheader{1}{11. ตติยนวก}
\proseitem{184}{อิธ ปน ภิกฺขเว ภิกฺขุ สมฺพหุลา สํฆาทิเสสา อาปตฺติโย อาปชฺชติ ปริมาณมฺปิ อปริมาณมฺปิ ฯเปฯ ววตฺถิตมฺปิ สมฺภินฺนมฺปิ.\csromanspacer{} โส สํฆํ ตาสํ อาปตฺตีนํ สโมธานปริวาสํ ยาจติ, ตสฺส สํโฆ ตาสํ \csromanfolio{181}อาปตฺตีนํ สโมธานปริวาสํ เทติ, โส ปริวสนฺโต อนฺตรา สมฺพหุลา สํฆาทิเสสา อาปตฺติโย อาปชฺชติ ปริมาณาโย อปฺปฏิจฺฉนฺนาโย.\csromanspacer{} โส สํฆํ อนฺตรา-อาปตฺตีนํ มูลายปฏิกสฺสนํ ยาจติ, ตํ สํโฆ อนฺตรา-อาปตฺตีนํ มูลาย ปฏิกสฺสติ อธมฺมิเกน กมฺเมน กุปฺเปน อฏฺฐานารเหน, อธมฺเมน สโมธานปริวาสํ เทติ, โส “ปริวสามี”ติ มญฺญมาโน อนฺตรา สมฺพหุลา สํฆาทิเสสา อาปตฺติโย อาปชฺชติ ปริมาณาโย อปฺปฏิจฺฉนฺนาโย.\csromanspacer{} โส ตสฺมิํ ภูมิยํ ฐิโต ปุริมา-อาปตฺตีนํ อนฺตรา-อาปตฺติโย สรติ, อปรา-อาปตฺตีนํ อนฺตรา-อาปตฺติโย สรติ, ตสฺส เอวํ โหติ “อหํ โข สมฺพหุลา สํฆาทิเสสา อาปตฺติโย อาปชฺชิํ ปริมาณมฺปิ อปริมาณมฺปิ ฯเปฯ ววตฺถิตมฺปิ สมฺภินฺนมฺปิ.\csromanspacer{} โสหํ สํฆํ ตาสํ อาปตฺตีนํ สโมธานปริวาสํ ยาจิํ, ตสฺส เม สํโฆ ตาสํ อาปตฺตีนํ สโมธานปริวาสํ อทาสิ, โสหํ ปริวสนฺโต อนฺตรา สมฺพหุลา สํฆาทิเสสา อาปตฺติโย อาปชฺชิํ ปริมาณาโย อปฺปฏิจฺฉนฺนาโย.\csromanspacer{} โสหํ สํฆํ อนฺตรา-อาปตฺตีนํ มูลายปฏิกสฺสนํ ยาจิํ, ตํ มํ สํโฆ อนฺตรา-อาปตฺตีนํ มูลาย ปฏิกสฺสิ อธมฺมิเกน กมฺเมน กุปฺเปน อฏฺฐานารเหน, อธมฺเมน สโมธานปริวาสํ อทาสิ, โสหํ ‘ปริวสามี’ติ มญฺญมาโน อนฺตรา สมฺพหุลา สํฆาทิเสสา อาปตฺติโย อาปชฺชิํ ปริมาณาโย อปฺปฏิจฺฉนฺนาโย.\csromanspacer{} โสหํ ตสฺมิํ ภูมิยํ ฐิโต ปุริมา-อาปตฺตีนํ อนฺตรา-อาปตฺติโย สรามิ, อปรา-อาปตฺตีนํ อนฺตรา-อาปตฺติโย สรามิ, ยํนูนาหํ สํฆํ ปุริมา-อาปตฺตีนํ อนฺตรา-อาปตฺตีนญฺจ อปรา-อาปตฺตีนํ อนฺตรา-อาปตฺตีนญฺจ มูลายปฏิกสฺสนํ ยาเจยฺยํ ธมฺมิเกน กมฺเมน อกุปฺเปน ฐานารเหน, ธมฺเมน สโมธานปริวาสํ, ธมฺเมน มานตฺตํ, ธมฺเมน อพฺภานนฺ”ติ.\csromanspacer{} โส สํฆํ ปุริมา-อาปตฺตีนํ อนฺตรา-อาปตฺตีนญฺจ อปรา-อาปตฺตีนํ อนฺตรา-อาปตฺตีนญฺจ มูลายปฏิกสฺสนํ ยาจติ ธมฺมิเกน กมฺเมน อกุปฺเปน ฐานารเหน, ธมฺเมน สโมธานปริวาสํ, ธมฺเมน มานตฺตํ, ธมฺเมน อพฺภานํ.\csromanspacer{} ตํ สํโฆ ปุริมา-อาปตฺตีนํ อนฺตรา-อาปตฺตีนญฺจ อปรา-อาปตฺตีนํ อนฺตรา-อาปตฺตีนญฺจ มูลาย ปฏิกสฺสติ ธมฺมิเกน กมฺเมน อกุปฺเปน ฐานารเหน, ธมฺเมน สโมธานปริวาสํ เทติ, ธมฺเมน มานตฺตํ เทติ, ธมฺเมน อพฺเภติ.\csromanspacer{} โส ภิกฺขเว ภิกฺขุ วิสุทฺโธ ตาหิ อาปตฺตีหิ\footnote{อยํ ปฐมวาโร สี-สฺยา-โปตฺถเกสุ ปริปุณฺโณ ทิสฺสติ.}. (1)}

\prose{\csromanfolio{182}อิธ ปน ภิกฺขเว ภิกฺขุ สมฺพหุลา สํฆาทิเสสา อาปตฺติโย อาปชฺชติ ปริมาณมฺปิ อปริมาณมฺปิ ฯเปฯ ววตฺถิตมฺปิ สมฺภินฺนมฺปิ.\csromanspacer{} โส สํฆํ ตาสํ อาปตฺตีนํ สโมธานปริวาสํ ยาจติ, ตสฺส สํโฆ ตาสํ อาปตฺตีนํ สโมธานปริวาสํ เทติ, โส ปริวสนฺโต อนฺตรา สมฺพหุลา สํฆาทิเสสา อาปตฺติโย อาปชฺชติ ปริมาณา ปฏิจฺฉนฺนาโย.\csromanspacer{} โส สํฆํ อนฺตรา-อาปตฺตีนํ มูลายปฏิกสฺสนํ ยาจติ, ตํ สํโฆ อนฺตรา-อาปตฺตีนํ มูลาย ปฏิกสฺสติ อธมฺมิเกน กมฺเมน กุปฺเปน อฏฺฐานารเหน, อธมฺเมน สโมธานปริวาสํ เทติ, โส “ปริวสฺสามี”ติ มญฺญมาโน อนฺตรา สมฺพหุลา สํฆาทิเสสา อาปตฺติโย อาปชฺชติ ปริมาณา ปฏิจฺฉนฺนาโย.\csromanspacer{} โส ตสฺมิํ ภูมิยํ ฐิโต ปุริมา-อาปตฺตีนํ\footnote{ปุริมานํ อาปตฺตีนํ (ก)} อนฺตรา-อาปตฺติโย สรติ, อปรา-อาปตฺตีนํ อนฺตรา-อาปตฺติโย สรติ, ตสฺส เอวํ โหติ “อหํ โข สมฺพหุลา สํฆาทิเสสา อาปตฺติโย อาปชฺชิํ ปริมาณมฺปิ อปริมาณมฺปิ ฯเปฯ ววตฺถิตมฺปิ สมฺภินฺนมฺปิ.\csromanspacer{} โสหํ สํฆํ ตาสํ อาปตฺตีนํ สโมธานปริวาสํ ยาจิํ, ตสฺส เม สํโฆ ตาสํ อาปตฺตีนํ สโมธานปริวาสํ อทาสิ, โสหํ ปริวสนฺโต อนฺตรา สมฺพหุลา สํฆาทิเสสา อาปตฺติโย อาปชฺชิํ ปริมาณา ปฏิจฺฉนฺนาโย.\csromanspacer{} โสหํ สํฆํ อนฺตรา-อาปตฺตีนํ มูลยปฏิกสฺสนํ ยาจิํ, ตํ มํ สํโฆ อนฺตรา-อาปตฺตีนํ มูลาย ปฏิกสฺสิ อธมฺมิเกน กมฺเมน กุปฺเปน อฏฺฐานารเหน, อธมฺเมน สโมธานปริวาสํ อทาสิ, โสหํ ‘ปริวสามี’ติ มญฺญมาโน อนฺตรา สมฺพหุลา สํฆาทิเสสา อาปตฺติโย อาปชฺชิํ ปริมาณา ปฏิจฺฉนฺนาโย โสหํ ตสฺมิํ ภูมิยํ ฐิโต ปุริมา-อาปตฺตีนํ อนฺตรา-อาปตฺติโย สรามิ, อปรา อาปตฺตีนํ อนฺตรา-อาปตฺติโย สรามิ, ยํนูนาหํ สํฆํ ปุริมา-อาปตฺตีนํ อนฺตรา อาปตฺตีนญฺจ อปรา อาปตฺตีนํ อนฺตรา อาปตฺตีนญฺจ มูลายปฏิกสฺสนํ ยาเจยฺยํ ธมฺมิเกน กมฺเมน อกุปฺเปน ฐานารเหน, ธมฺเมน สโมธานปริวาสํ, ธมฺเมน มานตฺตํ, ธมฺเมน อพฺภานนฺ”ติ.\csromanspacer{} โส สํฆํ ปุริมา-อาปตฺตีนํ อนฺตรา-อาปตฺตีนญฺจ อปรา-อาปตฺตีนํ อนฺตรา-อาปตฺตีนญฺจ มูลายปฏิกสฺสนํ ยาจติ ธมฺมิเกน กมฺเมน อกุปฺเปน ฐานารเหน, ธมฺเมน สโมธานปริวาสํ, ธมฺเมน มานตฺตํ, ธมฺเมน อพฺภานํ.\csromanspacer{} ตํ สํโฆ ปุริมา-อาปตฺตีนํ อนฺตรา-อาปตฺตีนญฺจ อปรา-อาปตฺตีนํ อนฺตรา-อาปตฺตีนญฺจ มูลาย ปฏิกสฺสติ \csromanfolio{183}ธมฺมิเกน กมฺเมน อกุปฺเปน ฐานารเหน, ธมฺเมน สโมธานปริวาสํ เทติ, ธมฺเมน มานตฺตํ เทติ, ธมฺเมน อพฺเภติ.\csromanspacer{} โส ภิกฺขเว ภิกฺขุ วิสุทฺโธ ตาหิ อาปตฺตีหิ. (2)}

\prose{อิธ ปน ภิกฺขเว ภิกฺขุ สมฺพหุลา สํฆาทิเสสา อาปตฺติโย อาปชฺชติ ปริมาณมฺปิ อปริมาณมฺปิ ฯเปฯ ววตฺถิตมฺปิ สมฺภินฺนมฺปิ.\csromanspacer{} โส สํฆํ ตาสํ อาปตฺตีนํ สโมธานปริวาสํ ยาจติ, ตสฺส สํโฆ ตาสํ อาปตฺตีนํ สโมธานปริวาสํ เทติ, โส ปริวสนฺโต อนฺตรา สมฺพหุลา สํฆาทิเสสา อาปตฺติโย อาปชฺชติ ปริมาณาปฏิจฺฉนฺนาโยปิ อปฺปฏิจฺฉนฺนาโยปิ.\csromanspacer{} โส สํฆํ อนฺตรา-อาปตฺตีนํ มูลายปฏิกสฺสนํ ยาจติ, ตํ สํโฆ อนฺตรา-อาปตฺตีนํ มูลาย ปฏิกสฺสติ อธมฺมิเกน กมฺเมน กุปฺเปน อฏฺฐานารเหน, อธมฺเมน สโมธานปริวาสํ เทติ, โส “ปริวสามี”ติ มญฺญมาโน อนฺตรา สมฺพหุลา สํฆาทิเสสา อาปตฺติโย อาปชฺชติ ปริมาณา ปฏิจฺฉนฺนาโยปิ อปฺปฏิจฺฉนาโยปิ.\csromanspacer{} โส ตสฺมิํ ภูมิยํ ฐิโต ปุริมา-อาปตฺตีนํ อนฺตรา-อาปตฺติโย สรติ, อปรา-อาปตฺตีนํ อนฺตรา-อาปตฺติโย สรติ, ตสฺส เอวํ โหติ “อหํ โข สมฺพหุลา สํฆาทิเสสา อาปตฺติโย อาปชฺชิํ ปริมาณมฺปิ อปริมาณมฺปิ ฯเปฯ ววตฺถิตมฺปิ สมฺภินฺนมฺปิ.\csromanspacer{} โสหํ สํฆํ ตาสํ อาปตฺตีนํ สโมธานปริวาสํ ยาจิํ, ตสฺส เม สํโฆ ตาสํ อาปตฺตีนํ สโมธานปริวาสํ อทาสิ, โสหํ ปริวสนฺโต อนฺตรา สมฺพหุลา สํฆาทิเสสา อาปตฺติโย อาปชฺชิํ ปริมาณา ปฏิจฺฉนฺนาโยปิ อปฺปฏิจฺฉนฺนาโยปิ.\csromanspacer{} โสหํ สํฆํ อนฺตรา-อาปตฺตีนํ มูลายปฏิกสฺสนํ ยาจิํ, ตํ มํ สํโฆ อนฺตรา-อาปตฺตีนํ มูลาย ปฏิกสฺสิ อธมฺมิเกน กมฺเมน กุปฺเปน อฏฺฐานารเหน, อธมฺเมน สโมธานปริวาสํ อทาสิ, โสหํ ‘ปริวสามี’ติ มญฺญมาโน อนฺตรา สมฺพหุลา สํฆาทิเสสา อาปตฺติโย อาปชฺชิํ ปริมาณา ปฏิจฺฉนฺนาโยปิ อปฺปฏิจฺฉนฺนาโยปิ.\csromanspacer{} โสหํ ตสฺมิํ ภูมิยํ ฐิโต ปุริมา-อาปตฺตีนํ อนฺตรา-อาปตฺติโย สรามิ, อปรา-อาปตฺตีนํ อนฺตรา-อาปตฺติโย สรามิ, ยํนูนาหํ สํฆํ ปุริมา อาปตฺตีนํ อนฺตรา อาปตฺตีนญฺจ อปรา อาปตฺตีนํ อนฺตรา อาปตฺตีนญฺจ มูลายปฏิกสฺสนํ ยาเจยฺยํ ธมฺมิเกน กมฺเมน อกุปฺเปน ฐานารเหน, ธมฺเมน สโมธานปริวาสํ, ธมฺเมน มานตฺตํ, ธมฺเมน อพฺภานนฺ”ติ.\csromanspacer{} โส สํฆํ ปุริมา-อาปตฺตีนํ อนฺตรา อาปตฺตีนญฺจ อปรา อาปตฺตีนํ อนฺตรา อาปตฺตินญฺจ \csromanfolio{184}มูลายปฏิกสฺสนํ ยาจติ, ธมฺมิเกน กมฺเมน อกุปฺเปน ฐานารเหน, ธมฺเมน สโมธานปริวาสํ, ธมฺเมน มานตฺตํ, ธมฺเมน อพฺภานํ.\csromanspacer{} ตํ สํโฆ ปุริมา-อาปตฺตีนํ อนฺตรา อาปตฺตีนญฺจ อปรา อาปตฺตีนํ อนฺตรา อาปตฺตีนญฺจ มูลาย ปฏิกสฺสติ ธมฺมิเกน กมฺเมน อกุปฺเปน ฐานารเหน, ธมฺเมน สโมธานปริวาสํ เทติ, ธมฺเมน มานตฺตํ เทติ, ธมฺเมน อพฺเภติ.\csromanspacer{} โส ภิกฺขเว ภิกฺขุ วิสุทฺโธ ตาหิ อาปตฺตีหิ. (3)}

\prose{อิธ ปน ภิกฺขเว ภิกฺขุ สมฺพหุลา สํฆาทิเสสา อาปตฺติโย อาปชฺชติ ปริมาณมฺปิ อปริมาณมฺปิ ฯเปฯ ววตฺถิตมฺปิ สมฺภินฺนมฺปิ.\csromanspacer{} โส สํฆํ ตาสํ อาปตฺตีนํ สโมธานปริวาสํ ยาจติ, ตสฺส สํโฆ ตาสํ อาปตฺตีนํ สโมธานปริวาสํ เทติ, โส ปริวสนฺโต อนฺตรา สมฺพหุลา สํฆาทิเสสา อาปตฺติโย อาปชฺชติ อปริมาณาโย อปฺปฏิจฺฉนฺนาโย ฯเปฯ อปริมาณาโย ปฏิจฺฉนฺนาโย ฯเปฯ อปริมาณาโย ปฏิจฺฉนาโยปิ อปฺปฏิจฺฉนฺนาโยปิ ฯเปฯ ปริมาณาโยปิ อปริมาณาโยปิ อปฺปฏิจฺฉนฺนาโย.\csromanspacer{} โส สํฆํ อนฺตรา- อาปตฺตีนํ มูลายปฏิกสฺสนํ ยาจติ.\csromanspacer{} ตํ สํโฆ อนฺตรา-อาปตฺตีนํ มูลาย ปฏิกสฺสิ อธมฺมิเกน กมฺเมน กุปฺเปน อฏฺฐานารเหน, อธมฺเมน สโมธานปริวาสํ เทติ, โส “ปริวสามี”ติ มญฺญมาโน ฯเปฯ.\csromanspacer{} ตํ สํโฆ ปุริมา-อาปตฺตีนํ อนฺตรา-อาปตฺตีนญฺจ อปรา-อาปตฺตีนํ อนฺตรา-อาปตฺตีนญฺจ มูลาย ปฏิกสฺสติ ธมฺมิเกน กมฺเมน อกุปฺเปน ฐานารเหน, ธมฺเมน สโมธานปริวาสํ เทติ, ธมฺเมน มานตฺตํ เทติ, ธมฺเมน อพฺเภติ.\csromanspacer{} โส ภิกฺขเว ภิกฺขุ วิสุทฺโธ ตาหิ อาปตฺตีหิ\footnote{อิเมปิ จตฺตาโร วารา สี-สฺยา-โปตฺถเกสุ ปริปุณฺณา ทิสฺสนฺติ.}. (4 - 7)}

\prose{อิธ ปน ภิกฺขเว ภิกฺขุ สมฺพหุลา สํฆาทิเสสา อาปตฺติโย อาปชฺชติ ปริมาณมฺปิ อปริมาณมฺปิ ฯเปฯ ววตฺถิตมฺปิ สมฺภินฺนมฺปิ.\csromanspacer{} โส สํฆํ ตาสํ อาปตฺตีนํ สโมธานปริวาสํ ยาจติ, ตสฺส สํโฆ ตาสํ อาปตฺตีนํ สโมธานปริวาสํ เทติ, โส ปริวสนฺโต อนฺตรา สมฺพหุลา สํฆาทิเสสา อาปตฺติโย อาปชฺชติ ปริมาณาโยปิ อปริมาณาโยปิ ปฏิจฺฉนฺนาโย.\csromanspacer{} โส สํฆํ อนฺตรา-อาปตฺตีนํ มูลายปฏิกสฺสนํ ยาจติ, ตํ สํโฆ อนฺตรา-อาปตฺตีนํ มูลาย ปฏิกสฺสติ อธมฺมิเกน กมฺเมน กุปฺเปน อฏฺฐานารเหน, อธมฺเมน สโมธานปริวาสํ เทติ, โส “ปริวสามี”ติ มญฺญมาโน อนฺตรา \csromanfolio{185}สํฆาทิเสสา อาปตฺติโย อาปชฺชติ ปริมาณาโยปิ อปริมาณาโยปิ ปฏิจฺฉนฺนาโย.\csromanspacer{} โส ตสฺมิํ ภูมิยํ ฐิโต ปุริมา-อาปตฺตีนํ อนฺตรา-อาปตฺติโย สรติ, อปรา อาปตฺตีนํ อนฺตรา อาปตฺติโย สรติ, ตสฺส เอวํ โหติ “อหํ โข สมฺพหุลา สํฆาทิเสสา อาปตฺติโย อาปชฺชิํ ปริมาณมฺปิ อปริมาณมฺปิ ฯเปฯ ววตฺถิตมฺปิ สมฺภินฺนมฺปิ.\csromanspacer{} โสหํ สํฆํ ตาสํ อาปตฺตีนํ สโมธานปริวาสํ ยาจิํ, ตสฺส เม สํโฆ ตาสํ อาปตฺตีนํ สโมธานปริวาสํ อทาสิ, โสหํ ปริวสนฺโต อนฺตรา สมฺพหุลา สํฆาทิเสสา อาปตฺติโย อาปชฺชิํ ปริมาณาโยปิ อปริมาณาโยปิ ปฏิจฺฉนฺนาโย.\csromanspacer{} โสหํ สํฆํ อนฺตรา-อาปตฺตีนํ มูลาย ปฏิกสฺสนํ ยาจิํ, ตํ มํ สํโฆ อนฺตรา-อาปตฺตีนํ มูลาย ปฏิกสฺสิ อธมฺมิเกน กมฺเมน กุปฺเปน อฏฺฐานารเหน, อธมฺเมน สโมธานปริวาสํ อทาสิ, โสหํ ‘ปริวสามี’ติ มญฺญมาโน อนฺตรา สมฺพหุลา สํฆาทิเสสา อาปตฺติโย อาปชฺชิํ ปริมาณาโยปิ อปริมาณาโยปิ ปฏิจฺฉนฺนาโย.\csromanspacer{} โสหํ ตสฺมิํ ภูมิยํ ฐิโต ปุริมา-อาปตฺตีนํ อนฺตรา อาปตฺติโย สรามิ, อปราปตฺตีนํ อนฺตรา-อาปตฺติโย สรามิ, ยํนูนาหํ สํฆํ ปุริมา-อาปตฺตีนํ อนฺตราปตฺตีนญฺจ อปรา อาปตฺตีนํ อนฺตรา อาปตฺตีนญฺจ มูลาย ปฏิกสฺสนํ ยาเจยฺยํ ธมฺมิเกน กมฺเมน อกุปฺเปน ฐานารเหน, ธมฺเมน สโมธานปริวาสํ, ธมฺเมน มานตฺตํ, ธมฺเมน อพฺภานนฺ”ติ.\csromanspacer{} โส สํฆํ ปุริมา-อาปตฺตีนํ อนฺตราปตฺตีนญฺจ อปรา อาปตฺตีนํ อนฺตรา อาปตฺตีนญฺจ มูลายปฏิกสฺสนํ ยาจติ, ธมฺมิเกน กมฺเมน อกุปฺเปน ฐานารเหน, ธมฺเมน สโมธานปริวาสํ, ธมฺเมน มานตฺตํ, ธมฺเมน อพฺภานํ.\csromanspacer{} ตํ สํโฆ ปุริมา-อาปตฺตีนํ อนฺตรา อาปตฺตีนญฺจ อปรา อาปตฺตีนํ อนฺตรา อาปตฺตีนญฺจ มูลาย ปฏิกสฺสติ ธมฺมิเกน กมฺเมน อกุปฺเปน ฐานารเหน, ธมฺเมน สโมธานปริวาสํ เทติ, ธมฺเมน มานตฺตํ เทติ, ธมฺเมน อพฺเภติ.\csromanspacer{} โส ภิกฺขเว ภิกฺขุ วิสุทฺโธ ตาหิ อาปตฺตีหิ. (8)}

\prose{อิธ ปน ภิกฺขเว ภิกฺขุ สมฺพหุลา สํฆาทิเสสา อาปตฺติโย อาปชฺชติ ปริมาณมฺปิ อปริมาณมฺปิ ฯเปฯ ววตฺถิตมฺปิ สมฺภินฺนมฺปิ.\csromanspacer{} โส สํฆํ ตาสํ อาปตฺตีนํ สโมธานปริวาสํ ยาจติ, ตสฺส สํโฆ ตาสํ อาปตฺตีนํ สโมธานปริวาสํ เทติ, โส ปริวสนฺโต อนฺตรา สมฺพหุลา สํฆาทิเสสา อาปตฺติโย อาปชฺชติ ปริมาณาโยปิ อปริมาณาโยปิ ปฏิจฺฉนฺนาโยปิ อปฺปฏิจฺฉนฺนาโยปิ.\csromanspacer{} โส สํฆํ \csromanfolio{186}อนฺตรา-อาปตฺตีนํ มูลายปฏิกสฺสนํ ยาจติ, ตํ สํโฆ อนฺตรา-อาปตฺตีนํ มูลาย ปฏิกสฺสติ อธมฺมิเกน กมฺเมน กุปฺเปน อฏฺฐานารเหน, อธมฺเมน สโมธานปริวาสํ เทติ, โส “ปริวสามี”ติ มญฺญมาโน อนฺตรา สมฺพหุลา สํฆาทิเสสา อาปตฺติโย อาปชฺชติ ปริมาณาโยปิ อปริมาณาโยปิ ปฏิจฺฉนฺนาโยปิ อปฺปฏิจฺฉนฺนาโยปิ.\csromanspacer{} โส ตสฺมิํ ภูมิยํ ฐิโต ปุริมา-อาปตฺตีนํ อนฺตรา-อาปตฺติโย สรติ.\csromanspacer{} อปรา อาปตฺตีนํ อนฺตรา-อาปตฺติโย สรติ, ตสฺส เอวํ โหติ “อหํ โข สมฺพหุลา สํฆาทิเสสา อาปตฺติโย อาปชฺชิํ ปริมาณมฺปิ อปริมาณมฺปิ ฯเปฯ ววตฺถิตมฺปิ สมฺภินฺนมฺปิ.\csromanspacer{} โสหํ สํฆํ ตาสํ อาปตฺตีนํ สโมธานปริวาสํ ยาจิํ, ตสฺส เม สํโฆ ตาสํ อาปตฺตีนํ สโมธานปริวาสํ อทาสิ, โสหํ ปริวสนฺโต อนฺตรา สมฺพหุลา สํฆาทิเสสา อาปตฺติโย อาปชฺชิํ ปริมาณาโยปิ อปริมาณาโยปิ ปฏิจฺฉนฺนาโยปิ อปฺปฏิจฺฉนฺนาโยปิ.\csromanspacer{} โสหํ สํฆํ อนฺตรา-อาปตฺตีนํ มูลายปฏิกสฺสนํ ยาจิํ, ตํ มํ สํโฆ อนฺตรา-อาปตฺตีนํ มูลาย ปฏิกสฺสิ อธมฺมิเกน กมฺเมน กุปฺเปน อฏฺฐานารเหน, อธมฺเมน สโมธานปริวาสํ อทาสิ, โสหํ ‘ปริวสามี’ติ มญฺญมาโน อนฺตรา สมฺพหุลา สํฆาทิเสสา อาปตฺติโย อาปชฺชิํ ปริมาณาโยปิ อปริมาณาโยปิ ปฏิจฺฉนฺนาโยปิ อปฺปฏิจฺฉนฺนาโยปิ.\csromanspacer{} โสหํ ตสฺมิํ ภูมิยํ ฐิโต ปุริมา-อาปตฺตีนํ อนฺตรา-อาปตฺติโย สรามิ, อปรา อาปตฺตีนํ อนฺตรา-อาปตฺติโย สรามิ, ยํนูนาหํ สํฆํ ปุริมา-อาปตฺตีนํ อนฺตรา อาปตฺตีนญฺจ อปรา อาปตฺตีนํ อนฺตรา อาปตฺตีนญฺจ มูลายปฏิกสฺสนํ ยาเจยฺยํ ธมฺมิเกน กมฺเมน อกุปฺเปน ฐานารเหน, ธมฺเมน สโมธานปริวาสํ, ธมฺเมน มานตฺตํ, ธมฺเมน อพฺภานนฺ”ติ.\csromanspacer{} โส สํฆํ ปุริมา-อาปตฺตีนํ อนฺตรา อาปตฺตีนญฺจ อปรา อาปตฺตีนํ อนฺตรา อาปตฺตีนญฺจ มูลายปฏิกสฺสนํ ยาจติ ธมฺมิเกน กมฺเมน อกุปฺเปน ฐานารเหน, ธมฺเมน สโมธานปริวาสํ, ธมฺเมน มานตฺตํ, ธมฺเมน อพฺภานํ.\csromanspacer{} ตํ สํโฆ ปุริมา-อาปตฺตีนํ อนฺตรา อาปตฺตีนญฺจ อปรา อาปตฺตีนํ อนฺตรา อาปตฺตีนญฺจ มูลาย ปฏิกสฺสติ ธมฺมิเกน กมฺเมน อกุปฺเปน ฐานารเหน, ธมฺเมน สโมธานปริวาสํ เทติ, ธมฺเมน มานตฺตํ เทติ, ธมฺเมน อพฺเภติ.\csromanspacer{} โส ภิกฺขเว ภิกฺขุ วิสุทฺโธ ตาหิ อาปตฺตีหิ. (9)}

\nitthitam{ตติยนวกํ นิฏฺฐิตํ.}
\vspace{\tipitakaparskip}
\csromancenter{สมุจฺจยกฺขนฺธโก ตติโย.}
\titlehead{\textbf{ตสฺสุทฺทานํ}}
\csromanmatikaanchor{mtk.100}
\csromantocmarkref{section}{อุทฺทานคาถา}{mtk.100}
\bookmark[dest={mtk.100},level=1]{อุทฺทานคาถา}
\csromangathameasure{อปฺปฏิจฺฉนฺนา เอกาห, ทฺวีห ตีห จตูห จ. \\ ปญฺจาหปกฺขทสนฺนํ, อาปตฺติมาห มหามุนิ. \\ สุทฺธนฺโต จ วิพฺภมนฺโต, ปริมาณมุขํ เทฺว ภิกฺขู. \\ ตตฺถ สญฺญิโน เทฺว ยถา, เวมติกา ตเถว จ. \\ มิสฺสกทิฏฺฐิโน เทฺว จ, อสุทฺธเกกทิฏฺฐิโน. \\ เทฺว เจว สุทฺธทิฏฺฐิโน. ตเถว จ เอโก ฉาเทติ, \\ อถ มกฺขมเตน จ. อุมฺมตฺตกเทสนญฺจ, \\ มูลา อฏฺฐารส วิสุทฺธโต. อาจริยานํ วิภชฺชปทานํ, \\ ตมฺพปณฺณิทีปปสาทกานํ. มหาวิหารวาสีนํ, \\ วาจนา สทฺธมฺมฏฺฐิติยาติ.}
\csromangathagroup{\csromangathastanza{\csromanfolio{187}อปฺปฏิจฺฉนฺนา เอกาห, ทฺวีห ตีห จตูห จ. \\ ปญฺจาหปกฺขทสนฺนํ, อาปตฺติมาห มหามุนิ.}\csromangathastanza{สุทฺธนฺโต จ วิพฺภมนฺโต, ปริมาณมุขํ เทฺว ภิกฺขู. \\ ตตฺถ สญฺญิโน เทฺว ยถา, เวมติกา ตเถว จ.}\csromangathastanza{มิสฺสกทิฏฺฐิโน เทฺว จ, อสุทฺธเกกทิฏฺฐิโน. \\ เทฺว เจว สุทฺธทิฏฺฐิโน. ตเถว จ เอโก ฉาเทติ,}\csromangathastanza{อถ มกฺขมเตน จ. อุมฺมตฺตกเทสนญฺจ, \\ มูลา อฏฺฐารส\footnote{ปนฺนรส (ก)} วิสุทฺธโต. อาจริยานํ วิภชฺชปทานํ\footnote{ปนฺนรส (ก)},}\csromangathastanza{ตมฺพปณฺณิทีปปสาทกานํ. มหาวิหารวาสีนํ, \\ วาจนา สทฺธมฺมฏฺฐิติยาติ.}}

\nitthitam{\textbf{สมุจฺจยกฺขนฺธโก นิฏฺฐิโต.}}
\csromanmatikaanchor{mtk.101}
\csromantocmarknopage{chapter}{4. สมถกฺขนฺธก}{mtk.101}
\bookmark[dest={mtk.101},level=0]{4. สมถกฺขนฺธก}
\chapterheadpageread{4. สมถกฺขนฺธก}
\csromanmatikaanchor{mtk.102}
\csromantocmarkref{section}{1. สมฺมุขาวินย}{mtk.102}
\bookmark[dest={mtk.102},level=1]{1. สมฺมุขาวินย}
\csromanheader{1}{1. สมฺมุขาวินย}
\proseitem{185}{\csromanfolio{188}เตน สมเยน พุทฺโธ ภควา สาวตฺถิยํ วิหรติ เชตวเน อนาถปิณฺฑิกสฺส อาราเม.\csromanspacer{} เตน โข ปน สมเยน ฉพฺพคฺคิยา ภิกฺขู อสมฺมุขีภูตานํ ภิกฺขูนํ กมฺมานิ กโรนฺติ ตชฺชนียมฺปิ นิยสฺสมฺปิ ปพฺพาชนียมฺปิ ปฏิสารณียมฺปิ อุกฺเขปนียมฺปิ.\csromanspacer{} เย เต ภิกฺขู อปฺปิจฺฉา ฯเปฯ เต อุชฺฌายนฺติ ขิยฺยนฺติ วิปาเจนฺติ “กถํ หิ นาม ฉพฺพคฺคิยา ภิกฺขู อสมฺมุขีภูตานํ ภิกฺขูนํ กมฺมานิ กริสฺสนฺติ ตชฺชนียมฺปิ นิยสฺสมฺปิ ปพฺพาชนียมฺปิ ปฏิสารณียมฺปิ อุกฺเขปณียมฺปี”ติ.\csromanspacer{} อถ โข เต ภิกฺขู ภควโต เอตมตฺถํ อาโรเจสุํ ฯเปฯ.\csromanspacer{} สจฺจํ กิร ภิกฺขเว ฉพฺพคฺคิยา ภิกฺขู อสมฺมุขีภูตานํ ภิกฺขูนํ กมฺมานิ กโรนฺติ ตชฺชนียมฺปิ นิยสฺสมฺปิ ปพฺพาชนียมฺปิ ปฏิสารณียมฺปิ อุกฺเขปนียมฺปีติ.\csromanspacer{} สจฺจํ ภควาติ.\csromanspacer{} วิครหิ พุทฺโธ ภควา อนนุจฺฉวิกํ ภิกฺขเว เตสํ โมฆปุริสานํ อนนุโลมิกํ อปฺปติรูปํ อสฺสามณกํ อกปฺปิยํ อกรณียํ.\csromanspacer{} กถํ หิ นาม เต ภิกฺขเว โมฆปุริสา อสมฺมุขีภูตานํ ภิกฺขูนํ กมฺมานิ กริสฺสนฺติ ตชฺชนียมฺปิ นิยสฺสมฺปิ ปพฺพาชนียมฺปิ ปฏิสารณียมฺปิ อุกฺเขปนียมฺปิ.\csromanspacer{} เนตํ ภิกฺขเว อปฺปสนฺนานํ วา ปสาทาย ฯเปฯ วิครหิตฺวา ฯเปฯ ธมฺมิํ กถํ กตฺวา ภิกฺขู อามนฺเตสิ\csromandash{}น ภิกฺขเว อสมฺมุขีภูตานํ ภิกฺขูนํ กมฺมํ กาตพฺพํ ตชฺชนียํ วา นิยสฺสํ วา ปพฺพาชนียํ วา ปฏิสารณียํ วา อุกฺเขปนียํ วา, โย กเรยฺย, อาปตฺติ ทุกฺกฏสฺส.}

\proseitem{186}{อธมฺมวาที ปุคฺคโล, อธมฺมวาที สมฺพหุลา, อธมฺมวาที สํโฆ, ธมฺมวาที ปุคฺคโล, ธมฺมวาที สมฺพหุลา, ธมฺมวาที สํโฆ.}

\csromanmatikaanchor{mtk.103}
\csromantocmarkref{subsection}{กณฺหปกฺขนวก}{mtk.103}
\bookmark[dest={mtk.103},level=2]{กณฺหปกฺขนวก}
\csromanheader{2}{\textbf{กณฺหปกฺขนวก}}
\proseitem{187}{อธมฺมวาที ปุคฺคโล ธมฺมวาทิํ ปุคฺคลํ สญฺญาเปติ นิชฺฌาเปติ เปกฺเขติ อนุเปกฺเขติ ทสฺเสติ อนุทสฺเสติ “อยํ ธมฺโม, อยํ วินโย, อิทํ สตฺถุสาสนํ, อิมํ คณฺหาหิ, อิมํ โรเจหี”ติ.\csromanspacer{} เอวญฺเจตํ อธิกรณํ วูปสมฺมติ, อธมฺเมน วูปสมฺมติ สมฺมุขาวินยปติรูปเกน. (1)}

\prose{อธมฺมวาที ปุคฺคโล ธมฺมวาที สมฺพหุเล สญฺญาเปติ นิชฺฌาเปติ เปกฺเขติ อนุเปกฺเขติ ทสฺเสติ อนุทสฺเสติ “อยํ ธมฺโม, อยํ \csromanfolio{189}วินโย, อิทํ สตฺถุสาสนํ, อิมํ คณฺหถ, อิมํ โรเจถา”ติ.\csromanspacer{} เอวญฺเจตํ อธิกรณํ วูปสมฺมติ, อธมฺเมน วูปสมฺมติ สมฺมุขาวินยปติรูปเกน. (2)}

\prose{อธมฺมวาที ปุคฺคโล ธมฺมวาทิํ สํฆํ สญฺญาเปติ นิชฺฌาเปติ เปกฺเขติ อนุเปกฺเขติ ทสฺเสติ อนุทสฺเสติ “อยํ ธมฺโม, อยํ วินโย, อิทํ สตฺถุสาสนํ, อิมํ คณฺหาหิ, อิมํ โรเจหี”ติ\footnote{อิมํ คณฺหถ, อิมํ โรเจถาติ (ก)}.\csromanspacer{} เอวญฺเจตํ อธิกรณํ วูปสมฺมติ, อธมฺเมน วูปสมฺมติ สมฺมุขาวินยปติรูปเกน. (3)}

\prose{อธมฺมวาที สมฺพหุลา ธมฺมวาทิํ ปุคฺคลํ สญฺญาเปนฺติ นิชฺฌาเปนฺติ เปกฺเขนฺติ อนุเปกฺเขนฺติ ทสฺเสนฺติ อนุทสฺเสนฺติ “อยํ ธมฺโม, อยํ วินโย, อิทํ สตฺถุสาสนํ, อิมํ คณฺหาหิ, อิมํ โรเจหี”ติ.\csromanspacer{} เอวญฺเจตํ อธิกรณํ วูปสมฺมติ, อธมฺเมน วูปสมฺมติ สมฺมุขาวินยปติรูปเกน. (4)}

\prose{อธมฺมวาที สมฺพหุลา ธมฺมวาที สมฺพหุเล สญฺญาเปนฺติ นิชฺฌาเปนฺติ เปกฺเขนฺติ อนุเปกฺเขนฺติ ทสฺเสนฺติ อนุทสฺเสนฺติ “อยํ ธมฺโม, อยํ วินโย, อิทํ สตฺถุสาสนํ, อิมํ คณฺหถ, อิมํ โรเจถา”ติ.\csromanspacer{} เอวญฺเจตํ อธิกรณํ วูปสมฺมติ, อธมฺเมน วูปสมฺมติ สมฺมุขาวินยปติรูปเกน. (5)}

\prose{อธมฺมวาที สมฺพหุลา ธมฺมวาทิํ สํฆํ สญฺญาเปนฺติ นิชฺฌาเปนฺติ เปกฺเขนฺติ อนุเปกฺเขนฺติ ทสฺเสนฺติ อนุทสฺเสนฺติ “อยํ ธมฺโม, อยํ วินโย, อิทํ สตฺถุสาสนํ, อิมํ คณฺหาหิ, อิมํ โรเจหี”ติ.\csromanspacer{} เอวญฺเจตํ อธิกรณํ วูปสมฺมติ, อธมฺเมน วูปสมฺมติ สมฺมุขาวินยปติรูปเกน. (6)}

\prose{อธมฺมวาที สํโฆ ธมฺมวาทิํ ปุคฺคลํ สญฺญาเปติ นิชฺฌาเปติ เปกฺเขติ อนุเปกฺเขติ ทสฺเสติ อนุทสฺเสติ “อยํ ธมฺโม, อยํ วินโย, อิทํ สตฺถุสาสนํ, อิมํ คณฺหาหิ, อิมํ โรเจหี”ติ.\csromanspacer{} เอวญฺเจตํ อธิกรณํ วูปสมฺมติ, อธมฺเมน วูปสมฺมติ สมฺมุขาวินยปติรูปเกน. (7)}

\prose{อธมฺมวาที สํโฆ ธมฺมวาที สมฺพหุเล สญฺญาเปติ นิชฺฌาเปติ เปกฺเขติ อนุเปกฺเขติ ทสฺเสติ อนุทสฺเสติ “อยํ ธมฺโม, อยํ วินโย, อิทํ สตฺถุสาสนํ, อิมํ คณฺหถ, อิมํ โรเจถา”ติ.\csromanspacer{} เอวญฺเจตํ อธิกรณํ วูปสมฺมติ, อธมฺเมน วูปสมฺมติ สมฺมุขาวินยปติรูปเกน. (8)}

\prose{อธมฺมวาที สํโฆ ธมฺมวาทิํ สํฆํ สญฺญาเปติ นิชฺฌาเปติ เปกฺเขติ, อนุเปกฺเขติ ทสฺเสติ อนุทสฺเสติ “อยํ ธมฺโม, อยํ วินโย, \csromanfolio{190}อิทํ สตฺถุสาสนํ, อิมํ คณฺหาหิ, อิมํ โรเจหี”ติ.\csromanspacer{} เอวญฺเจตํ อธิกรณํ วูปสมฺมติ, อธมฺเมน วูปสมฺมติ สมฺมุขาวินยปติรุปเกน. (9) กณฺหปกฺขนวกํ นิฏฺฐิตํ.}

\csromanmatikaanchor{mtk.104}
\csromantocmarkref{subsection}{สุกฺกปกฺขนวก}{mtk.104}
\bookmark[dest={mtk.104},level=2]{สุกฺกปกฺขนวก}
\csromanheader{2}{\textbf{สุกฺกปกฺขนวก}}
\proseitem{188}{ธมฺมวาที ปุคฺคโล อธมฺมวาทิํ ปุคฺคลํ สญฺญาเปติ นิชฺฌาเปติ เปกฺเขติ อนุเปกฺเขติ ทสฺเสติ อนุทสฺเสติ “อยํ ธมฺโม, อยํ วินโย, อิทํ สตฺถุสาสนํ, อิมํ คณฺหาหิ, อิมํ โรเจหี”ติ.\csromanspacer{} เอวญฺเจตํ อธิกรณํ วูปสมฺมติ, ธมฺเมน วูปสมฺมติ สมฺมุขาวินเยน. (1)}

\prose{ธมฺมวาที ปุคฺคโล อธมฺมวาที สมฺพหุเล สญฺญาเปติ นิชฺฌาเปติ เปกฺเขติ อนุเปกฺเขติ ทสฺเสติ อนุทสฺเสติ “อยํ ธมฺโม, อยํ วินโย, อิทํ สตฺถุสาสนํ, อิมํ คณฺหถ, อิมํ โรเจถา”ติ.\csromanspacer{} เอวญฺเจตํ อธิกรณํ วูปสมฺมติ, ธมฺเมน วูปสมฺมติ สมฺมุขาวินเยน. (2)}

\prose{ธมฺมวาที ปุคฺคโล อธมฺมวาทิํ สํฆํ สญฺญาเปติ นิชฺฌาเปติ เปกฺเขติ อนุเปกฺเขติ ทสฺเสติ อนุทสฺเสติ “อยํ ธมฺโม, อยํ วินโย, อิทํ สตฺถุสาสนํ, อิมํ คณฺหาหิ, อิมํ โรเจหี”ติ.\csromanspacer{} เอวญฺเจตํ อธิกรณํ วูปสมฺมติ, ธมฺเมน วูปสมฺมติ สมฺมุขาวินเยน. (3)}

\prose{ธมฺมวาที สมฺพหุลา อธมฺมวาทิํ ปุคฺคลํ สญฺญาเปนฺติ นิชฺฌาเปนฺติ เปกฺเขนฺติ อนุเปกฺเขนฺติ ทสฺเสนฺติ อนุทสฺเสนฺติ “อยํ ธมฺโม, อยํ วินโย, อิทํ สตฺถุสาสนํ, อิมํ คณฺหาหิ, อิมํ โรเจหี”ติ.\csromanspacer{} เอวญฺเจตํ อธิกรณํ วูปสมฺมติ ธมฺเมน วูปสมฺมติ สมฺมุขาวินเยน. (4)}

\prose{ธมฺมวาที สมฺพหุลา อธมฺมวาที สมฺพหุเล สญฺญาเปนฺติ นิชฺฌาเปนฺติ เปกฺเขนฺติ อนุเปกฺเขนฺติ ทสฺเสนฺติ อนุทสฺเสนฺติ “อยํ ธมฺโม, อยํ วินโย, อิทํ สตฺถุสาสนํ, อิมํ คณฺหถ, อิมํ โรเจถา”ติ.\csromanspacer{} เอวญฺเจตํ อธิกรณํ วูปสมฺมติ, ธมฺเมน วูปสมติ สมฺมุขาวินเยน. (5)}

\prose{ธมฺมวาที สมฺพหุลา อธมฺมวาทิํ สํฆํ สญฺญาเปนฺติ นิชฺฌาเปนฺติ เปกฺเขนฺติ อนุเปกฺเขนฺติ ทสฺเสนฺติ อนุทสฺเสนฺติ “อยํ ธมฺโม, อยํ วินโย, อิทํ \csromanfolio{191}สตฺถุสาสนํ, อิมํ คณฺหาหิ, อิมํ โรเจหี”ติ.\csromanspacer{} เอวญฺเจตํ อธิกรณํ วูปสมฺมติ, ธมฺเมน วูปสมฺมติ สมฺมุขาวินเยน. (6)}

\prose{ธมฺมวาที สํโฆ อธมฺมวาทิํ ปุคฺคลํ สญฺญาเปติ นิชฺฌาเปติ เปกฺเขติ อนุเปกฺเขติ ทสฺเสติ อนุทสฺเสติ “อยํ ธมฺโม, อยํ วินโย, อิทํ สตฺถุสาสนํ, อิมํ คณฺหาหิ, อิมํ โรเจหี”ติ.\csromanspacer{} เอวญฺเจตํ อธิกรณํ วูปสมฺมติ, ธมฺเมน วูปสมฺมติ สมฺมุขาวินเยน. (7)}

\prose{ธมฺมวาที สํโฆ อธมฺมวาที สมฺพหุเล สญฺญาเปติ นิชฺฌาเปติ เปกฺเขติ อนุเปกฺเขติ ทสฺเสติ อนุทสฺเสติ “อยํ ธมฺโม, อยํ วินโย, อิทํ สตฺถุสาสนํ, อิมํ คณฺหถ, อิมํ โรเจถา”ติ.\csromanspacer{} เอวญฺเจตํ อธิกรณํ วูปสมฺมติ, ธมฺเมน วูปสมฺมติ สมฺมุขาวินเยน. (8)}

\prose{ธมฺมวาที สํโฆ อธมฺมวาทิํ สํฆํ สญฺญาเปติ นิชฺฌาเปติ เปกฺเขติ อนุเปกฺเขติ ทสฺเสติ อนุทสฺเสติ “อยํ ธมฺโม, อยํ วินโย, อิทํ สตฺถุสาสนํ, อิมํ คณฺหาหิ, อิมํ โรเจหี”ติ.\csromanspacer{} เอวญฺเจตํ อธิกรณํ วูปสมฺมติ, ธมฺเมน วูปสมฺมติ สมฺมุขาวินเยนาติ. (9)}

\nitthitam{สุกฺกปกฺขนวกํ นิฏฺฐิตํ.}
\csromanmatikaanchor{mtk.105}
\csromantocmarkref{section}{2. สติวินย}{mtk.105}
\bookmark[dest={mtk.105},level=1]{2. สติวินย}
\csromanheader{1}{2. สติวินย}
\proseitem{189}{\csromansymbolfootnote{*}{อิทํ วตฺถุ วิ 1. 243 ปิฏฺฐาทีสุปิ อาคตํ.}เตน สมเยน พุทฺโธ ภควา ราชคเห วิหรติ เวฬุวเน กลนฺทกนิวาเป.\csromanspacer{} เตน โข ปน สมเยน อายสฺมตา ทพฺเพน มลฺลปุตฺเตน ชาติยา สตฺตวสฺเสน อรหตฺตํ สจฺฉิกตํ โหติ, ยํ กิญฺจิ สาวเกน ปตฺตพฺพํ, สพฺพํ เตน อนุปฺปตฺตํ โหติ, นตฺถิ จสฺส กิญฺจิ อุตฺตริ\footnote{อุตฺตริํ (สี)} กรณียํ, กตสฺส วา ปติจโย.\csromanspacer{} อถ โข อายสฺมโต ทพฺพสฺส มลฺลปุตฺตสฺส รโหคตสฺส ปฏิสลฺลีนสฺส เอวํ เจตสฺโส ปริวิตกฺโก อุทปาทิ “มยา โข ชาติยา สตฺตวสฺเสน อรหตฺตํ สจฺฉิกตํ, ยํ กิญฺจิ สาวเกน ปตฺตพฺพํ, สพฺพํ มยา อนุปฺปตฺตํ, นตฺถิ จ เม กิญฺจิ อุตฺตริกรณียํ, กตสฺส วา ปติจโย, กิํ นุ โข อหํ สํฆสฺส เวยฺยาวจฺจํ กเรยฺยนฺ”ติ.}

\prose{\csromanfolio{192}อถ โข อายสฺมโต ทพฺพสฺส มลฺลปุตฺตสฺส เอตทโหสิ “ยํนูนาหํ สํฆสฺส เสนาสนญฺจ ปญฺญเปยฺยํ ภตฺตานิ จ อุทฺทิเสยฺยนฺ”ติ.\csromanspacer{} อถ โข อายสฺมา ทพฺโพ มลฺลปุตฺโต สายนฺหสมยํ ปฏิสลฺลานา วุฏฺฐิโต เยน ภควา เตนุปสงฺกมิ, อุปสงฺกมิตฺวา ภควนฺตํ อภิวาเทตฺวา เอกมนฺตํ นิสีทิ, เอกมนฺตํ นิสินฺโน โข อายสฺมา ทพฺโพ มลฺลปุตฺโต ภควนฺตํ เอตทโวจ “อิธ มยฺหํ ภนฺเต รโหคตสฺส ปฏิสลฺลีนสฺส เอวํ เจตโส ปริวิตกฺโก อุทปาทิ ‘มยา โข ชาติยา สตฺตวสฺเสน อรหตฺตํ สจฺฉิกตํ, ยํ กิญฺจิ สาวเกน ปตฺตพฺพํ, สพฺพํ มยา อนุปฺปตฺตํ, นตฺถิ จ เม กิญฺจิ อุตฺตริกรณียํ, กตสฺส วา ปติจโย, กิํ นุ โข อหํ สํฆสฺส เวยฺยาวจฺจํ กเรยฺยนฺ’ติ.\csromanspacer{} ตสฺส มยฺหํ ภนฺเต เอตทโหสิ ‘ยํนูนาหํ สํฆสฺส เสนาสนญฺจ ปญฺญเปยฺยํ ภตฺตานิ จ อุทฺทิเสยฺยนฺ’ติ, อิจฺฉามหํ ภนฺเต สํฆสฺส เสนาสนญฺจ ปญฺญาเปตุํ ภตฺตานิ จ อุทฺทิสิตุนฺ”ติ.\csromanspacer{} สาธุ สาธุ ทพฺพ, เตน หิ ตฺวํ ทพฺพ สํฆสฺส เสนาสนญฺจ ปญฺญเปติ ภตฺตานิ จ อุทฺทิสาหีติ\footnote{อุทฺทิสาติ (วิ 1. 244 ปิฏฺเฐ.)}. “เอวํ ภนฺเต”ติ โข อายสฺมา ทพฺโพ มลฺลปุตฺโต ภควโต ปจฺจสฺโสสิ.\csromanspacer{} อถ โข ภควา เอตสฺมิํ นิทาเน เอตสฺมิํ ปกรเณ ธมฺมิํ กถํ กตฺวา ภิกฺขู อามนฺเตสิ\csromandash{}เตน หิ ภิกฺขเว สํโฆ ทพฺพํ มลฺลปุตฺตํ เสนาสนปญฺญาปกญฺจ ภตฺตุทฺเทสกญฺจ สมฺมนฺนตุ.\csromanspacer{} เอวญฺจ ปน ภิกฺขเว สมฺมนฺนิตพฺโพ, ปฐมํ ทพฺโพ มลฺลปุตฺโต ยาจิตพฺโพ\footnote{ทพฺโพ ยจิตพฺโพ (สฺยา, ก)}, ยาจิตฺวา พฺยตฺเตน ภิกฺขุนา ปฏิพเลน สํโฆ ญาเปตพฺโพ\csromandash{}}

\proseitem{190}{“สุณาตุ เม ภนฺเต สํโฆ, ยทิ สํฆสฺส ปตฺตกลฺลํ, สํโฆ อายสฺมนฺตํ ทพฺพํ มลฺลปุตฺตํ เสนาสนปญฺญาปกญฺจ ภตฺตุทฺเทสกญฺจ สมฺมนฺเนยฺย, เอสา ญตฺติ.}

\prose{สุณาตุ เม ภนฺเต สํโฆ, สํโฆ อายสฺมนฺตํ ทพฺพํ มลฺลปุตฺตํ เสนาสนปญฺญาปกญฺจ ภตฺตุทฺเทสกญฺจ สมฺมนฺนติ, ยสฺสายสฺมโต ขมติ อายสฺมโต ทพฺพสฺส มลฺลปุตฺตสฺส เสนาสนปญฺญาปกสฺส จ ภุตฺตุทฺเทสกสฺส จ สมฺมุติ, โส ตุณฺหสฺส, ยสฺส นกฺขมติ, โส ภาเสยฺย.}

\prose{สมฺมโต สํเฆน อายสฺมา ทพฺโพ มลฺลปุตฺโต เสนาสนปญฺญาปโก จ ภตฺตุทฺเทสโก จ ขมติ สํฆสฺส, ตสฺมา ตุณฺหี, เอวเมตํ ธารยามี”ติ.}

\proseitem{191}{\csromanfolio{193}สมฺมโต จ ปนายสฺมา ทพฺโพ มลฺลปุตฺโต สภาคานํ ภิกฺขูนํ เอกชฺฌํ เสนาสนํ ปญฺญเปติ, เย เต ภิกฺขู สุตฺตนฺติกา, เตสํ เอกชฺฌํ เสนาสนํ ปญฺญเปติ “เต อญฺญมญฺญํ สุตฺตนฺตํ สงฺคายิสฺสนฺตี”ติ.\csromanspacer{} เย เต ภิกฺขู วินยธรา, เตสํ เอกชฺฌํ เสนาสนํ ปญฺญเปติ “เต อญฺญมญฺญํ วินยํ วินิจฺฉินิสฺสนฺตี’ติ.\csromanspacer{} เย เต ภิกฺขู ธมฺมกถิกา, เตสํ เอกชฺฌํ เสนาสนํ ปญฺญเปติ “เต อญฺญมญฺญํ ธมฺมํ สากจฺฉิสฺสนฺตี”ติ.\csromanspacer{} เย เต ภิกฺขู ฌายิโน, เตสํ เอกชฺฌํ เสนาสนํ ปญฺญเปติ “เต อญฺญมญฺญํ น พฺยาพาธิสฺสนฺตี”ติ.\csromanspacer{} เย เต ภิกฺขู ติรจฺฉานกถิกา กายทฬฺหิพหุลา\footnote{กายทฑฺฒิพหุลา (สี)} วิหรนฺติ, เตสมฺปิ เอกชฺฌํ เสนาสนํ ปญฺญเปติ “อิมายปิเม อายสฺมนฺโต รติยา อจฺฉิสฺสนฺตี”ติ.\csromanspacer{} เยปิ เต ภิกฺขู วิกาเล อาคจฺฉนฺติ, เตสมฺปิ เตโชธาตุํ สมาปชฺชิตฺวา เตเนว อาโลเกน เสนาสนํ ปญฺญเปติ, อปิสุ ภิกฺขู สญฺจิจฺจ วิกาเล อาคจฺฉนฺติ “มยํ อายสฺมโต ทพฺพสฺส มลฺลปุตฺตสฺส อิทฺธิปาฏิหาริยํ ปสฺสิสฺสามา”ติ.}

\prose{เต อายสฺมนฺตํ ทพฺพํ มลฺลปุตฺตํ อุปสงฺกมิตฺวา เอวํ วทนฺติ “อมฺหากํ อาวุโส ทพฺพ เสนาสนํ ปญฺญเปหี”ติ.\csromanspacer{} เต อายสฺมา ทพฺโพ มลฺลปุตฺโต เอวํ วเทติ “กตฺถ อายสฺมนฺตา อิจฺฉนฺติ, กตฺถ ปญฺญเปมี”ติ.\csromanspacer{} เต สญฺจิจฺจ ทูเร อปทิสนฺติ “อมฺหากํ อาวุโส ทพฺพ คิชฺฌกูเฏ ปพฺพเต เสนาสนํ ปญฺญเปหิ, อมฺหากํ อาวุโส โจรปปาเต เสนาสนํ ปญฺญเปหิ, อมฺหากํ อาวุโส อิสิคิลิปสฺเส กาฬสิลายํ เสนาสนํ ปญฺญเปหิ, อมฺหากํ อาวุโส เวภารปสฺเส สตฺตปณฺณิคุหายํ เสนาสนํ ปญฺญเปหิ, อมฺหากํ อาวุโส สีตวเน สปฺปโสณฺฑิกปพฺภาเร เสนาสนํ ปญฺญเปหิ, อมฺหากํ อาวุโส โคตมกกนฺทรายํ\footnote{โคมฏกนฺทรายํ (สฺยา, กํ)} เสนาสนํ ปญฺญเปหิ, อมฺหากํ อาวุโส ตินฺทุกกนฺทรายํ เสนาสนํ ปญฺญเปหิ, อมฺหากํ อาวุโส ตโปทกนฺทรายํ\footnote{กโปตกนฺทรายํ (ก)} เสนาสนํ ปญฺญเปหิ, อมฺหากํ อาวุโส ตโปทาราเม เสนาสนํ ปญฺญเปหิ, อมฺหากํ อาวุโส ชีวกมฺพวเน เสนาสนํ ปญฺญเปหิ, อมฺหากํ อาวุโส มทฺทกุจฺฉิมฺหิ มิคทาเย เสนาสนํ ปญฺญเปหี”ติ.}

\prose{\csromanfolio{194}เตสํ อายสฺมา ทพฺโพ มลฺลปุตฺโต เตโชธาตุํ สมาปชฺชิตฺวา องฺคุลิยา ชลมานาย ปุรโต ปุรโต คจฺฉติ.\csromanspacer{} เตปิ เตเนว อาโลเกน อายสฺมโต ทพฺพสฺส มลฺลปุตฺตสฺส ปิฏฺฐิโต ปิฏฺฐิโต คจฺฉนฺติ.\csromanspacer{} เตสํ อายสฺมา ทพฺโพ มลฺลปุตฺโต เอวํ เสนาสนํ ปญฺญเปติ “อยํ มญฺโจ, อิทํ ปีฐํ, อยํ ภิสิ, อิทํ พิพฺโพหนํ\footnote{พิมฺโพหนํ (สี, สฺยา, กํ)}, อิทํ วจฺจฏฺฐานํ, อิทํ ปสฺสาวฏฺฐานํ, อิทํ ปานียํ, อิทํ ปริโภชนียํ, อยํ กตฺตรทณฺโฑ, อิทํ สํฆสฺส กติกสณฺฐานํ, อิมํ กาลํ ปวิสิตพฺพํ, อิมํ กาลํ นิกฺขมิตพฺพนฺ”ติ.\csromanspacer{} เตสํ อายสฺมา ทพฺโพ มลฺลปุตฺโต เอวํ เสนาสนํ ปญฺญเปตฺวา ปุนเทว เวฬุวนํ ปจฺจาคจฺฉติ.}

\proseitem{192}{เตน โข ปน สมเยน เมตฺติยภูมชกา\footnote{เมตฺติยภุมฺมชกา (สี, สฺยา, กํ)} ภิกฺขู นวกา เจว โหนฺติ อปฺปปุญฺญา จ, ยานิ สํฆสฺส ลามกานิ เสนาสนานิ, ตานิ เตสํ ปาปุณนฺติ ลามกานิ จ ภตฺตานิ.\csromanspacer{} เตน โข ปน สมเยน ราชคเห มนุสฺสา อิจฺฉนฺติ เถรานํ ภิกฺขูนํ อภิสงฺขาริกํ ปิณฺฑปาตํ ทาตุํ สปฺปิมฺปิ เตลมฺปิ อุตฺตริภงฺคมฺปิ.\csromanspacer{} เมตฺติยภูมชกานํ ปน ภิกฺขูนํ ปากติกํ เทนฺติ ยถารนฺธํ\footnote{ยถารทฺธํ (สฺยา)} กณาชกํ พิลงฺคทุติยํ, เต ปจฺฉาภตฺตํ ปิณฺฑปาตปฺปฏิกฺกนฺตา เถเร ภิกฺขู ปุจฺฉนฺติ “ตุมฺหากํ อาวุโส ภตฺตคฺเค กิํ อโหสิ, ตุมฺหากํ กิํ อโหสี”ติ\footnote{กิํ นาโหสิ (สฺยา, กํ)}.\csromanspacer{} เอกจฺเจ เถรา เอวํ วทนฺติ “อมฺหากํ อาวุโส สปฺปิ อโหสิ, เตลํ อโหสิ, อุตฺตริภงฺคํ อโหสี”ติ.\csromanspacer{} เมตฺติยภูมชกา ปน ภิกฺขู เอวํ วทนฺติ “อมฺหากํ อาวุโส น กิญฺจิ อโหสิ, ปากติกํ ยถารนฺธํ กณาชกํ พิลงฺคทุติยนฺ’ติ.}

\prose{เตน โข ปน สมเยน กลฺยาณภตฺติโก คหปติ สํฆสฺส จตุกฺกภตฺตํ เทติ นิจฺจภตฺตํ, โส ภตฺตคฺเค สปุตฺตทาโร อุปติฏฺฐิตฺวา ปริวิสติ, อญฺเญ โอทเนน ปุจฺฉนฺติ, อญฺเญ สูเปน ปุจฺฉนฺติ, อญฺเญ เตเลน ปุจฺฉนฺติ, อญฺเญ อุตฺตริภงฺเคน ปุจฺฉนฺติ.\csromanspacer{} เตน โข ปน สมเยน กลฺยาณภตฺติกสฺส คหปติโน ภตฺตํ สฺวาตนาย เมตฺติยภูมชกานํ ภิกฺขูนํ อุทฺทิฏฺฐํ โหติ.\csromanspacer{} อถ โข กลฺยาณภตฺติโก คหปติ อารามํ อคมาสิ เกนจิเทว กรณีเยน.\csromanspacer{} โส เยนยสฺมา ทพฺโพ มลฺลปุตฺโต เตนุปสงฺกมิ, อุปสงฺกมิตฺวา อายสฺมนฺตํ ทพฺพํ มลฺลปุตฺตํ อภิวาเทตฺวา \csromanfolio{195}เอกมนฺตํ นิสีทิ, เอกมนฺตํ นิสินฺนํ โข กลฺยาณภตฺติกํ คหปติํ อายสฺมา ทพฺโพ มลฺลปุตฺโต ธมฺมิยา กถาย สนฺทสฺเสสิ สมาทเปสิ สมุตฺเตเชสิ สมฺปหํเสสิ.\csromanspacer{} อถ โข กลฺยาณภตฺติโก คหปติ อายสฺมตา ทพฺเพน มลฺลปุตฺเตน ธมฺมิยา กถาย สนฺทสฺสิโต สมาทปิโต สมุตฺเตชิโต สมฺปหํสิโต อายสฺมนฺตํ ทพฺพํ มลฺลปุตฺตํ เอตทโวจ “กสฺส ภนฺเต อมฺหากํ ฆเร สฺวาตนาย ภตฺตํ อุทฺทิฏฺฐนฺ”ติ.\csromanspacer{} เมตฺติยภูมชกานํ โข คหปติ ภิกฺขูนํ ตุมฺหากํ ฆเร สฺวาตนาย ภตฺตํ อุทฺทิฏฺฐนฺติ.\csromanspacer{} อถ โข กลฺยาณภตฺติโก คหปติ อนตฺตมโน อโหสิ “กถํ หิ นาม ปาปภิกฺขู อมฺหากํ ฆเร ภุญฺชิสฺสนฺตี”ติ ฆรํ คนฺตฺวา ทาสิํ อาณาเปสิ “เย เช เสฺว ภตฺติกา อาคจฺฉนฺติ, เต โกฏฺฐเก อาสนํ ปญฺญเปตฺวา กณาชเกน พิลงฺคทุติเยน ปริวิสา”ติ. “เอวํ อยฺยา”ติ โข สา ทาสี กลฺยาณภตฺติกสฺส คหปติโน ปจฺจสฺโสสิ.}

\prose{อถ โข เมตฺติยภูมชกา ภิกฺขู “หิยฺโย โข อาวุโส อมฺหากํ กลฺยาณภตฺติกสฺส คหปติโน ภตฺตํ อุทฺทิฏฺฐํ, เสฺว อเมฺห กลฺยาณภตฺติโก คหปติ สปุตฺตทาโร อุปติฏฺฐิตฺวา ปริวิสิสฺสติ, อญฺเญ โอทเนน ปุจฺฉิสฺสนฺติ, อญฺเญ สูเปน ปุจฺฉิสฺสนฺติ, อญฺเญ เตเลน ปุจฺฉิสฺสนฺติ, อญฺเญ อุตฺตริภงฺเคน ปุจฺฉิสฺสนฺตี”ติ เต เตเนว โสมนสฺเสน น จิตฺตรูปํ รตฺติยา สุปิํสุ.\csromanspacer{} อถ โข เมตฺติยภูมชกา ภิกฺขู ปุพฺพณฺหสมยํ นิวาเสตฺวา ปตฺตจีวรมาทาย เยน กลฺยาณภตฺติกสฺส คหปติโน นิเวสนํ เตนุปสงฺกมิํสุ.\csromanspacer{} อทฺทสา โข สา ทาสี เมตฺติยภูมชเก ภิกฺขู ทูรโตว อาคจฺฉนฺเต, ทิสฺวาน โกฏฺฐเก อาสนํ ปญฺญเปตฺวา เมตฺติยภูมชเก ภิกฺขู เอตทโวจ “นิสีทถ ภนฺเต”ติ.\csromanspacer{} อถ โข เมตฺติยภูมชกานํ ภิกฺขูนํ เอตทโหสิ “นิสฺสํสยํ โข น ตาว ภตฺตํ สิทฺธํ ภวิสฺสติ, ยถา มยํ โกฏฺฐเก นิสีทาปิยามา”ติ\footnote{นิสีทาเปยฺยามาติ (ก)}.\csromanspacer{} อถ โข สา ทาสี กณาชเกน\footnote{กาณาชเกน (สฺยา, กํ)} พิลงฺคทุติเยน อุปคญฺฉิ “ภุญฺชถ ภนฺเต”ติ.\csromanspacer{} มยํ โข ภคินิ นิจฺจภตฺติกาติ.\csromanspacer{} ชานามิ อยฺยา “นิจฺจภตฺติกา”ติ, อปิ จาหํ หิยฺโยว คหปตินา อาณตฺตา “เย เช เสฺว ภตฺติกา อาคจฺฉนฺติ, เต โกฏฺฐเก อาสนํ ปญฺญเปตฺวา กณาชเกน พิลงฺคทุติเยน ปริวิสา”ติ, ภุญฺชถ ภนฺเตติ.\csromanspacer{} อถ โข เมตฺติยภูมชกา ภิกฺขู “หิยฺโย โข อาวุโส กลฺยาณภตฺติโก คหปติ อารามํ อคมาสิ ทพฺพสฺส \csromanfolio{196}มลฺลปุตฺตสฺส สนฺติเก, นิสฺสํสยํ โข มยํ ทพฺเพน มลฺลปุตฺเตน คหปติโน อนฺตเร ปริภินฺนา”ติ\footnote{สนฺติเก ปริภินฺนาติ (สฺยา, กํ)} เต เตเนว โทมนสฺเสน น จิตฺตรูปํ ภุญฺชิํสุ.\csromanspacer{} อถ โข เมตฺติยภูมชกา ภิกฺขู ปจฺฉาภตฺตํ ปิณฺฑปาตปฺปฏิกฺกนฺตา อารามํ คนฺตฺวา ปตฺตจีวรํ ปฏิสาเมตฺวา พหารามโกฏฺฐเก สํฆาฏิปลฺลตฺถิกาย นิสีทิํสุ ตุณฺหีภูตา มงฺกุภูตา ปตฺตกฺขนฺธา อโธมุขา ปชฺฌายนฺตา อปฺปฏิภานา.}

\prose{อถ โข เมตฺติยา ภิกฺขุนี เยน เมตฺติยภูมชกา ภิกฺขู เตนุปสงฺกมิ, อุปสงฺกมิตฺวา เมตฺติยภูมชเก ภิกฺขู เอตทโวจ “วนฺทามิ อยฺยา”ติ.\csromanspacer{} เอวํ วุตฺเต เมตฺติยภูมชกา ภิกฺขู นาลปิํสุ.\csromanspacer{} ทุติยมฺปิ โข ฯเปฯ.\csromanspacer{} ตติยมฺปิ โข เมตฺติยา ภิกฺขุนี เมตฺติยภูมชเก ภิกฺขู เอตทโวจ “วนฺทามิ อยฺยา”ติ.\csromanspacer{} ตติยมฺปิ โข เมตฺติยภูมชกา ภิกฺขู นาลปิํสุ.\csromanspacer{} กฺยาหํ อยฺยานํ อปรชฺฌามิ, กิสฺส มํ อยฺยา นาลปนฺตีติ.\csromanspacer{} ตถา หิ ปน ตฺวํ ภคินิ อเมฺห ทพฺเพน มลฺลปุตฺเตน วิเหฐิยมาเน อชฺฌุเปกฺขสีติ.\csromanspacer{} กฺยาหํ อยฺยา กโรมีติ.\csromanspacer{} สเจ โข ตฺวํ ภคินิ อิจฺเฉยฺยาสิ, อชฺเชว ภควา อายสฺมนฺตํ ทพฺพํ มลฺลปุตฺตํ นาสาเปยฺยาติ.\csromanspacer{} กฺยาหํ อยฺยา กโรมิ, กิํ มยา สกฺกา กาตุนฺติ.\csromanspacer{} เอหิ ตฺวํ ภคินิ, เยน ภควา เตนุปสงฺกม, อุปสงฺกมิตฺวา ภควนฺตํ เอวํ วเทหิ “อิทํ ภนฺเต นจฺฉนฺนํ นปฺปติรูปํ, ยายํ ภนฺเต ทิสา อภยา อนีติกา อนุปทฺทวา, สายํ ทิสา สภยา ส-อีติกา ส-อุปทฺทวา, ยโต นิวาตํ ตโต สวาตํ\footnote{ตโต ปวาตํ (สี, สฺยา, กํ)}, อุทกํ มญฺเญ อาทิตฺตํ, อยฺเยนมฺหิ ทพฺเพน มลฺลปุตฺเตน ทูสิตา”ติ. “เอวํ อยฺยา”ติ โข เมตฺติยา ภิกฺขุนี เมตฺติยภูมชกานํ ภิกฺขูนํ ปฏิสฺสุตฺวา\footnote{ปฏิสฺสุณิตฺวา (สฺยา, กํ)} เยน ภควา เตนุปสงฺกมิ, อุปสงฺกมิตฺวา ภควนฺตํ อภิวาเทตฺวา เอกมนฺตํ อฏฺฐาสิ, เอกมนฺตํ ฐิตา โข เมตฺติยา\footnote{สา เมตฺติยา (สฺยา, ก)} ภิกฺขุนี ภควนฺตํ เอตทโวจ “อิทํ ภนฺเต นจฺฉนฺนํ นปฺปติรูปํ, ยายํ ภนฺเต ทิสา อภยา อนีติกา อนุปทฺทวา, สายํ ทิสา สภยา ส-อีติกา ส-อุปทฺทวา, ยโต นิวาตํ ตโต สวาตํ, อุทกํ มญฺเญ อาทิตฺตํ, อยฺเยนมฺหิ ทพฺเพน มลฺลปุตฺเตน ทูสิตา”ติ.}

\proseitem{193}{อถ โข ภควา เอตสฺมิํ นิทาเน เอตสฺมิํ ปกรเณ ภิกฺขุสํฆํ สนฺนิปาตาเปตฺวา อายสฺมนฺตํ ทพฺพํ มลฺลปุตฺตํ ปฏิปุจฺฉิ “สรสิ ตฺวํ ทพฺพ เอวรูปํ \csromanfolio{197}กตฺตา, ยถายํ ภิกฺขุนี อาหา”ติ.\csromanspacer{} ยถา มํ ภนฺเต ภควา ชานาตีติ.\csromanspacer{} ทุติยมฺปิ โข ภควา อายสฺมนฺตํ ทพฺพํ มลฺลปุตฺตํ เอตทโวจ “สรสิ ตฺวํ ทพฺพ เอวรูปํ กตฺตา, ยถายํ ภิกฺขุนี อาหา”ติ.\csromanspacer{} ยถา มํ ภนฺเต ภควา ชานาตีติ.\csromanspacer{} ตติยมฺปิ โข ภควา อายสฺมนฺตํ ทพฺพํ มลฺลปุตฺตํ เอตทโวจ “สรสิ ตฺวํ ทพฺพ เอวรูปํ กตฺตา, ยถายํ ภิกฺขุนี อาหา”ติ.\csromanspacer{} ยถา มํ ภนฺเต ภควา ชานาตีติ.\csromanspacer{} น โข ทพฺพ ทพฺพา เอวํ นิพฺเพเฐนฺติ, สเจ ตยา กตํ ‘กตนฺ’ติ วเทหิ, สเจ อกตํ ‘อกตนฺ’ติ วเทหีติ.\csromanspacer{} ยโตหํ ภนฺเต ชาโต, นาภิชานามิ สุปินนฺเตนปิ เมถุนํ ธมฺมํ ปฏิเสวิตา, ปเคว ชาคโรติ.\csromanspacer{} อถ โข ภควา ภิกฺขู อามนฺเตสิ “เตน หิ ภิกฺขเว เมตฺติยํ ภิกฺขุนิํ นาเสถ, อิเม จ ภิกฺขู อนุยุญฺชถา”ติ.\csromanspacer{} อิทํ วตฺวา ภควา อุฏฺฐายาสนา วิหารํ ปาวิสิ.}

\prose{อถ โข เต ภิกฺขู เมตฺติยํ ภิกฺขุนิํ นาเสสุํ.\csromanspacer{} อถ โข เมตฺติยภูมชกา ภิกฺขู เต ภิกฺขู เอตทโวจุํ “มาวุโส เมตฺติยํ ภิกฺขุนิํ นาเสถ, น สา กิญฺจิ อปรชฺฌติ, อเมฺหหิ สา อุสฺสาหิตา กุปิเตหิ อนตฺตมเนหิ จาวนาธิปฺปาเยหี”ติ.\csromanspacer{} กิํ ปน ตุเมฺห อาวุโส อายสฺมนฺตํ ทพฺพํ มลฺลปุตฺตํ อมูลิกาย สีลวิปตฺติยา อนุทฺธํเสถาติ.\csromanspacer{} เอวมาวุโสติ.\csromanspacer{} เย เต ภิกฺขู อปฺปิจฺฉา ฯเปฯ เต อุชฺฌายนฺติ ขิยฺยนฺติ วิปาเจนฺติ “กถํ หิ นาม เมตฺติยภูมชกา ภิกฺขู อายสฺมนฺตํ ทพฺพํ มลฺลปุตฺตํ อมูลิกาย สีลวิปตฺติยา อนุทฺธํเสสฺสนฺตี”ติ.\csromanspacer{} อถ โข เต ภิกฺขู ภควโต เอตมตฺถํ อาโรเจสุํ ฯเปฯ.\csromanspacer{} สจฺจํ กิร ภิกฺขเว เมตฺติยภูมชกา ภิกฺขู ทพฺพํ มลฺลปุตฺตํ อมูลิกาย สีลวิปตฺติยา อนุทฺธํเสนฺตีติ.\csromanspacer{} สจฺจํ ภควาติ ฯเปฯ วิครหิตฺวา ฯเปฯ ธมฺมิํ กถํ กตฺวา ภิกฺขู อามนฺเตสิ\csromandash{}}

\prose{เตน หิ ภิกฺขเว สํโฆ ทพฺพสฺส มลฺลปุตฺตสฺส สติเวปุลฺลปฺปตฺตสฺส สติวินยํ เทตุ.\csromanspacer{} เอวญฺจ ปน ภิกฺขเว ทาตพฺโพ, เตน ภิกฺขเว ทพฺเพน มลฺลปุตฺเตน สํฆํ อุปสงฺกมิตฺวา เอกํสํ อุตฺตราสงฺคํ กริตฺวา วุฑฺฒานํ ภิกฺขูนํ ปาเท วนฺทิตฺวา อุกฺกุฏิกํ นิสีทิตฺวา อญฺชลิํ ปคฺคเหตฺวา เอวมสฺส วจนีโย “อิเม มํ ภนฺเต เมตฺติยภูมชกา ภิกฺขู อมูลิกาย สีลวิปตฺติยา อนุทฺธํเสนฺติ, โสหํ ภนฺเต สติเวปุลฺลปฺปตฺโต สํฆํ สติวินยํ ยาจามี”ติ.\csromanspacer{} ทุติยมฺปิ ยาจิตพฺโพ.\csromanspacer{} ตติยมฺปิ ยาจิตพฺโพ “อิเม มํ ภนฺเต เมตฺติยภูมชกา ภิกฺขู อมูลิกาย สีลวิปตฺติยา อนุทฺธํเสนฺติ, \csromanfolio{198}โสหํ\footnote{โสหํ ภนฺเต (ก)} สติเวปุลฺลปฺปตฺโต ตติยมฺปิ ภนฺเต สํฆํ สติวินยํ ยาจามี”ติ.\csromanspacer{} พฺยตฺเตน ภิกฺขุนา ปฏิพเลน สํโฆ ญาเปตพฺโพ\csromandash{}}

\proseitem{194}{“สุณาตุ เม ภนฺเต สํโฆ, อิเม เมตฺติยภูมชกา ภิกฺขู อายสฺมนฺตํ ทพฺพํ มลฺลปุตฺตํ อมูลิกาย สีลวิปตฺติยา อนุทฺธํเสนฺติ, อายสฺมา ทพฺโพ มลฺลปุตฺโต สติเวปุลฺลปฺปตฺโต สํฆํ สติวินยํ ยาจติ, ยทิ สํฆสฺส ปตฺตกลฺลํ, สํโฆ อายสฺมโต ทพฺพสฺส มลฺลปุตฺตสฺส สติเวปุลฺลปฺปตฺตสฺส สติวินยํ ทเทยฺย, เอสา ญตฺติ.}

\prose{สุณาตุ เม ภนฺเต สํโฆ, อิเม เมตฺติยภูมชกา ภิกฺขู อายสฺมนฺตํ ทพฺพํ มลฺลปุตฺตํ อมูลิกาย สีลวิปตฺติยา อนุทฺธํเสนฺติ, อายสฺมา ทพฺโพ มลฺลปุตฺโต สติเวปุลฺลปฺปตฺโต สํฆํ สติวินยํ ยาจติ, สํโฆ อายสฺมโต ทพฺพสฺส มลฺลปุตฺตสฺส สติเวปุลฺลปฺปตฺตสฺส สติวินยํ เทติ, ยสฺสายสฺมโต ขมติ อายสฺมโต ทพฺพสฺส มลฺลปุตฺตสฺส สติเวปุลฺลปฺปตฺตสฺส สติวินยสฺส ทานํ, โส ตุณฺหสฺส, ยสฺส นกฺขมติ, โส ภาเสยฺย.}

\prose{ทุติยมฺปิ เอตมตฺถํ วทามิ ฯเปฯ.\csromanspacer{} ตติยมฺปิ เอตมตฺถํ วทามิ.\csromanspacer{} สุณาตุ เม ภนฺเต สํโฆ อิเม เมตฺติยภูมชกา ภิกฺขู อายสฺมนฺตํ ทพฺพํ มลฺลปุตฺตํ อมูลิกาย สีลวิปตฺติยา อนุทฺธํเสนฺติ, อายสฺมา ทพฺโพ มลฺลปุตฺโต สติเวปุลฺลปฺปตฺโต สํฆํ สติวินยํ ยาจติ, สํโฆ อายสฺมโต ทพฺพสฺส มลฺลปุตฺตสฺส สติเวปุลฺลปฺปตฺตสฺส สติวินยํ เทติ, ยสฺสายสฺมโต ขมติ อายสฺมโต ทพฺพสฺส มลฺลปุตฺตสฺส สติเวปุลฺลปฺปตฺตสฺส สติวินยสฺส ทานํ, โส ตุณฺหสฺส, ยสฺส นกฺขมติ, โส ภาเสยฺย.}

\prose{ทินฺโน สํเฆน อายสฺมโต ทพฺพสฺส มลฺลปุตฺตสฺส สติเวปุลฺลปฺปตฺตสฺส สติวินโย, ขมติ สํฆสฺส, ตสฺมา ตุณฺหี, เอวเมตํ ธารยามี”ติ.}

\proseitem{195}{ปญฺจิมานิ ภิกฺขเว ธมฺมิกานิ สติวินยสฺส ทานานิ.\csromanspacer{} สุทฺโธ โหติ ภิกฺขุ อนาปตฺติโก, อนุวทนฺติ จ นํ, ยาจติ จ, ตสฺส สํโฆ สติวินยํ เทติ ธมฺเมน สมคฺเคน.\csromanspacer{} อิมานิ โข ภิกฺขเว ปญฺจ ธมฺมิกานิ สติวินยสฺส ทานานีติ.}

\csromanmatikaanchor{mtk.106}
\csromantocmarkref{section}{3. อมูฬฺหวินย}{mtk.106}
\bookmark[dest={mtk.106},level=1]{3. อมูฬฺหวินย}
\csromanheader{1}{3. อมูฬฺหวินย}
\proseitem{196}{\csromanfolio{199}เตน โข ปน สมเยน คคฺโค ภิกฺขุ อุมฺมตฺตโก โหติ จิตฺตวิปริยาสกโต, เตน อุมฺมตฺตเกน จิตฺตวิปริยาสกเตน พหุํ อสฺสามณกํ อชฺฌาจิณฺณํ โหติ ภาสิตปริกฺกนฺตํ.\csromanspacer{} ภิกฺขู คคฺคํ ภิกฺขุํ อุมฺมตฺตเกน จิตฺตวิปริยาสกเตน อชฺฌาจิณฺเณน อาปตฺติยา โจเทนฺติ “สรตายสฺมา เอวรูปิํ อาปตฺติํ อาปชฺชิตา”ติ.\csromanspacer{} โส เอวํ วเทติ “อหํ โข อาวุโส อุมฺมตฺตโก อโหสิํ จิตฺตวิปริยาสกโต, เตน เม อุมฺมตฺตเกน จิตฺตวิปริยาสกเตน พหุํ อสฺสามณกํ อชฺฌาจิณฺณํ ภาสิตปริกฺกนฺตํ, นาหํ ตํ สรามิ, มูเฬฺหน เม เอตํ กตนฺ”ติ.\csromanspacer{} เอวมฺปิ นํ วุจฺจมานา โจเทนฺเตว “สรตายสฺมา เอวรูปิํ อาปตฺติํ อาปชฺชิตา”ติ.\csromanspacer{} เย เต ภิกฺขู อปฺปิจฺฉา ฯเปฯ เต อุชฺฌายนฺติ ขิยฺยนฺติ วิปาเจนฺติ “กถํ หิ นาม ภิกฺขู คคฺคํ ภิกฺขุํ อุมฺมตฺตเกน จิตฺตวิปริยาสกเตน อชฺฌาจิณฺเณน อาปตฺติยา โจเทสฺสนฺติ ‘สรตายสฺมา เอวรูปิํ อาปตฺติํ อาปชฺชิตา’ติ.\csromanspacer{} โส เอวํ วเทติ ‘อหํ โข อาวุโส อุมฺมตฺตโก อโหสิํ จิตฺตวิปริยาสกโต, เตน เม อุมฺมตฺตเกน จิตฺตวิปริยาสกเตน พหุํ อสฺสามณกํ อชฺฌาจิณฺณํ ภาสิตปริกฺกนฺตํ, นาหํ ตํ สรามิ, มูเฬฺหน เม เอตํ กตนฺ’ติ.\csromanspacer{} เอวมฺปิ นํ วุจฺจมานา โจเทนฺเตว สรตายสฺมา เอวรูปิํ อาปตฺติํ อาปชฺชิตา”ติ.\csromanspacer{} อถ โข เต ภิกฺขู ภควโต เอตมตฺถํ อาโรเจสุํ ฯเปฯ.\csromanspacer{} สจฺจํ กิร ภิกฺขเว ฯเปฯ.\csromanspacer{} สจฺจํ ภควาติ ฯเปฯ วิครหิตฺวา ฯเปฯ ธมฺมิํ กถํ กตฺวา ภิกฺขู อามนฺเตสิ\csromandash{}}

\prose{เตน หิ ภิกฺขเว สํโฆ คคฺคสฺส ภิกฺขุโน อมูฬฺหสฺส อมูฬฺหวินยํ เทตุ.\csromanspacer{} เอวญฺจ ปน ภิกฺขเว ทาตพฺโพ, เตน ภิกฺขเว วคฺเคน ภิกฺขุนา สํฆํ อุปสงฺกมิตฺวา เอกํสํ อุตฺตราสงฺคํ กริตฺวา วุฑฺฒานํ ภิกฺขูนํ ปาเท วนฺทิตฺวา อุกฺกุฏิกํ นิสีทิตฺวา อญฺชลิํ ปคฺคเหตฺวา เอวมสฺส วจนีโย “อหํ ภนฺเต อุมฺมตฺตโก อโหสิํ จิตฺตวิปริยาสกโต, เตน เม อุมฺมตฺตเกน จิตฺตวิปริยาสกเตน พหุํ อสฺสามณกํ อชฺฌาจิณฺณํ ภาสิตปริกฺกนฺตํ, มํ ภิกฺขู อุมฺมตฺตเกน จิตฺตวิปริยาสกเตน อชฺฌาจิณฺเณน อาปตฺติยา โจเทนฺติ ‘สรตายสฺมา เอวรูปิํ อาปตฺติํ อาปชฺชิตา’ติ, ตฺยาหํ เอวํ วทามิ ‘อหํ โข อาวุโส อุมฺมตฺตโก อโหสิํ จิตฺตวิปริยาสกโต, เตน เม อุมฺมตฺตเกน จิตฺตวิปริยาสกเตน พหุํ อสฺสามณกํ อชฺฌาจิณฺณํ ภาสิตปริกฺกนฺตํ, นาหํ ตํ สรามิ, มูเฬฺหน เม เอตํ กตนฺ’ติ, เอวมฺปิ มํ วุจฺจมานา \csromanfolio{200}โจเทนฺเตว ‘สรตายสฺมา เอวรูปิํ อาปตฺติํ อาปชฺชิตา’ติ, โสหํ ภนฺเต อมูโฬฺห สํฆํ อมูฬฺหวินยํ ยาจามี”ติ.\csromanspacer{} ทุติยมฺปิ ยาจิตพฺโพ.\csromanspacer{} ตติยมฺปิ ยาจิตพฺโพ “อหํ ภนฺเต อุมฺมตฺตโก อโหสิํ จิตฺตวิปริยาสกโต, เตน เม อุมฺมตฺตเกน จิตฺตวิปริยาสกเตน พหุํ อสฺสามณกํ อชฺฌาจิณฺณํ ภาสิตปริกฺกนฺตํ, มํ ภิกฺขู อุมฺมตฺตเกน จิตฺตวิปริยาสกเตน อชฺฌาจิณฺเณน อาปตฺติยา โจเทนฺติ ‘สรตายสฺมา เอวรูปิํ อาปตฺติํ อาปชฺชิตา’ติ, ตฺยาหํ เอวํ วทามิ ‘อหํ โข อาวุโส อุมฺมตฺตโก อโหสิํ จิตฺตวิปริยาสกโต, เตน เม อุมฺมตฺตเกน จิตฺตวิปริยาสกเตน พหุํ อสฺสามณกํ อชฺฌาจิณฺณํ ภาสิตปริกฺกนฺตํ, นาหํ ตํ สรามิ, มูเฬฺหน เม เอตํ กตนฺ’ติ, เอวมฺปิ มํ วุจฺจมานา โจเทนฺเตว ‘สรตายสฺมา เอวรูปิํ อาปตฺติํ อาปชฺชิตา’ติ, โสหํ อมูโฬฺห\footnote{โสหํ ภนฺเต อมูโฬฺห (ก)} ตติยมฺปิ ภนฺเต สํฆํ อมูฬฺหวินยํ ยาจามี”ติ.\csromanspacer{} พฺยตฺเตน ภิกฺขุนา ปฏิพเลน สํโฆ ญาเปตพฺโพ\csromandash{}}

\proseitem{197}{“สุณาตุ เม ภนฺเต สํโฆ, อยํ คคฺโค ภิกฺขุ อุมฺมตฺตโก อโหสิ จิตฺตวิปริยาสกโต, เตน อุมฺมตฺตเกน จิตฺตวิปริยาสกเตน พหุํ อสฺสามณกํ อชฺฌาจิณฺณํ ภาสิตปริกฺกนฺตํ, ภิกฺขู คคฺคํ ภิกฺขุํ อุมฺมตฺตเกน จิตฺตวิปริยาสกเตน อชฺฌาจิณฺเณน อาปตฺติยา โจเทนฺติ ‘สรตายสฺมา เอวรูปิํ อาปตฺติํ อาปชฺชิตา’ติ, โส เอวํ วเทติ ‘อหํ โข อาวุโส อุมฺมตฺตโก อโหสิํ จิตฺตวิปริยาสกโต, เตน เม อุมฺมตฺตเกน จิตฺตวิปริยาสกเตน พหุํ อสฺสามณกํ อชฺฌาจิณฺณํ ภาสิตปริกฺกนฺตํ, นาหํ ตํ สรามิ, มูเฬฺหน เม เอตํ กตนฺ’ติ, เอวมฺปิ นํ วุจฺจมานา โจเทนฺเตว ‘สรตายสฺมา เอวรูปิํ อาปตฺติํ อาปชฺชิตา’ติ, โส อมูโฬฺห สํฆํ อมูฬฺหวินยํ ยาจติ, ยทิ สํฆสฺส ปตฺตกลฺลํ, สํโฆ คคฺคสฺส ภิกฺขุโน อมูฬฺหสฺส อมูฬฺหวินยํ ทเทยฺย, เอสา ญตฺติ.}

\prose{สุณาตุ เม ภนฺเต สํโฆ, อยํ คคฺโค ภิกฺขุ อุมฺมตฺตโก อโหสิ จิตฺตวิปริยาสกโต, เตน อุมฺมตฺตเกน จิตฺตวิปริยาสกเตน พหุํ อสฺสามณกํ อชฺฌาจิณฺณํ ภาสิตปริกฺกนฺตํ, ภิกฺขู คคฺคํ ภิกฺขุํ อุมฺมตฺตเกน จิตฺตวิปริยาสกเตน อชฺฌาจิณฺเณน อาปตฺติยา โจเทนฺติ ‘สรตายสฺมา เอวรูปิํ อาปตฺติํ อาปชฺชิตา’ติ, โส เอวํ วเทติ, ‘อหํ โข อาวุโส อุมฺมตฺตโก อโหสิํ จิตฺตวิปริยาสกโต, เตน เม อุมฺมตฺตเกน จิตฺตวิปริยาสกเตน พหุํ อสฺสามณกํ \csromanfolio{201}อชฺฌาจิณฺณํ ภาสิตปริกฺกนฺตํ, นาหํ ตํ สรามิ, มูเฬฺหน เม เอตํ กตนฺ’ติ, เอวมฺปิ นํ วุจฺจมานา โจเทนฺเตว ‘สรตายสฺมา เอวรูปิํ อาปตฺติํ อาปชฺชิตา’ติ, โส อมูโฬฺห สํฆํ อมูฬฺหวินยํ ยาจติ, สํโฆ คคฺคสฺส ภิกฺขุโน อมูฬฺหสฺส อมูฬฺหวินยํ เทติ, ยสฺสายสฺมโต ขมติ คคฺคสฺส ภิกฺขุโน อมูฬฺหสฺส อมูฬฺหวินยสฺส ทานํ, โส ตุณฺหสฺส, ยสฺส นกฺขมติ, โส ภาเสยฺย.}

\prose{ทุติยมฺปิ เอตมตฺถํ วทามิ ฯเปฯ.\csromanspacer{} ตติยมฺปิ เอตมตฺถํ วทามิ ฯเปฯ.}

\prose{ทินฺโน สํเฆน คคฺคสฺส ภิกฺขุโน อมูฬฺหสฺส อมูฬฺหวินโย, ขมติ สํฆสฺส, ตสฺมา ตุณฺหี, เอวเมตํ ธารยามี”ติ.}

\proseitem{198}{ตีณิมานิ ภิกฺขเว อธมฺมิกานิ อมูฬฺหวินยสฺส ทานานิ, ตีณิ ธมฺมิกานิ.\csromanspacer{} กตมานิ ตีณิ อธมฺมิกานิ อมูฬฺหวินยสฺส ทานานิ.}

\prose{อิธ ปน ภิกฺขเว ภิกฺขุ อาปตฺติํ อาปนฺโน โหติ, ตเมนํ โจเทติ สํโฆ วา สมฺพหุลา วา เอกปุคฺคโล วา “สรตายสฺมา เอวรูปิํ อาปตฺติํ อาปชฺชิตา”ติ, โส สรมาโนว เอวํ วเทติ “น โข อหํ อาวุโส สรามิ เอวรูปิํ อาปตฺติํ อาปชฺชิตา”ติ, ตสฺส สํโฆ อมูฬฺหวินยํ เทติ, อธมฺมิกํ อมูฬฺหวินยสฺส ทานํ.}

\prose{อิธ ปน ภิกฺขเว ภิกฺขุ อาปตฺติํ อาปนฺโน โหติ, ตเมนํ โจเทติ สํโฆ วา สมฺพหุลา วา เอกปุคฺคโล วา “สรตายสฺมา เอวรูปิํ อาปตฺติํ อาปชฺชิตา”ติ, โส สรมาโนว เอวํ วเทติ “สรามิ โข อหํ อาวุโส ยถาสุปินนฺเตนา”ติ, ตสฺส สํโฆ อมูฬฺหวินยํ เทติ, อธมฺมิกํ อมูฬฺหวินยสฺส ทานํ.}

\prose{อิธ ปน ภิกฺขเว ภิกฺขุ อาปตฺติํ อาปนฺโน โหติ, ตเมนํ โจเทติ สํโฆ วา สมฺพหุลา วา เอกปุคฺคโล วา “สรตายสฺมา เอวรูปิํ อาปตฺติํ อาปชฺชิตา”ติ, โส อนุมฺมตฺตโก อุมฺมตฺตกาลยํ กโรติ “อหมฺปิ โข เอวํ กโรมิ, ตุเมฺหปิ เอวํ กโรถ, มยฺหมฺปิ เอตํ กปฺปติ, ตุมฺหากมฺเปตํ กปฺปตี”ติ, ตสฺส สํโฆ อมูฬฺหวินยํ เทติ, อธมฺมิกํ อมูฬฺหวินยสฺส ทานํ.\csromanspacer{} อิมานิ ตีณิ อธมฺมิกานิ อมูฬฺหวินยสฺส ทานานิ.}

\csromanmatikaanchor{mtk.107}
\csromantocmarkref{section}{4. ปฏิญฺญาตกรณ}{mtk.107}
\bookmark[dest={mtk.107},level=1]{4. ปฏิญฺญาตกรณ}
\proseitem{199}{\csromanfolio{202}กตมานิ ตีณิ ธมฺมิกานิ อมูฬฺหวินยสฺส ทานานิ.}

\prose{อิธ ปน ภิกฺขเว ภิกฺขุ อุมฺมตฺตโก โหติ จิตฺตวิปริยาสกโต, เตน อุมฺมตฺตเกน จิตฺตวิปริยาสกเตน พหุํ อสฺสามณกํ อชฺฌาจิณฺณํ โหติ ภาสิตปริกฺกนฺตํ, ตเมนํ โจเทติ สํโฆ วา สมฺพหุลา วา เอกปุคฺคโล วา “สรตายสฺมา เอวรูปิํ อาปตฺติํ อาปชฺชิตา”ติ, โส อสฺสรมาโนว เอวํ วเทติ “น โข อหํ อาวุโส สรามิ เอวรูปิํ อาปตฺติํ อาปชฺชิตา”ติ, ตสฺส สํโฆ อมูฬฺหวินยํ เทติ, ธมฺมิกํ อมูฬฺหวินยสฺส ทานํ.}

\prose{อิธ ปน ภิกฺขเว ภิกฺขุ อุมฺมตฺตโก โหติ จิตฺตวิปริยาสกโต, เตน อุมฺมตฺตเกน จิตฺตวิปริยาสกเตน พหุํ อสฺสามณกํ อชฺฌาจิณฺณํ โหติ ภาสิตปริกฺกนฺตํ, ตเมนํ โจเทติ สํโฆ วา สมฺพหุลา วา เอกปุคฺคโล วา “สรตายสฺมา เอวรูปิํ อาปตฺติํ อาปชฺชิตา”ติ, โส อสฺสรมาโนว เอวํ วเทติ “สรามิ โข อหํ อาวุโส ยถา สุปินนฺเตนา”ติ, ตสฺส สํโฆ อมูฬฺหวินยํ เทติ, ธมฺมิกํ อมูฬฺหวินยสฺส ทานํ.}

\prose{อิธ ปน ภิกฺขเว ภิกฺขุ อุมฺมตฺตโก โหติ จิตฺตวิปริยาสกโต, เตน อุมฺมตฺตเกน จิตฺตวิปริยาสกเตน พหุํ อสฺสามณกํ อชฺฌาจิณฺณํ โหติ ภาสิตปริกฺกนฺตํ, ตเมนํ โจเทติ สํโฆ วา สมฺพหุลา วา เอกปุคฺคโล วา “สรตายสฺมา เอวรูปิํ อาปตฺติํ อาปชฺชิตา”ติ, โส อุมฺมตฺตโก อุมฺมตฺตกาลยํ กโรติ “อหมฺปิ เอวํ กโรมิ, ตุเมฺหปิ เอวํ กโรถ, มยฺหมฺปิ เอตํ กปฺปติ, ตุมฺหากมฺเปตํ กปฺปตี”ติ, ตสฺส สํโฆ อมูฬฺหวินยํ เทติ, ธมฺมิกํ อมูฬฺหวินยสฺส ทานํ.\csromanspacer{} อิมานิ ตีณิ ธมฺมิกานิ อมูฬฺหวินยสฺส ทานานีติ.}

\csromanheader{2}{4. ปฏิญฺญาตกรณา}
\proseitem{200}{เตน โข ปน สมเยน ฉพฺพคฺคิยา ภิกฺขู อปฺปฏิญฺญาย ภิกฺขูนํ กมฺมานิ กโรนฺติ ตชฺชนียมฺปิ นิยสฺสมฺปิ ปพฺพาชนียมฺปิ ปฏิสารณียมฺปิ อุกฺเขปนียมฺปิ.\csromanspacer{} เย เต ภิกฺขู อปฺปิจฺฉา ฯเปฯ เต อุชฺฌายนฺติ ขิยฺยนฺติ วิปาเจนฺติ “กถํ หิ นาม ฉพฺพคฺคิยา ภิกฺขู อปฺปฏิญฺญาย ภิกฺขูนํ กมฺมานิ กริสฺสนฺติ ตชฺชนียมฺปิ นิยสฺสมฺปิ ปพฺพาชนียมฺปิ ปฏิสารณียมฺปิ อุกฺเขปนียมฺปี”ติ.\csromanspacer{} อถ \csromanfolio{203}โข เต ภิกฺขู ภควโต เอตมตฺถํ อาโรเจสุํ ฯเปฯ.\csromanspacer{} สจฺจํ กิร ภิกฺขเว ฯเปฯ.\csromanspacer{} สจฺจํ ภควาติ ฯเปฯ วิครหิตฺวา ฯเปฯ ธมฺมิํ กถํ กตฺวา ภิกฺขู อามนฺเตสิ\csromandash{}น ภิกฺขเว อปฺปฏิญฺญาย ภิกฺขูนํ กมฺมํ กาตพฺพํ ตชฺชนียํ วา นิยสฺสํ วา ปพฺพาชนียํ วา ปฏิสารณียํ วา อุกฺเขปนียํ วา, โย กเรยฺย, อาปตฺติ ทุกฺกฏสฺส.}

\proseitem{201}{เอวํ โข ภิกฺขเว อธมฺมิกํ โหติ ปฏิญฺญาตกรณํ, เอวํ ธมฺมิกํ.\csromanspacer{} กถญฺจ ภิกฺขเว อธมฺมิกํ โหติ ปฏิญฺญาตกรณํ.}

\prose{ภิกฺขุ ปาราชิกํ อชฺฌาปนฺโน โหติ, ตเมนํ โจเทติ สํโฆ วา สมฺพหุลา วา เอกปุคฺคโล วา “ปาราชิกํ อายสฺมา อชฺฌาปนฺโน”ติ.\csromanspacer{} โส เอวํ วเทติ “น โข อหํ อาวุโส ปาราชิกํ อชฺฌาปนฺโน สํฆาทิเสสํ อชฺฌาปนฺโน”ติ.\csromanspacer{} ตํ สํโฆ สํฆาทิเสเสน กาเรติ, อธมฺมิกํ ปฏิญฺญาตกรณํ.}

\prose{ภิกฺขุ ปาราชิกํ อชฺฌาปนฺโน โหติ, ตเมนํ โจเทติ สํโฆ วา สมฺพหุลา วา เอกปุคฺคโล วา “ปาราชิกํ อายสฺมา อชฺฌาปนฺโน”ติ.\csromanspacer{} โส เอวํ วเทติ “น โข อหํ อาวุโส ปาราชิกํ อชฺฌาปนฺโน, ถุลฺลจฺจยํ ฯเปฯ ปาจิตฺติยํ ฯเปฯ ปาฏิเทสนียํ ฯเปฯ ทุกฺกฏํ ฯเปฯ ทุพฺภาสิตํ อชฺฌาปนฺโน”ติ.\csromanspacer{} ตํ สํโฆ ทุพฺภาสิเตน กาเรติ, อธมฺมิกํ ปฏิญฺญาตกรณํ.}

\prose{ภิกฺขุ สํฆาทิเสสํ ฯเปฯ ถุลฺลจฺจยํ ฯเปฯ ปาจิตฺติยํ ฯเปฯ ปาฏิเทสนียํ ฯเปฯ ทุกฺกฏํ ฯเปฯ ทุพฺภาสิตํ อชฺฌาปนฺโน โหติ, ตเมนํ โจเทติ สํโฆ วา สมฺพหุลา วา เอกปุคฺคโล วา “ทุพฺภาสิตํ อายสฺมา อชฺฌาปนฺโน”ติ.\csromanspacer{} โส เอวํ วเทติ “น โข อหํ อาวุโส ทุพฺภาสิตํ อชฺฌาปนฺโน ปาราชิกํ อชฺฌาปนฺโน”ติ.\csromanspacer{} ตํ สํโฆ ปาราชิเกน กาเรติ, อธมฺมิกํ ปฏิญฺญาตกรณํ.}

\prose{ภิกฺขุ ทุพฺภาสิตํ อชฺฌาปนฺโน โหติ, ตเมนํ โจเทติ สํโฆ วา สมฺพหุลา วา เอกปุคฺคโล วา “ทุพฺภาสิตํ อายสฺมา อชฺฌาปนฺโน”ติ.\csromanspacer{} โส เอวํ วเทติ “น โข อหํ อาวุโส ทุพฺภาสิตํ อชฺฌาปนฺโน สํฆาทิเสสํ ฯเปฯ ถุลฺลจฺจยํ ฯเปฯ ปาจิตฺติยํ ฯเปฯ ปาฏิเทสนียํ ฯเปฯ ทุกฺกฏํ อชฺฌาปนฺโน”ติ.\csromanspacer{} ตํ สํโฆ ทุกฺกเฏน กาเรติ, อธมฺมิกํ ปฏิญฺญาตกรณํ.\csromanspacer{} เอวํ โข ภิกฺขเว อธมฺมิกํ โหติ ปฏิญฺญาตกรณํ.}

\prose{\csromanfolio{204}กถญฺจ ภิกฺขเว ธมฺมิกํ โหติ ปฏิญฺญาตกรณํ.\csromanspacer{} ภิกฺขุ ปาราชิกํ อชฺฌาปนฺโน โหติ, ตเมนํ โจเทติ สํโฆ วา สมฺพหุลา วา เอกปุคฺคโล วา “ปาราชิกํ อายสฺมา อชฺฌาปนฺโน”ติ.\csromanspacer{} โส เอวํ วเทติ “อาม อาวุโส ปาราชิกํ อชฺฌาปนฺโน”ติ.\csromanspacer{} ตํ สํโฆ ปาราชิเกน กาเรติ, ธมฺมิกํ ปฏิญฺญาตกรณํ.}

\prose{ภิกฺขุ สํฆาทิเสสํ ฯเปฯ ถุลฺลจฺจยํ ฯเปฯ ปาจิตฺติยํ ฯเปฯ ปาฏิเทสนียํ ฯเปฯ ทุกฺกฏํ ฯเปฯ ทุพฺภาสิตํ อชฺฌาปนฺโน โหติ, ตเมนํ โจเทติ สํโฆ วา สมฺพหุลา วา เอกปุคฺคโล วา “ทุพฺภาสิตํ อายสฺมา อชฺฌาปนฺโน”ติ.\csromanspacer{} โส เอวํ วเทติ “อาม อาวุโส ทุพฺภาสิตํ อชฺฌาปนฺโน”ติ.\csromanspacer{} ตํ สํโฆ ทุพฺภาสิเตน กาเรติ, ธมฺมิกํ ปฏิญฺญาตกรณํ.\csromanspacer{} เอวํ โข ภิกฺขเว ธมฺมิกํ โหติ ปฏิญฺญาตกรณนฺติ.}

\csromanmatikaanchor{mtk.108}
\csromantocmarkref{section}{5. เยภุยฺยสิกา}{mtk.108}
\bookmark[dest={mtk.108},level=1]{5. เยภุยฺยสิกา}
\csromanheader{1}{5. เยภุยฺยสิกา}
\proseitem{202}{เตน โข ปน สมเยน ภิกฺขู สํฆมชฺเฌ ภณฺฑนชาตา กลหชาตา วิวาทาปนฺนา อญฺญมญฺญํ มุขสตฺตีหิ วิตุทนฺตา วิหรนฺติ, น สกฺโกนฺติ ตํ อธิกรณํ วูปสเมตุํ.\csromanspacer{} ภควโต เอตมตฺถํ อาโรเจสุํ.\csromanspacer{} อนุชานามิ ภิกฺขเว เอวรูปํ อธิกรณํ เยภุยฺยสิกาย วูปสเมตุํ.\csromanspacer{} ปญฺจหงฺเคหิ สมนฺนาคโต ภิกฺขุ สลากคฺคาหาปโก สมฺมนฺนิตพฺโพ, โย น ฉนฺทาคติํ คจฺเฉยฺย, น โทสาคติํ คจฺเฉยฺย, น โมหาคติํ คจฺเฉยฺย, น ภยาคติํ คจฺเฉยฺย, คติตาคหิตญฺจ ชาเนยฺย.\csromanspacer{} เอวญฺจ ปน ภิกฺขเว สมฺมนฺนิตพฺโพ, ปฐมํ ภิกฺขุ ยาจิตพฺโพ, ยาจิตฺวา พฺยตฺเตน ภิกฺขุนา ปฏิพเลน สํโฆ ญาเปตพฺโพ\csromandash{}}

\proseitem{203}{“สุณาตุ เม ภนฺเต สํโฆ, ยทิ สํฆสฺส ปตฺตกลฺลํ, สํโฆ อิตฺถนฺนามํ ภิกฺขุํ สลากคฺคาหาปกํ สมฺมนฺเนยฺย, เอสา ญตฺติ.}

\prose{สุณาตุ เม ภนฺเต สํโฆ, สํโฆ อิตฺถนฺนามํ ภิกฺขุํ สลากคฺคาหาปกํ สมฺมนฺนติ, ยสฺสายสฺมโต ขมติ อิตฺถนฺนามสฺส ภิกฺขุโน สลากคฺคาหาปกสฺส สมฺมุติ, โส ตุณฺหสฺส, ยสฺส นกฺขมติ, โส ภาเสยฺย.}

\prose{สมฺมโต สํเฆน อิตฺถนฺนาโม ภิกฺขุ สลากคฺคาหาปโก, ขมติ สํฆสฺส, ตสฺมา ตุณฺหี, เอวเมตํ ธารยามี”ติ.}

\proseitem{204}{\csromanfolio{205}ทสยิเม ภิกฺขเว อธมฺมิกา สลากคฺคาหา, ทส ธมฺมิกา\footnote{ทส ธมฺมิกา สลากคฺคาหา (ก)}.\csromanspacer{} กตเม ทส อธมฺมิกา สลากคฺคาหา.\csromanspacer{} โอรมตฺตกญฺจ อธิกรณํ โหติ, น จ คติคตํ โหติ, น จ สริตสาริตํ โหติ, ชานาติ “อธมฺมวาที พหุตรา”ติ, อปฺเปว นาม อธมฺมวาที พหุตรา อสฺสูติ, ชานาติ “สํโฆ ภิชฺชิสฺสตี”ติ, อปฺเปว นาม สํโฆ ภิชฺเชยฺยาติ, อธมฺเมน คณฺหนฺติ, วคฺคา คณฺหนฺติ, น จ ยถาทิฏฺฐิยา คณฺหนฺติ.\csromanspacer{} อิเม ทส อธมฺมิกา สลากคฺคาหา.}

\prose{กตเม ทส ธมฺมิกา สลากคฺคาหา.\csromanspacer{} น จ โอรมตฺตกํ อธิกรณํ โหติ, คติคตญฺจ โหติ, สริตสาริตญฺจ โหติ, ชานาติ “ธมฺมวาที พหุตรา”ติ, อปฺเปว นาม ธมฺมวาที พหุตรา อสฺสูติ, ชานาติ “สํโฆ น ภิชฺชิสฺสตี”ติ, อปฺเปว นาม สํโฆ น ภิชฺเชยฺยาติ, ธมฺเมน คณฺหนฺติ, สมคฺคา คณฺหนฺติ, ยถาทิฏฺฐิยา จ คณฺหนฺติ.\csromanspacer{} อิเม ทส ธมฺมิกา สลากคฺคาหาติ.}

\csromanmatikaanchor{mtk.109}
\csromantocmarkref{section}{6. ตสฺสปาปิยสิกา}{mtk.109}
\bookmark[dest={mtk.109},level=1]{6. ตสฺสปาปิยสิกา}
\csromanheader{1}{6. ตสฺสปาปิยสิกา}
\proseitem{205}{เตน โข ปน สมเยน อุปวาโฬ ภิกฺขุ สํฆมชฺเฌ อาปตฺติยา อนุยุญฺชิยมาโน อวชานิตฺวา ปฏิชานาติ, ปฏิชานิตฺวา อวชานาติ, อญฺเญนญฺญํ ปฏิจรติ, สมฺปชานมุสา ภาสติ.\csromanspacer{} เย เต ภิกฺขู อปฺปิจฺฉา ฯเปฯ เต อุชฺฌายนฺติ ขิยฺยนฺติ วิปาเจนฺติ “กถํ หิ นาม อุปวาโฬ ภิกฺขุ สํฆมชฺเฌ อาปตฺติยา อนุยุญฺชิยมาโน อวชานิตฺวา ปฏิชานิสฺสติ, ปฏิชานิตฺวา อวชานิสฺสติ, อญฺเญนญฺญํ ปฏิจริสฺสติ, สมฺปชานมุสา ภาสิสฺสตี”ติ.\csromanspacer{} อถ โข เต ภิกฺขู ภควโต เอตมตฺถํ อาโรเจสุํ ฯเปฯ.\csromanspacer{} สจฺจํ กิร ภิกฺขเว ฯเปฯ.\csromanspacer{} สจฺจํ ภควาติ ฯเปฯ วิครหิตฺวา ฯเปฯ ธมฺมิํ กถํ กตฺวา ภิกฺขู อามนฺเตสิ\csromandash{} เตน หิ ภิกฺขเว สํโฆ อุปวาฬสฺส ภิกฺขุโน ตสฺสปาปิยสิกากมฺมํ กโรตุ.\csromanspacer{} เอวญฺจ ปน ภิกฺขเว กาตพฺพํ, ปฐมํ อุปวาโฬ ภิกฺขุ โจเทตพฺโพ, โจเทตฺวา สาเรตพฺโพ, สาเรตฺวา อาปตฺติํ อาโรเปตพฺโพ, อาปตฺติํ อาโรเปตฺวา พฺยตฺเตน ภิกฺขุนา ปฏิพเลน สํโฆ ญาเปตพฺโพ\csromandash{}}

\proseitem{206}{“สุณาตุ เม ภนฺเต สํโฆ, อยํ อุปวาโฬ ภิกฺขุ สํฆมชฺเฌ อาปตฺติยา อนุยุญฺชิยมาโน อวชานิตฺวา ปฏิชานาติ, ปฏิชานิตฺวา อวชานาติ, อญฺเญนญฺญํ ปฏิจรติ, สมฺปชานมุสา ภาสติ, \csromanfolio{206}ยทิ สํฆสฺส ปตฺตกลฺลํ, สํโฆ อุปวาฬสฺส ภิกฺขุโน ตสฺสปาปิยสิกากมฺมํ กเรยฺย, เอสา ญตฺติ.}

\prose{สุณาตุ เม ภนฺเต สํโฆ, อยํ อุปวาโฬ ภิกฺขุ สํฆมชฺเฌ อาปตฺติยา อนุยุญฺชิยมาโน อวชานิตฺวา ปฏิชานาติ, ปฏิชานิตฺวา อวชานาติ, อญฺเญนญฺญํ ปฏิจรติ, สมฺปชานมุสา ภาสติ, สํโฆ อุปวาฬสฺส ภิกฺขุโน ตสฺสปาปิยสิกากมฺมํ กโรติ, ยสฺสายสฺมโต ขมติ อุปวาฬสฺส ภิกฺขุโน ตสฺสปาปิยสิกากมฺมสฺส กรณํ, โส ตุณฺหสฺส, ยสฺส นกฺขมติ, โส ภาเสยฺย.}

\prose{ทุติยมฺปิ เอตมตฺถํ วทามิ ฯเปฯ.\csromanspacer{} ตติยมฺปิ เอตมตฺถํ วทามิ ฯเปฯ.}

\prose{กตํ สํเฆน อุปวาฬสฺส ภิกฺขุโน ตสฺสปาปิยสิกากมฺมํ, ขมติ สํฆสฺส, ตสฺมา ตุณฺหี, เอวเมตํ ธารยามี”ติ.}

\proseitem{207}{ปญฺจิมานิ ภิกฺขเว ธมฺมิกานิ ตสฺสปาปิยสิกากมฺมสฺส กรณานิ.\csromanspacer{} อสุจิ จ โหติ, อลชฺชี จ, สานุวาโท จ, ตสฺส สํโฆ ตสฺสปาปิยสิกากมฺมํ กโรติ ธมฺเมน สมคฺเคน.\csromanspacer{} อิมานิ โข ภิกฺขเว ปญฺจ ธมฺมิกานิ ตสฺสปาปิยสิกากมฺมสฺส กรณานิ.}

\csromanmatikaanchor{mtk.110}
\csromantocmarkref{subsection}{อธมฺมกมฺมทฺวาทสก}{mtk.110}
\bookmark[dest={mtk.110},level=2]{อธมฺมกมฺมทฺวาทสก}
\csromanheader{2}{\textbf{อธมฺมกมฺมทฺวาทสก}}
\proseitem{208}{ตีหิ ภิกฺขเว องฺเคหิ สมนฺนาคตํ ตสฺสปาปิยสิกากมฺมํ อธมฺมกมฺมญฺจ โหติ อวินยกมฺมญฺจ ทุวูปสนฺตญฺจ.\csromanspacer{} อสมฺมุขา กตํ โหติ, อปฺปฏิปุจฺฉากตํ โหติ, อปฺปฏิญฺญาย กตํ โหติ ฯเปฯ.\csromanspacer{} อธมฺเมน กตํ โหติ, วคฺเคน กตํ โหติ.\csromanspacer{} อิเมหิ โข ภิกฺขเว ตีหงฺเคหิ สมนฺนาคตํ ตสฺสปาปิยสิกากมฺมํ อธมฺมกมฺมญฺจ โหติ อวินยกมฺมญฺจ ทุวูปสนฺตญฺจ.}

\csromanmatikaanchor{mtk.111}
\csromantocmarkref{subsection}{ธมฺมกมฺมทฺวาทสก}{mtk.111}
\bookmark[dest={mtk.111},level=2]{ธมฺมกมฺมทฺวาทสก}
\csromanheader{2}{\textbf{ธมฺมกมฺมทฺวาทสก}}
\proseitem{209}{ตีหิ ภิกฺขเว องฺเคหิ สมนฺนาคตํ ตสฺสปาปิยสิกากมฺมํ ธมฺมกมฺมญฺจ โหติ วินยกมฺมญฺจ สุวูปสนฺตญฺจ.\csromanspacer{} สมฺมุขา กตํ โหติ, ปฏิปุจฺฉากตํ โหติ, ปฏิญฺญาย กตํ โหติ ฯเปฯ.\csromanspacer{} ธมฺเมน กตํ โหติ, สมคฺเคน กตํ โหติ.\csromanspacer{} อิเมหิ โข ภิกฺขเว ตีหงฺเคหิ สมนฺนาคตํ ตสฺสปาปิยสิกากมฺมํ ธมฺมกมฺมญฺจ โหติ วินยกมฺมญฺจ สุวูปสนฺตญฺจ.}

\csromanmatikaanchor{mtk.112}
\csromantocmarkref{subsection}{อากงฺขมานฉกฺก}{mtk.112}
\bookmark[dest={mtk.112},level=2]{อากงฺขมานฉกฺก}
\csromanheader{2}{\textbf{อากงฺขมานฉกฺก}}
\proseitem{210}{\csromanfolio{207}ตีหิ ภิกฺขเว องฺเคหิ สมนฺนาคตสฺส ภิกฺขุโน อากงฺขมาโน สํโฆ ตสฺสปาปิยสิกากมฺมํ กเรยฺย.\csromanspacer{} ภณฺฑนการโก โหติ กลหการโก วิวาทการโก ภสฺสการโก สํเฆ อธิกรณการโก, พาโล โหติ อพฺยตฺโต อาปตฺติพหุโล อนปทาโน, คิหิสํสฏฺโฐ วิหรติ อนนุโลมิเกหิ คิหิสํสคฺเคหิ.\csromanspacer{} อิเมหิ โข ภิกฺขเว ตีหงฺเคหิ สมนฺนาคตสฺส ภิกฺขุโน อากงฺขมาโน สํโฆ ตสฺสปาปิยสิกากมฺมํ กเรยฺย. (1)}

\prose{อปเรหิปิ ภิกฺขเว ตีหงฺเคหิ สมนฺนาคตสฺส ภิกฺขุโน อากงฺขมาโน สํโฆ ตสฺสปาปิยสิกากมฺมํ กเรยฺย.\csromanspacer{} อธิสีเล สีลวิปนฺโน โหติ, อชฺฌาจาเร อาจารวิปนฺโน โหติ, อติทิฏฺฐิยา ทิฏฺฐิวิปนฺโน โหติ.\csromanspacer{} อิเมหิ โข ภิกฺขเว ฯเปฯ. (2)}

\prose{อปเรหิปิ ภิกฺขเว ฯเปฯ.\csromanspacer{} พุทฺธสฺส อวณฺณํ ภาสติ, ธมฺมสฺส อวณฺณํ ภาสติ, สํฆสฺส อวณฺณํ ภาสติ.\csromanspacer{} อิเมหิ โข ภิกฺขเว ตีหงฺเคหิ สมนฺนาคตสฺส ภิกฺขุโน อากงฺขมาโน สํโฆ ตสฺสปาปิยสิกากมฺมํ กเรยฺย. (3)}

\prose{ติณฺณํ ภิกฺขเว ภิกฺขูนํ อากงฺขมาโน สํโฆ ตสฺสปาปิยสิกากมฺมํ กเรยฺย.\csromanspacer{} เอโก ภณฺฑนการโก โหติ ฯเปฯ สํเฆ อธิกรณการโก, เอโก พาโล โหติ อพฺยตฺโต อาปตฺติพหุโล อนปทาโน, เอโก คิหิสํสฏฺโฐ วิหรติ อนนุโลมิเกหิ คิหิสํสคฺเคหิ.\csromanspacer{} อิเมสํ โข ภิกฺขเว ติณฺณํ ภิกฺขูนํ อากงฺขมาโน สํโฆ ตสฺสปาปิยสิกากมฺมํ กเรยฺย. (4)}

\prose{อปเรสมฺปิ ภิกฺขเว ติณฺณํ ภิกฺขูนํ อากงฺขมาโน สํโฆ ตสฺสปาปิยสิกากมฺมํ กเรยฺย.\csromanspacer{} เอโก อธิสีเล สีลวิปนฺโน โหติ, เอโก อชฺฌาจาเร อาจารวิปนฺโน โหติ, เอโก อติทิฏฺฐิยา ทิฏฺฐิวิปนฺโน โหติ.\csromanspacer{} อิเมสํ โข ภิกฺขเว ฯเปฯ. (5)}

\prose{อปเรสมฺปิ ภิกฺขเว ฯเปฯ.\csromanspacer{} เอโก พุทฺธสฺส อวณฺณํ ภาสติ, เอโก ธมฺมสฺส อวณฺณํ ภาสติ, เอโก สํฆสฺส อวณฺณํ ภาสติ.\csromanspacer{} อิเมสํ โข ภิกฺขเว ติณฺณํ ภิกฺขูนํ อากงฺขมาโน สํโฆ ตสฺสปาปิยสิกากมฺมํ กเรยฺย. (6)}

\csromanmatikaanchor{mtk.113}
\csromantocmarkref{subsection}{อฏฺฐารสวตฺต}{mtk.113}
\bookmark[dest={mtk.113},level=2]{อฏฺฐารสวตฺต}
\csromanheader{2}{\textbf{อฏฺฐารสวตฺต}}
\proseitem{211}{\csromanfolio{208}ตสฺสปาปิยสิกากมฺมกเตน ภิกฺขเว ภิกฺขุนา สมฺมา วตฺติตพฺพํ.\csromanspacer{} ตตฺรายํ สมฺมาวตฺตนา\csromandash{}}

\prose{น อุปสมฺปาเทตพฺพํ, น นิสฺสโย ทาตพฺโพ, น สามเณโร อุปฏฺฐาเปตพฺโพ, น ภิกฺขุโนวาทกสมฺมุติ สาทิตพฺพา, สมฺมเตนปิ ภิกฺขุนิโย น โอวทิตพฺพา ฯเปฯ น ภิกฺขูหิ สมฺปโยเชตพฺพนฺติ.\csromanspacer{} อถ โข สํโฆ อุปวาฬสฺส ภิกฺขุโน ตสฺสปาปิยสิกากมฺมํ อกาสิ.}

\csromanmatikaanchor{mtk.114}
\csromantocmarkref{section}{7. ติณวตฺถารก}{mtk.114}
\bookmark[dest={mtk.114},level=1]{7. ติณวตฺถารก}
\csromanheader{1}{7. ติณวตฺถารก}
\proseitem{212}{เตน โข ปน สมเยน ภิกฺขูนํ ภณฺฑนชาตานํ กลหชาตานํ วิวาทาปนฺนานํ วิหรตํ พหุํ อสฺสามณกํ อชฺฌาจิณฺณํ โหติ ภาสิตปริกฺกนฺตํ.\csromanspacer{} อถ โข เตสํ ภิกฺขูนํ เอตทโหสิ “อมฺหากํ โข ภณฺฑนชาตานํ กลหชาตานํ วิเวทาปนฺนานํ วิหรตํ พหุํ อสฺสามณกํ อชฺฌาจิณฺณํ ภาสิตปริกฺกนฺตํ, สเจ มยํ อิมาหิ อาปตฺตีหิ อญฺญมญฺญํ กาเรสฺสาม, สิยาปิ ตํ อธิกรณํ กกฺขฬตฺตาย วาฬตฺตาย\footnote{กกฺขฬตาย วาฬตาย (สฺยา, กํ)} เภทาย สํวตฺเตยฺย, กถํ นุ โข อเมฺหหิ ปฏิปชฺชิตพฺพนฺ”ติ.\csromanspacer{} ภควโต เอตมตฺถํ อาโรเจสุํ.}

\prose{อิธ ปน ภิกฺขเว ภิกฺขูนํ ภณฺฑนชาตานํ กลหชาตานํ วิวาทาปนฺนานํ วิหรตํ พหุํ อสฺสามณกํ อชฺฌาจิณฺณํ โหติ ภาสิตปริกฺกนฺตํ, ตตฺร เจ ภิกฺขูนํ\footnote{ตตฺร เจ ภิกฺขเว ภิกฺขูนํ (สฺยา)} เอวํ โหติ “อมฺหากํ โข ภณฺฑนชาตานํ กลหชาตานํ วิวาทาปนฺนานํ วิหรตํ พหุํ อสฺสามณกํ อชฺฌาจิณฺณํ ภาสิตปริกฺกนฺตํ, สเจ มยํ อิมาหิ อาปตฺตีหิ อญฺญมญฺญํ กาเรสฺสาม, สิยาปิ ตํ อธิกรณํ กกฺขฬตฺตาย วาฬตฺตาย เภทาย สํวตฺเตยฺยา”ติ.\csromanspacer{} อนุชานามิ ภิกฺขเว เอวรูปํ อธิกรณํ ติณวตฺถารเกน วูปสเมตุํ.\csromanspacer{} เอวญฺจ ปน ภิกฺขเว วูปสเมตพฺพํ, สพฺเพเหว เอกชฺฌํ สนฺนิปติตพฺพํ, สนฺนิปติตฺวา พฺยตฺเตน ภิกฺขุนา ปฏิพเลน สํโฆ ญาเปตพฺโพ\csromandash{}}

\prose{“สุณาตุ เม ภนฺเต สํโฆ, อมฺหากํ ภณฺฑนชาตานํ กลหชาตานํ วิวาทาปนฺนานํ วิหรตํ พหุํ อสฺสามณกํ อชฺฌาจิณฺณํ ภาสิตปริกฺกนฺตํ, สเจ มยํ อิมาหิ อาปตฺตีหิ อญฺญมญฺญํ กาเรสฺสาม, สิยาปิ \csromanfolio{209}ตํ อธิกรณํ กกฺขฬตฺตาย วาฬตฺตาย เภทาย สํวตฺเตยฺย, ยทิ สํฆสฺส ปตฺตกลฺลํ, สํโฆ อิมํ อธิกรณํ ติณวตฺถารเกน วูปสเมยฺย ฐเปตฺวา ถุลฺลวชฺชํ ฐเปตฺวา คิหิปฺปฏิสํยุตฺตนฺ”ติ.\csromanspacer{} เอกโตปกฺขิกานํ ภิกฺขูนํ พฺยตฺเตน ภิกฺขุนา ปฏิพเลน สโก ปกฺโข ญาเปตพฺโพ\csromandash{}}

\prose{“สุณนฺตุ เม อายสฺมนฺตา, อมฺหากํ ภณฺฑนชาตานํ กลหชาตานํ วิวาทาปนฺนานํ วิหรตํ พหุํ อสฺสามณกํ อชฺฌาจิณฺณํ ภาสิตปริกฺกนฺตํ, สเจ มยํ อิมาหิ อาปตฺตีหิ อญฺญมญฺญํ กาเรสฺสาม, สิยาปิ ตํ อธิกรณํ กกฺขฬตฺตาย วาฬตฺตาย เภทาย สํวตฺเตยฺย, ยทายสฺมนฺตานํ ปตฺตกลฺลํ, อหํ ยา เจว อายสฺมนฺตานํ อาปตฺติ, ยา จ อตฺตโน อาปตฺติ, อายสฺมนฺตานญฺเจว อตฺถาย อตฺตโน จ อตฺถาย สํฆมชฺเฌ ติณวตฺถารเกน เทเสยฺยํ ฐเปตฺวา ถุลฺลวชฺชํ, ฐเปตฺวา คิหิปฺปฏิสํยุตฺตนฺ”ติ.}

\prose{อถาปเรสํ เอกโตปกฺขิกานํ ภิกฺขูนํ พฺยตฺเตน ภิกฺขุนา ปฏิพเลน สโก ปกฺโข ญาเปตพฺโพ\csromandash{}}

\prose{“สุณนฺตุ เม อายสฺมนฺตา, อมฺหากํ ภณฺฑนชาตานํ กลหชาตานํ วิวาทาปนฺนานํ วิหรตํ พหุํ อสฺสามณกํ อชฺฌาจิณฺณํ ภาสิตปริกฺกนฺตํ, สเจ มยํ อิมาหิ อาปตฺตีหิ อญฺญมญฺญํ กาเรสฺสาม, สิยาปิ ตํ อธิกรณํ กกฺขฬตฺตาย วาฬตฺตาย เภทาย สํวตฺเตยฺย, ยทายสฺมนฺตานํ ปตฺตกลฺลํ อหํ ยา เจว อายสฺมนฺตานํ อาปตฺติ, ยา จ อตฺตโน อาปตฺติ, อายสฺมนฺตานญฺเจว อตฺถาย อตฺตโน จ อตฺถาย สํฆมชฺเฌ ติณวตฺถารเกน เทเสยฺยํ ฐเปตฺวา ถุลฺลวชฺชํ ฐเปตฺวา คิหิปฺปฏิสํยุตฺตนฺ”ติ.}

\proseitem{213}{อถาปเรสํ เอกโตปกฺขิกานํ ภิกฺขูนํ พฺยตฺเตน ภิกฺขุนา ปฏิพเลน สํโฆ ญาเปตพฺโพ\csromandash{}}

\prose{“สุณาตุ เม ภนฺเต สํโฆ, อมฺหากํ ภณฺฑนชาตานํ กลหชาตานํ วิวาทาปนฺนานํ วิหรตํ พหุํ อสฺสามณกํ อชฺฌาจิณฺณํ ภาสิตปริกฺกนฺตํ, สเจ มยํ อิมาหิ อาปตฺตีหิ อญฺญมญฺญํ กาเรสฺสาม, สิยาปิ ตํ อธิกรณํ กกฺขฬตฺตาย วาฬตฺตาย เภทาย สํวตฺเตยฺย, ยทิ สํฆสฺส ปตฺตกลฺลํ อหํ ยา เจว อิเมสํ อายสฺมนฺตานํ อาปตฺติ, ยา จ อตฺตโน อาปตฺติ, อิเมสญฺเจว อายสฺมนฺตานํ อตฺถาย อตฺตโน \csromanfolio{210}จ อตฺถาย สํฆมชฺเฌ ติณวตฺถารเกน เทเสยฺยํ ฐเปตฺวา ถุลฺลวชฺชํ, ฐเปตฺวา คิหิปฺปฏิสํยุตฺตํ, เอสา ญตฺติ.}

\prose{สุณาตุ เม ภนฺเต สํโฆ, อมฺหากํ ภณฺฑนชาตานํ กลหชาตานํ วิวาทาปนฺนานํ วิหรตํ พหุํ อสฺสามณกํ อชฺฌาจิณฺณํ ภาสิตปริกฺกนฺตํ, สเจ มยํ อิมาหิ อาปตฺตีหิ อญฺญมญฺญํ กาเรสฺสาม, สิยาปิ ตํ อธิกรณํ กกฺขฬตฺตาย วาฬตฺตาย เภทาย สํวตฺเตยฺย, อหํ ยา เจว อิเมสํ อายสฺมนฺตานํ อาปตฺติ, ยา จ อตฺตโน อาปตฺติ, อิเมสญฺเจว อายสฺมนฺตานํ อตฺถาย อตฺตโน จ อตฺถาย สํฆมชฺเฌ ติณวตฺถารเกน เทเสมิ ฐเปตฺวา ถุลฺลวชฺชํ, ฐเปตฺวา คิหิปฺปฏิสํยุตฺตํ, ยสฺสายสฺมโต ขมติ อมฺหากํ อิมาสํ อาปตฺตีนํ สํฆมชฺเฌ ติณวตฺถารเกน เทสนา ฐเปตฺวา ถุลฺลวชฺชํ, ฐเปตฺวา คิหิปฺปฏิสํยุตฺตํ, โส ตุณฺหสฺส, ยสฺส นกฺขมติ, โส ภาเสยฺย.}

\prose{เทสิตา อมฺหากํ อิมา อาปตฺติโย สํฆมชฺเฌ ติณวตฺถารเกน ฐเปตฺวา ถุลฺลวชฺชํ, ฐเปตฺวา คิหิปฺปฏิสํยุตฺตํ, ขมติ สํฆสฺส, ตสฺมา ตุณฺหี, เอวเมตํ ธารยามี”ติ.}

\proseitem{214}{อถาปเรสํ เอกโตปกฺขิกานํ ภิกฺขูนํ พฺยตฺเตน ภิกฺขุนา ปฏิพเลน สํโฆ ญาเปตพฺโพ\csromandash{}}

\prose{“สุณาตุ เม ภนฺเต สํโฆ, อมฺหากํ ภณฺฑนชาตานํ กลหชาตานํ วิวาทาปนฺนานํ วิหรตํ พหุํ อสฺสามณกํ อชฺฌาจิณฺณํ ภาสิตปริกฺกนฺตํ, สเจ มยํ อิมาหิ อาปตฺตีหิ อญฺญมญฺญํ กาเรสฺสาม, สิยาปิ ตํ อธิกรณํ กกฺขฬตฺตาย วาฬตฺตาย เภทาย สํวตฺเตยฺย, ยทิ สํฆสฺส ปตฺตกลฺลํ, อหํ ยา เจว อิเมสํ อายสฺมนฺตานํ อาปตฺติ, ยา จ อตฺตโน อาปตฺติ, อิเมสญฺเจว อายสฺมนฺตานํ อตฺถาย อตฺตโน จ อตฺถาย สํฆมชฺเฌ ติณวตฺถารเกน เทเสยฺยํ ฐเปตฺวา ถุลฺลวชฺชํ, ฐเปตฺวา คิหิปฺปฏิสํยุตฺตํ, เอสา ญตฺติ.}

\prose{สุณาตุ เม ภนฺเต สํโฆ, อมฺหากํ ภณฺฑนชาตานํ กลหชาตานํ วิวาทาปนฺนานํ วิหรตํ พหุํ อสฺสามณกํ อชฺฌาจิณฺณํ ภาสิตปริกฺกนฺตํ, สเจ มยํ อิมาหิ อาปตฺตีหิ อญฺญมญฺญํ กาเรสฺสาม, สิยาปิ ตํ อธิกรณํ กกฺขฬตฺตาย วาฬตฺตาย เภทาย สํวตฺเตยฺย, อหํ ยา เจว อิเมสํ อายสฺมนฺตานํ อาปตฺติ, ยา จ อตฺตโน อาปตฺติ, อิเมสญฺเจว \csromanfolio{211}อายสฺมนฺตานํ อตฺถาย อตฺตโน จ อตฺถาย สํฆมชฺเฌ ติณวตฺถารเกน เทเสมิ ฐเปตฺวา ถุลฺลวชฺชํ, ฐเปตฺวา คิหิปฺปฏิสํยุตฺตํ, ยสฺสายสฺมโต ขมติ อมฺหากํ อิมาสํ อาปตฺตีนํ สํฆมชฺเฌ ติณวตฺถารเกน เทสนา ฐเปตฺวา ถุลฺลวชฺชํ, ฐเปตฺวา คิหิปฺปฏิสํยุตฺตํ, โส ตุณฺหสฺส, ยสฺส นกฺขมติ, โส ภาเสยฺย.}

\prose{เทสิตา อมฺหากํ อิมา อาปตฺติโย สํฆมชฺเฌ ติณวตฺถารเกน ฐเปตฺวา ถุลฺลวชฺชํ, ฐเปตฺวา คิหิปฺปฏิสํยุตฺตํ, ขมติ สํฆสฺส, ตสฺมา ตุณฺหี, เอวเมตํ ธารยามี”ติ.}

\prose{เอวญฺจ ปน ภิกฺขเว เต ภิกฺขู ตาหิ อาปตฺตีหิ วุฏฺฐิตา โหนฺติ ฐเปตฺวา ถุลฺลวชฺชํ, ฐเปตฺวา คิหิปฺปฏิสํยุตฺตํ ฐเปตฺวา ทิฏฺฐาวิกมฺมํ ฐเปตฺวา เย น ตตฺถ โหนฺตีติ.}

\csromanmatikaanchor{mtk.115}
\csromantocmarkref{section}{8. อธิกรณ}{mtk.115}
\bookmark[dest={mtk.115},level=1]{8. อธิกรณ}
\csromanheader{1}{8. อธิกรณ}
\proseitem{215}{เตน โข ปน สมเยน (ภิกฺขูปิ ภิกฺขูหิ วิวทนฺติ)\footnote{( ) นตฺถิ (สี, สฺยา, กํ)} ภิกฺขูปิ ภิกฺขุนีหิ วิวทนฺติ, ภิกฺขุนิโยปิ ภิกฺขูหิ วิวทนฺติ, ฉนฺโนปิ ภิกฺขุ ภิกฺขุนีนํ อนุปขชฺช ภิกฺขูหิ สทฺธิํ วิวทติ, ภิกฺขุนีนํ ปกฺขํ คาเหติ.\csromanspacer{} เย เต ภิกฺขู อปฺปิจฺฉา ฯเปฯ เต อุชฺฌายนฺติ ขิยฺยนฺติ วิปาเจนฺติ “กถํ หิ นาม ฉนฺโน ภิกฺขุ ภิกฺขุนีนํ อนุปขชฺช ภิกฺขูหิ สทฺธิํ วิวทิสฺสติ, ภิกฺขุนีนํ ปกฺขํ คาเหสฺสตี”ติ.\csromanspacer{} อถ โข เต ภิกฺขู ภควโต เอตมตฺถํ อาโรเจสุํ ฯเปฯ.\csromanspacer{} สจฺจํ กิร ภิกฺขเว ฯเปฯ.\csromanspacer{} สจฺจํ ภควาติ ฯเปฯ วิครหิตฺวา ฯเปฯ ธมฺมิํ กถํ กตฺวา ภิกฺขู อามนฺเตสิ\csromandash{}}

\prose{\csromansymbolfootnote{*}{วิ 5. 169 ปิฏฺเฐปิ. 2. กิจฺจาธิกรณญฺจ (ก)}จตฺตาริมานิ ภิกฺขเว อธิกรณานิ วิวาทาธิกรณํ อนุวาทาธิกรณํ อาปตฺตาธิกรณํ กิจฺจาธิกรณํ2.}

\prose{ตตฺถ กตมํ วิวาทาธิกรณํ. + อิธ ปน ภิกฺขเว ภิกฺขู\footnote{ภิกฺขู ภิกฺขุํ (ก)} วิวทนฺติ ธมฺโมติ วา อธมฺโมติ วา, วินโยติ วา อวินโยติ วา, ภาสิตํ ลปิตํ ตถาคเตนาติ วา อภาสิตํ อลปิตํ ตถาคเตนาติ วา, อาจิณฺณํ ตถาคเตนาติ วา อนาจิณฺณํ ตถาคเตนาติ วา, ปญฺญตฺตํ ตถาคเตนาติ วา อปฺปญฺญตฺตํ ตถาคเตนาติ วา, อาปตฺตีติ วา อนาปตฺตีติ \csromanfolio{212}วา, ลหุกา อาปตฺตีติ วา ครุกา อาปตฺตีติ วา, สาวเสสา อาปตฺตีติ วา อนวเสสา อาปตฺตีติ วา, ทุฏฺฐุลฺลา อาปตฺตีติ วา อทุฏฺฐุลฺลา อาปตฺตีติ วา.\csromanspacer{} ยํ ตตฺถ ภณฺฑนํ กลโห วิคฺคโห วิวาโท นานาวาโท อญฺญถาวาโท วิปจฺจตาย โวหาโร เมธคํ, อิทํ วุจฺจติ วิวาทาธิกรณํ.}

\prose{ตตฺถ กตมํ อนุวาทาธิกรณํ. * อิธ ปน ภิกฺขเว ภิกฺขู ภิกฺขุํ อนุวทนฺติ สีลวิปตฺติยา วา อาจารวิปตฺติยา วา ทิฏฺฐิวิปตฺติยา วา อาชีววิปตฺติยา วา.\csromanspacer{} โย ตตฺถ อนุวาโท อนุวทนา อนุลฺลปนา อนุภณนา อนุสมฺปวงฺกตา อพฺภุสฺสหนตา อนุพลปฺปทานํ, อิทํ วุจฺจติ อนุวาทาธิกรณํ.}

\prose{ตตฺถ กตมํ อาปตฺตาธิกรณํ.\csromanspacer{} ปญฺจปิ อาปตฺติกฺขนฺธา อาปตฺตาธิกรณํ, สตฺตปิ อาปตฺติกฺขนฺธา อาปตฺตาธิกรณํ, อิทํ วุจฺจติ อาปตฺตาธิกรณํ.}

\prose{ตตฺถ กตมํ กิจฺจาธิกรณํ. + ยา สํฆสฺส กิจฺจยตา กรณียตา อปโลกนกมฺมํ ญตฺติกมฺมํ ญตฺติทุติยกมฺมํ ญตฺติจตุตฺถกมฺมํ, อิทํ วุจฺจติ กิจฺจาธิกรณํ.}

\proseitem{216}{วิวาทาธิกรณสฺส กิํ มูลํ.\csromanspacer{} ฉ วิวาทมูลานิ วิวาทาหิกรณสฺส มูลํ, ตีณิปิ อกุสลมูลนิ, วิวาทาธิกรณสฺส มูลํ, ตีณิปิ กุสลมูลานิ วิวาทาธิกรณสฺส มูลํ. ++ กตมานิ ฉ วิวาทมูลานิ วิวาทาธิกรณสฺส มูลํ.\csromanspacer{} อิธ ปน ภิกฺขเว ภิกฺขุ โกธโน โหติ อุปนาหี, โย โส ภิกฺขเว ภิกฺขุ โกธโน โหติ อุปนาหี, โส สตฺถริปิ อคารโว วิหรติ อปฺปติสฺโส, ธมฺเมปิ อคารโว วิหรติ อปฺปติสฺโส, สํเฆปิ อคารโว วิหรติ อปฺปติสฺโส, สิกฺขายปิ น ปริปูรการี โหติ, โย โส ภิกฺขเว ภิกฺขุ สตฺถริปิ อคารโว วิหรติ อปฺปติสฺโส, ธมฺเมปิ ฯเปฯ สํเฆปิ ฯเปฯ สิกฺขายปิ น ปริปูรการี, โส สํเฆ วิวาทํ ชเนติ, โย โหติ วิวาโท พหุชนาหิตาย พหุชนาสุขาย พหุโน ชนสฺส อนตฺถาย อหิตาย ทุกฺขาย เทวมนุสฺสานํ, เอวรูปญฺเจ ตุเมฺห ภิกฺขเว วิวาทมูลํ อชฺฌตฺตํ วา พหิทฺธา วา สมนุปสฺเสยฺยาถ, ตตฺร ตุเมฺห ภิกฺขเว ตสฺเสว ปาปกสฺส วิวาทมูลสฺส ปหานาย วายเมยฺยาถ, เอวรูปญฺเจ ตุเมฺห ภิกฺขเว วิวาทมูลํ อชฺฌตฺตํ วา พหิทฺธา วา น สมนุปสฺเสยฺยาถ, ตตฺร ตุเมฺห ภิกฺขเว ตสฺเสว ปาปกสฺส \csromanfolio{213}วิวาทมูลสฺส อายติํ อนวสฺสวาย ปฏิปชฺเชยฺยาถ, เอวเมตสฺส ปาปกสฺส วิวาทมูลสฺส ปหานํ โหติ, เอวเมตสฺส ปาปกสฺส วิวาทมูลสฺส อายติํ อนวสฺสโว โหติ.}

\prose{\csromansymbolfootnote{*}{วิ 5. 166 ปิฏฺเฐปิ.}ปุน จปรํ ภิกฺขเว ภิกฺขุ มกฺขี โหติ ปฬาสี ฯเปฯ อิสฺสุกี โหติ มจฺฉรี.\csromanspacer{} สโฐ โหติ มายาวี.\csromanspacer{} ปาปิจฺโฉ โหติ มิจฺฉาทิฏฺฐี.\csromanspacer{} สนฺทิฏฺฐิปรามาสี โหติ อาธานคฺคาหี ทุปฺปฏินิสฺสคฺคี, โย โส ภิกฺขเว ภิกฺขุ สนฺทิฏฺฐิปรามาสี โหติ อาธานคฺคาหี ทุปฺปฏินิสฺสคฺคี, โส สตฺถริปิ อคารโว วิหรติ อปฺปติสฺโส, ธมฺเมปิ อคารโว วิหรติ อปฺปติสฺโส, สํเฆปิ อคารโว วิหรติ อปฺปติสฺโส, สิกฺขายปิ น ปริปูรการี โหติ, โย โส ภิกฺขเว ภิกฺขุ สตฺถริปิ อคารโว วิหรติ อปฺปติสฺโส, ธมฺเมปิ ฯเปฯ สํเฆปิ ฯเปฯ สิกฺขายปิ น ปริปูรการี, โส สํเฆ วิวาทํ ชเนติ, โย โหติ วิวาโท พหุชนาหิตาย พหุชนาสุขาย พหุโน ชนสฺส อนตฺถาย อหิตาย ทุกฺขาย เทวมนุสฺสานํ, เอวรูปญฺเจ ตุเมฺห ภิกฺขเว วิวาทมูลํ อชฺฌตฺตํ วา พหิทฺธา วา สมนุปสฺเสยฺยาถ, ตตฺร ตุเมฺห ภิกฺขเว ตสฺเสว ปาปกสฺส วิวาทมูลสฺส ปหานาย วายเมยฺยาถ, เอวรูปญฺเจ ตุเมฺห ภิกฺขเว วิวาทมูลํ อชฺฌตฺตํ วา พหิทฺธา วา น สมนุปสฺเสยฺยาถ, ตตฺร ตุเมฺห ภิกฺขเว ตสฺเสว ปาปกสฺส วิวาทมูลสฺส อายติํ อนวสฺสวาย ปฏิปชฺเชยฺยาถ, เอวเมตสฺส ปาปกสฺส วิวาทมูลสฺส ปหานํ โหติ, เอวเมตสฺส ปาปกสฺส วิวาทมูลสฺส อายติํ อนวสฺสโว โหติ.\csromanspacer{} อิมานิ ฉ วิวาทมูลานิ วิวาทาธิกรณสฺส มูลํ.}

\prose{กตมานิ ตีณิ อกุสลมูลานิ วิวาทาธิกรณสฺส มูลํ.\csromanspacer{} อิธ ปน ภิกฺขเว ภิกฺขู ลุทฺธจิตฺตา วิวทนฺติ.\csromanspacer{} ทุฏฺฐจิตฺตา วิวทนฺติ.\csromanspacer{} มูฬฺหจิตฺตา วิวทนฺติ ธมฺโมติ วา อธมฺโมติ วา, วินโยติ วา อวินโยติ วา, ภาสิตํ ลปิตํ ตถาคเตนาติ วา อภาสิตํ อลปิตํ ตถาคเตนาติ วา, อาจิณฺณํ ตถาคเตนาติ วา อนาจิณฺณํ ตถาคเตนาติ วา, ปญฺญตฺตํ ตถาคเตนาติ วา อปฺปญฺญตฺตํ ตถาคเตนาติ วา, อาปตฺตีติ วา อนาปตฺตีติ วา, ลหุกา อาปตฺตีติ วา ครุกา อาปตฺตีติ วา, สาวเสสา อาปตฺตีติ วา อนวเสสา อาปตฺตีติ วา, ทุฏฺฐุลฺลา อาปตฺตีติ วา อทุฏฺฐุลฺลา อาปตฺตีติ วา.\csromanspacer{} อิมานิ ตีณิ อกุสลมูลานิ วิวาทาธิกรณสฺส มูลํ.}

\prose{\csromanfolio{214}กตมานิ ตีณิ กุสลมูลานิ วิวาทาธิกรณสฺส มูลํ.\csromanspacer{} อิธ ปน ภิกฺขเว ภิกฺขู อลุทฺธจิตฺตา วิวทนฺติ.\csromanspacer{} อทุฏฺฐจิตฺตา วิวทนฺติ.\csromanspacer{} อมูฬฺหจิตฺตา วิวทนฺติ ธมฺโมติ วา อธมฺโมติ วา ฯเปฯ ทุฏฺฐุลฺลา อาปตฺตีติ วา อทุฏฺฐุลฺลา อาปตฺตีติ วา.\csromanspacer{} อิมานิ ตีณี กุสลมูลานิ วิวาทาธิกรสฺส มูลํ.}

\proseitem{217}{\csromansymbolfootnote{*}{วิ 5. 166 ปิฏฺเฐปิ.}อนุวาทาธิกรณสฺส กิํ มูลํ.\csromanspacer{} ฉ อนุวาทมูลานิ อนุวาทาธิกรณสฺส มูลํ, ตีณิปิ อกุสลมูลานิ อนุวาทาธิกรณสฺส มูลํ, ตีณิปิ กุสลมูลานิ อนุวาทาธิกรณสฺส มูลํ, กาโยปิ อนุวาทาธิกรณสฺส มูลํ, วาจาปิ อนุวาทาธิกรณสฺส มูลํ.\csromanspacer{} กตมานิ ฉ อนุวาทมูลานิ อนุวาทาธิกรณสฺส มูลํ.\csromanspacer{} อิธ ปน ภิกฺขเว ภิกฺขุ โกธโน โหติ อุปนาหี, โย โส ภิกฺขเว ภิกฺขุ โกธโน โหติ อุปนาหี, โส สตฺถริปิ อคารโว วิหรติ อปฺปติสฺโส, ธมฺเมปิ อคารโว วิหรติ อปฺปติสฺโส, สํเฆปิ อคารโว วิหรติ อปฺปติสฺโส, สิกฺขายปิ น ปริปูรการี โหติ, โย โส ภิกฺขเว ภิกฺขุ สตฺถริปิ อคารโว วิหรติ อปฺปติสฺโส, ธมฺเมปิ ฯเปฯ สํเฆปิ ฯเปฯ สิกฺขายปิ น ปริปูรการี, โส สํเฆ อนุวาทํ ชเนติ, โย โหติ อนุวาโท พหุชนาหิตาย พหุชนาสุขาย พหุโน ชนสฺส อนตฺถาย อหิตาย ทุกฺขาย เทวมนุสฺสานํ, เอวรูปญฺเจ ตุเมฺห ภิกฺขเว อนุวาทมูลํ อชฺฌตฺตํ วา พหิทฺธา วา สมนุปสฺเสยฺยาถ, ตตฺร ตุเมฺห ภิกฺขเว ตสฺเสว ปาปกสฺส อนุวาทมูลสฺส ปหานาย วายเมยฺยาถ, เอวรูปญฺเจ ตุเมฺห ภิกฺขเว อนุวาทมูลํ อชฺฌตฺตํ วา พหิทฺธา วา น สมนุปสฺเสยฺยาถ, ตตฺร ตุเมฺห ภิกฺขเว ตสฺเสว ปาปกสฺส อนุวาทมูลสฺส อายติํ อนวสฺสวาย ปฏิปชฺเชยฺยาถ, เอวเมตสฺส ปาปกสฺส อนุวาทมูลสฺส ปหานํ โหติ, เอวเมตสฺส ปาปกสฺส อนุวาทมูลสฺส อายติํ อนวสฺสโว โหติ.}

\prose{ปุน จปรํ ภิกฺขเว ภิกฺขุ มกฺขี โหติ ปฬาสี ฯเปฯ อิสฺสุกี โหติ มจฺฉรี.\csromanspacer{} สโฐ โหติ มายาวี.\csromanspacer{} ปาปิจฺโฉ โหติ มิจฺฉาทิฏฺฐี.\csromanspacer{} สนฺทิฏฺฐิปรามาสี โหติ อาธานคฺคาหี ทุปฺปฏินิสฺสคฺคี, โย โส ภิกฺขเว ภิกฺขุ สนฺทิฏฺฐิปรามาสี โหติ อาธานคฺคาหี ทุปฺปฏินิสฺสคฺคี, โส สตฺถริปิ อคารโว วิหรติ อปฺปติสฺโส, ธมฺเมปิ อคารโว วิหรติ อปฺปติสฺโส, สํเฆปิ อคารโว วิหรติ อปฺปติสฺโส, สิกฺขายปิ น ปริปูรการี โหติ, โย โส ภิกฺขเว ภิกฺขุ สตฺถริปิ อคารโว วิหรติ \csromanfolio{215}อปฺปติสฺโส, ธมฺเมปิ ฯเปฯ สํเฆปิ ฯเปฯ สิกฺขายปิ น ปริปูรการี, โส สํเฆ อนุวาทํ ชเนติ, โย โหติ อนุวาโท พหุชนาหิตาย พหุชนาสุขาย พหุโน ชนสฺส อนตฺถาย อหิตาย ทุกฺขาย เทวมนุสฺสานํ.\csromanspacer{} เอวรูปญฺเจ ตุเมฺห ภิกฺขเว อนุวาทมูลํ อชฺฌตฺตํ วา พหิทฺธา วา สมนุปสฺเสยฺยาถ, ตตฺร ตุเมฺห ภิกฺขเว ตสฺเสว ปาปกสฺส อนุวาทมูลสฺส ปหานาย วายเมยฺยาถ, เอวรุปญฺเจ ตุเมฺห ภิกฺขเว อนุวาทมูลํ อชฺฌตฺตํ วา พหิทฺธา วา น สมนุปสฺเสยฺยาถ, ตตฺร ตุเมฺห ภิกฺขเว ตสฺเสว ปาปกสฺส อนุวาทมูลสฺส อายติํ อนวสฺสวาย ปฏิปชฺเชยฺยาถ, เอวเมตสฺส ปาปกสฺส อนุวาทมูลสฺส ปหานํ โหติ, เอวเมตสฺส ปาปกสฺส อนุวาทมูลสฺส อายติํ อนวสฺสโว โหติ.\csromanspacer{} อิมานิ ฉ อนุวาทมูลานิ อนุวาทาธิกรณสฺส มูลํ.}

\prose{กตมานิ ตีณิ อกุสลมูลานิ อนุวาทาธิกรณสฺส มูลํ.\csromanspacer{} อิธ ปน ภิกฺขเว ภิกฺขู ภิกฺขุํ ลุทฺธจิตฺตา อนุวทนฺติ.\csromanspacer{} ทุฏฺฐจิตฺตา อนุวทนฺติ.\csromanspacer{} มูฬฺหจิตฺตา อนุวทนฺติ สีลวิปตฺติยา วา อาจารวิปตฺติยา วา ทิฏฺฐิวิปตฺติยา วา อาชีววิปตฺติยา วา.\csromanspacer{} อิมานิ ตีณิ อกุสลมูลานิ อนุวาทาธิกรณสฺส มูลํ.}

\prose{กตมานิ ตีณิ กุสลมูลานิ อนุวาทาธิกรณสฺส มูลํ.\csromanspacer{} อิธ ปน ภิกฺขเว ภิกฺขู ภิกฺขุํ อลุทฺธจิตฺตา อนุวทนฺติ.\csromanspacer{} อทุฏฺฐจิตฺตา อนุวทนฺติ.\csromanspacer{} อมูฬฺหจิตฺตา อนุวทนฺติ สีลวิปตฺติยา วา อาจารวิปตฺติยา วา ทิฏฺฐิวิปตฺติยา วา อาชีววิปตฺติยา วา.\csromanspacer{} อิมานิ ตีณิ กุสลมูลานิ อนุวาทาธิกรณสฺส มูลํ.}

\prose{กตโม กาโย\footnote{กตโม จ กาโย (สฺยา, กํ)} อนุวาทาธิกรณสฺส มูลํ.\csromanspacer{} อิเธกจฺโจ ทุพฺพณฺโณ โหติ ทุทฺทสฺสิโก โอโกฏิมโก พหฺวาพาโธ กาโณ วา กุณี วา ขญฺโช วา ปกฺขหโต วา, เยน นํ อนุวทนฺติ, อยํ กาโย อนุวาทาธิกรณสฺส มูลํ.}

\prose{กตมา วาจา อนุวาทาธิกรณสฺส มูลํ.\csromanspacer{} อิเธกจฺโจ ทุพฺพโจ โหติ มมฺมโน เอฬคลวาโจ, ยาย นํ อนุวทนฺติ, อยํ วาจา อนุวาทาธิกรณสฺส มูลํ.}

\proseitem{218}{อาปตฺตาธิกรณสฺส กิํ มูลํ.\csromanspacer{} ฉ อาปตฺติสมุฏฺฐานา อาปตฺตาธิกรณสฺส มูลํ.\csromanspacer{} อตฺถาปตฺติ กายโต สมุฏฺฐาติ น วาจโต \csromanfolio{216}น จิตฺตโต, อตฺถาปตฺติ วาจโต สมุฏฺฐาติ น กายโต น จิตฺตโต, อตฺถาปตฺติ กายโต จ วาจโต จ สมุฏฺฐาติ น จิตฺตโต, อตฺถาปตฺติ กายโต จ จิตฺตโต จ สมุฏฺฐาติ น วาจโต, อตฺถาปตฺติ วาจโต จ จิตฺตโต จ สมุฏฺฐาติ น กายโต, อตฺถาปตฺติ กายโต จ วาจโต จ จิตฺตโต จ สมุฏฺฐาติ.\csromanspacer{} อิเม ฉ อาปตฺติสมุฏฺฐานา อาปตฺตาธิกรณสฺส มูลํ.}

\proseitem{219}{กิจฺจาธิกรณสฺส กิํ มูลํ.\csromanspacer{} กิจฺจาธิกรณสฺส เอกํ มูลํ สํโฆ.}

\proseitem{220}{วิวาทาธิกรณํ กุสลํ อกุสลํ อพฺยากตํ? วิวาทาธิกรณํ สิยา กุสลํ สิยา อกุสลํ สิยา อพฺยากตํ.\csromanspacer{} ตตฺถ กตมํ วิวาทาธิกรณํ กุสลํ.\csromanspacer{} อิธ ปน ภิกฺขเว ภิกฺขู กุสลจิตฺตา วิวทนฺติ ธมฺโมติ วา อธมฺโมติ วา ฯเปฯ ทุฏฺฐุลฺลา อาปตฺตีติ วา อทุฏฺฐุลฺลา อาปตฺตีติ วา.\csromanspacer{} ยํ ตตฺถ ภณฺฑนํ กลโห วิคฺคโห วิวาโท นานาวาโท อญฺญถาวาโท วิปจฺจตาย โวหาโร เมธคํ, อิทํ วุจฺจติ วิวาทาธิกรณํ กุสลํ.}

\prose{ตตฺถ กตมํ วิวาทาธิกรณํ อกุสลํ.\csromanspacer{} อิธ ปน ภิกฺขเว ภิกฺขู อกุสลจิตฺตา วิวทนฺติ ธมฺโมติ วา อธมฺโมติ วา ฯเปฯ ทุฏฺฐุลฺลา อาปตฺตีติ วา อทุฏฺฐุลฺลา อาปตฺตีติ วา.\csromanspacer{} ยํ ตตฺถ ภณฺฑนํ กลโห วิคฺคโห วิวาโท นานาวาโท อญฺญถาวาโท วิปจฺจตาย โวหาโร เมธคํ, อิทํ วุจฺจติ วิวาทาธิกรณํ อกุสลํ.}

\prose{ตตฺถ กตมํ วิวาทาธิกรณํ อพฺยากตํ.\csromanspacer{} อิธ ปน ภิกฺขเว ภิกฺขู อพฺยากตจิตฺตา วิวทนฺติ ธมฺโมติ วา อธมฺโมติ วา ฯเปฯ ทุฏฺฐุลฺลา อาปตฺตีติ วา อทุฏฺฐุลฺลา อาปตฺตีติ วา.\csromanspacer{} ยํ ตตฺถ ภณฺฑนํ กลโห วิคฺคโห วิวาโท นานาวาโท อญฺญถาวาโท วิปจฺจตาย โวหาโร เมธคํ, อิทํ วุจฺจติ วิวาทาธิกรณํ อพฺยากตํ.}

\proseitem{221}{อนุวาทาธิกรณํ กุสลํ อกุสลํ อพฺยากตํ? อนุวาทาธิกรณํ สิยา กุสลํ สิยา อกุสลํ สิยา อพฺยากตํ.\csromanspacer{} ตตฺถ กตมํ อนุวาทาธิกรณํ กุสลํ.\csromanspacer{} อิธ ปน ภิกฺขเว ภิกฺขู ภิกฺขุํ กุสลจิตฺตา อนุวทนฺติ สีลวิปตฺติยา วา อาจารวิปตฺติยา วา ทิฏฺฐิวิปตฺติยา วา อาชีววิปตฺติยา วา.\csromanspacer{} โย ตตฺถ อนุวาโท อนุวทนา อนุลฺลปนา \csromanfolio{217}อนุภณนา อนุสมฺปวงฺกตา อพฺภุสฺสหนตา อนุพลปฺปทานํ, อิทํ วุจฺจติ อนุวาทาธิกรณํ กุสลํ.}

\prose{ตตฺถ กตมํ อนุวาทาธิกรณํ อกุสลํ.\csromanspacer{} อิธ ปน ภิกฺขเว ภิกฺขู ภิกฺขุํ อกุสลจิตฺตา อนุวทนฺติ สีลวิปตฺติยา วา อาจารวิปตฺติยา วา ทิฏฺฐิวิปตฺติยา วา อาชีววิปตฺติยา วา.\csromanspacer{} โย ตตฺถ อนุวาโท อนุวทนา อนุลฺลปนา อนุภณนา อนุสมฺปวงฺกตา อพฺภุสฺสหนตา อนุพลปฺปทานํ, อิทํ วุจฺจติ อนุวาทาธิกรณํ อกุสลํ.}

\prose{ตตฺถ กตมํ อนุวาทาธิกรณํ อพฺยากตํ.\csromanspacer{} อิธ ปน ภิกฺขเว ภิกฺขู ภิกฺขุํ อพฺยากตจิตฺตา อนุวทนฺติ สีลวิปตฺติยา วา อาจารวิปตฺติยา วา ทิฏฺฐิวิปตฺติยา วา อาชีววิปตฺติยา วา.\csromanspacer{} โย ตตฺถ อนุวาโท อนุวทนา อนุลฺลปนา อนุภณนา อนุสมฺปวงฺกตา อพฺภุสฺสหนตา อนุพลปฺปทานํ, อิทํ วุจฺจติ อนุวาทาธิกรณํ อพฺยากตํ.}

\proseitem{222}{อาปตฺตาธิกรณํ กุสลํ\footnote{อิทํ ปทํ เกสุจิ นตฺถิ.} อกุสลํ อพฺยากตํ? อาปตฺตาธิกรณํ สิยา อกุสลํ สิยา อพฺยากตํ, นตฺถิ อาปตฺตาธิกรณํ กุสลํ.\csromanspacer{} ตตฺถ กตมํ อาปตฺตาธิกรณํ อกุสลํ.\csromanspacer{} ยํ ชานนฺโต สญฺชานนฺโต เจจฺจ อภิวิตริตฺวา วิติกฺกโม, อิทํ วุจฺจติ อาปตฺตาธิกรณํ อกุสลํ.}

\prose{ตตฺถ กตมํ อาปตฺตาธิกรณํ อพฺยากตํ.\csromanspacer{} ยํ อชานนฺโต อสญฺชานนฺโต อเจจฺจ อนภิวิตริตฺวา วีติกฺกโม, อิทํ วุจฺจติ อาปตฺตาธิกรณํ อพฺยากตํ.}

\proseitem{223}{กิจฺจาธิกรณํ กุสลํ อกุสลํ อพฺยากตํ? กิจฺจาธิกรณํ สิยา กุสลํ สิยา อกุสลํ สิยา อพฺยากตํ.\csromanspacer{} ตตฺถ กตมํ กิจฺจาธิกรณํ กุสลํ.\csromanspacer{} ยํ สํโฆ กุสลจิตฺโต กมฺมํ กโรติ อปโลกนกมฺมํ ญตฺติกมฺมํ ญตฺติทุติยกมฺมํ ญตฺติจตุตฺถกมฺมํ, อิทํ วุจฺจติ กิจฺจาธิกรณํ กุสลํ.}

\prose{ตตฺถ กตมํ กิจฺจาธิกรณํ อกุสลํ.\csromanspacer{} ยํ สํโฆ อกุสลจิตฺโต กมฺมํ กโรติ อปโลกนกมฺมํ ญตฺติกมฺมํ ญตฺติทุติยกมฺมํ ญตฺติจตุตฺถกมฺมํ, อิทํ วุจฺจติ กิจฺจาธิกรณํ อกุสลํ.}

\prose{\csromanfolio{218}ตตฺถ กตมํ กิจฺจาธิกรณํ อพฺยากตํ.\csromanspacer{} ยํ สํโฆ อพฺยากตจิตฺโต กมฺมํ กโรติ อปโลกนกมฺมํ ญตฺติกมฺมํ ญตฺติทุติยกมฺมํ ญตฺติจตุตฺถกมฺมํ, อิทํ วุจฺจติ กิจฺจาธิกรณํ อพฺยากตํ.}

\proseitem{224}{\csromansymbolfootnote{*}{วิ 5. 275 ปิฏฺเฐปิ.}วิวาโท วิวาทาธิกรณํ, วิวาโท โน อธิกรณํ, อธิกรณํ โน วิวาโท, อธิกรณญฺเจว วิวาโท จ.\csromanspacer{} สิยา วิวาโท วิวาทาธิกรณํ, สิยา วิวาโท โน อธิกรณํ, สิยา อธิกรณํ โน วิวาโท, สิยา อธิกรณญฺเจว วิวาโท จ.}

\prose{ตตฺถ กตโม วิวาโท วิวาทาธิกรณํ.\csromanspacer{} อิธ ปน ภิกฺขเว ภิกฺขู วิวทนฺติ ธมฺโมติ วา อธมฺโมติ วา ฯเปฯ ทุฏฺฐุลฺลา อาปตฺตีติ วา อทุฏฺฐุลฺลา อาปตฺตีติ วา.\csromanspacer{} ยํ ตตฺถ ภณฺฑนํ กลโห วิคฺคโห วิวาโท นานาวาโท อญฺญถาวาโท วิปจฺจตาย โวหาโร เมธคํ, อยํ วิวาโท วิวาทามิกรณํ.}

\prose{ตตฺถ กตโม วิวาโท โน อธิกรณํ.\csromanspacer{} มาตาปิ ปุตฺเตน วิวทติ, ปุตฺโตปิ มาตรา วิวทติ, ปิตาปิ ปุตฺเตน วิวทติ, ปุตฺโตปิ ปิตรา วิวทติ, ภาตาปิ ภาตรา วิวทติ, ภาตาปิ ภคินิยา วิวทติ, ภคินีปิ ภาตรา วิวทติ, สหาโยปิ สหาเยน วิวทติ, อยํ วิวาโท โน อธิกรณํ.}

\prose{ตตฺถ กตมํ อธิกรณํ โน วิวาโท.\csromanspacer{} อนุวาทาธิกรณํ อาปตฺตาธิกรณํ กิจฺจาธิกรณํ, อิทํ อธิกรณํ โน วิวาโท.}

\prose{ตตฺถ กตมํ อธิกรณญฺเจว วิวาโท จ.\csromanspacer{} วิวาทาธิกรณํ อธิกรณญฺเจว วิวาโท จ.}

\proseitem{225}{\csromansymbolfootnote{+}{วิ 5. 276 ปิฏฺเฐปิ.}อนุวาโท อนุวาทาธิกรณํ, อนุวาโท โน อธิกรณํ, อธิกรณํ โน อนุวาโท, อธิกรณญฺเจว อนุวาโท จ.\csromanspacer{} สิยา อนุวาโท อนุวาทาธิกรณํ, สิยา อนุวาโท โน อธิกรณํ, สิยา อธิกรณํ โน อนุวาโท, สิยา อธิกรณญฺเจว อนุวาโท จ.}

\prose{ตตฺถ กตโม อนุวาโท อนุวาทาธิกรณํ.\csromanspacer{} อิธ ปน ภิกฺขเว ภิกฺขู ภิกฺขุํ อนุวทนฺติ สีลวิปตฺติยา วา อาจารวิปตฺติยา วา ทิฏฺฐิวิปตฺติยา วา อาชีววิปตฺติยา วา.\csromanspacer{} โย ตตฺถ อนุวาโท อนุวทนา อนุลฺลปนา \csromanfolio{219}อนุภณนา อนุสมฺปวงฺกตา อพฺภุสฺสหนตา อนุพลปฺปทานํ, อยํ อนุวาโท อนุวาทาธิกรณํ.}

\prose{ตตฺถ กตโม อนุวาโท โน อธิกรณํ.\csromanspacer{} มาตาปิ ปุตฺตํ อนุวทติ, ปุตฺโตปิ มาตรํ อนุวทติ, ปิตาปิ ปุตฺตํ อนุวทติ, ปุตฺโตปิ ปิตรํ อนุวทติ, ภาตาปิ ภาตรํ อนุวทติ, ภาตาปิ ภคินิํ อนุวทติ, ภคินีปิ ภาตรํ อนุวทติ, สหาโยปิ สหายํ อนุวทติ, อยํ อนุวาโท โน อธิกรณํ.}

\prose{ตตฺถ กตมํ อธิกรณํ โน อนุวาโท.\csromanspacer{} อาปตฺตาธิกรณํ กิจฺจาธิกรณํ วิเวทาธิกรณํ, อิทํ อธิกรณํ โน อนุวาโท.}

\prose{ตตฺถ กตมํ อธิกรณญฺเจว อนุวาโท จ.\csromanspacer{} อนุวาทาธิกรณํ อธิกรณญฺเจว อนุวาโท จ.}

\proseitem{226}{\csromansymbolfootnote{*}{วิ 5. 277 ปิฏฺเฐปิ.}อาปตฺติ อาปตฺตาธิกรณํ, อาปตฺติ โน อธิกรณํ, อธิกรณํ โน อาปตฺติ, อธิกรณญฺเจว อาปตฺติ จ.\csromanspacer{} สิยา อาปตฺติ อาปตฺตาธิกรณํ, สิยา อาปตฺติ โน อธิกรณํ, สิยา อธิกรณํ โน อาปตฺติ, สิยา อธิกรณญฺเจว อาปตฺติ จ.\csromanspacer{} ตตฺถ กตมํ อาปตฺติ อาปตฺตาธิกรณํ.\csromanspacer{} ปญฺจปิ อาปตฺติกฺขนฺธา อาปตฺตาธิกรณํ, สตฺตปิ อาปตฺถิกฺขนฺธา อาปตฺตาธิกรณํ, อยํ อาปตฺติ อาปตฺตาธิกรณํ.}

\prose{ตตฺถ กตมํ อาปตฺติ โน อธิกรณํ.\csromanspacer{} โสตาปตฺติ สมาปตฺติ, อยํ อาปตฺติ โน อธิกรณํ.}

\prose{ตตฺถ กตมํ อธิกรณํ โน อาปตฺติ.\csromanspacer{} กิจฺจาธิกรณํ วิวาทาธิกรณํ อนุวาทาธิกรณํ, อิทํ อธิกรณํ โน อาปตฺติ.}

\prose{ตตฺถ กตมํ อธิกรณญฺเจว อาปตฺติ จ.\csromanspacer{} อาปตฺตาธิกรณํ อธิกรณญฺเจว อาปตฺติ จ.}

\proseitem{227}{\csromansymbolmark{*}กิจฺจํ กิจฺจาธิกรณํ, กิจฺจํ โน อธิกรณํ, อธิกรณํ โน กิจฺจํ, อธิกรณญฺเจว กิจฺจญฺจ.\csromanspacer{} สิยา กิจฺจํ กิจฺจาธิกรณํ, สิยา กิจฺจํ โน อธิกรณํ, สิยา อธิกรณํ โน กิจฺจํ, สิยา อธิกรณญฺเจว กิจฺจญฺจ.}

\prose{\csromanfolio{220}ตตฺถ กตมํ กิจฺจํ กิจฺจาธิกรณํ.\csromanspacer{} ยา สํฆสฺส กิจฺจยตา กรณียตา อปโลกนกมฺมํ ญตฺติกมฺมํ ญตฺติทุติยกมฺมํ ญตฺติจตุตฺถกมฺมํ, อิทํ กิจฺจํ กิจฺจาธิกรณํ.}

\prose{ตตฺถ กตมํ กิจฺจํ โน อธิกรณํ.\csromanspacer{} อาจริยกิจฺจํ อุปชฺฌายกิจฺจํ สมานุปชฺฌายกิจฺจํ สมานาจริยกิจฺจํ, อิทํ กิจฺจํ โน อธิกรณํ.}

\prose{ตตฺถ กตมํ อธิกรณํ โน กิจฺจํ, วิวาทาธิกรณํ อนุวาทาธิกรณํ อาปตฺตาธิกรณํ, อิทํ อธิกรณํ โน กิจฺจํ.}

\prose{ตตฺถ กตมํ อธิกรณญฺเจว กิจฺจญฺจ.\csromanspacer{} กิจฺจาธิกรณํ อธิกรณญฺเจว กิจฺจญฺจ.}

\csromanmatikaanchor{mtk.116}
\csromantocmarkref{section}{9. อธิกรณวูปสมนสมถ}{mtk.116}
\bookmark[dest={mtk.116},level=1]{9. อธิกรณวูปสมนสมถ}
\csromanheader{1}{9. อธิกรณวูปสมนสมถ}
\csromanmatikaanchor{mtk.117}
\csromantocmarkref{subsection}{สมฺมุขาวินย}{mtk.117}
\bookmark[dest={mtk.117},level=2]{สมฺมุขาวินย}
\csromanheader{2}{\textbf{สมฺมุขาวินย}}
\proseitem{228}{\csromansymbolfootnote{*}{วิ 5. 183, 195 ปิฏฺฐาทีสุปิ.}วิวาทาธิกรณํ กติหิ สมเถหิ สมฺมติ.\csromanspacer{} วิวาทาธิกรณํ ทฺวีหิ สมเถหิ สมฺมติ สมฺมุขาวินเยน จ เยภุยฺยสิกาย จ.\csromanspacer{} สิยา วิวาทาธิกรณํ เอกํ สมถํ อนาคมฺม เยภุยฺยสิกํ เอเกน สมเถน สเมยฺย สมฺมุขาวินเยนาติ? สิยาติสฺส วจนียํ.\csromanspacer{} ยถา กถํ วิย.\csromanspacer{} อิธ ปน ภิกฺขเว ภิกฺขู วิวทนฺติ ธมฺโมติ วา อธมฺโมติ วา, วินโยติ วา อวินโยติ วา, ภาสิตํ ลปิตํ ตถาคเตนาติ วา อภาสิตํ อลปิตํ ตถาคเตนาติ วา, อาจิณฺณํ ตถาคเตนาติ วา อนาจิณฺณํ ตถาคเตนาติ วา, ปญฺญตฺตํ ตถาคเตนาติ วา อปญฺญตฺตํ ตถาคเตนาติ วา, อาปตฺตีติ วา อนาปตฺตีติ วา, ลหุกา อาปตฺตีติ วา ครุกา อาปตฺตีติ วา, สาวเสสา อาปตฺตีติ วา อนวเสสา อาปตฺตีติ วา, ทุฏฺฐุลฺลา อาปตฺตีติ วา อทุฏฺฐุลฺลา อาปตฺตีติ วา.}

\prose{เต เจ ภิกฺขเว ภิกฺขู สกฺโกนฺติ ตํ อธิกรณํ วูปสเมตุํ, อิทํ วุจฺจติ ภิกฺขเว อธิกรณํ วุปสนฺตํ.\csromanspacer{} เกน วูปสนฺตํ, สมฺมุขาวินเยน.\csromanspacer{} กิญฺจ ตตฺถ สมฺมุขาวินยสฺมิํ, สํฆสมฺมุขตา ธมฺมสมฺมุขตา วินยสมฺมุขตา ปุคฺคลสมฺมุขตา.\csromanspacer{} กา จ ตตฺถ สํฆสมฺมุขตา, ยาวติกา ภิกฺขู กมฺมปฺปตฺตา, เต อาคตา โหนฺติ, ฉนฺทารหานํ ฉนฺโท อาหโฏ โหติ, สมฺมุขีภูตา น ปฏิกฺโกสนฺติ, อยํ ตตฺถ สํฆสมฺมุขตา.\csromanspacer{} กา \csromanfolio{221}จ ตตฺถ ธมฺมสมฺมุขตา วินยสมฺมุขตา, เยน ธมฺเมน เยน วินเยน เยน สตฺถุสาสเนน ตํ อธิกรณํ วูปสมฺมติ, อยํ ตตฺถ ธมฺมสมฺมุขตา วินยสมฺมุขตา.\csromanspacer{} กา จ ตตฺถ ปุคฺคลสมฺมุขตา, โย จ วิวทติ, เยน จ วิวทติ, อุโภ อตฺถปจฺจตฺถิกา สมฺมุขีภูตา โหนฺติ, อยํ ตตฺถ ปุคฺคลสมฺมุขตา.\csromanspacer{} เอวํ วูปสนฺตํ เจ ภิกฺขเว อธิกรณํ การโก อุกฺโกเฏติ, อุกฺโกฏนกํ ปาจิตฺติยํ.\csromanspacer{} ฉนฺททายโก ขียติ, ขียนกํ ปาจิตฺติยํ.}

\proseitem{229}{เต เจ ภิกฺขเว ภิกฺขู น สกฺโกนฺติ ตํ อธิกรณํ ตสฺมิํ อาวาเส วูปสเมตุํ, เตหิ ภิกฺขเว ภิกฺขูหิ ยสฺมิํ อาวาเส สมฺพหุลา\footnote{พหุตรา (สี, สฺยา)} ภิกฺขู, โส อาวาโส คนฺตพฺโพ.\csromanspacer{} เต เจ ภิกฺขเว ภิกฺขู ตํ อาวาสํ คจฺฉนฺตา อนฺตรามคฺเค สกฺโกนฺติ ตํ อธิกรณํ วูปสเมตุํ.\csromanspacer{} อิทํ วุจฺจติ ภิกฺขเว อธิกรณํ วูปสนฺตํ.\csromanspacer{} เกน วูปสนฺตํ, สมฺมุขาวินเยน.\csromanspacer{} กิญฺจ ตตฺถ สมฺมุขาวินยสฺมิํ, สํฆสมฺมุขตา ธมฺมสมฺมุขตา วินยสมฺมุขตา ปุคฺคลสมฺมุขตา.\csromanspacer{} กา จ ตตฺถ สํฆสมฺมุขตา, ยาวติกา ภิกฺขู กมฺมปฺปตฺตา, เต อาคตา โหนฺติ, ฉนฺทารหานํ ฉนฺโท อาหโฏ โหติ, สมฺมุขีภูตา น ปฏิกฺโกสนฺติ, อยํ ตตฺถ สํฆสมฺมุขตา.\csromanspacer{} กา จ ตตฺถ ธมฺมสมฺมุขตา วินยสมฺมุขตา, เยน ธมฺเมน เยน วินเยน เยน สตฺถุสาสเนน ตํ อธิกรณํ วูปสมฺมติ, อยํ ตตฺถ ธมฺมสมฺมุขตา วินยสมฺมุขตา.\csromanspacer{} กา จ ตตฺถ ปุคฺคลสมฺมุขตา, โย จ วิวทติ, เยน จ วิวทติ, อุโภ อตฺถปจฺจตฺถิกา สมฺมุขีภูตา โหนฺติ, อยํ ตตฺถ ปุคฺคลสมฺมุขตา.\csromanspacer{} เอวํ วูปสนฺตํ เจ ภิกฺขเว อธิกรณํ การโก อุกฺโกเฏติ, อุกฺโกฏนกํ ปาจิตฺติยํ.\csromanspacer{} ฉนฺททายโก ขียติ, ขียนกํ ปาจิตฺติยํ.}

\proseitem{230}{เต เจ ภิกฺขเว ภิกฺขู ตํ อาวาสํ อาคจฺฉนฺตา อนฺตรามคฺเค น สกฺโกนฺติ ตํ อธิกรณํ วูปสเมตุํ, เตหิ ภิกฺขเว ภิกฺขูหิ ตํ อาวาสํ คนฺตฺวา อาวาสิกา ภิกฺขู เอวมสฺสุ วจนียา “อิทํ โข อาวุโส อธิกรณํ เอวํ ชาตํ เอวํ สมุปฺปนฺนํ, สาธายสฺมนฺตา อิมํ อธิกรณํ วูปสเมนฺตุ ธมฺเมน วินเยน สตฺถุสาสเนน, ยถยิทํ อธิกรณํ สุวูปสนฺตํ อสฺสา”ติ.}

\prose{\csromanfolio{222}สเจ ภิกฺขเว อาวาสิกา ภิกฺขู วุฑฺฒตรา โหนฺติ, อาคนฺตุกา ภิกฺขู นวกตรา, เตหิ ภิกฺขเว อาวาสิเกหิ ภิกฺขูหิ อาคนฺตุกา ภิกฺขู เอวมสฺสุ วจนียา “อิงฺฆ ตุเมฺห อายสฺมนฺโต มุหุตฺตํ เอกมนฺตํ โหถ, ยาว มยํ มนฺเตมา”ติ.\csromanspacer{} สเจ ปน ภิกฺขเว อาวาสิกา ภิกฺขู นวกตรา โหนฺติ, อาคนฺตุกา ภิกฺขู วุฑฺฒุตรา, เตหิ ภิกฺขเว อาวาสิเกหิ ภิกฺขูหิ อาคนฺตุกา ภิกฺขู เอวมสฺสุ วจนียา “เตน หิ ตุเมฺห อายสฺมนฺโต มุหุตฺตํ อิเธว ตาว โหถ, ยาว มยํ มนฺเตมา”ติ.}

\prose{สเจ ปน ภิกฺขเว อาวาสิกานํ ภิกฺขูนํ มนฺตยมานานํ เอวํ โหติ “น มยํ สกฺโกม อิมํ อธิกรณํ วูปสเมตุํ ธมฺเมน วินเยน สตฺถุสาสเนนา”ติ, น ตํ อธิกรณํ อาวาสิเกหิ ภิกฺขูหิ สมฺปฏิจฺฉิตพฺพํ.\csromanspacer{} สเจ ปน ภิกฺขเว อาวาสิกานํ ภิกฺขูนํ มนฺตยมานานํ เอวํ โหติ “สกฺโกม มยํ อิมํ อธิกรณํ วูปสเมตุํ ธมฺเมน วินเยน สตฺถุสาสเนนา”ติ, เตหิ ภิกฺขเว อาวาสิเกหิ ภิกฺขูหิ อาคนฺตุกา ภิกฺขู เอวมสฺสุ วจนียา “สเจ ตุเมฺห อายสฺมนฺโต อมฺหากํ อิมํ อธิกรณํ ยถาชาตํ ยถาสมุปฺปนฺนํ อาโรเจสฺสถ, ยถา จ มยํ อิมํ อธิกรณํ วูปสเมสฺสาม ธมฺเมน วินเยน สตฺถุสาสเนน, ตถาสุ วูปสนฺตํ ภวิสฺสติ, เอวํ มยํ อิมํ อธิกรณํ สมฺปฏิจฺฉิสฺสาม.\csromanspacer{} โน เจ ตุเมฺห อายสฺมนฺโต อมฺหากํ อิมํ อธิกรณํ ยถาชาตํ ยถาสมุปฺปนฺนํ อาโรเจสฺสถ, ยถา จ มยํ อิมํ อธิกรณํ วูปสเมสฺสาม ธมฺเมน วินเยน สตฺถุสาสเนน, ตถา น สุวูปสนฺตํ ภวิสฺสติ, น มยํ อิมํ อธิกรณํ สมฺปฏิจฺฉิสฺสามา”ติ.\csromanspacer{} เอวํ สุปริคฺคหิตํ โข ภิกฺขเว กตฺวา อาวาสิเกหิ ภิกฺขูหิ ตํ อธิกรณํ สมฺปฏิจฺฉิตพฺพํ.}

\prose{เตหิ ภิกฺขเว อาคนฺตุเกหิ ภิกฺขูหิ อาวาสิกา ภิกฺขู เอวมสฺสุ วจนียา “ยถาชาตํ ยถาสมุปฺปนฺนํ มยํ อิมํ อธิกรณํ อายสฺมนฺตานํ อาโรเจสฺสาม.\csromanspacer{} สเจ อายสฺมนฺตา สกฺโกนฺติ เอตฺตเกน วา เอตฺตเกน วา อนฺตเรน อิมํ อธิกรณํ วูปสเมตุํ ธมฺเมน วินเยน สตฺถุสาสเนน, ตถา สุวูปสนฺตํ ภวิสฺสติ, เอวํ มยํ อิมํ อธิกรณํ อายสฺมนฺตานํ นิยฺยาเทสฺสาม.\csromanspacer{} โน เจ อายสฺมนฺตา สกฺโกนฺติ เอตฺตเกน วา เอตฺตเกน วา อนฺตเรน อิมํ อธิกรณํ วูปสเมตุํ ธมฺเมน วินเยน สตฺถุสาสเนน, ตถา น สุวูปสนฺตํ ภวิสฺสติ, น มยํ อิมํ อธิกรณํ อายสฺมนฺตานํ \csromanfolio{223}นิยฺยาเทสฺสาม, มยเมว อิมสฺส อธิกรณสฺส สามิโน ภวิสฺสามา”ติ.\csromanspacer{} เอวํ สุปริคฺคหิตํ โข ภิกฺขเว กตฺวา อาคนฺตุเกหิ ภิกฺขูหิ ตํ อธิกรณํ อาวาสิกานํ ภิกฺขูนํ นิยฺยาเทตพฺพํ.}

\prose{เต เจ ภิกฺขเว ภิกฺขู สกฺโกนฺติ ตํ อธิกรณํ วูปสเมตุํ, อิทํ วุจฺจติ ภิกฺขเว อธิกรณํ วูปสนฺตํ.\csromanspacer{} เกน วูปสนฺตํ, สมฺมุขาวินเยน.\csromanspacer{} กิญฺจ ตตฺถ สมฺมุขาวินยสฺมิํ, สํฆสมฺมุขตา ธมฺมสมฺมุขตา วินยสมฺมุขตา ปุคฺคลสมฺมุขตา ฯเปฯ.\csromanspacer{} เอวํ วูปสนฺตํ เจ ภิกฺขเว อธิกรณํ การโก อุกฺโกเฏติ, อุกฺโกฏนกํ ปาจิตฺติยํ.\csromanspacer{} ฉนฺททายโก ขียติ, ขียนกํ ปาจิตฺติยํ.}

\csromanmatikaanchor{mtk.118}
\csromantocmarkref{subsection}{อุพฺพาหิกายวูปสมน}{mtk.118}
\bookmark[dest={mtk.118},level=2]{อุพฺพาหิกายวูปสมน}
\csromanheader{2}{\textbf{อุพฺพาหิกายวูปสมน}}
\proseitem{231}{เตหิ เจ ภิกฺขเว ภิกฺขูหิ ตสฺมิํ อธิกรเณ วินิจฺฉิยมาเน อนนฺตานิ\footnote{อนคฺคานิ (สี)} เจว ภสฺสานิ ชายนฺติ, น เจกสฺส\footnote{น เจตสฺส (ก)} ภาสิตสฺส อตฺโถ วิญฺญายติ, อนุชานามิ ภิกฺขเว เอวรูปํ อธิกรณํ อุพฺพาหิกาย วูปสเมตุํ.\csromanspacer{} ทสหงฺเคหิ สมนฺนาคโต ภิกฺขุ อุพฺพาหิกาย สมฺมนฺนิตพฺโพ.\csromanspacer{} สีลวา โหติ ปาติโมกฺขสํวรสํวุโต วิหรติ อาจารโคจรสมฺปนฺโน อณุมตฺเตสุ วชฺเชสุ ภยทสฺสาวี สมาทาย สิกฺขติ สิกฺขาปเทสุ, พหุสฺสุโต โหติ สุตธโร สุตสนฺนิจโย, เย เต ธมฺมา อาทิกลฺยาณา มชฺเฌกลฺยาณา ปริโยสานกลฺยาณา สาตฺถํ สพฺยญฺชนํ เกวลปริปุณฺณํ ปริสุทฺธํ พฺรหฺมจริยํ อภิวทนฺติ, ตถารูปสฺส ธมฺมา พหุสฺสุตา โหนฺติ ธาตา\footnote{ธตา (สี, สฺยา)} วจสา ปริจิตา มนสานุเปกฺขิตา ทิฏฺฐิยา สุปฺปฏิวิทฺธา, อุภยานิ โข ปนสฺส ปาติโมกฺขานิ วิตฺถาเรน สฺวาคตานิ โหนฺติ สุวิภตฺตานิ สุปฺปวตฺตีนิ สุวินิจฺฉิตานิ สุตฺตโส อนุพฺยญฺชนโส, วินเย โข ปน ฐิโต\footnote{เฉโก (ก)} โหติ อสํหีโร, ปฏิพโล โหติ อุโภ อตฺถปจฺจตฺถิเก อสฺสาเสตุํ สญฺญาเปตุํ นิชฺฌาเปตุํ เปกฺเขตุํ ปสาเทตุํ, อธิกรณสมุปฺปาทวูปสมนกุสโล โหติ, อธิกรณํ ชานาติ, อธิกรณสมุทยํ ชานาติ, อธิกรณนิโรธํ ชานาติ, อธิกรณนิโรธคามินิปฏิปทํ ชานาติ.\csromanspacer{} อนุชานามิ ภิกฺขเว อิเมหิ ทสหงฺเคหิ สมนฺนาคตํ ภิกฺขุํ อุพฺพาหิกาย สมฺมนฺนิตุํ. \csromanfolio{224}เอวญฺจ ปน ภิกฺขเว สมฺมนฺนิตพฺโพ, ปฐมํ ภิกฺขุ ยาจิตพฺโพ, ยาจิตฺวา พฺยตฺเตน ภิกฺขุนา ปฏิพเลน สํโฆ ญาเปตพฺโพ\csromandash{}}

\proseitem{232}{“สุณาตุ เม ภนฺเต สํโฆ, อมฺหากํ อิมสฺมิํ อธิกรเณ วินิจฺฉิยมาเน อนนฺตานิ เจว ภสฺสานิ ชายนฺติ, น เจกสฺส ภาสิตสฺส อตฺโถ วิญฺญายติ, ยทิ สํฆสฺส ปตฺตกลฺลํ, สํโฆ อิตฺถนฺนามญฺจ อิตฺถนฺนามญฺจ ภิกฺขุํ สมฺมนฺเนยฺย อุพฺพาหิกาย อิมํ อธิกรณํ วูปสเมตุํ, เอสา ญตฺติ.}

\prose{สุณาตุ เม ภนฺเต สํโฆ, อมฺหากํ อิมสฺมิํ อธิกรเณ วินิจฺฉิยมาเน อนนฺตานิ เจว ภสฺสานิ ชายนฺติ, น เจกสฺส ภาสิตสฺส อตฺโถ วิญฺญายติ, สํโฆ อิตฺถนฺนามญฺจ อิตฺถนฺนามญฺจ ภิกฺขุํ สมฺมนฺนติ อุพฺพาหิกาย อิมํ อธิกรณํ วูปสเมตุํ, ยสฺสายสฺมโต ขมติ อิตฺถนฺนามสฺส จ อิตฺถนฺนามสฺส จ ภิกฺขูโน สมฺมุติ อุพฺพาหิกาย อิมํ อธิกรณํ วูปสเมตุํ, โส ตุณฺหสฺส, ยสฺส นกฺขมติ, โส ภาเสยฺย.}

\prose{สมฺมโต สํเฆน อิตฺถนฺนาโม จ อิตฺถนฺนาโม จ ภิกฺขุ อุพฺพาหิกาย อิมํ อธิกรณํ วูปสเมตุํ, ขมติ สํฆสฺส, ตสฺมา ตุณฺหี, เอวเมตํ ธารยามี”ติ.}

\prose{เต เจ ภิกฺขเว ภิกฺขู สกฺโกนฺติ ตํ อธิกรณํ อุพฺพาหิกาย วูปสเมตุํ, อิทํ วุจฺจติ ภิกฺขเว อธิกรณํ วูปสนฺตํ.\csromanspacer{} เกน วูปสนฺตํ, สมฺมุขาวินเยน.\csromanspacer{} กิญฺจ ตตฺถ สมฺมุขาวินยสฺมิํ, ธมฺมสมฺมุขตา วินยสมฺมุขตา ปุคฺคลสมฺมุขตา ฯเปฯ.\csromanspacer{} เอวํ วูปสนฺตํ เจ ภิกฺขเว อธิกรณํ การโก อุกฺโกเฏติ, อุกฺโกฏนกํ ปาจิตฺติยํ.}

\proseitem{233}{เตหิ เจ ภิกฺขเว ภิกฺขูหิ ตสฺมิํ อธิกรเณ วินิจฺฉิยมาเน ตตฺราสฺส ภิกฺขุ ธมฺมกถิโก, ตสฺส เนว สุตฺตํ อาคตํ โหติ, โน สุตฺตวิภงฺโค, โส อตฺถํ อสลฺลกฺเขนฺโต พฺยญฺชนจฺฉายาย อตฺถํ ปฏิพาหติ.\csromanspacer{} พฺยตฺเตน ภิกฺขุนา ปฏิพเลน เต ภิกฺขู ญาเปตพฺพา\csromandash{}}

\prose{“สุณนฺตุ เม อายสฺมนฺตา, อยํ อิตฺถนฺนาโม ภิกฺขู ธมฺมกถิโก, อิมสฺส เนว สุตฺตํ อาคตํ โหติ, โน สุตฺตวิภงฺโค, โส อตฺถํ อสลฺลกฺเขนฺโต พฺยญฺชนจฺฉายาย อตฺถํ ปฏิพาหติ, ยทายสฺมนฺตานํ \csromanfolio{225}ปตฺตกลฺลํ, อิตฺถนฺนามํ ภิกฺขุํ วุฏฺฐาเปตฺวา อวเสสา อิมํ อธิกรณํ วูปสเมยฺยามา”ติ.}

\prose{เต เจ ภิกฺขเว ภิกฺขู ตํ ภิกฺขุํ วุฏฺฐาเปตฺวา สกฺโกนฺติ ตํ อธิกรณํ วูปสเมตุํ, อิทํ วุจฺจติ ภิกฺขเว อธิกรณํ วูปสนฺตํ.\csromanspacer{} เกน วูปสนฺตํ, สมฺมุขาวินเยน.\csromanspacer{} กิญฺจ ตตฺถ สมฺมุขาวินยสฺมิํ, ธมฺมสมฺมุขตา วินยสมฺมุขตา ปุคฺคลสมฺมุขตา ฯเปฯ.\csromanspacer{} เอวํ วูปสนฺตํ เจ ภิกฺขเว อธิกรณํ กรโก อุกฺโกเฏติ, อุกฺโกฏนกํ ปาจิตฺติยํ.}

\prose{เตหิ เจ ภิกฺขเว ภิกฺขูหิ ตสฺมิํ อธิกรเณ วินิจฺฉิยมาเน ตตฺราสฺส ภิกฺขุ ธมฺมกถิโก, ตสฺส สุตฺตํ หิ โข อาคตํ โหติ, โน สุตฺตวิภงฺโค, โส อตฺถํ อสลฺลกฺเขนฺโต พฺยญฺชนจฺฉายาย อตฺถํ ปฏิพาหติ.\csromanspacer{} พฺยตฺเตน ภิกฺขูนา ปฏิพเลน เต ภิกฺขู ญาเปตพฺพา\csromandash{}}

\prose{“สุณนฺตุ เม อายสฺมนฺตา, อยํ อิตฺถนฺนาโม ภิกฺขุ ธมฺมกถิโก, อิมสฺส สุตฺตํ หิ โข อาคตํ โหติ, โน สุตฺตวิภงฺโค, โส อตฺถํ อสลฺลกฺเขนฺโต พฺยญฺชนจฺฉายาย อตฺถํ ปฏิพาหติ, ยทายสฺมนฺตานํ ปตฺตกลฺลํ, อิตฺถนฺนามํ ภิกฺขุํ วุฏฺฐาเปตฺวา อวเสสา อิมํ อธิกรณํ วูปสเมยฺยามา”ติ.}

\prose{เต เจ ภิกฺขเว ภิกฺขู ตํ ภิกฺขุํ วุฏฺฐาเปตฺวา สกฺโกนฺติ ตํ อธิกรณํ วูปสเมตุํ, อิทํ วุจฺจติ ภิกฺขเว อธิกรณํ วูปสนฺตํ.\csromanspacer{} เกน วูปสนฺตํ, สมฺมุขาวินเยน.\csromanspacer{} กิญฺจ ตตฺถ สมฺมุขาวินยสฺมิํ, ธมฺมสมฺมุขตา วินยสมฺมุขตา ปุคฺคลสมฺมุขตา ฯเปฯ.\csromanspacer{} เอวํ วูปสนฺตํ เจ ภิกฺขเว อธิกรณํ การโก อุกฺโกเฏติ, อุกฺโกฏนกํ ปาจิตฺติยํ.}

\csromanmatikaanchor{mtk.119}
\csromantocmarkref{subsection}{เยภุยฺยสิกาวินย}{mtk.119}
\bookmark[dest={mtk.119},level=2]{เยภุยฺยสิกาวินย}
\csromanheader{2}{\textbf{เยภุยฺยสิกาวินย}}
\proseitem{234}{เต เจ ภิกฺขเว ภิกฺขู น สกฺโกนฺติ ตํ อธิกรณํ อุพฺพาหิกาย วูปสเมตุํ.\csromanspacer{} เตหิ ภิกฺขเว ภิกฺขูหิ ตํ อธิกรณํ สํฆสฺส นิยฺยาเทตพฺพํ “น มยํ\footnote{น จ มยํ (ก)} ภนฺเต สกฺโกม อิมํ อธิกรณํ อุพฺพาหิกาย วูปสเมตุํ, สํโฆว อิมํ อธิกรณํ วูปสเมตู”ติ.\csromanspacer{} อนุชานามิ ภิกฺขเว เอวรูปํ อธิกรณํ เยภุยฺยสิกาย วูปสเมตุํ.\csromanspacer{} ปญฺจหงฺเคหิ สมนฺนาคโต ภิกฺขุ สลากคฺคาหาปโก สมฺมนฺนิตพฺโพ, โย น ฉนฺทาคติํ คจฺเฉยฺย. \csromanfolio{226}น โทสาคติํ คจฺเฉยฺย, น โมหาคติํ คจฺเฉยฺย, น ภยาคติํ คจฺเฉยฺย, คหิตาคหิตญฺจ ชาเนยฺย ฯเปฯ.\csromanspacer{} เอวญฺจ ปน ภิกฺขเว สมฺมนฺนิตพฺโพ, ปฐมํ ภิกฺขุ ยาจิตพฺโพ, ยาจิตฺวา พฺยตฺเตน ภิกฺขุนา ปฏิพเลน สํโฆ ญาเปตพฺโพ\csromandash{}}

\prose{“สุณาตุ เม ภนฺเต สํโฆ, ยทิ สํฆสฺส ปตฺตกลฺลํ, สํโฆ อิตฺถนฺนามํ ภิกฺขุํ สลากคฺคาหาปกํ สมฺมนฺเนยฺย, เอสา ญตฺติ.}

\prose{สุณาตุ เม ภนฺเต สํโฆ, สํโฆ อิตฺถนฺนามํ ภิกฺขุํ สลากคฺคาหาปกํ สมฺมนฺนติ, ยสฺสายสฺมโต ขมติ อิตฺถนฺนามสฺส ภิกฺขุโน สลากคฺคาหาปกสฺส สมฺมุติ, โส ตุณฺหสฺส, ยสฺส นกฺขมติ, โส ภาเสยฺย.}

\prose{สมฺมโต สํเฆน อิตฺถนฺนาโม ภิกฺขุ สลากคฺคาหาปโก, ขมติ สํฆสฺส, ตสฺมา ตุณฺหี, เอวเมตํ ธารยามี”ติ.}

\prose{เตน สลากคฺคาหาปเกน ภิกฺขุนา สลากา คาเหตพฺพา, ยถา พหุตรา ภิกฺขู ธมฺมวาทิโน วทนฺติ, ตถา ตํ อธิกรณํ วูปสเมตพฺพํ, อิทํ วุจฺจติ ภิกฺขเว อธิกรณํ วูปสนฺตํ.\csromanspacer{} เกน วูปสนฺตํ, สมฺมุขาวินเยน จ เยภุยฺยสิกาย จ.\csromanspacer{} กิญฺจ ตตฺถ สมฺมุขาวินยสฺมิํ, สํฆสมฺมุขตา ธมฺมสมฺมุขตา วินยสมฺมุขตา ปุคฺคลสมฺมุขตา.\csromanspacer{} กา จ ตตฺถ สํฆสมฺมุขตา, ยาวติกา ภิกฺขู กมฺมปฺปตฺตา, เต อาคตา โหนฺติ, ฉนฺทารหานํ ฉนฺโท อาหโฏ โหติ, สมฺมุขีภูตา น ปฏิกฺโกสนฺติ, อยํ ตตฺถ สํฆสมฺมุขตา.\csromanspacer{} กา จ ตตฺถ ธมฺมสมฺมุขตา วินยสมฺมุขตา, เยน ธมฺเมน เยน วินเยน เยน สตฺถุสาสเนน ตํ อธิกรณํ วูปสมฺมติ, อยํ ตตฺถ ธมฺมสมฺมุขตา วินยสมฺมุขตา.\csromanspacer{} กา จ ตตฺถ ปุคฺคลสมฺมุขตา, โย จ วิวทติ, เยน จ วิวทติ, อุโภ อตฺถปจฺจตฺถิกา สมฺมุขีภูตา โหนฺติ, อยํ ตตฺถ ปุคฺคลสมฺมุขตา.\csromanspacer{} กา จ ตตฺถ เยภุยฺยสิกาย, ยา เยภุยฺยสิกากมฺมสฺส กิริยา กรณํ อุปคมนํ อชฺฌุปคมนํ อธิวาสนา อปฺปฏิกฺโกสนา, อยํ ตตฺถ เยภุยฺยสิกาย.\csromanspacer{} เอวํ วูปสนฺตํ เจ ภิกฺขเว อธิกรณํ การโก อุกฺโกเฏติ, อุกฺโกฏนกํ ปาจิตฺติยํ.}

\csromangathameasure{ฉนฺททายโก ขียติ, ขียนกํ ปาจิตฺติยนฺติ.}
\csromangathagroup{\csromangathastanza{ฉนฺททายโก ขียติ, ขียนกํ ปาจิตฺติยนฺติ.}}

\csromanmatikaanchor{mtk.120}
\csromantocmarkref{subsection}{ติวิธสลากคฺคาห}{mtk.120}
\bookmark[dest={mtk.120},level=2]{ติวิธสลากคฺคาห}
\csromanheader{2}{\textbf{ติวิธสลากคฺคาห}}
\proseitem{253}{\csromanfolio{227}เตน โข ปน สมเยน สาวตฺถิยา เอวํ ชาตํ เอวํ สมุปฺปนฺนํ อธิกรณํ โหติ.\csromanspacer{} อถ โข เต ภิกฺขู อสนฺตุฏฺฐา สาวตฺถิยา สํฆสฺส อธิกรณวูปสมเนน.\csromanspacer{} อโสสุํ โข “อมุกสฺมิํ กิร อาวาเส สมฺพหุลา เถรา วิหรนฺติ พหุสฺสุตา อาคตาคมา ธมฺมธรา วินยธรา มาติกาธรา ปณฺฑิตา วิยตฺตา เมธาวิโน ลชฺชิโน กุกฺกุจฺจกา สิกฺขากามา, เต เจ เถรา อิมํ อธิกรณํ วูปสเมยฺยุํ ธมฺเมน วินเยน สตฺถุสาสเนน, เอวมิทํ อธิกรณํ สุวูปสนฺตํ อสฺสา”ติ.\csromanspacer{} อถ โข เต ภิกฺขู ตํ อาวาสํ คนฺตฺวา เต เถเร เอตทโวจุํ “อิทํ ภนฺเต อธิกรณํ เอวํ ชาตํ เอวํ สมุปฺปนฺนํ, สาธุ ภนฺเต เถรา อิมํ อธิกรณํ วูปสเมนฺตุ ธมฺเมน วินเยน สตฺถุสาสเนน, ยถยิทํ อธิกรณํ สุวูปสนฺตํ อสฺสา”ติ.\csromanspacer{} อถ โข เต เถรา “ยถา สาวตฺถิยา สํเฆน อธิกรณํ วูปสมิตํ, ตถา สุวูปสนฺตนฺ”ติ\footnote{ยถา สุวูปสนฺตํ (สี, สฺยา)} ตถา ตํ อธิกรณํ วูปสเมสุํ.}

\prose{อถ โข เต ภิกฺขู อสนฺตุฏฺฐา สาวตฺถิยา สํฆสฺส อธิกรณวูปสมเนน อสนฺตุฏฺฐา สมฺพหุลานํ เถรานํ อธิกรณวูปสมเนน.\csromanspacer{} อสฺโสสุํ โข “อมุกสฺมิํ กิร อาวาเส ตโย เถรา วิหรนฺติ ฯเปฯ เทฺว เถรา วิหรนฺติ ฯเปฯ เอโก เถโร วิหรติ พหุสฺสุโต อาคตาคโม ธมฺมธโร วินยธโร มาติกาธโร ปณฺฑิโต วิยตฺโต เมธาวี ลชฺชี กุกฺกุจฺจโก สิกฺขากาโม, โส เจ เถโร อิมํ อธิกรณํ วูปสเมยฺย ธมฺเมน วินเยน สตฺถุสาสเนน, เอวมิทํ อธิกรณํ สุวูปสนฺตํ อสฺสา”ติ.\csromanspacer{} อถ โข เต ภิกฺขู ตํ อาวาสํ คนฺตฺวา ตํ เถรํ เอตทโวจุํ “อิทํ ภนฺเต อธิกรณํ เอวํ ชาตํ เอวํ สมุปฺปนฺนํ, สาธุ ภนฺเต เถโร อิมํ อธิกรณํ วูปสเมตุ ธมฺเมน วินเยน สตฺถุสาสเนน, ยถยิทํ อธิกรณํ สุวูปสนฺตํ อสฺสา”ติ.\csromanspacer{} อถ โข โส เถโร “ยถา สาวตฺถิยา สํเฆน อธิกรณํ วูปสมิตํ ยถา สมฺพหุเลหิ เถเรหิ อธิกรณํ วูปสมิตํ ยถา ตีหิ เถเรหิ อธิกรณํ วูปสมิตํ ยถา ทฺวีหิ เถเรหิ อธิกรณํ วูปสมิตํ, ตถา สุวูปสนฺตนฺ”ติ ตถา ตํ อธิกรณํ วูปสเมสิ.}

\prose{\csromanfolio{228}อถ โข เต ภิกฺขู อสนฺตุฏฺฐา สาวตฺถิยา สํฆสฺส อธิกรณวูปสมเนน อสนฺตุฏฺฐา สมฺพหุลานํ เถรานํ อธิกรณวูปสมเนน อสนฺตุฏฺฐา ติณฺณํ เถเรนํ อธิกรณวูปสมเนน อสนฺตุฏฺฐา ทฺวินฺนํ เถรานํ อธิกรณวูปสมเนน อสนฺตุฏฺฐา เอกสฺส เถรสฺส อธิกรณวูปสมเนน เยน ภควา เตนุปสงฺกมิํสุ, อุปสงฺกมิตฺวา ภควโต เอตมตฺถํ อาโรเจสุํ ฯเปฯ.\csromanspacer{} นิหตเมตํ ภิกฺขเว อธิกรณํ สนฺตํ วูปสนฺตํ สุวูปสนฺตํ.\csromanspacer{} อนุชานามิ ภิกฺขเว เตสํ ภิกฺขูนํ สญฺญตฺติยา ตโย สลากคฺคาเห คุฬฺหกํ สกณฺณชปฺปกํ วิวฏกํ.}

\prose{กถญฺจ ภิกฺขเว คุฬฺหโก สลากคฺคาโห โหติ.\csromanspacer{} เตน สลากคฺคาหาปเกน ภิกฺขุนา สลากาโย วณฺณาวณฺณาโย กตฺวา เอกเมโก ภิกฺขุ อุปสงฺกมิตฺวา เอวมสฺส วจนีโย “อยํ เอวํวาทิสฺส สลากา, อยํ เอวํวาทิสฺส สลากา, ยํ อิจฺฉสิ, ตํ คณฺหาหี”ติ.\csromanspacer{} คหิเต วตฺตพฺโพ “มา จ กสฺสจิ ทสฺเสหี”ติ.\csromanspacer{} สเจ ชานาติ “อธมฺมวาที พหุตรา”ติ, “ทุคฺคโห”ติ ปจฺจุกฺกฑฺฒิตพฺพํ.\csromanspacer{} สเจ ชานาติ “ธมฺมวาที พหุตรา”ติ, “สุคฺคโห”ติ สาเวตพฺพํ.\csromanspacer{} เอวํ โข ภิกฺขเว คูฬฺหโก สลากคฺคาโห โหติ.}

\prose{กถญฺจ ภิกฺขเว สกณฺณชปฺปโก สลากคฺคาโห โหติ.\csromanspacer{} เตน สลากคฺคาหาปเกน ภิกฺขุนา เอกเมกสฺส ภิกฺขุโน อุปกณฺณเก อาโรเจตพฺพํ “อยํ เอวํวาทิสฺส สลากา, อยํ เอวํวาทิสฺส สลากา, ยํ อิจฺฉสิ, ตํ คณฺหาหี”ติ.\csromanspacer{} คหิเต วตฺตพฺโพ “มา จ กสฺสจิ อาโรเจหี”ติ.\csromanspacer{} สเจ ชานาติ “อธมฺมวาที พหุตรา”ติ, “ทุคฺคโห”ติ ปจฺจุกฺกฑฺฒิตพฺพํ.\csromanspacer{} สเจ ชานาติ “ธมฺมวาที พหุตรา”ติ, “สุคฺคโห”ติ สาเวตพฺพํ.\csromanspacer{} เอวํ โข ภิกฺขเว สกณฺณชปฺปโก สลากคฺคาโห โหติ.}

\prose{กถญฺจ ภิกฺขเว วิวฏโก สลากคฺคาโห โหติ.\csromanspacer{} สเจ ชานาติ “ธมฺมวาที พหุตรา”ติ, วิสฺสฏฺเฐเนว วิวเฏน คาเหตพฺโพ.\csromanspacer{} เอวํ โข ภิกฺขเว วิวฏโก สลากคฺคาโห โหติ.\csromanspacer{} อิเม โข ภิกฺขเว ตโย สลากคฺคาหาติ.}

\csromanmatikaanchor{mtk.121}
\csromantocmarkref{subsection}{สติวินย}{mtk.121}
\bookmark[dest={mtk.121},level=2]{สติวินย}
\csromanheader{2}{\textbf{สติวินย}}
\proseitem{236}{\csromanfolio{229}\csromansymbolfootnote{*}{วิ 5. 184, 195 ปิฏฺฐาทีสุปิ.}อนุวาทาธิกรณํ กติหิ สมเถหิ สมฺมติ.\csromanspacer{} อนุวาทาธิกรณํ จตูหิ สมเถหิ สมฺมติ สมฺมุขาวินเยน จ สติวินเยน จ อมูฬฺหวินเยน จ ตสฺสปาปิยสิกาย จ.\csromanspacer{} สิยา อนุวาทาธิกรณํ เทฺว สมเถ อนาคมฺม อมูฬฺหวินยญฺจ ตสฺสปาปิยสิกญฺจ ทฺวีหิ สมเถหิ สมฺเมยฺย สมฺมุขาวินเยน จ สติวินเยน จาติ? สิยาติสฺส วจนียํ.\csromanspacer{} ยถา กถํ วิย.\csromanspacer{} อิธ ปน ภิกฺขเว ภิกฺขู ภิกฺขุํ อมูลิกาย สีลวิปตฺติยา อนุทฺธํเสนฺติ, ตสฺส โข ภิกฺขเว\footnote{ตสฺส โข ตํ ภิกฺขเว (สฺยา, ก)} ภิกฺขุโน สติเวปุลฺลปฺปตฺตสฺส สติวินโย ทาตพฺโพ.\csromanspacer{} เอวญฺจ ปน ภิกฺขเว ทาตพฺโพ,}

\prose{เตน ภิกฺขเว ภิกฺขุนา สํฆํ อุปสงฺกมิตฺวา เอกํสํ อุตฺตราสงฺคํ กริตฺวา ฯเปฯ เอวมสฺส วจนีโย “มํ ภนฺเต ภิกฺขู อมูลิกาย สีลวิปตฺติยา อนุทฺธํเสนฺติ, โสหํ ภนฺเต สติเวปุลฺลปฺปตฺโต สํฆํ สติวินยํ ยาจามี”ติ.\csromanspacer{} ทุติยมฺปิ ยาจิตพฺโพ.\csromanspacer{} ตติยมฺปิ ยาจิตพฺโพ.\csromanspacer{} พฺยตฺเตน ภิกฺขุนา ปฏิพเลน สํโฆ ญาเปตพฺโพ\csromandash{}}

\prose{“สุณาตุ เม ภนฺเต สํโฆ, ภิกฺขู อิตฺถนฺนามํ ภิกฺขุํ อมูลิกาย สีลวิปตฺติยา อนุทฺธํเสนฺติ, โส สติเวปุลฺลปฺปตฺโต สํฆํ สติวินยํ ยาจติ, ยทิ สํฆสฺส ปตฺตกลฺลํ, สํโฆ อิตฺถนฺนามสฺส ภิกฺขุโน สติเวปุลฺลปฺปตฺตสฺส สติวินยํ ทเทยฺย, เอสา ญตฺติ.}

\prose{สุณาตุ เม ภนฺเต สํโฆ, ภิกฺขู อิตฺถนฺนามํ ภิกฺขุํ อมูลิกาย สีลวิปตฺติยา อนุทฺธํเสนฺติ, โส สติเวปุลฺลปฺปตฺโต สํฆํ สติวินยํ ยาจติ, สํโฆ อิตฺถนฺนามสฺส ภิกฺขุโน สติเวปุลฺลปฺปตฺตสฺส สติวินยํ เทติ, ยสฺสายสฺมโต ขมติ อิตฺถนฺนามสฺส ภิกฺขุโน สติเวปุลฺลปฺปตฺตสฺส สติวินยสฺส ทานํ, โส ตุณฺหสฺส, ยสฺส นกฺขมติ, โส ภาเสยฺย.}

\prose{ทุติยมฺปิ เอตมตฺถํ วทามิ ฯเปฯ.\csromanspacer{} ตติยมฺปิ เอตมตฺถํ วทามิ ฯเปฯ.}

\prose{ทินฺโน สํเฆน อิตฺถนฺนามสฺส ภิกฺขุโน สติเวปุลฺลปฺปตฺตสฺส สติวินโย, ขมติ สํฆสฺส, ตสฺมา ตุณฺหี, เอวเมตํ ธารยามี”ติ.}

\prose{อิทํ วุจฺจติ ภิกฺขเว อธิกรณํ วูปสนฺตํ.\csromanspacer{} เกน วูปสนฺตํ, สมฺมุขาวินเยน จ สติวินเยน จ.\csromanspacer{} กิญฺจ ตตฺถ สมฺมุขาวินยสฺมิํ, สํฆสมฺมุขตา \csromanfolio{230}ธมฺมสมฺมุขตา วินยสมฺมุขตา ปุคฺคลสมฺมุขตา ฯเปฯ.\csromanspacer{} กา จ ตตฺถ ปุคฺคลสมฺมุขตา, โย จ อนุวทติ, ยญฺจ อนุวทติ, อุโภ สมฺมุขีภูตา โหนฺติ, อยํ ตตฺถ ปุคฺคลสมฺมุขตา.\csromanspacer{} กิญฺจ ตตฺถ สติวินยสฺมิํ, ยา สติวินยสฺส กมฺมสฺส กิริยา กรณํ อุปคมนํ อชฺฌุปคมนํ อธิวาสนา อปฺปฏิกฺโกสนา, อิทํ ตตฺถ สติวินยสฺมิํ.\csromanspacer{} เอวํ วูปสนฺตํ เจ ภิกฺขเว อธิกรณํ การโก อุกฺโกเฏติ, อุกฺโกฏนกํ ปาจิตฺติยํ.\csromanspacer{} ฉนฺททายโก ขียติ, ขียนกํ ปาจิตฺติยํ.}

\csromanmatikaanchor{mtk.122}
\csromantocmarkref{subsection}{อมูฬฺหวินย}{mtk.122}
\bookmark[dest={mtk.122},level=2]{อมูฬฺหวินย}
\csromanheader{2}{\textbf{อมูฬฺหวินย}}
\proseitem{237}{สิยา อนุวาทาธิกรณํ เทฺว สมเถ อนาคมฺม สติวินยญฺจ ตสฺสปาปิยสิกญฺจ ทฺวีหิ สมเถหิ สมฺเมยฺย สมฺมุขาวินเยน จ อมูฬวินเยน จาติ? สิยาติสฺส วจนียํ.\csromanspacer{} ยถา กถํ วิย.\csromanspacer{} อิธ ปน ภิกฺขเว ภิกฺขุ อุมฺมตฺตโก โหติ จิตฺตวิปริยาสกโต, เตน อุมฺมตฺตเกน จิตฺตวิปริยาสกเตน พหุํ อสฺสามณกํ อชฺฌาจิณฺณํ โหติ ภาสิตปริกฺกนฺตํ, ตํ ภิกฺขู อุมฺมตฺตเกน จิตฺตวิปริยาสกเตน อชฺฌาจิณฺเณน อาปตฺติยา โจเทนฺติ “สรตายสฺมา เอวรูปิํ อาปตฺติํ อาปชฺชิตา”ติ.\csromanspacer{} โส เอวํ วเทติ “อหํ โข อาวุโส อุมฺมตฺตโก อโหสิํ จิตฺตวิปริยาสกโต, เตน เม อุมฺมตฺตเกน จิตฺตวิปริยาสกเตน พหุํ อสฺสามณกํ อชฺฌาจิณฺณํ ภาสิตปริกฺกนฺตํ, นาหํ ตํ สรามิ, มูเฬฺหน เม เอตํ กตนฺ”ติ.\csromanspacer{} เอวมฺปิ นํ วุจฺจมานา โจเทนฺเตว “สรตายสฺมา เอวรูปิํ อาปตฺติํ อาปชฺชิตา”ติ.\csromanspacer{} ตสฺส โข ภิกฺขเว ภิกฺขุโน อมูฬฺหสฺส อมูฬฺหวินโย ทาตพฺโพ.\csromanspacer{} เอวญฺจ ปน ภิกฺขเว ทาตพฺโพ.}

\prose{เตน ภิกฺขเว ภิกฺขุนา สํฆํ อุปสงฺกมิตฺวา เอกํสํ อุตฺตราสงฺคํ กริตฺวา ฯเปฯ เอวมสฺส วจนีโย “อหํ ภนฺเต อุมฺมตฺตโก อโหสิํ จิตฺตวิปริยาสกโต, เตน เม อุมฺมตฺตเกน จิตฺตวิปริยาสกเตน พหุํ อสฺสามณกํ อชฺฌาจิณฺณํ ภาสิตปริกฺกนฺตํ, มํ ภิกฺขู อุมฺมตฺตเกน จิตฺตวิปริยาสกเตน อชฺฌาจิณฺเณน อาปตฺติยา โจเทนฺติ “สรตายสฺมา เอวรูปิํ อาปตฺติํ อาปชฺชิตา”ติ, ตฺยาหํ เอวํ วทามิ ‘อหํ โข อาวุโส อุมฺมตฺตโก อโหสิํ จิตฺตวิปริยาสกโต, เตน เม อุมฺมตฺตเกน จิตฺตวิปริยาสกเตน พหุํ อสฺสามณกํ อชฺฌาจิณฺณํ ภาสิตปริกฺกนฺตํ, นาหํ ตํ สรามิ, มูเฬฺหน เม เอตํ กตนฺ’ติ, เอวมฺปิ มํ วุจฺจมาโน โจเทนฺเตว \csromanfolio{231}“สรตายสฺมา เอวรูปิํ อาปตฺติํ อาปชฺชิตา’ติ, โสหํ ภนฺเต อมูโฬฺห สํฆํ อมูฬฺหวินยํ ยาจามี”ติ.\csromanspacer{} ทุติยมฺปิ ยาจิตพฺโพ.\csromanspacer{} ตติยมฺปิ ยาจิตพฺโพ.\csromanspacer{} พฺยตฺเตน ภิกฺขุนา ปฏิพเลน สํโฆ ญาเปตพฺโพ\csromandash{}}

\prose{“สุณาตุ เม ภนฺเต สํโฆ, อยํ อิตฺถนฺนาโม ภิกฺขุ อุมฺมตฺตโก อโหสิ จิตฺตวิปริยาสกโต, เตน อุมฺมตฺตเกน จิตฺตวิปริยาสกเตน พหุํ อสฺสามณกํ อชฺฌาจิณฺณํ ภาสิตปริกฺกนฺตํ, ตํ ภิกฺขู อุมฺมตฺตเกน จิตฺตวิปริยาสกเตน อชฺฌาจิณฺเณน อาปตฺติยา โจเทนฺติ ‘สรตายสฺมา เอวรูปิํ อาปตฺติํ อาปชฺชิตา’ติ, โส เอวํ วเทติ ‘อหํ โข อาวุโส อุมฺมตฺตโก อโหสิํ จิตฺตวิปริยาสกโต, เตน เม อุมฺมตฺตเกน จิตฺตวิปริยาสกเตน พหุํ อสฺสามณกํ อชฺฌาจิณฺณํ ภาสิตปริกฺกนฺตํ, นาหํ ตํ สรามิ, มูเฬฺหน เม เอตํ กตนฺ’ติ เอวมฺปิ นํ วุจฺจมานา โจเทนฺเตว ‘สรตายสฺมา เอวรูปิํ อาปตฺติํ อาปชฺชิตา’ติ.\csromanspacer{} โส อมูโฬฺห สํฆํ อมูฬฺหวินยํ ยาจติ, ยทิ สํฆสฺส ปตฺตกลฺลํ, สํโฆ อิตฺถนฺนามสฺส ภิกฺขุโน อมูฬฺหสฺส อมูฬฺหวินยํ ทเทยฺย, เอสา ญตฺติ.}

\prose{สุณาตุ เม ภนฺเต สํโฆ, อยํ อิตฺถนฺนาโม ภิกฺขุ อุมฺมตฺตโก อโหสิ จิตฺตวิปริยาสกโต, เตน อุมฺมตฺตเกน จิตฺตวิปริยาสกเตน พหุํ อสฺสามณกํ อชฺฌาจิณฺณํ ภาสิตปริกฺกนฺตํ, ตํ ภิกฺขู อุมฺมตฺตเกน จิตฺตวิปริยาสกเตน อชฺฌาจิณฺเณน อาปตฺติยา โจเทนฺติ ‘สรตายสฺมา เอวรูปิํ อาปตฺติํ อาปชฺชิตา’ติ, โส เอวํ วเทติ ‘อหํ โข อาวุโส อุมฺมตฺตโก อโหสิํ จิตฺตวิปริยาสกโต, เตน เม อุมฺมตฺตเกน จิตฺตวิปริยาสกเตน พหุํ อสฺสามณกํ อชฺฌาจิณฺณํ ภาสิตปริกฺกนฺตํ, นาหํ ตํ สรามิ, มูเฬฺหน เม เอตํ กตนฺ’ติ, เอวมฺปิ นํ วุจฺจมานา โจเทนฺเตว ‘สรตายสฺมา เอวรูปิํ อาปตฺติํ อาปชฺชิตา’ติ, โส อมูโฬฺห สํฆํ อมูฬฺหวินยํ ยาจติ, สํโฆ อิตฺถนฺนามสฺส ภิกฺขุโน อมูฬฺหสฺส อมูฬฺหวินยํ เทติ, ยสฺสายสฺมโต ขมติ อิตฺถนฺนามสฺส ภิกฺขุโน อมูฬฺหสฺส อมูฬฺหวินยสฺส ทานํ, โส ตุณฺหสฺส, ยสฺส นกฺขมติ, โส ภาเสยฺย.}

\prose{ทุติยมฺปิ เอตมตฺถํ วทามิ ฯเปฯ.\csromanspacer{} ตติยมฺปิ เอตมตฺถํ วทามิ ฯเปฯ.}

\prose{ทินฺโน สํเฆน อิตฺถนฺนามสฺส ภิกฺขุโน อมูฬฺหสฺส อมูฬฺหวินโย, ขมติ สํฆสฺส, ตสฺมา ตุณฺหี, เอวเมตํ ธารยามี”ติ.}

\prose{\csromanfolio{232}อิทํ วุจฺจติ ภิกฺขเว อธิกรณํ วูปสนฺตํ.\csromanspacer{} เกน วูปสนฺตํ, สมฺมุขาวินเยน จ อมูฬฺหวินเยน จ.\csromanspacer{} กิญฺจ ตตฺถ สมฺมุขาวินยสฺมิํ, สํฆสมฺมุขตา ธมฺมสมฺมุขตา วินยสมฺมุขตา ปุคฺคลสมฺมุขตา ฯเปฯ.\csromanspacer{} กิญฺจ ตตฺถ อมูฬฺหวินยสฺมิํ, ยา อมูฬฺหวินยสฺส กมฺมสฺส กิริยา กรณํ อุปคมนํ อชฺฌุปคมนํ อธิวาสนา อปฺปฏิกฺโกสนา, อิทํ ตตฺถ อมูฬฺหวินยสฺมิํ.\csromanspacer{} เอวํ วูปสนฺตํ เจ ภิกฺขเว อธิกรณํ การโก อุกฺโกเฏติ, อุกฺโกฏนกํ ปาจิตฺติยํ.\csromanspacer{} ฉนฺททายโก ขียติ, ขียนกํ ปาจิตฺติยํ.}

\csromanmatikaanchor{mtk.123}
\csromantocmarkref{subsection}{ตสฺสปาปิยสิกาวินย}{mtk.123}
\bookmark[dest={mtk.123},level=2]{ตสฺสปาปิยสิกาวินย}
\csromanheader{2}{\textbf{ตสฺสปาปิยสิกาวินย}}
\proseitem{238}{สิยา อนุวาทาธิกรณํ เทฺว สมเถ อนาคมฺม สติวินยญฺจ อมูฬฺหวินยญฺจ ทฺวีหิ สมเถหิ สมฺเมยฺย สมฺมุขาวินเยน จ ตสฺสปาปิยสิกาย จาติ? สิยาติสฺส วจนียํ.\csromanspacer{} ยถา กถํ วิย.\csromanspacer{} อิธ ปน ภิกฺขเว ภิกฺขุ ภิกฺขุํ สํฆมชฺเฌ ครุกาย อาปตฺติยา โจเทติ “สรตายสฺมา เอวรูปิํ ครุกํ อาปตฺติํ อาปชฺชิตา ปาราชิกํ วา ปาราชิกสามนฺตํ วา”ติ.\csromanspacer{} โส เอวํ วเทติ “น โข อหํ อาวุโส สรามิ เอวรูปิํ ครุกํ อาปตฺติํ อาปชฺชิตา ปาราชิกํ วา ปาราชิกสามนฺตํ วา”ติ.\csromanspacer{} ตเมนํ โส นิพฺเพเฐนฺตํ อติเวเฐติ “อิงฺฆายสฺมา สาธุกเมว ชานาติ, ยทิ สรสิ เอวรูปิํ ครุกํ อาปตฺติํ อาปชฺชิตา ปาราชิกํ วา ปาราชิกสามนฺตํ วา”ติ.\csromanspacer{} โส เอวํ วเทติ “น โข อหํ อาวุโส สรามิ เอวรูปิํ ครุกํ อาปตฺติํ อาปชฺชิตา ปาราชิกํ วา ปาราชิกสามนฺตํ วา, สรามิ จ โข อหํ อาวุโส เอวรูปิํ อปฺปมตฺติกํ อาปตฺติํ อาปชฺชิตา”ติ.\csromanspacer{} ตเมนํ โส นิพฺเพเฐนฺตํ อติเวเฐติ “อิงฺฆายสฺมา สาธุกเมว ชานาหิ, ยทิ สรสิ เอวรูปิํ ครุกํ อาปตฺติํ อาปชฺชิตา ปาราชิกํ วา ปาราชิกสามนฺตํ วา”ติ.\csromanspacer{} โส เอวํ วเทติ “อิมํ หิ นามาหํ อาวุโส อปฺปมตฺติกํ อาปตฺติํ อาปชฺชิตฺวา อปุฏฺโฐ ปฏิชานิสฺสามิ, กิํ ปนาหํ เอวรูปิํ ครุกํ อาปตฺติํ อาปชฺชีตฺวา ปาราชิกํ วา ปาราชิกสามนฺตํ วา ปุฏฺโฐ น ปฏิชานิสฺสามี”ติ.\csromanspacer{} โส เอวํ วเทติ “อิมํ หิ นาม ตฺวํ อาวุโส อปฺปมตฺติกํ อาปตฺติํ อาปชฺชิตฺวา อปุฏฺโฐ น ปฏิชานิสฺสสิ, กิํ ปน ตฺวํ เอวรูปิํ ครุกํ อาปตฺติํ อาปชฺชิตฺวา ปาราชิกํ วา ปาราชิกสามนฺตํ วา อปุฏฺโฐ ปฏิชานิสฺสสิ, อิงฺฆายสฺมา สาธุกเมว ชานาหิ, ยทิ สรสิ เอวรูปิํ ครุกํ อาปตฺติํ อาปชฺชิตา ปาราชิกํ วา ปาราชิกสามนฺตํ วา”ติ.\csromanspacer{} โส เอวํ วเทสิ “สรามิ โข อหํ อาวุโส \csromanfolio{233}เอวรูปิํ ครุกํ อาปตฺติํ อาปชฺชิตา ปาราชิกํ วา ปาราชิกสามนฺตํ วา, ทวา เม เอตํ วุตฺตํ, รวา เม เอตํ วุตฺตํ, นาหํ ตํ สรามิ เอวรูปิํ ครุกํ อาปตฺติํ อาปชฺชิตา ปาราชิกํ วา ปาราชิกสามนฺตํ วา”ติ.\csromanspacer{} ตสฺส โข ภิกฺขเว\footnote{ตสฺส โข ตํ ภิกฺขเว (ก), ตสฺส เขฺวตํ ภิกฺขเว (สฺยา)} ภิกฺขุโน ตสฺสปาปิยสิกากมฺมํ กาตพฺพํ.\csromanspacer{} เอวญฺจ ปน ภิกฺขเว กาตพฺพํ, พฺยตฺเตน ภิกฺขุนา ปฏิพเลน สํโฆ ญาเปตพฺโพ\csromandash{}}

\prose{“สุณาตุ เม ภนฺเต สํโฆ, อยํ อิตฺถนฺนาโม ภิกฺขุ สํฆมชฺเฌ ครุกาย อาปตฺติยา อนุยุญฺชิยมาโน อวชานิตฺวา ปฏิชานาติ, ปฏิชานิตฺวา อวชานาติ, อญฺเญนาญฺญํ ปฏิจรติ, สมฺปชานมุสา ภาสติ, ยทิ สํฆสฺส ปตฺตกลฺลํ, สํโฆ อิตฺถนฺนามสฺส ภิกฺขุโน ตสฺสปาปิยสิกากมฺมํ กเรยฺย, เอสา ญตฺติ.}

\prose{สุณาตุ เม ภนฺเต สํโฆ, อยํ อิตฺถนฺนาโม ภิกฺขุ สํฆมชฺเฌ ครุกาย อาปตฺติยา อนุยุญฺชิยมาโน อวชานิตฺวา ปฏิชานาติ ปฏิชานิตฺวา อวชานาติ, อญฺเญนาญฺญํ ปฏิจรติ, สมฺปชานมุสา ภาสติ.\csromanspacer{} สํโฆ อิตฺถนฺนามสฺส ภิกฺขุโน ตสฺสปาปิยสิกากมฺมํ กโรติ, ยสฺสายสฺมโต ขมติ อิตฺถนฺนามสฺส ภิกฺขุโน ตสฺสปาปิยสิกากมฺมสฺส กรณํ, โส ตุณฺหสฺส, ยสฺส นกฺขมติ, โส ภาเสยฺย.}

\prose{ทุติยมฺปิ เอตมตฺถํ วทามิ ฯเปฯ.\csromanspacer{} ตติยมฺปิ เอตมตฺถํ วทามิ ฯเปฯ.}

\prose{กตํ สํเฆน อิตฺถนฺนามสฺส ภิกฺขุโน ตสฺสปาปิยสิกากมฺมํ, ขมติ สํฆสฺส, ตสฺมา ตุณฺหี, เอวเมตํ ธารยามี”ติ.}

\prose{อิทํ วุจฺจติ ภิกฺขเว อธิกรณํ วูปสนฺตํ.\csromanspacer{} เกน วูปสนฺตํ, สมฺมุขาวินเยน จ ตสฺสปาปิยสิกาย จ.\csromanspacer{} กิญฺจ ตตฺถ สมฺมุขาวินยสฺมิํ, สํฆสมฺมุขตา ธมฺมสมฺมุขตา วินยสมฺมุขตา ฯเปฯ.\csromanspacer{} กา จ ตตฺถ ตสฺสปาปิยสิกาย, ยา ตสฺสปาปิยสิกากมฺมสฺส กิริยา กรณํ อุปคมนํ อชฺฌุปคมนํ อธิวาสนา อปฺปฏิกฺโกสนา, อยํ ตตฺถ ตสฺสปาปิยสิกาย.\csromanspacer{} เอวํ วูปสนฺตํ เจ ภิกฺขเว อธิกรณํ การโก อุกฺโกเฏติ, อุกฺโกฏนกํ ปาจิตฺติยํ.\csromanspacer{} ฉนฺททายโก ขียติ, ขียนกํ ปาจิตฺติยํ.}

\csromanmatikaanchor{mtk.124}
\csromantocmarkref{subsection}{ปฏิญฺญาตกรณ}{mtk.124}
\bookmark[dest={mtk.124},level=2]{ปฏิญฺญาตกรณ}
\csromanheader{2}{\textbf{ปฏิญฺญาตกรณ}}
\proseitem{239}{\csromansymbolfootnote{*}{วิ 5. 184, 195, 198 ปิฏฺฐาทีสุปิ.}อาปตฺตาธิกรณํ กติหิ สมเถหิ สมฺมติ.\csromanspacer{} อาปตฺตาธิกรณํ ตีหิ สมเถหิ สมฺมติ สมฺมุขาวินเยน จ ปฏิญฺญาตกรเณน จ \csromanfolio{234}ติณวตฺถารเกน จ.\csromanspacer{} สิยา อาปตฺตาธิกรณํ เอกํ สมถํ อนาคมฺม ติณวตฺถารกํ ทฺวีหิ สมเถหิ สมฺเมยฺย สมฺมุขาวินเยน จ ปฏิญฺญาตกรเณน จาติ? สิยาติสฺส วจนียํ.\csromanspacer{} ยถา กถํ วิย.\csromanspacer{} อิธ ปน ภิกฺขเว ภิกฺขุ ลหุกํ อาปตฺติํ อาปนฺโน โหติ, เตน ภิกฺขเว ภิกฺขุโน เอกํ ภิกฺขุํ อุปสงฺกมิตฺวา เอกํสํ อุตฺตราสงฺคํ กริตฺวา อุกฺกุฏิกํ นิสีทิตฺวา อญฺชลิํ ปคฺคเหตฺวา เอวมสฺส วจนีโย “อหํ อาวุโส อิตฺถนฺนามํ อาปตฺติํ อาปนฺโน, ตํ ปฏิเทเสมี”ติ.\csromanspacer{} เตน วตฺตพฺโพ “ปสฺสสี”ติ. “อาม ปสฺสามี”ติ. “อายติํ สํวเรยฺยาสี”ติ.}

\prose{อิทํ วุจฺจติ ภิกฺขเว อธิกรณํ วูปสนฺตํ.\csromanspacer{} เกน วูปสนฺตํ, สมฺมุขาวินเยน จ ปฏิญฺญาตกรเณน จ.\csromanspacer{} กิญฺจ ตตฺถ สมฺมุขาวินยสฺมิํ, ธมฺมสมฺมุขตา วินยสมฺมุขตา ปุคฺคลสมฺมุขตา ฯเปฯ.\csromanspacer{} กา จ ตตฺถ ปุคฺคลสมฺมุขตา, โย จ เทเสติ, ยสฺส จ เทเสติ, อุโภ สมฺมุขีภูตา โหนฺติ.\csromanspacer{} อยํ ตตฺถ ปุคฺคลสมฺมุขตา.\csromanspacer{} กิญฺจ ตตฺถ ปฏิญฺญาตกรณสฺมิํ, ยา ปฏิญฺญาตกรณสฺส กมฺมสฺส กิริยา กรณํ อุปคมนํ อชฺฌุปคมนํ อธิวาสนา อปฺปฏิกฺโกสนา, อิทํ ตตฺถ ปฏิญฺญาตกรณสฺมิํ.\csromanspacer{} เอวํ วูปสนฺตํ เจ ภิกฺขเว อธิกรณํ ปฏิคฺคาหโก อุกฺโกเฏติ, อุกฺโกฏนกํ ปาจิตฺติยํ.\csromanspacer{} เอวญฺเจตํ ลเภถ, อิจฺเจตํ กุสลํ.\csromanspacer{} โน เจ ลเภถ, เตน ภิกฺขเว ภิกฺขุนา สมฺพหุเล ภิกฺขู อุปสงฺกมิตฺวา เอกํสํ อุตฺตราสงฺคํ กริตฺวา วุฑฺฒานํ ภิกฺขูนํ ปาเท วนฺทิตฺวา อุกฺกุฏิกํ นิสีทิตฺวา อญฺชลิํ ปคฺคเหตฺวา เอวมสฺสุ วจนียา “อหํ ภนฺเต อิตฺถนฺนามํ อาปตฺติํ อาปนฺโน,ตํ ปฏิเทเสมี”ติ.\csromanspacer{} พฺยตฺเตน ภิกฺขุนา ปฏิพเลน เต ภิกฺขู ญาเปตพฺพา\csromandash{}}

\prose{“สุณนฺตุ เม อายสฺมนฺตา, อยํ อิตฺถนฺนาโม ภิกฺขุ อาปตฺติํ สรติ วิวรติ อุตฺตานิํ กโรติ เทเสติ, ยทายสฺมนฺตานํ ปตฺตกลฺลํ, อหํ อิตฺถนฺนามสฺส ภิกฺขุโน อาปตฺติํ ปฏิคฺคเณฺหยฺยนฺ”ติ.\csromanspacer{} เตน วตฺตพฺโพ “ปสฺสสี”ติ. “อาม ปสฺสามี”ติ. “อายติํ สํวเรยฺยาสี”ติ.}

\prose{อิทํ วุจฺจติ ภิกฺขเว อธิกรณํ วูปสนฺตํ.\csromanspacer{} เกน วูปสนฺตํ, สมฺมุขาวินเยน จ ปฏิญฺญาตกรเณน จ.\csromanspacer{} กิญฺจ ตตฺถ สมฺมุขาวินยสฺมิํ, ธมฺมสมฺมุขตา วินยสมฺมุขตา ปุคฺคลสมฺมุขตา ฯเปฯ.\csromanspacer{} กา จ ตตฺถ ปุคฺคลสมฺมุขตา, โย จ เทเสติ, ยสฺส จ เทเสติ, อุโภ สมฺมุขีภูตา \csromanfolio{235}โหนฺติ.\csromanspacer{} อยํ ตตฺถ ปุคฺคลสมฺมุขตา.\csromanspacer{} กิญฺจ ตตฺถ ปฏิญฺญาตกรณสฺมิํ, ยา ปฏิญฺญาตกรณสฺส กมฺมสฺส กิริยา กรณํ อุปคมนํ อชฺฌุปคมนํ อธิวาสนา อปฺปฏิกฺโกสนา, อิทํ ตตฺถ ปฏิญฺญาตกรณสฺมิํ.\csromanspacer{} เอวํ วูปสนฺตํ เจ ภิกฺขเว อธิกรณํ ปฏิคฺคาหโก อุกฺโกเฏติ, อุกฺโกฏนกํ ปาจิตฺติยํ.\csromanspacer{} เอวญฺเจตํ ลเภถ, อิจฺเจตํ กุสลํ.\csromanspacer{} โน เจ ลเภถ, เตน ภิกฺขเว ภิกฺขุนา สํฆํ อุปสงฺกมิตฺวา เอกํสํ อุตฺตราสงฺคํ กริตฺวา วุฑฺฒานํ ภิกฺขูนํ ปาเท วนฺทิตฺวา อุกฺกุฏิกํ นิสีทิตฺวา อญฺชลิํ ปคฺคเหตฺวา เอวมสฺส วจนีโย “อหํ ภนฺเต อิตฺถนฺนามํ อาปตฺติํ อาปนฺโน, ตํ ปฏิเทเสมี”ติ.\csromanspacer{} พฺยตฺเตน ภิกฺขุนา ปฏิพเลน สํโฆ ญาเปตพฺโพ\csromandash{}}

\prose{“สุณาตุ เม ภนฺเต สํโฆ, อยํ อิตฺถนฺนาโม ภิกฺขุ อาปตฺติํ สรติ วิวรติ อุตฺตานิํ กโรติ เทเสติ, ยทิ สํฆสฺส ปตฺตกลฺลํ, อหํ อิตฺถนฺนามสฺส ภิกฺขุโน อาปตฺติํ ปฏิคฺคเณฺหยฺยนฺ”ติ.\csromanspacer{} เตน วตฺตพฺโพ “ปสฺสสี”ติ. “อาม ปสฺสามี”ติ. “อายติํ สํวเรยฺยาสี”ติ.}

\prose{อิทํ วุจฺจติ ภิกฺขเว อธิกรณํ วูปสนฺตํ.\csromanspacer{} เกน วูปสนฺตํ, สมฺมุขาวินเยน จ ปฏิญฺญาตกรเณน จ.\csromanspacer{} กิญฺจ ตตฺถ สมฺมุขาวินยสฺมิํ, สํฆสมฺมุขตา ธมฺมสมฺมุขตา วินยสมฺมุขตา ปุคฺคลสมฺมุขตา ฯเปฯ.\csromanspacer{} เอวํ วูปสนฺตํ เจ ภิกฺขเว อธิกรณํ ปฏิคฺคาหโก อุกฺโกเฏติ, อุกฺโกฏนกํ ปาจิตฺติยํ.\csromanspacer{} ฉนฺททายโก ขียติ, ขียนกํ ปาจิตฺติยํ.}

\csromanmatikaanchor{mtk.125}
\csromantocmarkref{subsection}{ติณวตฺถารก}{mtk.125}
\bookmark[dest={mtk.125},level=2]{ติณวตฺถารก}
\csromanheader{2}{\textbf{ติณวตฺถารก}}
\proseitem{240}{สิยา อาปตฺตาธิกรณํ เอกํ สมถํ อนาคมฺม ปฏิญฺญาตกรณํ ทฺวีหิ สมเถหิ สมฺเมยฺย สมฺมุขาวินเยน จ ติณวตฺถารเกน จาติ? สิยาติสฺส วจนียํ.\csromanspacer{} ยถา กถํ วิย.\csromanspacer{} อิธ ปน ภิกฺขเว ภิกฺขูนํ ภณฺฑนชาตานํ กลหชาตานํ วิวาทาปนฺนานํ วิหรตํ พหุํ อสฺสามณกํ อชฺฌาจิณฺณํ โหติ ภาสิตปริกฺกนฺตํ.\csromanspacer{} ตตฺร เจ ภิกฺขูนํ เอวํ โหติ “อมฺหากํ โข ภณฺฑนชาตานํ กลหชาตานํ วิเวทาปนฺนานํ วิหรตํ พหุํ อสฺสามณกํ อชฺฌาจิณฺณํ ภาสิตปริกฺกนฺตํ, สเจ มยํ อิมาหิ อาปตฺตีหิ อญฺญมญฺญํ กาเรสฺสาม, สิยาปิ ตํ อธิกรณํ กกฺขฬตฺตาย วาฬตฺตาย เภทาย สํวตฺเตยฺยา”ติ.\csromanspacer{} อนุชานามิ ภิกฺขเว เอวรูปํ อธิกรณํ ติณวตฺถารเกน วูปสเมตุํ.\csromanspacer{} เอวญฺจ ปน ภิกฺขเว วูปสเมตพฺพํ, \csromanfolio{236}สพฺเพเหว เอกชฺฌํ สนฺนิปติตพฺพํ, สนฺนิปติตฺวา พฺยตฺเตน ภิกฺขุนา ปฏิพเลน สํโฆ ญาเปตพฺโพ\csromandash{}}

\prose{“สุณาตุ เม ภนฺเต สํโฆ, อมฺหากํ ภณฺฑนชาตานํ กลหชาตานํ วิวาทาปนฺนานํ วิหรตํ พหุํ อสฺสามณกํ อชฺฌาจิณฺณํ ภาสิตปริกฺกนฺตํ, สเจ มยํ อิมาหิ อาปตฺตีหิ อญฺญมญฺญํ กาเรสฺสาม, สิยาปิ ตํ อธิกรณํ กกฺขฬตฺตาย วาฬตฺตาย เภทาย สํวตฺเตยฺย, ยทิ สํฆสฺส ปตฺตกลฺลํ, สํโฆ อิมํ อธิกรณํ ติณวตฺถารเกน วูปสเมยฺย ฐเปตฺวา ถุลฺลวชฺชํ, ฐเปตฺวา คิหิปฺปฏิสํยุตฺตนฺ”ติ.\csromanspacer{} เอกโตปกฺขิกานํ ภิกฺขูนํ พฺยตฺเตน ภิกฺขุนา ปฏิพเลน สโก ปกฺโข ญาเปตพฺโพ\csromandash{}}

\prose{“สุณนฺตุ เม อายสฺมนฺตา, อมฺหากํ ภณฺฑนชาตานํ กลหชาตานํ วิวาทาปนฺนานํ วิหรตํ พหุํ อสฺสามณกํ อชฺฌาจิณฺณํ ภาสิตปริกฺกนฺตํ, สเจ มยํ อิมาหิ อาปตฺตีหิ อญฺญมญฺญํ กาเรสฺสาม, สิยาปิ ตํ อธิกรณํ กกฺขฬตฺตาย วาฬตฺตาย เภทาย สํวตฺเตยฺย, ยทายสฺมนฺตานํ ปตฺตกลฺลํ, อหํ ยา เจว อายสฺมนฺตานํ อาปตฺติ, ยา จ อตฺตโน อาปตฺติ, อายสฺมนฺตานญฺเจว อตฺถาย อตฺตโน จ อตฺถาย สํฆมชฺเฌ ติณวตฺถารเกน เทเสยฺยํ ฐเปตฺวา ถุลฺลวชฺชํ, ฐเปตฺวา คิหิปฺปฏิสํยุตฺตนฺ”ติ.}

\proseitem{241}{อถาปเรสํ เอกโตปกฺขิกานํ ภิกฺขูนํ พฺยตฺเตน ภิกฺขุนา ปฏิพเลน สโก ปกฺโข ญาเปตพฺโพ\csromandash{}}

\prose{“สุณนฺตุ เม อายสฺมนฺตา, อมฺหากํ ภณฺฑนชาตานํ กลหชาตานํ วิวาทาปนฺนานํ วิหรตํ พหุํ อสฺสามณกํ อชฺฌาจิณฺณํ ภาสิตปริกฺกนฺตํ, สเจ มยํ อิมาหิ อาปตฺตีหิ อญฺญมญฺญํ กาเรสฺสาม, สิยาปิ ตํ อธิกรณํ กกฺขฬตฺตาย วาฬตฺตาย เภทาย สํวตฺเตยฺย, ยทายสฺมนฺตานํ ปตฺตกลฺลํ, อหํ ยา เจว อายสฺมนฺตานํ อาปตฺติ, ยา จ อตฺตโน อาปตฺติ, อายสฺมนฺตานญฺเจว อตฺถาย อตฺตโน จ อตฺถาย สํฆมชฺเฌ ติณวตฺถารเกน เทเสยฺยํ ฐเปตฺวา ถุลฺลวชฺชํ, ฐเปตฺวา คิหิปฺปฏิสํยุตฺตนฺ”ติ.}

\proseitem{242}{อถาปเรสํ เอกโตปกฺขิกานํ ภิกฺขูนํ พฺยตฺเตน ภิกฺขุนา ปฏิพเลน สํโฆ ญาเปตพฺโพ\csromandash{}}

\prose{\csromanfolio{237}“สุณาตุ เม ภนฺเต สํโฆ, อมฺหากํ ภณฺฑนชาตานํ กลหชาตานํ วิวาทาปนฺนานํ วิหรตํ พหุํ อสฺสามณกํ อชฺฌาจิณฺณํ ภาสิตปริกฺกนฺตํ, สเจ มยํ อิมาหิ อาปตฺตีหิ อญฺญมญฺญํ กาเรสฺสาม, สิยาปิ ตํ อธิกรณํ กกฺขฬตฺตาย วาฬตฺตาย เภทาย สํวตฺเตยฺย, ยทิ สํฆสฺส ปตฺตกลฺลํ, อหํ ยา เจว อิเมสํ อายสฺมนฺตานํ อาปตฺติ, ยา จ อตฺตโน อาปตฺติ, อิเมสญฺเจว อายสฺมนฺตานํ อตฺถาย อตฺตโน จ อตฺถาย สํฆมชฺเฌ ติณวตฺถารเกน เทเสยฺยํ ฐเปตฺวา ถุลฺลวชฺชํ, ฐเปตฺวา คิหิปฺปฏิสํยุตฺตํ, เอสา ญตฺติ.}

\prose{สุณาตุ เม ภนฺเต สํโฆ, อมฺหากํ ภณฺฑนชาตานํ กลหชาตานํ วิวาทาปนฺนานํ วิหรตํ พหุํ อสฺสามณกํ อชฺฌาจิณฺณํ ภาสิตปริกฺกนฺตํ, สเจ มยํ อิมาหิ อาปตฺตีหิ อญฺญมญฺญํ กาเรสฺสาม, สิยาปิ ตํ อธิกรณํ กกฺขฬตฺตาย วาฬตฺตาย เภทาย สํวตฺเตยฺย, อหํ ยา เจว อิเมสํ อายสฺมนฺตานํ อาปตฺติ, ยา จ อตฺตโน อาปตฺติ, อิเมสญฺเจจ อายสฺมนฺตานํ อตฺถาย อตฺตโน จ อตฺถาย สํฆมชฺเฌ ติณวตฺถารเกน เทเสมิ ฐเปตฺวา ถุลฺลวชฺชํ, ฐเปตฺวา คิหิปฺปฏิสํยุตฺตํ, ยสฺสายสฺมโต ขมติ อมฺหากํ อิมาสํ อาปตฺตีนํ สํฆมชฺเฌ ติณวตฺถารเกน เทสนา ฐเปตฺวา ถุลฺลวชฺชํ, ฐเปตฺวา คิหิปฺปฏิสํยุตฺตํ, โส ตุณฺหสฺส, ยสฺส นกฺขมติ, โส ภาเสยฺย.}

\prose{เทสิตา อมฺหากํ อิมา อาปตฺติโย สํฆมชฺเฌ ติณวตฺถารเกน ฐเปตฺวา ถุลฺลวชฺชํ, ฐเปตฺวา คิหิปฺปฏิสํยุตฺตํ, ขมติ สํฆสฺส, ตสฺมา ตุณฺหี, เอวเมตํ ธารยามี”ติ.}

\prose{อถาปเรสํ ฯเปฯ เอวเมตํ ธารยามีติ.}

\prose{อิทํ วุจฺจติ ภิกฺขเว อธิกรณํ วูปสนฺตํ.\csromanspacer{} เกน วูปสนฺตํ, สมฺมุขาวินเยน จ ติณวตฺถารเกน จ.\csromanspacer{} กิญฺจ ตตฺถ สมฺมุขาวินยสฺมิํ, สํฆสมฺมุขตา ธมฺมสมฺมุขตา วินยสมฺมุขตา ปุคฺคลสมฺมุขตา.}

\prose{กา จ ตตฺถ สํฆสมฺมุขตา, ยาวติกา ภิกฺขู กมฺมปฺปตฺตา, เต อาคตา โหนฺติ, ฉนฺทารหานํ ฉนฺโท อาหโฏ โหติ สมฺมุขีภูตา น ปฏิกฺโกสนฺติ, อยํ ตตฺถ สํฆสมฺมุขตา.}

\prose{\csromanfolio{238}กา จ ตตฺถ ธมฺมสมฺมุขตา วินยสมฺมุขตา, เยน ธมฺเมน เยน วินเยน เยน สตฺถุสาสเนน ตํ อธิกรณํ วูปสมฺมติ, อยํ ตตฺถ ธมฺมสมฺมุขตา วินยสมฺมุขตา.\csromanspacer{} กา จ ตตฺถ ปุคฺคลสมฺมุขตา, โย จ เทเสติ, ยสฺส จ เทเสติ, อุโภ สมฺมุขีภูตา โหติ, อยํ ตตฺถ ปุคฺคลสมฺมุขตา.}

\prose{กิญฺจ ตตฺถ ติณวตฺถารกสฺมิํ, ยา ติณวตฺถารกสฺส กมฺมสฺส กิริยา กรณํ อุปคมนํ อชฺฌุปคมนํ อธิวาสนา อปฺปฏิกฺโกสนา, อิทํ ตตฺถ ติณวตฺถารกสฺมิํ.\csromanspacer{} เอวํ วูปสนฺตํ เจ ภิกฺขเว อธิกรณํ ปฏิคฺคาหโก อุกฺโกเฏติ, อุกฺโกฏนกํ ปาจิตฺติยํ.\csromanspacer{} ฉนฺททายโก ขียติ, ขียนกํ ปาจิตฺติยํ.}

\prose{\csromansymbolfootnote{*}{วิ 5. 185, 195, 198 ปิฏฺฐาทีสุปิ.}กิจฺจาธิกรณํ กติหิ สมเถหิ สมฺมติ.\csromanspacer{} กิจฺจาธิกรณํ เอเกน สมเถน สมฺมติ สมฺมุขาวินเยนาติ.}

\nitthitam{\textbf{สมถกฺขนฺธโก นิฏฺฐิโต จตุตฺโถ.}}
\csromanmatikaanchor{mtk.126}
\csromantocmarknopage{chapter}{5. ขุทฺทกวตฺถุกฺขนฺธก}{mtk.126}
\bookmark[dest={mtk.126},level=0]{5. ขุทฺทกวตฺถุกฺขนฺธก}
\chapterheadpageread{5. ขุทฺทกวตฺถุกฺขนฺธก}
\csromanmatikaanchor{mtk.127}
\csromantocmarkref{section}{ขุทฺทกวตฺถูนิ}{mtk.127}
\bookmark[dest={mtk.127},level=1]{ขุทฺทกวตฺถูนิ}
\csromanheader{1}{\textbf{ขุทฺทกวตฺถูนิ}}
\proseitem{243}{\csromanfolio{239}เตน สมเยน พุทฺโธ ภควา ราชคเห วิหรติ เวฬุวเน กลนฺทกนิวาเป.\csromanspacer{} เตน โข ปน สมเยน ฉพฺพคฺคิยา ภิกฺขู นหายมานา รุกฺเข กายํ อุคฺฆํเสนฺติ, อูรุมฺปิ พาหุมฺปิ อุรมฺปิ ปิฏฺฐิมฺปิ.\csromanspacer{} มนุสฺสา อุชฺฌายนฺติ ขิยฺยนฺติ วิปาเจนฺติ “กถํ หิ นาม สมณา สกฺยปุตฺติยา นหายมานา รุกฺเข กายํ อุคฺฆํเสสฺสนฺติ, อูรุมฺปิ พาหุมฺปิ อุรมฺปิ ปิฏฺฐิมฺปิ เสยฺยถาปิ มลฺลมุฏฺฐิกา คามโมทฺทวา”ติ\footnote{คามโปทฺทวา (สี), คามปูตวา (สฺยา)}.\csromanspacer{} อสฺโสสุํ โข ภิกฺขู เตสํ มนุสฺสานํ อุชฺฌายนฺตานํ ขิยฺยนฺตานํ วิปาเจนฺตานํ.\csromanspacer{} อถ โข เต ภิกฺขู ภควโต เอตมตฺถํ อาโรเจสุํ.\csromanspacer{} อถ โข ภควา เอตสฺมิํ นิทาเน เอตสฺมิํ ปกรเณ ภิกฺขุสํฆํ สนฺนิปาตาเปตฺวา ภิกฺขู ปฏิปุจฺฉิ “สจฺจํ กิร ภิกฺขเว ฉพฺพคฺคิยา ภิกฺขู นหายมานา รุกฺเข กายํ อุคฺฆํเสนฺติ, อูรุํปิ พาหุํปิ อุรํปิ ปิฏฺฐิํปี”ติ.\csromanspacer{} สจฺจํ ภควาติ.\csromanspacer{} วิครหิ พุทฺโธ ภควา อนนุจฺฉวิกํ ภิกฺขเว เตสํ โมฆปุริสานํ อนนุโลมิกํ อปฺปติรูปํ อสฺสามณกํ อกปฺปิยํ อกรณียํ, กถํ หิ นาม เต ภิกฺขเว โมฆปุริสา นหายมานา รุกฺเข กายํ อุคฺฆํเสสฺสนฺติ, อูรุมฺปิ พาหุมฺปิ อุรมฺปิ ปิฏฺฐิมฺปิ.\csromanspacer{} เนตํ ภิกฺขเว อปฺปสนฺนานํ วา ปสาทาย ฯเปฯ วิครหิตฺวา ฯเปฯ ธมฺมิํ กถํ กตฺวา ภิกฺขู อามนฺเตสิ “น ภิกฺขเว นหายมาเนน ภิกฺขุนา รุกฺเข กาโย อุคฺฆํเสตพฺโพ, โย อุคฺฆํเสยฺย, อาปตฺติ ทุกฺกฏสฺสา”ติ.}

\prose{เตน โข ปน สมเยน ฉพฺพคฺคิยา ภิกฺขู นหายมานา ถมฺเภ กายํ อุคฺฆํเสนฺติ, อูรุมฺปิ พาหุมฺปิ อุรมฺปิ ปิฏฺฐิมฺปิ.\csromanspacer{} มนุสฺสา อุชฺฌายนฺติ ขิยฺยนฺติ วิปาเจนฺติ “กถํ หิ นาม สมณา สกฺยปุตฺติยา นหายมานา ถมฺเภ กายํ อุคฺฆํเสสฺสนฺติ, อูรุมฺปิ พาหุมฺปิ อุรมฺปิ ปิฏฺฐิมฺปิ เสยฺยถาปิ มลฺลมุฏฺฐิกา คามโมทฺทวา”ติ.\csromanspacer{} อสฺโสสุํ โข ภิกฺขู เตสํ มนุสฺสานํ อุชฺฌายนฺตานํ ขิยฺยนฺตานํ วิปาเจนฺตานํ.\csromanspacer{} อถ โข เต ภิกฺขู ภควโต เอตมตฺถํ อาโรเจสุํ ฯเปฯ.\csromanspacer{} สจฺจํ ภควาติ ฯเปฯ วิครหิตฺวา ฯเปฯ ธมฺมิํ กถํ กตฺวา ภิกฺขู อามนฺเตสิ “น ภิกฺขเว นหายมาเนน ภิกฺขุนา ถมฺเภ กาโย อุคฺฆํเสตพฺโพ, โย อุคฺฆํเสยฺย, อาปตฺติ ทุกฺกฏสฺสา”ติ.}

\prose{\csromanfolio{240}เตน โข ปน สมเยน ฉพฺพคฺคิยา ภิกฺขู นหายมานา กุฏฺเฏ\footnote{กุฑฺเฑ (สี, สฺยา)} กายํ อุคฺฆํเสนฺติ, อูรุมฺปิ พาหุมฺปิ อุรมฺปิ ปิฏฺฐิมฺปิ.\csromanspacer{} มนุสฺสา อุชฺฌายนฺติ ขิยฺยนฺติ วิปาเจนฺติ “กถํ หิ นาม สมณา สกฺยปุตฺติยา นหายมานา กุฏฺเฏ กายํ อุคฺฆํเสสฺสนฺติ, อูรุมฺปิพาหุมฺปิ อุรมฺปิ ปิฏฺฐิมฺปิ เสยฺยถาปิ มลฺลมุฏฺฐิกา คามโมทฺทวา”ติ ฯเปฯ น ภิกฺขเว นหายมาเนน ภิกฺขุนา กุฏฺเฏ กาโย อุคฺฆํเสตพฺโพ, โย อุคฺฆํเสยฺย, อาปตฺติ ทุกฺกฏสฺสาติ.}

\prose{เตน โข ปน สมเยน ฉพฺพคฺคิยา ภิกฺขู อฏฺฏาเน\footnote{อฏฺฐาเน (สี, สฺยา)} นหายนฺติ.\csromanspacer{} มนุสฺสา อุชฺฌายนฺติ ขิยฺยนฺติ วิปาเจนฺติ “ฯเปฯ เสยฺยถาปิ คิหี กามโภคิโน”ติ.\csromanspacer{} อสฺโสสุํ โข ภิกฺขู เตสํ มนุสฺสานํ อุชฺฌายนฺตานํ ขิยฺยนฺตานํ วิปาเจนฺตานํ.\csromanspacer{} อถ โข เต ภิกฺขู ภควโต เอตมตฺถํ อาโรเจสุํ ฯเปฯ.\csromanspacer{} สจฺจํ ภควาติ ฯเปฯ วิครหิตฺวา ฯเปฯ ธมฺมิํ กถํ กตฺวา ภิกฺขู อามนฺเตสิ “น ภิกฺขเว อฏฺฏาเน นหายิตพฺพํ, โย นหาเยยฺย, อาปตฺติ ทุกฺกฏสฺสา”ติ.}

\prose{เตน โข ปน สมเยน ฉพฺพคฺคิยา ภิกฺขู คนฺธพฺพหตฺถเกน นหายนฺติ.\csromanspacer{} มนุสฺสา อุชฺฌายนฺติ ขิยฺยนฺติ วิปาเจนฺติ “ฯเปฯ เสยฺยถาปิ คิหี กามโภคิโน”ติ.\csromanspacer{} อสฺโสสุํ โข ภิกฺขู เตสํ มนุสฺสานํ อุชฺฌายนฺตานํ ขิยฺยนฺตานํ วิปาเจนฺตานํ.\csromanspacer{} อถ โข เต ภิกฺขู ภควโต เอตมตฺถํ อาโรเจสุํ ฯเปฯ น ภิกฺขเว คนฺธพฺพหตฺถเกน นหายิตพฺพํ, โย นหาเยยฺย, อาปตฺติ ทุกฺกฏสฺสาติ.}

\prose{เตน โข ปน สมเยน ฉพฺพคฺคิยา ภิกฺขู กุรุวินฺทกสุตฺติยา นหายนฺติ.\csromanspacer{} มนุสฺสา อุชฺฌายนฺติ ขิยฺยนฺติ วิปาเจนฺติ “ฯเปฯ เสยฺยถาปิ คิหี กามโภคิโน”ติ.\csromanspacer{} ภควโต เอตมตฺถํ อาโรเจสุํ.\csromanspacer{} น ภิกฺขเว กุรุวินฺทกสุตฺติยา นหายิตพฺพํ, โย นหาเยยฺย, อาปตฺติ ทุกฺกฏสฺสาติ.}

\prose{เตน โข ปน สมเยน ฉพฺพคฺคิยา ภิกฺขู วิคฺคยฺห ปริกมฺมํ การาเปนฺติ.\csromanspacer{} มนุสฺสา อุชฺฌายนฺติ ขิยฺยนฺติ วิปาเจนฺติ “ฯเปฯ เสยฺยถาปิ คิหี กามโภคิโน”ติ.\csromanspacer{} ภควโต เอตมตฺถํ อาโรเจสุํ.\csromanspacer{} น ภิกฺขเว วิคฺคยฺห ปริกมฺมํ การาเปตพฺพํ, โย การาเปยฺย, อาปตฺติ ทุกฺกฏสฺสาติ.}

\prose{\csromanfolio{241}เตน โข ปน สมเยน ฉพฺพคฺคิยา ภิกฺขู มลฺลเกน นหายนฺติ.\csromanspacer{} มนุสฺสา อุชฺฌายนฺติ ขิยฺยนฺติ วิปาเจนฺติ “ฯเปฯ เสยฺยถาปิ คิหี กามโภคิโน”ติ.\csromanspacer{} ภควโต เอตมตฺถํ อาโรเจสุํ.\csromanspacer{} น ภิกฺขเว มลฺลเกน นหายิตพฺพํ, โย นหาเยยฺย, อาปตฺติ ทุกฺกฏสฺสาติ.}

\proseitem{244}{เตน โข ปน สมเยน อญฺญตรสฺส ภิกฺขุโน กจฺฉุโรคาพาโธ โหติ, น ตสฺส วินา มลฺลเกน ผาสุ โหติ.\csromanspacer{} ภควโต เอตมตฺถํ อาโรเจสุํ.\csromanspacer{} อนุชานามิ ภิกฺขเว คิลานสฺส อกตมลฺลกนฺติ.}

\prose{เตน โข ปน สมเยน อญฺญตโร ภิกฺขุ ชราทุพฺพโล นหายมาโน น สกฺโกติ อตฺตโน กายํ อุคฺฆํเสตุํ.\csromanspacer{} ภควโต เอตมตฺถํ อาโรเจสุํ.\csromanspacer{} อนุชานามิ ภิกฺขเว อุกฺกาสิกนฺติ.}

\prose{เตน โข ปน สมเยน ภิกฺขู ปิฏฺฐิปริกมฺมํ กาตุํ กุกฺกุจฺจายนฺติ.\csromanspacer{} ภควโต เอตมตฺถํ อาโรเจสุํ.\csromanspacer{} อนุชานามิ ภิกฺขเว ปุถุปาณิกนฺติ.}

\proseitem{245}{เตน โข ปน สมเยน ฉพฺพคฺคิยา ภิกฺขู วลฺลิกํ ธาเรนฺติ ฯเปฯ ปามงฺคํ ธาเรนฺติ ฯเปฯ กณฺฐสุตฺตกํ ธาเรนฺติ ฯเปฯ กฏิสุตฺตกํ ธาเรนฺติ ฯเปฯ โอวฏฺฏิกํ ธาเรนฺติ ฯเปฯ กายุรํ ธาเรนฺติ ฯเปฯ หตฺถาภรณํ ธาเรนฺติ ฯเปฯ องฺคุลิมุทฺทิกํ ธาเรนฺติ.\csromanspacer{} มนุสฺสา อุชฺฌายนฺติ ขิยฺยนฺติ วิปาเจนฺติ “เสยฺยถาปิ คิหี กามโภคิโน”ติ.\csromanspacer{} อสฺโสสุํ โข ภิกฺขู เตสํ มนุสฺสานํ อุชฺฌายนฺตานํ ขิยฺยนฺตานํ วิปาเจนฺตานํ.\csromanspacer{} อถ โข เต ภิกฺขู ภควโต เอตมตฺถํ อาโรเจสุํ.\csromanspacer{} สจฺจํ กิร ภิกฺขเว ฉพฺพคฺคิยา ภิกฺขู วลฺลิกํ ธาเรนฺติ ฯเปฯ ปามงฺคํ ธาเรนฺติ.\csromanspacer{} กณฺฐสุตฺตกํ ธาเรนฺติ.\csromanspacer{} กฏิสุตฺตกํ ธาเรนฺติ.\csromanspacer{} โอวฏฺฏิกํ ธาเรนฺติ.\csromanspacer{} กายุรํ ธาเรนฺติ.\csromanspacer{} หตฺถาภรณํ ธาเรนฺติ.\csromanspacer{} องฺคุลิมุทฺทิกํ ธาเรนฺตีติ.\csromanspacer{} สจฺจํ ภควาติ ฯเปฯ วิครหิตฺวา ฯเปฯ ธมฺมิํ กถํ กตฺวา ภิกฺขู อามนฺเตสิ “น ภิกฺขเว วลฺลิกา ธาเรตพฺพา ฯเปฯ น ปามงฺโค ธาเรตพฺโพ.\csromanspacer{} น กณฺฐสุตฺตกํ ธาเรตพฺพํ.\csromanspacer{} น กฏิสุตฺตกํ ธาเรตพฺพํ.\csromanspacer{} น โอวฏฺฏิกํ ธาเรตพฺพํ.\csromanspacer{} น กายุรํ ธาเรตพฺพํ.\csromanspacer{} น หตฺถาภรณํ ธาเรตพฺพํ.\csromanspacer{} น องฺคุลิมุทฺทิกา ธาเรตพฺพา, โย ธาเรยฺย, อาปตฺติ ทุกฺกฏสฺสา”ติ.}

\proseitem{246}{เตน โข ปน สมเยน ฉพฺพคฺคิยา ภิกฺขู ทีเฆ เกเส ธาเรนฺติ.\csromanspacer{} มนุสฺสา อุชฺฌายนฺติ ขิยฺยนฺติ วิปาเจนฺติ “เสยฺยถาปิ คิหี กามโภคิโน”ติ.\csromanspacer{} ภควโต เอตมตฺถํ อาโรเจสุํ.\csromanspacer{} น ภิกฺขเว ทีฆา \csromanfolio{242}เกสา ธาเรตพฺพา, โย ธาเรยฺย, อาปตฺติ ทุกฺกฏสฺส.\csromanspacer{} อนุชานามิ ภิกฺขเว ทุมาสิกํ วา ทุวงฺคุลํ วาติ.}

\prose{เตน โข ปน สมเยน ฉพฺพคฺคิยา ภิกฺขู โกจฺเฉน เกเส โอสณฺเฐนฺติ\footnote{โอสเณฺหนฺติ (สี, สฺยา)} ฯเปฯ ผณเกน เกเส โอสณฺเฐนฺติ.\csromanspacer{} หตฺถผณเกน เกเส โอสณฺฐนฺติ.\csromanspacer{} สิตฺถเตลเกน เกเส โอสณฺเฐนฺติ.\csromanspacer{} อุทกเตลเกน เกเส โอสณฺเฐนฺติ.\csromanspacer{} มนุสฺสา อุชฺฌายนฺติ ขิยฺยนฺติ วิปาเจนฺติ “เสยฺยถาปิ คิหี กามโภคิโน”ติ.\csromanspacer{} ภควโต เอตมตฺถํ อาโรเจสุํ.\csromanspacer{} น ภิกฺขเว โกจฺเฉน เกสา โอสณฺเฐตพฺพา ฯเปฯ.\csromanspacer{} น สิตฺถเตลเกน เกสา โอสณฺเฐตพฺพา.\csromanspacer{} น อุทกเตลเกน เกสา โอสณฺเฐตพฺพา, โย โอสณฺเฐยฺย, อาปตฺติ ทุกฺกฏสฺสาติ.}

\proseitem{247}{เตน โข ปน สมเยน ฉพฺพคฺคิยา ภิกฺขู อาทาเสปิ อุทกปตฺเตปิ มุขนิมิตฺตํ โอโลเกนฺติ.\csromanspacer{} มนุสฺสา อุชฺฌายนฺติ ขิยฺยนฺติ วิปาเจนฺติ “เสยฺยถาปิ คิหี กามโภคิโน”ติ.\csromanspacer{} ภควโต เอตมตฺถํ อาโรเจสุํ.\csromanspacer{} น ภิกฺขเว อาทาเส วา อุทกปตฺเต วา มุขนิมิตฺตํ โอโลเกตพฺพํ, โย โอโลเกยฺย, อาปตฺติ ทุกฺกฏสฺสาติ.}

\prose{เตน โข ปน สมเยน อญฺญตรสฺส ภิกฺขุโน มุเข วโณ โหติ, โส ภิกฺขู ปุจฺฉิ “กีทิโส เม อาวุโส วโณ”ติ.\csromanspacer{} ภิกฺขู เอวมาหํสุ “เอทิโส เต อาวุโส วโณ”ติ.\csromanspacer{} โส น สทฺทหติ.\csromanspacer{} ภควโต เอตมตฺถํ อาโรเจสุํ.\csromanspacer{} อนุชานามี ภิกฺขเว อาพาธปฺปจฺจยา อาทาเส วา อุทกปตฺเต วา มุขนิมิตฺตํ โอโลเกตุนฺติ.}

\prose{เตน โข ปน สมเยน ฉพฺพคฺคิยา ภิกฺขู มุขํ อาลิมฺปนฺติ ฯเปฯ มุขํ อุมฺมทฺเทนฺติ.\csromanspacer{} มุขํ จุณฺเณนฺติ.\csromanspacer{} มโนสิลิกาย มุขํ ลญฺเฉนฺติ.\csromanspacer{} องฺคราคํ กโรนฺติ.\csromanspacer{} มุขราคํ กโรนฺติ.\csromanspacer{} องฺคราคมุขราคํ กโรนฺติ.\csromanspacer{} มนุสฺสา อุชฺฌายนฺติ ขิยฺยนฺติ วิปาเจนฺติ “เสยฺยถาปิ คิหี กามโภคิโน”ติ.\csromanspacer{} ภควโต เอตมตฺถํ อาโรเจสุํ.\csromanspacer{} น ภิกฺขเว มุขํ อาลิมฺปิตพฺพํ ฯเปฯ น มุขํ อุมฺมทฺทิตพฺพํ.\csromanspacer{} น มุขํ จุณฺเณตพฺพํ.\csromanspacer{} น มโนสิลิกาย มุขํ ลญฺเฉ ตพฺพํ.\csromanspacer{} น องฺคราโค กาตพฺโพ.\csromanspacer{} น มุขราโค กาตพฺโพ.\csromanspacer{} น องฺคราคมุขราโค กาตพฺโพ, โย กเรยฺย, อาปตฺติ ทุกฺกฏสฺสาติ.}

\prose{\csromanfolio{243}เตน โข ปน สมเยน อญฺญตรสฺส ภิกฺขุโน จกฺขุโรคาพาโธ โหติ.\csromanspacer{} ภควโต เอตมตฺถํ อาโรเจสุํ.\csromanspacer{} อนุชานามิ ภิกฺขเว อาพาธปฺปจฺจยา มุขํ อาลิมฺปิตุนฺติ.}

\proseitem{248}{เตน โข ปน สมเยน ราชคเห คิรคฺคสมชฺโช\footnote{สมชฺชา (อภิธานคนฺเถสุ)} โหติ, ฉพฺพคฺคิยา ภิกฺขู คิรคฺคสมชฺชํ ทสฺสนาย อคมํสุ.\csromanspacer{} มนุสฺสา อุชฺฌายนฺติ ขิยฺยนฺติ วิปาเจนฺติ “กถํ หิ นาม สมณา สกฺยปุตฺติยา นจฺจมฺปิ คีตมฺปิ วาทิตมฺปิ ทสฺสนาย คจฺฉิสฺสนฺติ เสยฺยถาปิ คิหี กามโภคิโน”ติ.\csromanspacer{} ภควโต เอตมตฺถํ อาโรเจสุํ.\csromanspacer{} น ภิกฺขเว นจฺจํ วา คีตํ วา วาทิตํ วา ทสฺสนาย คนฺตพฺพํ, โย คจฺเฉยฺย, อาปตฺติ ทุกฺกฏสฺสาติ.}

\proseitem{249}{เตน โข ปน สมเยน ฉพฺพคฺคิยา ภิกฺขู อายตเกน คีตสฺสเรน ธมฺมํ คายนฺติ.\csromanspacer{} มนุสฺสา อุชฺฌายนฺติ ขิยฺยนฺติ วิปาเจนฺติ “ยเถว\footnote{ยถา จ (ก)} มยํ คายาม, เอวเมวิเม สมณา สกฺยปุตฺติยา อายตเกน คีตสฺสเรน ธมฺมํ คายนฺตี”ติ.\csromanspacer{} อสฺโสสุํ โข ภิกฺขู เตสํ มนุสฺสานํ อุชฺฌายนฺตานํ ขิยฺยนฺตานํ วิปาเจนฺตานํ.\csromanspacer{} เย เต ภิกฺขู อปฺปิจฺฉา ฯเปฯ เต อุชฺฌายนฺติ ขิยฺยนฺติ วิปาเจนฺติ “กถํ หิ นาม ฉพฺพคฺคิยา ภิกฺขู อายตเกน คีตสฺสเรน ธมฺมํ คายิสฺสนฺตี”ติ.\csromanspacer{} อถ โข เต ภิกฺขู ภควโต เอตมตฺถํ อาโรเจสุํ.\csromanspacer{} สจฺจํ กิร ภิกฺขเว ฯเปฯ.\csromanspacer{} สจฺจํ ภควาติ ฯเปฯ ธมฺมิํ กถํ กตฺวา ภิกฺขู อามนฺเตสิ “ปญฺจิเม ภิกฺขเว อาทีนวา อายตเกน คีตสฺสเรน ธมฺมํ คายนฺตสฺส.\csromanspacer{} อตฺตนาปิ ตสฺมิํ สเร สารชฺชติ, ปเรปิ ตสฺมิํ สเร สารชฺชนฺติ, คหปติกาปิ อุชฺฌายนฺติ, สรกุตฺติมฺปิ นิกามยมานสฺส สมาธิสฺส ภงฺโค โหติ, ปจฺฉิมา ชนตา ทิฏฺฐานุคติํ อาปชฺชติ.\csromanspacer{} อิเม โข ภิกฺขเว ปญฺจ อาทีนวา อายตเกน คีตสฺสเรน ธมฺมํ คายนฺตสฺส.\csromanspacer{} น ภิกฺขเว อายตเกน คีตสฺสเรน ธมฺโม คายิตพฺโพ, โย คาเยยฺย, อาปตฺติ ทุกฺกฏสฺสา”ติ.}

\prose{เตน โข ปน สมเยน ภิกฺขู สรภญฺเญ กุกฺกุจฺจายนฺติ.\csromanspacer{} ภควโต เอตมตฺถํ อาโรเจสุํ.\csromanspacer{} อนุชานามิ ภิกฺขเว สรภญฺญนฺติ.}

\prose{เตน โข ปน สมเยน ฉพฺพคฺคิยา ภิกฺขู พาหิรโลมิํ\footnote{พาหิยโลมิํ (ก)} อุณฺณิํ ธาเรนฺติ.\csromanspacer{} มนุสฺสา อุชฺฌายนฺติ ขิยฺยนฺติ วิปาเจนฺติ “เสยฺยถาปิ คิหี \csromanfolio{244}กามโภคิโน”ติ.\csromanspacer{} ภควโต เอตมตฺถํ อาโรเจสุํ.\csromanspacer{} น ภิกฺขเว พาหิยโลมิ อุณฺณิ ธาเรตพฺพา, โย ธาเรยฺย, อาปตฺติ ทุกฺกฏสฺสาติ.}

\proseitem{250}{เตน โข ปน สมเยน รญฺโญ มาคธสฺส เสนิยสฺส พิมฺพิสารสฺส อาราเม อมฺพา ผลิโน โหนฺติ.\csromanspacer{} รญฺโญ มาคเธน เสนิเยน พิมฺพิสาเรน อนุญฺญาตํ โหติ “ยถาสุขํ อยฺยา อมฺพํ ปริภุญฺชนฺตู”ติ.\csromanspacer{} ฉพฺพคฺคิยา ภิกฺขู ตรุณญฺเญว อมฺพํ ปาตาเปตฺวา ปริภุญฺชิํสุ.\csromanspacer{} รญฺโญ จ มาคธสฺส เสนิยสฺส พิมฺพิสารสฺส อมฺเพน อตฺโถ โหติ.\csromanspacer{} อถ โข ราชา มาคโธ เสนิโย พิมฺพิสาโร มนุสฺเส อาณาเปสิ “คจฺฉถ ภเณ, อารามํ คนฺตฺวา อมฺพํ อาหรถา”ติ. “เอวํ เทวา”ติ โข เต มนุสฺสา รญฺโญ มาคธสฺส เสนิยสฺส พิมฺพิสารสฺส ปฏิสฺสุตฺวา อารามํ คนฺตฺวา อารามปาลํ เอตทโวจุํ “เทวสฺส ภเณ อมฺเพน อตฺโถ, อมฺพํ เทถา”ติ.\csromanspacer{} นตฺถายฺยา อมฺพํ, ตรุณญฺเญว อมฺพํ ปาตาเปตฺวา ภิกฺขู ปริภุญฺชิํสูติ.\csromanspacer{} อถ โข เต มนุสฺสา รญฺโญ มาคธสฺส เสนิยสฺส พิมฺพิสารสฺส เอตมตฺถํ อาโรเจสุํ.\csromanspacer{} สุปริภุตฺตํ ภเณ อยฺเยหิ อมฺพํ, อปิ จ ภควตา มตฺตา วณฺณิตาติ.\csromanspacer{} มนุสฺสา อุชฺฌายนฺติ ขิยฺยนฺติ วิปาเจนฺติ “กถํ หิ นาม สมณา สกฺยปุตฺติยา น มตฺตํ ชานิตฺวา รญฺโญ อมฺพํ ปริภุญฺชิสฺสนฺตี”ติ.\csromanspacer{} อสฺโสสุํ โข ภิกฺขู เตสํ มนุสฺสานํ อุชฺฌายนฺตานํ ขิยฺยนฺตานํ วิปาเจนฺตานํ.\csromanspacer{} อถ โข เต ภิกฺขู ภควโต เอตมตฺถํ อาโรเจสุํ.\csromanspacer{} น ภิกฺขเว อมฺพํ ปริภุญฺชิตพฺพํ, โย ปริภุญฺเชยฺย, อาปตฺติ ทุกฺกฏสฺสาติ.}

\prose{เตน โข ปน สมเยน อญฺญตรสฺส ปูคสฺส สํฆภตฺตํ โหติ, สูเป อมฺพเปสิกาโย ปกฺขิตฺตา โหนฺติ, ภิกฺขู กุกฺกุจฺจายนฺตา นปฺปฏิคฺคณฺหนฺติ - ป-.\csromanspacer{} ปฏิคฺคณฺหถ ภิกฺขเว ปริภุญฺชถ, อนุชานามิ ภิกฺขเว อมฺพเปสิกนฺติ.}

\prose{เตน โข ปน สมเยน อญฺญตรสฺส ปูคสฺส สํฆภตฺตํ โหติ, เต น ปริยาปุณิํสุ อมฺพเปสิกํ กาตุํ, ภตฺตคฺเค สกเลเหว อมฺเพหิ เทนฺติ, ภิกฺขู กุกฺกุจฺจายนฺตา นปฺปฏิคฺคณฺหนฺติ.\csromanspacer{} ปฏิคฺคณฺหถ ภิกฺขเว ปริภุญฺชถ, อนุชานามิ ภิกฺขเว ปญฺจหิ สมณกปฺเปหิ ผลํ ปริภุญฺชิตุํ, \csromanfolio{245}อคฺคิปริจิตํ สตฺถปริจิตํ นขปริวิตํ อพีชํ นิพฺพตฺตพีชญฺเญว\footnote{นิพฺพฏฺฏพีชํ (สี, สฺยา)} ปญฺจมํ.\csromanspacer{} อนาชานามิ ภิกฺขเว อิเมหิ ปญฺจหิ สมณกปฺเปหิ ผลํ ปริภุญฺชิตุนฺติ.}

\proseitem{251}{เตน โข ปน สมเยน อญฺญตโร ภิกฺขุ อหินา ทฏฺโฐ กาลงฺกโต โหติ.\csromanspacer{} ภควโต เอตมตฺถํ อาโรเจสุํ.\csromanspacer{} น หิ นูน โส ภิกฺขเว ภิกฺขุ อิมานิ จตฺตาริ อหิราชกุลานิ เมตฺเตน จิตฺเตน ผริ.\csromanspacer{} สเจ หิ โส ภิกฺขเว ภิกฺขุ อิมานิ จตฺตาริ อหิราชกุลานิ เมตฺเตน จิตฺเตน ผเรยฺย, น หิ โส ภิกฺขเว ภิกฺขุ อหินา ทฏฺโฐ กาลํ กเรยฺย.\csromanspacer{} กตมานิ จตฺตาริ อหิราชกุลานิ, วิรูปกฺขํ อหิราชกุลํ, เอราปถํ อหิราชกุลํ, ฉพฺยาปุตฺตํ อหิราชกุลํ, กณฺหา โคตมํ อหิราชกุลํ.\csromanspacer{} น หิ นูน โส ภิกฺขเว ภิกฺขุ อิมานิ จตฺตาริ อหิราชกุลานิ เมตฺเตน จิตฺเตน ผริ.\csromanspacer{} สเจ หิ โส ภิกฺขเว ภิกฺขุ อิมานิ จตฺตาริ อหิราชกุลานิ เมตฺเตน จิตฺเตน ผเรยฺย, น หิ โส ภิกฺขเว ภิกฺขุ อหินา ทฏฺโฐ กาลํ กเรยฺย.\csromanspacer{} อนุชานามิ ภิกฺขเว อิมานิ จตฺตาริ อหิราชกุลานิ เมตฺเตน จิตฺเตน ผริตุํ อตฺตคุตฺติยา อตฺตรกฺขาย อตฺตปริตฺตํ กาตุํ.\csromanspacer{} เอวญฺจ ปน ภิกฺขเว กาตพฺพํ\csromandash{}}

\prose{\csromansymbolfootnote{*}{ขุ 5. 53; อํ 1. 384 ปิฏฺเฐสุปิ อิทํ วตฺถุ อาคตํ. 2. สริํสปานิ (สี, สฺยา)}“วิรูปกฺเขหิ เม เมตฺตํ, เมตฺตํ เอราปเถหิ เม.}

\prose{ฉพฺยาปุตฺเตหิ เม เมตฺตํ, เมตฺตํ กณฺหาโคตมเกหิ จ.}

\prose{อปาทเกหิ เม เมตฺตํ, เมตฺตํ ทฺวิปาทเกหิ เม.}

\prose{จตุปฺปเทหิ เม เมตฺตํ, เมตฺตํ พหุปฺปเทหิ เม.}

\prose{มา มํ อปาทโก หิํสิ, มา มํ หิํสิ ทฺวิปาทโก.}

\prose{มา มํ จตุปฺปโท หิํสิ, มา มํ หิํสิ พหุปฺปโท.}

\prose{สพฺเพ สตฺตา สพฺเพ ปาณา, สพฺเพ ภูตา จ เกวลา.}

\prose{สพฺเพ ภทฺรานิ ปสฺสนฺตุ, มา กญฺจิ ปาปมาคมา.}

\prose{อปฺปมาโณ พุทฺโธ, อปฺปมาโณ ธมฺโม.}

\prose{อปฺปมาโณ สํโฆ, ปมาณวนฺตานิ สรีสปานิ2.}

\prose{อหิ วิจฺฉิกา สตปที อุณฺณนาภิ สรพู มูสิกา.}

\prose{กตา เม รกฺขา กตํ เม ปริตฺตํ, ปฏิกฺกมนฺตุ ภูตานิ.}

\prose{โสหํ นโม ภควโต, นโม สตฺตนฺนํ สมฺมาสมฺพุทฺธานนฺ”ติ.}

\prose{\csromanfolio{246}เตน โข ปน สมเยน อญฺญตโร ภิกฺขุ อนภิรติยา ปีฬิโต อตฺตโน องฺคชาตํ ฉินฺทิ.\csromanspacer{} ภควโต เอตมตฺถํ อาโรเจสุํ.\csromanspacer{} อญฺญมฺหิ โส ภิกฺขเว โมฆปุริโส เฉตพฺพมฺหิ อญฺญํ ฉินฺทิ, น ภิกฺขเว อตฺตโน องฺคชาตํ เฉตพฺพํ, โย ฉินฺเทยฺย, อาปตฺติ ถุลฺลจฺจยสฺสาติ.}

\proseitem{252}{เตน โข ปน สมเยน ราชคหกสฺส เสฏฺฐิสฺส มหคฺฆสฺส จนฺทนสฺส\footnote{จนฺทนสารสฺส (สี, สฺยา)} จนฺทนคณฺฐิ อุปฺปนฺนา โหติ.\csromanspacer{} อถ โข ราชคหกสฺส เสฏฺฐิสฺส เอตทโหสิ “ยํนูนาหํ อิมาย จนฺทนคณฺฐิยา ปตฺตํ เลขาเปยฺยํ, เลขญฺจ เม ปริโภคํ ภวิสฺสติ, ปตฺตญฺจ ทานํ ทสฺสามี”ติ.\csromanspacer{} อถ โข ราชคหโก เสฏฺฐิ ตาย จนฺทนคณฺฐิยา ปตฺตํ เลขาเปตฺวา สิกฺกาย อุฑฺฑิตฺวา\footnote{วาหิตฺวา (สี)} เวฬคฺเค อาลคฺเคตฺวา เวฬุปรมฺปราย พนฺธิตฺวา\footnote{วาหิตฺวา (สฺยา)} เอวมาห “โย สมโณ วา พฺราหฺมโณ วา อรหา เจว อิทฺธิมา จ, ทินฺนํเยว ปตฺตํ, โอหรตู”ติ.\csromanspacer{} อถ โข ปูรโณ กสฺสโป เยน ราชคหโก เสฏฺฐิ เตนุปสงฺกมิ, อุปสงฺกมิตฺวา ราชคหกํ เสฏฺฐิํ เอตทโวจ “อหํ หิ คหปติ อรหา เจว อิทฺธิมา จ, เทหิ เม ปตฺตนฺ”ติ.\csromanspacer{} สเจ ภนฺเต อายสฺมา อรหา เจว อิทฺธิมา จ, ทินฺนํเยว ปตฺตํ, โอหรตูติ.\csromanspacer{} อถ โข มกฺขลิ โคสาโล.\csromanspacer{} อชิโต เกสกมฺพโล.\csromanspacer{} ปกุโธ กจฺจายโน.\csromanspacer{} สญฺจโย เพลฏฺฐปุตฺโต\footnote{สญฺชโย เพลฺลฏฺฐิปุตฺโต (สี)}.\csromanspacer{} นิคณฺโฐ นาฏปุตฺโต\footnote{นาถปุตฺโต (สี)} เยน ราชคหโก เสฏฺฐิ เตนุปสงฺกมิ, อุปสงฺกมิตฺวา ราชคหกํ เสฏฺฐิํ เอตทโวจ “อหํ หิ คหปติ อรหา เจว อิทฺธิมา จ, เทหิ เม ปตฺตนฺ”ติ.\csromanspacer{} สเจ ภนฺเต อายสฺมา อรหา เจว อิทฺธิมา จ, ทินฺนํเยว ปตฺตํ โอหรตูติ.}

\prose{เตน โข ปน สมเยน อายสฺมา จ มหาโมคฺคลฺลาโน อายสฺมา จ ปิณฺโฑลภารทฺวาโช ปุพฺพณฺหสมยํ นิวาเสตฺวา ปตฺตจีวรมาทาย ราชคหํ ปิณฺฑายา ปวิสิํสุ.\csromanspacer{} อถ โข อายสฺมา ปิณฺโฑลภารทฺวาโช อายสฺมนฺตํ มหาโมคฺคลฺลานํ เอตทโวจ “อายสฺมา โข มหาโมคฺคลฺลาโน อรหา เจว อิทฺธิมา จ, คจฺฉาวุโส โมคฺคลฺลาน, เอตํ ปตฺตํ โอหร, ตุเยฺหโส ปตฺโต”ติ.\csromanspacer{} อายสฺมา โข ภารทฺวาโช อรหา เจว อิทฺธิมา จ, คจฺฉาวุโส ภารทฺวาช เอตํ ปตฺตํ โอหร, ตุเยฺหโส ปตฺโตติ.\csromanspacer{} อถ โข อายสฺมา ปิณฺโฑลภารทฺวาโช เวหาสํ อพฺภุคฺคนฺตฺวา ตํ ปตฺตํ คเหตฺวา ติกฺขตฺตุํ ราชคหํ อนุปริยายิ.}

\prose{\csromanfolio{247}เตน โข ปน สมเยน ราชคหโก เสฏฺฐิ สปุตฺตทาโร สเก นิเวสเน ฐิโต โหติ ปญฺชลิโก นมสฺสมาโน “อิเธว ภนฺเต อยฺโย ภารทฺวาโช อมฺหากํ นิเวสเน ปติฏฺฐาตู”ติ.\csromanspacer{} อถ โข อายสฺมา ปิณฺโฑลภารทฺวาโช ราชคหกสฺส เสฏฺฐิสฺส นิเวสเน ปติฏฺฐาสิ.\csromanspacer{} อถ โข ราชคหโก เสฏฺฐิ อายสฺมโต ปิณฺโฑลภารทฺวาชสฺส หตฺถโต ปตฺตํ คเหตฺวา มหคฺฆสฺส ขาทนียสฺส ปูเรตฺวา อายสฺมโต ปิณฺโฑลภารทฺวาชสฺส อทาสิ.\csromanspacer{} อถ โข อายสฺมา ปิณฺโฑลภารทฺวาโช ตํ ปตฺตํ คเหตฺวา อารามํ อคมาสิ.\csromanspacer{} อสฺโสสุํ โข มนุสฺสา “อยฺเยน กิร ปิณฺโฑลภารทฺวาเชน ราชคหกสฺส เสฏฺฐิสฺส ปตฺโต โอหาริโต”ติ.\csromanspacer{} เต จ มนุสฺสา อุจฺจาสทฺทา มหาสทฺทา อายสฺมนฺตํ ปิณฺโฑลภารทฺวาชํ ปิฏฺฐิโต ปิฏฺฐิโต อนุพนฺธิํสุ.}

\prose{อสฺโสสิ โข ภควา อุจฺจาสทฺทํ มหาสทฺทํ, สุตฺวาน อายสฺมนฺตํ อานนฺทํ อามนฺเตสิ “กิํ นุ โข โส อานนฺท อุจฺจาสทฺโท มหาสทฺโท”ติ.\csromanspacer{} อายสฺมตา ภนฺเต ปิณฺโฑลภารทฺวาเชน ราชคหกสฺส เสฏฺฐิสฺส ปตฺโต โอหาริโต, อสฺโสสุํ โข ภนฺเต มนุสฺสา อยฺเยน กิร ปิณฺโฑลภารทฺวาเชน ราชคหกสฺส เสฏฺฐิสฺส ปตฺโต โอหาริโตติ, เต จ ภนฺเต มนุสฺสา อุจฺจาสทฺทา มหาสทฺทา อายสฺมนฺตํ ปิณฺโฑลภารทฺวาชํ ปิฏฺฐิโต ปิฏฺฐิโต อนุพนฺธา, โส เอโส ภนฺเต ภควา อุจฺจาสทฺโท มหาสทฺโทติ.\csromanspacer{} อถ โข ภควา เอตสฺมิํ นิทาเน เอตสฺมิํ ปกรเณ ภิกฺขุสํฆํ สนฺนิปาตาเปตฺวา อายสฺมนฺตํ ปิณฺโฑลภารทฺวาชํ ปฏิปุจฺฉิ “สจฺจํ กิร ตยา ภารทฺวาช ราชคหกสฺส เสฏฺฐิสฺส ปตฺโต โอหาริโต”ติ.\csromanspacer{} สจฺจํ ภควาติ.\csromanspacer{} วิครหิ พุทฺโธ ภควา อนนุจฺฉวิกํ ภารทฺวาช อนนุโลมิกํ อปฺปติรูปํ อสฺสามณกํ อกปฺปิยํ อกรณียํ, กถํ หิ นาม ตฺวํ ภารทฺวาช ฉวสฺส ทารุปกฺกสฺส การณา คิหีนํ อุตฺตริมนุสฺสธมฺมํ อิทฺธิปาฏิหาริยํ ทสฺเสสฺสสิ, เสยฺยถาปิ ภารทฺวาช มาตุคาโม ฉวสฺส มาสกรูปสฺส การณา โกปินํ ทสฺเสติ, เอวเมว โข ตยา ภารทฺวาช ฉวสฺส ทารุปตฺตสฺส การณา คิหีนํ อุตฺตริมนุสฺสธมฺมํ อิทฺธิปาฏิหาริยํ ทสฺสิตํ.\csromanspacer{} เนตํ ภารทฺวาช อปฺปสนฺนานํ วา ปสาทาย - ป- วิครหิตฺวา ฯเปฯ ธมฺมิํ กถํ กตฺวา ภิกฺขู อามนฺเตสิ “น ภิกฺขเว คิหีนํ อุตฺตริมนุสฺสธมฺมํ อิทฺธิปาฏิหาริยํ ทสฺเสตพฺพํ, โย ทสฺเสยฺย, อาปตฺติ ทุกฺกฏสฺส.\csromanspacer{} ภินฺทเถตํ ภิกฺขเว ทารุปตฺตํ, สกลิกํ สกลิกํ \csromanfolio{248}กตฺวา ภิกฺขูนํ อญฺชนุปปิสนํ เทถ.\csromanspacer{} น จ ภิกฺขเว ทารุปตฺโต ธาเรตพฺโพ, โย ธาเรยฺย, อาปตฺติ ทุกฺกฏสฺสา”ติ.}

\prose{เตน โข ปน สมเยน ฉพฺพคฺคิยา ภิกฺขู อุจฺจาวเจ ปตฺเต ธาเรนฺติ โสวณฺณมยํ รูปิยมยํ.\csromanspacer{} มนุสฺสา อุชฺฌายนฺติ ขิยฺยนฺติ วิปาเจนฺติ “เสยฺยถาปิ คิหี กามโภคิโน”ติ.\csromanspacer{} ภควโต เอตมตฺถํ อาโรเจสุํ.\csromanspacer{} น ภิกฺขเว โสวณฺณมโย ปตฺโต ธาเรตพฺโพ ฯเปฯ.\csromanspacer{} น รูปิยมโย ปตฺโต ธาเรตพฺโพ.\csromanspacer{} น มณิมโย ปตฺโต ธาเรตพฺโพ.\csromanspacer{} น เวฬุริยมโย ปตฺโต ธาเรตพฺโพ.\csromanspacer{} น ผลิกมโย ปตฺโต ธาเรตพฺโพ.\csromanspacer{} น กํสมโย ปตฺโต ธาเรตพฺโพ.\csromanspacer{} น กาจมโย ปตฺโต ธาเรตพฺโพ.\csromanspacer{} น ติปุมโย ปตฺโต ธาเรตพฺโพ.\csromanspacer{} น สีสมโย ปตฺโต ธาเรตพฺโพ.\csromanspacer{} น ตมฺพโลหมโย ปตฺโต ธาเรตพฺโพ, โย ธาเรยฺย, อาปตฺติ ทุกฺกฏสฺส.\csromanspacer{} อนุชานามิ ภิกฺขเว เทฺว ปตฺเต อโยปตฺตํ มตฺติกาปตฺตนฺติ.}

\proseitem{253}{เตน โข ปน สมเยน ปตฺตมูลํ ฆํสิยติ.\csromanspacer{} ภควโต เอตมตฺถํ อาโรเจสุํ.\csromanspacer{} อนุชานามิ ภิกฺขเว ปตฺตมณฺฑลนฺติ.}

\prose{เตน โข ปน สมเยน ฉพฺพคฺคิยา ภิกฺขู อุจฺจาวจานิ ปตฺตมณฺฑลานิ ธาเรนฺติ โสวณฺณมยํ รูปิยมยํ.\csromanspacer{} มนุสฺสา อุชฺฌายนฺติ ขิยฺยนฺติ วิปาเจนฺติ “เสยฺยถาปิ คิหี กามโภคิโน”ติ, ภควโต เอตมตฺถํ อาโรเจสุํ.\csromanspacer{} น ภิกฺขเว อุจฺจาวจานิ ปตฺตมณฺฑลานิ ธาเรตพฺพานิ, โย ธาเรยฺย, อาปตฺติ ทุกฺกฏสฺส.\csromanspacer{} อนุชานามิ ภิกฺขเว เทฺว ปตฺตมณฺฑลานิ ติปุมยํ สีสมยนฺติ.\csromanspacer{} พหลานิ มณฺฑลานิ น อจฺฉุปิยนฺติ.\csromanspacer{} ภควโต เอตมตฺถํ อาโรเจสุํ.\csromanspacer{} อนุชานามิ ภิกฺขเว ลิขิตุนฺติ.\csromanspacer{} วลี\footnote{จลี (ก)} โหนฺติ.\csromanspacer{} ภควโต เอตมตฺถํ อาโรเจสุํ.\csromanspacer{} อนุชานามิ ภิกฺขเว มกรทนฺตกํ ฉินฺทิตุนฺติ.}

\prose{เตน โข ปน สมเยน ฉพฺพคฺคิยา ภิกฺขู จิตฺรานิ ปตฺตมณฺฑลานิ ธาเรนฺติ รูปกากิณฺณานิ ภิตฺติกมฺมกตานิ, ตานิ รถิกายปิ ทสฺเสนฺตา อาหิณฺฑนฺติ.\csromanspacer{} มนุสฺสา อุชฺฌายนฺติ ขิยฺยนฺติ วิปาเจนฺติ “เสยฺยถาปิ คิหี กามโภคิโน”ติ.\csromanspacer{} ภควโต เอตมตฺถํ อาโรเจสุํ.\csromanspacer{} น ภิกฺขเว จิตฺรานิ ปตฺตมณฺฑลานิ ธาเรตพฺพานิ รูปกากิณฺณานิ ภิตฺติกมฺมกตานิ, โย ธาเรยฺย, อาปตฺติ ทุกฺกฏสฺส.\csromanspacer{} อนุชานามิ ภิกฺขเว ปกติมณฺฑลนฺติ.}

\proseitem{254}{\csromanfolio{249}เตน โข ปน สมเยน ภิกฺขู โสทกํ ปตฺตํ ปฏิสาเมนฺติ, ปตฺโต ทุสฺสติ.\csromanspacer{} ภควโต เอตมตฺถํ อาโรเจสุํ.\csromanspacer{} น ภิกฺขเว โสทโก ปตฺโต ปฏิสาเมตพฺโพ, โย ปฏิสาเมยฺย, อาปตฺติ ทุกฺกฏสฺส.\csromanspacer{} อนุชานามิ ภิกฺขเว โอตาเปตฺวา ปตฺตํ ปฏิสาเมตุนฺติ.}

\prose{เตน โข ปน สมเยน ภิกฺขู โสทกํ\footnote{ส-อุทกํ (ก)} ปตฺตํ โอตาเปนฺติ, ปตฺโต ทุคฺคนฺโธ โหติ.\csromanspacer{} ภควโต เอตมตฺถํ อาโรเจสุํ.\csromanspacer{} น ภิกฺขเว โสทโก ปตฺโต โอตาเปตพฺโพ, โย โอตาเปยฺย, อาปตฺติ ทุกฺกฏสฺส.\csromanspacer{} อนุชานามิ ภิกฺขเว โวทกํ กตฺวา โอตาเปตฺวา ปตฺตํ ปฏิสาเมตุนฺติ.}

\prose{เตน โข ปน สมเยน ภิกฺขู อุเณฺห ปตฺตํ นิทหนฺติ, ปตฺตสฺส วณฺโณ ทุสฺสติ.\csromanspacer{} ภควโต เอตมตฺถํ อาโรเจสุํ.\csromanspacer{} น ภิกฺขเว อุเณฺห ปตฺโต นิทหิตพฺโพ, โย นิทเหยฺย, อาปตฺติ ทุกฺกฏสฺส.\csromanspacer{} อนุชานามิ ภิกฺขเว มุหุตฺตํ อุเณฺห โอตาเปตฺวา ปตฺตํ ปฏิสาเมตุนฺติ.}

\prose{เตน โข ปน สมเยน สมฺพหุลา ปตฺตา อชฺโฌกาเส อนาธารา นิกฺขิตฺตา โหนฺติ, วาตมณฺฑลิกาย อาวฏฺเฏตฺวา ปตฺตา ภิชฺชิํสุ.\csromanspacer{} ภควโต เอตมตฺถํ อาโรเจสุํ.\csromanspacer{} อนุชานามิ ภิกฺขเว ปตฺตาธารกนฺติ.}

\prose{เตน โข ปน สมเยน ภิกฺขู มิฑฺฒนฺเต ปตฺตํ นิกฺขิปนฺติ, ปริปติตฺวา\footnote{ปริวฏฺฏิตฺวา (สฺยา)} ปตฺโต ภิชฺชติ.\csromanspacer{} ภควโต เอตมตฺถํ อาโรเจสุํ.\csromanspacer{} น ภิกฺขเว มิฑฺฒนฺเต ปตฺโต นิกฺขิปิตพฺโพ, โย นิกฺขิเปยฺย, อาปตฺติ ทุกฺกฏสฺสาติ.}

\prose{เตน โข ปน สมเยน ภิกฺขู ปริภณฺฑนฺเต ปตฺตํ นิกฺขิปนฺติ, ปริปติตฺวา ปตฺโต ภิชฺชติ.\csromanspacer{} ภควโต เอตมตฺถํ อาโรเจสุํ.\csromanspacer{} น ภิกฺขเว ปริภณฺฑนฺเต ปตฺโต นิกฺขิปิตพฺโพ, โย นิกฺขิเปยฺย, อาปตฺติ ทุกฺกฏสฺสาติ.}

\prose{เตน โข ปน สมเยน ภิกฺขู ฉมาย ปตฺตํ นิกฺกุชฺชนฺติ, โอฏฺโฐ ฆํสิยติ.\csromanspacer{} ภควโต เอตมตฺถํ อาโรเจสุํ.\csromanspacer{} อนุชานามิ ภิกฺขเว ติณสนฺถารกนฺติ.\csromanspacer{} ติณสนฺถารโก อุปจิกาหิ ขชฺชติ.\csromanspacer{} ภควโต เอตมตฺตํ อาโรเจสุํ.\csromanspacer{} อนุชานามิ ภิกฺขเว โจฬกนฺติ.\csromanspacer{} โจฬกํ อุปจิกาหิ ขชฺชติ.\csromanspacer{} ภควโต เอตมตฺถํ อาโรเจสุํ.\csromanspacer{} อนุชานามิ ภิกฺขเว ปตฺตมาฬกนฺติ.\csromanspacer{} ปตฺตมาฬโก ปริปติตฺวา ปตฺโต ภิชฺชติ.\csromanspacer{} ภควโต \csromanfolio{250}เอตมตฺถํ อาโรเจสุํ.\csromanspacer{} อนาชานามิ ภิกฺขเว ปตฺตกุณฺโฑลิกนฺติ.\csromanspacer{} ปตฺตกุณฺโฑลิกาย ปตฺโต ฆํสิยติ.\csromanspacer{} ภควโต เอตมตฺถํ อาโรเจสุํ.\csromanspacer{} อนุชานามิ ภิกฺขเว ปตฺตถวิกนฺติ.\csromanspacer{} อํสพทฺธโก น โหติ.\csromanspacer{} ภควโต เอตมตฺถํ อาโรเจสุํ.\csromanspacer{} อนุชานามิ ภิกฺขเว อํสพทฺธกํ พนฺธนสุตฺตกนฺติ.}

\prose{เตน โข ปน สมเยน ภิกฺขู ภิตฺติขิเลปิ นาคทนฺตเกปิ ปตฺตํ ลคฺเคนฺติ, ปริปติตฺวา ปตฺโต ภิชฺชติ.\csromanspacer{} ภควโต เอตมตฺถํ อาโรเจสุํ.\csromanspacer{} น ภิกฺขเว ปตฺโต ลคฺเคตพฺโพ, โย ลคฺเคยฺย, อาปตฺติ ทุกฺกฏสฺสาติ.}

\prose{เตน โข ปน สมเยน ภิกฺขู มญฺจ ปตฺตํ นิกฺขิปนฺติ, สติสมฺโมสา นิสีทนฺตา โอตฺถริตฺวา ปตฺตํ ภินฺเทนฺติ.\csromanspacer{} ภควโต เอตมตฺถํ อาโรเจสุํ.\csromanspacer{} น ภิกฺขเว มญฺเจ ปตฺโต นิกฺขิปิตพฺโพ, โย นิกฺขิเปยฺย, อาปตฺติ ทุกฺกฏสฺสาติ.}

\prose{เตน โข ปน สมเยน ภิกฺขู ปีเฐ ปตฺตํ นิกฺขิปนฺติ, สติสมฺโมสา นิสีทนฺตา โอตฺถริตฺวา ปตฺตํ ภินฺเทนฺติ.\csromanspacer{} ภควโต เอตมตฺถํ อาโรเจสุํ.\csromanspacer{} น ภิกฺขเว ปีเฐ ปตฺโต นิกฺขิปิตพฺโพ, โย นิกฺขิเปยฺย, อาปตฺติ ทุกฺกฏสฺสาติ.}

\prose{เตน โข ปน สมเยน ภิกฺขู องฺเก ปตฺตํ นิกฺขิปนฺติ, สติสมฺโมสา อุฏฺฐหนฺติ, ปริปติตฺวา ปตฺโต ภิชฺชติ.\csromanspacer{} ภควโต เอตมตฺถํ อาโรเจสุํ.\csromanspacer{} น ภิกฺขเว องฺเก ปตฺโต นิกฺขิปิตพฺโพ, โย นิกฺขิเปยฺย, อาปตฺติ ทุกฺกฏสฺสาติ.}

\prose{เตน โข ปน สมเยน ภิกฺขู ฉตฺเต ปตฺตํ นิกฺขิปนฺติ, วาตมณฺฑลิกาย ฉตฺตํ อุกฺขิปิยติ, ปริปติตฺวา ปตฺโต ภิชฺชติ.\csromanspacer{} ภควโต เอตมตฺถํ อาโรเจสุํ.\csromanspacer{} น ภิกฺขเว ฉตฺเต ปตฺโต นิกฺขิปิตพฺโพ, โย นิกฺขิเปยฺย, อาปตฺติ ทุกฺกฏสฺสาติ.}

\proseitem{255}{เตน โข ปน สมเยน ภิกฺขู ปตฺตหตฺถา กวาฏํ ปณาเมนฺติ, กวาโฏ อาวฏฺฏิตฺวา ปตฺโต ภิชฺชติ.\csromanspacer{} ภควโต เอตมตฺถํ อาโรเจสุํ.\csromanspacer{} น ภิกฺขเว ปตฺตหตฺเถน กวาฏํ ปณาเมตพฺพํ\footnote{กวาโฏ ปณาเมตพฺโพ (ก)}, โย ปณาเมยฺย, อาปตฺติ ทุกฺกฏสฺสาติ.}

\prose{\csromanfolio{251}เตน โข ปน สมเยน ภิกฺขู ตุมฺพกฏาเห ปิณฺฑาย จรนฺติ.\csromanspacer{} มนุสฺสา อุชฺฌายนฺติ ขิยฺยนฺติ วิปาเจนฺติ “เสยฺยถาปิ ติตฺถิยา”ติ.\csromanspacer{} ภควโต เอกมตฺถํ อาโรเจสุํ.\csromanspacer{} น ภิกฺขเว ตุมฺพกฏาเห ปิณฺฑาย จริตพฺพํ, โย จเรยฺย, อาปตฺติ ทุกฺกฏสฺสาติ.}

\prose{เตน โข ปน สมเยน ภิกฺขู ฆฏิกฏาเห\footnote{ฆฏิกฏาเหน (สฺยา)} ปิณฺฑาย จรนฺติ.\csromanspacer{} มนุสฺสา อุชฺฌายนฺติ ขิยฺยนฺติ วิปาเจนฺติ “เสยฺยถาปิ ติตฺถิยา”ติ.\csromanspacer{} ภควโต เอตมตฺถํ อาโรเจสุํ.\csromanspacer{} น ภิกฺขเว ฆฏิกฏาเห ปิณฺฑาย จริตพฺพํ, โย จเรยฺย, อาปตฺติ ทุกฺกฏสฺสาติ.}

\prose{เตน โข ปน สมเยน อญฺญตโร ภิกฺขุ สพฺพปํสุกูลิโก โหติ, โส ฉวสีสสฺส ปตฺตํ ธาเรติ, อญฺญตรา อิตฺถี ปสฺสิตฺวา ภีตา วิสฺสรมกาสิ “อภุํ เม ปิสาโจ วตายนฺ”ติ.\csromanspacer{} มนุสฺสา อุชฺฌายนฺติ ขิยฺยนฺติ วิปาเจนฺติ “กยํ หิ นาม สมณา สกฺยปุตฺติยา ฉวสีสสฺส ปตฺตํ ธาเรสฺสนฺติ, เสยฺยถาปิ ปิสาจิลฺลิกา”ติ.\csromanspacer{} ภควโต เอตมตฺถํ อาโรเจสุํ.\csromanspacer{} น ภิกฺขเว ฉวสีสสฺส ปตฺโต ธาเรตพฺโพ, โย ธาเรยฺย, อาปตฺติ ทุกฺกฏสฺส.\csromanspacer{} น จ ภิกฺขเว สพฺพปํสุกูลิเกน สวิตพฺพํ, โย ภเวยฺย, อาปตฺติ ทุกฺกฏสฺสาติ.}

\prose{เตน โข ปน สมเยน ภิกฺขู จลกานิปิ อฏฺฐิกานิปิ อุจฺฉิฏฺโฐทกมฺปิ ปตฺเตน นีหรนฺติ.\csromanspacer{} มนุสฺสา อุชฺฌายนฺติ ขิยฺยนฺติ วิปาเจนฺติ “ยสฺมิํ เยวิเม สมณา สกฺยปุตฺติยา ภุญฺชนฺติ, โสว เนสํ ปฏิคฺคโห”ติ.\csromanspacer{} ภควโต เอตมตฺถํ อาโรเจสุํ.\csromanspacer{} น ภิกฺขเว จลกานิ วา อฏฺฐิกานิ วา อุจฺฉิฏฺโฐทกํ วา ปตฺเตน นีหริตพฺพํ, โย นีหเรยฺย, อาปตฺติ ทุกฺกฏสฺส.\csromanspacer{} อนุชานามิ ภิกฺขเว ปฏิคฺคหนฺติ.}

\proseitem{256}{เตน โข ปน สมเยน ภิกฺขู หตฺเถน วิปฺผาเฬตฺวา จีวรํ สิพฺเพนฺติ, จีวรํ วิโลมิกํ โหติ.\csromanspacer{} ภควโต เอตมตฺถํ อาโรเจสุํ.\csromanspacer{} อนุชานามิ ภิกฺขเว สตฺถกํ นมตกนฺติ.}

\prose{เตน โข ปน สมเยน สํฆสฺส ทณฺฑสตฺถกํ อุปฺปนฺนํ โหติ.\csromanspacer{} ภควโต เอตมตฺถํ อาโรเจสุํ.\csromanspacer{} อนุชานามิ ภิกฺขเว ทณฺฑสตฺถกนฺติ.}

\prose{\csromanfolio{252}เตน โข ปน สมเยน ฉพฺพคฺคิยา ภิกฺขู อุจฺจาวเจ สตฺถกทณฺเฑ ธาเรนฺติ โสวณฺณมยํ รูปิยมยํ.\csromanspacer{} มนุสฺสา อุชฺฌายนฺติ ขิยฺยนฺติ วิปาเจนฺติ “เสยฺยถาปิ คิหี กามโภคิโน”ติ.\csromanspacer{} ภควโต เอตมตฺถํ อาโรเจสุํ.\csromanspacer{} น ภิกฺขเว อุจฺจาวจา สตฺถกทณฺฑา ธาเรตพฺพา, โย ธาเรยฺย, อาปตฺติ ทุกฺกฏสฺส.\csromanspacer{} อนุชานามิ ภิกฺขเว อฏฺฐิมยํ ทนฺตมยํ วิสาณมยํ นฬมยํ เวฬุมยํ กฏฺฐมยํ ชตุมยํ ผลมยํ โลหมยํ สงฺขนาภิมยนฺติ.}

\prose{เตน โข ปน สมเยน ภิกฺขู กุกฺกุฏปตฺเตนปิ เวฬุเปสิกายปิ จีวรํ สิพฺเพนฺติ, จีวรํ ทุสฺสิพฺพิตํ โหติ.\csromanspacer{} ภควโต เอตมตฺถํ อาโรเจสุํ.\csromanspacer{} อนุชานามิ ภิกฺขเว สูจินฺติ.\csromanspacer{} สูจิโย กณฺณกิตาโย โหนฺติ ฯเปฯ.\csromanspacer{} อนุชานามิ ภิกฺขเว สูจินาฬิกนฺติ.\csromanspacer{} สูจินาฬิกายปิ กณฺณกิตาโย โหนฺติ ฯเปฯ.\csromanspacer{} อนุชานามิ ภิกฺขเว กิณฺเณน ปุเรตุนฺติ.\csromanspacer{} กิณฺเณปิ กณฺณกิตาโย โหนฺติ ฯเปฯ.\csromanspacer{} อนุชานามิ ภิกฺขเว สตฺตุยา ปุเรตุนฺติ.\csromanspacer{} สตฺตุยาปิ กณฺณกิตาโย โหนฺติ ฯเปฯ.\csromanspacer{} อนุชานามิ ภิกฺขเว สริตกนฺติ.\csromanspacer{} สริตเกปิ กณฺณกิตาโย โหนฺติ ฯเปฯ.\csromanspacer{} อนุชานามิ ภิกฺขเว มธุสิตฺถเกน สาเรตุนฺติ.\csromanspacer{} สริตกํ ปริภิชฺชติ ฯเปฯ.\csromanspacer{} อนุชานามิ ภิกฺขเว สริตกสิปาฏิกนฺติ.}

\prose{เตน โข ปน สมเยน ภิกฺขู ตตฺถ ตตฺถ ขิลํ นิกฺขณิตฺวา สมฺพนฺธิตฺวา จีวรํ สิพฺเพนฺติ, จีวรํ วิกณฺณํ โหติ.\csromanspacer{} ภควโต เอตมตฺถํ อาโรเจสุํ.\csromanspacer{} อนุชานามิ ภิกฺขเว กถินํ กถินรชฺชุํ\footnote{กฐินํ กฐินรชฺชุํ (สี, สฺยา)}, ตตฺถ ตตฺถ โอพนฺธิตฺวา จีวรํ สิพฺเพตุนฺติ.\csromanspacer{} วิสเม กถินํ ปตฺถรนฺติ, กถินํ ปริภิชฺชติ ฯเปฯ.\csromanspacer{} น ภิกฺขเว วิสเม กถินํ ปตฺถริตพฺพํ, โย ปตฺถเรยฺย, อาปตฺติ ทุกฺกฏสฺสาติ.}

\prose{ฉมาย กถินํ ปตฺถรนฺติ, กถินํ ปํสุกิตํ โหติ ฯเปฯ.\csromanspacer{} อนุชานามิ ภิกฺขเว ติณสนฺถารกนฺติ.\csromanspacer{} กถินสฺส อนฺโต ชีรติ ฯเปฯ.\csromanspacer{} อนุชานามิ ภิกฺขเว อนุวาตํ ปริภณฺฑํ อาโรเปตุนฺติ.\csromanspacer{} กถินํ นปฺปโหติ ฯเปฯ.\csromanspacer{} อนุชานามิ ภิกฺขเว ทณฺฑกถินํ พิทลกํ สลากํ วินนฺธนรชฺชุํ วินนฺธนสุตฺตํ วินนฺธิตฺวา จีวรํ สิพฺเพตุนฺติ.\csromanspacer{} สุตฺตนฺตริกาโย วิสมา โหนฺติ ฯเปฯ. \csromanfolio{253}อนุชานามิ ภิกฺขเว กฬิมฺภกนฺติ.\csromanspacer{} สุตฺตา วงฺกา โหนฺติ ฯเปฯ.\csromanspacer{} อนุชานามิ ภิกฺขเว โมฆสุตฺตกนฺติ.}

\prose{เตน โข ปน สมเยน ภิกฺขู อโธเตหิ ปาเทหิ กถินํ อกฺกมนฺติ กถินํ ทุสฺสติ.\csromanspacer{} ภควโต เอตมตฺถํ อาโรเจสุํ.\csromanspacer{} น ภิกฺขเว อโธเตหิ ปาเทหิ กถินํ อกฺกมิตพฺพํ, โย อกฺกเมยฺย อาปตฺติ ทุกฺกฏสฺสาติ.}

\prose{เตน โข ปน สมเยน ภิกฺขู อลฺเลหิ ปาเทหิ กถินํ อกฺกมนฺติ, กถินํ ทุสฺสติ.\csromanspacer{} ภควโต เอตมตฺถํ อาโรเจสุํ.\csromanspacer{} น ภิกฺขเว อลฺเลหิ ปาเทหิ กถินํ อกฺกมิตพฺพํ, โย อกฺกเมยฺย, อาปตฺติ ทุกฺกฏสฺสาติ.}

\prose{เตน โข ปน สมเยน ภิกฺขู ส-อุปาหนา กถินํ อกฺกมนฺติ, กถินํ ทุสฺสติ.\csromanspacer{} ภควโต เอตมตฺถํ อาโรเจสุํ.\csromanspacer{} น ภิกฺขเว ส-อุปาหเนน กถินํ อกฺกมิตพฺพํ, โย อกฺกเมยฺย, อาปตฺติ ทุกฺกฏสฺสาติ.}

\prose{เตน โข ปน สมเยน ภิกฺขู จีวรํ สิพฺพนฺตา องฺคุลิยา ปฏิคฺคณฺหนฺติ, องฺคุลิโย ทุกฺขา โหนฺติ.\csromanspacer{} ภควโต เอตมตฺถํ อาโรเจสุํ.\csromanspacer{} อนุชานามิ ภิกฺขเว ปฏิคฺคหนฺติ.}

\proseitem{257}{เตน โข ปน สมเยน ฉพฺพคฺคิยา ภิกฺขู อุจฺจาวเจ ปฏิคฺคเห ธาเรนฺติ โสวณฺณมยํ รูปิยมยํ.\csromanspacer{} มนุสฺสา อุชฺฌายนฺติ ขิยฺยนฺติ วิปาเจนฺติ “เสยฺยถาปิ คิหี กามโภคิโน”ติ.\csromanspacer{} ภควโต เอตมตฺถํ อาโรเจสุํ.\csromanspacer{} น ภิกฺขเว อุจฺจาวจา ปฏิคฺคหา ธาเรตพฺพา, โย ธาเรยฺย, อาปตฺติ ทุกฺกฏสฺส.\csromanspacer{} อนุชามานิ ภิกฺขเว อฏฺฐิมยํ ฯเปฯ สงฺขนาภิมยนฺติ.}

\prose{เตน โข ปน สมเยน สูจิโยปิ สตฺถกาปิ ปฏิคฺคหาปิ นสฺสนฺติ.\csromanspacer{} ภควโต เอตมตฺถํ อาโรเจสุํ.\csromanspacer{} อนุชานามิ ภิกฺขเว อาเวสนวิตฺถกนฺติ.\csromanspacer{} อาเวสนวิตฺถเก สเมกุลา โหนฺติ.\csromanspacer{} ภควโต เอตมตฺถํ อาโรเจสุํ.\csromanspacer{} อนุชานามิ ภิกฺขเว ปฏิคฺคหถวิกนฺติ.\csromanspacer{} อํสพทฺธโก น โหติ.\csromanspacer{} ภควโต เอตมตฺถํ อาโรเจสุํ.\csromanspacer{} อนุชานามิ ภิกฺขเว อํสพทฺธกํ พนฺธนสุตฺตกนฺติ.}

\prose{เตน โข ปน สมเยน ภิกฺขู อพฺโภกาเส จีวรํ สิพฺพนฺตา สีเตนปิ อุเณฺหนปิ กิลมนฺติ.\csromanspacer{} ภควโต เอตมตฺถํ อาโรเจสุํ.\csromanspacer{} อนุชานามิ ภิกฺขเว กถินสาลํ กถินมณฺฑปนฺติ.\csromanspacer{} กถินสาลา \csromanfolio{254}นีจวตฺถุกา โหติ, อุทเกน โอตฺถริยฺยติ.\csromanspacer{} ภควโต เอตมตฺถํ อาโรเจสุํ.\csromanspacer{} อนุชานามิ ภิกฺขเว อุจฺจวตฺถุกํ กาตุนฺติ.\csromanspacer{} จโย ปริปตติ ฯเปฯ.\csromanspacer{} อนุชานามิ ภิกฺขเว จินิตุํ ตโย จเย อิฏฺฐกาจยํ สิลาจยํ ทารุจยนฺติ.\csromanspacer{} อาโรหนฺตา วิหญฺญนฺติ ฯเปฯ.\csromanspacer{} อนุชานามิ ภิกฺขเว ตโย โสปาเน อิฏฺฐกาโสปานํ สิลาโสปานํ ทารุโสปานนฺติ.\csromanspacer{} อาโรหนฺตา ปริปตนฺติ ฯเปฯ.\csromanspacer{} อนุชานามิ ภิกฺขเว อาลมฺพนพาหนฺติ.\csromanspacer{} กถินสาลาย ติณจุณฺณํ ปริปตติ ฯเปฯ.\csromanspacer{} อนุชานามิ ภิกฺขเว โอคุมฺเผตฺวา\footnote{โอคุมฺเพตฺวา (สี, สฺยา)} อุลฺลิตฺตาวลิตฺตํ กาตุํ, เสตวณฺณํ กาฬวณฺณํ เครุกปริกมฺมํ มาลากมฺมํ ลตากมฺมํ มกรทนฺตกํ ปญฺจปฏิกํ จีวรวํสํ จีวรรชฺชุกนฺติ.}

\prose{เตน โข ปน สมเยน ภิกฺขู จีวรํ สิพฺเพตฺวา ตตฺเถว กถินํ อุชฺฌิตฺวา ปกฺกมนฺติ, อุนฺทูเรหิปิ อุปจิกาหิปิ ขชฺชติ.\csromanspacer{} ภควโต เอตมตฺถํ อาโรเจสุํ.\csromanspacer{} อนุชานามิ ภิกฺขเว กถินํ สงฺฆริตุนฺติ.\csromanspacer{} กถินํ ปริภิชฺชติ ฯเปฯ.\csromanspacer{} อนุชานามิ ภิกฺขเว โคฆํสิกาย กถินํ สงฺฆริตุนฺติ.\csromanspacer{} กถินํ วินิเวฐิยติ ฯเปฯ.\csromanspacer{} อนุชานามิ ภิกฺขเว พนฺธนรชฺชุนฺติ.}

\prose{เตน โข ปน สมเยน ภิกฺขู กุฏฺเฏปิ ถมฺเภปิ กถินํ อุสฺสาเปตฺวา ปกฺกมนฺติ, ปริปติตฺวา กถินํ ภิชฺชติ.\csromanspacer{} ภควโต เอตมตฺถํ อาโรเจสุํ.\csromanspacer{} อนุชานามิ ภิกฺขเว ภิตฺติขิเล วา นาคทนฺเต วา ลคฺเคตุนฺติ.}

\proseitem{258}{อถ โข ภควา ราชคเห ยถาภิรนฺตํ วิหริตฺวา เยน เวสาลี เตน จาริกํ ปกฺกามิ.\csromanspacer{} เตน โข ปน สมเยน ภิกฺขู สูจิกมฺปิ สตฺถกมฺปิ เภสชฺชมฺปิ ปตฺเตน อาทาย คจฺฉนฺติ.\csromanspacer{} ภควโต เอตมตฺถํ อาโรเจสุํ.\csromanspacer{} อนุชานามิ ภิกฺขเว เภสชฺชตฺถวิกนฺติ.\csromanspacer{} อํสพทฺธโก น โหติ ฯเปฯ.\csromanspacer{} อนุชานามิ ภิกฺขเว อํสพทฺธกํ พนฺธนสุตฺตกนฺติ.}

\prose{เตน โข ปน สมเยน อญฺญตโร ภิกฺขุ อุปาหนาโย กายพนฺธเนน พนฺธิตฺวา คามํ ปิณฺฑาย ปาวิสิ, อญฺญตโร อุปาสโก ตํ ภิกฺขุํ อภิวาเทนฺโต อุปาหนาโย สีเสน ฆฏฺเฏติ, โส ภิกฺขุ มงฺกุ อโหสิ.\csromanspacer{} อถ โข โส ภิกฺขุ อารามํ คนฺตฺวา ภิกฺขูนํ เอตมตฺถํ อาโรเจสิ.\csromanspacer{} ภิกฺขู ภควโต เอตมตฺถํ อาโรเจสุํ.\csromanspacer{} อนุชานามิ ภิกฺขเว อุปาหนตฺถวิกนฺติ.\csromanspacer{} อํสพทฺธโก น โหติ ฯเปฯ.\csromanspacer{} อนุชานามิ ภิกฺขเว อํสพทฺธกํ พนฺธนสุตฺตกนฺติ.}

\prose{\csromanfolio{255}เตน โข ปน สมเยน อนฺตรามคฺเค อุทกํ อกปฺปิยํ โหติ, ปริสฺสาวนํ น โหติ.\csromanspacer{} ภควโต เอตมตฺถํ อาโรเจสุํ.\csromanspacer{} อนุชานามิ ภิกฺขเว ปริสฺสาวนนฺติ.\csromanspacer{} โจฬกํ นปฺปโหติ ฯเปฯ.\csromanspacer{} อนุชานามิ ภิกฺขเว กฏจฺฉุปริสฺสาวนนฺติ.\csromanspacer{} โจฬกํ นปฺปโหติ ฯเปฯ.\csromanspacer{} อนุชานามิ ภิกฺขเว ธมฺมกรณนฺติ\footnote{ธมฺมกรกํ (สี, สฺยา), ธมกรณํ (ก)}.}

\proseitem{259}{เตน โข ปน สมเยน เทฺว ภิกฺขู โกสเลสุ ชนปเท อทฺธานมคฺคปฺปฏิปนฺนา โหนฺติ, เอโก ภิกฺขุํ อนาจารํ อาจรติ, ทุติโย ภิกฺขุ ตํ ภิกฺขุํ เอตทโวจ “มา อาวุโส เอวรูปํ อกาสิ, เนตํ กปฺปตี”ติ.\csromanspacer{} โส ตสฺมิํ อุปนนฺธิ.\csromanspacer{} อถ โข โส ภิกฺขุ ปิปาสาย ปีฬิโต อุปนทฺธํ ภิกฺขุํ เอตทโวจ “เทหิ เม อาวุโส ปริสฺสาวนํ, ปานียํ ปิวิสฺสามี”ติ.\csromanspacer{} อุปนทฺโธ ภิกฺขุ น อทาสิ, โส ภิกฺขุ ปิปาสาย ปีฬิโต กาลมกาสิ.\csromanspacer{} อถ โข โส ภิกฺขุ อารามํ คนฺตฺวา ภิกฺขูนํ เอตมตฺถํ อาโรเจสิ.\csromanspacer{} กิํ ปน ตฺวํ อาวุโส ปริสฺสาวนํ ยาจิยมาโน น อทาสีติ.\csromanspacer{} เอวมาวุโสติ.\csromanspacer{} เย เต ภิกฺขู อปฺปิจฺฉา ฯเปฯ เต อุชฺฌายนฺติ ขิยฺยนฺติ วิปาเจนฺติ “กถํ หิ นาม ภิกฺขุ ปริสฺสาวนํ ยาจิยมาโน น ทสฺสตี”ติ.\csromanspacer{} อถ โข เต ภิกฺขู ภควโต เอตมตฺถํ อาโรเจสุํ.\csromanspacer{} อถ โข ภควา เอตสฺมิํ นิทาเน เอตสฺมิํ ปกรเณ ภิกฺขุสํฆํ สนฺนิปาตาเปตฺวา ตํ ภิกฺขุํ ปฏิปุจฺฉิ “สจฺจํ กิร ตฺวํ ภิกฺขู ปริสฺสาวนํ ยาจิยมาโน น อทาสี”ติ.\csromanspacer{} สจฺจํ ภควาติ.\csromanspacer{} วิครหิ พุทฺโธ ภควา อนนุจฺฉวิกํ โมฆปุริส อนนุโลมิกํ อปฺปติรูปํ อสฺสามณกํ อกปฺปิยํ อกรณียํ, กถํ หิ นาม ตฺวํ โมฆปุริส ปริสฺสาวนํ ยาจิยมาโน น ทสฺสสิ.\csromanspacer{} เนตํ โมฆปุริส อปฺปสนฺนานํ วา ปสาทาย ฯเปฯ วิครหิตฺวา ฯเปฯ ธมฺมิํ กถํ กตฺวา ภิกฺขู อามนฺเตสิ “น ภิกฺขเว อทฺธานมคฺคปฺปฏิปนฺเนน ภิกฺขุนา ปริสฺสาวนํ ยาจิยมาเนน น ทาตพฺพํ, โย น ทเทยฺย, อาปตฺติ ทุกฺกฏสฺส.\csromanspacer{} น จ ภิกฺขเว อปริสฺสาวนเกน อทฺธาโน ปฏิปชฺชิตพฺโพ, โย ปฏิปชฺเชยฺย, อาปตฺติ ทุกฺกฏสฺส.\csromanspacer{} สเจ น โหติ ปริสฺสาวนํ วา ธมฺมกรโณ วา, สํฆาฏิกณฺโณปิ อธิฏฺฐาตพฺโพ อิมินา ปริสฺสาเวตฺวา ปิวิสฺสามี”ติ.\csromanspacer{} อถ โข ภควา อนุปุพฺเพน จาริกํ จรมาโน เยน เวสาลี ตทวสริ, ตตฺร สุทํ ภควา เวสาลิยํ วิหรติ มหาวเน \csromanfolio{256}กูฏาคารสาลายํ.\csromanspacer{} เตน โข ปน สมเยน ภิกฺขู นวกมฺมํ กโรนฺติ, ปริสฺสาวนํ น สมฺมติ.\csromanspacer{} ภควโต เอตมตฺถํ อาโรเจสุํ.\csromanspacer{} อนุชานามิ ภิกฺขเว ทณฺฑปริสฺสาวนนฺติ.\csromanspacer{} ทณฺฑปริสฺสาวนํ น สมฺมติ.\csromanspacer{} ภควโต เอตมตฺถํ อาโรเจสุํ.\csromanspacer{} อนุชานามิ ภิกฺขเว โอตฺถรกนฺติ. เตน โข ปน สมเยน ภิกฺขู มกเสหิ อุพฺพาฬฺหา โหนฺติ.\csromanspacer{} ภควโต เอตมตฺถํ อาโรเจสุํ.\csromanspacer{} อนุชานามิ ภิกฺขเว มกสกุฏิกนฺติ.}

\proseitem{260}{เตน โข ปน สมเยน เวสาลิยํ ปณีตานํ ภตฺตานํ สตฺตปฏิปาฏิ อฏฺฐิตา โหติ.\csromanspacer{} ภิกฺขู ปณีตานิ โภชนานิ ภุญฺชิตฺวา อภิสนฺนกายา โหนฺติ พหฺวาพาธา\footnote{พวฺหาพาธา (สี)}.\csromanspacer{} อถ โข ชีวโก โกมารภจฺโจ เวสาลิํ อคมาสิ เกนจิเทว กรณีเยน.\csromanspacer{} อทฺทสา โข ชีวโก โกมารภจฺโจ ภิกฺขู อภิสนฺนกาเย พหฺวาพาเธ, ทิสฺวาน เยน ภควา เตนุปสงฺกมิ, อุปสงฺกมิตฺวา ภควนฺตํ อภิวาเทตฺวา เอกมนฺตํ นิสีทิ, เอกมนฺตํ นิสินฺโน โข ชีวโก โกมารภจฺโจ ภควนฺตํ เอตทโวจ “เอตรหิ ภนฺเต ภิกฺขู อภิสนฺนกายา พหฺวาพาธา, สาธุ ภนฺเต ภควา ภิกฺขูนํ จงฺกมญฺจ ชนฺตาฆรญฺจ อนุชานาตุ, เอวํ ภิกฺขู อปฺปาพาธา ภวิสฺสนฺตี”ติ.\csromanspacer{} อถ โข ภควา ชีวกํ โกมารภจฺจํ ธมฺมิยา กถาย สนฺทสฺเสสิ สมาทเปสิ สมุตฺเตเชสิ สมฺปหํเสสิ.\csromanspacer{} อถ โข ชีวโก โกมารภจฺโจ ภควตา ธมฺมิกา กถาย สนฺทสฺสิโต สมาทปิโต สมุตฺเตชิโต สมฺปหํสิโต อุฏฺฐายาสนา ภควนฺตํ อภิวาเทตฺวา ปทกฺขิณํ กตฺวา ปกฺกามิ.\csromanspacer{} อถ โข ภควา เอตสฺมิํ นิทาเน เอตสฺมิํ ปกรเณ ธมฺมิํ กถํ กตฺวา ภิกฺขู อามนฺเตสิ “อนุชานามิ ภิกฺขเว จงฺกมญฺจ ชนฺตาฆรญฺจา”ติ.}

\prose{เตน โข ปน สมเยน ภิกฺขู วิสเม จงฺกเม จงฺกมนฺติ, ปาทา ทุกฺขา โหนฺติ.\csromanspacer{} ภควโต เอตมตฺถํ อาโรเจสุํ.\csromanspacer{} อนุชานามิ ภิกฺขเว สมํ กาตุนฺติ.\csromanspacer{} จงฺกโม นีจวตฺถุโก โหติ, อุทเกน โอตฺถริยฺยติ ฯเปฯ.\csromanspacer{} อนุชานามิ ภิกฺขเว อุจฺจวตฺถุตํ กาตุนฺติ.\csromanspacer{} จโย ปริปตติ ฯเปฯ.\csromanspacer{} อนุชานามิ ภิกฺขเว จินิตุํ ตโย จเย อิฏฺฐกาจยํ สิลาจยํ ทารุจยนฺติ.\csromanspacer{} อาโรหนฺตา วิหญฺญนฺติ ฯเปฯ.\csromanspacer{} อนุชานามิ ภิกฺขเว ตโย โสปาเน อิฏฺฐกาโสปานํ สิลาโสปานํ ทารุโสปานนฺติ.\csromanspacer{} อาโร หนฺตา ปริปตนฺติ ฯเปฯ.\csromanspacer{} อนุชานามิ ภิกฺขเว อาลมฺพนพาหนฺติ.}

\prose{\csromanfolio{257}เตน โข ปน สมเยน ภิกฺขู จงฺกเม จงฺกมนฺตา ปริปตนฺติ.\csromanspacer{} ภควโต เอตมตฺถํ อาโรเจสุํ.\csromanspacer{} อนุชานามิ ภิกฺขเว จงฺกมนเวทิกนฺติ.}

\prose{เตน โข ปน สมเยน ภิกฺขู อชฺโฌกาเส จงฺกมนฺตา สีเตนปิ อุเณฺหนปิ กิลมนฺติ.\csromanspacer{} ภควโต เอตมตฺถํ อาโรเจสุํ.\csromanspacer{} อนุชานามิ ภิกฺขเว จงฺกมนสาลนฺติ.\csromanspacer{} จงฺกมนสาลายํ ติณจุณฺณํ ปริปตติ ฯเปฯ.\csromanspacer{} อนุชานามิ ภิกฺขเว โอคุมฺเผตฺวา อุลฺลิตฺตาวลิตฺตํ กาตุํ เสตวณฺณํ กาฬวณฺณํ เครุกปริกมฺมํ มาลากมฺมํ ลตากมฺมํ มกรทนฺตกํ ปญฺจปฏิกํ จีวรวํสํ จีวรรชฺชุนฺติ.\csromanspacer{} ชนฺตาฆรํ นีจวตฺถุกํ โหติ, อุทเกน โอตฺถริยฺยติ ฯเปฯ.\csromanspacer{} อนุชานามิ ภิกฺขเว อุจฺจวตฺถุกํ กาตุนฺติ.\csromanspacer{} จโย ปริปตติ ฯเปฯ.\csromanspacer{} อนุชานามิ ภิกฺขเว จินิตุํ ตโย จเย อิฏฺฐกาจยํ สิลาจยํ ทารุจยนฺติ.\csromanspacer{} อาโรหนฺตา วิหญฺญนฺติ ฯเปฯ.\csromanspacer{} อนุชานามิ ภิกฺขเว ตโย โสปาเน อิฏฺฐกาโสปานํ สิลาโสปานํ ทารุโสปานนฺติ.\csromanspacer{} อาโรหนฺตา ปริปตนฺติ ฯเปฯ.\csromanspacer{} อนุชานามิ ภิกฺขเว อาลมฺพนพาหนฺติ.\csromanspacer{} ชนฺตาฆรสฺส กวาฏํ น โหติ ฯเปฯ.\csromanspacer{} อนุชานามิ ภิกฺขเว กวาฏํ ปิฏฺฐสํฆาฏํ\footnote{ปีฐสํฆาฏํ (ก)} อุทุกฺขลิกํ อุตฺตรปาสกํ อคฺคฬวฏฺฏิํ กปิสีสกํ สูจิกํ ฆฏิกํ ตาฬจฺฉิทฺทํ อาวิญฺฉนฉิทฺทํ อาวิญฺฉนรชฺชุนฺติ.\csromanspacer{} ชนฺตาฆรสฺส กุฏฺฏปาโท ชีรติ ฯเปฯ.\csromanspacer{} อนุชานามิ ภิกฺขเว มณฺฑลิกํ กาตุนฺติ.\csromanspacer{} ชนฺตาฆรสฺส ธูมเนตฺตํ น โหติ ฯเปฯ.\csromanspacer{} อนุชานามิ ภิกฺขเว ธูมเนตฺตนฺติ.}

\prose{เตน โข ปน สมเยน ภิกฺขู ขุทฺทเก ชนฺตาฆเร มชฺเฌ อคฺคิฏฺฐานํ กโรนฺติ, อุปจาโร น โหติ.\csromanspacer{} ภควโต เอตมตฺถํ อาโรเจสุํ.\csromanspacer{} อนุชานามิ ภิกฺขเว ขุทฺทเก ชนฺตาฆเร เอกมนฺตํ อคฺคิฏฺฐานํ กาตุํ มหลฺลเก มชฺเฌติ.\csromanspacer{} ชนฺตาฆเร อคฺคิ มุขํ ฑหติ ฯเปฯ.\csromanspacer{} อนุชานามิ ภิกฺขเว มุขมตฺติกนฺติ.\csromanspacer{} หตฺเถ มตฺติกํ เตเมนฺติ ฯเปฯ.\csromanspacer{} อนุชานามิ ภิกฺขเว มตฺติกาโทณิกนฺติ.\csromanspacer{} มตฺติกา ทุคฺคนฺธา โหติ ฯเปฯ.\csromanspacer{} อนุชานามิ ภิกฺขเว วาเสตุนฺติ.\csromanspacer{} ชนฺตาฆเร อคฺคิ กายํ ฑหติ ฯเปฯ.\csromanspacer{} อนุชานามิ ภิกฺขเว อุทกํ อติหริตุนฺติ.\csromanspacer{} ปาติยาปิ ปตฺเตนปิ อุทกํ อติหรนฺติ ฯเปฯ.\csromanspacer{} อนุชานามิ ภิกฺขเว อุทกฏฺฐานํ อุทกสราวกนฺติ.\csromanspacer{} ชนฺตาฆรํ ติณจฺฉทนํ น เสเทติ\footnote{ติณจฺฉาทเนน ฉาเทติ (ก)} ฯเปฯ.\csromanspacer{} อนุชานามิ ภิกฺขเว โอคุมฺเผตฺวา อุลฺลิตฺตาวลิตฺตํ กาตุนฺติ. \csromanfolio{258}ชนฺตาฆรํ จิกฺขลฺลํ โหติ ฯเปฯ.\csromanspacer{} อนุชานามิ ภิกฺขเว สนฺถริตุํ ตโย สนฺถเร อิฏฺฐกาสนฺถรํ สิลาสนฺถรํ ทารุสนฺถรนฺติ.\csromanspacer{} จิกฺขลฺลํเยว โหติ ฯเปฯ.\csromanspacer{} อนุชานามิ ภิกฺขเว โธวิตุนฺติ.\csromanspacer{} อุทกํ สนฺติฏฺฐติ ฯเปฯ.\csromanspacer{} อนุชานามิ ภิกฺขเว อุทกนิทฺธมนนฺติ.}

\prose{เตน โข ปน สมเยน ภิกฺขู ชนฺตาฆเร ฉมาย นิสีทนฺติ, คตฺตานิ กณฺฑูวนฺติ.\csromanspacer{} ภควโต เอตมตฺถํ อาโรเจสุํ.\csromanspacer{} อนุชานามิ ภิกฺขเว ชนฺตาฆรปีฐนฺติ.\csromanspacer{} เตน โข ปน สมเยน ชนฺตาฆรํ อปริกฺขิตฺตํ โหติ ฯเปฯ.\csromanspacer{} อนุชานามิ ภิกฺขเว ปริกฺขิปิตุํ ตโย ปากาเร อิฏฺฐกาปาการํ สิลาปาการํ ทารุปาการนฺติ.\csromanspacer{} โกฏฺฐโก น โหติ ฯเปฯ.\csromanspacer{} อนุชานามิ ภิกฺขเว โกฏฺฐกนฺติ.\csromanspacer{} โกฏฺฐโก นีจวตฺถุโก โหติ, อุทเกน โอตฺถริยฺยติ ฯเปฯ.\csromanspacer{} อนุชานามิ ภิกฺขเว อุจฺจวตฺถุกํ กาตุนฺติ.\csromanspacer{} จโย ปริปตติ ฯเปฯ.\csromanspacer{} อนุชานามิ ภิกฺขเว จินิตุํ ตโย จเย อิฏฺฐกาจยํ สิลาจยํ ทารุจยนฺติ.\csromanspacer{} อาโรหนฺตา วิหญฺญนฺติ ฯเปฯ.\csromanspacer{} อนุชานามิ ภิกฺขเว ตโย โสปาเน อิฏฺฐกาโสปานํ สิลาโสปานํ ทารุโสปานนฺติ.\csromanspacer{} อาโรหนฺตา ปริปตนฺติ ฯเปฯ.\csromanspacer{} อนุชานามิ ภิกฺขเว อาลมฺพนพาหนฺติ.\csromanspacer{} โกฏฺฐกสฺส กวาฏํ น โหติ ฯเปฯ.\csromanspacer{} อนุชานามิ ภิกฺขเว กวาฏํ ปิฏฺฐสํฆาฏํ อุทุกฺขลิกํ อุตฺตรปาสกํ อคฺคฬวฏฺฏิํ กปิสีสกํ สูจิกํ ฆฏิกํ ตาฬจฺฉิทฺทํ อาวิญฺฉนฉทฺทํ อาวิญฺฉนรชฺชุนฺติ.\csromanspacer{} โกฏฺฐเก ติณจุณฺณํ ปริปตติ ฯเปฯ.\csromanspacer{} อนุชานามิ ภิกฺขเว โอคุมฺเผตฺวา อุลฺลิตฺตาวลิตฺตํ กาตุํ กาฬวณฺณํ เครุกปริกมฺมํ มาลากมฺมํ ลเตกมฺมํ มกรทนฺตกํ ปญฺจปฏิกนฺติ.\csromanspacer{} ปริเวณํ จิกฺขลฺลํ โหติ ฯเปฯ.\csromanspacer{} อนุชานามิ ภิกฺขเว มรุมฺพํ อุปกิริตุนฺติ.\csromanspacer{} น ปริยาปุณนฺติ ฯเปฯ.\csromanspacer{} อนุชานามิ ภิกฺขเว ปทรสิลํ นิกฺขิปิตุนฺติ.\csromanspacer{} อุทกํ สนฺติฏฺฐติ ฯเปฯ.\csromanspacer{} อนุชานามิ ภิกฺขเว อุทกนิทฺธมนนฺติ.}

\proseitem{261}{เตน โข ปน สมเยน ภิกฺขู นคฺคา นคฺคํ อภิวาเทนฺติ ฯเปฯ.\csromanspacer{} นคฺคา นคฺคํ อภิวาทาเปนฺติ.\csromanspacer{} นคฺคา นคฺคสฺส ปริกมฺมํ กโรนฺติ.\csromanspacer{} นคฺคา นคฺคสฺส ปริกมฺมํ การาเปนฺติ.\csromanspacer{} นคฺคา นคฺคสฺส เทนฺติ.\csromanspacer{} นคฺคา ปฏิคฺคณฺหนฺติ.\csromanspacer{} นคฺคา ขาทนฺติ.\csromanspacer{} นคฺคา ภุญฺชนฺติ.\csromanspacer{} นคฺคา สายนฺติ.\csromanspacer{} นคฺคา ปิวนฺติ.\csromanspacer{} ภควโต เอตมตฺตํ อาโรเจสุํ.\csromanspacer{} น ภิกฺขเว นคฺเคน\footnote{อิทํ ปทํ กณฺฑจิ นตฺถิ.} นคฺโค อภิวาเทตพฺโพ ฯเปฯ.\csromanspacer{} น นคฺเคน อภิวาเทตพฺพํ.\csromanspacer{} น นคฺเคน1 นคฺโค อภิวาทาเปตพฺโพ.\csromanspacer{} น นคฺเคน อภิวาทาเปตพฺพํ.\csromanspacer{} น นคฺเคน นคฺคสฺส \csromanfolio{259}ปริกมฺมํ กาตพฺพํ.\csromanspacer{} น นคฺเคน นคฺคสฺส ปริกมฺมํ การาเปตพฺพํ.\csromanspacer{} น นคฺเคน นคฺคสฺส ทาตพฺพํ.\csromanspacer{} น นคฺเคน ปฏิคฺคเหตพฺพํ.\csromanspacer{} น นคฺเคน ขาทิตพฺพํ.\csromanspacer{} น นคฺเคน ภุญฺชิตพฺพํ.\csromanspacer{} น นคฺเคน สายิตพฺพํ.\csromanspacer{} น นคฺเคน ปาตพฺพํ, โย ปิเวยฺย, อาปตฺติ ทุกฺกฏสฺสาติ.}

\prose{เตน โข ปน สมเยน ภิกฺขู ชนฺตาฆเร ฉมาย จีวรํ นิกฺขิปนฺติ, จีวรํ ปํสุกิตํ โหติ.\csromanspacer{} ภควโต เอตมตฺถํ อาโรเจสุํ.\csromanspacer{} อนุชานามิ ภิกฺขเว จีวรวํสํ จีวรรชฺชุนฺติ.\csromanspacer{} เทเว วสฺสนฺเต จีวรํ โอวสฺสติ ฯเปฯ.\csromanspacer{} อนุชานามิ ภิกฺขเว ชนฺตาฆรสาลนฺติ.\csromanspacer{} ชนฺตาฆรสาลา นีจวตฺถุกา โหติ, อุทเกน โอตฺถริยฺยติ ฯเปฯ.\csromanspacer{} อนุชานามิ ภิกฺขเว อุจฺจวตฺถุกํ กาตุนฺติ.\csromanspacer{} จโย ปริปตติ ฯเปฯ.\csromanspacer{} อนุชานามิ ภิกฺขเว จินิตุํ ฯเปฯ อาโรหนฺตา วิหญฺญนฺติ ฯเปฯ อาโรหนฺตา ปริปตนฺติ ฯเปฯ.\csromanspacer{} อนุชานามิ ภิกฺขเว อาลมฺพนพาหนฺติ.\csromanspacer{} ชนฺตาฆรสาลาย ติณจุณฺณํ ปริปตติ ฯเปฯ.\csromanspacer{} อนุชานามิ ภิกฺขเว โอคุมฺเผตฺวา อุลฺลิตฺตาวลิตฺตํ กาตุํ ฯเปฯ จีวรวํสํ จีวรรชฺชุนฺติ.}

\prose{เตน โข ปน สมเยน ภิกฺขู ชนฺตาฆเรปิ อุทเกปิ ปริกมฺมํ กาตุํ กุกฺกุจฺจายนฺติ.\csromanspacer{} ภควโต เอตมตฺถํ อาโรเจสุํ.\csromanspacer{} อนุชานามิ ภิกฺขเว ติสฺโส ปฏิจฺฉาทิโย ชนฺตาฆรปฏิจฺฉาทิํ อุทกปฏิจฺฉาทิํ วตฺถปฏิจฺฉาทินฺติ.}

\prose{เตน โข ปน สมเยน ชนฺตาฆเร อุทกํ น โหติ.\csromanspacer{} ภควโต เอตมตฺถํ อาโรเจสุํ.\csromanspacer{} อนุชานามิ ภิกฺขเว อุทปานนฺติ.\csromanspacer{} อุทปานสฺส กูลํ ลุชฺชติ ฯเปฯ.\csromanspacer{} อนุชานามิ ภิกฺขเว จินิตุํ ตโย จเย อิฏฺฐกาจยํ สิลาจยํ ทารุจยนฺติ.\csromanspacer{} อุทปาโน นีจวตฺถุโก โหติ, อุทเกน โอตฺถริยฺยติ ฯเปฯ.\csromanspacer{} อนุชานามิ ภิกฺขเว อุจฺจวตฺถุกํ กาตุนฺติ.\csromanspacer{} จโย ปริปตติ ฯเปฯ อาโรหนฺตา วิหญฺญนฺติ ฯเปฯ อาโรหนฺตา ปริปตนฺติ ฯเปฯ.\csromanspacer{} อนุชานามิ ภิกฺขเว อาลมฺพนพาหนฺติ.}

\proseitem{262}{เตน โข ปน สมเยน ภิกฺขู วลฺลิกายปิ กายพนฺธเนนปิ อุทกํ วาเหนฺติ\footnote{วาหนฺติ (ก)}.\csromanspacer{} ภควโต เอตมตฺถํ อาโรเจสุํ.\csromanspacer{} อนุชานามิ ภิกฺขเว อุทกวาหนรชฺชุนฺติ.\csromanspacer{} หตฺถา ทุกฺขา โหนฺติ ฯเปฯ.\csromanspacer{} อนุชานามิ ภิกฺขเว ตุลํ กรกฏกํ จกฺกวฏฺฏกนฺติ.\csromanspacer{} ภาชนา พหู ภิชฺชนฺติ ฯเปฯ.\csromanspacer{} อนุชานามิ ภิกฺขเว ตโย วารเก โลหวารกํ ทารุวารกํ จมฺมกฺขณฺฑนฺติ.}

\prose{\csromanfolio{260}เตน โข ปน สมเยน ภิกฺขู อชฺโฌกาเส อุทกํ วาเหนฺตา สีเตนปิ อุเณฺหนปิ กิลมนฺติ.\csromanspacer{} ภควโต เอตมตฺถํ อาโรเจสุํ.\csromanspacer{} อนุชานามิ ภิกฺขเว อุทปานสาลนฺติ.\csromanspacer{} อุทปานสาลาย ติณจุณฺณํ ปริปตติ ฯเปฯ.\csromanspacer{} อนุชานามิ ภิกฺขเว โอคุมฺเผตฺวา อุลฺลิตฺตาวลิตฺตํ กาตุํ เสตวณฺณํ กาฬวณฺณํ เครุกปริกมฺมํ มาลากมฺมํ ลตากมฺมํ มกรทนฺตกํ ปญฺจปฏิกํ จีวรวํสํ จีวรรชฺชุนฺติ.\csromanspacer{} อุทปาโน อปารุโต โหติ, ติณจุณฺเณหิปิ ปํสุเกหิปิ โอกิริยฺยติ ฯเปฯ.\csromanspacer{} อนุชานามิ ภิกฺขเว อปิธานนฺติ.\csromanspacer{} อุทกภาชนํ น สํวิชฺชติ ฯเปฯ.\csromanspacer{} อนุชานามิ ภิกฺขเว อุทกโทณิํ อุทกกฏาหนฺติ.}

\prose{เตน โข ปน สมเยน ภิกฺขู อาราเม ตหํ ตหํ นหายนฺติ, อาราโม จิกฺขลฺโล โหติ.\csromanspacer{} ภควโต เอตมตฺถํ อาโรเจสุํ.\csromanspacer{} อนุชานามิ ภิกฺขเว จนฺทนิกนฺติ.\csromanspacer{} จนฺทนิกา ปากฏา โหติ.\csromanspacer{} ภิกฺขู หิริยนฺติ นหายิตุํ ฯเปฯ.\csromanspacer{} อนุชานามิ ภิกฺขเว ปริกฺขิปิตุํ ตโย ปากาเร อิฏฺฐกาปาการํ สิลาปาการํ ทารุปาการนฺติ.\csromanspacer{} จนฺทนิกา จิกฺขลฺลา โหติ ฯเปฯ.\csromanspacer{} อนุชานามิ ภิกฺขเว สนฺตริตุํ ตโย สนฺถเร อิฏฺฐกาสนฺถรํ สิลาสนฺถรํ ทารุสนฺถรนฺติ.\csromanspacer{} อุทกํ สนฺติฏฺฐติ ฯเปฯ.\csromanspacer{} อนุชานามิ ภิกฺขเว อุกกนิทฺธมนนฺติ.}

\prose{เตน โข ปน สมเยน ภิกฺขูนํ คตฺตานิ สีติคตานิ โหนฺติ, ภควโต เอตมตฺถํ อาโรเจสุํ ฯเปฯ.\csromanspacer{} อนุชานามิ ภิกฺขเว อุทกปุญฺฉนิํ โจฬเกนปิ ปจฺจุทฺธริตุนฺติ.}

\proseitem{263}{เตน โข ปน สมเยน อญฺญตโร อุปาสโก สํฆสฺส อตฺถาย โปกฺขรณิํ กาเรตุกาโม โหติ ภควโต เอกมตฺถํ อาโรเจสุํ ฯเปฯ.\csromanspacer{} อนุชานามิ ภิกฺขเว โปกฺขรณินฺติ.\csromanspacer{} โปกฺขรณิยา กูลํ ลุชฺชติ ฯเปฯ.\csromanspacer{} อนุชานามิ ภิกฺขเว จินิตุํ ตโย จเย อิฏฺฐกาจยํ สิลาจยํ ทารุจยนฺติ.\csromanspacer{} อาโรหนฺตา วิหญฺญนฺติ ฯเปฯ.\csromanspacer{} อนุชานามิ ภิกฺขเว ตโย โสปาเน อิฏฺฐกาโสปานํ สิลาโสปานํ ทารุโสปานนฺติ.\csromanspacer{} อาโรหนฺตา ปริปตนฺติ ฯเปฯ.\csromanspacer{} อนุชานามิ ภิกฺขเว อาลมฺพนพาหนฺติ.\csromanspacer{} โปกฺขรณิยา อุทกํ ปุราณํ โหติ ฯเปฯ.\csromanspacer{} อนุชานามิ ภิกฺขเว อุทกมาติกํ อุทกนิทฺธมนนฺติ.}

\prose{\csromanfolio{261}เตน โข ปน สมเยน อญฺญตโร ภิกฺขุ สํฆสฺส อตฺถาย นิลฺเลขํ ชนฺตาฆรํ กตฺตุกาโม โหติ.\csromanspacer{} ภควโต เอตมตฺถํ อาโรเจสุํ ฯเปฯ.\csromanspacer{} อนุชานามิ ภิกฺขเว นิลฺเลขํ ชนฺตาฆรนฺติ.}

\prose{เตน โข ปน สมเยน ฉพฺพคฺคิยา ภิกฺขู จาตุมาสํ นิสีทเนน วิปฺปวสนฺติ.\csromanspacer{} ภควโต เอตมตฺถํ อาโรเจสุํ ฯเปฯ.\csromanspacer{} น ภิกฺขเว จาตุมาสํ นิสีทเนน วิปฺปวสิตพฺพํ.\csromanspacer{} โย วิปฺปวเสยฺย, อาปตฺติ ทุกฺกฏสฺสาติ.}

\proseitem{264}{เตน โข ปน สมเยน ฉพฺพคฺคิยา ภิกฺขู ปุปฺผาภิกิณฺเณสุ สยเนสุ สยนฺติ, มนุสฺสา วิหารจาริกํ อาหิณฺฑนฺตา ปสฺสิตฺวา อุชฺฌายนฺติ ขิยฺยนฺติ วิปาเจนฺติ “เสยฺยถาปิ คิหี กามโภคิโน”ติ.\csromanspacer{} ภควโต เอตมตฺถํ อาโรเจสุํ.\csromanspacer{} น ภิกฺขเว ปุปฺผาภิกิณฺเณสุ สยเนสุ สยิตพฺพํ, โย สเยยฺย, อาปตฺติ ทุกฺกฏสฺสาติ.}

\prose{เตน โข ปน สมเยน มนุสฺสา คนฺธมฺปิ มาลมฺปิ อาทาย อารามํ อาคจฺฉนฺติ, ภิกฺขู กุกฺกุจฺจายนฺตา น ปฏิคฺคณฺหนฺติ.\csromanspacer{} ภควโต เอตมตฺถํ อาโรเจสุํ.\csromanspacer{} อนุชานามิ ภิกฺขเว คนฺธํ คเหตฺวา กวาเฏ ปญฺจงฺคุลิกํ ทาตุํ, ปุปฺผํ คเหตฺวา วิหาเร เอกมนฺตํ นิกฺขิปิตุนฺติ.}

\prose{เตน โข ปน สมเยน สํฆสฺส นมตกํ อุปฺปนฺนํ โหติ.\csromanspacer{} ภควโต เอตมตฺถํ อาโรเจสุํ.\csromanspacer{} อนุชานามิ ภิกฺขเว นมตกนฺติ.\csromanspacer{} อถ โข ภิกฺขูนํ เอตทโหสิ “นมตกํ อธิฏฺฐาตพฺพํ นุ โข อุทาหุ วิกปฺเปตพฺพนฺ”ติ ฯเปฯ.\csromanspacer{} น ภิกฺขเว นมตกํ อธิฏฺฐาตพฺพํ, น วิกปฺเปตพฺพนฺติ.}

\prose{เตน โข ปน สมเยน ฉพฺพคฺคิยา ภิกฺขู อาสิตฺตกูปธาเน ภุญฺชนฺติ.\csromanspacer{} มนุสฺสา อุชฺฌายนฺติ ขิยฺยนฺติ วิปาเจนฺติ “เสยฺยถาปิ คิหี กามโภคิโน”ติ.\csromanspacer{} ภควโต เอตมตฺถํ อาโรเจสุํ ฯเปฯ.\csromanspacer{} น ภิกฺขเว อาสิตฺตกูปธาเน ภุญฺชิตพฺพํ, โย ภุญฺเชยฺย, อาปตฺติ ทุกฺกฏสฺสาติ.}

\prose{เตน โข ปน สมเยน อญฺญตโร ภิกฺขุ คิลาโน โหติ, โส ภุญฺชมาโน น สกฺโกติ หตฺเถน ปตฺตํ สนฺธาเรตุํ.\csromanspacer{} ภควโต เอตมตฺถํ อาโรเจสุํ.\csromanspacer{} อนุชานามิ ภิกฺขเว มโฬริกนฺติ.}

\prose{เตน โข ปน สมเยน ฉพฺพคฺคิยา ภิกฺขู เอกภาชเนปิ ภุญฺชนฺติ ฯเปฯ เอกถาลเกปิ ปิวนฺติ.\csromanspacer{} เอกมญฺเจปิ ตุวฏฺเฏนฺติ.\csromanspacer{} เอกตฺถรณาปิ ตุวฏฺเฏนฺติ.\csromanspacer{} เอกปาวุรณาปิ ตุวฏฺเฏนฺติ.\csromanspacer{} เอกตฺถรณปาวุรณาปิ ตุวฏฺเฏนฺติ.\csromanspacer{} มนุสฺสา \csromanfolio{262}อุชฺฌายนฺติ ขิยฺยนฺติ วิปาเจนฺติ “เสยฺยถาปิ คิหี กามโภคิโน”ติ.\csromanspacer{} ภควโต เอตมตฺถํ อาโรเจสุํ.\csromanspacer{} น ภิกฺขเว เอกภาชเน ภุญฺชิตพฺพํ ฯเปฯ น เอกถาลเก ปาตพฺพํ.\csromanspacer{} น เอกมญฺเจ ตุวฏฺฏิตพฺพํ.\csromanspacer{} น เอกตฺถรณา\footnote{น เอกตฺถรเณ.} ตุวฏฺฏิตพฺพํ.\csromanspacer{} น เอกปาวุรณา\footnote{น เอกปาวุรเณ.} ตุวฏฺฏิตพฺพํ.\csromanspacer{} น เอกตฺถรณปาวุรณา\footnote{น เอกตฺถรณปาวุรเณ (สฺยา)} ตุวฏฺฏิตพฺพํ, โย ตุวฏฺเฏยฺย, อาปตฺติ ทุกฺกฏสฺสาติ.}

\proseitem{265}{เตน โข ปน สมเยน วฑฺโฒ ลิจฺฉวี เมตฺติยภูมชกานํ ภิกฺขูนํ สหาโย โหติ.\csromanspacer{} อถ โข วฑฺโฒ ลิจฺฉวี เยน เมตฺติยภูมชกา ภิกฺขู เตนุปสงฺกมิ, อุปสงฺกมิตฺวา เมตฺติยภูมชเก ภิกฺขู เอตทโวจ “วนฺทามิ อยฺยา”ติ.\csromanspacer{} เอวํ วุตฺเต เมตฺติยภูมชกา ภิกฺขู นาลปิํสุ.\csromanspacer{} ทุติยมฺปิ โข วฑฺโฒ ลิจฺฉวี เมตฺติยภูมชเก ภิกฺขู เอตทโวจ “วนฺทามิ อยฺยา”ติ.\csromanspacer{} ทุติยมฺปิ โข เมตฺติยภูมชกา ภิกฺขู นาลปิํสุ.\csromanspacer{} ตติยมฺปิ โข วฑฺโฒ ลิจฺฉวี เมตฺติยภูมชเก ภิกฺขู เอตทโวจ “วนฺทามิ อยฺยา”ติ.\csromanspacer{} ตติยมฺปิ โข เมตฺติยภูมชกา ภิกฺขู นาลปิํสุ.\csromanspacer{} กฺยาหํ อยฺยานํ อปรชฺฌามิ, กิสฺส มํ อยฺยา นาลปนฺตีติ.\csromanspacer{} ตถา หิ ปน ตฺวํ อาวุโส วฑฺฒ อเมฺห ทพฺเพน มลฺลปุตฺเตน วิเหฐิยมาเน อชฺฌุเปกฺขสีติ.\csromanspacer{} กฺยาหํ อยฺยา กโรมีติ.\csromanspacer{} สเจ โข ตฺวํ อาวุโส วฑฺฒ อิจฺเฉยฺยาสิ, อชฺเชว ภควา อายสฺมนฺตํ ทพฺพํ มลฺลปุตฺตํ นาสาเปยฺยาติ.\csromanspacer{} กฺยาหํ อยฺยา กโรมิ, กิํ มยา สกฺกา กาตุนฺติ.\csromanspacer{} เอหิ ตฺวํ อาวุโส วฑฺฒ เยน ภควา เตนุปสงฺกม, อุปสงฺกมิตฺวา ภควนฺตํ เอวํ วเทหิ “อิทํ ภนฺเต นจฺฉนฺนํ นปฺปติรูปํ, ยายํ ภนฺเต ทิสา อภยา อนีติกา อนุปทฺทวา, สายํ ทิสา สภยา ส-อีติกา ส-อุปทฺทวา, ยโต นิวาตํ, ตโต สวาตํ, อุทกํ มญฺเญ อาทิตฺตํ, อยฺเยน เม ทพฺเพน มลฺลปุตฺเตน ปชาปติ ทูสิตา”ติ. “เอวํ อยฺยา”ติ โข วฑฺโฒ ลิจฺฉวี เมตฺติยภูมชกานํ ภิกฺขูนํ ปฏิสฺสุตฺวา เยน ภควา เตนุปสงฺกมิ, อุปสงฺกมิตฺวา ภควนฺตํ อภิวาเทตฺวา เอกมนฺตํ นิสีทิ, เอกมนฺตํ นิสินฺโน โข วฑฺโฒ ลิจฺฉวี ภควนฺตํ เอตทโวจ “อิทํ ภนฺเต นจฺฉนฺนํ นปฺปติรูปํ, ยายํ ภนฺเต ทิสา อภยา อนีติกา อนุปทฺทวา, สายํ ทิสา สภยา ส-อีติกา ส-อุปทฺทวา, ยโต นิวาตํ, ตโต สวาตํ, อุทกํ มญฺเญ อาทิตฺตํ, อยฺเยน เม ทพฺเพน มลฺลปุตฺเตน ปชาปติ ทูสิตา”ติ.}

\prose{\csromanfolio{263}อถ โข ภควา เอตสฺมิํ นิทาเน เอตสฺมิํ ปกรเณ ภิกฺขุสํฆํ สนฺนิปาตาเปตฺวา อายสฺมนฺตํ ทพฺพํ มลฺลปุตฺตํ ปฏิปุจฺฉิ “สรสิ ตฺวํ ทพฺพ เอวรูปํ กตฺตา ยถายํ วฑฺโฒ อาหา”ติ ฯเปฯ.\csromanspacer{} ยถา มํ ภนฺเต ภควา ชานาตีติ.\csromanspacer{} ทุติยมฺปิ โข ภควา ฯเปฯ.\csromanspacer{} ตติยมฺปิ โข ภควา อายสฺมนฺตํ ทพฺพํ มลฺลปุตฺตํ เอตทโวจ “สรสิ ตฺวํ ทพฺพ เอวรูปํ กตฺตา ยถายํ วฑฺโฒ อาหา”ติ.\csromanspacer{} ยถา มํ ภนฺเต ภควา ชานาตี’ติ.\csromanspacer{} น โข ทพฺพ ทพฺพา เอวํ นิพฺเพเฐนฺติ, สเจ ตยา กตํ กตนฺติ วเทหิ, สเจ อกตํ อกตนฺติ วเทหีติ.\csromanspacer{} ยโต อหํ ภนฺเต ชาโต นาภิชานามิ สุปินนฺเตนปิ เมถุนํ ธมฺมํ ปฏิเสวิตา, ปเคว ชาคโรติ.\csromanspacer{} อถ โข ภควา ภิกฺขู อามนฺเตสิ\csromandash{}เตน หิ ภิกฺขเว สํโฆ วฑฺฒสฺส ลิจฺฉวิสฺส ปตฺตํ นิกฺกุชฺชตุ, อสมฺโภคํ สํเฆน กโรตุ.}

\prose{อฏฺฐหิ ภิกฺขเว องฺเคหิ สมนฺนาคตสฺส อุปาสกสฺส ปตฺโต นิกฺกุชฺชิตพฺโพ.\csromanspacer{} ภิกฺขูนํ อลาภาย ปริสกฺกติ, ภิกฺขูนํ อนตฺถาย ปริสกฺกติ, ภิกฺขูนํ อวาสาย\footnote{อนาวาสาย (สฺยา)} ปริสกฺกติ, ภิกฺขู อกฺโกสติ ปริภาสติ, ภิกฺขู ภิกฺขูหิ เภเทติ, พุทฺธสฺส อวณฺณํ ภาสติ, ธมฺมสฺส อวณฺณํ ภาสติ, สํฆสฺส อวณฺณํ ภาสติ.\csromanspacer{} อนุชานามิ ภิกฺขเว อิเมหิ อฏฺฐหงฺเคหิ สมนฺนาคตสฺส อุปาสกสฺส ปตฺตํ นิกฺกุชฺชิตุํ.\csromanspacer{} เอวญฺจ ปน ภิกฺขเว นิกฺกุชฺชิตพฺโพ, พฺยตฺเตน ภิกฺขุนา ปฏิพเลน สํโฆ ญาเปตพฺโพ\csromandash{}}

\proseitem{266}{“สุณาตุ เม ภนฺเต สํโฆ, วฑฺโฒ ลิจฺฉวี อายสฺมนฺตํ ทพฺพํ มลฺลปุตฺตํ อมูลิกาย สีลวิปตฺติยา อนุทฺธํเสติ, ยทิ สํฆสฺส ปตฺตกลฺลํ, สํโฆ วฑฺฒสฺส ลิจฺฉวิสฺส ปตฺตํ นิกฺกุชฺเชยฺย, อสมฺโภคํ สํเฆน กเรยฺย, เอสา ญตฺติ.}

\prose{สุณาตุ เม ภนฺเต สํโฆ, วฑฺโฒ ลิจฺฉวี อายสฺมนฺตํ ทพฺพํ มลฺลปุตฺตํ อมูลิกาย สีลวิปตฺติยา อนุทฺธํเสติ, สํโฆ วฑฺฒสฺส ลิจฺฉวิสฺส ปตฺตํ นิกฺกุชฺชติ, อสมฺโภคํ สํเฆน กโรติ, ยสฺสายสฺมโต ขมติ วฑฺฒสฺส ลิจฺฉวิสฺส ปตฺตสฺส นิกฺกุชฺชนา, อสมฺโภคํ สํเฆน กรณํ, โส ตุณฺหสฺส, ยสฺส นกฺขมติ, โส ภาเสยฺย.}

\prose{นิกฺกุชฺชิโต สํเฆน วฑฺฒสฺส ลิจฺฉวิสฺส ปตฺโต อสมฺโภโค สํเฆน, ขมติ สํฆสฺส, ตสฺมา หุณฺหี, เอวเมตํ ธารยามี”ติ.}

\prose{\csromanfolio{264}อถ โข อายสฺมา อานนฺโท ปุพฺพณฺหสมยํ นิวาเสตฺวา ปตฺตจีวร มาทาย เยน วฑฺฒสฺส ลิจฺฉวิสฺส นิเวสนํ เตนุปสงฺกมิ, อุปสงฺกมิตฺวา วฑฺฒํ ลิจฺฉวิํ เอตทโวจ “สํเฆน เต อาวุโส วฑฺฒ ปตฺโต นิกฺกุชฺชิโต อสมฺโภโคสิ สํเฆนา”ติ.\csromanspacer{} อถ โข วฑฺโฒ ลิจฺฉวี “สํเฆน กิร เม ปตฺโต นิกฺกุชฺชิโต อสมฺโภโคมฺหิ กิร สํเฆนา”ติ ตตฺเถว มุจฺฉิโต ปปโต.\csromanspacer{} อถ โข วฑฺฒสฺส ลิจฺฉวิสฺส มิตฺตามจฺจา ญาติสาโลหิตา วฑฺฒํ ลิจฺฉวิํ เอตทโวจุํ “อลํ อาวุโส วฑฺฒ, มา โสจิ มา ปริเทวิ, มยํ ภควนฺตํ ปสาเทสฺสาม ภิกฺขุสํฆญฺจา”ติ.}

\prose{อถ โข วฑฺโฒ ลิจฺฉวี สปุตฺตทาโร สมิตฺตามจฺโจ สญาติสาโลหิโต อลฺลวตฺโถ อลฺลเกโส เยน ภควา เตนุปสงฺกมิ, อุปสงฺกมิตฺวา ภควโต ปาเทสุ สิรสา นิปติตฺวา ภควนฺตํ เอตทโวจ “อจฺจโย มํ ภนฺเต อจฺจคมา ยถาพาลํ ยถามูฬฺหํ ยถา-อกุสลํ, โยหํ อยฺยํ ทพฺพํ มลฺลปุตฺตํ อมูลิกาย สีลวิปตฺติยา อนุทฺธํเสสิํ, ตสฺส เม ภนฺเต ภควา อจฺจยํ อจฺจยโต ปฏิคฺคณฺหาตุ อายติํ สํวรายา”ติ.\csromanspacer{} ตคฺฆ ตฺวํ อาวุโส วฑฺฒ อจฺจโย อจฺจคมา ยถาพาลํ ยถามูฬฺหํ ยถา-อกุสลํ, ยํ ตฺวํ ทพฺพํ มลฺลปุตฺตํ อมูลิกาย สีลวิปตฺติยา อนุทฺธํเสสิ, ยโต จ โข ตฺวํ อาวุโส วฑฺฒ อจฺจยํ อจฺจยโต ทิสฺวา ยถาธมฺมํ ปฏิกโรสิ, ตํ เต มยํ ปฏิคฺคณฺหาม, วุฑฺฒิเหสา อาวุโส วฑฺฒ อริยสฺส วินเย โย อจฺจยํ อจฺจยโต ทิสฺวา ยถาธมฺมํ ปฏิกโรติ, อายติํ สํวรํ อาปชฺชตีติ.\csromanspacer{} อถ โข ภควา ภิกฺขู อามนฺเตสิ\csromandash{}เตน หิ ภิกฺขเว สํโฆ วฑฺฒสฺส ลิจฺฉวิสฺส ปตฺตํ อุกฺกุชฺชตุ, สมฺโภคํ สํเฆน กโรตุ.}

\prose{อฏฺฐหิ ภิกฺขเว องฺเคหิ สมนฺนาคตสฺส อุปาสกสฺส ปตฺโต อุกฺกุชฺชิตพฺโพ.\csromanspacer{} น ภิกฺขูนํ อลาภาย ปริสกฺกติ, น ภิกฺขูนํ อนตฺถาย ปริสกฺกติ, น ภิกฺขูนํ อวาสาย ปริสกฺกติ, น ภิกฺขู อกฺโกสติ ปริภาสติ, น ภิกฺขู ภิกฺขูหิ เภเทติ, น พุทฺธสฺส อวณฺณํ ภาสติ, น ธมฺมสฺส อวณฺณํ ภาสติ, น สํฆสฺส อวณฺณํ ภาสติ.\csromanspacer{} อนุชานามิ ภิกฺขเว อิเมหิ อฏฺฐหงฺเคหิ สมนฺนาคตสฺส อุปาสกสฺส ปตฺตํ อุกฺกุชฺชิตุํ.\csromanspacer{} เอวญฺจ ปน ภิกฺขเว อุกฺกุชฺชิตพฺโพ.\csromanspacer{} เตน ภิกฺขเว วฑฺเฒน ลิจฺฉวินา สํฆํ อุปสงฺกมิตฺวา เอกํสํ อุตฺตราสงฺคํ กริตฺวา ภิกฺขูนํ ปาเท วนฺทิตฺวา อุกฺกุฏิกํ นิสีทิตฺวา อญฺชลิํ ปคฺคเหตฺวา เอวมสฺส วจนีโย “สํเฆน เม ภนฺเต \csromanfolio{265}ปตฺโต นิกฺกุชฺชิโต, อสมฺโภโคมฺหิ สํเฆน, โสหํ ภนฺเต สมฺมา วตฺตามิ, โลมํ ปาเตมิ, เนตฺถารํ วตฺตามิ, สํฆํ ปตฺตุกฺกุชฺชนํ ยาจามี”ติ.\csromanspacer{} ทุติยมฺปิ ยาจิตพฺโพ.\csromanspacer{} ตติยมฺปิ ยาจิตพฺโพ.\csromanspacer{} พฺยตฺเตน ภิกฺขุนา ปฏิพเลน สํโฆ ญาเปตพฺโพ\csromandash{}}

\proseitem{267}{“สุณาตุ เม ภนฺเต สํโฆ, สํเฆน วฑฺฒสฺส ลิจฺฉวิสฺส ปตฺโต นิกฺกุชฺชิโต, อสมฺโภโค สํเฆน, โส สมฺมา วตฺตติ, โลมํ ปาเตติ, เนตฺถารํ วตฺตติ, สํฆํ ปตฺตุกฺกุชฺชนํ ยาจติ, ยทิ สํฆสฺส ปตฺตกลฺลํ สํโฆ วฑฺฒสฺส ลิจฺฉวิสฺส ปตฺตํ อุกฺกุชฺเชยฺย, สมฺโภคํ สํเฆน กเรยฺย, เอสา ญตฺติ.}

\prose{สุณาตุ เม ภนฺเต สํโฆ, สํเฆน วฑฺฒสฺส ลิจฺฉวิสฺส ปตฺโต นิกฺกุชฺชิโต, อสมฺโภโค สํเฆน, โส สมฺมา วตฺตติ, โลมํ ปาเตติ, เนตฺถารํ วตฺตติ, สํฆํ ปตฺตุกฺกุชฺชนํ ยาจติ, สํโฆ วฑฺฒสฺส ลิจฺฉวิสฺส ปตฺตํ อุกฺกุชฺชติ, สมฺโภคํ สํเฆน กโรติ, ยสฺสายสฺมโต ขมติ วฑฺฒสฺส ลิจฺฉวิสฺส ปตฺตสฺส อุกฺกุชฺชนา, สมฺโภคํ สํเฆน กรณํ, โส ตุณฺหสฺส, ยสฺส นกฺขมติ, โส ภาเสยฺย.}

\prose{อุกฺกุชฺชิโต สํเฆน วฑฺฒสฺส ลิจฺฉวิสฺส ปตฺโต สมฺโภโค สํเฆน, ขมติ สํฆสฺส, ตสฺมา ตุณฺหี, เอวเมตํ ธารยามี”ติ.}

\proseitem{268}{อถ โข ภควา เวสาลิยํ ยถาภิรนฺตํ วิหริตฺวา เยน ภคฺคา เตน จาริกํ ปกฺกามิ, อนุปุพฺเพน จาริกํ จรมาโน เยน ภคฺคา ตทวสริ, ตตฺร สุทํ ภควา ภคฺเคสุ วิหรติ สุสุมารคิเร\footnote{สุํสูมารคิเร (สี, สฺยา), สํสุมารคิเร (ก)} เภสกฬวเน มิคทาเย.\csromanspacer{} เตน โข ปน สมเยน โพธิสฺส ราชกุมารสฺส โกกนโท\footnote{โกกนุโท (ก)} นาม ปาสาโท อจิรการิโต โหติ อนชฺฌาวุตฺโถ สมเณน วา พฺราหฺมเณน วา เกนจิ วา มนุสฺสภูเตน.\csromanspacer{} อถ โข โพธิ ราช กุมาโร สญฺจิกาปุตฺตํ มาณวํ อามนฺเตสิ\csromandash{}เอหิ ตฺวํ สมฺม สญฺจิกาปุตฺต, เยน ภควา เตนุปสงฺกม, อุปสงฺกมิตฺวา มม วจเนน ภควโต ปาเท สิรสา วนฺท, อปฺปาพาธํ อปฺปาตงฺกํ ลหุฏฺฐานํ พลํ ผาสุวิหารํ ปุจฺฉ “โพธิ ภนฺเต ราชกุมาโร ภควโต ปาเท สิรสา วนฺทติ, อปฺปาพาธํ \csromanfolio{266}อปฺปาตงฺกํ ลหุฏฺฐานํ พลํ ผาสุวิหารํ ปุจฺฉตี”ติ.\csromanspacer{} เอวญฺจ วเทหิ “อธิวาเสตุ กิร ภนฺเต ภควา โพธิสฺส ราชกุมารสฺส สฺวาตนาย ภตฺตํ สทฺธิํ ภิกฺขุสํเฆนา”ติ. “เอวํ โภ”ติ โข สญฺจิกาปุตฺโต มาณโว โพธิสฺส ราชกุมารสฺส ปฏิสฺสุตฺวา เยน ภควา เตนุปสงฺกมิ, อุปสงฺกมิตฺวา ภควตา สทฺธิํ สมฺโมทิ, สมฺโมทนียํ กถํ สารณียํ วีติสาเรตฺวา เอกมนฺตํ นิสีทิ, เอกมนฺตํ นิสินฺโน โข สญฺจิกาปุตฺโต มาณโว ภควนฺตํ เอตทโวจ\csromandash{}โพธิ โข ราชกุมาโร โภโต โคตมสฺส ปาเท สิรสา วนฺทติ, อปฺปาพาธํ อปฺปาตงฺกํ ลหุฏฺฐานํ พลํ ผาสุวิหารํ ปุจฺฉติ, เอวญฺจ วเทติ “อธิวาเสตุ กิร ภวํ โคตโม โพธิสฺส ราชกุมารสฺส สฺวาตนาย ภตฺตํ สทฺธิํ ภิกฺขุสํเฆนา”ติ.\csromanspacer{} อธิวาเสสิ ภควา ตุณฺหีภาเวน.\csromanspacer{} อถ โข สญฺจิกาปุตฺโต มาณโว ภควโต อธิวาสนํ วิทิตฺวา อุฏฺฐายาสนา เยน โพธิ ราชกุมาโร เตนุปสงฺกมิ, อุปสงฺกมิตฺวา โพธิํ ราชกุมารํ เอตทโวจ “อโวจุมฺห โข มยํ โภโต วจเนน ตํ ภวนฺตํ โคตมํ ‘โพธิ โข ราชกุมาโร โภโต โคตมสฺส ปาเท สิรสา วนฺทติ, อปฺปาพาธํ อปฺปาตงฺกํ ลหุฏฺฐานํ พลํ ผาสุวิหารํ ปุจฺฉติ, เอวญฺจ วเทติ อธิวาเสตุ กิร ภวํ โคตโม โพธิสฺส ราชกุมารสฺส สฺวาตนาย ภตฺตํ สทฺธิํ ภิกฺขุสํเฆนา’ติ, อธิวุตฺถญฺจ ปน สมเณน โคตเมนา”ติ.}

\prose{อถ โข โพธิ ราชกุมาโร ตสฺสา รตฺติยา อจฺจเยน ปณีตํ ขาทนียํ โภชนียํ ปฏิยาทาเปตฺวา โกกนทญฺจ ปาสาทํ โอทาเตหิ ทุสฺเสหิ สนฺตราเปตฺวา ยาว ปจฺฉิมโสปานกเฬวรา สญฺจิกาปุตฺตํ มาณวํ อามนฺเตสิ\csromandash{}เอหิ ตฺวํ สมฺม สญฺจิกาปุตฺต, เยน ภควา เตนุปสงฺกม, อุปสงฺกมิตฺวา ภควโต กาลํ อาโรเจหิ “กาโล ภนฺเต, นิฏฺฐิตํ ภตฺตนฺ”ติ. “เอวํ โภ”ติ โข สญฺจิกาปุตฺโต มาณโว โพธิสฺส ราชกุมารสฺส ปฏิสฺสุตฺวา เยน ภควา เตนุปสงฺกมิ, อุปสงฺกมิตฺวา ภควโต กาลํ อาโรเจสิ “กาโล โภ โคตม, นิฏฺฐิตํ ภตฺตนฺ”ติ.}

\prose{อถ โข ภควา ปุพฺพณฺหสมยํ นิวาเสตฺวา ปตฺตจีวรทาย เยน โพธิสฺส ราชกุมารสฺส นิเวสนํ เตนุปสงฺกมิ.\csromanspacer{} เตน โข ปน สมเยน โพธิ ราชกุมาโร พหิทฺวารโกฏฺฐเก ฐิโต โหติ ภควนฺตํ อาคมยมาโน.\csromanspacer{} อทฺทสา โข โพธิ ราชกุมาโร ภควนฺตํ \csromanfolio{267}ทูรโตว อาคจฺฉนฺตํ, ทิสฺวาน ตโต ปจฺจุคฺคนฺตฺวา ภควนฺตํ อภิวาเทตฺวา ปุเรกฺขตฺวา เยน โกกนโท ปาสาโท เตนุปสงฺกมิ.\csromanspacer{} อถ โข ภควา ปจฺฉิมโสปานกเฬวรํ นิสฺสาย อฏฺฐาสิ.\csromanspacer{} อถ โข โพธิ ราชกุมาโร ภควนฺตํ เอตทโวจ “อกฺกมตุ ภนฺเต ภควา ทุสฺสานิ, อกฺกมตุ สุคโต ทุสฺสานิ, ยํ มม อสฺส ทีฆรตฺตํ หิตาย สุขายา”ติ.\csromanspacer{} เอวํ วุตฺเต ภควา ตุณฺหี อโหสิ.\csromanspacer{} ทุติยมฺปิ โข ฯเปฯ.\csromanspacer{} ตติยมฺปิ โข โพธิ ราชกุมาโร ภควนฺตํ เอตทโวจ “อกฺกมตุ ภนฺเต ภควา ทุสฺสานิ, อกฺกมตุ สุคโต ทุสฺสานิ, ยํ มม อสฺส ทีฆรตฺตํ หิตาย สุขายา”ติ.\csromanspacer{} อถ โข ภควา อายสฺมนฺตํ อานนฺทํ อปโลเกสิ.\csromanspacer{} อถ โข อายสฺมา อานนฺโท โพธิํ ราชกุมารํ เอตทโวจ “สํหรนฺตุ ราชกุมาร ทุสฺสานิ, น ภควา เจลปฏิกํ\footnote{เจลปตฺติกํ (สี)} อกฺกมิสฺสติ, ปจฺฉิมํ ชนตํ ตถาคโต อนุกมฺปตี”ติ.}

\prose{อถ โข โพธิ ราชกุมาโร ทุสฺสานิ สํหราเปตฺวา อุปริโกกนเท ปาสาเท อาสนํ ปญฺญเปสิ.\csromanspacer{} อถ โข ภควา โกกนทํ ปาสาทํ อภิรุหิตฺวา ปญฺญตฺเต อาสเน นิสีทิ สทฺธิํ ภิกฺขุสํเฆน.\csromanspacer{} อถ โข โพธิ ราชกุมาโร พุทฺธปฺปมุขํ ภิกฺขุสํฆํ ปณีเตน ขาทนีเยน โภชนีเยน สหตฺถา สนฺตปฺเปตฺวา สมฺปวาเรตฺวา ภควนฺตํ ภุตฺตาวิํ โอนีตปตฺตปาณิํ เอกมนฺตํ นิสีทิ, เอกมนฺตํ นิสินฺนํ โข โพธิํ ราชกุมารํ ภควา ธมฺมิยา กถาย สนฺทสฺเสตฺวา สมาทเปตฺวา สมุตฺเตเชตฺวา สมฺปหํเสตฺวา อุฏฺฐายาสนา ปกฺกามิ.\csromanspacer{} อถ โข ภควา เอตสฺมิํ นิทาเน เอตสฺมิํ ปกรเณ ธมฺมิํ กถํ กตฺวา ภิกฺขู อามนฺเตสิ “น ภิกฺขเว เจลปฺปฏิกา อกฺกมิตพฺพา, โย อกฺกเมยฺย, อาปตฺติ ทุกฺกฏสฺสา”ติ.}

\prose{เตน โข ปน สมเยน อญฺญตรา อิตฺตี อปคตคพฺภา ภิกฺขู นิมนฺเตตฺวา ทุสฺสํ ปญฺญเปตฺวา เอตทโวจ “อกฺกมถ ภนฺเต ทุสฺสนฺ”ติ.\csromanspacer{} ภิกฺขู กุกฺกุจฺจายนฺตา น อกฺกมนฺติ.\csromanspacer{} อกฺกมถ ภนฺเต ทุสฺสํ มงฺคลตฺถายาติ.\csromanspacer{} ภิกฺขู กุกฺกุจฺจายนฺตา น อกฺกมิํสุ.\csromanspacer{} อถ โข สา อิตฺถี อุชฺฌายติ ขิยฺยติ วิปาเจติ “กถํ หิ นาม อยฺยา มงฺคลตฺถาย ยาจิยมานา เจลปฺปฏิกํ น อกฺกมิสฺสนฺตี”ติ.\csromanspacer{} อโสสุํ โข ภิกฺขู ตสฺสา อิตฺถิยา อุชฺฌายนฺติยา ขิยฺยนฺติยา วิปาเจนฺติยา.\csromanspacer{} อถ โข เต ภิกฺขู ภควโต \csromanfolio{268}เอตมตฺถํ อาโรเจสุํ.\csromanspacer{} คิหี ภิกฺขเว มงฺคลิกา.\csromanspacer{} อนุชานามิ ภิกฺขเว คิหีนํ มงฺคลตฺถาย ยาจิยมาเนน เจลปฺปฏิกํ อกฺกมิตุนฺติ.}

\prose{เตน โข ปน สมเยน ภิกฺขู โธตปาทกํ อกฺกมิตุํ กุกฺกุจฺจายนฺติ.\csromanspacer{} ภควโต เอตมตฺถํ อาโรเจสุํ.\csromanspacer{} อนุชานามิ ภิกฺขเว โธตปาทกํ อกฺกมิตุนฺติ.}

\nitthitam{ทุติยภาณวาโร นิฏฺฐิโต.}
\proseitem{269}{อถ โข ภควา ภคฺเคสุ ยถาภิรนฺตํ วิหริตฺวา เยน สาวตฺถิ เตน จาริกํ ปกฺกามิ, อนุปุพฺเพน จาริกํ จรมาโน เยน สาวตฺถิ ตทวสริ, ตตฺร สุทํ ภควา สาวตฺถิยํ วิหรติ เชตวเน อนาถปิณฺฑิกสฺส อาราเม.\csromanspacer{} อถ โข วิสาขา มิคารมาตา ฆฏกญฺจ กตกญฺจ สมฺมชฺชนิญฺจ อาทาย เยน ภควา เตนุปสงฺกมิ, อุปสงฺกมิตฺวา ภควนฺตํ อภิวาเทตฺวา เอกมนฺตํ นิสีทิ, เอกมนฺตํ นิสินฺนา โข วิสาขา มิคารมาตา ภควนฺตํ เอตทโวจ “ปฏิคฺคณฺหาตุ เม ภนฺเต ภควา ฆฏกญฺจ กตกญฺจ สมฺมชฺชนิญฺจ, ยํ มม อสฺส ทีฆรตฺตํ หิตาย สุขายา”ติ.\csromanspacer{} ปฏิคฺคเหสิ ภควา ฆฏกญฺจ สมฺมชฺชนิญฺจ, น ภควา กตกํ ปฏิคฺคเหสิ.\csromanspacer{} อถ โข ภควา วิสาขํ มิคารมาตรํ ธมฺมิยา กถาย สนฺทสฺเสสิ สมาทเปสิ สมุตฺเตเชสิ สมฺปหํเสสิ.\csromanspacer{} อถ โข วิสาขา มิคารมาตา ภควตา ธมฺมิยา กถาย สนฺทสฺสิตา สมาทปิตา สมุตฺเตชิตา สมฺปหํสิตา อุฏฺฐายาสนา ภควนฺตํ อภิวาเทตฺวา ปทกฺขิณํ กตฺวา ปกฺกามิ.}

\prose{อถ โข ภควา เอตสฺมิํ นิทาเน เอตสฺมิํ ปกรเณ ธมฺมิํ กถํ กตฺวา ภิกฺขู อามนฺเตสิ “อนุชานามิ ภิกฺขเว ฆฏกญฺจ สมฺมชฺชนิญฺจ.\csromanspacer{} น ภิกฺขเว กตกํ ปริภุญฺชิตพฺพํ, โย ปริภุญฺเชยฺย, อาปตฺติ ทุกฺกฏสฺส.\csromanspacer{} อนุชานามิ ภิกฺขเว ติสฺโส ปาทฆํสนิโย สกฺขรํ กถลํ สมุทฺทเผณกนฺ”ติ.}

\prose{อถ โข วิสาขา มิคารมาตา วิธูปนญฺจ ตาลวณฺฏญฺจ\footnote{ตาลวณฺฑญฺจ (ก)} อาทาย เยน ภควา เตนุปสงฺกมิ, อุปสงฺกมิตฺวา ภควนฺตํ อภิวาเทตฺวา เอกมนฺตํ นิสีทิ, เอกมนฺตํ นิสินฺนา โข วิสาขา มิคารมาตา ภควนฺตํ นิสินฺนา โข วิสาขา มิคารมาตา ภควนฺตํ เอตทโวจ “ปฏิคฺคณฺหาตุ เม ภนฺเต ภควา วิธูปนญฺจ ตาลวณฺฏญฺจ, ยํ มม อสฺส ทีฆรตฺตํ หิตาย สุขายา”ติ.\csromanspacer{} ปฏิคฺคเหสิ ภควา วิธูปนญฺจ ตาลวณฺฏญฺจ.}

\prose{\csromanfolio{269}อถ โข ภควา วิสาขํ มิคารมาตรํ ธมฺมิยา กถาย สนฺทสฺเสสิ สมาทเปสิ สมุตฺเตเชสิ สมฺปหํเสสิ ฯเปฯ ปทกฺขิณํ กตฺวา ปกฺกามิ.\csromanspacer{} อถ โข ภควา เอตสฺมิํ นิทาเน เอตสฺมิํ ปกรเณ ธมฺมิํ กถํ กตฺวา ภิกฺขู อามนฺเตสิ “อนุชานามิ ภิกฺขเว วิธูปนญฺจ ตาลวณฺฏญฺจา”ติ.}

\prose{เตน โข ปน สมเยน สํฆสฺส มกสพีชนี อุปฺปนฺนา โหติ.\csromanspacer{} ภควโต เอตมตฺถํ อาโรเจสุํ.\csromanspacer{} อนุชานามิ ภิกฺขเว มกสพีชนินฺติ.\csromanspacer{} จามริพีชนี อุปฺปนฺนา โหติ.\csromanspacer{} ภควโต เอตมตฺถํ อาโรเจสุํ.\csromanspacer{} น ภิกฺขเว จามริพีชนี ธาเรตพฺพา, โย ธาเรยฺย, อาปตฺติ ทุกฺกฏสฺส.\csromanspacer{} อนุชานามิ ภิกฺขเว ติสฺโส พีชนิโย วากมยํ อุสีรมยํ โมรปิญฺฉมยนฺติ.}

\proseitem{270}{เตน โข ปน สมเยน สํฆสฺส ฉตฺตํ อุปฺปนฺนํ โหติ.\csromanspacer{} ภควโต เอตมตฺถํ อาโรเจสุํ.\csromanspacer{} อนุชานามิ ภิกฺขเว ฉตฺตนฺติ.}

\prose{เตน โข ปน สมเยน ฉพฺพคฺคิยา ภิกฺขู ฉตฺตปฺปคฺคหิตา\footnote{ฉตฺตํ มคฺคเหตฺวา (ก)} อาหิณฺฑนฺติ.\csromanspacer{} เตน โข ปน สมเยน อญฺญตโร อุปาสโก สมฺพหุเลหิ อาชีวกสาวเกหิ สทฺธิํ อุยฺยานํ อคมาสิ.\csromanspacer{} อทฺทสาสุํ โข เต อาชีวกสาวกา ฉพฺพคฺคิเย ภิกฺขู ทูรโตว ฉตฺตปฺปคฺคหิเต อาคจฺฉนฺเต, ทิสฺวาน ตํ อุปาสกํ เอตทโวจุํ “เอเต โข อยฺยา\footnote{อยฺโย (ก)} ตุมฺหากํ ภทนฺตา ฉตฺตปฺปคฺคหิตา อาคจฺฉนฺติ เสยฺยถาปิ คณกมหามตฺตา”ติ.\csromanspacer{} นายฺยา เอเต ภิกฺขู ปริพฺพาชกาติ. “ภิกฺขู น ภิกฺขู”ติ อพฺภุตํ อกํสุ.\csromanspacer{} อถ โข โส อุปาสโก อุปคเต สญฺชานิตฺวา อุชฺฌายติ ขิยฺยติ วิปาเจติ “กถํ หิ นาม ภทนฺตา ฉตฺตปฺปคฺคหิตา อาหิณฺฑิสฺสนฺตี”ติ.\csromanspacer{} อสฺโสสุํ โข ภิกฺขู ตสฺส อุปาสกสฺส อุชฺฌายนฺตสฺส ขิยฺยนฺตสฺส วิปาเจนฺตสฺส.\csromanspacer{} อถ โข เต ภิกฺขู ภควโต เอตมตฺถํ อาโรเจสุํ.\csromanspacer{} สจฺจํ กิร ภิกฺขเว ฯเปฯ.\csromanspacer{} สจฺจํ ภควาติ ฯเปฯ วิครหิตฺวา ฯเปฯ ธมฺมิํ กถํ กตฺวา ภิกฺขู อามนฺเตสิ “น ภิกฺขเว ฉตฺตํ ธาเรตพฺพํ, โย ธาเรยฺย, อาปตฺติ ทุกฺกฏสฺสา”ติ.}

\prose{เตน โข ปน สมเยน อญฺญตโร ภิกฺขุ คิลาโน โหติ, ตสฺส ภิกฺขุโน วินา ฉตฺตํ น ผาสุ โหติ.\csromanspacer{} ภควโต เอตมตฺถํ อาโรเจสุํ.\csromanspacer{} อนุชานามิ ภิกฺขเว คิลานสฺส ฉตฺตํ ธาเรตุนฺติ.}

\prose{\csromanfolio{270}เตน โข ปน สมเยน ภิกฺขู “คิลานสฺเสว ภควตา ฉตฺตํ อนุญฺญาตํ โน อคิลานสฺสา”ติ อาราเม อารามูปจาเร ฉตฺตํ ธาเรตุํ กุกฺกุจฺจายนฺติ.\csromanspacer{} ภควโต เอตมตฺถํ อาโรเจสุํ.\csromanspacer{} อนุชานามิ ภิกฺขเว อคิลาเนนปิ อาราเม อารามูปจาเร ฉตฺตํ ธาเรตุนฺติ.}

\prose{เตน โข ปน สมเยน อญฺญตโร ภิกฺขุ สิกฺกาย ปตฺตํ อุฏฺฏิตฺวา ทณฺเฑ อาลคฺคิตฺวา วิกาเล อญฺญตเรน คามทฺวาเรน อติกฺกมติ.\csromanspacer{} มนุสฺสา “เอสยฺโย โจโร คจฺฉติ, อสิสฺส วิชฺโชตลตี”ติ อนุปติตฺวา คเหตฺวา สญฺชานิตฺวา มุญฺจิํสุ.\csromanspacer{} อถ โข โส ภิกฺขุ อารามํ คนฺตฺวา ภิกฺขูนํ เอตมตฺถํ อาโรเจสิ.\csromanspacer{} กิํ ปน ตฺวํ อาวุโส ทณฺฑสิกฺกํ ธาเรสีติ.\csromanspacer{} เอวมาวุโสติ.\csromanspacer{} เย เต ภิกฺขู อปฺปิจฺฉา ฯเปฯ เต อุชฺฌายนฺติ ขิยฺยนฺติ วิปาเจนฺติ “กถํ หิ นาม ภิกฺขุ ทณฺฑสิกฺกํ ธาเรสฺสสี”ติ.\csromanspacer{} อถ โข เต ภิกฺขู ภควโต เอตมตฺถํ อาโรเจสุํ ฯเปฯ.\csromanspacer{} สจฺจํ ภควาติ ฯเปฯ วิครหิตฺวา ฯเปฯ ธมฺมิํ กถํ กตฺวา ภิกฺขู อามนฺเตสิ “น ภิกฺขเว ทณฺฑสิกฺกา ธาเรตพฺพา, โย ธาเรยฺย, อาปตฺติ ทุกฺกฏสฺสา”ติ.}

\prose{เตน โข ปน สมเยน อญฺญตโร ภิกฺขุ คิลาโน โหติ, น สกฺโกติ วินา ทณฺเฑน อาหิณฺฑิตุํ.\csromanspacer{} ภควโต เอตมตฺถํ อาโรเจสุํ.\csromanspacer{} อนุชานามิ ภิกฺขเว คิลานสฺส ภิกฺขุโน ทณฺฑสมฺมุติํ ทาตุํ.\csromanspacer{} เอวญฺจ ปน ภิกฺขเว ทาตพฺพา, เตน คิลาเนน ภิกฺขุนา สํฆํ อุปสงฺกมิตฺวา เอกํสํ อุตฺตราสงฺคํ กริตฺวา วุฑฺฒานํ ภิกฺขูนํ ปาเท วนฺทิตฺวา อุกฺกุฏิกํ นิสีทิตฺวา อญฺชลิํ ปคฺคเหตฺวา เอวมสฺส วจนีโย “อหํ ภนฺเต คิลาโน น สกฺโกมิ วินา ทณฺเฑน อาหิณฺฑิตุํ, โสหํ ภนฺเต สํฆํ ทณฺฑสมฺมุติํ ยาจมี”ติ.\csromanspacer{} ทุติยมฺปิ ยาจิตพฺพา.\csromanspacer{} ตติยมฺปิ ยาจิตพฺพา.\csromanspacer{} พฺยตฺเตน ภิกฺขุนา ปฏิพเลน สํโฆ ญาเปตพฺโพ\csromandash{}}

\prose{“สุณาตุ เม ภนฺเต สํโฆ, อยํ อิตฺถนฺนาโม ภิกฺขุ คิลาโน น สกฺโกติ วินา ทณฺเฑน อาหิณฺฑิตุํ, โส สํฆํ ทณฺฑสมฺมุติํ ยาจติ, ยทิ สํฆสฺส ปตฺตกลฺลํ, สํโฆ อิตฺถนฺนามสฺส ภิกฺขุโน ทณฺฑสมฺมุติํ ทเทยฺย, เอสา ญตฺติ.}

\prose{สุณาตุ เม ภนฺเต สํโฆ, อยํ อิตฺถนฺนาโม ภิกฺขุ คิลาโน น สกฺโกติ วิเน ทณฺเฑน อาหิณฺฑิตุํ, โส สํฆํ ทณฺฑสมฺมุติํ ยาจติ, \csromanfolio{271}สํโฆ อิตฺถนฺนามสฺส ภิกฺขุโน ทณฺฑสมฺมุติํ เทติ, ยสฺสายสฺมโต ขมติ อิตฺถนฺนามสฺส ภิกฺขุโน ทณฺฑสมฺมุติยา ทานํ, โส ตุณฺหสฺส, ยสฺส นกฺขมติ, โส ภาเสยฺย.}

\prose{ทินฺนา สํเฆน อิตฺถนฺนามสฺส ภิกฺขุโน ทณฺฑสมฺมุติ, ขมติ สํฆสฺส, ตสฺมา ตุณฺหี, เอวเมตํ ธารยามี”ติ.}

\proseitem{271}{เตน โข ปน สมเยน อญฺญตโร ภิกฺขุ คิลาโน โหติ, น สกฺโกติ วินา สิกฺกาย ปตฺตํ ปริหริตุํ.\csromanspacer{} ภควโต เอตมตฺถํ อาโรเจสุํ.\csromanspacer{} อนุชานามิ ภิกฺขเว คิลานสฺส ภิกฺขุโน สิกฺกาสมฺมุติํ ทาตุํ.\csromanspacer{} เอวญฺจ ปน ภิกฺขเว ทาตพฺพา, เตน คิลาเนน ภิกฺขุนา สํฆํ อุปสงฺกมิตฺวา เอกํสํ อุตฺตราสงฺคํ กริตฺวา วุฑฺฒานํ ภิกฺขูนํ ปาเท วนฺทิตฺวา อุกฺกุฏิกํ นิสีทิตฺวา อญฺชลิํ ปคฺคเหตฺวา เอวมสฺส วจนีโย “อหํ ภนฺเต คิลาโน น สกฺโกมิ วินา สิกฺกาย ปตฺตํ ปริหริตุํ, โสหํ ภนฺเต สํฆํ สิกฺกาสมฺมุติํ ยาจามี”ติ.\csromanspacer{} ทุติยมฺปิ ยาจิตพฺพา.\csromanspacer{} ตติยมฺปิ ยาจิตพฺพา.\csromanspacer{} พฺยตฺเตน ภิกฺขุนา ปฏิพเลน สํโฆ ญาเปตพฺโพ\csromandash{}}

\prose{“สุณาตุ เม ภนฺเต สํโฆ, อยํ อิตฺถนฺนาโม ภิกฺขุ คิลาโน น สกฺโกติ วินา สิกฺกาย ปตฺตํ ปริหริตุํ, โส สํฆํ สิกฺกาสมฺมุติํ ยาจติ, ยทิ สํฆสฺส ปตฺตกลฺลํ, สํโฆ อิตฺถนฺนามสฺส ภิกฺขุโน สิกฺกาสมฺมุติํ ทเทยฺย, เอสา ญตฺติ.}

\prose{สุณาตุ เม ภนฺเต สํโฆ, อยํ อิตฺถนฺนาโม ภิกฺขุ คิลาโน น สกฺโกติ วินา สิกฺกาย ปตฺตํ ปริหริตุํ, โส สํฆํ สิกฺกาสมฺมุติํ ยาจติ, สํโฆ อิตฺถนฺนามสฺส ภิกฺขุโน สิกฺกาสมฺมุติํ เทติ, ยสฺสายสฺมโต ขมติ อิตฺถนฺนามสฺส ภิกฺขุโน สิกฺกาสมฺมุติยา ทานํ, โส ตุณฺหสฺส, ยสฺส นกฺขมติ, โส ภาเสยฺย.}

\prose{ทินฺนา สํเฆน อิตฺถนฺนามสฺส ภิกฺขุโน สิกฺกาสมฺมุติ, ขมติ สํฆสฺส, ตสฺมา ตุณฺหี, เอวเมตํ ธารยามี”ติ.}

\proseitem{272}{เตน โข ปน สมเยน อญฺญตโร ภิกฺขุ คิลาโน โหติ, น สกฺโกติ วินา ทณฺเฑน อาหิณฺฑิตุํ, น สกฺโกติ วินา สิกฺกาย ปตฺตํ ปริหริตุํ.\csromanspacer{} ภควโต เอตมตฺถํ อาโรเจสุํ.\csromanspacer{} อนุชานามิ ภิกฺขเว คิลานสฺส ภิกฺขุโน ทณฺฑสิกฺกาสมฺมุติํ ทาตุํ.\csromanspacer{} เอวญฺจ ปน ภิกฺขเว \csromanfolio{272}ทาตพฺพา, เตน คิลาเนน ภิกฺขุนา สํฆํ อุปสงฺกมิตฺวา เอกํสํ อุตฺตราสงฺคํ กริตฺวา วุฑฺฒานํ ภิกฺขูนํ ปาเท วนฺทิตฺวา อุกฺกุฏิกํ นิสีทิตฺวา อญฺจลิํ ปคฺคเหตฺวา เอวมสฺส วจนีโย “อหํ ภนฺเต คิลาโน น สกฺโกมิ วินา ทณฺเฑน อาหิณฺฑิตุํ, น สกฺโกมิ วินา สิกฺกาย ปตฺตํ ปริหริตุํ, โสหํ ภนฺเต สํฆํ ทณฺฑสิกฺกาสมฺมุติํ ยาจามี”ติ.\csromanspacer{} ทุติยมฺปิ ยาจิตพฺพา.\csromanspacer{} ตติยมฺปิ ยาจิตพฺพา.\csromanspacer{} พฺยตฺเตน ภิกฺขุนา ปฏิพเลน สํโฆ ญาเปตพฺโพ\csromandash{}}

\prose{“สุณาตุ เม ภนฺเต สํโฆ, อยํ อิตฺถนฺนาโม ภิกฺขุ คิลาโน น สกฺโกติ วินา ทณฺเฑน อาหิณฺฑิตุํ, น สกฺโกติ วินา สิกฺกาย ปตฺตํ ปริหริตุํ, โส สํฆํ ทณฺฑสิกฺกาสมฺมุติํ ยาจติ, ยทิ สํฆสฺส ปตฺตกลฺลํ, สํโฆ อิตฺถนฺนามสฺส ภิกฺขุโน ทณฺฑสิกฺกาสมฺมุติํ ทเทยฺย, เอสา ญตฺติ.}

\prose{สุณาตุ เม ภนฺเต สํโฆ, อยํ อิตฺถนฺนาโม ภิกฺขุ คิลาโน น สกฺโกติ วินา ทณฺเฑน อาหิณฺฑิตุํ, น สกฺโกติ วินา สิกฺกาย ปตฺตํ ปริหริตุํ, โส สํฆํ ทณฺฑสิกฺกาสมฺมุติํ ยาจติ, สํโฆ อิตฺถนฺนามสฺส ภิกฺขุโน ทณฺฑสิกฺกาสมฺมุติํ เทติ, ยสฺสายสฺมโต ขมติ อิตฺถนฺนามสฺส ภิกฺขุโน ทณฺฑสิกฺกาสมฺมุติยา ทานํ, โส ตุณฺหสฺส, ยสฺส นกฺขมติ, โส ภาเสยฺย.}

\prose{ทินฺนา สํเฆน อิตฺถนฺนามสฺส ภิกฺขุโน ทณฺฑสิกฺกาสมฺมุติ, ขมติ สํฆสฺส, ตสฺมา ตุณฺหี, เอวเมตํ ธารยามี”ติ.}

\proseitem{273}{เตน โข ปน สมเยน อญฺญตโร ภิกฺขุ โรมนฺถโก\footnote{เอรามฏฺฐโก (ก)} โหติ, โส โรมนฺถิตฺวา โรมนฺถิตฺวา อชฺโฌหรติ.\csromanspacer{} ภิกฺขู อุชฺฌายนฺติ ขิยฺยนฺติ วิปาเจนฺติ “วิกาลายํ\footnote{กถํ หิ นาม วิกาลายํ (ก)} ภิกฺขุ โภชนํ ภุญฺชตี”ติ.\csromanspacer{} ภควโต เอตมตฺถํ อาโรเจสุํ.\csromanspacer{} เอโส ภิกฺขเว ภิกฺขุ อจิรํโคโยนิยา จุโต.\csromanspacer{} อนุชานามิ ภิกฺขเว โรมนฺถกสฺส โรมนฺถนํ.\csromanspacer{} น จ ภิกฺขเว พหิมุขทฺวารํ นีหริตฺวา อชฺโฌหริตพฺพํ, โย อชฺโฌหเรยฺย, ยถาธมฺโม กาเรตพฺโพติ.}

\prose{\csromanfolio{273}เตน โข ปน สมเยน อญฺญตรสฺส ปูคสฺส สํฆภตฺตํ โหติ, ภตฺตคฺเค พหุสิตฺถานิ ปกิริยิํสุ\footnote{วิปฺปกิรียิํสุ (สี), ปริกิริํสุ (สฺยา)}.\csromanspacer{} มนุสฺสา อุชฺฌายนฺติ ขิยฺยนฺติ วิปาเจนฺติ “กถํ หิ นาม สมณา สกฺยปุตฺติยา โอทเน ทิยฺยมาเน น สกฺกจฺจํ ปฏิคฺคเหสฺสนฺติ, เอกเมกํ สิตฺถํ กมฺมสเตน นิฏฺฐายตี”ติ.\csromanspacer{} อสฺโสสุํ โข ภิกฺขู เตสํ มนุสฺสานํ อุชฺฌายนฺตานํ ขิยฺยนฺตานํ วิปาเจนฺตานํ.\csromanspacer{} อถ โข เต ภิกฺขู ภควโต เอตมตฺถํ อาโรเจสุํ.\csromanspacer{} อนุชานามิ ภิกฺขเว ยํ ทิยฺยมานํ ปตติ, ตํ สามํ คเหตฺวา ปริภุญฺชิตุํ, ปริจฺจตฺตํ ตํ ภิกฺขเว ทายเกหีติ.}

\proseitem{274}{เตน โข ปน สมเยน อญฺญตโร ภิกฺขุ ทีเฆหิ นเขหิ ปิณฺฑาย จรติ, อญฺญตรา อิตฺถี ปสฺสิตฺวา ตํ ภิกฺขุํ เอตทโวจ “เอหิ ภนฺเต เมถุนํ ธมฺมํ ปฏิเสวา”ติ.\csromanspacer{} อลํ ภคินิ, เนตํ กปฺปตีติ.\csromanspacer{} สเจ โข ตฺวํ ภนฺเต นปฺปฏิเสวิสฺสสิ, อิทานาหํ อตฺตโน นเขหิ คตฺตานิ วิลิขิตฺวา กุปฺปํ กริสฺสามิ “อยํ มํ ภิกฺขุ วิปฺปกโรตี”ติ.\csromanspacer{} ปชานาหิ ตฺวํ ภคินีติ.\csromanspacer{} อถ โข สา อิตฺถี อตฺตโน นเขหิ คตฺตานิ วิลิขิตฺวา กุปฺปํ อกาสิ “อยํ มํ ภิกฺขุ วิปฺปกโรตี”ติ.\csromanspacer{} มนุสฺสา อุปธาวิตฺวา ตํ ภิกฺขุํ อคฺคเหสุํ.\csromanspacer{} อทฺทสาสุํ โข เต มนุสฺสา ตสฺสา อิตฺถิยา นเข ฉวิมฺปิ โลหิตมฺปิ.\csromanspacer{} ทิสฺวาน “อิมิสฺสาเยว อิตฺถิยา อิทํ กมฺมํ, อการโก ภิกฺขู”ติ ตํ ภิกฺขุํ มุญฺจิํสุ.\csromanspacer{} อถ โข โส ภิกฺขุ อารามํ คนฺตฺวา ภิกฺขูนํ เอตมตฺถํ อาโรเจสิ.\csromanspacer{} กิํ ปน ตฺวํ อาวุโส ทีเฆ นเข ธาเรสีติ.\csromanspacer{} เอวมาวุโสติ.\csromanspacer{} เย เต ภิกฺขู อปฺปิจฺฉา ฯเปฯ เต อุชฺฌายนฺติ ขิยฺยนฺติ วิปาเจนฺติ “กถํ หิ นาม ภิกฺขุ ทีเฆ นเข ธาเรสฺสสี”ติ.\csromanspacer{} อถ โข เต ภิกฺขู ภควโต เอตมตฺถํ อาโรเจสุํ.\csromanspacer{} น ภิกฺขเว ทีฆา นขา ธาเรตพฺพา, โย ธาเรยฺย, อาปตฺติ ทุกฺกฏสฺสาติ.}

\prose{เตน โข ปน สมเยน ภิกฺขู นเขนปิ นขํ ฉินฺทนฺติ, มุเขนปิ นขํ ฉินฺทนฺติ, กุฏฺเฏปิ ฆํสนฺติ, องฺคุลิโย ทุกฺขา โหนฺติ.\csromanspacer{} ภควโต เอตมตฺถํ อาโรเจสุํ.\csromanspacer{} อนุชานามิ ภิกฺขเว นขจฺเฉทนนฺติ.\csromanspacer{} สโลหิตํ นขํ ฉินฺทนฺติ, องฺคุลิโย ทุกฺขา โหนฺติ ฯเปฯ.\csromanspacer{} อนุชานามิ ภิกฺขเว มํสปฺปมาเณน นขํ ฉินฺทิตุนฺติ.}

\prose{\csromanfolio{274}เตน โข ปน สมเยน ฉพฺพคฺคิยา ภิกฺขู วีสติมฏฺฐํ\footnote{วีสติมฏฺฏํ (สี)} การาเปนฺติ.\csromanspacer{} มนุสฺสา อุชฺฌายนฺติ ขิยฺยนฺติ วิปาเจนฺติ “เสยฺยถาปิ คิหี กาม โภคิโน”ติ.\csromanspacer{} ภควโต เอตมตฺถํ อาโรเจสุํ.\csromanspacer{} น ภิกฺขเว วีสติมฏฺฐํ การาเปตพฺพํ, โย การาเปยฺย, อาปตฺติ ทุกฺกฏสฺส.\csromanspacer{} อนุชานามิ ภิกฺขเว มลมตฺตํ อปกฑฺฒิตุนฺติ.}

\proseitem{257}{เตน โข ปน สมเยน ภิกฺขูนํ เกสา ทีฆา โหนฺติ.\csromanspacer{} ภควโต เอตมตฺถํ อาโรเจสุํ.\csromanspacer{} อุสฺสหนฺติ ปน ภิกฺขเว ภิกฺขู อญฺญมญฺญํ เกเส โอโรเปตุนฺติ.\csromanspacer{} อุสฺสหนฺติ ภควาติ.\csromanspacer{} อถ โข ภควา เอตสฺมิํ นิทาเน เอตสฺมิํ ปกรเณ ภิกฺขุสํฆํ สนฺนิปาตาเปตฺวา ฯเปฯ ภิกฺขู อามนฺเตสิ “อนุชานามิ ภิกฺขเว ขุรํ ขุรสิลํ ขุรสิปาฏิกํ นมตกํ สพฺพํ ขุรภณฺฑนฺ”ติ.}

\prose{เตน โข ปน สมเยน ฉพฺพคฺคิยา ภิกฺขู มสฺสุํ กปฺปาเปนฺติ ฯเปฯ มสฺสุํ วฑฺฒาเปนฺติ.\csromanspacer{} โคโลมิกํ การาเปนฺติ.\csromanspacer{} จตุรสฺสกํ การาเปนฺติ.\csromanspacer{} ปริมุขํ การาเปนฺติ.\csromanspacer{} อฑฺฒทุกํ\footnote{อฑฺฒุรกํ (สี), อฑฺฒรุกํ (สฺยา)} การาเปนฺติ.\csromanspacer{} ทาฐิกํ ฐเปนฺติ.\csromanspacer{} สมฺพาเธ โลมํ สํหราเปนฺติ.\csromanspacer{} มนุสฺสา อุชฺฌายนฺติ ขิยฺยนฺติ วิปาเจนฺติ “เสยฺยถาปิ คิหี กามโภคิโน”ติ.\csromanspacer{} ภควโต เอตมตฺถํ อาโรเจสุํ.\csromanspacer{} น ภิกฺขเว มสฺสุ กปฺปาเปตพฺพํ ฯเปฯ น มสฺสุ วฑฺฒาเปตพฺพํ.\csromanspacer{} น โคโลมิกํ การาเปตพฺพํ.\csromanspacer{} น จตุรสฺสกํ การาเปตพฺพํ.\csromanspacer{} น ปริมุขํ การาเปตพฺพํ.\csromanspacer{} น อฑฺฒทุกํ การาเปตพฺพํ.\csromanspacer{} น ทาฐิกา ฐเปตพฺพา.\csromanspacer{} น สมฺพาเธ โลมํ สํหราเปตพฺพํ, โย สํหราเปยฺย, อาปตฺติ ทุกฺกฏสฺสาติ.}

\prose{เตน โข ปน สมเยน อญฺญตรสฺส ภิกฺขุโน สมฺพาเธ วโณ โหติ, เภสชฺชํ น สนฺติฏฺฐติ.\csromanspacer{} ภควโต เอตมตฺถํ อาโรเจสุํ.\csromanspacer{} อนุชานามิ ภิกฺขเว อาพาธปฺปจฺจยา สมฺพาเธ โลมํ สํหราเปตุนฺติ.}

\prose{เตน โข ปน สมเยน ฉพฺพคฺคิยา ภิกฺขู กตฺตริกาย เกเส เฉทาเปนฺติ.\csromanspacer{} มนุสฺสา อุชฺฌายนฺติ ขิยฺยนฺติ วิปาเจนฺติ “เสยฺยถาปิ คิหี กามโภคิโน”ติ.\csromanspacer{} ภควโต เอตมตฺถํ อาโรเจสุํ.\csromanspacer{} น ภิกฺขเว กตฺตริกาย เกสา เฉทาเปตพฺพา, โย เฉทาเปยฺย, อาปตฺติ ทุกฺกฏสฺสาติ.}

\prose{\csromanfolio{275}เตน โข ปน สมเยน อญฺญตรสฺส ภิกฺขุโน สีเส วโณ โหติ, น สกฺโกติ ขุเรน เกเส โอโรเปตุํ.\csromanspacer{} ภควโต เอตมตฺถํ อาโรเจสุํ.\csromanspacer{} อนุชานามิ ภิกฺขเว อาพาธปฺปจฺจยา กตฺตริกาย เกเส เฉทาเปตุนฺติ.}

\prose{เตน โข ปน สมเยน ภิกฺขู ทีฆานิ นาสิกาโลมานิ ธาเรนฺติ.\csromanspacer{} มนุสฺสา อุชฺฌายนฺติ ขิยฺยนฺติ วิปาเจนฺติ “เสยฺยถาปิ ปิสาจิลฺลิกา”ติ.\csromanspacer{} ภควโต เอตมตฺถํ อาโรเจสุํ.\csromanspacer{} น ภิกฺขเว ทีฆํ นาสิกาโลมํ ธาเรตพฺพํ, โย ธาเรยฺย, อาปตฺติ ทุกฺกฏสฺสาติ.}

\prose{เตน โข ปน สมเยน ภิกฺขู สกฺขริกายปิ มธุสิตฺถเกนปิ นาสิกาโลมํ คาหาเปนฺติ, นาสิกา ทุกฺขา โหนฺติ.\csromanspacer{} ภควโต เอตมตฺถํ อาโรเจสุํ.\csromanspacer{} อนุชานามิ ภิกฺขเว สณฺฑาสนฺติ.}

\prose{เตน โข ปน สมเยน ฉพฺพคฺคิยา ภิกฺขู ปลิตํ คาหาเปนฺติ.\csromanspacer{} มนุสฺสา อุชฺฌายนฺติ ขิยฺยนฺติ วิปาเจนฺติ “เสยฺยถาปิ คิหี กามโภคิโน”ติ.\csromanspacer{} ภควโต เอตมตฺถํ อาโรเจสุํ.\csromanspacer{} น ภิกฺขเว ปลิตํ คาหาเปตพฺพํ, โย คาหาเปยฺย, อาปตฺติ ทุกฺกฏสฺส.}

\proseitem{276}{เตน โข ปน สมเยน อญฺญตรสฺส ภิกฺขุโน กณฺณคูถเกหิ กณฺณา ถกิตา โหนฺติ.\csromanspacer{} ภควโต เอตมตฺถํ อาโรเจสุํ.\csromanspacer{} อนุชานามิ ภิกฺขเว กณฺณมลหรณินฺติ.}

\prose{เตน โข ปน สมเยน ฉพฺพคฺคิยา ภิกฺขู อุจฺจาวจา กณฺณมลหรณิโย ธาเรนฺติ โสวณฺณมยํ รูปิยมยํ.\csromanspacer{} มนุสฺสา อุชฺฌายนฺติ ขิยฺยนฺติ วิปาเจนฺติ “เสยฺยถาปิ คิหี กามโภคิโน”ติ.\csromanspacer{} ภควโต เอตมตฺถํ อาโรเจสุํ.\csromanspacer{} น ภิกฺขเว อุจฺจาวจา กณฺณมลหรณิโย ธาเรตพฺพา, โย ธาเรยฺย, อาปตฺติ ทุกฺกฏสฺสาติ.\csromanspacer{} อนุชานามิ ภิกฺขเว อฏฺฐิมยํ ทนฺตมยํ วิสาณมยํ นฬมยํ เวฬุมยํ กฏฺฐมยํ ชตุมยํ ผลมยํ โลหมยํ สงฺขนาภิมยนฺติ.}

\proseitem{277}{เตน โข ปน สมเยน ฉพฺพคฺคิยา ภิกฺขู พหุํ โลหภณฺฑํ กํสภณฺฑํ สนฺนิจยํ กโรนฺติ.\csromanspacer{} มนุสฺสา วิหารจาริกํ อาหิณฺฑนฺตา ปสฺสิตฺวา อุชฺฌายนฺติ ขิยฺยนฺติ วิปาเจนฺติ “กถํ หิ นาม สมณา สกฺยปุตฺติยา พหุํ โลหภณฺฑํ กํสภณฺฑํ สนฺนิจยํ กริสฺสนฺติ \csromanfolio{276}เสยฺยถาปิ กํสปตฺถริกา”ติ.\csromanspacer{} ภควโต เอตมตฺถํ อาโรเจสุํ.\csromanspacer{} น ภิกฺขเว พหุํ โลหภณฺฑํ กํสภณฺฑํ สนฺนิจโย กาตพฺโพ, โย กเรยฺย, อาปตฺติ ทุกฺกฏสฺสาติ.}

\prose{เตน โข ปน สมเยน ภิกฺขู อญฺชนิมฺปิ อญฺชนิสลากมฺปิ กณฺณมลหรณิมฺปิ พนฺธนมตฺตมฺปิ กุกฺกุจฺจายนฺติ.\csromanspacer{} ภควโต เอตมตฺถํ อาโรเจสุํ.\csromanspacer{} อนุชานามิ ภิกฺขเว อญฺชนิํ อญฺชนิสลากํ กณฺณมลหรณิํ พนฺธนมตฺตนฺติ.}

\prose{เตน โข ปน สมเยน ฉพฺพคฺคิยา ภิกฺขู สํฆาฏิปลฺลตฺถิกาย นิสีทนฺติ, สํฆาฏิยา ปตฺตา ลุชฺชนฺติ.\csromanspacer{} ภควโต เอตมตฺถํ อาโรเจสุํ.\csromanspacer{} น ภิกฺขเว สํฆาฏิปลฺลตฺถิกาย นิสีทิตพฺพํ, โย นิสีเทยฺย, อาปตฺติ ทุกฺกฏสฺสาติ.}

\prose{เตน โข ปน สมเยน อญฺญตโร ภิกฺขุ คิลาโน โหติ, ตสฺส วินา อาโยเคน\footnote{อาโยคา (สี, สฺยา)} น ผาสุ โหติ.\csromanspacer{} ภควโต เอตมตฺถํ อาโรเจสุํ.\csromanspacer{} อนุชานามิ ภิกฺขเว อาโยคนฺติ.\csromanspacer{} อถ โข ภิกฺขูนํ เอตทโหสิ “กถํ นุ โข อาโยโค กาตพฺโพ”ติ.\csromanspacer{} ภควโต เอตมตฺถํ อาโรเจสุํ.\csromanspacer{} อนุชานามิ ภิกฺขเว ตนฺตกํ เวมํ กวฏํ\footnote{เวมกํ วฏฺฏํ (สี), เวมกํ วฏํ (สฺยา)} สลากํ สพฺพํ ตนฺตภณฺฑกนฺติ.}

\proseitem{278}{เตน โข ปน สมเยน อญฺญตโร ภิกฺขุ อกายพนฺธโน คามํ ปิณฺฑาย ปาวิสิ, ตสฺส รถิกาย อนฺตรวาสโก ปภสฺสิตฺถ.\csromanspacer{} มนุสฺสา อุกฺกุฏฺฐิมกํสุ.\csromanspacer{} โส ภิกฺขุ มงฺกุ อโหสิ.\csromanspacer{} อถ โข โส ภิกฺขุ อารามํ คนฺตา ภิกฺขูนํ เอตมตฺถํ อาโรเจสิ.\csromanspacer{} ภิกฺขู ภควโต เอตมตฺถํ อาโรเจสุํ.\csromanspacer{} น ภิกฺขเว อกายพนฺธเนน คาโม ปวิสิตพฺโพ, โย ปวิเสยฺย, อาปตฺติ ทุกฺกฏสฺส.\csromanspacer{} อนุชานามิ ภิกฺขเว กายพนฺธนนฺติ.}

\prose{เตน โข ปน สมเยน ฉพฺพคฺคิยา ภิกฺขู อุจฺจาวจานิ กายพนฺธนานิ ธาเรนฺติ กลาพุกํ เทฑฺฑุภกํ มุรชํ มทฺทวีณํ.\csromanspacer{} มนุสฺสา อุชฺฌายนฺติ ขิยฺยนฺติ วิปาเจนฺติ “เสยฺยถาปิ คิหี กามโภคิโน”ติ.\csromanspacer{} ภควโต เอตมตฺถํ อาโรเจสุํ.\csromanspacer{} น ภิกฺขเว อุจฺจาวจานิ กายพนฺธนานิ ธาเรตพฺพานิ กลาพุกํ เทฑฺฑุภกํ มุรชํ มทฺทวีณํ, โย ธาเรยฺย, อาปตฺติ ทุกฺกฏสฺส.\csromanspacer{} อนุชานามิ ภิกฺขเว เทฺว กายพนฺธนานิ ปฏฺฏิกํ สูกรนฺตกนฺติ.\csromanspacer{} กายพนฺธนสฺส \csromanfolio{277}ทสา ชีรนฺติ ฯเปฯ.\csromanspacer{} อนุชานามิ ภิกฺขเว มุรชํ มทฺทวีณนฺติ.\csromanspacer{} กายพนฺธนสฺส อนฺโต ชีรติ ฯเปฯ.\csromanspacer{} อนุชานามิ ภิกฺขเว โสภณํ คุณกนฺติ.\csromanspacer{} กายพนฺธนสฺส ปวนนฺโต ชีรติ ฯเปฯ.\csromanspacer{} อนุชานามิ ภิกฺขเว วิธนฺติ\footnote{วีถนฺติ (สี, สฺยา)}.}

\prose{เตน โข ปน สมเยน ฉพฺพคฺคิยา ภิกฺขู อุจฺจาวเจ วิเธ ธาเรนฺติ โสวณฺณมยํ รูปิยมยํ.\csromanspacer{} มนุสฺสา อุชฺฌายนฺติ ขิยฺยนฺติ วิปาเจนฺติ “เสยฺยถาปิ คิหี กามโภคิโน”ติ.\csromanspacer{} ภควโต เอตมตฺถํ อาโรเจสุํ.\csromanspacer{} น ภิกฺขเว อุจฺจาวจา วิธา ธาเรตพฺพา, โย ธาเรยฺย, อาปตฺติ ทุกฺกฏสฺส.\csromanspacer{} อนุชานามิ ภิกฺขเว อฏฺฐิมยํ ฯเปฯ สงฺขนาภิมยํ สุตฺตมยนฺติ.}

\proseitem{279}{เตน โข ปน สมเยน อายสฺมา อานนฺโท ลหุกา สํฆาฏิโย ปารุปิตฺวา คามํ ปิณฺฑาย ปาวิสิ, วาตมณฺฑลิกาย สํฆาฏิโย อุกฺขิปิยิํสุ.\csromanspacer{} อถ โข อายสฺมา อานนฺโท อารามํ คนฺตา ภิกฺขูนํ เอตมตฺถํ อาโรเจสิ.\csromanspacer{} ภิกฺขู ภควโต เอตมตฺถํ อาโรเจสุํ.\csromanspacer{} อนุชานามิ ภิกฺขเว คณฺฐิกํ ปาสกนฺติ.}

\prose{เตน โข ปน สมเยน ฉพฺพคฺคิยา ภิกฺขู อุจฺจาวจา คณฺฐิกาโย ธาเรนฺติ โสวณฺณมยํ รูปิยมยํ.\csromanspacer{} มนุสฺสา อุชฺฌายนฺติ ขิยฺยนฺติ วิปาเจนฺติ “เสยฺยถาปิ คิหี กามโภคิโน”ติ.\csromanspacer{} ภควโต เอตมตฺถํ อาโรเจสุํ.\csromanspacer{} น ภิกฺขเว อุจฺจาวจา คณฺฐิกา ธาเรตพฺพา, โย ธาเรยฺย, อาปตฺติ ทุกฺกฏสฺส.\csromanspacer{} อนุชานามิ ภิกฺขเว อฏฺฐิมยํ ทนฺตมยํ วิสาณมยํ นฬมยํ เวฬุมยํ กฏฺฐมยํ ชตุมยํ ผลมยํ โลหมยํ สงฺขนาภิมยํ สุตฺตมยนฺติ.}

\prose{เตน โข ปน สมเยน ภิกฺขู คณฺฐิกมฺปิ ปาสกมฺปิ จีวเร อปฺเปนฺติ.\csromanspacer{} จีวรํ ชีรติ.\csromanspacer{} ภควโต เอตมตฺถํ อาโรเจสุํ.\csromanspacer{} อนุชานามิ ภิกฺขเว คณฺฐิกผลกํ ปาสกผลกนฺติ.\csromanspacer{} คณฺฐิกผลกมฺปิ ปาสกผลกมฺปิ อนฺเต อปฺเปนฺติ, โกฏฺโฏ\footnote{โกโณ (สี, สฺยา)} วิวริยติ.\csromanspacer{} ภควโต เอตมตฺถํ อาโรเจสุํ.\csromanspacer{} อนุชานามิ ภิกฺขเว คณฺฐิกผลกํ อนฺเต อปฺเปตุํ, ปาสกผลกํ สตฺตงฺคุลํ วา อฏฺฐงฺคุลํ วา โอคาเหตฺวา อปฺเปตุนฺติ.}

\proseitem{280}{เตน โข ปน สมเยน ฉพฺพคฺคิยา ภิกฺขู คิหินิวตฺถํ นิวาเสนฺติ หตฺถิโสณฺฑกํ มจฺฉวาฬกํ จตุกณฺณกํ ตาลวณฺฏกํ สตวลิกํ.\csromanspacer{} มนุสฺสา \csromanfolio{278}อุชฺฌายนฺติ ขิยฺยนฺติ วิปาเจนฺติ “เสยฺยถาปิ คิหี กามโภคิโน”ติ.\csromanspacer{} ภควโต เอตมตฺถํ อาโรเจสุํ.\csromanspacer{} น ภิกฺขเว คิหินิวตฺถํ นิวาเสตพฺพํ หตฺถิโสณฺฑกํ มจฺฉวาฬกํ จตุกณฺณกํ ตาลวณฺฏกํ สตวลิกํ, โย นิวาเสยฺย, อาปตฺติ ทุกฺกฏสฺสาติ.}

\prose{เตน โข ปน สมเยน ฉพฺพคฺคิยา ภิกฺขู คิหิปารุตํ ปารุปนฺติ.\csromanspacer{} มนุสฺสา อุชฺฌายนฺติ ขิยฺยนฺติ วิปาเจนฺติ “เสยฺยถาปิ คิหี กามโภคิโน”ติ.\csromanspacer{} ภควโต เอตมตฺถํ อาโรเจสุํ.\csromanspacer{} น ภิกฺขเว คิหิปารุตํ ปารุปิตพฺพํ, โย ปารุเปยฺย, อาปตฺติ ทุกฺกฏสฺสาติ.}

\prose{เตน โข ปน สมเยน ฉพฺพคฺคิยา ภิกฺขู สํเวลฺลิยํ นิวาเสนฺติ.\csromanspacer{} มนุสฺสา อุชฺฌายนฺติ ขิยฺยนฺติ วิปาเจนฺติ “เสยฺยถาปิ รญฺโญ มุณฺฑวฏฺฏี”ติ.\csromanspacer{} ภควโต เอตมตฺถํ อาโรเจสุํ.\csromanspacer{} น ภิกฺขเว สํเวลฺลิยํ นิวาเสตพฺพํ, โย นิวาเสยฺย, อาปตฺติ ทุกฺกฏสฺสาติ.}

\proseitem{281}{เตน โข ปน สมเยน ฉพฺพคฺคิยา ภิกฺขู อุภโตกาชํ หรนฺติ.\csromanspacer{} มนุสฺสา อุชฺฌายนฺติ ขิยฺยนฺติ วิปาเจนฺติ “เสยฺยถาปิ รญฺโญ มุณฺฑวฏฺฏี”ติ.\csromanspacer{} ภควโต เอตมตฺถํ อาโรเจสุํ.\csromanspacer{} น ภิกฺขเว อุภโตกาชํ หริตพฺพํ, โย หเรยฺย, อาปตฺติ ทุกฺกฏสฺส.\csromanspacer{} อนุชานามิ ภิกฺขเว เอกโตกาชํ อนฺตรากาชํ สีสภารํ ขนฺธภารํ กฏิภารํ โอลมฺพกนฺติ.}

\proseitem{282}{เตน โข ปน สมเยน ภิกฺขู ทนฺตกฏฺฐํ น ขาทนฺติ, มุขํ ทุคฺคนฺธํ โหติ.\csromanspacer{} ภควโต เอตมตฺถํ อาโรเจสุํ.}

\prose{ปญฺจิเม ภิกฺขเว อาทีนวา ทนฺตกฏฺฐสฺส อขาทเน.\csromanspacer{} อจกฺขุสฺสํ, มุขํ ทุคฺคนฺธํ โหติ, รสหรณิโย น วิสุชฺฌนฺติ, ปิตฺตํ เสมฺหํ ภตฺตํ ปริโยนนฺธติ, ภตฺตมสฺส นจฺฉาเทติ.\csromanspacer{} อิเม โข ภิกฺขเว ปญฺจ อาทีนวา ทนฺตกฏฺฐสฺส อขาทเน.}

\prose{ปญฺจิเม ภิกฺขเว อานิสํสา ทนฺตกฏฺฐสฺส ขาทเน.\csromanspacer{} จกฺขุสฺสํ, มุขํ น ทุคฺคนฺธํ โหติ, รสหรณิโย วิสุชฺฌนฺติ, ปิตฺตํ เสมฺหํ ภตฺตํ น ปริโยนนฺธติ, ภตฺตมสฺส ฉาเทติ.\csromanspacer{} อิเม โข ภิกฺขเว ปญฺจ อานิสํสา ทนฺตกฏฺฐสฺส ขาทเน.\csromanspacer{} อนุชานามิ ภิกฺขเว ทนฺตกฏฺฐนฺติ.}

\prose{เตน โข ปน สมเยน ฉพฺพคฺคิยา ภิกฺขู ทีฆานิ ทนฺตกฏฺฐานิ ขาทนฺติ, เตเหว สามเณรํ อาโกเฏนฺติ.\csromanspacer{} ภควโต เอตมตฺถํ อาโรเจสุํ. \csromanfolio{279}น ภิกฺขเว ทีฆํ ทนฺตกฏฺฐํ ขาทิตพฺพํ, โย ขาเทยฺย, อาปตฺติ ทุกฺกฏสฺส.\csromanspacer{} อนุชานามิ ภิกฺขเว อฏฺฐงฺคุลปรมํ ทนฺตกฏฺฐํ, น จ เตน สามเณโร อาโกเฏตพฺโพ, โย อาโกเฏยฺย, อาปตฺติ ทุกฺกฏสฺสาติ.}

\prose{เตน โข ปน สมเยน อญฺญตรสฺส ภิกฺขุโน อติมฏาหกํ ทนฺตกฏฺฐํ ขาทนฺตสฺส กณฺเฐ วิลคฺคํ โหติ.\csromanspacer{} ภควโต เอตมตฺถํ อาโรเจสุํ.\csromanspacer{} น ภิกฺขเว อติมฏาหกํ ทนฺตกฏฺฐํ ขาทิตพฺพํ, โย ขาเทยฺย, อาปตฺติ ทุกฺกฏสฺส.\csromanspacer{} อนุชานามิ ภิกฺขเว จตุรงฺคุลปจฺฉิมํ\footnote{จตุรงฺคุลํ ปจฺฉิมํ (ก)} ทนฺตกฏฺฐนฺติ.}

\proseitem{283}{เตน โข ปน สมเยน ฉพฺพคฺคิยา ภิกฺขู ทายํ อาลิมฺเปนฺติ.\csromanspacer{} มนุสฺสา อุชฺฌายนฺติ ขิยฺยนฺติ วิปาเจนฺติ “เสยฺยถาปิ ทวฑาหกา”ติ.\csromanspacer{} ภควโต เอตมตฺถํ อาโรเจสุํ.\csromanspacer{} น ภิกฺขเว ทาโย อาลิมฺปิตพฺโพ, โย อาลิมฺเปยฺย, อาปตฺติ ทุกฺกฏสฺสาติ.}

\prose{เตน โข ปน สมเยน วิหารา ติณคหนา โหนฺติ, ทวฑาเห ฑยฺหมาเน วิหารา ฑยฺหนฺติ, ภิกฺขู กุกฺกุจฺจายนฺติ ปฏคฺคิํ ทาตุํ ปริตฺตํ กาตุํ.\csromanspacer{} ภควโต เอตมตฺถํ อาโรเจสุํ.\csromanspacer{} อนุชานามิ ภิกฺขเว ทวฑาเห ฑยฺหมาเน ปฏคฺคิํ ทาตุํ ปริตฺตํ กาตุนฺติ.}

\proseitem{284}{เตน โข ปน สมเยน ฉพฺพคฺคิยา ภิกฺขู รุกฺขํ อภิรุหนฺติ, รุกฺขา รุกฺขํ สงฺกมนฺติ.\csromanspacer{} มนุสฺสา อุชฺฌายนฺติ ขิยฺยนฺติ วิปาเจนฺติ “เสยฺยถาปิ มกฺกฏา”ติ.\csromanspacer{} ภควโต เอตมตฺถํ อาโรเจสุํ.\csromanspacer{} น ภิกฺขเว รุกฺโข อภิรุหิตพฺโพ, โย อภิรุเหยฺย, อาปตฺติ ทุกฺกฏสฺสาติ.}

\prose{เตน โข ปน สมเยน อญฺญตรสฺส ภิกฺขุโน โกสเลสุ ชนปเท สาวตฺถิํ คจฺฉนฺตสฺส อนฺตรามคฺเค หตฺถี ปริยุฏฺฐาติ.\csromanspacer{} อถ โข โส ภิกฺขุ รุกฺขมูลํ อุปธาวิตฺวา กุกฺกุจฺจายนฺโต รุกฺขํ น อภิรุหิ, โส หตฺถี อญฺเญน อคมาสิ.\csromanspacer{} อถ โข โส ภิกฺขุ สาวตฺถิํ คนฺตฺวา ภิกฺขูนํ เอตมตฺถํ อาโรเจสิ.\csromanspacer{} ภิกฺขู ภควโต เอตมตฺถํ อาโรเจสุํ.\csromanspacer{} อนุชานามิ ภิกฺขเว สติ กรณีเย โปริสํ รุกฺขํ อภิรุหิตุํ, อาปทาสุ ยาวทตฺถนฺติ.}

\proseitem{285}{\csromanfolio{280}เตน โข ปน สมเยน ยเมฬเกกุฏา นาม\footnote{ยเมฬุเตกุลา นาม (สี), เมฏฺฐโกกุฏฺฐา นาม (สฺยา)} ภิกฺขู เทฺว ภาติกา โหนฺติ พฺราหฺมณชาติกา กลฺยาณวาจา กลฺยาณวากฺกรณา, เต เยน ภควา เตนุปสงฺกมิํสุ, อุปสงฺกมิตฺวา ภควนฺตํ อภิวาเทตฺวา เอกมนฺตํ นิสีทิํสุ, เอกมนฺตํ นิสินฺนา โข เต ภิกฺขู ภควนฺตํ เอตทโวจุํ “เอตรหิ ภนฺเต ภิกฺขู นานานามา นานาโคตฺตา นานาชจฺจา นานากุลา ปพฺพชิตา, เต สกาย นิรุตฺติยา พุทฺธวจนํ ทูเสนฺติ, หนฺท มยํ ภนฺเต พุทฺธวจนํ ฉนฺทโส อาโรเปมา”ติ.\csromanspacer{} วิครหิ พุทฺโธ ภควา ฯเปฯ.\csromanspacer{} กถํ หิ นาม ตุเมฺห โมฆปุริสา เอวํ วกฺขถ “หนฺท มยํ ภนฺเต พุทฺธวจนํ ฉนฺทโส อาโรเปมา”ติ.\csromanspacer{} เนตํ โมฆปุริสา อปฺปสนฺนานํ วา ปสาทย ฯเปฯ วิครหิตฺวา ฯเปฯ ธมฺมิํ กถํ กตฺวา ภิกฺขู อามนฺเตสิ “น ภิกฺขเว พุทฺธวจนํ ฉนฺทโส อาโรเปตพฺพํ, โย อาโรเปยฺย, อาปตฺติ ทุกฺกฏสฺส.\csromanspacer{} อนุชานามิ ภิกฺขเว สกาย นิรุตฺติยา พุทฺธวจนํ ปริยาปุณิตุนฺ”ติ.}

\proseitem{286}{เตน โข ปน สมเยน ฉพฺพคฺคิยา ภิกฺขู โลกายตํ ปริยาปุณนฺติ.\csromanspacer{} มนุสฺสา อุชฺฌายนฺติ ขิยฺยนฺติ วิปาเจนฺติ “เสยฺยถาปิ คิหี กามโภคิโน”ติ.\csromanspacer{} อสฺโสสุํ โข ภิกฺขู เตสํ มนุสฺสานํ อุชฺฌายนฺตานํ ขิยฺยนฺตานํ วิปาเจนฺตานํ.\csromanspacer{} อถ โข เต ภิกฺขู ภควโต เอตมตฺถํ อาโรเจสุํ.\csromanspacer{} อปิ นุ โข ภิกฺขเว โลกายเต สารทสฺสาวี อิมสฺมิํ ธมฺมวินเย วุทฺธิํ วิรุฬฺหิํ เวปุลฺลํ อาปชฺเชยฺยาติ.\csromanspacer{} โนเหตํ ภนฺเต.\csromanspacer{} อิมสฺมิํ วา ปน ธมฺมวินเย สารทสฺสาวี โลกายตํ ปริยาปุเณยฺยาติ.\csromanspacer{} โนเหตํ ภนฺเต.\csromanspacer{} น ภิกฺขเว โลกายตํ ปริยาปุณิตพฺพํ, โย ปริยาปุเณยฺย, อาปตฺติ ทุกฺกฏสฺสาติ.}

\prose{เตน โข ปน สมเยน ฉพฺพคฺคิยา ภิกฺขู โลกายตํ วาเจนฺติ.\csromanspacer{} มนุสฺสา อุชฺฌายนฺติ ขิยฺยนฺติ วิปาเจนฺติ “เสยฺยถาปิ คิหี กามโภคิโน”ติ.\csromanspacer{} ภควโต เอตมตฺถํ อาโรเจสุํ.\csromanspacer{} น ภิกฺขเว โลกายตํ วาเจตพฺพํ, โย วาเจยฺย, อาปตฺติ ทุกฺกฏสฺสาติ.}

\proseitem{287}{เตน โข ปน สมเยน ฉพฺพคฺคิยา ภิกฺขู ติรจฺฉานวิชฺชํ ปริยาปุณนฺติ.\csromanspacer{} มนุสฺสา อุชฺฌายนฺติ ขิยฺยนฺติ วิปาเจนฺติ “เสยฺยถาปิ คิหี \csromanfolio{281}กามโภคิโน”ติ.\csromanspacer{} ภควโต เอตมตฺถํ อาโรเจสุํ.\csromanspacer{} น ภิกฺขเว ติรจฺฉานวิชฺชา ปริยาปุณิตพฺพา, โย ปริยาปุเณยฺย, อาปตฺติ ทุกฺกฏสฺสาติ.}

\prose{เตน โข ปน สมเยน ฉพฺพคฺคิยา ภิกฺขู ติรจฺฉานวิชฺชํ วาเจนฺติ.\csromanspacer{} มนุสฺสา อุชฺฌายนฺติ ขิยฺยนฺติ วิปาเจนฺติ “เสยฺยถาปิ คิหี กามโภคิโน”ติ.\csromanspacer{} ภควโต เอตมตฺถํ อาโรเจสุํ.\csromanspacer{} น ภิกฺขเว ติรจฺฉานวิชฺชา วาเจตพฺพา, โย วาเจยฺย, อาปตฺติ ทุกฺกฏสฺสาติ.}

\proseitem{288}{เตน โข ปน สมเยน ภควา มหติยา ปริสาย ปริวุโต ธมฺมํ เทเสนฺโต ขิปิ.\csromanspacer{} ภิกฺขู “ชีวตุ ภนฺเต ภควา, ชีวตุ สุคโต”ติ อุจฺจาสทฺทํ มหาสทฺทํ อกํสุ, เตน สทฺเทน ธมฺมกถา อนฺตรา อโหสิ.\csromanspacer{} อถ โข ภควา ภิกฺขู อามนฺเตสิ “อปิ นุ โข ภิกฺขเว ขิปิเต ‘ชีวา’ติ วุตฺโต\footnote{วุตฺเต (ก)} ตปฺปจฺจยา ชีเวยฺย วา มเรยฺย วา”ติ.\csromanspacer{} โนเหตํ ภนฺเต.\csromanspacer{} น ภิกฺขเว ขิปิเต “ชีวา”ติ วตฺตพฺโพ, โย วเทยฺย, อาปตฺติ ทุกฺกฏสฺสาติ.}

\prose{เตน โข ปน สมเยน มนุสฺสา ภิกฺขูนํ ขิปิเต “ชีวถ ภนฺเต”ติ วทนฺติ, ภิกฺขู กุกฺกุจฺจายนฺตา นาลปนฺติ.\csromanspacer{} มนุสฺสา อุชฺฌายนฺติ ขิยฺยนฺติ วิปาเจนฺติ “กถํ หิ นาม สมณา สกฺยปุตฺติยา ‘ชีวถ ภนฺเต’ติ วุจฺจมานา นาลปิสฺสนฺตี”ติ.\csromanspacer{} ภควโต เอตมตฺถํ อาโรเจสุํ.\csromanspacer{} คิหี ภิกฺขเว มงฺคลิกา.\csromanspacer{} อนุชานามิ ภิกฺขเว คิหีนํ “ชีวถ ภนฺเต”ติ วุจฺจมาเนน “จิรํ ชีวา”ติ วตฺตุนฺติ.}

\proseitem{289}{เตน โข ปน สมเยน ภควา มหติยา ปริสาย ปริวุโต ธมฺมํ เทเสนฺโต นิสินฺโน โหติ, อญฺญตเรน ภิกฺขุนา ลสุณํ ขายิตํ โหติ.\csromanspacer{} โส “มา ภิกฺขู พฺยาพาธิํสู”ติ เอกมนฺตํ นิสีทิ.\csromanspacer{} อทฺทสา โข ภควา ตํ ภิกฺขุํ เอกมนฺตํ นิสินฺนํ, ทิสฺวาน ภิกฺขู อามนฺเตสิ “กิํ นุ โข โส ภิกฺขเว ภิกฺขุ เอกมนฺตํ นิสินฺโน”ติ.\csromanspacer{} เอเตน ภนฺเต ภิกฺขุนา ลสุณํ ขายิตํ, โส “มา ภิกฺขู พฺยาพาธิํสู”ติ เอกมนฺตํ นิสินฺโนติ.\csromanspacer{} อปิ นุ โข ภิกฺขเว\footnote{ภิกฺขเว ภิกฺขุนา (สฺยา)} ตํ ขาทิตพฺพํ, ยํ ขาทิตฺวา เอวรูปาย ธมฺมกถาย ปริพาหิโย อสฺสาติ.\csromanspacer{} โนเหตํ ภนฺเต.\csromanspacer{} น ภิกฺขเว ลสุณํ ขาทิตพฺพํ, โย ขาเทยฺย, อาปตฺติ ทุกฺกฏสฺสาติ.}

\prose{\csromanfolio{282}เตน โข ปน สมเยน อายสฺมโต สาริปุตฺตสฺส อุทรวาตาพาโธ โหติ.\csromanspacer{} อถ โข อายสฺมา มหาโมคฺคลฺลาโน เยนายสฺมา สาริปุตฺโต เตนุปสงฺกมิ, อุปสงฺกมิตฺวา อายสฺมนฺตํ สาริปุตฺตํ เอตทโวจ “ปุพฺเพ เต อาวุโส สาริปุตฺต อุทรวาตาพาโธ เกน ผาสุ โหตี”ติ.\csromanspacer{} ลสุเณน เม อาวุโสติ\footnote{อถ โข อายสฺมา สาริปุตฺโต อายสฺมนฺตํ มหาโมคฺคลฺลานํ เอตทโวจ “ปุพฺเพ โข เม อาวุโส โมคฺคลฺลาน อุทรวาตาพาโธ ลสุเณสุ ผาสุ โหตี”ติ (ก)}.\csromanspacer{} ภควโต เอตมตฺถํ อาโรเจสุํ.\csromanspacer{} อนุชานามิ ภิกฺขเว อาพาธปฺปจฺจยา ลสุณํ ขาทิตุนฺติ.}

\proseitem{290}{เตน โข ปน สมเยน ภิกฺขู อาราเม ตหํ ตหํ ปสฺสาวํ กโรนฺติ, อาราโม ทุสฺสติ.\csromanspacer{} ภควโต เอตมตฺถํ อาโรเจสุํ.\csromanspacer{} อนุชานามิ ภิกฺขเว เอกมนฺตํ ปสฺสาวํ กาตุนฺติ.\csromanspacer{} อาราโม ทุคฺคนฺโธ โหติ ฯเปฯ.\csromanspacer{} อนุชานามิ ภิกฺขเว ปสฺสาวกุมฺภินฺติ.\csromanspacer{} ทุกฺขํ นิสินฺนา ปสฺสาวํ กโรนฺติ ฯเปฯ.\csromanspacer{} อนุชานามิ ภิกฺขเว ปสฺสาวปาทุกนฺติ.\csromanspacer{} ปสฺสาวปาทุกา ปากฏา โหนฺติ, ภิกฺขู หิริยนฺติ ปสฺสาวํ กาตุํ ฯเปฯ.\csromanspacer{} อนุชานามิ ภิกฺขเว ปริกฺขิปิตุํ ตโย ปากาเร อิฏฺฐกาปาการํ สิลาปาการํ ทารุปาการนฺติ.\csromanspacer{} ปสฺสาวกุมฺภี อปารุตา ทุคฺคนฺธา โหติ ฯเปฯ.\csromanspacer{} อนุชานามิ ภิกฺขเว อปิธานนฺติ.}

\proseitem{291}{เตน โข ปน สมเยน ภิกฺขู อาราเม ตหํ ตหํ วจฺจํ กโรนฺติ, อาราโม ทุสฺสติ.\csromanspacer{} ภควโต เอตมตฺถํ อาโรเจสุํ.\csromanspacer{} อนุชานามิ ภิกฺขเว เอกมนฺตํ วจฺจํ กาตุนฺติ.\csromanspacer{} อาราโม ทุคฺคนฺโธ โหติ ฯเปฯ.\csromanspacer{} อนุชานามิ ภิกฺขเว วจฺจกูปนฺติ.\csromanspacer{} วจฺจกูปสฺส กูลํ ลุชฺชติ ฯเปฯ.\csromanspacer{} อนุชานามิ ภิกฺขเว จินิตุํ ตโย จเย อิฏฺฐกาจยํ สิลาจยํ ทารุจยนฺติ.\csromanspacer{} วจฺจกูโป นีจวตฺถุโก โหติ, อุทเกน โอตฺถริยฺยติ ฯเปฯ.\csromanspacer{} อนุชานามิ ภิกฺขเว อุจฺจวตฺถุกํ กาตุนฺติ.\csromanspacer{} จโย ปริปตติ ฯเปฯ.\csromanspacer{} อนุชานามิ ภิกฺขเว จินิตุํ ตโย จเย อิฏฺฐกาจยํ สิลาจยํ ทารุจยนฺติ.\csromanspacer{} อาโรหนฺตา วิหญฺญนฺติ ฯเปฯ.\csromanspacer{} อนุชานามิ ภิกฺขเว ตโย โสปาเน อิฏฺฐกาโสปานํ สิลาโสปานํ ทารุโสปานนฺติ.\csromanspacer{} อาโรหนฺตา ปริปตนฺติ ฯเปฯ.\csromanspacer{} อนุชานามิ ภิกฺขเว อาลมฺพนพาหนฺติ.\csromanspacer{} อนฺเต นิสินฺนา วจฺจํ กโรนฺตา ปริปตนฺติ ฯเปฯ.\csromanspacer{} อนุชานามิ ภิกฺขเว สนฺถริตฺวา มชฺเฌ ฉิทฺทํ กตฺวา วจฺจํ กาตุนฺติ.\csromanspacer{} ทุกฺขํ นิสินฺนา วจฺจํ กโรนฺติ ฯเปฯ.\csromanspacer{} อนุชานามิ ภิกฺขเว วจฺจปาทุกนฺติ.\csromanspacer{} พหิทฺธา \csromanfolio{283}ปสฺสาวํ กโรนฺติ ฯเปฯ.\csromanspacer{} อนุชานามิ ภิกฺขเว ปสฺสาวโทณิกนฺติ.\csromanspacer{} อวเลขนกฏฺฐํ น โหติ ฯเปฯ.\csromanspacer{} อนุชานามิ ภิกฺขเว อวเลขนกฏฺฐนฺติ.\csromanspacer{} อวเลขนปิฐโร น โหติ ฯเปฯ.\csromanspacer{} อนุชานามิ ภิกฺขเว อวเลขนปิฐรนฺติ.\csromanspacer{} วจฺจกูโป อปารุโต ทุคฺคนฺโธ โหติ ฯเปฯ.\csromanspacer{} อนุชานามิ ภิกฺขเว อปิธานนฺติ.\csromanspacer{} อชฺโฌกาเส วจฺจํ กโรนฺตา สีเตนปิ อุเณฺหนปิ กิลมนฺติ ฯเปฯ.\csromanspacer{} อนุชานามิ ภิกฺขเว วจฺจกุฏินฺติ.\csromanspacer{} วจฺจกุฏิยา กวาฏํ น โหติ ฯเปฯ.\csromanspacer{} อนุชานามิ ภิกฺขเว กวาฏํ ปิฏฺฐสํฆาฏํ อุทุกฺขลิกํ อุตฺตรปาสกํ อคฺคฬวฏฺฏิํ กปิสีสกํ สูจิกํ ฆฏิกํ ตาฬจฺฉิทฺทํ อาวิญฺฉนจฺฉิทฺทํ อาวิญฺฉนรชฺชุนฺติ.\csromanspacer{} วจฺจกุฏิยา ติณจุณฺณํ ปริปตติ ฯเปฯ.\csromanspacer{} อนุชานามิ ภิกฺขเว โอคุมฺเผตฺวา อุลฺลิตฺตาวลิตฺตํ กาตุํ เสตวณฺณํ กาฬวณฺณํ เครุกปริกมฺมํ มาลากมฺมํ ลตากมฺมํ มกรทนฺตกํ ปญฺจปฏิกํ จีวรวํสํ จีวรรชฺชุนฺติ.}

\proseitem{292}{เตน โข ปน สมเยน อญฺญตโร ภิกฺขุ ชราทุพฺพโล วจฺจํ กตฺวา วุฏฺฐหนฺโต ปริปตติ.\csromanspacer{} ภควโต เอตมตฺถํ อาโรเจสุํ.\csromanspacer{} อนุชานามิ ภิกฺขเว โอลมฺพกนฺติ.\csromanspacer{} วจฺจกุฏิ อปริกฺขิตฺตา โหติ ฯเปฯ.\csromanspacer{} อนุชานามิ ภิกฺขเว ปริกฺขิปิตุํ ตโย ปากาเร อิฏฺฐกาปาการํ สิลาปาการํ ทารุปาการนฺติ.\csromanspacer{} โกฏฺฐโก น โหติ ฯเปฯ. (อนุชานามิ ภิกฺขเว โกฏฺฐกนฺติ.\csromanspacer{} โกฏฺฐกสฺส กวาฏํ น โหติ.)\footnote{(อนุชานามิ ภิกฺขเว โกฏฺฐกนฺติ. โกฏฺฐโก นีจวตฺถุโก โหติ ฯเปฯ. อนุชานามิ ภิกฺขเว อุจฺจวตฺถุกํ กาตุนฺติ. จโย ปริปตติ ฯเปฯ. อนุชานามิ ภิกฺขเว จินิตุํ ตโย จเย อิฏฺฐกาจยํ สิลาจยํ ทารุจยนฺติ. อาโรหนฺตา วิหญฺญนฺติ ฯเปฯ. อนุชานามิ ภิกฺขเว ตโย โสปาเน อิฏฺฐกาโสปานํ สิลาโสปานํ ทารุโสปานนฺติ. อาโรหนฺตา ปริปตนฺติ ฯเปฯ. อนุชานามิ ภิกฺขเว อาลมฺพนพาหนฺติ. โกฏฺฐกสฺส กวาฏํ น โหติ.) (สฺยา, กํ)} ฯเปฯ.\csromanspacer{} อนุชานามิ ภิกฺขเว กวาฏํ ปิฏฺฐสํฆาฏํ อุทุกฺขลิกํ อุตฺตรปาสกํ อคฺคฬวฏฺฏิํ กปิสีสกํ สูจิกํ ฆฏิกํ ตาฬจฺฉิทฺทํ อาวิญฺฉนจฺฉิทฺทํ อาวิญฺฉนรชฺชุนฺติ.\csromanspacer{} โกฏฺฐเก ติณจุณฺณํ ปริปตติ ฯเปฯ.\csromanspacer{} อนุชานามิ ภิกฺขเว โอคุมฺเผตฺวา อุลฺลิตฺตาวลิตฺตํ กาตุํ เสตวณฺณํ กาฬวณฺณํ เครุกปริกมฺมํ มาลากมฺมํ ลตากมฺมํ มกรทนฺตกํ ปญฺจปฏิกนฺติ.\csromanspacer{} ปริเวณํ จิกฺขลฺลํ โหติ ฯเปฯ.\csromanspacer{} อนุชานามิ ภิกฺขเว มรุมฺพํ ปกิริตุนฺติ.\csromanspacer{} น ปริยาปุณนฺติ ฯเปฯ.\csromanspacer{} อนุชานามิ ภิกฺขเว ปทรสิลํ\footnote{ปฏฺฏสิลํ (ก)} นิกฺขิปิตุนฺติ.\csromanspacer{} อุทกํ สนฺติฏฺฐติ ฯเปฯ.\csromanspacer{} อนุชานามิ ภิกฺขเว อุทกนิทฺธมนนฺติ.\csromanspacer{} อาจมนกุมฺภี น โหติ ฯเปฯ.\csromanspacer{} อนุชานามิ ภิกฺขเว อาจมนกุมฺภินฺติ. \csromanfolio{284}อาจมนสราวโก น โหติ ฯเปฯ.\csromanspacer{} อนุชานามิ ภิกฺขเว อาจมนสราวกนฺติ.\csromanspacer{} ทุกฺขํ นิสินฺนา อาจเมนฺติ ฯเปฯ.\csromanspacer{} อนุชานามิ ภิกฺขเว อาจมน ปาทุกนฺติ.\csromanspacer{} อาจมนปาทุกา ปากฏา โหนฺติ, ภิกฺขู หิริยนฺติ อาจเมตุํ ฯเปฯ.\csromanspacer{} อนุชานามิ ภิกฺขเว ปริกฺขิปิตุํ ตโย ปากาเร อิฏฺฐกาปาการํ สิลาปาการํ ทารุปาการนฺติ.\csromanspacer{} อาจมนกุมฺภี อปารุตา โหติ, ติณจุณฺเณหิปิ ปํสุเกหิปิ โอกิริยฺยติ ฯเปฯ.\csromanspacer{} อนุชานามิ ภิกฺขเว อปิธานนฺติ.}

\csromanmatikaanchor{mtk.128}
\csromantocmarkref{section}{อุทฺทานคาถา}{mtk.128}
\bookmark[dest={mtk.128},level=1]{อุทฺทานคาถา}
\proseitem{293}{เตน โข ปน สมเยน ฉพฺพคฺคิยา ภิกฺขู เอวรูปํ อนาจารํ อาจรนฺติ, มาลาวจฺฉํ โรเปนฺติปิ โรปาเปนฺติปิ, สิญฺจนฺติปิ สิญฺจาเปนฺติปิ, โอจินนฺติปิ โอจินาเปนฺติปิ, คนฺเถนฺติปิ คนฺถาเปนฺติปิ, เอกโตวณฺฏิกมาลํ กโรนฺติปิ การาเปนฺติปิ, อุภโตวณฺฏิกมาลํ กโรนฺติปิ การาเปนฺติปิ, มญฺชริกํ กโรนฺติปิ การาเปนฺติปิ, วิธูติกํ กโรนฺติปิ การาเปนฺติปิ, วฏํสกํ กโรนฺติปิ การาเปนฺติปิ, อาเวฬํ กโรนฺติปิ การาเปนฺติปิ, อุรจฺฉทํ กโรนฺติปิ การาเปนฺติปิ.\csromanspacer{} เต กุลิตฺถีนํ กุลธีตานํ กุลกุมารีนํ กุลสุณฺหานํ กุลทาสีนํ เอกโตวณฺฏิกมาลํ หรนฺติปิ หราเปนฺติปิ, อุภโตวณฺฏิกมาลํ หรนฺติปิ หราเปนฺติปิ, มญฺชริกํ หรนฺติปิ หราเปนฺติปิ, วิธูติกํ หรนฺติปิ หราเปนฺติปิ, วฏํสกํ หรนฺติปิ หราเปนฺติปิ, อาเวฬํ หรนฺติปิ หาราเปนฺติปิ, อุรจฺฉทํ หรนฺติปิ หราเปนฺติปิ.\csromanspacer{} เต กุลิตฺถีหิ กุลธีตาหิ กุลกุมารีหิ กุลสุณฺหาหิ กุลทาสีหิ สทฺธิํ เอกภาชเนปิ ภุญฺชนฺติ, เอกถาลเกปิ ปิวนฺติ, เอกาสเนปิ นิสีทนฺติ, เอกมญฺเจปิ ตุวฏฺเฏนฺติ, เอกตฺถรณาปิ ตุวฏฺเฏนฺติ, เอกปาวุรณาปิ ตุวฏฺเฏนฺติ, เอกตฺถรณปาวุรณาปิ ตุวฏฺเฏนฺติ, วิกาเลปิ ภุญฺชนฺติ, มชฺชมฺปิ ปิวนฺติ, มาลาคนฺธวิเลปนมฺปิ ธาเรนฺติ, นจฺจนฺติปิ คายนฺติปิ วาเทนฺติปิ ลาเสนฺติปิ, นจฺจนฺติยาปิ นจฺจนฺติ นจฺจนฺติยาปิ คายนฺติ นจฺจนฺติยาปิ วาเทนฺติ นจฺจนฺติยาปิ ลาเสนฺติ ฯเปฯ.\csromanspacer{} ลาเสนฺติยาปิ นจฺจนฺติ ลาเสนฺติยาปิ คายนฺติ ลาเสนฺติยาปิ วาเทนฺติ ลาเสนฺติยาปิ ลาเสนฺติ, อฏฺฐปเทปิ กีฬนฺติ, ทสปเทปิ กีฬนฺติ, อากาเสปิ กีฬนฺติ, ปริหารปเถปิ กีฬนฺติ, สนฺติกายปิ กีฬนฺติ, ขลิกายปิ กีฬนฺติ, ฆฏิกายปิ กีฬนฺติ, สลากหตฺเถนปิ กีฬนฺติ, อกฺเขนปิ กีฬนฺติ, ปงฺคจีเรนปิ กีฬนฺติ, วงฺกเกนปิ กีฬนฺติ, โมกฺขจิกายปิ กีฬนฺติ, จิงฺคุลเกนปิ กีฬนฺติ, ปตฺตาฬฺหเกนปิ กีฬนฺติ, รถเกนปิ กีฬนฺติ, ธนุเกนปิ กีฬนฺติ, อกฺขริกายปิ กีฬนฺติ, มเนสิกายปิ กีฬนฺติ, ยถาวชฺเชนปิ กีฬนฺติ, หตฺถิสฺมิมฺปิ สิกฺขนฺติ, อสฺสสฺมิมฺปิ สิกฺขนฺติ, รถสฺมิมฺปิ สิกฺขนฺติ, \csromanfolio{285}ธนุสฺมิมฺปิ สิกฺขนฺติ, ถรุสฺมิมฺปิ สิกฺขนฺติ, หตฺถิสฺสปิ ปุรโต ธาวนฺติ, อสฺสสฺสปิ ปุรโต ธาวนฺติ, รถสฺสปิ ปุรโต ธาวนฺติปิ อาธาวนฺติปิ, อุสฺเสเฬนฺติปิ, อปฺโผเฏนฺติปิ, นิพฺพุชฺฌนฺติปิ, มุฏฺฐีหิปิ ยุชฺฌนฺติ, รงฺคมชฺเฌปิ สํฆาฏิํ ปตฺถริตฺวา นจฺจกิํ เอวํ วทนฺติ “อิธ ภคินิ นจฺจสฺสู”ติ, นลาฏิกมฺปิ เทนฺติ, วิวิธมฺปิ อนาจารํ อาจรนฺติ.\csromanspacer{} ภควโต เอตมตฺถํ อาโรเจสุํ ฯเปฯ.\csromanspacer{} น ภิกฺขเว วิวิธํ อนาจารํ อาจริตพฺพํ, โย อาจเรยฺย, ยถาธมฺโม กาเรตพฺโพติ.}

\prose{เตน โข ปน สมเยน อายสฺมนฺเต อุรุเวลกสฺสเป ปพฺพชิเต สํฆสฺส พหุํ โลหภณฺฑํ ทารุภณฺฑํ มตฺติกาภณฺฑํ อุปฺปนฺนํ โหติ.\csromanspacer{} อถ โข ภิกฺขูนํ เอตทโหสิ “กิํ นุ โข ภควตา โลหภณฺฑํ อนุญฺญาตํ, กิํ อนนุญฺญาตํ, กิํ ทารุภณฺฑํ อนุญฺญาตํ, กิํ อนนุญฺญาตํ, กิํ มตฺติกาภณฺฑํ อนุญฺญาตํ, กิํ อนนุญฺญาตนฺ”ติ.\csromanspacer{} ภควโต เอตมตฺถํ อาโรเจสุํ.\csromanspacer{} อถ โข ภควา เอตสฺมิํ นิทาเน เอตสฺมิํ ปกรเณ ธมฺมิํ กถํ กตฺวา ภิกฺขู อามนฺเตสิ “อนุชานามิ ภิกฺขเว ฐเปตฺวา ปหรณิํ สพฺพํ โลหภณฺฑํ, ฐเปตฺวา อาสนฺทิํ ปลฺลงฺกํ ทารุปตฺตํ ทารุปาทุกํ สพฺพํ ทารุภณฺฑํ, ฐเปตฺวา กตกญฺจ กุมฺภการิกญฺจ สพฺพํ มตฺติกาภณฺฑนฺ”ติ.}

\vspace{\tipitakaparskip}
\csromancenter{ขุทฺทกวตฺถุกฺขนฺธโก ปญฺจโม.}
\titlehead{\textbf{ตสฺสุทฺทานํ}}
\csromangathameasure{รุกฺเข ถมฺเภ จ กุฏฺเฏ จ, อฏฺฏาเน คนฺธสุตฺติยา. \\ วิคยฺห มลฺลโก กจฺฉุ, ชรา จ ปุถุปาณิกา. \\ วลฺลิกาปิ จ ปามงฺโค, กณฺฐสุตฺตํ นธารเย. \\ กฏิ โอวฏฺฏิ กายุรํ, หตฺถาภรณมุทฺทิกา. \\ ทีเฆ โกจฺเฉ ผเณ หตฺเถ, สิตฺถา อุทกเตลเก. \\ อาทาสุทปตฺตวณา, อาเลโปมฺมทฺทจุณฺณนา. \\ ลญฺเฉนฺติ องฺคราคญฺจ, มุขราคํ ตทูภยํ. \\ จกฺขุโรคํ คิรคฺคญฺจ, อายตํ สรพาหิรํ. \\ อมฺพเปสิสกเลหิ, อหิจฺฉินฺทิ จ จนฺทนํ. \\ อุจฺจาวจา ปตฺตมูลา, สุวณฺโณ พหลา วลี. \\ จิตฺรา ทุสฺสติ ทุคฺคนฺโธ, อุเณฺห ภิชฺชิํสุ มิฑฺฒิยา. \\ ปริภณฺฑํ ติณํ โจฬํ, มาฬํ กุณฺโฑลิกาย จ. \\ ถวิกา อํสพทฺธญฺจ, ตถา พนฺธนสุตฺตกา. \\ ขิเล มญฺเจ จ ปีเฐ จ, องฺเก ฉตฺเต ปณามนา. \\ ตุมฺพฆฏิฉวสีสํ, จลกานิ ปฏิคฺคโห. \\ วิปฺผาลิทณฺฑโสวณฺณํ, ปตฺเต เปสิ จ นาฬิกา. \\ กิณฺณสตฺตุ สริตญฺจ, มธุสิตฺถํ สิปาฏิกํ. \\ วิกณฺณํ พนฺธิวิสมํ, ฉมาชิรปโหติ จ. \\ กฬิมฺภํ โมฆสุตฺตญฺจ, อโธตลฺลํ อุปาหนา. \\ องฺคุเล ปฏิคฺคหญฺจ, วิตฺถกํ ถวิกพทฺธกา. \\ อชฺโฌกาเส นีจวตฺถุ, จโย จาปิ วิหญฺญเร. \\ ปริปติ ติณจุณฺณํ, อุลฺลิตฺต-อวลิตฺตกํ. \\ เสตํ กาฬกวณฺณญฺจ, ปริกมฺมญฺจ เครุกํ. \\ มาลากมฺมํ ลตากมฺมํ, มกรทนฺตกปาฏิกํ. \\ จีวรวํสํ รชฺชุญฺจ, อนุญฺญาสิ วินายโก. \\ อุชฺฌิตฺวา ปกฺกมนฺติ จ, กถินํ ปริภิชฺชติ. \\ วินิเวฐิยติ กุฏฺเฏปิ, ปตฺเตนาทาย คจฺฉเร. \\ ถวิกา พนฺธสุตฺตญฺจ, พนฺธิตฺวา จ อุปาหนา. \\ อุปาหนตฺถวิกญฺจ, อํสพทฺธญฺจ สุตฺตกํ. \\ อุทกากปฺปิยํ มคฺเค, ปริสฺสาวนโจฬกํ. \\ ธมฺมกรณํ เทฺว ภิกฺขู, เวสาลิํ อคมา มุนิ. \\ ทณฺฑํ โอตฺถรกํ ตตฺถ, อนุญฺญาสิ ปริสฺสาวนํ. \\ มกเสหิ ปณีเตน, พหฺวาพาธา จ ชีวโก. \\ จงฺกมนชนฺตาฆรํ, วิสเม นีจวตฺถุโก. \\ ตโย จเย วิหญฺญนฺติ, โสปานาลมฺพเวทิตํ. \\ อชฺโฌกาเส ติณจุณฺณํ, อุลฺลิตฺต-อวลิตฺตกํ. \\ เสตกํ กาฬวณฺณญฺจ, ปริกมฺมญฺจ เครุกํ. \\ มาลากมฺมํ ลตากมฺมํ, มกรทนฺตกปาฏิกํ. \\ วํสํ จีวรรชฺชุญฺจ, อุจฺจํ จ วตฺถุกํ กเร. \\ จโย โสปานพาหญฺจ, กวาฏํ ปิฏฺฐสํฆาฏํ. \\ อุทุกฺขลุตฺตรปาสกํ, วฏฺฏิญฺจ กปิสีสกํ. \\ สูจิฆฏิตาฬจฺฉิทฺทํ, อาวิญฺฉนญฺจ รชฺชุกํ. \\ มณฺฑลํ ธูมเนตฺตญฺจ, มชฺเฌ จ มุขมตฺติกา. \\ โทณิ ทุคฺคนฺธา ทหติ, อุทกฏฺฐานํ สราวกํ. \\ น เสเทติ จ จิกฺขลฺลํ, โธวิ นิทฺธมนํ กเร. \\ ปีฐญฺจ โกฏฺฐเก กมฺมํ, มรุมฺพา สิลา นิทฺธมนํ. \\ นคฺคา ฉมาย วสฺสนฺเต, ปฏิจฺฉาที ตโย ตหิํ. \\ อุทปานํ ลุชฺชติ นีจํ, วลฺลิยา กายพนฺธเน. \\ ตุลํ กฏกฏํ จกฺกํ, พหู ภิชฺชนฺติ ภาชนา. \\ โลหทารุจมฺมขณฺฑํ, สาลาติณาปิธานิ จ. \\ โทณิจนฺทนิ ปาการํ, จิกฺขลฺลํ นิทฺธมเนน จ. \\ สีติคตํ โปกฺขรณิํ, ปุราณญฺจ นิลฺเลขณํ. \\ จาตุมาสํ สยนฺติ จ, นมตกญฺจ นธิฏฺฐเห. \\ อาสิตฺตกํ มโฬริกํ, ภุญฺชนฺเตกํ ตุวฏฺเฏยฺยุํ. \\ วฑฺโฒ โพธิ น อกฺกมิ, ฆฏํ กตกสมฺมชฺชนิ. \\ สกฺขรํ กถลญฺเจว, เผณกํ ปาทฆํสนี. \\ วิธูปนํ ตาลวณฺฏํ, มกสญฺจาปิ จามรี. \\ ฉตฺตํ วินา จ อาราเม, ตโย สิกฺกาย สมฺมุติ. \\ โรมสิตฺถา นขา ทีฆา, ฉินฺทนฺตงฺคุลิกา ทุกฺขา. \\ สโลหิตํ ปมาณญฺจ, วีสติ ทีฆเกสตา. \\ ขุรํ สิลํ สิปาฏิกํ, นมตกํ ขุรภณฺฑกํ. \\ มสฺสุํ กปฺเปนฺติ วฑฺเฒนฺติ, โคโลมิจตุรสฺสกํ. \\ ปริมุขํ อฑฺฒทุกญฺจ, ทาฐิสมฺพาธสํหเร. \\ อาพาธา กตฺตริวโณ, ทีฆํ สกฺขริกาย จ. \\ ปลิตํ ถกิตํ อุจฺจา, โลหภณฺฑญฺชนี สห. \\ ปลฺลตฺถิกญฺจ อาโยโค, วฏํ สลากพนฺธนํ. \\ กลาพุกํ เทฑฺฑุภกํ, มุรชํ มทฺทวีณกํ. \\ ปฏฺฏิกา สูกรนฺตญฺจ, ทสา มุรชเวณิกา. \\ อนฺโต โสภํ คุณกญฺจ, ปวนนฺโตปิ ชีรติ. \\ คณฺฐิกา อุจฺจาวจญฺจ, ผลกนฺเตปิ โอคาเห. \\ คิหิวตฺถํ หตฺถิโสณฺฑํ, มจฺฉกํ จตุกณฺณกํ. \\ ตาลวณฺฏํ สตวลิ, คิหิปารุตปารุปํ. \\ สํเวลฺลิ อุภโตกาชํ, ทนฺตกฏฺฐํ อาโกฏเน. \\ กณฺเฐ วิลคฺคํ ทายญฺจ, ปฏคฺคิ รุกฺขหตฺถินา. \\ ยเมเฬ โลกายตกํ, ปริยาปุณิํสุ วาจยุํ. \\ ติรจฺฉานกถา วิชฺฌา, ขิปิ มงฺคลํ ขาทิ จ. \\ วาตาพาโธ ทุสฺสติ จ, ทุคฺคนฺโธ ทุกฺขปาทุกา. \\ หิริยนฺติ ปารุทุคฺคนฺโธ, ตหํ ตหํ กโรนฺติ จ. \\ ทุคฺคนฺโธ กูปํ ลุชฺชนฺติ, อุจฺจวตฺถุ จเยน จ. \\ โสปานาลมฺพนพาหา, อนฺเต ทุกฺขญฺจ ปาทุกา. \\ พหิทฺธา โทณิ กฏฺฐญฺจ, ปิฐโร จ อปารุโต. \\ วจฺจกุฏิํ กวาฏญฺจ, ปิฏฺฐสํฆาฏเมว จ. \\ อุทุกฺขลุตฺตรปาโส, วฏฺฏิญฺจ กปิสีสกํ. \\ สูจิฆฏิตาฬจฺฉิทฺทํ, อาวิญฺฉนจฺฉิทฺทเมว จ. \\ รชฺชุ อุลฺลิตฺตาวลิตฺตํ, เสตวณฺณญฺจ กาฬกํ. \\ มาลากมฺมํ ลตากมฺมํ, มกรํ ปญฺจปฏิกํ. \\ จีวรวํสํ รชฺชุญฺจ, ชราทุพฺพลปาการํ. \\ โกฏฺฐเก จาปิ ตเถว, มรุมฺพํ ปทรสิลา. \\ สนฺติฏฺฐติ นิทฺธมนํ, กุมฺภิญฺจาปิ สราวกํ. \\ ทุกฺขํ หิริ อปิธานํ, อนาจารญฺจ อาจรุํ. \\ โลหภณฺฑํ อนุญฺญาสิ, ฐปยิตฺวา ปหรณิํ. \\ ฐปยิตฺวาสนฺทิปลฺลงฺกํ, ทารุปตฺตญฺจ ปาทุกํ. \\ สพฺพํ ทารุมยํ ภณฺฑํ, อนุญฺญาสิ มหามุนิ. \\ กตกํ กุมฺภการญฺจ, ฐปยิตฺวา ตถาคโต. \\ สพฺพมฺปิ มตฺติกาภณฺฑํ, อนุญฺญาสิ อนุกมฺปโก. \\ ยสฺส วตฺถุสฺส นิทฺเทโส, ปุริเมน ยทิ สมํ. \\ ตมฺปิ สํขิตฺตํ อุทฺทาเน, นยโต ตํ วิชานิยา. \\ เอวํ ทสสตา วตฺถู, วินเย ขุทฺทกวตฺถุเก. \\ สทฺธมฺมฏฺฐิติโก เจว, เปสลานญฺจนุคฺคโห. \\ สุสิกฺขิโต วินยธโร, หิตจิตฺโต สุเปสโล. \\ ปทีปกรโณ ธีโร, ปูชารโห พหุสฺสุโตติ.}
\csromangathagroup{\csromangathastanza{รุกฺเข ถมฺเภ จ กุฏฺเฏ จ, อฏฺฏาเน คนฺธสุตฺติยา. \\ วิคยฺห มลฺลโก กจฺฉุ, ชรา จ ปุถุปาณิกา.}\csromangathastanza{วลฺลิกาปิ จ ปามงฺโค, กณฺฐสุตฺตํ นธารเย. \\ กฏิ โอวฏฺฏิ กายุรํ, หตฺถาภรณมุทฺทิกา.}\csromangathastanza{ทีเฆ โกจฺเฉ ผเณ หตฺเถ, สิตฺถา อุทกเตลเก. \\ อาทาสุทปตฺตวณา, อาเลโปมฺมทฺทจุณฺณนา.}\csromangathastanza{ลญฺเฉนฺติ องฺคราคญฺจ, มุขราคํ ตทูภยํ. \\ จกฺขุโรคํ คิรคฺคญฺจ, อายตํ สรพาหิรํ.}\csromangathastanza{อมฺพเปสิสกเลหิ, อหิจฺฉินฺทิ จ จนฺทนํ. \\ อุจฺจาวจา ปตฺตมูลา, สุวณฺโณ พหลา วลี.}\csromangathastanza{\csromanfolio{286}จิตฺรา ทุสฺสติ ทุคฺคนฺโธ, อุเณฺห ภิชฺชิํสุ มิฑฺฒิยา. \\ ปริภณฺฑํ ติณํ โจฬํ, มาฬํ กุณฺโฑลิกาย จ.}\csromangathastanza{ถวิกา อํสพทฺธญฺจ, ตถา พนฺธนสุตฺตกา. \\ ขิเล มญฺเจ จ ปีเฐ จ, องฺเก ฉตฺเต ปณามนา.}\csromangathastanza{ตุมฺพฆฏิฉวสีสํ, จลกานิ ปฏิคฺคโห. \\ วิปฺผาลิทณฺฑโสวณฺณํ, ปตฺเต เปสิ จ นาฬิกา.}\csromangathastanza{กิณฺณสตฺตุ สริตญฺจ, มธุสิตฺถํ สิปาฏิกํ. \\ วิกณฺณํ พนฺธิวิสมํ, ฉมาชิรปโหติ จ.}\csromangathastanza{กฬิมฺภํ โมฆสุตฺตญฺจ, อโธตลฺลํ อุปาหนา. \\ องฺคุเล ปฏิคฺคหญฺจ, วิตฺถกํ ถวิกพทฺธกา.}\csromangathastanza{อชฺโฌกาเส นีจวตฺถุ, จโย จาปิ วิหญฺญเร. \\ ปริปติ ติณจุณฺณํ, อุลฺลิตฺต-อวลิตฺตกํ.}\csromangathastanza{เสตํ กาฬกวณฺณญฺจ, ปริกมฺมญฺจ เครุกํ. \\ มาลากมฺมํ ลตากมฺมํ, มกรทนฺตกปาฏิกํ.}\csromangathastanza{จีวรวํสํ รชฺชุญฺจ, อนุญฺญาสิ วินายโก. \\ อุชฺฌิตฺวา ปกฺกมนฺติ จ, กถินํ ปริภิชฺชติ.}\csromangathastanza{วินิเวฐิยติ กุฏฺเฏปิ, ปตฺเตนาทาย คจฺฉเร. \\ ถวิกา พนฺธสุตฺตญฺจ, พนฺธิตฺวา จ อุปาหนา.}\csromangathastanza{อุปาหนตฺถวิกญฺจ, อํสพทฺธญฺจ สุตฺตกํ. \\ อุทกากปฺปิยํ มคฺเค, ปริสฺสาวนโจฬกํ.}\csromangathastanza{ธมฺมกรณํ เทฺว ภิกฺขู, เวสาลิํ อคมา มุนิ. \\ ทณฺฑํ โอตฺถรกํ ตตฺถ, อนุญฺญาสิ ปริสฺสาวนํ.}\csromangathastanza{มกเสหิ ปณีเตน, พหฺวาพาธา จ ชีวโก. \\ จงฺกมนชนฺตาฆรํ, วิสเม นีจวตฺถุโก.}\csromangathastanza{ตโย จเย วิหญฺญนฺติ, โสปานาลมฺพเวทิตํ. \\ อชฺโฌกาเส ติณจุณฺณํ, อุลฺลิตฺต-อวลิตฺตกํ.}\csromangathastanza{\csromanfolio{287}เสตกํ กาฬวณฺณญฺจ, ปริกมฺมญฺจ เครุกํ. \\ มาลากมฺมํ ลตากมฺมํ, มกรทนฺตกปาฏิกํ.}\csromangathastanza{วํสํ จีวรรชฺชุญฺจ, อุจฺจํ จ วตฺถุกํ กเร. \\ จโย โสปานพาหญฺจ, กวาฏํ ปิฏฺฐสํฆาฏํ.}\csromangathastanza{อุทุกฺขลุตฺตรปาสกํ, วฏฺฏิญฺจ กปิสีสกํ. \\ สูจิฆฏิตาฬจฺฉิทฺทํ, อาวิญฺฉนญฺจ รชฺชุกํ.}\csromangathastanza{มณฺฑลํ ธูมเนตฺตญฺจ, มชฺเฌ จ มุขมตฺติกา. \\ โทณิ ทุคฺคนฺธา ทหติ, อุทกฏฺฐานํ สราวกํ.}\csromangathastanza{น เสเทติ จ จิกฺขลฺลํ, โธวิ นิทฺธมนํ กเร. \\ ปีฐญฺจ โกฏฺฐเก กมฺมํ, มรุมฺพา สิลา นิทฺธมนํ.}\csromangathastanza{นคฺคา ฉมาย วสฺสนฺเต, ปฏิจฺฉาที ตโย ตหิํ. \\ อุทปานํ ลุชฺชติ นีจํ, วลฺลิยา กายพนฺธเน.}\csromangathastanza{ตุลํ กฏกฏํ จกฺกํ, พหู ภิชฺชนฺติ ภาชนา. \\ โลหทารุจมฺมขณฺฑํ, สาลาติณาปิธานิ จ.}\csromangathastanza{โทณิจนฺทนิ ปาการํ, จิกฺขลฺลํ นิทฺธมเนน จ. \\ สีติคตํ โปกฺขรณิํ, ปุราณญฺจ นิลฺเลขณํ.}\csromangathastanza{จาตุมาสํ สยนฺติ จ, นมตกญฺจ\footnote{คนฺธปุปฺผํ (สฺยา)} นธิฏฺฐเห. \\ อาสิตฺตกํ มโฬริกํ, ภุญฺชนฺเตกํ ตุวฏฺเฏยฺยุํ.}\csromangathastanza{วฑฺโฒ โพธิ น อกฺกมิ, ฆฏํ กตกสมฺมชฺชนิ. \\ สกฺขรํ กถลญฺเจว, เผณกํ ปาทฆํสนี.}\csromangathastanza{วิธูปนํ ตาลวณฺฏํ, มกสญฺจาปิ จามรี. \\ ฉตฺตํ วินา จ อาราเม, ตโย สิกฺกาย สมฺมุติ.}\csromangathastanza{โรมสิตฺถา นขา ทีฆา, ฉินฺทนฺตงฺคุลิกา ทุกฺขา. \\ สโลหิตํ ปมาณญฺจ, วีสติ ทีฆเกสตา.}\csromangathastanza{ขุรํ สิลํ สิปาฏิกํ, นมตกํ ขุรภณฺฑกํ. \\ มสฺสุํ กปฺเปนฺติ วฑฺเฒนฺติ, โคโลมิจตุรสฺสกํ.}\csromangathastanza{\csromanfolio{288}ปริมุขํ อฑฺฒทุกญฺจ, ทาฐิสมฺพาธสํหเร. \\ อาพาธา กตฺตริวโณ, ทีฆํ สกฺขริกาย จ.}\csromangathastanza{ปลิตํ ถกิตํ อุจฺจา, โลหภณฺฑญฺชนี สห. \\ ปลฺลตฺถิกญฺจ อาโยโค, วฏํ สลากพนฺธนํ.}\csromangathastanza{กลาพุกํ เทฑฺฑุภกํ, มุรชํ มทฺทวีณกํ. \\ ปฏฺฏิกา สูกรนฺตญฺจ, ทสา มุรชเวณิกา.}\csromangathastanza{อนฺโต โสภํ คุณกญฺจ, ปวนนฺโตปิ ชีรติ. \\ คณฺฐิกา อุจฺจาวจญฺจ, ผลกนฺเตปิ โอคาเห.}\csromangathastanza{คิหิวตฺถํ หตฺถิโสณฺฑํ, มจฺฉกํ จตุกณฺณกํ. \\ ตาลวณฺฏํ สตวลิ, คิหิปารุตปารุปํ.}\csromangathastanza{สํเวลฺลิ อุภโตกาชํ, ทนฺตกฏฺฐํ อาโกฏเน. \\ กณฺเฐ วิลคฺคํ ทายญฺจ, ปฏคฺคิ รุกฺขหตฺถินา.}\csromangathastanza{ยเมเฬ โลกายตกํ, ปริยาปุณิํสุ วาจยุํ. \\ ติรจฺฉานกถา วิชฺฌา, ขิปิ มงฺคลํ ขาทิ จ.}\csromangathastanza{วาตาพาโธ ทุสฺสติ จ, ทุคฺคนฺโธ ทุกฺขปาทุกา. \\ หิริยนฺติ ปารุทุคฺคนฺโธ, ตหํ ตหํ กโรนฺติ จ.}\csromangathastanza{ทุคฺคนฺโธ กูปํ ลุชฺชนฺติ, อุจฺจวตฺถุ จเยน จ. \\ โสปานาลมฺพนพาหา, อนฺเต ทุกฺขญฺจ ปาทุกา.}\csromangathastanza{พหิทฺธา โทณิ กฏฺฐญฺจ, ปิฐโร จ อปารุโต. \\ วจฺจกุฏิํ กวาฏญฺจ, ปิฏฺฐสํฆาฏเมว จ.}\csromangathastanza{อุทุกฺขลุตฺตรปาโส, วฏฺฏิญฺจ กปิสีสกํ. \\ สูจิฆฏิตาฬจฺฉิทฺทํ, อาวิญฺฉนจฺฉิทฺทเมว จ.}\csromangathastanza{รชฺชุ อุลฺลิตฺตาวลิตฺตํ, เสตวณฺณญฺจ กาฬกํ. \\ มาลากมฺมํ ลตากมฺมํ, มกรํ ปญฺจปฏิกํ.}\csromangathastanza{จีวรวํสํ รชฺชุญฺจ, ชราทุพฺพลปาการํ. \\ โกฏฺฐเก จาปิ ตเถว, มรุมฺพํ ปทรสิลา.}\csromangathastanza{สนฺติฏฺฐติ นิทฺธมนํ, กุมฺภิญฺจาปิ สราวกํ. \\ ทุกฺขํ หิริ อปิธานํ, อนาจารญฺจ อาจรุํ.}\csromangathastanza{\csromanfolio{289}โลหภณฺฑํ อนุญฺญาสิ, ฐปยิตฺวา ปหรณิํ. \\ ฐปยิตฺวาสนฺทิปลฺลงฺกํ, ทารุปตฺตญฺจ ปาทุกํ.}\csromangathastanza{สพฺพํ ทารุมยํ ภณฺฑํ, อนุญฺญาสิ มหามุนิ. \\ กตกํ กุมฺภการญฺจ, ฐปยิตฺวา ตถาคโต.}\csromangathastanza{สพฺพมฺปิ มตฺติกาภณฺฑํ, อนุญฺญาสิ อนุกมฺปโก. \\ ยสฺส วตฺถุสฺส นิทฺเทโส, ปุริเมน ยทิ สมํ.}\csromangathastanza{ตมฺปิ สํขิตฺตํ อุทฺทาเน, นยโต ตํ วิชานิยา. \\ เอวํ ทสสตา วตฺถู, วินเย ขุทฺทกวตฺถุเก.}\csromangathastanza{สทฺธมฺมฏฺฐิติโก เจว, เปสลานญฺจนุคฺคโห. \\ สุสิกฺขิโต วินยธโร, หิตจิตฺโต สุเปสโล. \\ ปทีปกรโณ ธีโร, ปูชารโห พหุสฺสุโตติ.}}

\nitthitam{\textbf{ขุทฺทกวตฺถุกฺขนฺธโก นิฏฺฐิโต.}}
\csromanmatikaanchor{mtk.129}
\csromantocmarknopage{chapter}{6. เสนาสนกฺขนฺธก}{mtk.129}
\bookmark[dest={mtk.129},level=0]{6. เสนาสนกฺขนฺธก}
\chapterheadpageread{6. เสนาสนกฺขนฺธก}
\csromanmatikaanchor{mtk.130}
\csromantocmarkref{section}{1. ปฐมภาณวาร}{mtk.130}
\bookmark[dest={mtk.130},level=1]{1. ปฐมภาณวาร}
\csromanheader{1}{1. ปฐมภาณวาร}
\csromanmatikaanchor{mtk.131}
\csromantocmarkref{subsection}{วิหารานุชานน}{mtk.131}
\bookmark[dest={mtk.131},level=2]{วิหารานุชานน}
\csromanheader{2}{\textbf{วิหารานุชานน}}
\proseitem{294}{\csromanfolio{290}เตน สมเยน พุทฺโธ ภควา ราชคเห วิหรติ เวฬุวเน กลนฺทกนิวาเป.\csromanspacer{} เตน โข ปน สมเยน ภควตา ภิกฺขูนํ เสนาสนํ อปญฺญตฺตํ โหติ.\csromanspacer{} เต จ ภิกฺขู ตหํ ตหํ วิหรนฺติ อรญฺเญ รุกฺขมูเล ปพฺพเต กนฺทรายํ คิริคุหายํ สุสาเน วนปตฺเถ อชฺโฌกาเส ปลาลปุญฺเช, เต กาลสฺเสว ตโต ตโต อุปนิกฺขมนฺติ อรญฺญา รุกฺขมูลา ปพฺพตา กนฺทรา คิริคุหา สุสานา วนปตฺถา อชฺโฌกาสา ปลาลปุญฺชา ปาสาทิเกน อภิกฺกนฺเตน ปฏิกฺกนฺเตน อาโลกิเตน วิโลกิเตน สมิญฺชิเตน ปสาริเตน โอกฺขิตฺตจกฺขู อิริยาปถสมฺปนฺนา.\csromanspacer{} เตน โข ปน สมเยน ราชคหโก เสฏฺฐี\footnote{เสฏฺฐิ (ก)} กาลสฺเสว อุยฺยานํ อคมาสิ.\csromanspacer{} อทฺทสา โข ราชคหโก เสฏฺฐี เต ภิกฺขู กาลสฺเสว ตโต ตโต อุปนิกฺขมนฺเต อรญฺญา รุกฺขมูลา ปพฺพตา กนฺทรา คิริคุหา สุสานา วนปตฺถา อชฺโฌกาสา ปลาลปุญฺชา ปาสาทิเกน อภิกฺกนฺเตน ปฏิกฺกนฺเตน อาโลกิเตน วิโลกิเตน สมิญฺชิเตน ปสาริเตน โอกฺขิตฺตจกฺขู อิริยาปถสมฺปนฺเน, ทิสฺวานสฺส จิตฺตํ ปสีทิ.\csromanspacer{} อถ โข ราชคหโก เสฏฺฐี เยน เต ภิกฺขู เตนุปสงฺกมิ, อุปสงฺกมิตฺวา เต ภิกฺขู เอตทโวจ “สจาหํ ภนฺเต วิหาเร การาเปยฺยํ, วเสยฺยาถ เม วิหาเรสู”ติ.\csromanspacer{} น โข คหปติ ภควตา วิหารา อนุญฺญาตาติ.\csromanspacer{} เตน หิ ภนฺเต ภควนฺตํ ปฏิปุจฺฉิตฺวา มม อาโรเจยฺยาถาติ. “เอวํ คหปตี”ติ โข เต ภิกฺขู ราชคหกสฺส เสฏฺฐิสฺส ปฏิสฺสุตฺวา เยน ภควา เตนุปสงฺกมิํสุ, อุปสงฺกมิตฺวา ภควนฺตํ อภิวาเทตฺวา เอกมนฺตํ นิสีทิํสุ.\csromanspacer{} เอกมนฺตํ นิสินฺนา โข เต ภิกฺขู ภควนฺตํ เอตทโวจุํ “ราชคหโก ภนฺเต เสฏฺฐี วิหาเร การาเปตุกาโม, กถํ นุ โข ภนฺเต อเมฺหหิ\footnote{ภนฺเต (สี, ก)} ปฏิปชฺชิตพฺพนฺ”ติ.\csromanspacer{} อถ โข ภควา เอตสฺมิํ นิทาเน เอตสฺมิํ ปกรเณ ธมฺมิํ กถํ กตฺวา ภิกฺขู อามนฺเตสิ “อนุชานามิ ภิกฺขเว ปญฺจ เลณานิ\footnote{ปญฺจ เสนาสนานิ (สฺยา)} วิหารํ อฑฺฒโยคํ ปาสาทํ หมฺมิยํ คุหนฺ”ติ.}

\prose{\csromanfolio{291}อถ โข เต ภิกฺขู เยน ราชคหโก เสฏฺฐี เตนุปสงฺกมิํสุ, อุปสงฺกมิตฺวา ราชคหกํ เสฏฺฐิํ เอตทโวจุํ “อนุญฺญาตา โข คหปติ ภควตา วิหารา, ยสฺสทานิ กาลํ มญฺญสี”ติ.\csromanspacer{} อถ โข ราชคหโก เสฏฺฐี เอกาเหเนว สฏฺฐิวิหาเร ปติฏฺฐาเปสิ.\csromanspacer{} อถ โข ราชคหโก เสฏฺฐี เต สฏฺฐิวิหาเร ปริโยสาเปตฺวา เยน ภควา เตนุปสงฺกมิ, อุปสงฺกมิตฺวา ภควนฺตํ อภิวาเทตฺวา เอกมนฺตํ นิสีทิ.\csromanspacer{} เอกมนฺตํ นิสินฺโน โข ราชคหโก เสฏฺฐี ภควนฺตํ เอตทโวจ “อธิวาเสตุ เม ภนฺเต ภควา สฺวาตนาย ภตฺตํ สทฺธิํ ภิกฺขุสํเฆนา”ติ.\csromanspacer{} อธิวาเสสิ ภควา ตุณฺหีภาเวน.\csromanspacer{} อถ โข ราชคหโก เสฏฺฐี ภควโต อธิวาสนํ วิทิตฺวา อุฏฺฐายาสนา ภควนฺตํ อภิวาเทตฺวา ปทกฺขิณํ กตฺวา ปกฺกามิ.}

\prose{อถ โข ราชคหโก เสฏฺฐี ตสฺสา รตฺติยา อจฺจเยน ปณีตํ ขาทนียํ โภชนียํ ปฏิยาทาเปตฺวา ภควโต กาลํ อาโรจาเปสิ “กาโล ภนฺเต, นิฏฺฐิตํ ภตฺตนฺ”ติ.\csromanspacer{} อถ โข ภควา ปุพฺพณฺหสมยํ นิวาเสตฺวา ปตฺตจีวรมาทาย เยน ราชคหกสฺส เสฏฺฐิสฺส นิเวสนํ เตนุปสงฺกมิ, อุปสงฺกมิตฺวา ปญฺญตฺเต อาสเน นิสีทิ สทฺธิํ ภิกฺขุสํเฆน.\csromanspacer{} อถ โข ราชคหโก เสฏฺฐี พุทฺธปฺปมุขํ ภิกฺขุสํฆํ ปณีเตน ขาทนีเยน โภชนีเยน สหตฺถา สนฺตปฺเปตฺวา สมฺปวาเรตฺวา ภควนฺตํ ภุตฺตาวิํ โอนีตปตฺตปาณิํ เอกมนฺตํ นิสีทิ.\csromanspacer{} เอกมนฺตํ นิสินฺโน โข ราชคหโก เสฏฺฐี ภควนฺตํ เอตทโวจ “เอเต เม ภนฺเต สฏฺฐิวิหารา ปุญฺญตฺถิเกน สคฺคตฺถิเกน การาปิตา, กถาหํ ภนฺเต เตสุ วิหาเรสุ ปฏิปชฺชามี”ติ.\csromanspacer{} เตน หิ ตฺวํ คหปติ เต สฏฺฐิวิหาเร อาคตานาคตสฺส จาตุทฺทิสสฺส สํฆสฺส ปติฏฺฐาเปหีติ. “เอวํ ภนฺเต”ติ โข ราชคหโก เสฏฺฐี ภควโต ปฏิสฺสุตฺวา เต สฏฺฐิวิหาเร อาคตานาคตสฺส จาตุทฺทิสสฺส สํฆสฺส ปติฏฺฐาเปสิ.}

\csromanheader{2}{295. อถ โข ภควา ราชคหกํ เสฏฺฐิํ อิมาหิ คาถาหิ อนุโมทิ\csromandash{}}
\csromangathameasure{“สีตํ อุณฺหํ ปฏิหนฺติ, ตโต วาฬมิคานิ จ. \\ สรีสเป จ มกเส, สิสิเร จาปิ วุฏฺฐิโย. \\ ตโต วาตาตโป โฆโร, สญฺชาโต ปฏิหญฺญติ. \\ เลณตฺถญฺจ สุขตฺถญฺจ, ฌายิตุญฺจ วิปสฺสิตุํ. \\ วิหารทานํ สํฆสฺส, อคฺคํ พุทฺเธน วณฺณิตํ. \\ ตสฺมา หิ ปณฺฑิโต โปโส, สมฺปสฺสํ อตฺถมตฺตโน. \\ วิหาเร การเย รมฺเม, วาสเยตฺถ พหุสฺสุเต. \\ เตสํ อนฺนญฺจ ปานญฺจ, วตฺถเสนาสนานิ จ. \\ ทเทยฺย อุชุภูเตสุ, วิปฺปสนฺเนน เจตสา. \\ เต ตสฺส ธมฺมํ เทเสนฺติ, สพฺพทุกฺขาปนูทนํ. \\ ยํ โส ธมฺมํ อิธญฺญาย, ปรินิพฺพาติ อนาสโว”ติ. \\ อถ โข ภควา ราชคหกํ เสฏฺฐิํ อิมาหิ คาถาหิ อนุโมทิตฺวา อุฏฺฐายาสนา ปกฺกามิ.}
\csromangathagroup{\csromangathastanza{“สีตํ อุณฺหํ ปฏิหนฺติ\footnote{ปฏิหนติ (ก)}, ตโต วาฬมิคานิ จ. \\ สรีสเป จ มกเส, สิสิเร จาปิ วุฏฺฐิโย.}\csromangathastanza{\csromanfolio{292}ตโต วาตาตโป โฆโร, สญฺชาโต\footnote{วาตาตเป โฆเร, สญฺชเต (ก, สทฺทนีติ)} ปฏิหญฺญติ. \\ เลณตฺถญฺจ สุขตฺถญฺจ, ฌายิตุญฺจ วิปสฺสิตุํ.}\csromangathastanza{วิหารทานํ สํฆสฺส, อคฺคํ พุทฺเธน\footnote{พุทฺเธหิ (สฺยา)} วณฺณิตํ. \\ ตสฺมา หิ ปณฺฑิโต โปโส, สมฺปสฺสํ อตฺถมตฺตโน.}\csromangathastanza{วิหาเร การเย รมฺเม, วาสเยตฺถ พหุสฺสุเต. \\ เตสํ อนฺนญฺจ ปานญฺจ, วตฺถเสนาสนานิ จ.}\csromangathastanza{ทเทยฺย อุชุภูเตสุ, วิปฺปสนฺเนน เจตสา. \\ เต ตสฺส ธมฺมํ เทเสนฺติ, สพฺพทุกฺขาปนูทนํ.}\csromangathastanza{ยํ โส ธมฺมํ อิธญฺญาย, ปรินิพฺพาติ อนาสโว”ติ. \\ อถ โข ภควา ราชคหกํ เสฏฺฐิํ อิมาหิ คาถาหิ อนุโมทิตฺวา อุฏฺฐายาสนา ปกฺกามิ.}}

\proseitem{296}{อสฺโสสุํ โข มนุสฺสา “ภควโต กิร วิหารา อนุญฺญาตา”ติ สกฺกจฺจํ\footnote{เต สกฺกจฺจํ (สฺยา, กํ)} วิหาเร การาเปนฺติ, เต วิหารา อกวาฏกา โหนฺติ, อหีปิ วิจฺฉิกาปิ สตปทิโยปิ ปวิสนฺติ.\csromanspacer{} ภควโต เอตมตฺถํ อาโรเจสุํ.\csromanspacer{} อนุชานามิ ภิกฺขเว กวาฏนฺติ.\csromanspacer{} ภิตฺติฉิทฺทํ กริตฺวา วลฺลิยาปิ รชฺชุยาปิ กวาฏํ พนฺธนฺติ, อุนฺทูเรหิปิ อุปจิกาหิปิ ขชฺชนฺติ, ขยิตพนฺธนานิ กวาฏานิ ปตนฺติ.\csromanspacer{} ภควโต เอตมตฺถํ อาโรเจสุํ.\csromanspacer{} อนุชานามิ ภิกฺขเว ปิฏฺฐสํฆาฏํ อุทุกฺขลิกํ อุตฺตรปาสกนฺติ.\csromanspacer{} กวาฏา น ผุสียนฺติ.\csromanspacer{} ภควโต เอตมตฺถํ อาโรเจสุํ.\csromanspacer{} อนุชานามิ ภิกฺขเว อาวิญฺฉนจฺฉิทฺทํ อาวิญฺฉนรชฺชุนฺติ.\csromanspacer{} กวาฏา น ถกิยนฺติ.\csromanspacer{} ภควโต เอตมตฺถํ อาโรเจสุํ.\csromanspacer{} อนุชานามิ ภิกฺขเว อคฺคฬวฏฺฏิํ กปิสีสกํ สูจิกํ ฆฏิกนฺติ.}

\prose{เตน โข ปน สมเยน ภิกฺขู น สกฺโกนฺติ กวาฏํ อปาปุริตุํ.\csromanspacer{} ภควโต เอตมตฺถํ อาโรเจสุํ.\csromanspacer{} อนุชานามิ ภิกฺขเว ตาฬจฺฉิทฺทํ, ตีณี ตาฬานิ โลหตาฬํ กฏฺฐตาฬํ วิสาณตาฬนฺติ, เยหิ\footnote{เยปิ (สี, สฺยา)} เต อุคฺฆาเฏตฺวา ปวิสนฺติ\footnote{วิสาณตาฬํ, เยหิ เต อุคฺฆาเฏตฺวา ปวิสนฺตีติ (ก)}.\csromanspacer{} วิหารา อคุตฺตา โหนฺติ.\csromanspacer{} ภควโต เอตมตฺถํ อาโรเจสุํ.\csromanspacer{} อนุชานามิ ภิกฺขเว ยนฺตกํ สูจิกนฺติ.}

\prose{\csromanfolio{293}เตน โข ปน สมเยน วิหารา ติณจฺฉทนา โหนฺติ, สีตกาเล สีตา, อุณฺหกาเล อุณฺหา.\csromanspacer{} ภควโต เอตมตฺถํ อาโรเจสุํ.\csromanspacer{} อนุชานามิ ภิกฺขเว โอคุมฺเผตฺวา อุลฺลิตฺตาวลิตฺตํ กาตุนฺติ.}

\prose{เตน โข ปน สมเยน วิหารา อวาตปานกา โหนฺติ อจกฺขุสฺสา ทุคฺคนฺธา.\csromanspacer{} ภควโต เอตมตฺถํ อาโรเจสุํ.\csromanspacer{} อนุชานามิ ภิกฺขเว ตีณิ วาตปานานิ เวทิกาวาตปานํ ชาลวาตปานํ สลากวาตปานนฺติ.\csromanspacer{} วาตปานนฺตริกาย กาฬกาปิ วคฺคุลิโยปิ ปวิสนฺติ.\csromanspacer{} ภควโต เอตมตฺถํอาโรเจสุํ.\csromanspacer{} อนุชานามิ ภิกฺขเว วาตปานจกฺกลิกนฺติ.\csromanspacer{} จกฺกลิกนฺตริกายปิ กาฬกาปิ วคฺคุลิโยปิ ปวิสนฺติ.\csromanspacer{} ภควโต เอตมตฺถํ อาโรเจสุํ.\csromanspacer{} อนุชานามิ ภิกฺขเว วาตปานกวาฏกํ วาตปานภิสิกนฺติ.}

\prose{เตน โข ปนสมเยน ภิกฺขู ฉมายํ สยนฺติ, คตฺตานิปิ จีวรานิปิ ปํสุกิตานิ โหนฺติ.\csromanspacer{} ภควโต เอตมตฺถํ อาโรเจสุํ.\csromanspacer{} อนุชานามิ ภิกฺขเว ติณสนฺถารกนฺติ.\csromanspacer{} ติณสนฺถารโก อุนฺทูเรหิปิ อุปจิกาหิปิ ขชฺชติ.\csromanspacer{} ภควโต เอตมตฺถํ อาโรเจสุํ.\csromanspacer{} อนุชานามิ ภิกฺขเว มิฑฺฒินฺติ\footnote{มีฒนฺติ (สี), มิฒินฺติ (สฺยา)}.\csromanspacer{} มิฑฺฒิยา คตฺตานิ ทุกฺขา โหนฺติ.\csromanspacer{} ภควโต เอตมตฺถํ อาโรเจสุํ.\csromanspacer{} อนุชานามิ ภิกฺขเว พิทลมญฺจกนฺติ.}

\csromanmatikaanchor{mtk.132}
\csromantocmarkref{subsection}{มญฺจปีฐาทิ-อนุชานน}{mtk.132}
\bookmark[dest={mtk.132},level=2]{มญฺจปีฐาทิ-อนุชานน}
\csromanheader{2}{\textbf{มญฺจปีฐาทิ-อนุชานน}}
\proseitem{297}{เตน โข ปน สมเยน สํฆสฺส โสสานิโก มสารโก มญฺโจ อุปฺปนฺโน โหติ.\csromanspacer{} ภควโต เอตมตฺถํ อาโรเจสุํ.\csromanspacer{} อนุชานามิ ภิกฺขเว มสารกํ มญฺจนฺติ.\csromanspacer{} มสารกํ ปีฐํ อุปฺปนฺนํ โหติ.\csromanspacer{} ภควโต เอตมตฺถํ อาโรเจสุํ.\csromanspacer{} อนุชานามิ ภิกฺขเว มสารกํ ปีฐนฺติ.}

\prose{เตน โข ปน สมเยน สํฆสฺส โสสานิโก พุนฺทิกาพทฺโธ มญฺโจ อุปฺปนฺโน โหติ.\csromanspacer{} ภควโต เอตมตฺถํ อาโรเจสุํ.\csromanspacer{} อนุชานามิ ภิกฺขเว พุนฺทิกาพทฺธํ มญฺจนฺติ.\csromanspacer{} พุนฺทิกาพทฺธํ ปีฐํ อุปฺปนฺนํ โหติ.\csromanspacer{} ภควโต เอตมตฺถํ อาโรเจสุํ.\csromanspacer{} อนุชานามิ ภิกฺขเว พุนฺทิกาพทฺธํ ปิฐนฺติ.}

\prose{เตน โข ปน สมเยน สํฆสฺส โสสานิโก กุฬีรปาทโก มญฺโจ อุปฺปนฺโน โหติ.\csromanspacer{} ภควโต เอตมตฺถํ อาโรเจสุํ.\csromanspacer{} อนุชานามิ \csromanfolio{294}ภิกฺขเว กุฬีรปาทกํ มญฺจนฺติ.\csromanspacer{} กุฬีรปาทกํ ปีฐํ อุปฺปนฺนํ โหติ.\csromanspacer{} ภควโต เอตมตฺถํ อาโรเจสุํ.\csromanspacer{} อนุชานามิ ภิกฺขเว กุฬีรปาทกํ ปีฐนฺติ.}

\prose{เตน โข ปน สมเยน สํฆสฺส โสสานิโก อาหจฺจปาทโก มญฺโจ อุปฺปนฺโน โหติ.\csromanspacer{} ภควโต เอตมตฺถํ อาโรเจสุํ.\csromanspacer{} อนุชานามิ ภิกฺขเว อาหจฺจปาทกํ มญฺจนฺติ.\csromanspacer{} อาหาจฺจปาทกํ ปีฐํ อุปฺปนฺนํ โหติ.\csromanspacer{} ภควโต เอตมตฺถํ อาโรเจสุํ.\csromanspacer{} อนุชานามิ ภิกฺขเว อาหจฺจปาทกํ ปีฐนฺติ.}

\prose{เตน โข ปน สมเยน สํฆสฺส อาสนฺทิโก อุปฺปนฺโน โหติ.\csromanspacer{} ภควโต เอตมตฺถํ อาโรเจสุํ.\csromanspacer{} อนุชานามิ ภิกฺขเว อาสนฺทิกนฺติ.\csromanspacer{} อุจฺจโก อาสนฺทิโก อุปฺปนฺโน โหติ.\csromanspacer{} ภควโต เอตมตฺถํ อาโรเจสุํ.\csromanspacer{} อนุชานามิ ภิกฺขเว อุจฺจกมฺปิ อาสนฺทิกนฺติ.\csromanspacer{} สตฺตงฺโค อุปฺปนฺโน โหติ.\csromanspacer{} ภควโต เอตมตฺถํ อาโรเจสุํ.\csromanspacer{} อนุชานามิ ภิกฺขเว สตฺตงฺคนฺติ.\csromanspacer{} อุจฺจโก สตฺตงฺโค อุปฺปนฺโน โหติ.\csromanspacer{} ภควโต เอตมตฺถํ อาโรเจสุํ.\csromanspacer{} อนุชานามิ ภิกฺขเว อุจฺจกมฺปิ สตฺตงฺคนฺติ.\csromanspacer{} ภทฺทปีฐํ อุปฺปนฺนํ โหติ.\csromanspacer{} ภควโต เอตมตฺถํ อาโรเจสุํ.\csromanspacer{} อนุชานามิ ภิกฺขเว ภทฺทปีฐนฺติ.\csromanspacer{} ปีฐิกา อุปฺปนฺนา โหติ.\csromanspacer{} ภควโต เอตมตฺถํ อาโรเจสุํ.\csromanspacer{} อนุชานามิ ภิกฺขเว ปีฐิกนฺติ.\csromanspacer{} เอฬกปาทกํ ปีฐํ อุปฺปนฺนํ โหติ.\csromanspacer{} ภควโต เอตมตฺถํ อาโรเจสุํ.\csromanspacer{} อนุชานามิ ภิกฺขเว เอฬกปาทกํ ปีฐนฺติ.\csromanspacer{} อามลกวฏฺฏิกํ\footnote{อามลกวณฺฏิกํ (สฺยา)} ปีฐํ อุปฺปนฺนํ โหติ.\csromanspacer{} ภควโต เอตมตฺถํ อาโรเจสุํ.\csromanspacer{} อนุชานามิ ภิกฺขเว อามลกวฏฺฏิกํ ปีฐนฺติ.\csromanspacer{} ผลกํ อุปฺปนฺนํ โหติ.\csromanspacer{} ภควโต เอตมตฺถํ อาโรเจสุํ.\csromanspacer{} อนุชานามิ ภิกฺขเว ผลกนฺติ.\csromanspacer{} โกจฺฉํ อุปฺปนฺนํ โหติ.\csromanspacer{} ภควโต เอตมตฺถํ อาโรเจสุํ.\csromanspacer{} อนุชานามิ ภิกฺขเว โกจฺฉนฺติ.\csromanspacer{} ปลาลปีฐํ อุปฺปนฺนํ โหติ.\csromanspacer{} ภควโต เอตมตฺถํ อาโรเจสุํ.\csromanspacer{} อนุชานามิ ภิกฺขเว ปลาลปีฐนฺติ.}

\prose{เตน โข ปน สมเยน ฉพฺพคฺคิยา ภิกฺขู อุจฺเจ มญฺเจ สยนฺติ.\csromanspacer{} มนุสฺสา วิหารจาริกํ อาหิณฺฑนฺตา ปสฺสิตฺวา อุชฺฌายนฺติ ขิยฺยนฺติ วิปาเจนฺติ “เสยฺยถาปิ คิหี กามโภคิโน”ติ.\csromanspacer{} ภควโต เอตมตฺถํ อาโรเจสุํ.\csromanspacer{} น ภิกฺขเว อุจฺเจ มญฺเจ สยิตพฺพํ, โย สเยยฺย, อาปตฺติ ทุกฺกฏสฺสาติ.}

\prose{\csromanfolio{295}เตน โข ปน สมเยน อญฺญตโร ภิกฺขุ นีเจ มญฺเจ สยนฺโต อหินา ทฏฺโฐ โหติ.\csromanspacer{} ภควโต เอตมตฺถํ อาโรเจสุํ.\csromanspacer{} อนุชานามิ ภิกฺขเว มญฺจปฏิปาทกนฺติ.}

\prose{เตน โข ปน สมเยน ฉพฺพคฺคิยา ภิกฺขู อุจฺเจ มญฺจปฏิปาทเก ธาเรนฺติ, สห มญฺจปฏิปาทเกหิ ปเวเธนฺติ.\csromanspacer{} มนุสฺสา วิหารจาริกํ อาหิณฺฑนฺตา ปสฺสิตฺวา อุชฺฌายนฺติ ขิยฺยนฺติ วิปาเจนฺติ “เสยฺยถาปิ คิหี กามโภคิโน”ติ.\csromanspacer{} ภควโต เอตมตฺถํ อาโรเจสุํ.\csromanspacer{} น ภิกฺขเว อุจฺจา มญฺจปฏิปาทกา ธาเรตพฺพา, โย ธาเรยฺย, อาปตฺติ ทุกฺกฏสฺส.\csromanspacer{} อนุชานามิ ภิกฺขเว อฏฺฐงฺคุลปรมํ มญฺจปฏิปาทกนฺติ.}

\prose{เตน โข ปน สมเยน สํฆสฺส สุตฺตํ อุปฺปนฺนํ โหติ.\csromanspacer{} ภควโต เอตมตฺถํ อาโรเจสุํ.\csromanspacer{} อนุชานามิ ภิกฺขเว สุตฺตํ มญฺจํ เวเฐตุนฺติ\footnote{เวตุนฺติ (สี)}.\csromanspacer{} องฺคานิ พหุสุตฺตํ ปริยาทิยนฺติ.\csromanspacer{} ภควโต เอตมตฺถํ อาโรเจสุํ.\csromanspacer{} อนุชานามิ ภิกฺขเว องฺเค วิชฺฌิตฺวา อฏฺฐปทกํ เวเฐตุนฺติ.\csromanspacer{} โจฬกํ อุปฺปนฺนํ โหติ.\csromanspacer{} ภควโต เอตมตฺถํ อาโรเจสุํ.\csromanspacer{} อนุชานามิ ภิกฺขเว จิมิลิกํ กาตุนฺติ.\csromanspacer{} ตูลิกา อุปฺปนฺนาโหติ.\csromanspacer{} ภควโต เอตมตฺถํ อาโรเจสุํ.\csromanspacer{} อนุชานามิ ภิกฺขเว วิชเฏตฺวา พิพฺโพหนํ\footnote{พิมฺโพหนํ (สี, สฺยา, พิมฺพ + โอหนํ)} กาตุํ, ตีณิ ตูลานิ รุกฺขตูลํ ลตาตูลํ โปฏกิตูลนฺติ.}

\prose{เตน โข ปน สมเยน ฉพฺพคฺคิยา ภิกฺขู อทฺธกายิกานิ\footnote{อฑฺฒกายิกานิ (ก)} พิพฺโพหนานิ ธาเรนฺติ.\csromanspacer{} มนุสฺสา วิหารจาริกํ อาหิณฺฑนฺตา ปสฺสิตฺวา อุชฺฌายนฺติ ขิยฺยนฺติ วิปาเจนฺติ “เสยฺยถาปิ คิหี กามโภคิโน”ติ.\csromanspacer{} ภควโต เอตมตฺถํ อาโรเจสุํ.\csromanspacer{} น ภิกฺขเว อฑฺฒกายิกานิ พิพฺโพหนานิ ธาเรตพฺพานิ, โย ธาเรยฺย, อาปตฺติ ทุกฺกฏสฺส.\csromanspacer{} อนุชานามิ ภิกฺขเว สีสปฺปมาณํ พิพฺโพหนํ กาตุนฺติ.}

\prose{เตน โข ปน สมเยน ราชคเห คิรคฺคสมชฺโช โหติ, มนุสฺสา มหามตฺตานํ อตฺถาย ภิสิโย ปฏิยาเทนฺติ อุณฺณภิสิํ โจฬภิสิํ วากภิสิํ ติณภิสิํ ปณฺณภิสิํ, เต วีติวตฺเต สมชฺเช ฉวิํ อุปฺปาเฏตฺวา หรนฺติ.\csromanspacer{} อทฺทสาสุํ โข ภิกฺขู สมชฺชฏฺฐาเน พหุํ อุณฺณมฺปิ โจฬกมฺปิ วากมฺปิ ติณมฺปิ ปณฺณมฺปิ ฉฏฺฏิตํ, ทิสฺวาน ภควโต เอตมตฺถํ \csromanfolio{296}อาโรเจสุํ.\csromanspacer{} อนุชานามิ ภิกฺขเว ปญฺจ ภิสิโย อุณฺณภิสิํ โจฬภิสิํ วากภิสิํ ติณภิสิํ ปณฺณภิสินฺติ.}

\prose{เตน โข ปน สมเยน สํฆสฺส เสนาสนปริกฺขาริกํ ทุสฺสํ อุปฺปนฺนํ โหติ.\csromanspacer{} ภควโต เอตมตฺถํ อาโรเจสุํ.\csromanspacer{} อนุชานามิ ภิกฺขเว ภิสิํ โอนนฺธิตุนฺติ\footnote{โอนทฺธิตุํ (สฺยา)}.}

\prose{เตน โข ปน สมเยน ภิกฺขู มญฺจภิสิํ ปีเฐ สนฺถรนฺติ, ปีฐภิสิํ มญฺเจ สนฺถรนฺติ, ภิสิโย ปริภิชฺชนฺติ.\csromanspacer{} ภควโต เอตมตฺถํ อาโรเจสุํ.\csromanspacer{} อนุชานามิ ภิกฺขเว โอนทฺธมญฺจํ\footnote{โอนนฺธมญฺจ (ก) เอวมุปริปิ.} โอนทฺธปีฐนฺติ.\csromanspacer{} อุลฺโลกํ อกริตฺวา สนฺถรนฺติ, เหฏฺฐโต นิปตนฺติ\footnote{นิปฏนฺติ (ก), นิปฺผฏนฺติ (สี), นิปฺปาเฏนฺติ (สฺยา)} ฯเปฯ.\csromanspacer{} อนุชานามิ ภิกฺขเว อุลฺโลกํ กริตฺวา สนฺถริตฺวา ภิสิํ โอนนฺธิตุนฺติ.\csromanspacer{} ฉวิํ อุปฺปาเฏตฺวา หรนฺติ ฯเปฯ.\csromanspacer{} อนุชานามิ ภิกฺขเว โผสิตุนฺติ.\csromanspacer{} หรนฺติเยว ฯเปฯ.\csromanspacer{} อนุชานามิ ภิกฺขเว ภตฺติกมฺมนฺติ.\csromanspacer{} หรนฺติเยว ฯเปฯ.\csromanspacer{} อนุชานามิ ภิกฺขเว หตฺถภตฺติกมฺมนฺติ.\csromanspacer{} หรนฺติเยว ฯเปฯ.\csromanspacer{} อนุชานามิ ภิกฺขเว หตฺถภตฺตินฺติ.}

\csromanmatikaanchor{mtk.133}
\csromantocmarkref{subsection}{เสตวณฺณาทิ-อนุชานน}{mtk.133}
\bookmark[dest={mtk.133},level=2]{เสตวณฺณาทิ-อนุชานน}
\csromanheader{2}{\textbf{เสตวณฺณาทิ-อนุชานน}}
\proseitem{298}{เตน โข ปน สมเยน ติตฺถิยานํ เสยฺยาโย เสตวณฺณา โหนฺติ, กาฬวณฺณกตา ภูมิ, เครุกปริกมฺมกตา ภิตฺติ, พหู มนุสฺสา เสยฺยาเปกฺขกา คจฺฉนฺติ.\csromanspacer{} ภควโต เอตมตฺถํ อาโรเจสุํ.\csromanspacer{} อนุชานามิ ภิกฺขเว วิหาเร เสตวณฺณํ กาฬวณฺณํ เครุกปริกมฺมนฺติ.}

\prose{เตน โข ปน สมเยน ผรุสาย ภิตฺติยา เสตวณฺโณ น นิปตติ.\csromanspacer{} ภควโต เอตมตฺถํ อาโรเจสุํ.\csromanspacer{} อนุชานามิ ภิกฺขเว ถุสปิณฺฑํ ทตฺวา ปาณิกาย ปฏิพาเหตฺวา เสตวณฺณํ นิปาเตตุนฺติ.\csromanspacer{} เสตวณฺโณ อนิพนฺธนีโย โหติ ฯเปฯ.\csromanspacer{} อนุชานามิ ภิกฺขเว สณฺหมตฺติกํ ทตฺวา ปาณิกาย ปฏิพาเหตฺวา เสตวณฺณํ นิปาเตตุนฺติ.\csromanspacer{} เสตวณฺโณ อนิพนฺธนีโย โหติ ฯเปฯ.\csromanspacer{} อนุชานามิ ภิกฺขเว อิกฺกาสํ ปิฏฺฐมทฺทนฺติ.}

\prose{เตน โข ปน สมเยน ผรุสาย ภิตฺติยา เครุกา น นิปตติ.\csromanspacer{} ภควโต เอตมตฺถํ อาโรเจสุํ.\csromanspacer{} อนุชานามิ ภิกฺขเว ถุสปิณฺฑํ ทตฺวา ปาณิกาย ปฏิพาเหตฺวา เครุกํ นิปาเตตุนฺติ.\csromanspacer{} เครุกา อนิพนฺธนียา \csromanfolio{297}โหติ ฯเปฯ.\csromanspacer{} อนุชานามิ ภิกฺขเว กุณฺฑกมตฺติกํ ทตฺวา ปาณิกาย ปฏิพาเหตฺวา เครุกํ นิปาเตตุนฺติ.\csromanspacer{} เครุกา อนิพนฺธนียา โหติ ฯเปฯ.\csromanspacer{} อนุชานามิ ภิกฺขเว สาสปกุฏฺฏํ สิตฺถเตลกนฺติ.\csromanspacer{} อจฺจุสฺสนฺนํ โหติ ฯเปฯ.\csromanspacer{} อนุชานามิ ภิกฺขเว โจฬเกน ปจฺจุทฺธริตุนฺติ.}

\prose{เตน โข ปน สมเยน ผรุสาย ภูมิยา กาฬวณฺโณ น นิปตติ.\csromanspacer{} ภควโต เอตมตฺถํ อาโรเจสุํ.\csromanspacer{} อนุชานามิ ภิกฺขเว ถุสปิณฺฑํ ทตฺวา ปาณิกาย ปฏิพาเหตฺวา กาฬวณฺณํ นิปาเตตุนฺติ.\csromanspacer{} กาฬวณฺโณ อนิพนฺธนีโย โหติ ฯเปฯ.\csromanspacer{} อนุชานามิ ภิกฺขเว คณฺฑุมตฺติกํ ทตฺวา ปาณิกาย ปฏิพาเหตฺวา กาฬวณฺณํ นิปาเตตุนฺติ.\csromanspacer{} กาฬวณฺโณ อนิพนฺธนีโย โหติ ฯเปฯ.\csromanspacer{} อนุชานามิ ภิกฺขเว อิกฺกาสํ กสาวนฺติ.}

\csromanmatikaanchor{mtk.134}
\csromantocmarkref{subsection}{ปฏิภานจิตฺตปฏิกฺเขป}{mtk.134}
\bookmark[dest={mtk.134},level=2]{ปฏิภานจิตฺตปฏิกฺเขป}
\csromanheader{2}{\textbf{ปฏิภานจิตฺตปฏิกฺเขป}}
\proseitem{299}{เตน โข ปน สมเยน ฉพฺพคฺคิยา ภิกฺขู วิหาเร ปฏิภานจิตฺตํ การาเปนฺติ อิตฺถิรูปกํ ปุริสรูปกํ.\csromanspacer{} มนุสฺสา วิหารจาริกํ อาหิณฺฑนฺตา ปสฺสิตฺวา อุชฺฌายนฺติ ขิยฺยนฺติ วิปาเจนฺติ “เสยฺยถาปิ คิหี กามโภคิโน”ติ.\csromanspacer{} ภควโต เอตมตฺถํ อาโรเจสุํ.\csromanspacer{} น ภิกฺขเว ปฏิภานจิตฺตํ การาเปตพฺพํ อิตฺถิรูปกํ ปุริสรูปกํ, โย การาเปยฺย, อาปตฺติ ทุกฺกฏสฺส.\csromanspacer{} อนุชานามิ ภิกฺขเว มาลากมฺมํ ลตากมฺมํ มกรทนฺตกํ ปญฺจปฏิกนฺติ.}

\csromanmatikaanchor{mtk.135}
\csromantocmarkref{subsection}{อิฏฺฐกาจยาทิ-อนุชานน}{mtk.135}
\bookmark[dest={mtk.135},level=2]{อิฏฺฐกาจยาทิ-อนุชานน}
\csromanheader{2}{\textbf{อิฏฺฐกาจยาทิ-อนุชานน}}
\proseitem{300}{เตน โข ปน สมเยน วิหารา นีจวตฺถุกา โหนฺติ, อุทเกน โอตฺถริยฺยนฺติ.\csromanspacer{} ภควโต เอตมตฺถํ อาโรเจสุํ.\csromanspacer{} อนุชานามิ ภิกฺขเว อุจฺจวตฺถุกํ กาตุนฺติ.\csromanspacer{} จโย ปริปตติ ฯเปฯ.\csromanspacer{} อนุชานามิ ภิกฺขเว จินิตุํ ตโย จเย อิฏฺฐกาจยํ สิลาจยํ ทารุจยนฺติ.\csromanspacer{} อาโรหนฺตา วิหญฺญนฺติ ฯเปฯ.\csromanspacer{} อนุชานามิ ภิกฺขเว ตโย โสปาเน อิฏฺฐกาโสปานํ สิลาโสปานํ ทารุโสปานนฺติ.\csromanspacer{} อาโรหนฺตา ปริปตนฺติ ฯเปฯ.\csromanspacer{} อนุชานามิ ภิกฺขเว อาลมฺพนพาหนฺติ.}

\prose{เตน โข ปน สมเยน วิหารา อาฬกมนฺทา โหนฺติ, ภิกฺขู หิริยนฺติ นิปชฺชิตุํ.\csromanspacer{} ภควโต เอตมตฺถํ อาโรเจสุํ.\csromanspacer{} อนุชานามิ ภิกฺขเว ติโรกรณินฺติ.\csromanspacer{} ติโรกรณิํ อุกฺขิปิตฺวา โอโลเกนฺติ ฯเปฯ.\csromanspacer{} อนุชานามิ \csromanfolio{298}ภิกฺขเว อฑฺฒกุฏฺฏกนฺติ, อฑฺฒกุฏฺฏกา อุปริโต โอโลเกนฺติ ฯเปฯ.\csromanspacer{} อนุชานามิ ภิกฺขเว ตโย คพฺเภ สิวิกาคพฺภํ นาฬิกาคพฺภํ หมฺมิยคพฺภนฺติ.}

\prose{เตน โข ปน สมเยน ภิกฺขู ขุทฺทเก วิหาเร มชฺเฌ คพฺภํ กโรนฺติ, อุปจาโร น โหติ.\csromanspacer{} ภควโต เอตมตฺถํ อาโรเจสุํ.\csromanspacer{} อนุชานามิ ภิกฺขเว ขุทฺทเก วิหาเร เอกมนฺตํ คพฺภํ กาตุํ, มหลฺลเก มชฺเฌติ.}

\prose{เตน โข ปน สมเยน วิหารสฺส กุฏฺฏปาโท ชีรติ.\csromanspacer{} ภควโต เอตมตฺถํ อาโรเจสุํ.\csromanspacer{} อนุชานามิ ภิกฺขเว กุลงฺกปาทกนฺติ\footnote{กุฬุงฺกปาทกนฺติ (สี)}.\csromanspacer{} วิหารสฺส กุฏฺโฏ โอวสฺสติ ฯเปฯ.\csromanspacer{} อนุชานามิ ภิกฺขเว ปริตฺตาณกิฏิกํ อุทฺทสุธนฺติ\footnote{อุทฺธาสุธนฺติ (สฺยา)}.}

\prose{เตน โข ปน สมเยน อญฺญตรสฺส ภิกฺขุโน ติณจฺฉทนา อหิ ขนฺเธ ปตติ, โส ภีโต วิสฺสรมกาสิ.\csromanspacer{} ภิกฺขู อุปธาวิตฺวา ตํ ภิกฺขุํ เอตทโวจุํ “กิสฺส ตฺวํ อาวุโส วิสฺสรมกาสี”ติ.\csromanspacer{} อถ โข โส ภิกฺขูนํ เอตมตฺถํ อาโรเจสิ.\csromanspacer{} ภิกฺขู ภควโต เอตมตฺถํ อาโรเจสุํ.\csromanspacer{} อนุชานามิ ภิกฺขเว วิตานนฺติ.}

\prose{เตน โข ปน สมเยน ภิกฺขู มญฺจปาเทปิ ปีฐปาเทปิ ถวิกาโย ลคฺเคนฺติ, อุนฺทูเรหิปิ อุปจิกาหิปิ ขชฺชนฺติ.\csromanspacer{} ภควโต เอตมตฺถํ อาโรเจสุํ.\csromanspacer{} อนุชานามิ ภิกฺขเว ภิตฺติขิลํ นาคทนฺตกนฺติ.}

\prose{เตน โข ปน สมเยน ภิกฺขู มญฺเจปิ ปีเฐปิ จีวรํ นิกฺขิปนฺติ, จีวรํ ปริภิชฺชิ.\csromanspacer{} ภควโต เอตมตฺถํ อาโรเจสุํ.\csromanspacer{} อนุชานามิ ภิกฺขเว วิหาเร จีวรวํสํ จีวรรชฺชุนฺติ.}

\prose{เตน โข ปน สมเยน วิหารา อนาฬินฺทกา โหนฺติ อปฺปฏิสฺสรณา.\csromanspacer{} ภควโต เอตมตฺถํ อาโรเจสุํ.\csromanspacer{} อนุชานามิ ภิกฺขเว อาฬินฺทํ ปฆนํ ปกุฏฺฏํ\footnote{ปกุฑฺฑํ (สี)} โอสารกนฺติ.\csromanspacer{} อาฬินฺทา ปากฏา โหนฺติ, ภิกฺขู หิริยนฺติ นิปชฺชิตุํ ฯเปฯ.\csromanspacer{} อนุชานามิ ภิกฺขเว สํสรณกิฏิกํ อุคฺฆาฏนกิฏิกนฺติ.}

\csromanmatikaanchor{mtk.136}
\csromantocmarkref{subsection}{อุปฏฺฐานสาลา-อนุชานน}{mtk.136}
\bookmark[dest={mtk.136},level=2]{อุปฏฺฐานสาลา-อนุชานน}
\csromanheader{2}{\textbf{อุปฏฺฐานสาลา-อนุชานน}}
\proseitem{301}{เตน โข ปน สมเยน ภิกฺขู อชฺโฌกาเส ภตฺตวิสฺสคฺคํ กโรนฺตา สีเตนปิ อุเณฺหนปิ กิลมนฺติ.\csromanspacer{} ภควโต เอตมตฺถํ อาโรเจสุํ.\csromanspacer{} อนุชานามิ ภิกฺขเว อุปฏฺฐานสาลา นีจวตฺถุกา \csromanfolio{299}โหติ, อุทเกน โอตฺถริยฺยติ ฯเปฯ.\csromanspacer{} อนุชานามิ ภิกฺขเว อุจฺจวตฺถุกํ กาตุนฺติ.\csromanspacer{} จโย ปริปตติ ฯเปฯ.\csromanspacer{} อนุชานามิ ภิกฺขเว จินิตุํ ตโย จเย อิฏฺฐกาจยํ สิลาจยํ ทารุจยนฺติ.\csromanspacer{} อาโรหนฺตา วิหญฺญนฺติ ฯเปฯ.\csromanspacer{} อนุชานามิ ภิกฺขเว ตโย โสปาเน อิฏฺฐกาโสปานํ สิลาโสปานํ ทารุโสปานนฺติ.\csromanspacer{} อาโรหนฺตา ปริปตนฺติ ฯเปฯ.\csromanspacer{} อนุชานามิ ภิกฺขเว อาลมฺพนพาหนฺติ.\csromanspacer{} อุปฏฺฐานสาลาย ติณจุณฺณํ ปริปตติ ฯเปฯ.\csromanspacer{} อนุชานามิ ภิกฺขเว โอคุมฺเผตฺวา อุลฺลิตฺตาวลิตฺตํ กาตุํ เสตวณฺณํ กาฬวณฺณํ เครุกปริกมฺมํ มาลากมฺมํ ลตากมฺมํ มกรทนฺตกํ ปญฺจปฏิกํ จีวรวํสํ จีวรรชฺชุนฺติ.}

\prose{เตน โข ปน สมเยน ภิกฺขู อชฺโฌกาเส ฉมาย จีวรํ ปตฺถรนฺติ, จีวรํ ปํสุกิตํ โหติ.\csromanspacer{} ภควโต เอตมตฺถํ อาโรเจสุํ.\csromanspacer{} อนุชานามิ ภิกฺขเว อชฺโฌกาเส จีวรวํสํ จีวรรชฺชุนฺติ.\csromanspacer{} ปานียํ โอตปฺปติ ฯเปฯ.\csromanspacer{} อนุชานามิ ภิกฺขเว ปานิยสาลํ ปานียมณฺฑปนฺติ.\csromanspacer{} ปานียสาลา นีจวตฺถุกา โหติ, อุทเกน โอตฺถริยฺยติ ฯเปฯ.\csromanspacer{} อนุชานามิ ภิกฺขเว อุจฺจวตฺถุกํ กาตุนฺติ.\csromanspacer{} จโย ปริปตติ ฯเปฯ.\csromanspacer{} อนุชานามิ ภิกฺขเว จินิตุํ ตโย จเย อิฏฺฐกาจยํ สิลาจยํ ทารุจยนฺติ.\csromanspacer{} อาโรหนฺตา วิหญฺญนฺติ ฯเปฯ.\csromanspacer{} อนุชานามิ ภิกฺขเว ตโย โสปาเน อิฏฺฐกาโสปานํ สิลาโสปานํ ทารุโสปานนฺติ.\csromanspacer{} อาโรหนฺตา ปริปตนฺติ ฯเปฯ.\csromanspacer{} อนุชานามิ ภิกฺขเว อาลมฺพนพาหนฺติ.\csromanspacer{} ปานียสาลาย ติณจุณฺณํ ปริปตติ ฯเปฯ.\csromanspacer{} อนุชานามิ ภิกฺขเว โอคุมฺเผตฺวา อุลฺลิตฺตาวลิตฺตํ กาตุํ เสตวณฺณํ กาฬวณฺณํ เครุกปริกมฺมํ มาลากมฺมํ ลตากมฺมํ มกรทนฺตกํ ปญฺจปฏิกํ จีวรวํสํ จีวรรชฺชุนฺติ.\csromanspacer{} ปานียภาชนํ น สํวิชฺชติ ฯเปฯ.\csromanspacer{} อนุชานามิ ภิกฺขเว ปานิยสงฺขํ ปานิยสราวกนฺติ.}

\csromanmatikaanchor{mtk.137}
\csromantocmarkref{subsection}{ปาการาทิ-อนุชานน}{mtk.137}
\bookmark[dest={mtk.137},level=2]{ปาการาทิ-อนุชานน}
\csromanheader{2}{\textbf{ปาการาทิ-อนุชานน}}
\proseitem{302}{เตน โข ปน สมเยน วิหารา อปริกฺขิตฺตา หนฺติ.\csromanspacer{} ภควโต เอตมตฺถํ อาโรเจสุํ.\csromanspacer{} อนุชานามิ ภิกฺขเว ปริกฺขิปิตุํ ตโย ปากาเร อิฏฺฐกาปาการํ สิลาปาการํ ทารุปาการนฺติ.\csromanspacer{} โกฏฺฐโก น โหติ ฯเปฯ.\csromanspacer{} อนุชานามิ ภิกฺขเว โกฏฺฐกนฺติ.\csromanspacer{} โกฏฺฐโก นีจวตฺถุโก โหติ, อุทเกน โอตฺถริยฺยติ ฯเปฯ.\csromanspacer{} อนุชานามิ ภิกฺขเว อุจฺจวตฺถุกํ กาตุนฺติ.\csromanspacer{} โกฏฺฐกสฺส กวาฏํ น โหติ ฯเปฯ.\csromanspacer{} อนุชานามิ ภิกฺขเว กวาฏํ ปิฏฺฐสํฆาฏํ อุทุกฺขลิกํ อุตฺตรปาสกํ อคฺคฬวฏฺฏิํ กปิสีสกํ สูจิกํ ฆฏิกํ \csromanfolio{300}ตาฬจฺฉิทฺทํ อาวิญฺฉนจฺฉิทํ อาวิญฺฉนรชฺชุนฺติ.\csromanspacer{} โกฏฺฐเก ติณจุณฺณํ ปริปตติ ฯเปฯ.\csromanspacer{} อนุชานามิ ภิกฺขเว โอคุมฺเผตฺวา อุลฺลิตฺตาวลิตฺตํ กาตุํ เสตวณฺณํ กาฬวณฺณํ เครุกปริกมฺมํ มาลากมฺมํ ลตากมฺมํ มกรทนฺตกํ ปญฺจปฏิกนฺติ.}

\prose{เตน โข ปน สมเยน ปริเวณํ จิกฺขลฺลํ โหติ.\csromanspacer{} ภควโต เอตมตฺถํ อาโรเจสุํ.\csromanspacer{} อนุชานามิ ภิกฺขเว มรุมฺพํ อุปกิริตุนฺติ.\csromanspacer{} น ปริยาปุณนฺติ ฯเปฯ.\csromanspacer{} อนุชานามิ ภิกฺขเว ปทรสิลํ นิกฺขิปิตุนฺติ.\csromanspacer{} อุทกํ สนฺติฏฺฐติ ฯเปฯ.\csromanspacer{} อนุชานามิ ภิกฺขเว อุทกนิทฺธมนนฺติ.}

\prose{เตน โข ปน สมเยน ภิกฺขู ปริเวเณ ตหํ ตหํ อคฺคิฏฺฐานํ กโรนฺติ, ปริเวณํ อุกฺลาปํ โหติ.\csromanspacer{} ภควโต เอตมตฺถํ อาโรเจสุํ.\csromanspacer{} อนุชานามิ ภิกฺขเว เอกมนฺตํ อคฺคิสาลํ กาตุนฺติ.\csromanspacer{} อคฺคิสาลา นีจวตฺถุกา โหติ, อุทเกน โอตฺถริยฺยติ ฯเปฯ.\csromanspacer{} อนุชานามิ ภิกฺขเว อุจฺจวตฺถุกํ กาตุนฺติ.\csromanspacer{} จโย ปริปตติ ฯเปฯ.\csromanspacer{} อนุชานามิ ภิกฺขเว จินิตุํ ตโย จเย อิฏฺฐกาจยํ สิลาจยํ ทารุจยนฺติ.\csromanspacer{} อาโรหนฺตา วิหญฺญนฺติ ฯเปฯ.\csromanspacer{} อนุชานามิ ภิกฺขเว ตโย โสปาเน อิฏฺฐกาโสปานํ สิลาโสปานํ ทารุโสปานนฺติ.\csromanspacer{} อาโรหนฺตา ปริปตนฺติ ฯเปฯ.\csromanspacer{} อนุชานามิ ภิกฺขเว อาลมฺพนพาหนฺติ.\csromanspacer{} อคฺคิสาลาย กวาฏํ น โหติ ฯเปฯ.\csromanspacer{} อนุชานามิ ภิกฺขเว กวาฏํ ปิฏฺฐสํฆาฏํ อุทุกฺขลิกํ อุตฺตรปาสกํ อคฺคฬวฏฺฏิํ กปิสีสกํ สูจิกํ ฆฏิกํ ตาฬจฺฉิทฺทํ อาวิญฺฉนจฺฉิทฺทํ อาวิญฺฉนรชฺชุนฺติ.\csromanspacer{} อคฺคิสาลาย ติณจุณฺณํ ปริปตติ ฯเปฯ.\csromanspacer{} อนุชานามิ ภิกฺขเว โอคุมฺเผตฺวา อุลฺลิตฺตาวลิตฺตํ กาตุํ เสตวณฺณํ กาฬวณฺณํ เครุกปริกมฺมํ มาลากมฺมํ ลตากมฺมํ มกรทนฺตกํ ปญฺจปฏิกํ ปญฺจปฏิกํ จีวรวํสํ จีวรรชฺชุนฺติ.}

\csromanmatikaanchor{mtk.138}
\csromantocmarkref{subsection}{อารามปริกฺเขป-อนุชานน}{mtk.138}
\bookmark[dest={mtk.138},level=2]{อารามปริกฺเขป-อนุชานน}
\csromanheader{2}{\textbf{อารามปริกฺเขป-อนุชานน}}
\proseitem{303}{เตน โข ปน สมเยน อาราโม อปริกฺขิตฺโต โหติ, อชกาปิ ปสุกาปิ อุปโรเป วิเหเฐนฺติ.\csromanspacer{} ภควโต เอตมตฺถํ อาโรเจสุํ.\csromanspacer{} อนุชานามิ ภิกฺขเว ปริกฺขิปิตุํ ตโย วาเฏ เวฬุวาฏํ กณฺฑกวาฏํ\footnote{วเฏ เวฬุวฏํ กณฺฑกวฏํ (สฺยา)} ปริกฺขนฺติ.\csromanspacer{} โกฏฺฐโก น โหติ.\csromanspacer{} ตเถว อชกาปิ ปสุกาปิ อุปโรเป วิเหเฐนฺติ ฯเปฯ.\csromanspacer{} อนุชานามิ ภิกฺขเว โกฏฺฐกํ อเปสิํ ยมกกวาฏํ โตรณํ ปลิฆนฺติ.\csromanspacer{} โกฏฺฐเก ติณจุณฺณํ ปริปตติ ฯเปฯ. \csromanfolio{301}อนุชานามิ ภิกฺขเว โอคุมฺเผตฺวา อุลฺลิตฺตาวลิตฺตํ กาตุํ เสตวณฺณํ กาฬวณฺณํ เครุกปริกมฺมํ มาลากมฺมํ ลตากมฺมํ มกรทนฺตกํ ปญฺจปฏิกนฺติ.\csromanspacer{} อาราโม จิกฺขลฺโล โหติ ฯเปฯ.\csromanspacer{} อนุชานามิ ภิกฺขเว มรุมฺพํ อุปกิริตุนฺติ.\csromanspacer{} น ปริยาปุณนฺติ ฯเปฯ.\csromanspacer{} อนุชานามิ ภิกฺขเว ปทรสิลํ นิกฺขิปิตุนฺติ.\csromanspacer{} อุทกํ สนฺติฏฺฐติ ฯเปฯ.\csromanspacer{} อนุชานามิ ภิกฺขเว อุทกนิทฺธมนนฺติ.}

\prose{เตน โข ปน สมเยน ราชา มาคโธ เสนิโย พิมฺพิสาโร สํฆสฺส อตฺถาย สุธามตฺติกาเลปนํ ปาสาทํ กาเรตุกาโม โหติ.\csromanspacer{} อถ โข ภิกฺขูนํ เอตทโหสิ “กิํ นุ โข ภควตา ฉทนํ อนุญฺญาตํ, กิํ อนนุญฺญาตนฺ”ติ.\csromanspacer{} ภควโต เอตมตฺถํ อาโรเจสุํ.\csromanspacer{} อนุชานามิ ภิกฺขเว ปญฺจ ฉทนานิ อิฏฺฐกาฉทนํ สิลาฉทนํ สุธาฉทนํ ติณจฺฉทนํ ปณฺณจฺฉทนนฺติ.}

\nitthitam{ปฐมภาณวาโร นิฏฺฐิโต.}
\csromanmatikaanchor{mtk.139}
\csromantocmarkref{section}{2. ทุติยภาณวาร}{mtk.139}
\bookmark[dest={mtk.139},level=1]{2. ทุติยภาณวาร}
\csromanheader{1}{2. ทุติยภาณวาร}
\csromanmatikaanchor{mtk.140}
\csromantocmarkref{section}{อนาถปิณฺฑิกวตฺถุ}{mtk.140}
\bookmark[dest={mtk.140},level=1]{อนาถปิณฺฑิกวตฺถุ}
\csromanheader{1}{\textbf{อนาถปิณฺฑิกวตฺถุ}}
\proseitem{304}{เตน โข ปน สมเยน อนาถปิณฺฑิโก คหปติ ราชคหกสฺส เสฏฺฐิสฺส ภคินิปติโก โหติ.\csromanspacer{} อถ โข อนาถปิณฺฑิโก คหปติ ราชคหํ อคมาสิ เกนจิเทว กรณีเยน.\csromanspacer{} เตน โข ปน สมเยน ราชคหเกน เสฏฺฐินา สฺวาตนาย พุทฺธปฺปมุโขสํโฆ นิมนฺติโต โหติ.\csromanspacer{} อถ โข ราชคหโก เสฏฺฐี ทาเส จ กมฺมกาเร\footnote{กมฺมกเร (สี, สฺยา)} จ อาณาเปสิ “เตน หิ ภเณ กาลสฺเสว อุฏฺฐาย ยาคุโย ปจถ, ภตฺตานิ ปจถ, สูปานิ สมฺปาเทถ, อุตฺตริภงฺคานิ สมฺปาเทถา”ติ.\csromanspacer{} อถ โข อนาถปิณฺฑิกสฺส คหปติสฺส เอตทโหสิ “ปุพฺเพ ขฺวายํ คหปติ มยิ อาคเต สพฺพกิจฺจานิ นิกฺขิปิตฺวา มมญฺเญว สทฺธิํ ปฏิสมฺโมทติ, โสทานายํ วิกฺขิตฺตรูโป ทาเส จ กมฺมกาเร จ อาณาเปสิ ‘เตน หิ ภเณ กาลสฺเสว อุฏฺฐาย ยาคุโย ปจถ, ภตฺตานิ ปจถ, สูปานิ สมฺปาเทถ, อุตฺตริภงฺคานิ สมฺปาเทถา’ติ.\csromanspacer{} กิํ นุ โข อิมสฺส คหปติสฺส อาวาโห วา ภวิสฺสกิ, วิวาโห วา ภวิสฺสติ, มหายญฺโญ วา \csromanfolio{302}ปจฺจุปฏฺฐิโต, ราชา วา มาคโธ เสนิโย พิมฺพิสาโร นิมนฺติโต สฺวาตนาย สทฺธิํ พลกาเยนา”ติ.}

\prose{อถ โข ราชคหโก เสฏฺฐี ทาเส จ กมฺมกาเร จ อาณาเปตฺวา เยน อนาถปิณฺฑิโก คหปติ เตนุปสงฺกมิ, อุปสงฺกมิตฺวา อนาถปิณฺฑิเกน คหปตินา สทฺธิํ ปฏิสมฺโมทิตฺวา เอกมนฺตํ นิสีทิ.\csromanspacer{} เอกมนฺตํ นิสินฺนํ โข ราชคหกํ เสฏฺฐิํ อนาถปิณฺฑิโก คหปติ เอตทโวจ “ปุพฺเพ โข ตฺวํ คหปติ มยิ อาคเต สพฺพกิจฺจานิ นิกฺขิปิตฺวา มมญฺเญว สทฺธิํ ปฏิสมฺโมทสิ, โสทานิ ตฺวํ วิกฺขิตฺตรูโป ทาเส จ กมฺมกาเร จ อาณาเปสิ ‘เตน หิ ภเณ กาลสฺเสว อุฏฺฐาย ยาคุโย ปจถ, ภตฺตานิ ปจถ, สูปานิ สมฺปาเทถ, อุตฺตริภงฺคานิ สมฺปาเทถา’ติ.\csromanspacer{} กิํ นุ โข เต คหปติ อาวาโห วา ภวิสฺสติ, วิวาโห วา ภวิสฺสติ, มหายญฺโญ วา ปจฺจุปฏฺฐิโต, ราชา วา มาคโธ เสนิโย พิมฺพิสาโร นิมนฺติโต สฺวาตนาย สทฺธิํ พลกาเยนา”ติ.\csromanspacer{} น เม คหปติ อาวาโห วา ภวิสฺสติ, นาปิ วิวาโห วา ภวิสฺสติ.\csromanspacer{} นาปิ ราชา วา มาคโธ เสนิโย พิมฺพิสาโร นิมนฺติโต สฺวาตนาย สทฺธิํ พลกาเยน.\csromanspacer{} อปิ จ เม มหายญฺโญ ปจฺจุปฏฺฐิโต สฺวาตนาย พุทฺธปฺปมุโข สํโฆ นิมนฺติโตติ. “พุทฺโธ”ติ ตฺวํ คหปติ วเทสีติ. “พุทฺโธ”ตฺยาหํ คหปติ วทามีติ. “พุทฺโธ”ติ ตฺวํ คหปติ วเทสีติ. “พุทฺโธ”ตฺยาหํ คหปติ วทามีติ. “พุทฺโธ”ติ ตฺวํ คหปติ วเทสีติ. “พุทฺโธ”ตฺยาหํ คหปติ วทามีติ.\csromanspacer{} โฆโสปิ โข เอโส คหปติ ทุลฺลโภ โลกสฺมิํ, ยทิทํ “พุทฺโธ พุทฺโธ”ติ.\csromanspacer{} สกฺกา นุ โข คหปติ อิมํ กาลํ ตํ ภควนฺตํ ทสฺสนาย อุปสงฺกมิตุํ อรหนฺตํ สมฺมาสมฺพุทฺธนฺติ.\csromanspacer{} อกาโล โข คหปติ อิมํ กาลํ ตํ ภควนฺตํ ทสฺสนาย อุปสงฺกมิตุํ อรหนฺตํ สมฺมาสมฺพุทฺธํ, เสฺวทานิ ตฺวํ กาเลน ตํ ภควนฺตํ ทสฺสนาย อุปสงฺกมิสฺสสิ อรหนฺตํ สมฺมาสมฺพุทฺธนฺติ.\csromanspacer{} อถ โข อนาถปิณฺฑิโก คหปติ “เสฺวทานาหํ กาเลน ตํ ภควนฺตํ ทสฺสนาย อุปสงฺกมิสฺสามิ อรหนฺตํ สมฺมาสมฺพุทฺธนฺ”ติ พุทฺธคตาย สติยา นิปชฺชิตฺวา รตฺติยา สุทฺทํ ติกฺขตฺตุํ วุฏฺฐาสิ ปภาตํ มญฺญมาโน.}

\proseitem{305}{\csromansymbolfootnote{*}{สํ 1. 213 ปิฏฺเฐปิ.}อถ โข อนาถปิณฺฑิโก คหปติ เยน สิวกทฺวารํ\footnote{สีวทฺวารํ (สี), สีตวนทฺวารํ (สฺยา)} เตนุปสงฺกมิ.\csromanspacer{} อมนุสฺสา ทฺวารํ วิวริํสุ.\csromanspacer{} อถ โข อนาถปิณฺฑิกสฺส \csromanfolio{303}คหปติสฺส นครมฺหา นิกฺขนฺตสฺส อาโลโก อนฺตรธายิ, อนฺธกาโร ปาตุรโหสิ, ภยํ ฉมฺภิตตฺตํ โลมหํโส อุทปาทิ, ตโตว ปุน นิวตฺติตุกาโม อโหสิ.\csromanspacer{} อถ โข สิวโก\footnote{สีวโก (สี, สฺยา)} ยกฺโข อนฺตรหิโต สทฺทมนุสฺสเวสิ\csromandash{}}

\csromangathameasure{“สตํ หตฺถี สตํ อสฺสา, สตํ อสฺสตรีรถา. \\ สตํ กญฺญาสหสฺสานิ, อามุกฺกมณิกุณฺฑลา. \\ เอกสฺส ปทวีติหารสฺส, \\ กลํ นาคฺฆนฺติ โสฬสิํ. \\ อภิกฺกม คหปติ อภิกฺกม คหปติ, \\ อภิกฺกนฺตํ เต เสยฺโย โน ปฏิกฺกนฺตนฺ”ติ.}
\csromangathagroup{\csromangathastanza{“สตํ หตฺถี สตํ อสฺสา, สตํ อสฺสตรีรถา. \\ สตํ กญฺญาสหสฺสานิ, อามุกฺกมณิกุณฺฑลา.}\csromangathastanza{เอกสฺส ปทวีติหารสฺส, \\ กลํ นาคฺฆนฺติ โสฬสิํ\footnote{โสฬสินฺติ (สี, ก)}. \\ อภิกฺกม คหปติ อภิกฺกม คหปติ, \\ อภิกฺกนฺตํ เต เสยฺโย โน ปฏิกฺกนฺตนฺ”ติ.}}

\prose{อถ โข อนาถปิณฺฑิกสฺส คหปติสฺส อนฺธกาโร อนฺตรธายิ, อาโลโก ปาตุรโหสิ, ยํ อโหสิ ภยํ ฉมฺภิตตฺตํ โลมหํโส, โส ปฏิปฺปสฺสมฺภิ.\csromanspacer{} ทุติยมฺปิ โข ฯเปฯ.\csromanspacer{} ตติยมฺปิ โข อนาถปิณฺฑิกสฺส คหปติสฺส อาโลโก อนฺตรธายิ, อนฺธกาโร ปาตุรโหสิ, ภยํ ฉมฺภิตตฺตํ โลมหํโส อุทปาทิ, ตโตว ปุน นิวตฺติตุกาโม อโหสิ.\csromanspacer{} ตติยมฺปิ โข สิวโก ยกฺโข อนฺตรหิโต สทฺทมนุสฺสาเวสิ\csromandash{}}

\csromangathameasure{“สตํ หตฺถี สตํ อสฺสา, สตํ อสฺสตรีรถา. \\ สตํ กญฺญาสหสฺสานิ, อามุกฺกมณิกุณฺฑลา. \\ เอกสฺส ปทวีติหารสฺส, \\ กลํ นาคฺฆนฺติ โสฬสิํ. \\ อภิกฺกม คหปติ อภิกฺกม คหปติ, \\ อภิกฺกนฺตํ เต เสยฺโย โน ปฏิกฺกนฺตนฺ”ติ.}
\csromangathagroup{\csromangathastanza{“สตํ หตฺถี สตํ อสฺสา, สตํ อสฺสตรีรถา. \\ สตํ กญฺญาสหสฺสานิ, อามุกฺกมณิกุณฺฑลา.}\csromangathastanza{เอกสฺส ปทวีติหารสฺส, \\ กลํ นาคฺฆนฺติ โสฬสิํ. \\ อภิกฺกม คหปติ อภิกฺกม คหปติ, \\ อภิกฺกนฺตํ เต เสยฺโย โน ปฏิกฺกนฺตนฺ”ติ.}}

\prose{ตติยมฺปิ โข อนาถปิณฺฑิกสฺส คหปติสฺส อนฺธกาโร อนฺตรธายิ, อาโลโก ปาตุรโหสิ, ยํ อโหสิ ภยํ ฉมฺภิตตฺตํ โลมหํโส, โส ปฏิปฺปสฺสมฺภิ.\csromanspacer{} อถ โข อนาถปิณฺฑิโก คหปติ เยน สีตวนํ เตนุปสงฺกมิ.\csromanspacer{} เตน โข ปน สมเยน ภควา รตฺติยา ปจฺจูสสมยํ ปจฺจุฏฺฐาย อชฺโฌกาเส จงฺกมติ.\csromanspacer{} อทฺทสา โข ภควา อนาถปิณฺฑิกํ คหปติํ ทูรโตว อาคจฺฉนฺตํ, ทิสฺวาน จงฺกมา โอโรหิตฺวา ปญฺญตฺเต อาสเน นิสีทิ, นิสชฺช โข ภควา อนาถปิณฺฑิกํ \csromanfolio{304}คหปติํ เอตทโวจ “เอหิ สุทตฺตา”ติ.\csromanspacer{} อถ โข อนาถปิณฺฑิโก คหปติ “นาเมน มํ ภควา อาลปตี”ติ หฏฺโฐ อุทคฺโค เยน ภควา เตนุปสงฺกมิ, อุปสงฺกมิตฺวา ภควโต ปาเทสุ สิรสา นิปติตฺวา ภควนฺตํ เอตทโวจ ‘กจฺจิ ภนฺเต ภควา สุขํ สยิตฺถา”ติ.}

\csromangathameasure{*สพฺพทา เว สุขํ เสติ, พฺราหฺมโณ ปรินิพฺพุโต. \\ โย น ลิมฺปติ กาเมสุ, สีติภูโต นิรูปธิ. \\ สพฺพา อาสตฺติโย เฉตฺวา, วิเนยฺย หทเย ทรํ. \\ อุปสนฺโต สุขํ เสติ, สนฺติํ ปปฺปุยฺย เจตสาติ.}
\csromangathagroup{\csromangathastanza{\csromansymbolfootnote{*}{สํ 1. 214 ปิฏฺเฐปิ.}สพฺพทา เว สุขํ เสติ, พฺราหฺมโณ ปรินิพฺพุโต. \\ โย น ลิมฺปติ กาเมสุ, สีติภูโต นิรูปธิ.}\csromangathastanza{สพฺพา อาสตฺติโย เฉตฺวา, วิเนยฺย หทเย ทรํ. \\ อุปสนฺโต สุขํ เสติ, สนฺติํ ปปฺปุยฺย เจตสาติ\footnote{เจตโสติ (สี, สฺยา)}.}}

\prose{อถ โข ภควา อนาถปิณฺฑิกสฺส คหปติสฺส อนุปุพฺพิํ กถํ\footnote{อานุปุพฺพิกถํ (สี)} กเถสิ, เสยฺยถิทํ, ทานกถํ สีลกถํ สคฺคกถํ กามานํ อาทีนวํ โอการํ สํกิเลสํ เนกฺขมฺเม อานิสํสํ ปกาเสสิ.\csromanspacer{} ยทา ภควา อญฺญาสิ อนาถปิณฺฑิกํ คหปติํ กลฺลจิตฺตํ มุทุจิตฺตํ วินีวรณจิตฺตํ อุทคฺคจิตฺตํ ปสนฺนจิตฺตํ, อถ ยา พุทฺธานํ สามุกฺกํสิกา ธมฺมเทสนา, ตํ ปกาเสสิ ทุกฺขํ สมุทยํ นิโรธํ มคฺคํ, เสยฺยถาปิ นาม สุทฺธํ วตฺถํ อปคตกาฬกํ สมฺมเทว รชนํ ปฏิคฺคเณฺหยฺย, เอวเมว อนาถปิณฺฑิกสฺส คหปติสฺส ตสฺมิํเยว อาสเน วิรชํ วีตมลํ ธมฺมจกฺขุํ อุทปาทิ “ยํ กิญฺจิ สมุทยธมฺมํ สพฺพํ ตํ นิโรธธมฺมนฺ”ติ.\csromanspacer{} อถ โข อนาถปิณฺฑิโก คหปติ ทิฏฺฐธมฺโม ปตฺตธมฺโม วิทิตธมฺโม ปริโยคาฬฺหธมฺโม ติณฺณวิจิกิจฺโฉ วิคตกถํกโถ เวสารชฺชปฺปตฺโต อปรปฺปจฺจโย สตฺถุสาสเน ภควนฺตํ เอตทโวจ “อภิกฺกนฺตํ ภนฺเต, อภิกฺกนฺตํ ภนฺเต, เสยฺยถาปิ ภนฺเต นิกฺกุชฺชิตํ วา อุกฺกุชฺเชยฺย, ปฏิจฺฉนฺนํ วา วิวเรยฺย, มูฬฺหสฺส วา มคฺคํ อาจิกฺเขยฺย, อนฺธกาเร วา เตลปชฺโชตํ ธาเรยฺย ‘จกฺขุมนฺโต รูปานิ ทกฺขนฺตี’ติ, เอวเมวํ ภควตา อเนกปริยาเยน ธมฺโม ปกาสิโต.\csromanspacer{} เอสาหํ ภนฺเต ภควนฺตํ สรณํ คจฺฉามิ ธมฺมญฺจ ภิกฺขุสํฆญฺจ, อุปาสกํ มํ ภควา ธาเรตุ อชฺชตคฺเค ปาณุเปตํ สรณํ คตํ, อธิวาเสตุ จ เม ภนฺเต ภควา สฺวาตนาย ภตฺตํ สทฺธิํ ภิกฺขุสํเฆนาติ.\csromanspacer{} อธิวาเสสิ ภควา ตุณฺหีภาเวน.}

\prose{อถ โข อนาถปิณฺฑิโก คหปติ ภควโต อธิวาสนํ วิทิตฺวา อุฏฺฐายาสนา ภควนฺตํ อภิวาเทตฺวา ปทกฺขิณํ กตฺวา ปกฺกามิ.\csromanspacer{} อสฺโสสิ \csromanfolio{305}โข ราชคหโก เสฏฺฐี “อนาถปิณฺฑิเกน กิร คหปตินา สฺวาตนาย พุทฺธปฺปมุโข สํโฆ นิมนฺติโต”ติ.}

\proseitem{306}{อถ โข ราชคหโก เสฏฺฐี อนาถปิณฺฑิกํคหปติํ เอตทโวจ “ตยา กิร คหปติ สฺวาตนาย พุทฺธปฺปมุโข สํโฆ นิมนฺติโต, ตฺวญฺจาสิ อาคนฺตุโก, เทมิ เต คหปติ, เวยฺยายิกํ, เยน ตฺวํ พุทฺธปฺปมุขสฺส สํฆสฺส ภตฺตํ กเรยฺยาสี”ติ.\csromanspacer{} อลํ คหปติ, อตฺถิ เม เวยฺยายิกํ, เยนาหํ พุทฺธปฺปมุขสฺส สํฆสฺส ภตฺตํ กริสฺสามีติ.}

\prose{อสฺโสสิ โข ราชคหโก เนคโม “อนาถปิณฺฑิเกน กิร คหปตินา สฺวาตนาย พุทฺธปฺปมุโข สํโฆ นิมนฺติโต”ติ.\csromanspacer{} อถ โข ราชคหโก เนคโม อนาถปิณฺฑิกํ คหปติํ เอตทโวจ “ตยา กิร คหปติ สฺวาตนาย พุทฺธปฺปมุโข สํโฆ นิมนฺติโต, ตฺวญฺจาสิ อาคนฺตุโก, เทมิ เก คหปติ เวยฺยายิกํ, เยน ตฺวํ พุทฺธปฺปมุขสฺส สํฆสฺส ภตฺตํ กเรยฺยาสี”ติ.\csromanspacer{} อลํ อยฺย, อตฺถิ เม เวยฺยายิกํ, เยนาหํ พุทฺธปฺปมุขสฺส สํฆสฺส ภตฺตํ กริสฺสามีติ.}

\prose{อสฺโสสิ โขราชา มาคโธ เสนิโย พิมฺพิสาโร “อนาถปิณฺฑิเกน กิร คหปตินา สฺวาตนาย พุทฺธปฺปมุโข สํโฆ นิมนฺติโต”ติ.\csromanspacer{} อถ โข ราชา มาคโธ เสนิโย พิมฺพิสาโร อนาถปิณฺฑิกํ คหปติํ เอตทโวจ “ตยา กิร คหปติ สฺวาตนาย พุทฺธปฺปมุโข สํโฆ นิมนฺติโต, ตฺวญฺจาสิ อาคนฺตุโก, เทมิ เต คหปติ เวยฺยายิกํ, เยน ตฺวํ พุทฺธปฺปมุขสฺส สํฆสฺส ภตฺตํ กเรยฺยาสี”ติ.\csromanspacer{} อลํ เทว, อตฺถิ เม เวยฺยายิกํ, เยนาหํ พุทฺธปฺปมุขสฺส สํฆสฺส ภตฺตํ กริสฺสมีติ.}

\prose{อถ โข อนาถปิณฺฑิโก คหปติ ตสฺสา รตฺติยา อจฺจเยน ราชคหกสฺส เสฏฺฐิสฺส นิเวสเน ปณีตํ ขาทนียํ โภชนิยํ ปฏิยาทาเปตฺวา ภควโต กาลํ อาโรจาเปสิ “กาโล ภนฺเต นิฏฺฐิตํ ภตฺตนฺ”ติ.\csromanspacer{} อถ โข ภควา ปุพฺพณฺหสมยํ นิวาเสตฺวา ปตฺตจีวรมาทาย เยน ราชคหกสฺส เสฏฺฐิสฺส นิเวสนํ เตนุปสงฺกมิ, อุปสงฺกมิตฺวา ปญฺญตฺเต อาสเน นิสีทิ สทฺธิํ ภิกฺขุสํเฆน.\csromanspacer{} อถ โข อนาถปิณฺฑิโก คหปติ พุทฺธปฺปมุขํ ภิกฺขุสํฆํ ปณีเตน ขาทนีเยน โภชนีเยน สหตฺถา สนฺตปฺเปตฺวา สมฺปวาเรตฺวา ภควนฺตํ ภุตฺตาวิํ โอนีตปตฺตวาณิํ เอกมนฺตํ \csromanfolio{306}นิสีทิ, เอกมนฺตํ นิสินฺโน โข อนาถปิณฺฑิโก คหปติ ภควนฺตํ เอตทโวจ “อธิวาเสตุ เม ภนฺเต ภควา สาวตฺถิยํ วสฺสาวาสํ สทฺธิํ ภิกฺขุสํเฆนา”ติ.\csromanspacer{} สุญฺญาคาเร โข คหปติ ตถาคตา อภิรมนฺตีติ.\csromanspacer{} อญฺญาตํ ภควา, อญฺญาตํ สุคตาติ.\csromanspacer{} อถ โข ภควา อนาถปิณฺฑิกํ คหปติํ ธมฺมิยา กถาย สนฺทสฺเสตฺวา สมาทเปตฺวา สมุตฺเตเชตฺวา สมฺปหํเสตฺวา อุฏฺฐายาสนา ปกฺกามิ.}

\proseitem{307}{เตน โข ปน สมเยน อนาถปิณฺฑิโก คหปติ พหุมิตฺโต โหติ พหุสหาโย อาเทยฺยวาโจ.\csromanspacer{} อถ โข อนาถปิณฺฑิโก คหปติ ราชคเห ตํ กรณียํ ตีเรตฺวา เยน สาวตฺถิ เตน ปกฺกามิ.\csromanspacer{} อถ โข อนาถปิณฺฑิโก คหปติ อนฺตรามคฺเค มนุสฺเส อาณาเปสิ “อาราเม อยฺยา กโรถ, วิหาเร ปติฏฺฐาเปถ, ทานานิ ปฏฺฐเปถ, พุทฺโธ โลเก อุปฺปนฺโน, โส จ มยา ภควา นิมนฺติโต อิมินา มคฺเคน อาคจฺฉิสฺสตี”ติ.\csromanspacer{} อถ โข เต มนุสฺสา อนาถปิณฺฑิเกน คหปตินา อุยฺโยชิตา อาราเม อกํสุ, วิหาเร ปติฏฺฐาเปสุํ, ทานานิ ปฏฺฐเปสุํ.}

\prose{อถ โข นาถปิณฺฑิโก คหปติ สาวตฺถิํ คนฺตา สมนฺตา สาวตฺถิํ อนุวิโลเกสิ “กตฺถ นุ โข ภควา วิหเรยฺย, ยํ อสฺส คามโต เนว อติทูเร น อจฺจาสนฺเน คมนาคมนสมฺปนฺนํ อตฺถิกานํ อตฺถิกานํ มนุสฺสานํ อภิกฺกมนียํ ทิวา อปฺปากิณฺณํ รตฺติํ อปฺปสทฺทํ อปฺปนิคฺโฆสํ วิชนวาตํ มนุสฺสราหสฺเสยฺยกํ ปฏิสลฺลานสารุปฺปนฺ”ติ.}

\prose{อทฺทสา โข อนาถปิณฺฑิโก คหปติ เชตสฺส กุมารสฺส\footnote{ราชกุมารสฺส (สี, สฺยา, กํ)} อุยฺยานํ คามโต เนว อติทูเร น อจฺจาสนฺเน คมนาคมนสมฺปนฺนํ อตฺถิกานํ อตฺถิกานํ มนุสฺสานํ อภิกฺกมนียํ ทิวา อปฺปากิณฺณํ รตฺติํ อปฺปสทฺทํ อปฺปนิคฺโฆสํ วิชนวาตํ มนุสฺสราหสฺเสยฺยกํ ปฏิสลฺลานสารุปฺปํ, ทิสฺวาน เยน เชโต กุมาโร เตนุปสงฺกมิ, อุปสงฺกมิตฺวา เชตํ กุมารํ เอตทโวจ “เทหิ เม อยฺยปุตฺต อุยฺยานํ อารามํ กาตุนฺ”ติ\footnote{เกตุํ (วชีรพุทฺธิฏีกายํ)}.\csromanspacer{} อเทยฺโย คหปติ อาราโม อปิ โกฏิสนฺถเรนาติ.\csromanspacer{} คหิโต อยฺยปุตฺต อาราโมติ.\csromanspacer{} น คหปติ คหิโต อาราโมติ. “คหิโต น คหิโต”ติ โวหาริเก มหามตฺเต ปุจฺฉิํสุ.\csromanspacer{} มหามตฺตา \csromanfolio{307}เอวมาหํสุ “ยโต ตยา อยฺยปุตฺต อคฺโฆ กโต, คหิโต อาราโม”ติ.\csromanspacer{} อถ โข อนาถปิณฺฑิโก คหปติ สกเฏหิ หิรญฺญํ นิพฺพาหาเปตฺวา เชตวนํ โกฏิสนฺถรํ สนฺถราเปสิ, สกิํ นีหฏํ หิรญฺญํ โถกสฺส โอกาสสฺส โกฏฺฐกสามนฺตา นปฺปโหติ.\csromanspacer{} อถ โข อนาถปิณฺฑิโก คหปติ มนุสฺเส อาณาเปสิ “คจฺฉถ ภเณ, หิรญฺญํ อาหรถ, อิมํ โอกาสํ สนฺถริสฺสามา”ติ.}

\prose{อถ โข เชตสฺส กุมารสฺส เอตทโหสิ “น โข อิทํ โอรกํ ภวิสฺสติ, ยถายํ คหปติ ตาว พหุํ หิรญฺญํ ปริจฺจชตี”ติ อนาถปิณฺฑิกํ คหปติํ เอตทโวจ “อลํ คหปติ, มา ตํ โอกาสํ สนฺถราเปสิ, เทหิ เม เอตํ โอกาสํ, มเมตํ ทานํ ภวิสฺสตี”ติ.\csromanspacer{} อถ โข อนาถปิณฺฑิโก คหปติ “อยํ โข เชโต กุมาโร อภิญฺญาโก ญาตมนุสฺโส, มหตฺถิโก โข ปน เอวรูปานํ ญาตมนุสฺสานํ อิมสฺมิํ ธมฺมวินเย ปสาโท”ติ ตํ โอกาสํ เชตสฺส กุมารสฺส ปาทาสิ\footnote{อทาสิ (สฺยา)}.\csromanspacer{} อถ โข เชโต กุมาโร ตสฺมิํ โอกาเส โกฏฺฐกํ มาเปสิ. อถ โข อนาถปิณฺฑิโก คหปติ เชตวเน วิหาเร การาเปสิ, ปริเวณานิ การาเปสิ, โกฏฺฐเก การาเปสิ, อุปฏฺฐานสาลาโย การาเปสิ, อคฺคิสาลาโย การาเปสิ, กปฺปิยกุฏิโย การาเปสิ, วจฺจกุฏิโย การาเปสิ, จงฺกเม การาเปสิ, จงฺกมนสาลาโย การาเปสิ, อุทปาเน การาเปสิ, อุทปานสาลาโย การาเปสิ, ชนฺตาฆเร การาเปสิ, ชนฺตาฆรสาลาโย การาเปสิ, โปกฺขรณิโย การาเปสิ, มณฺฑเป การาเปสิ.}

\csromanmatikaanchor{mtk.141}
\csromantocmarkref{subsection}{นวกมฺมทาน}{mtk.141}
\bookmark[dest={mtk.141},level=2]{นวกมฺมทาน}
\csromanheader{2}{\textbf{นวกมฺมทาน}}
\proseitem{308}{อถ โข ภควา ราชคเห ยถาภิรนฺตํ วิหริตฺวา เยน เวสาลี เตน จาริกํ ปกฺกามิ.\csromanspacer{} อนุปุพฺเพน จาริกํ จรมาโน เยน เวสาลี ตทวสริ, ตตฺร สุทํ ภควา เวสาลิยํ วิหรติ มหาวเน กูฏาคารสาลายํ.\csromanspacer{} เตน โข ปน สมเยน มนุสฺสา สกฺกจฺจํ นวกมฺมํ กโรนฺติ, เยปิ ภิกฺขู นวกมฺมํ อธิฏฺเฐนฺติ, เตปิ สกฺกจฺจํ อุปฏฺเฐนฺติ. \csromanfolio{308}จีวรปิณฺฑปาตเสนาสนคิลานปฺปจฺจยเภสชฺชปริกฺขาเรน.\csromanspacer{} อถ โข อญฺญตรสฺส ทลิทฺทสฺส ตุนฺนวายสฺส เอตทโหสิ “น โข อิทํ โอรกํ ภวิสฺสติ, ยถยิเม มนุสฺสา สกฺกจฺจํ นวกมฺมํ กโรนฺติ, ยํนูนาหมฺปิ นวกมฺมํ กเรยฺยนฺ”ติ.\csromanspacer{} อถ โข โส ทลิทฺโท ตุนฺนวาโย สามํ จิกฺขลฺลํ มทฺทิตฺวา อิฏฺฐกาโย จินิตฺวา กุฏฺฏํ อุฏฺฐาเปสิ, เตน อกุสลเกน จิตา วงฺกา ภิตฺติ ปริปติ.\csromanspacer{} ทุติยมฺปิ โข ฯเปฯ.\csromanspacer{} ตติยมฺปิ โข โส ทลิทฺโท ตุนฺนวาโย สามํ จิกฺขลฺลํ มทฺทิตฺวา อิฏฺฐกาโย จินิตฺวา กุฏฺฏํ อุฏฺฐาเปสิ, เตน อกุสลเกน จิตา วงฺกา ภิตฺติ ปริปติ.\csromanspacer{} อถ โข โส ทลิทฺโท ตุนฺนวาโย อุชฺฌายติ ขิยฺยติ วิปาเจติ “เย อิเมสํ สมณานํ สกฺยปุตฺติยานํ เทนฺติ จีวรปิณฺฑปาตเสนาสนคิลานปฺปจฺจยเภสชฺชปริกฺขารํ, เต อิเม โอวทนฺติ อนุสาสนฺติ, เตสญฺจ นวกมฺมํ อธิฏฺเฐนฺติ, อหํ ปนมฺหิ ทลิทฺโท, น มํ โกจิ โอวทติ วา อนุสาสติ วา, นวกมฺมํ วา อธิฏฺเฐตี”ติ.\csromanspacer{} อสฺโสสุํ โข ภิกฺขู ตสฺส ทลิทฺทสฺส ตุนฺนวายสฺส อุชฺฌายนฺตสฺส ขิยฺยนฺตสฺส วิปาเจนฺตสฺส.\csromanspacer{} อถ โข เต ภิกฺขู ภควโต เอตมตฺถํ อาโรเจสุํ.\csromanspacer{} อถ โข ภควา เอตสฺมิํ นิทาเน เอตสฺมิํ ปกรเณ ธมฺมิํ กถํ กตฺวา ภิกฺขู อามนฺเตสิ\csromandash{}อนุชานามิ ภิกฺขเว นวกมฺมํ ทาตุํ, นวกมฺมิโก ภิกฺขเว ภิกฺขุ อุสฺสุกฺกํ อาปชฺชิสฺสติ “กินฺติ นุ โข วิหาโร ขิปฺปํ ปริโยสานํ คจฺเฉยฺยา”ติ, ขณฺฑํ ผุลฺลํ ปฏิสงฺขริสฺสติ.\csromanspacer{} เอวญฺจ ปน ภิกฺขเว ทาตพฺพํ, ปฐมํ ภิกฺขุ ยาจิตพฺโพ, ยาจิตฺวา พฺยตฺเตน ภิกฺขุนา ปฏิพเลน สํโฆ ญาเปตพฺโพ⎯}

\proseitem{309}{“สุณาตุ เม ภนฺเต สํโฆ, ยทิ สํฆสฺส ปตฺตกลฺลํ, สํโฆ อิตฺถนฺนามสฺส คหปติโน วิหารํ อิตฺถนฺนามสฺส ภิกฺขุโน นวกมฺมํ ทเทยฺย, เอสา ญตฺติ.}

\prose{สุณาตุ เม ภนฺเต สํโฆ, สํโฆ อิตฺถนฺนามสฺส คหปติโน วิหารํ อิตฺถนฺนามสฺส ภิกฺขุโน นวกมฺมํ เทติ, ยสฺสายสฺมโต ขมติ อิตฺถนฺนามสฺส คหปติโน วิหารํ อิตฺถนฺนามสฺส ภิกฺขุโน นวกมฺมสฺส ทานํ, โส ตุณฺหสฺส, ยสฺส นกฺขมติ, โส ภาเสยฺย.}

\prose{ทินฺโน สํเฆน อิตฺถนฺนามสฺส คหปติโน วิหาโร อิตฺถนฺนามสฺส ภิกฺขุโน นวกมฺมํ, ขมติ สํฆสฺส, ตสฺมา ตุณฺหี, เอวเมตํ ธารยามี”ติ.}

\csromanmatikaanchor{mtk.142}
\csromantocmarkref{subsection}{อคฺคาสนาทิ-อนุชานน}{mtk.142}
\bookmark[dest={mtk.142},level=2]{อคฺคาสนาทิ-อนุชานน}
\csromanheader{2}{\textbf{อคฺคาสนาทิ-อนุชานน}}
\proseitem{310}{\csromanfolio{309}อถ โข ภควา เวสาลิยํ ยถาภิรนฺตํ วิหริตฺวา เยน สาวตฺถิ เตน จาริกํ ปกฺกามิ.\csromanspacer{} เตน โข ปน สมเยน ฉพฺพคฺคิยานํ ภิกฺขูนํ อนฺเตวาสิกา ภิกฺขู พุทฺธปฺปมุขสฺส สํฆสฺส ปุรโต ปุรโต คนฺตฺวา วิหาเร ปริคฺคณฺหนฺติ เสยฺยาโย ปริคฺคณฺหนฺติ “อิทํ อมฺหากํ อุปชฺฌายานํ ภวิสฺสติ, อิทํ อมฺหากํ อาจริยานํ ภวิสฺสติ, อิทํ อมฺหากํ ภวิสฺสตี”ติ.}

\prose{อถ โข อายสฺมา สาริปุตฺโต พุทฺธปฺปมุขสฺส สํฆสฺส ปิฏฺฐิโต ปิฏฺฐิโต คนฺตฺวา วิหาเรสุ ปริคฺคหิเตสุ เสยฺยาสุ ปริคฺคหิตาสุ เสยฺยํ อลภมาโน อญฺญตรสฺมิํ รุกฺขมูเล นิสีทิ.\csromanspacer{} อถ โข ภควา รตฺติยา ปจฺจูสสมยํ ปจฺจุฏฺฐาย อุกฺกาสิ.\csromanspacer{} อายสฺมาปิ สาริปุตฺโต อุกฺกาสิ.\csromanspacer{} โก เอตฺถาติ.\csromanspacer{} อหํ ภควา สาริปุตฺโตติ.\csromanspacer{} กิสฺส ตฺวํ สาริปุตฺต อิธ นิสินฺโนติ.\csromanspacer{} อถ โข อายสฺมา สาริปุตฺโต ภควโต เอตมตฺถํ อาโรเจสิ.\csromanspacer{} อถ โข ภควา เอตสฺมิํ นิทาเน เอตสฺมิํ ปกรเณ ภิกฺขุสํฆํ สนฺนิปาตาเปตฺวา ภิกฺขู ปฏิปุจฺฉิ\csromandash{}สจฺจํ กิร ภิกฺขเว ฉพฺพคฺคิยานํ ภิกฺขูนํ อนฺเตวาสิกา ภิกฺขู พุทฺธปฺปมุขสฺส สํฆสฺส ปุรโต ปุรโต คนฺตฺวา วิหาเร ปริคฺคณฺหนฺติ เสยฺยาโย ปริคฺคณฺหนฺติ “อิทํ อมฺหากํ อุปชฺฌายานํ ภวิสฺสติ, อิทํ อมฺหากํ อาจริยานํ ภวิสฺสติ, อิทํ อมฺหากํ ภวิสฺสตี”ติ.\csromanspacer{} สจฺจํ ภควาติ.\csromanspacer{} วิครหิ พุทฺโธ ภควา ฯเปฯ.\csromanspacer{} กถํ หิ นาม เต ภิกฺขเว โมฆปุริสา พุทฺธปฺปมุขสฺส สํฆสฺส ปุรโต ปุรโต คนฺตฺวา วิหาเร ปริคฺคเหสฺสนฺติ เสยฺยาโย ปริคฺคเหสฺสนฺติ “อิทํ อมฺหากํ อุปชฺฌายานํ ภวิสฺสติ, อิทํ อมฺหากํ อาจริยานํ ภวิสฺสติ, อิทํ อมฺหากํ ภวิสฺสตี”ติ.\csromanspacer{} เนตํ ภิกฺขเว อปฺปสนฺนานํ วา ปสาทาย ฯเปฯ วิครหิตฺวา ฯเปฯ ธมฺมิํ กถํ กตฺวา ภิกฺขู อามนฺเตสิ “โก ภิกฺขเว อรหติ อคฺคาสนํ อคฺโคทกํ อคฺคปิณฺฑนฺ”ติ.}

\prose{เอกจฺเจ ภิกฺขู เอวมาหํสุ “โย ภควา ขตฺติยกุลา ปพฺพชิโต, โส อรหติ อคฺคาสนํ อคฺโคทกํ อคฺคปิณฺฑนฺ”ติ.\csromanspacer{} เอกจฺเจ ภิกฺขู เอวมาหํสุ “โย ภควา พฺราหฺมณกุลา ปพฺพชิโต, โส อรหติ อคฺคาสนํ อคฺโคทกํ อคฺคปิณฺฑนฺ”ติ.\csromanspacer{} เอกจฺเจ ภิกฺขู เอวมาหํสุ “โย ภควา คหปติกุลา ปพฺพชิโต, โส อรหติ อคฺคาสนํ อคฺโคทกํ อคฺคปิณฺฑนฺ”ติ. \csromanfolio{310}เอกจฺเจ ภิกฺขู เอวมาหํสุ “โย ภควา สุตฺตนฺติโก, โส อรหติ อคฺคาสนํ อคฺโคทกํ อคฺคปิณฺฑนฺ”ติ.\csromanspacer{} เอกจฺเจ ภิกฺขู เอวมาหํสุ “โย ภควา วินยธโร, โส อรหติ อคฺคาสนํ อคฺโคทกํ อคฺคปิณฺฑนฺ”ติ.\csromanspacer{} เอกจฺเจ ภิกฺขู เอวมาหํสุ “โย ภควา ธมฺมกถิโก, โส อรหติ อคฺคาสนํ อคฺโคทกํ อคฺคปิณฺฑนฺ”ติ.\csromanspacer{} เอกจฺเจ ภิกฺขู เอวมาหํสุ “โย ภควา ปฐมสฺส ฌานสฺส ลาภี, โส อรหติ อคฺคาสนํ อคฺโคทกํ อคฺคปิณฺฑนฺ”ติ.\csromanspacer{} เอกจฺเจ ภิกฺขู เอวมาหํสุ “โย ภควา ทุติยสฺส ฌานสฺส ลาภี, โส อรหติ อคฺคาสนํ อคฺโคทกํ อคฺคปิณฺฑนฺ”ติ.\csromanspacer{} เอกจฺเจ ภิกฺขู เอวมาหํสุ “โย ภควา ตติยสฺส ฌานสฺส ลาภี, โส อรหติ อคฺคาสนํ อคฺโคทกํ อคฺคปิณฺฑนฺ”ติ.\csromanspacer{} เอกจฺเจ ภิกฺขู เอวมาหํสุ “โย ภควา จตุตฺถสฺส ฌานสฺส ลาภี, โส อรหติ อคฺคาสนํ อคฺโคทกํ อคฺคปิณฺฑนฺ”ติ.\csromanspacer{} เอกจฺเจ ภิกฺขู เอวมาหํสุ “โย ภควา โสตาปนฺโน, โส อรหติ อคฺคาสนํ อคฺโคทกํ อคฺคปิณฺฑนฺ”ติ.\csromanspacer{} เอกจฺเจ ภิกฺขู เอวมาหํสุ “โย ภควา สกทาคามี ฯเปฯ.\csromanspacer{} โย ภควา อนาคามี ฯเปฯ.\csromanspacer{} โย ภควา อรหา, โส อรหติ อคฺคาสนํ อคฺโคทกํ อคฺคปิณฺฑนฺ”ติ.\csromanspacer{} เอกจฺเจ ภิกฺขู เอวมาหํสุ “โย ภควา เตวิชฺโช, โส อรหติ อคฺคาสนํ อคฺโคทกํ อคฺคปิณฺฑนฺ”ติ.\csromanspacer{} เอกจฺเจ ภิกฺขู เอวมาหํสุ “โย ภควา ฉฬภิญฺโญ, โส อรหติ อคฺคาสนํ อคฺโคทกํ อคฺคปิณฺฑนฺ”ติ.}

\proseitem{311}{อถ โข ภควา ภิกฺขู อามนฺเตสิ “ภูตปุพฺพํ ภิกฺขเว หิมวนฺตปเทเส\footnote{หิมวนฺตสฺส (สี, สฺยา)} มหานิโคฺรโธ อโหสิ, ตํ ตโย สหายา อุปนิสฺสาย วิหริํสุ ติตฺติโร จ มกฺกโฏ จ หตฺถินาโค จ, เต อญฺญมญฺญํ อคารวา อปฺปติสฺสา อสภาควุตฺติกา วิหรนฺติ.\csromanspacer{} อถ โข ภิกฺขเว เตสํ สหายานํ เอตทโหสิ “อโห นูน มยํ ชาเนยฺยาม, ยํ อมฺหากํ ชาติยา มหนฺตตรํ, ตํ มยํ สกฺกเรยฺยาม ครุํ กเรยฺยาม มาเนยฺยาม ปูเชยฺยาม, ตสฺส จ มยํ โอวาเท ติฏฺเฐยฺยามา”ติ.}

\prose{อถ โข ภิกฺขเว ติตฺติโร จ มกฺกโฏ จ หตฺถินาคํ ปุจฺฉิํสุ “ตฺวํ สมฺม กิํ โปราณํ สรสี”ติ.\csromanspacer{} ยทาหํ สมฺมา โปโต โหมิ, อิมํ นิโคฺรธํ \csromanfolio{311}อนฺตรา สตฺถีนํ\footnote{อนฺตราสตฺถิกํ (สี)} กริตฺวา อติกฺกมามิ, อคฺคงฺกุรกํ เม อุทรํ ฉุปติ, อิมาหํ สมฺมา โปราณํ สรามีติ.}

\prose{อถ โข ภิกฺขเว ติตฺติโร จ หตฺถินาโค จ มกฺกฏํ ปุจฺฉิํสุ “ตฺวํ สมฺม กิํ โปราณํ สรสี”ติ.\csromanspacer{} ยทาหํ สมฺมา ฉาโป โหมิ, ฉมายํ นิสีทิตฺวา อิมสฺส นิโคฺรธสฺส อคฺคงฺกุรกํ ขาทามิ, อิมาหํ สมฺมา โปราณํ สรามีติ.}

\prose{อถ โข ภิกฺขเว มกฺกโฏ จ หตฺถินาโค จ ติตฺติรํ ปุจฺฉิํสุ “ตฺวํ สมฺม กิํ โปราณํ สรสี”ติ.\csromanspacer{} อมุกสฺมิํ สมฺมา โอกาเส มหานิโคฺรโธ อโหสิ, ตโต อหํ ผลํ ภกฺขิตฺวา อิมสฺมิํ โอกาเส วจฺจํ อกาสิํ, ตสฺสายํ นิโคฺรโธ ชาโต, ตทาหํ สมฺมา ชาติยา มหนฺตตโรติ.}

\prose{อถ โข ภิกฺขเว มกฺกโฏ จ หตฺถิ นาโค จ ติตฺติรํ เอตทโวจุํ “ตฺวํ สมฺม อมฺหากํ ชาติยา มหนฺตตโร, ตํ มยํ สกฺกริสฺสาม ครุํ กริสฺสาม มาเนสฺสาม ปูเชสฺสาม, ตุยฺหญฺจ มยํ โอวาเท ปติฏฺฐิสฺสามา”ติ.\csromanspacer{} อถ โข ภิกฺขเว ติตฺติโร มกฺกฏญฺจ หตฺถินาคญฺจ ปญฺจสุ สีเลสุ สมาทเปสิ, อตฺตนา จ ปญฺจสุ สีเลสุ สมาทาย วตฺตติ, เต อญฺญมญฺญํ สคารวา สปฺปติสฺสา สภาควุตฺติกา วิหริตฺวา กายสฺส เภทา ปรํ มรณา สุคติํ สคฺคํ โลกํ อุปปชฺชิํสุ.\csromanspacer{} เอวํ โข ตํ ภิกฺขเว ติตฺติริยํ นาม พฺรหฺมจริยํ อโหสิ.}

\csromangathameasure{*เย วุฑฺฒมปจายนฺติ, นรา ธมฺมสฺส โกวิทา. \\ ทิฏฺเฐ ธมฺเม จ ปาสํสา, สมฺปราเย จ สุคฺคตีติ.}
\csromangathagroup{\csromangathastanza{\csromansymbolfootnote{*}{ขุ 5. 9 ปิฏฺเฐ ชาตเก.}เย วุฑฺฒมปจายนฺติ, นรา ธมฺมสฺส โกวิทา. \\ ทิฏฺเฐ ธมฺเม จ ปาสํสา, สมฺปราเย จ สุคฺคตีติ.}}

\prose{เต หิ นาม ภิกฺขเว ติรจฺฉานคตา ปาณา อญฺญมญฺญํ สคารวา สปฺปติสฺสา สภาควุตฺติกา วิหริสฺสนฺติ.\csromanspacer{} อิธ โข ตํ ภิกฺขเว โสเภถ, ยํ ตุเมฺห เอวํ สฺวากฺขาเต ธมฺมวินเย ปพฺพชิตา สมานา อญฺญมญฺญํ อคารวา อปฺปติสฺสา อสภาควุตฺติกา วิหเรยฺยาถ.\csromanspacer{} เนตํ ภิกฺขเว อปฺปสนฺนานํ วา ปสาทาย ฯเปฯ วิครหิตฺวา ฯเปฯ ธมฺมิํ กถํ กตฺวา ภิกฺขู อามนฺเตสิ “อนุชานามิ ภิกฺขเว ยถาวุฑฺฒํ อภิวาทนํ ปจฺจุฏฺฐานํ อญฺชลิกมฺมํ สามีจิกมฺมํ อคฺคาสนํ อคฺโคทกํ อคฺคปิณฺฑํ, น จ ภิกฺขเว สํฆิกํ ยถาวุฑฺฒํ ปฏิพาหิตพฺพํ, โย ปฏิพาเหยฺย,อาปตฺติ ทุกฺกฏสฺสา”ติ.}

\csromanmatikaanchor{mtk.143}
\csromantocmarkref{subsection}{อวนฺทิยาทิปุคฺคล}{mtk.143}
\bookmark[dest={mtk.143},level=2]{อวนฺทิยาทิปุคฺคล}
\csromanheader{2}{\textbf{อวนฺทิยาทิปุคฺคล}}
\proseitem{312}{\csromanfolio{312}ทสยิเม ภิกฺขเว อวนฺทิยา.\csromanspacer{} ปุเร อุปสมฺปนฺเนน ปจฺฉา อุปสมฺปนฺโน อวนฺทิโย, อนุปสมฺปนฺโน อวนฺทิโย, นานาสํวาสโก วุฑฺฒตโร อธมฺมวาที อวนฺทิโย, มาตุคาโม อวนฺทิโย, ปณฺฑโก อวนฺทิโย, ปาริวาสิโก อวนฺทิโย, มูลายปฏิกสฺสนารโห อวนฺทิโย, มานตฺตารโห อวนฺทิโย, มานตฺตจาริโก อวนฺทิโย, อพฺภานารโห อวนฺทิโย.\csromanspacer{} อิเม โข ภิกฺขเว ทส อวนฺทิยา.}

\prose{ตโยเม ภิกฺขเว วนฺทิยา.\csromanspacer{} ปจฺฉา-อุปสมฺปนฺเนน ปุเร-อุปสมฺปนฺโน วนฺทิโย, นานาสํวาสโก วุฑฺฒตโร ธมฺมวาที วนฺทิโย, สเทวเก ภิกฺขเว โลเก สมารเก สพฺรหฺมเก สสฺสมณพฺราหฺมณิยา ปชาย สเทวมนุสฺสาย ตถาคโต อรหํ สมฺมาสมฺพุทฺโธ วนฺทิโย.\csromanspacer{} อิเม โข ภิกฺขเว ตโย วนฺทิยาติ.}

\csromanmatikaanchor{mtk.144}
\csromantocmarkref{subsection}{อาสนปฺปฏิพาหนปฏิกฺเขป}{mtk.144}
\bookmark[dest={mtk.144},level=2]{อาสนปฺปฏิพาหนปฏิกฺเขป}
\csromanheader{2}{\textbf{อาสนปฺปฏิพาหนปฏิกฺเขป}}
\proseitem{313}{เตน โข ปน สมเยน มนุสฺสา สํฆํ อุทฺทิสฺส มณฺฑเป ปฏิยาเทนฺติ, สนฺถเร ปฏิยาเทนฺติ, โอกาเส ปฏิยาเทนฺติ, ฉพฺพคฺคิยานํ ภิกฺขูนํ อนฺเตวาสิกา ภิกฺขู “สํฆิกญฺเญว ภควตา ยถาวุฑฺฒํ อนุญฺญาตํ, โน อุทฺทิสฺสกตนฺ”ติ พุทฺธปฺปมุขสฺส สํฆสฺส ปุรโต ปุรโต คนฺตฺวา มณฺฑเปปิ ปริคฺคณฺหนฺติ สนฺตเรปิ ปริคฺคณฺหนฺติ โอกาเสปิ ปริคฺคณฺหนฺติ “อิทํ อมฺหากํ อุปชฺฌายานํ ภวิสฺสติ, อิทํ อมฺหากํ อาจริยานํ ภวิสฺสติ, อิทํ อมฺหากํ ภวิสฺสตี”ติ.\csromanspacer{} อถ โข อายสฺมา สาริปุตฺโต พุทฺธปฺปมุขสฺส สํฆสฺส ปิฏฺฐิโต ปิฏฺฐิโต คนฺตฺวา มณฺฑเปสุ ปริคฺคหิเตสุ สนฺถเรสุ ปริคฺคหิเตสุ โอกาเสสุ ปริคฺคหิเตสุ โอกาสํ อลภมาโน อญฺญตรสฺมิํ รุกฺขมูเล นิสีทิ.\csromanspacer{} อถ โข ภควา รตฺติยา ปจฺจูสสมยํ ปจฺจุฏฺฐาย อุกฺกาสิ.\csromanspacer{} อายสฺมาปิ สาริปุตฺโต อุกฺกาสิ.\csromanspacer{} โก เอตฺถาติ.\csromanspacer{} อหํ ภควา สาริปุตฺโตติ.\csromanspacer{} กิสฺส ตฺวํ สาริปุตฺต อิธ นิสินฺโนติ.\csromanspacer{} อถ โข อายสฺมา สาริปุตฺโต ภควโต เอตมตฺถํ อาโรเจสิ.\csromanspacer{} อถ โข ภควา เอตสฺมิํ นิทาเน เอตสฺมิํ ปกรเณ ภิกฺขุสํฆํ สนฺนิปาตาเปตฺวา ภิกฺขู ปฏิปุจฺฉิ “สจฺจํ กิร ภิกฺขเว ฉพฺพคฺคิยานํ ภิกฺขูนํ อนฺเตวาสิกา ภิกฺขู ‘สํฆิกญฺเญว ภควตา ยถาวุฑฺฒํ อนุญฺญาตํ, โน อุทฺทิสฺสกตนฺ’ติ พุทฺธปฺปมุขสฺส สํฆสฺส ปุรโต ปุรโต \csromanfolio{313}คนฺตฺวา มณฺฑเป ปริคฺคณฺหนฺติ, สนฺถเร ปริคฺคณฺหนฺติ, โอกาเส ปริคฺคณฺหนฺติ, อิทํ อมฺหากํ อุปชฺฌายานํ ภวิสฺสติ, อิทํ อมฺหากํ อาจริยานํ ภวิสฺสติ, อิทํ อมฺหากํ ภวิสฺสตี”ติ.\csromanspacer{} สจฺจํ ภควาติ ฯเปฯ วิครหิตฺวา ฯเปฯ ธมฺมิํ กถํ กตฺวา ภิกฺขู อามนฺเตสิ “น ภิกฺขเว อุทฺทิสฺสกตมฺปิ ยถาวุฑฺฒํ ปฏิพาหิตพฺพํ, โย ปฏิพาเหยฺย, อาปตฺติ ทุกฺกฏสฺสา”ติ.}

\csromanmatikaanchor{mtk.145}
\csromantocmarkref{subsection}{คิหิวิกต-อนุชานน}{mtk.145}
\bookmark[dest={mtk.145},level=2]{คิหิวิกต-อนุชานน}
\csromanheader{2}{\textbf{คิหิวิกต-อนุชานน}}
\proseitem{314}{เตน โข ปน สมเยน มนุสฺสา ภตฺตคฺเค อนฺตรฆเร อุจฺจาสยนมหาสยนานิ ปญฺญเปนฺติ.\csromanspacer{} เสยฺยถิทํ, อาสนฺทิํ ปลฺลงฺกํ โคนกํ จิตฺตกํ ปฏิกํ ปฏลิกํ ตูลิกํ วิกติกํ อุทฺทโลมิํ เอกนฺตโลมิํ กฏฺฏิสฺสํ โกเสยฺยํ\footnote{โกเสยฺยํ กมฺพลํ (สี, สฺยา)} กุตฺตกํ หตฺถตฺถรํ อสฺสตฺถรํ รถตฺถรํ อชินปเวณิํ กทลิมิคปวรปจฺจตฺถรณํ ส- อุตฺตรจฺฉทํ อุภโตโลหิตกูปธานํ.\csromanspacer{} ภิกฺขู กุกฺกุจฺจายนฺตา นาภินิสีทนฺติ.\csromanspacer{} ภควโต เอตมตฺถํ อาโรเจสุํ.\csromanspacer{} อนุชานามิ ภิกฺขเว ฐเปตฺวา ตีณิ อาสนฺทิํ ปลฺลงฺกํ ตูลิกํ คิหิวิกตํ\footnote{คิหิวิกฏํ (สี, ก), อวเสสํ คิหิวิกฏํ (สฺยา)} อภินิสีทิตุํ, นเตฺวว อภินิปชฺชิตุนฺติ.}

\prose{เตน โข ปน สมเยน มนุสฺสา ภตฺตคฺเค อนฺตรฆเร ตูโลนทฺธํ มญฺจํปิ ปีฐํปิ ปญฺญเปนฺติ.\csromanspacer{} ภิกฺขู กุกฺกุจฺจายนฺตา นาภินิสีทนฺติ.\csromanspacer{} ภควโต เอตมตฺถํ อาโรเจสุํ.\csromanspacer{} อนุชานามิ ภิกฺขเว คิหิวิกตํ อภินิสีทิตุํ, นเตฺวว อภินิปชฺชิตุนฺติ.}

\csromanmatikaanchor{mtk.146}
\csromantocmarkref{subsection}{เชตวนวิหารานุโมทนา}{mtk.146}
\bookmark[dest={mtk.146},level=2]{เชตวนวิหารานุโมทนา}
\csromanheader{2}{\textbf{เชตวนวิหารานุโมทนา}}
\proseitem{315}{อถ โข ภควา อนุปุพฺเพน จาริกํ จรมาโน เยน สาวตฺถิ ตทวสริ, ตตฺร สุทํ ภควา สาวตฺถิยํ วิหรติ เชตวเน อนาถปิณฺฑิกสฺส อาราเม.\csromanspacer{} อถ โข อนาถปิณฺฑิโก คหปติ เยน ภควา เตนุปสงฺกมิ, อุปสงฺกมิตฺวา ภควนฺตํ อภิวาเทตฺวา เอกมนฺตํ นิสีทิ, เอกมนฺตํ นิสินฺโน โข อนาถปิณฺฑิโก คหปติ ภควนฺตํ เอตทโวจ “อธิวาเสตุ เม ภนฺเต ภควา สฺวาตนาย ภตฺตํ สทฺธิํ ภิกฺขุสํเฆนา”ติ.\csromanspacer{} อธิวาเสสิ ภควา ตุณฺหีภาเวน.\csromanspacer{} อถ โข อนาถปิณฺฑิโก คหปติ ภควโต อธิวาสนํ วิทิตฺวา อุฏฺฐายาสนา \csromanfolio{314}ภควนฺตํ อภิวาเทตฺวา ปทกฺขิณํ กตฺวา ปกฺกามิ.\csromanspacer{} อถ โข อนาถปิณฺฑิโก คหปติ ตสฺสา รตฺติยา อจฺจเยน ปณีตํ ขาทนียํ โภชนียํ ปฏิยาทาเปตฺวา ภควโต กาลํ อาโรจาเปสิ “กาโล ภนฺเต, นิฏฺฐิตํ ภตฺตนฺ”ติ.\csromanspacer{} อถ โข ภควา ปุพฺพณฺหสมยํ นิวาเสตฺวา ปตฺตจีวรมาทาย เยน อนาถปิณฺฑิกสฺส คหปติสฺส นิเวสนํ เตนุปสงฺกมิ, อุปสงฺกมิตฺวา ปญฺญตฺเต อาสเน นิสีทิ สทฺธิํ ภิกฺขุสํเฆน.\csromanspacer{} อถ โข อนาถปิณฺฑิโก คหปติ พุทฺธปฺปมุขํ ภิกฺขุสํฆํ ปณีเตน ขาทนีเยน โภชนีเยน สหตฺถา สนฺตปฺเปตฺวา สมฺปวาเรตฺวา ภควนฺตํ ภุตฺตาวิํ โอนีตปตฺตปาณิํ เอกมนฺตํ นิสีทิ, เอกมนฺตํ นิสินฺโน โข อนาถปิณฺฑิโก คหปติ ภควนฺตํ เอตทโวจ “กถาหํ ภนฺเต เชตวเน ปฏิปชฺชามี”ติ.\csromanspacer{} เตน หิ ตฺวํ คหปติ เชตวนํ อาคตานาคตสฺส จาตุทฺทิสสฺส สํฆสฺส ปติฏฺฐเปหีติ. “เอวํ ภนฺเต”ติ โข อนาถปิณฺฑิโก คหปติ ภควโต ปฏิสฺสุตฺวา เชตวนํ อาคตานาคตสฺส จาตุทฺทิสสฺส สํฆสฺส ปติฏฺฐาเปสิ.}

\prose{อถ โข ภควา อนาถปิณฺฑิกํ คหปติํ อิมาหิ คาถาหิ อนุโมทิ\csromandash{}}

\csromangathameasure{“สีตํ อุณฺหํ ปฏิหนฺติ, ตโต วาฬมิคานิ จ. \\ สริสเป จ มกเส, สิสิเร จาปิ วุฏฺฐิโย. \\ ตโต วาตาตโป โฆโร, สญฺชาโต ปฏิหญฺญติ. \\ เลณตฺถญฺจ สุขตฺถญฺจ, ฌายิตุญฺจ วิปสฺสิตุํ. \\ วิหารทานํ สํฆสฺส, อคฺคํ พุทฺเธน วณฺณิตํ. \\ ตสฺมา หิ ปณฺฑิโต โปโส, สมฺปสฺสํ อตฺถมตฺตโน. \\ วิหาเร การเย รมฺเม, วาสเยตฺถ พหุสฺสุเต. \\ เตสํ อนฺนญฺจ ปานญฺจ, วตฺถเสนาสนานิ จ. \\ ทเทยฺย อุชุภูเตสุ, วิปฺปสนฺเนน เจตสา. \\ เต ตสฺส ธมฺมํ เทเสนฺติ, สพฺพทุกฺขาปนูทนํ. \\ ยํ โส ธมฺมํ อิธญฺญาย, ปรินิพฺพาติ อนาสโว”ติ.}
\csromangathagroup{\csromangathastanza{“สีตํ อุณฺหํ ปฏิหนฺติ, ตโต วาฬมิคานิ จ. \\ สริสเป จ มกเส, สิสิเร จาปิ วุฏฺฐิโย.}\csromangathastanza{ตโต วาตาตโป โฆโร, สญฺชาโต ปฏิหญฺญติ. \\ เลณตฺถญฺจ สุขตฺถญฺจ, ฌายิตุญฺจ วิปสฺสิตุํ.}\csromangathastanza{วิหารทานํ สํฆสฺส, อคฺคํ พุทฺเธน วณฺณิตํ. \\ ตสฺมา หิ ปณฺฑิโต โปโส, สมฺปสฺสํ อตฺถมตฺตโน.}\csromangathastanza{วิหาเร การเย รมฺเม, วาสเยตฺถ พหุสฺสุเต. \\ เตสํ อนฺนญฺจ ปานญฺจ, วตฺถเสนาสนานิ จ.}\csromangathastanza{ทเทยฺย อุชุภูเตสุ, วิปฺปสนฺเนน เจตสา. \\ เต ตสฺส ธมฺมํ เทเสนฺติ, สพฺพทุกฺขาปนูทนํ. \\ ยํ โส ธมฺมํ อิธญฺญาย, ปรินิพฺพาติ อนาสโว”ติ.}}

\prose{อถ โข ภควา อนาถปิณฺฑิกํ คหปติํ อิมาหิ คาถาหิ อนุโมทิตฺวา อุฏฺฐายาสนา ปกฺกามิ.}

\csromanmatikaanchor{mtk.147}
\csromantocmarkref{subsection}{อาสนปฺปฏิพาหนาทิ}{mtk.147}
\bookmark[dest={mtk.147},level=2]{อาสนปฺปฏิพาหนาทิ}
\csromanheader{2}{\textbf{อาสนปฺปฏิพาหนาทิ}}
\proseitem{316}{\csromanfolio{315}เตน โข ปน สมเยน อญฺญตรสฺส อาชีวกสาวกสฺส มหามตฺตสฺส สํฆภตฺตํ โหติ.\csromanspacer{} อายสฺมา อุปนนฺโท สกฺยปุตฺโต ปจฺฉา อาคนฺตฺวา วิปฺปกตโภชนํ อานนฺตริกํ ภิกฺขุํ วุฏฺฐาเปสิ, ภตฺตคฺคํ โกลาหลํ อโหสิ.\csromanspacer{} อถ โข โส มหามตฺโต อุชฺฌายติ ขิยฺยติ วิปาเจติ “กถํ หิ นาม สมณา สกฺยปุตฺติยา ปจฺฉา อาคนฺตฺวา วิปฺปกตโภชนํ อานนฺตริกํ ภิกฺขุํ วุฏฺฐาเปสฺสนฺติ, ภตฺตคฺคํ โกลาหลํ อโหสิ, นนุ นาม ลพฺภา อญฺญตฺราปิ นิสินฺเนน ยาวทตฺถํ ภุญฺชิตุนฺ”ติ.\csromanspacer{} อสฺโสสุํ โข ภิกฺขู ตสฺส มหามตฺตสฺส อุชฺฌายนฺตสฺส ขิยฺยนฺตสฺส วิปาเจนฺตสฺส.\csromanspacer{} เย เต ภิกฺขู อปฺปิจฺฉา ฯเปฯ เต อุชฺฌายนฺติ ขิยฺยนฺติ วิปาเจนฺติ “กถํ หิ นาม อายสฺมา อุปนนฺโท สกฺยปุตฺโต ปจฺฉา อาคนฺตฺวา วิปฺปกตโภชนํ อานนฺตริกํ ภิกฺขุํ วุฏฺฐาเปสฺสติ, ภตฺตคฺคํ โกลาหลํ อโหสี”ติ.\csromanspacer{} อถ โข เต ภิกฺขู ภควโต เอตมตฺถํ อาโรเจสุํ.\csromanspacer{} สจฺจํ กิร ตฺวํ อุปนนฺท ปจฺฉา อาคนฺตฺวา วิปฺปกตโภชนํ อานนฺตริกํ ภิกฺขุํ วุฏฺฐาเปสิ, ภตฺตคฺคํ โกลาหลํ อโหสีติ.\csromanspacer{} สจฺจํ ภควาติ.\csromanspacer{} วิครหิ พุทฺโธ ภควา ฯเปฯ.\csromanspacer{} กถํ หิ นาม ตฺวํ โมฆปุริส ปจฺฉา อาคนฺตฺวา วิปฺปกตโภชนํ อานนฺตริกํ ภิกฺขุํ วุฏฺฐาเปสฺสสิ, ภตฺตคฺคํ โกลาหลํ อโหสิ.\csromanspacer{} เนตํ โมฆปุริส อปฺปสนฺนานํ วา ปสาทาย ฯเปฯ วิครหิตฺวา ฯเปฯ ธมฺมิํ กถํ กตฺวา ภิกฺขู อามนฺเตสิ “น ภิกฺขเว วิปฺปกตโภชโน\footnote{โภชโน อานนฺตริโก (สฺยา)} ภิกฺขู วุฏฺฐาเปตพฺโพ, โย วุฏฺฐาเปยฺย, อาปตฺติ ทุกฺกฏสฺส.\csromanspacer{} สเจ วุฏฺฐาเปติ, ปวาริโต จ โหติ, ‘คจฺฉ อุทกํ อาหรา’ติ วตฺตพฺโพ.\csromanspacer{} เอวญฺเจตํ ลเภถ, อิจฺเจตํ กุสลํ.\csromanspacer{} โน เจ ลเภถ, สาธุกํ สิตฺถานิ คิลิตฺวา วุฑฺฒตรสฺส ภิกฺขุโน อาสนํ ทาตพฺพํ.\csromanspacer{} น เตฺววาหํ ภิกฺขเว ‘เกนจิ ปริยาเยน วุฑฺฒตรสฺส ภิกฺขุโน อาสนํ ปฏิพาหิตพฺพนฺ’ติ วทามิ.\csromanspacer{} โย ปฏิพาเหยฺย, อาปตฺติ ทุกฺกฏสฺสา”ติ.}

\prose{เตน โข ปน สมเยน ฉพฺพคฺคิยา ภิกฺขู คิลาเน ภิกฺขู วุฏฺฐาเปนฺติ.\csromanspacer{} คิลานา เอวํ วเทนฺติ “น มยํ อาวุโส สกฺโกม วุฏฺฐาตุํ, คิลานามฺหา”ติ. “มยํ อายสฺมนฺเต วุฏฺฐาเปสฺสามา”ติ ปริคฺคเหตฺวา วุฏฺฐาเปตฺวา ฐิตเก มุญฺจนฺติ, คิลานา มุจฺฉิตา ปปตนฺติ.\csromanspacer{} ภควโต เอตมตฺถํ \csromanfolio{316}อาโรเจสุํ.\csromanspacer{} น ภิกฺขเว คิลาโน วุฏฺฐาเปตพฺโพ, โย วุฏฺฐาเปยฺย, อาปตฺติ ทุกฺกฏสฺสาติ.}

\prose{เตน โข ปน สมเยน ฉพฺพคฺคิยา ภิกฺขู “คิลานา มยมฺหา อวุฏฺฐาปนียา”ติ วรเสยฺยาโย ปลิพุนฺเธนฺติ\footnote{ปลิพุนฺธนฺติ (ก)}.\csromanspacer{} ภควโต เอตมตฺถํ อาโรเจสุํ.\csromanspacer{} อนุชานามิ ภิกฺขเว คิลานสฺส ปติรูปํ เสยฺยํ ทาตุนฺติ.}

\prose{เตน โข ปน สมเยน ฉพฺพคฺคิยา ภิกฺขู เลสกปฺเปน เสนาสนํ ปฏิพาหนฺติ.\csromanspacer{} ภควโต เอตมตฺถํ อาโรเจสุํ.\csromanspacer{} น ภิกฺขเว เลสกปฺเปน เสนาสนํ ปฏิพาหิตพฺพํ, โย ปฏิพาเหยฺย, อาปตฺติ ทุกฺกฏสฺสาติ.}

\prose{\csromansymbolfootnote{*}{วิ 2. 63 ปิฏฺเฐปิ.}เตน โข ปน สมเยน สตฺตรสวคฺคิยา ภิกฺขู อญฺญตรํ ปจฺจนฺติมํ มหาวิหารํ ปฏิสงฺขโรนฺติ “อิธ มยํ วสฺสํ วสิสฺสามา”ติ.\csromanspacer{} อทฺทสํสุ\footnote{อทฺทสาสุํ (ก)} โข ฉพฺพคฺคิยา ภิกฺขู สตฺตรสวคฺคิเย ภิกฺขู วิหารํ ปฏิสงฺขโรนฺเต, ทิสฺวาน เอวมาหํสุ “อิเม อาวุโส สตฺตรสวคฺคิยา ภิกฺขู วิหารํ\footnote{อญฺญตรํ วิหารํ (ก)} ปฏิสงฺขโรนฺติ, หนฺท เน วุฏฺฐาเปสฺสามา”ติ.\csromanspacer{} เอกจฺเจ เอวมาหํสุ “อาคเมถาวุโส ยาว ปฏิสงฺขโรนฺติ, ปฏิสงฺขเต วุฏฺฐาเปสฺสามา”ติ.\csromanspacer{} อถ โข ฉพฺพคฺคิยา ภิกฺขู สตฺตรสวคฺคิเย ภิกฺขู เอตทโวจุํ “อุฏฺเฐถาวุโส อมฺหากํ วิหาโร ปาปุณาตี”ติ.\csromanspacer{} นนุ อาวุโส ปฏิกจฺเจว อาจิกฺขิตพฺพํ, มยญฺจญฺญํ ปฏิสงฺขเรยฺยามาติ.\csromanspacer{} นนุ อาวุโส สํฆิโก วิหาโรติ.\csromanspacer{} อามาวุโส สํฆิโก วิหาโรติ.\csromanspacer{} อุฏฺเฐถาวุโส, อมฺหากํ วิหาโร ปาปุณาตีติ.\csromanspacer{} มหลฺลโก อาวุโส วิหาโร, ตุเมฺหปิ วสถ, มยมฺปิ วสิสฺสามาติ. “อุฏฺเฐถาวุโส, อมฺหากํ วิหาโร ปาปุณาตี”ติ กุปิตา อนตฺตมนา คีวายํ คเหตฺวา นิกฺกฑฺฒนฺติ.\csromanspacer{} เต นิกฺกฑฺฒิยมานา โรทนฺติ.\csromanspacer{} ภิกฺขู เอวมาหํสุ “กิสฺส ตุเมฺห อาวุโส โรทถา”ติ.\csromanspacer{} อิเม อาวุโส ฉพฺพคฺคิยา ภิกฺขู กุปิตา อนตฺตมนา อเมฺห สํฆิกา วิหารา นิกฺกฑฺฒนฺตีติ.\csromanspacer{} เย เต ภิกฺขู อปฺปิจฺฉา ฯเปฯ เต อุชฺฌายนฺติ ขิยฺยนฺติ วิปาเจนฺติ “กถํ หิ นาม ฉพฺพคฺคิยา ภิกฺขู กุปิตา อนตฺตมนา ภิกฺขู สํฆิกา วิหารา นิกฺกฑฺฒิสฺสนฺตี”ติ.\csromanspacer{} อถ โข เต ภิกฺขู ภควโต เอตมตฺถํ อาโรเจสุํ ฯเปฯ.\csromanspacer{} สจฺจํ กิร ตุเมฺห ภิกฺขเว กุปิตา อนตฺตมนา สํฆิกา วิหารา ภิกฺขู นิกฺกฑฺฒถาติ.\csromanspacer{} สจฺจํ ภควาติ ฯเปฯ วิครหิตฺวา ฯเปฯ ธมฺมิํ กถํ กตฺวา ภิกฺขู อามนฺเตสิ “น ภิกฺขเว \csromanfolio{317}กุปิเตน อนตฺตมเนน ภิกฺขู สํฆิกา วิหารา นิกฺกฑฺฒิตพฺโพ, โย นิกฺกฑฺเฒยฺย, ยถาธมฺโม กเรตพฺโพ.\csromanspacer{} อนุชานามิ ภิกฺขเว เสนาสนํ คาเหตุนฺ”ติ.}

\csromanmatikaanchor{mtk.148}
\csromantocmarkref{subsection}{เสนาสนคฺคาหาปกสมฺมุติ}{mtk.148}
\bookmark[dest={mtk.148},level=2]{เสนาสนคฺคาหาปกสมฺมุติ}
\csromanheader{2}{\textbf{เสนาสนคฺคาหาปกสมฺมุติ}}
\proseitem{317}{อถ โข ภิกฺขูนํ เอตทโหสิ “เกน นุ โข เสนาสนํ คาเหตพฺพนฺ”ติ.\csromanspacer{} ภควโต เอตมตฺถํ อาโรเจสุํ ฯเปฯ.\csromanspacer{} อนุชานามิ ภิกฺขเว ปญฺจหงฺเคหิ สมนฺนาคตํ ภิกฺขุํ เสนาสนคฺคาหาปกํ สมฺมนฺนิตุํ.\csromanspacer{} โย น ฉนฺทาคติํ คจฺเฉยฺย, น โทสาคติํ คจฺเฉยฺย, น โมหาคติํ คจฺเฉยฺย, น ภยาคติํ คจฺเฉยฺย, คหิตาคหิตญฺจ ชาเนยฺย.\csromanspacer{} เอวญฺจ ปน ภิกฺขเว สมฺมนฺนิตพฺโพ.}

\prose{ปฐมํ ภิกฺขู ยาจิตพฺโพ, ยาจิตฺวา พฺยตฺเตน ภิกฺขุนา ปฏิพเลน สํโฆ ญาเปตพฺโพ\csromandash{}}

\prose{“สุณาตุ เม ภนฺเต สํโฆ, ยทิ สํฆสฺส ปตฺตกลฺลํ, สํโฆ อิตฺถนฺนามํ ภิกฺขุํ เสนาสนคฺคาหาปกํ สมฺมนฺเนยฺย, เอสา ญตฺติ.}

\prose{สุณาตุ เม ภนฺเต สํโฆ, สํโฆ อิตฺถนฺนามํ ภิกฺขุํ เสนาสนคฺคาหาปกํ สมฺมนฺนติ, ยสฺสายสฺมโต ขมติ อิตฺถนฺนามสฺส ภิกฺขุโน เสนาสนคฺคาหาปกสฺส สมฺมุติ, โส ตุณฺหสฺส, ยสฺส นกฺขมติ, โส ภาเสยฺย.}

\prose{สมฺมโต สํเฆน อิตฺถนฺนาโม ภิกฺขุ เสนาสนคฺคาหาปโก, ขมติ สํฆสฺส, ตสฺมา ตุณฺหี, เอวเมตํ ธารยามี”ติ.}

\proseitem{318}{อถ โข เสนาสนคฺคาหาปกานํ ภิกฺขุนํ เอตทโหสิ “กถํ นุ โข เสนาสนํ คาเหตพฺพนฺ”ติ.\csromanspacer{} ภควโต เอตมตฺถํ อาโรเจสุํ.\csromanspacer{} อนุชานามิ ภิกฺขเว ปฐมํ ภิกฺขู คเณตุํ, ภิกฺขู คเณตฺวา เสยฺยา คเณตุํ, เสยฺยา คเณตฺวา เสยฺยคฺเคน คาเหตุนฺติ.\csromanspacer{} เสยฺยคฺเคน คาเหนฺตา เสยฺยา อุสฺสารยิํสุ ฯเปฯ.\csromanspacer{} อนุชานามิ ภิกฺขเว วิหารคฺเคน คาเหตุนฺติ.\csromanspacer{} วิหารคฺเคน คาเหนฺตา วิหารา อุสฺสารยิํสุ ฯเปฯ.\csromanspacer{} อนุชานามิ ภิกฺขเว ปริเวณคฺเคน คาเหตุนฺติ.\csromanspacer{} ปริเวณคฺเคน คาเหนฺตา ปริเวณา อุสฺสารยิํสุ ฯเปฯ.\csromanspacer{} อนุชานามิ ภิกฺขเว อนุภาคมฺปิ ทาตุํ, คหิเต อนุภาเค อญฺโญ ภิกฺขุ อาคจฺฉติ, อกามา น ทาตพฺโพติ.}

\prose{\csromanfolio{318}เตน โข ปน สมเยน ภิกฺขู นิสฺสีเม ฐิตสฺส เสนาสนํ คาเหนฺติ.\csromanspacer{} ภควโต เอตมตฺถํ อาโรเจสุํ.\csromanspacer{} น ภิกฺขเว นิสฺสีเม ฐิตสฺสํ เสนาสนํ คาเหตพฺพํ, โย คาเหยฺย, อาปตฺติ ทุกฺกฏสฺสาติ.}

\prose{เตน โข ปน สมเยน ภิกฺขู เสนาสนํ คเหตฺวา สพฺพกาลํ ปฏิพาหนฺติ.\csromanspacer{} ภควโต เอตมตฺถํ อาโรเจสุํ.\csromanspacer{} น ภิกฺขเว เสนาสนํ คเหตฺวา สพฺพกาลํ ปฏิพาหิตพฺพํ, โย ปฏิพาเหยฺย, อาปตฺติ ทุกฺกฏสฺส.\csromanspacer{} อนุชานามิ ภิกฺขเว วสฺสานํ เตมาสํ ปฏิพาหิตุํ, อุตุกาลํ ปน น ปฏิพาหิตุนฺติ.}

\prose{อถ โข ภิกฺขูนํ เอตทโหสิ “กติ นุ โข เสนาสนคฺคาหา”ติ.\csromanspacer{} ภควโต เอตมตฺถํ อาโรเจสุํ.\csromanspacer{} ตโย เม ภิกฺขเว เสนาสนคฺคาหา, ปุริมโก ปจฺฉิมโก อนฺตรามุตฺตโก.\csromanspacer{} อปรชฺชุคตาย อาสาฬฺหิยา ปุริมโก คาเหตพฺโพ, มาสคตาย อาสาฬฺหิยา ปจฺฉิมโก คาเหตพฺโพ, อปรชฺชุคตาย ปวารณาย อายติํ วสฺสาวาสตฺถาย อนฺตรามุตฺตโก คาเหตพฺโพ.\csromanspacer{} อิเม โข ภิกฺขเว ตโย เสนาสนคฺคาหาติ.}

\nitthitam{ทุติยภาณวาโร นิฏฺฐิโต.}
\csromanmatikaanchor{mtk.149}
\csromantocmarkref{section}{3. ตติยภาณวาร}{mtk.149}
\bookmark[dest={mtk.149},level=1]{3. ตติยภาณวาร}
\csromanheader{1}{3. ตติยภาณวาร}
\proseitem{319}{เตน โข ปน สมเยน อายสฺมา อุปนนฺโท สกฺยปุตฺโต สาวตฺถิยํ เสนาสนํ คเหตฺวา อญฺญตรํ คามกาวาสํ อคมาสิ, ตตฺถปิ เสนาสนํ อคฺคเหสิ.\csromanspacer{} อถ โข เตสํ ภิกฺขูนํ เอตทโหสิ “อยํ อาวุโส อายสฺมา อุปนนฺโท สกฺยปุตฺโต ภณฺฑนการโก กลหการโก วิวาทการโก ภสฺสการโก สํเฆ อธิกรณการโก, สจายํ อิธ วสฺสํ วสิสฺสติ, สพฺเพว มยํ น ผาสุ ภวิสฺสาม, หนฺท นํ ปุจฺฉามา”ติ.\csromanspacer{} อถ โข เต ภิกฺขู อายสฺมนฺตํ อุปนนฺทํ สกฺยปุตฺตํ เอตทโวจุํ “นนุ ตยา อาวุโส อุปนนฺท สาวตฺถิยํ เสนาสนํ คหิตนฺ”ติ.\csromanspacer{} เอวมาวุโสติ.\csromanspacer{} กิํ ปน ตฺวํ อาวุโส อุปนนฺท เอโก เทฺว ปฏิพาหสีติ.\csromanspacer{} อิธทานาหํ อาวุโส มุญฺจามิ, ตตฺถ คณฺหามีติ.\csromanspacer{} เย เต ภิกฺขู อปฺปิจฺฉา ฯเปฯ เต อุชฺฌายนฺติ ขิยฺยนฺติ วิปาเจนฺติ “กถํ หิ นาม อายสฺมา อุปนนฺโท \csromanfolio{319}สกฺยปุตฺโต เอโก เทฺว ปฏิพาเหสฺสตี”ติ.\csromanspacer{} ภควโต เอตมตฺถํ อาโรเจสุํ.\csromanspacer{} สจฺจํ กิร ตฺวํ อุปนนฺท เอโก เทฺว ปฏิพาหสีติ.\csromanspacer{} สจฺจํ ภควาติ.\csromanspacer{} วิครหิ พุทฺโธ ภควา ฯเปฯ.\csromanspacer{} กถํ หิ นาม ตฺวํ โมฆปุริส เอโก เทฺว ปฏิพาหิสฺสสิ, ตตฺถ ตยา โมฆปุริส คหิตํ อิธ มุตฺตํ, อิธ ตยา คหิตํ ตตฺร มุตฺตํ, เอวํ โข ตฺวํ โมฆปุริส อุภยตฺถ ปริพาหิโร.\csromanspacer{} เนตํ โมฆปุริส อปฺปสนฺนานํ วา ปสาทาย ฯเปฯ วิครหิตฺวา ฯเปฯ ธมฺมิํ กถํ กตฺวา ภิกฺขู อามนฺเตสิ “น ภิกฺขเว เอเกน เทฺว ปฏิพาหิตพฺพา, โย ปฏิพาเหยฺย, อาปตฺติ ทุกฺกฏสฺสา”ติ.}

\proseitem{320}{\csromansymbolfootnote{*}{วิ 2. 187 ปิฏฺเฐปิ.}เตน โข ปน สมเยน ภควา ภิกฺขูนํ อเนกปริยาเยน วินยกถํ กเถสิ, วินยสฺส วณฺณํ ภาสติ, วินยปริยตฺติยา วณฺณํ ภาสติ, อาทิสฺส อาทิสฺส อายสฺมโต อุปาลิสฺส วณฺณํ ภาสติ.\csromanspacer{} ภิกฺขูนํ เอตทโหสิ “ภควา โข อเนกปริยาเยน วินยกถํ กเถติ, วินยสฺส วณฺณํ ภาสติ, วินยปริยตฺติยา วณฺณํ ภาสติ, อาทิสฺส อาทิสฺส อายสฺมโต อุปาลิสฺส วณฺณํ ภาสติ, หนฺท มยํ อาวุโส อายสฺมโต อุปาลิสฺส สนฺติเก วินยํ ปริยาปุณามา”ติ.\csromanspacer{} เต'ธ\footnote{เต จ (สฺยา, ก)} พหู ภิกฺขู เถรา จ นวา จ มชฺฌิมา จ อายสฺมโต อุปาลิสฺส สนฺติเก วินยํ ปริยาปุณนฺติ, อายสฺมา อุปาลิ ฐิตโกว อุทฺทิสติ เถรานํ ภิกฺขูนํ คารเวน, เถราปิ ภิกฺขู ฐิตกาว อุทฺทิสาเปนฺติ ธมฺมคารเวน, ตตฺถ เถรา เจว ภิกฺขู กิลมนฺติ, อายสฺมา จ อุปาลิ กิลมติ.\csromanspacer{} ภควโต เอตมตฺถํ อาโรเจสุํ.\csromanspacer{} อนุชานามิ ภิกฺขเว นวเกน ภิกฺขุนา อุทฺทิสนฺเตน สมเก วา อาสเน นิสีทิตุํ อุจฺจตเร วา ธมฺมคารเวน, เถเรน ภิกฺขุนา อุทฺทิสาเปนฺเตน สมเก วา อาสเน นิสีทิตุํ นีจตเร วา ธมฺมคารเวนาติ.}

\prose{เตน โข ปน สมเยน พหู ภิกฺขู อายสฺมโต อุปาลิสฺส สนฺติเก ฐิตกา อุทฺเทสํ ปฏิมาเนนฺตา กิลมนฺติ.\csromanspacer{} ภควโต เอตมตฺถํ อาโรเจสุํ.\csromanspacer{} อนุชานามิ ภิกฺขเว สมานาสนิเกหิ สห นิสีทิตุนฺติ.\csromanspacer{} อถ โข ภิกฺขูนํ เอตทโหสิ “กิตฺตาวตา นุ โข สมานาสนิโก โหตี”ติ.\csromanspacer{} ภควโต เอตมตฺถํ อาโรเจสุํ.\csromanspacer{} อนุชานามิ ภิกฺขเว ติวสฺสนฺตเรน สห นิสีทิตุนฺติ.}

\prose{\csromanfolio{320}เตน โข ปน สมเยน สมฺพหุลา ภิกฺขู สมานาสนิกา มญฺเจ\footnote{เอกมญฺเจ (สฺยา)} นิสีทิตฺวา มญฺจํ ภินฺทิํสุ, ปีเฐ\footnote{เอกปีเฐ (สฺยา)} นิสีทิตฺวา ปีฐํ ภินฺทิํสุ.\csromanspacer{} ภควโต เอตมตฺถํ อาโรเจสุํ.\csromanspacer{} อนุชานามิ ภิกฺขเว ติวคฺคสฺส มญฺจํ ติวคฺคสฺส ปีฐนฺติ.\csromanspacer{} ติวคฺโคปิ มญฺเจ นิสีทิตฺวา มญฺจํ ภินฺทิ.\csromanspacer{} ปีเฐ นิสีทิตฺวา ปีฐํ ภินฺทิ ฯเปฯ.\csromanspacer{} อนุชานามิ ภิกฺขเว ทุวคฺคสฺส มญฺจํ ทุวคฺคสฺส ปีฐนฺติ.}

\prose{เตน โข ปน สมเยน ภิกฺขู อสมานาสนิเกหิ สห ทีฆาสเน นิสีทิตุํ กุกฺกุจฺจายนฺติ.\csromanspacer{} ภควโต เอตมตฺถํ อาโรเจสุํ.\csromanspacer{} อนุชานามิ ภิกฺขเว ฐเปตฺวา ปณฺฑกํ มาตุคามํ อุภโตพฺยญฺชนกํ อสมานาสนิเกหิ สห ทีฆาสเน นิสีทิตุนฺติ.\csromanspacer{} อถ โข ภิกฺขูนํ เอตทโหสิ “กิตฺตกํ ปจฺฉิมํ นุ โข ทีฆาสนํ โหตี”ติ.\csromanspacer{} ภควโต เอตมตฺถํ อาโรเจสุํ.\csromanspacer{} อนุชานามิ ภิกฺขเว ยํ ติณฺณํ ปโหติ, เอตฺตกํ ปจฺฉิมํ ทีฆาสนนฺติ.}

\prose{เตน โข ปน สมเย วิสาขา มิคารมาตา สํฆสฺส อตฺถาย สาฬินฺทํ ปาสาทํ การาเปตุกามา โหติ หตฺถินขกํ.\csromanspacer{} อถ โข ภิกฺขูนํ เอตทโหสิ “กิํ นุ โข ภควตา ปาสาทปริโภโค อนุญฺญาโต, กิํ อนนุญฺญาโต”ติ.\csromanspacer{} ภควโต เอตมตฺถํ อาโรเจสุํ.\csromanspacer{} อนุชานามิ ภิกฺขเว สพฺพํ ปาสาทปริโภคนฺติ.}

\prose{เตน โข ปน สมเยน รญฺโญ ปเสนทิสฺส โกสลสฺส อยฺยิกา กาลงฺกตา โหติ, ตสฺส กาลงฺกิริยาย สํฆสฺส พหุํ อกปฺปิยภณฺฑํ อุปนฺนํ โหติ, เสยฺยถิทํ, อาสนฺทิ ปลฺลงฺโก โคนโก จิตฺตโก ปฏิกา ปฏลิกา ตูลิกา วิกติกา อุทฺทโลมี เอกนฺตลิมี กฏฺฏิสฺสํ โกเสยฺยํ กุตฺตกํ หตฺถตฺถรํ อสฺสตฺถรํ รถตฺถรํ อชินปฺปเวณิ กทลิมิคปฺปวรปจฺจตฺถรณํ ส-อุตฺตรจฺฉทํ อุภโตโลหิตกูปธานํ.\csromanspacer{} ภควโต เอตมตฺถํ อาโรเจสุํ.\csromanspacer{} อนุชานามิ ภิกฺขเว อาสนฺทิยา ปาเท ฉินฺทิตฺวา ปริภุญฺชิตุํ.\csromanspacer{} ปลฺลงฺกสฺส วาเฬ ภินฺทิตฺวา ปริภุญฺชิตุํ.\csromanspacer{} ตูลิกํ วิชเฏตฺวา พิพฺโพหนํ กาตุํ.\csromanspacer{} อวเสสํ ภูมตฺถรณํ\footnote{ภุมฺมตฺถรณํ (สี, สฺยา)} กาตุนฺติ.}

\csromanmatikaanchor{mtk.150}
\csromantocmarkref{subsection}{อวิสฺสชฺชิยวตฺถุ}{mtk.150}
\bookmark[dest={mtk.150},level=2]{อวิสฺสชฺชิยวตฺถุ}
\csromanheader{2}{\textbf{อวิสฺสชฺชิยวตฺถุ}}
\proseitem{321}{\csromanfolio{321}เตน โข ปน สมเยน สาวตฺถิยา อวิทูเร อญฺญตรสฺมิํ คามกาวาเส อาวาสิกา ภิกฺขู อุปทฺทุตา โหนฺติ อาคนฺตุกคมิกานํ ภิกฺขูนํ เสนาสนํ ปญฺญเปนฺตา.\csromanspacer{} อถ โข เตสํ ภิกฺขูนํ เอตทโหสิ “เอตรหิ โข มยํ อาวุโส อุปทฺทุตา อาคนฺตุกคมิกานํ ภิกฺขูนํ เสนาสนํ ปญฺญเปนฺตา, หนฺท มยํ อาวุโส สพฺพํ สํฆิกํ เสนาสนํ เอกสฺส เทม, ตสฺส สนฺตกํ ปริภุญฺชิสฺสามา”ติ.\csromanspacer{} เต สพฺพํ สํฆิกํ เสนาสนํ เอกสฺส อทํสุ.\csromanspacer{} อาคนฺตุกา ภิกฺขู เต ภิกฺขู เอตทโวจุํ “อมฺหากํ อาวุโส เสนาสนํ ปญฺญเปถา”ติ.\csromanspacer{} นตฺถาวุโส สํฆิกํ เสนาสนํ, สพฺพํ อมฺหาหิ เอกสฺส ทินฺนนฺติ.\csromanspacer{} กิํ ปน ตุเมฺห อาวุโส สํฆิกํ เสนาสนํ วิสฺสชฺเชถาติ.\csromanspacer{} เอวมาวุโสติ.\csromanspacer{} เย เต ภิกฺขู อปฺปิจฺฉา ฯเปฯ เต อุชฺฌายนฺติ ขิยฺยนฺติ วิปาเจนฺติ “กถํ หิ นาม ภิกฺขู สํฆิกํ เสนาสนํ วิสฺสชฺเชสฺสนฺตี”ติ\footnote{วิสฺสชฺชิสฺสนฺตีติ (ก)}.\csromanspacer{} อถ โข เต ภิกฺขู ภควโต เอตมตฺถํ อาโรเจสุํ ฯเปฯ.\csromanspacer{} สจฺจํ กิร ภิกฺขเว ภิกฺขู สํฆิกํ เสนาสนํ วิสฺสชฺเชนฺตีติ.\csromanspacer{} สจฺจํ ภควาติ.\csromanspacer{} วิครหิ พุทฺโธ ภควา ฯเปฯ.\csromanspacer{} กถํ หิ นาม เต ภิกฺขเว โมฆปุริสา สํฆิกํ เสนาสนํ วิสฺสชฺเชสฺสนฺติ.\csromanspacer{} เนตํ ภิกฺขเว อปฺปสนฺนานํ วา ปสาทย ฯเปฯ วิครหิตฺวา ฯเปฯ ธมฺมิํ กถํ กตฺวา ภิกฺขู อามนฺเตสิ\csromandash{}}

\prose{ปญฺจิมานิ ภิกฺขเว อวิสฺสชฺชิยานิ, น วิสฺสชฺเชตพฺพานิ\footnote{น วิสฺสชฺชิตพฺพานิ (ก)} สํเฆน วา คเณน วา ปุคฺคเลน วา, วิสฺสชฺชิตานิปิ อวิสฺสชฺชิตานิ โหนฺติ, โย วิสฺสชฺเชยฺย, อาปตฺติ ถุลฺลจฺจยสฺส.\csromanspacer{} กตมานิ ปญฺจ.\csromanspacer{} อาราโม, อารามวตฺถุ, อิทํ ปฐมํ อวิสฺสชฺชิยํ, น วิสฺสชฺเชตพฺพํ สํเฆน วา คเณน วา ปุคฺคเลน วา วิสฺสชฺชิตมฺปิ อวิสฺสชฺชิตํ โหติ, โย วิสฺสชฺเชยฺย, อาปตฺติ ถุลฺลจฺจยสฺส.}

\prose{วิหาโร, วิหารวตฺถุ, อิทํ ทุติยํ อวิสฺสชฺชิยํ, น วิสฺสชฺเชตพฺพํ, สํเฆน วา คเณน วา ปุคฺคเลน วา, วิสฺสชฺชิตมฺปิ อวิสฺสชฺชิตํ โหติ, โย วิสฺสชฺเชยฺย, อาปตฺติ ถุลฺลจฺจยสฺส.}

\prose{มญฺโจ, ปีฐํ, ภิสิ, พิพฺโพหนํ, อิทํ ตติยํ อวิสฺสชฺชิยํ, น วิสฺสชฺเชตพฺพํ สํเฆน วา คเณน วา ปุคฺคเลน วา, วิสฺสชฺชิตมฺปิ อวิสฺสชฺชิตํ โหติ, โย วิสฺสชฺเชยฺย, อาปตฺติ ถุลฺลจฺจยสฺส.}

\prose{\csromanfolio{322}โลหกุมฺภี, โลหภาณกํ, โลหวารโก, โลหกฏาหํ, วาสิ, ปรสุ\footnote{ผรสุ (สี, สฺยา, ก)}, กุฐารี\footnote{กุธารี (ก)}, กุทาโล, นิขาทนํ, อิทํ จตุตฺถํ อวิสฺสชฺชิยํ, น วิสฺสชฺเชตพฺพํ สํเฆน วา คเณน วา ปุคฺคเลน วา, วิสฺสชฺชิตมฺปิ อวิสฺสชฺชิตํ โหติ, โย วิสฺสชฺเชยฺย, อาปตฺติ ถุลฺลจฺจยสฺส.}

\prose{วลฺลิ, เวฬุ, มุญฺชํ, ปพฺพชํ\footnote{พพฺพชํ (สี)}, มตฺติกา, ทารุภณฺฑํ, มตฺติกาภณฺฑํ, อิทํ ปญฺจมํ อวิสฺสชฺชิยํ, น วิสฺสชฺเชตพฺพํ สํเฆน วา คเณน วา ปุคฺคเลน วา, วิสฺสชฺชิตมฺปิ อวิสฺสชฺชิตํ โหติ, โย วิสฺสชฺเชยฺย, อาปตฺติ ถุลฺลจฺจยสฺส.\csromanspacer{} อิมานิ โข ภิกฺขเว ปญฺจ อวิสฺสชฺชิยานิ, น วิสฺสชฺเชตพฺพานิ สํเฆน วา คเณน วา ปุคฺคเลน วา, วิสฺสชฺชิตานิปิ อวิสฺสชฺชิตานิ โหนฺติ, โย วิสฺสชฺเชยฺย, อาปตฺติ ถุลฺลจฺจยสฺสาติ.}

\csromanmatikaanchor{mtk.151}
\csromantocmarkref{subsection}{อเวภงฺคิยวตฺถุ}{mtk.151}
\bookmark[dest={mtk.151},level=2]{อเวภงฺคิยวตฺถุ}
\csromanheader{2}{\textbf{อเวภงฺคิยวตฺถุ}}
\proseitem{322}{อถ โข ภควา สาวตฺถิยํ ยถาภิรนฺตํ วิหริตฺวา เยน กีฏาคิริ เตน จาริกํ ปกฺกามิ มหตา ภิกฺขุสํเฆน สทฺธิํ ปญฺจมตฺเตหิ ภิกฺขุสเตหิ สาริปุตฺตโมคฺคลฺลาเนหิ จ.\csromanspacer{} อสฺโสสุํ โข อสฺสชิปุนพฺพสุกา ภิกฺขู “ภควา กิร กีฏาคิริํ อาคจฺฉติ มหตา ภิกฺขุสํเฆน สทฺธิํ ปญฺจมตฺเตหิ ภิกฺขุสเตหิ สาริปุตฺตโมคฺคลฺลาเนหิ จ.\csromanspacer{} หนฺท มยํ อาวุโส สพฺพํ สํฆิกํ เสนาสนํ ภาเชม, ปาปิจฺฉา สาริปุตฺตโมคฺคลฺลานา ปาปิกานํ อิจฺฉานํ วสํ คตา, น มยํ เตสํ เสนาสนํ ปญฺญเปสฺสามา”ติ.\csromanspacer{} เต สพฺพํ สํฆิกํ เสนาสนํ ภาเชสุํ.\csromanspacer{} อถ โข ภควา อนุปุพฺเพน จาริกํ จรมาโน เยน กีฏาคิริ ตทวสริ.\csromanspacer{} อถ โข ภควา สมฺพหุเล ภิกฺขู อามนฺเตสิ\csromandash{}คจฺฉถ ตุเมฺห ภิกฺขเว, อสฺสชิปุนพฺพสุเก ภิกฺขู อุปสงฺกมิตฺวา เอวํ วเทถ “ภควา อาวุโส อาคจฺฉติ มหตา ภิกฺขุสํเฆน สทฺธิํ ปญฺจมตฺเตหิ ภิกฺขุสเตหิ สาริปุตฺตโมคฺคลฺลาเนหิ จ, ภควโต จ อาวุโส เสนาสนํ ปญฺญเปถ ภิกฺขุสํฆสฺส จ สาริปุตฺตโมคฺคลฺลานานญฺจา”ติ. “เอวํ ภนฺเต”ติ โข เต ภิกฺขู ภควโต ปฏิสฺสุตฺวา เยน อสฺสชิปุนพฺพสุกา ภิกฺขู เตนุปสงฺกมิํสุ, อุปสงฺกมิตฺวา อสฺสชิปุนพฺพสุเก ภิกฺขู เอตทโวจุํ “ภควา อาวุโส อาคจฺฉติ มหตา ภิกฺขุสํเฆน สทฺธิํ ปญฺจมตฺเตหิ ภิกฺขุสเตหิ สาริปุตฺตโมคฺคลฺลาเนหิ จ, ภควโต จ อาวุโส \csromanfolio{323}เสนาสนํ ปญฺญเปถ ภิกฺขุสํฆสฺส จ สาริปุตฺตโมคฺคลฺลานานญฺจา”ติ.\csromanspacer{} นตฺถาวุโส สํฆิกํ เสนาสนํ, สพฺพํ อเมฺหหิ ภาชิตํ, สฺวาคตํ อาวุโส ภควโต, ยสฺมิํ วิหาเร ภควา อิจฺฉิสฺสติ, ตสฺมิํ วิหาเร วสิสฺสติ, ปาปิจฺฉา สาริปุตฺตโมคฺคลฺลานา ปาปิกานํ อิจฺฉานํ วสํ คตา, น มยํ เตสํ เสนสนํ ปญฺญเปสฺสามาติ.\csromanspacer{} กิํ ปน ตุเมฺห อาวุโส สํฆิกํ เสนาสนํ ภาชิตฺถาติ.\csromanspacer{} เอวมาวุโสติ.\csromanspacer{} เย เต ภิกฺขู อปฺปิจฺฉา ฯเปฯ เต อุชฺฌายนฺติ ขิยฺยนฺติ วิปาเจนฺติ “กถํ หิ นาม อสฺสชิปุนพฺพสุกา ภิกฺขู สํฆิกํ เสนาสนํ ภาเชสฺสนฺตี”ติ.\csromanspacer{} อถ โข เต ภิกฺขู ภควโต เอตมตฺถํ อาโรเจสุํ ฯเปฯ.\csromanspacer{} สจฺจํ กิร ภิกฺขเว ฯเปฯ.\csromanspacer{} สจฺจํ ภควาติ.\csromanspacer{} วิครหิ พุทฺโธ ภควา ฯเปฯ.\csromanspacer{} กถํ หิ นาม เต ภิกฺขเว โมฆปุริสา สํฆิกํ เสนาสนํ ภาเชสฺสนฺติ.\csromanspacer{} เนตํ ภิกฺขเว อปฺปสนฺนานํ วา ปสาทาย ฯเปฯ วิครหิตฺวา ฯเปฯ ธมฺมิํ กถํ กตฺวา ภิกฺขู อามนฺเตสิ\csromandash{}}

\prose{ปญฺจิมานิ ภิกฺขเว อเวภงฺคิยานิ\footnote{อเวภงฺคิกานิ (ก)}, น วิภชิตพฺพานิ สํเฆน วา คเณน วา ปุคฺคเลน วา, วิภตฺตานิปิ อวิภตฺตานิ โหนฺติ, โย วิภเชยฺย, อาปตฺติ ถุลฺลจฺจยสฺส.\csromanspacer{} กตมานิ ปญฺจ.\csromanspacer{} อาราโม, อารามวตฺถุ, อิทํ ปฐมํ อเวภงฺคิยํ, น วิภชิตพฺพํ สํเฆน วา คเณน วา ปุคฺคเลน วา, วิภตฺตมฺปิ อวิภตฺตํ โหติ, โย วิภเชยฺย, อาปตฺติ ถุลฺลจฺจยสฺส.}

\prose{วิหาโร, วิหารวตฺถุ, อิทํ ทุติยํ อเวภงฺคิยํ, น วิภชิตพฺพํ สํเฆน วา คเณน วา ปุคฺคเลน วา, วิภตฺตมฺปิ อวิภตฺตํ โหติ, โย วิภเชยฺย, อาปตฺติ ถุลฺลจฺจยสฺส.}

\prose{มญฺโจ, ปีฐํ, ภิสิ, พิพฺโพหนํ, อิทํ ตติยํ อเวงฺคิยํ, น วิภชิตพฺพํ สํเฆน วา คเณน วา ปุคฺคเลน วา, วิภตฺตมฺปิ อวิภตฺตํ โหติ, โย วิภเชยฺย, อาปตฺติ ถุลฺลจฺจยสฺส.}

\prose{โลหกุมฺภี, โลหภาณกํ, โลหวารโก, โลหกฏาหํ, วาสี, ปรสุ, กุฐารี, กุทาโล, นิขาทนํ, อิทํ จตุตฺถํ อเวภงฺคิยํ, น วิภชิตพฺพํ สํเฆน วา คเณน วา ปุคฺคเลน วา, วิภตฺตมฺปิ อวิภตฺตํ โหติ, โย วิภเชยฺย, อาปตฺติ ถุลฺลจฺจยสฺส.}

\prose{วลฺลี, เวฬุ, มุญฺชํ, ปพฺพชํ, ติณํ, มตฺติกา, ทารุภณฺฑํ, มตฺติกาภณฺฑํ, อิทํ ปญฺจมํ อเวภงฺคิยํ, น วิภชิตพฺพํ สํเฆน วา คเณน วา ปุคฺคเลน วา, \csromanfolio{324}วิภตฺตํปิ อวิภตฺตํ โหติ, โย วิภเชยฺย, อาปตฺติ ถุลฺลจฺจยสฺส.\csromanspacer{} อิมานิ โข ภิกฺขเว ปญฺจ ภิกฺขเว ปญฺจ อเวภงฺคิยานิ, น วิภชิตพฺพานิ สํเฆน วา คเณน วา ปุคฺคเลน วา, วิภตฺตานิปิ อวิภตฺตานิ โหนฺติ, โย วิภเชยฺย, อาปตฺติ ถุลฺลจฺจยสฺสาติ.}

\csromanmatikaanchor{mtk.152}
\csromantocmarkref{subsection}{นวกมฺมทานกถา}{mtk.152}
\bookmark[dest={mtk.152},level=2]{นวกมฺมทานกถา}
\csromanheader{2}{\textbf{นวกมฺมทานกถา}}
\proseitem{323}{อถ โข ภควา กีฏาคิริสฺมิํ ยถาภิรนฺตํ วิหริตฺวา เยน อาฬวี เตน จาริกํ ปกฺกามิ.\csromanspacer{} อนุปุพฺเพน จาริกํ จรมาโน เยน อาฬวี ตทวสริ, ตตฺร สุทํ ภควา อาฬวิยํ วิหรติ อคฺคาฬเว เจติเย.\csromanspacer{} เตน โข ปน สมเยน อาฬวกา\footnote{อาฬวิกา (สฺยา, ก)} ภิกฺขู เอวรุปานิ นวกมฺมานิ เทนฺติ, ปิณฺฑนิกฺเขปนมตฺเตนปิ นวกมฺมํ เทนฺติ, กุฏฺฏเลปนมตฺเตนปิ นวกมฺมํ เทนฺติ, ทฺวารฏฺฐปนมตฺเตนปิ นวกมฺมํ เทนฺติ, อคฺคฬวฏฺฏิกรณมตฺเตนปิ นวกมฺมํ เทนฺติ, อาโลกสนฺธิกรณมตฺเตนปิ นวกมฺมํ เทนฺติ, เสตวณฺณกรณมตฺเตนปิ นวกมฺมํ เทนฺติ, กาฬวณฺณกรณมตฺเตนปิ นวกมฺมํ เทนฺติ, เครุกปริกมฺมกรณ- มตฺเตนปิ นวกมฺมํ เทนฺติ, ฉาทนมตฺเตนปิ นวกมฺมํ เทนฺติ, พนฺธนมตฺเตนปิ นวกมฺมํ เทนฺติ, ภณฺฑิกาฏฺฐปนมตฺเตนปิ นวกมฺมํ เทนฺติ, ขณฺฑผุลฺลปฏิสงฺขรณมตฺเตนปิ นวกมฺมํ เทนฺติ, ปริภณฺฑกรณมตฺเตนปิ นวกมฺมํ เทนฺติ, วีสติวสฺสิกมฺปิ นวกมฺมํ เทนฺติ, ติํสวสฺสิกมฺปิ นวกมฺมํ เทนฺติ, ยาวชีวิกมฺปิ นวกมฺมํ เทนฺติ, ธูมกาลิกมฺปิ ปริโยสิตํ วิหารํ นวกมฺมํ เทนฺติ.\csromanspacer{} เย เต ภิกฺขู อปฺปิจฺฉา ฯเปฯ เต อุชฺฌายนฺติ ขิยฺยนฺติ วิปาเจนฺติ “กถํ หิ นาม อาฬวกา ภิกฺขู เอวรูปานิ นวกมฺมานิ ทสฺสนฺติ, ปิณฺฑนิกฺเขปนมตฺเตนปิ นวกมฺมํ ทสฺสนฺติ, กุฏฺฏเลปนมตฺเตนปิ.\csromanspacer{} ทฺวารฏฺฐปนมตฺเตนปิ.\csromanspacer{} อคฺคฬวฏฺฏิกรณมตฺเตนปิ.\csromanspacer{} อาโลกสนฺธิกรณ- มตฺเตนปิ.\csromanspacer{} เสตวณฺณกรณมตฺเตนปิ.\csromanspacer{} กาฬวณฺณกรณมตฺเตนปิ.\csromanspacer{} เครุกปริกมฺมกรณมตฺเตนปิ.\csromanspacer{} ฉาทนมตฺเตนปิ.\csromanspacer{} พนฺธนมตฺเตนปิ.\csromanspacer{} ภณฺฑิกาฏฺฐปนมตฺเตนปิ.\csromanspacer{} ขณฺฑผุลฺลปฏิสงฺขรณมตฺเตนปิ.\csromanspacer{} ปริภณฺฑกรณมตฺเตนปิ.\csromanspacer{} วีสติวสฺสิกมฺปิ.\csromanspacer{} ติํสวสฺสิกมฺปิ.\csromanspacer{} ยาวชีวิกมฺปิ.\csromanspacer{} ธูมกาลิ กมฺปิ ปริโยสิตํ วิหารํ นวกมฺมํ ทสฺสนฺตี”ติ.\csromanspacer{} ภควโต เอตมตฺถํ อาโรเจสุํ ฯเปฯ.\csromanspacer{} สจฺจํ กิร ภิกฺขเว ฯเปฯ.\csromanspacer{} สจฺจํ ภควาติ ฯเปฯ วิครหิตฺวา ฯเปฯ ธมฺมิํ กถํ กตฺวา ภิกฺขู อามนฺเตสิ “น ภิกฺขเว ปิณฺฑนิกฺเขปนมตฺเตน นวกมฺมํ ทาตพฺพํ ฯเปฯ.\csromanspacer{} น กุฏฺฏเลปนมตฺเตน นวกมฺมํ ทาตพฺพํ. \csromanfolio{325}น ทฺวารฏฺฐปนมตฺเตน นวกมฺมํ ทาตพฺพํ.\csromanspacer{} น อคฺคฬวฏฺฏิกรณ- มตฺเตน นวกมฺมํ ทาตพฺพํ.\csromanspacer{} น อาโลกสนฺธิกรณมตฺเตน นวกมฺมํ ทาตพฺพํ.\csromanspacer{} น เสตวณฺณกรณมตฺเตน นวกมฺมํ ทาตพฺพํ.\csromanspacer{} น กาฬวณฺณกรณมตฺเตน นวกมฺมํ ทาตพฺพํ.\csromanspacer{} น เครุกปริกมฺมกรณมตฺเตน นวกมฺมํ ทาตพฺพํ.\csromanspacer{} น ฉาทนมตฺเตน นวกมฺมํ ทาตพฺพํ.\csromanspacer{} น พนฺธนมตฺเตน นวกมฺมํ ทาตพฺพํ.\csromanspacer{} น ภณฺฑิกาฏฺฐปนมตฺเตน นวกมฺมํ ทาตพฺพํ.\csromanspacer{} น ขณฺฑผุลฺลปฏิสงฺขรณมตฺเตน นวกมฺมํ ทาตพฺพํ.\csromanspacer{} น ปริภณฺฑกรณมตฺเตน นวกมฺมํ ทาตพฺพํ.\csromanspacer{} น วีสติวสฺสิกํ นวกมฺมํ ทาตพฺพํ.\csromanspacer{} น ติํสวสฺสิกํ นวกมฺมํ ทาตพฺพํ.\csromanspacer{} น ยาวชีวิกํ นวกมฺมํ ทาตพฺพํ.\csromanspacer{} น ธูมกาลิกมฺปิ ปริโยสิตํ วิหารํ นวกมฺมํ ทาตพฺพํ, โย ทเทยฺย, อาปตฺติ ทุกฺกฏสฺส.\csromanspacer{} อนุชานามิ ภิกฺขเว อกตํ วา วิปฺปกตํ วา นวกมฺมํ ทาตุํ, ขุทฺทเก วิหาเร กมฺมํ โอโลเกตฺวา ฉปฺปญฺจวสฺสิกํ นวกมฺมํ ทาตุํ, อฑฺฒโยเค กมฺมํ โอโลเกตฺวา สตฺตฏฺฐวสฺสิกํ นวกมฺมํ ทาตุํ, มหลฺลเก วิหาเร ปาสาเท วา กมฺมํ โอโลเกตฺวา ทสทฺวาทสวสฺสิกํ นวกมฺมํ ทาตุนฺ”ติ.}

\prose{เตน โข ปน สมเยน ภิกฺขู สพฺเพ วิหาเร นวกมฺมํ เทนฺติ.\csromanspacer{} ภควโต เอตมตฺถํ อาโรเจสุํ.\csromanspacer{} น ภิกฺขเว สพฺเพ วิหาเร นวกมฺมํ ทาตพฺพํ, โย ทเทยฺย, อาปตฺติ ทุกฺกฏสฺสาติ.}

\prose{เตน โข ปน สมเยน ภิกฺขู เอกสฺส เทฺว เทนฺติ.\csromanspacer{} ภควโต เอตมตฺถํ อาโรเจสุํ.\csromanspacer{} น ภิกฺขเว เอกสฺส เทฺว ทาตพฺพา, โย ทเทยฺย, อาปตฺติ ทุกฺกฏสฺสาติ.}

\prose{เตน โข ปน สมเยน ภิกฺขู นวกมฺมํ คเหตฺวา อญฺญํ วาเสนฺติ.\csromanspacer{} ภควโต เอตมตฺถํ อาโรเจสุํ.\csromanspacer{} น ภิกฺขเว นวกมฺมํ คเหตฺวา อญฺโญ วาเสตพฺโพ, โย วาเสยฺย, อาปตฺติ ทุกฺกฏสฺสาติ.}

\prose{เตน โข ปน สมเยน ภิกฺขู นวกมฺมํ คเหตฺวา สํฆิกํ ปฏิพาหนฺติ.\csromanspacer{} ภควโต เอตมตฺถํ อาโรเจสุํ.\csromanspacer{} น ภิกฺขเว นวกมฺมํ คเหตฺวา สํฆิกํ ปฏิพาหิตพฺพํ, โย ปฏิพาเหยฺย, อาปตฺติ ทุกฺกฏสฺส.\csromanspacer{} อนุชานามิ ภิกฺขเว เอกํ วรเสยฺยํ คเหตุนฺติ.}

\prose{เตน โข ปน สมเยน ภิกฺขู นิสฺสีเม ฐิตสฺส นวกมฺมํ เทนฺติ.\csromanspacer{} ภควโต เอตมตฺถํ อาโรเจสุํ.\csromanspacer{} น ภิกฺขเว นิสฺสีเม ฐิตสฺส นวกมฺมํ ทาตพฺพํ, โย ทเทยฺย อาปตฺติ ทุกฺกฏสฺสาติ.}

\prose{\csromanfolio{326}เตน โข ปน สมเยน ภิกฺขู นวกมฺมํ คเหตฺวา สพฺพกาลํ ปฏิพาหนฺติ.\csromanspacer{} ภควโต เอตมตฺถํ อาโรเจสุํ.\csromanspacer{} น ภิกฺขเว นวกมฺมํ คเหตฺวา สพฺพกาลํ ปฏิพาหิตพฺพํ, โย ปฏิพาเหยฺย, อาปตฺติ ทุกฺกฏสฺส.\csromanspacer{} อนุชานามิ ภิกฺขเว วสฺสานํ เตมาสํ ปฏิพาหิตุํ.\csromanspacer{} อุตุกาลํ ปน น ปฏิพาหิตุนฺติ.}

\prose{เตน โข ปน สมเยน ภิกฺขู นวกมฺมํ คเหตฺวา ปกฺกมนฺติปิ วิพฺภมนฺติปิ, กาลํปิ กโรนฺติ, สามเณราปิ ปฏิชานนฺติ, สิกฺขํ ปจฺจกฺขาตกาปิ ปฏิชานนฺติ, อนฺติมวตฺถุํ อชฺฌาปนฺนกาปิ ปฏิชานนฺติ, อุมฺมตฺตกาปิ ปฏิชานนฺติ, ขิตฺตจิตฺตาปิ ปฏิชานนฺติ, เวทนาฏฺฏาปิ ปฏิชานนฺติ, อาปตฺติยา อทสฺสเน อุกฺขิตฺตกาปิ ปฏิชานนฺติ, อาปตฺติยา อปฺปฏิกมฺเม อุกฺขิตฺตกาปิ ปฏิชานนฺติ, ปาปิกาย ทิฏฺฐิยา อปฺปฏินิสฺสคฺเค อุกฺขิตฺตกาปิ ปฏิชานนฺติ, ปณฺฑกาปิ ปฏิชานนฺติ, เถยฺยสํวาสกาปิ ปฏิชานนฺติ, ติตฺถิยปกฺกนฺตกาปิ ปฏิชานนฺติ, ติรจฺฉานคตาปิ ปฏิชานนฺติ, มาตุฆาตกาปิ ปฏิชานนฺติ, ปิตุฆาตกาปิ ปฏิชานนฺติ, อรหนฺตฆาตกาปิ ปฏิชานนฺติ, ภิกฺขุนิทูสกาปิ ปฏิชานนฺติ, สํฆเภทกาปิ ปฏิชานนฺติ, โลหิตุปฺปาทกาปิ ปฏิชานนฺติ, อุภโตพฺยญฺชนกาปิ ปฏิชานนฺติ.\csromanspacer{} ภควโต เอตมตฺถํ อาโรเจสุํ.}

\prose{อิธ ปน ภิกฺขเว ภิกฺขุ นวกมฺมํ คเหตฺวา ปกฺกมติ, “มา สํฆสฺส หายี”ติ อญฺญสฺส ทาตพฺพํ.}

\prose{อิธ ปน ภิกฺขเว ภิกฺขุ นวกมฺมํ คเหตฺวา วิพฺภมติ ฯเปฯ กาลํ กโรติ.\csromanspacer{} สามเณโร ปฏิชานาติ.\csromanspacer{} สิกฺขํ ปจฺจกฺขาตโก ปฏิชานาติ.\csromanspacer{} อนฺติมวตฺถุํ อชฺฌาปนฺนโก ปฏิชานาติ.\csromanspacer{} อุมฺมตฺตโก ปฏิชานาติ.\csromanspacer{} ขิตฺตจิตฺโต ปฏิชานาติ.\csromanspacer{} เวทนาฏฺโฏ ปฏิชานาติ.\csromanspacer{} อาปตฺติยา อทสฺสเน อุกฺขิตฺตโก ปฏิชานาติ.\csromanspacer{} อาปตฺติยา อปฺปฏิกมฺเม อุกฺขิตฺตโก ปฏิชานาติ.\csromanspacer{} ปาปิกาย ทิฏฺฐิยา อปฺปฏินิสฺสคฺเค อุกฺขิตฺตโก ปฏิชานาติ.\csromanspacer{} ปณฺฑโก ปฏิชานาติ.\csromanspacer{} เถยฺยสํวาสโก ปฏิชานาติ.\csromanspacer{} ติตฺถิยปกฺกนฺตโก ปฏิชานาติ.\csromanspacer{} ติรจฺฉานคโต ปฏิชานาติ.\csromanspacer{} มาตุฆาตโก ปฏิชานาติ.\csromanspacer{} ปิตุฆาตโก ปฏิชานาติ.\csromanspacer{} อรหนฺตฆาตโก ปฏิชานาติ.\csromanspacer{} ภิกฺขุนิทูสโก ปฏิชานาติ.\csromanspacer{} สํฆเภทโก ปฏิชานาติ.\csromanspacer{} โลหิตุปฺปาทโก ปฏิชานาติ.\csromanspacer{} อุภโตพฺยญฺชนโก ปฏิชานาติ, “มา สํฆสฺส หายี”ติ อญฺญสฺส ทาตพฺพํ.}

\prose{\csromanfolio{327}อิธ ปน ภิกฺขเว ภิกฺขุ นวกมฺมํ คเหตฺวา วิปฺปกเต ปกฺกมติ, “มา สํฆสฺส หายี”ติ อญฺญสฺส ทาตพฺพํ.}

\prose{อิธ ปน ภิกฺขเว ภิกฺขุ นวกมฺมํ คเหตฺวา วิปฺปกเต วิพฺภมติ ฯเปฯ.\csromanspacer{} อุภโตพฺยญฺชนโก ปฏิชานาติ, “มา สํฆสฺส หายี”ติ อญฺญสฺส ทาตพฺพํ.}

\prose{อิธ ปน ภิกฺขเว ภิกฺขุ นวกมฺมํ คเหตฺวา ปริโยสิเต ปกฺกมติ, ตสฺเสเวตํ.}

\prose{อิธ ปน ภิกฺขเว ภิกฺขุ นวกมฺมํ คเหตฺวา ปริโยสิเต วิพฺภมติ ฯเปฯ กาลํ กโรติ.\csromanspacer{} สามเณโร ปฏิชานาติ.\csromanspacer{} สิกฺขํ ปจฺจกฺขาตโก ปฏิชานาติ.\csromanspacer{} อนฺติมวตฺถุํ อชฺฌาปนฺนโก ปฏิชานาติ.\csromanspacer{} สํโฆ สามี.}

\prose{อิธ ปน ภิกฺขเว ภิกฺขุ นวกมฺมํ คเหตฺวา ปริโยสิเต อุมฺมตฺตโก ปฏิชานาติ.\csromanspacer{} ขิตฺตจิตฺโต ปฏิชานาติ.\csromanspacer{} เวทนาฏฺโฏ ปฏิชานาติ.\csromanspacer{} อาปตฺติยา อทสฺสเน อุกฺขิตฺตโก ปฏิชานาติ.\csromanspacer{} อาปตฺติยา อปฺปฏิกมฺเม อุกฺขิตฺตโก ปฏิชานาติ.\csromanspacer{} ปาปิกาย ทิฏฺฐิยา อปฺปฏินิสฺสคฺเค อุกฺขิตฺตโก ปฏิชานาติ, ตสฺเสเวตํ.}

\prose{อิธ ปน ภิกฺขเว ภิกฺขุ นวกมฺมํ คเหตฺวา ปริโยสิเต ปณฺฑโก ปฏิชานาติ.\csromanspacer{} เถยฺยสํวาสโก ปฏิชานาติ.\csromanspacer{} ติตฺถิยปกฺกนฺตโก ปฏิชานาติ.\csromanspacer{} ติรจฺฉานคโต ปฏิชานาติ.\csromanspacer{} มาตุฆาตโก ปฏิชานาติ.\csromanspacer{} ปิตุฆาตโก ปฏิชานาติ.\csromanspacer{} อรหนฺตฆาตโก ปฏิชานาติ.\csromanspacer{} ภิกฺขุนีทูสโก ปฏิชานาติ.\csromanspacer{} สํฆเภทโก ปฏิชานาติ.\csromanspacer{} โลหิตุปฺปาทโก ปฏิชานาติ.\csromanspacer{} อุภโตพฺยญฺชนโก ปฏิชานาติ, สํโฆ สามี.}

\csromanmatikaanchor{mtk.153}
\csromantocmarkref{subsection}{อญฺญตฺรปริโภคปฏิกฺเขปาทิ}{mtk.153}
\bookmark[dest={mtk.153},level=2]{อญฺญตฺรปริโภคปฏิกฺเขปาทิ}
\csromanheader{2}{\textbf{อญฺญตฺรปริโภคปฏิกฺเขปาทิ}}
\proseitem{324}{\csromansymbolfootnote{*}{วิ 1. 83 ปิฏฺเฐปิ.}เตน โข ปน สมเยน ภิกฺขู อญฺญตรสฺส อุปาสกสฺส วิหารปริโภคํ เสนาสนํ อญฺญตฺร ปริภุญฺชนฺติ.\csromanspacer{} อถ โข โส อุปาสโก อุชฺฌายติ ขิยฺยติ วิปาเจติ “กถํ หิ นาม ภทนฺตา อญฺญตฺร ปริโภคํ อญฺญตฺร ปริภุญฺชิสฺสนฺตี”ติ.\csromanspacer{} ภควโต เอตมตฺถํ อาโรเจสุํ.\csromanspacer{} น ภิกฺขเว อญฺญตร ปริโภโค อญฺญตฺร ปริภุญฺชิตพฺโพ, โย ปริภุญฺเชยฺย, อาปตฺติ ทุกฺกฏสฺสาติ.}

\prose{\csromanfolio{328}\csromansymbolfootnote{*}{วิ 1. 83 ปิฏฺเฐปิ.}เตน โข ปน สมเยน ภิกฺขู อุโปสถคฺคมฺปิ สนฺนิสชฺชมฺปิ หริตุํ กุกฺกุจฺจายนฺตา ฉมาย นิสีทนฺติ, คตฺตานิปิ จีวรานิปิ ปํสุกิตานิ โหนฺติ.\csromanspacer{} ภควโต เอตมตฺถํ อาโรเจสุํ.\csromanspacer{} อนุชานามิ ภิกฺขเว ตาวกาลิกํ หริตุนฺติ.}

\prose{เตน โข ปน สมเยน สํฆสฺส มหาวิหาโร อุนฺทฺริยติ, ภิกฺขู กุกฺกุจฺจายนฺตา เสนาสนํ นาติหรนฺติ\footnote{นาภิหรนฺติ (ก)}.\csromanspacer{} ภควโต เอตมตฺถํ อาโรเจสุํ.\csromanspacer{} อนุชานามิ ภิกฺขเว คุตฺตตฺถาย หริตุนฺติ.}

\prose{เตน โข ปน สมเยน สํฆสฺส เสนาสนปริกฺขาริโก มหคฺโฆ กมฺพโล อุปฺปนฺโน โหติ.\csromanspacer{} ภควโต เอตมตฺถํ อาโรเจสุํ.\csromanspacer{} อนุชานามิ ภิกฺขเว ผาติกมฺมตฺถาย ปริวตฺเตตุนฺติ.}

\prose{เตน โข ปน สมเยน สํฆสฺส เสนาสนปริกฺขาริกํ มฆคฺฆํ ทุสฺสํ อุปฺปนฺนํ โหติ.\csromanspacer{} ภควโต เอตมตฺถํ อาโรเจสุํ.\csromanspacer{} อนุชานามิ ภิกฺขเว ผาติกมฺมตฺถาย ปริวตฺเตตุนฺติ.}

\prose{เตน โข ปน สมเยน สํฆสฺส อจฺฉจมฺมํ อุปฺปนฺนํ โหติ.\csromanspacer{} ภควโต เอตมตฺถํ อาโรเจสุํ.\csromanspacer{} อนุชานามิ ภิกฺขเว ปาทปุจฺฉนิํ กาตุนฺติ.}

\prose{เตน โข ปน สมเยน สํฆสฺส จกฺกลิกํ อุปฺปนฺนํ โหติ.\csromanspacer{} ภควโต เอตมตฺถํ อาโรเจสุํ.\csromanspacer{} อนุชานามิ ภิกฺขเว ปาทปุญฺฉนิํ กาตุนฺติ.}

\prose{เตน โข ปน สมเยน สํฆสฺส โจฬกํ อุปฺปนฺนํ โหติ.\csromanspacer{} ภควโต เอตมตฺถํ อาโรเจสุํ.\csromanspacer{} อนุชานามิ ภิกฺขเว ปาทปุญฺฉนิํ กาตุนฺติ.}

\prose{เตน โข ปน สมเยน ภิกฺขู อโธเตหิ ปาเทหิ เสนาสนํ อกฺกมนฺติ, เสนาสนํ ทุสฺสติ.\csromanspacer{} ภควโต เอตมตฺถํ อาโรเจสุํ.\csromanspacer{} น ภิกฺขเว อโธเตหิ ปาเทหิ เสนาสนํ อกฺกมิตพฺพํ, โย อกฺกเมยฺย, อาปตฺติ ทุกฺกฏสฺสาติ.}

\prose{เตน โข ปน สมเยน ภิกฺขู อลฺเลหิ ปาเทหิ เสนาสนํ อกฺกมนฺติ, เสนาสนํ ทุสฺสติ.\csromanspacer{} ภควโต เอตมตฺถํ อาโรเจสุํ.\csromanspacer{} น ภิกฺขเว อลฺเลหิ ปาเทหิ เสนาสนํ อกฺกมิตพฺพํ, โย อกฺกเมยฺย, อาปตฺติ ทุกฺกฏสฺสาติ.}

\prose{\csromanfolio{329}เตน โข ปน สมเยน ภิกฺขู ส-อุปาหนา เสนาสนํ อกฺกมนฺติ, เสนาสนํ ทุสฺสติ.\csromanspacer{} ภควโต เอตมตฺถํ อาโรเจสุํ.\csromanspacer{} น ภิกฺขเว ส-อุปาหเนน เสนาสนํ อกฺกมิตพฺพํ, โย อกฺกเมยฺย, อาปตฺติ ทุกฺกฏสฺสาติ.}

\prose{เตน โข ปน สมเยน ภิกฺขู ปริกมฺมกตาย ภูมิยา นิฏฺฐุภนฺติ, วณฺโณ ทุสฺสติ.\csromanspacer{} ภควโต เอตมตฺถํ อาโรเจสุํ.\csromanspacer{} น ภิกฺขเว ปริกมฺมกตาย ภูมิยา นิฏฺฐุภิตพฺพํ, โย นิฏฺฐุเภยฺย, อาปตฺติ ทุกฺกฏสฺส.\csromanspacer{} อนุชานามิ ภิกฺขเว เขฬมลฺลกนฺติ.}

\prose{เตน โข ปน สมเยน มญฺจปาทาปิ ปีฐปาทาปิ ปริกมฺมกตํ ภูมิํ วิลิขนฺติ.\csromanspacer{} ภควโต เอตมตฺถํ อาโรเจสุํ.\csromanspacer{} อนุชานามิ ภิกฺขเว โจฬเกน ปลิเวเฐตุนฺติ.}

\prose{เตน โข ปน สมเยน ภิกฺขู ปริกมฺมกตํ ภิตฺติํ อปสฺเสนฺติ, วณฺโณ ทุสฺสติ.\csromanspacer{} ภควโต เอตมตฺถํ อาโรเจสุํ.\csromanspacer{} น ภิกฺขเว ปริกมฺมกตา ภิตฺติ อปสฺเสตพฺพา, โย อปสฺเสยฺย, อาปตฺติ ทุกฺกฏสฺส.\csromanspacer{} อนุชานามิ ภิกฺขเว อปสฺเสนผลกนฺติ.\csromanspacer{} อปสฺเสนผลกํ เหฏฺฐโต ภูมิํ วิลิขติ อุปริโต ภิตฺติํ จ.\csromanspacer{} ภควโต เอตมตฺถํ อาโรเจสุํ.\csromanspacer{} อนุชานามิ ภิกฺขเว เหฏฺฐโต จ อุปริโต จ โจฬเกน ปลิเวเฐตุนฺติ.}

\prose{เตน โข ปน สมเยน ภิกฺขู โธตปาทกา นิปชฺชิตุํ กุกฺกุจฺจายนฺติ.\csromanspacer{} ภควโต เอตมตฺถํ อาโรเจสุํ.\csromanspacer{} อนุชานามิ ภิกฺขเว ปจฺจตฺถริตฺวา นิปชฺชิตุนฺติ.}

\csromanmatikaanchor{mtk.154}
\csromantocmarkref{subsection}{สํฆภตฺตาทิ-อนุชานน}{mtk.154}
\bookmark[dest={mtk.154},level=2]{สํฆภตฺตาทิ-อนุชานน}
\csromanheader{2}{\textbf{สํฆภตฺตาทิ-อนุชานน}}
\proseitem{325}{อถ โข ภควา อาฬวิยํ ยถาภิรนฺตํ วิหริตฺวา เยน ราชคหํ เตน จาริกํ ปกฺกามิ.\csromanspacer{} อนุปุพฺเพน จาริกํ จรมาโน เยน ราชคหํ ตทวสริ, ตตฺร สุทํ ภควา ราชคเห วิหรติ เวฬุวเน กลนฺทกนิวาเป.\csromanspacer{} เตน โข ปน สมเยน ราชคหํ ทุพฺภิกฺขํ โหติ, มนุสฺสา น สกฺโกนฺติ สํฆภตฺตํ กาตุํ, อิจฺฉนฺติ อุทฺเทสภตฺตํ นิมนฺตนํ สลากภตฺตํ ปกฺขิกํ อุโปสถิกํ ปาฏิปทิกํ กาตุํ.\csromanspacer{} ภควโต เอตมตฺถํ อาโรเจสุํ.\csromanspacer{} อนุชานามิ ภิกฺขเว สํฆภตฺตํ อุทฺเทสภตฺตํ นิมนฺตนํ สลากภตฺตํ ปกฺขิกํ อุโปสถิกํ ปาฏิปทิกนฺติ.}

\csromanmatikaanchor{mtk.155}
\csromantocmarkref{subsection}{ภตฺตุทฺเทสกสมฺมุติ}{mtk.155}
\bookmark[dest={mtk.155},level=2]{ภตฺตุทฺเทสกสมฺมุติ}
\csromanheader{2}{\textbf{ภตฺตุทฺเทสกสมฺมุติ}}
\csromanmatikaanchor{mtk.156}
\csromantocmarkref{subsection}{เสนาสนปญฺญาปกาทิสมฺมุติ}{mtk.156}
\bookmark[dest={mtk.156},level=2]{เสนาสนปญฺญาปกาทิสมฺมุติ}
\proseitem{326}{\csromanfolio{330}เตน โข ปน สมเยน ฉพฺพคฺคิยา ภิกฺขู อตฺตโน วรภตฺตานิ คเหตฺวา ลามกานิ ภตฺตานิ ภิกฺขูนํ เทนฺติ.\csromanspacer{} ภควโต เอตมตฺถํ อาโรเจสุํ.\csromanspacer{} อนุชานามิ ภิกฺขเว ปญฺจหงฺเคหิ สมนฺนาคตํ ภิกฺขุํ ภตฺตุทฺเทสกํ สมฺมนฺนิตุํ, โย น ฉนฺทาคติํ คจฺเฉยฺย, น โทสาคติํ คจฺเฉยฺย, น โมหาคติํ คจฺเฉยฺย, น ภยาคติํ คจฺเฉยฺย, อุทฺทิฏฺฐานุทฺทิฏฺฐญฺจ ชาเนยฺย.\csromanspacer{} เอวญฺจ ปน ภิกฺขเว สมฺมนฺนิตพฺโพ, ปฐมํ ภิกฺขุ ยาจิตพฺโพ, ยาจิตฺวา พฺยตฺเตน ภิกฺขุนา ปฏิพเลน สํโฆ ญาเปตพฺโพ\csromandash{}}

\prose{“สุณาตุ เม ภนฺเต สํโฆ, ยทิ สํฆสฺส ปตฺตกลฺลํ, สํโฆ อิตฺถนฺนามํ ภิกฺขุํ ภตฺตุทฺเทสกํ สมฺมนฺเนยฺย, เอสา ญตฺติ.}

\prose{สุณาตุ เม ภนฺเต สํโฆ, สํโฆ อิตฺถนฺนามํ ภิกฺขุํ ภตฺตุทฺเทสกํ สมฺมนฺนติ, ยสฺสายสฺมโต ขมติ อิตฺถนฺนามสฺส ภิกฺขุโน ภตฺตุทฺเทสกสฺส สมฺมุติ, โส ตุณฺหสฺส, ยสฺส นกฺขมติ, โส ภาเสยฺย.}

\prose{สมฺมโต สํเฆน อิตฺถนฺนาโม ภิกฺขุ ภตฺตุทฺเทสโก, ขมติ สํฆสฺส, ตสฺมา ตุณฺหี, เอวเมตํ ธารยามี”ติ.}

\prose{อถ โข ภตฺตุทฺเทสกานํ ภิกฺขูนํ เอตทโหสิ “กถํ นุ โข ภตฺตํ อุทฺทิสิตพฺพนฺ”ติ.\csromanspacer{} ภควโต เอตมตฺถํ อาโรเจสุํ.\csromanspacer{} อนุชานามิ ภิกฺขเว สลากาย วา ปฏฺฏิกาย วา\footnote{ปฏิกาย วา (สฺยา)} อุปนิพนฺธิตฺวา โอปุญฺชิตฺวา ภตฺตํ อุทฺทิสิตุนฺติ.}

\csromanheader{2}{\textbf{เสนาสน ปญฺญาปกาทิสมฺมุติ}}
\proseitem{327}{เตน โข ปน สมเยน สํฆสฺส เสนาสนปญฺญาปโก น โหติ ฯเปฯ ภณฺฑาคาริโก น โหติ ฯเปฯ จีวรปฺปฏิคฺคาหโก น โหติ ฯเปฯ จีวรภาชโก น โหติ ฯเปฯ ยาคุภาชโก น โหติ ฯเปฯ ผลภาชโก น โหติ ฯเปฯ ขชฺชกภาชโก น โหติ, ขชฺชกํ อภาชิยมานํ นสฺสติ.\csromanspacer{} ภควโต เอตมตฺถํ อาโรเจสุํ.\csromanspacer{} อนุชานามิ ภิกฺขเว ปญฺจหงฺเคหิ สมนฺนาคตํ ภิกฺขุํ ขชฺชกภาชกํ สมฺมนฺนิตุํ, โย น ฉนฺทาคติํ คจฺเฉยฺย, น โทสาคติํ คจฺเฉยฺย, น โมหาคติํ คจฺเฉยฺย, น ภยาคติํ \csromanfolio{331}คจฺเฉยฺย, ภาชิตาภาชิตญฺจ ชาเนยฺย.\csromanspacer{} เอวญฺจ ปน ภิกฺขเว สมฺมนฺนิตพฺโพ, ปฐมํ ภิกฺขุ ยาจิตพฺโพ, ยาจิตฺวา พฺยตฺเตน ภิกฺขุนา ปฏิพเลน สํโฆ ญาเปตพฺโพ\csromandash{}}

\prose{“สุณาตุ เม ภนฺเต สํโฆ, ยทิ สํฆสฺส ปตฺตกลฺลํ, สํโฆ อิตฺถนฺนามํ ภิกฺขุํ ขชฺชกภาชกํ สมฺมนฺเนยฺย, เอสา ญตฺติ.}

\prose{สุณาตุ เม ภนฺเต สํโฆ, สํโฆ อิตฺถนฺนามํ ภิกฺกุํ ขชฺชกภาชกํ สมฺมนฺนติ, ยสฺสายสฺมโต ขมติ อิตฺถนฺนามสฺส ภิกฺขุโน ขชฺชกภาชกสฺส สมฺมุติ, โส ตุณฺหสฺส, ยสฺส นกฺขมติ, โส ภาเสยฺย.}

\prose{สมฺมโต สํเฆน อิตฺถนฺนาโม ภิกฺขุ ขชฺชกภาชโก, ขมติ สํฆสฺส, ตสฺมา ตุณฺหี, เอวเมตํ ธารยามี”ติ.}

\csromanmatikaanchor{mtk.157}
\csromantocmarkref{subsection}{อปฺปมตฺตกวิสฺสชฺชกสมฺมุติ}{mtk.157}
\bookmark[dest={mtk.157},level=2]{อปฺปมตฺตกวิสฺสชฺชกสมฺมุติ}
\csromanheader{2}{\textbf{อปฺปมตฺตกวิสฺสชฺชกสมฺมุติ}}
\proseitem{328}{เตน โข ปน สมเยน สํฆสฺส ภณฺฑาคาเร อปฺปมตฺตโก ปริกฺขาโร อุปฺปนฺโน\footnote{อุสฺสนฺโน (สฺยา)} โหติ.\csromanspacer{} ภควโต เอตมตฺถํ อาโรเจสุํ.\csromanspacer{} อนุชานามิ ภิกฺขเว ปญฺจหงฺเคหิ สมนฺนาคตํ ภิกฺขุํ อปฺปมตฺตกวิสฺสชฺชกํ สมฺมนฺนิตุํ, โย น ฉนฺทาคติํ คจฺเฉยฺย, น โทสาคติํ คจฺเฉยฺย, น โมหาคติํ คจฺเฉยฺย, น ภยาคติํ คจฺเฉยฺย, วิสฺสชฺชิตาวิสฺสชฺชิตญฺจ ชาเนยฺย.\csromanspacer{} เอวญฺจ ปน ภิกฺขเว สมฺมนฺนิตพฺโพ, ปฐมํ ภิกฺขุ ยาจิตพฺโพ, ยาจิตฺวา พฺยตฺเตน ภิกฺขุนา ปฏิพเลน สํโฆ ญาเปตพฺโพ\csromandash{}}

\prose{“สุณาตุ เม ภนฺเต สํโฆ, ยทิ สํฆสฺส ปตฺตกลฺลํ, สํโฆ อิตฺถนฺนามํ ภิกฺขุํ อปฺปมตฺตกวิสฺสชฺชกํ สมฺมนฺเนยฺย, เอสา ญตฺติ.}

\prose{สุณาตุ เม ภนฺเต สํโฆ, สํโฆ อิตฺถนฺนามํ ภิกฺขุํ อปฺปมตฺตกวิสฺสชฺชกํ สมฺมนฺนติ, ยสฺสายสฺมโต ขมติ อิตฺถนฺนามสฺส ภิกฺขุโน อปฺปมตฺตกวิสฺสชฺชกสฺส สมฺมุติ, โส ตุณฺหสฺส, ยสฺส นกฺขมติ, โส ภาเสยฺย.}

\prose{สมฺมโต สํเฆน อิตฺถนฺนาโม ภิกฺขุ อปฺปมตฺตกวิสฺสชฺชโก, ขมติ สํฆสฺส, ตสฺมา ตุณฺหี, เอวเมตํ ธารยามี”ติ.}

\prose{\csromanfolio{332}เตน อปฺปมตฺตกวิสฺสชฺชเกน ภิกฺขุนา เอกา\footnote{เอเตกา (สี)} สูจิ ทาตพฺพา, สตฺถกํ ทาตพฺพํ, อุปาหนา ทาตพฺพา, กายพนฺธนํ ทาตพฺพํ, อํสพทฺธโก ทาตพฺโพ, ปริสฺสาวนํ ทาตพฺพํ, ธมฺมกรโณ ทาตพฺโพ, กุสิ ทาตพฺพา, อฑฺฒกุสิ ทาตพฺพา, มณฺฑลํ ทาตพฺพํ, อฑฺฒมณฺฑลํ ทาตพฺพํ, อนุวาโต ทาตพฺโพ, ปริภณฺฑํ ทาตพฺพํ.\csromanspacer{} สเจ โหติ สํฆสฺส สปฺปิ วา เตลํ วา มธุ วา ผาณิตํ วา, สกิํ ปฏิสายิตุํ ทาตพฺพํ.\csromanspacer{} สเจ ปุนปิ อตฺโถ โหติ, ปุนปิ ทาตพฺพํ.}

\csromanmatikaanchor{mtk.158}
\csromantocmarkref{subsection}{สาฏิยคฺคาหาปกาทิสมฺมุติ}{mtk.158}
\bookmark[dest={mtk.158},level=2]{สาฏิยคฺคาหาปกาทิสมฺมุติ}
\csromanheader{2}{\textbf{สาฏิยคฺคาหาปกาทิสมฺมุติ}}
\proseitem{329}{เตน โข ปน สมเยน สํฆสฺส สาฏิยคฺคาหาปโก น โหติ ฯเปฯ ปตฺตคฺคาหาปโก น โหติ ฯเปฯ อารามิกเปสโก น โหติ ฯเปฯ สามเณรเปสโก น โหติ, สามเณรา อเปสิยมานา กมฺมํ น กโรนฺติ.\csromanspacer{} ภควโต เอตมตฺถํ อาโรเจสุํ.\csromanspacer{} อนุชานามิ ภิกฺขเว ปญฺจหงฺเคหิ สมนฺนาคตํ ภิกฺขุํ สามเณรเปสกํ สมฺมนฺนิตุํ, โย น ฉนฺทาคติํ คจฺเฉยฺย, น โทสาคติํ คจฺเฉยฺย, น โมหาคติํ คจฺเฉยฺย, น ภยาคติํ คจฺเฉยฺย, เปสิตาเปสิตญฺจ ชาเนยฺย.\csromanspacer{} เอวญฺจ ปน ภิกฺขเว สมฺมนฺนิตพฺโพ, ปฐมํ ภิกฺขุ ยาจิตพฺโพ, ยาจิตฺวา พฺยตฺเตน ภิกฺขุนา ปฏิพเลน สํโฆ ญาเปตพฺโพ\csromandash{}}

\prose{“สุณาตุ เม ภนฺเต สํโฆ, ยทิ สํฆสฺส ปตฺตกลฺลํ, สํโฆ อิตฺถนฺนามํ ภิกฺขุํ สามเณรเปสกํ สมฺมนฺเนยฺย, เอสา ญตฺติ.}

\prose{สุณาตุ เม ภนฺเต สํโฆ, สํโฆ อิตฺถนฺนามํ ภิกฺขุํ สามเณรเปสกํ สมฺมนฺนติ, ยสฺสายสฺมโต ขมติ อิตฺถนฺนามสฺส ภิกฺขุโน สามเณรเปสกสฺส สมฺมุติ, โส ตุณฺหสฺส, ยสฺส นกฺขมติ, โส ภาเสยฺย.}

\prose{สมฺมโต สํเฆน อิตฺถนฺนาโม ภิกฺขุ สามเณรเปสโก, ขมติ สํฆสฺส, ตสฺมา ตุณฺหี, เอวเมตํ ธารยามี”ติ.}

\nitthitam{ตติยภาณวาโร นิฏฺฐิโต.}
\vspace{\tipitakaparskip}
\csromancenter{เสนาสนกฺขนฺธโก ฉฏฺโฐ.}
\titlehead{\textbf{ตสฺสุทฺทานํ}}
\csromanmatikaanchor{mtk.159}
\csromantocmarkref{section}{อุทฺทานคาถา}{mtk.159}
\bookmark[dest={mtk.159},level=1]{อุทฺทานคาถา}
\csromangathameasure{วิหารํ พุทฺธเสฏฺเฐน, อปญฺญตฺตํ ตทา อหุ. \\ ตหํ ตหํ นิกฺขมนฺติ, วาสา เต ชินสาวกา. \\ เสฏฺฐี คหปติ ทิสฺวา, ภิกฺขูนํ อิทมพฺรวิ. \\ การาเปยฺยํ วเสยฺยาถ, ปฏิปุจฺฉิํสุ นายกํ. \\ วิหารํ อฑฺฒโยคญฺจ, ปาสาทํ หมฺมิยํ คุหํ. \\ ปญฺจเลณํ อนุญฺญาสิ, วิหาเร เสฏฺฐิ การยิ. \\ ชโน วิหารํ กาเรติ, อกวาฏํ อสํวุตํ. \\ กวาฏํ ปิฏฺฐสํฆาฏํ, อุทุกฺขลญฺจ อุตฺตริ. \\ อาวิญฺฉนจฺฉิทฺทํ รชฺชุํ, วฏฺฏิญฺจ กปิสีสกํ. \\ สูจิฆฏิตาฬจฺฉิทฺทํ, โลหกฏฺฐวิสาณกํ. \\ ยนฺตกํ สูจิกญฺเจว, ฉทนํ อุลฺลิตฺตาวลิตฺตํ. \\ เวทิชาลสลากญฺจ, จกฺกลิ สนฺถเรน จ. \\ มิฑฺฒิ พิทลมญฺจญฺจ, โสสานิกมสารโก. \\ พุนฺทิกุฬิรปาทญฺจ, อาหจฺจาสนฺทิ อุจฺจเก. \\ สตฺตงฺโค จ ภทฺทปีฐํ, ปีฐเกฬกปาทกํ. \\ อามลาผลกา โกจฺฉา, ปลาลปีฐเมว จ. \\ อุจฺจาหิปฏิปาทกา, อฏฺฐงฺคุลิ จ ปาทกา. \\ สุตฺตํ อฏฺฐปทํ โจฬํ, ตูลิกํ อฑฺฒกายิกํ. \\ คิรคฺโค ภิสิโย จาปิ, ทุสฺสํ เสนาสนํปิ จ. \\ โอนทฺธํ เหฏฺฐา ปตติ, อุปฺปาเฏตฺวา หรนฺติ จ. \\ ภตฺติญฺจ หตฺถภตฺติญฺจ, อนุญฺญาสิ ตถาคโต. \\ ติตฺถิยา วิหาเร จาปิ, ถุสํ สณฺหญฺจ มตฺติกา. \\ อิกฺกาสํ ปาณิกํ กุณฺฑํ, สาสปํ สิตฺถเตลกํ. \\ อุสฺสนฺเน ปจฺจุทฺธริตุํ, ปรุสํ คณฺฑุมตฺติกํ. \\ อิกฺกาสํ ปฏิภานญฺจ, นีจา จโย จ อารุหํ. \\ ปริปตนฺติ อาฬกา, อฑฺฒกุฏฺฏํ ตโย ปุน. \\ ขุทฺทเก กุฏฺฏปาโท จ, โอวสฺสติ สรํ ขิลํ. \\ จีวรวํสํ รชฺชุญฺจ, อาฬินฺทํ กิฏิเกน จ. \\ อาลมฺพนํ ติณจุณฺณํ, เหฏฺฐามคฺเค นยํ กเร. \\ อชฺโฌกาเส โอตปฺปติ, สาลํ เหฏฺฐา จ ภาชนํ. \\ วิหาโร โกฏฺฐโก เจว, ปริเวณคฺคิสาลกํ. \\ อาราเม จ ปุน โกฏฺเฐ, เหฏฺฐญฺเญว นยํ กเร. \\ สุธํ อนาถปิณฺฑิ จ, สทฺโธ สีตวนํ อคา. \\ ทิฏฺฐธมฺโม นิมนฺเตสิ, สห สํเฆน นายกํ. \\ อาณาเปสนฺตรามคฺเค, อารามํ การยี คโณ. \\ เวสาลิยํ นวกมฺมํ, ปุรโต จ ปริคฺคหิ. \\ โก อรหติ ภตฺตคฺเค, ติตฺติรญฺจ อวนฺทิยา. \\ ปริคฺคหิตนฺตรฆรา, ตูโล สาวตฺถิ โอสริ. \\ ปติฏฺฐาเปสิ อารามํ, ภตฺตคฺเค จ โกลาหลํ. \\ คิลานา วรเสยฺยา จ, เลสา สตฺตรสา ตหิํ. \\ เกน นุ โข กถํ นุ โข, วิหารคฺเคน ภาชยิ. \\ ปริเวณํ อนุภาคญฺจ, อกามา ภาคํ โน ทเท. \\ นิสฺสีมํ สพฺพกาลญฺจ, คาหา เสนาสเน ตโย. \\ อุปนนฺโท จ วณฺเณสิ, ฐิตกา สมกาสนา. \\ สมานาสนิกา ภินฺทิํสุ, ติวคฺคา จ ทุวคฺคิกํ. \\ อสมานาสนิกา ทีฆํ, สาฬินฺทํ ปริภุญฺชิตุํ. \\ อยฺยิกา จ อวิทูเร, ภาชิตญฺจ กีฏาคิเร. \\ อาฬวี ปิณฺฑกกุฏฺเฏหิ, ทฺวาร-อคฺคฬวฏฺฏิกา. \\ อาโลกเสตกาฬญฺจ, เครุฉาทนพนฺธนา. \\ ภณฺฑิขณฺฑปริภณฺฑํ, วีส ติํสา จ กาลิกา. \\ โอสิเต อกตํ วิปฺปํ, ขุทฺเท ฉปฺปญฺจวสฺสิกํ. \\ อฑฺฒโยเค จ สตฺตฏฺฐ, มหลฺเล ทส ทฺวาทส. \\ สพฺพํ วิหารํ เอกสฺส, อญฺญํ วาเสนฺติ สํฆิกํ. \\ นิสฺสีมํ สพฺพกาลญฺจ, ปกฺกมิ วิพฺภมนฺติ จ. \\ กาลญฺจ สามเณรญฺจ, สิกฺขาปจฺจกฺข-อนฺติมํ. \\ อุมฺมตฺตขิตฺตจิตฺตา จ, เวทนาปตฺติทสฺสนา. \\ อปฺปฏิกมฺมทิฏฺฐิยา, ปณฺฑกา เถยฺยติตฺถิยา. \\ ติรจฺฉานมาตุปิตุ, อรหนฺตา จ ทูสกา. \\ เภทกา โลหิตุปฺปาทา, อุภโต จาปิ พฺยญฺชนกา. \\ มา สํฆสฺส ปริหายิ, กมฺมํ อญฺญสฺส ทาตเว. \\ วิปฺปกเต จ อญฺญสฺส, กเต ตสฺเสว ปกฺกเม. \\ วิพฺภมติ กาลงฺกโต, สามเณโร จ ชายติ. \\ ปจฺจกฺขาโต จ สิกฺขาย, อนฺติมชฺฌาปนฺนโก ยทิ. \\ สํโฆว สามิโก โหติ, อุมฺมตฺตขิตฺตเวทนา. \\ อทสฺสนาปฺปฏิกมฺเม, ทิฏฺฐิ ตสฺเสว โหติ ตํ. \\ ปณฺฑโก เถยฺยติตฺถี จ, ติรจฺฉานมาตุเปตฺติกํ. \\ ฆาตโก ทูสโก จาปิ, เภทโลหิตพฺยญฺชนา. \\ ปฏิชานาติ ยทิ โส, สํโฆว โหติ สามิโก. \\ หรนฺตญฺญตฺร กุกฺกุจฺจํ, อุนฺทฺริยติ จ กมฺพลํ. \\ ทุสฺสํ จ จมฺมจกฺกลี, โจฬกํ อกฺกมนฺติ จ. \\ อลฺลา อุปาหนานิฏฺฐุ, ลิขนฺติ อปสฺเสนฺติ จ. \\ อปสฺเสนํ ลิขเตว, โธตปจฺจตฺถเรน จ. \\ ราชคเห น สกฺโกนฺติ, ลามกํ ภตฺตุทฺเทสกํ. \\ กถํ นุ โข ปญฺญาปกํ, ภณฺฑาคาริกสมฺมุติ. \\ ปฏิคฺคาหภาชโก จาปิ, ยาคุ จ ผลภาชโก. \\ ขชฺชกภาชโก เจว, อปฺปมตฺตกวิสฺสชฺเช. \\ สาฏิยคฺคาหาปโก เจว, ตเถว ปตฺตคฺคาหโก. \\ อารามิกสามเณร, เปสกสฺส จ สมฺมุติ. \\ สพฺพาภิภู โลกวิทู, หิตจิตฺโต วินายโก.}
\csromangathagroup{\csromangathastanza{\csromanfolio{333}วิหารํ พุทฺธเสฏฺเฐน, อปญฺญตฺตํ ตทา อหุ. \\ ตหํ ตหํ นิกฺขมนฺติ, วาสา เต ชินสาวกา.}\csromangathastanza{เสฏฺฐี คหปติ ทิสฺวา, ภิกฺขูนํ อิทมพฺรวิ. \\ การาเปยฺยํ วเสยฺยาถ, ปฏิปุจฺฉิํสุ นายกํ.}\csromangathastanza{วิหารํ อฑฺฒโยคญฺจ, ปาสาทํ หมฺมิยํ คุหํ. \\ ปญฺจเลณํ อนุญฺญาสิ, วิหาเร เสฏฺฐิ การยิ.}\csromangathastanza{ชโน วิหารํ กาเรติ, อกวาฏํ อสํวุตํ. \\ กวาฏํ ปิฏฺฐสํฆาฏํ, อุทุกฺขลญฺจ อุตฺตริ.}\csromangathastanza{อาวิญฺฉนจฺฉิทฺทํ รชฺชุํ, วฏฺฏิญฺจ กปิสีสกํ. \\ สูจิฆฏิตาฬจฺฉิทฺทํ, โลหกฏฺฐวิสาณกํ.}\csromangathastanza{ยนฺตกํ สูจิกญฺเจว, ฉทนํ อุลฺลิตฺตาวลิตฺตํ. \\ เวทิชาลสลากญฺจ, จกฺกลิ สนฺถเรน จ.}\csromangathastanza{มิฑฺฒิ พิทลมญฺจญฺจ, โสสานิกมสารโก. \\ พุนฺทิกุฬิรปาทญฺจ, อาหจฺจาสนฺทิ อุจฺจเก.}\csromangathastanza{สตฺตงฺโค จ ภทฺทปีฐํ, ปีฐเกฬกปาทกํ. \\ อามลาผลกา โกจฺฉา, ปลาลปีฐเมว จ.}\csromangathastanza{อุจฺจาหิปฏิปาทกา, อฏฺฐงฺคุลิ จ ปาทกา. \\ สุตฺตํ อฏฺฐปทํ โจฬํ, ตูลิกํ อฑฺฒกายิกํ.}\csromangathastanza{คิรคฺโค ภิสิโย จาปิ, ทุสฺสํ เสนาสนํปิ จ. \\ โอนทฺธํ เหฏฺฐา ปตติ, อุปฺปาเฏตฺวา หรนฺติ จ.}\csromangathastanza{ภตฺติญฺจ หตฺถภตฺติญฺจ, อนุญฺญาสิ ตถาคโต. \\ ติตฺถิยา วิหาเร จาปิ, ถุสํ สณฺหญฺจ มตฺติกา.}\csromangathastanza{อิกฺกาสํ ปาณิกํ กุณฺฑํ, สาสปํ สิตฺถเตลกํ. \\ อุสฺสนฺเน ปจฺจุทฺธริตุํ, ปรุสํ คณฺฑุมตฺติกํ.}\csromangathastanza{อิกฺกาสํ ปฏิภานญฺจ, นีจา จโย จ อารุหํ. \\ ปริปตนฺติ อาฬกา, อฑฺฒกุฏฺฏํ ตโย ปุน.}\csromangathastanza{\csromanfolio{334}ขุทฺทเก กุฏฺฏปาโท จ, โอวสฺสติ สรํ ขิลํ. \\ จีวรวํสํ รชฺชุญฺจ, อาฬินฺทํ กิฏิเกน จ.}\csromangathastanza{อาลมฺพนํ ติณจุณฺณํ, เหฏฺฐามคฺเค นยํ กเร. \\ อชฺโฌกาเส โอตปฺปติ, สาลํ เหฏฺฐา จ ภาชนํ.}\csromangathastanza{วิหาโร โกฏฺฐโก เจว, ปริเวณคฺคิสาลกํ. \\ อาราเม จ ปุน โกฏฺเฐ, เหฏฺฐญฺเญว นยํ กเร.}\csromangathastanza{สุธํ อนาถปิณฺฑิ จ, สทฺโธ สีตวนํ อคา. \\ ทิฏฺฐธมฺโม นิมนฺเตสิ, สห สํเฆน นายกํ.}\csromangathastanza{อาณาเปสนฺตรามคฺเค, อารามํ การยี คโณ. \\ เวสาลิยํ นวกมฺมํ, ปุรโต จ ปริคฺคหิ.}\csromangathastanza{โก อรหติ ภตฺตคฺเค, ติตฺติรญฺจ อวนฺทิยา. \\ ปริคฺคหิตนฺตรฆรา, ตูโล สาวตฺถิ โอสริ.}\csromangathastanza{ปติฏฺฐาเปสิ อารามํ, ภตฺตคฺเค จ โกลาหลํ. \\ คิลานา วรเสยฺยา จ, เลสา สตฺตรสา ตหิํ.}\csromangathastanza{เกน นุ โข กถํ นุ โข, วิหารคฺเคน ภาชยิ. \\ ปริเวณํ อนุภาคญฺจ, อกามา ภาคํ โน ทเท.}\csromangathastanza{นิสฺสีมํ สพฺพกาลญฺจ, คาหา เสนาสเน ตโย. \\ อุปนนฺโท จ วณฺเณสิ, ฐิตกา สมกาสนา.}\csromangathastanza{สมานาสนิกา ภินฺทิํสุ, ติวคฺคา จ ทุวคฺคิกํ. \\ อสมานาสนิกา ทีฆํ, สาฬินฺทํ ปริภุญฺชิตุํ.}\csromangathastanza{อยฺยิกา จ อวิทูเร, ภาชิตญฺจ กีฏาคิเร. \\ อาฬวี ปิณฺฑกกุฏฺเฏหิ, ทฺวาร-อคฺคฬวฏฺฏิกา.}\csromangathastanza{อาโลกเสตกาฬญฺจ, เครุฉาทนพนฺธนา. \\ ภณฺฑิขณฺฑปริภณฺฑํ, วีส ติํสา จ กาลิกา.}\csromangathastanza{โอสิเต อกตํ วิปฺปํ, ขุทฺเท ฉปฺปญฺจวสฺสิกํ. \\ อฑฺฒโยเค จ สตฺตฏฺฐ, มหลฺเล ทส ทฺวาทส.}\csromangathastanza{สพฺพํ วิหารํ เอกสฺส, อญฺญํ วาเสนฺติ สํฆิกํ. \\ นิสฺสีมํ สพฺพกาลญฺจ, ปกฺกมิ วิพฺภมนฺติ จ.}\csromangathastanza{\csromanfolio{335}กาลญฺจ สามเณรญฺจ, สิกฺขาปจฺจกฺข-อนฺติมํ. \\ อุมฺมตฺตขิตฺตจิตฺตา จ, เวทนาปตฺติทสฺสนา.}\csromangathastanza{อปฺปฏิกมฺมทิฏฺฐิยา, ปณฺฑกา เถยฺยติตฺถิยา. \\ ติรจฺฉานมาตุปิตุ, อรหนฺตา จ ทูสกา.}\csromangathastanza{เภทกา โลหิตุปฺปาทา, อุภโต จาปิ พฺยญฺชนกา. \\ มา สํฆสฺส ปริหายิ, กมฺมํ อญฺญสฺส ทาตเว.}\csromangathastanza{วิปฺปกเต จ อญฺญสฺส, กเต ตสฺเสว ปกฺกเม. \\ วิพฺภมติ กาลงฺกโต, สามเณโร จ ชายติ.}\csromangathastanza{ปจฺจกฺขาโต จ สิกฺขาย, อนฺติมชฺฌาปนฺนโก ยทิ. \\ สํโฆว สามิโก โหติ, อุมฺมตฺตขิตฺตเวทนา.}\csromangathastanza{อทสฺสนาปฺปฏิกมฺเม, ทิฏฺฐิ ตสฺเสว โหติ ตํ. \\ ปณฺฑโก เถยฺยติตฺถี จ, ติรจฺฉานมาตุเปตฺติกํ.}\csromangathastanza{ฆาตโก ทูสโก จาปิ, เภทโลหิตพฺยญฺชนา. \\ ปฏิชานาติ ยทิ โส, สํโฆว โหติ สามิโก.}\csromangathastanza{หรนฺตญฺญตฺร กุกฺกุจฺจํ, อุนฺทฺริยติ จ กมฺพลํ. \\ ทุสฺสํ จ จมฺมจกฺกลี, โจฬกํ อกฺกมนฺติ จ.}\csromangathastanza{อลฺลา อุปาหนานิฏฺฐุ, ลิขนฺติ อปสฺเสนฺติ จ. \\ อปสฺเสนํ ลิขเตว, โธตปจฺจตฺถเรน จ.}\csromangathastanza{ราชคเห น สกฺโกนฺติ, ลามกํ ภตฺตุทฺเทสกํ. \\ กถํ นุ โข ปญฺญาปกํ, ภณฺฑาคาริกสมฺมุติ.}\csromangathastanza{ปฏิคฺคาหภาชโก จาปิ, ยาคุ จ ผลภาชโก. \\ ขชฺชกภาชโก เจว, อปฺปมตฺตกวิสฺสชฺเช.}\csromangathastanza{สาฏิยคฺคาหาปโก เจว, ตเถว ปตฺตคฺคาหโก. \\ อารามิกสามเณร, เปสกสฺส จ สมฺมุติ. \\ สพฺพาภิภู โลกวิทู, หิตจิตฺโต วินายโก.}}

\nitthitam{เลณตฺถญฺจ สุขตฺถญฺจ, ฌายิตุญฺจ วิปสฺสิตุนฺติ. เสนาสนกฺขนฺธโก นิฏฺฐิโต.}
\csromanmatikaanchor{mtk.160}
\csromantocmarknopage{chapter}{7. สํฆเภทกกฺขนฺธก}{mtk.160}
\bookmark[dest={mtk.160},level=0]{7. สํฆเภทกกฺขนฺธก}
\chapterheadpageread{7. สํฆเภทกกฺขนฺธก}
\csromanmatikaanchor{mtk.161}
\csromantocmarkref{section}{1. ปฐมภาณวาร}{mtk.161}
\bookmark[dest={mtk.161},level=1]{1. ปฐมภาณวาร}
\csromanheader{1}{1. ปฐมภาณวาร}
\csromanmatikaanchor{mtk.162}
\csromantocmarkref{subsection}{ฉสกฺยปพฺพชฺชากถา}{mtk.162}
\bookmark[dest={mtk.162},level=2]{ฉสกฺยปพฺพชฺชากถา}
\csromanheader{2}{\textbf{ฉสกฺยปพฺพชฺชากถา}}
\proseitem{330}{\csromanfolio{336}เตน สมเยน พุทฺโธ ภควา อนุปิยายํ วิหรติ อนุปิยํ นาม มลฺลานํ นิคโม.\csromanspacer{} เตน โข ปน สมเยน อภิญฺญาตา อภิญฺญาตา สกฺยกุมารา ภควนฺตํ ปพฺพชิตํ อนุปพฺพชนฺติ.\csromanspacer{} เตน โข ปน สมเยน มหานาโม จ สกฺโก อนุรุทฺโธ จ สกฺโก เทฺวภาติกา โหนฺติ.\csromanspacer{} อนุรุทฺโธ สกฺโก สุขุมาโล โหติ, ตสฺส ตโย ปาสาทา โหนฺติ เอโก เหมนฺติโก เอโก คิมฺหิโก เอโก วสฺสิโก, โส วสฺสิเก ปาสาเท จตฺตาโร มาเส\footnote{วสฺสิเก ปาสาเท วสฺสิเก จตฺตาโร มาเส (สี)} นิปฺปุริเสหิ ตูริเยหิ ปริจารยมาโน\footnote{ปริจาริยมาโน (ก)} น เหฏฺฐาปาสาทํ โอโรหติ.\csromanspacer{} อถ โข มหานามสฺส สกฺกสฺส เอตทโหสิ “เอตรหิ โข อภิญฺญาตา อภิญฺญาตา สกฺยกุมารา ภควนฺตํ ปพฺพชิตํ อนุปพฺพชนฺติ, อมฺหากญฺจ ปน กุลา นตฺถิ โกจิ อคารสฺมา อนคาริยํ ปพฺพชิโต, ยํนูนาหํ วา ปพฺพเชยฺยํ, อนุรุทฺโธ วา”ติ.\csromanspacer{} อถ โข มหานาโม สกฺโก เยน อนุรุทฺโธ สกฺโก เตนุปสงฺกมิ, อุปสงฺกมิตฺวา อนุรุทฺธํ สกฺกํ เอตทโวจ “เอตรหิ ตาต อนุรุทฺธ อภิญฺญาตา อภิญฺญาตา สกฺยกุมารา ภควนฺตํ ปพฺพชิตํ อนุปพฺพชนฺติ, อมฺหากญฺจ ปน กุลา นตฺถิ โกจิ อคารสฺมา อนคาริยํ ปพฺพชิโต, เตน หิ ตฺวํ วา ปพฺพช, อหํ วา ปพฺพชิสฺสามี”ติ, อหํ โข สุขุมาโล, นาหํ สกฺโกมิ อคารสฺมา อนคาริยํ ปพฺพชิตุํ, ตฺวํ ปพฺพชาหีติ.\csromanspacer{} เอหิ โข เต ตาต อนุรุทฺธ, ฆราวาสตฺถํ อนุสาสิสฺสามิ.\csromanspacer{} ปฐมํ เขตฺตํ กสาเปตพฺพํ, กสาเปตฺวา วปาเปตพฺพํ, วปาเปตฺวา อุทกํ อภิเนตพฺพํ, อุทกํ อภิเนตฺวา อุทกํ นินฺเนตพฺพํ, อุทกํ นินฺเนตฺวา นิทฺธาเปตพฺพํ, นิทฺธาเปตฺวา\footnote{นิฑฺฑเหตพฺพํ, นิฑฺฑตฺวา (สี)} ลวาเปตพฺพํ, ลวาเปตฺวา อุพฺพาหาเปตพฺพํ, อุพฺพาหาเปตฺวา ปุญฺชํ การาเปตพฺพํ, ปุญฺชํ การาเปตฺวา มทฺทาเปตพฺพํ, มทฺทาเปตฺวา ปลาลานิ อุทฺธราเปตพฺพานิ, ปลาลานิ อุทฺธราเปตฺวา ภุสิกา อุทฺธราเปตพฺพา, ภุสิกํ อุทฺธราเปตฺวา โอปุนาเปตพฺพํ\footnote{โอผุนาเปตพฺพํ (สฺยา, ก), โอผุณาเปตพฺพํ (โยชนา)}, โอปุนาเปตฺวา อติหราเปตพฺพํ, อติหราเปตฺวา \csromanfolio{337}อายติมฺปิ วสฺสํ เอวเมว กาตพฺพํ, อายติมฺปิ วสฺสํ เอวเมว กาตพฺพนฺติ.\csromanspacer{} น กมฺมา ขียนฺติ, น กมฺมานํ อนฺโต ปญฺญายติ, กทา กมฺมา ขียิสฺสนฺติ, กทา กมฺมานํ อนฺโต ปญฺญายิสฺสติ, กทา มยํ อปฺโปสฺสุกฺกา ปญฺจหิ กามคุเณหิ สมปฺปิตา สมงฺคีภูตา ปริจาเรสฺสามาติ.\csromanspacer{} น หิ ตาต อนุรุทฺธ กมฺมา ขียนฺติ, น กมฺมานํ อนฺโต ปญฺญายติ, อขีเณว กมฺเม ปิตโร จ ปิตามหา จ กาลงฺกตาติ.\csromanspacer{} เตน หิ ตฺวญฺเญว ฆราวาสตฺเถน อุปชานาหิ, อหํ อคารสฺมา อนคาริยํ ปพฺพชิสฺสามีติ.}

\prose{อถ โข อนุรุทฺโธ สกฺโก เยน มาตา เตนุปสงฺกมิ, อุปสงฺกมิตฺวา มาตรํ เอตทโวจ “อิจฺฉามหํ อมฺม อคารสฺมา อนคาริยํ ปพฺพชิตุํ, อนุชานาหิ มํ อคารสฺมา อนคาริยํ ปพฺพชฺชายา”ติ.\csromanspacer{} เอวํ วุตฺเต อนุรุทฺธสฺส สกฺกสฺส มาตา อนุรุทฺธํ สกฺกํ เอตทโวจ “ตุเมฺห โข เม ตาต อนุรุทฺธ เทฺว ปุตฺตา ปิยา มนาปา อปฺปฏิกูลา, มรเณนปิ โว อกามกา วินา ภวิสฺสามิ, กิํ ปนาหํ ตุเมฺห ชีวนฺเต อนุชานิสฺสามิ อคารสฺมา อนคาริยํ ปพฺพชฺชายา”ติ.\csromanspacer{} ทุติยมฺปิ โข ฯเปฯ.\csromanspacer{} ตติยมฺปิ โข อนุรุทฺโธ สกฺโก มาตรํ เอตทโวจ “อิจฺฉามหํ อมฺม อคารสฺมา อนคาริยํ ปพฺพชิตุํ, อนุชานาหิ มํ อคารสฺมา อนคาริยํ ปพฺพชฺชายา”ติ\footnote{อิมสฺส อนนฺตรํ เกสุจิ โปตฺถเกสุ เอวมฺปิ ปาโฐ ทิสฺสติ⎯ “เอวํ วุตฺเต อนุรุทฺธสฺส สกฺกสฺส มตา เอวมาห ‘สเจ ตาต อนุรุทฺธ ภทฺทิโย สกฺยราชา สกฺยานํ รชฺชํ กาเรติ อคารสฺมา อนคาริยํ ปพฺพชติ, เอวํ ตฺวมฺปิ ปพฺพชาหี’ติ”.}.\csromanspacer{} เตน โข ปน สมเยน ภทฺทิโย สกฺยราชา สกฺยานํ รชฺชํ กาเรสิ, โส จ อนุรุทฺธสฺส สกฺกสฺส สหาโย โหติ.\csromanspacer{} อถ โข อนุรุทฺธสฺส สกฺกสฺส มาตา “อยํ โข ภทฺทิโย สกฺยราชา สกฺยานํ รชฺชํ กาเรติ อนุรุทฺธสฺส สกฺกสฺส สหาโย, โส น อุสฺสหติ อคารสฺมา อนคาริยํ ปพฺพชิตุนฺ”ติ อนุรุทฺธํ สกฺกํ เอตทโวจ “สเจ ตาต อนุรุทฺธ ภทฺทิโย สกฺยราชา อคารสฺมา อนคาริยํ ปพฺพชติ, เอวํ ตฺวมฺปิ ปพฺพชาหี”ติ.\csromanspacer{} อถ โข อนุรุทฺโธ สกฺโก เยน ภทฺทิโย สกฺยราชา เตนุปสงฺกมิ, อุปสงฺกมิตฺวา ภทฺทิยํ สกฺยราชานํ เอตทโวจ “มม โข สมฺม ปพฺพชฺชา ตว ปฏิพทฺธา”ติ.\csromanspacer{} สเจ เต สมฺม ปพฺพชฺชา มม ปฏิพทฺธา วา อปฺปฏิพทฺธา วา สา โหตุ, อหํ ตยา, ยถา สุขํ ปพฺพชาหีติ.\csromanspacer{} เอหิ สมฺม อุโภ อคารสฺมา อนคาริยํ ปพฺพชิสฺสามาติ.\csromanspacer{} นาหํ สมฺม สกฺโกมิ อคารสฺมา อนคาริยํ ปพฺพชิตุนฺติ.\csromanspacer{} ยํ เต สกฺกา อญฺญํ \csromanfolio{338}มยา กาตุํ, กฺยาหํ\footnote{ตฺยาหํ (สี, สฺยา)} กริสฺสามิ, ตฺวํ ปพฺพชาหีติ.\csromanspacer{} มาตา โข มํ สมฺม เอวมาห “สเจ ตาต อนุรุทฺธ ภทฺทิโย สกฺยราชา อคารสฺมา อนคาริยํ ปพฺพชติ, เอวํ ตฺวมฺปิ ปพฺพชาหี”ติ.\csromanspacer{} ภาสิตา โข ปน เต สมฺม เอสา วาจา “สเจ เต สมฺม ปพฺพชฺชา มม ปฏิพทฺธา วา อปฺปฏิพทฺธา วา สา โหตุ, อหํ ตยา, ยถา สุขํ ปพฺพชาหี”ติ.\csromanspacer{} เอหิ สมฺม, อุโภ อคารสฺมา อนคาริยํ ปพฺพชิสฺสามาติ.}

\prose{เตน โข ปน สมเยน มนุสฺสา สจฺจวาทิโน โหนฺติ สจฺจปฏิญฺญา.\csromanspacer{} อถ โข ภทฺทิโย สกฺยราชา อนุรุทฺธํ สกฺกํ เอตทโวจ “อาคเมหิ สมฺม สตฺตวสฺสานิ, สตฺตนฺนํ วสฺสานํ อจฺจเยน อุโภ อคารสฺมา อนคาริยํ ปพฺพชิสฺสามา”ติ.\csromanspacer{} อติจิรํ สมฺม สตฺตวสฺสานิ, นาหํ สกฺโกมิ สตฺตวสฺสานิ อาคเมตุนฺติ.\csromanspacer{} อาคเมหิ สมฺม ฉวสฺสานิ ฯเปฯ ปญฺจวสฺสานิ.\csromanspacer{} จตฺตาริ วสฺสานิ.\csromanspacer{} ตีณิ วสฺสานิ.\csromanspacer{} เทฺว วสฺสานิ.\csromanspacer{} เอกํ วสฺสํ, เอกสฺส วสฺสสฺส อจฺจเยน อุโภ อคารสฺมา อนคาริยํ ปพฺพชิสฺสามาติ.\csromanspacer{} อติจิรํ สมฺม เอกวสฺสํ, นาหํ สกฺโกมิ เอกํ วสฺสํ อาคเมตุนฺติ.\csromanspacer{} อาคเมหิ สมฺม สตฺตมาเส, สตฺตนฺนํ มาสานํ อจฺจเยน อุโภปิ อคารสฺมา อนคาริยํ ปพฺพชิสฺสามาติ.\csromanspacer{} อติจิรํ สมฺม สตฺตมาสา, นาหํ สกฺโกมิ สตฺตมาเส อาคเมตุนฺติ.\csromanspacer{} อาคเมหิ สมฺม ฉ มาเส ฯเปฯ ปญฺจ มาเส.\csromanspacer{} จตฺตาโร มาเส.\csromanspacer{} ตโย มาเส.\csromanspacer{} เทฺว มาเส.\csromanspacer{} เอกํ มาสํ.\csromanspacer{} อฑฺฒมาสํ, อฑฺฒมาสสฺส อจฺจเยน อุโภปิ อคารสฺมา อนคาริยํ ปพฺพชิสฺสามาติ.\csromanspacer{} อติจิรํ สมฺม อฑฺฒมาโส, นาหํ สกฺโกมิ อฑฺฒมาสํ อาคเมตุนฺติ.\csromanspacer{} อาคเมหิ สมฺม สตฺตาหํ, ยาวาหํ ปุตฺเต จ ภาตโร จ รชฺชํ นิยฺยาเทมีติ.\csromanspacer{} น จิรํ สมฺม สตฺตาโห อาคเมสฺสามีติ.}

\proseitem{331}{อถ โข ภทฺทิโย จ สกฺยราชา อนุรุทฺโธ จ อานนฺโท จ ภคุ จ กิมิโล จ เทวทตฺโต จ อุปาลิกปฺปเกน สตฺตมา ยถา ปุเร จตุรงฺคินิยา เสนาย อุยฺยานภูมิํ นิยฺยนฺติ, เอวเมว จตุรงฺคินิยา เสนาย นิยฺยิํสุ.\csromanspacer{} เต ทูรํ คนฺตฺวา เสนํ นิวตฺตาเปตฺวา ปรวิสยํ โอกฺกมิตฺวา อาภรณํ โอมุญฺจิตฺวา อุตฺตราสงฺเคน ภณฺฑิกํ พนฺธิตฺวา อุปาลิํ กปฺปกํ เอตทโวจุํ “หนฺท ภเณ อุปาลิ นิวตฺตสฺสุ, อลํ เต เอตฺตกํ ชีวิกายา”ติ.\csromanspacer{} อถ โข อุปาลิสฺส กปฺปกสฺส นิวตฺตนฺตสฺส เอตทโหสิ “จณฺฑา โข สากิยา, ‘อิมินา กุมารา นิปฺปาติตา’ติ ฆาตาเปยฺยุมฺปิ มํ. \csromanfolio{339}อิเม หิ นาม สกฺยกุมารา อคารสฺมา อนคาริยํ ปพฺพชิสฺสนฺติ, กิมงฺคํ\footnote{กิมงฺค (สี)} ปนาหนฺ’ติ ภณฺฑิกํ มุญฺจิตฺวา ตํ ภณฺฑํ รุกฺเข อาลคฺเคตฺวา “โย ปสฺสติ, ทินฺนํเยว หรตู”ติ วตฺวา เยน เต สกฺยกุมารา เตนุปสงฺกมิ, อทฺทสาสุํ โข เต สกฺยกุมารา อุปาลิํ กปฺปกํ ทูรโตว อาคจฺฉนฺตํ, ทิสฺวาน อุปาลิํ กปฺปกํ เอตทโวจุํ “กิสฺส ภเณ อุปาลิ นิวตฺเตสี”ติ.\csromanspacer{} อิธ เม อยฺยปุตฺตา นิวตฺตนฺตสฺส เอตทโหสิ “จณฺฑา โข สากิยา, ‘อิมินา กุมารา นิปฺปาติตา’ติ ฆาตาเปยฺยุมฺปิ มํ.\csromanspacer{} อิเม หิ นาม สกฺยกุมารา อคารสฺมา อนคาริยํ ปพฺพชิสฺสนฺติ, กิมงฺคํ ปนาหนฺ”ติ.\csromanspacer{} โส โข อหํ อยฺยปุตฺตา ภณฺฑิกํ มุญฺจิตฺวา ตํ ภณฺฑํ รุกฺเข อาลคฺเคตฺวา “โย ปสฺสติ, ทินฺนํเยว หรตู”ติ วตฺวา ตโตมฺหิ ปฏินิวตฺโตติ.\csromanspacer{} สุฏฺฐุ ภเณ อุปาลิ อกาสิ ยมฺปิ น นิวตฺโต\footnote{ยํ นิวตฺโต (สี), ยํ ปน นิวตฺโต (สฺยา)}.\csromanspacer{} จณฺฑา โข สากิยา “อิมินา กุมารา นิปฺปาติตา”ติ ฆาตาเปยฺยุมฺปิ ตนฺติ.}

\prose{อถ โข สกฺยกุมารา อุปาลิํ กปฺปกํ อาทาย เยน ภควา เตนุปสงฺกมิํสุ, อุปสงฺกมิตฺวา ภควนฺตํ อภิวาเทตฺวา เอกมนฺตํ นิสีทิํสุ.\csromanspacer{} เอกมนฺตํ นิสินฺนา โข เต สกฺยกุมารา ภควนฺตํ เอตทโวจุํ “มยํ ภนฺเต สากิยา นาม มานสฺสิโน, อยํ ภนฺเต อุปาลิ กปฺปโก อมฺหากํ ทีฆรตฺตํ ปริจารโก, อิมํ ภควา ปฐมํ ปพฺพาเชตุ, อิมสฺส มยํ อภิวาทนปจฺจุฏฺฐาน-อญฺชลิกมฺมสามีจิกมฺมํ กริสฺสาม, เอวํ อมฺหากํ สากิยานํ สากิยมาโน นิมฺมานายิสฺสตี”ติ\footnote{นิมฺมาทยิสฺสตีติ (สี), นิมฺมนิยิสฺสตีติ (สฺยา)}.}

\prose{อถ โข ภควา อุปาลิํ กปฺปกํ ปฐมํ ปพฺพาเชสิ, ปจฺฉา เต สกฺยกุมาเร.\csromanspacer{} อถ โข อายสฺมา ภทฺทิโย เตเนว อนฺตรวสฺเสน ติสฺโส วิชฺชา สจฺฉากาสิ, อายสฺมา อนุรุทฺโธ ทิพฺพจกฺขุํ อุปฺปาเทสิ, อายสฺมา อานนฺโท โสตาปตฺติผลํ สจฺฉากาสิ, เทวทตฺโต โปถุชฺชนิกํ อภินิปฺผาเทสิ.}

\proseitem{332}{เตน โข ปน สมเยน อายสฺมา ภทฺทิโย อรญฺญคโตปิ รุกฺขมูลคโตปิ สุญฺญาคารคโตปิ อภิกฺขณํ อุทานํ อุทาเนสิ “อโห สุขํ, อโห สุขนฺ”ติ.\csromanspacer{} อถ โข สมฺพหุลา ภิกฺขู เยน ภควา เตนุปสงฺกมิํสุ, อุปสงฺกมิตฺวา ภควนฺตํ อภิวาเทตฺวา เอกมนฺตํ นิสีทิํสุ. \csromanfolio{340}เอกมนฺตํ นิสินฺนา โข เต ภิกฺขู ภควนฺตํ เอตทโวจุํ “อายสฺมา ภนฺเต ภทฺทิโย อรญฺญคโตปิ รุกฺขมูลคโตปิ สุญฺญาคารคโตปิ อภิกฺขณํ อุทานํ อุทาเนสิ ‘อโห สุขํ, อโห สุขนฺ’ติ, นิสฺสํสยํ โข ภนฺเต อายสฺมา ภทฺทิโย อนภิรโตว พฺรหฺมจริยํ จรติ, ตํเยว วา ปุริมํ รชฺชสุขํ สมนุสฺสรนฺโต อรญฺญคโตปิ รุกฺขมูลคโตปิ สุญฺญาคารคโตปิ อภิกฺขณํ อุทานํ อุทาเนสิ ‘อโห สุขํ, อโห สุขนฺ’ติ”.}

\prose{อถ โข ภควา อญฺญตรํ ภิกฺขุํ อามนฺเตสิ\csromandash{}เอหิ ตฺวํ ภิกฺขุ มม วจเนน ภทฺทิยํ ภิกฺขุํ อามนฺเตหิ “สตฺถา ตํ อาวุโส ภทฺทิย อามนฺเตตี”ติ. “เอวํ ภนฺเต”ติ โข โส ภิกฺขุ ภควโต ปฏิสฺสุตฺวา เยนายสฺมา ภทฺทิโย เตนุปสงฺกมิ, อุปสงฺกมิตฺวา อายสฺมนฺตํ ภทฺทิยํ เอตทโวจ “สตฺถา ตํ อาวุโส ภทฺทิย อามนฺเตตี”ติ. “เอวมาวุโส”ติ โข อายสฺมา ภทฺทิโย ตสฺส ภิกฺขุโน ปฏิสฺสุตฺวา เยน ภควา เตนุปสงฺกมิ, อุปสงฺกมิตฺวา ภควนฺตํ อภิวาเทตฺวา เอกมนฺตํ นิสีทิ.\csromanspacer{} เอกมนฺตํ นิสินฺนํ โข อายสฺมนฺตํ ภทฺทิยํ ภควา เอตทโวจ “สจฺจํ กิร ตฺวํ ภทฺทิย อรญฺญคโตปิ รุกฺขมูลคโตปิ สุญฺญาคารคโตปิ อภิกฺขณํ อุทานํ อุทาเนสิ “อโห สุขํ อโห สุขนฺ”ติ.\csromanspacer{} เอวํ ภนฺเตติ.\csromanspacer{} กิํ ปน ตฺวํ ภทฺทิย อตฺถวสํ สมฺปสฺสมาโน อรญฺญคโตปิ รุกฺขมูลคโตปิ สุญฺญาคารคโตปิ อภิกฺขณํ อุทานํ อุทาเนสิ “อโห สุขํ อโห สุขนฺ”ติ.\csromanspacer{} ปุพฺเพ เม ภนฺเต รญฺโญ สโตปิ อนฺโตปิ อนฺเตปุเร รกฺขา สุสํวิหิตา โหติ, พหิปิ อนฺเตปุเร รกฺขา สุสํวิหิตา โหติ, อนฺโตปิ นคเร รกฺขา สุสํวิหิตา โหติ, พหิปิ นคเร รกฺขา สุสํวิหิตา โหติ, อนฺโตปิ ชนปเท รกฺขา สุสํวิหิตา โหติ, พหิปิ ชนปเท รกฺขา สุสํวิหิตา โหติ.\csromanspacer{} โส โข อหํ ภนฺเต เอวํ รกฺขิโตปิ โคปิโตปิ สนฺโต ภีโต อุพฺพิคฺโค อุสฺสงฺกี อุตฺรสฺโต วิหรามิ.\csromanspacer{} เอตรหิ โข ปน อหํ เอโก ภนฺเต อรญฺญคโตปิ รุกฺขมูลคโตปิ สุญฺญาคารคโตปิ อภีโต อนุพฺพิคฺโค อนุสฺสงฺกี อนุตฺรสฺโต อปฺโปสฺสุกฺโก ปนฺนโลโม ปรทตฺตวุตฺโต มิคภูเตน เจตสา วิหรามีติ อิมํ โข อหํ ภนฺเต อตฺถวสํ สมฺปสฺสมาโน อรญฺญคโตปิ รุกฺขมูลคโตปิ สุญฺญาคารคโตปิ อภิกฺขณํ อุทานํ อุทาเนมิ “อโห สุขํ อโห สุขนฺ”ติ.\csromanspacer{} อถ โข ภควา เอตมตฺถํ วิทิตฺวา ตายํ เวลายํ อิมํ อุทานํ อุทาเนสิ\csromandash{}}

\csromangathameasure{*“ยสฺสนฺตรโต น สนฺติ โกปา, \\ อิติ ภวาภวตญฺจ วีติวตฺโต. \\ ตํ วิคตภยํ สุขิํ อโสกํ, \\ เทวา นานุภวนฺติ ทสฺสนายา”ติ.}
\csromangathagroup{\csromangathastanza{\csromanfolio{341}\csromansymbolfootnote{*}{ขุ 1. 100 ปิฏฺเฐ อุทาเนปิ.}“ยสฺสนฺตรโต น สนฺติ โกปา, \\ อิติ ภวาภวตญฺจ วีติวตฺโต. \\ ตํ วิคตภยํ สุขิํ อโสกํ, \\ เทวา นานุภวนฺติ ทสฺสนายา”ติ.}}

\csromanmatikaanchor{mtk.163}
\csromantocmarkref{subsection}{เทวทตฺตวตฺถุ}{mtk.163}
\bookmark[dest={mtk.163},level=2]{เทวทตฺตวตฺถุ}
\csromanheader{2}{\textbf{เทวทตฺตวตฺถุ}}
\proseitem{33}{อถ โข ภควา อนุปิยายํ ยถาภิรนฺตํ วิหริตฺวา เยน โกสมฺพี เตน จาริกํ ปกฺกามิ, อนุปุพฺเพน จาริกํ จรมาโน เยน โกสมฺพี ตทวสริ, ตตฺร สุทํ ภควา โกสมฺพิยํ วิหรติ โฆสิตาราเม.\csromanspacer{} อถ โข เทวทตฺตสฺส รโหคตสฺส ปฏิสลฺลีนสฺส เอวํ เจตโส ปริวิตกฺโก อุทปาทิ “กํ นุ โข อหํ ปสาเทยฺยํ, ยสฺมิํ เม ปสนฺเน พหุลาภสกฺกาโร อุปฺปชฺเชยฺยา”ติ.\csromanspacer{} อถ โข เทวทตฺตสฺส เอตทโหสิ “อยํ โข อชาตสตฺตุ กุมาโร ตรุโณ เจว อายติํ ภทฺโท จ, ยํนุนาหํ อชาตสตฺตุํ กุมารํ ปสาเทยฺยํ, ตสฺมิํ เม ปสนฺเน พหุลาภสกฺกาโร อุปฺปชฺชิสฺสตี”ติ.}

\prose{อถ โข เทวทตฺโต เสนาสนํ สํสาเมตฺวา ปตฺตจีวรมาทาย เยน ราชคหํ เตน ปกฺกามิ, อนุปุพฺเพน เยน ราชคหํ ตทวสริ.\csromanspacer{} อถ โข เทวทตฺโต สกวณฺณํ ปฏิสํหริตฺวา กุมารกวณฺณํ อภินิมฺมินิตฺวา อหิเมขลิกาย อชาตสตฺตุสฺส กุมารสฺส อุจฺฉงฺเค\footnote{อุจฺจงฺเก (สฺยา)} ปาตุรโหสิ.\csromanspacer{} อถ โข อชาตสตฺตุ กุมาโร ภีโต อโหสิ อุพฺพิคฺโค อุสฺสงฺกี อุตฺรสฺโต.\csromanspacer{} อถ โข เทวทตฺโต อชาตสตฺตุํ กุมารํ เอตทโวจ “ภายสิ มํ ตฺวํ กุมารา”ติ.\csromanspacer{} อาม ภายามิ, โกสิ ตฺวนฺติ.\csromanspacer{} อหํ เทวทตฺโตติ.\csromanspacer{} สเจ โข ตฺวํ ภนฺเต อยฺโย เทวทตฺโต อิงฺฆ สเกเนว วณฺเณน ปาตุภวสฺสูติ.\csromanspacer{} อถ โข เทวทตฺโต กุมารกวณฺณํ ปฏิสํหริตฺวา สํฆาฏิปตฺตจีวรธโร อชาตสตฺตุสฺส กุมารสฺส ปุรโต อฏฺฐาสิ.\csromanspacer{} อถ โข อชาตสตฺตุ กุมาโร เทวทตฺตสฺส อิมินา อิทฺธิปาฏิหาริเยน อภิปฺปสนฺโน ปญฺจหิ รถสเตหิ สายํ ปาตํ อุปฏฺฐานํ คจฺฉติ, ปญฺจ จ ถาลิปากสตานิ ภตฺตาภิหาโร อภิหรียติ.\csromanspacer{} อถ โข เทวทตฺตสฺส ลาภสกฺการสิโลเกน อภิภูตสฺส ปริยาทินฺนจิตฺตสฺส เอวรูปํ \csromanfolio{342}อิจฺฉาคตํ อุปฺปชฺชิ “อหํ ภิกฺขุสํฆํ ปริหริสฺสามี”ติ.\csromanspacer{} สห จิตฺตุปฺปาทาว เทวทตฺโต ตสฺสา อิทฺธิยา ปริหายิ.}

\prose{\csromansymbolfootnote{*}{อํ 2. 107 ปิฏฺเฐปิ.}เตน โข ปน สมเยน กกุโธ นาม โกฬิยปุตฺโต อายสฺมโต มหาโมคฺคลฺลานสฺส อุปฏฺฐาโก อธุนา กาลงฺกโต อญฺญตรํ มโนมยํ กายํ อุปปนฺโน.\csromanspacer{} ตสฺส เอวรูโป อตฺตภาวปฺปฏิลาโภ โหติ เสยฺยถาปิ นาม เทฺว วา ตีณิ วา มาคธกานิ\footnote{มาคธิกานิ (สฺยา)} คามกฺเขตฺตานิ.\csromanspacer{} โส เตน อตฺตภาวปฺปฏิลาเภน เนว อตฺตานํ น ปรํ พฺยาพาเธติ.\csromanspacer{} อถ โข กกุโธ เทวปุตฺโต เยนายสฺมา มหาโมคฺคลฺลาโน เตนุปสงฺกมิ, อุปสงฺกมิตฺวา อายสฺมนฺตํ มหาโมคฺคลฺลานํ อภิวาเทตฺวา เอกมนฺตํ อฏฺฐาสิ.\csromanspacer{} เอกมนฺตํ ฐิโต โข กกุโธ เทวปุตฺโต อายสฺมนฺตํ มหาโมคฺคลฺลานํ เอตทโวจ “เทวทตฺตสฺส ภนฺเต ลาภสกฺการสิโลเกน อภิภูตสฺส ปริยาทินฺนจิตฺตสฺส\footnote{ปริยาทิณฺณจิตฺตสฺส (ก)} เอวรูปํ อิจฺฉาคตํ อุปฺปชฺชิ ‘อหํ ภิกฺขุสํฆํ ปริหริสฺสามี”ติ.\csromanspacer{} สห จิตฺตุปฺปาทาว ภนฺเต เทวทตฺโต ตสฺสา อิทฺธิยา ปริหีโน”ติ.\csromanspacer{} อิทมโวจ กกุโธ เทวปุตฺโต, อิทํ วตฺวา อายสฺมนฺตํ มหาโมคฺคลฺลานํ อภิวาเทตฺวา ปทกฺขิณํ กตฺวา ตตฺเถว อนฺตรธายิ.}

\prose{อถ โข อายสฺมา มหาโมคฺคลฺลาโน เยน ภควา เตนุปสงฺกมิ, อุปสงฺกมิตฺวา ภควนฺตํ อภิวาเทตฺวา เอกมนฺตํ นิสีทิ.\csromanspacer{} เอกมนฺตํ นิสินฺโน โข อายสฺมา มหาโมคฺคลฺลาโน ภควนฺตํ เอตทโวจ “กกุโธ นาม ภนฺเต โกฬิยปุตฺโต มม อุปฏฺฐาโก อธุนา กาลงฺกโต อญฺญตรํ มโนมยํ กายํ อุปปนฺโน, ตสฺส เอวรูโป อตฺตภาวปฺปฏิลาโภ เสยฺยถาปิ นาม เทฺว วา ตีณิ วา มาคธกานิ คามกฺเขตฺตานิ.\csromanspacer{} โส เตน อตฺตภาวปฺปฏิลาเภน เนว อตฺตานํ น ปรํ พฺยาพาเธติ.\csromanspacer{} อถ โข ภนฺเต กกุโธ เทวปุตฺโต เยนาหํ เตนุปสงฺกมิ, อุปสงฺกมิตฺวา มํ อภิวาเทตฺวา เอกมนฺตํ อฏฺฐาสิ.\csromanspacer{} เอกมนฺตํ ฐิโต โข ภนฺเต กกุโธ เทวปุตฺโต มํ เอตทโวจ “เทวทตฺตสฺส ภนฺเต ลาภสกฺการสิโลเกน อภิภูตสฺส ปริยาทินฺนจิตฺตสฺส เอวรูปํ อิจฺฉาคตํ อุปฺปชฺชิ ‘อหํ ภิกฺขุสํฆํ ปริหริสฺสามี’ติ, สห จิตฺตุปฺปาทาว ภนฺเต เทวทตฺโต ตสฺสา อิทฺธิยา ปริหีโนติ.\csromanspacer{} อิทมโวจ ภนฺเต กกุโธ เทวปุตฺโต, อิทํ วตฺวา มํ อภิวาเทตฺวา ปทกฺขิณํ กตฺวา ตตฺเถว อนฺตรธายี”ติ.}

\prose{\csromanfolio{343}กิํ ปน เต โมคฺคลฺลาน กกุโธ เทวปุตฺโต เจตสา เจโต ปริจฺจ วิทิโต, ยํ กิญฺจิ กกุโธ เทวปุตฺโต ภาสติ, สพฺพํ ตํ ตเถว โหติ, โน อญฺญถาติ.\csromanspacer{} เจตสา เจโต ปริจฺจ วิทิโต จ เม ภนฺเต กกุโธ เทวปุตฺโต, ยํ กิญฺจิ กกุโธ เทวปุตฺโต ภาสติ, สพฺพํ ตํ ตเถว โหติ, โน อญฺญถาติ.\csromanspacer{} รกฺขสฺเสตํ โมคฺคลฺลาน วาจํ, รกฺขสฺเสตํ โมคฺคลฺลาน วาจํ, อิทานิ โส โมฆปุริโส อตฺตนาว อตฺตานํ ปาตุกริสฺสติ.}

\csromanmatikaanchor{mtk.164}
\csromantocmarkref{subsection}{ปญฺจสตฺถุกถา}{mtk.164}
\bookmark[dest={mtk.164},level=2]{ปญฺจสตฺถุกถา}
\csromanheader{2}{\textbf{ปญฺจสตฺถุกถา}}
\proseitem{334}{\csromansymbolfootnote{*}{อํ 2. 108 ปิฏฺเฐปิ.}ปญฺจิเม โมคฺคลฺลาน สตฺถาโร สนฺโต สํวิชฺชมานา โลกสฺมิํ.\csromanspacer{} กตเม ปญฺจ.\csromanspacer{} อิธ โมคฺคลฺลาน เอกจฺโจ สตฺถา อปริสุทฺธสีโล สมาโน “ปริสุทฺธสีโลมฺหี”ติ ปฏิชานาติ “ปริสุทฺธํ เม สีลํ ปริโยทาตํ อสํกิลิฏฺฐนฺ”ติ จ.\csromanspacer{} ตเมนํ สาวกา เอวํ ชานนฺติ “อยํ โข ภวํ สตฺถา อปริสุทฺธสีโล สมาโน ‘ปริสุทฺธสีโลมฺหี’ติ ปฏิชานาติ ‘ปริสุทฺธํ เม สีลํ ปริโยทาตํ อสํกิลิฏฺฐนฺ’ติ จ.\csromanspacer{} มยญฺเจว โข ปน คิหีนํ อาโรเจยฺยาม ‘นาสฺสสฺส มนาปํ, ยํ โข ปนสฺส อมนาปํ กถํ นํ มยํ เตน สมุทาจเรยฺยาม.\csromanspacer{} สมฺมนฺนติ โข ปน จีวรปิณฺฑปาตเสนาสนคิลานปฺปจฺจยเภสชฺชปริกฺขาเรน.\csromanspacer{} ยํ ตุโม กริสฺสติ, ตุโมว เตน ปญฺญายิสฺสตี”ติ.\csromanspacer{} เอวรูปํ โข โมคฺคลฺลาน สตฺถารํ สาวกา สีลโต รกฺขนฺติ, เอวรูโป จ ปน สตฺถา สาวเกหิ สีลโต รกฺขํ ปจฺจาสีสติ\footnote{ปจฺจาสิํสติ (สี, สฺยา)}.}

\prose{ปุน จปรํ โมคฺคลฺลาน อิเธกจฺโจ สตฺถา อปริสุทฺธาชีโว สมาโน “ปริสุทฺธาชีโวมฺหี”ติ ปฏิชานาติ “ปริสุทฺโธ เม อาชีโว ปริโยทาโต อสํกิลิฏฺโฐ”ติ จ.\csromanspacer{} ตเมนํ สาวกา เอวํ ชานนฺติ “อยํ โข ภวํ สตฺถา อปริสุทฺธาชีโว สมาโน ‘ปริสุทฺธาชีโวมฺหี’ติ ปฏิชานาติ ‘ปริสุทฺโธ เม อาชีโว ปริโยทาโต อสํกิลิฏฺโฐ’ติ จ.\csromanspacer{} มยญฺเจว โข ปน คิหีนํ อาโรเจยฺยาม ‘นาสฺสสฺส มนาปํ, ยํ โข ปนสฺส อมนาปํ, กถํ นํ มยํ เตน สมุทาจเรยฺยาม.\csromanspacer{} สมฺมนฺนติ โข ปน จีวรปิณฺฑปาตเสนาสนคิลานปฺปจฺจยเภสชฺชปริกฺขาเรน.\csromanspacer{} ยํ \csromanfolio{344}ตุโม กริสฺสติ, ตุโมว เตน ปญฺญายิสฺสตี”ติ.\csromanspacer{} เอวรูปํ โข โมคฺคลฺลาน สตฺถารํ สาวกา อาชีวโต รกฺขนฺติ, เอวรูโป จ ปน สตฺถา สาวเกหิ อาชีวโต รกฺขํ ปจฺจาสีสติ.}

\prose{ปุน จปรํ โมคฺคลฺลาน อิเธกจฺโจ สตฺถา อปริสุทฺธธมฺมเทสโน สมาโน “ปริสุทฺธธมฺมเทสโนมฺหี”ติ ปฏิชานาติ “ปริสุทฺธา เม ธมฺมเทสนา ปริโยทาตา อสํกิลิฏฺฐา”ติ จ.\csromanspacer{} ตเมนํ สาวกา เอวํ ชานนฺติ “อยํ โข ภวํ สตฺถา อปริสุทฺธธมฺมเทสโน สมาโน ‘ปริสุทฺธธมฺมเทสโนมฺหี’ติ ปฏิชานาติ ‘ปริสุทฺธา เม ธมฺมเทสนา ปริโยทาตา อสํกิลิฏฺฐา’ติ จ.\csromanspacer{} มยญฺเจว โข ปน คิหีนํ อาโรเจยฺยาม ‘นาสฺสสฺส มนาปํ, ยํ โข ปนสฺส อมนาปํ.\csromanspacer{} กถํ นํ มยํ เตน สมุทาจเรยฺยาม.\csromanspacer{} สมฺมนฺนติ โข ปน จีวรปิณฺฑปาตเสนาสนคิลานปฺปจฺจยเภสชฺชปริกฺขาเรน.\csromanspacer{} ยํ ตุโม กริสฺสติ, ตุโมว เตน ปญฺญายิสฺสาตี”ติ.\csromanspacer{} เอวรูปํ โข โมคฺคลฺลาน สตฺถารํ สาวกา ธมฺมเทสนโต รกฺขนฺติ, เอวรูโป จ ปน สตฺถา สาวเกหิ ธมฺมเทสนโต รกฺขํ ปจฺจาสีสติ.}

\prose{ปุน จปรํ โมคฺคลฺลาน อิเธกจฺโจ สตฺถา อปริสุทฺธเวยฺยากรโณ สมาโน “ปริสุทฺธเวยฺยากรโณมฺหี”ติ ปฏิชานาติ “ปริสุทฺธํ เม เวยฺยากรณํ ปริโยทาตํ อสํกิลิฏฺฐนฺ”ติ จ.\csromanspacer{} ตเมนํ สาวกา เอวํ ชานนฺติ “อยํ โข ภวํ สตฺถา อปริสุทฺธเวยฺยากรโณ สมาโน ‘ปริสุทฺธเวยฺยากรโณมฺหี’ติ ปฏิชานาติ ‘ปริสุทฺธํ เม เวยฺยากรณํ ปริโยทาตํ อสํกิลิฏฺฐนฺ’ติ จ.\csromanspacer{} มยญฺเจว โข ปน คิหีนํ อาโรเจยฺยาม ‘นาสฺสสฺส มนาปํ, ยํ โข ปนสฺส อมนาปํ, กถํ นํ มยํ เตน สมุทาจเรยฺยาม.\csromanspacer{} สมฺมนฺนติ โข ปน จีวรปิณฺฑปาตเสนาสนคิลานปฺปจฺจยเภสชฺชปริกฺขาเรน.\csromanspacer{} ยํ ตุโม กริสฺสติ, ตุโมว เตน ปญฺญายิสฺสตี”ติ.\csromanspacer{} เอวรูปํ โข โมคฺคลฺลาน สตฺถารํ สาวกา เวยฺยากรณโต รกฺขนฺติ, เอวรูโป จ ปน สตฺถา สาวเกหิ เวยฺยากรณโต รกฺขํ ปจฺจาสีสติ.}

\prose{ปุน จปรํ โมคฺคลฺลาน อิเธกจฺโจ สตฺถา อปริสุทฺธญาณทสฺสโน สมาโน “ปริสุทฺธญาณทสฺสโนมฺหี”ติ ปฏิชานาติ “ปริสุทฺธํ เม ญาณทสฺสนํ ปริโยทาตํ อสํกิลิฏฺฐนฺ”ติ จ.\csromanspacer{} ตเมนํ สาวกา เอวํ ชานนฺติ “อยํ โข ภวํ สตฺถา อปริสุทฺธญาณทสฺสโน สมาโน \csromanfolio{345}‘ปริสุทฺธญาณทสฺสโนมฺหี’ติ ปฏิชานาติ ‘ปริสุทฺธํ เม ญาณทสฺสนํ ปริโยทาตํ อสํกิลิฏฺฐนฺ’ติ จ.\csromanspacer{} มยญฺเจว โข ปน คิหีนํ อาโรเจยฺยาม ‘นาสฺสสฺส มนาปํ, ยํ โข ปนสฺส อมนาปํ, กถํ นํ มยํ เตน สมุทาจเรยฺยาม.\csromanspacer{} สมฺมนฺนติ โข ปน จีวรปิณฺฑปาตเสนาสนคิลานปฺปจฺจยเภสชฺชปริกฺขาเรน.\csromanspacer{} ยํ ตุโม กริสฺสติ, ตุโมว เตน ปญฺญายิสฺสตี”ติ.\csromanspacer{} เอวรูปํ โข โมคฺคลฺลาน สตฺถารํ สาวกา ญาณทสฺสนโต รกฺขนฺติ, เอวรูโป จ ปน สตฺถา สาวเกหิ ญาณทสฺสนโต รกฺขํ ปจฺจาสีสตีติ.\csromanspacer{} อิเม โข โมคฺคลฺลาน ปญฺจ สตฺถาโร สนฺโต สํวิชฺชมานา โลกสฺมิํ.}

\prose{อหํ โข ปน โมคฺคลฺลาน ปริสุทฺธสีโล สมาโน “ปริสุทฺธสีโลมฺหี”ติ ปฏิชานามิ “ปริสุทฺธํ เม สีลํ ปริโยทาตํ อสํกิลิฏฺฐนฺ”ติ จ, น จ มํ สาวกา สีลโต รกฺขนฺติ, น จาหํ สาวเกหิ สีลโต รกฺขํ ปจฺจาสีสามิ.\csromanspacer{} ปริสุทฺธาชีโว สมาโน ฯเปฯ.\csromanspacer{} ปริสุทฺธธมฺมเทสโน สมาโน ฯเปฯ.\csromanspacer{} ปริสุทฺธเวยฺยากรโณ สมาโน ฯเปฯ.\csromanspacer{} ปริสุทฺธญาณทสฺสโน สมาโน “ปริสุทฺธญาณทสฺสโนมฺหี”ติ ปฏิชานามิ “ปริสุทฺธํ เม ญาณทสฺสนํ ปริโยทาตํ อสํกิลิฏฺฐนฺ”ติ จ.\csromanspacer{} น จ มํ สาวกา ญาณทสฺสนโต รกฺขนฺติ, น จาหํ สาวเกหิ ญาณทสฺสนโต รกฺขํ ปจฺจาสีสามีติ.}

\proseitem{335}{อถ โข ภควา โกสมฺพิยํ ยถาภิรนฺตํ วิหริตฺวา เยน ราชคหํ เตน จาริกํ ปกฺกามิ, อนุปุพฺเพน จาริกํ จรมาโน เยน ราชคหํ ตทวสริ, ตตฺร สุทํ ภควา ราชคเห วิหรติ เวฬุวเน กลนฺทกนิวาเป.\csromanspacer{} อถ โข สมฺพหุลา ภิกฺขู เยน ภควา เตนุปสงฺกมิํสุ, อุปสงฺกมิตฺวา ภควนฺตํ อภิวาเทตฺวา เอกมนฺตํ นิสีทิํสุ.\csromanspacer{} เอกมนฺตํ นิสินฺนา โข เต ภิกฺขู ภควนฺตํ เอตทโวจุํ “เทวทตฺตสฺส ภนฺเต อชาตสตฺตุ กุมาโร ปญฺจหิ รถสเตหิ สายํ ปาตํ อุปฏฺฐานํ คจฺฉติ, ปญฺจ จ ถาลิปากสตานิ ภตฺตาภิหาโร อภิหรียตี”ติ.\csromanspacer{} มา ภิกฺขเว เทวทตฺตสฺส ลาภสกฺการสิโลกํ ปิหยิตฺถ, ยาวกีวญฺจ ภิกฺขเว เทวทตฺตสฺส อชาตสตฺตุ กุมาโร ปญฺจหิ รถสเตหิ สายํ ปาตํ อุปฏฺฐานํ คมิสฺสติ, ปญฺจ จ ถาลิปากสตานิ ภตฺตาภิหาโร อภิหรียิสฺสติ, หานิเยว ภิกฺขเว เทวทตฺตสฺส ปาฏิกงฺขา กุสเลสุ ธมฺเมสุ, โน วุฑฺฒิ.}

\prose{\csromanfolio{346}เสยฺยถาปิ ภิกฺขเว จณฺฑสฺส กุกฺกุรสฺส นาสาย ปิตฺตํ ภินฺเทยฺยุํ, เอวญฺหิ โส ภิกฺขเว กุกฺกุโร ภิยฺโยโสมตฺตาย จณฺฑตโร อสฺส, เอวเมว โข ภิกฺขเว ยาวกีวญฺจ เทวทตฺตสฺส อชาตสตฺตุ กุมาโร ปญฺจหิ รถสเตหิ สายํ ปาตํ อุปฏฺฐานํ คมิสฺสติ, ปญฺจ จ ถาลิปากสตานิ ภตฺตาภิหาโร อภิหรียิสฺสติ.\csromanspacer{} หานิเยว ภิกฺขเว เทวทตฺตสฺส ปาฏิกงฺขา กุสเลสุ ธมฺเมสุ, โน วุฑฺฒิ.}

\prose{\csromansymbolfootnote{*}{สํ 1. 436; อํ 1. 385 ปิฏฺเฐสุปิ.}อตฺตวธาย ภิกฺขเว เทวทตฺตส ลาภสกฺการสิโลโก อุทปาทิ, ปราภวาย เทวทตฺตสฺส ลาภสกฺการสิโลโก อุทปาทิ.}

\prose{เสยฺยถาปิ ภิกฺขเว กทลี อตฺตวธาย ผลํ เทติ, ปราภวาย ผลํ เทติ, เอวเมว โข ภิกฺขเว อตฺตวธาย เทวทตฺตสฺส ลาภสกฺการสิโลโก อุทปาทิ, ปราภวาย เทวทตฺตสฺส ลาภสกฺการสิโลโก อุทปาทิ.}

\prose{เสยฺยถาปิ ภิกฺขเว เวฬุ อตฺตวธาย ผลํ เทติ, ปราภวาย ผลํ เทติ, เอวเมว โข ภิกฺขเว อตฺตวธาย เทวทตฺตสฺส ลาภสกฺการสิโลโก อุทปาทิ, ปราภวาย เทวทตฺตสฺส ลาภสกฺการสิโลโก อุทปาทิ.}

\prose{เสยฺยถาปิ ภิกฺขเว นโฬ อตฺถวธาย ผลํ เทติ, ปราภวาย ผลํ เทติ, เอวเมว โข ภิกฺขเว อตฺตวธาย เทวทตฺตสฺส ลาภสกฺการสิโลโก อุทปาทิ, ปราภวาย เทวทตฺตสฺส ลาภสกฺการสิโลโก อุทปาทิ.}

\prose{เสยฺยถาปิ ภิกฺขเว อสฺสตรี อตฺตวธาย คพฺภํ คณฺหาติ, ปราภวาย คพฺภํ คณฺหาติ, เอวเมว โข ภิกฺขเว อตฺตวธาย เทวทตฺตสฺส ลาภสกฺการสิโลโก อุทปาทิ, ปราภวาย เทวทตฺตสฺส ลาภสกฺการสิโลโก อุทปาทีติ.}

\csromangathameasure{+ผลํ เว กทลิํ หนฺติ, ผลํ เวฬุํ ผลํ นฬํ. \\ สกฺกาโร กาปุริสํ หนฺติ, คพฺโภ อสฺสตริํ ยถาติ.}
\csromangathagroup{\csromangathastanza{\csromansymbolfootnote{+}{สํ 1. 156; ขุ 10. 109 ปิฏฺเฐสุปิ.}ผลํ เว กทลิํ หนฺติ, ผลํ เวฬุํ ผลํ นฬํ. \\ สกฺกาโร กาปุริสํ หนฺติ, คพฺโภ อสฺสตริํ ยถาติ.}}

\nitthitam{ปฐมภาณวาโร นิฏฺฐิโต.}
\csromanmatikaanchor{mtk.165}
\csromantocmarkref{section}{2. ทุติยภาณวาร}{mtk.165}
\bookmark[dest={mtk.165},level=1]{2. ทุติยภาณวาร}
\csromanheader{1}{2. ทุติยภาณวาร}
\csromanmatikaanchor{mtk.166}
\csromantocmarkref{subsection}{ปกาสนียกมฺม}{mtk.166}
\bookmark[dest={mtk.166},level=2]{ปกาสนียกมฺม}
\csromanheader{2}{\textbf{ปกาสนียกมฺม}}
\proseitem{336}{\csromanfolio{347}เตน โข ปน สมเยน ภควา มหติยา ปริสาย ปริวุโต ธมฺมํ เทเสนฺโต นิสินฺโน โหติ สราชิกาย ปริสาย.\csromanspacer{} อถ โข เทวทตฺโต อุฏฺฐายาสนา เอกํสํ อุตฺตราสงฺคํ กริตฺวา เยน ภควา เตนญฺชลิํ ปณาเมตฺวา ภควนฺตํ เอตทโวจ “ชิณฺโณ ทานิ ภนฺเต ภควา วุฑฺโฒ มหลฺลโก อทฺธคโต วโย-อนุปฺปตฺโต, อปฺโปสฺสุกฺโก ทานิ ภนฺเต ภควา ทิฏฺฐธมฺมสุขวิหารํ อนุยุตฺโต วิหรตุ, มมํ ภิกฺขุสํฆํ นิสฺสชฺชตุ, อหํ ภิกฺขุสํฆํ ปริหริสฺสามี”ติ.\csromanspacer{} อลํ เทวทตฺต, มา เต รุจฺจิ ภิกฺขุสํฆํ ปริหริตุนฺติ.\csromanspacer{} ทุติยมฺปิ โข เทวทตฺโต ฯเปฯ.\csromanspacer{} ตติยมฺปิ โข เทวทตฺโต ภควนฺตํ เอตทโวจ “ชิณฺโณ ทานิ ภนฺเต ภควา วุฑฺโฒ มหลฺลโก อทฺธคโต วโย-อนุปฺปตฺโต, อปฺโปสฺสุกฺโก ทานิ ภนฺเต ภควา ทิฏฺฐธมฺมสุขวิหารํ อนุยุตฺโต วิหรตุ, มมํ ภิกฺขุสํฆํ นิสฺสชฺชตุ, อหํ ภิกฺขุสํฆํ ปริหริสฺสามี”ติ.\csromanspacer{} สาริปุตฺตโมคฺคลฺลนานมฺปิ โข อหํ เทวทตฺต ภิกฺขุสํฆํ น นิสฺสชฺเชยฺยํ, กิํ ปน ตุยฺหํ ฉวสฺส เขฬาสกสฺสาติ.\csromanspacer{} อถ โข เทวทตฺโต “สราชิกาย มํ ภควา ปริสาย เขฬาสกวาเทน อปสาเทติ, สาริปุตฺตโมคฺคลฺลาเนว อุกฺกํสตี”ติ กุปิโต อนตฺตมโน ภควนฺตํ อภิวาเทตฺวา ปทกฺขิณํ กตฺวา ปกฺกามิ.\csromanspacer{} อยญฺจรหิ เทวทตฺตสฺส ภควติ ปฐโม อาฆาโต อโหสิ.}

\prose{อถ โข ภควา ภิกฺขู อามนฺเตสิ\csromandash{}เตน หิ ภิกฺขเว สํโฆ เทวทตฺตสฺส ราชคเห ปกาสนียํ กมฺมํ กโรตุ “ปุพฺเพ เทวทตฺตสฺส อญฺญา ปกติ อโหสิ, อิทานิ อญฺญา ปกติ, ยํ เทวทตฺโต กเรยฺย กาเยน วาจาย, น เตน พุทฺโธ วา ธมฺโม วา สํโฆ วา ทฏฺฐพฺโพ, เทวทตฺโตว เตน ทฏฺฐพฺโพ”ติ.\csromanspacer{} เอวญฺจ ปน ภิกฺขเว กาตพฺพํ, พฺยตฺเตน ภิกฺขุนา ปฏิพเลน สํโฆ ญาเปตพฺโพ\csromandash{}}

\proseitem{337}{“สุณาตุ เม ภนฺเต สํโฆ, ยทิ สํฆสฺส ปตฺตกลฺลํ, สํโฆ เทวทตฺตสฺส ราชคเห ปกาสนียํ กมฺมํ กเรยฺย ‘ปุพฺเพ เทวทตฺถสฺส อญฺญา ปกติ อโหสิ, อิทานิ อญฺญา ปกติ, ยํ เทวทตฺโต กเรยฺย \csromanfolio{348}กาเยน วาจายาน เตน พุทฺโธ วา ธมฺโม วา สํโฆ วา ทฏฺฐพฺโพ, เทวทตฺโตว เตน ทฏฺฐพฺโพ’ติ, เอสา ญตฺติ.}

\prose{สุณาตุ เม ภนฺเต สํโฆ, สํโฆ เทวทตฺตสฺส ราชคเห ปกาสนียํ กมฺมํ กโรติ ‘ปุพฺเพ เทวทตฺตสฺส อญฺญา ปกติ อโหสิ, อิทานิ อญฺญา ปกติ, ยํ เทวทตฺโต กเรยฺย กาเยน วาจาย, น เตน พุทฺโธ วา ธมฺโม วา สํโฆ วา ทฏฺฐพฺโพ, เทวทตฺโตว เตน ทฏฺฐพฺโพ’ติ, ยสฺสายสฺมโต ขมติ เทวทตฺตสฺส ราชคเห ปกาสนียสฺส กมฺมสฺส กรณํ ‘ปุพฺเพ เทวทตฺตสฺส อญฺญา ปกติ อโหสิ, อิทานิ อญฺญา ปกติ, ยํ เทวทตฺโต กเรยฺย กาเยน วาจาย, น เตน พุทฺโธ วา ธมฺโม วา สํโฆ วา ทฏฺฐพฺโพ, เทวทตฺโตว เตน ทฏฺฐพฺโพ’ติ, โส ตุณฺหสฺส, ยสฺส นกฺขมติ, โส ภาเสยฺย.}

\prose{กถํ สํเฆน เทวทตฺตสฺส ราชคเห ปกาสนียํ กมฺมํ ‘ปุพฺเพ เทวทตฺตสฺส อญฺญา ปกติ อโหสิ, อิทานิ อญฺญา ปกติ, ยํ เทวทตฺโต กเรยฺย กาเยน วาจาย, น เตน พุทฺโธ วา ธมฺโม วา สํโฆ วา ทฏฺฐพฺโพ, เทวทตฺโตว เตน ทฏฺฐพฺโพ’ติ, ขมติ สํฆสฺส, ตสฺมา ตุณฺหี, เอวเมตํ ธารยามี”ติ.}

\proseitem{338}{อถ โข ภควา อายสฺมนฺตํ สาริปุตฺตํ อามนฺเตสิ “เตน หิ ตฺวํ สาริปุตฺต เทวทตฺตํ ราชคเห ปกาเสหี”ติ.\csromanspacer{} ปุพฺเพ มยา ภนฺเต เทวทตฺตสฺส ราชคเห วณฺโณ ภาสิโต “มหิทฺธิโก โคธิปุตฺโต, มหานุภาโว โคธิปุตฺโต”ติ.\csromanspacer{} กถาหํ ภนฺเต เทวทตฺตํ ราชคเห ปกาเสมีติ.\csromanspacer{} นนุ ตยา สาริปุตฺต ภูโตเยว เทวทตฺตสฺส ราชคเห วณฺโณ ภาสิโต “มหิทฺธิโก โคธิปุตฺโต, มหานุภาโว โคธิปุตฺโต”ติ.\csromanspacer{} เอวํ ภนฺเตติ.\csromanspacer{} เอวเมว โข ตฺวํ สาริปุตฺต ภูตํเยว เทวทตฺตํ ราชคเห ปกาเสหีติ. “เอวํ ภนฺเต”ติ โข อายสฺมา สาริปุตฺโต ภควโต ปจฺจสฺโสสิ.}

\prose{อถ โข ภควา ภิกฺขู อามนฺเตสิ “เตน หิ ภิกฺขเว สํโฆ สาริปุตฺตํ สมฺมนฺนตุ เทวทตฺตํ ราชคเห ปกาเสตุํ ‘ปุพฺเพ เทวทตฺตสฺส อญฺญา ปกติ อโหสิ, อิทานิ อญฺญา ปกติ, ยํ เทวทตฺโต กเรยฺย กาเยน วาจาย, น เตน พุทฺโธ วา ธมฺโม วา สํโฆ วา ทฏฺฐพฺโพ, เทวทตฺโตว เตน ทฏฺฐพฺโพ’ติ”.\csromanspacer{} เอวญฺจ ปน ภิกฺขเว สมฺมนฺนิตพฺโพ, ปฐมํ \csromanfolio{349}สาริปุตฺโต ยาจิตพฺโพ, ยาจิตฺวา พฺยตฺเตน ภิกฺขุนา ปฏิพเลน สํโฆ ญาเปตพฺโพ\csromandash{}}

\prose{“สุณาตุ เม ภนฺเต สํโฆ, ยทิ สํฆสฺส ปตฺตกลฺลํ, สํโฆ อายสฺมนฺตํ สาริปุตฺตํ สมฺมนฺเนยฺย เทวทตฺตํ ราชคเห ปกาเสตุํ ‘ปุพฺเพ เทวทตฺตสฺส อญฺญา ปกติ อโหสิ, อิทานิ อญฺญา ปกติ, ยํ เทวทตฺโต กเรยฺย กาเยน วาจาย, น เตน พุทฺโธ วา ธมฺโม วา สํโฆ วา ทฏฺฐพฺโพ, เทวทตฺโตว เตน ทฏฺฐพฺโพ’ติ, เอสา ญตฺติ.}

\prose{สุณาตุ เม ภนฺเต สํโฆ, สํโฆ อายสฺมนฺตํ สาริปุตฺตํ สมฺมนฺนติ เทวทตฺตํ ราชคเห ปกาเสตุํ ‘ปุพฺเพ เทวทตฺตส อญฺญา ปกติ อโหสิ, อิทานิ อญฺญา ปกติ, ยํ เทวทตฺโต กเรยฺย กาเยน วาจาย, น เตน พุทฺโธ วา ธมฺโม วา สํโฆ วา ทฏฺฐพฺโพ, เทวทตฺโตว เตน ทฏฺฐพฺโพ’ติ, ยสฺสายสฺมโต ขมติ อายสฺมโต สาริปุตฺตสฺส สมฺมุติ เทวทตฺตํ ราชคเห ปกาเสตุํ ‘ปุพฺเพ เทวทตฺตสฺส อญฺญา ปกติ อโหสิ, อิทานิ อญฺญา ปกติ, ยํ เทวทตฺโต กเรยฺย กาเยน วาจาย, น เตน พุทฺโธ วา ธมฺโม วา สํโฆ วา ทฏฺฐพฺโพ, เทวทตฺโตว เตน ทฏฺฐพฺโพ’ติ.\csromanspacer{} โส ตุณฺหสฺส, ยสฺส นกฺขมติ, โส ภาเสยฺย.}

\prose{สมฺมโต สํเฆน อายสฺมา สาริปุตฺโต เทวทตฺตํ ราชคเห ปกาเสตุํ ‘ปุพฺเพ เทวทตฺตสฺส อญฺญา ปกติ อโหสิ, อิทานิ อญฺญา ปกติ, ยา เทวทตฺโต กเรยฺย กาเยน วาจาย, น เตน พุทฺโธ วา ธมฺโม วา สํโฆ วา ทฏฺฐพฺโพ, เทวทตฺโตว เตน ทฏฺฐพฺโพ’ติ, ขมติ สํฆสฺส, ตสฺมา ตุณฺหี, เอวเมตํ ธารยามี”ติ.}

\prose{สมฺมโต จ อายสฺมา สาริปุตฺโต สมฺพหุเลหิ ภิกฺขูหิ สทฺธิํ ราชคหํ ปวิสิตฺวา เทวทตฺตํ ราชคเห ปกาเสสิ “ปุพฺเพ เทวทตฺตสฺส อญฺญา ปกติ อโหสิ, อิทานิ อญฺญา ปกติ, ยํ เทวทตฺโต กเรยฺย กาเยน วาจาย, น เตน พุทฺโธ วา ธมฺโม วา สํโฆ วา ทฏฺฐพฺโพ, เทวทตฺโตว เตน ทฏฺฐพฺโพ”ติ.\csromanspacer{} ตตฺถ เย เต มนุสฺสา อสฺสทฺธา อปฺปสนฺนา ทุพฺพุทฺธิโน, เต เอวมาหํสุ “อุสูยกา อิเม สมณา สกฺยปุตฺติยา เทวทตฺตสฺส ลาภสกฺการํ อุสูยนฺตี”ติ.\csromanspacer{} เย ปน เต มนุสฺสา สทฺธา ปสนฺนา ปณฺฑิตา พฺยตฺตา พุทฺธิมนฺโต, เต เอวมาหํสุ “น โข อิทํ โอรกํ ภวิสฺสติ, ยถา ภควา เทวทตฺตํ ราชคเห ปกาสาเปตี”ติ.}

\csromanmatikaanchor{mtk.167}
\csromantocmarkref{subsection}{อชาตสตฺตุกุมารวตฺถุ}{mtk.167}
\bookmark[dest={mtk.167},level=2]{อชาตสตฺตุกุมารวตฺถุ}
\csromanheader{2}{\textbf{อชาตสตฺตุกุมารวตฺถุ}}
\proseitem{339}{\csromanfolio{350}อถ โข เทวทตฺโต เยน อชาตสตฺตุ กุมาโร เตนุปสงฺกมิ, อุปสงฺกมิตฺวา อชาตสตฺตุํ กุมารํ เอตทโวจ “ปุพฺเพ โข กุมาร มนุสฺสา ทีฆายุกา, เอตรหิ อปฺปายุกา.\csromanspacer{} ฐานํ โข ปเนตํ วิชฺชติ, ยํ ตฺวํ กุมาโรว สมาโน กาลํ กเรยฺยาสิ, เตน หิ ตฺวํ กุมาร ปิตรํ หนฺตฺวา ราชา โหติ, อหํ ภควนฺตํ หนฺตฺวา พุทฺโธ ภวิสฺสามี”ติ.}

\prose{อถ โข อชาตสตฺตุ กุมาโร “อยฺโย โข เทวทตฺโต มหิทฺธิโก มหานุภาโว, ชาเนยฺยาสิ อยฺโย เทวทตฺโต”ติ อูรุยา โปตฺถนิกํ พนฺธิตฺวา ทิวา ทิวสฺส\footnote{ทิวา ทิวสสฺส (ก)} ภีโต อุพฺพิคฺโค อุสฺสงฺกี อุตฺรสฺโต สหสา อนฺเตปุรํ ปาวิสิ.\csromanspacer{} อทฺทสาสุํ โข อนฺเตปุเร อุปจารกา มหามตฺตา อชาตสตฺตุํ กุมารํ ทิวา ทิวสฺส ภีตํ อุพฺพิคฺคํ อุสฺสงฺกิํ อุตฺรสฺตํ สหสา อนฺเตปุรํ ปวิสนฺตํ, ทิสฺวาน อคฺคเหสุํ.\csromanspacer{} เต วิจินนฺตา อูรุยา โปตฺถนิกํ พทฺธํ\footnote{พนฺธํ (ก)} ทิสฺวาน อชาตสตฺตุํ กุมารํ เอตทโวจุํ “กิํ ตฺวํ กุมาร กตฺตุกาโมสี”ติ.\csromanspacer{} ปิตรมฺหี หนฺตุกาโมติ.\csromanspacer{} เกนาสิ อุสฺสาหิโตติ.\csromanspacer{} อยฺเยน เทวทตฺเตนาติ.\csromanspacer{} เอกจฺเจ มหามตฺตา เอวํ มติํ อกํสุ “กุมาโร จ หนฺตพฺโพ เทวทตฺโต จ, สพฺเพ จ ภิกฺขู หนฺตพฺพา”ติ.\csromanspacer{} เอกจฺเจ มหามตฺตา เอวํ มติํ อกํสุ “น ภิกฺขู หนฺตพฺพา, น ภิกฺขู กิญฺจิ อปรชฺฌนฺติ, กุมาโร จ หนฺตพฺโพ เทวทตฺโต จา”ติ.\csromanspacer{} เอกจฺเจ มหามตฺตา เอวํ มติํ อกํสุ “น กุมาโร จ หนฺตพฺโพ น เทวทตฺโต, น ภิกฺขู หนฺตพฺพา, รญฺโญ อาโรเจตพฺพํ, ยถา ราชา วกฺขติ ตถา กริสฺสามา”ติ.}

\prose{อถ โข เต มหามตฺตา อชาตสตฺตุํ กุมารํ อาทาย เยน ราชา มาคโธ เสนิโย พิมฺพิสาโร เตนุปสงฺกมิํสุ, อุปสงฺกมิตฺวา รญฺโญ มาคธสฺส เสนิยสฺส พิมฺพิสารสฺส เอตมตฺถํ อาโรเจสุํ.\csromanspacer{} กถํ ภเณ มหามตฺเตหิ มติ กตาติ.\csromanspacer{} เอกจฺเจ เทว มหามตฺตา เอวํ มติํ อกํสุ “กุมาโร จ หนฺตพฺโพ เทวทตฺโต จ, สพฺเพ จ ภิกฺขู หนฺตพฺพา”ติ.\csromanspacer{} เอกจฺเจ มหามตฺตา เอวํ มติํ อกํสุ “น ภิกฺขู หนฺตพฺพา, น ภิกฺขู กิญฺจิ อปรชฺฌนฺติ, กุมาโร จ หนฺตพฺโพ เทวทตฺโต จา”ติ.\csromanspacer{} เอกจฺเจ มหามตฺตา เอวํ มติํ อกํสุ “น กุมาโร จ หนฺตพฺโพ, น เทวทตฺโต, น ภิกฺขู หนฺตพฺพา, \csromanfolio{351}รญฺโญ อาโรเจตพฺพํ, ยถา ราชา วกฺขติ ตถา กริสฺสามา”ติ.\csromanspacer{} กิํ ภเณ กริสฺสติ พุทฺโธ วา ธมฺโม วา สํโฆ วา, นนุ ภควตา ปฏิกจฺเจว เทวทตฺโต ราชคเห ปกาสาปิโต “ปุพฺเพ เทวทตฺตสฺส อญฺญา ปกติ อโหสิ, อิทานิ อญฺญา ปกติ, ยํ เทวทตฺโต กเรยฺย กาเยน วาจาย, น เตน พุทฺโธ วา ธมฺโม วา สํโฆ วา ทฏฺฐพฺโพ, เทวทตฺโตว เตน ทฏฺฐพฺโพ”ติ.\csromanspacer{} ตตฺถ เย เต มหามตฺตา เอวํ มติํ อกํสุ “กุมาโร จ หนฺตพฺโพ เทวทตฺโต จ, สพฺเพ จ ภิกฺขู หนฺตพฺพา”ติ.\csromanspacer{} เต อฏฺฐาเน อกาสิ.\csromanspacer{} เย เต มหามตฺตา เอวํ มติํ อกํสุ “น ภิกฺขู หนฺตพฺพา, น ภิกฺขู กิญฺจิ อปรชฺฌนฺติ, กุมาโร จ หนฺตพฺโพ เทวทตฺโต จา”ติ.\csromanspacer{} เต นีเจ ฐาเน ฐเปสิ.\csromanspacer{} เย เต มหามตฺตา เอวํ มติํ อกํสุ “น กุมาโร จ หนฺตพฺโพ น เทวทตฺโต, น ภิกฺขู หนฺตพฺพา, รญฺโญ อาโรเจตพฺพํ, ยถา ราชา วกฺขติ ตถา กริสฺสามา”ติ.\csromanspacer{} เต อุจฺเจ ฐาเน ฐเปสิ.\csromanspacer{} อถ โข ราชา มาคโธ เสนิโย พิมฺพิสาโร อชาตสตฺตุํ กุมารํ เอตทโวจ “กิสฺส มํ ตฺวํ กุมาร หนฺตุกาโมสี”ติ.\csromanspacer{} รชฺเชนามฺหิ เทว อตฺถิโกติ.\csromanspacer{} สเจ โข ตฺวํ กุมาร รชฺเชน อตฺถิโก, “เอตํ เต รชฺชนฺ”ติ อชาตสตฺตุสฺส กุมารสฺส รชฺชํ นิยฺยาเทสิ.}

\csromanmatikaanchor{mtk.168}
\csromantocmarkref{subsection}{อภิมารเปสน}{mtk.168}
\bookmark[dest={mtk.168},level=2]{อภิมารเปสน}
\csromanheader{2}{\textbf{อภิมารเปสน}}
\proseitem{340}{อถ โข เทวทตฺโต เยน อชาตสตฺตุ กุมาโร เตนุปสงฺกมิ, อุปสงฺกมิตฺวา อชาตสตฺตุํ กามารํ เอตทโวจ “ปุริเส มหาราช อาณาเปหิ, เย สมณํ โคตมํ ชีวิตา โวโรเปสฺสนฺตี”ติ.\csromanspacer{} อถ โข อชาตสตฺตุ กุมาโร มนุสฺเส อาณาเปสิ “ยถา ภเณ อยฺโย เทวทตฺโต อาห, ตถา กเรยฺยาถา”ติ.\csromanspacer{} อถ โข เทวทตฺโต เอกํ ปุริสํ อาณาเปสิ “คจฺฉาวุโส, อมุกสฺมิํ โอกาเส สมโณ โคตโม วิหรติ, ตํ ชีวิตา โวโรเปตฺวา อิมินา มคฺเคน อาคจฺฉา”ติ.\csromanspacer{} ตสฺมิํ มคฺเค เทฺว ปุริเส ฐเปสิ “โย อิมินา มคฺเคน เอโก ปุริโส อาคจฺฉติ, ตํ ชีวิตา โวโรเปตฺวา อิมินา มคฺเคน อาคจฺฉถา”ติ.\csromanspacer{} ตสฺมิํ มคฺเค จตฺตาโร ปุริเส ฐเปสิ “เย อิมินา มคฺเคน เทฺว ปุริสา อาคจฺฉนฺติ, เต ชีวิตา โวโรเปตฺวา อิมินา มคฺเคน อาคจฺฉถา”ติ.\csromanspacer{} ตสฺมิํ มคฺเค อฏฺฐ ปุริเส ฐเปสิ “เย อิมินา มคฺเคน จตฺตาโร ปุริสา อาคจฺฉนฺติ, เต ชีวิตา โวโรเปตฺวา อิมินา \csromanfolio{352}มคฺเคน อาคจฺฉถา”ติ.\csromanspacer{} ตสฺมิํ มคฺเค โสฬส ปุริเส ฐเปสิ “เย อิมินา มคฺเคน อฏฺฐ ปุริสา อาคจฺฉนฺติ, เต ชีวิตา โวโรเปตฺวา อาคจฺฉถา”ติ.}

\prose{อถ โข โส เอโก ปุริโส อสิจมฺมํ คเหตฺวา ธนุกลาปํ สนฺนยฺหิตฺวา เยน ภควา เตนุปสงฺกมิ, อุปสงฺกมิตฺวา ภควโต อวิทูเร ภีโต อุพฺพิคฺโค อุสฺสงฺกี อุตฺรสฺโต ปตฺถทฺเธน กาเยน อฏฺฐาสิ.\csromanspacer{} อทฺทสา โข ภควา ตํ ปุริสํ ภีตํ อุพฺพิคฺคํ อุสฺสงฺกิํ อุตฺรสฺตํ ปตฺถทฺเธน กาเยน ฐิตํ, ทิสฺวาน ตํ ปุริสํ เอตทโวจ “เอหาวุโส, มา ภายี”ติ.\csromanspacer{} อถ โข โส ปุริโส อสิจมฺมํ เอกมนฺตํ กริตฺวา ธนุกลาปํ นิกฺขิปิตฺวา เยน ภควา เตนุปสงฺกมิ, อุปสงฺกมิตฺวา ภควโต ปาเทสุ สิรสา นิปติตฺวา ภควนฺตํ เอตทโวจ “อจฺจโย มํ ภนฺเต อจฺจคมา ยถาพาลํ ยถามูฬฺหํ ยถา-อกุสลํ, โยหํ ทุฏฺฐจิตฺโต วธกจิตฺโต อิธูปสงฺกนฺโต, ตสฺส เม ภนฺเต ภควา อจฺจยํ อจฺจยโต ปฏิคฺคณฺหาตุ อายติํ สํวรายา”ติ.\csromanspacer{} ตคฺฆ ตฺวํ อาวุโส อจฺจโย อจฺจคมา ยถาพาลํ ยถามูฬฺหํ ยถา-อกุสลํ, ยํ ตฺวํ ทุฏฺฐจิตฺโต วธกจิตฺโต อิธูปสงฺกนฺโต, ยโต จ โข ตฺวํ อาวุโส อจฺจยํ อจฺจยโต ทิสฺวา ยถาธมฺมํ ปฏิกโรสิ, ตํ เต มยํ ปฏิคฺคณฺหาม, วุทฺธิ เหสา อาวุโส อริยสฺส วินเย, โย อจฺจยํ อจฺจยโต ทิสฺวา ยถาธมฺมํ ปฏิกโรติ, อายติํ สํวรํ อาปชฺชตีติ.}

\prose{อถ โข ภควา ตสฺส ปุริสสฺส อนุปุพฺพิํ กถํ กถสิ, เสยฺยถิทํ ทานกถํ สีลกถํ สคฺคกถํ กามานํ อาทีนวํ โอการํ สํกิเลสํ เนกฺขมฺเม อานิสํสํ ปกาเสสิ.\csromanspacer{} ยทา ตํ ภควา อญฺญาสิ กลฺลจิตฺตํ มุทุจิตฺตํ วินีวรณจิตฺตํ อุทคฺคจิตฺตํ ปสนฺนจิตฺตํ, อถ ยา พุทฺธานํ สามุกฺกํสิกา ธมฺมเทสนา, ตํ ปกาเสสิ ทุกฺขํ สมุทยํ นิโรธํ มคฺคํ, เสยฺยถาปิ นาม สุทฺธํ วตฺถํ อปคตกาฬกํ สมฺมเทว รชนํ ปฏิคฺคเหยฺย, เอวเมว ตสฺส ปุริสสฺส ตสฺมิํเยว อาสเน วิรชํ วีตมลํ ธมฺมจกฺขุํ อุทปาทิ “ยํ กิญฺจิ สมุทยธมฺมํ, สพฺพํ ตํ นิโรธธมฺมนฺ”ติ.\csromanspacer{} อถ โข โส ปุริโส ทิฏฺฐธมฺโม ปตฺตธมฺโม วิทิตธมฺโม ปริโยคาฬฺหธมฺโม ติณฺณวิจิกิจฺโฉ วิคตกถํกโถ เวสารชฺชปฺปตฺโต อปรปฺปจฺจโย สตฺถุสาสเน ภควนฺตํ เอตทโวจ “อภิกฺกนฺตํ ภนฺเต, อภิกฺกนฺตํ ภนฺเต, เสยฺยถาปิ ภนฺเต นิกฺกุชฺชิตํ วา อุกฺกุชฺเชยฺย, ปฏิจฺฉนฺนํ วา วิวเรยฺย, \csromanfolio{353}มูฬสฺส วา มคฺคํ อาจิกฺเขยฺย, อนฺธกาเร วา เตลปชฺโชตํ ธาเรยฺย ‘จกฺขุมนฺโต รูปานิ ทกฺขนฺตี’ติ.\csromanspacer{} เอวเมวํ ภควตา อเนกปริยาเยน ธมฺโม ปกาสิโต.\csromanspacer{} เอสาหํ ภนฺเต ภควนฺตํ สรณํ คจฺฉามิ ธมฺมญฺจ ภิกฺขุสํฆญฺจ, อุปาสกํ มํ ภควา ธาเรตุ อชฺชตคฺเค ปาณุเปตํ สรณํ คตนฺ”ติ.\csromanspacer{} อถ โข ภควา ตํ ปุริสํ เอตทโวจ “มา โข ตฺวํ อาวุโส อิมินา มคฺเคน คจฺฉ, อิมินา มคฺเคน คจฺฉาหี”ติ, อญฺเญน มคฺเคน อุยฺโยเชสิ.}

\prose{อถ โข เต เทฺว ปุริสา “กิํ นุ โข โส เอโก ปุริโส จิเรน อาคจฺฉตี”ติ ปฏิปถํ คจฺฉนฺตา อทฺทสาสุํ ภควนฺตํ อญฺญตรสฺมิํ รุกฺขมูเล นิสินฺนํ, ทิสฺวาน เยน ภควา เตนุสงฺกมิํสุ, อุปสงฺกมิตฺวา ภควนฺตํ อภิวาเทตฺวา เอกมนฺตํ นิสีทิํสุ.\csromanspacer{} เตสํ ภควา อนุปุพฺพิํ กถํ กเถสิ ฯเปฯ อปรปฺปจฺจยา สตฺถุสาสเน ภควนฺตํ เอตทโวจุํ “อภิกฺกนฺตํ ภนฺเต ฯเปฯ อุปาสเก โน ภควา ธาเรตุ อชฺชตคฺเค ปาณุเปตํ สรณํ คเต”ติ.\csromanspacer{} อถ โข ภควา เต ปุริเส เอตทโวจ “มา โข ตุเมฺห อาวุโส อิมินา มคฺเคน คจฺฉิตฺถ, อิมินา มคฺเคน คจฺฉถา, อญฺเญน มคฺเคน อุยฺโยเชสิ.}

\prose{อถ โข เต จตฺตาโร ปุริสา ฯเปฯ.\csromanspacer{} อถ โข เต อฏฺฐปุริสา ฯเปฯ.\csromanspacer{} อถ โข เต โสฬส ปุริสา “กิํ นุ โข เต อฏฺฐ ปุริสา จิเรน อาคจฺฉนฺตี”ติ ปฏิปถํ คจฺฉนฺตา อทฺทสาสุํ ภควนฺตํ อญฺญตรสฺมิํ รุกฺขมูเล นิสินฺนํ, ทิสฺวาน เยน ภควา เตนุปสงฺกมิํสุ, อุปสงฺกมิตฺวา ภควนฺตํ อภิวาเทตฺวา เอกมนฺตํ นิสีทิํสุ.\csromanspacer{} เตสํ ภควา อนุปุพฺพิํ กถํ กเถสิ, เสยฺยถิทํ, ทานกถํ ฯเปฯ อปรปฺปจฺจยา สตฺถุสาสเน ภควนฺตํ เอตทโวจุํ “อภิกฺกนฺตํ ภนฺเต ฯเปฯ อุปาสเก โน ภควา ธาเรตุ อชฺชตคฺเค ปาณุเปตํ สรณํ คเต”ติ.\csromanspacer{} อถ โข ภควา เต ปุริเส เอตทโวจ “มา โข ตุเมฺห อาวุโส อิมินา มคฺเคน คจฺฉิตฺถ, อิมินา มคฺเคน คจฺฉถา”ติ, อญฺเญน มคฺเคน อุยฺโยเชสิ.}

\prose{อถ โข โส เอโก ปุริโส เยน เทวทตฺโต เตนุปสงฺกมิ, อุปสงฺกมิตฺวา เทวทตฺตํ เอตทโวจ “นาหํ ภนฺเต สกฺโกมิ ตํ ภควนฺตํ ชีวิตา โวโรเปตุํ, มหิทฺธิโก โส ภควา มหานุภาโว”ติ.\csromanspacer{} อลํ อาวุโส, มา ตฺวํ สมณํ โคตมํ ชีวิตา โวโรเปสิ, อหเมว สมณํ โคตมํ ชีวิตา โวโรเปสฺสามีติ.}

\csromanmatikaanchor{mtk.169}
\csromantocmarkref{subsection}{โลหิตุปฺปาทกกมฺม}{mtk.169}
\bookmark[dest={mtk.169},level=2]{โลหิตุปฺปาทกกมฺม}
\csromanheader{2}{\textbf{โลหิตุปฺปาทกกมฺม}}
\proseitem{341}{\csromanfolio{354}เตน โข ปนสมเยน ภควา คิชฺฌกูฏสฺส ปพฺพตสฺส ฉายายํ จงฺกมติ.\csromanspacer{} อถ โข เทวทตฺโต คิชฺฌกูฏํ ปพฺพตํ อารุหิตฺวา มหติํ สิลํ ปวิชฺฌิ “อิมาย สมณํ โคตมํ ชีวิตา โวโรเปสฺสามี”ติ.\csromanspacer{} เทฺว ปพฺพตกูฏา สมาคนฺตฺวา ตํ สิลํ สมฺปฏิจฺฉิํสุ, ตโต ปปติกา อุปฺปติตฺวา ภควโต ปาเท รุหิรํ อุปฺปาเทสิ.\csromanspacer{} อถ โข ภควา อุทฺธํ อุลฺโลเกตฺวา เทวทตฺตํ เอตทโวจ “พหุํ ตยา โมฆปุริส อปุญฺญํ ปสุตํ, ยํ ตฺวํ ทุฏฺฐจิตฺโต วธกจิตฺโต ตถาคตสฺส รุหิรํ อุปฺปาเทสี”ติ.\csromanspacer{} อถ โข ภควา ภิกฺขู อามนฺเตสิ “อิทํ ภิกฺขเว เทวทตฺเตน ปฐมํ อานนฺตริยํ กมฺมํ อุปจิตํ, ยํทุฏฺฐจิตฺเตน วธกจิตฺเตน ตถาคตสฺส รุหิรํ อุปฺปาทิตนฺ”ติ.}

\prose{อสฺโสสุํ โข ภิกฺขู “เทวทตฺเตน กิร ภควโต วโธ ปยุตฺโต”ติ.\csromanspacer{} เต จ ภิกฺขู ภควโต วิหารสฺส ปริโต ปริโต จงฺกมนฺติ อุจฺจาสทฺทา มหาสทฺทา สชฺฌายํ กโรนฺตา ภควโต รกฺขาวรณคุตฺติยา.\csromanspacer{} อสฺโสสิ โข ภควา อุจฺจาสทฺทํ มหาสทฺทํ สชฺฌายสทฺทํ, สุตฺวาน อายสฺมนฺตํ อานนฺทํ อามนฺเตสิ “กิํ นุ โข โส อานนฺท อุจฺจาสทฺโท มหาสทฺโท สชฺฌายสทฺโท”ติ.\csromanspacer{} อสฺโสสุํ โข ภนฺเต ภิกฺขู “เทวทตฺเตน กิร ภควโต วโธ ปยุตฺโต”ติ.\csromanspacer{} เต จ\footnote{เตธ (สี)} ภนฺเต ภิกฺขู ภควโต วิหารสฺส ปริโต ปริโต จงฺกมนฺติ อุจฺจาสทฺทา มหาสทฺทา สชฺฌายํ กโรนฺตา ภควโต รกฺขาวรณคุตฺติยา, โส เอโส ภควา อุจฺจาสทฺโท มหาสทฺโท สชฺฌายสทฺโทติ.\csromanspacer{} เตน หานนฺท มม วจเนน เต ภิกฺขู อามนฺเตหิ “สตฺถา อายสฺมนฺเต อามนฺเตตี”ติ. “เอวํ ภนฺเต”ติ โข อายสฺมา อานนฺโท ภควโต ปฏิสฺสุตฺวา เยน เต ภิกฺขู เตนุปสงฺกมิ, อุปสงฺกมิตฺวา เต ภิกฺขู เอตทโวจ “สตฺถา อายสฺมนฺเต อามนฺเตตี”ติ. “เอวมาวุโส”ติ โข เต ภิกฺขู อายสฺมโต อานนฺทสฺส ปฏิสฺสุตฺวา เยน ภควา เตนุปสงฺกมิํสุ, อุปสงฺกมิตฺวา ภควนฺตํ อภิวาเทตฺวา เอกมนฺตํ นิสีทิํสุ.\csromanspacer{} เอกมนฺตํ นิสินฺเน โข เต ภิกฺขู ภควา เอตทโวจ\csromandash{}}

\prose{\csromanfolio{355}อฏฺฐานเมตํ ภิกฺขเว อนวกาโส, ยํ ปรูปกฺกเมน ตถาคตํ ชีวิตา โวโรเปยฺย, อนุปกฺกเมน ภิกฺขเว ตถาคตา ปรินิพฺพายนฺติ.}

\prose{\csromansymbolfootnote{*}{วิ 4. 343; อํ 2. 108 ปิฏฺเฐสุปิ.}ปญฺจิเม ภิกฺขเว สตฺถาโร สนฺโต สํวิชฺชมานา โลกสฺมิํ, กตเม ปญฺจ, อิธ ภิกฺขเว เอกจฺโจ สตฺถา อปริสุทฺธสีโล สมาโน “ปริสุทฺธสีโลมฺหี”ติ ปฏิชานาติ “ปริสุทฺธํ เม สีลํ ปริโยทาตํ อสํกิลิฏฺฐนฺ”ติ จ.\csromanspacer{} ตเมนํ สาวกา เอวํ ชานนฺติ “อยํ โข ภวํ สตฺถา อปริสุทฺธสีโล สมาโน “ปริสุทฺธสีโลมฺหี”ติ ปฏิชานาติ “ปริสุทฺธํ เม สีลํ ปริโยทาตํ อสํกิลิฏฺฐนฺ”ติ จ.\csromanspacer{} มยญฺเจว โข ปน คิหีนํ อาโรเจยฺยาม “นาสฺสสฺส มนาปํ, ยํ โข ปนสฺส อมนาปํ, กถํ นํ มยํ เตน สมุทาจเรยฺยาม.\csromanspacer{} สมฺมนฺนติ โข ปน จีวรปิณฺฑปาตเสนาสนคิลานปฺปจฺจยเภสชฺชปริกฺขาเรน.\csromanspacer{} ยํ ตุโม กริสฺสติ, ตุโมว เตน ปญฺญายิสฺสตี”ติ.\csromanspacer{} เอวรูปํ โข ภิกฺขเว สตฺถารํ สาวกา สีลโต รกฺขนฺติ, เอวรูโป จ ปน สตฺถา สาวเกหิ สีลโต รกฺขํ ปจฺจาสีสติ.}

\prose{ปุน จปรํ ภิกฺขเว อิเธกจฺโจ สตฺถา อปริสุทฺธ-อาชีโว สมาโน ฯเปฯ อปริสุทฺธธมฺมเทสโน สมาโน ฯเปฯ อปริสุทฺธเวยฺยากรโณ สมาโน ฯเปฯ อปริสุทฺธญาณทสฺสโน สมาโน “ปริสุทฺธญาณทสฺสโนมฺหี”ติ ปฏิชานาติ “ปริสุทฺธํ เม ญาณทสฺสนํ ปริโยทาตํ อสํกิลิฏฺฐนฺ”ติ จ.\csromanspacer{} ตเมนํ สาวกา เอวํ ชานนฺติ อยํ โข ภวํ สตฺถา อปริสุทฺธญาณทสฺสโน สมโน “ปริสุทฺธญาณทสฺสโนมฺหี”ติ ปฏิชานาติ “ปริสุทฺธํ เม ญาณทสฺสนํ ปริโยทาตํ อสํกิลิฏฺฐนฺ”ติ จ.\csromanspacer{} มยญฺเจว โข ปน คิหีนํ อาโรเจยฺยาม “นาสฺสสฺส มนาปํ, ยํ โข ปนสฺส อมนาปํ, กถํ นํ มยํ เตน สมุทาจเรยฺยาม, สมฺมนฺนติ โข ปน จีวรปิณฺฑปาตเสนาสนคิลานปฺปจฺจยเภสชฺชปริกฺขาเรน.\csromanspacer{} ยํ ตุโม กริสฺสติ, ตุโมว เตน ปญฺญายิสฺสตี”ติ.\csromanspacer{} เอวรูปํ โข ภิกฺขเว สตฺถารํ สาวกา ญาณทสฺสนโต รกฺขนฺติ, เอวรูโป จ ปน สตฺถา สาวเกหิ ญาณทสฺสนโต รกฺขํ ปจฺจาสีสติ.\csromanspacer{} อิเม โข ภิกฺขเว ปญฺจ สตฺถาโร สนฺโต สํวิชฺชมานา โลกสฺมิํ.\csromanspacer{} อหํ โข ปน ภิกฺขเว ปริสุทฺธสีโล สมาโน “ปริสุทฺธสีโลมฺหี”ติ ปฏิชานามิ “ปริสุทฺธํ เม สีลํ ปริโยทาตํ อสํกิลิฏฺฐนฺ”ติ จ, น จ มํ สาวกา สีลโต รกฺขนฺติ, น จาหํ สาวเกหิ สีลโต รกฺขํ ปจฺจาสีสามิ.\csromanspacer{} อหํ โข \csromanfolio{356}ปน ภิกฺขเว ปริสุทฺธาชีโว สมาโน ฯเปฯ ปริสุทฺธธมฺมเทสโน สมาโน ฯเปฯ ปริสุทฺธเวยฺยากรโณ สมาโน ฯเปฯ ปริสุทฺธญาณทสฺสโน สมาโน “ปริสุทฺธญาณทสฺสโนมฺหี”ติ ปฏิชานามิ “ปริสุทฺธํ เม ญาณทสฺสนํ ปริโยทาตํ อสํกิลิฏฺฐนฺ”ติ จ, น จ มํ สาวกา ญาณทสฺสนโต รกฺขนฺติ, น จาหํ สาวเกหิ ญาณทสฺสนโต รกฺขํ ปจฺจาสีสามิ.\csromanspacer{} อฏฺฐานเมตํ ภิกฺขเว อนวกาโส, ยํ ปรูปกฺกเมน ตถาคตํ ชีวิตา โวโรเปยฺย, อนุปกฺกเมน ภิกฺขเว ตถาคตา ปรินิพฺพายนฺติ.\csromanspacer{} คจฺฉถ ตุเมฺห ภิกฺขเว ยถาวิหารํ, อรกฺขิยา ภิกฺขเว ตถาคตาติ.}

\csromanmatikaanchor{mtk.170}
\csromantocmarkref{subsection}{นาฬาคิริเปสน}{mtk.170}
\bookmark[dest={mtk.170},level=2]{นาฬาคิริเปสน}
\csromanheader{2}{\textbf{นาฬาคิริเปสน}}
\proseitem{342}{เตน โข ปน สมเยน ราชคเห นาฬาคิริ นาม หตฺถี จณฺโฑ โหติ มนุสฺสฆาตโก.\csromanspacer{} อถ โข เทวทตฺโต ราชคหํ ปวิสิตฺวา หตฺถิสาลํ คนฺตฺวา หตฺถิภณฺเฑ เอตทโวจ “มยํ โข ภเณ ราชญาตกา นาม ปฏิพลา นีจฏฺฐานิยํ อุจฺจฏฺฐาเน ฐเปตุํ, ภตฺตมฺปิ เวตนมฺปิ วฑฺฒาเปตุํ.\csromanspacer{} เตน หิ ภเณ ยทา สมโณ โคตโม อิมํ รจฺฉํ ปฏิปนฺโน โหติ, ตทา อิมํ นาฬาคิริํ หตฺถิํ มุญฺเจตฺวา อิมํ รจฺฉํ ปฏิปาเทถา”ติ. “เอวํ ภนฺเต”ติ โข เต หตฺถิภณฺฑา เทวทตฺตสฺส ปจฺจสฺโสสุํ.\csromanspacer{} อถ โข ภควา ปุพฺพณฺหสมยํ นิวาเสตฺวา ปตฺตจีวรมาทาย สมฺพหุเลหิ ภิกฺขูหิ สทฺธิํ รชคหํ ปิณฺฑาย ปาวิสิ.\csromanspacer{} อถ โข ภควา ตํ รจฺฉํ ปฏิปชฺชิ.\csromanspacer{} อทฺทสาสุํ โข เต หตฺถิภณฺฑา ภควนฺตํ ตํ รจฺฉํ ปฏิปนฺนํ, ทิสฺวาน นาฬาคิริํ หตฺถิํ มุญฺจิตฺวา ตํ รจฺฉํ ปฏิปาเทสุํ.\csromanspacer{} อทฺทสา โข นาฬาคิริ หตฺถี ภควนฺตํ ทูรโตว อาคจฺฉนฺตํ, ทิสฺวาน โสณฺฑํ อุสฺสาเปตฺวา ปหฏฺฐกณฺณวาโล เยน ภควา เตน อภิธาวิ.\csromanspacer{} อทฺทสาสุํ โข เต ภิกฺขู นาฬาคิริํ หตฺถิํ ทูรโตว อาคจฺฉนฺตํ, ทิสฺวาน ภควนฺตํ เอตทโวจุํ “อยํ ภนฺเต นาฬาคิริ หตฺถี จณฺโฑ มนุสฺสฆาตโก อิมํ รจฺฉํ ปฏิปนฺโน, ปฏิกฺกมตุ ภนฺเต ภควา, ปฏิกฺกมตุ สุคโต”ติ.\csromanspacer{} อาคจฺฉถ ภิกฺขเว, มา ภายิตฺถ.\csromanspacer{} อฏฺฐานเมตํ ภิกฺขเว อนวกาโส, ยํ ปรูปกฺกเมน ตถาคตํ ชีวิตา โวโรเปยฺย, อนุปกฺกเมน ภิกฺขเว ตถาคตา ปรินิพฺพายนฺตีติ.\csromanspacer{} ทุติยมฺปิ โข เต ภิกฺขู ฯเปฯ.\csromanspacer{} ตติยมฺปิ โข เต ภิกฺขู ภควนฺตํ เอตทโวจุํ “อยํ ภนฺเต นาฬาคิริ หตฺถี จณฺโฑ มนุสฺสฆาตโก อิมํ รจฺฉํ ปฏิปนฺโน, ปฏิกฺกมตุ ภนฺเต ภควา, ปฏิกฺกมตุ สุคโต”ติ.\csromanspacer{} อาคจฺฉถ ภิกฺขเว, มา ภายิตฺถ. \csromanfolio{357}อฏฺฐานเมตํ ภิกฺขเว อนวกาโส, ยํ ปรูปกฺกเมน ตถาคตํ ชีวิตา โวโรเปยฺย, อนุปกฺกเมน ภิกฺขเว ตถาคตา ปรินิพฺพายนฺตีติ.}

\prose{เตน โข ปน สมเยน มนุสฺสา ปาสาเทสุปิ หมฺมิเยสุปิ ฉทเนสุปิ อารุฬฺหา อจฺฉนฺติ.\csromanspacer{} ตตฺถ เย เต มนุสฺสา อสฺสทฺธา อปฺปสนฺนา ทุพฺพุทฺธิโน, เต เอวมาหํสุ “อภิรูโป วต โภ\footnote{อภิรูโป วต โภ โคตโม (สฺยา, กํ)} มหาสมโณ นาเคน วิเหฐียิสฺสตี”ติ.\csromanspacer{} เย ปน เต มนุสฺสาสทฺธา ปสนฺนา ปณฺฑิตา พฺยตฺตา พุทฺธิมนฺโต, เต เอวมาหํสุ “นจิรสฺสํ วต โภ นาโค นาเคน สงฺคาเมสฺสตี”ติ.\csromanspacer{} อถ โข ภควา นาฬาคิริํ หตฺถิํ เมตฺเตน จิตฺเตน ผริ.\csromanspacer{} อถ โข นาฬาคิริ หตฺถี ภควโต\footnote{ภควตา (สี)} เมตฺเตน จิตฺเตน ผุฏฺโฐ\footnote{ผุโฏ (ก)} โสณฺฑํ โอโรเปตฺวา เยน ภควา เตนุปสงฺกมิ, อุปสงฺกมิตฺวา ภควโต ปุรโต อฏฺฐาสิ.\csromanspacer{} อถ โข ภควา ทกฺขิเณน หตฺเถน นาฬาคิริสฺส หตฺถิสฺส กุมฺภํ ปรามสนฺโต นาฬาคิริํ หตฺถิํ อิมาหิ คาถาหิ อชฺฌภาสิ\csromandash{}}

\csromangathameasure{“มา กุญฺชร นาคมาสโท, ทุกฺขํ หิ กุญฺชร นาคมาสโท. \\ น หิ นาคหตสฺส กุญฺชร, สุคติ โหติ อิโต ปรํ ยโต. \\ มา จ มโท มา จ ปมาโท, น หิ ปมตฺตา สุคติํ วชนฺติ เต. \\ ตฺวญฺเญว ตถา กริสฺสสิ, เยน ตฺวํ สุคติํ คมิสฺสสี”ติ.}
\csromangathagroup{\csromangathastanza{“มา กุญฺชร นาคมาสโท, ทุกฺขํ หิ กุญฺชร นาคมาสโท. \\ น หิ นาคหตสฺส กุญฺชร, สุคติ โหติ อิโต ปรํ ยโต.}\csromangathastanza{มา จ มโท มา จ ปมาโท, น หิ ปมตฺตา สุคติํ วชนฺติ เต. \\ ตฺวญฺเญว ตถา กริสฺสสิ, เยน ตฺวํ สุคติํ คมิสฺสสี”ติ.}}

\prose{อถ โข นาฬาคิริ หตฺถี โสณฺฑาย ภควโต ปาทปํสูนิ คเหตฺวา อุปริมุทฺธนิ อากิริตฺวา ปฏิกุฏิโยว\footnote{ปฏิกุฏิโต ปฏิสกฺกิ. (สี, สฺยา)} โอสกฺกิ.\csromanspacer{} ยาว ภควนฺตํ อทฺทกฺขิ.\csromanspacer{} อถ โข นาฬาคิริ หตฺถี หตฺถิสาลํ คนฺตฺวา สเก ฐาเน อฏฺฐาสิ, ตถา ทนฺโต จ ปน นาฬาคิริ หตฺถี อโหสิ.\csromanspacer{} เตน โข ปน สมเยน มนุสฺสา อิมํ คาถํ คายนฺติ\csromandash{}}

\prose{\csromansymbolfootnote{*}{ม 2. 308; ขุ 2. 335 ปิฏฺเฐสุปิ.}“ทณฺเฑเนเก ทมยนฺติ, องฺกุเสหิ กสาหิ จ.}

\csromangathameasure{อทณฺเฑน อสตฺเถน, นาโค ทนฺโต มเหสินา”ติ.}
\csromangathagroup{\csromangathastanza{อทณฺเฑน อสตฺเถน, นาโค ทนฺโต มเหสินา”ติ.}}

\prose{มนุสฺสา อุชฺฌายนฺติ ขิยฺยนฺติ วิปาเจนฺติ “ยาว ปาโป อยํ เทวทตฺโต อลกฺขิโก, ยตฺร หิ นาม สมณสฺส โคตมสฺส เอวํมหิทฺธิกสฺส เอวํมหานุภาวสฺส วธาย ปรกฺกมิสฺสตี”ติ.\csromanspacer{} เทวทตฺตสฺส ลาภสกฺกาโร ปริหายิ, ภควโต จ ลาภสกฺกาโร อภิวฑฺฒิ.}

\csromanmatikaanchor{mtk.171}
\csromantocmarkref{subsection}{ปญฺจวตฺถุยาจนกถา}{mtk.171}
\bookmark[dest={mtk.171},level=2]{ปญฺจวตฺถุยาจนกถา}
\csromanheader{2}{\textbf{ปญฺจวตฺถุยาจนกถา}}
\proseitem{343}{\csromanfolio{358}\csromansymbolfootnote{*}{วิ 2. 97 ปิฏฺเฐปิ.}เตน โข ปน สมเยน เทวทตฺโต ปริหีนลาภสกฺกาโร สปริโส กุเลสุ วิญฺญาเปตฺวา ภุญฺชติ.\csromanspacer{} มนุสฺสา อุชฺฌายนฺติ ขิยฺยนฺติ วิปาจินฺติ “กถํ หิ นาม สมณา สกฺยปุตฺติยา กุเลสุ วิญฺญาเปตฺวา ภุญฺชิสฺสนฺติ, กสฺส สมฺปนฺนํ น มนาปํ, กสฺส สาทุํ น รุจฺจตี”ติ.\csromanspacer{} อสฺโสสุํ โข ภิกฺขู เตสํ มนุสฺสานํ อุชฺฌายนฺตานํ ขิยฺยนฺตานํ วิปาเจนฺตานํ.\csromanspacer{} เย เต ภิกฺขู อปฺปิจฺฉา ฯเปฯ เต อุชฺฌายนฺติ ขิยฺยนฺติ วิปาเจนฺติ “กถํ หิ นาม เทวทตฺโต สปริโส กุเลสุ วิญฺญาเปตฺวา วิญฺญาเปตฺวา ภุญฺชิสฺสตี”ติ.\csromanspacer{} ภควโต เอตมตฺถํ อาโรเจสุํ ฯเปฯ.\csromanspacer{} สจฺจํ กิร ตฺวํ เทวทตฺต สปริโส กุเลสุ วิญฺญาเปตฺวา วิญฺญาเปตฺวา ภุญฺชสีติ.\csromanspacer{} สจฺจํ ภควาติ ฯเปฯ วิครหิตฺวา ฯเปฯ ธมฺมิํ กถํ กตฺวา ภิกฺขู อามนฺเตสิ “เตน หิ ภิกฺขเว ภิกฺขูนํ กุเลสุ ติกโภชนํ ปญฺญเปสฺสามิ ตโย อตฺถวเส ปฏิจฺจ, ทุมฺมงฺกูนํ ปุคฺคลานํ นิคฺคหาย, เปสลานํ ภิกฺขูนํ ผาสุวิหาราย ‘มา ปาปิจฺฉา ปกฺขํ นิสฺสาย สํฆํ ภินฺเทยฺยุนฺ’ติ, กุลานุทฺทยาย\footnote{กุลานุทยตาย (สี, สฺยา)} จ.\csromanspacer{} คณโภชเน ยถาธมฺโม กาเรตพฺโพ”ติ.}

\prose{\csromansymbolfootnote{+}{วิ 1. 263 ปิฏฺเฐปิ.}อถ โข เทวทตฺโต เยน โกกาลิโก กฏโมทกติสฺสโก\footnote{กุฏโมรกติสฺสโก (สี, สฺยา)} ขณฺฑเทวิยา ปุตฺโต สมุทฺททตฺโต เตนุปสงฺกมิ, อุปสงฺกมิตฺวา โกกาลิกํ กฏโมทกติสฺสกํ ขณฺฑเทวิยา ปุตฺตํ สมุทฺททตฺตํ เอตทโวจ “เอถ มยํ อาวุโส สมณสฺส โคตมสฺส สํฆเภทํ กริสฺสาม จกฺกเภทนฺ”ติ.\csromanspacer{} เอวํ วุตฺเต โกกาลิโก เทวทตฺตํ เอตทโวจ “สมโณ โข อาวุโส โคตโม มหิทฺธิโก มหานุภาโว, กถํ มยํ สมณสฺส โคตมสฺส สํฆเภทํ กริสฺสาม จกฺกเภทนฺ”ติ.\csromanspacer{} เอถ มยํ อาวุโส สมณํ โคตมํ อุปสงฺกมิตฺวา ปญฺจ วตฺถูนิ ยาจิสฺสาม “ภควา ภนฺเต อเนกปริยาเยน อปฺปิจฺฉสฺส สนฺตุฏฺฐสฺส สลฺเลขสฺส ธุตสฺส ปาสาทิกสฺส อปจยสฺส วีริยารมฺภสฺส วณฺณวาที, อิมานิ ภนฺเต ปญฺจ วตฺถูนิ อเนกปริยาเยน อปฺปิจฺฉตาย สนฺตุฏฺฐิยา สลฺเลขาย ธุตตาย ปาสาทิกตาย อปจยาย วีริยารมฺภาย สํวตฺตนฺติ, สาธุ ภนฺเต ภิกฺขู ยาวชีวํ อารญฺญิกา อสฺสุ, โย \csromanfolio{359}คามนฺตํ โอสเรยฺย, วชฺชํ นํ ผุเสยฺย.\csromanspacer{} ยาวชีวํ ปิณฺฑปาติกา อสฺสุ, โย นิมนฺตนํ สาทิเยยฺย, วชฺชํ นํ ผุเสยฺย.\csromanspacer{} ยาวชีวํ ปํสุกูลิกา อสฺสุ, โย คหปติจีวรํ สาทิเยยฺย, วชฺชํ นํ ผุเสยฺย.\csromanspacer{} ยาวชีวํ รุกฺขมูลิกา อสฺสุ, โย ฉนฺนํ อุปคจฺเฉยฺย, วชฺชํ นํ ผุเสยฺย.\csromanspacer{} ยาวชีวํ มจฺฉมํสํ น ขาเทยฺยุํ, โย มจฺฉมํสํ ขาเทยฺย, วชฺชํ นํ ผุเสยฺยา”ติ.\csromanspacer{} อิมานิ ปญฺจ วตฺถูนิ สมโณ โคตโม นานุชานิสฺสติ.\csromanspacer{} เต มยํ อิเมหิ ปญฺจหิ วตฺถูหิ ชนํ สญฺญาเปสฺสามาติ.\csromanspacer{} สกฺกา โข อาวุโส อิเมหิ ปญฺจหิ วตฺถูหิ สมณสฺส โคตมสฺส สํฆเภโท กาตุํ จกฺกเภโท, ลูขปฺปสนฺนา หิ อาวุโส มนุสฺสาติ.}

\prose{อถ โข เทวทตฺโต สปริโส เยน ภควา เตนุปสงฺกมิ, อุปสงฺกมิตฺวา ภควนฺตํ อภิวาเทตฺวา เอกมนฺตํ นิสีทิ, เอกมนฺตํ นิสินฺโน โข เทวทตฺโต ภควนฺตํ เอตทโวจ “ภควา ภนฺเต อเนกปริยาเยน อปฺปิจฺฉสฺส สนฺตุฏฺฐสฺส สลฺเลขสฺส ธุตสฺส ปาสาทิกสฺส อปจยสฺส วีริยารมฺภสฺส วณฺณวาที, อิมานิ ภนฺเต ปญฺจ วตฺถูนิ อเนกปริยาเยน อปฺปิจฺฉตาย สนฺตุฏฺฐิยา สลฺเลขาย ธุตตาย ปาสาทิกตาย อปจยาย วีริยารมฺภาย สํวตฺตนฺติ, สาธุ ภนฺเต ภิกฺขู ยาวชีวํ อารญฺญิกา อสฺสุ, โย คามนฺตํ โอสเรยฺย, วชฺชํ นํ ผุเสยฺย.\csromanspacer{} ยาวชีวํ ปิณฺฑปาติกา อสฺสุ, โย นิมนฺตนํ สาทิเยยฺย, วชฺชํ นํ ผุเสยฺย.\csromanspacer{} ยาวชีวํ ปํสุกูลิกา อสฺสุ, โย คหปติจีวรํ สาทิเยยฺย, วชฺชํ นํ ผุเสยฺย.\csromanspacer{} ยาวชีวํ รุกฺขมูลิกา อสฺสุ, โย ฉนฺนํ อุปคจฺเฉยฺย, วชฺชํ นํ ผุเสยฺย.\csromanspacer{} ยาวชีวํ มจฺฉมํสํ น ขาเทยฺยุํ, โย มจฺฉมํสํ ขาเทยฺย, วชฺชํ นํ ผุเสยฺยา”ติ.\csromanspacer{} อลํ เทวทตฺต, โย อิจฺฉติ อารญฺญิโก โหตุ, โย อิจฺฉติ คามนฺเต วิหรตุ, โย อิจฺฉติ ปิณฺฑปาติโก โหตุ, โย อิจฺฉติ นิมนฺตนํ สาทิยตุ, โย อิจฺฉติ ปํสุกูลิโก โหตุ, โย อิจฺฉติ คหปติจีวรํ สาทิยตุ, อฏฺฐมาเส โข มยา เทวทตฺต รุกฺขมูลเสนาสนํ อนุญฺญาตํ, ติโกฏิปริสุทฺธํ มจฺฉมํสํ อทิฏฺฐํ อสฺสุตํ อปริสงฺกิตนฺติ.\csromanspacer{} อถ โข เทวทตฺโต “น ภควา อิมานิ ปญฺจ วตฺถูนิ อนุชานาตี”ติ หฏฺโฐ อุทคฺโค สปริโส อุฏฺฐายาสนา ภควนฺตํ อภิวาเทตฺวา ปทกฺขิณํ กตฺวา ปกฺกามิ.}

\prose{อถ โข เทวทตฺโต สปริโส ราชคหํ ปวิสิตฺวา ปญฺจหิ วตฺถูหิ ชนํ สญฺญาเปสิ “มยํ อาวุโส สมณํ โคตมํ อุปสงฺกมิตฺวา ปญฺจ \csromanfolio{360}วตฺถูนิ ยาจิมฺหา ‘ภควา ภนฺเต อเนกปริยาเยน อปฺปิจฺฉสฺส ฯเปฯ วีริยารมฺภสฺส วณฺณวาที, อิมานิ ภนฺเต ปญฺจ วตฺถูนิ อเนกปริยาเยน อปฺปิจฺฉตาย ฯเปฯ วีริยารมฺภาย สํวตฺตนฺติ, สาธุ ภนฺเต ภิกฺขู ยาวชีวํ อารญฺญิกา อสฺสุ, โย คามนฺตํ โอสเรยฺย, วชฺชํ นํ ผุเสยฺย ฯเปฯ.\csromanspacer{} ยาวชีวํ มจฺฉมํสํ น ขาเทยฺยุํ, โย มจฺฉมํสํ ขาเทยฺย, วชฺชํ นํ ผุเสยฺยา’ติ.\csromanspacer{} อิมานิ ปญฺจ วตฺถูนิ สมโณ โคตโม นานุชานาติ, เต มยํ อิเมหิ ปญฺจหิ วตฺถูหิ สมาทาย วตฺตามา”ติ.}

\prose{ตตฺถ เย เต มนุสฺสา อสฺสทฺธา อปฺปสนฺนา ทุพฺพุทฺธิโน, เต เอวมาหํสุ “อิเม โข สมณา สกฺยปุตฺติยา ธุตา สลฺเลขวุตฺติโน, สมโณ ปน โคตโม พาหุลฺลิโก พาหุลฺลาย เจเตตี”ติ.\csromanspacer{} เย ปน เต มนุสฺสา สทฺธา ปสนฺนา ปณฺฑิตา พฺยตฺตา พุทฺธิมนฺโต, เต อุชฺฌายนฺติ ขิยฺยนฺติ วิปาเจนฺติ “กถํ หิ นาม เทวทตฺโต ภควโต สํฆเภทาย ปรกฺกมิสฺสติ จกฺกเภทายา”ติ.\csromanspacer{} อสฺโสสุํ โข ภิกฺขู เตสํ มนุสฺสานํ อุชฺฌายนฺตานํ ขิยฺยนฺตานํ วิปาเจนฺตานํ.\csromanspacer{} เย เต ภิกฺขู อปฺปิจฺฉา ฯเปฯ เต อุชฺฌายนฺติ ขิยฺยนฺติ วิปาเจนฺติ “กถํ หิ นาม เทวทตฺโต สํฆเภทาย ปรกฺกมิสฺสติ จกฺกเภทายา”ติ.\csromanspacer{} อถ โข เต ภิกฺขู ภควโต เอตมตฺถํ อาโรเจสุํ ฯเปฯ.\csromanspacer{} สจฺจํ กิร ตฺวํ เทวทตฺต สํฆเภทาย ปรกฺกมสิ จกฺกเภทายาติ.\csromanspacer{} สจฺจํ ภควาติ.\csromanspacer{} อลํ เทวทตฺต, มา เต รุจฺจิ สํฆเภโท, ครุโก โข เทวทตฺต สํฆเภโท, โย โข เทวทตฺต สมคฺคํ สํฆํ ภินฺทติ, กปฺปฏฺฐิกํ\footnote{กปฺปฏฺฐิติกํ (สฺยา)} กิพฺพิสํ ปสวติ, กปฺปํ นิรยมฺหิ ปจฺจติ.\csromanspacer{} โย จ โข เทวทตฺต ภินฺนํ สํฆํ สมคฺคํ กโรติ, พฺรหฺมํ ปุญฺญํ ปสวติ, กปฺปํ สคฺคมฺหิ โมทติ, อลํ เทวทตฺต, มา เต รุจฺจิ สํฆเภโท, ครุโก โข เทวทตฺต สํฆเภโทติ.}

\prose{\csromansymbolfootnote{*}{ขุ 1. 149 ปิฏฺเฐ อุทาเนปิ.}อถ โข อายสฺมา อานนฺโท ปุพฺพณฺหสมยํ นิวาเสตฺวา ปตฺตจีวรมาทาย ราชคหํ ปิณฺฑาย ปาวิสิ.\csromanspacer{} อทฺทสา โข เทวทตฺโต อายสฺมนฺตํ อานนฺทํ ราชคเห ปิณฺฑาย จรนฺตํ, ทิสฺวาน เยนายสฺมา อานนฺโท เตนุปสงฺกมิ, อุปสงฺกมิตฺวา อายสฺมนฺตํ อานนฺทํ เอตทโวจ “อชฺชตคฺเคทานาหํ อาวุโส อานนฺท อญฺญเตฺรว ภควตา, อญฺญเตฺรว ภิกฺขุสํฆา อุโปสถํ กริสฺสามิ, สํฆกมฺมํ กริสฺสามี”ติ.}

\csromanmatikaanchor{mtk.172}
\csromanmatikaanchor{mtk.173}
\csromantocmarkref{section}{3. ตติยภาณวาร}{mtk.172}
\bookmark[dest={mtk.172},level=1]{3. ตติยภาณวาร}
\csromantocmarkref{subsection}{สํฆเภทกถา}{mtk.173}
\bookmark[dest={mtk.173},level=2]{สํฆเภทกถา}
\prose{\csromanfolio{361}อถ โข อายสฺมา อานนฺโท ราชคเห ปิณฺฑาย จริตฺวา ปจฺฉาภตฺตํ ปิณฺฑปาตปฺปฏิกฺกนฺโต เยน ภควา เตนุปสงฺกมิ, อุปสงฺกมิตฺวา ภควนฺตํ อภิวาเทตฺวา เอกมนฺตํ นิสีทิ, เอกมนฺตํ นิสินฺโน โข อายสฺมา อานนฺโท ภควนฺตํ เอตทโวจ “อิธาหํ ภนฺเต ปุพฺพณฺหสมยํ นิวาเสตฺวา ปตฺตจีวรมาทาย ราชคหํ ปิณฺฑาย ปาวิสิํ, อทฺทสา โข มํ ภนฺเต เทวทตฺโต ราชคเห ปิณฺฑาย จรนฺตํ, ทิสฺวาน เยนาหํ เตนุปสงฺกมิ, อุปสงฺกมิตฺวา มํ เอตทโวจ ‘อชฺชตคฺเคทานาหํ อาวุโส อานนฺท อญฺญเตฺรว ภควตา, อญฺญเตฺรว ภิกฺขุสํฆา อุโปสถํ กริสฺสามิ, สํฆกมฺมํ กริสฺสามี”ติ, อชฺชตคฺเค ภนฺเต เทวทตฺโต สํฆํ ภินฺทิสฺสตี”ติ.\csromanspacer{} อถ โข ภควา เอตมตฺถํ วิทิตฺวา ตายํ เวลายํ อิมํ อุทานํ อุทาเนสิ\csromandash{}}

\csromangathameasure{*“สุกรํ สาธุนา สาธุํ, สาธุํ ปาเปน ทุกฺกรํ. \\ ปาปํ ปาเปน สุกรํ, ปาปมริเยหิ ทุกฺกรนฺ”ติ.}
\csromangathagroup{\csromangathastanza{\csromansymbolfootnote{*}{ขุ 1. 150 ปิฏฺเฐ อุทาเนปิ.}“สุกรํ สาธุนา สาธุํ, สาธุํ ปาเปน ทุกฺกรํ. \\ ปาปํ ปาเปน สุกรํ, ปาปมริเยหิ ทุกฺกรนฺ”ติ.}}

\nitthitam{ทุติยภาณวาโร นิฏฺฐิโต.}
\csromanheader{1}{3. ตติยภาณวาร}
\csromanheader{2}{\textbf{สํฆเภทถา}}
\proseitem{344}{อถ โข เทวทตฺโต ตทหุโปสเถ อุฏฺฐายาสนา สลากํ คาเหสิ “มยํ อาวุโส สมณํ โคตมํ อุปสงฺกมิตฺวา ปญฺจ วตฺถูนิ ยาจิมฺหา ‘ภควา ภนฺเต อเนกปริยาเยน อปฺปิจฺฉสฺส ฯเปฯ วีริยารมฺภสฺส วณฺณวาที, อิมานิ ภนฺเต ปญฺจ วตฺถูนิ อเนกปริยาเยน อปฺปิจฺฉตาย ฯเปฯ วีริยารมฺภาย สํวตฺตนฺติ, สาธุ ภนฺเต ภิกฺขู ยาวชีวํ อารญฺญิกา อสฺสุ, โย คามนฺตํ โอสเรยฺย, วชฺชํ นํ ผุเสยฺย ฯเปฯ.\csromanspacer{} ยาวชีวํ มจฺฉมํสํ น ขาเทยฺยุํ, โย มจฺฉมํสํ ขาเทยฺย, วชฺชํ นํ ผุเสยฺยา’ติ.\csromanspacer{} อิมานิ ปญฺจ วตฺถูนิ สมโณ โคตโม นานุชานาติ, เต มยํ อิเมหิ ปญฺจหิ วตฺถูหิ สมาทาย วตฺตาม, ยสฺสายสฺมโต อิมานิ ปญฺจ วตฺถูนิ ขมนฺติ, โส สลากํ คณฺหาตู”ติ.}

\prose{เตน โข ปน สมเยน เวสาลิกา วชฺชิปุตฺตกา ปญฺจมตฺตานิ ภิกฺขุสตานิ นวกา เจว โหนฺติ อปฺปกตญฺญุโน จ, เต “อยํ \csromanfolio{362}ธมฺโม, อยํ วินโย, อิทํ สตฺถุสาสนนฺ”ติ สลากํ คณฺหิํสุ.\csromanspacer{} อถ โข เทวทตฺโต สํฆํ ภินฺทิตฺวา ปญฺจมตฺตานิ ภิกฺขุสตานิ อาทาย เยน คยาสีสํ เตน ปกฺกามิ.\csromanspacer{} อถ โข สาริปุตฺตโมคฺคลฺลานา เยน ภควา เตนุปสงฺกมิํสุ, อุปสงฺกมิตฺวา ภควนฺตํ อภิวาเทตฺวา เอกมนฺตํ นิสีทิํสุ.\csromanspacer{} เอกมนฺตํ นิสินฺโน โข อายสฺมา สาริปุตฺโต ภควนฺตํ เอตทโวจ “เทวทตฺโต ภนฺเต สํฆํ ภินฺทิตฺวา ปญฺจมตฺตานิ ภิกฺขุสตานิ อาทาย เยน คยาสีสํ, เตน ปกฺกนฺโต”ติ.\csromanspacer{} น หิ นาม ตุมฺหากํ สาริปุตฺตา เตสุ นวเกสุ ภิกฺขูสุ การุญฺญมฺปิ ภวิสฺสติ, คจฺฉถ ตุเมฺห สาริปุตฺตา, ปุรา เต ภิกฺขู อนยพฺยสนํ อาปชฺชนฺตีติ. “เอวํ ภนฺเต”ติ โข สาริปุตฺตโมคฺคลฺลานา ภควโต ปฏิสฺสุตฺวา อุฏฺฐายาสนา ภควนฺตํ อภิวาเทตฺวา ปทกฺขิณํ กตฺวา เยน คยาสีสํ เตนุปสงฺกมิํสุ.}

\prose{เตน โข ปน สมเยน อญฺญตโร ภิกฺขุ ภควโต อวิทูเร โรทมาโน ฐิโต โหติ.\csromanspacer{} อถ โข ภควา ตํ ภิกฺขุํ เอตทโวจ “กิสฺส ตฺวํ ภิกฺขุ โรทสี”ติ.\csromanspacer{} เยปิ เต ภนฺเต ภควโต อคฺคสาวกา สาริปุตฺตโมคลฺลานา, เตปิ เทวทตฺตสฺส สนฺติเก คจฺฉนฺติ เทวทตฺตสฺส ธมฺมํ โรเจนฺตาติ.\csromanspacer{} อฏฺฐานเมตํ ภิกฺขุ อนวกาโส, ยํ สาริปุตฺตโมคฺคลฺลานา เทวทตฺตสฺส ธมฺมํ โรเจยฺยุํ.\csromanspacer{} อปิ จ เต คตา ภิกฺขูนํ สญฺญตฺติยาติ\footnote{ภิกฺขุสญฺญตฺติยาติ (สี, สฺยา), ภิกฺขู สญฺญตฺติยา (ก)}.}

\proseitem{345}{เตน โข ปน สมเยน เทวทตฺโต มหติยา ปริสาย ปริวุโต ธมฺมํ เทเสนฺโต นิสินฺโน โหติ.\csromanspacer{} อทฺทสา โข เทวทตฺโต สาริปุตฺตโมคฺคลฺลาเน ทูรโตว อาคจฺฉนฺเต, ทิสฺวาน ภิกฺขู อามนฺเตสิ “ปสฺสถ ภิกฺขเว, ยาว สฺวากฺขาโต มยํ ธมฺโม, เยปิ เต สมณสฺส โคตมสฺส อคฺคสาวกา สาริปุตฺตโมคฺคลฺลานา, เตปิ มม สนฺติเก อาคจฺฉนฺติ มม ธมฺมํ โรเจนฺตา”ติ.\csromanspacer{} เอวํ วุตฺเต โกกาลิโก เทวทตฺตํ เอตทโวจ “มา อาวุโส เทวทตฺต สาริปุตฺตโมคฺคลฺลาเน วิสฺสาสิ, ปาปิจฺฉา สาริปุตฺตโมคฺคลฺลานา ปาปิกานํ อิจฺฉานํ วสํ คตา”ติ.\csromanspacer{} อลํ อาวุโส, สฺวาคตํ เตสํ, ยโต เม ธมฺมํ โรเจนฺตีติ.}

\prose{อถ โข เทวทตฺโต อายสฺมนฺตํ สาริปุตฺตํ อุปฑฺฒาสเนน นิมนฺเตสิ “เอหาวุโส สาริปุตฺต, อิธ นิสีทาหี”ติ. “อลํ อาวุโส”ติ โข \csromanfolio{363}อายสฺมา สาริปุตฺโต อญฺญตรํ อาสนํ คเหตฺวา เอกมนฺตํ นิสีทิ อายสฺมาปิ โข มหาโมคฺคลฺลาโน อญฺญตรํ อาสนํ คเหตฺวา เอกมนฺตํ นิสีทิ.\csromanspacer{} อถ โข เทวทตฺโต พหุเทว รตฺติํ ภิกฺขู ธมฺมิยา กถาย สนฺทสฺเสตฺวา สมาทเปตฺวา สมุตฺเตเชตฺวา สมฺปหํเสตฺวา อายสฺมนฺตํ สาริปุตฺตํ อชฺเฌสิ “วิคตถินมิทฺโธ โข อาวุโส สาริปุตฺต ภิกฺขุสํโฆ, ปฏิภาตุ ตํ อาวุโส สาริปุตฺต ภิกฺขูนํ ธมฺมี กถา, ปิฏฺฐิ เม อาคิลายติ, ตมหํ อายมิสฺสามี”ติ. “เอวมาวุโส”ติ โข อายสฺมา สาริปุตฺโต เทวทตฺตสฺส ปจฺจสฺโสสิ.\csromanspacer{} อถ โข เทวทตฺโต จตุคฺคุณํ สํฆาฏิํ ปญฺญเปตฺวา ทกฺขิเณน ปสฺเสน เสยฺยํ กปฺเปสิ.\csromanspacer{} ตสฺส กิลมนฺตสฺส มุฏฺฐสฺสติสฺส อสมฺปชานสฺส มุหุตฺตเกเนว นิทฺทา โอกฺกมิ.}

\prose{อถ โข อายสฺมา สาริปุตฺโต อาเทสนาปาฏิหาริยานุสาสนิยา ภิกฺขู ธมฺมิยา กถาย โอวทิ อนุสาสิ, อายสฺมา มหาโมคฺคลฺลาโน อิทฺธิปาฏิหาริยานุสาสนิยา ภิกฺขู ธมฺมิยา กถาย โอวทิ อนุสาสิ.\csromanspacer{} อถ โข เตสํ ภิกฺขูนํ อายสฺมตา สาริปุตฺเตน อาเทสนาปาฏิหาริยานุสาสนิยา, อายสฺมตา จ มหาโมคฺคลฺลาเนน อิทฺธิปาฏิหาริยานุสาสนิยา โอวทิยมานานํ อนุสาสิยมานานํ วิรชํ วีตมลํ ธมฺมจกฺขุํ อุทปาทิ “ยํ กิญฺจิ สมุทยธมฺมํ, สพฺพํ ตํ นิโรธธมฺมนฺ”ติ.}

\prose{อถ โข อายสฺมา สาริปุตฺโต ภิกฺขู อามนฺเตสิ “คจฺฉาม มยํ อาวุโส ภควโต สนฺติเก, โย ตสฺส ภควโต ธมฺมํ โรเจสิ, โส อาคจฺฉตู”ติ.\csromanspacer{} อถ โข สาริปุตฺตโมคฺคลฺลานา ตานิ ปญฺจ ภิกฺขุสตานิ อาทาย เยน เวฬุวนํ เตนุปสงฺกมิํสุ.\csromanspacer{} อถ โข โกกาลิโก เทวทตฺตํ อุฏฺฐาเปสิ, อุฏฺเฐหิ อาวุโส เทวทตฺต, นีตา เต ภิกฺขู สาริปุตฺตโมคฺคลฺลาเนหิ, นนุ ตฺวํ อาวุโส เทวทตฺต มยา วุตฺโต “มา อาวุโส เทวทตฺต สาริปุตฺตโมคฺคลฺลาเน วิสฺสาสิ, ปาปิจฺฉา สาริปุตฺตโมคฺคลฺลานา ปาปิกานํ อิจฺฉานํ วสํ คตา”ติ.\csromanspacer{} อถ โข เทวทตฺตสฺส ตตฺเถว อุณฺหํ โลหิตํ มุขโต อุคฺคญฺฉิ. อถ โข สาริปุตฺตโมคฺคลฺลานา เยน ภควา เตนุปสงฺกมิํสุ, อุปสงฺกมิตฺวา ภควนฺตํ อภิวาเทตฺวา เอกมนฺตํ นิสีทิํสุ.\csromanspacer{} เอกมนฺตํ นิสินฺโน โข อายสฺมา สาริปุตฺโต ภควนฺตํ เอตทโวจ “สาธุ \csromanfolio{364}ภนฺเต เภทกานุวตฺตกา ภิกฺขู ปุน อุปสมฺปชฺเชยฺยุนฺ”ติ.\csromanspacer{} อลํ สาริปุตฺต, มา เต รุจฺจิ เภทกานุวตฺตกานํ ภิกฺขูนํ ปุน อุปสมฺปทา.\csromanspacer{} เตน หิ ตฺวํ สาริปุตฺต เภทกานุวตฺตเก ภิกฺขู ถุลฺลจฺจยํ เทสาเปหิ.\csromanspacer{} กถํ ปน เต สาริปุตฺต เทวทตฺโต ปฏิปชฺชีติ.\csromanspacer{} ยเถว ภนฺเต ภควา พหุเทว รตฺติํ ภิกฺขู ธมฺมิยา กถาย สนฺทสฺเสตฺวา สมาทเปตฺวา สมุตฺเตเชตฺวา สมฺปหํเสตฺวา มํ อชฺเฌสติ “วิคตถินมิทฺโธ โข สาริปุตฺต ภิกฺขุสํโฆ, ปฏิภาตุ ตํ สาริปุตฺต ภิกฺขูนํ ธมฺมี กถา, ปิฏฺฐิ เม อาคิลายติ, ตมหํ อายมิสฺสามี”ติ.\csromanspacer{} เอวเมว โข ภนฺเต เทวทตฺโต ปฏิปชฺชีติ.}

\proseitem{346}{อถ โข ภควา ภิกฺขู อามนฺเตสิ “ภูตปุพฺพํ ภิกฺขเว อรญฺญายตเน มหาสรสี, ตํ นาคา อุปนิสฺสาย วิหริํสุ, เต ตํ สรสิํ โอคาเหตฺวา โสณฺฑาย ภิสมุฬาลํ อพฺพุหิตฺวา สุวิกฺขาลิตํ วิกฺขาเลตฺวา อกทฺทมํ สํขาทิตฺวา อชฺโฌหรนฺติ.\csromanspacer{} เตสํ ตํ วณฺณาย เจว โหติ พลาย จ, น จ ตโตนิทานํ มรณํ วา นิคจฺฉนฺติ, มรณมตฺตํ วา ทุกฺขํ.\csromanspacer{} เตสํเยว โข ปน ภิกฺขเว มหานาคานํ อนุสิกฺขมานา ตรุณา ภิงฺกจฺฉาปา เต ตํ สรสิํ โอคาเหตฺวา โสณฺฑาย ภิสมุฬาลํ อพฺพุหิตฺวา น สุวิกฺขาลิตํ วิกฺขาลิตฺวา สกทฺทมํ สํขาทิตฺวา อชฺโฌหรนฺติ.\csromanspacer{} เตสํ ตํ เนว วณฺณาย โหติ, น พลาย, ตโตนิทานญฺจ มรณํ วา นิคจฺฉนฺติ, มรณมตฺตํ วา ทุกฺขํ.\csromanspacer{} เอวเมว โข ภิกฺขเว เทวทตฺโต มมานุกฺรุพฺพํ\footnote{มมานุกุพฺพํ (สี, สฺยา)} กปโณ มริสฺสตี”ติ.}

\csromangathameasure{มหาวราหสฺส มหิํ วิกฺรุพฺพโต, \\ ภิสํ ฆสานสฺส นทีสุ ชคฺคโต. \\ ภิงฺโกว ปงฺกํ อภิภกฺขยิตฺวา, \\ มมานุกฺรุพฺพํ กปโณ มริสฺสตีติ.}
\csromangathagroup{\csromangathastanza{มหาวราหสฺส มหิํ วิกฺรุพฺพโต\footnote{วิกุพฺพโต (สี, สฺยา)}, \\ ภิสํ ฆสานสฺส\footnote{วิกุพฺพโต (สี, สฺยา)} นทีสุ ชคฺคโต. \\ ภิงฺโกว ปงฺกํ อภิภกฺขยิตฺวา, \\ มมานุกฺรุพฺพํ กปโณ มริสฺสตีติ.}}

\proseitem{347}{\csromansymbolfootnote{*}{อํ 3. 37 ปิฏฺเฐปิ.}อฏฺฐหิ ภิกฺขเว องฺเคหิ สมนฺนาคโต ภิกฺขุ ทูเตยฺยํ คนฺตุมรหติ.\csromanspacer{} กตเมหิ อฏฺฐหิ.\csromanspacer{} อิธ ภิกฺขเว ภิกฺขุ โสตา จ โหติ, สาเวตา จ, อุคฺคเหตา จ, ธาเรตา จ วิญฺญาตา จ, วิญฺญาเปตา จ, กุสโล จ สหิตาสหิตสฺส, โน จ กลหการโก.\csromanspacer{} อิเมหิ โข ภิกฺขเว อฏฺฐหงฺเคหิ สมนฺนาคโต ภิกฺขุ ทูเตยฺยํ คนฺตุมรหติ.}

\prose{\csromanfolio{365}\csromansymbolfootnote{*}{อํ 3. 37 ปิฏฺเฐปิ. 1. พฺยาธติ (สี, สฺยา) 2. อขาตา (ก) + อํ 3. 10 ปิฏฺเฐปิ.}อฏฺฐหิ ภิกฺขเว องฺเคหิ สมนฺนาคโต สาริปุตฺโต ทูเตยฺยํ คนฺตุมรหติ.\csromanspacer{} กตเมหิ อฏฺฐหิ.\csromanspacer{} อิธ ภิกฺขเว สาริปุตฺโต โสตา จ โหติ, สาเวตา จ, อุคฺคเหตา จ, ธาเรตา จ, วิญฺญาตา จ, วิญฺญาเปตา จ, กุสโล จ สหิตาสหิตสฺส, โน จ กลหการโก.\csromanspacer{} อิเมหิ โข ภิกฺขเว อฏฺฐหงฺเคหิ สมนฺนาคโต สาริปุตฺโต ทูเตยฺยํ คนฺตุมรหตีติ.}

\csromangathameasure{*โย เว น พฺยถติ1 ปตฺวา, ปริสํ อุคฺควาทินิํ. \\ น จ หาเปติ วจนํ, น จ ฉาเทติ สาสนํ. \\ อสนฺทิทฺโธ จ อกฺขาติ2, ปุจฺฉิโต จ น กุปฺปติ. \\ ส เว ตาทิสโก ภิกฺขุ, ทูเตยฺยํ คนฺตุมรหตีติ.}
\csromangathagroup{\csromangathastanza{\csromansymbolmark{*}โย เว น พฺยถติ1 ปตฺวา, ปริสํ อุคฺควาทินิํ. \\ น จ หาเปติ วจนํ, น จ ฉาเทติ สาสนํ.}\csromangathastanza{อสนฺทิทฺโธ จ อกฺขาติ2, ปุจฺฉิโต จ น กุปฺปติ. \\ ส เว ตาทิสโก ภิกฺขุ, ทูเตยฺยํ คนฺตุมรหตีติ.}}

\proseitem{348}{\csromansymbolmark{+}อฏฺฐหิ ภิกฺขเว อสทฺธมฺเมหิ อภิภูโต ปริยาทินฺนจิตฺโต เทวทตฺโต อาปายิโก เนรยิโก กปฺปฏฺโฐ อเตกิจฺโฉ.\csromanspacer{} กตเมหิ อฏฺฐหิ.\csromanspacer{} ลาเภน ภิกฺขเว อภิภูโต ปริยาทินฺนจิตฺโต เทวทตฺโต อาปายิโก เนรยิโก กปฺปฏฺโฐ อเตกิจฺโฉ.\csromanspacer{} อลาเภน ภิกฺขเว ฯเปฯ.\csromanspacer{} ยเสน ภิกฺขเว ฯเปฯ.\csromanspacer{} อยเสน ภิกฺขเว ฯเปฯ.\csromanspacer{} สกฺกาเรน ภิกฺขเว ฯเปฯ.\csromanspacer{} อสกฺกาเรน ภิกฺขเว ฯเปฯ.\csromanspacer{} ปาปิจฺฉตาย ภิกฺขเว ฯเปฯ.\csromanspacer{} ปาปมิตฺตตาย ภิกฺขเว อภิภูโต ปริยาทินฺนจิตฺโต เทวทตฺโต อาปายิโก เนรยิโก กปฺปฏฺโฐ อเตกิจฺโฉ.\csromanspacer{} อิเมหิ โข ภิกฺขเว อฏฺฐหิ อสทฺธมฺเมหิ อภิภูโต ปริยาทินฺนจิตฺโต เทวทตฺโต อาปายิโก เนรยิโก กปฺปฏฺโฐ อเตกิจฺโฉ.}

\proseitem{349}{สาธุ ภิกฺขเว ภิกฺขุ อุปฺปนฺนํ ลาภํ อภิภุยฺย อภิภุยฺย วิหเรยฺย.\csromanspacer{} อุปฺปนฺนํ อลาภํ ฯเปฯ.\csromanspacer{} อุปฺปนฺนํ ยสํ.\csromanspacer{} อุปฺปนฺนํ อยสํ.\csromanspacer{} อุปฺปนฺนํ สกฺการํ.\csromanspacer{} อุปฺปนฺนํ อสกฺการํ.\csromanspacer{} อุปฺปนฺนํ ปาปิจฺฉตํ.\csromanspacer{} อุปฺปนฺนํ ปาปมิตฺตตํ อภิภุยฺย อภิภุยฺย วิหเรยฺย.\csromanspacer{} กถญฺจ ภิกฺขเว ภิกฺขุ อตฺถวสํ ปฏิจฺจ อุปฺปนฺนํ ลาภํ อภิภุยฺย อภิภุยฺย วิหเรยฺย.\csromanspacer{} อุปฺปนฺนํ อลาภํ ฯเปฯ.\csromanspacer{} อุปฺปนฺนํ ยสํ.\csromanspacer{} อุปฺปนฺนํ อยสํ.\csromanspacer{} อุปฺปนฺนํ สกฺการํ.\csromanspacer{} อุปฺปนฺนํ อสกฺการํ.\csromanspacer{} อุปฺปนฺนํ ปาปิจฺฉตํ.\csromanspacer{} อุปฺปนฺนํ ปาปมิตฺตตํ อภิภุยฺย อภิภุยฺย วิหเรยฺย.\csromanspacer{} ยํ หิสฺส ภิกฺขเว อุปฺปนฺนํ ลาภํ อนภิภุยฺย วิหรโต อุปฺปชฺเชยฺยุํ อาสวา วิฆาตปริฬาหา, อุปฺปนฺนํ ลาภํ อภิภุยฺย \csromanfolio{366}อภิภุยฺย วิหรโต เอวํสเต อาสวา วิฆาตปริฬาหา น โหนฺติ.\csromanspacer{} ยํ หิสฺส ภิกฺขเว อุปฺปนฺนํ อลาภํ ฯเปฯ อุปฺปนฺนํ ยสํ.\csromanspacer{} อุปฺปนฺนํ อยสํ.\csromanspacer{} อุปฺปนฺนํ สกฺการํ.\csromanspacer{} อุปฺปนฺนํ อสกฺการํ.\csromanspacer{} อุปฺปนฺนํ ปาปิจฺฉตํ.\csromanspacer{} อุปฺปนฺนํ ปาปมิตฺตตํ อนภิภุยฺย วิหรโต อุปฺปชฺเชยฺยุํ อาสวา วิฆาตปริฬาหา, อุปฺปนฺนํ ปาปมิตฺตตํ อภิภุยฺย อภิภุยฺย วิหรโต เอวํสเต อาสวา วิฆาตปริฬาหา น โหนฺติ.\csromanspacer{} อิทํ โข ภิกฺขเว อตฺถวสํ ปฏิจฺจ อุปฺปนฺนํ ลาภํ อภิภุยฺย อภิภุยฺย วิหเรยฺย.\csromanspacer{} อุปฺปนฺนํ อลาภํ ฯเปฯ อุปฺปนฺนํ ยสํ.\csromanspacer{} อุปฺปนฺนํ อยสํ.\csromanspacer{} อุปฺปนฺนํ สกฺการํ.\csromanspacer{} อุปฺปนฺนํ อสกฺการํ.\csromanspacer{} อุปฺปนฺนํ ปาปิจฺฉตํ.\csromanspacer{} อุปฺปนฺนํ ปาปมิตฺตตํ อภิภุยฺย อภิภุยฺย วิหเรยฺย.\csromanspacer{} ตสฺมา ติห ภิกฺขเว อุปฺปนฺนํ ลาภํ อภิภุยฺย อภิภุยฺย วิหริสฺสาม.\csromanspacer{} อุปฺปนฺนํ อลาภํ ฯเปฯ อุปฺปนฺนํ ยสํ.\csromanspacer{} อุปฺปนฺนํ อยสํ.\csromanspacer{} อุปฺปนฺนํ สกฺการํ.\csromanspacer{} อุปฺปนฺนํ อสกฺการํ.\csromanspacer{} อุปฺปนฺนํ ปาปิจฺฉตํ.\csromanspacer{} อุปฺปนฺนํ ปาปมิตฺตตํ อภิภุยฺย อภิภุยฺย วิหริสฺสามาติ.\csromanspacer{} เอวญฺหิ โว ภิกฺขเว สิกฺขิตพฺพนฺติ.}

\proseitem{350}{ตีหิ ภิกฺขเว อสทฺธมฺเมหิ อภิภูโต ปริยาทินฺนจิตฺโต เทวทตฺโต อาปายิโก เนรยิโก กปฺปฏฺโฐ อเตกิจฺโฉ.\csromanspacer{} กตเมหิ ตีหิ.\csromanspacer{} ปาปิจฺฉตา, ปาปมิตฺตตา, โอรมตฺตเกน วิเสสาธิคเมน อนฺตรา โวสานํ อาปาทิ.\csromanspacer{} อิเมหิ โข ภิกฺขเว ตีหิ อสทฺธมฺเมหิ อภิภูโต ปริยาทินฺนจิตฺโต เทวทตฺโต อาปายิโก เนรยิโก กปฺปฏฺโฐ อเตกิจฺโฉติ.}

\csromangathameasure{มา ชาตุ โกจิ โลกสฺมิํ, ปาปิจฺโฉ อุทปชฺชถ. \\ ตทมินาปิ ชานาถ, ปาปิจฺฉานํ ยถาคติ. \\ ปณฺฑิโตติ สมญฺญาโต, ภาวิตตฺโตติ สมฺมโต. \\ ชลํว ยสสา อฏฺฐา, เทวทตฺโตติ เม สุตํ. \\ โส ปมาทํ อนุจิณฺโณ, อาสชฺช นํ ตถาคตํ. \\ อวีจินิรยํ ปตฺโต, จตุทฺวารํ ภยานกํ. \\ อทุฏฺฐสฺส หิ โย ทุพฺเภ, ปาปกมฺมํ อกฺรุพฺพโต. \\ ตเมว ปาปํ ผุสติ, ทุฏฺฐจิตฺตํ อนาทรํ. \\ สมุทฺทํ วิสกุมฺเภน, โย มญฺเญยฺย ปทูสิตุํ. \\ น โส เตน ปทูเสยฺย, เภสฺมา หิ อุทธี มหา. \\ เอวเมว ตถาคตํ, โย วาเทนุปหิํสติ. \\ สมคฺคตํ สนฺตจิตฺตํ, วาโท ตมฺหิ น รูหติ. \\ ตาทิสํ มิตฺตํ กฺรุพฺเพถ, ตญฺจ เสเวถ ปณฺฑิโต. \\ ยสฺส มคฺคานุโค ภิกฺขุ, ขยํ ทุกฺขสฺส ปาปุเณติ.}
\csromangathagroup{\csromangathastanza{มา ชาตุ โกจิ โลกสฺมิํ, ปาปิจฺโฉ อุทปชฺชถ. \\ ตทมินาปิ ชานาถ, ปาปิจฺฉานํ ยถาคติ.}\csromangathastanza{ปณฺฑิโตติ สมญฺญาโต, ภาวิตตฺโตติ สมฺมโต. \\ ชลํว ยสสา อฏฺฐา, เทวทตฺโตติ เม สุตํ.}\csromangathastanza{โส ปมาทํ อนุจิณฺโณ, อาสชฺช นํ ตถาคตํ. \\ อวีจินิรยํ ปตฺโต, จตุทฺวารํ ภยานกํ.}\csromangathastanza{อทุฏฺฐสฺส หิ โย ทุพฺเภ, ปาปกมฺมํ อกฺรุพฺพโต. \\ ตเมว ปาปํ ผุสติ, ทุฏฺฐจิตฺตํ อนาทรํ.}\csromangathastanza{สมุทฺทํ วิสกุมฺเภน, โย มญฺเญยฺย ปทูสิตุํ\footnote{ปทุสฺสิตุํ (ก)}. \\ น โส เตน ปทูเสยฺย, เภสฺมา หิ อุทธี มหา.}\csromangathastanza{\csromanfolio{367}เอวเมว ตถาคตํ, โย วาเทนุปหิํสติ. \\ สมคฺคตํ\footnote{สมฺมาคตํ (สี), สมคตํ (สฺยา)} สนฺตจิตฺตํ, วาโท ตมฺหิ น รูหติ.}\csromangathastanza{ตาทิสํ มิตฺตํ กฺรุพฺเพถ\footnote{กุพฺเพถ (สี, สฺยา)}, ตญฺจ เสเวถ ปณฺฑิโต. \\ ยสฺส มคฺคานุโค ภิกฺขุ, ขยํ ทุกฺขสฺส ปาปุเณติ.}}

\csromanmatikaanchor{mtk.174}
\csromantocmarkref{subsection}{อุปาลิปญฺหา}{mtk.174}
\bookmark[dest={mtk.174},level=2]{อุปาลิปญฺหา}
\csromanheader{2}{\textbf{อุปาลิปญฺหา}}
\proseitem{351}{อถ โข อายสฺมา อุปาลิ เยน ภควา เตนุปสงฺกมิ, อุปสงฺกมิตฺวา ภควนฺตํ อภิวาเทตฺวา เอกมนฺตํ นิสีทิ, เอกมนฺตํ นิสินฺโน โข อายสฺมา อุปาลิ ภควนฺตํ เอตทโวจ “สํฆราชิ สํฆราชี”ติ ภนฺเต วุจฺจติ, กิตฺตาวตา นุ โข ภนฺเต สํฆราชิ โหติ, โน จ สํฆเภโท.\csromanspacer{} กิตฺตาวตา จ ปน สํฆราชิ เจว โหติ สํฆเภโท จา”ติ.}

\prose{เอกโต อุปาลิ เอโก โหติ, เอกโต เทฺว, จตุตฺโถ อนุสฺสาเวติ, สลากํ คาเหติ “อยํ ธมฺโม, อยํ วินโย, อิทํ สตฺถุสาสนํ, อิมํ คณฺหถ, อิมํ โรเจถา”ติ.\csromanspacer{} เอวมฺปิ โข อุปาลิ สํฆราชิ โหติ, โน จ สํฆเภโท.\csromanspacer{} เอกโต อุปาลิ เทฺว โหนฺติ, เอกโต เทฺว.\csromanspacer{} ปญฺจโม อนุสฺสาเวติ, สลากํ คาเหติ “อยํ ธมฺโม, อยํ วินโย, อิทํ สตฺถุสาสนํ, อิมํ คณฺหถ, อิมํ โรเจถา”ติ.\csromanspacer{} เอวมฺปิ โข อุปาลิ สํฆราชิ โหติ, โน จ สํฆเภโท.\csromanspacer{} เอกโต อุปาลิ เทฺว โหนฺติ.\csromanspacer{} เอกโต ตโย, ฉฏฺโฐ อนุสฺสาเวติ, สลากํ คาเหติ “อยํ ธมฺโม, อยํ วินโย, อิทํ สตฺถุสาสนํ, อิมํ คณฺหถ, อิมํ โรเจถา”ติ.\csromanspacer{} เอวมฺปิ โข อุปาลิ สํฆราชิ โหติ, โน จ สํฆเภโท.\csromanspacer{} เอกโต อุปาลิ ตโย โหนฺติ.\csromanspacer{} เอกโต ตโย, สตฺตโม อนุสฺสาเวติ, สลากํ คาเหติ “อยํ ธมฺโม, อยํ วินโย, อิทํ สตฺถุสาสนํ, อิมํ คณฺหถ, อิมํ โรเจถา”ติ.\csromanspacer{} เอวมฺปิ โข อุปาลิ สํฆราชิ โหติ, โน จ สํฆเภโท.\csromanspacer{} เอกโต อุปาลิ ตโย โหนฺติ.\csromanspacer{} เอกโต จตฺตาโร, อฏฺฐโม อนุสฺสาเวติ, สลากํ คาเหติ “อยํ ธมฺโม, อยํ วินโย, อิทํ สตฺถุสาสนํ, อิมํ คณฺหถ, อิมํ โรเจถา”ติ.\csromanspacer{} เอวมฺปิ โข อุปาลิ สํฆราชิ โหติ, โน จ สํฆเภโท.\csromanspacer{} เอกโต อุปาลิ จตฺตาโร โหนฺติ, เอกโต จตฺตาโร, นวโม อนุสฺสาเวติ, สลากํ คาเหติ “อยํ ธมฺโม, อยํ วินโย, อิทํ สตฺถุสาสนํ, อิมํ คณฺหถ, อิมํ โรเจถา”ติ. \csromanfolio{368}เอวํ โข อุปาลิ สํฆราชิ เจว โหติ สํฆเภโท จ.\csromanspacer{} นวนฺนํ วา อุปาลิ อติเรกนวนฺนํ วา สํฆราชิ เจว โหติ สํฆเภโท จ.\csromanspacer{} น โข อุปาลิ ภิกฺขุนี สํฆํ ภินฺทติ, อปิ จ เภทาย ปรกฺกมติ.\csromanspacer{} น สิกฺขมานา สํฆํ ภินฺทติ ฯเปฯ.\csromanspacer{} น สามเณโร สํฆํ ภินฺทติ.\csromanspacer{} น สามเณรี สํฆํ ภินฺทติ.\csromanspacer{} น อุปาสโก สํฆํ ภินฺทติ.\csromanspacer{} น อุปาสิกา สํฆํ ภินฺทติ, อปิ จ เภทาย ปรกฺกมติ.\csromanspacer{} ภิกฺขุ โข อุปาลิ ปกตตฺโต สมานสํวาสโก สมานสีมายํ ฐิโต สํฆํ ภินฺทตีติ.}

\proseitem{352}{\csromansymbolfootnote{*}{อํ 3. 314 ปิฏฺเฐทีสุปิ.}“สํฆเภโท สํฆเภโท”ติ ภนฺเต วุจฺจติ, กิตฺตาวตา นุ โข ภนฺเต สํโฆ ภินฺโน โหตีติ.}

\prose{อิธุปาลิ ภิกฺขู อธมฺมํ “ธมฺโม”ติ ทีเปนฺติ, ธมฺมํ “อธมฺโม”ติ ทีเปนฺติ, อวินยํ “วินโย”ติ ทีเปนฺติ, วินยํ “อวินโย”ติ ทีเปนฺติ, อภาสิตํ อลปิตํ ตถาคเตน “ภาสิตํ ลปิตํ ตถาคเตนา”ติ ทีเปนฺติ, ภาสิตํ ลปิตํ ตถาคเตน “อภาสิตํ อลปิตํ ตถาคเตนา”ติ ทีเปนฺติ, อนาจิณฺณํ ตถาคเตน “อาจิณฺณํ ตถาคเตนา”ติ ทีเปนฺติ, อาจิณฺณํ ตถาคเตน “อนาจิณฺณํ ตถาคเตนา”ติ ทีเปนฺติ, อปญฺญตฺตํ ตถาคเตน “ปญฺญตฺตํ ตถาคเตนา”ติ ทีเปนฺติ, ปญฺญตฺตํ ตถาคเตน “อปญฺญตฺตํ ตถาคเตนา”ติ ทีเปนฺติ, อนาปตฺติํ “อาปตฺตี”ติ ทีเปนฺติ, อาปตฺติํ “อานาปตฺตี”ติ ทีเปนฺติ, ลหุกํ อาปตฺติํ “ครุกา อาปตฺตี”ติ ทีเปนฺติ, ครุกํ อาปตฺติํ “ลหุกา อาปตฺตี”ติ ทีเปนฺติ, สาวเสสํ อาปตฺติํ “อนวเสสา อาปตฺตี”ติ ทีเปนฺติ, อนวเสสํ อาปตฺติํ “สาวเสสา อาปตฺตี”ติ ทีเปนฺติ, ทุฏฺฐุลฺลํ อาปตฺติํ “อทุฏฺฐุลฺลา อาปตฺตี”ติ ทีเปนฺติ, อทุฏฺฐุลฺลํ อาปตฺติํ “ทุฏฺฐุลฺลา อาปตฺตี”ติ ทีเปนฺติ.\csromanspacer{} เต อิเมหิ อฏฺฐารสหิ วตฺถูหิ อปกสฺสนฺติ, อวปกสฺสนฺติ, อาเวนิํ\footnote{อาเวณิ (สี), อาเวณิกํ (สฺยา)} อุโปสถํ กโรนฺติ, อาเวนิํ ปวารณํ กโรนฺติ, อาเวนิํ สํฆกมฺมํ กโรนฺติ.\csromanspacer{} เอตฺตาวตา โข อุปาลิ สํโฆ ภินฺโน โหตีติ.}

\proseitem{353}{\csromansymbolmark{*}“สํฆสามคฺคี สํฆสามคฺคี”ติ ภนฺเต วุจฺจติ, กิตฺตาวตา นุ โข ภนฺเต สํโฆ สมคฺโค โหตีติ.\csromanspacer{} อิธุปาลิ ภิกฺขู อธมฺมํ “อธมฺโม”ติ ทีเปนฺติ, ธมฺมํ “ธมฺโม”ติ ทีเปนฺติ, อวินยํ “อวินโย”ติ ทีเปนฺติ, วินยํ \csromanfolio{369}“วินโย”ติ ทีเปนฺติ, อภาสิตํ อลปิตํ ตถาคเตน “อภาสิตํ อลปิตํ ตถาคเตนา”ติ ทีเปนฺติ, ภาสิตํ ลปิตํ ตถาคเตน “ภาสิตํ ลปิตํ ตถาคเตนา”ติ ทีเปนฺติ, อนาจิณฺณํ ตถาคเตน “อนาจิณฺณํ ตถาคเตนา”ติ ทีเปนฺติ, อาจิณฺณํ ตถาคเตน “อาจิณฺณํ ตถาคเตนา”ติ ทีเปนฺติ, อปญฺญตฺตํ ตถาคเตน “อปญฺญตฺตํ ตถาคเตนา”ติ ทีเปนฺติ, ปญฺญตฺตํ ตถาคเตน “ปญฺญตฺตํ ตถาคเตนา”ติ ทีเปนฺติ, อนาปตฺติํ “อนาปตฺตี”ติ ทีเปนฺติ, อาปตฺติํ “อาปตฺตี”ติ ทีเปนฺติ, ลหุกํ อาปตฺติํ “ลหุกา อาปตฺตี”ติ ทีเปนฺติ, ครุกํ อาปตฺติํ “ครุกา อาปตฺตี”ติ ทีเปนฺติ, สาวเสสํ อาปตฺติํ “สาวเสสา อาปตฺตี”ติ ทีเปนฺติ, อนวเสสํ อาปตฺติํ “อนวเสสา อาปตฺตี”ติ ทีเปนฺติ, ทุฏฺฐุลฺลํ อาปตฺติํ “ทุฏฺฐุลฺลา อาปตฺตี”ติ ทีเปนฺติ, อทุฏฺฐุลฺลํ อาปตฺติํ “อทุฏฺฐุลฺลา อาปตฺตี”ติ ทีเปนฺติ.\csromanspacer{} เต อิเมหิ อฏฺฐารสหิ วตฺถูหิ น อปกสฺสนฺติ, อวปกสฺสนฺติ, น อาเวนิํ อุโปสถํ กโรนฺติ, น อาเวนิํ ปวารณํ กโรนฺติ, น อาเวนิํ สํฆกมฺมํ กโรนฺติ.\csromanspacer{} เอตฺตาวตา โข อุปาลิ สํโฆ สมคฺโค โหตีติ.}

\proseitem{345}{สมคฺคํ ปน ภนฺเต สํฆํ ภนฺทิตฺวา กิํ โส ปสวตีติ.\csromanspacer{} สมคฺคํ โข อุปาลิ สํฆํ ภินฺทิตฺวา กปฺปฏฺฐิกํ กิพฺพิสํ ปสวติ, กปฺปํ นิรยมฺหิ ปจฺจตีติ.}

\prose{\csromansymbolfootnote{*}{ขุ 1. 203; อํ 3. 315, 316 ปิฏฺเฐสุปิ.}อาปายิโก เนรยิโก, กปฺปฏฺโฐ สํฆเภทโก.}

\csromangathameasure{วคฺครโต อธมฺมฏฺโฐ, โยคกฺเขมา ปธํสติ. \\ สํฆํ สมคฺคํ ภินฺทิตฺวา, กปฺปํ นิรยมฺหิ ปจฺจตีติ.}
\csromangathagroup{\csromangathastanza{วคฺครโต อธมฺมฏฺโฐ, โยคกฺเขมา ปธํสติ. \\ สํฆํ สมคฺคํ ภินฺทิตฺวา, กปฺปํ นิรยมฺหิ ปจฺจตีติ.}}

\prose{ภินฺนํ ปน ภนฺเต สํฆํ สมคฺคํ กตฺวา กิํ โส ปสวตีติ.\csromanspacer{} ภินฺนํ โข อุปาลิ สํฆํ สมคฺคํ กตฺวา พฺรหฺมํ ปุญฺญํ ปสวติ, กปฺปํ สคฺคมฺหิ โมทตีติ.}

\prose{\csromansymbolmark{*}สุขา สํฆสฺส สามคฺคี, สมคฺคานญฺจ อนุคฺคโห.}

\csromangathameasure{สมคฺครโต ธมฺมฏฺโฐ, โยคกฺเขมา น ธํสติ. \\ สํฆํ สมคฺคํ กตฺวาน, กปฺปํ สคฺคมฺหิ โมทตีติ.}
\csromangathagroup{\csromangathastanza{สมคฺครโต ธมฺมฏฺโฐ, โยคกฺเขมา น ธํสติ. \\ สํฆํ สมคฺคํ กตฺวาน, กปฺปํ สคฺคมฺหิ โมทตีติ.}}

\proseitem{355}{\csromansymbolfootnote{+}{วิ 5. 350 ปิฏฺเฐปิ.}สิยา นุ โข ภนฺเต สํฆเภทโก อาปายิโก เนรยิโก กปฺปฏฺโฐ อเตกิจฺโฉติ.\csromanspacer{} สิยา อุปาลิ สํฆเภทโก อาปายิโก เนรยิโก กปฺปฏฺโฐ อเตกิจฺโฉติ.}

\prose{\csromanfolio{370}\csromansymbolfootnote{*}{วิ 5. 350 ปิฏฺเฐปิ.}สิยา\footnote{สิยา นุ โข (สฺยา, กํ)} ปน ภนฺเต สํฆเภทโก น อาปายิโก น เนรยิโก น กปฺปฏฺโฐ น อเตกิจฺโฉติ.\csromanspacer{} สิยา อุปาลิ สํฆเภทโก น อาปายิกา น เนรยิโก น กปฺปฏฺโฐ น อเตกิจฺโฉติ.}

\prose{กตโม ปน ภนฺเต สํฆเภทโก อาปายิโก เนรยิโก กปฺปฏฺโฐ อเตกิจฺโฉติ.\csromanspacer{} อิธุปาลิ ภิกฺขุ อธมฺมํ “ธมฺโม”ติ ทีเปติ.\csromanspacer{} ตสฺมิํ อธมฺมทิฏฺฐิ เภเท อธมฺมทิฏฺฐิ วินิธาย ทิฏฺฐิํ วินิธาย ขนฺติํ วินิธาย รุจิํ วินิธาย ภาวํ อนุสฺสาเวติ, สลากํ คาเหติ “อยํ ธมฺโม, อยํ วินโย, อิทํ สตฺถุสาสนํ, อิมํ คณฺหถ, อิมํ โรเจตา”ติ.\csromanspacer{} อยมฺปิ โข อุปาลิ สํฆเภทโก อาปายิโก เนรยิโก กปฺปฏฺโฐ อเตกิจฺโฉ.}

\prose{ปุน จปรํ อุปาลิ ภิกฺขุ อธมฺมํ “ธมฺโม”ติ ทีเปติ, ตสฺมิํ อธมฺมทิฏฺฐิ เภเท ธมฺมทิฏฺฐิ วินิธาย ทิฏฺฐิํ วินิธาย ขนฺติํ วินิธาย รุจิํ วินิธาย ภาวํ อนุสฺสาเวติ, สลากํ คาเหติ “อยํ ธมฺโม, อยํ วินโย, อิทํ สตฺถุสาสนํ, อิมํ คณฺหถ, อิมํ โรเจถา”ติ.\csromanspacer{} อยมฺปิ โข อุปาลิ สํฆเภทโก อาปายิโก เนรยิโก กปฺปฏฺโฐ อเตกิจฺโฉ.}

\prose{ปุน จปรํ อุปาลิ ภิกฺขุ อธมฺมํ “ธมฺโม”ติ ทีเปติ, ตสฺมิํ อธมฺมทิฏฺฐิ เภเท เวมติโก วินิธาย ทิฏฺฐิํ วินิธาย ขนฺติํ วินิธาย รุจิํ วินิธาย ภาวํ อนุสฺสาเวติ, สลากํ คาเหติ “อยํ ธมฺโม, อยํ วินโย, อิทํ สตฺถุสาสนํ, อิมํ คณฺหถ, อิมํ โรเจถา”ติ.\csromanspacer{} อยมฺปิ โข อุปาลิ สํฆเภทโก อาปายิโก เนรยิโก กปฺปฏฺโฐ อเตกิจฺโฉ.}

\prose{ปุน จปรํ อุปาลิ ภิกฺขุ อธมฺมํ “ธมฺโม”ติ ทีเปติ, ตสฺมิํ ธมฺมทิฏฺฐิ เภเท อธมฺมทิฏฺฐิ ฯเปฯ (ตสฺมิํ ธมฺมทิฏฺฐิ เภเท ธมฺมทิฏฺฐิ)\footnote{( ) สฺยาม-โปตฺถเก นตฺถิ, วิมติวิโนทนีฏีกาย สเมติ.}.\csromanspacer{} ตสฺมิํ ธมฺมทิฏฺฐิ เภเท เวมติโก.\csromanspacer{} ตสฺมิํ เวมติโก เภเท อธมฺมทิฏฺฐิ.\csromanspacer{} ตสฺมิํ เวมติโก เภเท ธมฺมทิฏฺฐิ.\csromanspacer{} ตสฺมิํ เวมติโก เภเท เวมติโก วินิธาย ทิฏฺฐิํ วินิธาย ขนฺติํ วินิธาย รุจิํ วินิธาย ภาวํ อนุสฺสาเวติ, สลากํ คาเหติ “อยํ ธมฺโม, อยํ วินโย, อิทํ สตฺถุสาสนํ, อิมํ คณฺหถ, อิมํ โรเจถา”ติ.\csromanspacer{} อยมฺปิ โข อุปาลิ สํฆเภทโก อาปายิโก เนรยิโก กปฺปฏฺโฐ อเตกิจฺโฉ.}

\prose{\csromanfolio{371}ปุน จปรํ อุปาลิ ภิกฺขุ ธมฺมํ “อธมฺโม”ติ ทีเปติ ฯเปฯ อวินยํ “วินโย”ติ ทีเปติ.\csromanspacer{} วินยํ “อวินโย”ติ ทีเปติ.\csromanspacer{} อภายิตํ อลปิตํ ตถาคเตน “ภาสิตํ ลปิตํ ตถาคเตนา”ติ ทีเปติ.\csromanspacer{} ภาสิตํ ลปิตํ ตถาคเตน “อภาสิตํ อลปิตํ ตถาคเตนา”ติ ทีเปติ.\csromanspacer{} อนาจิณฺณํ ตถาคเตน “อาจิณฺณํ ตถาคเตนา”ติ ทีเปติ.\csromanspacer{} อาจิณฺณํ ตถาคเตน “อนาจิณฺณํ ตถาคเตนา”ติ ทีเปติ.\csromanspacer{} อปญฺญตฺตํ ตถาคเตน “ปญฺญตฺตํ ตถาคเตนา”ติ ทีเปติ.\csromanspacer{} ปญฺญตฺตํ ตถาคเตน “อปญฺญตฺตํ ตถาคเตนา”ติ ทีเปติ.\csromanspacer{} อนาปตฺติํ “อาปตฺตี”ติ ทีเปติ.\csromanspacer{} อาปตฺติํ “อนาปตฺตี”ติ ทีเปติ.\csromanspacer{} ลหุกํ อาปตฺติํ “ครุกา อาปตฺตี”ติ ทีเปติ.\csromanspacer{} ครุกํ อาปตฺติํ “ลหุกา อาปตฺตี”ติ ทีเปติ.\csromanspacer{} สาวเสสํ อาปตฺติํ “อนวเสสา อาปตฺตี”ติ ทีเปติ.\csromanspacer{} อนวเสสํ อาปตฺติํ “สาวเสสา อาปตฺตี”ติ ทีเปติ.\csromanspacer{} ทุฏฺฐุลฺลํ อาปตฺติํ “อทุฏฺฐุลฺลา อาปตฺตี”ติ ทีเปติ.\csromanspacer{} อทุฏฺฐุลฺลํ อาปตฺติํ “ทุฏฺฐุลฺลา อาปตฺตี”ติ ทีเปติ.\csromanspacer{} ตสฺมิํ อธมฺมทิฏฺฐิ เภเท อธมฺมทิฏฺฐิ ฯเปฯ.\csromanspacer{} ตสฺมิํ อธมฺมทิฏฺฐิ เภเท ธมฺมทิฏฺฐิ.\csromanspacer{} ตสฺมิํ อธมฺมทิฏฺฐิ เภเท เวมติโก.\csromanspacer{} ตสฺมิํ ธมฺมทิฏฺฐิ เภเท อธมฺมทิฏฺฐิ. (ตสฺมิํ ธมฺมทิฏฺฐิ เภเท ธมฺมทิฏฺฐิ)\footnote{( ) สฺยาม-โปตฺถเก นตฺถิ.}.\csromanspacer{} ตสฺมิํ ธมฺมทิฏฺฐิ เภเท เวมติโก.\csromanspacer{} ตสฺมิํ เวมติโก เภเท อธมฺมทิฏฺฐิ.\csromanspacer{} ตสฺมิํ เวมติโก เภเท ธมฺมทิฏฺฐิ.\csromanspacer{} ตสฺมิํ เวมติโก เภเท เวมติโก วินิธาย ทิฏฺฐิํ วินิธาย ขนฺติํ วินิธาย รุจิํ วินิธาย ภาวํ อนุสฺสาเวติ, สลากํ คาเหติ “อยํ ธมฺโม, อยํ วินโย, อิทํ สตฺถุสาสนํ, อิมํ คณฺหถ, อิมํ โรเจถา”ติ.\csromanspacer{} อยมฺปิ โข อุปาลิ สํฆเภทโก อาปายิโก เนรยิโก กปฺปฏฺโฐ อเตกิจฺโฉติ.}

\prose{กตโม ปน ภนฺเต สํฆเภทโก น อาปายิโก น เนรยิโก น กปฺปฏฺโฐ น อเตกิจฺโฉติ.\csromanspacer{} อิธุปาลิ ภิกฺขุ อธมฺมํ “ธมฺโม”ติ ทีเปติ, ตสฺมิํ ธมฺมทิฏฺฐิ เภเท ธมฺมทิฏฺฐิ อวินิธาย ทิฏฺฐิํ อวินิธาย ขนฺติํ อวินิธาย รุจิํ อวินิธาย ภาวํ อนุสฺสาเวติ, สลากํ คาเหติ “อยํ ธมฺโม, อยํ วินโย, อิทํ สตฺถุสาสนํ, อิมํ คณฺหถ, อิมํ โรเจถา”ติ.\csromanspacer{} อยมฺปิ โข อุปาลิ สํฆเภทโก น อาปายิโก น เนรยิโก น กปฺปฏฺโฐ น อเตกิจฺโฉ.}

\csromanmatikaanchor{mtk.175}
\csromantocmarkref{section}{อุทฺทานคาถา}{mtk.175}
\bookmark[dest={mtk.175},level=1]{อุทฺทานคาถา}
\prose{\csromanfolio{372}ปุน จปรํ อุปาลิ ภิกฺขุ ธมฺมํ “อธมฺโม”ติ ทีเปติ ฯเปฯ ทุฏฺฐุลฺลํ อาปตฺติํ “อทุฏฺฐุลฺลา อาปตฺตี”ติ ทีเปติ, ตสฺมิํ ธมฺมทิฏฺฐิ เภเท ธมฺมทิฏฺฐิ อวินิธาย ทิฏฺฐิํ อวินิธาย ขนฺติํ อวินิธาย รุจิํ อวินิธาย ภาวํ อนุสฺสาเวติ, สลากํ คาเหติ “อยํ ธมฺโม, อยํ วินโย, อิทํ สตฺถุสาสนํ, อิมํ คณฺหถ, อิมํ โรเจถา”ติ.\csromanspacer{} อยมฺปิ โข อุปาลิ สํฆเภทโก น อาปายิโก น เนรยิโก น กปฺปฏฺโฐ น อเตกิจฺโฉติ.}

\nitthitam{ตติยภาณวาโร นิฏฺฐิโต.}
\vspace{\tipitakaparskip}
\csromancenter{สํฆเภทกกฺขนฺธโก สตฺตโม.}
\titlehead{\textbf{ตสฺสุทฺทานํ}}
\csromangathameasure{อนุปิเย อภิญฺญาตา, สุขุมาโล น อิจฺฉติ. \\ กสา วปา อภิ นินฺเน, นิทฺธา ลาเว จ อุพฺพเห. \\ ปุญฺชมทฺทปลาลญฺจ, ภุส-โอผุณนีหเร. \\ อายติมฺปิ น ขียนฺติ, ปิตโร จ ปิตามหา. \\ ภทฺทิโย อนุรุทฺโธ จ, อานนฺโท ภคุ กิมิโล. \\ สกฺยมาโน จ โกสมฺพิํ, ปริหายิ กกุเธน จ. \\ ปกาเสสิ ปิตุโน จ, ปุริเส สิลํ นาฬาคิริํ. \\ ติกปญฺจครุโก โข, ภินฺทิ ถุลฺลจฺจเยน จ. \\ ตโย อฏฺฐ ปุน ตีณิ, ราชิ เภทา สิยา นุ โขติ.}
\csromangathagroup{\csromangathastanza{อนุปิเย อภิญฺญาตา, สุขุมาโล น อิจฺฉติ. \\ กสา วปา อภิ นินฺเน, นิทฺธา ลาเว จ อุพฺพเห.}\csromangathastanza{ปุญฺชมทฺทปลาลญฺจ, ภุส-โอผุณนีหเร. \\ อายติมฺปิ น ขียนฺติ, ปิตโร จ ปิตามหา.}\csromangathastanza{ภทฺทิโย อนุรุทฺโธ จ, อานนฺโท ภคุ กิมิโล. \\ สกฺยมาโน จ โกสมฺพิํ, ปริหายิ กกุเธน จ.}\csromangathastanza{ปกาเสสิ ปิตุโน จ, ปุริเส สิลํ นาฬาคิริํ. \\ ติกปญฺจครุโก โข, ภินฺทิ ถุลฺลจฺจเยน จ. \\ ตโย อฏฺฐ ปุน ตีณิ, ราชิ เภทา สิยา นุ โขติ.}}

\nitthitam{\textbf{สํฆเภทกกฺขนฺธโก นิฏฺฐิโต.}}
\csromanmatikaanchor{mtk.176}
\csromantocmarknopage{chapter}{8. วตฺตกฺขนฺธก}{mtk.176}
\bookmark[dest={mtk.176},level=0]{8. วตฺตกฺขนฺธก}
\chapterheadpageread{8. วตฺตกฺขนฺธก}
\csromanmatikaanchor{mtk.177}
\csromantocmarkref{section}{1. อาคนฺตุกวตฺตกถา}{mtk.177}
\bookmark[dest={mtk.177},level=1]{1. อาคนฺตุกวตฺตกถา}
\csromanheader{1}{1. อาคนฺตุกวตฺตกถา}
\proseitem{356}{\csromanfolio{373}เตน สมเยน พุทฺโธ ภควา สาวตฺถิยํ วิหรติ เชตวเน อนาถปิณฺฑิกสฺส อาราเม.\csromanspacer{} เตน โข ปน สมเยน อาคนฺตุกา ภิกฺขู ส-อุปาหนาปิ อารามํ ปวิสนฺติ, ฉตฺตปคฺคหิตาปิ อารามํ ปวิสนฺติ, โอคุณฺฐิตาปิ อารามํ ปวิสนฺติ, สีเสปิ จีวรํ กริตฺวา อารามํ ปวิสนฺติ, ปานีเยนปิ ปาเท โธวนฺติ, วุฑฺฒตเรปิ อาวาสิเก ภิกฺขู น อภิวาเทนฺติ, นปิ เสนาสนํ ปุจฺฉนฺติ, อญฺญตโรปิ อาคนฺตุโก ภิกฺขุ อนชฺฌาวุฏฺฐํ วิหารํ ฆฏิกํ อุคฺฆาเฏตฺวา กวาฏํ ปณาเมตฺวา สหสา ปาวิสิ, ตสฺส อุปริปิฏฺฐิโต\footnote{อุปริปิฏฺฐโต (?)} อหิ ขนฺเธ ปปติ, โส ภีโต วิสฺสรมกาสิ, ภิกฺขู อุปธาวิตฺวา ตํ ภิกฺขุํ เอตทโวจุํ “กิสฺส ตฺวํ อาวุโส วิสฺสรมกาสี”ติ.\csromanspacer{} อถ โข โส ภิกฺขุ ภิกฺขูนํ เอตมตฺถํ อาโรเจสิ.\csromanspacer{} เย เต ภิกฺขู อปฺปิจฺฉา ฯเปฯ เต อุชฺฌายนฺติ ขิยฺยนฺติ วิปาเจนฺติ “กถํ หิ นาม อาคนฺตุกา ภิกฺขู ส-อุปาหนาปิ อารามํ ปวิสิสฺสนฺติ, ฉตฺตปคฺคหิตาปิ อารามํ ปวิสิสฺสนฺติ, โอคุณฺฐิตาปิ อารามํ ปวิสิสฺสนฺติ, สีเสปิ จีวรํ กริตฺวา อารามํ ปวิสิสฺสนฺติ, ปานีเยนปิ ปาเท โธวิสฺสนฺติ, วุฑฺฒตเรปิ อาวาสิเก ภิกฺขู น อภิวาเทสฺสนฺติ, นปิ เสนาสนํ ปุจฺฉิสฺสนฺตี”ติ.\csromanspacer{} อถ โข เต ภิกฺขู ภควโต เอตมตฺถํ อาโรเจสุํ ฯเปฯ.\csromanspacer{} สจฺจํ กิร ภิกฺขเว อาคนฺตุกา ภิกฺขู ส-อุปาหนาปิ อารามํ ปวิสนฺติ, ฉตฺตปคฺคหิตาปิ อารามํ ปวิสนฺติ, โอคุณฺฐิตาปิ อารามํ ปวิสนฺติ, สีเสปิ จีวรํ กริตฺวา อารามํ ปวิสนฺติ, ปานีเยนปิ ปาเท โธวนฺติ, วุฑฺฒตเรปิ อาวาสิเก ภิกฺขู น อภิวาเทนฺติ, นปิ เสนาสนํ ปุจฺฉนฺตีติ.\csromanspacer{} สจฺจํ ภควาติ.\csromanspacer{} วิครหิ พุทฺโธ ภควา ฯเปฯ.\csromanspacer{} กถํ หิ นาม ภิกฺขเว อาคนฺตุกา ภิกฺขู ส-อุปาหนาปิ อารามํ ปวิสิสฺสนฺติ, ฉตฺตปคฺคหิตาปิ อารามํ ปวิสิสฺสนฺติ, โอคุณฺฐิตาปิ อารามํ ปวิสิสฺสนฺติ, สีเสปิ จีวรํ กริตฺวา อารามํ ปวิสิสฺสนฺติ, ปานีเยนปิ ปาเท โธวิสฺสนฺติ, วุฑฺฒตเรปิ อาวาสิเก ภิกฺขู น อภิวาเทสฺสนฺติ, นปิ เสนาสนํ ปุจฺฉิสฺสนฺติ, เนตํ ภิกฺขเว อปฺปสนฺนานํ วา ปสาทาย ฯเปฯ วิครหิตฺวา ฯเปฯ ธมฺมิํ กถํ กตฺวา ภิกฺขู อามนฺเตสิ\csromandash{}}

\proseitem{357}{\csromanfolio{374}เตน หิ ภิกฺขเว อาคนฺตุกานํ ภิกฺขูนํ วตฺตํ ปญฺญเปสฺสามิ, ยถา อาคนฺตุเกหิ ภิกฺขูหิ สมฺมา วตฺติตพฺพํ.\csromanspacer{} อาคนฺตุเกน ภิกฺขเว ภิกฺขุนา “อิทานิ อารามํ ปวิสิสฺสามี”ติ อุปาหนา โอมุญฺจิตฺวา นีจํ กตฺว ปปฺโผเฏตฺวา คเหตฺวา ฉตฺตํ อปนาเมตฺวา สีสํ วิวริตฺวา สีเส จีวรํ\footnote{วิวริตฺวา จีวรํ (ก)} ขนฺเธ กตฺวา สาธุกํ อตรมาเนน อาราโม ปวิสิตพฺโพ, อารามํ ปวิสนฺเตน สลฺลกฺเขตพฺพํ “กตฺถ อาวาสิกา ภิกฺขู ปฏิกฺกมนฺตี”ติ.\csromanspacer{} ยตฺถ อาวาสิกา ภิกฺขู ปฏิกฺกมนฺติ อุปฏฺฐานสาลาย วา มณฺฑเป วา รุกฺขมูเล วา, ตตฺถ คนฺตฺวา เอกมนฺตํ ปตฺโต นิกฺขิปิตพฺโพ, เอกมนฺตํ จีวรํ นิกฺขิปิตพฺพํ, ปติรูปํ อาสนํ คเหตฺวา นิสีทิตพฺพํ, ปานียํ ปุจฺฉิตพฺพํ, ปริโภชนียํ ปุจฺฉิตพฺพํ “กตมํ ปานียํ, กตมํ ปริโภชนียนฺ”ติ.\csromanspacer{} สเจ ปานีเยน อตฺโถ โหติ, ปานียํ คเหตฺวา ปาตพฺพํ, สเจ ปริโภชนีเยน อตฺโถ โหติ, ปริโภชนียํ คเหตฺวา ปาทา โธวิตพฺพา, ปาเท โธวนฺเตน เอเกน หตฺเถน อุทกํ อาสิญฺจิตพฺพํ, เอเกน หตฺเถน ปาทา โธวิตพฺพา, น เตเนว หตฺเถน อุทกํ อาสิญฺจิตพฺพํ\footnote{เยน หตฺเถน อุทกํ อาสิญฺจิตพฺพํ (สฺยา)}, น เตเนว หตฺเถน ปาทา โธวิตพฺพา, อุปาหนาปุญฺฉนโจฬกํ ปุจฺฉิตฺวา อุปาหนา ปุญฺฉิตพฺพา, อุปาหนา ปุญฺฉนฺเตน ปฐมํ สุกฺเขน โจฬเกน ปุญฺฉิตพฺพา, ปจฺฉา อลฺเลน, อุปาหนาปุญฺฉนโจฬกํ โธวิตฺวา\footnote{ปีเฬตฺวา (สฺยา)} เอกมนฺตํ วิสฺสชฺเชตพฺพํ.}

\prose{สเจ อาวาสิโก ภิกฺขุ วุฑฺโฒ โหติ, อภิวาเทตพฺโพ, สเจ นวโก โหติ, อภิวาทาเปตพฺโพ.\csromanspacer{} เสนาสนํ ปุจฺฉิตพฺพํ “กตมํ เม เสนาสนํ ปาปุณาตี”ติ, อชฺฌาวุฏฺฐํ วา อนชฺฌาวุฏฺฐํ วา ปุจฺฉิตพฺพํ, โคจโร ปุจฺฉิตพฺโพ, อโคจโร ปุจฺฉิตพฺโพ, เสกฺขสมฺมตานิ\footnote{เสขสมฺมตานิ (ก)} กุลานิ ปุจฺฉิตพฺพานิ, วจฺจฏฺฐานํ ปุจฺฉิตพฺพํ, ปสฺสาวฏฺฐานํ ปุจฺฉิตพฺพํ, ปานียํ ปุจฺฉิตพฺพํ, ปริโภชนียํ ปุจฺฉิตพฺพํ, กตฺตรทณฺโฑ ปุจฺฉิตพฺโพ, สํฆสฺส กติกสณฺฐานํ ปุจฺฉิตพฺพํ “กํ กาลํ ปวิสิตพฺพํ, กํ กาลํ นิกฺขมิตพฺพนฺ”ติ.\csromanspacer{} สเจ วิหาโร อนชฺฌาวุฏฺโฐ โหติ, กวาฏํ อาโกเฏตฺวา มุหุตฺตํ อาคเมตฺวา ฆฏิกํ อุคฺฆาเฏตฺวา กวาฏํ ปณาเมตฺวา พหิ ฐิเตน นิลฺโลเกตพฺโพ.}

\prose{\csromanfolio{375}สเจ วิหาโร อุกฺลาโป โหติ, มญฺเจ วา มญฺโจ อาโรปิโต โหติ, ปีเฐ วา ปีฐํ อาโรปิตํ โหติ, เสนาสนํ อุปริ ปุญฺชีกตํ\footnote{ปุญฺชกิตํ (ก)} โหติ, สเจ อุสฺสหติ โสเธตพฺโพ, * วิหารํ โสเธนฺเตน ปฐมํ ภูมตฺถรณํ นีหริตฺวา เอกมนฺตํ นิกฺขิปิตพฺพํ, มญฺจปฏิปาทกา นีหริตฺวา เอกมนฺตํ นิกฺขิปิตพฺพา, ภิสิพิพฺโพหนํ นีหริตฺวา เอกมนฺตํ นิกฺขิปิตพฺพํ, นิสีทนปจฺจตฺถรณํ นีหริตฺวา เอกมนฺตํ นิกฺขิปิตพฺพํ, มญฺโจ นีจํ กตฺวา สาธุกํ อปฺปฏิฆํสนฺเตน อสํฆฏฺเฏนฺเตน กวาฏปิฏฺฐํ นีหริตฺวา เอกมนฺตํ นิกฺขิปิตพฺโพ, ปีฐํ นีจํ กตฺวา สาธุกํ อปฺปฏิฆํสนฺเตน อสํฆฏฺเฏนฺเตน กวาฏปิฏฺฐํ นีหริตฺวา เอกมนฺตํ นิกฺขิปิตพฺพํ, เขฬมลฺลโก นีหริตฺวา เอกมนฺตํ นิกฺขิปิตพฺโพ, อปสฺเสนผลกํ นีหริตฺวา เอกมนฺตํ นิกฺขิปิตพฺพํ.\csromanspacer{} สเจ วิหาเร สนฺตานกํ โหติ, อุลฺโลกา ปฐมํ โอหาเรตพฺพํ, อาโลกสนฺธิกณฺณภาคา ปมชฺชิตพฺพา.\csromanspacer{} สเจ เครุกปริกมฺมกตา ภิตฺติ กณฺณกิตา โหติ, โจฬกํ เตเมตฺวา ปีเฬตฺวา ปมชฺชิตพฺพา.\csromanspacer{} สเจ กาฬวณฺณกตา ภูมิ กณฺณกิตา โหติ, โจฬกํ เตเมตฺวา ปีเฬตฺวา ปมชฺชิตพฺพา.\csromanspacer{} สเจ อกตา โหติ ภูมิ, อุทเกน ปริปฺโผสิตฺวา สมฺมชฺชิตพฺพา “มา วิหาโร รเชน อุหญฺญี”ติ\footnote{อูหญฺญีติ (สี, สฺยา)}, สงฺการํ วิจินิตฺวา เอกมนฺตํ ฉฑฺเฑตพฺพํ.}

\prose{\csromansymbolfootnote{*}{วิ 3. 60, 65 ปิฏฺฐาทีสุปิ (โถกํ วิสทิสํ)}ภูมตฺถรณํ โอตาเปตฺวา โสเธตฺวา ปปฺโผเฏตฺวา อติหริตฺวา ยถาฐาเน\footnote{ยถาปญฺญตฺตํ (สี, สฺยา), ยถาภาคํ (ก)} ปญฺญเปตพฺพํ, มญฺจปฏิปาทกา โอตาเปตฺวา ปมชฺชิตฺวา อติหริตฺวา ยถาฐาเน\footnote{ยถาภาคํ (สฺยา, ก)} ฐเปตพฺพา, มญฺโจ โอตาเปตฺวา โสเธตฺวา ปปฺโผเฏตฺวา นีจํ กตฺวา สาธุกํ อปฺปฏิฆํสนฺเตน อสํฆฏฺเฏนฺเตน กวาฏปิฏฺฐํ อติหริตฺวา ยถาฐาเน4 ปญฺญเปตพฺโพ, ปีฐํ โอตาเปตฺวา โสเธตฺวา ปปฺโผเฏตฺวา นีจํ กตฺวา สาธุกํ อปฺปฏิฆํสนฺเตน อสํฆฏฺเฏนฺเตน กวาฏปิฏฺฐํ อติหริตฺวา ยถาฐาเน4 ปญฺญเปตพฺพํ, ภิสิพิพฺโพหนํ โอตาเปตฺวา โสเธตฺวา ปปฺโผเฏตฺวา อติหริตฺวา ยถาภาคํ ปญฺญเปตพฺพํ, นิสีทนปจฺจตฺถรณํ โอตาเปตฺวา โสเธตฺวา ปปฺโผเฏตฺวา อติหริตฺวา ยถาภาคํ ปญฺญเปตพฺพํ, เขฬมลฺลโก โอตาเปตฺวา ปมชฺชิตฺวา อติหริตฺวา ยถาภาคํ ฐเปตพฺโพ, อปสฺเสนผลกํ โอตาเปตฺวา ปมชฺชิตฺวา อติหริตฺวา ยถาภาคํ ฐเปตพฺพํ.\csromanspacer{} ปตฺตจีวรํ นิกฺขิปิตพฺพํ, ปตฺตํ \csromanfolio{376}นิกฺขิปนฺเตน เอเกน หตฺเถน ปตฺตํ คเตตฺวา เอเกน หตฺเถน เหฏฺฐามญฺจํ วา เหฏฺฐาปีฐํ วา ปรามสิตฺวา ปตฺโต นิกฺขิปิตพฺโพ, น จ อนนฺตรหิตาย ภูมิยา ปตฺโต นิกฺขิปิตพฺโพ.\csromanspacer{} จีวรํ นิกฺขิปนฺเตน เอเกน หตฺเถน จีวรํ คเหตฺวา เอเกน หตฺเถน จีวรวํสํ วา จีวรรชฺชุํ วา ปมชฺชิตฺวา ปารโต อนฺตํ โอรโต โภคํ กตฺวา จีวรํ นิกฺขิปิตพฺพํ.}

\prose{สเจ ปุรตฺถิมา สรชา วาตา วายนฺติ, ปุรตฺถิมา วาตปานา ถเกตพฺพา.\csromanspacer{} สเจ ปจฺฉิมา สรชา วาตา วายนฺติ, ปจฺฉิมา วาตปานา ตเกตพฺพา.\csromanspacer{} สเจ อุตฺตรา สรชา วาตา วายนฺติ, อุตฺตรา วาตปานา ถเกตพฺพา.\csromanspacer{} สเจ ทกฺขิณา สรชา วาตา วายนฺติ, ทกฺขิณา วาตปานา ถเกตพฺพา.\csromanspacer{} สเจ สีตกาโล โหติ, ทิวา วาตปานา วิวริตพฺพา, รตฺติํ ถเกตพฺพา.\csromanspacer{} สเจ อุณฺหกาโล โหติ, ทิวา วาตปานา ถเกตพฺพา, รตฺติํ วิวริตพฺพา.}

\prose{สเจ ปริเวณํ อุกฺลาปํ โหติ, ปริเวณํ สมฺมชฺชิตพฺพํ.\csromanspacer{} สเจ โกฏฺฐโก อุกฺลาโป โหติ, โกฏฺฐโก สมฺมชฺชิตพฺโพ.\csromanspacer{} สเจ อุปฏฺฐานสาลา อุกฺลาปา โหติ, อุปฏฺฐานสาลา สมฺมชฺชิตพฺพา.\csromanspacer{} สเจ อคฺคิสาลา อุกฺลาปา โหติ, อคฺคิสาลา สมฺมชฺชิตพฺพา.\csromanspacer{} สเจ วจฺจกุฏิ อุกฺลาปา โหติ, วจฺจกุฏิ สมฺมชฺชิตพฺพา.\csromanspacer{} สเจ ปานียํ น โหติ, ปานียํ อุปฏฺฐาเปตพฺพํ.\csromanspacer{} สเจ ปริโภชนียํ น โหติ, ปริโภชนียํ อุปฏฺฐาเปตพฺพํ.\csromanspacer{} สเจ อาจมนกุมฺภิยา อุทกํ น โหติ, อาจมนกุมฺภิยา อุทกํ อาสิญฺจิตพฺพํ.\csromanspacer{} อิทํ โข ภิกฺขเว อาคนฺตุกานํ ภิกฺขูนํ วตฺตํ, ยถา อาคนฺตุเกหิ ภิกฺขูหิ สมฺมา วตฺติตพฺพนฺติ.}

\csromanmatikaanchor{mtk.178}
\csromantocmarkref{section}{2. อาวาสิกวตฺตกถา}{mtk.178}
\bookmark[dest={mtk.178},level=1]{2. อาวาสิกวตฺตกถา}
\csromanheader{1}{2. อาวาสิกวตฺตกถา}
\proseitem{358}{เตน โข ปน สมเยน อาวาสิกา ภิกฺขู อาคนฺตุเก ภิกฺขู ทิสฺวา เนว อาสนํ ปญฺญเปนฺติ, น ปาโททกํ ปาทปีฐํ ปาทกถลิกํ อุปนีกฺขิปนฺติ, น ปจฺจุคฺคนฺตฺวา ปตฺตจีวรํ ปฏิคฺคณฺหนฺติ, น ปานีเยน ปุจฺฉนฺติ\footnote{น ปานีเยน ปุจฺฉนฺติ, น ปริโภชนีเยน ปุจฺฉนฺติ (สฺยา, กํ)}, น วุฑฺฒตเรปิ อาคนฺตุเก ภิกฺขู อภิวาเทนฺติ, น เสนาสนํ ปญฺญเปนฺติ, เย เต ภิกฺขู อปฺปิจฺฉา ฯเปฯ เต อุชฺฌายนฺติ ขิยฺยนฺติ วิปาเจนฺติ “กถํ หิ นาม อาวาสิกา ภิกฺขู อาคนฺตุเก ภิกฺขู ทิสฺวา เนว อาสนํ ปญฺญเปสฺสนฺติ, น ปาโททกํ ปาทปีฐํ ปาทกถลิกํ อุปนิกฺขิปิสฺสนฺติ, \csromanfolio{377}น ปจฺจุคฺคนฺตฺวา ปตฺตจีวรํ ปฏิคฺคหิสฺสนฺติ, น ปานีเยน ปุจฺฉิสฺสนฺติ.\csromanspacer{} วุฑฺฒตเรปิ อาคนฺตุเก ภิกฺขู น อภิวาเทสฺสนฺติ, น เสนาสนํ ปญฺญเปสฺสนฺตี”ติ.\csromanspacer{} อถ โข เต ภิกฺขู ภควโต เอตมตฺถํ อาโรเจสุํ ฯเปฯ.\csromanspacer{} สจฺจํ กิร ภิกฺขเว ฯเปฯ.\csromanspacer{} สจฺจํ ภควาติ ฯเปฯ วิครหิตฺวา ฯเปฯ ธมฺมิํ กถํ กตฺวา ภิกฺขู อามนฺเตสิ\csromandash{}}

\proseitem{359}{เตน หิ ภิกฺขเว อาวาสิกานํ ภิกฺขูนํ วตฺตํ ปญฺญเปสฺสามิ, ยถา อาวาสิเกหิ ภิกฺขูหิ สมฺมา วตฺติตพฺพํ.\csromanspacer{} อาวาสิเกน ภิกฺขเว ภิกฺขุนา อาคนฺตุกํ ภิกฺขุํ วุฑฺฒตรํ ทิสฺวา อาสนํ ปญฺญเปตพฺพํ, ปาโททกํ ปาทปีฐํ ปาทกถลิกํ อุปนิกฺขิปิตพฺพํ, ปจฺจุคฺคนฺตฺวา ปตฺตจีวรํ ปฏิคฺคเหตพฺพํ, ปานีเยน ปุจฺฉิตพฺโพ\footnote{ปานีเยน ปุจฺฉิตพฺโพ, ปริโภชนีเยน ปุจฺฉิตพฺโพ (สฺยา)}, สเจ อุสฺสหติ, อุปาหนา ปุญฺฉิตพฺพา, อุปาหนา ปุญฺฉนฺเตน ปฐมํ สุกฺเขน โจฬเกน ปุญฺฉิตพฺพา, ปจฺฉา อลฺเลน, อุปาหนาปุญฺฉนโจฬกํ โธวิตฺวา\footnote{โธวิตฺวา ปีเฬตฺวา (สฺยา)} เอกมนฺตํ วิสฺสชฺเชตพฺพํ.}

\prose{อาคนฺตุโก ภิกฺขุ วุฑฺฒตโร อภิวาเทตพฺโพ, เสนาสนํ ปญฺญเปตพฺพํ “เอตํ เต เสนาสนํ ปาปุณาตี”ติ.\csromanspacer{} อชฺฌาวุฏฺฐํ วา อนชฺฌาวุฏฺฐํ วา อาจิกฺขิตพฺพํ, โคจโร อาจิกฺขิตพฺโพ, อาโคจโร อาจิกฺขิตพฺโพ, เสกฺขสมฺมตานิ กุลานิ อาจิกฺขิตพฺพานิ, วจฺจฏฺฐานํ อาจิกฺขิตพฺพํ, ปสฺสาวฏฺฐานํ อาจิกฺขิตพฺพํ, ปานียํ อาจิกฺขิตพฺพํ, ปริโภชนียํ อาจิกฺขิตพฺพํ, กตฺตรทณฺโฑ อาจิกฺขิตพฺโพ, สํฆสฺส กติกสณฺฐานํ อาจิกฺขิตพฺพํ “อิมํ กาลํ ปวิสิตพฺพํ, อิมํ กาลํ นิกฺขมิตพฺพนฺ”ติ.}

\prose{สเจ นวโก โหติ, นิสินฺนเกเนว อาจิกฺขิตพฺพํ “อตฺร ปตฺตํ นิกฺขิปาหิ, อตฺร จีวรํ นิกฺขิปาหิ, อิทํ อาสนํ นิสีทาหี”ติ, ปานียํ อาจิกฺขิตพฺพํ, ปริโภชนียํ อาจิกฺขิตพฺพํ, อุปาหนาปุญฺฉนโจฬกํ อาจิกฺขิตพฺพํ, อาคนฺตุโก ภิกฺขุ นวโก อภิวาทาเปตพฺโพ, เสนาสนํ อาจิกฺขิตพฺพํ “เอตํ เต เสนาสนํ ปาปุณาตี”ติ, อชฺฌาวุฏฺฐํ วา อนชฺฌาวุฏฺฐํ วา อาจิกฺขิตพฺพํ, โคจโร อาจิกฺขิตพฺโพ, อโคจโร อาจิกฺขิตพฺโพ, เสกฺขสมฺมตานิ กุลานิ อาจิกฺขิตพฺพานิ, วจฺจฏฺฐานํ อาจิกฺขิตพฺพํ, ปสฺสาวฏฺฐานํ อาจิกฺขิตพฺพํ, ปานียํ อาจิกฺขิตพฺพํ, ปริโภชนียํ อาจิกฺขิตพฺพํ, กตฺตรทณฺโฑ อาจิกฺขิตพฺโพ, สํฆสฺส กติกสณฺฐานํ อาจิกฺขิตพฺพํ “อิมํ กาลํ ปวิสิตพฺพํ, อิมํ กาลํ นิกฺขมิตพฺพนฺ”ติ.\csromanspacer{} อิทํ โข \csromanfolio{378}ภิกฺขเว อาวาสิกานํ ภิกฺขูนํ วตฺตํ, ยถา อาวาสิเกหิ ภิกฺขูหิ สมฺมา วตฺติตพฺพนฺติ.}

\csromanmatikaanchor{mtk.179}
\csromantocmarkref{section}{3. คมิกวตฺตกถา}{mtk.179}
\bookmark[dest={mtk.179},level=1]{3. คมิกวตฺตกถา}
\csromanheader{1}{3. คมิกวตฺตกถา}
\proseitem{360}{เตน โข ปน สมเยน คมิกา ภิกฺขู ทารุภณฺฑํ มตฺติกาภณฺฑํ อปฺปฏิสาเมตฺวา ทฺวารวาตปานํ วิวริตฺวา เสนาสนํ อนาปุจฺฉา ปกฺกมนฺติ, ทารุภณฺฑํ มตฺติกาภณฺฑํ นสฺสติ, เสนาสนํ อคุตฺตํ โหติ.\csromanspacer{} เย เต ภิกฺขู อปฺปิจฺฉา ฯเปฯ เต อุชฺฌายนฺติ ขิยฺยนฺติ วิปาเจนฺติ “กถํ หิ นาม คมิกา ภิกฺขู ทารุภณฺฑํ มตฺติกาภณฺฑํ อปฺปฏิสาเมตฺวา ทฺวารวาตปานํ วิวริตฺวา เสนาสนํ อนาปุจฺฉา ปกฺกมิสฺสนฺติ, ทารุภณฺฑํ มตฺติกาภณฺฑํ นสฺสติ, เสนาสนํ อคุตฺตํ โหตี”ติ.\csromanspacer{} อถ โข เต ภิกฺขู ภควโต เอตมตฺถํ อาโรเจสุํ ฯเปฯ.\csromanspacer{} สจฺจํ กิร ภิกฺขเว ฯเปฯ.\csromanspacer{} สจฺจํ ภควาติ ฯเปฯ วิครหิตฺวา ฯเปฯ ธมฺมิํ กถํ กตฺวา ภิกฺขู อามนฺเตสิ\csromandash{}}

\proseitem{361}{เตน หิ ภิกฺขเว คมิกานํ ภิกฺขูนํ วตฺตํ ปญฺญเปสฺสามิ, ยถา คมิเกหิ ภิกฺขูหิ สมฺมา วตฺติตพฺพํ.\csromanspacer{} คมิเกน ภิกฺขเว ภิกฺขุนา ทารุภณฺฑํ มตฺติกาภณฺฑํ ปฏิสาเมตฺวา ทฺวารวาตปานํ ถเกตฺวา เสนาสนํ อาปุจฺฉา ปกฺกมิตพฺพํ\footnote{อาปุจฺฉิตพฺพํ (สฺยา)}.\csromanspacer{} สเจ ภิกฺขุ น โหติ, สามเณโร อาปุจฺฉิตพฺโพ.\csromanspacer{} สเจ สามเณโร น โหติ, อารามิโก อาปุจฺฉิตพฺโพ.\csromanspacer{} สเจ อารามิโก น โหติ, อุปาสโก อาปุจฺฉิตพฺโพ.\csromanspacer{} สเจ น โหติ ภิกฺขุ วา สามเณโร วา อารามิโก วา อุปาสโก วา, จตูสุ ปาสาเณสุ มญฺจํ ปญฺญเปตฺวา มญฺเจ มญฺจํ อาโรเปตฺวา ปีเฐ ปีฐํ อาโรเปตฺวา เสนาสนํ อุปาริ ปุญฺชํ กริตฺวา ทารุภณฺฑํ มตฺติกาภณฺฑํ ปฏิสาเมตฺวา ทฺวารวาตปานํ ถเกตฺวา ปกฺกมิตพฺพํ.\csromanspacer{} สเจ วิหาโร โอวสฺสติ, สเจ อุสฺสหติ, ฉาเทตพฺโพ, อุสฺสุกฺกํ วา กาตพฺพํ “กินฺติ นุ โข วิหาโร ฉาทิเยถา”ติ, เอวญฺเจตํ ลเภถ, อิจฺเจตํ กุสลํ.\csromanspacer{} โน เจ ลเภถ, โย เทโส อโนวสฺสโก โหติ, ตตฺถ จตูสุ ปาสาเณสุ มญฺจํ ปญฺญเปตฺวา มญฺเจ มญฺจํ อาโรเปตฺวา ปีเฐ ปีฐํ อาโรเปตฺวา เสนาสนํ อุปริ ปุญฺชํ กริตฺวา ทารุภณฺฑํ มตฺติกาภณฺฑํ ปฏิสาเมตฺวา ทฺวารวาตปานํ ถเกตฺวา ปกฺกมิตพฺพํ.\csromanspacer{} สเจ สพฺโพ วิหาโร โอวสฺสติ, สเจ อุสฺสหติ, เสนาสนํ คามํ อติหริตพฺพํ, อุสฺสุกฺกํ วา \csromanfolio{379}กาตพฺพํ “กินฺติ นุ โข เสนาสนํ คามํ อติหริเยถา”ติ, เอวญฺเจตํ ลเภถ, อิจฺเจตํ กุสลํ.\csromanspacer{} โน เจ ลเภถ, อชฺโฌกาเส จตูสุ ปาสาเณสุ มญฺจํ ปญฺญเปตฺวา มญฺเจ มญฺจํ อาโรเปตฺวา ปีเฐ ปีฐํ อาโรเปตฺวา เสนาสนํ อุปริ ปุญฺชํ กริตฺวา ทารุภณฺฑํ มตฺติกาภณฺฑํ ปฏิสาเมตฺวา ติเณน วา ปณฺเณน วา ปฏิจฺฉาเทตฺวา ปกฺกมิตพฺพํ “อปฺเปว นาม องฺคานิปิ เสเสยฺยุนฺ”ติ.\csromanspacer{} อิทํ โข ภิกฺขเว คมิกานํ ภิกฺขูนํ วตฺตํ, ยถา คมิเกหิ ภิกฺขูหิ สมฺมา วตฺติตพฺพนฺติ.}

\csromanmatikaanchor{mtk.180}
\csromantocmarkref{section}{4. อนุโมทนวตฺตกถา}{mtk.180}
\bookmark[dest={mtk.180},level=1]{4. อนุโมทนวตฺตกถา}
\csromanheader{1}{4. อนุโมทนวตฺตกถา}
\proseitem{362}{เตน โข ปน สมเยน ภิกฺขู ภตฺตคฺเค น อนุโมทนฺติ.\csromanspacer{} มนุสฺสา อุชฺฌายนฺติ ขิยฺยนฺติ วิปาเจนฺติ “กถํ หิ นาม สมณา สกฺยปุตฺติยา ภตฺตคฺเค น อนุโมทิสฺสนฺตี”ติ.\csromanspacer{} อสฺโสสุํ โข ภิกฺขู เตสํ มนุสฺสานํ อุชฺฌายนฺตานํ ขิยฺยนฺตานํ วิปาเจนฺตานํ.\csromanspacer{} อถ โข เต ภิกฺขู ภควโต เอตมตฺถํ อาโรเจสุํ.\csromanspacer{} อถ โข ภควา เอตสฺมิํ นิทาเน เอตสฺมิํ ปกรเณ ธมฺมิํ กถํ กตฺวา ภิกฺขู อามนฺเตสิ “อนุชานามิ ภิกฺขเว ภตฺตคฺเค อนุโมทิตุนฺ”ติ.\csromanspacer{} อถ โข เตสํ ภิกฺขูนํ เอตทโหสิ “เกน นุ โข ภตฺตคฺเค อนุโมทิตพฺพนฺ”ติ.\csromanspacer{} ภควโต เอตมตฺถํ อาโรเจสุํ.\csromanspacer{} อถ โข ภควา เอตสฺมิํ นิทาเน เอตสฺมิํ ปกรเณ ธมฺมิํ กถํ กตฺวา ภิกฺขู อามนฺเตสิ “อนุชานามิ ภิกฺขเว เถเรน ภิกฺขุนา ภตฺตคฺเค อนุโมทิตุนฺ”ติ.}

\prose{เตน โข ปน สมเยน อญฺญตรสฺส ปูคสฺส สํฆภตฺตํ โหติ, อายสฺมา สาริปุตฺโต สํฆตฺเถโร โหติ, ภิกฺขู “ภควตา อนุญฺญาตํ เถเรน ภิกฺขุนา ภตฺตคฺเค อนุโมทิตุนฺ”ติ อายสฺมนฺตํ สาริปุตฺตํ เอกกํ โอหาย ปกฺกมิํสุ.\csromanspacer{} อถ โข อายสฺมา สาริปุตฺโต เต มนุสฺเส ปฏิสมฺโมทิตฺวา ปจฺฉา เอกโก อคมาสิ.\csromanspacer{} อทฺทสา โข ภควา อายสฺมนฺตํ สาริปุตฺตํ ทูรโตว เอกกํ อาคจฺฉนฺตํ, ทิสฺวาน อายสฺมนฺตํ สาริปุตฺตํ เอตทโวจ “กจฺจิ สาริปุตฺต ภตฺตํ อิทฺธํ อโหสี”ติ.\csromanspacer{} อิทฺธํ โข ภนฺเต ภตฺตํ อโหสิ, อปิ จ มํ ภิกฺขู เอกกํ โอหาย ปกฺกนฺตาติ.\csromanspacer{} อถ โข ภควา เอตสฺมิํ นีทาเน เอตสฺมิํ ปกรเณ ธมฺมิํ กถํ กตฺวา ภิกฺขู อามนฺเตสิ “อนุชานามิ ภิกฺขเว ภตฺตคฺเค จตูหิ ปญฺจหิ เถรานุเถเรหิ ภิกฺขูหิ อาคเมตุนฺ”ติ.}

\prose{\csromanfolio{380}เตน โข ปน สมเยน อญฺญตโร เถโร ภตฺตคฺเค วจฺจิโต อาคเมสิ, โส วจฺจํ สทฺธาเรตุํ อสกฺโกนฺโต มุจฺฉิโต ปปติ.\csromanspacer{} ภควโต เอตมตฺถํ อาโรเจสุํ.\csromanspacer{} อนุชานามิ ภิกฺขเว สติ กรณีเย อานนฺตริกํ ภิกฺขุํ อาปุจฺฉิตฺวา คนฺตุนฺติ.}

\csromanmatikaanchor{mtk.181}
\csromantocmarkref{section}{5. ภตฺตคฺควตฺตกถา}{mtk.181}
\bookmark[dest={mtk.181},level=1]{5. ภตฺตคฺควตฺตกถา}
\csromanheader{1}{5. ภตฺตคฺควตฺตกถา}
\proseitem{363}{เตน โข ปน สมเยน ฉพฺพคฺคิยา ภิกฺขู ทุนฺนิวตฺถา ทุปฺปารุตา อนากปฺปสมฺปนฺนา ภตฺตคฺคํ คจฺฉนฺติ, โวกฺกมฺมปิ เถรานํ ภิกฺขูนํ ปุรโต ปุรโต คจฺฉนฺติ, เถเรปิ ภิกฺขู อนุปขชฺช นิสีทนฺติ, นเวปิ ภิกฺขู อาสเนน ปฏิพาหนฺติ, สํฆาฏิมฺปิ โอตฺถริตฺวา อนฺตรฆเร นิสีทนฺติ.\csromanspacer{} เย เต ภิกฺขู อปฺปิจฺฉา ฯเปฯ เต อุชฺฌายนฺติ ขิยฺยนฺติ วิปาเจนฺติ “กถํ หิ นาม ฉพฺพคฺคิยา ภิกฺขู ทุนฺนิวตฺถา ทุปฺปารุตา อนากปฺปสมฺปนฺนา ภตฺตคฺคํ คจฺฉิสฺสนฺติ, โวกฺกมฺมปิ เถรานํ ภิกฺขูนํ ปุรโต ปุรโต คจฺฉิสฺสนฺติ, เถเรปิ ภิกฺขู อนุปขชฺช นิสีทิสฺสนฺติ, นเวปิ ภิกฺขู อาสเนนปิ ปฏิพาหิสฺสนฺติ, สํฆาฏิมฺปิ โอตฺถริตฺวา อนฺตรฆเร นิสีทิสฺสนฺตี”ติ.\csromanspacer{} อถ โข เต ภิกฺขู ภควโต เอตมตฺถํ อาโรเจสุํ ฯเปฯ.\csromanspacer{} สจฺจํ กิร ภิกฺขเว ฉพฺพคฺคิยา ภิกฺขู ทุนฺนิวตฺถา ทุปฺปารุตา อนากปฺปสมฺปนฺนา ภตฺตคฺคํ คจฺฉนฺติ, โวกฺกมฺมปิ เถรานํ ภิกฺขูนํ ปุรโต ปุรโต คจฺฉนฺติ, เถเรปิ ภิกฺขู อนุปขชฺช นิสีทนฺติ, นเวปิ ภิกฺขู อาสเนน ปฏิพาหนฺติ, สํฆาฏิมฺปิ โอตฺถริตฺวา อนฺตรฆเร นิสีทนฺตีติ.\csromanspacer{} สจฺจํ ภควาติ ฯเปฯ วิครหิตฺวา ฯเปฯ ธมฺมิํ กถํ กตฺวา ภิกฺขู อามนฺเตสิ\csromandash{}}

\proseitem{364}{เตน หิ ภิกฺขเว ภิกฺขูนํ ภตฺตคฺควตฺตํ ปญฺญเปสฺสามิ, ยถา ภิกฺขูหิ ภตฺตคฺเค สมฺมา วตฺติตพฺพํ.\csromanspacer{} สเจ อาราเม กาโล อาโรจิโต โหติ, ติมณฺฑลํ ปฏิจฺฉาเทนฺเตน ปริมณฺฑลํ นิวาเสตฺวา กายพนฺธนํ พนฺธิตฺวา สคุณํ กตฺวา สํฆาฏิโย ปารุปิตฺวา คณฺฐิกํ ปฏิมุญฺจิตฺวา โธวิตฺวา ปตฺตํ คเหตฺวา สาธุกํ อตรมาเนน คาโม ปวิสิตพฺโพ.}

\prose{น โวกฺกมฺม เถรานํ ภิกฺขูนํ ปุรโต ปุรโต คนฺตพฺพํ, สุปฺปฏิจฺฉนฺเนน อนฺตรฆเร คนฺตพฺพํ, สุสํวุเตน อนฺตรฆเร คนฺตพฺพํ, โอกฺขิตฺตจกฺขุนา อนฺตรฆเร คนฺตพฺพํ, น อุกฺขิตฺตกาย อนฺตรฆเร คนฺตพฺพํ, น อุชฺชคฺฆิกาย อนฺตรฆเร คนฺตพฺพํ, อปฺปสทฺเทน อนฺตรฆเร คนฺตพฺพํ, น กายปฺปจาลกํ อนฺตรฆเร \csromanfolio{381}คนฺตพฺพํ, น พาหุปฺปจาลกํ อนฺตรฆเร คนฺตพฺพํ, น สีสปฺปจาลกํ อนฺตรฆเร คนฺตพฺพํ, น ขมฺภกเตน อนฺตรฆเร คนฺตพฺพํ, น โอคุณฺฐิเตน อนฺตรฆเร คนฺตพฺพํ, น อุกฺกุฏิกาย อนฺตรฆเร คนฺตพฺพํ.}

\prose{สุปฺปฏิจฺฉนฺเนน อนฺตรฆเร นิสีทิตพฺพํ, สุสํวุเตน อนฺตรฆเร นิสีทิตพฺพํ โอกฺขิตฺตจกฺขุนา อนฺตรฆเร นิสีทิตพฺพํ, น อุกฺขิตฺตกาย อนฺตรฆเร นิสีทิตพฺพํ, น อุชฺชคฺฆิกาย อนฺตรฆเร นิสีทิตพฺพํ, อปฺปสทฺเทน อนฺตรฆเร นิสีทิตพฺพํ, น กายปฺปจาลกํ อนฺตรฆเร นิสีทิตพฺพํ, น พาหุปฺปจาลกํ อนฺตรฆเร นิสีทิตพฺพํ, น สีสปฺปจาลกํ อนฺตรฆเร นิสีทิตพฺพํ, น ขมฺภกเตน อนฺตรฆเร นิสีทิตพฺพํ, น โอคุณฺฐิเตน อนฺตรฆเร นิสีทิตพฺพํ, น ปลฺลตฺถิกาย อนฺตรฆเร นิสีทิตพฺพํ, น เถเร ภิกฺขู อนุปขชฺช นิสีทิตพฺพํ, น นวา ภิกฺขู อาสเนน ปฏิพาหิตพฺพา, น สํฆาฏิํ โอตฺถริตฺวา อนฺตรฆเร นิสีทิตพฺพํ.}

\prose{อุทเก ทิยฺยมาเน อุโภหิ หตฺเถหิ ปตฺตํ ปฏิคฺคเหตฺวา อุทกํ ปฏิคฺคเหตพฺพํ, นีจํ กตฺวา สาธุกํ อปฺปฏิฆํสนฺเตน ปตฺโต โธวิตพฺโพ.\csromanspacer{} สเจ อุทกปฺปฏิคฺคาหโก โหติ, นีจํ กตฺวา อุทกปฺปฏิคฺคเห อุทกํ อาสิญฺจิตพฺพํ “มา อุทกปฺปฏิคฺคาหโก อุทเกน โอสิญฺจิ\footnote{โอสิญฺจิยฺยี (ก)}, มา สามนฺตา ภิกฺขู อุทเกน โอสิญฺจิํสุ\footnote{โอสิญฺจิยฺยิํสุ (ก)}, มา สํฆาฏิ อุทเกน โอสิญฺจี”ติ.\csromanspacer{} สเจ อุทกปฺปฏิคฺคาหโก น โหติ, นีจํ กตฺวา ฉมาย อุทกํ อาสิญฺจิตพฺพํ “มา สามนฺตา ภิกฺขู อุทเกน โอสิญฺจิํสุ, มา สํฆาฏิ อุทเกน โอสิญฺจี”ติ.}

\prose{โอทเน ทิยฺยมาเน อุโภหิ หตฺเถหิ ปตฺตํ ปฏิคฺคเหตฺวา โอทโน ปฏิคฺคเหตพฺโพ, สูปสฺส โอกาโส กาตพฺโพ.\csromanspacer{} สเจ โหติ สปฺปิ วา เตลํ วา อุตฺตริภงฺคํ วา, เถเรน วตฺตพฺโพ “สพฺเพสํ สมกํ สมฺปาเทหี”ติ.\csromanspacer{} สกฺกจฺจํ ปิณฺฑปาโต ปฏิคฺคเหตพฺโพ, ปตฺตสญฺญินา ปิณฺฑปาโต ปฏิคฺคเหตพฺโพ, สมสูปโก ปิณฺฑปาโต ปฏิคฺคเหตพฺโพ, สมติตฺติโก ปิณฺฑปาโต ปฏิคฺคเหตพฺโพ.}

\prose{น ตาว เถเรน ภุญฺชิตพฺพํ, ยาว น สพฺเพสํ โอทโน สมฺปตฺโต โหติ, สกฺกจฺจํ ปิณฺฑปาโต ภุญฺชิตพฺโพ, ปตฺตสญฺญินา ปิณฺฑปาโต ภุญฺชิตพฺโพ, สปทานํ ปิณฺฑปาโต ภุญฺชิตพฺโพ, สมสูปโก ปิณฺฑปาโต \csromanfolio{382}ภุญฺชิตพฺโพ, น ถูปกโต โอมทฺทิตฺวา ปิณฺฑปาโต ภุญฺชิตพฺโพ, น สูปํ วา พฺยญฺชนํ วา โอทเนน ปฏิจฺฉาเทตพฺพํ ภิยฺโยกมฺยตํ อุปาทาย, น สูปํ วา โอทนํ วา อคิลาเนน อตฺตโน อตฺถาย วิญฺญาเปตฺวา ภุญฺชิตพฺพํ, น อุชฺฌานสญฺญินา ปเรสํ ปตฺโต โอโลเกตพฺโพ, นาติมหนฺโต กพโฬ กาตพฺโพ, ปฏิมณฺฑโล อาโลโป กาตพฺโพ, น อนาหเฏ กพเฬ มุขทฺวารํ วิวริตพฺพํ, น ภุญฺชมาเนน สพฺโพ หตฺโถ มุเข ปกฺขิปิตพฺโพ, น สกพเฬน มุเขน พฺยาหริตพฺพํ, น ปิณฺฑุกฺเขปกํ ภุญฺชิตพฺพํ, น กพฬาวจฺเฉทกํ ภุญฺชิตพฺพํ, น อวคณฺฑการกํ ภุญฺชิตพฺพํ, น หตฺถนิทฺธุนกํ ภุญฺชิตพฺพํ, น สิตฺถาวการกํ ภุญฺชิตพฺพํ, น ชีวฺหานิจฺฉารกํ ภุญฺชิตพฺพํ, น จปุจปุการกํ ภุญฺชิตพฺพํ, น สุรุสุรุการกํ ภุญฺชิตพฺพํ, น หตฺถนิลฺเลหกํ ภุญฺชิติพฺพํ, น ปตฺตนิลฺเลหกํ ภุญฺชิตพฺพํ, น โอฏฺฐนิลฺเลหกํ ภุญฺชิตพฺพํ.}

\prose{น สามิเสน หตฺเถน ปานียถาลโก ปฏิคฺคเหตพฺโพ, น ตาว เถเรน อุทกํ ปฏิคฺคเหตพฺพํ, ยาว น สพฺเพว ภุตฺตาวิโน โหนฺติ, อุทเก ทิยฺยมาเน อุโภหิ หตฺเถหิ ปตฺตํ ปฏิคฺคเหตฺวา อุทกํ ปฏิคฺคเหตพฺพํ, นีจํ กตฺวา สาธุกํ อปฺปฏิฆํสนฺเตน ปตฺโต โธวิตพฺโพ.\csromanspacer{} สเจ อุทกปฺปฏิคฺคาหโก โหติ, นีจํ กตฺวา อุทกปฺปฏิคฺคเห อุทกํ อาสิญฺจิตพฺพํ “มา อุทกปฺปฏิคฺคาหโก อุทเกน โอสิญฺจิ, มา สามนฺตา ภิกฺขู อุทเกน โอสิญฺจิํสุ, มา สํฆาฏิ อุทเกน โอสิญฺจี”ติ.\csromanspacer{} สเจ อุทกปฺปฏิคฺคาหโก น โหติ, นีจํ กตฺวา ฉมาย อุทกํ อาสิญฺจิตพฺพํ “มา สามนฺตา ภิกฺขู อุทเกน โอสิญฺจิํสุ, มา สํฆาฏิ อุทเกน โอสิญฺจี”ติ, น สสิตฺถกํ ปตฺตโธวนํ อนฺตรฆเร ฉฑฺเฑตพฺพํ.}

\prose{นิวตฺตนฺเตน นวเกหิ ภิกฺขูหิ ปฐมตรํ นิวตฺติตพฺพํ, ปจฺฉา เถเรหิ, สุปฺปฏิจฺฉนฺเนน อนฺตรฆเร คนฺตพฺพํ, สุสํวุเตน อนฺตรฆเร คนฺตพฺพํ, โอกฺขิตฺตจกฺขุนา อนฺตรฆเร คนฺตพฺพํ, น อุกฺขิตฺตกาย อนฺตรฆเร คนฺตพฺพํ, น อุชฺชคฺฆิกาย อนฺตรฆเร คนฺตพฺพํ, อปฺปสทฺเทน อนฺตรฆเร คนฺตพฺพํ, น กายปฺปจาลกํ อนฺตรฆเร คนฺตพฺพํ, น พาหุปฺปจาลกํ อนฺตรฆเร คนฺตพฺพํ, น สีสปฺปจาลกํ อนฺตรฆเร คนฺตพฺพํ, น ขมฺภกเตน อนฺตรฆเร คนฺตพฺพํ, น โอคุณฺฐิเตน อนฺตรฆเร คนฺตพฺพํ, น อุกฺกุฏิกาย อนฺตรฆเร คนฺตพฺพํ.\csromanspacer{} อิทํ โข ภิกฺขเว ภิกฺขูนํ ภตฺตคฺควตฺตํ, ยถา ภิกฺขูหิ ภตฺตคฺเค สมฺมา วตฺติตพฺพนฺติ.}

\nitthitam{ปฐมภาณวาโร นิฏฺฐิโต.}
\csromanmatikaanchor{mtk.182}
\csromantocmarkref{section}{6. ปิณฺฑจาริกวตฺตกถา}{mtk.182}
\bookmark[dest={mtk.182},level=1]{6. ปิณฺฑจาริกวตฺตกถา}
\csromanheader{1}{6. ปิณฺฑจาริกวตฺตกถา}
\proseitem{365}{\csromanfolio{383}เตน โข ปน สมเยน ปิณฺฑจาริกา ภิกฺขู ทุนฺนิวตฺถา ทุปฺปารุตา อนากปฺปสมฺปนฺนา ปิณฺฑาย จรนฺติ, อสลฺลกฺเขตฺวาปิ นิเวสนํ ปวิสนฺติ, อสลฺลกฺเขตฺวาปิ นิกฺขมนฺติ, อติสหสาปิ ปวิสนฺติ, อติสหสาปิ นิกฺขมนฺติ, อติทูเรปิ ติฏฺฐนฺติ, อจฺจาสนฺเนปิ ติฏฺฐนฺติ, อติจิรมฺปิ ติฏฺฐนฺติ, อติลหุมฺปิ นิวตฺตนฺติ, อญฺญตโรปิ ปิณฺฑจาริโก ภิกฺขุ อสลฺลกฺเขตฺวา นิเวสนํ ปาวิสิ, โส จ ทฺวารํ มญฺญมาโน อญฺญตรํ โอวรกํ ปาวิสิ, ตสฺมิํปิ โอวรเก อิตฺถี นคฺคา อุตฺตานา นิปนฺนา โหติ.\csromanspacer{} อทฺทสา โข โส ภิกฺขุ ตํ อิตฺถิํ นคฺคํ อุตฺตานํ นิปนฺนํ, ทิสฺวาน “นยิทํ ทฺวารํ, โอวรกํ อิทนฺ”ติ ตมฺหา โอวรกา นิกฺขมิ.\csromanspacer{} อทฺทสา โข ตสฺสา อิตฺถิยา สามิโก ตํ อิตฺถิํ นคฺคํ อุตฺตานํ นิปนฺนํ, ทิสฺวาน “อิมินา เม ภิกฺขุนา ปชาปติ ทูสิตา”ติ ตํ ภิกฺขุํ คเหตฺวา อาโกเฏสิ.\csromanspacer{} อถ โข สา อิตฺถี เตน สทฺเทน ปฏิพุชฺฌิตฺวา ตํ ปุริสํ เอตทโวจ “กิสฺส ตฺวํ อยฺย อิมํ ภิกฺขุํ อาโกเฏสี”ติ.\csromanspacer{} อิมินาสิ ตฺวํ ภิกฺขุนา ทูสิตาติ. “นาหํ อยฺย อิมินา ภิกฺขุนา ทูสิตา, อการโก โส ภิกฺขู”ติ ตํ ภิกฺขุํ มุญฺจาเปสิ.\csromanspacer{} อถ โข โส ภิกฺขุ อารามํ คนฺตฺวา ภิกฺขูนํ เอตมตฺถํ อาโรเจสิ.\csromanspacer{} เย เต ภิกฺกู อปฺปิจฺฉา ฯเปฯ เต อุชฺฌายนฺติ ขิยฺยนฺติ วิปาเจนฺติ “กถํ หิ นาม ปิณฺฑจาริกา ภิกฺขู ทุนฺนิวตฺถา ทุปฺปารุตา อนากปฺปสมฺปนฺนา ปิณฺฑาย จริสฺสนฺติ, อสลฺลกฺเขตฺวาปิ นิเวสนํ ปวิสิสฺสนฺติ, อสลฺลกฺเขตฺวาปิ นิกฺขมิสฺสนฺติ, อติสหสาปิ ปวิสิสฺสนฺติ, อติสหสาปิ นิกฺขมิสฺสนฺติ, อติทูเรปิ ติฏฺฐิสฺสนฺติ, อจฺจาสนฺเนปิ ติฏฺฐิสฺสนฺติ, อติจิรมฺปิ ติฏฺฐิสฺสนฺติ, อติลหุมฺปิ นิวตฺติสฺสนฺตี”ติ.\csromanspacer{} อถ โข เต ภิกฺขุ ภควโต เอตมตฺถํ อาโรเจสุํ ฯเปฯ.\csromanspacer{} สจฺจํ กิร ภิกฺขเว ฯเปฯ.\csromanspacer{} สจฺจํ ภควาติ ฯเปฯ วิครหิตฺวา ฯเปฯ ธมฺมิํ กถํ กตฺวา ภิกฺขู อามนฺเตสิ\csromandash{}}

\proseitem{366}{เตน หิ ภิกฺขเว ปิณฺฑจาริกานํ ภิกฺขูนํ วตฺตํ ปญฺญเปสฺสามิ, ยถา ปิณฺฑจาริเกหิ ภิกฺขูหิ สมฺมา วตฺติตพฺพํ.\csromanspacer{} ปิณฺฑจาริเกน ภิกฺขเว ภิกฺขุนา “อิทานิ คามํ ปวิสิสฺสามี”ติ ติมณฺฑลํ ปฏิจฺฉาเทนฺเตน ปริมณฺฑลํ นิวาเสตฺวา กายพนฺธนํ พนฺธิตฺวา สคุณํ กตฺวา สํฆาฏิโย ปารุปิตฺวา คณฺฐิตํ ปฏิมุญฺจิตฺวา โธวิตฺวา ปตฺตํ คเหตฺวา สาธุกํ อตรมาเนน คาโม ปวิสิตพฺโพ.}

\prose{\csromanfolio{384}สุปฺปฏิจฺฉนฺเนน อนฺตรฆเร คนฺตพฺพํ, สุสํวุเตน อนฺตเรฆเร คนฺตพฺพํ, โอกฺขิตฺตจกฺขุนา อนฺตรฆเร คนฺตพฺพํ, น อุกฺขิตฺตกาย อนฺตรฆเร คนฺตพฺพํ, น อุชฺชคฺฆิกาย อนฺตรฆเร คนฺตพฺพํ, อปฺปสทฺเทน อนฺตรฆเร คนฺตพฺพํ, น กายปฺปจาลกํ อนฺตรฆเร คนฺตพฺพํ, น พาหุปฺปจาลกํ อนฺตรฆเร คนฺตพฺพํ, น สีสปฺปจาลกํ อนฺตรฆเร คนฺตพฺพํ, น ขมฺภกเตน อนฺตรฆเร คนฺตพฺพํ, น โอคุณฺฐิเตน อนฺตรฆเร คนฺตพฺพํ, น อุกฺกุฏิกาย อนฺตรฆเร คนฺตพฺพํ.}

\prose{นิเวสนํ ปวิสนฺเตน สลฺเลกฺเขตพฺพํ “อิมินา ปวิสิสฺสมิ, อิมินา นิกฺขมิสฺสามี”ติ, นาติสหสา ปวิสิตพฺพํ, นาติสหสา นิกฺขมิตพฺพํ, นาติทูเร ฐาตพฺพํ, นาจฺจาสนฺเน ฐาตพฺพํ, นาติจิรํ ฐาตพฺพํ, นาติลหุํ นิวตฺติตพฺพํ, ฐิตเกน สลฺลกฺเขตพฺพํ “ภิกฺขํ ทาตุกามา วา อทาตุกามา วา”ติ.\csromanspacer{} สเจ กมฺมํ วา นิกฺขิปติ, อาสนา วา วุฏฺฐาติ, กฏจฺฉุํ วา ปรามสติ, ภาชนํ วา ปรามสติ, ฐเปติ\footnote{ฐาเปติ (ก)} วา, “ทาตุกามสฺสา”ติ\footnote{ทาตุกามิยาติ (สฺยา), ทาตุกามา วิยาติ (สี)} ฐาตพฺพํ, ภิกฺขาย ทิยฺยมานาย วาเมน หตฺเถน สํฆาฏิํ อุจฺจาเรตฺวา ทกฺขิเณน หตฺเถน ปตฺตํ ปณาเมตฺวา อุโภหิ หตฺเถหิ ปตฺตํ ปฏิคฺคเหตฺวา ภิกฺขา ปฏิคฺคเหตพฺพา, น จ ภิกฺขาทายิกาย มุขํ อุลฺโลเกตพฺพํ\footnote{โอโลเกตพฺพํ (สฺยา)}, สลฺลกฺเขตพฺพํ “สูปํ ทาตุกามา วา อทาตุกามา วา”ติ.\csromanspacer{} สเจ กฏจฺฉุํ วา ปรามสติ, ภาชนํ วา ปรามสติ, ฐเปติ วา, “ทาตุกามสฺสา”ติ ฐาตพฺพํ, ภิกฺขาย ทินฺนาย สํฆาฏิยา ปตฺตํ ปฏิจฺฉาเทตฺวา สาธุกํ อตรมาเนน นิวตฺติตพฺพํ.}

\prose{สุปฺปฏิจฺฉนฺเนน อนฺตรฆเร คนฺตพฺพํ, สุสํวุเตน อนฺตรฆเร คนฺตพฺพํ, โอกฺขิตฺตจกฺขุนา อนฺตรฆเร คนฺตพฺพํ, น อุกฺขิตฺตกาย อนฺตรฆเร คนฺตพฺพํ, น อุชฺชคฺฆิกาย อนฺตรฆเร คนฺตพฺพํ, อปฺปสทฺเทน อนฺตรฆเร คนฺตพฺพํ, น กายปฺปจาลกํ อนฺตรฆเร คนฺตพฺพํ, น พาหุปฺปจาลกํ อนฺตรฆเร คนฺตพฺพํ, น สีสปฺปจาลกํ อนฺตรฆเร คนฺตพฺพํ, น ขมฺภกเตน อนฺตรฆเร คนฺตพฺพํ, น โอคุณฺฐิเตน อนฺตรฆเร คนฺตพฺพํ, น อุกฺกุฏิกาย อนฺตรฆเร คนฺตพฺพํ, โย ปฐมํ คามโต ปิณฺฑาย ปฏิกฺกมติ, เตน อาสนํ ปญฺญเปตพฺพํ, ปาโททกํ ปาทปีฐํ ปาทกถลิกํ อุปนิกฺขิปิตพฺพํ, อวกฺการปาติ โธวิตฺวา อุปฏฺฐาเปตพฺพา, ปานียํ ปริโภชนียํ อุปฏฺฐาเปตพฺพํ, โย ปจฺฉา คามโต ปิณฺฑาย ปฏิกฺกมติ, สเจ โหติ ภุตฺตาวเสโส, สเจ อากงฺขติ, \csromanfolio{385}ภุญฺชิตพฺพํ, โน เจ อากงฺขติ, อปฺปหริเต วา ฉฑฺเฑตพฺพํ, อปฺปาณเก วา อุทเก โอปิลาเปตพฺพํ, เตน อาสนํ อุทฺธริตพฺพํ, ปาโททกํ ปาทปีฐํ ปาทกถลิกํ ปฏิสาเมตพฺพํ, อวกฺการปาติ โธวิตฺวา ปฏิสาเมตพฺพา, ปานียํ ปริโภชนียํ ปฏิสาเมตพฺพํ, ภตฺตคฺคํ สมฺมชฺชิตพฺพํ, โย ปสฺสติ ปานียฆฏํ วา ปริโภชนียฆฏํ วา วจฺจฆฏํ วา ริตฺตํ ตุจฺฉํ, เตน อุปฏฺฐาเปตพฺพํ.\csromanspacer{} สจสฺส โหติ, อวิสยฺหํ, หตฺถวิกาเรน ทุติยํ อามนฺเตตฺวา หตฺถวิลงฺฆเกน อุปฏฺฐาเปตพฺพํ, น จ ตปฺปจฺจยา วาจา ภินฺทิตพฺพา.\csromanspacer{} อิทํ โข ภิกฺขเว ปิณฺฑจาริกานํ ภิกฺขูนํ วตฺตํ, ยถา ปิณฺฑจาริเกหิ ภิกฺขูหิ สมฺมา วตฺติตพฺพนฺติ.}

\csromanmatikaanchor{mtk.183}
\csromantocmarkref{section}{7. อารญฺญิกวตฺตกถา}{mtk.183}
\bookmark[dest={mtk.183},level=1]{7. อารญฺญิกวตฺตกถา}
\csromanheader{1}{7. อารญฺญิกวตฺตกถา}
\proseitem{367}{เตน โข ปน สมเยน สมฺพหุลา ภิกฺขู อรญฺเญ วิหรนฺติ, เต เนว ปานียํ อุปฏฺฐาเปนฺติ, น ปริโภชนียํ อุปฏฺฐาเปนฺติ, น อคฺคิํ อุปฏฺฐาเปนฺติ, น อรณิสหิตํ อุปฏฺฐาเปนฺติ, น นกฺขตฺตปทานิ ชานนฺติ, น ทิสาภาคํ ชานนฺติ.\csromanspacer{} โจรา ตตฺถ คนฺตฺวา เต ภิกฺขู เอตทโวจุํ “อตฺถิ ภนฺเต ปานียนฺ”ติ.\csromanspacer{} นตฺถาวุโสติ.\csromanspacer{} อตฺถิ ภนฺเต ปริโภชนียนฺติ.\csromanspacer{} นตฺตาวุโสติ.\csromanspacer{} อตฺถิ ภนฺเต อคฺคีติ.\csromanspacer{} นตฺถาวุโสติ.\csromanspacer{} อตฺถิ ภนฺเต อรณิสหิตนฺติ.\csromanspacer{} นตฺถาวุโสติ. ( )\footnote{(อตฺถิ ภนฺเต นกฺขตฺตปทานีติ, น ชานาม อาวุโสติ, อตฺถิ ภนฺเต ทิสาภาคนฺติ, น ชานาม อาวุโสติ.) สี, วิมติฏีกาย ปน สเมติ.} เกนชฺช ภนฺเต ยุตฺตนฺติ.\csromanspacer{} น โข มยํ อาวุโส ชานามาติ.\csromanspacer{} กตมายํ ภนฺเต ทิสาติ.\csromanspacer{} น โข มยํ อาวุโส ชานามาติ.\csromanspacer{} อถ โข เต โจรา “เนวิเมสํ ปานียํ อตฺถิ, น ปริโภชนียํ อตฺถิ, น อคฺคิ อตฺถิ, น อรณิสหิตํ อตฺถิ, น นกฺขตฺตปทานิ ชานนฺติ, น ทิสาภาคํ ชานนฺติ, โจรา อิเม, นยิเม ภิกฺขู”ติ อาโกเฏตฺวา ปกฺกมิํสุ.\csromanspacer{} อถ โข เต ภิกฺขู ภิกฺขูนํ เอตมตฺถํ อาโรเจสุํ.\csromanspacer{} ภิกฺขู ภควโต เอตมตฺถํ อาโรเจสุํ.\csromanspacer{} อถ โข ภควา เอตสฺมิํ นิทาเน เอตสฺมิํ ปกรเณ ธมฺมิํ กถํ กตฺวา ภิกฺขู อามนฺเตสิ\csromandash{}}

\proseitem{368}{เตน หิ ภิกฺขเว อารญฺญิกานํ ภิกฺขูนํ วตฺตํ ปญฺญเปสฺสามิ, ยถา อารญฺญิเกหิ ภิกฺขูหิ สมฺมา วตฺติตพฺพํ.\csromanspacer{} อารญฺญิเกน ภิกฺขเว ภิกฺขุนา กาลสฺเสว อุฏฺฐาย ปตฺตํ ถวิกาย ปกฺขิปิตฺวา อํเส \csromanfolio{386}อาลคฺเคตฺวา จีวรํ ขนฺเธ กริตฺวา อุปาหนา อาโรหิตฺวา ทารุภณฺฑํ มตฺติกาภณฺฑํ ปฏิสาเมตฺวา ทฺวารวาตปานํ ถเกตฺวา เสนาสนา โอตริตพฺพํ. “อิทานิ คามํ ปวิสิสฺสามี”ติ อุปาหนา โอมุญฺจิตฺวา นีจํ กตฺวา ปปฺโผเฏตฺวา ถวิกาย ปกฺขิปิตฺวา อํเส อาลคฺเคตฺวา ติมณฺฑลํ ปฏิจฺฉาเทนฺเตน ปริมณฺฑลํ นิวาเสตฺวา กายพนฺธนํ พนฺธิตฺวา สคุณํ กตฺวา สํฆาฏิโย ปารุปิตฺวา คณฺฐิกํ ปฏิมุญฺจิตฺวา โธวิตฺวา ปตฺตํ คเหตฺวา สาธุกํ อตรมาเนน คาโม ปวิสิตพฺโพ, สุปฺปฏิจฺฉนฺเนน อนฺตรฆเร คนฺตพฺพํ ฯเปฯ น ขมฺภกเตน อนฺตรฆเร คนฺตพฺพํ, น โอคุณฺฐิเตน อนฺตรฆเร คนฺตพฺพํ, น อุกฺกุฏิกาย อนฺตรฆเร คนฺตพฺพํ.}

\prose{นิเวสนํ ปวิสนฺเตน สลฺลกฺเขตพฺพํ “อิมินา ปวิสิสฺสามิ, อิมินา นิกฺขมิสฺสามี”ติ.\csromanspacer{} นาติสหสา ปวิสิตพฺพํ, นาติสหสา นิกฺขมิตพฺพํ, นาติทูเร ฐาตพฺพํ, นาจฺจาสนฺเน ฐาตพฺพํ, นาติจิรํ ฐาตพฺพํ, นาติลหุํ นิวตฺติตพฺพํ, ฐิตเกน สลฺลกฺเขตพฺพํ “ภิกฺขํ ทาตุกามา วา อทาตุกามาวา”ติ.\csromanspacer{} สเจ กมฺมํ วา นิกฺขิปติ, อาสนา วา วุฏฺฐาติ, กฏจฺฉุํ วา ปรามสติ, ภาชนํ วา ปรามสติ, ฐเปติ วา, “ทาตุกามสฺสา”ติ ฐาตพฺพํ, ภิกฺขาย ทิยฺยมานาย วาเมน หตฺเถน สํฆาฏิํ อุจฺจาเรตฺวา ทกฺขิเณน หตฺเถน ปตฺตํ ปณาเมตฺวา อุโภหิ หตฺเถหิ ปตฺตํ ปฏิคฺคเหตฺวา ภิกฺขา ปฏิคฺคเหตพฺพา, น จ ภิกฺขาทายิกาย มุขํ อุลฺโลเกตพฺพํ, สลฺลกฺเขตพฺพํ “สูปํ ทาตุกามา วา อทาตุกามา วา”ติ.\csromanspacer{} สเจ กฏจฺฉุ วา ปรามสติ, ภาชนํ วา ปรามสติ, ฐเปติ วา, “ทาตุกามสฺสา”ติ ฐาตพฺพํ.\csromanspacer{} ภิกฺขาย ทินฺนาย สํฆาฏิยา ปตฺตํ ปฏิจฺฉาเทตฺวา สาธุกํ อตรมาเนน นิวตฺติตพฺพํ.}

\prose{สุปฺปฏิจฺฉนฺเนน อนฺตรฆเร คนฺตพฺพํ ฯเปฯ น อุกฺกุฏิกาย อนฺตรฆเร คนฺตพฺพํ.\csromanspacer{} คามโต นิกฺขมิตฺวา ปตฺตํ ถวิกาย ปกฺขิปิตฺวา อํเส อาลคฺเคตฺวา จีวรํ สงฺฆริตฺวา สีเส กริตฺวา อุปาหนา อาโรหิตฺวา คนฺตพฺพํ.}

\prose{อารญฺญิเกน ภิกฺขเว ภิกฺขุนา ปานียํ อุปฏฺฐาเปตพฺพํ, ปริโภชนียํ อุปฏฺฐาเปตพฺพํ, อคฺคิ อุปฏฺฐาเปตพฺโพ, อรณิสหิตํ อุปฏฺฐาเปตพฺพํ, กตฺตรทณฺโฑ อุปฏฺฐาเปตพฺโพ, นกฺขตฺตปทานิ อุคฺคเหตพฺพานิ สกลานิ วา เอกเทสานิ วา, ทิสากุสเลน ภวิตพฺพํ.\csromanspacer{} อิทํ โข ภิกฺขเว อารญฺญิกานํ ภิกฺขูนํ วตฺตํ, ยถา อารญฺญิเกหิ ภิกฺขูหิ สมฺมา วตฺติตพฺพนฺติ.}

\csromanmatikaanchor{mtk.184}
\csromantocmarkref{section}{8. เสนาสนวตฺตกถา}{mtk.184}
\bookmark[dest={mtk.184},level=1]{8. เสนาสนวตฺตกถา}
\csromanheader{1}{8. เสนาสนวตฺตกถา}
\proseitem{369}{\csromanfolio{387}เตน โข ปน สมเยน สมฺพหุลา ภิกฺขู อชฺโฌกาเส จีวรกมฺมํ กโรนฺติ.\csromanspacer{} ฉพฺพคฺคิยา ภิกฺขู ปฏิวาเต องฺคเณ\footnote{ปงฺคเณ (สี, สฺยา)} เสนาสนํ ปปฺโผเฏสุํ, ภิกฺขู รเชน โอกิริํสุ.\csromanspacer{} เย เต ภิกฺขู อปฺปิจฺฉา ฯเปฯ เต อุชฺฌายนฺติ ขิยฺยนฺติ วิปาเจนฺติ “กถํ หิ นาม ฉพฺพคฺคิยา ภิกฺขู ปฏิวาเต องฺคเณ เสนาสนํ ปปฺโผเฏสฺสนฺติ, ภิกฺขู รเชน โอกิริํสู”ติ.\csromanspacer{} อถ โข เต ภิกฺขู ภควโต เอตมตฺถํ อาโรเจสุํ ฯเปฯ.\csromanspacer{} สจฺจํ กิร ภิกฺขเว ฉพฺพคฺคิยา ภิกฺขู ปฏิวาเต องฺคเณ เสนาสนํ ปปฺโผเฏนฺติ, ภิกฺขู รเชน โอกิริํสูติ.\csromanspacer{} สจฺจํ ภควาติ ฯเปฯ วิครหิตฺวา ฯเปฯ ธมฺมิํ กตฺวา ภิกฺขู อามนฺเตสิ\csromandash{}}

\proseitem{370}{เตน หิ ภิกฺขเว ภิกฺขูนํ เสนาสนวตฺตํ ปญฺญเปสฺสามิ, ยถา ภิกฺขูหิ เสนาสเน สมฺมา วตฺติตพฺพํ.\csromanspacer{} ยสฺมิํ วิหาเร วิหรติ, สเจ โส วิหาโร อุกฺลาโป โหติ, สเจ อุสฺสหติ โสเธตพฺโพ, * วิหารํ โสเธนฺเตน ปฐมํ ปตฺตจีวรํ นีหริตฺวา เอกมนฺตํ นิกฺขิปิตพฺพํ, นิสีทนปจฺจตฺถรณํ นีหริตฺวา เอกมนฺตํ นิกฺขิปิตพฺพํ, ภิสิพิพฺโพหนํ นีหริตฺวา เอกมนฺตํ นิกฺขิปิตพฺพํ, มญฺโจ นีจํ กตฺวา สาธุกํ อปฺปฏิฆสนฺเตน อสํฆฏฺเฏนฺเตน กวาฏปิฏฺฐํ นีหริตฺวา เอกมนฺตํ นิกฺขิปิตพฺโพ, ปีฐํ นีจํ กตฺวา สาธุกํ อปฺปฏิฆํสนฺเตน อสํฆฏฺเฏนฺเตน กวาฏปิฏฺฐํ นีหริตฺวา เอกมนฺตํ นิกฺขิปิตพฺพํ, มญฺจปฏิปาทกา นีหริตฺวา เอกมนฺตํ นิกฺขิปิตพฺพา, เขฬมลฺลโก นีหริตฺวา เอกมนฺตํ นิกฺขิปิตพฺโพ, อปสฺเสนผลกํ นีหริตฺวา เอกมนฺตํ นิกฺขิปิตพฺพํ, ภูมตฺถรณํ ยถาปญฺญตฺตํ สลฺลกฺเขตฺวา นีหริตฺวา เอกมนฺตํ นิกฺขิปิตพฺพํ.\csromanspacer{} สเจ วิหาเร สนฺตานกํ โหติ, อุลฺโลกา ปฐมํ โอหาเรตพฺพํ, อาโลกสนฺธิกณฺณภาคา ปมชฺชิตพฺพา.\csromanspacer{} สเจ เครุกปริกมฺมกตา ภิตฺติ กณฺณกิตา โหติ, โจฬกํ เตเมตฺวา ปีเฬตฺวา ปมชฺชิตพฺพา.\csromanspacer{} สเจ กาฬวณฺณกตา ภูมิ กณฺณกิตา โหติ, โจฬกํ เตเมตฺวา ปีเฬตฺวา ปมชฺชิตพฺพา.\csromanspacer{} สเจ อกตา โหติ ภูมิ, อุทเกน ปริปฺโผสิตฺวา ปริปฺโผสิตฺวา สมฺมชฺชิตพฺพา “มา วิหาโร รเชน อุหญฺญี”ติ, สงฺการํ วิจินิตฺวา เอกมนฺตํ ฉฑฺเฑตพฺพํ.}

\prose{\csromanfolio{388}น ภิกฺขุสามนฺตา เสนาสนํ ปปฺโผเฏตพฺพํ, น วิหารสามนฺตา เสนาสนํ ปปฺโผเฏตพฺพํ, น ปานียสามนฺตา เสนาสนํ ปปฺโผเฏตพฺพํ, น ปริโภชนียสามนฺตา เสนาสนํ ปปฺโผเฏตพฺพํ, น ปฏิวาเต องฺคเณ เสนาสนํ ปปฺโผเฏตพฺพํ, อโธวาเต เสนาสนํ ปปฺโผเฏตพฺพํ.}

\prose{ภูมตฺถรณํ เอกมนฺตํ โอตาเปตฺวา โสเธตฺวา ปปฺโผเฏตฺวา อติหริตฺวา ยถาปญฺญตฺตํ, ปญฺญเปตพฺพํ, มญฺจปฏิปาทกา เอกมนฺตํ โอตาเปตฺวา ปมชฺชิตฺวา อติหริตฺวา ยถาฐาเน ฐเปตพฺพา, มญฺโจ เอกมนฺตํ โอตาเปตฺวา โสเธตฺวา ปปฺโผเฏตฺวา นีจํ กตฺวา สาธุกํ อปฺปฏิฆํสนฺเตน อสํฆฏฺเฏนฺเตน กวาฏปิฏฺฐํ อติหริตฺวา ยถาปญฺญตฺตํ ปญฺญเปตพฺโพ, ปีฐํ เอกมนฺตํ โอตาเปตฺวา โสเธตฺวา ปปฺโผเฏตฺวา นีจํ กตฺวา สาธุกํ อปฺปฏิฆํสนฺเตน อสํฆฏฺเฏนฺเตน กวาฏปิฏฺฐํ อติหริตฺวา ยถาปญฺญตฺตํ ปญฺญเปตพฺพํ, ภิสิพิพฺโพหนํ เอกมนฺตํ โอตาเปตฺวา โสเธตฺวา ปปฺโผเฏตฺวา อติหริตฺวา ยถาปญฺญตฺตํ ปญฺญเปตพฺพํ, นิสีทนปจฺจตฺถรณํ เอกมนฺตํ โอตาเปตฺวา โสเธตฺวา ปปฺโผเฏตฺวา อติหริตฺวา ยถาปญฺญตฺตํ ปญฺญเปตพฺพํ, เขฬมลฺลโก เอกมนฺตํ โอตาเปตฺวา ปมชฺชิตฺวา อติหริตฺวา ยถาฐาเน ฐเปตพฺโพ, อปสฺเสนผลกํ เอกมนฺตํ โอตาเปตฺวา ปมชฺชิตฺวา อติหริตฺวา ยถาฐาเน ฐเปตพฺพํ.\csromanspacer{} ปตฺตจีวรํ นิกฺขิปิตพฺพํ, ปตฺตํ นิกฺขิปนฺเตน เอเกน หตฺเถน ปตฺตํ คเหตฺวา เอเกน หตฺเถน เหฏฺฐามญฺจํ วา เหฏฺฐาปีฐํ วา ปรามสิตฺวา ปตฺโต นิกฺขิปิตพฺโพ, น จ อนนฺตรหิตาย ภูมิยา ปตฺโต นิกฺขิปิตพฺโพ.\csromanspacer{} จีวรํ นิกฺขิปนฺเตน เอเกน หตฺเถน จีวรํ คเหตฺวา เอเกน หตฺเถน จีวรวํสํ วา จีวรรชฺชุํ วา ปมชฺชิตฺวา ปารโต อนฺตํ โอรโต โภคํ กตฺวา จีวรํ นิกฺขิปิตพฺพํ.}

\prose{สเจ ปุรตฺถิมา สรชา วาตา วายนฺติ, ปุรตฺถิมา วาตปานา ถเกตพฺพา.\csromanspacer{} สเจ ปจฺฉิมา สรชา วาตา วายนฺติ, ปจฺฉิมา วาตปานา ถเกตพฺพา.\csromanspacer{} สเจ อุตฺตรา สรชา วาตา วายนฺติ, อุตฺตรา วาตปานา ถเกตพฺพา.\csromanspacer{} สเจ ทกฺขิณา สรชา วาตา วายนฺติ, ทกฺขิณา วาตปานา ถเกตพฺพา.\csromanspacer{} สเจ สีตกาโล โหติ, ทิวา วาตปานา วิวริตพฺพา รตฺติํ ถเกตพฺพา.\csromanspacer{} สเจ อุณฺหกาโล โหติ, ทิวา วาตปานา ถเกตพฺพา, รตฺติํ วิวริตพฺพา.}

\prose{\csromanfolio{389}สเจ ปริเวณํ อุกฺลาปํ โหติ, ปริเวณํ สมฺมชฺชิตพฺพํ.\csromanspacer{} สเจ โกฏฺฐโก อุกฺลาโป โหติ, โกฏฺฐโก สมฺมชฺชิตพฺโพ.\csromanspacer{} สเจ อุปฏฺฐานสาลา อุกฺลาปา โหติ, อุปฏฺฐานสาลา สมฺมชฺชิตพฺโพ.\csromanspacer{} สเจ อคฺคิสาลา อุกฺลาปา โหติ, อคฺคิสาลา สมฺมชฺชิตพฺพา.\csromanspacer{} สเจ วจฺจกุฏิ อุกฺลาปา โหติ, วจฺจกุฏิ สมฺมชฺชิตพฺพา.\csromanspacer{} สเจ ปานียํ น โหติ, ปานียํ อุปฏฺฐาเปตพฺพํ.\csromanspacer{} สเจ ปริโภชนียํ น โหติ, ปริโภชนียํ อุปฏฺฐาเปตพฺพํ.\csromanspacer{} สเจ อาจมนกุมฺภิยา อุทกํ น โหติ, อาจมนกุมฺภิยา อุทกํ อาสิญฺจิตพฺพํ.}

\prose{สเจ วุฑฺเฒน สทฺธิํ เอกวิหาเร วิหรติ, น วุฑฺฒํ อนาปุจฺฉา อุทฺเทโส ทาตพฺโพ, น ปริปุจฺฉา ทาตพฺพา, น สชฺฌาโย กาตพฺโพ, น ธมฺโม ภาสิตพฺโพ, น ปทีโป กาตพฺโพ, น ปทีโป วิชฺฌาเปตพฺโพ, น วาตปานา วิวริตพฺพา, น วาตปานา ถเกตพฺพา.\csromanspacer{} สเจ วุฑฺเฒน สทฺธิํ เอกจงฺกเม จงฺกมติ, เยน วุฑฺโฒ เตน ปริวตฺติตพฺพํ, น จ วุฑฺโฒ สํฆาฏิกณฺเณน ฆฏฺเฏตพฺโพ.\csromanspacer{} อิทํ โข ภิกฺขเว ภิกฺขูนํ เสนาสนวตฺตํ, ยถา ภิกฺขูหิ เสนาสเน สมฺมา วตฺติตพฺพนฺติ.}

\csromanmatikaanchor{mtk.185}
\csromantocmarkref{section}{9. ชนฺตาฆรวตฺตกถา}{mtk.185}
\bookmark[dest={mtk.185},level=1]{9. ชนฺตาฆรวตฺตกถา}
\csromanheader{1}{9. ชนฺตาฆรวตฺตกถา}
\csromanmatikaanchor{mtk.186}
\csromantocmarkref{section}{10. วจฺจกุฏิวตฺตกถา}{mtk.186}
\bookmark[dest={mtk.186},level=1]{10. วจฺจกุฏิวตฺตกถา}
\proseitem{371}{เตน โข ปน สมเยน ฉพฺพคฺคิยา ภิกฺขู ชนฺตาฆเร เถเรหิ ภิกฺขูหิ นิวาริยมานา อนาทริยํ ปฏิจฺจ ปหูตํ กฏฺฐํ อาโรเปตฺวา อคฺคิํ ทตฺวา ทฺวารํ ถเกตฺวา ทฺวาเร นิสีทนฺติ, ภิกฺขู\footnote{เถรา จ ภิกฺขู (สฺยา, กํ)} อุณฺหาภิตตฺตา ทฺวารํ อลภมานา มุจฺฉิตา ปปตนฺติ.\csromanspacer{} เย เต ภิกฺขู อปฺปิจฺฉา ฯเปฯ เต อุชฺฌายนฺติ ขิยฺยนฺติ วิปาเจนฺติ “กถํ หิ นาม ฉพฺพคฺคิยา ภิกฺขู ชนฺตาฆเร เถเรหิ ภิกฺขูหิ นิวาริยมานา อนาทริยํ ปฏิจฺจ ปหูตํ กฏฺฐํ อาโรเปตฺวา อคฺคิํ ทตฺวา ทฺวารํ ถเกตฺวา ทฺวาเร นิสีทิสฺสนฺติ, ภิกฺขู อุณฺหาภิตตฺตา ทฺวารํ อลภมานา มุจฺฉิตา ปปตนฺตี”ติ.\csromanspacer{} อถ โข เต ภิกฺขู ภควโต เอตมตฺถํ อาโรเจสุํ ฯเปฯ.\csromanspacer{} สจฺจํ กิร ภิกฺขเว ฉพฺพคฺคิยา ภิกฺขู ชนฺตาฆเร เถเรหิ ภิกฺขูหิ นิวาริยมานา อนาทริยํ ปฏิจฺจ ปหูตํ กฏฺฐํ อาโรเปตฺวา อคฺคิํ ทตฺวา ทฺวารํ ถเกตฺวา ทฺวาเร นิสีทนฺติ, ภิกฺขู อุณฺหาภิตตฺตา ทฺวารํ อลภมานา มุจฺฉิตา ปปตนฺตีติ.\csromanspacer{} สจฺจํ ภควาติ ฯเปฯ วิครหิตฺวา ฯเปฯ ธมฺมิํ กถํ กตฺวา ภิกฺขู อามนฺเตสิ\csromandash{} \csromanfolio{390}น ภิกฺขเว ชนฺตาฆเร เถเรน ภิกฺขุนา นิวาริยมาเนน อนาทริยํ ปฏิจฺจ ปหูตํ กฏฺฐํ อาโรเปตฺวา อคฺคิ ทาตพฺโพ, โย ทเทยฺย, อาปตฺติ ทุกฺกฏสฺส.\csromanspacer{} น ภิกฺขเว ทฺวารํ ถเกตฺวา ทฺวาเร นิสีทิตพฺพํ, โย นิสีเทยฺย, อาปตฺติ ทุกฺกฏสฺส.}

\proseitem{372}{เตน หิ ภิกฺขเว ภิกฺขูนํ ชนฺตาฆรวตฺตํ ปญฺญเปสฺสามิ, ยถา ภิกฺขูหิ ชนฺตาฆเร สมฺมา วตฺติตพฺพํ.\csromanspacer{} โย ปฐมํ ชนฺตาฆรํ คจฺฉติ, สเจ ฉาริกา อุสฺสนฺนา โหติ, ฉาริกา ฉฑฺเฑตพฺพา.\csromanspacer{} สเจ ชนฺตาฆรํ อุกฺลาปํ โหติ, ชนฺตาฆรํ สมฺมชฺชิตพฺพํ.\csromanspacer{} สเจ ปริภณฺฑํ อุกฺลาปํ โหติ, ปริภณฺฑํ สมฺมชฺชิตพฺพํ.\csromanspacer{} สเจ ปริเวณํ อุกฺลาปํ โหติ, ปริเวณํ สมฺมชฺชิตพฺพํ.\csromanspacer{} สเจ โกฏฺฐโก อุกฺลาโป โหติ, โกฏฺฐโก สมฺมชฺชิตพฺโพ.\csromanspacer{} สเจ ชนฺตาฆรสาลา อุกฺลาปา โหติ, ชนฺตาฆรสาลา สมฺมชฺชิตพฺพา.}

\prose{จุณฺณํ สนฺเนตพฺพํ, มตฺติกา เตเมตพฺพา, อุทกโทณิกาย อุทกํ อาสิญฺจิตพฺพํ, ชนฺตาฆรํ ปวิสนฺเตน มตฺติกาย มุขํ มกฺเขตฺวา ปุรโต จ ปจฺฉโต จ ปฏิจฺฉาเทตฺวา ชนฺตาฆรํ ปวิสิตพฺพํ, น เถเร ภิกฺขู อนุปขชฺช นิสีทิตพฺพํ, น นวา ภิกฺขู อาสเนน ปฏิพาหิตพฺพา.\csromanspacer{} สเจ อุสฺสหติ, ชนฺตาฆเร เถรานํ ภิกฺขุนํ ปริกมฺมํ กาตพฺพํ, ชนฺตาฆรา นิกฺขมนฺเตน ชนฺตาฆรปีฐํ อาทาย ปุรโต จ ปจฺฉโต จ ปฏิจฺฉาเทตฺวา ชนฺตาฆรา นิกฺขมิตพฺพํ.\csromanspacer{} สเจ อุสฺสหติ, อุทเกปิ เถรานํ ภิกฺขูนํ ปริกมฺมํ กาตพฺพํ, น เถรานํ ภิกฺขูนํ ปุรโตปิ นหายิตพฺพํ, น อุปริโตปิ นหายิตพฺพํ, นหาเตน อุตฺตรนฺเตน โอตรนฺตานํ มคฺโค ทาตพฺโพ.\csromanspacer{} โย ปจฺฉา ชนฺตาฆรา นิกฺขมติ, สเจ ชนฺตาฆรํ จิกฺขลฺลํ โหติ, โธวิตพฺพํ, มตฺติกาโทณิกํ โธวิตฺวา ชนฺตาฆรปีฐํ ปฏิสาเมตฺวา อคฺคิํ วิชฺฌาเปตฺวา ทฺวารํ ถเกตฺวา ปกฺกมิตพฺพํ.\csromanspacer{} อิทํ โข ติกฺขเว ภิกฺขูนํ ชนฺตาฆรวตฺตํ, ยถา ภิกฺขูหิ ชนฺตาฆเร สมฺมา วตฺติตพฺพนฺติ.}

\csromanheader{2}{10. วจฺจกุฏิวตฺตตถา}
\proseitem{373}{เตน โข ปน สมเยน อญฺญตโร ภิกฺขุ พฺราหฺมณชาติโก วจฺจํ กตฺวา น อิจฺฉติ อาจเมตุํ “โก อิมํ วสลํ ทุคฺคนฺธํ \csromanfolio{391}อามสิสฺสตี”ติ\footnote{อาจมิสฺสตีติ (ก)}.\csromanspacer{} ตสฺส วจฺจมคฺเค กิมิ สณฺฐาติ.\csromanspacer{} อถ โข โส ภิกฺขุ ภิกฺขูนํ เอตมตฺถํ อาโรเจสิ.\csromanspacer{} กิํ ปน ตฺวํ อาวุโส วจฺจํ กตฺวา น อาจเมสีติ.\csromanspacer{} เอวมาวุโสติ.\csromanspacer{} เย เต ภิกฺขู อปฺปิจฺฉา ฯเปฯ เต อุชฺฌายนฺติ ขิยฺยนฺติ วิปาเจนฺติ “กถํ หิ นาม ภิกฺขุ วจฺจํ กตฺวา น อาจเมสฺสตี”ติ.\csromanspacer{} อถ โข เต ภิกฺขู ภควโต เอตมตฺถํ อาโรเจสุํ ฯเปฯ.\csromanspacer{} สจฺจํ กิร ตฺวํ ภิกฺขุ วจฺจํ กตฺวา น อาจเมสีติ.\csromanspacer{} สจฺจํ ภควาติ ฯเปฯ วิครหิตฺวา ฯเปฯ ธมฺมิํ กถํ กตฺวา ภิกฺขู อามนฺเตสิ “น ภิกฺขเว วจฺจํ กตฺวา สติ อุทเก นาจเมตพฺพํ, โย นาจเมยฺย, อาปตฺติ ทุกฺกฏสฺสา”ติ.}

\prose{เตน โข ปน สมเยน ภิกฺขู วจฺจกุฏิยา ยถาวุฑฺฒํ วจฺจํ กโรนฺติ, นวกา ภิกฺขู ปฐมตรํ อาคนฺตฺวา วจฺจิตา อาคเมนฺติ, เต วจฺจํ สนฺธาเรนฺตา มุจฺฉิตา ปปตนฺติ.\csromanspacer{} ภควโต เอตมตฺถํ อาโรเจสุํ ฯเปฯ.\csromanspacer{} สจฺจํ กิร ภิกฺขเว ฯเปฯ.\csromanspacer{} สจฺจํ ภควาติ ฯเปฯ.\csromanspacer{} น ภิกฺขเว วจฺจกุฏิยา ยถาวุฑฺฒํ วจฺโจ กาตพฺโพ, โย กเรยฺย อาปตฺติ ทุกฺกฏสฺส.\csromanspacer{} อนุชานามิ ภิกฺขเว อาคตปฏิปาฏิยา วจฺจํ กาตุนฺติ.}

\prose{เตน โข ปน สมเยน ฉพฺพคฺคิยา ภิกฺขู อติสหสาปิ วจฺจกุฏิํ ปวิสนฺติ, อุพฺภชิตฺวาปิ\footnote{อุพฺภุชฺชิตฺวาปิ (สี), อุพฺภุชิตฺวา (สฺยา)} ปวิสนฺติ, นิตฺถุนนฺตาปิ วจฺจํ กโรนฺติ, ทนฺตกฏฺฐํ ขาทนฺตาปิ วจฺจํ กโรนฺติ, พหิทฺธาปิ วจฺจโทณิกาย วจฺจํ กโรนฺติ, พหิทฺธาปิ ปสฺสาวโทณิกาย ปสฺสาวํ กโรนฺติ, ปสฺสาวโทณิกายปิ เขฬํ กโรนฺติ, ผรุเสนปิ กฏฺเฐน อวเลขนฺติ, อวเลขนกฏฺฐมฺปิ วจฺจกูปมฺหิ ปาเตนฺติ, อติสหสาปิ นิกฺขมนฺติ, อุพฺภชิตฺวาปิ นิกฺขมนฺติ, จปุจปุการกมฺปิ อาจเมนฺติ, อาจมนสราวเกปิ อุทกํ เสเสนฺติ.\csromanspacer{} เย เต ภิกฺขู อปฺปิจฺฉา ฯเปฯ เต อุชฺฌายนฺติ ขิยฺยนฺติ วิปาเจนฺติ “กถํ หิ นาม ฉพฺพคฺคิยา ภิกฺขู อติสหสาปิ วจฺจกุฏิํ ปวิสิสฺสนฺติ, อุพฺภชิตฺวาปิ ปวิสิสฺสนฺติ, นิตฺถุนนฺตาปิ วจฺจํ กริสฺสนฺติ, ทนฺตกฏฺฐํ ขาทนฺตาปิ วจฺจํ กริสฺสนฺติ, พหิทฺธาวิ วจฺจโทณิกาย วจฺจํ กริสฺสนฺติ, พหิทฺธาปิ ปสฺสาวโทณิกาย ปสฺสาวํ กริสฺสนฺติ, ปสฺสาวโทณิกายปิ เขฬํ กริสฺสนฺติ, ผรุเสนปิ กฏฺเฐน อวเลขิสฺสนฺติ, อวเลขนกฏฺฐมฺปิ วจฺจกูปมฺหิ ปาเตสฺสนฺติ, อติสหสาปิ นิกฺขมิสฺสนฺติ, อุพฺภชิตฺวาปิ นิกฺขมิสฺสนฺติ, จปุจปุการกมฺปิ อาจเมสฺสนฺติ, อาจมนสราวเกปิ อุทกํ เสเสสฺสนฺตี”ติ.\csromanspacer{} อถ โข เต ภิกฺขู ภควโต \csromanfolio{392}เอตมตฺถํ อาโรเจสุํ ฯเปฯ.\csromanspacer{} สจฺจํ กิร ภิกฺขเว ฯเปฯ.\csromanspacer{} สจฺจํ ภควาติ ฯเปฯ วิครหิตฺวา ฯเปฯ ธมฺมิํ กถํ กตฺวา ภิกฺขู อามนฺเตสิ\csromandash{}}

\proseitem{374}{เตน หิ ภิกฺขเว ภิกฺขูนํ วจฺจกุฏิวตฺตํ ปญฺญเปสฺสามิ, ยถา ภิกฺขูหิ วจฺจกุฏิยา สมฺมา วตฺติตพฺพํ.\csromanspacer{} โย วจฺจกุฏิํ คจฺฉติ, เตน พหิ ฐิเตน\footnote{พหิ ฐิเตน (สี, ก)} อุกฺกาสิตพฺพํ, อนฺโต นิสินฺเนนปิ อุกฺกาสิตพฺพํ, จีวรวํเส วา จีวรรชฺชุยา วา จีวรํ นิกฺขิปิตฺวา สาธุกํ อตรมาเนน วจฺจกุฏิ ปวิสิตพฺพา, นาติสหสา ปวิสิตพฺพา, น อุพฺภชิตฺวา ปวิสิตพฺพา, วจฺจปาทุกาย ฐิเตน อุพฺภชิตพฺพํ, น นิตฺถุนนฺเตน วจฺโจ กาตพฺโพ, น ทนฺตกฏฺฐํ ขาทนฺเตน วจฺโจ กาตพฺโพ, น พหิทฺธา วจฺจโทณิกาย วจฺโจ กาตพฺโพ, น พหิทฺธา ปสฺสาวโทณิกาย ปสฺสาโว กาตพฺโพ, น ปสฺสาวโทณิกาย เขโฬ กาตพฺโพ, น ผรุเสน กฏฺเฐน อวเลขิตพฺพํ, น อวเลขนกฏฺฐํ วจฺจกูปมฺหิ ปาเตตพฺพํ, วจฺจปาทุกาย ฐิเตน ปฏิจฺฉาเทตพฺพํ, นาติสหสา นิกฺขมิตพฺพํ, น อุพฺภชิตฺวา นิกฺขมิตพฺพํ, อาจมนปาทุกาย ฐิเตน อุพฺภชิตพฺพํ, น จปุจปุการกํ อาจเมตพฺพํ, น อาจมนสราวเก อุทกํ เสเสตพฺพํ, อาจมนปาทุกาย ฐิเตน ปฏิจฺฉาเทตพฺพํ.}

\prose{สเจ วจฺจกุฏิ อุหตา\footnote{อูหตา (สี, สฺยา)} โหติ, โธวิตพฺพา.\csromanspacer{} สเจ อวเลขนปิธโร ปูโร โหติ, อวเลขนกฏฺฐํ ฉฑฺเฑตพฺพํ.\csromanspacer{} สเจ วจฺจกุฏิ อุกฺลาปา โหติ, วจฺจกุฏิ สมฺมชฺชิตพฺพา.\csromanspacer{} สเจ ปริภณฺฑํ อุกฺลาปํ โหติ, ปริภณฺฑํ สมฺมชฺชิตพฺพํ.\csromanspacer{} สเจ ปริเวณํ อุกฺลาปํ โหติ, ปริเวณํ สมฺมชฺชิตพฺพํ.\csromanspacer{} สเจ โกฏฺฐโก อุกฺลาโป โหติ, โกฏฺฐโก สมฺมชฺชิตพฺโพ.\csromanspacer{} สเจ อาจมนกุมฺภิยา อุทกํ น โหติ, อาจมนกุมฺภิยา อุทกํ อาสิญฺจิตพฺพํ.\csromanspacer{} อิทํ โข ภิกฺขเว ภิกฺขูนํ วจฺจกุฏิวตฺตํ, ยถา ภิกฺขูหิ วจฺจกุฏิยา สมฺมา วตฺติตพฺพนฺติ.}

\csromanmatikaanchor{mtk.187}
\csromantocmarkref{section}{11. อุปชฺฌายวตฺตกถา}{mtk.187}
\bookmark[dest={mtk.187},level=1]{11. อุปชฺฌายวตฺตกถา}
\csromanheader{1}{11. อุปชฺฌายวตฺตกถา}
\proseitem{375}{เตน โข ปน สมเยน สทฺธิวิหาริกา อุปชฺฌาเยสุ น สมฺมา วตฺตนฺติ.\csromanspacer{} เย เต ภิกฺขู อปฺปิจฺฉา ฯเปฯ เต อุชฺฌายนฺติ ขิยฺยนฺติ วิปาเจนฺติ “กถํ หิ นาม สทฺธิวิหาริกา อุปชฺฌาเยสุ น สมฺมา วตฺติสฺสนฺตี”ติ.\csromanspacer{} อถ โข เต ภิกฺขู ภควโต เอตมตฺถํ อาโรเจสุํ ฯเปฯ.\csromanspacer{} สจฺจํ กิร ภิกฺขเว สทฺธิวิหาริกา อุปชฺฌาเยสุ น สมฺมา วตฺตนฺตีติ.\csromanspacer{} สจฺจํ ภควาติ.\csromanspacer{} วิครหิ \csromanfolio{393}พุทฺโธ ภควา ฯเปฯ.\csromanspacer{} กถํ หิ นาม ภิกฺขเว สทฺธิวิหาริกา อุปชฺฌาเยสุ น สมฺมา วตฺติสฺสนฺติ.\csromanspacer{} เนตํ ภิกฺขเว อปฺปสนฺนานํ วา ปสาทาย ฯเปฯ วิครหิตฺวา ฯเปฯ ธมฺมิํ กตํ กตฺวา ภิกฺขู อามนฺเตสิ\csromandash{}}

\proseitem{376}{เตน หิ ภิกฺขเว สทฺธิวิหาริกานํ อุปชฺฌาเยสุ วตฺตํ ปญฺญเปสฺสามิ, ยถา สทฺธิวิหาริเกหิ อุปชฺฌาเยสุ สมฺมา วตฺติตพฺพํ. * สทฺธิวิหาริเกน ภิกฺขเว อุปชฺฌายมฺหิ สมฺมา วตฺติตพฺพํ.\csromanspacer{} ตตฺรายํ สมฺมาวตฺตนา\csromandash{}}

\prose{กาลสฺเสว อุฏฺฐาย อุปาหนา โอมุญฺจิตฺวา เอกํสํ อุตฺตราสงฺคํ กริตฺวา ทนฺตกฏฺฐํ ทาตพฺพํ, มุโขทกํ ทาตพฺพํ อาสนํ ปญฺญเปตพฺพํ.\csromanspacer{} สเจ ยาคุ โหติ, ภาชนํ โธวิตฺวา ยาคุ อุปนาเมตพฺพา, ยาคุํ ปีตสฺส อุทกํ ทตฺวา ภาชนํ ปฏิคฺคเหตฺวา นีจํ กตฺวา สาธุกํ อปฺปฏิฆํสนฺเตน โธวิตฺวา ปฏิสาเมตพฺพํ, อุปชฺฌายมฺหิ วุฏฺฐิเต อาสนํ อุทฺธริตพฺพํ.\csromanspacer{} สเจ โส เทโส อุกฺลาโป โหติ, โส เทโส สมฺมชฺชิตพฺโพ.}

\prose{สเจ อุปชฺฌาโย คามํ ปวิสิตุกาโม โหติ, นิวาสนํ ทาตพฺพํ, ปฏินิวาสนํ ปฏิคฺคเหตพฺพํ, กายพนฺธนํ ทาตพฺพํ, สคุณํ กตฺวา สํฆาฏิโย ทาตพฺพา, โธวิตฺวา ปตฺโต โสทโก\footnote{ส-อุทโก (ก)} ทาตพฺโพ.\csromanspacer{} สเจ อุปชฺฌาโย ปจฺฉาสมณํ อากงฺขติ, ติมณฺฑลํ ปฏิจฺฉาเทนฺเตน ปริมณฺฑลํ นิวาเสตฺวา กายพนฺธนํ พนฺธิตฺวา สคุณํ กตฺวา สํฆาฏิโย ปารุปิตฺวา คณฺฐิกํ ปฏิมุญฺจิตฺวา โธวิตฺวา ปตฺตํ คเหตฺวา อุปชฺฌายสฺส ปจฺฉาสมเณน โหตพฺพํ, นาติทูเร คนฺตพฺพํ, นาจฺจาสนฺเน คนฺตพฺพํ, ปตฺตปริยาปนฺนํ ปฏิคฺคเหตพฺพํ, น อุปชฺฌายสฺส ภณมานสฺส อนฺตรนฺตรา กถา โอปาเตตพฺพา, อุปชฺฌาโย อาปตฺติสามนฺตา ภณมาโน นิวาเรตพฺโพ.}

\prose{นิวตฺตนฺเตน ปฐมตรํ อาคนฺตฺวา อาสนํ ปญฺญเปตพฺพํ, ปาโททกํ ปาทปีฐํ ปาทกถลิกํ อุปนิกฺขิปิตพฺพํ, ปจฺจุคฺคนฺตฺวา ปตฺตจีวรํ ปฏิคฺคเหตพฺพํ ปฏินิวาสนํ ทาตพฺพํ, นิวาสนํ ปฏิคฺคเหตพฺพํ.\csromanspacer{} สเจ จีวรํ สินฺนํ โหติ.\csromanspacer{} มุหุตฺตํ อุเณฺห โอตาเปตพฺพํ, น จ อุเณฺห จีวรํ นิทหิตพฺพํ, จีวรํ สงฺฆริตพฺพํ, จีวรํ สงฺฆรนฺเตน จตุรงฺคุลํ กณฺณํ อุสฺสาเรตฺวา จีวรํ สงฺฆริตพฺพํ “มา มชฺเฌ ภงฺโค อโหสี”ติ, โอโภเค กายพนฺธนํ กาตพฺพํ.}

\prose{\csromanfolio{394}สเจ ปิณฺฑปาโต โหติ, อุปชฺฌาโย จ ภุญฺชิตุกาโม โหติ, อุทกํ ทตฺวา ปิณฺฑปาโต อุปนาเมตพฺโพ, อุปชฺฌาโย ปานีเยน ปุจฺฉิตพฺโพ, ภุตฺตาวิสฺส อุทกํ ทตฺวา ปตฺตํ ปฏิคฺคเหตฺวา นีจํ กตฺวา สาธุกํ อปฺปฏิฆํสนฺเตน โธวิตฺวา โวทกํ กตฺวา มุหุตฺตํ อุเณฺห โอตาเปตพฺโพ, น จ อุเณฺห ปตฺโต นิทหิตพฺโพ.\csromanspacer{} ปตฺตจีวรํ นิกฺขิปิตพฺพํ ปตฺตํ นิกฺขิปนฺเตน เอเกน หตฺเถน ปตฺตํ คเหตฺวา เอเกน หตฺเถน เหฏฺฐามญฺจํ วา เหฏฺฐาปีฐํ วา ปรามสิตฺวา ปตฺโต นิกฺขิปิตพฺโพ, น จ อนนฺตรหิตาย ภูมิยา ปตฺโต นิกฺขิปิตพฺโพ.\csromanspacer{} จีวรํ นิกฺขิปนฺเตน เอเกน หตฺเถน จีวรํ คเหตฺวา เอเกน หตฺเถน จีวรวํสํ วา จีวรรชฺชุํ วา ปมชฺชิตฺวา ปารโต อนฺตํ โอรโต โภคํ กตฺวา จีวรํ นิกฺขิปิตพฺพํ.\csromanspacer{} อุปชฺฌายมฺหิ วุฏฺฐิเต อาสนํ อุทฺธริตพฺพํ, ปาโททกํ ปาทปีฐํ ปาทกถลิกํ ปฏิสาเมตพฺพํ.\csromanspacer{} สเจ โส เทโส อุกฺลาโป โหติ, โสเทโส สมฺมชฺชิตพฺโพ.}

\prose{สเจ อุปชฺฌาโย นหายิตุกาโม โหติ, นหานํ ปฏิยาเทตพฺพํ.\csromanspacer{} สเจ สีเตน อตฺโถ โหติ, สีตํ ปฏิยาเทตพฺพํ.\csromanspacer{} สเจ อุเณฺหน อตฺโถ โหติ, อุณฺหํ ปฏิยาเทตพฺพํ.}

\prose{สเจ อุปชฺฌาโย ชนฺตาฆรํ ปวิสิตุกาโม โหติ, จุณฺณํ สนฺเนตพฺพํ, มตฺติกา เตเมตพฺพา, ชนฺตาฆรปีฐํ อาทาย อุปชฺฌายสฺส ปิฏฺฐิโต ปิฏฺฐิโต คนฺตฺวา ชนฺตาฆรปีฐํ ทตฺวา จีวรํ ปฏิคฺคเหตฺวา เอกมนฺตํ นิกฺขิปิตพฺพํ, จุณฺณํ ทาตพฺพํ, มตฺติกา ทาตพฺพา.\csromanspacer{} สเจ อุสฺสหติ, ชนฺตาฆรํ ปวิสิตพฺพํ, ชนฺตาฆรํ ปวิสนฺเตน มตฺติกาย มุขํ มกฺเขตฺวา ปุรโต จ ปจฺฉโต จ ปฏิจฺฉาเทตฺวา ชนฺตาฆรํ ปวิสิตพฺพํ, น เถเร ภิกฺขู อนุปขชฺช นิสีทิตพฺพํ, น นวา ภิกฺขู อาสเนน ปฏิพาหิตพฺพา, ชนฺตาฆเร อุปชฺฌายสฺส ปริกมฺมํ กาตพฺพํ, ชนฺตาฆรา นิกฺขมนฺเตน ชนฺตาฆรปีฐํ อาทาย ปุรโต จ ปจฺฉโต จ ปฏิจฺฉาเทตฺวา ชนฺตาฆรา นิกฺขมิตพฺพํ.}

\prose{อุทเกปิ อุปชฺฌายสฺส ปริกมฺมํ กาตพฺพํ, นหาเตน ปฐมตรํ อุตฺตริตฺวา อตฺตโน คตฺตํ โวทกํ กตฺวา นิวาเสตฺวา อุปชฺฌายสฺส คตฺตโต อุทกํ ปมชฺชิตพฺพํ, นิวาสนํ ทาตพฺพํ, สํฆาฏิ ทาตพฺโพ, ชนฺตาฆรปีฐํ อาทาย ปฐมตรํ อาคนฺตฺวา อาสนํ ปญฺญเปตพฺพํ, ปาโททกํ ปาทปิฐํ ปาทกถลิกํ อุปนิกฺขิปิตพฺพํ, อุปชฺฌาโย ปานีเยน ปุจฺฉิตพฺโพ.\csromanspacer{} สเจ อุทฺทิสาเปตุกาโม โหติ อุทฺทิสิตพฺโพ.\csromanspacer{} สเจ ปริปุจฺฉิตุกาโม โหติ ปริปุจฺฉิตพฺโพ.}

\prose{\csromanfolio{395}ยสฺมิํ วิหาเร อุปชฺฌาโย วิหรติ, สเจ โส วิหาโร อุกฺลาโป โหติ, สเจ อุสฺสหติ, โสเธตพฺโพ, วิหารํ โสเธนฺเตน ปฐมํ ปตฺตจีวรํ นีหริตฺวา เอกมนฺตํ นิกฺขิปิตพฺพํ, นิสีทนปจฺจตฺถรณํ นีหริตฺวา เอกมนฺตํ นิกฺขิปิตพฺพํ, ภิสิพิพฺโพหนํ นีหริตฺวา เอกมนฺตํ นิกฺขิปิตพฺพํ, มญฺโจ นีจํ กตฺวา สาธุกํ อปฺปฏิฆํสนฺเตน อสํฆฏฺเฏนฺเตน กวาฏปิฏฺฐํ นีหริตฺวา เอกมนฺตํ นิกฺขิปิตพฺโพ, ปีฐํ นีจํ กตฺวา สาธุกํ อปฺปฏิฆํสนฺเตน อสํฆฏฺเฏนฺเตน กวาฏปิฏฺฐํ นีหริตฺวา เอกมนฺตํ นิกฺขิปิตพฺพํ, มญฺจปฏิปาทกา นีหริตฺวา เอกมนฺตํ นิกฺขิปิตพฺพา, เขฬมลฺลโก นีหริตฺวา เอกมนฺตํ นิกฺขิปิตพฺโพ, อปสฺเสนผลกํ นีหริตฺวา เอกมนฺตํ นิกฺขิปิตพฺพํ ภูมตฺถรณํ ยถาปญฺญตฺตํ สลฺลกฺเขตฺวา นีหริตฺวา เอกมนฺตํ นิกฺขิปิตพฺพํ.\csromanspacer{} สเจ วิหาเร สนฺตานกํ โหติ, อุลฺโลกา ปฐมํ โอหาเรตพฺพํ, อาโลกสนฺธิกณฺณภาคา ปมชฺชิตพฺพา.\csromanspacer{} สเจ เครุกปริกมฺมกตา ภิตฺติ กณฺณกิตา โหติ, โจฬกํ เตเมตฺวา ปีเฬตฺวา ปมชฺชิตพฺพา.\csromanspacer{} สเจ กาฬวณฺณกตา ภูมิ กณฺณกิตา โหติ, โจฬกํ เตเมตฺวา ปีเฬตฺวา ปมชฺชิตพฺพา.\csromanspacer{} สเจ อกตา โหติ ภูมิ, อุทเกน ปริปฺโผสิตฺวา ปริปฺโผสิตฺวา สมฺมชฺชิตพฺพา “มา วิหาโร รเชน อุหญฺญี”ติ, สงฺการํ วิจินิตฺวา เอกมนฺตํ ฉฑฺเฑตพฺพํ.}

\prose{ภูมตฺถรณํ โอตาเปตฺวา โสเธตฺวา ปปฺโผเฏตฺวา อติหริตฺวา ยถาปญฺญตฺตํ ปญฺญเปตพฺพํ, มญฺจปฏิปาทกา โอตาเปตฺวา ปมชฺชิตฺวา อติหริตฺวา ยถาฐาเน ฐเปตพฺพา, มญฺโจ โอตาเปตฺวา โสเธตฺวา ปปฺโผเฏตฺวา นีจํ กตฺวา สาธุกํ อปฺปฏิฆํสนฺเตน อสํฆฏฺเฏนฺเตน กวาฏปิฏฺฐํ อติหริตฺวา ยถาปญฺญตฺตํ ปญฺญเปตพฺโพ, ปีฐํ โอตาเปตฺวา โสเธตฺวา ปปฺโผเฏตฺวา นีจํ กตฺวา สาธุกํ อปฺปฏิฆํสนฺเตน อสํฆฏฺเฏนฺเตน กวาฏปิฏฺฐํ อติหริตฺวา ยถาปญฺญตฺตํ ปญฺญเปตพฺพํ, ภิสิพิพฺโพหนํ โอตาเปตฺวา โสเธตฺวา ปปฺโผเฏตฺวา อติหริตฺวา ยถาปญฺญตฺตํ ปญฺญเปตพฺพํ, นิสีทนปจฺจตฺถรณํ โอตาเปตฺวา โสเธตฺวา ปปฺโผเฏตฺวา อติหริตฺวา ยถาปญฺญตฺตํ ปญฺญเปตพฺพํ, เขฬมลฺลโก โอตาเปตฺวา ปมชฺชิตฺวา อติหริตฺวา ยถาฐาเน ฐเปตพฺโพ, อปสฺเสนผลกํ โอตาเปตฺวา ปมชฺชิตฺวา อติหริตฺวา ยถาฐาเน ฐเปตพฺพํ, ปตฺตจีวรํ นิกฺขิปิตพฺพํ, ปตฺตํ นิกฺขิปนฺเตน เอเกน หตฺเถน ปตฺตํ คเหตฺวา เอเกน หตฺเถน เหฏฺฐามญฺจํ วา เหฏฺฐาปีฐํ วา ปรามสิตฺวา \csromanfolio{396}ปตฺโต นิกฺขิปิตพฺโพ, น จ อนนฺตรหิตาย ภูมิยา ปตฺโต นิกฺขิปิตพฺโพ, จีวรํ นิกฺขิปนฺเตน เอเกน หเถน จีวรํ คเหตฺวา เอเกน หตฺเถน จีวรํวํสํ วา จีวรรชฺชุํ วา ปมชฺชิตฺวา ปารโต อนฺตํ โอรโต โภคํ กตฺวา จีวรํ นิกฺขิปิตพฺพํ.}

\prose{สเจ ปุรตฺถิมา สรชา วาตา วายนฺติ, ปุรตฺถิมา วาตปานา ถเกตพฺพา.\csromanspacer{} สเจ ปจฺฉิมา สรชา วาตา วายนฺติ, ปจฺฉิมา วาตปานา ถเกตพฺพา.\csromanspacer{} สเจ อุตฺตรา สรชา วาตา วายนฺติ, อุตฺตรา วาตปานา ถเกตพฺพา.\csromanspacer{} สเจ ทกฺขิณา สรชา วาตา วายนฺติ, ทกฺขิณา วาตปานา ถเกตพฺพา.\csromanspacer{} สเจ สีตกาโล โหติ, ทิวา วาตปานา วิวริตพฺพา, รตฺติํ ถเกตพฺพา.\csromanspacer{} สเจ อุณฺหกาโล โหติ, ทิวา วา ตปานา ถเกตพฺพา, รตฺติํ วิวริตพฺพา.}

\prose{สเจ ปริเวณํ อุกฺลาปํ โหติ, ปริเวณํ สมฺมชฺชิตพฺพํ.\csromanspacer{} สเจ โกฏฺฐโก อุกฺลาโป โหติ, โกฏฺฐโก สมฺมชฺชิตพฺโพ.\csromanspacer{} สเจ อุปฏฺฐานสาลา อุกฺลาปา โหติ, อุปฏฺฐานสาลา สมฺมชฺชิตพฺพา.\csromanspacer{} สเจ อคฺคิสาลา อุกฺลาปา โหติ, อคฺคิสาลา สมฺมชฺชิตพฺพา.\csromanspacer{} สเจ วจฺจกุฏิ อุกฺลาปา โหติ, วจฺจกุฏิ สมฺมชฺชิตพฺพา.\csromanspacer{} สเจ ปานียํ น โหติ, ปานียํ อุปฏฺฐาเปตพฺพํ.\csromanspacer{} สเจ ปริโภชนียํ น โหติ, ปริโภชนียํ อุปฏฺฐาเปตพฺพํ.\csromanspacer{} สเจ อาจมนกุมฺภิยา อุทกํ น โหติ, อาจมนกุมฺภิยา อุทกํ อาสิญฺจิตพฺพํ.}

\prose{สเจ อุปชฺฌายสฺส อนภิรติ อุปฺปนฺนา โหติ, สทฺธิวิหาริเกน วูปกาเสตพฺโพ วูปกาสาเปตพฺโพ ธมฺมกถา วาสฺส กาตพฺพา.\csromanspacer{} สเจ อุปชฺฌายสฺส กุกฺกุจฺจํ อุปฺปนฺนํ โหติ, สทฺธิวิหาริเกน วิโนเทตพฺพํ วิโนทาเปตพฺพํ ธมฺมกถา วาสฺส กาตพฺพา.\csromanspacer{} สเจ อุปชฺฌายสฺส ทิฏฺฐิคตํ อุปฺปนฺนํ โหติ, สทฺธิวิหาริเกน วิเวเจตพฺพํ วิเวจาเปตพฺพํ ธมฺมกถา วาสฺส กาตพฺพา.\csromanspacer{} สเจ อุปชฺฌาโย ครุธมฺมํ อชฺฌาปนฺโน โหติ ปริวาสารโห, สทฺธิวิหาริเกน อุสฺสุกฺกํ กาตพฺพํ “กินฺติ นุ โข สํโฆ อุปชฺฌายสฺส ปริวาสํ ทเทยฺยา”ติ.\csromanspacer{} สเจ อุปชฺฌาโย มูลายปฏิกสฺสนารโห โหติ, สทฺธิวิหาริเกน อุสฺสุกฺกํ กาตพฺพํ “กินฺติ นุ โข สํโฆ อุปชฺฌายํ มูลาย ปฏิกสฺเสยฺยา”ติ.\csromanspacer{} สเจ อุปชฺฌาโย มานตฺตารโห โหติ, สทฺธิวิหาริเกน อุสฺสุกฺกํ กาตพฺพํ “กินฺติ นุ โข สํโฆ อุปชฺฌายสฺส มานตฺตํ ทเทยฺยา”ติ.\csromanspacer{} สเจ อุปชฺฌาโย อพฺภานารโห โหติ, สทฺธิวิหาริเกน อุสฺสุกฺกํ กาตพฺพํ “กินฺติ นุ \csromanfolio{397}โข สํโฆ อุปชฺฌายํ อพฺเภยฺยา”ติ.\csromanspacer{} สเจ สํโฆ อุปชฺฌายสฺส กมฺมํ กตฺตุกาโม โหติ ตชฺชนียํ วา นิยสฺสํ วา ปพฺพาชนียํ วา ปฏิสารณียํ วา อุกฺเขปนียํ วา, สทฺธิวิหาริเกน อุสฺสุกฺกํ กาตพฺพํ “กินฺติ นุ โข สํโฆ อุปชฺฌายสฺส กมฺมํ น กเรยฺย, ลหุกาย วา ปริณาเมยฺยา”ติ.\csromanspacer{} กตํ วา ปนสฺส โหติ สํเฆน กมฺมํ ตชฺชนียํ วา นิยสฺสํ วา ปพฺพาชนียํ วา ปฏิสารณียํ วา อุกฺเขปนียํ วา, สทฺธิวิหาริเกน อุสฺสุกฺกํ กาตพฺพํ “กินฺติ นุ โข อุปชฺฌาโย สมฺมา วตฺเตยฺย, โลมํ ปาเตยฺย, เนตฺถารํ วตฺเตยฺย, สํโฆ ตํ กมฺมํ ปฏิปฺปสฺสมฺเภยฺยา”ติ.}

\prose{สเจ อุปชฺฌายสฺส จีวรํ โธวิตพฺพํ โหติ, สทฺธิวิหาริเกน โธวิตพฺพํ, อุสฺสุกฺกํ วา กาตพฺพํ “กินฺติ นุ โข อุปชฺฌายสฺส จีวรํ โธวิเยถา”ติ.\csromanspacer{} สเจ อุปชฺฌายสฺส จีวรํ กาตพฺพํ โหติ, สทฺธิวิหาริเกน กาตพฺพํ, อุสฺสุกฺกํ วา กาตพฺพํ “กินฺติ นุ โข อุปชฺฌายสฺส จีวรํ กริเยถา”ติ.\csromanspacer{} สเจ อุปชฺฌายสฺส รชนา ปจิตพฺพา โหติ, สทฺธิวิหาริเกน ปจิตพฺพา, อุสฺสุกฺกํ วา กาตพฺพํ “กินฺติ นุ โข อุปชฺฌายสฺส รชนํ ปจิเยถา”ติ.\csromanspacer{} สเจ อุปชฺฌายสฺส จีวรํ รชิตพฺพํ\footnote{รเชตพฺพํ (สฺยา)} โหติ, สทฺธิวิหาริเกน รชิตพฺพํ, อุสฺสุกฺกํ วา กาตพฺพํ “กินฺติ นุ โข อุปชฺฌายสฺส จีวรํ รชิเยถา”ติ.\csromanspacer{} จีวรํ รชนฺเตน\footnote{รเชนฺเตน (สฺยา)} สาธุกํ สมฺปริวตฺตกํ สมฺปริวตฺตกํ รชิตพฺพํ, น จ อจฺฉินฺเน เถเว ปกฺกมิตพฺพํ.}

\prose{น อุปชฺฌายํ อนาปุจฺฉา เอกจฺจสฺส ปตฺโต ทาตพฺโพ, น เอกจฺจสฺส ปตฺโต ปฏิคฺคเหตพฺโพ, น เอกจฺจสฺส จีวรํ ทาตพฺพํ, น เอกจฺจสฺส จีวรํ ปฏิคฺคเหตพฺพํ, น เอกจฺจสฺส ปริกฺขาโร ทาตพฺโพ, น เอกจฺจสฺส ปริกฺขาโร ปฏิคฺคเหตพฺโพ, น เอกจฺจสฺส เกสา เฉเทตพฺพา\footnote{เฉตฺตพฺพา (สี), เฉทิตพฺพา (ก)}, น เอกจฺเจน เกสา เฉทาเปตพฺพา, น เอกจฺจสฺส ปริกมฺมํ กาตพฺพํ, น เอกจฺเจน ปริกมฺมํ การาเปตพฺพํ, น เอกจฺจสฺส เวยฺยาวจฺโจ\footnote{เวยฺยาวจฺจํ (ก)} กาตพฺโพ, น เอกจฺเจน เวยฺยาวจฺโจ การาเปตพฺโพ, น เอกจฺจสฺส ปจฺฉาสมเณน โหตพฺพํ, น เอกจฺโจ ปจฺฉาสมโณ อาทาตพฺโพ, น เอกจฺจสฺส ปิณฺฑปาโต นีหริตพฺโพ, น เอกจฺเจน ปิณฺฑปาโต นีหราเปตพฺโพ, น อุปชฺฌายํ อนาปุจฺฉา คาโม ปวิสิตพฺโพ, น สุสานํ คนฺตพฺพํ, น ทิสา ปกฺกมิตพฺพา.\csromanspacer{} สเจ อุปชฺฌาโย คิลาโน โหติ, ยาวชีวํ \csromanfolio{398}อุปฏฺฐาตพฺโพ, วุฏฺฐานมสฺส อาคเมตพฺพํ.\csromanspacer{} อิทํ โข ภิกฺขเว สทฺธิวิหาริกานํ อุปชฺฌาเยสุ วตฺตํ, ยถา สทฺธิวิหาริเกหิ อุปชฺฌาเยสุ สมฺมา วตฺติตพฺพนฺติ.}

\csromanmatikaanchor{mtk.188}
\csromantocmarkref{section}{12. สทฺธิวิหาริกวตฺตกถา}{mtk.188}
\bookmark[dest={mtk.188},level=1]{12. สทฺธิวิหาริกวตฺตกถา}
\csromanheader{1}{12. สทฺธิวิหาริกวตฺตกถา}
\proseitem{377}{เตน โข ปน สมเยน อุปชฺฌายา สทฺธิวิหาริเกสุ น สมฺมา วตฺตนฺติ.\csromanspacer{} เย เต ภิกฺขู อปฺปิจฺฉา ฯเปฯ เต อุชฺฌายนฺติ ขิยฺยนฺติ วิปาเจนฺติ “กถํ หิ นาม อุปชฺฌายา สทฺธิวิหาริเกสุ น สมฺมา วตฺติสฺสนฺตี”ติ.\csromanspacer{} อถ โข เต ภิกฺขู ภควโต เอตมตฺถํ อาโรเจสุํ ฯเปฯ.\csromanspacer{} สจฺจํ กิร ภิกฺขเว อุปชฺฌายา สทฺธิวิหาริเกสุ น สมฺมา วตฺตนฺตีติ.\csromanspacer{} สจฺจํ ภควาติ ฯเปฯ วิครหิตฺวา ฯเปฯ ธมฺมิํ กถํ กตฺวา ภิกฺขู อามนฺเตสิ\csromandash{}}

\proseitem{378}{เตน หิ ภิกฺขเว อุปชฺฌายานํ สทฺธิวิหาริเกสุ วตฺตํ ปญฺญเปสฺสามิ, ยถา อุปชฺฌาเยหิ สทฺธิวิหาริเกสุ สมฺมา วตฺติตพฺพํ. * อุปชฺฌาเยน ภิกฺขเว สทฺธิวิหาริกมฺหิ สมฺมา วตฺติตพฺพํ, ตตฺรายํ สมฺมาวตฺตนา\csromandash{}}

\prose{อุปชฺฌาเยน ภิกฺขเว สทฺธิวิหาริโก สงฺคเหตพฺโพ อนุคฺคเหตพฺโพ อุทฺเทเสน ปริปุจฺฉาย โอวาเทน อนุสาสนิยา.\csromanspacer{} สเจ อุปชฺฌายสฺส ปตฺโต โหติ สทฺธิวิหาริกสฺส ปตฺโต น โหติ, อุปชฺฌาเยน สทฺธิวิหาริกสฺส ปตฺโต ทาตพฺโพ, อุสฺสุกฺกํ วา กาตพฺพํ “กินฺติ นุ โข สทฺธิวิหาริกสฺส ปตฺโต อุปฺปชฺชิเยถา”ติ.\csromanspacer{} สเจ อุปชฺฌายสฺส จีวรํ โหติ สทฺธิวิหาริกสฺส จีวรํ น โหติ, อุปชฺฌาเยน สทฺธิวิหาริกสฺส จีวรํ ทาตพฺพํ, อุสฺสุกฺกํ วา กาตพฺพํ “กินฺติ นุ โข สทฺธิวิหาริกสฺส จีวรํ อุปฺปชฺชิเยถา”ติ.\csromanspacer{} สเจ อุปชฺฌายสฺส ปริกฺขาโร โหติ สทฺธิวิหาริกสฺส ปริกฺขาโร น โหติ, อุปชฺฌาเยน สทฺธิวิหาริกสฺส ปริกฺขาโร ทาตพฺโพ, อุสฺสุกฺกํ วา กาตพฺพํ “กินฺติ นุ โข สทฺธิวิหาริกสฺส ปริกฺขาโร อุปฺปชฺชิเยถา”ติ.}

\prose{สเจ สทฺธิวิหาริโก คิลาโน โหติ, กาลสฺเสว อุฏฺฐาย ทนฺตกฏฺฐํ ทาตพฺพํ, มุโขทกํ ทาตพฺพํ, อาสนํ ปญฺญเปตพฺพํ.\csromanspacer{} สเจ ยาคุ โหติ, ภาชนํ โธวิตฺวา ยาคุ อุปนาเมตพฺพา, ยาคุํ ปีตสฺส \csromanfolio{399}อุทกํ ทตฺวา ภาชนํ ปฏิคฺคเหตฺวา นีจํ กตฺวา สาธุกํ อปฺปฏิฆํสนฺเตน โธวิตฺวา ปฏิสาเมตพฺพํ สทฺธิวิหาริกมฺหิ วุฏฺฐิเต อาสนํ อุทฺธริตพฺพํ.\csromanspacer{} สเจ โส เทโส อุกฺลาโป โหติ, โส เทโส สมฺมชฺชิตพฺโพ.}

\prose{สเจ สทฺธิวิหาริโก คามํ ปวิสิตุกาโม โหติ, นิวาสนํ ทาตพฺพํ, ปฏินิวาสนํ ปฏิคฺคเหตพฺพํ, กายพนฺธนํ ทาตพฺพํ, สคุณํ กตฺวา สํฆาฏิโย ทาตพฺพา, โธวิตฺวา ปตฺโต โสทโก ทาตพฺโพ.}

\prose{“เอตฺตาวตา นิวตฺติสฺสตี”ติ อาสนํ ปญฺญเปตพฺพํ, ปาโททกํ ปาทปีฐํ ปาทกถลิกํ อุปนิกฺขิปิตพฺพํ.\csromanspacer{} ปจฺจุคฺคนฺตฺวา ปตฺตจีวรํ ปฏิคฺคเหตพฺพํ ปฏินิวาสนํ ทาตพฺพํ, นิวาสนํ ปฏิคฺคเหตพฺพํ.\csromanspacer{} สเจ จีวรํ สินฺนํ โหติ, มุหุตฺตํ อุเณฺห โอตาเปตพฺพํ, น จ อุเณฺห จีวรํ นิทหิตพฺพํ, จีวรํ สงฺฆริตพฺพํ, จีวรํ สงฺฆรนฺเตน จตุรงฺคุลํ กณฺณํ อุสฺสาเรตฺวา จีวรํ สงฺฆริตพฺพํ “มา มชฺเฌ ภงฺโค อโหสี”ติ, โอโภเค กายพนฺธนํ กาตพฺพํ.}

\prose{สเจ ปิณฺฑปาโต โหติ สทฺธิวิหาริโก จ ภุญฺชิตุกาโม โหติ, อุทกํ ทตฺวา ปิณฺฑปาโต อุปนาเมตพฺโพ, สทฺธิวิหาริโกปานีเยน ปุจฺฉิตพฺโพ, ภุตฺตาวิสฺส อุทกํ ทตฺวา ปตฺตํ ปฏิคฺคเหตฺวา นีจํ กตฺวา สาธุกํ อปฺปฏิฆํสนฺเตน โธวิตฺวา โวทกํ กตฺวา มุหุตฺตํ อุเณฺห โอตาเปตพฺโพ, น จ อุเณฺห ปตฺโต นิทหิตพฺโพ.\csromanspacer{} ปตฺตจีวรํ นิกฺขิปิตพฺพํ, ปตฺตํ นิกฺขิปนฺเตน เอเกน หตฺเถน ปตฺตํ คเหตฺวา เอเกน หตฺเถน เหฏฺฐามญฺจํ วา เหฏฺฐาปีฐํ วา ปรามสิตฺวา ปตฺโต นิกฺขิปิตพฺโพ, น จ อนนฺตรหิตาย ภูมิยา ปตฺโต นิกฺขิปิตพฺโพ.\csromanspacer{} จีวรํ นิกฺขิปนฺเตน เอเกน หตฺเถน จีวรํ คเหตฺวา เอเกน หตฺเถน จีวรวํสํ วา จีวรรชฺชุํ วา ปมชฺชิตฺวา ปารโต อนฺตํ โอรโต โภคํ กตฺวา จีวรํ นิกฺขิปิตพฺพํ.\csromanspacer{} สทฺธิวิหาริกมฺหิ วุฏฺฐิเต อาสนํ อุทฺธริตพฺพํ, ปาโททกํ ปาทปีฐํ ปาทกถลิกํ ปฏิสาเมตพฺพํ.\csromanspacer{} สเจ โส เทโส อุกฺลาโป โหติ, โส เทโส สมฺมชฺชิตพฺโพ.}

\prose{สเจ สทฺธิวิหาริโก นหายิตุกาโม โหติ, นหานํ ปฏิยาเทตพฺพํ.\csromanspacer{} สเจ สีเตน อตฺโถ โหติ, สีตํ ปฏิยาเทตพฺพํ.\csromanspacer{} สเจ อุเณฺหน อตฺโถ โหติ, อุณฺหํ ปฏิยาเทตพฺพํ.}

\prose{\csromanfolio{400}สเจ สทฺธิวิหาริโก ชนฺตาฆรํ ปวิสิตุกาโม โหติ, จุณฺณํ สนฺเนตพฺพํ, มตฺติกา เตเมตพฺพา, ชนฺตาฆรปีฐํ อาทาย\footnote{อาทาย สทฺธิวิหาริกสฺส ปิฏฺฐิโต ปิฏฺฐิโต (ก)} คนฺตฺวา ชนฺตาฆรปีฐํ ทตฺวา จีวรํ ปฏิคฺคเหตฺวา เอกมนฺตํ นิกฺขิปิตพฺพํ, จุณฺณํ ทาตพฺพํ, มตฺติกา ทาตพฺพา.\csromanspacer{} สเจ อุสฺสหติ, ชนฺตาฆรํ ปวิสิตพฺพํ, ชนฺตาฆรํ ปวิสนฺเตน มตฺติกาย มุขํ มกฺเขตฺวา ปุรโต จ ปจฺฉโต จ ปฏิจฺฉาเทตฺวา ชนฺตาฆรํ ปวิสิตพฺพํ, น เถเร ภิกฺขู อนุปขชฺช นิสีทิตพฺพํ, น นวา ภิกฺขู อาสเนน ปฏิพาหิตพฺพา, ชนฺตาฆเร สทฺธิวิหาริกสฺส ปริกมฺมํ กาตพฺพํ, ชนฺตาฆรา นิกฺขมนฺเตน ชนฺตาฆรปีฐํ อาทาย ปุรโต จ ปจฺฉโต จ ปฏิจฺฉาเทตฺวา ชนฺตาฆรา นิกฺขมิตพฺพํ.}

\prose{อุทเกปิ สทฺธิวิหาริกสฺส ปริกมฺมํ กาตพฺพํ, นหาเตน ปฐมตรํ อุตฺตริตฺวา อตฺตโน คตฺตํ โวทกํ กตฺวา นิวาเสตฺวา สทฺธิวิหาริกสฺส คตฺตโต อุทกํ ปมชฺชิตพฺพํ, นิวาสนํ ทาตพฺพํ, สํฆาฏิ ทาตพฺพา, ชนฺตาฆรปีฐํ อาทาย ปฐมตรํ อาคนฺตฺวา อาสนํ ปญฺญเปตพฺพํ, ปาโททกํ ปาทปีฐํ ปาทกถลิกํ อุปนิกฺขิปิตพฺพํ, สทฺธิวิหาริโก ปานีเยน ปุจฺฉิตพฺโพ.}

\prose{ยสฺมิํ วิหาเร สทฺธิวิหาริโก วิหรติ, สเจ โส วิหาโร อุกฺลาโป โหติ, สเจ อุสฺสหติ, โสเธตพฺโพ, วิหารํ โสเธนฺเตน ปฐมํ ปตฺตจีวรํ นีหริตฺวา เอกมนฺตํ นิกฺขิปิตพฺพํ ฯเปฯ.\csromanspacer{} สเจ อาจมนกุมฺภิยา อุทกํ น โหติ, อาจมนกุมฺภิยา อุทกํ อาสิญฺจิตพฺพํ.}

\prose{สเจ สทฺธิวิหาริกสฺส อนภิรติ อุปฺปนฺนา โหติ, อุปชฺฌาเยน วูปกาเสตพฺโพ วูปกาสาเปตพฺโพ ธมฺมกถา วาสฺส กาตพฺพา.\csromanspacer{} สเจ สทฺธิวิหาริกสฺส กุกฺกุจฺจํ อุปฺปนฺนํ โหติ, อุปชฺฌาเยน วิโนเทตพฺพํ วิโนทาเปตพฺพํ, ธมฺมกถา วาสฺส กาตพฺพา.\csromanspacer{} สเจ สทฺธิวิหาริกสฺส ทิฏฺฐิคตํ อุปฺปนฺนํ โหติ, อุปชฺฌาเยน วิเวเจตพฺพํ วิเวจาเปตพฺพํ ธมฺมกถา วาสฺส กาตพฺพา.\csromanspacer{} สเจ สทฺธิวิหาริโก ครุธมฺมํ อชฺฌาปนฺโน โหติ ปริวาสารโห, อุปชฺฌาเยน อุสฺสุกฺกํ กาตพฺพํ “กินฺติ นุ โข สํโฆ สทฺธิวิหาริกสฺส ปริวาสํ ทเทยฺยาติ.\csromanspacer{} สเจ สทฺธิวิหาริโก มูลายปฏิกสฺสนารโห โหติ, อุปชฺฌาเยน อุสฺสุกฺกํ กาตพฺพํ “กินฺติ นุ โข สํโฆ สทฺธิวิหาริกํ มูลาย ปฏิกสฺเสยฺยา”ติ.\csromanspacer{} สเจ สทฺธิวิหาริโก มานตฺตารโห โหติ, อุปชฺฌาเยน อุสฺสุกฺกํ กาตพฺพํ “กินฺติ นุ โข \csromanfolio{401}สํโฆ สทฺธิวิหาริกสฺส มานตฺตํ ทเทยฺยา”ติ.\csromanspacer{} สเจ สทฺธิวิหาริโก อพฺภานารโห โหติ, อุปชฺฌาเยน อุสฺสุกฺกํ กาตพฺพํ “กินฺติ นุ โข สํโฆ สทฺธิวิหาริกํ อพฺเภยฺยา”ติ.\csromanspacer{} สเจ สํโฆ สทฺธิวิหาริกสฺส กมฺมํ กตฺตุกาโม โหติ ตชฺชนียํ วา นิยสฺสํ วา ปพฺพาชนียํ วา ปฏิสารณียํ วา อุกฺเขปนียํ วา, อุปชฺฌาเยน อุสฺสุกฺกํ กาตพฺพํ “กินฺติ นุ โข สํโฆ สทฺธิวิหาริกสฺส กมฺมํ น กเรยฺย, ลหุกาย วา ปริณาเมยฺยา”ติ.\csromanspacer{} กตํ วา ปนสฺส โหติ สํเฆน กมฺมํ ตชฺชนียํ วา นิยสฺสํ วา ปพฺพาชนียํ วา ปฏิสารณียํ วา อุกฺเขปนียํ วา, อุปชฺฌาเยน อุสฺสุกฺกํ กาตพฺพํ “กินฺติ นุ โข สทฺธิวิหาริโก สมฺมา วตฺเตยฺย, โลมํ ปาเตยฺย, เนตฺถารํ วตฺเตยฺย, สํโฆ ตํ กมฺมํ ปฏิปฺปสฺสมฺเภยา”ติ.}

\prose{สเจ สทฺธิวิหาริกสฺส จีวรํ โธวิตพฺพํ โหติ, อุปชฺฌาเยน อาจิกฺขิตพฺพํ “เอวํ โธเวยฺยาสี”ติ, อุสฺสุกฺกํ วา กาตพฺพํ “กินฺติ นุ โข สทฺธิวิหาริกสฺส จีวรํ โธวิเยถา”ติ.\csromanspacer{} สเจ สทฺธิวิหาริกสฺส จีวรํ กาตพฺพํ โหติ, อุปชฺฌาเยน อาจิกฺขิตพฺพํ “เอวํ กเรยฺยาสี”ติ, อุสฺสุกฺกํ วา กาตพฺพํ “กินฺติ นุ โข สทฺธิวิหาริกสฺส จีวรํ กริเยถา”ติ.\csromanspacer{} สเจ สทฺธิวิหาริกสฺส รชนํ ปจิตพฺพํ โหติ, อุปชฺฌาเยน อาจิกฺขิตพฺพํ “เอวํ ปเจยฺยาสี”ติ, อุสฺสุกฺกํ วา กาตพฺพํ “กินฺติ นุ โข สทฺธิวิหาริกสฺส รชนํ ปจิเยถา”ติ.\csromanspacer{} สเจ สทฺธิวิหาริกสฺส จีวรํ รชิตพฺพํ โหติ, อุปชฺฌาเยน อาจิกฺขิตพฺพํ “เอวํ รเชยฺยาสี”ติ, อุสฺสุกฺกํ วา กาตพฺพํ “กินฺติ นุ โข สทฺธิวิหาริกสฺส จีวรํ รชิเยถา”ติ.\csromanspacer{} จีวรํ รชนฺเตน สาธุกํ สมฺปริวตฺตกํ สมฺปริวตฺตกํ รชิตพฺพํ, น จ อจฺฉินฺเน เถเว ปกฺกมิตพฺพํ, สเจ สทฺธิวิหาริโก คิลาโน โหติ, ยาวชีวํ อุปฏฺฐาตพฺโพ, วุฏฺฐานมสฺส อาคเมตพฺพํ.\csromanspacer{} อิทํ โข ภิกฺขเว อุปชฺฌายานํ สทฺธิวิหาริเกสุ วตฺตํ, ยถา อุปชฺฌาเยหิ สทฺธิวิหาริเกสุ สมฺมา วตฺติตพฺพนฺติ.}

\nitthitam{ทุติยภาณวาโร นิฏฺฐิโต.}
\csromanmatikaanchor{mtk.189}
\csromantocmarkref{section}{13. อาจริยวตฺตกถา}{mtk.189}
\bookmark[dest={mtk.189},level=1]{13. อาจริยวตฺตกถา}
\csromanheader{1}{13. อาจริยวตฺตกถา}
\proseitem{379}{เตน โข ปน สมเยน อนฺเตวาสิกา อาจริเยสุ น สมฺมา วตฺตนฺติ.\csromanspacer{} เย เต ภิกฺขู อปฺปิจฺฉา ฯเปฯ เต อุชฺฌายนฺติ ขิยฺยนฺติ วิปาเจนฺติ “กถํ หิ นาม อนฺเตวาสิกา อาจริเยสุ น สมฺมา วตฺติสฺสนฺตี”ติ. \csromanfolio{402}อถ โข เต ภิกฺขู ภควโต เอตมตฺถํ อาโรเจสุํ ฯเปฯ.\csromanspacer{} สจฺจํ กิร ภิกฺขเว อนฺเตวาสิกา อาจริเยสุ น สมฺมา วตฺตนฺตีติ.\csromanspacer{} สจฺจํ ภควาติ ฯเปฯ วิครหิตฺวา ฯเปฯ ธมฺมิํ กถํ กตฺวา ภิกฺขู อามนฺเตสิ\csromandash{}}

\proseitem{380}{เตน หิ ภิกฺขเว อนฺเตวาสิกานํ อาจริเยสุ วตฺตํ ปญฺญเปสฺสามิ, ยถา อนฺเตวาสิเกหิ อาจริเยสุ สมฺมา วตฺติตพฺพํ. * อนฺเตวาสิเกน ภิกฺขเว อาจริยมฺหิ สมฺมา วตฺติตพฺพํ, ตตฺรายํ สมฺมาวตฺตนา\csromandash{}}

\prose{กาลสฺเสว อุฏฺฐาย อุปาหนา โอมุญฺจิตฺวา เอกํสํ อุตฺตราสงฺคํ กริตฺวา ทนฺตกฏฺฐํ ทาตพฺพํ, มุโขทกํ ทาตพฺพํ, อาสนํ ปญฺญเปตพฺพํ.\csromanspacer{} สเจ ยาคุ โหติ, ภาชนํ โธวิตฺวา ยาคุ อุปนาเมตพฺพา, ยาคุํ ปีตสฺส อุทกํ ทตฺวา ภาชนํ ปฏิคฺคเหตฺวา นีจํ กตฺวา สาธุกํ อปฺปฏิฆํสนฺเตน โธวิตฺวา ปฏิสาเมตพฺพํ, อาจริยมฺหิ วุฏฺฐิเต อาสนํ อุทฺธริตพฺพํ.\csromanspacer{} สเจ โส เทโส อุกฺลาโป โหติ, โส เทโส สมฺมชฺชิตพฺโพ.}

\prose{สเจ อาจริโย คามํ ปวิสิตุกาโม โหติ, นิวาสนํ ทาตพฺพํ, ปฏินิวาสนํ ปฏิคฺคเหตพฺพํ, กายพนฺธนํ ทาตพฺพํ, สคุณํ กตฺวา สํฆาฏิโย ทาตพฺพา, โธวิตฺวา ปตฺโต โสทโก ทาตพฺโพ.\csromanspacer{} สเจ อาจริโย ปจฺฉาสมณํ อากงฺขติ, ติมณฺฑลํ ปฏิจฺฉาเทนฺเตน ปริมณฺฑลํ นิวาเสตฺวา กายพนฺธนํ พนฺธิตฺวา สคุณํ กตฺวา สํฆาฏิโย ปารุปิตฺวา คณฺฐิกํ ปฏิมุญฺจิตฺวา โธวิตฺวา ปตฺตํ คเหตฺวา อาจริยสฺส ปจฺฉาสมเณน โหตพฺพํ, นาติทูเร คนฺตพฺพํ, นาจฺจาสนฺเน คนฺตพฺพํ, ปตฺตปริยาปนฺนํ ปฏิคฺคเหตพฺพํ, น อาจริยสฺส ภณมานสฺส อนฺตรนฺตรา กถา โอปาเตตพฺพา, อาจริโย อาปตฺติสามนฺตา ภณมาโน นิวาเรตพฺโพ.}

\prose{นิวตฺตนฺเตน ปฐมตรํ อาคนฺตฺวา อาสนํ ปญฺญเปตพฺพํ, ปาโททกํ ปาทปีฐํ ปาทกถลิกํ อุปนิกฺขิปิตพฺพํ, ปจฺจุคฺคนฺตฺวา ปตฺตจีวรํ ปฏิคฺคเหตพฺพํ, ปฏินิวาสนํ ทาตพฺพํ, นิวาสนํ ปฏิคฺคเหตพฺพํ.\csromanspacer{} สเจ จีวรํ สินฺนํ โหติ, มุหุตฺตํ อุเณฺห โอตาเปตพฺพํ, น จ อุเณฺห จีวรํ นิทหิตพฺพํ, จีวรํ สงฺฆริตพฺพํ, จีวรํ สงฺฆรนฺเตน จตุรงฺคุลํ กณฺณํ อุสฺสาเรตฺวา จีวรํ สงฺฆริตพฺพํ “มา มชฺเฌ ภงฺโค อโหสี”ติ, โอโภเค กายพนฺธนํ กาตพฺพํ.}

\prose{สเจ ปิณฺฑปาโต โหติ, อาจริโย จ ภุญฺชิตุกาโม โหติ, อุทกํ ทตฺวา ปิณฺฑปาโต อุปนาเมตพฺโพ, อาจริโย ปานีเยน \csromanfolio{403}ปุจฺฉิตพฺโพ, ภุตฺตาวิสฺส อุทกํ ทตฺวา ปตฺตํ ปฏิคฺคเหตฺวา นีจํ กตฺวา สาธุกํ อปฺปฏิฆํสนฺเตน โธวิตฺวา โวทกํ กตฺวา มุหุตฺตํ อุเณฺห โอตาเปตพฺโพ, น จ อุเณฺห ปตฺโต นิทหิตพฺโพ.\csromanspacer{} ปตฺตจีวรํ นิกฺขิปิตพฺพํ, ปตฺตํ นิกฺขิปนฺเตน เอเกน หตฺเถน ปตฺตํ คเหตฺวา เอเกน หตฺเถน เหฏฺฐามญฺจํ วา เหฏฺฐาปีฐํ วา ปรามสิตฺวา ปตฺโต นิกฺขิปิตพฺโพ, น จ อนนฺตรหิตาย ภูมิยา ปตฺโต นิกฺขิปิตพฺโพ.\csromanspacer{} จีวรํ นิกฺขิปนฺเตน เอเกน หตฺเถน จีวรํ คเหตฺวา เอเกน หตฺเถน จีวรวํสํ วา จีวรรชฺชุํ วา ปมชิตฺวา ปารโต อนฺตํ โอรโต โภคํ กตฺวา จีวรํ นิกฺขิปิตพฺพํ.\csromanspacer{} อาจริยมฺหิ วุฏฺฐิเต อาสนํ อุทฺธริตพฺพํ, ปาโททกํ ปาทปีฐํ ปาทกถลิกํ ปฏิสาเมตพฺพํ.\csromanspacer{} สเจ โส เทโส อุกฺลาโป โหติ, โส เทโส สมฺมชฺชิตพฺโพ.}

\prose{สเจ อาจริโย นหายิตุกาโม โหติ, นหานํ ปฏิยาเทตพฺพํ.\csromanspacer{} สเจ สีเตน อตฺโถ โหติ, สีตํ ปฏิยาเทตพฺพํ.\csromanspacer{} สเจ อุเณฺหน อตฺโถ โหติ, อุณฺหํ ปฏิยาเทตพฺพํ.}

\prose{สเจ อาจริโย ชนฺตาฆรํ ปวิสิตุกาโม โหติ, จุณฺณํ สนฺเนตพฺพํ, มตฺติกา เตเมตพฺพา, ชนฺตาฆรปีฐํ อาทาย อาจริยสฺส ปิฏฺฐิโต ปิฏฺฐิโต คนฺตฺวา ชนฺตาฆรปีฐํ ทตฺวา จีวรํ ปฏิคฺคเหตฺวา เอกมนฺตํ นิกฺขิปิตพฺพํ, จุณฺณํ ทาตพฺพํ, มตฺติกา ทาตพฺพา.\csromanspacer{} สเจ อุสฺสหติ, ชนฺตาฆรํ ปวิสิตพฺพํ, ชนฺตาฆรํ ปวิสนฺเตน มตฺติกาย มุขํ มกฺเขตฺวา ปุรโต จ ปจฺฉโต จ ปฏิจฺฉาเทตฺวา ชนฺตาฆรํ ปวิสิตพฺพํ, น เถเร ภิกฺขู อนุปขชฺช นิสีทิตพฺพํ, น นวา ภิกฺขู อาสเนน ปฏิพาหิตพฺพา, ชนฺตาฆเร อาจริยสฺส ปริกมฺมํ กาตพฺพํ, ชนฺตาฆรา นิกฺขมนฺเตน ชนฺตาฆรปีฐํ อาทาย ปุรโต จ ปจฺฉโต จ ปฏิจฺฉาเทตฺวา ชนฺตาฆรา นิกฺขมิตพฺพํ.}

\prose{อุทเกปิ อาจริยสฺส ปริกมฺมํ กาตพฺพํ, นหาเตน ปฐมตรํ อุตฺตริตฺวา อตฺตโน คตฺตํ โวทกํ กตฺวา นิวาเสตฺวา อาจริยสฺส คตฺตโต อุทกํ ปมชฺชิตพฺพํ, นิวาสนํ ทาตพฺพํ, สํฆาฏิ ทาตพฺพา, ชนฺตาฆรปีฐํ อาทาย ปฐมตรํ อาคนฺตฺวา อาสนํ ปญฺญเปตพฺพํ, ปาโททกํ ปาทปีฐํ ปาทกถลิกํ อุปนิกฺขิปิตพฺพํ, อาจริโย ปานีเยน ปุจฺฉิตพฺโพ.\csromanspacer{} สเจ อุทฺทิสาเปตุกาโม โหติ อุทฺทิสิตพฺโพ, สเจ ปริปุจฺฉิตุกาโม โหติ ปริปุจฺฉิตพฺโพ.}

\prose{\csromanfolio{404}ยสฺมิํ วิหาเร อาจริโย วิหรติ, สเจ โส วิหาโร อุกฺลาโป โหติ, สเจ อุสฺสหติ, โสเธตพฺโพ.\csromanspacer{} วิหารํ โสเธนฺเตน ปฐมํ ปตฺตจีวรํ นีหริตฺวา เอกมนฺตํ นิกฺขิปิตพฺพํ, นิสีทนปจฺจตฺถรณํ นีหริตฺวา เอกมนฺตํ นิกฺขิปิตพฺพํ, ภิสิพิพฺโพหนํ นีหริตฺวา เอกมนฺตํ นิกฺขิปิตพฺพํ, มญฺโจ นีจํ กตฺวา สาธุกํ อปฺปฏิฆํสนฺเตน อสํฆฏฺเฏนฺเตน กวาฏปิฏฺฐํ นีหริตฺวา เอกมนฺตํ นิกฺขิปิตพฺโพ, ปีฐํ นีจํ กตฺวา สาธุกํ อปฺปฏิฆํสนฺเตน อสํฆฏฺเฏนฺเตน กวาฏปิฏฺฐํ นีหริตฺวา เอกมนฺตํ นิกฺขิปิตพฺพํ, มญฺจปฏิปาทกา นีหริตฺวา เอกมนฺตํ นิกฺขิปิตพฺพา, เขฬมลฺลโก นีหริตฺวา เอกมนฺตํ นิกฺขิปิตพฺโพ, อปสฺเสนผลกํ นีหริตฺวา เอกมนฺตํ นิกฺขิปิตพฺพํ, ภูมตฺถรณํ ยถาปญฺญตฺตํ สลฺลกฺเขตฺวา นีหริตฺวา เอกมนฺตํ นิกฺขิปิตพฺพํ.\csromanspacer{} สเจ วิหาเร สนฺตานกํ โหติ, อุลฺโลกา ปฐมํ โอหาเรตพฺพํ, อาโลกสนฺธิกณฺณภาคา ปมชฺชิตพฺพา.\csromanspacer{} สเจ เครุกปริกมฺมกตา ภิตฺติ กณฺณกิตา โหติ, โจฬกํ เตเมตฺวา ปีเฬตฺวา ปมชฺชิตพฺพา.\csromanspacer{} สเจ กาฬวณฺณกตา ภูมิ กณฺณกิตา โหติ, โจฬกํ เตเมตฺวา ปีเฬตฺวา ปมชฺชิตพฺพา.\csromanspacer{} สเจ อกตา โหติ ภูมิ, อุทเกน ปริปฺโผสิตฺวา ปริปฺโผสิตฺวา สมฺมชฺชิตพฺพา “มา วิหาโร รเชน อุหญฺญี”ติ.\csromanspacer{} สงฺการํ วิจินิตฺวา เอกมนฺตํ ฉฑฺเฑตพฺพํ.}

\prose{ภูมตฺถรณํ โอตาเปตฺวา โสเธตฺวา ปปฺโผเฏตฺวา อติหริตฺวา ยถาปญฺญตฺตํ ปญฺญเปตพฺพํ, มญฺจปฏิปาทกา โอตาเปตฺวา ปมชฺชิตฺวา อติหริตฺวา ยถาฐาเน ฐเปตพฺพา, มญฺโจ โอตาเปตฺวา ปปฺโผเฏตฺวา นีจํ กตฺวา สาธุกํ อปฺปฏิฆํสนฺเตน อสํฆฏฺเฏนฺเตน กวาฏปิฏฺฐํ อติหริตฺวา ยถาปญฺญตฺตํ ปญฺญเปตพฺโพ, ปีฐํ โอตาเปตฺวา โสเธตฺวา ปปฺโผเฏตฺวา นีจํ กตฺวา สาธุกํ อปฺปฏิฆํสนฺเตน อสํฆฏฺเฏนฺเตน กวาฏปิฏฺฐํ อติหริตฺวา ยถาปญฺญตฺตํ ปญฺญเปตพฺพํ, ภิสิพิพฺโพหนํ โอตาเปตฺวา โสเธตฺวา ปปฺโผเฏตฺวา อติหริตฺวา ยถาปญฺญตฺตํ ปญฺญเปตพฺพํ, นิสีทนปจฺจตฺถรณํ โอตาเปตฺวา โสเธตฺวา ปปฺโผเฏตฺวา อติหริตฺวา ยถาปญฺญตฺตํ ปญฺญเปตพฺพํ, เขฬมลฺลโก โอตาเปตฺวา ปมชฺชิตฺวา อติหริตฺวา ยถาฐาเน ฐเปตพฺโพ, อปสฺเสนผลกํ โอตาเปตฺวา ปมชฺชิตฺวา อติหริตฺวา ยถาฐาเน ฐเปตพฺพํ.\csromanspacer{} ปตฺตจีวรํ นิกฺขิปิตพฺพํ, ปตฺตํ นิกฺขิปนฺเตน เอเกน หตฺเถน ปตฺตํ คเหตฺวา เอเกน หตฺเถน เหฏฺฐามญฺจํ วา เหฏฺฐาปีฐํ วา ปรามสิตฺวา ปตฺโต นิกฺขิปิตพฺโพ, น จ อนนฺตรหิตาย \csromanfolio{405}ภูมิยา ปตฺโต นิกฺขิปิตพฺโพ.\csromanspacer{} จีวรํ นิกฺขิปนฺเตน เอเกน หตฺเถน จีวรํ คเหตฺวา เอเกน หตฺเถน จีวรวํสํ วา จีวรรชฺชุํ วา ปมชฺชิตฺวา ปารโต อนฺตํ โอรโต โภคํ กตฺวา จีวรํ นิกฺขิปิตพฺพํ.}

\prose{สเจ ปุรตฺถิมา สรชา วาตา วายนฺติ, ปุรตฺถิมา วาตปานา ถเกตพฺพา.\csromanspacer{} สเจ ปจฺฉิมา สรชา วาตา วายนฺติ, ปจฺฉิมา วาตปานา ถเกตพฺพา.\csromanspacer{} สเจ อุตฺตรา สรชา วาตา วายนฺติ, อุตฺตรา วาตปานา ถเกตพฺพา.\csromanspacer{} สเจ ทกฺขิณา สรชา วาตา วายนฺติ, ทกฺขิณา วาตปานา ถเกตพฺพา.\csromanspacer{} สเจ สีตกาโล โหติ, ทิวา วาตปานา วิวริตพฺพา, รตฺติํ ถเกตพฺพา.\csromanspacer{} สเจ อุณฺหกาโล โหติ, ทิวา วาตปานา ถเกตพฺพา, รตฺติํ วิวริตพฺพา.}

\prose{สเจ ปริเวณํ อุกฺลาปํ โหติ, ปริเวณํ สมฺมชฺชิตพฺพํ.\csromanspacer{} สเจ โกฏฺฐโก อุกฺลาโป โหติ, โกฏฺฐโก สมฺมชฺชิตพฺโพ.\csromanspacer{} สเจ อุปฏฺฐานสาลา อุกฺลาปา โหติ, อุปฏฺฐานสาลา สมฺมชฺชิตพฺพา.\csromanspacer{} สเจ อคฺคิสาลา อุกฺลาปา โหติ, อคฺคิสาลา สมฺมชฺชิตพฺพา.\csromanspacer{} สเจ วจฺจกุฏิ อุกฺลาปา โหติ, วจฺจกุฏิ สมฺมชฺชิตพฺพา.\csromanspacer{} สเจ ปานียํ น โหติ, ปานียํ อุปฏฺฐาเปตพฺพํ, สเจ ปริโภชนียํ น โหติ, ปริโภชนียํ อุปฏฺฐาเปตพฺพํ.\csromanspacer{} สเจ อาจมนกุมฺภิยา อุทกํ น โหติ, อาจมนกุมฺภิยา อุทกํ อาสิญฺจิตพฺพํ.}

\prose{สเจ อาจริยสฺส อนภิรติ อุปฺปนฺนา โหติ, อนฺเตวาสิเกน วูปกาเสตพฺโพ วูปกาสาเปตพฺโพ ธมฺมกถา วาสฺส กาตพฺพา.\csromanspacer{} สเจ อาจริยสฺส กุกฺกุจฺจํ อุปฺปนฺนํ โหติ, อนฺเตวาสิเกน วิโนเทตพฺพํ วิโนทาเปตพฺพํ ธมฺมกถา วาสฺส กาตพฺพา.\csromanspacer{} สเจ อาจริยสฺส ทิฏฺฐิคตํ อุปฺปนฺนํ โหติ, อนฺเตวาสิเกน วิเวเจตพฺพํ วิเวจาเปตพฺพํ ธมฺมกถา วาสฺส กาตพฺพา.\csromanspacer{} สเจ อาจริโย ครุธมฺมํ อชฺฌาปนฺโน โหติ ปริวาสารโห, อนฺเตวาสิเกน อุสฺสุกฺกํ กาตพฺพํ “กินฺติ นุ โข สํโฆ อาจริยสฺส ปริวาสํ ทเทยฺยา”ติ.\csromanspacer{} สเจ อาจริโย มูลายปฏิกสฺสนารโห โหติ, อนฺเตวาสิเกน อุสฺสุกฺกํ กาตพฺพํ “กินฺติ นุ โข สํโฆ อาจริยํ มูลาย ปฏิกสฺเสยฺยา”ติ.\csromanspacer{} สเจ อาจริโย มานตฺตารโห โหติ, อนฺเตวาสิเกน อุสฺสุกฺกํ กาตพฺพํ “กินฺติ นุ โข สํโฆ อาจริยสฺส มานตฺตํ ทเทยฺยา”ติ.\csromanspacer{} สเจ อาจริโย อพฺภานารโห โหติ, อนฺเตวาสิเกน อุสฺสุกฺกํ กาตพฺพํ “กินฺติ นุ โข สํโฆ อาจริยํ อพฺเภยฺยา”ติ.\csromanspacer{} สเจ สํโฆ \csromanfolio{406}อาจริยสฺส กมฺมํ กตฺตุกาโม โหติ ตชฺชนียํ วา นิยสฺสํ วา ปพฺพาชนียํ วา ปฏิสารณียํ วา อุกฺเขปนียํ วา, อนฺเตวาสิเกน อุสฺสุกฺกํ กาตพฺพํ “กินฺติ นุ โข สํโฆ อาจริยสฺส กมฺมํ น กเรยฺย ลหุกาย วา ปริณาเมยฺยา”ติ.\csromanspacer{} กตํ วา ปนสฺส โหติ สํเฆน กมฺมํ ตชฺชนียํ วา นิยสฺสํ วา ปพฺพาชนียํ วา ปฏิสารณียํ อุกฺเขปนียํ วา, อนฺเตวาสิเกน อุสฺสุกฺกํ กาตพฺพํ “กินฺติ นุ โข อาจริโย สมฺมา วตฺเตยฺย, โลมํ ปาเตยฺย, เนตฺถารํ วตฺเตยฺย, สํโฆ ตํ กมฺมํ ปฏิปฺปสฺสมฺเภยฺยา”ติ.}

\prose{สเจ อาจริยสฺส จีวรํ โธวิตพฺพํ โหติ, อนฺเตวาสิเกน โธวิตพฺพํ, อุสฺสุกฺกํ วา กาตพฺพํ “กินฺติ นุ โข อาจริยสฺส จีวรํ โธวิเยถา”ติ.\csromanspacer{} สเจ อาจริยสฺส จีวรํ กาตพฺพํ โหติ, อนฺเตวาสิเกน กาตพฺพํ, อุสฺสุกฺกํ วา กาตพฺพํ “กินฺติ นุ โข อาจริยสฺส จีวรํ กริเยถา”ติ.\csromanspacer{} สเจ อาจริยสฺส รชนํ ปจิตพฺพํ โหติ, อนฺเตวาสิเกน ปจิตพฺพํ, อุสฺสุกฺกํ วา กาตพฺพํ “กินฺติ นุ โข อาจริยสฺส รชนํ ปจิเยถา”ติ.\csromanspacer{} สเจ อาจริยสฺส จีวรํ รชิตพฺพํ โหติ, อนฺเตวาสิเกน รชิตพฺพํ, อุสฺสุกฺกํ วา กาตพฺพํ “กินฺติ นุ โข อาจริยสฺส จีวรํ รชิเยถา”ติ.\csromanspacer{} จีวรํ รชนฺเตน สาธุกํ สมฺปริวตฺตกํ สมฺปริวตฺตกํ รชิตพฺพํ, น จ อจฺฉินฺเน เถว ปกฺกมิตพฺพํ.}

\prose{น อาจริยํ อนาปุจฺฉา เอกจฺจสฺส ปตฺโต ทาตพฺโพ, น เอกจฺจสฺส ปตฺโต ปฏิคฺคเหตพฺโพ, น เอกจฺจสฺส จีวรํ ทาตพฺพํ, น เอกจฺจสฺส จีวรํ ปฏิคฺคเหตพฺพํ, น เอกจฺจสฺส ปริกฺขาโร ทาตพฺโพ, น เอกจฺจสฺส ปริกฺขาโร ปฏิคฺคเหตพฺโพ, น เอกจฺจสฺส เกสา เฉทิตพฺพา, น เอกจฺเจน เกสา เฉทาเปตพฺพา, น เอกจฺจสฺส ปริกมฺมํ กาตพฺพํ, น เอกจฺเจน ปริกมฺมํ การาเปตพฺพํ, น เอกจฺจสฺส เวยฺยาวจฺโจ กาตพฺโพ, น เอกจฺเจน เวยฺยาวจฺโจ การาเปตพฺโพ, น เอกจฺจสฺส ปจฺฉาสมเณน โหตพฺพํ, น เอกจฺโจ ปจฺฉาสมโณ อาทาตพฺโพ, น เอกจฺจสฺส ปิณฺฑปาโต นีหริตพฺโพ, น เอกจฺเจน ปิณฺฑปาโต นีหราเปตพฺโพ, น อาจริยํ อนาปุจฺฉา คาโม ปวิสิตพฺโพ, น สุสานํ คนฺตพฺพํ, น ทิสา ปกฺกมิตพฺพา, สเจ อาจริโย คิลาโน โหติ, ยาวชีวํ \csromanfolio{407}อุปฏฺฐาตพฺโพ, วุฏฺฐานมสฺส อาคเมตพฺพํ.\csromanspacer{} อิทํ โข ภิกฺขเว อนฺเตวาสิกานํ อาจริเยสุ วตฺตํ, ยถา อนฺเตวาสิเกหิ อาจริเยสุ สมฺมา วตฺติตพฺพนฺติ.}

\csromanmatikaanchor{mtk.190}
\csromantocmarkref{section}{14. อนฺเตวาสิกวตฺตกถา}{mtk.190}
\bookmark[dest={mtk.190},level=1]{14. อนฺเตวาสิกวตฺตกถา}
\csromanheader{1}{14. อนฺเตวาสิกวตฺตกถา}
\proseitem{381}{เตน โข ปน สมเยน อาจริยา อนฺเตวาสิเกสุ น สมฺมา วตฺตนฺติ.\csromanspacer{} เย เต ภิกฺขู อปฺปิจฺฉา ฯเปฯ เต อุชฺฌายนฺติ ขิยฺยนฺติ วิปาเจนฺติ “กถํ หิ นาม อาจริยา อนฺเตวาสิเกสุ น สมฺมา วตฺติสฺสนฺตี”ติ.\csromanspacer{} อถ โข เต ภิกฺขู ภควโต เอตมตฺถํ อาโรเจสุํ.\csromanspacer{} อถ โข ภควา เอตสฺมิํ นิทาเน เอตสฺมิํ ปกรเณ ภิกฺขุสํฆํ สนฺนิปาตาเปตฺวา ภิกฺขู ปฏิปุจฺฉิ “สจฺจํ กิร ภิกฺขเว อาจริยา อนฺเตวาสิเกสุ น สมฺมา วตฺตนฺตี”ติ.\csromanspacer{} สจฺจํ ภควาติ ฯเปฯ วิครหิตฺวา ฯเปฯ ธมฺมิํ กถํ กตฺวา ภิกฺขู อามนฺเตสิ\csromandash{}}

\proseitem{382}{เตน หิ ภิกฺขเว อาจริยานํ อนฺเตวาสิเกสุ วตฺตํ ปญฺญเปสฺสามิ, ยถา อาจาริเยหิ อนฺเตวาสิเกสุ สมฺมา วตฺติตพฺพํ. * อาจริเยน ภิกฺขเว อนฺเตวาสิกมฺหิ สมฺมา วตฺติตพฺพํ.\csromanspacer{} ตตฺรายํ สมฺมาวตฺตนา\csromandash{}}

\prose{อาจริเยน ภิกฺขเว อนฺเตวาสิโก สงฺคเหตพฺโพ อนุคฺคเหตพฺโพ อุทฺเทเสน ปริปุจฺฉาย โอวาเทน อนุสาสนิยา, สเจ อาจริยสฺส ปตฺโต โหติ อนฺเตวาสิกสฺส ปตฺโต น โหติ, อาจริเยน อนฺเตวาสิกสฺส ปตฺโต ทาตพฺโพ, อุสฺสุกฺกํ วา กาตพฺพํ “กินฺติ นุ โข อนฺเตวาสิกสฺส ปตฺโต อุปฺปชฺชิเยถา”ติ.\csromanspacer{} สเจ อาจริยสฺส จีวรํ โหติ อนฺเตวาสิกสฺส จีวรํ น โหติ, อาจริเยน อนฺเตวาสิกสฺส จีวรํ ทาตพฺพํ, อุสฺสุกฺกํ วา กาตพฺพํ “กินฺติ นุ โข อนฺเตวาสิกสฺส จีวรํ อุปฺปชฺชิเยถา”ติ.\csromanspacer{} สเจ อาจริยสฺส ปริกฺขาโร โหติ อนฺเตวาสิกสฺส ปริกฺขาโร น โหติ, อาจริเยน อนฺเตวาสิกสฺส ปริกฺขาโร ทาตพฺโพ, อุสฺสุกฺกํ วา กาตพฺพํ “กินฺติ นุ โข อนฺเตวาสิกสฺส ปริกฺขาโร อุปฺปชฺชิเยถา”ติ.}

\prose{สเจ อนฺเตวาสิโก คิลาโน โหติ, กลาสฺเสว อุฏฺฐาย ทนฺตกฏฺฐํ ทาตพฺพํ, มุโขทกํ ทาตพฺพํ, อาสนํ ปญฺญเปตพฺพํ, สเจ ยาคุ \csromanfolio{408}โหติ, ภาชนํ โธวิตฺวา ยาคุ อุปนาเมตพฺพา, ยาคุํ ปีตสฺส อุทกํ ทตฺวา ภาชนํ ปฏิคฺคเหตฺวา นีจํ กตฺวา สาธุกํ อปฺปฏิฆํสนฺเตน โธวิตฺวา ปฏิสาเมตพฺพํ, อนฺเตวาสิกมฺหิ วุฏฺฐิเต อาสนํ อุทฺธริตพฺพํ, สเจ โส เทโส อุกฺลาโป โหติ, โส เทโส สมฺมชฺชิตพฺโพ.}

\prose{สเจ อนฺเตวาสิโก คามํ ปวิสิตุกาโม โหติ, นิวาสนํ ทาตพฺพํ, ปฏินิวาสนํ ปฏิคฺคเหตพฺพํ, กายพนฺธนํ ทาตพฺพํ, สคุณํ กตฺวา สํฆาฏิโย ทาตพฺพา, โธวิตฺวา ปตฺโต โสทโก ทาตพฺโพ.}

\prose{“เอตฺตาวตา นิวตฺติสฺสตี”ติ อาสนํ ปญฺญเปตพฺพํ, ปาโททกํ ปาทปีฐํ ปาทกถลิกํ อุปนิกฺขิปิตพฺพํ, ปจฺจุคฺคนฺตฺวา ปตฺตจีวรํ ปฏิคฺคเหตพฺพํ, ปฏินิวาสนํ ทาตพฺพํ, นิวาสนํ ปฏิคฺคเหตพฺพํ, สเจ จีวรํ สินฺนํ โหติ, มุหุตฺตํ อุเณฺห โอตาเปตพฺพํ, น จ อุเณฺห จีวรํ นิทหิตพฺพํ, จีวรํ สงฺฆริตพฺพํ, จีวรํ สงฺฆรนฺเตน จตุรงฺคุลํ กณฺณํ อุสฺสาเรตฺวา จีวรํ สงฺฆริตพฺพํ “มา มชฺเฌ ภงฺโค อโหสี”ติ.\csromanspacer{} โอโภเค กายพนฺธนํ กาตพฺพํ.}

\prose{สเจ ปิณฺฑปาโต โหติ อนฺเตวาสิโก จ ภุญฺชิตุกาโม โหติ, อุทกํ ทตฺวา ปิณฺฑปาโต อุปนาเมตพฺโพ, อนฺเตวาสิโก ปานีเยน ปุจฺฉิตพฺโพ, ภุตฺตาวิสฺส อุทกํ ทตฺวา ปตฺตํ ปฏิคฺคเหตฺวา นีจํ กตฺวา สาธุกํ อปฺปฏิฆํสนฺเตน โธวิตฺวา โวทกํ กตฺวา มุหุตฺตํ อุเณฺห โอตาเปตพฺโพ, น จ อุเณฺห ปตฺโต นิทหิตพฺโพ.\csromanspacer{} ปตฺตจีวรํ นิกฺขิปิตพฺพํ, ปตฺตํ นิกฺขิปนฺเตน ฯเปฯ จีวรํ นิกฺขิปนฺเตน ฯเปฯ ปารโต อนฺตํ โอรโต โภคํ กตฺวา จีวรํ นิกฺขิปิตพฺพํ.\csromanspacer{} อนฺเตวาสิกมฺหิ อุฏฺฐิเต อาสนํ อุทฺธริตพฺพํ, ปาโททกํ ปาทปีฐํ ปาทกถลิกํ ปฏิสาเมตพฺพํ, สเจ โส เทโส อุกฺลาโป โหติ, โส เทโส สมฺมชฺชิตพฺโพ.}

\prose{สเจ อนฺเตวาสิโก นหายิตุกาโม โหติ, นหานํ ปฏิยาเทตพฺพํ, สเจ สีเตน อตฺโถ โหติ, สีตํ ปฏิยาเทตพฺพํ, สเจ อุเณฺหน อตฺโถ โหติ, อุณฺหํ ปฏิยาเทตพฺพํ.}

\prose{สเจ อนฺเตวาสิโก ชนฺตาฆรํ ปวิสิตุกาโม โหติ, จุณฺณํ สนฺเนตพฺพํ, มตฺติกา เตเมตพฺพา, ชนฺตาฆรปีฐํ อาทาย\footnote{อาทาย อนฺเตวาสิกสฺส ปิฏฺฐิโต ปิฏฺฐิโต (ก)} คนฺตฺวา ชนฺตาฆรปีฐํ ทตฺวา จีวรํ ปฏิคฺคเหตฺวา เอกมนฺตํ นิกฺขิปิตพฺพํ, จุณฺณํ ทาตพฺพํ, มตฺติกา ทาตพฺพา, สเจ อุสฺสหติ, ชนฺตาฆรํ ปวิสิตพฺพํ, ชนฺตาฆรํ \csromanfolio{409}ปวิสนฺเตน มตฺติกาย มุขํ มกฺเขตฺวา ปุรโต จ ปจฺฉโต จ ปฏิจฺฉาเทตฺวา ชนฺตาฆรํ ปวิสิตพฺพํ, น เถเร ภิกฺขู อนุปขชฺช นิสีทิตพฺพํ, น นวา ภิกฺขู อาสเนน ปฏิพาหิตพฺพา, ชนฺตาฆเร อนฺเตวาสิกสฺส ปริกมฺมํ กาตพฺพํ, ชนฺตาฆรา นิกฺขมนฺเตน ชนฺตาฆรปีฐํ อาทาย ปุรโต จ ปจฺฉโต จ ปฏิจฺฉาเทตฺวา ชนฺตาฆรา นิกฺขมิตพฺพํ.}

\prose{อุทเกปิ อนฺเตวาสิกสฺส ปริกมฺมํ กาตพฺพํ, นหาเตน ปฐมตรํ อุตฺตริตฺวา อตฺตโน คตฺตํ โวทกํ กตฺวา นิวาเสตฺวา อนฺเตวาสิกสฺส คตฺตโต อุทกํ ปมชฺชิตพฺพํ, นิวาสนํ ทาตพฺพํ, สํฆาฏิ ทาตพฺพา, ชนฺตาฆรปีฐํ อาทาย ปฐมตรํ อาคนฺตฺวา อาสนํ ปญฺญเปตพฺพํ, ปาโททกํ ปาทปีฐํ ปาทกถลิกํ อุปนิกฺขิปิตพฺพํ, อนฺเตวาสิโก ปานีเยน ปุจฺฉิตพฺโพ.}

\prose{ยสฺมิํ วิหาเร อนฺเตวาสิโก วิหรติ, สเจ โส วิหาโร อุกฺลาโป โหติ, สเจ อุสฺสหติ, โสเธตพฺโพ.\csromanspacer{} วิหารํ โสเธนฺเตน ปฐมํ ปตฺตจีวรํ นีหริตฺวา เอกมนฺตํ นิกฺขิปิตพฺพํ ฯเปฯ อาจมนกุมฺภิยา อุทกํ น โหติ, อาจมนกุมฺภิยา อุทกํ อาสิญฺจิตพฺพํ.}

\csromanmatikaanchor{mtk.191}
\csromantocmarkref{section}{อุทฺทานคาถา}{mtk.191}
\bookmark[dest={mtk.191},level=1]{อุทฺทานคาถา}
\prose{สเจ อนฺเตวาสิกสฺส อนภิรติ อุปฺปนฺนา โหติ, อาจริเยน วูปกาเสตพฺโพ วูปกาสาเปตพฺโพ ธมฺมกถา วาสฺส กาตพฺพา.\csromanspacer{} สเจ อนฺเตวาสิกสฺส กุกฺกุจฺจํ อุปฺปนฺนํ โหติ, อาจริเยน วิโนเทตพฺพํ วิโนทาเปตพฺพํ ธมฺมกถา วาสฺส กาตพฺพา.\csromanspacer{} สเจ อนฺเตวาสิกสฺส ทิฏฺฐิคตํ อุปฺปนฺนํ โหติ, อาจริเยน วิเวเจตพฺพํ วิเวจาเปตพฺพํ ธมฺมกถา วาสฺส กาตพฺพา.\csromanspacer{} สเจ อนฺเตวาสิโก ครุธมฺมํ อชฺฌาปนฺโน โหติ ปริวาสารโห, อาจริเยน อุสฺสุกฺกํ กาตพฺพํ “กินฺติ นุ โข สํโฆ อนฺเตวาสิกสฺส ปริวาสํ ทเทยฺยา”ติ.\csromanspacer{} สเจ อนฺเตวาสิโก มูลายปฏิกสฺสนารโห โหติ, อาจริเยน อุสฺสุกฺกํ กาตพฺพํ “กินฺติ นุ โข สํโฆ อนฺเตวาสิกํ มูลาย ปฏิกสฺเสยฺยา”ติ.\csromanspacer{} สเจ อนฺเตวาสิโก มานตฺตารโห โหติ, อาจริเยน อุสฺสุกฺกํ กาตพฺพํ “กินฺติ นุ โข สํโฆ อนฺเตวาสิกสฺส มานตฺตํ ทเทยฺยา”ติ.\csromanspacer{} สเจ อนฺเตวาสิโก อพฺภานารโห โหติ, อาจริเยน อุสฺสุกฺกํ กาตพฺพํ “กินฺติ นุ โข สํโฆ อนฺเตวาสิกํ อพฺเภยฺยา”ติ.\csromanspacer{} สเจ สํโฆ อนฺเตวาสิกสฺส กมฺมํ กตฺตุกาโม โหติ ตชฺชนียํ วา นิยสฺสํ วา ปพฺพาชนียํ วา ปฏิสารณียํ วา อุกฺเขปนียํ วา, อาจริเยน อุสฺสุกฺกํ กาตพฺพํ “กินฺติ นุ โข สํโฆ อนฺเตวาสิกสฺส กมฺมํ น \csromanfolio{410}กเรยฺย ลหุกาย วา ปริณาเมยฺยา”ติ.\csromanspacer{} กตํ วา ปนสฺส โหติ สํเฆน กมฺมํ ตชฺชนียํ วา นิยสฺสํ วา ปพฺพาชนียํ วา ปฏิสารณียํ วา อุกฺเขปนียํ วา, อาจริเยน อุสฺสุกฺกํ กาตพฺพํ “กินฺติ นุ โข อนฺเตวาสิโก สมฺมา วตฺเตยฺย, โลมํ ปาเตยฺย, เนตฺถารํ วตฺเตยฺย, สํโฆ ตํ กมฺมํ ปฏิปฺปสฺสมฺเภยฺยา”ติ.}

\prose{สเจ อนฺเตวาสิกสฺส จีวรํ โธวิตพฺพํ โหติ, อาจริเยน อาจิกฺขิตพฺพํ “เอวํ โธเวยฺยาสี”ติ, อุสฺสุกฺกํ วา กาตพฺพํ “กินฺติ นุ โข อนฺเตวาสิกสฺส จีวรํ โธวิเยถา”ติ.\csromanspacer{} สเจ อนฺเตวาสิกสฺส จีวรํ กาตพฺพํ โหติ, อาจริเยน อาจิกฺขิตพฺพํ “เอวํ กเรยฺยาสี”ติ, อุสฺสุกฺกํ วา กาตพฺพํ “กินฺติ นุ โข อนฺเตวาสิกสฺส จีวรํ กริเยถา”ติ.\csromanspacer{} สเจ อนฺเตวาสิกสฺส รชนํ ปจิตพฺพํ โหติ, อาจริเยน อาจิกฺขิตพฺพํ “กฺวํ ปเจยฺยาสี”ติ, อุสฺสุกฺกํ วา กาตพฺพํ “กินฺติ นุ โข อนฺเตวาสิกสฺส รชนํ ปจิเยถา”ติ.\csromanspacer{} สเจ อนฺเตวาสิกสฺส จีวรํ รชิตพฺพํ โหติ, อาจริเยน อาจิกฺขิตพฺพํ “เอวํ รเชยฺยาสี”ติ, อุสฺสุกฺกํ วา กาตพฺพํ “กินฺติ นุ โข อนฺเตวาสิกสฺส จีวรํ รชิเยถา”ติ.\csromanspacer{} จีวรํ รชนฺเตน สาธุกํ สมฺปริวตฺตกํ สมฺปริวตฺตกํ รชิตพฺพํ, น จ อจฺฉินฺเน เถเว ปกฺกมิตพฺพํ.\csromanspacer{} สเจ อนฺเตวาสิโก คิลาโน โหติ, ยาวชีวํ อุปฏฺฐาตพฺโพ, วุฏฺฐานมสฺส อาคเมตพฺพํ.\csromanspacer{} อิทํ โข ภิกฺขเว อาจริยานํ อนฺเตวาสิเกสุ วตฺตํ, ยถา อาจริเยหิ อนฺเตวาสิเกสุ สมฺมา วตฺติตพฺพนฺติ.}

\vspace{\tipitakaparskip}
\csromancenter{วตฺตกฺขนฺธโก อฏฺฐโม.}
\vspace{0.5\baselineskip}
\csromangathameasure{อิมมฺหิ ขนฺธเก วตฺถู เอกูนวีสติ, วตฺตา จุทฺทส.}
\csromangathagroup{\csromangathastanza{อิมมฺหิ ขนฺธเก วตฺถู เอกูนวีสติ, วตฺตา จุทฺทส.}}

\titlehead{\textbf{ตสฺสุทฺทานํ}}
\csromangathameasure{ส-อุปาหนา ฉตฺตา จ โอคุณฺฐิ สีสํ ปานียํ. นาภิวาเท น ปุจฺฉนฺติ, \\ อหิ อุชฺฌนฺติ เปสลา. โอมุญฺจิ ฉตฺตํ ขนฺเธ จ, \\ อตรญฺจ ปฏิกฺกมํ. ปตฺตจีวรํ นิกฺขิปา, \\ ปติรูปญฺจ ปุจฺฉิตา. อาสิญฺเจยฺย โธวิเตน, \\ สุกฺเขนลฺเลนุปาหนา. วุฑฺโฒ นวโก ปุจฺเฉยฺย, \\ อชฺฌาวุฏฺฐญฺจ โคจรา. เสกฺขา วจฺจา ปานี ปริ, \\ กตฺตรํ กติกํ ตโต. กาลํ มุหุตฺตํ อุกฺลาโป, \\ ภูมตฺถรณํ นีหเร. ปฏิปาโท ภิสิพิพฺโพ, \\ มญฺจปีฐญฺจ มลฺลกํ. อปสฺเสนุลฺโลกกณฺณา, \\ เครุกา กาฬ อกตา. สงฺการญฺจ ภูมตฺถรณํ, \\ ปฏิปาทกํ มญฺจปีฐํ. ภิสิ นิสีทนมฺปิ, \\ มลฺลกํ อปสฺเสน จ. ปตฺตจีวรํ ภูมิ จ, \\ ปารนฺตํ โอรโต โภคํ. ปุรตฺถิมา ปจฺฉิมา จ, \\ อุตฺตรา อถ ทกฺขิณา. สีตุเณฺห จ ทิวารตฺติํ, \\ ปริเวณญฺจ โกฏฺฐโก. อุปฏฺฐานคฺคิ สาลา จ, \\ วตฺตํ วจฺจกุฏีสุ จ. ปานี ปริโภชนิยา, \\ กุมฺภิ อาจมเนสุ จ. อโนปเมน ปญฺญตฺตํ, \\ วตฺตํ อาคนฺตุเกหิเม. เนวาสนํ น อุทกํ, \\ น ปจฺจุ น จ ปานิยํ. นาภิวาเท นปญฺญเป, \\ อุชฺฌายนฺติ จ เปสลา. วุฑฺฒาสนญฺจ อุทกํ, \\ ปจฺจุคฺคนฺตฺวา จ ปานิยํ. อุปาหเน เอกมนฺตํ, \\ อภิวาเท จ ปญฺญเป. วุตฺถํ โคจรเสกฺโข จ, \\ ฐานํ ปานียโภชนํ. กตฺตรํ กติกํ กาลํ, \\ นวกสฺส นิสินฺนเก. อภิวาทเย อาจิกฺเข, \\ ยถา เหฏฺฐา ตถา นเย. นิทฺทิฏฺฐํ สตฺถวาเหน, \\ วตฺตํ อาวาสิเกหิเม. คมิกา ทารุมตฺติ จ, \\ วิวริตฺวา น ปุจฺฉิย. นสฺสนฺติ จ อคุตฺตญฺจ, \\ อุชฺฌายนฺติ จ เปสลา. ปฏิสาเมตฺวา ถเกตฺวา, \\ อาปุจฺฉิตฺวาว ปกฺกเม. ภิกฺขุ วา สามเณโร วา, \\ อารามิโก อุปาสโก. ปาสาณเกสุ จ ปุญฺชํ, \\ ปฏิสาเมถเกยฺย จ. สเจ อุสฺสหติ อุสฺสุกฺกํ, \\ อโนวสฺเส ตเถว จ. สพฺโพ โอวสฺสติ คามํ, \\ อชฺโฌกาเส ตเถว จ. อปฺเปวงฺคานิ เสเสยฺยุํ, \\ วตฺตํ คมิกภิกฺขุนา. นานุโมทนฺติ เถเรน, \\ โอหาย จตุปญฺจหิ. วจฺจิโต มุจฺฉิโต อาสิ, \\ วตฺตานุโมทเนสุเม. ฉพฺพคฺคิยา ทุนฺนิวตฺถา, \\ อโถปิ จ ทุปฺปารุกา. อนากปฺปา จ โวกฺกมฺม, \\ เถเร จ อนุปขชฺชเน. นเว ภิกฺขู จ สํฆาฏิ, \\ อุชฺฌายนฺติ จ เปสลา. ติมณฺฑลํ นิวาเสตฺวา, \\ กายสคุณคณฺฐิกา. น โวกฺกมฺม ปฏิจฺฉนฺนํ, \\ สุสํวุโตกฺขิตฺตจกฺขุ. อุกฺขิตฺโตชฺชคฺฆิกาสทฺโท, \\ ตโย เจว ปจาลนา. ขมฺโภคุณฺฐิ-อุกฺกุฏิกา, \\ ปฏิจฺฉนฺนํ สุสํวุโต. โอกฺขิตฺตุกฺขิตฺต-อุชฺชคฺฆิ, \\ อปฺปสทฺโท ตโย จลา. ขมฺโภคุณฺฐิปลฺลตฺถิ จ, \\ อนุปขชฺช นาสเน. โอตฺถริตฺวาน อุทเก, \\ นีจํ กตฺวาน สิญฺจิยา. ปฏิ สามนฺตา สํฆาฏิ, \\ โอทเน จ ปฏิคฺคเห. สูปํ อุตฺตริภงฺเคน, \\ สพฺเพสํ สมติตฺถิ จ. สกฺกจฺจํ, \\ ปตฺตสญฺญี จ, สปทานญฺจ สูปกํ. น ถูปโต ปฏิจฺฉาเท, \\ วิญฺญตฺตุชฺฌานสญฺญินา. มหนฺตมณฺฑลทฺวารํ, \\ สพฺพหตฺโถ น พฺยาหเร. อุกฺเขโป เฉทนาคณฺฑ, \\ ธุนํ สิตฺถาวการกํ. ชิวฺหานิจฺฉารกญฺเจว, \\ จปุจปุ สุรุสุรุ. หตฺถปตฺโตฏฺฐนิลฺเลหํ, \\ สามิเสน ปฏิคฺคเห. ยาว น สพฺเพ อุทเก, \\ นีจํ กตฺวาน สิญฺจิยํ. ปฏิ สามนฺตา สํฆาฏิ, \\ นีจํ กตฺวา ฉมาย จ. สสิตฺถกํ นิวตฺตนฺเต, \\ สุปฺปฏิจฺฉนฺนมุกฺกุฏิ. ธมฺมราเชน ปญฺญตฺตํ, \\ อิทํ ภตฺตคฺควตฺตนํ. ทุนฺนิวตฺถา อนากปฺปา, \\ อสลฺลกฺเขตฺวา จ สหสา. ทูเร อจฺจ จิรํ ลหุํ, \\ ตเถว ปิณฺฑจาริโก. ปฏิจฺฉนฺโนว คจฺเฉยฺย, \\ สุสํวุโตกฺขิตฺตจกฺขุ. อุกฺขิตฺโตชฺชคฺฆิกาสทฺโท, \\ ตโย เจว ปจาลนา. ขมฺโภคุณฺฐิ-อุกฺกุฏิกา, \\ สลฺลกฺเขตฺวา จ สหสา. ทูเร อจฺจ จิรํ ลหุํ, \\ อาสนกํ กฏจฺฉุกา. ภาชนํ วา ฐเปติ จ, \\ อุจฺจาเรตฺวา ปณาเมตฺวา. ปฏิคฺคเห น อุลฺโลเก, \\ สูเปสุปิ ตเถว ตํ. ภิกฺขุ สํฆาฏิยา ฉาเท, \\ ปฏิจฺฉนฺเนว คจฺฉิยํ. สํวุโตกฺขิตฺตจกฺขุ จ, \\ อุกฺขิตฺโตชฺชคฺฆิกาย จ. อปฺปสทฺโท ตโย จาลา, \\ ขมฺโภคุณฺฐิก-อุกฺกุฏิ. ปฐมาสนวกฺการ, \\ ปานิยํ ปริโภชนี. ปจฺฉากงฺขติ ภุญฺเชยฺย, \\ โอปิลาเปยฺย อุทฺธเร. ปฏิสาเมยฺย สมฺมชฺเช, \\ ริตฺตํ ตุจฺฉํ อุปฏฺฐเป. หตฺถวิกาเร ภินฺเทยฺย, \\ วตฺติทํ ปิณฺฑจาริเก. ปานี ปริ อคฺคิรณิ, \\ นกฺขตฺตทิสโจรา จ. “สพฺพํ นตฺถี”ติ โกฏฺเฏตฺวา, \\ ปตฺตํเส จีวรํ ตโต. อิทานิ อํเส ลคฺเคตฺวา, \\ ติมณฺฑลํ ปริมณฺฑลํ. ยถา ปิณฺฑจาริวตฺตํ, \\ นเย อารญฺญเกสุปิ. ปตฺตํเส จีวรํ สีเส, \\ อาโรหิตฺวา จ ปานิยํ. ปริโภชนิยํ อคฺคิ, \\ อรณี จาปิ กตฺตรํ. นกฺขตฺตํ สปฺปเทสํ วา, \\ ทิสาปิ กุสโล ภเว. สตฺตุตฺตเมน ปญฺญตฺตํ, \\ วตฺตํ อรญฺญเกสุเม. อชฺโฌกาเส โอกิริํสุ, \\ อุชฺฌายนฺติ จ เปสลา. สเจ วิหาโร อุกฺลาโป, \\ ปฐมํ ปตฺตจีวรํ. ภิสิพิพฺโพหนํ มญฺจํ, \\ ปีฐญฺจ เขฬมลฺลกํ. อปสฺเสนุลฺโลกกณฺณา, \\ เครุกา กาฬ อกตา. สงฺการํ ภิกฺขุสามนฺตา, \\ เสนาวิหารปานิยํ. ปริโภชนสามนฺตา, \\ ปฏิวาเต จ องฺคเณ. อโธวาเต อตฺถรณํ, \\ ปฏิปาทกมญฺโจ จ. ปีฐํ ภิสิ นิสีทนํ, \\ มลฺลกํ อปสฺเสน จ. ปตฺตจีวรํ ภูมิ จ, \\ ปารนฺตํ โอรโต โภคํ. ปุรตฺถิมา จ ปจฺฉิมา, \\ อุตฺตรา อถ ทกฺขิณา. สีตุเณฺห จ ทิวา รตฺติํ, \\ ปริเวณญฺจ โกฏฺฐโก. อุปฏฺฐานคฺคิสาลา จ, \\ วจฺจกุฏี จ ปานิยํ. อาจมกุมฺภิ วุฑฺเฒ จ, \\ อุทฺเทสปุจฺฉนา สชฺฌา. ธมฺโม ปทีปํ วิชฺฌาเป, \\ น วิวเร นปิ ถเก. เยน วุฑฺโฒ ปริวตฺติ, \\ กณฺเณนปิ น ฆฏฺเฏเย. ปญฺญเปสิ มหาวีโร, \\ วตฺตํ เสนาสเนสุ ตํ. นิวาริยมานา ทฺวารํ, \\ มุจฺฉิตุชฺฌนฺติ เปสลา. ฉาริกํ ฉฑฺฑเย ชนฺตา, \\ ปริภณฺฑํ ตเถว จ. ปริเวณํ โกฏฺฐโก สาลา, \\ จุณฺณมตฺติกโทณิกา. มุขํ ปุรโต น เถเร, \\ น นเว อุสฺสหติ สเจ. ปุรโต อุปริมคฺโค, \\ จิกฺขลฺลํ มตฺติ ปีฐกํ. วิชฺฌาเปตฺวา ถเกตฺวา จ, \\ วตฺตํ ชนฺตาฆเรสุเม. นาจเมติ ยถาวุฑฺฒุํ, \\ ปฏิปาฏิ จ สหสา. อุพฺภชิ นิตฺถุโน กฏฺฐํ, \\ วจฺจํ ปสฺสาว เขฬกํ. ผรุสา กูป สหสา, \\ อุพฺภชิ จปุ เสเสน. พหิ อนฺโต จ อุกฺกาเส, \\ รชฺชุ อตรมานญฺจ. สหสา อุพฺภชิ ฐิเต, \\ นิตฺถุเน กฏฺฐ วจฺจญฺจ. ปสฺสาว เขฬ ผรุสา, \\ กูปญฺจ วจฺจปาทุเก. นติสหสา อุพฺภชิ, \\ ปาทุกาย จปุจปุ. น เสสเย ปฏิจฺฉาเท, \\ อุหตปิธเรน จ. วจฺจกุฏี ปริภณฺฑํ, \\ ปริเวณญฺจ โกฏฺฐโก. อาจมเน จ อุทกํ, \\ วตฺตํ วจฺจกุฏีสุเม. อุปาหนา ทนฺตกฏฺฐํ, \\ มุโขทกญฺจ อาสนํ. ยาคุ อุทกํ โธวิตฺวา, \\ อุทฺธารุกฺลาป คาม จ. นิวาสนา กายพนฺธา, \\ สคุณํ ปตฺตโสทกํ. ปจฺฉา ติมณฺฑโล เจว, \\ ปริมณฺฑล พนฺธนํ. สคุณํ โธวิตฺวา ปจฺฉา, \\ นาติทูเร ปฏิคฺคเห. ภณมานสฺส อาปตฺติ, \\ ปฐมาคนฺตฺวาน อาสนํ. อุทกํ ปีฐกถลิ, \\ ปจฺจุคฺคนฺตฺวา นิวาสนํ. โอตาเป นิทหิ ภงฺโค, \\ โอโภเค ภุญฺชิตุ นเม. ปานียํ อุทกํ นีจํ, \\ มุหุตฺตํ น จ นิทเห. ปตฺตจีวรํ ภูมิ จ, \\ ปารนฺตํ โอรโต โภคํ. อุทฺธเร ปฏิสาเม จ, \\ อุกฺลาโป จ นหายิตุํ. สีตํ อุณฺหํ ชนฺตาฆรํ, \\ จุณฺณํ มตฺติก ปิฏฺฐิโต. ปีฐญฺจ จีวรํ จุณฺณํ, \\ มตฺติกุสฺสหติ มุขํ. ปุรโต เถเร นเว จ, \\ ปริกมฺมญฺจ นิกฺขเม. ปุรโต อุทเก นฺหาเต, \\ นิวาเสตฺวา อุปชฺฌายํ. นิวาสนญฺจ สํฆาฏิ, \\ ปีฐกํ อาสเนน จ. ปาโท ปีฐํ กถลิญฺจ, \\ ปานียุทฺเทสปุจฺฉนา. อุกฺลาปํ สุโสเธยฺย, \\ ปฐมํ ปตฺตจีวรํ. นิสีทนปจฺจตฺถรณํ, \\ ภิสิ พิพฺโพหนานิ จ. มญฺโจ ปีฐํ ปฏิปาทํ, \\ มลฺลกํ อปสฺเสน จ. ภูม สนฺตาน อาโลก, \\ เครุกา กาฬ อกตา. ภูมตฺถรปฏิปาทา, \\ มญฺโจ ปีฐํ พิพฺโพหนํ. นิสีทถรณํ เขฬ, \\ อปสฺเส ปตฺตจีวรํ. ปุรตฺถิมา ปจฺฉิมา จ, \\ อุตฺตรา อถ ทกฺขิณา. สีตุณฺหญฺจ ทิวา รตฺติํ, \\ ปริเวณญฺจ โกฏฺฐโก. อุปฏฺฐานคฺคิสาลา จ, \\ วจฺจปานิยโภชนี. อาจมํ อนภิรติ, \\ กุกฺกุจฺจํ ทิฏฺฐิ จ ครุ. มูลมานตฺต-อพฺภานํ, \\ ตชฺชนียํ นิยสฺสกํ. ปพฺพาช ปฏิสารณี, \\ อุกฺเขปญฺจ กตํ ยทิ. โธเว กาตพฺพํ รชญฺจ, \\ รเช สมฺปริวตฺตกํ. ปตฺตญฺจ จีวรญฺจาปิ, \\ ปริกฺขารญฺจ เฉทนํ. ปริกมฺมํ เวยฺยาวจฺจํ, \\ ปจฺฉา ปิณฺฑํ ปวิสนํ. น สุสานํ ทิสา เจว, \\ ยาวชีวํ อุปฏฺฐเห. สทฺธิวิหาริเกเนตํ, \\ วตฺตุปชฺฌายเกสุเม. โอวาทสาสนุทฺเทสา, \\ ปุจฺฉา ปตฺตญฺจ จีวรํ. ปริกฺขาโร คิลาโน จ, \\ น ปจฺฉาสมโณ ภเว. อุปชฺฌาเยสุ เย วตฺตา, \\ เอวํ อาจริเยสุปิ. สทฺธิวิหาริเก วตฺตา, \\ ตเถว อนฺเตวาสิเก. อาคนฺตุเกสุ เย วตฺตา, \\ ปุน อาวาสิเกสุ จ. คมิกานุโมทนิกา, \\ ภตฺตคฺเค ปิณฺฑจาริเก. อารญฺญเกสุ ยํ วตฺตํ, \\ ยญฺจ เสนาสเนสุปิ. ชนฺตาฆเร วจฺจกุฏี, \\ อุปชฺฌา สทฺธิวิหาริเก. อาจริเยสุ ยํ วตฺตํ, \\ ตเถว อนฺเตวาสิเก. เอกูนวีสติ วตฺถู, \\ วตฺตา จุทฺทส ขนฺธเก. วตฺตํ อปริปูเรนฺโต, \\ น สีลํ ปริปูรติ. อสุทฺธสีโล ทุปฺปญฺโญ, \\ จิตฺเตกคฺคํ น วินฺทติ. วิกฺขิตฺตจิตฺโตเนกคฺโค, \\ สมฺมา ธมฺมํ น ปสฺสติ. อปสฺสมาโน สทฺธมฺมํ, \\ ทุกฺขา น ปริมุจฺจติ. ยํ วตฺตํ ปริปูเรนฺโต, \\ สีลมฺปิ ปริปูรติ. วิสุทฺธสีโล สปฺปญฺโญ, \\ จิตฺเตกคฺคมฺปิ วินฺทติ. อวิกฺขิตฺตจิตฺโต เอกคฺโค, \\ สมฺมา ธมฺมํ วิปสฺสติ. สมฺปสฺสมาโน สทฺธมฺมํ, \\ ทุกฺขา โส ปริมุจฺจติ. ตสฺมา หิ วตฺตํ ปูเรยฺย, \\ ชินปุตฺโต วิจฺจกฺขโณ. โอวาทํ พุทฺธเสฏฺฐสฺส, \\ ตโต นิพฺพานเมหิตีติ.}
\csromangathagroup{\csromangathastanza{ส-อุปาหนา ฉตฺตา จ โอคุณฺฐิ สีสํ ปานียํ. นาภิวาเท น ปุจฺฉนฺติ, \\ อหิ อุชฺฌนฺติ เปสลา. โอมุญฺจิ ฉตฺตํ ขนฺเธ จ,}\csromangathastanza{อตรญฺจ ปฏิกฺกมํ. ปตฺตจีวรํ นิกฺขิปา, \\ ปติรูปญฺจ ปุจฺฉิตา. อาสิญฺเจยฺย โธวิเตน,}\csromangathastanza{สุกฺเขนลฺเลนุปาหนา. วุฑฺโฒ นวโก ปุจฺเฉยฺย, \\ อชฺฌาวุฏฺฐญฺจ โคจรา. เสกฺขา วจฺจา ปานี ปริ,}\csromangathastanza{\csromanfolio{411}กตฺตรํ กติกํ ตโต. กาลํ มุหุตฺตํ อุกฺลาโป, \\ ภูมตฺถรณํ นีหเร. ปฏิปาโท ภิสิพิพฺโพ,}\csromangathastanza{มญฺจปีฐญฺจ มลฺลกํ. อปสฺเสนุลฺโลกกณฺณา, \\ เครุกา กาฬ อกตา. สงฺการญฺจ ภูมตฺถรณํ,}\csromangathastanza{ปฏิปาทกํ มญฺจปีฐํ. ภิสิ นิสีทนมฺปิ, \\ มลฺลกํ อปสฺเสน จ. ปตฺตจีวรํ ภูมิ จ,}\csromangathastanza{ปารนฺตํ โอรโต โภคํ. ปุรตฺถิมา ปจฺฉิมา จ, \\ อุตฺตรา อถ ทกฺขิณา. สีตุเณฺห จ ทิวารตฺติํ,}\csromangathastanza{ปริเวณญฺจ โกฏฺฐโก. อุปฏฺฐานคฺคิ สาลา จ, \\ วตฺตํ วจฺจกุฏีสุ จ. ปานี ปริโภชนิยา,}\csromangathastanza{กุมฺภิ อาจมเนสุ จ. อโนปเมน ปญฺญตฺตํ, \\ วตฺตํ อาคนฺตุเกหิเม\footnote{เว (ก, เอวมุปริปิ)}. เนวาสนํ น อุทกํ,}\csromangathastanza{น ปจฺจุ น จ ปานิยํ. นาภิวาเท นปญฺญเป, \\ อุชฺฌายนฺติ จ เปสลา. วุฑฺฒาสนญฺจ อุทกํ,}\csromangathastanza{ปจฺจุคฺคนฺตฺวา จ ปานิยํ. อุปาหเน เอกมนฺตํ, \\ อภิวาเท จ ปญฺญเป. วุตฺถํ โคจรเสกฺโข จ,}\csromangathastanza{ฐานํ ปานียโภชนํ. กตฺตรํ กติกํ กาลํ, \\ นวกสฺส นิสินฺนเก. อภิวาทเย อาจิกฺเข,}\csromangathastanza{ยถา เหฏฺฐา ตถา นเย. นิทฺทิฏฺฐํ สตฺถวาเหน, \\ วตฺตํ อาวาสิเกหิเม. คมิกา ทารุมตฺติ จ,}\csromangathastanza{วิวริตฺวา น ปุจฺฉิย. นสฺสนฺติ จ อคุตฺตญฺจ, \\ อุชฺฌายนฺติ จ เปสลา. ปฏิสาเมตฺวา ถเกตฺวา,}\csromangathastanza{อาปุจฺฉิตฺวาว ปกฺกเม. ภิกฺขุ วา สามเณโร วา, \\ อารามิโก อุปาสโก. ปาสาณเกสุ จ ปุญฺชํ,}\csromangathastanza{ปฏิสาเมถเกยฺย จ. สเจ อุสฺสหติ อุสฺสุกฺกํ, \\ อโนวสฺเส ตเถว จ. สพฺโพ โอวสฺสติ คามํ,}\csromangathastanza{\csromanfolio{412}อชฺโฌกาเส ตเถว จ. อปฺเปวงฺคานิ เสเสยฺยุํ, \\ วตฺตํ คมิกภิกฺขุนา. นานุโมทนฺติ เถเรน,}\csromangathastanza{โอหาย จตุปญฺจหิ. วจฺจิโต มุจฺฉิโต อาสิ, \\ วตฺตานุโมทเนสุเม. ฉพฺพคฺคิยา ทุนฺนิวตฺถา,}\csromangathastanza{อโถปิ จ ทุปฺปารุกา. อนากปฺปา จ โวกฺกมฺม, \\ เถเร จ อนุปขชฺชเน. นเว ภิกฺขู จ สํฆาฏิ,}\csromangathastanza{อุชฺฌายนฺติ จ เปสลา. ติมณฺฑลํ นิวาเสตฺวา, \\ กายสคุณคณฺฐิกา. น โวกฺกมฺม ปฏิจฺฉนฺนํ,}\csromangathastanza{สุสํวุโตกฺขิตฺตจกฺขุ. อุกฺขิตฺโตชฺชคฺฆิกาสทฺโท, \\ ตโย เจว ปจาลนา. ขมฺโภคุณฺฐิ-อุกฺกุฏิกา,}\csromangathastanza{ปฏิจฺฉนฺนํ สุสํวุโต. โอกฺขิตฺตุกฺขิตฺต-อุชฺชคฺฆิ, \\ อปฺปสทฺโท ตโย จลา. ขมฺโภคุณฺฐิปลฺลตฺถิ จ,}\csromangathastanza{อนุปขชฺช นาสเน. โอตฺถริตฺวาน อุทเก, \\ นีจํ กตฺวาน สิญฺจิยา. ปฏิ สามนฺตา สํฆาฏิ,}\csromangathastanza{โอทเน จ ปฏิคฺคเห. สูปํ อุตฺตริภงฺเคน, \\ สพฺเพสํ สมติตฺถิ จ. สกฺกจฺจํ,}\csromangathastanza{ปตฺตสญฺญี จ, สปทานญฺจ สูปกํ. น ถูปโต ปฏิจฺฉาเท, \\ วิญฺญตฺตุชฺฌานสญฺญินา. มหนฺตมณฺฑลทฺวารํ,}\csromangathastanza{สพฺพหตฺโถ น พฺยาหเร. อุกฺเขโป เฉทนาคณฺฑ, \\ ธุนํ สิตฺถาวการกํ. ชิวฺหานิจฺฉารกญฺเจว,}\csromangathastanza{จปุจปุ สุรุสุรุ. หตฺถปตฺโตฏฺฐนิลฺเลหํ, \\ สามิเสน ปฏิคฺคเห. ยาว น สพฺเพ อุทเก,}\csromangathastanza{นีจํ กตฺวาน สิญฺจิยํ. ปฏิ สามนฺตา สํฆาฏิ, \\ นีจํ กตฺวา ฉมาย จ. สสิตฺถกํ นิวตฺตนฺเต,}\csromangathastanza{สุปฺปฏิจฺฉนฺนมุกฺกุฏิ. ธมฺมราเชน ปญฺญตฺตํ, \\ อิทํ ภตฺตคฺควตฺตนํ. ทุนฺนิวตฺถา อนากปฺปา,}\csromangathastanza{\csromanfolio{413}อสลฺลกฺเขตฺวา จ สหสา. ทูเร อจฺจ จิรํ ลหุํ, \\ ตเถว ปิณฺฑจาริโก. ปฏิจฺฉนฺโนว คจฺเฉยฺย,}\csromangathastanza{สุสํวุโตกฺขิตฺตจกฺขุ. อุกฺขิตฺโตชฺชคฺฆิกาสทฺโท, \\ ตโย เจว ปจาลนา. ขมฺโภคุณฺฐิ-อุกฺกุฏิกา,}\csromangathastanza{สลฺลกฺเขตฺวา จ สหสา. ทูเร อจฺจ จิรํ ลหุํ, \\ อาสนกํ กฏจฺฉุกา. ภาชนํ วา ฐเปติ จ,}\csromangathastanza{อุจฺจาเรตฺวา ปณาเมตฺวา. ปฏิคฺคเห น อุลฺโลเก, \\ สูเปสุปิ ตเถว ตํ. ภิกฺขุ สํฆาฏิยา ฉาเท,}\csromangathastanza{ปฏิจฺฉนฺเนว คจฺฉิยํ. สํวุโตกฺขิตฺตจกฺขุ จ, \\ อุกฺขิตฺโตชฺชคฺฆิกาย จ. อปฺปสทฺโท ตโย จาลา,}\csromangathastanza{ขมฺโภคุณฺฐิก-อุกฺกุฏิ. ปฐมาสนวกฺการ, \\ ปานิยํ ปริโภชนี. ปจฺฉากงฺขติ ภุญฺเชยฺย,}\csromangathastanza{โอปิลาเปยฺย อุทฺธเร. ปฏิสาเมยฺย สมฺมชฺเช, \\ ริตฺตํ ตุจฺฉํ อุปฏฺฐเป. หตฺถวิกาเร ภินฺเทยฺย,}\csromangathastanza{วตฺติทํ ปิณฺฑจาริเก. ปานี ปริ อคฺคิรณิ, \\ นกฺขตฺตทิสโจรา จ. “สพฺพํ นตฺถี”ติ โกฏฺเฏตฺวา,}\csromangathastanza{ปตฺตํเส จีวรํ ตโต. อิทานิ อํเส ลคฺเคตฺวา, \\ ติมณฺฑลํ ปริมณฺฑลํ. ยถา ปิณฺฑจาริวตฺตํ,}\csromangathastanza{นเย อารญฺญเกสุปิ. ปตฺตํเส จีวรํ สีเส, \\ อาโรหิตฺวา จ ปานิยํ. ปริโภชนิยํ อคฺคิ,}\csromangathastanza{อรณี จาปิ กตฺตรํ. นกฺขตฺตํ สปฺปเทสํ วา, \\ ทิสาปิ กุสโล ภเว. สตฺตุตฺตเมน ปญฺญตฺตํ,}\csromangathastanza{วตฺตํ อรญฺญเกสุเม. อชฺโฌกาเส โอกิริํสุ, \\ อุชฺฌายนฺติ จ เปสลา. สเจ วิหาโร อุกฺลาโป,}\csromangathastanza{ปฐมํ ปตฺตจีวรํ. ภิสิพิพฺโพหนํ มญฺจํ, \\ ปีฐญฺจ เขฬมลฺลกํ. อปสฺเสนุลฺโลกกณฺณา,}\csromangathastanza{เครุกา กาฬ อกตา. สงฺการํ ภิกฺขุสามนฺตา, \\ เสนาวิหารปานิยํ. ปริโภชนสามนฺตา,}\csromangathastanza{\csromanfolio{414}ปฏิวาเต จ องฺคเณ. อโธวาเต อตฺถรณํ, \\ ปฏิปาทกมญฺโจ จ. ปีฐํ ภิสิ นิสีทนํ,}\csromangathastanza{มลฺลกํ อปสฺเสน จ. ปตฺตจีวรํ ภูมิ จ, \\ ปารนฺตํ โอรโต โภคํ. ปุรตฺถิมา จ ปจฺฉิมา,}\csromangathastanza{อุตฺตรา อถ ทกฺขิณา. สีตุเณฺห จ ทิวา รตฺติํ, \\ ปริเวณญฺจ โกฏฺฐโก. อุปฏฺฐานคฺคิสาลา จ,}\csromangathastanza{วจฺจกุฏี จ ปานิยํ. อาจมกุมฺภิ วุฑฺเฒ จ, \\ อุทฺเทสปุจฺฉนา สชฺฌา. ธมฺโม ปทีปํ วิชฺฌาเป,}\csromangathastanza{น วิวเร นปิ ถเก. เยน วุฑฺโฒ ปริวตฺติ, \\ กณฺเณนปิ น ฆฏฺเฏเย. ปญฺญเปสิ มหาวีโร,}\csromangathastanza{วตฺตํ เสนาสเนสุ ตํ. นิวาริยมานา ทฺวารํ, \\ มุจฺฉิตุชฺฌนฺติ เปสลา. ฉาริกํ ฉฑฺฑเย ชนฺตา,}\csromangathastanza{ปริภณฺฑํ ตเถว จ. ปริเวณํ โกฏฺฐโก สาลา, \\ จุณฺณมตฺติกโทณิกา. มุขํ ปุรโต น เถเร,}\csromangathastanza{น นเว อุสฺสหติ สเจ. ปุรโต อุปริมคฺโค, \\ จิกฺขลฺลํ มตฺติ ปีฐกํ. วิชฺฌาเปตฺวา ถเกตฺวา จ,}\csromangathastanza{วตฺตํ ชนฺตาฆเรสุเม. นาจเมติ ยถาวุฑฺฒุํ, \\ ปฏิปาฏิ จ สหสา. อุพฺภชิ นิตฺถุโน กฏฺฐํ,}\csromangathastanza{วจฺจํ ปสฺสาว เขฬกํ. ผรุสา กูป สหสา, \\ อุพฺภชิ จปุ เสเสน. พหิ อนฺโต จ อุกฺกาเส,}\csromangathastanza{รชฺชุ อตรมานญฺจ. สหสา อุพฺภชิ ฐิเต, \\ นิตฺถุเน กฏฺฐ วจฺจญฺจ. ปสฺสาว เขฬ ผรุสา,}\csromangathastanza{กูปญฺจ วจฺจปาทุเก. นติสหสา อุพฺภชิ, \\ ปาทุกาย จปุจปุ. น เสสเย ปฏิจฺฉาเท,}\csromangathastanza{อุหตปิธเรน จ. วจฺจกุฏี ปริภณฺฑํ, \\ ปริเวณญฺจ โกฏฺฐโก. อาจมเน จ อุทกํ,}\csromangathastanza{\csromanfolio{415}วตฺตํ วจฺจกุฏีสุเม. อุปาหนา ทนฺตกฏฺฐํ, \\ มุโขทกญฺจ อาสนํ. ยาคุ อุทกํ โธวิตฺวา,}\csromangathastanza{อุทฺธารุกฺลาป คาม จ. นิวาสนา กายพนฺธา, \\ สคุณํ ปตฺตโสทกํ. ปจฺฉา ติมณฺฑโล เจว,}\csromangathastanza{ปริมณฺฑล พนฺธนํ. สคุณํ โธวิตฺวา ปจฺฉา, \\ นาติทูเร ปฏิคฺคเห. ภณมานสฺส อาปตฺติ,}\csromangathastanza{ปฐมาคนฺตฺวาน อาสนํ. อุทกํ ปีฐกถลิ, \\ ปจฺจุคฺคนฺตฺวา นิวาสนํ. โอตาเป นิทหิ ภงฺโค,}\csromangathastanza{โอโภเค ภุญฺชิตุ นเม. ปานียํ อุทกํ นีจํ, \\ มุหุตฺตํ น จ นิทเห. ปตฺตจีวรํ ภูมิ จ,}\csromangathastanza{ปารนฺตํ โอรโต โภคํ. อุทฺธเร ปฏิสาเม จ, \\ อุกฺลาโป จ นหายิตุํ. สีตํ อุณฺหํ ชนฺตาฆรํ,}\csromangathastanza{จุณฺณํ มตฺติก ปิฏฺฐิโต. ปีฐญฺจ จีวรํ จุณฺณํ, \\ มตฺติกุสฺสหติ มุขํ. ปุรโต เถเร นเว จ,}\csromangathastanza{ปริกมฺมญฺจ นิกฺขเม. ปุรโต อุทเก นฺหาเต, \\ นิวาเสตฺวา อุปชฺฌายํ. นิวาสนญฺจ สํฆาฏิ,}\csromangathastanza{ปีฐกํ อาสเนน จ. ปาโท ปีฐํ กถลิญฺจ, \\ ปานียุทฺเทสปุจฺฉนา. อุกฺลาปํ สุโสเธยฺย,}\csromangathastanza{ปฐมํ ปตฺตจีวรํ. นิสีทนปจฺจตฺถรณํ, \\ ภิสิ พิพฺโพหนานิ จ. มญฺโจ ปีฐํ ปฏิปาทํ,}\csromangathastanza{มลฺลกํ อปสฺเสน จ. ภูม สนฺตาน อาโลก, \\ เครุกา กาฬ อกตา. ภูมตฺถรปฏิปาทา,}\csromangathastanza{มญฺโจ ปีฐํ พิพฺโพหนํ. นิสีทถรณํ เขฬ, \\ อปสฺเส ปตฺตจีวรํ. ปุรตฺถิมา ปจฺฉิมา จ,}\csromangathastanza{อุตฺตรา อถ ทกฺขิณา. สีตุณฺหญฺจ ทิวา รตฺติํ, \\ ปริเวณญฺจ โกฏฺฐโก. อุปฏฺฐานคฺคิสาลา จ,}\csromangathastanza{\csromanfolio{416}วจฺจปานิยโภชนี. อาจมํ อนภิรติ, \\ กุกฺกุจฺจํ ทิฏฺฐิ จ ครุ. มูลมานตฺต-อพฺภานํ,}\csromangathastanza{ตชฺชนียํ นิยสฺสกํ. ปพฺพาช ปฏิสารณี, \\ อุกฺเขปญฺจ กตํ ยทิ. โธเว กาตพฺพํ รชญฺจ,}\csromangathastanza{รเช สมฺปริวตฺตกํ. ปตฺตญฺจ จีวรญฺจาปิ, \\ ปริกฺขารญฺจ เฉทนํ. ปริกมฺมํ เวยฺยาวจฺจํ,}\csromangathastanza{ปจฺฉา ปิณฺฑํ ปวิสนํ. น สุสานํ ทิสา เจว, \\ ยาวชีวํ อุปฏฺฐเห. สทฺธิวิหาริเกเนตํ,}\csromangathastanza{วตฺตุปชฺฌายเกสุเม. โอวาทสาสนุทฺเทสา, \\ ปุจฺฉา ปตฺตญฺจ จีวรํ. ปริกฺขาโร คิลาโน จ,}\csromangathastanza{น ปจฺฉาสมโณ ภเว. อุปชฺฌาเยสุ เย วตฺตา, \\ เอวํ อาจริเยสุปิ. สทฺธิวิหาริเก วตฺตา,}\csromangathastanza{ตเถว อนฺเตวาสิเก. อาคนฺตุเกสุ เย วตฺตา, \\ ปุน อาวาสิเกสุ จ. คมิกานุโมทนิกา,}\csromangathastanza{ภตฺตคฺเค ปิณฺฑจาริเก. อารญฺญเกสุ ยํ วตฺตํ, \\ ยญฺจ เสนาสเนสุปิ. ชนฺตาฆเร วจฺจกุฏี,}\csromangathastanza{อุปชฺฌา สทฺธิวิหาริเก. อาจริเยสุ ยํ วตฺตํ, \\ ตเถว อนฺเตวาสิเก. เอกูนวีสติ วตฺถู,}\csromangathastanza{วตฺตา จุทฺทส ขนฺธเก. วตฺตํ อปริปูเรนฺโต, \\ น สีลํ ปริปูรติ. อสุทฺธสีโล ทุปฺปญฺโญ,}\csromangathastanza{จิตฺเตกคฺคํ น วินฺทติ. วิกฺขิตฺตจิตฺโตเนกคฺโค, \\ สมฺมา ธมฺมํ น ปสฺสติ. อปสฺสมาโน สทฺธมฺมํ,}\csromangathastanza{ทุกฺขา น ปริมุจฺจติ. ยํ วตฺตํ ปริปูเรนฺโต, \\ สีลมฺปิ ปริปูรติ. วิสุทฺธสีโล สปฺปญฺโญ,}\csromangathastanza{\csromanfolio{417}จิตฺเตกคฺคมฺปิ วินฺทติ. อวิกฺขิตฺตจิตฺโต เอกคฺโค, \\ สมฺมา ธมฺมํ วิปสฺสติ. สมฺปสฺสมาโน สทฺธมฺมํ,}\csromangathastanza{ทุกฺขา โส ปริมุจฺจติ. ตสฺมา หิ วตฺตํ ปูเรยฺย, \\ ชินปุตฺโต วิจฺจกฺขโณ. โอวาทํ พุทฺธเสฏฺฐสฺส, \\ ตโต นิพฺพานเมหิตีติ.}}

\nitthitam{\textbf{วตฺตกฺขนฺธโก นิฏฺฐิโต.}}
\csromanmatikaanchor{mtk.192}
\csromantocmarknopage{chapter}{9. ปาติโมกฺขฏฺฐปนกฺขนฺธก}{mtk.192}
\bookmark[dest={mtk.192},level=0]{9. ปาติโมกฺขฏฺฐปนกฺขนฺธก}
\chapterheadpageread{9. ปาติโมกฺขฏฺฐปนกฺขนฺธก}
\csromanmatikaanchor{mtk.193}
\csromantocmarkref{section}{1. ปาติโมกฺขุทฺเทสยาจน}{mtk.193}
\bookmark[dest={mtk.193},level=1]{1. ปาติโมกฺขุทฺเทสยาจน}
\csromanheader{1}{1. ปาติโมกฺขุทฺเทสยาจน}
\proseitem{383}{\csromanfolio{418}\csromansymbolfootnote{*}{ขุ 1. 138; อํ 3. 44 ปิฏฺเฐสุปิ.}เตน สมเยน พุทฺโธ ภควา สาวตฺถิยํ วิหรติ ปุพฺพาราเม มิคารมาตุ ปาสาเท.\csromanspacer{} เตน โข ปน สมเยน ภควา ตทหุโปสเถ ภิกฺขุสํฆปริวุโต นิสินฺโน โหติ.\csromanspacer{} อถ โข อายสฺมา อานนฺโท อภิกฺกนฺตาย รตฺติยา นิกฺขนฺเต ปฐเม ยาเม อุฏฺฐายาสนา เอกํสํ อุตฺตราสงฺคํ กริตฺวา เยน ภควา เตนญฺชลิํ ปณาเมตฺวา ภควนฺตํ เอตทโวจ “อภิกฺกนฺตา ภนฺเต รตฺติ, นิกฺขนฺโต ปฐโม ยาโม, จิรนิสินฺโน ภิกฺขุสํโฆ, อุทฺทิสตุ ภนฺเต ภควา ภิกฺขูนํ ปาติโมกฺขนฺ”ติ.\csromanspacer{} เอวํ วุตฺเต ภควา ตุณฺหี อโหสิ.\csromanspacer{} ทุติยมฺปิ โข อายสฺมา อานนฺโท อภิกฺกนฺตาย รตฺติยา นิกฺขนฺเต มชฺฌิเม ยาเม อุฏฺฐายาสนา เอกํสํ อุตฺตราสงฺคํ กริตฺวา เยน ภควา เตนญฺชลิํ ปณาเมตฺวา ภควนฺตํ เอตทโวจ “อภิกฺกนฺตา ภนฺเต รตฺติ, นิกฺขนฺโต มชฺฌิโม ยาโม, จิรนิสินฺโน ภิกฺขุสํโฆ, อุทฺทิสตุ ภนฺเต ภควา ภิกฺขูนํ ปาติโมกฺขนฺ”ติ.\csromanspacer{} ทุติยมฺปิ โข ภควา ตุณฺหี อโหสิ.\csromanspacer{} ตติยมฺปิ โข อายสฺมา อานนฺโท อภิกฺกนฺตาย รตฺติยา นิกฺขนฺเต ปจฺฉิเม ยาเม อุทฺธสฺเต อรุเณ นนฺทิมุขิยา รตฺติยา อุฏฺฐายาสนา เอกํสํ อุตฺตราสงฺคํ กริตฺวา เยน ภควา เตนญฺชลิํ ปณาเมตฺวา ภควนฺตํ เอตทโวจ “อภิกฺกนฺตา ภนฺเต รตฺติ, นิกฺขนฺโต ปจฺฉิโม ยาโม, อุทฺธสฺตํ อรุณํ, นนฺทิมุขิ รตฺติ, จิรนิสินฺโน ภิกฺขุสํโฆ, อุทฺทิสตุ ภนฺเต ภควา ภิกฺขูนํ ปาติโมกฺขนฺ”ติ.\csromanspacer{} อปริสุทฺธา อานนฺท ปริสาติ.}

\prose{อถ โข อายสฺมโต มหาโมคฺคลฺลานสฺส เอตทโหสิ “กํ นุ โข ภควา ปุคฺคลํ สนฺธาย เอวมาห ‘อปริสุทฺธา อานนฺท ปริสา’ติ”.\csromanspacer{} อถ โข อายสฺมา มหาโมคฺคลฺลาโน สพฺพาวนฺตํ ภิกฺขุสํฆํ เจตสา เจโต ปริจฺจ มนสากาสิ.\csromanspacer{} อทฺทสา โข อายสฺมา มหาโมคฺคลฺลาโน ตํ ปุคฺคลํ ทุสฺสีลํ ปาปธมฺมํ อสุจิสงฺกสฺสรสมาจารํ ปฏิจฺฉนฺนกมฺมนฺตํ อสฺสมณํ สมณปฏิญฺญํ อพฺรหฺมจาริํ พฺรหฺมจาริปฏิญฺญํ อนฺโตปูติํ อวสฺสุตํ กสมฺพุชาตํ\footnote{กสมฺพุกชาตํ (ก)} มชฺเฌ ภิกฺขุสํฆสฺส นิสินฺนํ, ทิสฺวาน เยน โส \csromanfolio{419}ปุคฺคโล เตนุปสงฺกมิ, อุปสงฺกมิตฺวา ตํ ปุคฺคลํ เอตทโวจ “อุฏฺเฐหิ อาวุโส, ทิฏฺโฐสิ ภควตา, นตฺถิ เต ภิกฺขูหิ สทฺธิํ สํวาโส”ติ.\csromanspacer{} เอวํ วุตฺเต โส ปุคฺคโล ตุณฺหี อโหสิ.\csromanspacer{} ทุติยมฺปิ โข อายสฺมา มหาโมคฺคลฺลาโน ตํ ปุคฺคลํ เอตทโวจ “อุฏฺเฐหิ อาวุโส, ทิฏฺโฐสิ ภควตา, นตฺถิ เต ภิกฺขูหิ สทฺธิํ สํวาโส”ติ.\csromanspacer{} ทุติยมฺปิ โข โส ปุคฺคโล ตุณฺหี อโหสิ.\csromanspacer{} ตติยมฺปิ โข อายสฺมา มหาโมคฺคลฺลาโน ตํ ปุคฺคลํ เอตทโวจ “อุฏฺเฐหิ อาวุโส, ทิฏฺโฐสิ ภควตา, นตฺถิ เต ภิกฺขูหิ สทฺธิํ สํวาโส”ติ.\csromanspacer{} ตติยมฺปิ โข โส ปุคฺคโล ตุณฺหี อโหสิ.\csromanspacer{} อถ โข อายสฺมา มหาโมคฺคลฺลาโน ตํ ปุคฺคลํ พาหายํ คเหตฺวา พหิทฺวารโกฏฺฐกา นิกฺขาเมตฺวา สูจิฆฏิกํ ทตฺวา เยน ภควา เตนุปสงฺกมิ, อุปสงฺกมิตฺวา ภควนฺตํ เอตทโวจ “นิกฺขามิโต โส ภนฺเต ปุคฺคโล มยา, สุทฺธา ปริสา อุทฺทิสตุ ภนฺเต ภควา ภิกฺขูนํ ปาติโมกฺขนฺ”ติ.}

\prose{อจฺฉริยํ โมคฺคลฺลาน, อพฺภุตํ โมคฺคลฺลาน, ยาว พาหาคหณาปิ นาม โส โมฆปุริโส อาคเมสฺสตีติ.\csromanspacer{} อถ โข ภควา ภิกฺขู อามนฺเตสิ\csromandash{}}

\csromanmatikaanchor{mtk.194}
\csromantocmarkref{section}{2. มหาสมุทฺเท อฏฺฐจฺฉริย}{mtk.194}
\bookmark[dest={mtk.194},level=1]{2. มหาสมุทฺเท อฏฺฐจฺฉริย}
\csromanheader{1}{2. มหาสมุทฺเท-อฏฺฐจฺฉริย}
\proseitem{384}{\csromansymbolfootnote{*}{ขุ 1. 139; อํ 3. 39 ปิฏฺเฐสุปิ.}อฏฺฐิเม ภิกฺขเว มหาสมุทฺเท อจฺฉริยา อพฺภุตา ธมฺมา, เย ทิสฺวา ทิสฺวา อสุรา มหาสมุทฺเท อภิรมนฺติ.\csromanspacer{} กตเม อฏฺฐ.}

\prose{มหาสมุทฺโท ภิกฺขเว อนุปุพฺพนินฺโน อนุปุพฺพโปโณ อนุปุพฺพปพฺภาโร น อายตเกเนว ปปาโต.\csromanspacer{} ยมฺปิ ภิกฺขเว มหาสมุทฺโท อนุปุพฺพนินฺโน อนุปุพฺพโปโณ อนุปุพฺพปพฺภาโร น อายตเกเนว ปปาโต, อยํ ภิกฺขเว มหาสมุทฺเท ปฐโม อจฺฉริโย อพฺภุโต ธมฺโม, ยํ ทิสฺวา ทิสฺวา อสุรา มหาสมุทฺเท อภิรมนฺติ.}

\prose{ปุน จปรํ ภิกฺขเว มหาสมุทฺโท ฐิตธมฺโม เวลํ นาติวตฺตติ, ยมฺปิ ภิกฺขเว มหาสมุทฺโท ฐิตธมฺโม เวลํ นาติวตฺตติ.\csromanspacer{} อยํ\footnote{อยมฺปิ (สฺยา)} ภิกฺขเว มหาสมุทฺเท ทุติโย อจฺฉริโย อพฺภุโต ธมฺโม, ยํ ทิสฺวา ทิสฺวา อสุรา มหาสมุทฺเท อภิรมนฺติ.}

\prose{\csromanfolio{420}ปุน จปรํ ภิกฺขเว มหาสมุทฺโท น มเตน กุณเปน สํวสติ.\csromanspacer{} ยํ โหติ มหาสมุทฺเท มตํ กุณปํ, ตํ ขิปฺปญฺเญว ตีรํ วาเหติ, ถลํ อุสฺสาเรติ.\csromanspacer{} ยมฺปิ ภิกฺขเว มหาสมุทฺโท น มเตน กุณเปน สํวสติ.\csromanspacer{} ยํ โหติ มหาสมุทฺเท มตํ กุณปํ, ตํ ขิปฺปญฺเญว ตีรํ วาเหติ, ถลํ อุสฺสาเรติ.\csromanspacer{} อยํ ภิกฺขเว มหาสมุทฺเท ตติโย อจฺฉริโย อพฺภุโต ธมฺโม, ยํ ทิสฺวา ทิสฺวา อสุรา มหาสมุทฺเท อภิรมนฺติ.}

\prose{ปุน จปรํ ภิกฺขเว ยา กาจิ มหานทิโย, เสยฺยถิทํ, คงฺคา ยมุนา อจิรวตี สรภู มหี, ตา มหาสมุทฺทํ ปตฺตา ชหนฺติ ปุริมานิ นามโคตฺตานิ, “มหาสมุทฺโท”เตฺวว สงฺขํ คจฺฉนฺติ.\csromanspacer{} ยมฺปิ ภิกฺขเว ยากาจิ มหานทิโย, เสยฺยถิทํ, คงฺคา ยมุนา อจิรวตี สรภู มหี, ตา มหาสมุทฺทํ ปตฺตา ชหนฺติ ปุริมานิ นามโคตฺตานิ, “มหาสมุทฺโท”เตฺวว สงฺขํ คจฺฉนฺติ, อยํ ภิกฺขเว มหาสมุทฺเท จตุตฺโถ อจฺฉริโย อพฺภุโต ธมฺโม, ยํ ทิสฺวา ทิสฺวา อสุรา มหาสมุทฺเท อภิรมนฺติ.}

\prose{ปุน จปรํ ภิกฺขเว ยา จ\footnote{ยากาจิ (สฺยา)} โลเก สวนฺติโย มหาสมุทฺทํ อปฺเปนฺติ, ยา จ อนฺตลิกฺขา ธารา ปปตนฺติ, น เตน มหาสมุทฺทสฺส อูนตฺตํ วา ปูรตฺตํ วา ปญฺญยติ.\csromanspacer{} ยมฺปิ ภิกฺขเว ยา จ โลเก สวนฺติโย มหาสมุทฺทํ อปฺเปนฺติ, ยา จ อนฺตลิกฺขา ธารา ปปตนฺติ, น เตน มหาสมุทฺทสฺส อูนตฺตํ วา ปูรตฺตํ วา ปญฺญายติ, อยํ ภิกฺขเว มหาสมุทฺเท ปญฺจโม อจฺฉริโย อพฺภุโต ธมฺโม, ยํ ทิสฺวา ทิสฺวา อสุรา มหาสมุทฺเท อภิรมนฺติ.}

\prose{ปุน จปรํ ภิกฺขเว มหาสมุทฺโท เอกรโส โลณรโส.\csromanspacer{} ยมฺปิ ภิกฺขเว มหาสมุทฺโท เอกรโส โลณรโส, อยํ ภิกฺขเว มหาสมุทฺเท ฉฏฺโฐ อจฺฉริโย อพฺภุโต ธมฺโม, ยํ ทิสฺวา ทิสฺวา อสุรา มหาสมุทฺเท อภิรมนฺติ.}

\prose{ปุน จปรํ ภิกฺขเว มหาสมุทฺโท พหุรตโน\footnote{ปหูตรตโน (ก)} อเนกรตโน, ตตฺริมานิ รตนานิ, เสยฺยถิทํ, มุตฺตา มณิ เวฬุริโย สงฺโข สิลา ปวาฬํ รชตํ ชาตรูปํ โลหิตโก มสารคลฺลํ.\csromanspacer{} ยมฺปิ ภิกฺขเว มหาสมุทฺโท พหุรตโน อเนกรตโน, ตตฺริมานิ รตนานิ เสยฺยถิทํ, มุตฺตา มณิ เวฬุริโย สงฺโข สิลา ปวาฬํ รชตํ ชตรูปํ \csromanfolio{421}โลหิตโก มสารคลฺลํ, อยํ ภิกฺขเว มหาสมุทฺเท สตฺตโม อจฺฉริโย อพฺภุโต ธมฺโม, ยํ ทิสฺวา ทิสฺวา อสุรา มหาสมุทฺเท อภิรมนฺติ.}

\prose{ปุน จปรํ ภิกฺขเว มหาสมุทฺโท มหตํ ภูตานํ อาวาโส, ตตฺริเม ภูตา, ติมิ ติมิงฺคโล ติมิติมิงฺคโล\footnote{ติมิรปิงฺคโล (สี), ติมิติมิงฺคโล มหาติมิงฺคโล (สฺยา, กํ)} อสุรา นาคา คนฺธพฺพา, สนฺติ มหาสมุทฺเท โยชนสติกาปิ อตฺตภาวา, ทฺวิโยชนสติกาปิ อตฺตภาวา, ติโยชนสติกาปิ อตฺตภาวา, จตุโยชนสติกาปิ อตฺตภาวา, ปญฺจโยชนสติกาปิ อตฺตภาวา.\csromanspacer{} ยมฺปิ ภิกฺขเว มหาสมุทฺโท มหตํ ภูตานํ อาวาโส, ตตฺริเม ภูตา, ติมิ ติมิงฺคโล ติมิติมิงฺคโล อสุรา นาคา คนฺธพฺพา, สนฺติ มหาสมุทฺเท โยชนสติกาปิ อตฺตภาวา ฯเปฯ ปญฺจโยชนสติกาปิ อตฺตภาวา, อยํ ภิกฺขเว มหาสมุทฺเท อฏฺฐโม อจฺฉริโย อพฺภุโต ธมฺโม, ยํ ทิสฺวา ทิสฺวา อสุรา มหาสมุทฺเท อภิรมนฺติ.\csromanspacer{} อิเม โข ภิกฺขเว มหาสมุทฺเท อฏฺฐ อจฺฉริยา อพฺภุตา ธมฺมา, เย ทิสฺวา ทิสฺวา อสุรา มหาสมุทฺเท อภิรมนฺติ.}

\csromanmatikaanchor{mtk.195}
\csromantocmarkref{section}{3. อิมสฺมิํ ธมฺมวินเย อฏฺฐจฺฉริย}{mtk.195}
\bookmark[dest={mtk.195},level=1]{3. อิมสฺมิํ ธมฺมวินเย อฏฺฐจฺฉริย}
\csromanheader{1}{3. อิมสฺมิํ ธมฺมวินเย อฏฺฐจฺฉริย}
\proseitem{385}{เอวเมว โข ภิกฺขเว อิมสฺมิํ ธมฺมวินเย อฏฺฐ อจฺฉริยา อพฺภุตา ธมฺมา, เย ทิสฺวา ทิสฺวา ภิกฺขู อิมสฺมิํ ธมฺมวินเย อภิรมนฺติ.\csromanspacer{} กตเม อฏฺฐ.}

\prose{เสยฺยถาปิ ภิกฺขเว มหาสมุทฺโท อนุปุพฺพนินฺโน อนุปุพฺพโปโณ อนุปุพฺพปพฺภาโร น อายตเกเนว ปปาโต, เอวเมว โข อภิกฺขเว อิมสฺมิํ ธมฺมวินเย อนุปุพฺพสิกฺขา อนุปุพฺพกิริยา อนุปุพฺพปฏิปทา น อายตเกเนว อญฺญาปฏิเวโธ.\csromanspacer{} ยมฺปิ ภิกฺขเว อิมสฺมิํ ธมฺมวินเย อนุปุพฺพสิกฺขา อนุปุพฺพกิริยา อนุปุพฺพปฏิปทา น อายตเกเนว อญฺญาปฏิเวโธ, อยํ ภิกฺขเว อิมสฺมิํ ธมฺมวินเย ปฐโม อจฺฉริโย อพฺภุโต ธมฺโม, ยํ ทิสฺวา ทิสฺวา ภิกฺขู อิมสฺมิํ ธมฺมวินเย อภิรมนฺติ.}

\prose{เสยฺยถาปิ ภิกฺขเว มหาสมุทฺโท ฐิตธมฺโม เวลํ นาติวตฺตติ เอวเมว โข ภิกฺขเว ยํ มยา มม สาวกานํ สิกฺขาปทํ ปญฺญตฺตํ, ตํ มม สาวกา ชีวิตเหตุปิ นาติกฺกมนฺติ.\csromanspacer{} ยมฺปิ ภิกฺขเว มยา มม สาวกานํ \csromanfolio{422}สิกฺขาปทํ ปญฺญตฺตํ, ตํ มม สาวกา ชีวิตเหตุปิ นาติกฺกมนฺติ, อยํ ภิกฺขเว อิมสฺมิํ ธมฺมวินเย ทุติโย อจฺฉริโย อพฺภุโต ธมฺโม, ยํ ทิสฺวา ทิสฺวา ภิกฺขู อิมสฺมิํ ธมฺมวินเย อภิรมนฺติ.}

\prose{เสยฺยถาปิ ภิกฺขเว มหาสมุทฺโท น มเตน กุณเปน สํวสติ, ยํ โหติ มหาสมุทฺเท มตํ กุณปํ, ตํ ขิปฺปเมว ตีรํ วาเหติ, ถลํ อุสฺสาเรติ, เอวเมว โข ภิกฺขเว โย โส ปุคฺคโล ทุสฺสีโล ปาปธมฺโม อสุจิสงฺกสฺสรสมาจาโร ปฏิจฺฉนฺนกมฺมนฺโต อสฺสมโณ สมณปฏิญฺโญ อพฺรหฺมจารี พฺรหฺมจาริปฏิญฺโญ อนฺโตปูติ อวสฺสุโต กสมฺพุชาโต, น เตน สํโฆ สํวสติ, ขิปฺปเมว นํ สนฺนิปติตฺวา อุกฺขิปติ, กิญฺจาปิ โข โส โหติ มชฺเฌ ภิกฺขุสํฆสฺส นิสินฺโน, อถ โข โส อารกาว สํฆมฺหา สํโฆ จ เตน.\csromanspacer{} ยมฺปิ ภิกฺขเว โย โส ปุคฺคโล ทุสฺสีโล ปาปธมฺโม อสุจิสงฺกสฺสรสมาจาโร ปฏิจฺฉนฺนกมฺมนฺโต อสฺสมโณ สมณปฏิญฺโญ อพฺรหฺมจารี พฺรหฺมจาริปฏิญฺโญ อนฺโตปูติ อวสฺสุโต กสมฺพุชาโต, น เตน สํโฆ สํวสติ, ขิปฺปเมว นํ สนฺนิปติตฺวา อุกฺขิปติ, กิญฺจาปิ โข โส โหติ มชฺเฌ ภิกฺขุสํฆสฺส นิสินฺโน, อถ โข โส อารกาว สํฆมฺหา สํโฆ จ เตน, อยํ ภิกฺขเว อิมสฺมิํ ธมฺมวินเย ตติโย อจฺฉริโย อพฺภุโต ธมฺโม, ยํ ทิสฺวา ทิสฺวา ภิกฺขู อิมสฺมิํ ธมฺมวินเย อภิรมนฺติ.}

\prose{เสยฺยถาปิ ภิกฺขเว ยา กาจิ มหานทิโย, เสยฺยถิทํ, คงฺคา ยมุนา อจิรวตี สรภู มหี, ตา มหาสมุทฺทํ ปตฺตา ชหนฺติ ปุริมานิ นามโคตฺตานิ, “มหาสมุทฺโท”เตฺวว สงฺขํ คจฺฉนฺติ, เอวเมว โข ภิกฺขเว จตฺตาโรเม วณฺณา ขตฺติยา พฺราหฺมณา เวสฺสา สุทฺทา, เต ตถาคตปฺปเวทิเต ธมฺมวินเย อคารสฺมา อนคาริยํ ปพฺพชิตฺวา\footnote{ปพฺพชิตา (สี)} ชหนฺติ ปุริมานิ นามโคตฺตานิ, “สมณา สกฺยปุตฺติยา”เตฺวว สงฺขํ คจฺฉนฺติ.\csromanspacer{} ยมฺปิ ภิกฺขเว จตฺตาโรเม วณฺณา ขตฺติยา พฺราหฺมณา เวสฺสา สุทฺทา, เต ตถาคตปฺปเวทิเต ธมฺมวินเย อคารสฺมา อนคาริยํ ปพฺพชิตฺวา ชหนฺติ ปุริมานิ นามโคตฺตานิ, “สมณา สกฺยปุตฺติยา”เตฺวว สงฺขํ คจฺฉนฺติ, อยํ ภิกฺขเว อิมสฺมิํ ธมฺมวินเย จตุตฺโถ อจฺฉริโย อพฺภุโต ธมฺโม, ยํ ทิสฺวา ทิสฺวา ภิกฺขู อิมสฺมิํ ธมฺมวินเย อภิรมนฺติ.}

\prose{\csromanfolio{423}เสยฺยถาปิ ภิกฺขเว ยา จ โลเก สวนฺติโย มหาสมุทฺทํ อปฺเปนฺติ, ยา จ อนฺตลิกฺขา ธารา ปปตนฺติ, น เตน มหาสมุทฺทสฺส อูนตฺตํ วา ปูรตฺตํ วา ปญฺญายติ, เอวเมว โข ภิกฺขเว พหู เจปิ ภิกฺขู อนุปาทิเสสาย นิพฺพานธาตุยา ปรินิพฺพายนฺติ, น เตน นิพฺพานธาตุยา อูนตฺตํ วา ปูรตฺตํ วา ปญฺญายติ.\csromanspacer{} ยมฺปิ ภิกฺขเว พหู เจปิ ภิกฺขู อนุปาทิเสสาย นิพฺพานธาตุยา ปรินิพฺพายนฺติ, น เตน นิพฺพานธาตุยา อูนตฺตํ วา ปูรตฺตํ วา ปญฺญายติ, อยํ ภิกฺขเว อิมสฺมิํ ธมฺมวินเย ปญฺจโม อจฺฉริโย อพฺภุโต ธมฺโม, ยํ ทิสฺวา ทิสฺวา ภิกฺขู อิมสฺมิํ ธมฺมวินเย อภิรมนฺติ.}

\prose{เสยฺยถาปิ ภิกฺขเว มหาสมุทฺโท เอกรโส โลณรโส, เอวเมว โข ภิกฺขเว อยํ ธมฺมวินโย เอกรโส วิมุตฺติรโส.\csromanspacer{} ยมฺปิ ภิกฺขเว อยํ ธมฺมวินโย เอกรโส วิมุตฺติรโส, อยํ ภิกฺขเว อิมสฺมิํ ธมฺมวินเย ฉฏฺโฐ อจฺฉริโย อพฺภุโต ธมฺโม, ยํ ทิสฺวา ทิสฺวา ภิกฺขู อิมสฺมิํ ธมฺมวินเย อภิรมนฺติ.}

\prose{เสยฺยถาปิ ภิกฺขเว มหาสมุทฺโท พหุรตโน อเนกรตโน, ตตฺริมานิ รตนานิ, เสยฺยถิทํ, มุตฺตา มณิ เวฬุริโย สงฺโข สิลา ปวาฬํ รชตํ ชาตรูปํ โลหิตโก มสารคลฺลํ, เอวเมว โข ภิกฺขเว อยํ ธมฺมวินโย พหุรตโน อเนกรตโน, ตตฺริมานิ รตนานิ, เสยฺยถิทํ, จตฺตาโร สติปฏฺฐานา จตฺตาโร สมฺมปฺปธานา จตฺตาโร อิทฺธิปาทา ปญฺจินฺทฺริยานิ ปญฺจ พลานิ สตฺต โพชฺฌงฺคา อริโย อฏฺฐงฺคิโก มคฺโค.\csromanspacer{} ยมฺปิ ภิกฺขเว อยํ ธมฺมวินโย พหุรตโน อเนกรตโน, ตตฺริมานิ รตนานิ, เสยฺยถิทํ, จตฺตาโร สติปฏฺฐานา ฯเปฯ อริโย อฏฺฐงฺคิโก มคฺโค, อยํ ภิกฺขเว อิมสฺมิํ ธมฺมวินเย สตฺตโม อจฺฉริโย อพฺภุโต ธมฺโม, ยํ ทิสฺวา ทิสฺวา ภิกฺขู อิมสฺมิํ ธมฺมวินเย อภิรมนฺติ.}

\prose{เสยฺยถาปิ ภิกฺขเว มหาสมุทฺโท มหตํ ภูตานํ อาวาโส, ตตฺริเม ภูตา, ติมิ ติมิงฺคโล ติมิติมิงฺคโล อสุรา นาคา คนฺธพฺพา, สนฺติ มหาสมุทฺเท โยชนสติกาปิ อตฺตภาวา, ทฺวิโยชนสติกาปิ อตฺตภาวา, ติโยชนสติกาปิ อตฺตภาวา, จตุโยชนสติกาปิ อตฺตภาวา, ปญฺจโยชนสติกาปิ อตฺตภาวา, เอวเมว โข ภิกฺขเว อยํ ธมฺมวินโย มหตํ ภูตานํ อาวาโส, ตตฺริเม ภูตา, โสตาปนฺโน, โสตาปตฺติผลสจฺฉิกิริยาย ปฏิปนฺโน, สกทาคามี, \csromanfolio{424}สกทาคามิผลสจฺฉิกิริยาย ปฏิปนฺโน, อนาคามี, อนาคามิผลสจฺฉิกิริยาย ปฏิปนฺโน อรหา, อรหตฺตผลสจฺฉิกิริยาย ปฏิปนฺโน.\csromanspacer{} ยมฺปิ ภิกฺขเว อยํ ธมฺมวินโย มหตํ ภูตานํ อาวาโส, ตตฺริเม ภูตา, โสตาปนฺโน, โสตาปตฺติผลสจฺฉิกิริยาย ปฏิปนฺโน ฯเปฯ อรหา, อรหตฺตผลสจฺฉิกิริยาย ปฏิปนฺโน, อยํ ภิกฺขเว อิมสฺมิํ ธมฺมวินเย อฏฺฐโม อจฺฉริโย อพฺภุโต ธมฺโม, ยํ ทิสฺวา ทิสฺวา ภิกฺขู อิมสฺมิํ ธมฺมวินเย อภิรมนฺติ.\csromanspacer{} อิเม โข ภิกฺขเว อิมสฺมิํ ธมฺมวินเย อฏฺฐ อจฺฉริยา อพฺภุตา ธมฺมา, เย ทิสฺวา ทิสฺวา ภิกฺขู อิมสฺมิํ ธมฺมวินเย อภิรมนฺตีติ.}

\prose{อถ โข ภควา เอตมตฺถํ วิทิตฺวา ตายํ เวลายํ อิมํ อุทานํ อุทาเนสิ\csromandash{}}

\csromangathameasure{*“ฉนฺนมติวสฺสติ, วิวฏํ นาติวสฺสติ. \\ ตสฺมา ฉนฺนํ วิวเรถ, เอวํ ตํ นาติวสฺสตี”ติ.}
\csromangathagroup{\csromangathastanza{\csromansymbolfootnote{*}{วิ 1. 44 ปิฏฺเฐ อุทาเนปิ.}“ฉนฺนมติวสฺสติ\footnote{สุฉนฺทมติวสฺสติ (ก)}, วิวฏํ นาติวสฺสติ. \\ ตสฺมา ฉนฺนํ วิวเรถ, เอวํ ตํ นาติวสฺสตี”ติ.}}

\csromanmatikaanchor{mtk.196}
\csromantocmarkref{section}{4. ปาติโมกฺขสวนารห}{mtk.196}
\bookmark[dest={mtk.196},level=1]{4. ปาติโมกฺขสวนารห}
\csromanheader{1}{4. ปาติโมกฺขสวนารห}
\proseitem{386}{อถ โข ภควา ภิกฺขู อามนฺเตสิ\csromandash{}นทานาหํ ภิกฺขเว อิโต ปรํ อุโปสถํ กริสฺสามิ, ปาติโมกฺขํ อุทฺทิสิสฺสามิ, ตุเมฺหวทานิ ภิกฺขเว อิโต ปรํ อุโปสถํ กเรยฺยาถ, ปาติโมกฺขํ อุทฺทิเสยฺยาถ.\csromanspacer{} อฏฺฐานเมตํ ภิกฺขเว อนวกาโส, ยํ ตถาคโต อปริสุทฺธาย ปริสาย อุโปสถํ กเรยฺย, ปาติโมกฺขํ อุทฺทิเสยฺย.\csromanspacer{} น จ ภิกฺขเว\footnote{น จ ภิกฺขเว ภิกฺขุนา (สฺยา, กํ)} สาปตฺติเกน ปาติโมกฺขํ โสตพฺพํ, โย สุเณยฺย, อาปตฺติ ทุกฺกฏสฺส.\csromanspacer{} อนุชานามิ ภิกฺขเว โย สาปตฺติโก ปาติโมกฺขํ สุณาติ, ตสฺส ปาติโมกฺขํ ฐเปตุํ.\csromanspacer{} เอวญฺจ ปน ภิกฺขเว ฐเปตพฺพํ, ตทหุโปสเถ จาตุทฺทเส วา ปนฺนรเส วา ตสฺมิํ ปุคฺคเล สมฺมุขีภูเต สํฆมชฺเฌ อุทาหริตพฺพํ\csromandash{}}

\prose{“สุณาตุ เม ภนฺเต สํโฆ, อิตฺถนฺนาโม ปุคฺคโล สาปตฺถิโก, ตสฺส ปาติโมกฺขํ ฐเปมิ ‘น ตสฺมิํ สมฺมุขีภูเต ปาติโมกฺขํ อุทฺทิสิตพฺพนฺ’ติ”.\csromanspacer{} ฐปิตํ โหติ ปาติโมกฺขนฺติ.}

\prose{\csromanfolio{425}เตน โข ปน สมเยน ฉพฺพคฺคิยา ภิกฺขู “นาเมฺห โกจิ ชานาตี”ติ สาปตฺติกาว ปาติโมกฺขํ สุณนฺติ, เถรา ภิกฺขู ปรจิตฺตวิทุโน ภิกฺขูนํ อาโรเจนฺติ “อิตฺถนฺนาโม จ อิตฺถนฺนาโม จ อาวุโส ฉพฺพคฺคิยา ภิกฺขู ‘นาเมฺห โกจิ ชานาตี’ติ สาปตฺติกาว ปาติโมกฺขํ สุณนฺตี”ติ.\csromanspacer{} อสฺโสสุํ โข ฉพฺพคฺคิยา ภิกฺขู เถรา กิร ภิกฺขู ปรจิตฺตวิทุโน อเมฺห ภิกฺขูนํ อาโรเจนฺติ “อิตฺถนฺนาโม จ อิตฺถนฺนาโม จ อาวุโส ฉพฺพคฺคิยา ภิกฺขู ‘นาเมฺห โกจิ ชานาตี’ติ สาปตฺติกาว ปาติโมกฺขํ สุณนฺตี”ติ.}

\prose{เต “ปุรมฺหากํ เปสลา ภิกฺขู ปาติโมกฺขํ ฐเปนฺตี”ติ ปฏิกจฺเจว สุทฺธานํ ภิกฺขูนํ อนาปตฺติกานํ อวตฺถุสฺมิํ อการเณ ปาติโมกฺขํ ฐเปนฺติ.\csromanspacer{} เย เต ภิกฺขู อปฺปิจฺฉา ฯเปฯ เต อุชฺฌายนฺติ ขิยฺยนฺติ วิปาเจนฺติ “กถํ หิ นาม ฉพฺพคฺคิยา ภิกฺขู สุทฺธานํ ภิกฺขูนํ อนาปตฺติกานํ อวตฺถุสฺมิํ อการเณ ปาติโมกฺขํ ฐเปสฺสนฺตี”ติ.\csromanspacer{} อถ โข เต ภิกฺขู ภควโต เอตมตฺถํ อาโรเจสุํ ฯเปฯ.\csromanspacer{} สจฺจํ กิร ภิกฺขเว ฉพฺพคฺคิยา ภิกฺขู สุทฺธานํ ภิกฺขูนํ อนาปตฺติกานํ อวตฺถุสฺมิํ อการเณ ปาติโมกฺขํ ฐเปนฺตีติ.\csromanspacer{} สจฺจํ ภควาติ ฯเปฯ วิครหิตฺวา ฯเปฯ ธมฺมิํ กถํ กตฺวา ภิกฺขู อามนฺเตสิ\csromandash{}น ภิกฺขเว สุทฺธานํ ภิกฺขูนํ อนาปตฺติกานํ อวตฺถุสฺมิํ อการเณ ปาติโมกฺขํ ฐเปตพฺพํ, โย ฐเปยฺย, อาปตฺติ ทุกฺกฏสฺส.}

\csromanmatikaanchor{mtk.197}
\csromantocmarkref{section}{5. ธมฺมิกาธมฺมิกปาติโมกฺขฏฺฐปน}{mtk.197}
\bookmark[dest={mtk.197},level=1]{5. ธมฺมิกาธมฺมิกปาติโมกฺขฏฺฐปน}
\csromanheader{1}{5. ธมฺมิกาธมฺมิกปาติโมกฺขฏฺฐปน}
\proseitem{387}{เอกํ ภิกฺขเว อธมฺมิกํ ปาติโมกฺขฏฺฐปนํ, เอกํ ธมฺมิกํ ปาติโมกฺขฏฺฐปนํ\footnote{เอกํ ธมฺมิกํ (สี, สฺยา)}.\csromanspacer{} เทฺว อธมฺมิกานิ ปาติโมกฺขฏฺฐปนานิ, เทฺว ธมฺมิกานิ.\csromanspacer{} ตีณิ อธมฺมิกานิ ปาติโมกฺขฏฺฐปนานิ, ตีณิ ธมฺมิกานิ.\csromanspacer{} จตฺตาริ อธมฺมิกานิ ปาติโมกฺขฏฺฐปนานิ, จตฺตาริ ธมฺมิกานิ.\csromanspacer{} ปญฺจ อธมฺมิกานิ ปาติโมกฺขฏฺฐปนานิ, ปญฺจ ธมฺมิกานิ.\csromanspacer{} ฉ อธมฺมิกานิ ปาติโมกฺขฏฺฐปนานิ, ฉ ธมฺมิกานิ.\csromanspacer{} สตฺต อธมฺมิกานิ ปาติโมกฺขฏฺฐปนานิ, สตฺต ธมฺมิกานิ.\csromanspacer{} อฏฺฐ อธมฺมิกานิ ปาติโมกฺขฏฺฐปนานิ, อฏฺฐ ธมฺมิกานิ.\csromanspacer{} นว อธมฺมิกานิ ปาติโมกฺขฏฺฐปนานิ, นว ธมฺมิกานิ.\csromanspacer{} ทส อธมฺมิกานิ ปาติโมกฺขฏฺฐปนานิ, ทส ธมฺมิกานิ.}

\prose{\csromanfolio{426}กตมํ เอกํ อธมฺมิกํ ปาติโมกฺขฏฺฐปนํ.\csromanspacer{} อมูลิกาย สีลวิปตฺติยา ปาติโมกฺขํ ฐเปติ.\csromanspacer{} อิทํ เอกํ อธมฺมิกํ ปาติโมกฺขฏฺฐปนํ.\csromanspacer{} กตมํ เอกํ ธมฺมิกํ ปาติโมกฺขฏฺฐปนํ.\csromanspacer{} สมูลิกาย สีลวิปตฺติยา ปาติโมกฺขํ ฐเปติ.\csromanspacer{} อิทํ เอกํ ธมฺมิกํ ปาติโมกฺขฏฺฐปนํ.}

\prose{กตมานิ เทฺว อธมฺมิกานิ ปาติโมกฺขฏฺฐปนานิ.\csromanspacer{} อมูลิกาย สีลวิปตฺติยา ปาติโมกฺขํ ฐเปติ, อมูลิกาย อาจารวิปตฺติยา ปาติโมกฺขํ ฐเปติ.\csromanspacer{} อิมานิ เทฺว อธมฺมิกานิ ปาติโมกฺขฏฺฐปนานิ.\csromanspacer{} กตมานิ เทฺว ธมฺมิกานิ ปาติโมกฺขฏฺฐปนานิ.\csromanspacer{} สมูลิกาย สีลวิปตฺติยา ปาติโมกฺขํ ฐเปติ, สมูลิกาย อาจารวิปตฺติยา ปาติโมกฺขํ ฐเปติ.\csromanspacer{} อิมานิ เทฺว ธมฺมิกานิ ปาติโมกฺขฏฺฐปนานิ.}

\prose{กตมานิ ตีณิ อธมฺมิกานิ ปาติโมกฺขฏฺฐปนานิ.\csromanspacer{} อมูลิกาย สีลวิปตฺติยา ปาติโมกฺขํ ฐเปติ, อมูลิกาย อาจารวิปตฺติยา ปาติโมกฺขํ ฐเปติ, อมูลิกาย ทิฏฺฐิวิปตฺติยา ปาติโมกฺขํ ฐเปติ.\csromanspacer{} อิมานิ ตีณิ อธมฺมิกานิ ปาติโมกฺขฏฺฐปนานิ.\csromanspacer{} กตมานิ ตีณิ ธมฺมิกานิ ปาติโมกฺขฏฺฐปนานิ.\csromanspacer{} สมูลิกาย สีลวิปตฺติยา ปาติโมกฺขํ ฐเปติ, สมูลิกาย อาจารวิปตฺติยา ปาติโมกฺขํ ฐเปติ, สมูลิกาย ทิฏฺฐิวิปตฺติยา ปาติโมกฺขํ ฐเปติ.\csromanspacer{} อิมานิ ตีณิ ธมฺมิกานิ ปาติโมกฺขฏฺฐปนานิ.}

\prose{กตมานิ จตฺตาริ อธมฺมิกานิ ปาติโมกฺขฏฺฐปนานิ.\csromanspacer{} อมูลิกาย สีลวิปตฺติยา ปาติโมกฺขํ ฐเปติ, อมูลิกาย อาจารวิปตฺติยา ปาติโมกฺขํ ฐเปติ, อมูลิกาย ทิฏฺฐิวิปตฺติยา ปาติโมกฺขํ ฐเปติ, อมูลิกาย อาชีววิปตฺติยา ปาติโมกฺขํ ฐเปติ.\csromanspacer{} อิมานิ จตฺตาริ อธมฺมิกานิ ปาติโมกฺขฏฺฐปนานิ.\csromanspacer{} กตมานิ จตฺตาริ ธมฺมิกานิ ปาติโมกฺขฏฺฐปนานิ.\csromanspacer{} สมูลิกาย สีลวิปตฺติยา ปาติโมกฺขํ ฐเปติ, สมูลิกาย อาจารวิปตฺติยา ปาติโมกฺขํ ฐเปติ, สมูลิกาย ทิฏฺฐิวิปตฺติยา ปาติโมกฺขํ ฐเปติ, สมูลิกาย อาชีววิปตฺติยา, ปาติโมกฺขํ ฐเปติ.\csromanspacer{} อิมานิ จตฺตาริ ธมฺมิกานิ ปาติโมกฺขฏฺฐปนานิ.}

\prose{กตมานิ ปญฺจ อธมฺมิกานิ ปาติโมกฺขฏฺฐปนานิ.\csromanspacer{} อมูลเกน ปาราชิเกน ปาติโมกฺขํ ฐเปติ, อมูลเกน สํฆาทิเสเสน ฯเปฯ อมูลเกน ปาจิตฺติเยน.\csromanspacer{} อมูลเกน ปาฏิเทสนีเยน.\csromanspacer{} อมูลเกน ทุกฺกเฏน \csromanfolio{427}ปาติโมกฺขํ ฐเปติ.\csromanspacer{} อิมานิ ปญฺจ อธมฺมิกานิ ปาติโมกฺขฏฺฐปนานิ.\csromanspacer{} กตมานิ ปญฺจ ธมฺมิกานิ ปาติโมกฺขฏฺฐปนานิ.\csromanspacer{} สมูลเกน ปาราชิเกน ปาติโมกฺขํ ฐเปติ, สมูลเกน สํฆาทิเสเสน ฯเปฯ ปาจิตฺติเยน.\csromanspacer{} ปาฏิเทสนีเยน.\csromanspacer{} ทุกฺกเฏน ปาติโมกฺขํ ฐเปติ.\csromanspacer{} อิมานิ ปญฺจ ธมฺมิกานิ ปาติโมกฺขฏฺฐปนานิ.}

\prose{กตมานิ ฉ อธมฺมิกานิ ปาติโมกฺขฏฺฐปนานิ.\csromanspacer{} อมูลิกาย สีลวิปตฺติยา ปาติโมกฺขํ ฐเปติ อกตาย, อมูลิกาย สีลวิปตฺติยา ปาติโมกฺขํ ฐเปติ กตาย, อมูลิกาย อาจารวิปตฺติยา ปาติโมกฺขํ ฐเปติ อกตาย, อมูลิกาย อาจารวิปตฺติยา ปาติโมกฺขํ ฐเปติ กตาย, อมูลิกาย ทิฏฺฐิวิปตฺติยา ปาติโมกฺขํ ฐเปติ อกตาย, อมูลิกาย ทิฏฺฐิวิปตฺติยา ปาติโมกฺขํ ฐเปติ กตาย.\csromanspacer{} อิมานิ ฉ อธมฺมิกานิ ปาติโมกฺขฏฺฐปนานิ.\csromanspacer{} กตมานิ ฉ ธมฺมิกานิ ปาติโมกฺขฏฺฐปนานิ.\csromanspacer{} สมูลิกาย สีลวิปตฺติยา ปาติโมกฺขํ ฐเปติ อกตาย, สมุลิกาย สีลวิปตฺติยา ปาติโมกฺขํ ฐเปติ กตาย, สมูลิกาย อาจารวิปตฺติยา ฯเปฯ ทิฏฺฐิวิปตฺติยา ปาติโมกฺขํ ฐเปติ อกตาย ฯเปฯ กตาย.\csromanspacer{} อิมานิ ฉ ธมฺมิกานิ ปาติโมกฺขฏฺฐปนานิ.}

\prose{กตมานิ สตฺต อธมฺมิกานิ ปาติโมกฺขฏฺฐปนานิ.\csromanspacer{} อมูลเกน ปาราชิเกน ปาติโมกฺขํ ฐเปติ, อมูลเกน สํฆาทิเสเสน ฯเปฯ ถุลฺลจฺจเยน.\csromanspacer{} ปาจิตฺติเยน.\csromanspacer{} ปาฏิเทสนีเยน.\csromanspacer{} ทุกฺกเฏน.\csromanspacer{} ทุพฺภาสิเตน ปาติโมกฺขํ ฐเปติ.\csromanspacer{} อิมานิ สตฺต อธมฺมิกานิ ปาติโมกฺขฏฺฐปนานิ.\csromanspacer{} กตมนิ สตฺต ธมฺมิกานิ ปาติโมกฺขฏฺฐปนานิ.\csromanspacer{} สมูลเกน ปาราชิเกน ปาติโมกฺขํ ฐเปติ, สมูลเกน สํฆาทิเสเสน ฯเปฯ ถุลฺลจฺจเยน.\csromanspacer{} ปาจิตฺติเยน.\csromanspacer{} ปาฏิเทสนีเยน.\csromanspacer{} ทุกฺกเฏน.\csromanspacer{} ทุพฺภาสิเตน ปาติโมกฺขํ ฐเปติ.\csromanspacer{} อิมานิ สตฺต ธมฺมิกานิ ปาติโมกฺขฏฺฐปนานิ.}

\prose{กตมานิ อฏฺฐ อธมฺมิกานิ ปาติโมกฺขฏฺฐปนานิ.\csromanspacer{} อมูลิกาย สีลวิปตฺติยา ปาติโมกฺขํ ฐเปติ อกตาย, อมูลิกาย สีลวิปตฺติยา ปาติโมกฺขํ ฐเปติ กตาย, อมูลิกาย อาจารวิปตฺติยา ฯเปฯ.\csromanspacer{} ทิฏฺฐิวิปตฺติยา.\csromanspacer{} อาชีววิปตฺติยา ปาติโมกฺขํ ฐเปติ อกตาย ฯเปฯ กตาย.\csromanspacer{} อิมานิ อฏฺฐ อธมฺมิกานิ ปาติโมกฺขฏฺฐปนานิ.\csromanspacer{} กตมานิ อฏฺฐ ธมฺมิกานิ ปาติโมกฺขฏฺฐปนานิ.\csromanspacer{} สมูลิกาย สีลวิปตฺติยา ปาติโมกฺขํ ฐเปติ อกตาย, สมูลิกาย สีลวิปตฺติยา ปาติโมกฺขํ \csromanfolio{428}ฐเปติ กตาย, สมูลิกาย อาจารวิปตฺติยา ฯเปฯ ทิฏฺฐิวิปตฺติยา.\csromanspacer{} อาชีววิปตฺติยา ปาติโมกฺขํ ฐเปติ อกตาย ฯเปฯ กตาย.\csromanspacer{} อิมานิ อฏฺฐ ธมฺมิกานิ ปาติโมกฺขฏฺฐปนานิ.}

\prose{กตมานิ นว อธมฺมิกานิ ปาติโมกฺขฏฺฐปนานิ.\csromanspacer{} อมูลิกาย สีลวิปตฺติยา ปาติโมกฺขํ ฐเปติ อกตาย, อมูลิกาย สีลวิปตฺติยา ปาติโมกฺขํ ฐเปติ กตาย, อมูลิกาย สีลวิปตฺติยา ปาติโมกฺขํ ฐเปติ กตากตาย, อมูลิกาย อาจารวิปตฺติยา ฯเปฯ ทิฏฺฐิวิปตฺติยา ปาติโมกฺขํ ฐเปติ อกตาย ฯเปฯ กตาย ฯเปฯ กตากตาย.\csromanspacer{} อิมานิ นว อธมฺมิกานิ ปาติโมกฺขฏฺฐปนานิ.\csromanspacer{} กตามานิ นว ธมฺมิกานิ ปาติโมกฺขฏฺฐปนานิ.\csromanspacer{} สมูลิกาย สีลวิปตฺติยา ปาติโมกฺขํ ฐเปติ อกตาย, สมูลิกาย สีลวิปตฺติยา ปาติโมกฺขํ ฐเปติ กตาย, สมูลิกาย, สีลวิปตฺติยา ปาติโมกฺขํ ฐเปติ กตากตาย, สมูลิกาย, อาจารวิปตฺติยา ฯเปฯ ทิฏฺฐิวิปตฺติยา ปาติโมกฺขํ ฐเปติ อกตาย ฯเปฯ กตาย ฯเปฯ กตากตาย.\csromanspacer{} อิมานิ นว ธมฺมิกานิ ปาติโมกฺขฏฺฐปนานิ.}

\prose{กตมานิ ทส อธมฺมิกานิ ปาติโมกฺขฏฺฐปนานิ.\csromanspacer{} น ปาราชิโก ตสฺสํ ปริสายํ นิสินฺโน โหติ, น ปาราชิกกถา วิปฺปกตา โหติ, น สิกฺขํ ปจฺจกฺขาตโก ตสฺสํ ปริสายํ นิสินฺโน โหติ, น สิกฺขํ ปจฺจกฺขาตกถา วิปฺปกตา โหติ, ธมฺมิกํ สามคฺคิํ อุเปติ, น ธมฺมิกํ สามคฺคิํ ปจฺจาทิยติ, น ธมฺมิกาย สามคฺคิยา ปจฺจาทานกถา วิปฺปกตา โหติ, น สีลวิปตฺติยา ทิฏฺฐสุตปริสงฺกิโต โหติ, น อาจารวิปตฺติยา ทิฏฺฐสุตปริสงฺกิโต โหติ, น ทิฏฺฐิวิปตฺติยา ทิฏฺฐสุตปริสงฺกิโต โหติ.\csromanspacer{} อิมานิ ทส อธมฺมิกานิ ปาติโมกฺขฏฺฐปนานิ.\csromanspacer{} กตมานิ ทส ธมฺมิกานิ ปาติโมกฺขฏฺฐปนานิ.\csromanspacer{} ปาราชิโก ตสฺสํ ปริสายํ นิสินฺโน โหติ, ปาราชิกกถา วิปฺปกตา โหติ, สิกฺขํ ปจฺจกฺขาตโก ตสฺสํ ปริสายํ นิสินฺโน โหติ, สิกฺขํ ปจฺจกฺขาตกถา วิปฺปกตา โหติ, ธมฺมิกํ สามคฺคิํ น อุเปติ, ธมฺมิกํ สามคฺคิํ ปจฺจาทิยติ, ธมฺมิกาย สามคฺคิยา ปจฺจาทานกถา วิปฺปกตา โหติ, สีลวิปตฺติยา ทิฏฺฐสุตปริสงฺกิโต โหติ, อาจารวิปตฺติยา ทิฏฺฐสุตปริสงฺกิโต โหติ, ทิฏฺฐิวิปตฺติยา ทิฏฺฐสุตปริสงฺกิโต โหติ.\csromanspacer{} อิมานิ ทส ธมฺมิกานิ ปาติโมกฺขฏฺฐปนานิ.}

\csromanmatikaanchor{mtk.198}
\csromantocmarkref{section}{6. ธมฺมิกปาติโมกฺขฏฺฐปน}{mtk.198}
\bookmark[dest={mtk.198},level=1]{6. ธมฺมิกปาติโมกฺขฏฺฐปน}
\csromanheader{1}{6. ธมฺมิกปาติโมกฺขฏฺฐปน}
\proseitem{388}{\csromanfolio{429}กถํ ปาราชิโก ตสฺสํ ปริสายํ นิสินฺโน โหติ.\csromanspacer{} อิธ ภิกฺขเว เยหิ อากาเรหิ เยหิ ลิงฺเคหิ เยหิ นิมิตฺเตหิ ปาราชิกสฺส ธมฺมสฺส อชฺฌาปตฺติ โหติ, เตหิ อากาเรหิ เตหิ ลิงฺเคหิ เตหิ นิมิตฺเตหิ ภิกฺขู ภิกฺขุํ ปสฺสติ ปาราชิกํ ธมฺมํ อชฺฌาปชฺชนฺตํ.\csromanspacer{} น เหว โข ภิกฺขุ ภิกฺขุํ ปสฺสติ ปาราชิกํ ธมฺมํ อชฺฌาปชฺชนฺตํ, อปิ จ อญฺโญ ภิกฺขุ ภิกฺขุสฺส อาโรเจติ “อิตฺถนฺนาโม อาวุโส ภิกฺขุ ปาราชิกํ ธมฺมํ อชฺฌาปนฺโน”ติ.\csromanspacer{} น เหว โข ภิกฺขุ ภิกฺขุํ ปสฺสติ ปาราชิกํ ธมฺมํ อชฺฌาปชฺชนฺตํ, นาปิ\footnote{นาปิ จ (ก)} อญฺโญ ภิกฺขุ ภิกฺขุสฺส อาโรเจติ “อิตฺถนฺนาโม อาวุโส ภิกฺขุ ปาราชิกํ ธมฺมํ อชฺฌาปนฺโน”ติ, อปิ จ โสว ภิกฺขุ ภิกฺขุสฺส อาโรเจติ “อหํ อาวุโส ปาราชิกํ ธมฺมํ อชฺฌาปนฺโน”ติ.\csromanspacer{} อากงฺขมาโน ภิกฺขเว ภิกฺขุ เตน ทิฏฺเฐน เตน สุเตน ตาย ปริสงฺกาย ตทหุโปสเถ จาตุทฺทเส วา ปนฺนรเส วา ตสฺมิํ ปุคฺคเล สมฺมุขีภูเต สํฆมชฺเฌ อุทาหเรยฺย\csromandash{}}

\prose{“สุณาตุ เม ภนฺเต สํโฆ, อิตฺถนฺนาโม ปุคฺคโล ปาราชิกํ ธมฺมํ อชฺฌาปนฺโน, ตสฺส ปาติโมกฺขํ ฐเปมิ ‘น ตสฺมิํ สมฺมุขีภูเต ปาติโมกฺขํ อุทฺทิสิตพฺพนฺ’ติ”.\csromanspacer{} ธมฺมิกํ ปาติโมกฺขฏฺฐปนํ.}

\proseitem{389}{ภิกฺขุสฺส ปาติโมกฺเข ฐปิเต ปริสา วุฏฺฐาติ ทสนฺนํ อนฺตรายานํ อญฺญตเรน\footnote{อญฺญตเรน อนฺตราเยน (สฺยา, กํ)} ราชนฺตราเยน วา โจรนฺตราเยน วา อคฺยนฺตราเยน วา อุทกนฺตราเยน วา มนุสฺสนฺตราเยน วา อมนุสฺสนฺตราเยน วา วาฬนฺตราเยน วา สรีสปนฺตราเยน วา ชีวิตนฺตราเยน วา พฺรหฺมจริยนฺตราเยน วา, อากงฺขมาโน ภิกฺขเว ภิกฺขุ ตสฺมิํ วา อาวาเส อญฺญสฺมิํ วา อาวาเส ตสฺมิํ ปุคฺคเล สมฺมุขีภูเต สํฆมชฺเฌ อุทาหเรยฺย⎯}

\prose{“สุณาตุ เม ภนฺเต สํโฆ, อิตฺถนฺนามสฺส ปุคฺคลสฺส ปาราชิกกถา วิปฺปกตา, ตํ วตฺถุ อวินิจฺฉิตํ, ยทิ สํฆสฺส ปตฺตกลฺลํ, สํโฆ ตํ วตฺถุํ วินิจฺฉิเนยฺยา”ติ.}

\prose{เอวญฺเจตํ ลเภถ, อิจฺเจตํ กุสลํ.\csromanspacer{} โน เจ ลเภถ, ตทหุโปสเถ จาตุทฺทเส วา ปนฺนรเส วา ตสฺมิํ ปุคฺคเล สมฺมุขีภูเต สํฆมชฺเฌ อุทาหริตพฺพํ\csromandash{}}

\prose{\csromanfolio{430}“สุณาตุ เม ภนฺเต สํโฆ, อิตฺถนฺนามสฺส ปุคฺคลสฺส ปาราชิกกถา วิปฺปกตา, ตํ วตฺถุ อวินิจฺฉิตํ, ตสฺส ปาติโมกฺขํ ฐเปมิ ‘น ตสฺมิํ สมฺมุขีภูเต ปาติโมกฺขํ อุทฺทิสิตพฺพนฺ’ติ”.\csromanspacer{} ธมฺมิกํ ปาติโมกฺขฏฺฐปนํ.}

\proseitem{390}{กถํ สิกฺขํ ปจฺจกฺขาตโก ตสฺสํ ปริสายํ นิสินฺโน โหติ.\csromanspacer{} อิธ ปน ภิกฺขเว ภิกฺขุ เยหิ อากาเรหิ เยหิ ลิงฺเคหิ เยหิ นิมิตฺเตหิ สิกฺขา ปจฺจกฺขาตา โหติ, เตหิ อากาเรหิ เตหิ ลิงฺเคหิ เตหิ นิมิตฺเตหิ ภิกฺขุ ภิกฺขุํ ปสฺสติ สิกฺขํ ปจฺจกฺขนฺตํ.\csromanspacer{} น เหว โข ภิกฺขุ ภิกฺขุํ ปสฺสติ สิกฺขํ ปจฺจกฺขนฺตํ, อปิ จ อญฺโญ ภิกฺขุ ภิกฺขุสฺส อาโรเจติ “อิตฺถนฺนาเมน อาวุโส ภิกฺขุนา สิกฺขา ปจฺจกฺขาตา”ติ.\csromanspacer{} น เหว โข ภิกฺขุ ภิกฺขุํ ปสฺสติ สิกฺขํ ปจฺจกฺขนฺตํ, นาปิ อญฺโญ ภิกฺขุ ภิกฺขุสฺส อาโรเจติ “อิตฺถนฺนาเมน อาวุโส ภิกฺขุนา สิกฺขา ปจฺจกฺขาตา”ติ.\csromanspacer{} อปิ จ โสว ภิกฺขุ ภิกฺขุสฺส อาโรเจติ “มยา อาวุโส สิกฺขา ปจฺจกฺขาตา”ติ.\csromanspacer{} อากงฺขมาโน ภิกฺขเว ภิกฺขุ เตน ทิฏฺเฐน เตน สุเตน ตาย ปริสงฺกาย ตทหุโปสเถ จาตุทฺทเส วา ปนฺนรเส วา ตสฺมิํ ปุคฺคเล สมฺมุขีภูเต สํฆมชฺเฌ อุทาหเรยฺย\csromandash{}}

\prose{“สุณาตุ เม ภนฺเต สํโฆ, อิตฺถนฺนาเมน ปุคฺคเลน สิกฺขา ปจฺจกฺขาตา, ตสฺส ปาติโมกฺขํ ฐเปมิ ‘น ตสฺมิํ สมฺมุขีภูเต ปาติโมกฺขํ อุทฺทิสิตพฺพนฺ’ติ”.\csromanspacer{} ธมฺมิกํ ปาติโมกฺขฏฺฐปนํ.}

\proseitem{391}{ภิกฺขุสฺส ปาติโมกฺเข ฐปิเต ปริสา วุฏฺฐาติ ทสนฺนํ อนฺตรายานํ อญฺญตเรน ราชนฺตราเยน วา ฯเปฯ พฺรหฺมจริยนฺตราเยน วา.\csromanspacer{} อากงฺขมาโน ภิกฺขเว ภิกฺขุ ตสฺมิํ วา อาวาเส อญฺญสฺมิํ วา อาวาเส ตสฺมิํ ปุคฺคเล สมฺมุขีภูเต สํฆมชฺเฌ อุทาหเรยฺย\csromandash{}}

\prose{“สุณาตุ เม ภนฺเต สํโฆ, อิตฺถนฺนามสฺส ปุคฺคลสฺส สิกฺขํ ปจฺจกฺขาตกถา\footnote{สิกฺขาปจฺจกฺขาตกถา (สี)} วิปฺปกตา, ตํ วตฺถุ อวินิจฺฉิตํ, ยทิ สํฆสฺส ปตฺตกลฺลํ, สํโฆ ตํ วตฺถุํ วินิจฺฉิเนยฺยา”ติ.}

\prose{เอวญฺเจตํ ลเภถ, อิจฺเจตํ กุสลํ.\csromanspacer{} โน เจ ลเภถ, ตทหุโปสเถ จาตุทฺทเส วา ปนฺนรเส วา ตสฺมิํ ปุคฺคเล สมฺมุขีภูเต สํฆมชฺเฌ อุทาหริตพฺพํ\csromandash{}}

\prose{\csromanfolio{431}“สุณาตุ เม ภนฺเต สํโฆ, อิตฺถนฺนามสฺส ปุคฺคลสฺส สิกฺขํ ปจฺจกฺขาตกถา วิปฺปกตา, ตํ วตฺถุ อวินิจฺฉิตํ, ตสฺส ปาติโมกฺขํ ฐเปมิ ‘น ตสฺมิํ สมฺมุขีภูเต ปาติโมกฺขํ อุทฺทิสิตพฺพนฺ’ติ”.\csromanspacer{} ธมฺมิกํ ปาติโมกฺขฏฺฐปนํ.}

\proseitem{392}{กถํ ธมฺมิกํ สามคฺคิํ น อุเปติ.\csromanspacer{} อิธ ปน ภิกฺขเว ภิกฺขุ เยหิ อากาเรหิ เยหิ ลิงฺเคหิ เยหิ นิมิตฺเตหิ ธมฺมิกาย สามคฺคิยานุปคมนํ โหติ, เตหิ อากาเรหิ เตหิ ลิงฺเคหิ เตหิ นิมิตฺเตหิ ภิกฺขุ ภิกฺขุํ ปสฺสติ ธมฺมิกํ สามคฺคิํ น อุเปนฺตํ.\csromanspacer{} น เหว โข ภิกฺขุ ภิกฺขุํ ปสฺสติ ธมฺมิกํ สามคฺคิํ น อุเปนฺตํ, อปิ จ อญฺโญ ภิกฺขุ ภิกฺขุสฺส อาโรเจติ “อิตฺถนฺนาโม อาวุโส ภิกฺขุ ธมฺมิกํ สามคฺคิํ น อุเปตี”ติ.\csromanspacer{} น เหว โข ภิกฺขุ ภิกฺขุํ ปสฺสติ ธมฺมิกํ สามคฺคิํ น อุเปนฺตํ, นาปิ อญฺโญ ภิกฺขุ ภิกฺขุสฺส อาโรเจติ “อิตฺถนฺนาโม อาวุโส ภิกฺขุ ธมฺมิกํ สามคฺคิํ น อุเปตี”ติ, อปิ จ โสว ภิกฺขุ ภิกฺขุสฺส อาโรเจติ “อหํ อาวุโส ธมฺมิกํ สามคฺคิํ น อุเปมี”ติ.\csromanspacer{} อากงฺขมาโน ภิกฺขเว ภิกฺขุ เตน ทิฏฺเฐน เตน สุเตน ตาย ปริสงฺกาย ตทหุโปสเถ จาตุทฺทเส วา ปนฺนรเส วา ตสฺมิํ ปุคฺคเล สมฺมุขีภูเต สํฆมชฺเฌ อุทาหเรยฺย\csromandash{}}

\prose{“สุณาตุ เม ภนฺเต สํโฆ, อิตฺถนฺนาโม ปุคฺคโล ธมฺมิกํ สามคฺคิํ น อุเปติ, ตสฺส ปาติโมกฺขํ ฐเปมิ ‘น ตสฺมิํ สมฺมุขีภูเต ปาติโมกฺขํ อุทฺทิสิตพฺพนฺ’ติ”.\csromanspacer{} ธมฺมิกํ ปาติโมกฺขฏฺฐปนํ.}

\proseitem{393}{กถํ ธมฺมิกํ สามคฺคิํ ปจฺจาทิยติ.\csromanspacer{} อิธ ภิกฺขเว ภิกฺขุ เยหิ อากาเรหิ เยหิ ลิงฺเคหิ เยหิ นิมิตฺเตหิ ธมฺมิกาย สามคฺคิยา ปจฺจาทานํ โหติ, เตหิ อากาเรหิ เตหิ ลิงฺเคหิ เตหิ นิมิตฺเตหิ ภิกฺขุ ภิกฺขุํ ปสฺสติ ธมฺมิกํ สามคฺคิํ ปจฺจาทิยนฺตํ.\csromanspacer{} น เหว โข ภิกฺขุ ภิกฺขุํ ปสฺสติ ธมฺมิกํ สามคฺคิํ ปจฺจาทิยนฺตํ, อปิ จ อญฺโญ ภิกฺขุ ภิกฺขุสฺส อาโรเจติ “อิตฺถนฺนาโม อาวุโส ภิกฺขุ ธมฺมิกํ สามคฺคิํ ปจฺจาทิยตี”ติ.\csromanspacer{} น เหว โข ภิกฺขุ ภิกฺขุํ ปสฺสติ ธมฺมิกํ สามคฺคิํ ปจฺจาทิยนฺตํ, นาปิ อญฺโญ ภิกฺขุ ภิกฺขุสฺส อาโรเจติ “อิตฺถนฺนาโม อาวุโส ภิกฺขุ ธมฺมิกํ สามคฺคิํ ปจฺจาทิยตี”ติ, อปิ จ โสว ภิกฺขุ ภิกฺขุสฺส อาโรเจติ “อหํ อาวุโส ธมฺมิกํ สามคฺคิํ \csromanfolio{432}ปจฺจาทิยามี”ติ.\csromanspacer{} อากงฺขมาโน ภิกฺขเว ภิกฺขุ เตน ทิฏฺเฐน เตน สุเตน ตาย ปริสงฺกาย ตทหุโปสเถ จาตุทฺทเส วา ปนฺนรเส วา ตสฺมิํ ปุคฺคเล สมฺมุขีภูเต สํฆมชฺเฌ อุทาหเรยฺย\csromandash{}}

\prose{“สุณาตุ เม ภนฺเต สํโฆ, อิตฺถนฺนโม ปุคฺคโล ธมฺมิกํ สามคฺคิํ ปจฺจาทิยติ, ตสฺส ปาติโมกฺขํ ฐเปมิ ‘น ตสฺมิํ สมฺมุขีภูเต ปาติโมกฺขํ อุทฺทิสิตพฺพนฺ’ติ”.\csromanspacer{} ธมฺมิกํ ปาติโมกฺขฏฺฐปนํ.}

\proseitem{394}{ภิกฺขุสฺส ปาติโมกฺเข ฐปิเต ปริสา วุฏฺฐาติ ทสนฺนํ อนฺตรายานํ อญฺญตเรน ราชนฺตราเยน วา ฯเปฯ พฺรหฺมจริยนฺตราเยน วา.\csromanspacer{} อากงฺขมาโน ภิกฺขเว ภิกฺขุ ตสฺมิํ วา อาวาเส อญฺญสฺมิํ วา อาวาเส ตสฺมิํ ปุคฺคเล สมฺมุขีภูเต สํฆมชฺเฌ อุทาหเรยฺย\csromandash{}}

\prose{“สุณาตุ เม ภนฺเต สํโฆ, อิตฺถนฺนามสฺส ปุคฺคลสฺส ธมฺมิกาย สามคฺคิยา ปจฺจาทานกถา วิปฺปกตา, ตํ วตฺถุ อวินิจฺฉิตํ, ยทิ สํฆสฺส ปตฺตกลฺลํ, สํโฆ ตํ วตฺถุํ วินิจฺฉิเนยฺยา”ติ.}

\prose{เอวญฺเจตํ ลเภถ, อิจฺเจตํ กุสลํ.\csromanspacer{} โน เจ ลเภถ, ตทหุโปสเถ จาตุทฺทเส วา ปนฺนรเส วา ตสฺมิํ ปุคฺคเล สมฺมุขีภูเต สํฆมชฺเฌ อุทาหริตพฺพํ\csromandash{}}

\prose{“สุณาตุ เม ภนฺเต สํโฆ, อิตฺถนฺนามสฺส ปุคฺคลสฺส ธมฺมิกาย สามคฺคิยา ปจฺจาทานกถา วิปฺปกตา, ตํ วตฺถุ อวินิจฺฉิตํ, ตสฺส ปาติโมกฺขํ ฐเปมิ ‘น ตสฺมิํ สมฺมุขีภูเต ปาติโมกฺขํ อุทฺทิสิตพฺพนฺ’ติ”.\csromanspacer{} ธมฺมิกํ ปาติโมกฺขฏฺฐปนํ.}

\proseitem{395}{กตํ สีลวิปตฺติยา ทิฏฺฐสุตปริสงฺกิโต โหติ.\csromanspacer{} อิธ ปน ภิกฺขเว ปน ภิกฺขุ เยหิ อากาเรหิ เยหิ ลิงฺเคหิ เยหิ นิมิตฺเตหิ สีลวิปตฺติยา ทิฏฺฐสุตปริสงฺกิโต โหติ, เตหิ อากาเรหิ เตหิ ลิงฺเคหิ เตหิ นิมิตฺเตหิ ภิกฺขุ ภิกฺขุํ ปสฺสติ สีลวิปตฺติยา ทิฏฺฐสุตปริสงฺกิตํ.\csromanspacer{} น เหว โข ภิกฺขุ ภิกฺขุํ ปสฺสติ สีลวิปตฺติยา ทิฏฺฐสุตปริสงฺกิตํ, อปิจ อญฺโญ ภิกฺขุ ภิกฺขุสฺส อาโรเจติ “อิตฺถนฺนาโม อาวุโส ภิกฺขุ สีลวิปตฺติยา ทิฏฺฐสุตปริสงฺกิโต”ติ.\csromanspacer{} น เหว โข ภิกฺขุ ภิกฺขุํ ปสฺสติ สีลวิปตฺติยา ทิฏฺฐสุตปริสงฺกิตํ, นาปิ อญฺโญ ภิกฺขุ ภิกฺขุสฺส อาโรเจติ “อิตฺถนฺนาโม อาวุโส \csromanfolio{433}ภิกฺขุ สีลวิปตฺติยา ทิฏฺฐสุตปริสงฺกิโต”ติ.\csromanspacer{} อปิ จ โสว ภิกฺขุ ภิกฺขุสฺส อาโรเจติ “อหํ อาวุโส สีลวิปตฺติยา ทิฏฺฐสุตปริสงฺกิโตมฺหี”ติ.\csromanspacer{} อากงฺขมาโน ภิกฺขเว ภิกฺขุ เตน ทิฏฺเฐน เตน สุเตน ตาย ปริสงฺกาย ตทหุโปสเถ จาตุทฺทเส วา ปนฺนรเส วา ตสฺมิํ ปุคฺคเล สมฺมุขีภูเต สํฆมชฺเฌ อุทาหเรยฺย\csromandash{}}

\prose{“สุณาตุ เม ภนฺเต สํโฆ, อิตฺถนฺนาโม ปุคฺคโล สีลวิปตฺติยา ทิฏฺฐสุตปริสงฺกิโต\footnote{ทิฏฺฐสุตปริสงฺกิโต โหติ (สฺยา, กํ)}, ตสฺส ปาติโมกฺขํ ฐเปมิ น ตสฺมิํ สมฺมุขีภูเต ปาติโมกฺขํ อุทฺทิสิตพฺพนฺ”ติ.\csromanspacer{} ธมฺมิกํ ปาติโมกฺขฏฺฐปนํ.}

\proseitem{396}{กถํ อาจารวิปตฺติยา ทิฏฺฐสุตปริสงฺกิโต โหติ.\csromanspacer{} อิธ ภิกฺขเว ภิกฺขุ เยหิ อากาเรหิ เยหิ ลิงฺเคหิ เยหิ นิมิตฺเตหิ อาจารวิปตฺติยา ทิฏฺฐสุตปริสงฺกิโต โหติ, เตหิ อากาเรหิ เตหิ ลิงฺเคหิ เตหิ นิมิตฺเตหิ ภิกฺขุ ภิกฺขุํ ปสฺสติ อาจารวิปตฺติยา ทิฏฺฐสุตปริสงฺกิตํ.\csromanspacer{} น เหว โข ภิกฺขุ ภิกฺขุํ ปสฺสติ อาจารวิปตฺติยา ทิฏฺฐสุตปริสงฺกิตํ, อปิ จ อญฺโญ ภิกฺขุ ภิกฺขุสฺส อาโรเจติ “อิตฺถนฺนาโม อาวุโส ภิกฺขุ อาจารวิปตฺติยา ทิฏฺฐสุตปริสงฺกิโต”ติ.\csromanspacer{} น เหว โข ภิกฺขุ ภิกฺขุํ ปสฺสติ อาจารวิปตฺติยา ทิฏฺฐสุตปริสงฺกิตํ, นาปิ อญฺโญ ภิกฺขุ ภิกฺขุสฺส อาโรเจติ “อิตฺถนฺนาโม อาวุโส ภิกฺขุ อาจารวิปตฺติยา ทิฏฺฐสุตปริสงฺกิโต”ติ.\csromanspacer{} อปิ จ โสว ภิกฺขุ ภิกฺขุสฺส อาโรเจติ “อหํ อาวุโส อาจารวิปตฺติยา ทิฏฺฐสุตปริสงฺกิโตมฺหี”ติ.\csromanspacer{} อากงฺขมาโน ภิกฺขเว ภิกฺขุ เตน ทิฏฺเฐน เตน สุเตน ตาย ปริสงฺกาย ตทหุโปสเถ จาตุทฺทเส วา ปนฺนรเส วา ตสฺมิํ ปุคฺคเล สมฺมุขีภูเต สํฆมชฺเฌ อุทาหเรยฺย\csromandash{}}

\prose{“สุณาตุ เม ภนฺเต สํโฆ, อิตฺถนฺนาโม ปุคฺคโล อาจารวิปตฺติยา ทิฏฺฐสุตปริสงฺกิโต, ตสฺส ปาติโมกฺขํ ฐเปมิ ‘น ตสฺมิํ สมฺมุขีภูเต ปาติโมกฺขํ อุทฺทิสิตพฺพนฺ’ติ”.\csromanspacer{} ธมฺมิกํ ปาติโมกฺขฏฺฐปนํ.}

\proseitem{397}{กถํ ทิฏฺฐิวิปตฺติยา ทิฏฺฐสุตปริสงฺกิโต โหติ.\csromanspacer{} อิธ ปน ภิกฺขเว ภิกฺขุ เยหิ อากาเรหิ เยหิ ลิงฺเคหิ เยหิ นิมิตฺเตหิ \csromanfolio{434}ทิฏฺฐิวิปตฺติยา ทิฏฺฐสุตปริสงฺกิโต โหติ, เตหิ อากาเรหิ เตหิ ลิงฺเคหิ เตหิ นิมิตฺเตหิ ภิกฺขุ ภิกฺขุํ ปสฺสติ ทิฏฺฐิวิปตฺติยา ทิฏฺฐสุตปริสงฺกิตํ.\csromanspacer{} น เหว โข ภิกฺขุ ภิกฺขุํ ปสฺสติ ทิฏฺฐิวิปตฺติยา ทิฏฺฐสุตปริสงฺกิตํ, อปิ จ อญฺโญ ภิกฺขุ ภิกฺขุสฺส อาโรเจติ “อิตฺถนฺนาโม อาวุโส ภิกฺขุ ทิฏฺฐิวิปตฺติยา ทิฏฺฐสุตปริสงฺกิโต”ติ.\csromanspacer{} น เหว โข ภิกฺขุ ภิกฺขุํ ปสฺสติ ทิฏฺฐิวิปตฺติยา ทิฏฺฐสุตปริสงฺกิตํ, นาปิ อญฺโญ ภิกฺขุ ภิกฺขุสฺส อาโรเจติ “อิตฺถนฺนาโม อาวุโส ภิกฺขุ ทิฏฺฐิวิปตฺติยา, ทิฏฺฐสุตปริสงฺกิโต”ติ, อปิ จ โสว ภิกฺขุ ภิกฺขุสฺส อาโรเจติ “อหํ อาวุโส ทิฏฺฐิวิปตฺติยา ทิฏฺฐสุตปริสงฺกิโตมฺหี”ติ.\csromanspacer{} อากงฺขมาโน ภิกฺขเว ภิกฺขุ เตน ทิฏฺเฐน เตน สุเตน ตาย ปริสงฺกาย ตทหุโปสเถ จาตุทฺทเส วา ปนฺนรเส วา ตสฺมิํ ปุคฺคเล สมฺมุขีภูเต สํฆมชฺเฌ อุทาหเรยฺย\csromandash{}}

\prose{“สุณาตุ เม ภนฺเต สํโฆ, อิตฺถนฺนาโม ปุคฺคโล ทิฏฺฐิวิปตฺติยา ทิฏฺฐสุตปริสงฺกิโต, ตสฺส ปาติโมกฺขํ ฐเปมิ ‘น ตสฺมิํ สมฺมุขีภูเต ปาติโมกฺขํ อุทฺทิสิตพฺพนฺ’ติ”.\csromanspacer{} ธมฺมิกํ ปาติโมกฺขฏฺฐปนํ.}

\vspace{\tipitakaparskip}
\csromancenter{อิมานิ ทส ธมฺมิกานิ ปาติโมกฺขฏฺฐปนานีติ.}
\nitthitam{ปฐมภาณวาโร นิฏฺฐิโต.}
\csromanmatikaanchor{mtk.199}
\csromantocmarkref{section}{7. อตฺตาทาน-องฺค}{mtk.199}
\bookmark[dest={mtk.199},level=1]{7. อตฺตาทาน-องฺค}
\csromanheader{1}{7. อตฺตาทาน-องฺค}
\proseitem{398}{อถ โข อายสฺมา อุปาลิ เยน ภควา เตนุปสงฺกมิ, อุปสงฺกมิตฺวา ภควนฺตํ อภิวาเทตฺวา เอกมนฺตํ นิสีทิ, เอกมนฺตํ เอกมนฺตํ นิสินฺโน โข อายสฺมา อุปาลิ ภควนฺตํ เอตทโวจ “อตฺตาทานํ อาทาตุกาเมน ภนฺเต ภิกฺขุนา กตมงฺคสมนฺนาคตํ\footnote{กตงฺคสมนฺนาคตํ (ก)} อตฺตาทานํ อาทาตพฺพนฺ”ติ.}

\prose{\csromansymbolfootnote{*}{วิ 5. 331 ปิฏฺเฐปิ.}อตฺตาทานํ อาทาตุกาเมน อุปาลิ ภิกฺขุนา ปญฺจงฺคสมนฺนาคตํ อตฺตาทานํ อาทาตพฺพํ, อตฺตาทานํ อาทาตุกาเมน อุปาลิ ภิกฺขุนา เอวํ ปจฺจเวกฺขิตพฺพํ “ยํ โข อหํ อิมํ อตฺตาทานํ อาทาตุกาโม, กาโล นุ โข อิมํ \csromanfolio{435}อตฺตาทานํ อาทาตุํ, อุทาหุ โน”ติ.\csromanspacer{} สเจ อุปาลิ ภิกฺขุ ปจฺจเวกฺขมาโน เอวํ ชานาติ “อกาโล อิมํ อตฺตาทานํ อาทาตุํ, โน กาโล”ติ.\csromanspacer{} น ตํ อุปาลิ อตฺตาทานํ อาทาตพฺพํ.}

\prose{สเจ ปนุปาลิ ภิกฺขุ ปจฺจเวกฺขมาโน เอวํ ชานาติ “กาโล อิมํ อตฺตาทานํ อาทาตุํ, โน อกาโล”ติ.\csromanspacer{} เตนุปาลิ ภิกฺขุนา อุตฺตริ ปจฺจเวกฺขิตพฺพํ “ยํ โข อหํ อิมํ อตฺตาทานํ อาทาตุกาโม, ภูตํ นุ โข อิทํ อตฺตาทานํ, อุทาหุ โน”ติ.\csromanspacer{} สเจ อุปาลิ ภิกฺขุ ปจฺจเวกฺขมาโน เอวํ ชานาติ “อภูตํ อิทํ อตฺตาทานํ, โน ภูตนฺ”ติ, น ตํ อุปาลิ อตฺตาทานํ อาทาตพฺพํ.}

\prose{สเจ ปนุปาลิ ภิกฺขุ ปจฺจเวกฺขมาโน เอวํ ชานาติ “ภูตํ อิทํ อตฺตาทานํ, โน อภูตนฺ”ติ, เตนุปาลิ ภิกฺขุนา อุตฺตริ ปจฺจเวกฺขิตพฺพํ “ยํ โข อหํ อิมํ อตฺตาทานํ อาทาตุกาโม, อตฺถสญฺหิตํ นุ โข อิทํ อตฺตาทานํ, อุทาหุ โน”ติ.\csromanspacer{} สเจ อุปาลิ ภิกฺขุ ปจฺจเวกฺขมาโน เอวํ ชานาติ “อนตฺถสญฺหิตํ อิทํ อตฺตาทานํ, โน อตฺถสญฺหิตนฺ”ติ, น ตํ อุปาลิ อตฺตาทานํ อาทาตพฺพํ.}

\prose{สเจ ปนุปาลิ ภิกฺขุ ปจฺจเวกฺขมาโน เอวํ ชานาติ “อตฺถสญฺหิตํ อิทํ อตฺตาทานํ, โน อนตฺถสญฺหิตนฺ”ติ, เตนุปาลิ ภิกฺขุนา อุตฺตริ ปจฺจเวกฺขิตพฺพํ “อิมํ โข อหํ อตฺตาทานํ อาทิยมาโน ลภิสฺสามิ สนฺทิฏฺเฐ สมฺภตฺเต ภิกฺขู ธมฺมโต วินยโต ปกฺเข, อุทาหุ โน”ติ.\csromanspacer{} สเจ อุปาลิ ภิกฺขุ ปจฺจเวกฺขมาโน เอวํ ชานาติ “อิมํ โข อหํ อตฺตาทานํ อาทิยมาโน น ลภิสฺสามิ สนฺทิฏฺเฐ สมฺเภตฺเต ภิกฺขู ธมฺมโต วินยโต ปกฺเข”ติ, น ตํ อุปาลิ อตฺตาทานํ อาทาตพฺพํ.}

\prose{สเจ ปนุปาลิ ภิกฺขุ ปจฺจเวกฺขมาโน เอวํ ชานาติ “อิมํ โข อหํ อตฺตาทานํ อาทิยมาโน ลภิสฺสามิ สนฺทิฏฺเฐ สมฺภตฺเต ภิกฺขู ธมฺมโต วินยโต ปกฺเข”ติ, เตนุปาลิ ภิกฺขุนา อุตฺตริ ปจฺจเวกฺขิตพฺพํ “อิมํ โข เม อตฺตาทานํ อาทิยโต ภวิสฺสติ สํฆสฺส ตโตนิทานํ ภณฺฑนํ กลโห วิคฺคโห วิวาโท สํฆเภโท สํฆราชิ สํฆววตฺถานํ สํฆนานากรณํ, อุทาหุ โน”ติ.\csromanspacer{} สเจ อุปาลิ ภิกฺขุ ปจฺจเวกฺขมาโน เอวํ ชานาติ “อิมํ โข เม อตฺตาทานํ อาทิยโต ภวิสฺสติ สํฆสฺส ตโตนิทานํ ภณฺฑํ กลโห วิคฺคโห วิวาโท สํฆเภโท สํฆราชิ \csromanfolio{436}สํฆววตฺถานํ สํฆนานากรณนฺ”ติ, น ตํ อุปาลิ อตฺตาทานํ อาทาตพฺพํ.\csromanspacer{} สเจ ปนุปาลิ ภิกฺขุ ปจฺจเวกฺขมาโน เอวํ ชานาติ “อิมํ โข เม อตฺตาทานํ อาทิยโต น ภวิสฺสติ สํฆสฺส ตโตนิทานํ ภณฺฑนํ กลโห วิคฺคโห วิวาโท สํฆเภโท สํฆราชิ สํฆววตฺถานํ สํฆนานากรณนฺ”ติ, อาทาตพฺพํ ตํ อุปาลิ อตฺตาทานํ.\csromanspacer{} เอวํ ปญฺจงฺคสมนฺนาคตํ โข อุปาลิ อตฺตาทานํ อาทินฺนํ ปจฺฉาปิ อวิปฺปฏิสารการกรํ ภวิสฺสตีติ.}

\csromanmatikaanchor{mtk.200}
\csromantocmarkref{section}{8. โจทเกนปจฺจเวกฺขิตพฺพธมฺม}{mtk.200}
\bookmark[dest={mtk.200},level=1]{8. โจทเกนปจฺจเวกฺขิตพฺพธมฺม}
\csromanheader{1}{8. โจทเกนปจฺจเวกฺขิตพฺพธมฺม}
\proseitem{399}{\csromansymbolfootnote{*}{วิ 5. 328 ปิฏฺเฐปิ. (อํ 3. 318 ปิฏฺเฐ ปน โถกํ วิสทิสํ)}โจทเกน ภนฺเต ภิกฺขุนา ปรํ โจเทตุกาเมน กติ ธมฺเม อชฺฌตฺตํ ปจฺจเวกฺขิตฺวา ปโร โจเทตพฺโพติ.\csromanspacer{} โจทเกน อุปาลิ ภิกฺขุนา ปรํ โจเทตุกาเมน ปญฺจ ธมฺเม อชฺฌตฺตํ ปจฺจเวกฺขิตฺวา ปโร โจเทตพฺโพ.}

\prose{โจทเกน อุปาลิ ภิกฺขุนา ปรํ โจเทตุกาเมน เอวํ ปจฺจเวกฺขิตพฺพํ “ปริสุทฺธกายสมาจาโร นุ โขมฺหิ, ปริสุทฺเธนมฺหิ กายสมาจาเรน สมนฺนาคโต อจฺฉิทฺเทน อปฺปฏิมํเสน, สํวิชฺชติ นุ โข เม เอโส ธมฺโม, อุทาหุ โน”ติ.\csromanspacer{} โน เจ อุปาลิ ภิกฺขุ ปริสุทฺธกายสมาจาโร โหติ ปริสุทฺเธน กายสมาจาเรน สมนฺนาคโต อจฺฉิทฺเทน อปฺปฏิมํเสน, ตสฺส ภวนฺติ วตฺตาโร “อิงฺฆ ตาว อายสฺมา กายิกํ สิกฺขสฺสู”ติ, อิติสฺส ภวนฺติ วตฺตาโร.}

\prose{ปุน จปรํ อุปาลิ โจทเกน ภิกฺขุนา ปรํ โจเทตุกาเมน เอวํ ปจฺจเวกฺขิตพฺพํ “ปริสุทฺธวจีสมาจาโร นุ โขมฺหิ, ปริสุทฺเธนมฺหิ วจีสมาจาเรน สมนฺนาคโต อจฺฉิทฺเทน อปฺปฏิมํเสน, สํวิชฺชติ นุ โข เม เอโส ธมฺโม, อุทาหุ โน”ติ.\csromanspacer{} โน เจ อุปาลิ ภิกฺขุ ปริสุทฺธวจีสมาจาโร โหติ ปริสุทฺเธน วจีสมาจาเรน สมนฺนาคโต อจฺฉิทฺเทน อปฺปฏิมํเสน, ตสฺส ภวนฺติ วตฺตาโร “อิงฺฆ ตาว อายสฺมา วาจสิกํ สิกฺขสฺสู”ติ, อิติสฺส ภวนฺติ วตฺตาโร.}

\prose{ปุน จปรํ อุปาลิ โจทเกน ภิกฺขุนา ปรํ โจเทตุกาเมน เอวํ ปจฺจเวกฺขิตพฺพํ “เมตฺตํ นุ โข เม จิตฺตํ ปจฺจุปฏฺฐิตํ สพฺรหฺมจารีสุ อนาฆาตํ, สํวิชฺชติ นุ โข เม เอโส ธมฺโม, อุทาหุ โน”ติ.\csromanspacer{} โน เจ อุปาลิ ภิกฺขุโน เมตฺตจิตฺตํ ปจฺจุปฏฺฐิตํ โหติ สพฺรหฺมจารีสุ อนาฆาตํ, \csromanfolio{437}ตสฺส ภวนฺติ วตฺตาโร “อิงฺฆ ตาว อายสฺมา สพฺรหฺมจารีสุ เมตฺตจิตฺตํ อุปฏฺฐาเปหี”ติ, อิติสฺส ภวนฺติ วตฺตาโร.}

\prose{ปุน จปรํ อุปาลิ โจทเกน ภิกฺขุนา ปรํ โจเทตุกาเมน เอวํ ปจฺจเวกฺขิตพฺพํ “พหุสฺสุโต นุ โขมฺหิ สุตธโร สุตสนฺนิจฺจโย, เย เต ธมฺมา อาทิกลฺยาณา มชฺเฌกลฺยาณา ปริโยสานกลฺยาณา สาตฺถํ สพฺยญฺชนํ เกวลปริปุณฺณํ ปริสุทฺธํ พฺรหฺมจริยํ อภิวทนฺติ, ตถารูปา เม ธมฺมา พหุสฺสุตา โหนฺติ ธาตา วจสา ปริจิตา มนสานุเปกฺขิตา ทิฏฺฐิยา สุปฺปฏิวิทฺธา, สํวิชฺชติ นุ โข เม เอโส ธมฺโม, อุทาหุ โน”ติ.\csromanspacer{} โน เจ อุปาลิ ภิกฺขุ พหุสฺสุโต โหติ สุตธโร สุตสนฺนิจฺจโย, เย เต ธมฺมา อาทิกลฺยาณา มชฺเฌกลฺยาณา ปริโยสานกลฺยาณา สาตฺถํ สพฺยญฺชนํ เกวลปริปุณฺณํ ปริสุทฺธํ พฺรหฺมจริยํ อภิวทนฺติ, ตถารูปสฺส ธมฺมา น พหุสฺสุตา โหนฺติ ธาตา วจสา ปริจิตา มนสานุเปกฺขิตา ทิฏฺฐิยา สุปฺปฏิวิทฺธา, ตสฺส ภวนฺติ วตฺตาโร “อิงฺฆ ตาว อายสฺมา อาคมํ ปริยาปุณสฺสู”ติ, อิติสฺส ภวนฺติ วตฺตาโร.}

\prose{ปุน จปรํ อุปาลิ โจทเกน ภิกฺขุนา ปรํ โจเทตุกาเมน เอวํ ปจฺจเวกฺขิตพฺพํ “อุภยานิ โข เม ปาติโมกฺขานิ วิตฺถาเรน สฺวาคตานิ โหนฺติ สุวิภตฺตานิ สุปฺปวตฺตีนิ สุวินิจฺฉิตานิ สุตฺตโส อนุพฺยญฺชนโส.\csromanspacer{} สํวิชฺชติ นุ โข เม เอโส ธมฺโม, อุทาหุ โน”ติ.\csromanspacer{} โน เจ อุปาลิ ภิกฺขุโน อุภยานิ ปาติโมกฺขานิ วิตฺถาเรน สฺวาคตานิ โหนฺติ สุวิภตฺตานิ สุปฺปวตฺตีนิ สุวินิจฺฉิตานิ สุตฺตโส อนุพฺยญฺชนโส, “อิทํ ปนาวุโส กตฺถ วุตฺตํ ภควตา”ติ อิติ ปุฏฺโฐ น สมฺปายติ\footnote{น สมฺปาทยติ (ก)}, ตสฺส ภวนฺติ วตฺตาโร “อิงฺฆ ตาว อายสฺมา วินยํ ปริยาปุณสฺสู”ติ, อิติสฺส ภวนฺติ วตฺตาโร.\csromanspacer{} โจทเกนุปาลิ ภิกฺขุนา ปรํ โจเทตุกาเมน อิเม ปญฺจ ธมฺเม อชฺฌตฺตํ ปจฺจเวกฺขิตฺวา ปโร โจเทตพฺโพติ.}

\csromanmatikaanchor{mtk.201}
\csromantocmarkref{section}{9. โจทเกน-อุปฏฺฐาเปตพฺพธมฺม}{mtk.201}
\bookmark[dest={mtk.201},level=1]{9. โจทเกน-อุปฏฺฐาเปตพฺพธมฺม}
\csromanheader{1}{9. โจทเกน-อุปฏฺฐาเปตพฺพธมฺม}
\proseitem{400}{\csromansymbolfootnote{*}{วิ 5. 330 ปิฏฺเฐปิ.}โจทเกน ภนฺเต ภิกฺขุนา ปรํ โจเทตุกาเมน กติ ธมฺเม อชฺฌตฺตํ อุปฏฺฐาเปตฺวา ปโร โจเทตพฺโพติ.\csromanspacer{} โจทเกนุปาลิ ภิกฺขุนา ปรํ โจเทตุกาเมน ปญฺจ ธมฺเม อชฺฌตฺตํ อุปฏฺฐาเปตฺวา ปโร โจเทตพฺโพ “กาเลน วกฺขามิ, โน อกาเลน, ภูเตน วกฺขามิ, โน อภูเตน, \csromanfolio{438}สเณฺหน วกฺขามิ, โน ผรุเสน, อตฺถสํหิเตน วกฺขามิ, โน อนตฺถสํหิเตน, เมตฺตจิตฺโต วกฺขามิ, โน โทสนฺตโร”ติ.\csromanspacer{} โจทเกนุปาลิ ภิกฺขุนา ปรํ โจเทตุกาเมน อิเม ปญฺจ ธมฺเม อชฺฌตฺตํ อุปฏฺฐาเปตฺวา ปโร โจเทตพฺโพติ.}

\csromanmatikaanchor{mtk.202}
\csromantocmarkref{section}{10. โจทกจุทิตกปฏิสํยุตฺตกถา}{mtk.202}
\bookmark[dest={mtk.202},level=1]{10. โจทกจุทิตกปฏิสํยุตฺตกถา}
\csromanheader{1}{10. โจทกจุทิตกปฏิสํยุตฺตกถา}
\proseitem{401}{อธมฺมโจทกสฺส ภนฺเต ภิกฺขุโน กติหากาเรหิ วิปฺปฏิสาโร อุปทหาตพฺโพติ.\csromanspacer{} อธมฺมโจทกสฺส อุปาลิ ภิกฺขุโน ปญฺจหากาเรหิ วิปฺปฏิสาโร อุปทหาตพฺโพ “อกาเลนายสฺมา โจเทสิ, โน กาเลน, อลํ เต วิปฺปฏิสาราย.\csromanspacer{} อภูเตนายสฺมา โจเทสิ, โน ภูเตน, อลํ เต วิปฺปฏิสาราย.\csromanspacer{} ผรุเสนายสฺมา โจเทสิ, โน สเณฺหน, อลํ เต วิปฺปฏิสาราย.\csromanspacer{} อนตฺถสํหิเตนายสฺมา โจเทสิ, โน อตฺถสํหิเตน, อลํ เต วิปฺปฏิสาราย.\csromanspacer{} โทสนฺตโร อายสฺมา โจเทสิ, โน เมตฺตจิตฺโต, อลํ เต วิปฺปฏิสารายา”ติ.\csromanspacer{} อธมฺมโจทกสฺส อุปาลิ ภิกฺขุโน อิเมหิ ปญฺจหากาเรหิ วิปฺปฏิสาโร อุปทหาตพฺโพ, ตํ กิสฺส เหตุ, ยถา น อญฺโญปิ ภิกฺขุ อภูเตน โจเทตพฺพํ มญฺเญยฺยาติ.}

\prose{อธมฺมจุทิตสฺส ปน ภนฺเต ภิกฺขุโน กติหากาเรหิ อวิปฺปฏิสาโร อุปทหาตพฺโพติ.\csromanspacer{} อธมฺมจุทิตสฺส อุปาลิ ภิกฺขุโน ปญฺจหากาเรหิ อวิปฺปฏิสาโร อุปทหาตพฺโพ “อกาเลนายสฺมา จุทิโต, โน กาเลน, อลํ เต อวิปฺปฏิสาราย.\csromanspacer{} อภูเตนายสฺมา จุทิโต, โน ภูเตน, อลํ เต อวิปฺปฏิสาราย.\csromanspacer{} ผรุเสนายสฺมา จุทิโต, โน สเณฺหน, อลํ เต อวิปฺปฏิสาราย.\csromanspacer{} อนตฺถสํหิเตนายสฺมา จุทิโต, โน อตฺถสํหิเตน, อลํ เต อวิปฺปฏิสาราย.\csromanspacer{} โทสนฺตเรนายสฺมา จุทิโต, โน เมตฺตจิตฺเตน, อลํ เต อวิปฺปฏิสารายา”ติ.\csromanspacer{} อธมฺมจุทิตส อุปาลิ ภิกฺขุโน อิเมหิ ปญฺจหากาเรหิ อวิปฺปฏิสาโร อุปทหาตพฺโพติ.}

\prose{ธมฺมโจทกสฺส ภนฺเต ภิกฺขุโน กติหากาเรหิ อวิปฺปฏิสาโร อุปทหาตพฺโพติ.\csromanspacer{} ธมฺมโจทกสฺส อุปาลิ ภิกฺขุโน ปญฺจหากาเรหิ อวิปฺปฏิสาโร อุปทหาตพฺโพ “กาเลนายสฺมา โจเทสิ, โน อกาเลน, อลํ เต อวิปฺปฏิสาราย.\csromanspacer{} ภูเตนายสฺมา โจเทสิ, โน \csromanfolio{439}อภูเตน, อลํ เต อวิปฺปฏิสาราย.\csromanspacer{} สเณฺหนายสฺมา โจเทสิ, โน ผรุเสน, อลํ เต อวิปฺปฏิสาราย.\csromanspacer{} อตฺถสํหิเตนายสฺมา โจเทสิ, โน อนตฺถสํหิเตน, อลํ เต อวิปฺปฏิสาราย.\csromanspacer{} เมตฺตจิตฺโต อายสฺมา โจเทสิ, โน โทสนฺตโร, อลํ เต อวิปฺปฏิสารายา”ติ.\csromanspacer{} ธมฺมโจทกสฺส อุปาลิ ภิกฺขุโน อิเมหิ ปญฺจหากาเรหิ อวิปฺปฏิสาโร อุปทหาตพฺโพ.\csromanspacer{} ตํ กิสฺส เหตุ, ยถา อญฺโญปิ ภิกฺขุ ภูเตน โจเทตพฺพํ มญฺเญยฺยาติ.}

\prose{ธมฺมจุทิตสฺส ปน ภนฺเต ภิกฺขุโน กติหากาเรหิ วิปฺปฏิสาโร อุปทหาตพฺโพติ.\csromanspacer{} ธมฺมจุทิตสฺส อุปาลิ ภิกฺขุโน ปญฺจหากาเรหิ วิปฺปฏิสาโร อุปทหาตพฺโพ “กาเลนายสฺมา จุทิโต, โน อกาเลน, อลํ เต วิปฺปฏิสาราย.\csromanspacer{} ภูเตนายสฺมา จุทิโต, โน อภูเตน, อลํ เต วิปฺปฏิสาราย.\csromanspacer{} สเณฺหนายสฺมา จุทิโต, โน ผรุเสน, อลํ เต วิปฺปฏิสาราย.\csromanspacer{} อตฺถสํหิเตนายสฺมา จุทิโต, โน อนตฺถสํหิเตน, อลํ เต วิปฺปฏิสาราย.\csromanspacer{} เมตฺตจิตฺเตนายสฺมา จุทิโต, โน โทสนฺตเรน, อลํ เต วิปฺปฏิสารายา”ติ.\csromanspacer{} ธมฺมจุทิตสฺส อุปาลิ ภิกฺขุโน อิเมหิ ปญฺจหากาเรหิ วิปฺปฏิสาโร อุปทหาตพฺโพติ.}

\prose{โจทเกน ภนฺเต ภิกฺขุนา ปรํ โจเทตุกาเมน กติ ธมฺเม อชฺฌตฺตํ มนสิ กริตฺวา ปโร โจเทตพฺโพติ.\csromanspacer{} โจทเกนุปาลิ ภิกฺขุนา ปรํ โจเทตุกาเมน ปญฺจ ธมฺเม อชฺฌตฺตํ มนสิ กริตฺวา ปโร โจเทตพฺโพ.\csromanspacer{} การุญฺญตา หิเตสิตา อนุกมฺปิตา\footnote{อนุกมฺปตา (ก)} อาปตฺติวุฏฺฐานตา วินยปุเรกฺขารตาติ.\csromanspacer{} โจทเกนุปาลิ ภิกฺขุนา ปรํ โจเทตุกาเมน อิเม ปญฺจ ธมฺเม อชฺฌตฺตํ มนสิ กริตฺวา ปโร โจเทตพฺโพติ.}

\prose{จุทิเตน ปน ภนฺเต ภิกฺขุนา กติสุ ธมฺเมสุ ปติฏฺฐาตพฺพนฺติ.\csromanspacer{} จุทิเตนุปาลิ ภิกฺขุนา ทฺวีสุ ธมฺเมสุ ปติฏฺฐาตพฺพํ สจฺเจ จ อกุปฺเป จาติ.}

\nitthitam{ทุติยภาณวาโร นิฏฺฐิโต.}
\vspace{\tipitakaparskip}
\csromancenter{ปาติโมกฺขฏฺฐปนกฺขนฺธโก นวโม.}
\csromancenter{อิมมฺหิ ขนฺธเก วตฺถู ติํส.}
\titlehead{\textbf{ตสฺสุทฺทานํ}}
\csromanmatikaanchor{mtk.203}
\csromantocmarkref{section}{อุทฺทานคาถา}{mtk.203}
\bookmark[dest={mtk.203},level=1]{อุทฺทานคาถา}
\csromangathameasure{อุโปสเถ ยาวติกํ, ปาปภิกฺขุ น นิกฺขมิ. \\ โมคฺคลฺลาเนน นิจฺฉุทฺโธ, อจฺเฉรา ชินสาสเน. \\ นินฺโนนุปุพฺพสิกฺขา จ, ฐิตธมฺโม นาติกฺกมฺม. \\ กุณปุกฺขิปติ สํโฆ, สวนฺติโย ชหนฺติ จ. \\ สวนฺติ ปรินิพฺพนฺติ, เอกรส วิมุตฺติ จ. \\ พหุ ธมฺมวินโยปิ, ภูตฏฺฐาริยปุคฺคลา. \\ สมุทฺทํ อุปมํ กตฺวา, วาเจสิ สาสเน คุณํ. \\ อุโปสเถ ปาติโมกฺขํ, น อเมฺห โกจิ ชานาติ. \\ ปฏิกจฺเจว อุชฺฌนฺติ, เอโก เทฺว ตีณิ จตฺตาริ. \\ ปญฺจ ฉ สตฺต อฏฺฐานิ, นวา จ ทสมานิ จ. \\ สีล อาจาร ทิฏฺฐิ จ, อาชีวํ จตุภาคิเก. \\ ปาราชิกญฺจ สํฆาทิ, ปาจิตฺติ ปาฏิเทสนิ. \\ ทุกฺกฏํ ปญฺจภาเคสุ, สีลาจารวิปตฺติ จ. \\ อกตาย กตาย จ, ฉภาเคสุ ยถาวิธิ. \\ ปาราชิกญฺจ สํฆาทิ, ถุลฺลํ ปาจิตฺติเยน จ. \\ ปาฏิเทสนิยญฺเจว, ทุกฺกฏญฺจ ทุพฺภาสิตํ. \\ สีลาจารวิปตฺติ จ, ทิฏฺฐิ-อาชีววิปตฺติ. \\ ยา จ อฏฺฐา กตากเต, เตเนตา สิลาจารทิฏฺฐิยา. \\ อกตาย กตายาปิ, กตากตายเมว จ. \\ เอวํ นววิธา วุตฺตา, ยถาภูเตน ญายโต. \\ ปาราชิโก วิปฺปกตา, ปจฺจกฺขาโต ตเถว จ. \\ อุเปติ ปจฺจาทิยติ, ปจฺจาทานกถา จ ยา. \\ สีลาจารวิปตฺติ จ, ตถา ทิฏฺฐิวิปตฺติยา. \\ ทิฏฺฐสุตปริสงฺกิตํ, ทสธา ตํ วิชานาถ. \\ ภิกฺขุ วิปสฺสติ ภิกฺขุํ, อญฺโญ จาโรจยาติ ตํ. \\ โส เยว ตสฺส อกฺขาติ, ปาติโมกฺขํ ฐเปติ โส. \\ วุฏฺฐาติ อนฺตราเยน, ราชโจรคฺคุทกา จ. \\ มนุสฺส-อมนุสฺสา จ, วาฬสรีสปา ชีวิพฺรหฺมํ. \\ ทสนฺนมญฺญตเรน, ตสฺมิํ อญฺญตเรสุ วา. \\ ธมฺมิกาธมฺมิกา เจว, ยถา มคฺเคน ชานาถ. \\ กาลภูตตฺถสํหิตํ, ลภิสฺสามิ ภวิสฺสติ. \\ กายวาจสิกา เมตฺตา, พาหุสจฺจํ อุภยานิ. \\ กาลภูเตน สเณฺหน, อตฺถเมตฺเตน โจทเย. \\ วิปฺปฏิสารธมฺเมน, ตถา วาจา วิโนทเย. \\ ธมฺมโจทจุทิตสฺส, วิโนเทติ วิปฺปฏิสาโร. \\ กรุณา หิตานุกมฺปิ, วุฏฺฐานปุเรกฺขารโต. \\ โจทกสฺส ปฏิปตฺติ, สมฺพุทฺเธน ปกาสิตา. \\ สจฺเจ เจว อกุปฺเป จ, จุทิตสฺเสว ธมฺมตาติ.}
\csromangathagroup{\csromangathastanza{\csromanfolio{440}อุโปสเถ ยาวติกํ, ปาปภิกฺขุ น นิกฺขมิ. \\ โมคฺคลฺลาเนน นิจฺฉุทฺโธ, อจฺเฉรา ชินสาสเน.}\csromangathastanza{นินฺโนนุปุพฺพสิกฺขา จ, ฐิตธมฺโม นาติกฺกมฺม. \\ กุณปุกฺขิปติ สํโฆ, สวนฺติโย ชหนฺติ จ.}\csromangathastanza{สวนฺติ ปรินิพฺพนฺติ, เอกรส วิมุตฺติ จ. \\ พหุ ธมฺมวินโยปิ, ภูตฏฺฐาริยปุคฺคลา.}\csromangathastanza{สมุทฺทํ อุปมํ กตฺวา, วาเจสิ\footnote{ฐาเปสิ (ก)} สาสเน คุณํ. \\ อุโปสเถ ปาติโมกฺขํ, น อเมฺห โกจิ ชานาติ.}\csromangathastanza{ปฏิกจฺเจว อุชฺฌนฺติ, เอโก เทฺว ตีณิ จตฺตาริ. \\ ปญฺจ ฉ สตฺต อฏฺฐานิ, นวา จ ทสมานิ จ.}\csromangathastanza{สีล อาจาร ทิฏฺฐิ จ, อาชีวํ จตุภาคิเก. \\ ปาราชิกญฺจ สํฆาทิ, ปาจิตฺติ ปาฏิเทสนิ.}\csromangathastanza{ทุกฺกฏํ ปญฺจภาเคสุ, สีลาจารวิปตฺติ จ. \\ อกตาย กตาย จ, ฉภาเคสุ ยถาวิธิ.}\csromangathastanza{ปาราชิกญฺจ สํฆาทิ, ถุลฺลํ ปาจิตฺติเยน จ. \\ ปาฏิเทสนิยญฺเจว, ทุกฺกฏญฺจ ทุพฺภาสิตํ.}\csromangathastanza{สีลาจารวิปตฺติ จ, ทิฏฺฐิ-อาชีววิปตฺติ. \\ ยา จ อฏฺฐา กตากเต, เตเนตา สิลาจารทิฏฺฐิยา.}\csromangathastanza{อกตาย กตายาปิ, กตากตายเมว จ. \\ เอวํ นววิธา วุตฺตา, ยถาภูเตน ญายโต.}\csromangathastanza{ปาราชิโก วิปฺปกตา, ปจฺจกฺขาโต ตเถว จ. \\ อุเปติ ปจฺจาทิยติ, ปจฺจาทานกถา จ ยา.}\csromangathastanza{สีลาจารวิปตฺติ จ, ตถา ทิฏฺฐิวิปตฺติยา. \\ ทิฏฺฐสุตปริสงฺกิตํ, ทสธา ตํ วิชานาถ.}\csromangathastanza{\csromanfolio{441}ภิกฺขุ วิปสฺสติ ภิกฺขุํ, อญฺโญ จาโรจยาติ ตํ. \\ โส เยว ตสฺส อกฺขาติ\footnote{วิปสฺสญฺโญ จาโรจติ. ตํ สุทฺเธว ตสฺส อกฺขาติ (ก)}, ปาติโมกฺขํ ฐเปติ โส.}\csromangathastanza{วุฏฺฐาติ อนฺตราเยน, ราชโจรคฺคุทกา จ. \\ มนุสฺส-อมนุสฺสา จ, วาฬสรีสปา ชีวิพฺรหฺมํ.}\csromangathastanza{ทสนฺนมญฺญตเรน, ตสฺมิํ อญฺญตเรสุ วา. \\ ธมฺมิกาธมฺมิกา เจว, ยถา มคฺเคน ชานาถ.}\csromangathastanza{กาลภูตตฺถสํหิตํ, ลภิสฺสามิ ภวิสฺสติ. \\ กายวาจสิกา เมตฺตา, พาหุสจฺจํ อุภยานิ.}\csromangathastanza{กาลภูเตน สเณฺหน, อตฺถเมตฺเตน โจทเย. \\ วิปฺปฏิสารธมฺเมน, ตถา วาจา\footnote{ตเถวาปิ (สฺยา)} วิโนทเย.}\csromangathastanza{ธมฺมโจทจุทิตสฺส, วิโนเทติ วิปฺปฏิสาโร. \\ กรุณา หิตานุกมฺปิ, วุฏฺฐานปุเรกฺขารโต.}\csromangathastanza{โจทกสฺส ปฏิปตฺติ, สมฺพุทฺเธน ปกาสิตา. \\ สจฺเจ เจว อกุปฺเป จ, จุทิตสฺเสว ธมฺมตาติ.}}

\nitthitam{\textbf{ปาติโมกฺขฏฺฐปนกฺขนฺธโก นิฏฺฐิโต.}}
\csromanmatikaanchor{mtk.204}
\csromantocmarknopage{chapter}{10. ภิกฺขุนิกฺขนฺธก}{mtk.204}
\bookmark[dest={mtk.204},level=0]{10. ภิกฺขุนิกฺขนฺธก}
\chapterheadpageread{10. ภิกฺขุนิกฺขนฺธก}
\csromanmatikaanchor{mtk.205}
\csromantocmarkref{section}{1. ปฐมภาณวาร}{mtk.205}
\bookmark[dest={mtk.205},level=1]{1. ปฐมภาณวาร}
\csromanheader{1}{1. ปฐมภาณวาร}
\csromanheader{2}{\textbf{มหาปชาปติ โคตมีวตฺถุ}}
\csromanmatikaanchor{mtk.206}
\csromantocmarkref{subsection}{มหาปชาปติโคตมีวตฺถุ}{mtk.206}
\bookmark[dest={mtk.206},level=2]{มหาปชาปติโคตมีวตฺถุ}
\proseitem{402}{\csromanfolio{442}\csromansymbolfootnote{*}{อํ 3. 101 ปิฏฺฐาทีสุปิ อิทํ วตฺถุ อาคตํ.}เตน สมเยน พุทฺโธ ภควา สกฺเกสุ วิหรติ กปิลวตฺถุสฺมิํ นิโคฺรธาราเม.\csromanspacer{} อถ โข มหาปชาปติ1โคตมี เยน ภควา เตนุปสงฺกมิ, อุปสงฺกมิตฺวา ภควนฺตํ อภิวาเทตฺวา เอกมนฺตํ อฏฺฐาสิ, เอกมนฺตํ ฐิตา โข มหาปชาปติ โคตมี ภควนฺตํ เอตทโวจ “สาธุ ภนฺเต, ลเภยฺย มาตุคาโม ตถาคตปฺปเวทิเต ธมฺมวินเย อคารสฺมา อนคาริยํ ปพฺพชฺชนฺ”ติ.\csromanspacer{} อลํ โคตมิ, มา เต รุจฺจิ มาตุคามสฺส ตถาคตปฺปเวทิเต ธมฺมวินเย อคารสฺมา อนคาริยํ ปพฺพชฺชาติ.\csromanspacer{} ทุติยมฺปิ โข ฯเปฯ.\csromanspacer{} ตติยมฺปิ โข มหาปชาปติ โคตมี ภควนฺตํ เอตทโวจ “สาธุ ภนฺเต, ลเภยฺย มาตุคาโม ตถาคตปฺปเวทิเต ธมฺมวินเย อคารสฺมา อนคาริยํ ปพฺพชฺชนฺ”ติ.\csromanspacer{} อลํ โคตมิ, มา เต รุจฺจิ มาตุคามสฺส ตถาคตปฺปเวทิเต ธมฺมวินเย อคารสฺมา อนคาริยํ ปพฺพชฺชาติ.\csromanspacer{} อถ โข มหาปชาปติ โคตมี “น ภควา อนุชานาติ มาตุคามสฺส ตถาคตปฺปเวทิเต ธมฺมวินเย อคารสฺมา อนคาริยํ ปพฺพชฺชนฺ”ติ ทุกฺขี ทุมฺมนา อสฺสุมุขี รุทมานา ภควนฺตํ อภิวาเทตฺวา ปทกฺขิณํ กตฺวา ปกฺกามิ.}

\prose{อถ โข ภควา กปิลวตฺถุสฺมิํ ยถาภิรนฺตํ วิหริตฺวา เยน เวสาลี เตน จาริกํ ปกฺกามิ, อนุปุพฺเพน จาริกํ จรมาโน เยน เวสาลี ตทวสริ, ตตฺร สุทํ ภควา เวสาลิยํ วิหรติ มหาวเน กูฏาคารสาลายํ.\csromanspacer{} อถ โข มหาปชาปติ โคตมี เกเส เฉทาเปตฺวา กาสายานิ วตฺถานิ อจฺฉาเทตฺวา สมฺพหุลาหิ สากิยานีหิ สทฺธิํ เยน เวสาลี เตน ปกฺกามิ.\csromanspacer{} อนุปุพฺเพน เยน เวสาลี มหาวนํ กูฏาคารสาลา เตนุปสงฺกมิ.\csromanspacer{} อถ โข มหาปชาปติ โคตมี สูเนหิ ปาเทหิ รโชกิณฺเณน คตฺเตน ทุกฺขี ทุมฺมนา อสฺสุมุขี รุทมานา พหิทฺวารโกฏฺฐเก อฏฺฐาสิ, อทฺทสา โข อายสฺมา อานนฺโท มหาปชาปติํ โคตมิํ สูเนหิ ปาเทหิ รโชกิณฺเณน คตฺเตน ทุกฺขิํ ทุมฺมนํ อสฺสุมุขิํ รุทมานํ พหิทฺวารโกฏฺฐเก ฐิตํ, ทิสฺวาน มหาปชาปติํ \csromanfolio{443}โคตมิํ เอตทโวจ “กิสฺส ตฺวํ โคตมิ สูเนหิ ปาเทหิ รโชกิณฺเณน คตฺเตน ทุกฺขี ทุมฺมนา อสฺสุมุขี รุทมานา พหิทฺวารโกฏฺฐเก ฐิตา”ติ.\csromanspacer{} ตถา หิ ปน ภนฺเต อานนฺท น ภควา อนุชานาติ มาตุคามสฺส ตถาคตปฺปเวทิเต ธมฺมวินเย อคารสฺมา อนคาริยํ ปพฺพชฺชนฺติ.\csromanspacer{} เตน หิ ตฺวํ โคตมิ มุหุตฺตํ อิเธว ตาว โหติ, ยาวาหํ ภควนฺตํ ยาจามิ มาตุคามสฺส ตถาคตปฺปเวทิเต ธมฺมวินเย อคารสฺมา อนคาริยํ ปพฺพชฺชนฺติ.}

\prose{อถ โข อายสฺมา อานนฺโท เยน ภควา เตนุปสงฺกมิ, อุปสงฺกมิตฺวา ภควนฺตํ อภิวาเทตฺวา เอกมนฺตํ นิสีทิ, เอกมนฺตํ นิสินฺโน โข อายสฺมา อานนฺโท ภควนฺตํ เอตทโวจ “เอสา ภนฺเต มหาปชาปติ โคตมี สูเนหิ ปาเทหิ รโชกิณฺเณน คตฺเตน ทุกฺขี ทุมฺมนา อสฺสุมุขี รุทมานา พหิทฺวารโกฏฺฐเก ฐิตา ‘น ภควา อนุชานาติ มาตุคามสฺส ตถาคตปฺปเวทิเต ธมฺมวินเย อคารสฺมา อนคาริยํ ปพฺพชฺชนฺ’ติ, สาธุ ภนฺเต, ลเภยฺย มาตุคาโม ตถาคตปฺปเวทิเต ธมฺมวินเย อคารสฺมา อนคาริยํ ปพฺพชฺชนฺ”ติ.\csromanspacer{} อลํ อานนฺท, มา เต รุจฺจิ มาตุคามสฺส ตถาคตปฺปเวทิเต ธมฺมวินเย อคารสฺมา อนคาริยํ ปพฺพชฺชาติ.\csromanspacer{} ทุติยมฺปิ โข ฯเปฯ.\csromanspacer{} ตติยมฺปิ โข อายสฺมา อานนฺโท ภควนฺตํ เอตทโวจ “สาธุ ภนฺเต, ลเภยฺย มาตุคาโม ตถาคตปฺปเวทิเต ธมฺมวินเย อคารสฺมา อนคาริยํ ปพฺพชฺชนฺ”ติ.\csromanspacer{} อลํ อานนฺท, มา เต รุจฺจิ มาตุคามสฺส ตถาคตปฺปเวทิเต ธมฺมวินเย อคารสฺมา อนคาริยํ ปพฺพชฺชาติ.}

\prose{อถ โข อายสฺมา อานนฺโท\footnote{อํ 3. 103 ปิฏฺเฐ ปน “อถ โข อายสฺมโต อานนฺทสฺส เอตทโหสี”ติ อาคตํ.} “น ภควา อนุชานาติ มาตุคามสฺส ตถาคตปฺปเวทิเต ธมฺมวินเย อคารสฺมา อนคาริยํ ปพฺพชฺชํ, ยํนูนาหํ อญฺเญนปิ ปริยาเยน ภควนฺตํ ยาเจยฺยํ มาตุคามสฺส ตถาคตปฺปเวทิเต ธมฺมวินเย อคารสฺมา อนคาริยํ ปพฺพชฺชนฺ”ติ.\csromanspacer{} อถ โข อายสฺมา อานนฺโท ภควนฺตํ เอตทโวจ “ภพฺโพ นุ โข ภนฺเต มาตุคาโม ตถาคตปฺปเวทิเต ธมฺมวินเย อคารสฺมา อนคาริยํ ปพฺพชิตฺวา โสตาปตฺติผลํ วา สกทาคามิผลํ\footnote{สกิทาคามิผลํ (สฺยา)} วา อนาคามิผลํ วา อรหตฺตผลํ วา สจฺฉิกาตุนฺ”ติ.\csromanspacer{} ภพฺโพ อานนฺท มาตุคาโม ตถาคตปฺปเวทิเต ธมฺมวินเย อคารสฺมา อนคาริยํ ปพฺพชิตฺวา โสตาปตฺติผลมฺปิ สกทาคามิผลมฺปิ \csromanfolio{444}อนาคามิผลมฺปิ อรหตฺตผลมฺปิ สจฺฉิกาตุนฺติ.\csromanspacer{} สเจ ภนฺเต ภพฺโพ มาตุคาโม ตถาคตปฺปเวทิเต ธมฺมวินเย อคารสฺมา อนคาริยํ ปพฺพชิตฺวา โสตาปตฺติผลมฺปิ สกทาคามิผลมฺปิ อนาคามิผลมฺปิ อรหตฺตผลมฺปิ สจฺฉิกาตุํ, พหูปการา ภนฺเต มหาปชาปติ โคตมี ภควโต มาตุจฺฉา อาปาทิกา โปสิกา ขีรสฺส ทายิกา ภควนฺตํ ชเนตฺติยา กาลงฺกตาย ถญฺญํ ปาเยสิ, สาธุ ภนฺเต, ลเภยฺย มาตุคาโม ตถาคตปฺปเวทิเต ธมฺมวินเย อคารสฺมา อนคาริยํ ปพฺพชฺชนฺติ.}

\csromanmatikaanchor{mtk.207}
\csromantocmarkref{subsection}{อฏฺฐครุธมฺม}{mtk.207}
\bookmark[dest={mtk.207},level=2]{อฏฺฐครุธมฺม}
\csromanheader{2}{\textbf{อฏฺฐครุธมฺม}}
\proseitem{403}{สเจ อานนฺท มหาปชาปติ โคตมี อฏฺฐ ครุธมฺเม ปฏิคฺคณฺหาติ, สาวสฺสา โหตุ อุปสมฺปทา.}

\prose{\csromansymbolfootnote{*}{วิ 2. 74 ปิฏฺเฐปิ.}วสฺสสตูปสมฺปนฺนาย ภิกฺขุนิยา ตทหุปสมฺปนฺนสฺส ภิกฺขุโน อภิวาทนํ ปจฺจุฏฺฐานํ อญฺชลิกมฺมํ สามีจิกมฺมํ กาตพฺพํ, อยมฺปิ ธมฺโม สกฺกตฺวา ครุกตฺวา\footnote{ครุํกตฺวา (ก)} มาเนตฺวา ปูเชตฺวา ยาวชีวํ อนติกฺกมนีโย. (1)}

\prose{น ภิกฺขุนิยา อภิกฺขุเก อาวาเส วสฺสํ วสิตพฺพํ, อยมฺปิ ธมฺโม สกฺกตฺวา ครุกตฺวา มาเนตฺวา ปูเชตฺวา ยาวชีวํ อนติกฺกมนีโย. (2)}

\prose{อนฺวทฺธมาสํ ภิกฺขุนิยา ภิกฺขุสํฆโต เทฺว ธมฺมา ปจฺจาสีสิตพฺพา อุโปสถปุจฺฉกญฺจ โอวาทูปสงฺกมนญฺจ, อยมฺปิ ธมฺโม สกฺกตฺวา ครุกตฺวา มาเนตฺวา ปูเชตฺวา ยาวชีวํ อนติกฺกมนีโย. (3)}

\prose{วสฺสํวุฏฺฐาย ภิกฺขุนิยา อุภโตสํเฆ ตีหิ ฐาเนหิ ปวาเรตพฺพํ ทิฏฺเฐน วา สุเตน วา ปริสงฺกาย วา, อยมฺปิ ธมฺโม สกฺกตฺวา ครุกตฺวา มาเนตฺวา ปูเชตฺวา ยาวชีวํ อนติกฺกมนีโย. (4)}

\prose{ครุธมฺมํ อชฺฌาปนฺนาย ภิกฺขุนิยา อุภโตสํเฆ ปกฺขมานตฺตํ จริตพฺพํ, อยมฺปิ ธมฺโม สกฺกตฺวา ครุกตฺวา มาเนตฺวา ปูเชตฺวา ยาวชีวํ อนติกฺกมนีโย. (5)}

\prose{เทฺว วสฺสานิ ฉสุ ธมฺเมสุ สิกฺขิตสิกฺขาย สิกฺขมานาย อุภโตสํเฆ อุปสมฺปทา ปริเยสิตพฺพา, อยมฺปิ ธมฺโม สกฺกตฺวา ครุกตฺวา มาเนตฺวา ปูเชตฺวา ยาวชีวํ อนติกฺกมนีโย. (6)}

\prose{\csromanfolio{445}น ภิกฺขุนิยา เกนจิ ปริยาเยน ภิกฺขุ อกฺโกสิตพฺโพ ปริภาสิตพฺโพ, อยมฺปิ ธมฺโม สกฺกตฺวา ครุกตฺวา มาเนตฺวา ปูเชตฺวา ยาวชีวํ อนติกฺกมนีโย. (7)}

\prose{อชฺชตคฺเค โอวโฏ ภิกฺขุนีนํ ภิกฺขูสุ วจนปโถ, อโนวโฏ ภิกฺขูนํ ภิกฺขุนีสุ วจนปโถ, อยมฺปิ ธมฺโม สกฺกตฺวา ครุกตฺวา มาเนตฺวา ปูเชตฺวา ยาวชีวํ อนติกฺกมนีโย. (8)}

\prose{สเจ อานนฺท มหาปชาปติ โคตมี อิเม อฏฺฐ ครุธมฺเม ปฏิคฺคณฺหาติ, สาวสฺสา โหตุ อุปสมฺปทาติ.}

\prose{อถ โข อายสฺมา อนนฺโท ภควโต สนฺติเก อฏฺฐ ครุธมฺเม อุคฺคเหตฺวา เยน มหาปชาปติ โคตมี เตนุปสงฺกมิ, อุปสงฺกมิตฺวา มหาปชาปติํ โคตมิํ เอตทโวจ “สเจ โข ตฺวํ โคตมิ อฏฺฐ ครุธมฺเม ปฏิคฺคเณฺหยฺยาสิ, สาว เต ภวิสฺสติ อุปสมฺปทา.}

\prose{วสฺสสตูปสมฺปนฺนาย ภิกฺขุนิยา ตทหุปสมฺปนฺนสฺส ภิกฺขุโน อภิวาทนํ ปจฺจุฏฺฐานํ อญฺชลิกมฺมํ สามีจิกมฺมํ กาตพฺพํ, อยมฺปิ ธมฺโม สกฺกตฺวา ครุกตฺวา มาเนตฺวา ปูเชตฺวา ยาวชีวํ อนติกฺกมนีโย.}

\prose{น ภิกฺขุนิยา อภิกฺขุเก อาวาเส วสฺสํ วสิตพฺพํ, อยมฺปิ ธมฺโม สกฺกตฺวา ครุกตฺวา มาเนตฺวา ปูเชตฺวา ยาวชีวํ อนติกฺกมนีโย.}

\prose{อนฺวทฺธมาสํ ภิกฺขุนิยา ภิกฺขุสํฆโต เทฺว ธมฺมา ปจฺจาสีสิตพฺพา อุโปสถปุจฺฉกญฺจ โอวาทูปสงฺกมนญฺจ, อยมฺปิ ธมฺโม สกฺกตฺวา ครุกตฺวา มาเนตฺวา ปูเชตฺวา ยาวชีวํ อนติกฺกมนีโย.}

\prose{วสฺสํวุฏฺฐาย ภิกฺขุนิยา อุภโตสํเฆ ตีหิ ฐาเนหิ ปวาเรตพฺพํ ทิฏฺเฐน วา สุเตน วา ปริสงฺกาย วา, อยมฺปิ ธมฺโม สกฺกตฺวา ครุกตฺวา มาเนตฺวา ปูเชตฺวา ยาวชีวํ อนติกฺกมนีโย.}

\prose{ครุธมฺมํ อชฺฌาปนฺนาย ภิกฺขุนิยา อุภโตสํเฆ ปกฺขมานตฺตํ จริตพฺพํ, อยมฺปิ ธมฺโม สกฺกตฺวา ครุกตฺวา มาเนตฺวา ปูเชตฺวา ยาวชีวํ อนติกฺกมนีโย.}

\prose{เทฺว วสฺสานิ ฉสุ ธมฺเมสุ สิกฺขิตสิกฺขาย สิกฺขมานาส อุภโตสํเฆ อุปสมฺปทา ปริเยสิตพฺพา, อยมฺปิ ธมฺโม สกฺกตฺวา ครุกตฺวา มาเนตฺวา ปูเชตฺวา ยาวชีวํ อนติกฺกมนีโย.}

\prose{\csromanfolio{446}น ภิกฺขุนิยา เกนจิ ปริยาเยน ภิกฺขุ อกฺโกสิตพฺโพ ปริภาสิตพฺโพ, อยมฺปิ ธมฺโม สกฺกตฺวา ครุกตฺวา มาเนตฺวา ปูเชตฺวา ยาวชีวํ อนติกฺกมนิโย.}

\prose{อชฺชตคฺเค โอวโฏ ภิกฺขุนีนํ ภิกฺขูสุ วจนปโถ, อโนวโฏ ภิกฺขูนํ ภิกฺขุนีสุ วจนปโถ, อยมฺปิ ธมฺโม สกฺกตฺวา ครุกตฺวา มาเนตฺวา ปูเชตฺวา ยาวชีวํ อนติกฺกมนีโย.}

\prose{สเจ โข ตฺวํ โคตมิ อิเม อฏฺฐ ครุธมฺเม ปฏิคฺคเณฺหยฺยาสิ, สาว เต ภวิสฺสติ อุปสมฺปทา”ติ.}

\prose{เสยฺยถาปิ ภนฺเต อานนฺท อิตฺถี วา ปุริโส วา ทหโร ยุวา มณฺฑนกชาติโก สีสํนหาโต อุปฺปลมาลํ วา วสฺสิกมาลํ วา อติมุตฺตกมาลํ\footnote{อธิมตฺตกมาลํ (สฺยา)} วา ลภิตฺวา อุโภหิ หตฺเถหิ ปฏิคฺคเหตฺวา อุตฺตมงฺเค สิรสฺมิํ ปติฏฺฐาเปยฺย, เอวเมว โข อหํ ภนฺเต อานนฺท อิเม อฏฺฐ ครุธมฺเม ปฏิคฺคณฺหามิ ยาวชีวํ อนติกฺกมนีเยติ.}

\prose{อถ โข อายสฺมา อานนฺโท เยน ภควา เตนุปสงฺกมิ, อุปสงฺกมิตฺวา ภควนฺตํ อภิวาเทตฺวา เอกมนฺตํ นิสีทิ, เอกมนฺตํ นิสินฺโน โข อายสฺมา อานนฺโท ภควนฺตํ เอตทโวจ “ปฏิคฺคหิตา ภนฺเต มหาปชาปติยา โคตมิยา อฏฺฐ ครุธมฺมา, อุปสมฺปนฺนา ภควโต มาตุจฺฉา”ติ.}

\prose{สเจ อานนฺท นาลภิสฺส มาตุคาโม ตถาคตปฺปเวทิเต ธมฺมวินเย อคารสฺมา อนคาริยํ ปพฺพชฺชํ, จิรฏฺฐิติกํ อานนฺท พฺรหฺมจริยํ อภวิสฺส, วสฺสสหสฺสํ สทฺธมฺโม ติฏฺเฐยฺย.\csromanspacer{} ยโต จ โข อานนฺท มาตุคาโม ตถาคตปฺปเวทิเต ธมฺมวินเย อคารสฺมา อนคาริยํ ปพฺพชิโต, น ทานิ อานนฺท พฺรหฺมจริยํ จิรฏฺฐิติกํ ภวิสฺสติ, ปญฺเจว ทานิ อานนฺท วสฺสสตานิ สทฺธมฺโม ฐสฺสติ.}

\prose{เสยฺยถาปิ อานนฺท ยานิ กานิจิ กุลานิ พหุตฺถิกานิ\footnote{พหุ-อิตฺถิกานิ (สี, สฺยา)} อปฺปปุริสกานิ, ตานิ สุปฺปธํสิยานิ โหนฺติ โจเรหิ กุมฺภเถนเกหิ, เอวเมว โข อานนฺท ยสฺมิํ ธมฺมวินเย ลภติ มาตุคาโม อคารสฺมา อนคาริยํ ปพฺพชฺชํ, น ตํ พฺรหฺมจริยํ จิรฏฺฐิติกํ โหติ.}

\prose{\csromanfolio{447}เสยฺยถาปิ อานนฺท สมฺปนฺเน สาลิกฺเขตฺเต เสตฏฺฏิกา นาม โรคชาติ นิปตติ, เอวํ ตํ สาลิกฺเขตฺตํ น จิรฏฺฐิติกํ โหติ, เอวเมว โข อานนฺท ยสฺมิํ ธมฺมวินเย ลภติ มาตุคาโม อคารสฺมา อนคาริยํ ปพฺพชฺชํ, น ตํ พฺรหฺมจริยํ จิรฏฺฐิติกํ โหติ.}

\prose{เสยฺยถาปิ อานนฺท สมฺปนฺเน อุจฺฉุกฺเขตฺเต มญฺชิฏฺฐิกา\footnote{มญฺเชฏฺฐิกา (สี, สฺยา)} นาม โรคชาติ นิปตติ, เอวํ ตํ อุจฺฉุกฺเขตฺตํ น จิรฏฺฐิติกํ โหติ, เอวเมว โข อานนฺท ยสฺมิํ ธมฺมวินเย ลภติ มาตุคาโม อคารสฺมา อนคาริยํ ปพฺพชฺชํ, น ตํ พฺรหฺมจริยํ จิรฏฺฐิติกํ โหติ.}

\prose{เสยฺยถาปิ อานนฺท ปุริโส มหโต ตฬากสฺส ปฏิกจฺเจว อาฬิํ พนฺเธยฺย ยาวเทว อุทกสฺส อนติกฺกมนาย, เอวเมว โข อานนฺท มยา ปฏิกจฺเจว ภิกฺขุนีนํ อฏฺฐ ครุธมฺมา ปญฺญตฺตา ยาวชีวํ อนติกฺกมนียาติ.}

\nitthitam{ภิกฺขุนีนํ อฏฺฐ ครุธมฺมา นิฏฺฐิตา.}
\csromanmatikaanchor{mtk.208}
\csromantocmarkref{subsection}{ภิกฺขุนี-อุปสมฺปทานุชานน}{mtk.208}
\bookmark[dest={mtk.208},level=2]{ภิกฺขุนี-อุปสมฺปทานุชานน}
\csromanheader{2}{\textbf{ภิกฺขุนี-อุปสมฺปทานุชานน}}
\proseitem{404}{อถ โข มหาปชาปติ โคตมี เยน ภควา เตนุปสงฺกมิ, อุปสงฺกมิตฺวา ภควนฺตํ อภิวาเทตฺวา เอกมนฺตํ อฏฺฐาสิ, เอกมนฺตํ ฐิตา โข มหาปชาปติ โคตมี ภควนฺตํ เอตทโวจ “กถาหํ ภนฺเต อิมาสุ สากิยานีสุ ปฏิปชฺชามี”ติ.\csromanspacer{} อถ โข ภควา มหาปชาปติํ โคตมิํ ธมฺมิยา กถาย สนฺทสฺเสสิ สมาทเปสิ สมุตฺเตเชสิ สมฺปหํเสสิ.\csromanspacer{} อถ โข มหาปชาปติ โคตมี ภควตา ธมฺมิยา กถาย สนฺทสฺสิตา สมาทปิตา สมุตฺเตชิตา สมฺปหํสิตา ภควนฺตํ อภิวาเทตฺวา ปทกฺขิณํ กตฺวา ปกฺกามิ.\csromanspacer{} อถ โข ภควา เอตสฺมิํ นิทาเน เอตสฺมิํ ปกรเณ ธมฺมิํ กถํ กตฺวา ภิกฺขู อามนฺเตสิ “อนุชานามิ ภิกฺขเว ภิกฺขูหิ ภิกฺขุนิโย อุปสมฺปาเทตุนฺ”ติ.}

\prose{อถ โข ตา ภิกฺขุนิโย มหาปชาปติํ โคตมิํ เอตทโวจุํ “อยฺยา อนุปสมฺปนฺนา, มยญฺจมฺหา อุปสมฺปนฺนา, เอวํ หิ ภควตา ปญฺญตฺตํ, ‘ภิกฺขูหิ ภิกฺขุนิโย อุปสมฺปาเทตพฺพา’ติ”.\csromanspacer{} อถ โข มหาปชาปติ \csromanfolio{448}โคตมี เยนายสฺมา อานนฺโท เตนุปสงฺกมิ, อุปสงฺกมิตฺวา อายสฺมนฺตํ อานนฺทํ อภิวาเทตฺวา เอกมนฺตํ อฏฺฐาสิ, เอกมนฺตํ ฐิตา โข มหาปชาปติ โคตมี อายสฺมนฺตํ อานนฺทํ เอตทโวจ “อิมา มํ ภนฺเต อานนฺท ภิกฺขุนิโย เอวมาหํสุ ‘อยฺยา อนุปสมฺปนฺนา, มยญฺจมฺหา อุปสมฺปนฺนา, เอวํ หิ ภควตา ปญฺญตฺตํ ภิกฺขูหิ ภิกฺขุนิโย อุปสมฺปาเทตพฺพา’ติ”.}

\prose{อถ โข อายสฺมา อานนฺโท เยน ภควา เตนุปสงฺกมิ, อุปสงฺกมิตฺวา ภควนฺตํ อภิวาเทตฺวา เอกมนฺตํ นิสีทิ, เอกมนฺตํ นิสินฺโน โข อายสฺมา อานนฺโท ภควนฺตํ เอตทโวจ “มหาปชาปติ ภนฺเต โคตมี เอวมาห ‘อิมา มํ ภนฺเต อานนฺท ภิกฺขุนิโย เอวมาหํสุ อยฺยา อนุปสมฺปนฺนา, มยญฺจมฺหา อุปสมฺปนฺนา, เอวํ หิ ภควตา ปญฺญตฺตํ ภิกฺขูหิ ภิกฺขุนิโย อุปสมฺปาเทตพฺพา’ติ”.}

\prose{ยทคฺเคน อานนฺท มหาปชาปติยา โคตมิยา อฏฺฐ ครุธมฺมา ปฏิคฺคหิตา, ตเทว สา\footnote{ตเทวสฺสา (ก)} อุปสมฺปนฺนาติ.}

\proseitem{405}{อถ โข มหาปชาปติ โคตมี เยนายสฺมา อานนฺโท เตนุปสงฺกมิ, อุปสงฺกมิตฺวา อายสฺมนฺตํ อานนฺทํ อภิวาเทตฺวา เอกมนฺตํ อฏฺฐาสิ, เอกมนฺตํ ฐิตา โข มหาปชาปติ โคตมี อายสฺมนฺตํ อานนฺทํ เอตทโวจ “เอกาหํ ภนฺเต อานนฺท ภควนฺตํ วรํ ยาจามิ, สาธุ ภนฺเต, ภควา อนุชาเนยฺย ภิกฺขูนญฺจ ภิกฺขุ-นีนญฺจ ยถาวุฑฺฒํ อภิวาทนํ ปจฺจุฏฺฐานํ อญฺชลิกมฺมํ สามีจิกมฺมนฺ”ติ.}

\prose{อถ โข อายสฺมา อานนฺโท เยน ภควา เตนุปสงฺกมิ, อุปสงฺกมิตฺวา ภควนฺตํ อภิวาเทตฺวา เอกมนฺตํ นิสีทิ, เอกมนฺตํ นิสินฺโน โข อายสฺมา อานนฺโท ภควนฺตํ เอตทโวจ “มหาปชาปติ ภนฺเต โคตมี เอวมาห ‘เอกาหํ ภนฺเต อานนฺท ภควนฺตํ วรํ ยาจามิ, สาธุ ภนฺเต ภควา อนุชาเนยฺย ภิกฺขูนญฺจ ภิกฺขุนีนญฺจ ยถาวุฑฺฒํ อภิวาทนํ ปจฺจุฏฺฐานํ อญฺชลิกมฺมํ สามีจิกมฺมนฺ’ติ”.}

\prose{อฏฺฐานเมตํ อานนฺท อนวกาโส, ยํ ตถาคโต อนุชาเนยฺย มาตุคามสฺส อภิวาทนํ ปจฺจุฏฺฐานํ อญฺชลิกมฺมํ สามีจิกมฺมํ, อิเมหิ นาม อานนฺท อญฺญติตฺถิยา ทุรกฺขาตธมฺมา มาตุคามสฺส อภิวาทนํ ปจฺจุฏฺฐานํ \csromanfolio{449}อญฺชลิกมฺมํ สามีจิกมฺมํ น กริสฺสนฺติ, กิมงฺคํ ปน ตถาคโต อนุชานิสฺสติ มาตุคามสฺส อภิวาทนํ ปจฺจุฏฺฐานํ อญฺชลิกมฺมํ สามีจิกมฺมนฺติ.}

\prose{อถ โข ภควา เอตสฺมิํ นิทาเน เอตสฺมิํ ปกรเณ ธมฺมิํ กถํ กตฺวา ภิกฺขู อามนฺเตสิ “น ภิกฺขเว มาตุคามสฺส อภิวาทนํ ปจฺจุฏฺฐานํ อญฺชลิกมฺมํ สามีจิกมฺมํ กาตพฺพํ, โย กเรยฺย, อาปตฺติ ทุกฺกฏสฺสา”ติ.}

\prose{อถ โข มหาปชาปติ โคตมี เยน ภควา เตนุปสงฺกมิ อุปสงฺกมิตฺวา ภควนฺตํ อภิวาเทตฺวา เอกมนฺตํ อฏฺฐาสิ, เอกมนฺตํ ฐิตา โข มหาปชาปติ โคตมี ภควนฺตํ เอตทโวจ “ยานิ ตานิ ภนฺเต ภิกฺขุนีนํ สิกฺขาปทานิ ภิกฺขูหิ สาธารณานิ, กถํ มยํ ภนฺเต เตสุ สิกฺขาปเทสุ ปฏิปชฺชามา”ติ.\csromanspacer{} ยานิ ตานิ โคตมิ ภิกฺขุนีนํ สิกฺขาปทานิ ภิกฺขูหิ สาธารณานิ, ยถา ภิกฺขู สิกฺขนฺติ ตถา เตสุ สิกฺขาปเทสุ สิกฺขถาติ.\csromanspacer{} ยานิ ปน ตานิ ภนฺเต ภิกฺขุนีนํ สิกฺขาปทานิ ภิกฺขูหิ อสาธารณานิ, กถํ มยํ ภนฺเต เตสุ สิกฺขาปเทสุ ปฏิปชฺชามาติ.\csromanspacer{} ยานิ ตานิ โคตมิ ภิกฺขุนีนํ สิกฺขาปทานิ ภิกฺขูหิ อสาธรณานิ, ยถาปญฺญตฺเตสุ สิกฺขาปเทสุ สิกฺขถาติ.}

\proseitem{406}{\csromansymbolfootnote{*}{อํ 3. 106 ปิฏฺเฐปิ.}อถ โข มหาปชาปติ โคตมี เยน ภควา เตนุปสงฺกมิ, อุปสงฺกมิตฺวา ภควนฺตํ อภิวาเทตฺวา เอกมนฺตํ อฏฺฐาสิ, เอกมนฺตํ ฐิตา โข มหาปชาปติ โคตมี ภควนฺตํ เอตทโวจ “สาธุ เม ภนฺเต ภควา สํขิตฺเตน ธมฺมํ เทเสตุ, ยมหํ ภควโต ธมฺมํ สุตฺวา เอกา วูปกฏฺฐา อปฺปมตฺตา อาตาปินี ปหิตตฺตา วิหเรยฺยนฺ”ติ.\csromanspacer{} เย โข ตฺวํ โคตมิ ธมฺเม ชาเนยฺยาสิ, อิเม ธมฺมา สราคาย สํวตฺตนฺติ, โน วิราคาย, สญฺโญคาย สํวตฺตนฺติ, โน วิสญฺโญคาย, อาจยาย สํวตฺตนฺติ, โน อปจยาย, มหิจฺฉตาย สํวตฺตนฺติ, โน อปฺปิจฺฉตาย, อสนฺตุฏฺฐิยา สํวตฺตนฺติ, โน สนฺตุฏฺฐิยา, สงฺคณิกาย สํวตฺตนฺติ, โน ปวิเวกาย, โกสชฺชาย สํวตฺตนิ, โน วีริยารมฺภาย, ทุพฺภรตาย สํวตฺตนฺติ, โน สุภรตาย, เอกํเสน โคตมิ ธาเรยฺยาสิ “เนโส ธมฺโม, เนโส วินโย, เนตํ สตฺถุสาสนนฺ”ติ.\csromanspacer{} เย จ โข ตฺวํ โคตมิ ธมฺเม ชาเนยฺยาสิ, อิเม ธมฺมา วิราคาย สํวตฺตนฺติ, โน สราคาย, วิสญฺโญคาย สํวตฺตนฺติ, โน สญฺโญคาย, อปจยาย สํวตฺตนฺติ, โน อาจยาย, อปฺปิจฺฉตาย สํวตฺตนฺติ, \csromanfolio{450}โน มหิจฺฉตาย, สนฺตุฏฺฐิยา สํวตฺตนฺติ, โน อสนฺตุฏฺฐิยา, ปวิเวกาย สํวตฺตนฺติ, โน สงฺคณิกาย, วีริยารมฺภาย สํวตฺตนฺติ, โน โกสชฺชาย, สุภรตาย สํวตฺตนฺติ, โน ทุพฺภรตาย, เอกํเสน โคตมิ ธาเรยฺยาสิ “เอโส ธมฺโม, เอโส วินโย, เอตํ สตฺถุสาสนนฺ”ติ.}

\proseitem{407}{เตน โข ปน สมเยน ภิกฺขุนีนํ ปาติโมกฺขํ น อุทฺทิสียติ.\csromanspacer{} ภควโต เอตมตฺถํ อาโรเจสุํ ฯเปฯ.\csromanspacer{} อนุชานามิ ภิกฺขเว ภิกฺขุนีนํ ปาติโมกฺขํ อุทฺทิสิตุนฺติ.\csromanspacer{} อถ โข ภิกฺขูนํ เอตทโหสิ “เกน นุ โข ภิกฺขุนีนํ ปาติโมกฺขํ อุทฺทิสิตพฺพนฺ”ติ.\csromanspacer{} ภควโต เอตมตฺถํ อาโรเจสุํ.\csromanspacer{} อนุชานามิ ภิกฺขเว ภิกฺขูหิ ภิกฺขุนีนํ ปาติโมกฺขํ อุทฺทิสิตุนฺติ.}

\prose{เตน โข ปน สมเยน ภิกฺขู ภิกฺขุนุปสฺสยํ อุปสงฺกมิตฺวา ภิกฺขุนีนํ ปาติโมกฺขํ อุทฺทิสนฺติ.\csromanspacer{} มนุสฺสา อุชฺฌายนฺติ ขิยฺยนฺติ วิปาเจนฺติ “ชายาโย อิมา อิเมสํ, ชาริโย อิมา อิเมสํ, อิทานิ อิเม อิมาหิ สทฺธิํ อภิรมิสฺสนฺตี”ติ.\csromanspacer{} ภควโต เอตมตฺถํ อาโรเจสุํ.\csromanspacer{} น ภิกฺขเว ภิกฺขูหิ ภิกฺขุนีนํ ปาติโมกฺขํ อุทฺทิสิตพฺพํ, โย อุทฺทิเสยฺย, อาปตฺติ ทุกฺกฏสฺส.\csromanspacer{} อนุชานามิ ภิกฺขเว ภิกฺขุนีหิ ภิกฺขุนีนํ ปาติโมกฺขํ อุทฺทิสิตุนฺติ.\csromanspacer{} ภิกฺขุนิโย น ชานนฺติ “เอวํ ปาติโมกฺขํ อุทฺทิสิตพฺพนฺ”ติ.\csromanspacer{} ภควโต เอตมตฺถํ อาโรเจสุํ.\csromanspacer{} อนุชานามิ ภิกฺขเว ภิกฺขูหิ ภิกฺขุนีนํ อาจิกฺขิตุํ “เอวํ ปาติโมกฺขํ อุทฺทิเสยฺยาถา”ติ.}

\proseitem{408}{เตน โข ปน สมเยน ภิกฺขุนิโย อาปตฺติํ น ปฏิกโรนฺติ.\csromanspacer{} ภควโต เอตมตฺถํ อาโรเจสุํ.\csromanspacer{} น ภิกฺขเว ภิกฺขุนิยา อาปตฺติ น ปฏิกาตพฺพา, ยา น ปฏิกเรยฺย, อาปตฺติ ทุกฺกฏสฺสาติ.\csromanspacer{} ภิกฺขุนิโย น ชานนฺติ “เอวํ อาปตฺติ ปฏิกาตพฺพา”ติ.\csromanspacer{} ภควโต เอตมตฺถํ อาโรเจสุํ.\csromanspacer{} อนุชานามิ ภิกฺขเว ภิกฺขูหิ ภิกฺขุนีนํ อาจิกฺขิตุํ “เอวํ อาปตฺติํ ปฏิกเรยฺยาถา”ติ.\csromanspacer{} อถ โข ภิกฺขูนํ เอตทโหสิ “เกน นุ โข ภิกฺขุนีนํ อาปตฺติ ปฏิคฺคเหตพฺพา”ติ.\csromanspacer{} ภควโต เอตมตฺถํ อาโรเจสุํ.\csromanspacer{} อนุชานามิ ภิกฺขเว ภิกฺขูหิ ภิกฺขุนีนํ อาปตฺติํ ปฏิคฺคเหตุนฺติ.}

\prose{เตน โข ปน สมเยน ภิกฺขุนิโย รถิกายปิ พฺยูเหปิ สิงฺฆาฏเกปิ ภิกฺขุํ ปสฺสิตฺวา ปตฺตํ ภูมิยํ นิกฺขิปิตฺวา เอกํสํ อุตฺตราสงฺคํ กริตฺวา \csromanfolio{451}อุกฺกุฏิกํ นิสีทิตฺวา อญฺชลิํ ปคฺคเหตฺวา อาปตฺติํ ปฏิกโรนฺติ.\csromanspacer{} มนุสฺสา อุชฺฌายนฺติ ขิยฺยนฺติ วิปาเจนฺติ “ชายาโย อิมา อิเมสํ, ชาริโย อิมา อิเมสํ, รตฺติํ วิมาเนตฺวา อิทานิ ขมาเปนฺตี”ติ.\csromanspacer{} ภควโต เอตมตฺถํ อาโรเจสุํ.\csromanspacer{} น ภิกฺขเว ภิกฺขูหิ ภิกฺขุนีนํ อาปตฺติ ปฏิคฺคเหตพฺพา, โย ปฏิคฺคเณฺหยฺย, อาปตฺติ ทุกฺกฏสฺส.\csromanspacer{} อนุชานามิ ภิกฺขเว ภิกฺขุนีหิ ภิกฺขุนีนํ อาปตฺติํ ปฏิคฺคเหตุนฺติ.\csromanspacer{} ภิกฺขุนิโย น ชานนฺติ “เอวํ อาปตฺติ ปฏิคฺคเหตพฺพา”ติ.\csromanspacer{} ภควโต เอตมตฺถํ อาโรเจสุํ.\csromanspacer{} อนุชานามิ ภิกฺขเว ภิกฺขูหิ ภิกฺขุนีนํ อาจิกฺขิตุํ “เอวํ อาปตฺติํ ปฏิคฺคเณฺหยฺยาถา”ติ.}

\proseitem{409}{เตน โข ปน สมเยน ภิกฺขุนีนํ กมฺมํ น กริยติ.\csromanspacer{} ภควโต เอตมตฺถํ อาโรเจสุํ.\csromanspacer{} อนุชานามิ ภิกฺขเว ภิกฺขุนีนํ กมฺมํ กาตุนฺติ.\csromanspacer{} อถ โข ภิกฺขูนํ เอตทโหสิ “เกน นุ โข ภิกฺขุนีนํ กมฺมํ กาตพฺพนฺ”ติ.\csromanspacer{} ภควโต เอตมตฺถํ อาโรเจสุํ.\csromanspacer{} อนุชานามิ ภิกฺขเว ภิกฺขูหิ ภิกฺขุนีนํ กมฺมํ กาตุนฺติ.}

\prose{เตน โข ปน สมเยน กตกมฺมา ภิกฺขุนิโย รถิกายปิ พฺยูเหปิ สิงฺฆาฏเกปิ ภิกฺขุํ ปสฺสิตฺวา ปตฺตํ ภูมิยํ นิกฺขิปิตฺวา เอกํสํ อุตฺตราสงฺคํ กริตฺวา อุกฺกุฏิกํ นิสีทิตฺวา อญฺชลิํ ปคฺคเหตฺวา ขมาเปนฺติ “เอวํ นูน กาตพฺพนฺ”ติ มญฺญมานา.\csromanspacer{} มนุสฺสา ตเถว อุชฺฌายนฺติ ขิยฺยนฺติ วิปาเจนฺติ “ชายาโย อิมา อิเมสํ, ชาริโย อิมา อิเมสํ, รตฺติํ วิมาเนตฺวา อิทานิ ขมาเปนฺตี”ติ.\csromanspacer{} ภควโต เอตมตฺถํ อาโรเจสุํ.\csromanspacer{} น ภิกฺขเว ภิกฺขูหิ ภิกฺขุนีนํ กมฺมํ กาตพฺพํ, โย กเรยฺย, อาปตฺติ ทุกฺกฏสฺส.\csromanspacer{} อนุชานามิ ภิกฺขเว ภิกฺขุนีหิ ภิกฺขุนีนํ กมฺมํ กาตุนฺติ.\csromanspacer{} ภิกฺขุนิโย น ชานนฺติ “เอวํ กมฺมํ กาตพฺพนฺ”ติ.\csromanspacer{} ภควโต เอตมตฺถํ อาโรเจสุํ.\csromanspacer{} อนุชานามิ ภิกฺขเว ภิกฺขูหิ ภิกฺขุนีนํ อาจิกฺขิตุํ “เอวํ กมฺมํ กเรยฺยาถา”ติ.}

\proseitem{410}{เตน โข ปน สมเยน ภิกฺขุนิโย สํฆมชฺเฌ ภณฺฑนชาตา กลหชาตา วิวาทาปนฺนา อญฺญมญฺญํ มุขสตฺตีหิ วิตุทนฺตา วิหรนฺติ, น สกฺโกนฺติ ตํ อธิกรณํ วูปสเมตุํ.\csromanspacer{} ภควโต เอตมตฺถํ อาโรเจสุํ.\csromanspacer{} อนุชานามิ ภิกฺขเว ภิกฺขูหิ ภิกฺขุนีนํ อธิกรณํ วูปสเมตุนฺติ.}

\prose{\csromanfolio{452}เตน โข ปน สมเยน ภิกฺขู ภิกฺขุนีนํ อธิกรณํ วูปสเมนฺติ, ตสฺมิํ โข ปน อธิกรเณ วินิจฺฉิยมาเน ทิสฺสนฺติ ภิกฺขุนิโย กมฺมปฺปตฺตาโยปิ อาปตฺติคามินิโยปิ, ภิกฺขุนิโย เอวมาหํสุ “สาธุ ภนฺเต, อยฺยาว ภิกฺขุนีนํ กมฺมํ กโรนฺตุ, อยฺยาว ภิกฺขุนีนํ อาปตฺติํ ปฏิคฺคณฺหนฺตุ, เอวํ หิ ภควตา ปญฺญตฺตํ ภิกฺขูหิ ภิกฺขุนีนํ อธิกรณํ วูปสเมตพฺพนฺ”ติ.\csromanspacer{} ภควโต เอตมตฺถํ อาโรเจสุํ.\csromanspacer{} อนุชานามิ ภิกฺขเว ภิกฺขูหิ ภิกฺขุนีนํ กมฺมํ อาโรเปตฺวา ภิกฺขุนีนํ นิยฺยาเทตุํ, ภิกฺขุนีหิ ภิกฺขุนีนํ กมฺมํ กาตุํ, ภิกฺขูหิ ภิกฺขุนีนํ อาปตฺติํ อาโรเปตฺวา ภิกฺขุนีนํ นิยฺยาเทตุํ, ภิกฺขุนีนํ อาปตฺติํ ปฏิคฺคเหตุนฺติ.}

\prose{เตน โข ปน สมเยน อุปฺปลวณฺณาย ภิกฺขุนิยา อนฺเตวาสินี ภิกฺขุนี สตฺต วสฺสานิ ภควนฺตํ อนุพนฺธา โหติ วินยํ ปริยาปุณนฺตี, ตสฺสา มุฏฺฐสฺสตินิยา คหิโต คหิโต มุสฺสติ.\csromanspacer{} อสฺโสสิ โข สา ภิกฺขุนี “ภควา กิร สาวตฺถิํ คนฺตุกาโม”ติ.\csromanspacer{} อถ โข ตสฺสา ภิกฺขุนิยา เอตทโหสี “อหํ โข สตฺต วสฺสานิ ภควนฺตํ อนุพนฺธิํ วินยํ ปริยาปุณนฺตี, ตสฺสา เม มุฏฺฐสฺสตินิยา คหิโต คหิโต มุสฺสติ, ทุกฺกรํ โข ปน มาตุคาเมน ยาวชีวํ สตฺถารํ อนุพนฺธิตุํ, กถํ นุ โข มยา ปฏิปชฺชิตพฺพนฺ”ติ.\csromanspacer{} อถ โข สา ภิกฺขุนี ภิกฺขุนีนํ เอตมตฺถํ อาโรเจสิ, ภิกฺขุนิโย ภิกฺขูนํ เอตมตฺถํ อาโรเจสุํ, ภิกฺขู ภควโต เอตมตฺถํ อาโรเจสุํ.\csromanspacer{} อนุชานามิ ภิกฺขเว ภิกฺขูหิ ภิกฺขุนีนํ วินยํ วาเจตุนฺติ.}

\nitthitam{ปฐมภาณวาโร นิฏฺฐิโต.}
\csromanmatikaanchor{mtk.209}
\csromantocmarkref{section}{2. ทุติยภาณวาร}{mtk.209}
\bookmark[dest={mtk.209},level=1]{2. ทุติยภาณวาร}
\csromanheader{1}{2. ทุติยภาณวาร}
\proseitem{411}{อถ โข ภควา เวสาลิยํ ยถาภิรนฺตํ วิหริตฺวา เยน สาวตฺถิ เตน จาริกํ ปกฺกามิ, อนุปุพฺเพน จาริกํ จรมาโน เยน สาวตฺถิ ตทวสริ, ตตฺร สุทํ ภควา สาวตฺถิยํ วิหรติ เชตวเน อนาถปิณฺฑิกสฺส อาราเม.\csromanspacer{} เตน โข ปน สมเยน ฉพฺพคฺคิยา ภิกฺขู ภิกฺขุนิโย กทฺทโมทเกน โอสิญฺจนฺติ “อปฺเปว นาม อเมฺหสุ สารชฺเชยฺยุนฺ”ติ.\csromanspacer{} ภควโต เอตมตฺถํ อาโรเจสุํ.\csromanspacer{} น ภิกฺขเว \csromanfolio{453}ภิกฺขุนา ภิกฺขุนิโย กทฺทโมทเกน โอสิญฺจิตพฺพา, โย โอสิญฺเจยฺย, อาปตฺติ ทุกฺกฏสฺส.\csromanspacer{} อนุชานามิ ภิกฺขเว ตสฺส ภิกฺขุโน ทณฺฑกมฺมํ กาตุนฺติ.\csromanspacer{} อถ โข ภิกฺขูนํ เอตทโหสิ “กิํ นุ โข ทณฺฑกมฺมํ กาตพฺพนฺ”ติ.\csromanspacer{} ภควโต เอตมตฺถํ อาโรเจสุํ.\csromanspacer{} อวนฺทิโย โส ภิกฺขเว ภิกฺขุ ภิกฺขุนิสํเฆน กาตพฺโพติ.}

\prose{เตน โข ปน สมเยน ฉพฺพคฺคิยา ภิกฺขู กายํ วิวริตฺวา ภิกฺขุนีนํ ทสฺเสนฺติ ฯเปฯ.\csromanspacer{} อูรุํ วิวริตฺวา ภิกฺขุนีนํ ทสฺเสนฺติ.\csromanspacer{} องฺคชาตํ วิวริตฺวา ภิกฺขุนีนํ ทสฺเสนฺติ.\csromanspacer{} ภิกฺขุนิโย โอภาเสนฺติ.\csromanspacer{} ภิกฺขุนีหิ สทฺธิํ สมฺปโยเชนฺติ “อปฺเปว นาม อเมฺหสุ สารชฺเชยฺยุนฺ”ติ.\csromanspacer{} ภควโต เอตมตฺถํ อาโรเจสุํ.\csromanspacer{} น ภิกฺขเว ภิกฺขุนา กาโย วิวริตฺวา ภิกฺขุนีนํ ทสฺเสตพฺโพ ฯเปฯ.\csromanspacer{} น อูรุ วิวริตฺวา ภิกฺขุนีนํ ทสฺเสตพฺโพ.\csromanspacer{} น องฺคชาตํ วิวริตฺวา ภิกฺขุนีนํ ทสฺเสตพฺพํ.\csromanspacer{} น ภิกฺขุนิโย โอภาสิตพฺพา.\csromanspacer{} น ภิกฺขุนีหิ สทฺธิํ สมฺปโยเชตพฺพํ, โย สมฺปโยเชยฺย, อาปตฺติ ทุกฺกฏสฺส.\csromanspacer{} อนุชานามิ ภิกฺขเว ตสฺส ภิกฺขุโน ทณฺฑกมฺมํ กาตุนฺติ.\csromanspacer{} อถ โข ภิกฺขูนํ เอตทโหสิ “กิํ นุ โข ทณฺฑกมฺมํ กาตพฺพนฺ”ติ.\csromanspacer{} ภควโต เอตมตฺถํ อาโรเจสุํ.\csromanspacer{} อวนฺทิโย โส ภิกฺขเว ภิกฺขุ ภิกฺขุนิสํเฆน กาตพฺโพติ.}

\prose{เตน โข ปน สมเยน ฉพฺพคฺคิยา ภิกฺขุนิโย ภิกฺขุํ กทฺทโมทเกน โอสิญฺจนฺติ “อปฺเปว นาม อเมฺหสุ สารชฺเชยฺยุนฺ”ติ.\csromanspacer{} ภควโต เอตมตฺถํ อาโรเจสุํ.\csromanspacer{} น ภิกฺขเว ภิกฺขุนิยา ภิกฺขุ กทฺทโมทเกน โอสิญฺจิตพฺโพ, ยา โอสิญฺเจยฺย, อาปตฺติ ทุกฺกฏสฺส.\csromanspacer{} อนุชานามิ ภิกฺขเว ตสฺสา ภิกฺขุนิยา ทณฺฑกมฺมํ กาตุนฺติ.\csromanspacer{} อถ โข ภิกฺขูนํ เอตทโหสิ “กิํ นุ โข ทณฺฑกมฺมํ กาตพฺพนฺ”ติ.\csromanspacer{} ภควโต เอตมตฺถํ อาโรเจสุํ.\csromanspacer{} อนุชานามิ ภิกฺขเว อาวรณํ กาตุนฺติ.\csromanspacer{} อาวรเณ กเต น อาทิยนฺติ.\csromanspacer{} ภควโต เอตมตฺถํ อาโรเจสุํ.\csromanspacer{} อนุชานามิ ภิกฺขเว โอวาทํ ฐเปตุนฺติ.}

\prose{เตน โข ปน สมเยน ฉพฺพคฺคิยา ภิกฺขุนิโย กายํ วิวริตฺวา ภิกฺขูนํ ทสฺเสนฺติ, ถนํ วิวริตฺวา ภิกฺขูนํ ทสฺเสนฺติ, อูรุํ วิวริตฺวา ภิกฺขูนํ ทสฺเสนฺติ, องฺคชาตํ วิวริตฺวา ภิกฺขูนํ ทสฺเสนฺติ, ภิกฺขู โอภาเสนฺติ, ภิกฺขูหิ สทฺธิํ สมฺปโยเชนฺติ “อปฺเปว นาม อเมฺหสุ สารชฺเชยฺยุนฺ”ติ.\csromanspacer{} ภควโต เอตมตฺถํ อาโรเจสุํ.\csromanspacer{} น ภิกฺขเว ภิกฺขุนิยา กาโย วิวริตฺวา ภิกฺขูนํ ทสฺเสตพฺโพ ฯเปฯ.\csromanspacer{} น ถโน \csromanfolio{454}วิวริตฺวา ภิกฺขูนํ ทสฺเสตพฺโพ.\csromanspacer{} น อูรุ วิวริตฺวา ภิกฺขูนํ ทสฺเสตพฺโพ.\csromanspacer{} น องฺคชาตํ วิวริตฺวา ภิกฺขูนํ ทสฺเสตพฺพํ.\csromanspacer{} น ภิกฺขู โอภาสิตพฺพา.\csromanspacer{} น ภิกฺขูหิ สทฺธิํ สมฺปโยเชตพฺพํ, ยา สมฺปโยเชยฺย, อาปตฺติ ทุกฺกฏสฺส.\csromanspacer{} อนุชานามิ ภิกฺขเว ตสฺสา ภิกฺขุนิยา ทณฺฑกมฺมํ กาตุนฺติ.\csromanspacer{} อถ โข ภิกฺขูนํ เอตทโหสิ “กิํ นุ โข ทณฺฑกมฺมํ กาตพฺพนฺ”ติ.\csromanspacer{} ภควโต เอตมตฺถํ อาโรเจสุํ.\csromanspacer{} อนุชานามิ ภิกฺขเว อาวรณํ กาตุนฺติ.\csromanspacer{} อาวรเณ กเต น อาทิยนฺติ.\csromanspacer{} ภควโต เอตมตฺถํ อาโรเจสุํ.\csromanspacer{} อนุชานามิ ภิกฺขเว โอวาทํ ฐเปตุนฺติ.\csromanspacer{} อถ โข ภิกฺขูนํ เอตทโหสิ “กปฺปติ นุ โข โอวาทฏฺฐปิตาย\footnote{โอวาทณฺฐปิตาย (สฺยา), โอวาทํฐปิตาย (ก)} ภิกฺขุนิยา สทฺธิํ อุโปสถํ กาตุํ, น นุ โข กปฺปตี”ติ.\csromanspacer{} ภควโต เอตมตฺถํ อาโรเจสุํ.\csromanspacer{} น ภิกฺขเว โอวาทฏฺฐปิตาย ภิกฺขุนิยา สทฺธิํ อุโปสโถ กาตพฺโพ, ยาว น ตํ อธิกรณํ วูปสนฺตํ โหตีติ.}

\proseitem{412}{เตน โข ปน สมเยน อายสฺมา อุทายี โอวาทํ ฐเปตฺวา จาริกํ ปกฺกามิ.\csromanspacer{} ภิกฺขุนิโย อุชฺฌายนฺติ ขิยฺยนฺติ วิปาเจนฺติ “กถํ หิ นาม อยฺโย อุทายี โอวาทํ ฐเปตฺวา จาริกํ ปกฺกมิสฺสตี”ติ.\csromanspacer{} ภควโต เอตมตฺถํ อาโรเจสุํ.\csromanspacer{} น ภิกฺขเว โอวาทํ ฐเปตฺวา จาริกา ปกฺกมิตพฺพา, โย ปกฺกเมยฺย, อาปตฺติ ทุกฺกฏสฺสาติ.}

\prose{เตน โข ปน สมเยน\footnote{เตน โข ปน สมเยน ภิกฺขู (สฺยา, กํ)} พาลา อพฺยตฺตา โอวาทํ ฐเปนฺติ.\csromanspacer{} ภควโต เอตมตฺถํ อาโรเจสุํ.\csromanspacer{} น ภิกฺขเว พาเลน อพฺยตฺเตน โอวาโท ฐเปตพฺโพ, โย ฐเปยฺย, อาปตฺติ ทุกฺกฏสฺสาติ.}

\prose{เตน โข ปน สมเยน ภิกฺขู อวตฺถุสฺมิํ อการเณ โอวาทํ ฐเปนฺติ.\csromanspacer{} ภควโต เอตมตฺถํ อาโรเจสุํ.\csromanspacer{} น ภิกฺขเว อวตฺถุสฺมิํ อการเณ โอวาโท ฐเปตพฺโพ, โย ฐเปยฺย, อาปตฺติ ทุกฺกฏสฺสาติ.}

\prose{เตน โข ปน สมเยน ภิกฺขู โอวาทํ ฐเปตฺวา วินิจฺฉยํ น เทนฺติ.\csromanspacer{} ภควโต เอตมตฺถํ อาโรเจสุํ.\csromanspacer{} น ภิกฺขเว โอวาทํ ฐเปตฺวา วินิจฺฉโย น ทาตพฺโพ, โย น ทเทยฺย, อาปตฺติ ทุกฺกฏสฺสาติ.}

\proseitem{413}{\csromanfolio{455}\csromansymbolfootnote{*}{วิ 2. 413 ปิฏฺเฐปิ.}เตน โข ปน สมเยน ภิกฺขุนิโย โอวาทํ น คจฺฉนฺติ.\csromanspacer{} ภควโต เอตมตฺถํ อาโรเจสุํ.\csromanspacer{} น ภิกฺขเว ภิกฺขุนิยา โอวาโท น คนฺตพฺโพ, ยา น คจฺเฉยฺย, ยถาธมฺโม กาเรตพฺโพติ.}

\prose{เตน โข ปน สมเยน สพฺโพ ภิกฺขุนิสํโฆ โอวาทํ คจฺฉติ.\csromanspacer{} มนุสฺสา อุชฺฌายนฺติ ขิยฺยนฺติ วิปาเจนฺติ “ชายาโย อิมา อิเมสํ, ชาริโย อิมา อิเมสํ, อิทานิ อิเม อิมาหิ สทฺธิํ อภิรมิสฺสนฺตี”ติ.\csromanspacer{} ภควโต เอตมตฺถํ อาโรเจสุํ.\csromanspacer{} น ภิกฺขเว สพฺเพน ภิกฺขุนิสํเฆน โอวาโท คนฺตพฺโพ, คจฺเฉยฺย เจ, อาปตฺติ ทุกฺกฏสฺส.\csromanspacer{} อนุชานามิ ภิกฺขเว จตูหิ ปญฺจหิ ภิกฺขุนีหิ โอวาทํ คนฺตุนฺติ.}

\prose{เตน โข ปน สมเยน จตสฺโส ปญฺจ ภิกฺขุนิโย โอวาทํ คจฺฉนฺติ.\csromanspacer{} มนุสฺสา ตเถว อุชฺฌายนฺติ ขิยฺยนฺติ วิปาเจนฺติ “ชายาโย อิมา อิเมสํ, ชาริโย อิมา อิเมสํ, อิทานิ อิเม อิมาหิ สทฺธิํ อภิรมิสฺสนฺตี”ติ.\csromanspacer{} ภควโต เอตมตฺถํ อาโรเจสุํ.\csromanspacer{} น ภิกฺขเว จตูหิ ปญฺจหิ ภิกฺขุนีหิ โอวาโท คนฺตพฺโพ, คจฺเฉยฺยุํ เจ, อาปตฺติ ทุกฺกฏสฺส.\csromanspacer{} อนุชานามิ ภิกฺขเว เทฺว ติสฺโส ภิกฺขุนิโย\footnote{ภิกฺขุนีหิ (ก)} โอวาทํ คนฺตุํ, เอกํ ภิกฺขุํ อุปสงฺกมิตฺวา เอกํสํ อุตฺตราสงฺคํ กริตฺวา ปาเท วนฺทิตฺวา อุกฺกุฏิกํ นิสีทิตฺวา อญฺชลิํ ปคฺคเหตฺวา เอวมสฺส วจนีโย “ภิกฺขุนิสํโฆ อยฺย ภิกฺขุสํฆสฺส ปาเท วนฺทติ, โอวทูปสงฺกมนญฺจ ยาจติ, ลภตุ กิร อยฺย ภิกฺขุนิสํโฆ โอวาทูปสงฺกมนนฺ”ติ.\csromanspacer{} เตน ภิกฺขุนา ปาติโมกฺขุทฺเทสโก อุปสงฺกมิตฺวา เอวมสฺส วจนีโย “ภิกฺขุนิสํโฆ ภนฺเต ภิกฺขุสํฆสฺส ปาเท วนฺทติ, โอวาทูปสงฺกมนญฺจ ยาจติ, ลภตุ กิร ภนฺเต ภิกฺขุนิสํโฆ โอวาทูปสงฺกมนนฺ”ติ.\csromanspacer{} ปาติโมกฺขุทฺเทสเกน วตฺตพฺโพ “อตฺถิ โกจิ ภิกฺขุ ภิกฺขุโนวาทโก สมฺมโต”ติ.\csromanspacer{} สเจ โหติ โกจิ ภิกฺขุ ภิกฺขุโนวาทโก สมฺมโต, ปาติโมกฺขุทฺเทสเกน วตฺตพฺโพ “อิตฺถนฺนาโม ภิกฺขุ ภิกฺขุโนวาทโก สมฺมโต, ตํ ภิกฺขุนิสํโฆ อุปสงฺกมตู”ติ.\csromanspacer{} สเจ น โหติ โกจิ ภิกฺขุ ภิกฺขุโนวาทโก สมฺมโต, ปาติโมกฺขุทฺเทสเกน วตฺตพฺโพ “โก อายสฺมา อุสฺสหติ ภิกฺขุนิโย โอวทิตุนฺ”ติ.\csromanspacer{} สเจ โกจิ อุสฺสหติ ภิกฺขุนิโย โอวทิตุํ, โส จ โหติ อฏฺฐหงฺเคหิ สมนฺนาคโต, สมฺมนฺนิตฺวา วตฺตพฺโพ “อิตฺถนฺนาโม \csromanfolio{456}ภิกฺขุ ภิกฺขุโนวาทโก สมฺมโต, ตํ ภิกฺขุนิสํโฆ อุปสงฺกมตู”ติ.\csromanspacer{} สเจ น โกจิ อุสฺสหภิ ภิกฺขุนิโย โอวทิตุํ, ปาติโมกฺขุทฺเทสเกน วตฺตพฺโพ “นตฺถิ โกจิ ภิกฺขุ ภิกฺขุโนวาทโก สมฺมโต, ปาสาทิเกน ภิกฺขุนิสํโฆ สมฺปาเทตู”ติ.}

\proseitem{414}{เตน โข ปน สมเยน ภิกฺขู โอวาทํ น คณฺหนฺติ.\csromanspacer{} ภควโต เอตมตฺถํ อาโรเจสุํ.\csromanspacer{} น ภิกฺขเว โอวาโท น คเหตพฺโพ, โย น คเณฺหยฺย, อาปตฺติ ทุกฺกฏสฺสาติ.}

\prose{เตน โข ปน สมเยน อญฺญตโร ภิกฺขุ พาโล โหติ, ตํ ภิกฺขุนิโย อุปสงฺกมิตฺวา เอตทโวจุํ “โอวาทํ อยฺย คณฺหาหี”ติ, อหํ หิ ภคินี พาโล, กถาหํ โอวาทํ คณฺหามีติ.\csromanspacer{} คณฺหาหยฺย โอวาทํ, เอวํ หิ ภควตา ปญฺญตฺตํ ภิกฺขูหิ ภิกฺขุนีนํ โอวาโท คเหตพฺโพ”ติ.\csromanspacer{} ภควโต เอตมตฺถํ อาโรเจสุํ.\csromanspacer{} อนุชานามิ ภิกฺขเว ฐเป-ตฺวา พาลํ อวเสเสหิ โอวาทํ คเหตุนฺติ.}

\prose{เตน โข ปน สมเยน อญฺญตโร ภิกฺขุ คิลาโน โหติ, ตํ ภิกฺขุนิโย อุปสงฺกมิตฺวา เอตทโวจุํ “โอวาทํ อยฺย คณฺหาหี”ติ.\csromanspacer{} อหํ หิ ภคินี คิลาโน, กถาหํ โอวาทํ คณฺหามีติ.\csromanspacer{} คณฺหาหยฺย โอวาทํ, เอวํ หิ ภควตา ปญฺญตฺตํ “ฐเปตฺวา พาลํ อวเสเสหิ โอวาโท คเหตพฺโพ”ติ.\csromanspacer{} ภควโต เอตมตฺถํ อาโรเจสุํ.\csromanspacer{} อนุชานามิ ภิกฺขเว ฐเปตฺวา พาลํ, ฐเปตฺวา คิลานํ อวเสเสหิ โอวาทํ คเหตุนฺติ.}

\prose{เตน โข ปน สมเยน อญฺญตโร ภิกฺขุ คมิโก โหติ, ตํ ภิกฺขุนิโย อุปสงฺกมิตฺวา เอตทโวจุํ “โอวาทํ อยฺย คณฺหาหี”ติ.\csromanspacer{} อหํ หิ ภคินี คมิโก, กถาหํ โอวาทํ คณฺหามีติ.\csromanspacer{} คณฺหาหยฺย โอวาทํ, เอวํ หิ ภควตา ปญฺญตฺตํ “ฐเปตฺวา พาลํ, ฐเปตฺวา คิลานํ อวเสเสหิ โอวาโท คเหตพฺโพ”ติ.\csromanspacer{} ภควโต เอตมตฺถํ อาโรเจสุํ.\csromanspacer{} อนุชานามิ ภิกฺขเว ฐเปตฺวา พาลํ, ฐเปตฺวา คิลานํ, ฐเปตฺวา คมิกํ อวเสเสหิ โอวาทํ คเหตุนฺติ.}

\prose{เตน โข ปน สมเยน อญฺญตโร ภิกฺขุ อรญฺเญ วิหรติ, ตํ ภิกฺขุนิโย อุปสงฺกมิตฺวา เอตทโวจุํ “โอวาทํ อยฺย คณฺหาหี”ติ.\csromanspacer{} อหํ หิ ภคินี อรญฺเญ วิหรามิ, กถาหํ โอวาทํ คณฺหามีติ. \csromanfolio{457}คณฺหาหยฺย โอวาทํ, เอวํ หิ ภควตา ปญฺญตฺตํ “ฐเปตฺวา พาลํ, ฐเปตฺวา คิลานํ, ฐเปตฺวา คมิกํ อวเสเสหิ โอวาโท คเหตพฺโพ”ติ.\csromanspacer{} ภควโต เอตมตฺถํ อาโรเจสุํ.\csromanspacer{} อนุชานามิ ภิกฺขเว อรญฺญิเกน ภิกฺขุนา โอวาทํ คเหตุํ, สงฺเกตญฺจ กาตุํ “อตฺร ปติหริสฺสามี”ติ.}

\proseitem{415}{เตน โข ปน สมเยน ภิกฺขู โอวาทํ คเหตฺวา น อาโรเจนฺติ.\csromanspacer{} ภควโต เอตมตฺถํ อาโรเจสุํ.\csromanspacer{} น ภิกฺขเว โอวาโท น อาโรเจตพฺโพ, โย น อาโรเจยฺย, อาปตฺติ ทุกฺกฏสฺสาติ.}

\prose{เตน โข ปน สมเยน ภิกฺขู โอวาทํ คเหตฺวา น ปจฺจาหรนฺติ.\csromanspacer{} ภควโต เอตมตฺถํ อาโรเจสุํ.\csromanspacer{} น ภิกฺขเว โอวาโท น ปจฺจาหริตพฺโพ, โย น ปจฺจาหเรยฺย, อาปตฺติ ทุกฺกฏสฺสาติ.}

\prose{เตน โข ปน สมเยน ภิกฺขุนิโย สงฺเกตํ น คจฺฉนฺติ.\csromanspacer{} ภควโต เอตมตฺถํ อาโรเจสุํ.\csromanspacer{} น ภิกฺขเว ภิกฺขุนิยา สงฺเกตํ น คนฺตพฺพํ, ยา น คจฺเฉยฺย, อาปตฺติ ทุกฺกฏสฺสาติ.}

\proseitem{416}{เตน โข ปน สมเยน ภิกฺขุนิโย ทีฆานิ กายพนฺธนานิ ธาเรนฺติ, เตเหว ผาสุกา นาเมนฺติ.\csromanspacer{} มนุสฺสา อุชฺฌายนฺติ ขิยฺยนฺติ วิปาเจนฺติ “ฯเปฯ เสยฺยถาปิ คิหินี กามโภคินิโย”ติ.\csromanspacer{} ภควโต เอตมตฺถํ อาโรเจสุํ.\csromanspacer{} น ภิกฺขเว ภิกฺขุนิยา ทีฆํ กายพนฺธนํ ธาเรตพฺพํ, ยา ธาเรยฺย, อาปตฺติ ทุกฺกฏสฺส.\csromanspacer{} อนุชานามิ ภิกฺขเว ภิกฺขุนิยา เอกปริยากตํ\footnote{เอกปริยายกตํ (สฺยา)} กายพนฺธนํ, น จ เตน ผาสุกา นาเมตพฺพา, ยา นาเมยฺย, อาปตฺติ ทุกฺกฏสฺสาติ.}

\prose{เตน โข ปน สมเยน ภิกฺขุนิโย วิลีเวน\footnote{วิลิเวน(ก)} ปฏฺเฏน ผาสุกา นาเมนฺติ ฯเปฯ จมฺมปฏฺเฏน ผาสุกา นาเมนฺติ.\csromanspacer{} ทุสฺสปฏฺเฏน ผาสุกา นาเมนฺติ.\csromanspacer{} ทุสฺสเวณิยา ผาสุกา นาเมนฺติ.\csromanspacer{} ทุสฺสวฏฺฏิยา ผาสุกา นาเมนฺติ.\csromanspacer{} โจฬปฏฺเฏน ผาสุกา นาเมนฺติ.\csromanspacer{} โจฬเวณิยา ผาสุกา นาเมนฺติ.\csromanspacer{} โจฬวฏฺฏิยา ผาสุกา นาเมนฺติ.\csromanspacer{} สุตฺตเวณิยา ผาสุกา นาเมนฺติ.\csromanspacer{} สุตฺตวฏฺฏิยา ผาสุกา นาเมนฺติ.\csromanspacer{} มนุสฺสา อุชฺฌายนฺติ ขิยฺยนฺติ วิปาเจนฺติ “ฯเปฯ เสยฺยถาปิ คิหินี กามโภคินิโย”ติ.\csromanspacer{} ภควโต เอตมตฺถํ \csromanfolio{458}อาโรเจสุํ.\csromanspacer{} น ภิกฺขเว ภิกฺขุนิยา วิลีเวน ปฏฺเฏน ผาสุกา นาเมตพฺพา ฯเปฯ.\csromanspacer{} น สุตฺตวฏฺฏิยา ผาสุกา นาเมตพฺพา, ยา นาเมยฺย อาปตฺติ ทุกฺกฏสฺสาติ.}

\prose{เตน โข ปน สมเยน ภิกฺขุนิโย อฏฺฐิลฺเลน ชฆนํฆํสาเปนฺติ ฯเปฯ โคหนุเกน ชฆนํ โกฏฺฏาเปนฺติ.\csromanspacer{} หตฺถํ โกฏฺฏาเปนฺติ.\csromanspacer{} หตฺถโกจฺฉํ โกฏฺฏาเปนฺติ, ปาทํ โกฏฺฏาเปนฺติ.\csromanspacer{} ปาทโกจฺฉํ โกฏฺฏาเปนฺติ.\csromanspacer{} อูรุํ โกฏฺฏาเปนฺติ.\csromanspacer{} มุขํ โกฏฺฏาเปนฺติ.\csromanspacer{} ทนฺตมํสํ โกฏฺฏาเปนฺติ.\csromanspacer{} มนุสฺสา อุชฺฌายนฺติ ขิยฺยนฺติ วิปาเจนฺติ “ฯเปฯ เสยฺยถาปิ คิหินี กามโภคินิโย”ติ.\csromanspacer{} ภควโต เอตมตฺถํ อาโรเจสุํ.\csromanspacer{} น ภิกฺขเว ภิกฺขุนิยา อฏฺฐิลฺเลน ชฆนํ ฆํสาเปตพฺพํ ฯเปฯ.\csromanspacer{} น ทนฺตมํสํ โกฏฺฏาเปตพฺพํ, ยา โกฏฺฏาเปยฺย, อาปตฺติ ทุกฺกฏสฺสาติ.}

\proseitem{417}{เตน โข ปน สมเยน ฉพฺพคฺคิยา ภิกฺขุนิโย มุขํ อาลิมฺปนฺติ ฯเปฯ.\csromanspacer{} มุขํ อุมฺมทฺเทนฺติ.\csromanspacer{} มุขํ จุณฺเณนฺติ.\csromanspacer{} มโนสิลิกาย มุขํ ลญฺเฉนฺติ.\csromanspacer{} องฺคราคํ กโรนฺติ.\csromanspacer{} มุขราคํ กโรนฺติ.\csromanspacer{} องฺคราคมุขราคํ กโรนฺติ.\csromanspacer{} มนุสฺสา อุชฺฌายนฺติ ขิยฺยนฺติ วิปาเจนฺติ “ฯเปฯ เสยฺยถาปิคิหินี กามโภคินิโย”ติ.\csromanspacer{} ภควโต เอตมตฺถํ อาโรเจสุํ.\csromanspacer{} น ภิกฺขเว ภิกฺขุนิยา มุขํ อาลิมฺปิตพฺพํ ฯเปฯ.\csromanspacer{} น มุขํ อุมฺมทฺทิตพฺพํ.\csromanspacer{} น มุขํ จุณฺเณตพฺพํ.\csromanspacer{} น มโนสิลิกาย มุขํ ลญฺฉิตพฺพํ.\csromanspacer{} น องฺคราโค กาตพฺโพ.\csromanspacer{} น มุขราโค กาตพฺโพ.\csromanspacer{} น องฺคราคมุขราโค กาตพฺโพ, ยา กเรยฺย, อาปตฺติ ทุกฺกฏสฺสาติ.}

\prose{เตน โข ปน สมเยน ฉพฺพคฺคิยา ภิกฺขุนิโย อวงฺคํ\footnote{อปางฺคํ (?)} กโรนฺติ ฯเปฯ.\csromanspacer{} วิเสสกํ กโรนฺติ.\csromanspacer{} โอโลกนเกน โอโลเกนฺติ.\csromanspacer{} สาโลเก ติฏฺฐนฺติ.\csromanspacer{} นจฺจํ\footnote{สนจฺจํ (สี, สฺยา), สมชฺชํ (ก)} การาเปนฺติ.\csromanspacer{} เวสิํ วุฏฺฐาเปนฺติ.\csromanspacer{} ปานาคารํ ฐเปนฺติ.\csromanspacer{} สูนํ ฐเปนฺติ.\csromanspacer{} อาปณํ ปสาเรนฺติ.\csromanspacer{} วฑฺฒิํ ปโยเชนฺติ.\csromanspacer{} วณิชฺชํ ปโยเชนฺติ.\csromanspacer{} ทาสํ อุปฏฺฐาเปนฺติ.\csromanspacer{} ทาสิํ อุปฏฺฐาเปนฺติ.\csromanspacer{} กมฺมการํ อุปฏฺฐาเปนฺติ.\csromanspacer{} กมฺมการิํ อุปฏฺฐาเปนฺติ.\csromanspacer{} ติรจฺฉานคตํ อุปฏฺฐาเปนฺติ.\csromanspacer{} หรีตกปตฺติกํ\footnote{หรีตกปณฺณิกํ (ก)} ปกิณนฺติ.\csromanspacer{} นมตกํ ธาเรนฺติ.\csromanspacer{} มนุสฺสา อุชฺฌายนฺติ ขิยฺยนฺติ วิปาเจนฺติ “ฯเปฯ เสยฺยถาปิ คิหินี กามโภคินิโย”ติ.\csromanspacer{} ภควโต เอตมตฺถํ อาโรเจสุํ.\csromanspacer{} น ภิกฺขเว ภิกฺขุนิยา อวงฺคํ กาตพฺพํ ฯเปฯ.\csromanspacer{} น วิเสสกํ กาตพฺพํ.\csromanspacer{} น \csromanfolio{459}โอโลกนเกน โอโลเกตพฺพํ.\csromanspacer{} น สาโลเก ฐาตพฺพํ.\csromanspacer{} น นจฺจํ การาเปตพฺพํ.\csromanspacer{} น เวสี วุฏฺฐาเปตพฺพา.\csromanspacer{} น ปานาคารํ ฐเปตพฺพํ.\csromanspacer{} น สูนา ฐเปตพฺพา.\csromanspacer{} น อาปโณ ปสาเรตพฺโพ.\csromanspacer{} น วฑฺฒิ ปโยเชตพฺพา.\csromanspacer{} น วณิชฺชา ปโยเชตพฺพา.\csromanspacer{} น ทาโส อุปฏฺฐาเปตพฺโพ.\csromanspacer{} น ทาสี อุปฏฺฐาเปตพฺพา.\csromanspacer{} น กมฺมกาโร อุปฏฺฐาเปตพฺโพ.\csromanspacer{} น กมฺมการี อุปฏฺฐาเปตพฺพา.\csromanspacer{} น ติรจฺฉานคโต อุปฏฺฐาเปตพฺโพ.\csromanspacer{} น หรีตกปกฺกิกํ ปกิณิตพฺพํ.\csromanspacer{} น นมตกํ ธาเรตพฺพํ, ยา ธาเรยฺย, อาปตฺติ ทุกฺกฏสฺสาติ.}

\proseitem{418}{เตน โข ปน สมเยน ฉพฺพคฺคิยา ภิกฺขุนิโย สพฺพนีลกานิ จีวรานิ ธาเรนฺติ ฯเปฯ.\csromanspacer{} สพฺพปีตกานิ จีวรานิ ธาเรนฺติ.\csromanspacer{} สพฺพโลหิตกานิ จีวรานิ ธาเรนฺติ.\csromanspacer{} สพฺพมญฺชิฏฺฐิกานิ จีวรานิ ธาเรนฺติ.\csromanspacer{} สพฺพกณฺหานิ จีวรานิ ธาเรนฺติ.\csromanspacer{} สพฺพมหารงฺครตฺตานิ จีวรานิ ธาเรนฺติ.\csromanspacer{} สพฺพมหานามรตฺตานิ จีวรานิ ธาเรนฺติ.\csromanspacer{} อจฺฉินฺนทสานิ จีวรานิ ธาเรนฺติ.\csromanspacer{} ทีฆทสานิ จีวรานิ ธาเรนฺติ.\csromanspacer{} ปุปฺผทสานิ จีวรานิ ธาเรนฺติ.\csromanspacer{} ผลทสานิ จีวรานิ ธาเรนฺติ.\csromanspacer{} กญฺจุกํ ธาเรนฺติ.\csromanspacer{} ติรีฏกํ ธาเรนฺติ.\csromanspacer{} มนุสฺสา อุชฺฌายนฺติ ขิยฺยนฺติ วิปาเจนฺติ “ฯเปฯ เสยฺยถาปิ คิหินี กามโภคินิโย”ติ.\csromanspacer{} ภควโต เอตมตฺถํ อาโรเจสุํ.\csromanspacer{} น ภิกฺขเว ภิกฺขุนิยา สพฺพนีลกานิ จีวรานิ ธาเรตพฺพานิ ฯเปฯ.\csromanspacer{} น ติรีฏกํ ธาเรตพฺพํ, ยา ธาเรยฺย, อาปตฺติ ทุกฺกฏสฺสาติ.}

\proseitem{419}{เตน โข ปน สมเยน อญฺญตรา ภิกฺขุนี กาลํ กโรนฺตี เอวมาห “มมจฺจเยน มยฺหํ ปริกฺขาโร สํฆสฺส โหตู”ติ.\csromanspacer{} ตตฺถ ภิกฺขู จ ภิกฺขุนิโย จ วิวทนฺติ “อมฺหากํ โหติ, อมฺหากํ โหตี”ติ.\csromanspacer{} ภควโต เอตมตฺถํ อาโรเจสุํ.\csromanspacer{} ภิกฺขุนี เจ ภิกฺขเว กาลํ กโรนฺตี เอวํ วเทยฺย “มมจฺจเยน มยฺหํ ปริกฺขาโร สํฆสฺส โหตู”ติ, อนิสฺสโร ตตฺถ ภิกฺขุสํโฆ, ภิกฺขุนิสํฆสฺเสเวตํ.\csromanspacer{} สิกฺขมานา เจ ภิกฺขเว ฯเปฯ.\csromanspacer{} สามเณรี เจ ภิกฺขเว กาลํ กโรนฺตี เอวํ วเทยฺย “มมจฺจเยน มยฺหํ ปริกฺขาโร สํฆสฺส โหตู”ติ, อนิสฺสโร ตตฺถ ภิกฺขุสํโฆ, ภิกฺขุนิสํฆสฺเสเวตํ.\csromanspacer{} ภิกฺขุ เจ ภิกฺขเว กาลํ กโรนฺโต เอวํ วเทยฺย “มมจฺจเยน มยฺหํ ปริกฺขาโร สํฆสฺส โหตู”ติ, อนิสฺสโร ตตฺถ ภิกฺขุนิสํโฆ, ภิกฺขุสํฆสฺเสเวตํ.\csromanspacer{} สามเณโร เจ ภิกฺขเว ฯเปฯ.\csromanspacer{} อุปาสโก เจ ภิกฺขเว ฯเปฯ.\csromanspacer{} อุปาสิกา เจ ภิกฺขเว ฯเปฯ.\csromanspacer{} อญฺโญ เจ ภิกฺขเว โกจิ กาลํ กโรนฺโต เอวํ วเทยฺย “มมจฺจเยน มยฺหํ \csromanfolio{460}ปริกฺขาโร สํฆสฺส โหตู”ติ, อนิสฺสโร ตตฺถ ภิกฺขุนิสํโฆ, ภิกฺขุสํฆสฺเสเวตนฺติ.}

\proseitem{420}{เตน โข ปน สมเยน อญฺญตรา อิตฺถี ปุราณมลฺลี ภิกฺขุนีสุ ปพฺพชิตา โหติ, สา รถิกาย ทุพฺพลกํ ภิกฺขุํ ปสฺสิตฺวา อํสกูเฏน ปหารํ ทตฺวา ปาเตสิ\footnote{ปวฏฺเฏสิ (สี)}.\csromanspacer{} ภิกฺขู อุชฺฌายนฺติ ขิยฺยนฺติ วิปาเจนฺติ “กถํ หิ นาม ภิกฺขุนี ภิกฺขุสฺส ปหารํ ทสฺสตี”ติ.\csromanspacer{} ภควโต เอตมตฺถํ อาโรเจสุํ.\csromanspacer{} น ภิกฺขเว ภิกฺขุนิยา ภิกฺขุสฺส ปหาโร ทาตพฺโพ, ยา ทเทยฺย, อาปตฺติ ทุกฺกฏสฺส.\csromanspacer{} อนุชานามิ ภิกฺขเว ภิกฺขุนิยา ภิกฺขุํ ปสฺสิตฺวา ทูรโตว โอกฺกมิตฺวา มคฺคํ ทาตุนฺติ.}

\prose{เตน โข ปน สมเยน อญฺญตรา อิตฺถี ปวุตฺถปติกา ชาเรน คพฺภินี โหติ, สา คพฺภํ ปาเตตฺวา กุลูปิกํ ภิกฺขุนิํ เอตทโวจ “หนฺทยฺเย อิมํ คพฺภํ ปตฺเตน นีหรา”ติ.\csromanspacer{} อถ โข สา ภิกฺขุนี ตํ คพฺภํ ปตฺเต ปกฺขิปิตฺวา สํฆาฏิยา ปฏิจฺฉาเทตฺวา อคมาสิ.\csromanspacer{} เตน โข ปน สมเยน อญฺญตเรน ปิณฺฑจาริเกน ภิกฺขุนา สมาทานํ กตํ โหติ “ยาหํ ปฐมํ ภิกฺขํ ลภิสฺสามิ, น ตํ อทตฺวา ภิกฺขุสฺส วา ภิกฺขุนิยา วา ปริภุญฺชิสฺสามี”ติ.\csromanspacer{} อถ โข โส ภิกฺขุ ตํ ภิกฺขุนิํ ปสฺสิตฺวา เอตทโวจ “หนฺท ภคินิ ภิกฺขํ ปฏิคฺคณฺหา”ติ.\csromanspacer{} อลํ อยฺยาติ.\csromanspacer{} ทุติยมฺปิ โข ฯเปฯ.\csromanspacer{} ตติยมฺปิ โข โส ภิกฺขุ ตํ ภิกฺขุนิํ เอตทโวจ “หนฺท ภคินิ ภิกฺขํ ปฏิคฺคณฺหา”ติ.\csromanspacer{} อลํ อยฺยาติ.\csromanspacer{} มยา โข ภคินิ สมาทานํ กตํ “ยาหํ ปฐมํ ภิกฺขํ ลภิสฺสามิ, น ตํ อทตฺวา ภิกฺขุสฺส วา ภิกฺขุนิยา วา ปริภุญฺชิสฺสามี”ติ.\csromanspacer{} หนฺท ภคินิ ภิกฺขํ ปฏิคฺคณฺหาติ.\csromanspacer{} อถ โข สา ภิกฺขุนี เตน ภิกฺขุนา นิปฺปีฬิยมานา นีหริตฺวา ปตฺตํ ทสฺเสสิ “ปสฺส อยฺย ปตฺเต คพฺภํ, มา จ กสฺสจิ อาโรเจสี”ติ.\csromanspacer{} อถ โข โส ภิกฺขุ อุชฺฌายติ ขิยฺยติ วิปาเจติ “กถํ หิ นาม ภิกฺขุนี ปตฺเตน คพฺภํ นีหริสฺสตี”ติ.\csromanspacer{} อถ โข โส ภิกฺขุ ภิกฺขูนํ เอตมตฺถํ อาโรเจสิ.\csromanspacer{} เย เต ภิกฺขู อปฺปิจฺฉา ฯเปฯ เต อุชฺฌายนฺติ ขิยฺยนฺติ วิปาเจนฺติ “กถํ หิ นาม ภิกฺขุนี ปตฺเตน คพฺภํ นีหริสฺสตี”ติ.\csromanspacer{} ภควโต เอตมตฺถํ อาโรเจสุํ.\csromanspacer{} น ภิกฺขเว ภิกฺขุนิยา ปตฺเตน คพฺโภ นีหริตพฺโพ, \csromanfolio{461}ยา นีหเรยฺย, อาปตฺติ ทุกฺกฏสฺส.\csromanspacer{} อนุชานามิ ภิกฺขเว ภิกฺขุนิยา ภิกฺขุํ ปสฺสิตฺวา นีหริตฺวา ปตฺตํ ทสฺเสตุนฺติ.}

\prose{เตน โข ปน สมเยน ฉพฺพคฺคิยา ภิกฺขุนิโย ภิกฺขุํ ปสฺสิตฺวา ปริวตฺเตตฺวา ปตฺตมูลํ ทสฺเสนฺติ.\csromanspacer{} ภิกฺขู อุชฺฌายนฺติ ขิยฺยนฺติ วิปาเจนฺติ “กถํ หิ นาม ฉพฺพคฺคิยา ภิกฺขุนิโย ภิกฺขุํ ปสฺสิตฺวา ปริวตฺเตตฺวา ปตฺตมูลํ ทสฺเสสฺสนฺตี”ติ.\csromanspacer{} อถ โข เต ภิกฺขู ภควโต เอตมตฺถํ อาโรเจสุํ.\csromanspacer{} น ภิกฺขเว ภิกฺขุนิยา ภิกฺขุํ ปสฺสิตฺวา ปริวตฺเตตฺวา ปตฺตมูลํ ทสฺเสตพฺพํ, ยา ทสฺเสยฺย, อาปตฺติ ทุกฺกฏสฺส.\csromanspacer{} อนุชานามิ ภิกฺขเว ภิกฺขุนิยา ภิกฺขุํ ปสฺสิตฺวา อุกฺกุชฺชิตฺวา ปตฺตํ ทสฺเสตุํ, ยญฺจ ปตฺเต อามิสํ โหติ, เตน จ ภิกฺขุ นิมนฺเตตพฺโพติ.}

\prose{เตน โข ปน สมเยน สาวตฺถิยํ รถิกาย ปุริสพฺยญฺชนํ ฉฑฺฑิตํ โหติ, ตํ ภิกฺขุนิโย สกฺกจฺจํ อุปนิชฺฌายิํสุ, มนุสฺสา อุกฺกุฏฺฐิํ อกํสุ, ตา ภิกฺขุนิโย มงฺกู อเหสุํ.\csromanspacer{} อถ โข ตา ภิกฺขุนิโย อุปสฺสยํ คนฺตฺวา ภิกฺขุนีนํ เอตมตฺถํ อาโรเจสุํ.\csromanspacer{} ยา ตา ภิกฺขุนิโย อปฺปิจฺฉา ฯเปฯ ตา อุชฺฌายนฺติ ขิยฺยนฺติ วิปาเจนฺติ “กถํ หิ นาม ภิกฺขุนิโย ปุริสพฺยญฺชนํ อุปนิชฺฌายิสฺสนฺตี”ติ.\csromanspacer{} อถ โข ตา ภิกฺขุนิโย ภิกฺขูนํ เอตมตฺถํ อาโรเจสุํ.\csromanspacer{} ภิกฺขู ภควโต เอตมตฺถํ อาโรเจสุํ.\csromanspacer{} น ภิกฺขเว ภิกฺขุนิยา ปุริสพฺยญฺชนํ อุปนิชฺฌายิตพฺพํ, ยา อุปนิชฺฌาเยยฺย, อาปตฺติ ทุกฺกฏสฺสาติ.}

\proseitem{421}{เตน โข ปน สมเยน มนุสฺสา ภิกฺขูนํ อามิสํ เทนฺติ, ภิกฺขู ภิกฺขุนีนํ เทนฺติ.\csromanspacer{} มนุสฺสา อุชฺฌายนฺติ ขิยฺยนฺติ วิปาเจนฺติ “กถํ หิ นาม ภทนฺตา อตฺตโน ปริโภคตฺถาย ทินฺนํ อญฺเญสํ ทสฺสนฺติ, มยมฺปิ น ชานาม ทานํ ทาตุนฺ”ติ.\csromanspacer{} ภควโต เอตมตฺถํ อาโรเจสุํ.\csromanspacer{} น ภิกฺขเว อตฺตโน ปริโภคตฺถาย ทินฺนํ อญฺเญสํ ทาตพฺพํ, โย ทเทยฺย, อาปตฺติ ทุกฺกฏสฺสาติ.}

\prose{เตน โข ปน สมเยน ภิกฺขูนํ อามิสํ อุสฺสนฺนํ โหติ.\csromanspacer{} ภควโต เอตมตฺถํ อาโรเจสุํ.\csromanspacer{} อนุชานามิ ภิกฺขเว สํฆสฺส ทาตุนฺติ.\csromanspacer{} พาฬฺหตรํ อุสฺสนฺนํ โหติ.\csromanspacer{} ภควโต เอตมตฺถํ อาโรเจสุํ.\csromanspacer{} อนุชานามิ ภิกฺขเว ปุคฺคลิกมฺปิ ทาตุนฺติ.}

\prose{\csromanfolio{462}เตน โข ปน สมเยน ภิกฺขูนํ สนฺนิธิกตํ อามิสํ อุสฺสนฺนํ โหติ.\csromanspacer{} ภควโต เอตมตฺถํ อาโรเจสุํ.\csromanspacer{} อนุชานามิ ภิกฺขเว ภิกฺขูนํ สนฺนิธิํ ภิกฺขุนีหิ\footnote{ภิกฺขุนีหิ ภิกฺขูหิ (สี)} ปฏิคฺคาหาเปตฺวา\footnote{ปฏิคฺคหาเปตฺวา (ก)} ปริภุญฺชิตุนฺติ.}

\prose{เตน โข ปน สมเยน มนุสฺสา ภิกฺขุนีนํ อามิสํ เทนฺติ, ภิกฺขุนิโย ภิกฺขูนํ เทนฺติ.\csromanspacer{} มนุสฺสา อุชฺฌายนฺติ ขิยฺยนฺติ วิปาเจนฺติ “กถํ หิ นาม ภิกฺขุนิโย อตฺตโน ปริโภคตฺถาย ทินฺนํ อญฺเญสํ ทสฺสนฺติ, มยมฺปิ น ชานาม ทานํ ทาตุนฺ”ติ.\csromanspacer{} ภควโต เอตมตฺถํ อาโรเจสุํ.\csromanspacer{} น ภิกฺขเว ภิกฺขุนิยา อตฺตโน ปริโภคตฺถาย ทินฺนํ อญฺเญสํ ทาตพฺพํ, ยา ทเทยฺย, อาปตฺติ ทุกฺกฏสฺสาติ.}

\prose{เตน โข ปน สมเยน ภิกฺขุนีนํ อามิสํ อุสฺสนฺนํ โหติ.\csromanspacer{} ภควโต เอตมตฺถํ อาโรเจสุํ.\csromanspacer{} อนุชานามิ ภิกฺขเว สํฆสฺส ทาตุนฺติ.\csromanspacer{} พาฬฺหตรํ อุสฺสนฺนํ โหติ.\csromanspacer{} ภควโต เอตมตฺถํ อาโรเจสุํ.\csromanspacer{} อนุชานามิ ภิกฺขเว ปุคฺคลิกมฺปิ ทาตุนฺติ.}

\prose{เตน โข ปน สมเยน ภิกฺขุนีนํ สนฺนิธิกตํ อามิสํ อุสฺสนฺนํ โหติ.\csromanspacer{} ภควโต เอตมตฺถํ อาโรเจสุํ.\csromanspacer{} อนุชานามิ ภิกฺขเว ภิกฺขุนีนํ สนฺนิธิํ ภิกฺขูหิ\footnote{ภิกฺขูหิ ภิกฺขุนีหิ (สี)} ปฏิคฺคาหาเปตฺวา ปริภุญฺชิตุนฺติ.}

\proseitem{442}{เตน โข ปน สมเยน ภิกฺขูนํ เสนาสนํ อุสฺสนฺนํ โหติ, ภิกฺขุนีนํ\footnote{ภิกฺขุนีนํ เสนาสนํ (สฺยา, กํ)} น โหติ.\csromanspacer{} ภิกฺขุนิโย ภิกฺขูนํ สนฺติเก ทูตํ ปาเหสุํ “สาธุ ภนฺเต, อยฺยา อมฺหากํ เสนาสนํ เทนฺตุ ตาวกาลิกนฺ”ติ.\csromanspacer{} ภควโต เอตมตฺถํ อาโรเจสุํ.\csromanspacer{} อนุชานามิ ภิกฺขเว ภิกฺขุนีนํ เสนาสนํ ทาตุํ ตาวกาลิกนฺติ.}

\prose{เตน โข ปน สมเยน อุตุนิโย ภิกฺขุนิโย โอนทฺธมญฺจํ โอนทฺธปีฐํ อภินิสีทนฺติปิ อภินิปชฺชนฺติปิ, เสนาสนํ โลหิเตน มกฺขิยฺยติ.\csromanspacer{} ภควโต เอตมตฺถํ อาโรเจสุํ.\csromanspacer{} น ภิกฺขเว ภิกฺขุนิยา โอนทฺธมญฺจํ โอนทฺธปีฐํ อภินิสีทิตพฺพํ อภินิปชฺชิตพฺพํ, ยา อภินิสีเทยฺย วา อภินิปชฺเชยฺย วา, อาปตฺติ ทุกฺกฏสฺส.\csromanspacer{} อนุชานามิ ภิกฺขเว อาวสถจีวรนฺติ.\csromanspacer{} อาวสถจีวรํ โลหิเตน มกฺขิยฺยติ.\csromanspacer{} ภควโต เอตมตฺถํ อาโรเจสุํ. \csromanfolio{463}อนุชานามิ ภิกฺขเว อาณิโจฬกนฺติ.\csromanspacer{} โจฬกํ นิปตติ.\csromanspacer{} ภควโต เอตมตฺถํ อาโรเจสุํ.\csromanspacer{} อนุชานามิ ภิกฺขเว สุตฺตเกน พนฺธิตฺวา อูรุยา พนฺธิตุนฺติ.\csromanspacer{} สุตฺตํ ฉิชฺชติ.\csromanspacer{} ภควโต เอตมตฺถํ อาโรเจสุํ.\csromanspacer{} อนุชานามิ ภิกฺขเว สํเวลฺลิยํ กฏิสุตฺตกนฺติ.}

\prose{เตน โข ปน สมเยน ฉพฺพคฺคิยา ภิกฺขุนิโย สพฺพกาลํ กฏิสุตฺตกํ ธาเรนฺติ.\csromanspacer{} มนุสฺสา อุชฺฌายนฺติ ขิยฺยนฺติ วิปาเจนฺติ “เสยฺยถาปิ คิหินี กามโภคินิโย”ติ.\csromanspacer{} ภควโต เอตมตฺถํ อาโรเจสุํ.\csromanspacer{} น ภิกฺขเว ภิกฺขุนิยา สพฺพกาลํ กฏิสุตฺตกํ ธาเรตพฺพํ, ยา ธาเรยฺย, อาปตฺติ ทุกฺกฏสฺส.\csromanspacer{} อนุชานามิ ภิกฺขเว อุตุนิยา กฏิสุตฺตกนฺติ.}

\nitthitam{ทุติยภาณวาโร นิฏฺฐิโต.}
\csromanmatikaanchor{mtk.210}
\csromantocmarkref{section}{3. ตติยภาณวาร}{mtk.210}
\bookmark[dest={mtk.210},level=1]{3. ตติยภาณวาร}
\csromanheader{1}{3. ตติยภาณวาร}
\proseitem{423}{เตน โข ปน สมเยน อุปสมฺปนฺนาโย ทิสฺสนฺติ อนิมิตฺตาปิ นิมิตฺตมตฺตาปิ อโลหิตาปิ ธุวโลหิตาปิ ธุวโจฬาปิ ปคฺฆรนฺตีปิ สิขรณีปิ อิตฺถิปณฺฑกาปิ เวปุริสิกาปิ สมฺภินฺนาปิ อุภโตพฺยญฺชนาปิ.\csromanspacer{} ภควโต เอตมตฺถํ อาโรเจสุํ.\csromanspacer{} อนุชานามิ ภิกฺขเว อุปสมฺปาเทนฺติยา จตุวีสติ อนฺตรายิเก ธมฺเม ปุจฺฉิตุํ.\csromanspacer{} เอวญฺจ ปน ภิกฺขเว ปุจฺฉิตพฺพา “นสิ อนิมิตฺตา, นสิ นิมิตฺตมตฺตา, นสิ อโลหิตา, นสิ ธุวโลหิตา, นสิ ธุวโจฬา, นสิ ปคฺฆรนฺตี, นสิ สิขรณี\footnote{สิขริณี (สี, สฺยา)}, นสิ อิตฺถิปณฺฑกา, นสิ เวปุริสิกา, นสิ สมฺภินฺนา, นสิ อุภโตพฺยญฺชนา.\csromanspacer{} สนฺติ เต เอวรูปา อาพาธา กุฏฺฐํ คณฺโฑ กิลาโส โสโส อปมาโร, มนุสฺสาสิ, อิตฺถีสิ, ภุชิสฺสาสิ, อณณาสิ, นสิ ราชภฏี, อนุญฺญาตาสิ มาตาปิตูหิ สามิเกน, ปริปุณฺณวีสติวสฺสาสิ, ปริปุณฺณํ เต ปตฺตจีวรํ, กินฺนามาสิ, กานามา เต ปวตฺตินี”ติ.}

\prose{เตน โข ปน สมเยน ภิกฺขู ภิกฺขุนีนํ อนฺตรายิเก ธมฺเม ปุจฺฉนฺติ, อุปสมฺปทาเปกฺขาโย วิตฺถายนฺติ, มงฺกู โหนฺติ, น สกฺโกนฺติ วิสฺสชฺเชตุํ.\csromanspacer{} ภควโต เอตมตฺถํ อาโรเจสุํ.\csromanspacer{} อนุชานามิ ภิกฺขเว เอกโต-อุปสมฺปนฺนาย ภิกฺขุนิสํเฆ วิสุทฺธาย ภิกฺขุสํเฆ อุปสมฺปาเทตุนฺติ.}

\prose{\csromanfolio{464}เตน โข ปน สมเยน ภิกฺขุนิโย อนนุสิฏฺฐา อุปสมฺปทาเปกฺขาโย อนฺตรายิเก ธมฺเม ปุจฺฉนฺติ, อุปสมฺปทาเปกฺขาโย วิตฺถายนฺติ, มงฺกู โหนฺติ, น สกฺโกนฺติ วิสฺสชฺเชตุํ.\csromanspacer{} ภควโต เอตมตฺถํ อาโรเจสุํ.\csromanspacer{} อนุชานามิ ภิกฺขเว ปฐมํ อนุสาสิตฺวา ปจฺฉา อนฺตรายิเก ธมฺเม ปุจฺฉิตุนฺติ.}

\prose{ตตฺเถว สํฆมชฺเฌ อนุสาสนฺติ, อุปสมฺปทาเปกฺขาโย ตเถว วิตฺถายนฺติ, มงฺกู โหนฺติ, น สกฺโกนฺติ วิสฺสชฺเชตุํ.\csromanspacer{} ภควโต เอตมตฺถํ อาโรเจสุํ.\csromanspacer{} อนุชานามิ ภิกฺขเว เอกมนฺตํ อนุสาสิตฺวา สํฆมชฺเฌ อนฺตรายิเก ธมฺเม ปุจฺฉิตุํ.\csromanspacer{} เอวญฺจ ปน ภิกฺขเว อนุสาสิตพฺพา\csromandash{}}

\proseitem{424}{ปฐมํ อุปชฺฌํ คาหาเปตพฺพา, อุปชฺฌํ คาหาเปตฺวา ปตฺตจีวรํ อาจิกฺขิตพฺพํ, อยํ เต ปตฺโต, อยํ สํฆาฏิ, อยํ อุตฺตราสงฺโค, อยํ อนฺตรวาสโก, อิทํ สํกจฺจิกํ\footnote{สงฺกจฺฉิกํ (สฺยา)}, อยํ อุทกสาฏิกา.\csromanspacer{} คจฺฉ อมุมฺหิ โอกาเส ติฏฺฐาหีติ.}

\prose{พาลา อพฺยตฺตา อนุสาสนฺติ, ทุรนุสิฏฺฐา อุปสมฺปทาเปกฺขาโย วิตฺถายนฺติ, มงฺกู โหนฺติ, น สกฺโกนฺติ วิสฺสชฺเชตุํ.\csromanspacer{} ภควโต เอตมตฺถํ อาโรเจสุํ.\csromanspacer{} น ภิกฺขเว พาลาย อพฺยตฺตาย อนุสาสิตพฺพา, ยา อนุสาเสยฺย, อาปตฺติ ทุกฺกฏสฺส.\csromanspacer{} อนุชานามิ ภิกฺขเว พฺยตฺตาย ภิกฺขุนิยา ปฏิพลาย อนุสาสิตุนฺติ.}

\prose{อสมฺมตา อนุสาสนฺติ.\csromanspacer{} ภควโต เอตมตฺถํ อาโรเจสุํ.\csromanspacer{} น ภิกฺขเว อสมฺมตาย อนุสาสิตพฺพา, ยา อนุสาเสยฺย, อาปตฺติ ทุกฺกฏสฺสา.\csromanspacer{} อนุชานามิ ภิกฺขเว สมฺมตาย อนุสาสิตุํ.\csromanspacer{} เอวญฺจ ปน ภิกฺขเว สมฺมนฺนิตพฺพา, อตฺตนา วา อตฺตานํ สมฺมนฺนิตพฺพํ, ปราย วา ปรา สมฺมนฺนิตพฺพา.}

\prose{กถญฺจ อตฺตนาว\footnote{อตฺตนา วา (ก)} อตฺตานํ สมฺมนฺนิตพฺพํ, พฺยตฺตาย ภิกฺขุนิยา ปฏิพลาย สํโฆ ญาเปตพฺโพ “สุณาตุ เม อยฺเย สํโฆ, อิตฺถนฺนามา อิตฺถนฺนามาย อยฺยาย อุปสมฺปทาเปกฺขา, ยทิ สํฆสฺส ปตฺตกลฺลํ, อหํ อิตฺถนฺนามา อิตฺถนฺนามํ อนุสาเสยฺยนฺ”ติ, เอวํ อตฺตนาว อตฺตานํ สมฺมนฺนิตพฺพํ.}

\prose{กถญฺจ ปราย\footnote{ปราย วา (ก)} ปรา สมฺมนฺนิตพฺพา, พฺยตฺตาย ภิกฺขุนิยา ปฏิพลาย สํโฆ ญาเปตพฺโพ “สุณาตุ เม อยฺเย สํโฆ, อิตฺถนฺนามา \csromanfolio{465}อิตฺถนฺนามาย อยฺยาย อุปสมฺปทาเปกฺขา, ยทิ สํฆสฺส ปตฺตกลฺลํ, อิตฺถนฺนามา อิตฺถนฺนามํ อนุสาเสยฺยา”ติ, เอวํ ปราย ปรา สมฺมนฺนิตพฺพา.}

\prose{ตาย สมฺมตาย ภิกฺขุนิยา อุปสมฺปทาเปกฺขํ อุปสงฺกมิตฺวา เอวมสฺส วจนียา “สุณสิ อิตฺถนฺนาเม อยํ เต สจฺจกาโล ภูตกาโล, ยํ ชาตํ ตํ สํฆมชฺเฌ ปุจฺฉนฺเต สนฺตํ ‘อตฺถี’ติ วตฺตพฺพํ, อสนฺตํ ‘นตฺถี’ติ วตฺตพฺพํ, มา โข วิตฺถายิ\footnote{วิตฺถาสิ (ก)}, มา โข มงฺกุ อโหสิ, เอวํ ตํ ปุจฺฉิสฺสนฺติ\csromandash{}นสิ อนิมิตฺตา, นสิ นิมิตฺตมตฺตา, นสิ อโลหิตา, นสิ ธุวโลหิตา, นสิ ธุวโจฬา, นสิ ปคฺฆรนฺตี, นสิ สิขรณี, นสิ อิตฺถิปณฺฑกา, นสิ เวปุริสิกา, นสิ สมฺภินฺนา, นสิ อุภโตพฺยญฺชนา.\csromanspacer{} สนฺติ เต เอวรูปา อาพาธา กุฏฺฐํ คณฺโฑ กิลาโส โสโส อปมาโร, มนุสฺสาสิ, อิตฺถีสิ, ภุชิสฺสาสิ, อณณาสิ, นสิ ราชภฏี, อนุญฺญาตาสิ มาตาปิตูหิ สามิเกน, ปริปุณฺณวีสติวสฺสาสิ, ปริปุณฺณํ เต ปตฺตจีวรํ กินฺนามาสิ, กานามา เต ปวตฺตินี”ติ.}

\prose{เอกโต อาคจฺฉนฺติ.\csromanspacer{} น เอกโต อาคนฺตพฺพํ, อนุสาสิกาย ปฐมตรํ อาคนฺตฺวา สํโฆ ญาเปตพฺโพ “สุณาตุ เม อยฺเย สํโฆ, อิตฺถนฺนามา อิตฺถนฺนามาย อยฺยาย อุปสมฺปทาเปกฺขา, อนุสิฏฺฐา สา มยา, ยทิ สํฆสฺส ปตฺตกลฺลํ, อิตฺถนฺนามา อาคจฺเฉยฺยา”ติ. "อาคจฺฉาหี”ติ วตฺตพฺพา.}

\prose{เอกํสํ อุตฺตราสงฺคํ การาเปตฺวา ภิกฺขุนีนํ ปาเท วนฺทาเปตฺวา อุกฺกุฏิกํ นิสีทาเปตฺวา อญฺชลิํ ปคฺคณฺหาเปตฺวา อุปสมฺปทํ ยาจาเปตพฺพา “สํฆํ อยฺเย อุปสมฺปทํ ยาจามิ, อุลฺลุมฺปตุ มํ อยฺเย สํโฆ อนุกมฺปํ อุปาทาย.\csromanspacer{} ทุติยมฺปิ อยฺเย สํฆํ อุปสมฺปทํ ยาจามิ, อุลฺลุมฺปตุ มํ อยฺเย สํโฆ อนุกมฺปํ อุปาทาย.\csromanspacer{} ตติยมฺปิ อยฺเย สํฆํ อุปสมฺปทํ ยาจามิ, อุลฺลุมฺปตุ มํ อยฺเย สํโฆ อนุกมฺปํ อุปาทายา”ติ.\csromanspacer{} พฺยตฺตาย ภิกฺขุนิยา ปฏิพลาย สํโฆ ญาเปตพฺโพ\csromandash{}}

\prose{“สุณาตุ เม อยฺเย สํโฆ, อยํ อิตฺถนฺนามา อิตฺถนฺนามาย อยฺยาย อุปสมฺปทาเปกฺขา, ยทิ สํฆสฺส ปตฺตกลฺลํ, อหํ อิตฺถนฺนามํ อนฺตรายิเก ธมฺเม ปุจฺเฉยฺยนฺ”ติ.}

\prose{\csromanfolio{466}สุณสิ อิตฺถนฺนาเม อยํ เต สจฺจกาโล ภูตกาโล, ยํชาตํ ตํ ปุจฺฉามิ, สนฺตํ ‘อตฺถี’ติ วตฺตพฺพํ, อสนฺตํ ‘นตฺถี’ติ วตฺตพฺพํ, นสิ อนิมิตฺตา ฯเปฯ กินฺนามาสิ, กานามา เต ปวตฺตินี”ติ.\csromanspacer{} พฺยตฺตาย ภิกฺขุนิยา ปฏิพลาย สํโฆ ญาเปตพฺโพ\csromandash{}}

\proseitem{425}{“สุณาตุ เม อยฺเย สํโฆ, อยํ อิตฺถนฺนามา อิตฺถนฺนามาย อยฺยาย อุปสมฺปทาเปกฺขา, ปริสุทฺธา อนฺตรายิเกหิ ธมฺเมหิ, ปริปุณฺณสฺสา ปตฺตจีวรํ, อิตฺถนฺนามา สํฆํ อุปสมฺปทํ ยาจติ อิตฺถนฺนามาย อยฺยาย ปวตฺตินิยา, ยทิ สํฆสฺส ปตฺตกลฺลํ, สํโฆ อิตฺถนฺนามํ อุปสมฺปาเทยฺย อิตฺถนฺนามาย อยฺยาย ปวตฺตินิยา, เอสา ญตฺติ.}

\prose{สุณาตุ เม อยฺเย สํโฆ, อยํ อิตฺถนฺนามา อิตฺถนฺนามาย อยฺยาย อุปสมฺปาทาเปกฺขา, ปริสุทฺธา อนฺตรายิเกหิ ธมฺเมหิ, ปริปุณฺณสฺสา ปตฺตจีวรํ, อิตฺถนฺนามา สํฆํ อุปสมฺปทํ ยาจติ อิตฺถนฺนามาย อยฺยาย ปวตฺตินิยา, สํโฆ อิตฺถนฺนามํ อุปสมฺปาเทติ อิตฺถนฺนามาย อยฺยาย ปวตฺตินิยา ยสฺสา อยฺยาย ขมติ อิตฺถนฺนามาย อุปสมฺปทา อิตฺถนฺนามาย อยฺยาย ปวตฺตินิยา, สา ตุณฺหสฺส, ยสฺสา นกฺขมติ, สา ภาเสยฺย.}

\prose{ทุติยมฺปิ เอตมตฺถํ วทามิ ฯเปฯ.\csromanspacer{} ตติยมฺปิ เอตมตฺถํ วทามิ.\csromanspacer{} สุณาตุ เม อยฺเย สํโฆ, อยํ อิตฺถนฺนามา อิตฺถนฺนามาย อยฺยาย อุปสมฺปทาเปกฺขา, ปริสุทฺธา อนฺตรายิเกหิ ธมฺเมหิ, ปริปุณฺณสฺสา ปตฺตจีวรํ, อิตฺถนฺนามา สํฆํ อุปสมฺปทํ ยาจติ อิตฺถนฺนามาย อยฺยาย ปวตฺตินิยา, สํโฆ อิตฺถนฺนามํ อุปสมฺปาเทติ อิตฺถนฺนามาย อยฺยาย ปวตฺตินิยา, ยสฺสา อยฺยาย ขมติ อิตฺถนฺนามาย อุปสมฺปทา อิตฺถนฺนามาย อยฺยาย ปวตฺตินิยา, สา ตุณฺหสฺส, ยสฺสา นกฺขมติ, สา ภาเสยฺย.}

\prose{อุปสมฺปนฺนา สํเฆน อิตฺถนฺนามา อิตฺถนฺนามาย อยฺยาย ปวตฺตินิยา, ขมติ สํฆสฺส, ตสฺมา ตุณฺหี, เอวเมตํ ธาเรยามี”ติ.}

\prose{ตาวเทว ตํ อาทาย ภิกฺขุสํฆํ อุปสงฺกมิตฺวา เอกํสํ อุตฺตราสงฺคํ การาเปตฺวา ภิกฺขูนํ ปาเท วนฺทาเปตฺวา อุกฺกุฏิกํ นิสีทาเปตฺวา อญฺชลิํ ปคฺคณฺหาเปตฺวา อุปสมฺปทํ ยาจาเปตพฺพา “อหํ อยฺยา อิตฺถนฺนามา อิตฺถนฺนามาย อยฺยาย อุปสมฺปทาเปกฺขา เอกโต-อุปสมฺปนฺนา ภิกฺขุนิสํเฆ วิสุทฺธา, สํฆํ อยฺยา อุปสมฺปทํ ยาจามิ, อุลฺลุมฺปตุมํ อยฺยา สํโฆ อนุกมฺปํ อุปาทาย.\csromanspacer{} อหํ อยฺยา อิตฺถนฺนามา อิตฺถนฺนามาย \csromanfolio{467}อยฺยาย อุปสมฺปทาเปกฺขา เอกโต-อุปสมฺปนฺนา ภิกฺขุนิสํเฆ วิสุทฺธา, ทุติยมฺปิ อยฺยา สํฆํ อุปสมฺปทํ ยาจามิ, อุลฺลุมฺปตุ มํ อยฺยา สํโฆ อนุกมฺปํ อุปาทาย.\csromanspacer{} อหํ อยฺยา อิตฺถนฺนามา อิตฺถนฺนามาย อยฺยาย อุปสมฺปทาเปกฺขา เอกโต-อุปสมฺปนฺนา ภิกฺขุนิสํเฆ วิสุทฺธา, ตติยมฺปิ อยฺยา สํฆํ อุปสมฺปทํ ยาจามิ, อุลฺลุมฺปตุ มํ อยฺยา สํโฆ อนุกมฺปํ อุปาทายา”ติ.\csromanspacer{} พฺยตฺเตน ภิกฺขุนา ปฏิพเลน สํโฆ ญาเปตพฺโพ\csromandash{}}

\prose{“สุณาตุ เม ภนฺเต สํโฆ, อยํ อิตฺถนฺนามา อิตฺถนฺนามาย อุปสมฺปทาเปกฺขา เอกโต-อุปสมฺปนฺนา ภิกฺขุนิสํเฆ วิสุทฺธา, อิตฺถนฺนามา สํฆํ อุปสมฺปทํ ยาจติ อิตฺถนฺนามาย ปวตฺตินิยา, ยทิ สํฆสฺส ปตฺตกลฺลํ, สํโฆ อิตฺถนฺนามํ อุปสมฺปาเทยฺย อิตฺถนฺนามาย ปวตฺตินิยา, เอสา ญตฺติ.}

\prose{สุณาตุ เม ภนฺเต สํโฆ, อยํ อิตฺถนฺนามา อิตฺถนฺนามาย อุปสมฺปทาเปกฺขา เอกโต-อุปสมฺปนฺนา ภิกฺขุนิสํเฆ วิสุทฺธา, อิตฺถนฺนามา สํฆํ อุปสมฺปทํ ยาจติ อิตฺถนฺนามาย ปวตฺตินิยา, สํโฆ อิตฺถนฺนามํ อุปสมฺปาเทติ อิตฺถนฺนามาย ปวตฺตินิยา, ยสฺสายสฺมโต ขมติ อิตฺถนฺนามาย อุปสมฺปทา อิตฺถนฺนามาย ปวตฺตินิยา, โส ตุณฺหสฺส, ยสฺส นกฺขมติ, โส ภาเสยฺย.}

\prose{ทุติยมฺปิ เอตมตฺถํ วทามิ ฯเปฯ.\csromanspacer{} ตติยมฺปิ เอตมตฺถํ วทามิ.\csromanspacer{} สุณาตุ เม ภนฺเต สํโฆ, อยํ อิตฺถนฺนามา อิตฺถนฺนามาย อุปสมฺปทาเปกฺขา เอกโต-อุปสมฺปนฺนา ภิกฺขุนิสํเฆ วิสุทฺธา, อิตฺถนฺนามา สํฆํ อุปสมฺปทํ ยาจติ อิตฺถนฺนามาย ปวตฺตินิยา, สํโฆ อิตฺถนฺนามํ อุปสมฺปาเทติ อิตฺถนฺนามาย ปวตฺตินิยา, ยสฺสายสฺมโต ขมติ อิตฺถนฺนามาย อุปสมฺปทา อิตฺถนฺนามาย ปวตฺตินิยา, โส ตุณฺหสฺส, ยสฺส นกฺขมติ, โส ภาเสยฺย.}

\prose{อุปสมฺปนฺนา สํเฆน อิตฺถนฺนามา อิตฺถนฺนามาย ปวตฺตินิยา, ขมติ สํฆสฺส, ตสฺมา ตุณฺหี, เอวเมตํ ธารยามี”ติ.}

\prose{ตาวเทว ฉายา เมตพฺพา, อุตุปฺปมาณํ อาจิกฺขิตพฺพํ, ทิวสภาโค อาจิกฺขิตพฺโพ, สงฺคีติ อาจิกฺขิตพฺพา, ภิกฺขุนิโย วตฺตพฺพา “อิมิสฺสา ตโย จ นิสฺสเย อฏฺฐ จ อกรณียานิ อาจิกฺเขยฺยาถา”ติ.}

\proseitem{426}{เตน โข ปน สมเยน ภิกฺขุนิโย ภตฺตคฺเค อาสนํ สํกสายนฺติโย กาลํ วีตินาเมสุํ.\csromanspacer{} ภควโต เอตมตฺถํ \csromanfolio{468}อาโรเจสุํ.\csromanspacer{} อนุชานามิ ภิกฺขเว อฏฺฐนฺนํ ภิกฺขุนีนํ ยถาวุฑฺฒํ อวเสสานํ ยถาคติกนฺติ.}

\prose{เตน โข ปน สมเยน ภิกฺขุนิโย “ภควตา อนุญฺญาตํ อฏฺฐนฺนํ ภิกฺขุนีนํ ยถาวุฑฺฒํ อวเสสานํ ยถาคติกนฺ”ติ สพฺพตฺถ อฏฺเฐว ภิกฺขุนิโย ยถาวุฑฺฒํ ปฏิพาหนฺติ อวเสสาโย ยถาคติกํ.\csromanspacer{} ภควโต เอตมตฺถํ อาโรเจสุํ.\csromanspacer{} อนุชานามิ ภิกฺขเว ภตฺตคฺเค อฏฺฐนฺนํ ภิกฺขุนีนํ ยถาวุฑฺฒํ อวเสสานํ ยถาคติกํ, อญฺญตฺถ สพฺพตฺถ ยถาวุฑฺฒํ\footnote{อญฺญตฺถ ยถาวุฑฺฒํ (สฺยา)} น ปฏิพาหิตพฺพํ, ยา ปฏิพาเหยฺย, อาปตฺติ ทุกฺกฏสฺสาติ.}

\proseitem{427}{\csromansymbolfootnote{*}{วิ 2. 412-3 ปิฏฺเฐสุปิ.}เตน โข ปน สมเยน ภิกฺขุนิโย น ปวาเรนฺติ.\csromanspacer{} ภควโต เอตมตฺถํ อาโรเจสุํ.\csromanspacer{} น ภิกฺขเว ภิกฺขุนิยา น ปวาเรตพฺพํ, ยา น ปวาเรยฺย, ยถาธมฺโม กาเรตพฺโพติ.}

\prose{เตน โข ปน สมเยน ภิกฺขุนิโย อตฺตนา ปวาเรตฺวา ภิกฺขุสํฆํ น ปเวเรนฺติ.\csromanspacer{} ภควโต เอตมตฺถํ อาโรเจสุํ.\csromanspacer{} น ภิกฺขเว ภิกฺขุนิยา อตฺตนา ปวาเรตฺวา ภิกฺขุสํโฆ น ปวาเรตพฺโพ, ยา น ปวาเรยฺย, ยถาธมฺโม กาเรตพฺโพติ.}

\prose{เตน โข ปน สมเยน ภิกฺขุนิโย ภิกฺขูหิ สทฺธิํ เอกโต ปวาเรนฺติโย โกลาหลํ อกํสุ.\csromanspacer{} ภควโต เอตมตฺถํ อาโรเจสุํ.\csromanspacer{} น ภิกฺขเว ภิกฺขุนิยา ภิกฺขูหิ สทฺธิํ เอกโต ปวาเรตพฺพํ, ยา ปวาเรยฺย, อาปตฺติ ทุกฺกฏสฺสาติ.}

\prose{เตน โข ปน สมเยน ภิกฺขุนิโย ปุเรภตฺตํ ปวาเรนฺติโย กาลํ วีตินาเมสุํ.\csromanspacer{} ภควโต เอตมตฺถํ อาโรเจสุํ.\csromanspacer{} อนุชานามิ ภิกฺขเว ปจฺฉาภตฺตํ\footnote{ภิกฺขเว ภิกฺขุนิยา ปจฺฉาภตฺตํ (สฺยา, กํ)} ปวาเรตุนฺติ.\csromanspacer{} ปจฺฉาภตฺตํ ปวาเรนฺติโย วิกาเล\footnote{วิกาลํ (ก)} อเหสุํ, ภควโต เอตมตฺถํ อาโรเจสุํ.\csromanspacer{} อนุชานามิ ภิกฺขเว อชฺชตนา ภิกฺขุนิสํฆํ\footnote{อชฺชตนา (สี), อชฺชตนาย (ก)} ปวาเรตฺวา อปรชฺชุ ภิกฺขุสํฆํ ปวาเรตุนฺติ.}

\prose{เตน โข ปน สมเยน สพฺโพ ภิกฺขุนิสํโฆ ปวาเรนฺโต โกลาหลํ อกาสิ.\csromanspacer{} ภควโต เอตมตฺถํ อาโรเจสุํ.\csromanspacer{} อนุชานามิ ภิกฺขเว เอกํ ภิกฺขุนิํ พฺยตฺตํ ปฏิพลํ สมฺมนฺนิตุํ ภิกฺขุนิสํฆสฺส อตฺถาย \csromanfolio{469}ภิกฺขุสํฆํ ปวาเรตุํ.\csromanspacer{} เอวญฺจ ปน ภิกฺขเว สมฺมนฺนิตพฺพา, ปฐมํ ภิกฺขุนี ยาจิตพฺพา, ยาจิตฺวา พฺยตฺตาย ภิกฺขุนิยา ปฏิพลาย สํโฆ ญาเปตพฺโพ\csromandash{}}

\prose{“สุณาตุ เม อยฺเย สํโฆ, ยทิ สํฆสฺส ปตฺตกลฺลํ, สํโฆ อิตฺถนฺนามํ ภิกฺขุนิํ สมฺมนฺเนยฺย ภิกฺขุนิสํฆสฺส อตฺถาย ภิกฺขุสํฆํ ปวาเรตุํ, เอสา ญตฺติ.}

\prose{สุณาตุ เม อยฺเย สํโฆ, สํโฆ อิตฺถนฺนามํ ภิกฺขุนิํ สมฺมนฺนติ ภิกฺขุนิสํฆสฺส อตฺถาย ภิกฺขุสํฆํ ปวาเรตุํ, ยสฺสา อยฺยาย ขมติ อิตฺถนฺนามาย ภิกฺขุนิยา สมฺมุติ ภิกฺขุนิสํฆสฺส อตฺถาย ภิกฺขุสํฆํ ปวาเรตุํ, สา ตุณฺหสฺส, ยสฺสา นกฺขมติ, สา ภาเสยฺย.}

\prose{สมฺมตา สํเฆน อิตฺถนฺนามา ภิกฺขุนี ภิกฺขุนิสํฆสฺส อตฺถาย ภิกฺขุสํฆํ ปวาเรตุํ, ขมติ สํฆสฺส, ตสฺมา ตุณฺหี, เอวเมตํ ธารยามี”ติ.}

\prose{ตาย สมฺมตาย ภิกฺขุนิยา ภิกฺขุนิสํฆํ อาทาย ภิกฺขุสํฆํ อุปสงฺกมิตฺวา เอกํสํ อุตฺตราสงฺคํ กริตฺวา ภิกฺขูนํ ปาเท วนฺทิตฺวา อุกฺกุฏิกํ นิสีทิตฺวา อญฺชลิํ ปคฺคเหตฺวา เอวมสฺส วจนีโย “ภิกฺขุนิสํโฆ อยฺยา ภิกฺขุสํฆํ ปวาเรติ ทิฏฺเฐน วา สุเตน วา ปริสงฺกาย วา, วทตุ\footnote{วทตุ มํ (ก)} อยฺยา ภิกฺขุสํโฆ ภิกฺขุนิสํฆํ อนุกมฺปํ อุปาทาย, ปสฺสนฺโต ปฏิกริสฺสติ, ทุติยมฺปิ อยฺยา ฯเปฯ.\csromanspacer{} ตติยมฺปิ อยฺยา ภิกฺขุนิสํโฆ ภิกฺขุสํฆํ ปวาเรติ ทิฏฺเฐน วา สุเตน วา ปริสงฺกาย วา, วทตุ อยฺยา ภิกฺขุสํโฆ ภิกฺขุนิสํฆํ อนุกมฺปํ อุปาทาย, ปสฺสนฺโต ปฏิกริสฺสตี”ติ.}

\proseitem{428}{เตน โข ปน สมเยน ภิกฺขุนิโย ภิกฺขูนํ อุโปสถํ ฐเปนฺติ.\csromanspacer{} ปวารณํ ฐเปนฺติ, สวจนียํ กโรนฺติ, อนุวาทํ ปฏฺฐเปนฺติ, โอกาสํ กาเรนฺติ, โจเทนฺติ, สาเรนฺติ.\csromanspacer{} ภควโต เอตมตฺถํ อาโรเจสุํ.\csromanspacer{} น ภิกฺขเว ภิกฺขุนิยา ภิกฺขุสฺส อุโปสโถ ฐเปตพฺโพ, ฐปิโตปิ อฏฺฐปิโต, ฐเปนฺติยา อาปตฺติ ทุกฺกฏสฺส.\csromanspacer{} น ปวารณา ฐเปตพฺพา, ฐปิตาปิ อฏฺฐปิตา, ฐเปนฺติยา อาปตฺติ ทุกฺกฏสฺส.\csromanspacer{} น สวจนียํ กาตพฺพํ, กตํปิ อกตํ, กโรนฺติยา อาปตฺติ ทุกฺกฏสฺส.\csromanspacer{} น อนุวาโท ปฏฺฐเปตพฺโพ, ปฏฺฐปิโตปิ อปฺปฏฺฐปิโต, ปฏฺฐเปนฺติยา อาปตฺติ ทุกฺกฏสฺส.\csromanspacer{} น โอกาโส \csromanfolio{470}กาเรตพฺโพ, การิโตปิ อการิโต, กาเรนฺติยา อาปตฺติ ทุกฺกฏสฺส.\csromanspacer{} น โจเทตพฺโพ, โจทิโตปิ อโจทิโต, โจเทนฺติยา อาปตฺติ ทุกฺกฏสฺส.\csromanspacer{} น สาเรตพฺโพ, สาริโตปิ อสาริโต, สาเรนฺติยา อาปตฺติ ทุกฺกฏสฺสาติ.}

\prose{เตน โข ปน สมเยน ภิกฺขู ภิกฺขุนีนํ อุโปสถํ ฐเปนฺติ, ปวารณํ ฐเปนฺติ, สวจนียํ กโรนฺติ, อนุวาทํ ปฏฺฐเปนฺติ, โอกาสํ กาเรนฺติ, โจเทนฺติ, สาเรนฺติ.\csromanspacer{} ภควโต เอตมตฺถํ อาโรเจสุํ.\csromanspacer{} อนุชานามิ ภิกฺขเว ภิกฺขุนา ภิกฺขุนิยา อุโปสถํ ฐเปตุํ, ฐปิโตปิ สุฏฺฐปิโต, ฐเปนฺตสฺส อนาปตฺติ.\csromanspacer{} ปวารณํ ฐเปตุํ, ฐปิตาปิ สุฏฺฐปิตา, ฐเปนฺตสฺส อนาปตฺติ.\csromanspacer{} สวจนียํ กาตุํ, กตมฺปิ สุกตํ, กโรนฺตสฺส อนาปตฺติ.\csromanspacer{} อนุวาทํ ปฏฺฐเปตุํ, ปฏฺฐปิโตปิ สุปฺปฏฺฐปิโต, ปฏฺฐเปนฺตสฺส อนาปตฺติ.\csromanspacer{} โอกาสํ กาเรตุํ, การิโตปิ สุการิโต, กาเรนฺตสฺส อนาปตฺติ.\csromanspacer{} โจเทตุํ, โจทิตาปิ สุโจทิตา, โจเทนฺตสฺส อนาปตฺติ.\csromanspacer{} สาเรตุํ, สาริตาปิ สุสาริตา, สาเรนฺตสฺส อนาปตฺตีติ.}

\proseitem{429}{\csromansymbolfootnote{*}{วิ 2. 452 ปิฏฺเฐปิ.}เตน โข ปน สมเยน ฉพฺพคฺคิยา ภิกฺขุนิโย ยาเนน ยายนฺติ, อิตฺถิยุตฺเตนปิ ปุริสนฺตเรน ปุริสยุตฺเตนปิ อิตฺถนฺตเรน.\csromanspacer{} มนุสฺสา อุชฺฌายนฺติ ขิยฺยนฺติ วิปาเจนฺติ “ฯเปฯ เสยฺยถาปิ คงฺคามหิยายา”ติ\footnote{คงฺคามหิยาติ (สี)}.\csromanspacer{} ภควโต เอตมตฺถํ อาโรเจสุํ.\csromanspacer{} น ภิกฺขเว ภิกฺขุนิยา ยาเนน ยายิตพฺพํ, ยา ยาเยยฺย, ยถาธมฺโม กาเรตพฺโพติ.}

\prose{เตน โข ปน สมเยน อญฺญตรา ภิกฺขุนี คิลานา โหติ, น สกฺโกติ ปทสา คนฺตุํ.\csromanspacer{} ภควโต เอตมตฺถํ อาโรเจสุํ.\csromanspacer{} อนุชานามิ ภิกฺขเว คิลานาย ยานนฺติ.\csromanspacer{} อถ โข ภิกฺขุนีนํ เอตทโหสิ “อิตฺถิยุตฺตํ นุ โข ปุริสยุตฺตํ นุ โข”ติ.\csromanspacer{} ภควโต เอตมตฺถํ อาโรเจสุํ.\csromanspacer{} อนุชานามิ ภิกฺขเว อิตฺถิยุตฺตํ ปุริสยุตฺตํ หตฺถวฏฺฏกนฺติ.}

\prose{เตน โข ปน สมเยน อญฺญตริสฺสา ภิกฺขุนิยา ยานุคฺฆาเตน พาฬฺหตรํ อผาสุ อโหสิ.\csromanspacer{} ภควโต เอตมตฺถํ อาโรเจสุํ.\csromanspacer{} อนุชานามิ ภิกฺขเว สิวิกํ ปาฏงฺกินฺติ.}

\proseitem{430}{\csromanfolio{471}เตน โข ปน สมเยน อฑฺฒกาสี คณิกา ภิกฺขุนีสุ ปพฺพชิตา โหติ, สา จ สาวตฺถิํ คนฺตุกามา โหติ “ภควโต สนฺติเก อุปสมฺปชฺชิสฺสามี”ติ.\csromanspacer{} อสฺโสสุํ โข ธุตฺตา “อฑฺฒกาสี กิร คณิกา สาวตฺถิํ คนฺตุกามา”ติ.\csromanspacer{} เต มคฺเค ปริยุฏฺฐิํสุ.\csromanspacer{} อสฺโสสิ โข อฑฺฒกาสี คณิกา “ธุตฺตา กิร มคฺเค ปริยุฏฺฐิตา”ติ.\csromanspacer{} ภควโต สนฺติเก ทูตํ ปาเหสิ “อหํ หิ อุปสมฺปชฺชิตุกามา, กถํ นุ โข มยา ปฏิปชฺชิตพฺพนฺ”ติ.\csromanspacer{} อถ โข ภควา เอตสฺมิํ นิทาเน เอตสฺมิํ ปกรเณ ธมฺมิํ กถํ กตฺวา ภิกฺขู อามนฺเตสิ “อนุชานามิ ภิกฺขเว ทูเตนปิ อุปสมฺปาเทตุนฺ”ติ.}

\prose{ภิกฺขุทูเตน อุปสมฺปาเทนฺติ.\csromanspacer{} ภควโต เอตมตฺถํ อาโรเจสุํ.\csromanspacer{} น ภิกฺขเว ภิกฺขุทูเตน อุปสมฺปาเทตพฺพา, โย อุปสมฺปาเทยฺย, อาปตฺติ ทุกฺกฏสฺสาติ.\csromanspacer{} สิกฺขามานทูเตน อุปสมฺปาเทนฺติ ฯเปฯ.\csromanspacer{} สามเณรทูเตน อุปสมฺปาเทนฺติ ฯเปฯ.\csromanspacer{} สามเณริทูเตน อุปสมฺปาเทนฺติ ฯเปฯ.\csromanspacer{} พาลาย อพฺยตฺตาย ทูเตน อุปสมฺปาเทนฺติ.\csromanspacer{} น ภิกฺขเว พาลาย อพฺยตฺตาย ทูเตน อุปสมฺปาเทตพฺพา, โย อุปสมฺปาเทยฺย, อาปตฺติ ทุกฺกฏสฺส.\csromanspacer{} อนุชานามิ ภิกฺขเว พฺยตฺตาย ภิกฺขุนิยา ปฏิพลาย ทูเตน อุปสมฺปาเทตุนฺติ.}

\prose{ตาย ทูตาย ภิกฺขุนิยา สํฆํ อุปสงฺกมิตฺวา เอกํสํ อุตฺตราสงฺคํ กริตฺวา ภิกฺขูนํ ปาเท วนฺทิตฺวา อุกฺกุฏิกํ นิสีทิตฺวา อญฺชลิํ ปคฺคเหตฺวา เอวมสฺส วจนีโย “อิตฺถนฺนามา อยฺยา อิตฺถนฺนามาย อยฺยาย อุปสมฺปทาเปกฺขา เอกโต-อุปสมฺปนฺนา ภิกฺขุนิสํเฆ วิสุทฺธา, สา เกนจิเทว อนฺตราเยน น อาคจฺฉติ, อิตฺถนฺนามา อยฺยา สํฆํ อุปสมฺปทํ ยาจติ, อุลฺลุมฺปตุ ตํ อยฺยา สํโฆ อนุกมฺปํ อุปาทาย.\csromanspacer{} อิตฺถนฺนามา อยฺยา อิตฺถนฺนามาย อยฺยาย อุปสมฺปทาเปกฺขา เอกโต-อุปสมฺปนฺนา ภิกฺขุนิสํเฆ วิสุทฺธา, สา เกนจิเทว อนฺตราเยน น อาคจฺฉติ, ทุติยมฺปิ อยฺยา อิตฺถนฺนามา สํฆํ อุปสมฺปทํ ยาจติ, อุลฺลุมฺปตุ ตํ อยฺยา สํโฆ อนุกมฺปํ อุปาทาย.\csromanspacer{} อิตฺถนฺนามา อยฺยา อิตฺถนฺนามาย อยฺยาย อุปสมฺปทาเปกฺขา เอกโต-อุปสมฺปนฺนา ภิกฺขุนิสํเฆ วิสุทฺธา, สา เกนจิเทว อนฺตราเยน น อาคจฺฉติ, ตติยมฺปิ อยฺยา อิตฺถนฺนามา สํฆํ อุปสมฺปทํ ยาจติ, อุลฺลุมฺปตุ ตํ อยฺยา สํโฆ อนุกมฺปํ อุปาทายา”ติ.\csromanspacer{} พฺยตฺเตน ภิกฺขุนา ปฏิพเลน สํโฆ ญาเปตพฺโพ\csromandash{}}

\prose{\csromanfolio{472}“สุณาตุ เม ภนฺเต สํโฆ, อิตฺถนฺนามา อิตฺถนฺนามาย อุปสมฺปทาเปกฺขา เอกโต-อุปสมฺปนฺนา ภิกฺขุนิสํเฆ วิสุทฺธา, สา เกนจิเทว อนฺตราเยน น อาคจฺฉติ, อิตฺถนฺนามา สํฆํ อุปสมฺปทํ ยาจติ อิตฺถนฺนามาย ปวตฺตินิยา, ยทิ สํฆสฺส ปตฺตกลฺลํ, สํโฆ อิตฺถนฺนามํ อุปสมฺปาเทยฺย อิตฺถนฺนามาย ปวตฺตินิยา, เอสา ญตฺติ.}

\prose{สุณาตุ เม ภนฺเต สํโฆ, อิตฺถนฺนามาย อุปสมฺปทาเปกฺขา เอกโต-อุปสมฺปนฺนา ภิกฺขุนิสํเฆ วิสุทฺธา, สา เกนจิเทว อนฺตราเยน น อาคจฺฉติ, อิตฺถนฺนามา สํฆํ อุปสมฺปทํ ยาจติ อิตฺถนฺนามาย ปวตฺตินิยา, สํโฆ อิตฺถนฺนามํ อุปสมฺปาเทติ อิตฺถนฺนามาย ปวตฺตินิยา, ยสฺสายสฺมโต ขมติ อิตฺถนฺนามาย อุปสมฺปทา อิตฺถนฺนามาย ปวตฺตินิยา, โส ตุณฺหสฺส, ยสฺส นกฺขมติ, โส ภาเสยฺย.}

\prose{ทุติยมฺปิ เอตมตฺถํ วทามิ ฯเปฯ.\csromanspacer{} ตติยมฺปิ เอตมตฺถํ วทามิ.\csromanspacer{} สุณาตุ เม ภนฺเต สํโฆ, อิตฺถนฺนามา อิตฺถนฺนามาย อุปสมฺปทาเปกฺขา เอกโต-อุปสมฺปนฺนา ภิกฺขุนิสํเฆ วิสุทฺธา, สา เกนจิเทว อนฺตราเยน น อาคจฺฉติ, อิตฺถนฺนามา สํฆํ อุปสมฺปทํ ยาจติ อิตฺถนฺนามาย ปวตฺตินิยา, สํโฆ อิตฺถนฺนามํ อุปสมฺปาเทติ อิตฺถนฺนามาย ปวตฺตินิยา, ยสฺสายสฺมโต ขมติ อิตฺถนฺนามาย อุปสมฺปทา อิตฺถนฺนามาย ปวตฺตินิยา, โส ตุณฺหสฺส, ยสฺส นกฺขมติ, โส ภาเสยฺย.}

\prose{อุปสมฺปนฺนา สํเฆน อิตฺถนฺนามา อิตฺถนฺนามาย ปวตฺตินิยา, ขมติ สํฆสฺส, ตสฺมา ตุณฺหี, เอวเมตํ ธารยามี”ติ.}

\prose{ตาวเทว ฉายา เมตพฺพา, อุตุปฺปมาณํ อาจิกฺขิตพฺพํ, ทิวสภาโค อาจิกฺขิตพฺโพ, สงฺคีติ อาจิกฺขิตพฺพา, ภิกฺขุนิโย วตฺตพฺพา “ตสฺสา ตโย จ นิสฺสเย อฏฺฐ จ อกรณียานิ อาจิกฺเขยฺยาถา”ติ.}

\proseitem{431}{เตน โข ปน สมเยน ภิกฺขุนิโย อรญฺเญ วิหรนฺติ, ธุตฺตา ทูเสนฺติ.\csromanspacer{} ภควโต เอตมตฺถํ อาโรเจสุํ.\csromanspacer{} น ภิกฺขเว ภิกฺขุนิยา อรญฺเญ วตฺถพฺพํ, ยา วเสยฺย, อาปตฺติ ทุกฺกฏสฺสาติ.}

\prose{เตน โข ปน สมเยน อญฺญตเรน อุปาสเกน ภิกฺขุนิสํฆสฺส อุโทสิโต ทินฺโน โหติ.\csromanspacer{} ภควโต เอตมตฺถํ อาโรเจสุํ.\csromanspacer{} อนุชานามิ ภิกฺขเว อุโทสิตนฺติ.\csromanspacer{} อุโทสิโต น สมฺมติ.\csromanspacer{} ภควโต \csromanfolio{473}เอตมตฺถํ อาโรเจสุํ.\csromanspacer{} อนุชานามิ ภิกฺขเว อุปสฺสยนฺติ.\csromanspacer{} อุปสฺสโย น สมฺมติ.\csromanspacer{} ภควโต เอตมตฺถํ อาโรเจสุํ.\csromanspacer{} อนุชานามิ ภิกฺขเว นวกมฺมนฺติ.\csromanspacer{} นวกมฺมํ น สมฺมติ.\csromanspacer{} ภควโต เอตมตฺถํ อาโรเจสุํ.\csromanspacer{} อนุชานามิ ภิกฺขเว ปุคฺคลิกมฺปิ กาตุนฺติ.}

\proseitem{432}{เตน โข ปน สมเยน อญฺญตรา อิตฺถี สนฺนิสินฺนคพฺภา ภิกฺขุนีสุ ปพฺพชิตา โหติ, ตสฺสา ปพฺพชิตาย คพฺโภ วุฏฺฐาติ.\csromanspacer{} อถ โข ตสฺสา ภิกฺขุนิยา เอตทโหสิ “กถํ นุ โข มยา อิมสฺมิํ ทารเก ปฏิปชฺชิตพฺพนฺ”ติ.\csromanspacer{} ภควโต เอตมตฺถํ อาโรเจสุํ.\csromanspacer{} อนุชานามิ ภิกฺขเว โปเสตุํ ยาว โส ทารโก วิญฺญุตํ ปาปุณาตีติ.}

\prose{อถ โข ตสฺสา ภิกฺขุนิยา เอตทโหสิ “มยา จ น ลพฺภา เอกิกาย วตฺถุํ, อญฺญาย จ ภิกฺขุนิยา น ลพฺภา ทารเกน สห วตฺถุํ, กถํ นุ โข มยา ปฏิปชฺชิตพฺพนฺ”ติ.\csromanspacer{} ภควโต เอตมตฺถํ อาโรเจสุํ.\csromanspacer{} อนุชานามิ ภิกฺขเว เอกํ ภิกฺขุนิํ สมฺมนฺนิตฺวา ตสฺสา ภิกฺขุนิยา ทุติยํ ทาตุํ.\csromanspacer{} เอวญฺจ ปน ภิกฺขเว สมฺมนฺนิตพฺพา, ปฐมํ ภิกฺขุนี ยาจิตพฺพา, ยาจิตฺวา พฺยตฺตาย ภิกฺขุนิยา ปฏิพลาย สํโฆ ญาเปตพฺโพ\csromandash{}}

\prose{“สุณาตุ เม อยฺเย สํโฆ, ยทิ สํฆสฺส ปตฺตกลฺลํ, สํโฆ อิตฺถนฺนามํ ภิกฺขุนิํ สมฺมนฺเนยฺย อิตฺถนฺนามาย ภิกฺขุนิยา ทุติยํ, เอสา ญตฺติ.}

\prose{สุณาตุ เม อยฺเย สํโฆ, สํโฆ อิตฺถนฺนามํ ภิกฺขุนิํ สมฺมนฺนติ อิตฺถนฺนามาย ภิกฺขุนิยา ทุติยํ, ยสฺสา อยฺยาย ขมติ อิตฺถนฺนามาย ภิกฺขุนิยา สมฺมุติ อิตฺถนฺนามาย ภิกฺขุนิยา ทุติยาย, สา ตุณฺหสฺส, ยสฺสา นกฺขมติ, สา ภาเสยฺย.}

\prose{สมฺมตา สํเฆน อิตฺถนฺนามา ภิกฺขุนี อิตฺถนฺนามาย ภิกฺขุนิยา ทุติยา, ขมติ สํฆสฺส, ตสฺมา ตุณฺหี, เอวเมตํ ธารยามี”ติ.}

\prose{อถ โข ตสฺสา ทุติยิกาย ภิกฺขุนิยา เอตทโหสิ “กถํ นุ โข มยา อิมสฺมิํ ทารเก ปฏิปชฺชิตพฺพนฺ”ติ.\csromanspacer{} ภควโต เอตมตฺถํ อาโรเจสุํ.\csromanspacer{} อนุชานามิ ภิกฺขเว ฐเปตฺวา สาคารํ ยถา อญฺญสฺมิํ ปุริเส ปฏิปชฺชนฺติ\footnote{ปฏิปชฺชติ (สฺยา)}, เอวํ ตสฺมิํ ทารเก ปฏิปชฺชิตุนฺติ.}

\proseitem{433}{\csromanfolio{474}เตน โข ปน สมเยน อญฺญตรา ภิกฺขุนี ครุธมฺมํ อชฺฌาปนฺนา โหติ มานตฺตจารินี.\csromanspacer{} อถ โข ตสฺสา ภิกฺขุนิยา เอตทโหสิ “มยา จ น ลพฺภา เอกิกาย วตฺถุํ, อญฺญาย จ ภิกฺขุนิยา น ลพฺภา สห มยา วตฺถุํ, กถํ นุ โข มยา ปฏิปชฺชิตพฺพนฺ”ติ.\csromanspacer{} ภควโต เอตมตฺถํ อาโรเจสุํ.\csromanspacer{} อนุชานามิ ภิกฺขเว เอกํ ภิกฺขุนิํ สมฺมนฺนิตฺวา ตสฺสา ภิกฺขุนิยา ทุติยํ ทาตุํ.\csromanspacer{} เอวญฺจ ปน ภิกฺขเว สมฺมนฺนิตพฺพา, ปฐมํ ภิกฺขุนี ยาจิตพฺพา, ยาจิตฺวา พฺยตฺตาย ภิกฺขุนิยา ปฏิพลาย สํโฆ ญาเปตพฺโพ\csromandash{}}

\prose{“สุณาตุ เม อยฺเย สํโฆ, ยทิ สํฆสฺส ปตฺตกลฺลํ, สํโฆ อิตฺถนฺนามํ ภิกฺขุนิํ สมฺมนฺเนยฺย อิตฺถนฺนามาย ภิกฺขุนิยา ทุติยํ, เอสา ญตฺติ.}

\prose{สุณาตุ เม อยฺเย สํโฆ, สํโฆ, สํโฆ อิตฺถนฺนามํ ภิกฺขุนิํ สมฺมนฺนติ อิตฺถนฺนามาย ภิกฺขุนิยา ทุติยํ, ยสฺสา อยฺยาย ขมติ อิตฺถนฺนามาย ภิกฺขุนิยา สมฺมุติ อิตฺถนฺนามาย ภิกฺขุนิยา ทุติยาย, สา ตุณฺหสฺส, ยสฺสา นกฺขมติ, สา ภาเสยฺย.}

\prose{สมฺมตา สํเฆน อิตฺถนฺนามา ภิกฺขุนี อิตฺถนฺนามาย ภิกฺขุนิยา ทุติยา, ขมติ สํฆสฺส, ตสฺมา ตุณฺหี, เอวเมตํ ธารยามี”ติ.}

\proseitem{434}{เตน โข ปน สมเยน อญฺญตรา ภิกฺขุนี สิกฺขํ ปจฺจกฺขาย วิพฺภมิ, สา ปุน ปจฺจาคนฺตฺวา ภิกฺขุนิโย อุปสมฺปทํ ยาจิ.\csromanspacer{} ภควโต เอตมตฺถํ อาโรเจสุํ.\csromanspacer{} น ภิกฺขเว ภิกฺขุนิยา สิกฺขาปจฺจกฺขานํ, ยเทว สา วิพฺภนฺตา, ตเทว สา อภิกฺขุนีติ.}

\prose{เตน โข ปน สมเยน อญฺญตรา ภิกฺขุนี สกาวาสา ติตฺถายตนํ สงฺกมิ, สา ปุน ปจฺจาคนฺตฺวา ภิกฺขุนิโย อุปสมฺปทํ ยาจิ.\csromanspacer{} ภควโต เอตมตฺถํ อาโรเจสุํ.\csromanspacer{} ยา สา ภิกฺขเว ภิกฺขุนี สกาวาสา ติตฺถายตนํ สงฺกนฺตา, สา อาคตา น อุปสมฺปาเทตพฺพาติ.}

\prose{เตน โข ปน สมเยน ภิกฺขุนิโย ปุริเสหิ อภิวาทนํ เกสจฺเฉทนํ นขจฺเฉทนํ วณปฺปฏิกมฺมํ กุกฺกุจฺจายนฺตา น สาทิยนฺติ.\csromanspacer{} ภควโต เอตมตฺถํ อาโรเจสุํ.\csromanspacer{} อนุชานามิ ภิกฺขเว สาทิตุนฺติ\footnote{สาทิยิตุํ (ก)}.}

\proseitem{435}{\csromanfolio{475}เตน โข ปน สมเยน ภิกฺขุนิโย ปลฺลงฺเกน นิสีทนฺติ ปณฺหิสมฺผสฺสํ สาทิยนฺตี\footnote{สาทิยนฺตา (ก)}.\csromanspacer{} ภควโต เอตมตฺถํ อาโรเจสุํ.\csromanspacer{} น ภิกฺขเว ภิกฺขุนิยา ปลฺลงฺเกน นิสีทิตพฺพํ, ยา นิสีเทยฺย, อาปตฺติ ทุกฺกฏสฺสาติ.}

\prose{เตน โข ปน สมเยน อญฺญตรา ภิกฺขุนี คิลานา โหติ, ตสฺสา วินา ปลฺลงฺเกน น ผาสุ โหติ.\csromanspacer{} ภควโต เอตมตฺถํ อาโรเจสุํ.\csromanspacer{} อนุชานามิ ภิกฺขเว ภิกฺขุนิยา อฑฺฒปลฺลงฺกนฺติ.}

\prose{เตน โข ปน สมเยน ภิกฺขุนิโย วจฺจกุฏิยา วจฺจํ กโรนฺติ, ฉพฺพคฺคิยา ภิกฺขุนิโย ตตฺเถว คพฺภํ ปาเตนฺติ.\csromanspacer{} ภควโต เอตมตฺถํ อาโรเจสุํ.\csromanspacer{} น ภิกฺขเว ภิกฺขุนิยา วจฺจกุฏิยา วจฺโจ กาตพฺโพ, ยา กเรยฺย, อาปตฺติ ทุกฺกฏสฺส.\csromanspacer{} อนุชานามิ ภิกฺขเว เหฏฺฐา วิวเฏ อุปริ ปฏิจฺฉนฺเน วจฺจํ กาตุนฺติ.}

\proseitem{436}{เตน โข ปน สมเยน ภิกฺขุนิโย จุณฺเณน นหายนฺติ.\csromanspacer{} มนุสฺสา อุชฺฌายนฺติ ขิยฺยนฺติ วิปาเจนฺติ. “ฯเปฯ เสยฺยถาปิ คิหินี กามโภคินิโย”ติ.\csromanspacer{} ภควโต เอตมตฺถํ อาโรเจสุํ.\csromanspacer{} น ภิกฺขเว ภิกฺขุนิยา จุณฺเณน นหายิตพฺพํ, ยา นหาเยยฺย, อาปตฺติ ทุกฺกฏสฺส.\csromanspacer{} อนุชานามิ ภิกฺขเว กุกฺกุสํ มตฺติกนฺติ.}

\prose{เตน โข ปน สมเยน ภิกฺขุนิโย วาสิตกาย มตฺติกาย นหายนฺติ.\csromanspacer{} มนุสฺสา อุชฺฌยนฺติ ขิยฺยนฺติ วิปาเจนฺติ. “ฯเปฯ เสยฺยถาปิ คิหินี กามโภคินิโย”ติ.\csromanspacer{} ภควโต เอตมตฺถํ อาโรเจสุํ.\csromanspacer{} น ภิกฺขเว ภิกฺขุนิยา วาสิตกาย มตฺติกาย นหายิตพฺพํ, ยา นหาเยยฺย, อาปตฺติ ทุกฺกฏสฺส.\csromanspacer{} อนุชานามิ ภิกฺขเว ปกติมตฺติกนฺติ.}

\prose{เตน โข ปน สมเยน ภิกฺขุนิโย ชนฺตาฆเร นหายนฺติโย โกลาหลํ อกํสุ.\csromanspacer{} ภควโต เอตมตฺถํ อาโรเจสุํ.\csromanspacer{} น ภิกฺขเว ภิกฺขุนิยา ชนฺตาฆเร นหายิตพฺพํ, ยา นหาเยยฺย, อาปตฺติ ทุกฺกฏสฺสาติ.}

\csromanmatikaanchor{mtk.211}
\csromantocmarkref{section}{อุทฺทานคาถา}{mtk.211}
\bookmark[dest={mtk.211},level=1]{อุทฺทานคาถา}
\prose{เตน โข ปน สมเยน ภิกฺขุนิโย ปฏิโสเต นหายนฺติ ธาราสมฺผสฺสํ สาทิยนฺติ.\csromanspacer{} ภควโต เอตมตฺถํ อาโรเจสุํ.\csromanspacer{} น ภิกฺขเว \csromanfolio{476}ภิกฺขุนิยา ปฏิโสเต นหายิตพฺพํ, ยา นหาเยยฺย, อาปตฺติ ทุกฺกฏสฺสาติ.}

\prose{เตน โข ปน สมเยน ภิกฺขุนิโย อติตฺเถ นหายนฺติ, ธุตฺตา ทูเสนฺติ.\csromanspacer{} ภควโต เอตมตฺถํ อาโรเจสุํ.\csromanspacer{} น ภิกฺขเว ภิกฺขุนิยา อติตฺเถ นหายิตพฺพํ, ยา นหาเยยฺย, อาปตฺติ ทุกฺกฏสฺสาติ.}

\prose{เตน โข ปน สมเยน ภิกฺขุนิโย ปุริสติตฺเถ นหายนฺติ.\csromanspacer{} มนุสฺสา อุชฺฌายนฺติ ขิยฺยนฺติ วิปาเจนฺติ. “ฯเปฯ เสยฺยถาปิ คิหินี กามโภคินิโย”ติ.\csromanspacer{} ภควโต เอตมตฺถํ อาโรเจสุํ.\csromanspacer{} น ภิกฺขเว ภิกฺขุนิยา ปุริสติตฺเถ นหายิตพฺพํ, ยา นหาเยยฺย, อาปตฺติ ทุกฺกฏสฺส.\csromanspacer{} อนุชานามิ ภิกฺขเว มหิลาติตฺเถ นหายิตุนฺติ.}

\nitthitam{ตติยภาณวาโร นิฏฺฐิโต.}
\vspace{\tipitakaparskip}
\csromancenter{ภิกฺขุนิกฺขนฺธโก ทสโม.}
\csromancenter{อิมสฺมิํ ขนฺธเก วตฺถู เอกสตํ.}
\titlehead{\textbf{ตสฺสุทฺทานํ}}
\csromangathameasure{ปพฺพชฺชํโคตมี ยาจิ, นานุญฺญาสิ ตถาคโต. \\ กปิลวตฺถุ เวสาลิํ, อคมาสิ วินายโก. \\ รโชกิณฺเณน โกฏฺฐเก, อานนฺทสฺส ปเวทยิ. \\ ภพฺโพติ นยโต ยาจิ, มาตาติ โปสิกาติ จ. \\ วสฺสสตํ ตทหุ จ, อภิกฺขุปจฺจาสีสนา. \\ ปวารณา ครุธมฺมา, เทฺว วสฺสา อนกฺโกสนา. \\ โอวโฏ จ อฏฺฐ ธมฺมา, ยาวชีวานุวตฺตนา. \\ ครุธมฺมปฏิคฺคาโห, สาวสฺสา อุปสมฺปทา. \\ วสฺสสหสฺสํ ปญฺเจว, กุมฺภเถนกเสตฏฺฏิ. \\ มญฺชิฏฺฐิก-อุปมาหิ, เอวํ สทฺธมฺมหิํสนา. \\ อาฬิํ พนฺเธยฺย ปา-เอว, ปุน สทฺธมฺมสณฺฐิติ. \\ อุปสมฺปาเทตุํ อยฺยา, ยถาวุฑฺฒาภิวาทนา. \\ น กริสฺสนฺติ กิเมว, สาธารณาสาธารณํ. \\ โอวาทํ ปาติโมกฺขญฺจ, เกน นุ โข อุปสฺสยํ. \\ น ชานนฺติ จ อาจิกฺขิ, น กโรนฺติ จ ภิกฺขุหิ. \\ ปฏิคฺคเหตุํ ภิกฺขูหิ, ภิกฺขุนีหิ ปฏิคฺคโห. \\ อาจิกฺขิ กมฺมํ ภิกฺขูหิ, อุชฺฌายนฺติ ภิกฺขุนีหิ วา. \\ อาจิกฺขิตุํ ภณฺฑนญฺจ, โรเปตฺวา อุปฺปลาย จ. \\ สาวตฺถิยา กทฺทโมท, อวนฺทิ กาย อูรุ จ. \\ องฺคชาตญฺจ โอภาสํ, สมฺปโยเชนฺติ วคฺคิกา. \\ อวนฺทิโย ทณฺฑกมฺมํ, ภิกฺขุนิโย ตถา ปุน. \\ อาวรณญฺจ โอวาทํ, กปฺปติ นุ โข ปกฺกมิ. \\ พาลา วตฺถุวินิจฺฉยา, โอวาทํ สํโฆ ปญฺจหิ. \\ ทุเว ติสฺโส น คณฺหนฺติ, พาลา คิลานคมิกํ. \\ อารญฺญิโก นาโรเจนฺติ, น ปจฺจาคจฺฉนฺติ จ. \\ ทีฆํ วิลีวจมฺมญฺจ, ทุสฺสา จ เวณิวฏฺฏิ จ. \\ โจฬเวณิ จ ปฏฺฏิ จ, สุตฺตเวณิ จ วฏฺฏิกา. \\ อฏฺฐิลฺลํ โคหนุเกน, หตฺถโกจฺฉํ ปาทํ ตถา. \\ อูรุํ มุขํ ทนฺตมํสํ, อาลิมฺโปมทฺทจุณฺณนา. \\ ลญฺเฉนฺติ องฺคราคญฺจ, มุขราคํ ตถา ทุเว. \\ อวงฺคํ วิเสโสโลโก, สาโลเกน นจฺเจน จ. \\ เวสี ปานาคารํ สูนํ, อาปณํ วฑฺฒิ วณิชฺชา. \\ ทาสํ ทาสิํ กมฺมกรํ, กมฺมการิํ อุปฏฺฐยฺยุํ. \\ ติรจฺฉานหรีตกิ, สํธารยนฺติ นมตกํ. \\ นีลํ ปีตํ โลหิตกํ, มญฺชิฏฺฐกณฺหจีวรา. \\ มหารงฺคมหานาม, อจฺฉินฺนา ทีฆเมว จ. \\ ปุปฺผผลกญฺจุกญฺจ, ติรีฏกญฺจ ธารยุํ. \\ ภิกฺขุนี สิกฺขมานาย, สามเณราย อจฺจเย. \\ นิยฺยาทิเต ปริกฺขาเร, ภิกฺขุนิโยว อิสฺสรา. \\ ภิกฺขุสฺส สามเณรสฺส, อุปาสกสฺสุปาสิกา. \\ อญฺเญสญฺจ ปริกฺขาเร, นิยฺยาเต ภิกฺขุ-อิสฺสรา. \\ มลฺลี คพฺภํ ปตฺตมูลํ, พฺยญฺชนํ อามิเสน จ. \\ อุสฺสนฺนญฺจ พาฬฺหตรํ, สนฺนิธิกตมามิสํ. \\ ภิกฺขูนํ ยาทิสํ โภฏฺฐํ, ภิกฺขุนีนํ ตถา กเร. \\ เสนาสนํ อุตุนิโย, มกฺขียติ ปฏาณิ จ. \\ ฉิชฺชนฺติ สพฺพกาลญฺจ, อนิมิตฺตาปิ ทิสฺสเร. \\ นิมิตฺตา โลหิตา เจว, ตเถว ธุวโลหิตา. \\ ธุวโจฬปคฺฆรนฺตี, สิขรณิตฺถิปณฺฑกา. \\ เวปุริสี จ สมฺภินฺนา, อุภโตพฺยญฺชนาปิ จ. \\ อนิมิตฺตาทิโต กตฺวา, ยาว อุภโตพฺยญฺชนา. \\ เอตํ เปยฺยาลโต เหฏฺฐา, กุฏฺฐํ คณฺโฑ กิลาโส จ โสสาปมาโร มานุสี, \\ อิตฺถีสิ ภุชิสฺสาสิ จ. อณณา น ราชภฏี, \\ อนุญฺญาตา จ วีสติ. ปริปุณฺณา จ กินฺนามา, \\ กานามา เต ปวตฺตินี. จตุวีสนฺตรายานํ, \\ ปุจฺฉิตฺวา อุปสมฺปทา. วิตฺถายนฺติ อนนุสิฏฺฐา, \\ สํฆมชฺเฌ ตเถว จ. อุปชฺฌาคาห สํฆาฏิ, \\ อุตฺตรนฺตรวาสโก. สํกจฺจุทกสาฏิ จ, \\ อาจิกฺขิตฺวาน เปสเย. พาลา อสมฺมเตกโต, \\ ยาเจ ปุจฺฉนฺตรายิกา. เอกโต-อุปสมฺปนฺนา, \\ ภิกฺขุสํเฆ ตถา ปุน. ฉายา อุตุ ทิวสา จ, \\ สงฺคีติ ตโย นิสฺสเย. อฏฺฐ อกรณียานิ, \\ กาลํ สพฺพตฺถ อฏฺเฐว. น ปวาเรนฺติ ภิกฺขุนี, \\ ภิกฺขุสํฆํ ตเถว จ. โกลาหลํ ปุเรภตฺตํ, \\ วิกาเล จ โกลาหลํ. อุโปสถํ ปวารณํ, \\ สวจนียานุวาทนํ. โอกาสํ โจเท สาเรนฺติ, \\ ปฏิกฺขิตฺตํ มเหสินา. ตเถว ภิกฺขุ ภิกฺขุนี, \\ อนุญฺญาตํ มเหสินา. ยานํ คิลานยุตฺตญฺจ, \\ ยานุคฺฆาตฑฺฒกาสิกา. ภิกฺขุ สิกฺขา สามเณร, \\ สามเณรี จ พาลาย. อรญฺเญ อุปาสเกน, \\ อุโทสิโต อุปสฺสยํ. น สมฺมติ นวกมฺมํ, \\ นิสินฺนคพฺภ-เอกิกา. สาคารญฺจ ครุธมฺมํ, \\ ปจฺจกฺขาย จ สงฺกมิ. อภิวาทนเกสา จ, \\ นขา จ วณกมฺมนา. ปลฺลงฺเกน คิลานา จ, \\ วจฺจํ จุณฺเณน วาสิตํ. ชนฺตาฆเร ปฏิโสเต, \\ อติตฺเถ ปุริเสน จ. มหาโคตมี อายาจิ, \\ อานนฺโท จาปิ โยนิโส. ปริสา จตสฺโส โหนฺติ, \\ ปพฺพชฺชา ชินสาสเน. สํเวคชนนตฺถาย, \\ สทฺธมฺมสฺส จ วุทฺธิยา. อาตุรสฺสาว เภสชฺชํ, \\ เอวํ พุทฺเธน เทสิตํ. เอวํ วินีตา สทฺธมฺเม, \\ มาตุคามาปิ อิตรา. ยายนฺติ อจฺจุตํ ฐานํ, \\ ยตฺถ คนฺตฺวา น โสจเรติ.}
\csromangathagroup{\csromangathastanza{ปพฺพชฺชํโคตมี ยาจิ, นานุญฺญาสิ ตถาคโต. \\ กปิลวตฺถุ เวสาลิํ, อคมาสิ วินายโก.}\csromangathastanza{รโชกิณฺเณน โกฏฺฐเก, อานนฺทสฺส ปเวทยิ. \\ ภพฺโพติ นยโต ยาจิ, มาตาติ โปสิกาติ จ.}\csromangathastanza{วสฺสสตํ ตทหุ จ, อภิกฺขุปจฺจาสีสนา. \\ ปวารณา ครุธมฺมา, เทฺว วสฺสา อนกฺโกสนา.}\csromangathastanza{โอวโฏ จ อฏฺฐ ธมฺมา, ยาวชีวานุวตฺตนา. \\ ครุธมฺมปฏิคฺคาโห, สาวสฺสา อุปสมฺปทา.}\csromangathastanza{วสฺสสหสฺสํ ปญฺเจว, กุมฺภเถนกเสตฏฺฏิ. \\ มญฺชิฏฺฐิก-อุปมาหิ, เอวํ สทฺธมฺมหิํสนา.}\csromangathastanza{อาฬิํ พนฺเธยฺย ปา-เอว, ปุน สทฺธมฺมสณฺฐิติ. \\ อุปสมฺปาเทตุํ อยฺยา, ยถาวุฑฺฒาภิวาทนา.}\csromangathastanza{\csromanfolio{477}น กริสฺสนฺติ กิเมว, สาธารณาสาธารณํ. \\ โอวาทํ ปาติโมกฺขญฺจ, เกน นุ โข อุปสฺสยํ.}\csromangathastanza{น ชานนฺติ จ อาจิกฺขิ, น กโรนฺติ จ ภิกฺขุหิ. \\ ปฏิคฺคเหตุํ ภิกฺขูหิ, ภิกฺขุนีหิ ปฏิคฺคโห.}\csromangathastanza{อาจิกฺขิ กมฺมํ ภิกฺขูหิ, อุชฺฌายนฺติ ภิกฺขุนีหิ วา. \\ อาจิกฺขิตุํ ภณฺฑนญฺจ, โรเปตฺวา อุปฺปลาย จ.}\csromangathastanza{สาวตฺถิยา กทฺทโมท, อวนฺทิ กาย อูรุ จ. \\ องฺคชาตญฺจ โอภาสํ, สมฺปโยเชนฺติ วคฺคิกา.}\csromangathastanza{อวนฺทิโย ทณฺฑกมฺมํ, ภิกฺขุนิโย ตถา ปุน. \\ อาวรณญฺจ โอวาทํ, กปฺปติ นุ โข ปกฺกมิ.}\csromangathastanza{พาลา วตฺถุวินิจฺฉยา, โอวาทํ สํโฆ ปญฺจหิ. \\ ทุเว ติสฺโส น คณฺหนฺติ, พาลา คิลานคมิกํ.}\csromangathastanza{อารญฺญิโก นาโรเจนฺติ, น ปจฺจาคจฺฉนฺติ จ. \\ ทีฆํ วิลีวจมฺมญฺจ, ทุสฺสา จ เวณิวฏฺฏิ จ.}\csromangathastanza{โจฬเวณิ จ ปฏฺฏิ จ, สุตฺตเวณิ จ วฏฺฏิกา. \\ อฏฺฐิลฺลํ โคหนุเกน, หตฺถโกจฺฉํ ปาทํ ตถา.}\csromangathastanza{อูรุํ มุขํ ทนฺตมํสํ, อาลิมฺโปมทฺทจุณฺณนา. \\ ลญฺเฉนฺติ องฺคราคญฺจ, มุขราคํ ตถา ทุเว.}\csromangathastanza{อวงฺคํ วิเสโสโลโก, สาโลเกน นจฺเจน จ\footnote{สาโลเก สนจฺเจน จ (สี), สาโลเกน สนจฺจนํ (สฺยา)}. \\ เวสี ปานาคารํ สูนํ, อาปณํ วฑฺฒิ วณิชฺชา.}\csromangathastanza{ทาสํ ทาสิํ กมฺมกรํ, กมฺมการิํ อุปฏฺฐยฺยุํ. \\ ติรจฺฉานหรีตกิ, สํธารยนฺติ นมตกํ.}\csromangathastanza{นีลํ ปีตํ โลหิตกํ, มญฺชิฏฺฐกณฺหจีวรา. \\ มหารงฺคมหานาม, อจฺฉินฺนา ทีฆเมว จ.}\csromangathastanza{ปุปฺผผลกญฺจุกญฺจ, ติรีฏกญฺจ ธารยุํ. \\ ภิกฺขุนี สิกฺขมานาย, สามเณราย อจฺจเย.}\csromangathastanza{\csromanfolio{478}นิยฺยาทิเต ปริกฺขาเร, ภิกฺขุนิโยว อิสฺสรา. \\ ภิกฺขุสฺส สามเณรสฺส, อุปาสกสฺสุปาสิกา.}\csromangathastanza{อญฺเญสญฺจ ปริกฺขาเร, นิยฺยาเต ภิกฺขุ-อิสฺสรา. \\ มลฺลี คพฺภํ ปตฺตมูลํ, พฺยญฺชนํ อามิเสน จ.}\csromangathastanza{อุสฺสนฺนญฺจ พาฬฺหตรํ, สนฺนิธิกตมามิสํ. \\ ภิกฺขูนํ ยาทิสํ โภฏฺฐํ\footnote{เหฏฺฐา (สี), เหฏฺฐํ (สฺยา), ภุตฺติ (ก)}, ภิกฺขุนีนํ ตถา กเร.}\csromangathastanza{เสนาสนํ อุตุนิโย, มกฺขียติ ปฏาณิ จ. \\ ฉิชฺชนฺติ สพฺพกาลญฺจ, อนิมิตฺตาปิ ทิสฺสเร.}\csromangathastanza{นิมิตฺตา โลหิตา เจว, ตเถว ธุวโลหิตา. \\ ธุวโจฬปคฺฆรนฺตี, สิขรณิตฺถิปณฺฑกา.}\csromangathastanza{เวปุริสี จ สมฺภินฺนา, อุภโตพฺยญฺชนาปิ จ. \\ อนิมิตฺตาทิโต กตฺวา, ยาว อุภโตพฺยญฺชนา.}\csromangathastanza{เอตํ เปยฺยาลโต เหฏฺฐา, กุฏฺฐํ คณฺโฑ กิลาโส จ โสสาปมาโร มานุสี, \\ อิตฺถีสิ ภุชิสฺสาสิ จ. อณณา น ราชภฏี,}\csromangathastanza{อนุญฺญาตา จ วีสติ. ปริปุณฺณา จ กินฺนามา, \\ กานามา เต ปวตฺตินี. จตุวีสนฺตรายานํ,}\csromangathastanza{ปุจฺฉิตฺวา อุปสมฺปทา. วิตฺถายนฺติ อนนุสิฏฺฐา, \\ สํฆมชฺเฌ ตเถว จ. อุปชฺฌาคาห สํฆาฏิ,}\csromangathastanza{อุตฺตรนฺตรวาสโก. สํกจฺจุทกสาฏิ จ, \\ อาจิกฺขิตฺวาน เปสเย. พาลา อสมฺมเตกโต,}\csromangathastanza{ยาเจ ปุจฺฉนฺตรายิกา. เอกโต-อุปสมฺปนฺนา, \\ ภิกฺขุสํเฆ ตถา ปุน. ฉายา อุตุ ทิวสา จ,}\csromangathastanza{สงฺคีติ ตโย นิสฺสเย. อฏฺฐ อกรณียานิ, \\ กาลํ สพฺพตฺถ อฏฺเฐว. น ปวาเรนฺติ ภิกฺขุนี,}\csromangathastanza{ภิกฺขุสํฆํ ตเถว จ. โกลาหลํ ปุเรภตฺตํ, \\ วิกาเล จ โกลาหลํ. อุโปสถํ ปวารณํ,}\csromangathastanza{\csromanfolio{479}สวจนียานุวาทนํ. โอกาสํ โจเท สาเรนฺติ, \\ ปฏิกฺขิตฺตํ มเหสินา. ตเถว ภิกฺขุ ภิกฺขุนี,}\csromangathastanza{อนุญฺญาตํ มเหสินา. ยานํ คิลานยุตฺตญฺจ, \\ ยานุคฺฆาตฑฺฒกาสิกา. ภิกฺขุ สิกฺขา สามเณร,}\csromangathastanza{สามเณรี จ พาลาย. อรญฺเญ อุปาสเกน, \\ อุโทสิโต อุปสฺสยํ. น สมฺมติ นวกมฺมํ,}\csromangathastanza{นิสินฺนคพฺภ-เอกิกา. สาคารญฺจ ครุธมฺมํ, \\ ปจฺจกฺขาย จ สงฺกมิ. อภิวาทนเกสา จ,}\csromangathastanza{นขา จ วณกมฺมนา. ปลฺลงฺเกน คิลานา จ, \\ วจฺจํ จุณฺเณน วาสิตํ. ชนฺตาฆเร ปฏิโสเต,}\csromangathastanza{อติตฺเถ ปุริเสน จ. มหาโคตมี อายาจิ, \\ อานนฺโท จาปิ โยนิโส. ปริสา จตสฺโส โหนฺติ,}\csromangathastanza{ปพฺพชฺชา ชินสาสเน. สํเวคชนนตฺถาย, \\ สทฺธมฺมสฺส จ วุทฺธิยา. อาตุรสฺสาว เภสชฺชํ,}\csromangathastanza{เอวํ พุทฺเธน เทสิตํ. เอวํ วินีตา สทฺธมฺเม, \\ มาตุคามาปิ อิตรา. ยายนฺติ\footnote{ตายนฺติ (สี, สฺยา)} อจฺจุตํ ฐานํ, \\ ยตฺถ คนฺตฺวา น โสจเรติ.}}

\nitthitam{\textbf{ภิกฺขุนิกฺขนฺธโก นิฏฺฐิโต.}}
\csromanmatikaanchor{mtk.212}
\csromantocmarknopage{chapter}{11. ปญฺจสติกกฺขนฺธก}{mtk.212}
\bookmark[dest={mtk.212},level=0]{11. ปญฺจสติกกฺขนฺธก}
\chapterheadpageread{11. ปญฺจสติกกฺขนฺธก}
\csromanmatikaanchor{mtk.213}
\csromantocmarkref{section}{1. สงฺคีตินิทาน}{mtk.213}
\bookmark[dest={mtk.213},level=1]{1. สงฺคีตินิทาน}
\csromanheader{1}{1. สงฺคีตินิทาน}
\proseitem{437}{\csromanfolio{480}อถ โข อายสฺมา มหากสฺสโป ภิกฺขู อามนฺเตสิ\csromandash{}เอกมิทาหํ อาวุโส สมยํ ปาวาย กุสินารํ อทฺธานมคฺคปฺปฏิปนฺโน มหตา ภิกฺขุสํเฆน สทฺธิํ ปญฺจมตฺเตหิ ภิกฺขุสเตหิ.\csromanspacer{} อถ ขฺวาหํ อาวุโส มคฺคา โอกฺกมฺม อญฺญตรสฺมิํ รุกฺขมูเล นิสีทิํ.}

\prose{\csromansymbolfootnote{*}{ที 2. 133 ปิฏฺเฐปิ.}เตน โข ปน สมเยน อญฺญตโร อาชีวโก กุสินาราย มนฺทารวปุปฺผํ คเหตฺวา ปาวํ อทฺธานมคฺคปฺปฏิปนฺโน โหติ.\csromanspacer{} อทฺทสํ โข อหํ อาวุโส ตํ อาชีวกํ ทูรโตว อาคจฺฉนฺตํ, ทิสฺวาน ตํ อาชีวกํ เอตทโวจํ “อปาวุโส อมฺหากํ สตฺถารํ ชานาสี”ติ.\csromanspacer{} อามาวุโส ชานามิ, อชฺช สตฺตาหปรินิพฺพุโต สมโณ โคตโม, ตโต เม อิทํ มนฺทารวปุปฺผํ คหิตนฺติ.\csromanspacer{} ตตฺราวุโส เย เต ภิกฺขู อวีตราคา, อปฺเปกจฺเจ พาหา ปคฺคยฺห กนฺทนฺติ, ฉินฺนปาตํ ปปตนฺติ, อาวฏฺฏนฺติ, วิวฏฺฏนฺติ, “อติขิปฺปํ ภควา ปรินิพฺพุโต, อติขิปฺปํ สุคโต ปรินิพฺพุโต, อติขิปฺปํ จกฺขุํโลเก อนฺตรหิตนฺ”ติ.\csromanspacer{} เย ปน เต ภิกฺขู วีตราคา, เต สตา สมฺปชานา อธิวาเสนฺติ “อนิจฺจา สงฺขารา, ตํ กุเตตฺถ ลพฺภา”ติ.}

\prose{อถ ขฺวาหํ อาวุโส เต ภิกฺขู เอตทโวจํ “อลํ อาวุโส มา โสจิตฺถ, มา ปริเทวิตฺถ, นเนฺวตํ อาวุโส ภควตา ปฏิกจฺเจว อกฺขาตํ สพฺเพเหว ปิเยหิ มนาเปหิ นานาภาโว วินาภาโว อญฺญถาภาโว, ตํ กุเตตฺถ อาวุโส ลพฺภา, ‘ยํ ตํ ชาตํ ภูตํ สงฺขตํ ปโลกธมฺมํ, ตํ วต มา ปลุชฺชี’ติ เนตํ ฐานํ วิชฺชตี”ติ.}

\prose{เตน โข ปนาวุโส สมเยน สุภทฺโท นาม วุฑฺฒปพฺพชิโต ตสฺสํ ปริสายํ นิสินฺโน โหติ.\csromanspacer{} อถ โข อาวุโส สุภทฺโท วุฑฺฒปพฺพชิโต เต ภิกฺขู เอตทโวจ “อลํ อาวุโส มา โสจิตฺถ, มา ปริเทวิตฺถ, สุมุตฺตา มยํ เตน มหาสมเณน, อุปทฺทุตา จ มยํ โหม ‘อิทํ โว กปฺปติ, อิทํ โว น กปฺปตี’ติ, อิทานิ ปน มยํ ยํ อิจฺฉิสฺสาม ตํ กริสฺสาม, ยํ น อิจฺฉิสฺสาม น ตํ กริสฺสามา”ติ. \csromanfolio{481}หนฺท มยํ อาวุโส ธมฺมญฺจ วินยญฺจ สงฺคายาม, ปุเร อธมฺโม ทิปฺปติ\footnote{ทิพฺพติ (ก)}, ธมฺโม ปฏิพาหิยฺยติ, ปุเร อวินโย ทิปฺปติ, วินโย ปฏิพาหิยฺยติ, ปุเร อธมฺมวาทิโน พลวนฺโต โหนฺติ, ธมฺมวาทิโน ทุพฺพลา โหนฺติ, ปุเร อวินยวาทิโน พลวนฺโต โหนฺติ, วินยวาทิโน ทุพฺพลา โหนฺตีติ.}

\prose{เตน หิ ภนฺเต เถโร ภิกฺขู อุจฺจินตูติ.\csromanspacer{} อถ โข อายสฺมา มหากสฺสโป เอเกนูนปญฺจ-อรหนฺตสตานิ อุจฺจินิ.\csromanspacer{} ภิกฺขู อายสฺมนฺตํ มหากสฺสปํ เอตทโวจุํ “อยํ ภนฺเต อายสฺมา อานนฺโท กิญฺจาปิ เสกฺโข, อภพฺโพ ฉนฺทา โทสา โมหา ภยา อคติํ คนฺตุํ, พหุ จ อเนน ภควโต สนฺติเก ธมฺโม จ วินโย จ ปริยตฺโต, เตน หิ ภนฺเต เถโร อายสฺมนฺตมฺปิ อานนฺทํ อุจฺจินตู”ติ.\csromanspacer{} อถ โข อายสฺมา มหากสฺสโป อายสฺมนฺตมฺปิ อานนฺทํ อุจฺจินิ.}

\prose{อถ โข เถรานํ ภิกฺขูนํ เอตทโหสิ “กตฺถ นุ โข มยํ ธมฺมญฺจ วินยญฺจ สงฺคาเยยฺยามา”ติ.\csromanspacer{} อถ โข เถรานํ ภิกฺขูนํ เอตทโหสิ “ราชคหํ โข มหาโคจรํ ปหูตเสนาสนํ, ยํนูน มยํ ราชคเห วสฺสํ วสนฺตา ธมฺมญฺจ วินยญฺจ สงฺคาเยยฺยาม, น อญฺเญ ภิกฺขู ราชคเห วสฺสํ อุปคจฺเฉยฺยุนฺติ.}

\csromanheader{2}{อถ โข อยสฺมา มหากสฺสโป สํฆํ ญาเปสิ\csromandash{}}
\proseitem{438}{“สุณาตุ เม อาวุโส สํโฆ, ยทิ สํฆสฺส ปตฺตกลฺลํ, สํโฆ อิมานิ ปญฺจ ภิกฺขุสตานิ สมฺมนฺเนยฺย ‘ราชคเห วสฺสํ วสนฺตานิ ธมฺมญฺจ วินยญฺจ สงฺคายิตุํ, น อญฺเญหิ ภิกฺขูหิ ราชคเห วสฺสํ วสิตพฺพนฺ’ติ, เอสา ญตฺติ.}

\prose{สุณาตุ เม อาวุโส สํโฆ, อิมานิ ปญฺจ ภิกฺขุสตานิ สมฺมนฺนติ ‘ราชคเห วสฺสํ วสนฺตานิ ธมฺมญฺจ วินยญฺจ สงฺคายิตุํ, น อญฺเญหิ ภิกฺขูหิ ราชคเห วสฺสํ วสิตพฺพนฺ’ติ, ยสฺสายสฺมโต ขมติ อิเมสํ ปญฺจนฺนํ ภิกฺขุสตานํ สมฺมุติ ‘ราชคเห วสฺสํ วสนฺตานํ ธมฺมญฺจ วินยญฺจ สงฺคายิตุํ น อญฺเญหิ ภิกฺขูหิ ราชคเห วสฺสํ วสิตพฺพนฺ’ติ, โส ตุณฺหสฺส, ยสฺส นกฺขมติ, โส ภาเสยฺย.}

\prose{\csromanfolio{482}สมฺมตานิ สํเฆน อิมานิ ปญฺจ ภิกฺขุสตานิ ‘ราชคเห วสฺสํ วสนฺตานิ ธมฺมญฺจ วินยญฺจ สงฺคายิตุํ น อญฺเญหิ ภิกฺขูหิ ราชคเห วสฺสํ วสิตพฺพนฺ’ติ, ขมติ สํฆสฺส, ตสฺมา ตุณฺหี, เอวเมตํ ธารยามี”ติ.}

\prose{อถ โข เถรา ภิกฺขู ราชคหํ อคมํสุ ธมฺมญฺจ วินยญฺจ สงฺคายิตุํ.\csromanspacer{} อถ โข เถรานํ ภิกฺขูนํ เอตทโหสิ “ภควตา โข อาวุโส ขณฺฑผุลฺลปฺปฏิสงฺขรณํ วณฺณิตํ, หนฺท มยํ อาวุโส ปฐมํ มาสํ ขณฺฑผุลฺลํ ปฏิสงฺขโรม, มชฺฌิมํ มาสํ สนฺนิปติตฺวา ธมฺมญฺจ วินยญฺจ สงฺคายิสฺสามา”ติ.}

\prose{อถ โข เถรา ภิกฺขู ปฐมํ มาสํ ขณฺฑผุลฺลํ ปฏิสงฺขาริํสุ.\csromanspacer{} อถ โข อายสฺมา อานนฺโท “เสฺว สนฺนิปาโต\footnote{สนฺนิปาโตติ (ก)}, น โข เมตํ ปติรูปํ, โยหํ เสกฺโข สมาโน สนฺนิปาตํ คจฺเฉยฺยนฺ”ติ พหุเทว รตฺติํ กายคตาย สติยา วีตินาเมตฺวา รตฺติยา ปจฺจูสสมยํ “นิปชฺชิสฺสามี”ติ กายํ อาวชฺเชสิ, อปฺปตฺตญฺจ สีสํ พิพฺโพหนํ, ภูมิโต จ ปาทา มุตฺตา, เอตสฺมิํ อนฺตเร อนุปาทาย อาสเวหิ จิตฺตํ วิมุจฺจิ.}

\proseitem{439}{อถ โข อายสฺมา อานนฺโท อรหา สมาโน สนฺนิปาตํ อคมาสิ.\csromanspacer{} อถ โข อายสฺมา มหากสฺสโป สํฆํ ญาเปสิ\csromandash{}}

\prose{“สุณาตุ เม อาวุโส สํโฆ, ยทิ สํฆสฺส ปตฺตกลฺลํ, อหํ อุปาลิํ วินยํ ปุจฺเฉยฺยนฺ”ติ.}

\csromanheader{2}{อายสฺมา อุปาลิ สํฆํ ญาเปสิ\csromandash{}}
\prose{“สุณาตุ เม ภนฺเต สํโฆ, ยทิ สํฆสฺส ปตฺตกลฺลํ, อหํ อายสฺมตา มหากสฺสเปน วินยํ ปุฏฺโฐ วิสฺสชฺเชยฺยนฺ”ติ.}

\prose{อถ โข อายสฺมา มหากสฺสโป อายสฺมนฺตํ อุปาลิํ เอตทโวจ “ปฐมํ อาวุโส อุปาลิ ปาราชิกํ กตฺถ ปญฺญตฺตนฺ”ติ.\csromanspacer{} เวสลิยํ ภนฺเตติ.\csromanspacer{} กํ อารพฺภาติ.\csromanspacer{} สุทินฺนํ กลนฺทปุตฺตํ อารพฺภาติ.\csromanspacer{} กิสฺมิํ วตฺถุสฺมินฺติ.\csromanspacer{} เมถุนธมฺเมติ.\csromanspacer{} อถ โข อายสฺมา มหากสฺสโป อายสฺมนฺตํ อุปาลิํ ปฐมสฺส ปาราชิกสฺส วตฺถุมฺปิ ปุจฺฉิ, นิทานมฺปิ ปุจฺฉิ, ปุคฺคลมฺปิ ปุจฺฉิ, ปญฺญตฺติมฺปิ ปุจฺฉิ, อนุปญฺญตฺติมฺปิ ปุจฺฉิ, อาปตฺติมฺปิ ปุจฺฉิ,}

\prose{\csromanfolio{483}อนาปตฺติมฺปิ ปุจฺฉิ.\csromanspacer{} ทุติยํ ปนาวุโส อุปาลิ ปาราชิกํ กตฺถ ปญฺญตฺตนฺติ.\csromanspacer{} ราชคเห ภนฺเตติ.\csromanspacer{} กํ อารพฺภาติ.\csromanspacer{} ธนิยํ กุมฺภการปุตฺตํ อารพฺภาติ.\csromanspacer{} กิสฺมิํ วตฺถุสฺมินฺติ.\csromanspacer{} อทินฺนาทาเนติ.\csromanspacer{} อถ โข อายสฺมา มหากสฺสโป อายสฺมนฺตํ อุปาลิํ ทุติยสฺส ปาราชิกสฺส วตฺถุมฺปิ ปุจฺฉิ, นิทานมฺปิ ปุจฺฉิ, ปุคฺคลมฺปิ ปุจฺฉิ, ปญฺญตฺติมฺปิ ปุจฺฉิ, อนุปญฺญตฺติมฺปิ ปุจฺฉิ, อาปตฺติมฺปิ ปุจฺฉิ, อนาปตฺติมฺปิ ปุจฺฉิ.\csromanspacer{} ตติยํ ปนาวุโส อุปาลิ ปาราชิกํ กตฺถ ปญฺญตฺตนฺติ.\csromanspacer{} เวสาลิยํ ภนฺเตติ.\csromanspacer{} กํ อารพฺภาติ.\csromanspacer{} สมฺพหุเล ภิกฺขู อารพฺภาติ.\csromanspacer{} กิสฺมิํ วตฺถุสฺมินฺติ.\csromanspacer{} มนุสฺสวิคฺคเหติ.\csromanspacer{} อถ โข อายสฺมา มหากสฺสโป อายสฺมนฺตํ อุปาลิํ ตติยสฺส ปาราชิกสฺส วตฺถุมฺปิ ปุจฺฉิ, นิทานมฺปิ ปุจฺฉิ, ปุคฺคลมฺปิ ปุจฺฉิ, ปญฺญตฺติมฺปิ ปุจฺฉิ, อนุปญฺญตฺติมฺปิ ปุจฺฉิ, อาปตฺติมฺปิ ปุจฺฉิ, อนาปตฺติมฺปิ ปุจฺฉิ.\csromanspacer{} จตุตฺถํ ปนาวุโส อุปาลิ ปาราชิกํ กตฺถ ปญฺญตฺตนฺติ.\csromanspacer{} เวสาลิยํ ภนฺเตติ.\csromanspacer{} กํ อารพฺภาติ.\csromanspacer{} วคฺคุมุทาตีริเย ภิกฺขู อารพฺภาติ.\csromanspacer{} กิสฺมิํ วตฺถุสฺมินฺติ.\csromanspacer{} อุตฺตริมนุสฺสธมฺเมติ.\csromanspacer{} อถ โข อายสฺมา มหากสฺสโป อายสฺมนฺตํ อุปาลิํ จตุตฺถสฺส ปาราชิกสฺส วตฺถุมฺปิ ปุจฺฉิ, นิทานมฺปิ ปุจฺฉิ, ปุคฺคลมฺปิ ปุจฺฉิ, ปญฺญตฺติมฺปิ ปุจฺฉิ, อนุปญฺญตฺติมฺปิ ปุจฺฉิ, อาปตฺติมฺปิ ปุจฺฉิ, อนาปตฺติมฺปิ ปุจฺฉิ.\csromanspacer{} เอเตเนว อุปาเยน อุภโตวิภงฺเค ปุจฺฉิ, ปุฏฺโฐ ปุฏฺโฐ อายสฺมา อุปาลิ วิสฺสชฺเชสิ.}

\csromanheader{2}{440. อถ โข อายสฺมา มหากสฺสโป สํฆํ ญาเปสิ\csromandash{}}
\prose{“สุณาตุ เม อาวุโส สํโฆ, ยทิ สํฆสฺส ปตฺตกลฺลํ, อหํ อานนฺทํ ธมฺมํ ปุจฺเฉยฺยนฺ”ติ.}

\csromanheader{2}{อายสฺมา อานนฺโท สํฆํ ญาเปสิ\csromandash{}}
\prose{“สุณาตุ เม ภนฺเต สํโฆ, ยทิ สํฆสฺส ปตฺตกลฺลํ, อหํ อายสฺมตา มหากสฺสเปน ธมฺมํ ปุฏฺโฐ วิสฺสชฺเชยฺยนฺ”ติ.}

\prose{อถ โข อายสฺมา มหากสฺสโป อายสฺมนฺตํ อานนฺทํ เอตทโวจ “พฺรหฺมชาลํ อาวุโส อานนฺท กตฺถ ภาสิตนฺ”ติ.\csromanspacer{} อนฺตรา จ ภนฺเต ราชคหํ, อนฺตรา จ นาฬนฺทํ ราชาคารเก อมฺพลฏฺฐิกายาติ.\csromanspacer{} กํ อารพฺภาติ.\csromanspacer{} สุปฺปิยญฺจ ปริพฺพาชกํ พฺรหฺมทตฺตญฺจ มาณวนฺติ.\csromanspacer{} อถ โข อายสฺมา มหากสฺสโป อายสฺมนฺตํ อานนฺทํ พฺรหฺมชาลสฺส นิทานมฺปิ ปุจฺฉิ, ปุคฺคลมฺปิ ปุจฺฉิ.\csromanspacer{} สามญฺญผลํ ปนาวุโส อานนฺท กตฺถ ภาสิตนฺติ.\csromanspacer{} ราชคเห ภนฺเต ชีวกมฺพวเนติ.\csromanspacer{} เกน สทฺธินฺติ.\csromanspacer{} อชาตสตฺตุนา}

\prose{\csromanfolio{484}เวเทหิปุตฺเตน สทฺธินฺติ.\csromanspacer{} อถ โข อายสฺมา มหากสฺสโป อายสฺมนฺตํ อานนฺทํ สามญฺญผลสฺส นิทานมฺปิ ปุจฺฉิ, ปุคฺคลมฺปิ ปุจฺฉิ.\csromanspacer{} เอเตเนว อุปาเยน ปญฺจปิ นิกาเย ปุจฺฉิ, ปุฏฺโฐ ปุฏฺโฐ อายสฺมา อานนฺโท วิสฺสชฺเชสิ.}

\csromanmatikaanchor{mtk.214}
\csromantocmarkref{section}{2. ขุทฺทานุขุทฺทกสิกฺขาปทกถา}{mtk.214}
\bookmark[dest={mtk.214},level=1]{2. ขุทฺทานุขุทฺทกสิกฺขาปทกถา}
\csromanheader{1}{2. ขุทฺทานุขุทฺทกสิกฺขาปทกถา}
\proseitem{441}{อถ โข อายสฺมา อานนฺโท เถเร ภิกฺขู เอตทโวจ “ภควา มํ ภนฺเต ปรินิพฺพานกาเล เอวมาห ‘อากงฺขมาโน อานนฺท สํโฆ มมจฺจเยน ขุทฺทานุขุทฺทกานิ สิกฺขาปทานิ สมูหเนยฺยา’ติ”.\csromanspacer{} ปุจฺฉิ ปุน ตฺวํ อาวุโส อานนฺท ภควนฺตํ “กตมานิ ปน ภนฺเต ขุทฺทานุขุทฺทกานิ สิกฺขาปทานี”ติ.\csromanspacer{} น โขหํ ภนฺเต ภควนฺตํ ปุจฺฉิํ “กตมานิ ปน ภนฺเต ขุทฺทานุขุทฺทกานิ สิกฺขาปทานี”ติ.\csromanspacer{} เอกจฺเจ เถรา เอวมาหํสุ “จตฺตาริ ปาราชิกานิ ฐเปตฺวา อวเสสานิ ขุทฺทานุขุทฺทกานิ สิกฺขาปทานี”ติ.\csromanspacer{} เอกจฺเจ เถรา เอวมาหํสุ “จตฺตาริ ปาราชิกานิ ฐเปตฺวา เตรส สํฆาทิเสเส ฐเปตฺวา อวเสสานิ ขุทฺทานุขุทฺทกานิ สิกฺขาปทานี”ติ.\csromanspacer{} เอกจฺเจ เถรา เอวมาหํสุ “จตฺตาริ ปาราชิกานิ ฐเปตฺวา เตรส สํฆาทิเสเส ฐเปตฺวา เทฺว อนิยเต ฐเปตฺวา อวเสสานิ ขุทฺทานุขุทฺทกานิ สิกฺขาปทานี”ติ.\csromanspacer{} เอกจฺเจ เถรา เอวมาหํสุ “จตฺตาริ ปาราชิกานิ ฐเปตฺวา เตรส สํฆาทิเสเส ฐเปตฺวา เทฺว อนิยเต ฐเปตฺวา ติํส นิสฺสคฺคิเย ปาจิตฺติเย ฐเปตฺวา อวเสสานิ ขุทฺทานุขุทฺทกานิ สิกฺขาปทานี”ติ.\csromanspacer{} เอกจฺเจ เถรา เอวมาหํสุ “จตฺตาริ ปาราชิกานิ ฐเปตฺวา เตรส สํฆาทิเสเส ฐเปตฺวา เทฺว อนิยเต ฐเปตฺวา ติํส นิสฺสคฺคิเย ปาจิตฺติเย ฐเปตฺวา เทฺวนวุติ ปาจิตฺติเย ฐเปตฺวา อวเสสานิ ขุทฺทานุขุทฺทกานิ สิกฺขาปทานี”ติ.\csromanspacer{} เอกจฺเจ เถรา เอวมาหํสุ “จตฺตาริ ปาราชิกานิ ฐเปตฺวา เตรส สํฆาทิเสเส ฐเปตฺวา เทฺว อนิยเต ฐเปตฺวา ติํส นิสฺสคฺคิเย ปาจิตฺติเย ฐเปตฺวา เทฺวนวุติ ปาจิตฺติเย ฐเปตฺวา จตฺตาโร ปาฏิเทสนีเย ฐเปตฺวา อวเสสานิ ขุทฺทานุขุทฺทกานิ สิกฺขาปทานี”ติ.}

\csromanheader{2}{442. อถ โข อายสฺมา มหากสฺสโป สํฆํ ญาเปสิ\csromandash{}}
\prose{“สุณาตุ เม อาวุโส สํโฆ, สนฺตมฺหากํ สิกฺขาปทานิ คิหิคตานิ คิหิโนปิ ชานนฺติ ‘อิทํ โว สมณานํ สกฺยปุตฺติยานํ กปฺปติ, อิทํ โว น กปฺปตี’ติ.\csromanspacer{} สเจ มยํ ขุทฺทานุขุทฺทกานิ สิกฺขาปทานิ สมูหนิสฺสาม,}

\prose{\csromanfolio{485}ภวิสฺสนฺติ วตฺตาโร ‘ธูมกาลิกํ สมเณน โคตเมน สาวกานํ สิกฺขาปทํ ปญฺญตฺตํ, ยาวิเมสํ สตฺถา อฏฺฐาสิ, ตาวิเม สิกฺขาปเทสุ สิกฺขิํสุ, ยโต อิเมสํ สตฺถา ปรินิพฺพุโต, น ทานิเม สิกฺขาปเทสุ สิกฺขนฺตี’ติ, ยทิ สํฆสฺส ปตฺตกลฺลํ, สํโฆ อปฺปญฺญตฺตํ นปฺปญฺญเปยฺย, ปญฺญตฺตํ น สมุจฺฉินฺเทยฺย, ยถาปญฺญตฺเตสุ สิกฺขาปเทสุ สมาทาย วตฺเตยฺย, เอสา ญตฺติ.}

\prose{สุณาตุ เม อาวุโส สํโฆ, สนฺตมฺหากํ สิกฺขาปทานิ คิหิคตานิ คิหิโนปิ ชานนฺติ ‘อิทํ โว สมณานํ สกฺยปุตฺติยานํ กปฺปติ, อิทํ โว น กปฺปตี”ติ, สเจ มยํ ขุทฺทานุขุทฺทกานิ สิกฺขาปทานิ สมูหนิสฺสาม, ภวิสฺสนฺติ วตฺตาโร ‘ธูมกาลิกํ สมเณน โคตเมน สาวกานํ สิกฺขาปทํ ปญฺญตฺตํ, ยาวิเมสํ สตฺถา อฏฺฐาสิ, ตาวิเม สิกฺขาปเทสุ สิกฺขิํสุ, ยโต อิเมสํ สตฺถา ปรินิพฺพุโต, น ทานิเม สิกฺขาปเทสุ สิกฺขนฺตี’ติ, สํโฆ อปฺปญฺญตฺตํ นปฺปญฺญเปติ, ปญฺญตฺตํ น สมุจฺฉินฺทติ, ยถาปญฺญตฺเตสุ สิกฺขาปเทสุ สมาทาย วตฺตติ, ยสฺสายสฺมโต ขมติ อปฺปญฺญตฺตสฺส อปฺปญฺญาปนา, ปญฺญตฺตสฺส อสมุจฺเฉโท, ยถาปญฺญตฺเตสุ สิกฺขาปเทสุ สมาทาย วตฺตนา, โส ตุณฺหสฺส, ยสฺส นกฺขมติ, โส ภาเสยฺย.}

\prose{สํโฆ อปฺปญฺญตฺตํ นปฺปญฺญเปติ, ปญฺญตฺตํ น สมุจฺฉินฺทติ, ยถาปญฺญตฺเตสุ สิกฺขาปเทสุ สมาทาย วตฺตติ, ขมติ สํฆสฺส, ตสฺมา ตุณฺหี, เอวเมตํ ธารยามี”ติ.}

\proseitem{443}{อถ โข เถรา ภิกฺขู อายสฺมนฺตํ อานนฺทํ เอตทโวจุํ “อิทํ เต อาวุโส อานนฺท ทุกฺกฏํ, ยํ ตฺวํ ภควนฺตํ น ปุจฺฉิ ‘กตมานิ ปน ภนฺเต ขุทฺทานุขุทฺทกานิ สิกฺขาปทานี’ติ, เทเสหิ ตํ ทุกฺกฏนฺ”ติ.\csromanspacer{} อหํ โข ภนฺเต อสฺสติยา ภควนฺตํ น ปุจฺฉิํ “กตมานิ ปน ภนฺเต ขุทฺทานุขุทฺทกานิ สิกฺขาปทานี”ติ, นาหํ ตํ ทุกฺกฏํ ปสฺสามิ, อปิ จายสฺมนฺตานํ สทฺธาย เทเสมิ ตํ ทุกฺกฏนฺติ.\csromanspacer{} อิทมฺปิ เต อาวุโส อานนฺท ทุกฺกฏํ, ยํ ตฺวํ ภควโต วสฺสิกสาฏิกํ อกฺกมิตฺวา สิพฺเพสิ, เทเสหิ ตํ ทุกฺกฏนฺติ.\csromanspacer{} อหํ โข ภนฺเต น อคารเวน ภควโต วสฺสิกสาฏิกํ อกฺกมิตฺวา สิพฺเพสิํ, นาหํ ตํ ทุกฺกฏํ ปสฺสามิ, อปิ จายสฺมนฺตานํ สทฺธาย เทเสมิ ตํ ทุกฺกฏนฺติ.\csromanspacer{} อิทมฺปิ เต อาวุโส อานนฺท ทุกฺกฏํ, ยํ ตฺวํ มาตุคาเมหิ \csromanfolio{486}ภควโต สรีรํ ปฐมํ วนฺทาเปสิ, ตาสํ โรทนฺตีนํ ภควโต สรีรํ อสฺสุเกน มกฺขิตํ, เทเสหิ ตํ ทุกฺกฏนฺติ.\csromanspacer{} อหํ โข ภนฺเต “มายิมาสํ\footnote{มายิมา (สี, สฺยา)} วิกาเล อเหสุนฺ”ติ มาตุคาเมหิ ภควโต สรีรํ ปฐมํ วนฺทาเปสิํ, นาหํ ตํ ทุกฺกฏํ ปสฺสามิ, อปิ จายสฺมนฺตานํ สทฺธาย เทเสมิ ตํ ทุกฺกฏนฺติ.\csromanspacer{} อิทมฺปิ เต อาวุโส อานนฺท ทุกฺกฏํ, ยํ ตฺวํ ภควโต โอฬาริเก นิมิตฺเต กยิรมาเน โอฬริเก โอภาเส กยิรมาเน น ภควนฺตํ ยาจิ “ติฏฺฐตุ ภควา กปฺปํ, ติฏฺฐตุ สุคโต กปฺปํ พหุชนหิตาย พหุชนสุขาย โลกานุกมฺปาย อตฺถาย หิตาย สุขาย เทวมนุสฺสานนฺ”ติ, เทเสหิ ตํ ทุกฺกฏนฺติ.\csromanspacer{} อหํ โข ภนฺเต มาเรน ปริยุฏฺฐิตจิตฺโต น ภควนฺตํ ยาจิํ “ติฏฺฐตุ ภควา กปฺปํ, ติฏฺฐตุ สุคโต กปฺปํ พหุชนหิตาย พหุชนสุขาย โลกานุกมฺปาย อตฺถาย หิตาย สุขาย เทวมนุสฺสานนฺ”ติ, นาหํ ตํ ทุกฺกฏํ ปสฺสามิ, อปิ จายสฺมนฺตานํ สทฺธาย เทเสมิ ตํ ทุกฺกฏนฺติ.\csromanspacer{} อิทมฺปิ เต อาวุโส อานนฺท ทุกฺกฏํ, ยํ ตฺวํ มาตุคามสฺส ตถาคตปฺปเวทิเต ธมฺมวินเย ปพฺพชฺชํ อุสฺสุกฺกํ อกาสิ, เทเสหิ ตํ ทุกฺกฏนฺติ.\csromanspacer{} อหํ โข ภนฺเต “อยํ มหาปชาปติ โคตมี ภควโต มาตุจฺฉา อาปาทิกา โปสิกา ขีรสฺส ทายิกา ภควนฺตํ ชเนตฺติยา กาลงฺกตาย ถญฺญํ ปาเยสี”ติ มาตุคามสฺส ตถาคตปฺปเวทิเต ธมฺมวินเย ปพฺพชฺชํ อุสฺสุกฺกํ อกาสิํ, นาหํ ตํ ทุกฺกฏํ ปสฺสามิ, อปิ จายสฺมนฺตานํ สทฺธาย เทเสมิ ตํ ทุกฺกฏนฺติ.}

\proseitem{444}{เตน โข ปน สมเยน อายสฺมา ปุราโณ ทกฺขิณาคิริสฺมิํ จาริกํ จรติ มหตา ภิกฺขุสํเฆน สทฺธิํ ปญฺจมตฺเตหิ ภิกฺขุสเตหิ.\csromanspacer{} อถ โข อายสฺมา ปุราโณ เถเรหิ ภิกฺขูหิ ธมฺเม จ วินเย จ สงฺคีเต ทกฺขิณาคิริสฺมิํ ยถาภิรนฺตํ วิหริตฺวา เยน ราชคหํ เยน เวฬุวนํ กลนฺทกนิวาโป เยน เถรา ภิกฺขู เตนุปสงฺกมิ, อุปสงฺกมิตฺวา เถเรหิ ภิกฺขูหิ สทฺธิํ ปฏิสมฺโมทิตฺวา เอกมนฺตํ นิสีทิ, เอกมนฺตํ นิสินฺนํ โข อายสฺมนฺตํ ปุราณํ เถรา ภิกฺขู เอตทโวจุํ “เถเรหิ อาวุโส ปุราณ ธมฺโม จ วินโย จ สงฺกีโต, อุเปหิ ตํ สงฺคีตินฺ”ติ.\csromanspacer{} สุสงฺคีตาวุโส เถเรหิ ธมฺโม จ วินโย จ, อปิ จ ยเถว มยา ภควโต สมฺมุขา สุตํ สมฺมุขา ปฏิคฺคหิตํ, ตเถวาหํ ธาเรสฺสามีติ.}

\csromanmatikaanchor{mtk.215}
\csromantocmarkref{section}{3. พฺรหฺมทณฺฑกถา}{mtk.215}
\bookmark[dest={mtk.215},level=1]{3. พฺรหฺมทณฺฑกถา}
\csromanheader{1}{3. พฺรหฺมทณฺฑกถา}
\proseitem{445}{\csromanfolio{487}อถ โข อายสฺมา อานนฺโท เถเร ภิกฺขู เอตทโวจ “ภควา มํ ภนฺเต ปรินิพฺพานกาเล เอวมาห ‘เตน หานนฺท สํโฆ มมจฺจเยน ฉนฺนสฺส ภิกฺขุโน พฺรหฺมทณฺฑํ อาณาเปตู’ติ”.\csromanspacer{} ปุจฺฉิ ปน ตฺวํ อาวุโส อานนฺท ภควนฺตํ “กตโม ปน ภนฺเต พฺรหฺมทณฺโฑ”ติ.\csromanspacer{} ปุจฺฉิํ โขหํ ภนฺเต ภควนฺตํ “กตโม ปน ภนฺเต พฺรหฺมทณฺโฑ”ติ.\csromanspacer{} ฉนฺโน อานนฺท ภิกฺขุ ยํ อิจฺเฉยฺย ตํ วเทยฺย, ภิกฺขูหิ ฉนฺโน ภิกฺขุ เนว วตฺตพฺโพ, น โอวทิตพฺโพ, นานุสาสิตพฺโพติ.\csromanspacer{} เตน หาวุโส อานนฺท ตฺวํเยว ฉนฺนสฺส ภิกฺขุโน พฺรหฺมทณฺฑํ อาณาเปหีติ.\csromanspacer{} กถาหํ ภนฺเต ฉนฺนสฺส ภิกฺขุโน พฺรหฺมทณฺฑํ อาณาเปมิ, จณฺโฑ โส ภิกฺขุ ผรุโสติ.\csromanspacer{} เตน หาวุโส อานนฺท พหุเกหิ ภิกฺขูหิ สทฺธิํ คจฺฉาหีติ. “เอวํ ภนฺเต”ติ โข อายสฺมา อานนฺโท เถรานํ ภิกฺขูนํ ปฏิสฺสุตฺวา มหตา ภิกฺขุสํเฆน สทฺธิํ ปญฺจมตฺเตหิ ภิกฺขุสเตหิ นาวาย อุชฺชวนิกาย โกสมฺพิํ อุชฺชวิ, นาวาย ปจฺโจโรหิตฺวา รญฺโญ อุเทนสฺส\footnote{อุเตนสฺส (ก)} อุยฺยานสฺส อวิทูเร อญฺญตรสฺมิํ รุกฺขมูเล นิสีทิ.\csromanspacer{} เตน โข ปน สมเยน ราชา อุเทโน อุยฺยาเน ปริจาเรสิ สทฺธิํ โอโรเธน.\csromanspacer{} อสฺโสสิ โข รญฺโญ อุเทนสฺส โอโรโธ “อมฺหากํ กิร อาจริโย อยฺโย อานนฺโท อุยฺยานสฺส อวิทูเร อญฺญตรสฺมิํ รุกฺขมูเล นิสินฺโน”ติ.\csromanspacer{} อถ โข รญฺโญ อุเทนสฺส โอโรโธ ราชานํ อุเทนํ เอตทโวจ “อมฺหากํ กิร เทว อาจริโย อยฺโย อานนฺเทว อุยฺยานสฺส อวิทูเร อญฺญตรสฺมิํ รุกฺขมูเล นิสินฺโน, อิจฺฉาม มยํ เทว อยฺยํ อานนฺทํ ปสฺสิตุนฺ”ติ.\csromanspacer{} เตน หิ ตุเมฺห สมณํ อานนฺทํ ปสฺสถาติ.}

\prose{อถ โข รญฺโญ อุเทนสฺส โอโรโธ เยนายสฺมา อานนฺโท เตนุปสงฺกมิ, อุปสงฺกมิตฺวา อายสฺมนฺตํ อานนฺทํ อภิวาเทตฺวา เอกมนฺตํ นิสีทิ.\csromanspacer{} เอกมนฺตํ นิสินฺนํ โข รญฺโญ อุเทนสฺส โอโรธํ อายสฺมา อานนฺโท ธมฺมิยา กถาย สนฺทสฺเสสิ สมาทเปสิ สมุตฺเตเชสิ สมฺปหํเสสิ.\csromanspacer{} อถ โข รญฺโญ อุเทนสฺส โอโรโธ อายสฺมตา อานนฺเทน ธมฺมิยา กถาย สนฺทสฺสิโต สมาทปิโต สมุตฺเตชิโต สมฺปหํสิโต อายสฺมโต อานนฺทสฺส ปญฺจ อุตฺตราสงฺคสตานิ ปาทาสิ. \csromanfolio{488}อถ โข รญฺโญ อุเทนสฺส โอโรโธ อายสฺมโต อานนฺทสฺส ภาสิตํ อภินนฺทิตฺวา อนุโมทิตฺวา อุฏฺฐายาสนา อายสฺมนฺตํ อานนฺทํ อภิวาเทตฺวา ปทกฺขิณํ กตฺวา เยน ราชา อุเทโน เตนุปสงฺกมิ.\csromanspacer{} อทฺทสา โข ราชา อุเทโน โอโรธํ ทูรโตว อาคจฺฉนฺตํ, ทิสฺวาน โอโรธํ เอตทโวจ “อปิ นุ โข ตุเมฺห สมณํ อานนฺทํ ปสฺสิตฺถา”ติ.\csromanspacer{} อปสฺสิมฺหา โข มยํ เทว อยฺยํ อานนฺทนฺติ.\csromanspacer{} อปิ นุ ตุเมฺห สมณสฺส อานนฺทสฺส กิญฺจิ อทตฺถาติ.\csromanspacer{} อทมฺหา โข มยํ เทว อยฺยสฺส อานนฺทสฺส ปญฺจ อุตฺตราสงฺคสตานีติ.\csromanspacer{} ราชา อุเทโน อุชฺฌายติ ขิยฺยติ วิปาเจติ “กถํ หิ นาม สมโณ อานนฺโท ตาว พหุํ จีวรํ ปฏิคฺคเหสฺสติ, ทุสฺสวาณิชฺชํ วา สมโณ อานนฺโท กริสฺสติ, ปคฺคาหิกสาลํ วา ปสาเรสฺสตี”ติ.}

\prose{อถ โข ราชา อุเทโน เยนายสฺมา อานนฺโท เตนุปสงฺกมิ, อุปสงฺกมิตฺวา อายสฺมตา อานนฺเทน สทฺธิํ สมฺโมทิ, สมฺโมทนียํ กถํ สารณียํ วีติสาเรตฺวา เอกมนฺตํ นิสีทิ, เอกมนฺตํ นิสินฺโน โข ราชา อุเทโน อายสฺมนฺตํ อานนฺทํ เอตทโวจ “อาคมา นุ ขฺวิธ โภ อานนฺท อมฺหากํ โอโรโธ”ติ.\csromanspacer{} อาคมาสิ โข เต อิธ มหาราช โอโรโธติ.\csromanspacer{} อปิ ปน โภโต อานนฺทสฺส กิญฺจิ อทาสีติ.\csromanspacer{} อทาสิ โข เม มหาราช ปญฺจ อุตฺตราสงฺคสตานีติ.\csromanspacer{} กิํ ปน ภวํ อานนฺโท ตาว พหุํ จีวรํ กริสฺสตีติ.\csromanspacer{} เย\footnote{เย ปน (ก)} เต มหาราช ภิกฺขู ทุพฺพลจีวรา, เตหิ สทฺธิํ สํวิภชิสฺสามีติ.\csromanspacer{} ยานิ โข ปน โภ อานนฺท โปราณกานิ ทุพฺพลจีวรานิ, ตานิ กถํ กริสฺสถาติ.\csromanspacer{} ตานิ มหาราช อุตฺตรตฺถรณํ กริสฺสามาติ.\csromanspacer{} ยานิ ปน โภ อานนฺท โปราณกานิ อุตฺตรตฺถรณานิ, ตานิ กถํ กริสฺสถาติ.\csromanspacer{} ตานิ มหาราช ภิสิจฺฉวิโย กริสฺสามาติ.\csromanspacer{} ยา ปน โภ อานนฺท โปราณกา ภิสิจฺฉวิโย, ตา กถํ กริสฺสถาติ.\csromanspacer{} ตา มหาราช ภูมตฺถรณํ กริสฺสามาติ.\csromanspacer{} ยานิ ปน โภ อานนฺท โปราณกานิ ภูมตฺถรณานิ, ตานิ กถํ กริสฺสถาติ.\csromanspacer{} ตานิ มหาราช ปาทปุญฺฉนิโย กริสฺสามาติ.\csromanspacer{} ยา ปน โภ อานนฺท โปราณกา ปาทปุญฺฉนิโย, ตา กถํ กริสฺสถาติ.\csromanspacer{} ตา มหาราช รโชหรณํ กริสฺสมาติ.\csromanspacer{} ยานิ ปน \csromanfolio{489}โภ อานนฺท โปราณกานิ รโชหรณานิ, ตานิ กถํ กริสฺสถาติ.\csromanspacer{} ตานิ มหาราช โกฏฺเฏตฺวา จิกฺขลฺเลน มทฺทิตฺวา ปริภณฺฑํ ลิมฺปิสฺสามาติ.}

\prose{อถ โข ราชา อุเทโน “สพฺเพวิเม สมณา สกฺยปุตฺติยา โยนิโส อุปเนนฺติ, น กุลวํ คเมนฺตี”ติ อายสฺมโต อานนฺทสฺส อญฺญานิปิ ปญฺจ ทุสฺสสตานิ ปาทาสิ.\csromanspacer{} อยญฺจรหิ อายสฺมโต อานนฺทสฺส ปฐมํ จีวรภิกฺขา อุปฺปชฺชิ จีวรสหสฺสํ.\csromanspacer{} อถ โข อายสฺมา อานนฺโท เยน โฆสิตาราโม เตนุปสงฺกมิ, อุปสงฺกมิตฺวา ปญฺญตฺเต อาสเน นิสีทิ.\csromanspacer{} อถ โข อายสฺมา ฉนฺโน เยนายสฺมา อานนฺโท เตนุปสงฺกมิ, อุปสงฺกมิตฺวา อายสฺมนฺตํ อานนฺทํ อภิวาเทตฺวา เอกมนฺตํ นิสีทิ, เอกมนฺตํ นิสินฺนํ โข อายสฺมนฺตํ ฉนฺนํ อายสฺมา อานนฺโท เอตทโวจ “สํเฆน เต อาวุโส ฉนฺน พฺรหฺมทณฺโฑ อาณาปิโต”ติ.}

\prose{กตโม ปน ภนฺเต อานนฺท พฺรหฺมทณฺโฑ อาณาปิโตติ.\csromanspacer{} ตฺวํ อาวุโส ฉนฺน ภิกฺขู ยํ อิจฺเฉยฺยาสิ ตํ วเทยฺยาสิ, ภิกฺขูหิ ตฺวํ เนว วตฺตพฺโพ, น โอวทิตพฺโพ, นานุสาสิตพฺโพติ. “นนฺวาหํ ภนฺเต อานนฺท หโต เอตฺตาวตา, ยโตหํ ภิกฺขูหิ เนว วตฺตพฺโพ, น โอวทิตพฺโพ, นานุสาสิตพฺโพ”ติ ตตฺเถว มุจฺฉิโต ปปโต.\csromanspacer{} อถ โข อายสฺมา ฉนฺโน พฺรหฺมทณฺเฑน อฏฺฏียมาโน หรายมาโน ชิคุจฺฉมาโน เอโก วูปกฏฺโฐ อปฺปมตฺโต อาตาปิ ปหิตตฺโต วิหรนฺโต นจิรสฺเสว ยสฺสตฺถาย กุลปุตฺตา สมฺมเทว อคารสฺมา อนคาริยํ ปพฺพชนฺติ, ตทนุตฺตรํ พฺรหฺมจริยปริโยสานํ ทิฏฺเฐว ธมฺเม สยํ อภิณฺญา สจฺฉิกตฺวา อุปสมฺปชฺช วิหาสิ, “ขีณา ชาติ, วุสิตํ พฺรหฺมจริยํ, กตํ กรณียํ, นาปรํ อิตฺถตฺตายา”ติ อพฺภญฺญาสิ, อญฺญตโร จ ปนายสฺมา ฉนฺโน อรหตํ อโหสิ.\csromanspacer{} อถ โข อายสฺมา ฉนฺโน อรหตฺตํ ปตฺโต เยนายสฺมา อานนฺโท เตนุปสงฺกมิ, อุปสงฺกมิตฺวา อายสฺมนฺตํ อานนฺทํ เอตทโวจ “ปฏิปฺปสฺสมฺเภหิ ทานิ เม ภนฺเต อานนฺท พฺรหฺมทณฺฑนฺ”ติ.\csromanspacer{} ยทคฺเคน ตยา อาวุโส ฉนฺน อรหตฺตํ สจฺฉิกตํ, ตทคฺเคน เต พฺรหฺมทณฺโฑ ปฏิปฺปสฺสทฺโธติ.\csromanspacer{} อิมาย โข ปน วินยสงฺคีติยา ปญฺจ ภิกฺขุสตานิ อนูนานิ อนธิกานิ อเหสุํ.\csromanspacer{} ตสฺมา อยํ วินยสงฺคีติ “ปญฺจสติกา”ติ วุจฺจตีติ.}

\vspace{\tipitakaparskip}
\csromancenter{ปญฺจสติกกฺขนฺธโก เอกาทสโม.}
\csromancenter{อิมมฺหิ ขนฺธเก วตฺถู เตวีสติ.}
\titlehead{\textbf{ตสฺสุทฺทานํ}}
\csromanmatikaanchor{mtk.216}
\csromantocmarkref{section}{อุทฺทานคาถา}{mtk.216}
\bookmark[dest={mtk.216},level=1]{อุทฺทานคาถา}
\csromangathameasure{ปรินิพฺพุเต สมฺพุทฺเธ, เถโร กสฺสปสวฺหโย. \\ อามนฺตยิ ภิกฺขุคณํ, สทฺธมฺมมนุปาลโก. \\ ปาวายทฺธานมคฺคมฺหิ, สุภทฺเทน ปเวทิตํ. \\ สงฺคายิสฺสาม สทฺธมฺมํ, อธมฺโม ปุเร ทิปฺปติ. \\ เอเกนูน ปญฺจสตํ, อานนฺทมฺปิ จ อุจฺจินิ. \\ ธมฺมวินยสงฺคีติํ, วสนฺโต คุหมุตฺตเม. \\ อุปาลิํ วินยํ ปุจฺฉิ, สุตฺตนฺตานนฺทปณฺฑิตํ. \\ ปิฏกํ ตีณิ สงฺคีติํ, อกํสุ ชินสาวกา. \\ ขุทฺทานุขุทฺทเก นานา, ยถาปญฺญตฺติวตฺตนา. \\ น ปุจฺฉิ อกฺกมิตฺวาน, วนฺทาเปสิ น ยาจิ จ. \\ ปพฺพชฺชํ มาตุคามสฺส, สทฺธาย ทุกฺกฏานิ เม. \\ ปุราโณ พฺรหฺมทณฺฑญฺจ, โอโรโธ อุเทเนน สห. \\ ตาว พหุ ทุพฺพลญฺจ, อุตฺตรตฺถรณา ภิสิ. \\ ภูมตฺถรณา ปุญฺฉนิโย, รโช จิกฺขลฺลมทฺทนา. \\ สหสฺสจีวรํ อุปฺปชฺชิ, ปฐมานนฺทสวฺหโย. \\ ตชฺชิโต พฺรหฺมทณฺเฑน, จตุสฺสจฺจํ อปาปุณิ. \\ วสีภูตา ปญฺจสตา, ตสฺมา “ปญฺจสตี” อิติ.}
\csromangathagroup{\csromangathastanza{\csromanfolio{490}ปรินิพฺพุเต สมฺพุทฺเธ, เถโร กสฺสปสวฺหโย. \\ อามนฺตยิ ภิกฺขุคณํ, สทฺธมฺมมนุปาลโก.}\csromangathastanza{ปาวายทฺธานมคฺคมฺหิ, สุภทฺเทน ปเวทิตํ. \\ สงฺคายิสฺสาม สทฺธมฺมํ, อธมฺโม ปุเร ทิปฺปติ.}\csromangathastanza{เอเกนูน ปญฺจสตํ, อานนฺทมฺปิ จ อุจฺจินิ. \\ ธมฺมวินยสงฺคีติํ, วสนฺโต คุหมุตฺตเม.}\csromangathastanza{อุปาลิํ วินยํ ปุจฺฉิ, สุตฺตนฺตานนฺทปณฺฑิตํ. \\ ปิฏกํ ตีณิ สงฺคีติํ, อกํสุ ชินสาวกา.}\csromangathastanza{ขุทฺทานุขุทฺทเก นานา, ยถาปญฺญตฺติวตฺตนา. \\ น ปุจฺฉิ อกฺกมิตฺวาน, วนฺทาเปสิ น ยาจิ จ.}\csromangathastanza{ปพฺพชฺชํ มาตุคามสฺส, สทฺธาย ทุกฺกฏานิ เม. \\ ปุราโณ พฺรหฺมทณฺฑญฺจ, โอโรโธ อุเทเนน สห.}\csromangathastanza{ตาว พหุ ทุพฺพลญฺจ, อุตฺตรตฺถรณา ภิสิ. \\ ภูมตฺถรณา ปุญฺฉนิโย, รโช จิกฺขลฺลมทฺทนา.}\csromangathastanza{สหสฺสจีวรํ อุปฺปชฺชิ, ปฐมานนฺทสวฺหโย. \\ ตชฺชิโต พฺรหฺมทณฺเฑน, จตุสฺสจฺจํ อปาปุณิ. \\ วสีภูตา ปญฺจสตา, ตสฺมา “ปญฺจสตี” อิติ.}}

\nitthitam{\textbf{ปญฺจสติกกฺขนฺธโก นิฏฺฐิโต.}}
\csromanmatikaanchor{mtk.217}
\csromantocmarknopage{chapter}{12. สตฺตสติกกฺขนฺธก}{mtk.217}
\bookmark[dest={mtk.217},level=0]{12. สตฺตสติกกฺขนฺธก}
\chapterheadpageread{12. สตฺตสติกกฺขนฺธก}
\csromanmatikaanchor{mtk.218}
\csromantocmarkref{section}{1. ปฐมภาณวาร}{mtk.218}
\bookmark[dest={mtk.218},level=1]{1. ปฐมภาณวาร}
\csromanheader{1}{1. ปฐมภาณวาร}
\proseitem{446}{\csromanfolio{491}เตน โข ปน สมเยน วสฺสสตปรินิพฺพุเต ภควติ เวสาลิกา วชฺชิปุตฺตกา ภิกฺขู เวสาลิยํ ทส วตฺถูนิ ทีเปนฺติ “กปฺปติ สิงฺคิโลณกปฺโป, กปฺปติ ทฺวงฺคุลกปฺโป, กปฺปติ คามนฺตรกปฺโป, กปฺปติ อาวาสกปฺโป, กปฺปติ อนุมติกปฺโป, กปฺปติ อาจิณฺณกปฺโป, กปฺปติ อมถิตกปฺโป, กปฺปติ ชโฬคิํ ปาตุํ, กปฺปติ อทสกํ นิสีทนํ, กปฺปติ ชาตรูปรชตนฺ”ติ.}

\prose{เตน โข ปน สมเยน อายสฺมา ยโส กากณฺฑกปุตฺโตวชฺชีสุ จาริกํ จรมาโน เยน เวสาลี ตทวสริ, ตตฺร สุทํ อายสฺมา ยโส กากณฺฑกปุตฺโตเวสาลิยํ วิหรติ มหาวเน กูฏาคารสาลายํ.\csromanspacer{} เตน โข ปน สมเยน เวสาลิกา วชฺชิปุตฺตกา ภิกฺขู ตทหุโปสเถ กํสปาติํ\footnote{กํสจฏิํ (สฺยา)} อุทเกน ปูเรตฺวา มชฺเฌ ภิกฺขุสํฆสฺส ฐเปตฺวา อาคตาคเต เวสาลิเก อุปาสเก เอวํ วทนฺติ “เทถาวุโส สํฆสฺส กหาปณมฺปิ อฑฺฒมฺปิ ปาทมฺปิ มาสกรูปมฺปิ, ภวิสฺสติ สํฆสฺส ปริกฺขาเรน กรณียนฺ”ติ.\csromanspacer{} เอวํ วุตฺเต อายสฺมา ยโส กากณฺฑกปุตฺโต เวสาลิเก อุปาสเก เอตทโวจ “มาวุโส อทตฺถ สํฆสฺส กหาปณมฺปิ อฑฺฒมฺปิ ปาทมฺปิ มาสกรูปมฺปิ, น กปฺปติ สมณานํ สกฺยปุตฺติยานํ ชาตรูปรชตํ, น สาทิยนฺติ สมณา สกฺยปุตฺติยา ชาตรูปรชตํ, น ปฏิคฺคณฺหนฺติ สมณา สกฺยปุตฺติยา ชาตรูปรชตํ, นิกฺขิตฺตมณิสุวณฺณา สมณา สกฺยปุตฺติยา, อเปตชาตรูปรชตา”ติ.\csromanspacer{} เอวมฺปิ โข เวสาลิกา อุปาสกา อายสฺมตา ยเสน กากณฺฑกปุตฺเตน วุจฺจมานา อทํสุเยว สํฆสฺส กหาปณมฺปิ อฑฺฒมฺปิ ปาทมฺปิ มาสกรูปมฺปิ.}

\prose{อถ โข เวสาลิกา วชฺชิปุตฺตกา ภิกฺขู ตสฺสา รตฺติยา อจฺจเยน ตํ หิรญฺญํ ภิกฺขคฺเคน\footnote{ภิกฺขุคฺเคน (สฺยา)} ปฏิวีสํ\footnote{ปฏิวิํสํ (สี), ปฏิวิสํ (สฺยา)} ฐเปตฺวา ภาเชสุํ.\csromanspacer{} อถ โข เวสาลิกา วชฺชิปุตฺตกา ภิกฺขู อายสฺมนฺตํ ยสํ กากณฺฑกปุตฺตํ เอตทโวจุํ “เอโส เต อาวุโส ยส หิรญฺญสฺส ปฏิวีโส”ติ.\csromanspacer{} นตฺถิ เม อาวุโส หิรญฺญสฺส ปฏิวีโส, นาหํ หิรญฺญํ สาทิยามีติ. \csromanfolio{492}อถ โข เวสาลิกา วชฺชิปุตฺตกา ภิกฺขู “อยํ อาวุโส ยโส กากณฺฑกปุตฺโต อุปาสเก สทฺเธ ปสนฺเน อกฺโกสติ ปริภาสติ อปฺปสาทํ กโรติ, หนฺทสฺส มยํ ปฏิสารณียกมฺมํ กโรมา”ติ, เต ตสฺส ปฏิสารณียกมฺมํ อกํสุ.\csromanspacer{} อถ โข อายสฺมา ยโส กากณฺฑกปุตฺโต เวสาลิเก วชฺชิปุตฺตเก ภิกฺขู เอตทโวจ “ภควโต อาวุโส ปญฺญตฺตํ ‘ปฏิสารณียกมฺมกตสฺส ภิกฺขุโน อนุทูโต ทาตพฺโพ’ติ, เทถ เม อาวุโส อนุทูตํ ภิกฺขุนฺ”ติ. อถ โข เวสาลิกา วชฺชิปุตฺตกา ภิกฺขู เอกํ ภิกฺขุํ สมฺมนฺนิตฺวา อายสฺมโต ยสสฺส กากณฺฑกปุตฺตสฺส อนุทูตํ อทํสุ.\csromanspacer{} อถ โข อายสฺมา ยโส กากณฺฑกปุตฺโต อนุทูเตน ภิกฺขุนา สทฺธิํ เวสาลิํ ปวิสิตฺวา เวสาลิเก อุปาสเก เอตทโวจ\csromandash{}อหํ กิรายสฺมนฺเต อุปาสเก สทฺเธ ปสนฺเน อกฺโกสามิ ปริภาสามิ อปฺปสาทํ กโรมิ, โยหํ อธมฺมํ “อธมฺโม”ติ วทามิ, ธมฺมํ “ธมฺโม”ติ วทามิ, อวินยํ “อวินโย”ติ วทามิ, วินยํ “วินโย”ติ วทามิ.}

\proseitem{447}{เอกมิทํ อาวุโส สมยํ ภควา สาวตฺถิยํ วิหรติ เชตวเน อนาถปิณฺฑิกสฺส อาราเม, ตตฺร โข อาวุโส ภควา ภิกฺขู อามนฺเตสิ\csromandash{} *จตฺตาโรเม ภิกฺขเว จนฺทิมสูริยานํ อุปกฺกิเลสา, เยหิ อุปกฺกิเลเสหิ อุปกฺกิลิฏฺฐา จนฺทิมสูริยา น ตปนฺติ, น ภาสนฺติ, น วิโรจนฺติ.\csromanspacer{} กตเม จตฺตาโร, อพฺภํ ภิกฺขเว จนฺทิมสูริยานํ อุปกฺกิเลโส, เยน อุปกฺกิเลเสน อุปกฺกิลิฏฺฐา จนฺทิมสูริยา น ตปนฺติ, น ภาสนฺติ, น วิโรจนฺติ.\csromanspacer{} มหิกา ภิกฺขเว จนฺทิมสูริยานํ อุปกฺกิเลโส, เยน อุปกฺกิเลเสน อุปกฺกิลิฏฺฐา จนฺทิมสูริยา น ตปนฺติ, น ภาสนฺติ, น วิโรจนฺติ.\csromanspacer{} ธูมรโช ภิกฺขเว จนฺทิมสูริยานํ อุปกฺกิเลโส, เยน อุปกฺกิเลเสน อุปกฺกิลิฏฺฐา จนฺทิมสูริยา น ตปนฺติ, น ภาสนฺติ, น วิโรจนฺติ.\csromanspacer{} ราหุ ภิกฺขเว อสุรินฺโท จนฺทิมสูริยานํ อุปกฺกิเลโส, เยน อุปกฺกิเลเสน อุปกฺกิลิฏฺฐา จนฺทิมสูริยา น ตปนฺติ, น ภาสนฺติ, น วิโรจนฺติ.\csromanspacer{} อิเม โข ภิกฺขเว จตฺตาโร จนฺทิมสูริยานํ อุปกฺกิเลสา, เยหิ อุปกฺกิเลเสหิ อุปกฺกิลิฏฺฐา จนฺทิมสูริยา น ตปนฺติ, น ภาสนฺติ, น วิโรจนฺติ.\csromanspacer{} เอวเมว โข ภิกฺขเว จตฺตาโรเม สมณพฺราหฺมณานํ อุปกฺกิเลสา, เยหิ อุปกฺกิเลเสหิ อุปกฺกิลิฏฺฐา เอเก สมณพฺราหฺมณา \csromanfolio{493}น ตปนฺติ, น ภาสนฺติ, น วิโรจนฺติ.\csromanspacer{} กตเม จตฺตาโร, สนฺติ ภิกฺขเว เอเก สมณพฺราหฺมณา สุรํ ปิวนฺติ, เมรยํ ปิวนฺติ, สุราเมรยปานา อปฺปฏิวิรตา, อยํ ภิกฺขเว ปฐโม สมณพฺราหฺมณานํ อุปกฺกิเลโส, เยน อุปกฺกิเลเสน อุปกฺกิลิฏฺฐา เอเก สมณพฺราหฺมณา น ตปนฺติ, น ภาสนฺติ, น วิโรจนฺติ.\csromanspacer{} ปุน จปรํ ภิกฺขเว เอเก สมณพฺรหฺมณา เมถุนํ ธมฺมํ ปฏิเสวนฺติ, เมถุนธมฺมา อปฺปฏิวิรตา, อยํ ภิกฺขเว ทุติโย สมณพฺราหฺมณานํ อุปกฺกิเลโส, เยน อุปกฺกิเลเสน อุปกฺกิลิฏฺฐา เอเก สมณพฺราหฺมณา น ตปนฺติ, น ภาสนฺติ, น วิโรจนฺติ.\csromanspacer{} ปุน จปรํ ภิกฺขเว เอเก สมณพฺราหฺมณา ชาตรูปรชตํ สาทิยนฺติ, ชาตรูปรชตปฺปฏิคฺคหณา อปฺปฏิวิรตา, อยํ ภิกฺขเว ตติโย สมณพฺราหฺมณานํ อุปกฺกิเลโส, เยน อุปกฺกิเลเสน อุปกฺกิลิฏฺฐา เอเก สมณพฺราหฺมณา น ตปนฺติ, น ภาสนฺติ, น วิโรจนฺติ.\csromanspacer{} ปุน จปรํ ภิกฺขเว เอเก สมณ พฺราหฺมณา มิจฺฉาชีเวน ชีวิตํ กปฺเปนฺติ, มิจฺฉาชีวา อปฺปฏิวิรตา, อยํ ภิกฺขเว จตุตฺโถ สมณพฺราหฺมณานํ อุปกฺกิเลโส, เยน อุปกฺกิเลเสน อุปกฺกิลิฏฺฐา เอเก สมณพฺราหฺมณา น ตปนฺติ, น ภาสนฺติ, น วิโรจนฺติ.\csromanspacer{} อิเม โข ภิกฺขเว จตฺตาโร สมณพฺราหฺมณานํ อุปกฺกิเลสา, เยหิ อุปกฺกิเลเสหิ อุปกฺกิลิฏฺฐา เอเก สมณพฺราหฺมณา น ตปนฺติ, น ภาสนฺติ, น วิโรจนฺติ.\csromanspacer{} อิทมโวจาวุโส ภควา, อิทํ วตฺวาน สุคโต อถาปรํ เอตทโวจ สตฺถา.}

\csromangathameasure{ราคโทสปริกฺลิฏฺฐา, เอเก สมณพฺราหฺมณา. \\ อวิชฺชานิวุฏา โปสา, ปิยรูปาภินนฺทิโน. \\ สุรํ ปิวนฺติ เมรยํ, ปฏิเสวนฺติ เมถุนํ. \\ รชตํ ชาตรูปญฺจ, สาทิยนฺติ อวิทฺทสู. \\ มิจฺฉาชีเวน ชีวนฺติ, เอเก สมณพฺราหฺมณา. \\ เอเต อุปกฺกิเลสา วุตฺตา, พุทฺเธนาทิจฺจพนฺธุนา. \\ เยหิ อุปกฺกิเลเสหิ อุปกฺกิลิฏฺฐา, เอเก สมณพฺราหฺมณา. \\ น ตปนฺติ น ภาสนฺติ, อสุทฺธา สรชา มคา. \\ อนฺธกาเรน โอนทฺธา, ตณฺหาทาสา สเนตฺติกา.}
\csromangathagroup{\csromangathastanza{ราคโทสปริกฺลิฏฺฐา, เอเก สมณพฺราหฺมณา. \\ อวิชฺชานิวุฏา\footnote{อวิชฺชานีวุตา (สฺยา)} โปสา, ปิยรูปาภินนฺทิโน.}\csromangathastanza{สุรํ ปิวนฺติ เมรยํ, ปฏิเสวนฺติ เมถุนํ. \\ รชตํ ชาตรูปญฺจ, สาทิยนฺติ อวิทฺทสู.}\csromangathastanza{มิจฺฉาชีเวน ชีวนฺติ, เอเก สมณพฺราหฺมณา. \\ เอเต อุปกฺกิเลสา วุตฺตา, พุทฺเธนาทิจฺจพนฺธุนา.}\csromangathastanza{เยหิ อุปกฺกิเลเสหิ อุปกฺกิลิฏฺฐา, เอเก สมณพฺราหฺมณา. \\ น ตปนฺติ น ภาสนฺติ, อสุทฺธา สรชา มคา. \\ อนฺธกาเรน โอนทฺธา, ตณฺหาทาสา สเนตฺติกา.}}

\prose{วฑฺเฒนฺติ กฏสิํ โฆรํ, อาทิยนฺติ ปุนพฺภวนฺติ}

\prose{\csromanfolio{494}เอวํวาที กิราหํ อายสฺมนฺเต อุปาสเก สทฺเธ ปสนฺเน อกฺโกสามิ ปริภาสามิ อปฺปสาทํ กโรมิ, โยหํ อธมฺมํ “อธมฺโม”ติ วทามิ, ธมฺมํ “ธมฺโม”ติ วทามิ, อวินยํ “อวินโย”ติ วทามิ, วินยํ “วินโย”ติ วทามิ.}

\proseitem{448}{\csromansymbolfootnote{*}{สํ 2. 509 ปิฏฺเฐปิ อิทํ วตฺถุ อาคตํ.}เอกมิทํ อาวุโส สมยํ ภควา ราชคเห วิหรติ เวฬุวเน กลนฺทกนิวาเป.\csromanspacer{} เตน โข ปน สมเยน ราชนฺเตปุเร ราชปริสายํ สนฺนิสินฺนานํ สนฺนิปติตานํ อยมนฺตรกถา อุทปาทิ “กปฺปติ สมณานํ สกฺยปุตฺติยานํ ชาตรูปรชตํ, สาทิยนฺติ สมณา สกฺยปุตฺติยา ชาตรูปรชตํ, ปฏิคฺคณฺหนฺติ สมณา สกฺยปุตฺติยา ชาตรูปรชตนฺ”ติ.\csromanspacer{} เตน โข ปนาวุโส สมเยน มณิจูฬโก คามณี ตสฺสํ ปริสายํ นิสินฺโน โหติ.\csromanspacer{} อถ โข อาวุโส มณิจูฬโก คามณี ตํ ปริสํ เอตทโวจ “มา อยฺยา เอวํ อวจุตฺถ, น กปฺปติ สมณานํ สกฺยปุตฺติยานํ ชาตรูปรชตํ, น สาทิยนฺติ สมณา สกฺยปุตฺติยา ชาตรูปรชตํ, น ปฏิคฺคณฺหนฺติ สมณา สกฺยปุตฺติยา ชาตรูปรชตํ นิกฺขิตฺตมณิสุวณฺณา สมณา สกฺยปุตฺติยา อเปตชาตรูปรชตา”ติ.\csromanspacer{} อสกฺขิ โข อาวุโส มณิจูฬโก คามณี ตํ ปริสํ สญฺญาเปตุํ.}

\prose{อถ โข อาวุโส มณิจูฬโก คามณี ตํ ปริสํ สญฺญาเปตฺวา เยน ภควา เตนุปสงฺกมิ, อุปสงฺกมิตฺวา ภควนฺตํ อภิวาเทตฺวา เอกมนฺตํ นิสีทิ, เอกมนฺตํ นิสินฺโน โข อาวุโส มณิจูฬโก คามณี ภควนฺตํ เอตทโวจ “อิธ ภนฺเต ราชนฺเตปุเร ราชปริสายํ สนฺนิสินฺนานํ สนฺนิปติตานํ อยมนฺตรกถา อุทปาทิ ‘กปฺปติ สมณานํ สกฺยปุตฺติยานํ ชาตรูปรชตํ, สาทิยนฺติ สมณา สกฺยปุตฺติยา ชาตรูปรชตํ, ปฏิคฺคณฺหนฺติ สมณา สกฺยปุตฺติยา ชาตรูปรชตนฺ’ติ, เอวํ วุตฺเต อหํ ภนฺเต ตํ ปริสํ เอตทโวจํ ‘มา อยฺยา เอวํ อวจุตฺถ, น กปฺปติ สมณานํ สกฺยปุตฺติยานํ ชาตรูปรชตํ, น สาทิยนฺติ สมณา สกฺยปุตฺติยา ชาตรูปรชตํ, น ปฏิคฺคณฺหนฺติ สมณา สกฺยปุตฺติยา ชาตรูปรชตํ, นิกฺขิตฺตมณิสุวณฺณา สมณา สกฺยปุตฺติยา อเปตชาตรูปรชาตา’ติ.\csromanspacer{} อสกฺขิํ โข อหํ ภนฺเต ตํ ปริสํ สญฺญาเปตุํ, กจฺจาหํ ภนฺเต เอวํ พฺยากรมาโน วุตฺตวาที เจว ภควโต โหมิ, น จ ภควนฺตํ อภูเตน อพฺภาจิกฺขามิ, ธมฺมสฺส จานุธมฺมํ พฺยากโรมิ, น จ โกจิ สหธมฺมิโก วาทานุวาโท คารยฺหํ ฐานํ อาคจฺฉตี”ติ.\csromanspacer{} ตคฺฆ ตฺวํ คามณิ เอวํ พฺยากรมาโน วุตฺตวาที เจว เม โหสิ, น จ มํ อภูเตน \csromanfolio{495}อพฺภาจิกฺขสิ, ธมฺมสฺส จานุธมฺมํ พฺยากโรสิ, น จโกจิ สหธมฺมิโก วาทานุวาโท คารยฺหํ ฐานํ อาคจฺฉติ, น หิ คามณิ กปฺปติ สมณานํ สกฺยปุตฺติยานํ ชาตรูปรชตํ, น สาทิยนฺติ สมณา สกฺยปุตฺติยา ชาตรูปรชตํ, น ปฏิคฺคณฺหนฺติ สมณา สกฺยปุตฺติยา ชาตรูปรชตํ, นิกฺขิตฺตมณิสุวณฺณา สมณา สกฺยปุตฺติยา อเปตชาตรูปรชตา.\csromanspacer{} ยสฺส โข คามณิ ชาตรูปรชตํ กปฺปติ, ปญฺจปิ ตสฺส กามคุณา กปฺปนฺติ.\csromanspacer{} ยสฺส ปญฺจ กามคุณา กปฺปนฺติ, เอกํเสเนตํ คามณิ ธาเรยฺยาสิ “อสฺสมณธมฺโม อสกฺยปุตฺติยธมฺโม”ติ.\csromanspacer{} อปิ จาหํ คามณิ เอวํ วทามิ “ติณํ ติณตฺถิเกน ปริเยสิตพฺพํ, ทารุ ทารุตฺถิเกน ปริเยสิตพฺพํ, สกฏํ สกฏตฺถิเกน ปริเยสิตพฺพํ, ปุริโส ปุริสตฺถิเกน ปริเยสิทฺตพฺโพ, น เตฺววาหํ คามณี เกนจิ ปริยาเยน ชาตรูปรชตํ ‘สาทิตพฺพํ ปริเยสิตพฺพนฺ’ติ วทามี”ติ.}

\prose{เอวํวาที กิราหํ อายสฺมนฺเต อุปาสเก สทฺเธ ปสนฺเน อกฺโกสามิ ปริภาสามิ อปฺปสาทํ กโรมิ, โยหํ อธมฺมํ “อธมฺโม”ติ วทามิ, ธมฺมํ “ธมฺโม”ติ วทามิ, อวินยํ “อวินโย”ติ วทามิ, วินยํ “วินโย”ติ วทามิ.}

\proseitem{449}{\csromansymbolfootnote{*}{วิ 1. 344 ปิฏฺเฐ.}เอกมิทํ อาวุโส สมยํ ภควา ราชคเห อายสฺมนฺตํ อุปนนฺทํ สกฺยปุตฺตํ อารพฺภ ชาตรูปรชตํ ปฏิกฺขิปิ, สิกฺขาปทญฺจ ปญฺญเปสิ, เอวํวาที กิราหํ อายสฺมนฺเต อุปาสเก สทฺเธ ปสนฺเน อกฺโกสามิ ปริภาสามิ อปฺปสาทํ กโรมิ, โยหํ อธมฺมํ “อธมฺโม”ติ วทามิ, ธมฺมํ “ธมฺโม”ติ วทามิ, อวินยํ “อวินโย”ติ วาทามิ, วินยํ “วินโย”ติ วทามีติ.}

\prose{เอวํ วุตฺเต เวสาลิกา อุปาสกา อายสฺมนฺตํ ยสํ กากณฺฑกปุตฺตํ เอตทโวจุํ “เอโกว ภนฺเต อยฺโย ยโส กากณฺฑกปุตฺโตสมโณ สกฺยปุตฺติโย, สพฺเพวิเม อสฺสมณา อสกฺยปุตฺติยา, วสตุ ภนฺเต อยฺโย ยโส กากณฺฑกปุตฺโตเวสาลิยํ, มยํ อยฺยสฺส ยสสฺส กากณฺฑกปุตฺตสฺส อุสฺสุกฺกํ กริสฺสาม จีวรปิณฺฑปาตเสนาสนคิลานปฺปจฺจยเภสชฺชปริกฺขารานนฺ”ติ.\csromanspacer{} อถ โข อายสฺมา ยโส กากณฺฑกปุตฺโต เวสาลิเก อุปาสเก สญฺญาเปตฺวา อนุทูเตน ภิกฺขุนา สทฺธิํ อารามํ อคมาสิ.}

\prose{อถ โข เวสาลิกา วชฺชิปุตฺตกา ภิกฺขู อนุทูตํ ภิกฺขุํ ปุจฺฉิํสุ “ขมาปิตาวุโส ยเสน กากณฺฑกปุตฺเตน เวสาลิกา \csromanfolio{496}อุปาสกา”ติ.\csromanspacer{} อุปาสเกหิ ปาปิกํ โน อาวุโส กตํ, เอโกว ยโส กากณฺฑกปุตฺโตสมโณ สกฺยปุตฺติโย กโต, สพฺเพว มยํ อสฺสมณา อสกฺยปุตฺติยา กตาติ.\csromanspacer{} อถ โข เวสาลิกา วชฺชิปุตฺตกา ภิกฺขู “อยํ อาวุโส ยโส กากณฺฑกปุตฺโต อเมฺหหิ อสมฺมโต คิหีนํ ปกาเสสิ, หนฺทสฺส มยํ อุกฺเขปนียกมฺมํ กโรมา”ติ, เต ตสฺส อุกฺเขปนียกมฺมํ กตฺตุกามา สนฺนิปติํสุ.\csromanspacer{} อถ โข อายสฺมา ยโส กากณฺฑกปุตฺโต เวหาสํ อพฺภุคฺคนฺตฺวา โกสมฺพิยํ ปจฺจุฏฺฐาสิ.}

\proseitem{450}{อถ โข อายสฺมา ยโส กากณฺฑกปุตฺโต ปาเวยฺยกานญฺจ\footnote{ปาเฐยฺยกานญฺจ (สฺยา)} อวนฺติทกฺขิณาปถกานญฺจ ภิกฺขูนํ สนฺติเก ทูตํ ปาเหสิ “อาคจฺฉนฺตุ อายสฺมนฺตา, อิมํ อธิกรณํ อาทิยิสฺสาม, ปุเร อธมฺโม ทิปฺปติ, ธมฺโม ปฏิพาหิยฺยติ, อวินโย ทิปฺปติ, วินโย ปฏิพาหิยฺยติ, ปุเร อธมฺมวาทิโน พลวนฺโต โหนฺติ, ธมฺมวาทิโน ทุพฺพลา โหนฺติ, อวินยวาทิโน พลวนฺโต โหนฺติ, วินยวาทิโน ทุพฺพลา โหนฺตี”ติ.}

\prose{เตน โข ปน สมเยน อายสฺมา สมฺภูโต สาณวาสี อโหคงฺเค ปพฺพเต ปฏิวสติ.\csromanspacer{} อถ โข อายสฺมา ยโส กากณฺฑกปุตฺโต เยน อโหคงฺโค ปพฺพโต เยนายสฺมา สมฺภูโต สาณวาสี เตนุปสงฺกมิ, อุปสงฺกมิตฺวา อายสฺมนฺตํ สมฺภูตํ สาณวาสิํ อภิวาเทตฺวา เอกมนฺตํ นิสีทิ, เอกมนฺตํ นิสินฺโน โข อายสฺมา ยโส กากณฺฑกปุตฺโต อายสฺมนฺตํ สมฺภูตํ สาณวาสิํ เอตทโวจ “อิเม ภนฺเต เวสาลิกา วชฺชิปุตฺตกา ภิกฺขู เวสาลิยํ ทส วตฺถูนิ ทีเปนฺติ ‘กปฺปติ สงฺคิโลณกปฺโป, กปฺปติ ทฺวงฺคุลกปฺโป, กปฺปติ คามนฺตรกปฺโป, กปฺปติ อาวาสกปฺโป, กปฺปติ อนุมติกปฺโป, กปฺปติ อาจิณฺณกปฺโป, กปฺปติ อมถิตกปฺโป, กปฺปติ ชโฬคิํ ปาตุํ, กปฺปติ อทสกํ นิสีทนํ, กปฺปติ ชาตรูปรชตนฺ’ติ, หนฺท มยํ ภนฺเต อิมํ อธิกรณํ อาทิยิสฺสาม, ปุเร อธมฺโม ทิปฺปติ, ธมฺโม ปฏิพาหิยฺยติ, อวินโย ทิปฺปติ, วินโย ปฏิคาหิยฺยติ, ปุเร อธมฺมวาทิโน พลวนฺโต โหนฺติ, ธมฺมวาทิโน ทุพฺพลา โหนฺติ, อวินยวาทิโน พลวนฺโต โหนฺติ, วินยวาทิโน ทุพฺพลา โหนฺตี”ติ.\csromanspacer{} เอวมาวุโสติ โข อายสฺมา สมฺภูโต สาณวาสี อายสฺมโต ยสสฺส กากณฺฑกปุตฺตสฺส ปจฺจสฺโสสิ. \csromanfolio{497}อถ โข สฏฺฐิมตฺตา ปาเวยฺยกา ภิกฺขู สพฺเพ อารญฺญิกา สพฺเพ ปิณฺณปาติกา สพฺเพ ปํสุกูลิกา สพฺเพ เตจีวริกา สพฺเพว อรหนฺโต อโหคงฺเค ปพฺพเต สนฺนิปติํสุ.\csromanspacer{} อฏฺฐาสีติมตฺตา อวนฺติทกฺขิณาปถกา ภิกฺขู อปฺเปกจฺเจ อารญฺญิกา อปฺเปกจฺเจ ปิณฺฑปาติกา อปฺเปกจฺเจ ปํสุกูลิกา อปฺเปกจฺเจ เตจีวริกา สพฺเพว อรหนฺโต อโหคงฺเค ปพฺพเต สนฺนิปติํสุ.\csromanspacer{} อถ โข เถรานํ ภิกฺขูนํ มนฺตยมานานํ เอตทโหสิ “อิทํ โข อธิกรณํ กกฺขฬญฺจ วาฬญฺจ, กํ นุ โข มยํ ปกฺขํ ลเภยฺยาม เยน มยํ อิมสฺมิํ อธิกรเณ พลวนฺตตรา อสฺสามา”ติ.}

\proseitem{451}{เตน โข ปน สมเยน อายสฺมา เรวโต โสเรยฺเย ปฏิวสติ พหุสฺสุโต อาคตาคโม ธมฺมธโร วินยธโร มาติกาธโร ปณฺฑิโต วิยตฺโต เมธาวี ลชฺชี กุกฺกุจฺจโก สิกฺขากาโม.\csromanspacer{} อถ โข เถรานํ ภิกฺขูนํ เอตทโหสิ “อยํ โข อายสฺมา เรวโต โสเรยฺเย ปฏิวสติ พหุสฺสุโต อาคตาคโม ธมฺมธโร วินยธโร มาติกาธโร ปณฺฑิโต วิยตฺโต เมธาวี ลชฺชี กุกฺกุจฺจโก สิกฺขากาโม, สเจ มยํ อายสฺมนฺตํ เรวตํ ปกฺขํ ลภิสฺสาม, เอวํ มยํ อิมสฺมิํ อธิกรเณ พลวนฺตตรา อสฺสามา”ติ.\csromanspacer{} อสฺโสสิ โข อายสฺมา เรวโต ทิพฺพาย โสตธาตุยา วิสุทฺธาย อติกฺกนฺตมานุสิกาย เถรานํ ภิกฺขูนํ มนฺตยมานานํ, สุตฺวานสฺส เอตทโหสิ “อิทํ โข อธิกรณํ กกฺขฬญฺจ วาฬญฺจ, น โข เมตํ ปติรูปํ, โยหํ เอวรูเป อธิกรเณ โอสกฺเกยฺยํ.\csromanspacer{} อิทานิ จ ปน เต ภิกฺขู อาคจฺฉิสฺสนฺติ, โสหํ เตหิ อากิณฺโณ น ผาสุ คมิสฺสามิ, ยํนูนาหํ ปฏิกจฺเจว คจฺเฉยฺยนฺ”ติ.\csromanspacer{} อถ โข อายสฺมา เรวโต โสเรยฺยา สงฺกสฺสํ อคมาสิ.}

\prose{อถ โข เถรา ภิกฺขู โสเรยฺยํ คนฺตฺวา ปุจฺฉิํสุ “กหํ อายสฺมา เรวโต”ติ.\csromanspacer{} เต เอวมาหํสุ “เอสายสฺมา เรวโต สงฺกสฺสํ คโต”ติ.\csromanspacer{} อถ โข อายสฺมา เรวโต สงฺกสฺสา กณฺณกุชฺชํ\footnote{กนฺนกุชฺชํ (สี)} อคมาสิ.\csromanspacer{} อถ โข เถรา ภิกฺขู สงฺกสฺสํ คนฺตฺวา ปุจฺฉิํสุ “กหํ อายสฺมา เรวโต”ติ.\csromanspacer{} เต เอวมาหํสุ “เอสายสฺมา เรวโต กณฺณกุชฺชํ คโต”ติ.\csromanspacer{} อถ โข อายสฺมา เรวโต กณฺณกุชฺชา \csromanfolio{498}อุทุมฺพรํ อคมาสิ.\csromanspacer{} อถ โข เถรา ภิกฺขู กณฺณกุชฺชํ คนฺตฺวา ปุจฺฉิํสุ “กหํ อายสฺมา เรวโต”ติ.\csromanspacer{} เต เอวมาหํสุ “เอสายสฺมา เรวโต อุทุมฺพรํ คโต”ติ.\csromanspacer{} อถ โข อายสฺมา เรวโต อุทุมฺพรา อคฺคฬปุรํ อคมาสิ.\csromanspacer{} อถ โข เถรา ภิกฺขู อุทุมฺพรํ คนฺตฺวา ปุจฺฉิํสุ “กหํ อายสฺมา เรวโต”ติ.\csromanspacer{} เต เอวมาหํสุ “เอสายสฺมา เรวโต อคฺคฬปุรํ คโต”ติ.\csromanspacer{} อถ โข อายสฺมา เรวโต อคฺคฬปุรา สหชาติํ\footnote{สหํชาติํ (ก)} อคมาสิ.\csromanspacer{} อถ โข เถรา ภิกฺขู อคฺคฬปุรํ คนฺตฺวา ปุจฺฉิํสุ “กหํ อายสฺมา เรวโต”ติ.\csromanspacer{} เต เอวมาหํสุ “เอสายสฺมา เรวโต สหชาติํ คโต”ติ.\csromanspacer{} อถ โข เถรา ภิกฺขู อายสฺมนฺตํ เรวตํ สหชาติยํ สมฺภาเวสุํ.}

\proseitem{452}{อถ โข อายสฺมา สมฺภูโต สาณวาสี อายสฺมนฺตํ ยสํ กากณฺฑกปุตฺตํ เอตทโวจ “อยํ อาวุโส อายสฺมา เรวโต พหุสฺสุโต อาคตาคโม ธมฺมธโร วินยธโร มาติกาธโร ปณฺฑิโต วิยตฺโต เมธาวี ลชฺชี กุกฺกุจฺจโก สิกฺขากาโม, สเจ มยํ อายสฺมนฺตํ เรวตํ ปญฺหํ ปุจฺฉิสฺสาม, ปฏิพโล อายสฺมา เรวโต เอเกเนว ปเญฺหน สกลมฺปิ รตฺติํ วีตินาเมตุํ, อิทานิ จ ปนายสฺมา เรวโต อนฺเตวาสิกํ\footnote{อนฺเตวาสิํ (สฺยา)} สรภาณกํ ภิกฺขุํ อชฺเฌสิสฺสติ, โส ตฺวํ ตสฺส ภิกฺขุโน สรภญฺญปริโยสาเน อายสฺมนฺตํ เรวตํ อุปสงฺกมิตฺวา อิมานิ ทส วตฺถูนิ ปุจฺเฉยฺยาสี”ติ. “เอวํ ภนฺเต”ติ โข อายสฺมา ยโส กากณฺฑกปุตฺโต อายสฺมโต สมฺภูตสฺส สาณวาสิสฺส ปจฺจสฺโสสิ.\csromanspacer{} อถ โข อายสฺมา เรวโต อนฺเตวาสิกํ สรภาณกํ ภิกฺขุํ อชฺเฌสิ.\csromanspacer{} อถ โข อายสฺมา ยโส กากณฺฑกปุตฺโตตสฺส ภิกฺขุโน สรภญฺญปริโยสาเน เยนายสฺมา เรวโต เตนุปสงฺกมิ, อุปสงฺกมิตฺวา อายสฺมนฺตํ เรวตํ อภิวาเทตฺวา เอกมนฺตํ นิสีทิ, เอกมนฺตํ นิสินฺโน โข อายสฺมา ยโส กากณฺฑกปุตฺโต อายสฺมนฺตํ เรวตํ เอตทโวจ “กปฺปติ ภนฺเต สิงฺคิโลณกปฺโป”ติ.\csromanspacer{} โก โส อาวุโส สิงฺคิโลณกปฺโปติ.\csromanspacer{} กปฺปติ ภนฺเต สิงฺคินา โลณํ ปริหริตุํ “ยตฺถ อโลณกํ ภวิสฺสติ, ตตฺถ ปริภุญฺชิสฺสามี”ติ.\csromanspacer{} นาวุโส กปฺปตีติ.\csromanspacer{} กปฺปติ ภนฺเต \csromanfolio{499}ทฺวงฺคุลกปฺโปติ.\csromanspacer{} โก โส อาวุโส ทฺวงฺคุลกปฺโปติ.\csromanspacer{} กปฺปติ ภนฺเต ทฺวงฺคุลาย ฉายาย วีติวตฺตาย วิกาเล โภชนํ ภุญฺชิตุนฺติ.\csromanspacer{} นาวุโส กปฺปตีติ.\csromanspacer{} กปฺปติ ภนฺเต คามนฺตรกปฺโปติ.\csromanspacer{} โก โส อาวุโส คามนฺตรกปฺโปติ.\csromanspacer{} กปฺปติ ภนฺเต “อิทานิ คามนฺตรํ คมิสฺสามี”ติ ภุตฺตาวินา ปวาริเตน อนติริตฺตํ โภชนํ ภุญฺชิตุนฺติ.\csromanspacer{} นาวุโส กปฺปตีติ.\csromanspacer{} กปฺปติ ภนฺเต อาวาสกปฺโปติ.\csromanspacer{} โก โส อาวุโส อาวาสกปฺโปติ.\csromanspacer{} กปฺปติ ภนฺเต สมฺพหุลา อาวาสา สมานสีมา นานุโปสถํ กาตุนฺติ.\csromanspacer{} นาวุโส กปฺปตีติ.\csromanspacer{} กปฺปติ ภนฺเต อนุมติกปฺโปติ.\csromanspacer{} โก โส อาวุโส อนุมติกปฺโปติ.\csromanspacer{} กปฺปติ ภนฺเต วคฺเคน สํเฆน กมฺมํ กาตุํ “อาคเต ภิกฺขู อนุมาเนสฺสามา”ติ\footnote{อนุชานิสฺสามาติ, อนุชาเนสฺสามาติ, อนุมติํ อาเนสฺสามาติ (ก)}.\csromanspacer{} นาวุโส กปฺปตีติ.\csromanspacer{} กปฺปติ ภนฺเต อาจิณฺณกปฺโปติ.\csromanspacer{} โก โส อาวุโส อาจิณฺณกปฺโปติ.\csromanspacer{} กปฺปติ ภนฺเต อิทํ เม อุปชฺฌาเยน อชฺฌาจิณฺณํ, อิทํ เม อาจริเยน อชฺฌาจิณฺณํ, ตํ อชฺฌาจริตุนฺติ.\csromanspacer{} อาจิณฺณกปฺโป โข อาวุโส เอกจฺโจ กปฺปติ, เอกจฺโจ น กปฺปตีติ.\csromanspacer{} กปฺปติ ภนฺเต อมถิตกปฺโปติ.\csromanspacer{} โก โส อาวุโส อมถิตกปฺโปติ.\csromanspacer{} กปฺปติ ภนฺเต ยํ ตํ ขีรํ ขีรภาวํ วิชหิตํ อสมฺปตฺตํ ทธิภาวํ, ตํ ภุตฺตาวินา ปวาริเตน อนติริตฺตํ ปาตุนฺติ.\csromanspacer{} นาวุโส กปฺปตีติ.\csromanspacer{} กปฺปติ ภนฺเต ชโฬคิํ ปาตุนฺติ.\csromanspacer{} กา สา อาวุโส ชโฬคีติ.\csromanspacer{} กปฺปติ ภนฺเต ยา สา สุรา อาสุตา อสมฺปตฺตา มชฺชภาวํ สา ปาตุนฺติ.\csromanspacer{} นาวุโส กปฺปตีติ.\csromanspacer{} กปฺปติ ภนฺเต อทสกํ นิสีทนนฺติ.\csromanspacer{} นาวุโส กปฺปตีติ.\csromanspacer{} กปฺปติ ภนฺเต ชาตรูปรชตนฺติ.\csromanspacer{} นาวุโส กปฺปตีติ.\csromanspacer{} อิเม ภนฺเต เวสาลิกา วชฺชิปุตฺตกา ภิกฺขู เวสาลิยํ อิมานิ ทส วตฺถูนิ ทีเปนฺติ, หนฺท มยํ ภนฺเต อิมํ อธิกรณํ อาทิยิสฺสาม, ปุเร อธมฺโม ทิปฺปติ, ธมฺโม ปฏิพาหิยฺยติ, อวินโย ทิปฺปติ, วินโย ปฏิพาหิยฺยติ, ปุเร อธมฺมวาทิโน พลวนฺโต โหนฺติ, ธมฺมวาทิโน ทุพฺพลา โหนฺติ, อวินยวาทิโน พลวนฺโต โหนฺติ, วินยวาทิโน ทุพฺพลา โหนฺตีติ. “เอวมาวุโส”ติ โข อายสฺมา เรวโต อายสฺมโต ยสสฺส กากณฺฑกปุตฺตสฺส ปจฺจสฺโสสิ.}

\nitthitam{ปฐมภาณวาโร นิฏฺฐิโต.}
\csromanmatikaanchor{mtk.219}
\csromantocmarkref{section}{2. ทุติยภาณวาร}{mtk.219}
\bookmark[dest={mtk.219},level=1]{2. ทุติยภาณวาร}
\csromanheader{1}{2. ทุติยภาณวาร}
\proseitem{453}{\csromanfolio{500}อสฺโสสุํ โข เวสาลิกา วชฺชิปุตฺตกา ภิกฺขู “ยโส กิร กากณฺฑกปุตฺโต อิทํ อธิกรณํ อาทิยิตุกาโม ปกฺขํ ปริเยสติ, ลภติ จ กิร ปกฺขนฺ”ติ.\csromanspacer{} อถ โข เวสาลิกานํ วชฺชิปุตฺตกานํ ภิกฺขูนํ เอตทโหสิ “อิทํ โข อธิกรณํ กกฺขฬญฺจ วาฬญฺจ, กํ นุ โข มยํ ปกฺขํ ลเภยฺยาม, เยน มยํ อิมสฺมิํ อธิกรเณ พลวนฺตตรา อสฺสามา”ติ.}

\prose{อถ โข เวสาลิกานํ วชฺชิปุตฺตกานํ ภิกฺขูนํ เอตทโหสิ “อยํ โข อายสฺมา เรวโต พหุสฺสุโต อาคตาคโม ธมฺมธโร วินยธโร มาติกาธโร ปณฺฑิโต วิยตฺโต เมธาวี ลชฺชี กุกฺกุจฺจโก สิกฺขากาโม, สเจ มยํ อายสฺมนฺตํ เรวตํ ปกฺขํ ลเภยฺยาม, เอวํ มยํ อิมสฺมิํ อธิกรเณ พลวนฺตตรา อสฺสามา”ติ.}

\prose{อถ โข เวสาลิกา วชฺชิปุตฺตก ภิกฺขู ปหูตํ สามณกํ ปริกฺขารํ ปฏิยาเทสุํ ปตฺตมฺปิ จีวรมฺปิ นิสีทนมฺปิ สูจิฆรมฺปิ กายพนฺธนมฺปิ ปริสฺสาวนมฺปิ ธมฺมกรณมฺปิ, อถ โข เวสาลิกา วชฺชิปุตฺตกา ภิกฺขู ตํ สามณกํ ปริกฺขารํ อาทาย นาวาย สหชาติํ อุชฺชวิํสุ, นาวาย ปจฺโจโรหิตฺวา อญฺญตรสฺมิํ รุกฺขมูเล ภตฺตวิสฺสคฺคํ กโรนฺติ.\csromanspacer{} อถ โข อายสฺมาโต สาฬฺหสฺส รโหคตสฺส ปฏิสลฺลีนสฺส เอวํ เจตโส ปริวิตกฺโก อุทปาทิ “เก นุ โข ธมฺมวาทิโน ปาจีนกา วา ภิกฺขู, ปาเวยฺยกา วา”ติ.\csromanspacer{} อถ โข อายสฺมโต สาฬฺหสฺส ธมฺมญฺจ วินยญฺจ เจตสา ปจฺจเวกฺขนฺตสฺส เอตทโหสิ อธมฺมวาทิโน ปาจีนกา ภิกฺขู, ธมฺมวาทิโน ปาเวยฺยกา\footnote{ปาเฐยฺยกา (สฺยา)} ภิกฺขู”ติ.}

\prose{อถ โข อญฺญตรา สุทฺธาวาสกายิกา เทวตา อายสฺมโต สาฬฺหสฺส เจตสา เจโตปริวิตกฺกมญฺญาย เสยฺยถาปิ นาม พลวาปุริโส สมิญฺชิตํ วา พาหํ ปสาเรยฺย ปสาริตํ วา พาหํ สมิญฺเชยฺย เอวเมว สุทฺธาวาเสสุ เทเวสุ อนฺตรหิตา อายสฺมโต สาฬฺหสฺส สมฺมุเข ปาตุรโหสิ.\csromanspacer{} อถ โข สา เทวตา อายสฺมนฺตํ สาฬฺหํ เอตทโวจ “สาธุ ภนฺเต สาฬฺห, อธมฺมวาที ปาจีนกา ภิกฺขู, ธมฺมวาที ปาเวยฺยกา ภิกฺขู, เตน หิ ภนฺเต สาฬฺห ยถาธมฺโม ตถา ติฏฺฐาหี”ติ.\csromanspacer{} ปุพฺเพปิ จาหํ เทวเต เอตรหิ จ ยถาธมฺโม ตถา \csromanfolio{501}ฐิโต, อปิ จาหํ น ตาว ทิฏฺฐิํ อาวิ กโรมิ “อปฺเปว นาม มํ อิมสฺมิํ อธิกรเณ สมฺมนฺเนยฺยา”ติ.}

\proseitem{454}{อถ โข เวสาลิกา วชฺชิปุตฺตกา ภิกฺขู ตํ สามณกํ ปริกฺขารํ อาทาย เยนายสฺมา เรวโต เตนุปสงฺกมิํสุ, อุปสงฺกมิตฺวา อายสฺมนฺตํ เรวตํ เอตทโวจุํ “ปฏิคฺคณฺหาตุ ภนฺเต เถโร สามณกํ ปริกฺขารํ ปตฺตมฺปิ จีวรมฺปิ นิสีทนมฺปิ สูจิฆรมฺปิ กายพนฺธนมฺปิ ปริสฺสาวนมฺปิ ธมฺมกรณมฺปี”ติ. “อลํ อาวุโส, ปริปุณฺณํ เม ปตฺตจีวรนฺ”ติ น อิจฺฉิ ปฏิคฺคเหตุํ.}

\prose{เตน โข ปน สมเยน อุตฺตโร นาม ภิกฺขุ วีสติวสฺโส อายสฺมโต เรวตสฺส อุปฏฺฐาโก โหติ.\csromanspacer{} อถ โข เวสาลิกา วชฺชิปุตฺตกา ภิกฺขู เยนายสฺมา อุตฺตโร เตนุปสงฺกมิํสุ, อุปสงฺกมิตฺวา อายสฺมนฺตํ อุตฺตรํ เอตทโวจุํ “ปฏิคฺคณฺหาตุ อายสฺมา อุตฺตโร สามณกํ ปริกฺขารํ ปตฺตมฺปิ จีวรมฺปิ นิสีทนมฺปิ สูจิฆรมฺปิ กายพนฺธนมฺปิ ปริสฺสาวนมฺปิ ธมฺมกรณมฺปี”ติ. “อลํ อาวุโส, ปริปุณฺณํ เม ปตฺตจีวรนฺ”ติ น อิจฺฉิ ปฏิคฺคเหตุํ.\csromanspacer{} มนุสฺสา โข อาวุโส อุตฺตร ภควโต สามณกํ ปริกฺขารํ อุปนาเมนฺติ, สเจ ภควา ปฏิคฺคณฺหาติ, เตเนว อตฺตมนา โหนฺติ.\csromanspacer{} โน เจ ภควา ปฏิคฺคณฺหาติ.\csromanspacer{} อายสฺมโต\footnote{อายสฺมโต จ (สฺยา)} อานนฺทสฺส อุปนาเมนฺติ “ปฏิคฺคณฺหาตุ ภนฺเต เถโร สามณกํ ปริกฺขารํ, ยถา ภควตา ปฏิคฺคหิโต, เอวเมว โส ภวิสฺสตี”ติ.\csromanspacer{} ปฏิคฺคณฺหาตุ อายสฺมา อุตฺตโร สามณกํ ปริกฺขารํ, ยถา เถเรน ปฏิคฺคหิโต, เอวเมว โส ภวิสฺสตีติ.\csromanspacer{} อถ โข อายสฺมา อุตฺตโร เวสาลิเกหิ วชฺชิปุตฺเตหิ ภิกฺขูหิ นิปฺปีฬิยมาโน เอกํ จีวรํ อคฺคเหสิ.\csromanspacer{} วเทยฺยาถ อาวุโส เยน อตฺโถติ.\csromanspacer{} เอตฺตกํ อายสฺมา อุตฺตโร เถรํ วเทตุ “เอตฺตกญฺจ ภนฺเต เถโร สํฆมชฺเฌ วเทตุ, ปุรตฺถิเมสุ ชนปเทสุ พุทฺธา ภควนฺโต อุปฺปชฺชนฺติ.\csromanspacer{} ธมฺมวาที ปาจีนกา ภิกฺขู, อธมฺมวาที ปาเวยฺยกา ภิกฺขู”ติ. “เอวมาวุโส”ติ โข อายสฺมา อุตฺตโร เวสาลิกานํ วชฺชิปุตฺตกานํ ภิกฺขูนํ ปฏิสฺสุตฺวา เยนายสฺมา เรวโต เตนุปสงฺกมิ, อุปสงฺกมิตฺวา อายสฺมนฺตํ เรวตํ เอตทโวจ “เอตฺตกํ ภนฺเต เถโร สํฆมชฺเฌ วเทตุ, ปุรตฺถิเมสุ ชนปเทสุ พุทฺธา ภควนฺโต \csromanfolio{502}อุปฺปชฺชนฺติ.\csromanspacer{} ธมฺมวาที ปาจีนกา ภิกฺขู, อธมฺมวาที ปาเวยฺยกา ภิกฺขู”ติ. “อธมฺเม มํ ตฺวํ ภิกฺขุ นิโยเชสี”ติ เถโร อายสฺมนฺตํ อุตฺตรํ ปณาเมสิ.}

\prose{อถ โข เวสาลิกา วชฺชิปุตฺตกา ภิกฺขู อายสฺมนฺตํ อุตฺตรํ เอตทโวจุํ “กิํ อาวุโส อุตฺตร เถโร อาหา”ติ.\csromanspacer{} ปาปิกํ โน อาวุโส กตํ, “อธมฺเม มํ ตฺวํ ภิกฺขุ นิโยเชสี”ติ เถโร มํ ปณาเมสีติ.\csromanspacer{} นนุ ตฺวํ อาวุโส\footnote{อาวุโส อุตฺตร (สฺยา, กํ)} วุฑฺโฒ วีสติวสฺโสสีติ.\csromanspacer{} อามาวุโส, อปิ จ มยํ ครุนิสฺสยํ คณฺหามาติ.}

\proseitem{455}{อถ โข สํโฆ ตํ อธิกรณํ วินิจฺฉินิตุกาโม สนฺนิปติ, อถ โข อายสฺมา เรวโต สํฆํ ญาเปสิ\csromandash{}}

\prose{“สุณาตุ เม อาวุโส สํโฆ, สเจ มยํ อิมํ อธิกรณํ อิธ วูปสเมสฺสาม, สิยาปิ มูลาทายกา\footnote{มูลทายกา (สี)} ภิกฺขู ปุนกมฺมาย อุกฺโกเฏยฺยุํ, ยทิ สํฆสฺส ปตฺตกลฺลํ, ยตฺเถวิมํ อธิกรณํ สมุปฺปนฺนํ, สํโฆ ตตฺเถวิมํ อธิกรณํ วูปสเมยฺยา”ติ.}

\prose{อถ โข เถรา ภิกฺขู เวสาลิํ อคมํสุ ตํ อธิกรณํ วินิจฺฉินิตุกามา.}

\prose{เตน โข ปน สมเยน สพฺพกามี นาม ปถพฺยา สํฆตฺเถโร วีสวสฺสสติโก อุปสมฺปทาย อายสฺมโต อานนฺทสฺส สทฺธิวิหาริโก เวสาลิยํ ปฏิวสติ.\csromanspacer{} อถ โข อายสฺมา เรวโต อายสฺมนฺตํ สมฺภูตํ สาณวาสิํ เอตทโวจ “อหํ อาวุโส ยสฺมิํ วิหาเร สพฺพกามี เถโร วิหรติ, ตํ วิหารํ อุปคจฺฉามิ, โส ตฺวํ กาลสฺเสว อายสฺมนฺตํ สพฺพกามิํ อุปสงฺกมิตฺวา อิมานิ ทส วตฺถูนิ ปุจฺเฉยฺยาสี”ติ.}

\prose{“เอวํ ภนฺเต”ติ โข อายสฺมา สมฺภูโต สาณวาสี อายสฺมโต เรวตสฺส ปจฺจสฺโสสิ.\csromanspacer{} อถ โข อายสฺมา เรวโต ยสฺมิํ วิหาเร สพฺพกามี เถโร วิหรติ ตํ วิหารํ อุปคจฺฉิ.\csromanspacer{} คพฺเภ อายสฺมโต สพฺพกามิสฺส เสนาสนํ ปญฺญตฺตํ โหติ, คพฺภปฺปมุเข อายสฺมโต เรวตสฺส.\csromanspacer{} อถ โข อายสฺมา เรวโต “อยํ เถโร มหลฺลโก, น นิปชฺชตี”ติ น เสยฺยํ กปฺเปสิ, อายสฺมา สพฺพกามี “อยํ ภิกฺขุ อาคนฺตุโก กิลนฺโต น นิปชฺชตี”ติ น เสยฺยํ กปฺเปสิ.\csromanspacer{} อถ โข \csromanfolio{503}อายสฺมา สพฺพกามี รตฺติยา ปจฺจูสสมยํ ปจฺจุฏฺฐาย อายสฺมนฺตํ เรวตํ เอตทโวจ “กตเมน ตฺวํ ภูมิ วิหาเรน เอตรหิ พหุลํ วิหรสี”ติ.\csromanspacer{} เมตฺตาวิหาเรน โข อหํ ภนฺเต เอตรหิ พหุลํ วิหรามีติ.\csromanspacer{} กุลฺลกวิหาเรน กิร ตฺวํ ภูมิ เอตรหิ พหุลํ วิหรสิ, กุลฺลกวิหาโร เอโส\footnote{เหโส (สฺยา)} ภูมิ, ยทิทํ เมตฺตาติ.\csromanspacer{} ปุพฺเพปิ เม ภนฺเต คิหิภูตสฺส อาจิณฺณา เมตฺตา, เตนาหํ เอตรหิปิ เมตฺตาวิหาเรน พหุลํ วิหรามิ, อปิ จ มยา จิรปฺปตฺตํ อรหตฺตนฺติ.\csromanspacer{} เถโร ปน ภนฺเต กตเมน วิหาเรน เอตรหิ พหุลํ วิหรตีติ.\csromanspacer{} สุญฺญตาวิหาเรน โข อหํ ภูมิ เอตรหิ พหุลํ วิหรามีติ.\csromanspacer{} มหาปุริสวิหาเรน กิร ภนฺเต เถโร เอตรหิ พหุลํ วิหรติ, มหาปุริสวิหาโร เอโส ภนฺเต, ยทิทํ สุญฺญตาติ.\csromanspacer{} ปุพฺเพปิ เม ภูมิ คิหิภูตสฺส อาจิณฺณา สุญฺญตา, เตนาหํ เอตรหิปิ สุญฺญาตาวิหาเรน พหุลํ วิหรามิ, อปิ จ มยา จิรปฺปตฺตํ อรหตฺตนฺติ.\csromanspacer{} อยญฺจรหิ เถเรนํ ภิกฺขูนํ อนฺตรากถา วิปฺปกตา, อถายสฺมา สมฺภูโต สาณวาสี ตสฺมิํ อนุปฺปตฺโต โหติ.\csromanspacer{} อถ โข อายสฺมา สมฺภูโต สาณวาสี เยนายสฺมา สพฺพกามี เตนุปสงฺกมิ, อุปสงฺกมิตฺวา อายสฺมนฺตํ สพฺพกามิํ อภิวาเทตฺวา เอกมนฺตํ นิสีทิ, เอกมนฺตํ นิสินฺโน โข อายสฺมา สมฺภูโต สาณวาสี อายสฺมนฺตํ สพฺพกามิํ เอตทโวจ “อิเม ภนฺเต เวสาลิกา วชฺชิปุตฺตกา ภิกฺขู เวสาลิยํ ทส วตฺถูนิ ทีเปนฺติ ‘กปฺปติ สิงฺคิโลณกปฺโป, กปฺปติ ทฺวงฺคุลกปฺโป, กปฺปติ คามนฺตรกปฺโป, กปฺปติ อาวาสกปฺโป, กปฺปติ อนุมติกปฺโป, กปฺปติ อาจิณฺณกปฺโป, กปฺปติ อมถิตกปฺโป, กปฺปติ ชโฬคิํ ปาตุํ, กปฺปติ อทสกํ นิสีทนํ, กปฺปติ ชาตรูปรชตนฺ”ติ.\csromanspacer{} เถเรน ภนฺเต อุปชฺฌายสฺส มูเล พหุ ธมฺโม จ วินโย จ ปริยตฺโต, เถเรสฺส ภนฺเต ธมฺมญฺจ วินยญฺจ ปจฺจเวกฺขนฺตสฺส กถํ โหติ “เก นุ โข ธมฺมวาทิโน, ปาจีนกา วา ภิกฺขู, ปาเวยฺยกา วา”ติ.\csromanspacer{} ตยาปิ โข อาวุโส อุปชฺฌายสฺส มูเล พหุ ธมฺโม จ วินโย จ ปริยตฺโต, ตุยฺหํ ปน อาวุโส ธมฺมญฺจ วินยญฺจ ปจฺจเวกฺขนฺตสฺส กถํ โหติ “เก นุ โข ธมฺมวาทิโน ปาจีนกา วา ภิกฺขู, ปาเวยฺยกา วา”ติ.\csromanspacer{} มยฺหํ โข ภนฺเต ธมฺมญฺจ วินยญฺจ ปจฺจเวกฺขนฺตสฺส เอวํ โหติ “อธมฺมวาที ปาจีนกา ภิกฺขู, ธมฺมวาที ปาเวยฺยกา ภิกฺขู”ติ.\csromanspacer{} อปิ จาหํ \csromanfolio{504}น ตาว ทิฏฺฐิํ อาวิ กโรมิ “อปฺเปว นาม มํ อิมสฺมิํ อธิกรเณ สมฺมนฺเนยฺยา”ติ.\csromanspacer{} มยฺหมฺปิ โข อาวุโส ธมฺมญฺจ วินยญฺจ ปจฺจเวกฺขนฺตสฺส เอวํ โหติ “อธมฺมวาที ปาจีนกา ภิกฺขู, ธมฺมวาที ปาเวยฺยกา ภิกฺขู”ติ.\csromanspacer{} อปิ จาหํ น ตาว ทิฏฺฐิํ อาวิ กโรมิ “อปฺเปว นาม มํ อิมสฺมิํ อธิกรเณ สมฺมนฺเนยฺยา”ติ.}

\proseitem{456}{อถ โข สํโฆ ตํ อธิกรณํ วินิจฺฉินิตุกาโม สนฺนิปติ, ตสฺมิํ โข ปน อธิกรเณ วินิจฺฉิยมาเน อนคฺคานิ เจว ภสฺสานิ ชายนฺติ, น เจกสฺส ภาสิตสฺส อตฺโถ วิญฺญายติ, อถ โข อายสฺมา เรวโต สํฆํ ญาเปสิ\csromandash{}}

\prose{“สุณาตุ เม ภนฺเต สํโฆ, อมฺหากํ อิมสฺมิํ อธิกรเณ วินิจฺฉิยมาเน อนคฺคานิ เจว ภสฺสานิ ชายนฺติ, น เจกสฺส ภาสิตสฺส อตฺโถ วิญฺญายติ, ยทิ สํฆสฺส ปตฺตกลฺลํ, สํโฆ อิมํ อธิกรณํ อุพฺพาหิกาย วูปสเมยฺยา”ติ.\csromanspacer{} สํโฆ จตฺตาโร ปาจีนเก ภิกฺขู จตฺตาโร ปาเวยฺยเก ภิกฺขู อุจฺจินิ, “ปาจีนกานํ ภิกฺขูนํ อายสฺมนฺตญฺจ สพฺพกามิํ อายสฺมนฺตญฺจ สาฬฺหํ อายสฺมนฺตญฺจ ขุชฺชโสภิตํ อายสฺมนฺตญฺจ วาสภคามิกํ, ปาเวยฺยกานํ ภิกฺขูนํ อายสฺมนฺตญฺจ เรวตํ อายสฺมนฺตญฺจ สมฺภูตํ สาณวาสิํ อายสฺมนฺตญฺจ ยสํ กากณฺฑกปุตฺตํ อายสฺมนฺตญฺจ สุมนนฺ”ติ.\csromanspacer{} อถ โข อายสฺมา เรวโต สํฆํ ญาเปสิ\csromandash{}}

\prose{“สุณาตุ เม ภนฺเต สํโฆ, อมฺหากํ อิมสฺมิํ อธิกรเณ วินิจฺฉิยมาเน อนคฺคานิ เจว ภสฺสานิ ชายนฺติ.\csromanspacer{} น เจกสฺส ภาสิตสฺส อตฺโถ วิญฺญายติ, ยทิ สํฆสฺส ปตฺตกลฺลํ, สํโฆ จตฺตาโร ปาจีนเก ภิกฺขู จตฺตาโร ปาเวยฺยเก ภิกฺขู สมฺมนฺเนยฺย อุพฺพาหิกาย อิมํ อธิกรณํ วูปสเมตุํ, เอสา ญตฺติ.}

\prose{สุณาตุ เม ภนฺเต สํโฆ, อมฺหากํ อิมสฺมิํ อธิกรเณ วินิจฺฉิยมาเน อนคฺคานิ เจว ภสฺสานิ ชายนฺติ, น เจกสฺส ภาสิตสฺส อตฺโถ วิญฺญายติ, สํโฆ จตฺตาโร ปาจีนเก ภิกฺขู จตฺตาโร ปาเวยฺยเก ภิกฺขู สมฺมนฺนติ อุพฺพาหิกาย อิมํ อธิกรณํ วูปสเมตุํ, ยสฺสายสฺมโต ขมติ จตุนฺนํ ปาจีนกานํ ภิกฺขูนํ จตุนฺนํ ปาเวยฺยกานํ ภิกฺขูนํ สมฺมุติ อุพฺพาหิกาย อิมํ อธิกรณํ วูปสเมตุํ, โส ตุณฺหสฺส, ยสฺส นกฺขมติ, โส ภาเสยฺย.}

\prose{\csromanfolio{505}สมฺมตา สํเฆน จตฺตาโร ปาจีนกา ภิกฺขู จตฺตาโร ปาเวยฺยกา ภิกฺขู อุพฺพาหิกาย อิมํ อธิกรณํ วูปสเมตุํ, ขมติ สํฆสฺส, ตสฺมา ตุณฺหี, เอวเมตํ ธารยามี”ติ.}

\prose{เตน โข ปน สมเยน อชิโต นาม ภิกฺขุ ทสวสฺโส สํฆสฺส ปาติโมกฺขุทฺเทสโก โหติ.\csromanspacer{} อถ โข สํโฆ อายสฺมนฺตมฺปิ อชิตํ สมฺมนฺนติ เถรานํ ภิกฺขูนํ อาสนปญฺญาปกํ.\csromanspacer{} อถ โข เถรานํ ภิกฺขูนํ เอตทโหสิ “กตฺถ นุ โข มยํ อิมํ อธิกรณํ วูปสเมยฺยามา”ติ.\csromanspacer{} อถ โข เถรานํ ภิกฺขูนํ เอตทโหสิ “อยํ โข วาลิการาโม รมณีโย อปฺปสทฺโท อปฺปนิคฺโฆโส, ยํนูน มยํ วาลิการาเม อิมํ อธิกรณํ วูปสเมยฺยามา”ติ.}

\proseitem{457}{อถ โข เถรา ภิกฺขู วาลิการามํ อคมํสุ ตํ อธิกรณํ วินิจฺฉินิตุกามา.\csromanspacer{} อถ โข อายสฺมา เรวโต สํฆํ ญาเปสิ\csromandash{}}

\prose{“สุณาตุ เม ภนฺเต สํโฆ, ยทิ สํฆสฺส ปตฺตกลฺลํ, อหํ อายสฺมนฺตํ สพฺพกามิํ วินยํ ปุจฺเฉยฺยนฺ”ติ.}

\csromanheader{2}{อายสฺมา สพฺพกามี สํฆํ ญาเปสิ\csromandash{}}
\prose{“สุณาตุ เม อาวุโส สํโฆ, ยทิ สํฆสฺส ปตฺตกลฺลํ, อหํ เรวเตน วินยํ ปุฏฺโฐ วิสฺสชฺเชยฺยนฺ”ติ.}

\prose{อถ โข อายสฺมา เรวโต อายสฺมนฺตํ สพฺพกามิํ เอตทโวจ “กปฺปติ ภนฺเต สิงฺคิโลณกปฺโป”ติ.\csromanspacer{} โก โส อาวุโส สิงฺคิโลณกปฺโปติ.\csromanspacer{} กปฺปติ ภนฺเต สิงฺคินา โลณํ ปริหริตุํ “ยตฺถ อโลณกํ ภวิสฺสติ, ตตฺถ ปริภุญฺชิสฺสามา”ติ.\csromanspacer{} นาวุโส กปฺปตีติ. * กตฺถ ปฏิกฺขิตฺตนฺติ.\csromanspacer{} สาวตฺถิยํ สุตฺตวิภงฺเคติ.\csromanspacer{} กิํ อาปชฺชตีติ.\csromanspacer{} สนฺนิธิการกโภชเน ปาจิตฺติยนฺติ.}

\prose{สุณาตุ เม ภนฺเต สํโฆ, อิทํ ปฐมํ วตฺถุ สํเฆน วินิจฺฉิตํ, อิติปิทํ วตฺถุ อุทฺธมฺมํ อุพฺพินยํ อปคตสตฺถุสาสนํ, อิทํ ปฐมํ สลากํ นิกฺขิปามิ. (1)}

\prose{กปฺปติ ภนฺเต ทฺวงฺคุลกปฺโปติ.\csromanspacer{} โก โส อาวุโส ทฺวงฺคุลกปฺโปติ.\csromanspacer{} กปฺปติ ภนฺเต ทฺวงฺคุลาย ฉายาย วีติวตฺตาย วิกาเล โภชนํ}

\prose{\csromanfolio{506}ภุญฺชิตุนฺติ.\csromanspacer{} นาวุโส กปฺปตีติ.\footnote{วิ 2. 114 ปิฏฺเฐ.} กตฺถ ปฏิกฺขิตฺตนฺติ.\csromanspacer{} ราชคเห สุตฺตวิภงฺเคติ.\csromanspacer{} กิํ อาปชฺชตีติ.\csromanspacer{} วิกาลโภชเน ปาจิตฺติยนฺติ.}

\prose{สุณาตุ เม ภนฺเต สํโฆ, อิทํ ทุติยํ วตฺถุ สํเฆน วินิจฺฉิตํ, อิติปิทํ วตฺถุ อุทฺธมฺมํ อุพฺพินยํ อปคตสตฺถุสาสนํ, อิทํ ทุติยํ สลากํ นิกฺขิปามิ. (2)}

\prose{กปฺปติ ภนฺเต คามนฺตรกปฺโปติ.\csromanspacer{} โก โส อาวุโส คามนฺตรกปฺโปติ.\csromanspacer{} กปฺปติ ภนฺเต “อิทานิ คามนฺตรํ คมิสฺสามี”ติ ภุตฺตาวินา ปวาริเตน อนติริตฺตํ โภชนํ ภุญฺชิตุนฺติ.\csromanspacer{} นาวุโส กปฺปตีติ.\footnote{วิ 2. 109 ปิฏฺเฐ.} กตฺถ ปฏิกฺขิตฺตนฺติ.\csromanspacer{} สาวตฺถิยํ สุตฺตวิภงฺเคติ.\csromanspacer{} กิํ อาปชฺชตีติ.\csromanspacer{} อนติริตฺตโภชเน ปาจิตฺติยนฺติ.}

\prose{สุณาตุ เม ภนฺเต สํโฆ, อิทํ ตติยํ วตฺถุ สํเฆน วินิจฺฉิตํ, อิติปิทํ วตฺถุ อุทฺธมฺมํ อุพฺพินยํ อปคตสตฺถุสาสนํ, อิทํ ตติยํ สลากํ นิกฺขิปามิ. (3)}

\prose{กปฺปติ ภนฺเต อาวาสกปฺโปติ.\csromanspacer{} โก โส อาวุโส อาวาสกปฺโปติ.\csromanspacer{} กปฺปติ ภนฺเต สมฺพหุลา อาวาสา สมานสีมานานุโปสถํ กาตุนฺติ.\csromanspacer{} นาวุโส กปฺปตีติ.\footnote{วิ 3. 146 ปิฏฺเฐ.} กตฺถ ปฏิกฺขิตฺตนฺติ.\csromanspacer{} ราชคเห อุโปสถสํยุตฺเตติ.\csromanspacer{} กิํ อาปชฺชตีติ.\csromanspacer{} วินยาติสาเร ทุกฺกฏนฺติ.}

\prose{สุณาตุ เม ภนฺเต สํโฆ, อิทํ จตุตฺถํ วตฺถุ สํเฆน วินิจฺฉิตํ, อิติปิทํ วตฺถุ อุทฺธมฺมํ อุพฺพินยํ อปคตสตฺถุสาสนํ, อิทํ จตุตฺถํ สลากํ นิกฺขิปามิ. (4)}

\prose{กปฺปติ ภนฺเต อนุมติกปฺโปติ.\csromanspacer{} โก โส อาวุโส อนุมติกปฺโปติ.\csromanspacer{} กปฺปติ ภนฺเต วคฺเคน สํเฆน กมฺมํ กาตุํ “อาคเต ภิกฺขู อนุมาเนสฺสามา”ติ.\csromanspacer{} นาวุโส กปฺปตีติ.\footnote{วิ 3. 435 ปิฏฺเฐ.} กตฺถ ปฏิกฺขิตฺตนฺติ.\csromanspacer{} จมฺเปยฺยเก วินยวตฺถุสฺมินฺติ.\csromanspacer{} กิํ อาปชฺชตีติ.\csromanspacer{} วินยาติสาเร ทุกฺกฏนฺติ.}

\prose{สุณาตุ เม ภนฺเต สํโฆ, อิทํ ปญฺจมํ วตฺถุ สํเฆน วินิจฺฉิตํ, อิติปิทํ วตฺถุ อุทฺธมฺมํ อุพฺพินยํ อปคตสตฺถุสาสนํ, อิทํ ปญฺจมํ สลากํ นิกฺขิปามิ. (5)}

\prose{กปฺปติ ภนฺเต อาจิณฺณกปฺโปติ.\csromanspacer{} โก โส อาวุโส อาจิณฺณกปฺโปติ.\csromanspacer{} กปฺปติ ภนฺเต อิทํ เม อุปชฺฌาเยน อชฺฌาจิณฺณํ, อิทํ เม}

\prose{\csromanfolio{507}อาจริเยน อชฺฌาจิณฺณํ, ตํ อชฺฌาจริตุนฺติ.\csromanspacer{} อาจิณฺณกปฺโป โข อาวุโส เอกจฺโจ กปฺปติ, เอกจฺโจ น กปฺปตีติ.}

\prose{สุณาตุ เม ภนฺเต สํโฆ, อิทํ ฉฏฺฐํ วตฺถุ สํเฆน วินิจฺฉิตํ, อิติปิทํ วตฺถุ อุทฺธมฺมํ อุพฺพินยํ อปคตสตฺถุสาสนํ, อิทํ ฉฏฺฐํ สลากํ นิกฺขิปามิ. (6)}

\prose{กปฺปติ ภนฺเต อมถิตกปฺโปติ.\csromanspacer{} โก โส อาวุโส อมถิตกปฺโปติ.\csromanspacer{} กปฺปติ ภนฺเต ยํ ตํ ขีรํ ขีรภาวํ วิชหิตํ อสมฺปตฺตํ ทธิภาวํ, ตํ ภุตฺตาวินา ปวาริเตน อนติริตฺตํ ปาตุนฺติ.\csromanspacer{} นาวุโส กปฺปตีติ.\footnote{วิ 2. 109 ปิฏฺเฐ.} กตฺถ ปฏิกฺขิตฺตนฺติ.\csromanspacer{} สาวตฺถิยํ สุตฺตวิภงฺเคติ.\csromanspacer{} กิํ อาปชฺชตีติ.\csromanspacer{} อนติริตฺตโภชเน ปาจิตฺติยนฺติ.}

\prose{สุณาตุ เม ภนฺเต สํโฆ, อิทํ สตฺตมํ วตฺถุ สํเฆน วินิจฺฉิตํ, อิติปิทํ วตฺถุ อุทฺธมฺมํ อุพฺพินยํ อปคตสตฺถุสาสนํ, อิทํ สตฺตมํ สลากํ นิกฺขิปามิ. (7)}

\prose{กปฺปติ ภนฺเต ชโฬคิํ ปาตุนฺติ.\csromanspacer{} กา สา อาวุโส ชโฬคีติ.\csromanspacer{} กปฺปติ ภนฺเต ยา สา สุรา อาสุตา อสมฺปตฺตา มชฺชภาวํ สา ปาตุนฺติ.\csromanspacer{} นาวุโส กปฺปตีติ.\footnote{วิ 2. 145 ปิฏฺเฐ. 3. วิ 2. 221 ปิฏฺเฐ.} กตฺถ ปฏิกฺขิตฺตนฺติ.\csromanspacer{} โกสมฺพิยํ สุตฺตวิภงฺเคติ.\csromanspacer{} กิํ อาปชฺชตีติ.\csromanspacer{} สุราเมรยปาเน ปาจิตฺติยนฺติ.}

\prose{สุณาตุ เม ภนฺเต สํโฆ, อิทํ อฏฺฐมํ วตฺถุ สํเฆน วินิจฺฉิตํ, อิติปิทํ วตฺถุ อุทฺธมฺมํ อุพฺพินยํ อปคตสตฺถุสาสนํ, อิทํ อฏฺฐมํ สลากํ นิกฺขิปามิ. (8)}

\prose{กปฺปติ ภนฺเต อทสกํ นิสีทนนฺติ.\csromanspacer{} นาวุโส กปฺปตีติ.3 กตฺถ ปฏิกฺขิตฺตนฺติ.\csromanspacer{} สาวตฺถิยํ สุตฺตวิภงฺเคติ.\csromanspacer{} กิํ อาปชฺชตีติ.\csromanspacer{} เฉทนเก ปาจิตฺติยนฺติ.}

\prose{สุณาตุ เม ภนฺเต สํโฆ, อิทํ นวมํ วตฺถุ สํเฆน วินิจฺฉิตํ, อิติปิทํ วตฺถุ อุทฺธมฺพํ อุพฺพินยํ อปคตสตฺถุสาสนํ, อิทํ นวมํ สลากํ นิกฺขิปามิ. (9)}

\prose{กปฺปติ ภนฺเต ชาตรูปรชตนฺติ.\csromanspacer{} นาวุโส กปฺปตีติ.\footnote{วิ 1. 344 ปิฏฺเฐ.} กตฺถ ปฏิกฺขิตฺตนฺติ.\csromanspacer{} ราชคเห สุตฺตวิภงฺเคติ.\csromanspacer{} กิํ อาปชฺชตีติ.\csromanspacer{} ชาตรูปรชตปฏิคฺคหเณ ปาจิตฺติยนฺติ.}

\csromanmatikaanchor{mtk.220}
\csromantocmarkref{section}{อุทฺทานคาถา}{mtk.220}
\bookmark[dest={mtk.220},level=1]{อุทฺทานคาถา}
\prose{\csromanfolio{508}สุณาตุ เม ภนฺเต สํโฆ, อิทํ ทสมํ วตฺถุ สํเฆน วินิจฺฉิตํ, อิติปิทํ วตฺถุ อุทฺธมฺมํ อุพฺพินยํ อปคตสตฺถุสาสนํ, อิทํ ทสมํ สลากํ นิกฺขิปามิ. (10)}

\prose{สุณาตุ เม ภนฺเต สํโฆ, อิมานิ ทส วตฺถูนิ สํเฆน วินิจฺฉิตานิ “อิติปิมานิ ทสวตฺถูนิ อุทฺธมฺมานิ อุพฺพินยานิ อปคตสตฺถุสาสนานี”ติ.}

\proseitem{458}{นิหตเมตํ อาวุโส อธิกรณํ สนฺตํ วูปสนฺตํ สุวูปสนฺตํ, อปิ จ มํ ตฺวํ อาวุโส สํฆมชฺเฌปิ อิมานิ ทส วตฺถูนิ ปุจฺเฉยฺยาสิ เตสํ ภิกฺขูนํ สญฺญตฺติยาติ.\csromanspacer{} อถ โข อายสฺมา เรวโต อายสฺมนฺตํ สพฺพกามิํ สํฆมชฺเฌปิ อิมานิ ทส วตฺถูนิ ปุจฺฉิ, ปุฏฺโฐ ปุฏฺโฐ อายสฺมา สพฺพกามี วิสฺสชฺเชสิ.\csromanspacer{} อิมายา โข ปน วินยสงฺคีติยา สตฺต ภิกฺขุสตานิ อนูนานิ อนธิกานิ อเหสุํ, ตสฺมายํ วินยสงฺคีติ “สตฺตสติกา”ติ วุจฺจตีติ.\csromanspacer{} สตฺตสติกกฺขนฺธโก ทฺวาทสโม.\csromanspacer{} อิมมฺหิ ขนฺธเก วตฺถู ปญฺจวีสติ.}

\titlehead{\textbf{ตสฺสุทฺทานํ}}
\csromangathameasure{ทส วตฺถูนิ ปูเรตฺวา, กมฺมํ ทูเตน ปาวิสิ. \\ จตฺตาโร ปุน รูปญฺจ, โกสมฺพิ จ ปาเวยฺยโก. \\ มคฺโค โสเรยฺยํ สงฺกสฺสํ, กณฺณกุชฺชํ อุทุมฺพรํ. \\ สหชาติ จ มชฺเฌสิ, อสฺโสสิ กํ นุ โข มยํ. \\ ปตฺตนาวาย อุชฺชวิ, รโหสิ อุปนามยํ. \\ ครุ สํโฆ จ เวสาลิํ, เมตฺตา สํโฆ อุพฺพาหิกาติ. สตฺตสติกกฺขนฺธโก นิฏฺฐิโต. จูฬวคฺโค นิฏฺฐิโต.}
\csromangathagroup{\csromangathastanza{ทส วตฺถูนิ ปูเรตฺวา, กมฺมํ ทูเตน ปาวิสิ. \\ จตฺตาโร ปุน รูปญฺจ, โกสมฺพิ จ ปาเวยฺยโก.}\csromangathastanza{มคฺโค โสเรยฺยํ สงฺกสฺสํ, กณฺณกุชฺชํ อุทุมฺพรํ. \\ สหชาติ จ มชฺเฌสิ, อสฺโสสิ กํ นุ โข มยํ.}\csromangathastanza{ปตฺตนาวาย อุชฺชวิ, รโหสิ อุปนามยํ\footnote{ทูรโตปิ อุทปาทิ (ก)}. \\ ครุ\footnote{ทูรโตปิ อุทปาทิ (ก)} สํโฆ จ เวสาลิํ, เมตฺตา สํโฆ อุพฺพาหิกาติ.\csromanspacer{} สตฺตสติกกฺขนฺธโก นิฏฺฐิโต.\csromanspacer{} จูฬวคฺโค\footnote{ทูรโตปิ อุทปาทิ (ก)} นิฏฺฐิโต.}}

\nitthitam{\textbf{จูฬวคฺคปาฬิ นิฏฺฐิตา.}}
\orphannote{วิ 4. 284 ปิฏฺฐาทีสุปิ.}

\orphannote{มูลายวิสุทฺธินวก (สี, สฺยา)}

\orphannote{อิธ ปน ภิกฺขเว ภิกฺขุ สมฺพหุลา สํฆาทิเสสา อาปตฺติโย อาปชฺชติ ปริมาณมฺปิ อปริมาณมฺปิ ฯเปฯ ววตฺถิตมฺปิ สมฺภินฺนมฺปิ. โส สํฆํ ตาสํ อาปตฺตีนํ สโมธานปริวาสํ ยาจติ, ตสฺส สํโฆ ตาสํ อาปตฺตีนํ สโมธานปริวาสํ เทติ, โส ปริวสนฺโต อนฺตรา สมฺพหุลา สํฆาทิเสสา อาปตฺติโย อาปชฺชติ ปริมาณาโย อปฺปฏิจฺฉนฺนาโย ฯเปฯ ปริมาณาโย ปฏิจฺฉนฺนาโย ฯเปฯ ปริมาณาโย}

\orphannote{มูลายวิสุทฺธินวก (สี, สฺยา) ปฏิจฺฉนฺนาโยปิ อปฺปฏิจฺฉนฺนาโยปิ. โส สํฆํ อนฺตรา-อาปตฺตีนํ มูลายปฏิกสฺสนํ ยาจติ, ตํ สํโฆ อนฺตรา-อาปตฺตีนํ มูลาย ปฏิกสฺสติ ธมฺมิเกน กมฺเมน อกุปฺเปน ฐานารเหน, ธมฺเมน สโมธานปริวาสํ เทติ, ธมฺเมน มานตฺตํ เทติ, ธมฺเมน อพฺเภติ. โส ภิกฺขเว ภิกฺขุ วิสุทฺโธ ตาหิ อาปตฺตีหิ. (1 - 3) อิธ ปน ภิกฺขเว ภิกฺขุ สมฺพหุลา สํฆาทิเสสา อาปตฺติโย อาปชฺชติ ปริมาณมฺปิ อปริมาณมฺปิ ฯเปฯ ววตฺถิตมฺปิ สมฺภินฺนมฺปิ. โส สํฆํ ตาสํ อาปตฺตีนํ สโมธานปริวาสํ ยาจติ, ตสฺส สํโฆ ตาสํ อาปตฺตีนํ สโมธานปริวาสํ เทติ, โส ปริวสนฺโต อนฺตรา สมฺพหุลา สํฆาทิเสสา อาปตฺติโย อาปชฺชติ อปริมาณาโย}

\orphannote{มูลายวิสุทฺธินวก (สี, สฺยา) อปฺปฏิจฺฉนฺนาโย ฯเปฯ อปริมาณาโย ปฏิจฺฉนฺนาโย ฯเปฯ อปริมาณาโย ปฏิจฺฉนฺนาโยปิ อปฺปฏิจฺฉนฺนาโยปิ. โส สํฆํ อนฺตรา-อาปตฺตีนํ มูลายปฏิกสฺสนํ ยาจติ, ตํ สํโฆ อนฺตรา-อาปตฺตีนํ มูลาย ปฏิกสฺสติ ธมฺมิเกน กมฺเมน อกุปฺเปน ฐานารเหน, ธมฺเมน สโมธานปริวาสํ เทติ, ธมฺเมน มานตฺตํ เทติ, ธมฺเมน อพฺเภติ. โส ภิกฺขเว ภิกฺขุ วิสุทฺโธ ตาหิ อาปตฺตีหิ. (4-6) อิธ ปน ภิกฺขเว ภิกฺขุ สมฺพหุลา สํฆาทิเสสา อาปตฺติโย อาปชฺชติ ปริมาณมฺปิ อปริมาณมฺปิ ฯเปฯ ววตฺถิตมฺปิ สมฺภินฺนมฺปิ. โส สํฆํ ตาสํ อาปตฺตีนํ สโมธานปริวาสํ ยาจติ, ตสฺส สํโฆ ตาสํ อาปตฺตีนํ สโมธานปริวาสํ เทติ, โส ปริวสนฺโต อนฺตรา สมฺพหุลา}

\orphannote{มูลายวิสุทฺธินวก (สี, สฺยา) สํฆาทิเสสา อาปตฺติโย อาปชฺชติ ปริมาณาโยปิ อปริมาณาโยปิ อปฺปฏิจฺฉนฺนาโย ฯเปฯ ปริมาณาโยปิ อปริมาณาโยปิ ปฏิจฺฉนฺนาโย ฯเปฯ ปริมาณาโยปิ อปริมาณาโยปิ ปฏิจฺฉนฺนาโยปิ อปฺปฏิจฺฉนฺนาโยปิ. โส สํฆํ อนฺตรา-อาปตฺตีนํ มูลายปฏิกสฺสนํ ยาจติ, ตํ สํโฆ อนฺตรา-อาปตฺตีนํ มูลาย ปฏิกสฺสติ ธมฺมิเกน กมฺเมน อกุปฺเปน ฐานารเหน, ธมฺเมน สโมธานปริวาสํ เทติ, ธมฺเมน มานตฺตํ เทติ, ธมฺเมน อพฺเภติ. โส ภิกฺขเว ภิกฺขุ วิสุทฺโธ ตาหิ อาปตฺตีหิ. (7-9) มูลาย วิสุทฺธินวกํ นิฏฺฐิตํ.}

\orphannote{วิ 5. 202 ปิฏฺเฐปิ.}

\orphannote{วิ 5. 203 ปิฏฺเฐปิ.}

\orphannote{วิ 5. 204 ปิฏฺเฐปิ. ++ วิ 5. 165; อํ 2. 294; ม 3. 33 ปิฏฺเฐสุปิ.}

\orphannote{วิ 3. 60, 65; วิ 4. 375 ปิฏฺฐาทีสุปิ.}

\orphannote{วิ 3. 58 ปิฏฺเฐปิ.}

\orphannote{วิ 3. 63 ปิฏฺเฐปิ.}

\orphannote{วิ 3. 77 ปิฏฺเฐปิ.}

\orphannote{วิ 3. 81 ปิฏฺเฐปิ.}

\orphannote{มหาปชาปตี (สี, สฺยา)}

\orphannote{อํ 1. 362 ปิฏฺเฐปิ.}

\orphannote{วิ 2. 117 ปิฏฺเฐ.}

